% Self-contained PE-CGG-HOUQFT article. The original O-QFT estimates
% and O-QED tests are retained with their action and domain identified.
% Builds with latexmk -pdf -interaction=nonstopmode -halt-on-error UV.tex
% The earlier preprint DOI and 45-result Lean snapshot are historical references.
\documentclass[final,5p,times,twocolumn,numbers,nopreprintline]{elsarticle}
\usepackage[T1]{fontenc}
\usepackage[utf8]{inputenc}
\usepackage[english]{babel}
\renewcommand{\ttdefault}{lmtt} % Outline font for DOI and URL text.
\usepackage{amsmath,amssymb,amsthm,mathtools,bm}
\usepackage{mathrsfs}
\usepackage{slashed}
\usepackage{microtype}
\usepackage{booktabs,tabularx,array}
\IfFileExists{flushend.sty}{\usepackage[nospread]{flushend}}{}
\usepackage{enumitem,needspace}
\usepackage{xcolor}
\usepackage[hidelinks]{hyperref}
\usepackage{bookmark}
\hypersetup{
pdftitle={Geo-Phase Suppression and Ultraviolet Finiteness in PE-CGG-HOUQFT: The Curvature Action, Raw O-QED, and the Induced Electromagnetic Sector},
 pdfauthor={Giovanni Joseph Chiappone},
 pdfsubject={Action-derived ultraviolet estimates, raw O-QED and the electromagnetic sector induced by the graded groupoid Cartan-matrix action},
 pdfkeywords={PE-CGG-HOUQFT, geo-phase suppression, O-QED, effective action, ultraviolet finiteness, BRST, vectorial pole}}
\hypersetup{pdfinfo={DOI={10.5281/zenodo.24407735}}}
\numberwithin{equation}{section}
\newtheorem{theorem}{Theorem}[section]
\newtheorem{proposition}[theorem]{Proposition}
\newtheorem{lemma}[theorem]{Lemma}
\newtheorem{corollary}[theorem]{Corollary}
\theoremstyle{definition}
\newtheorem{definition}[theorem]{Definition}
\newtheorem{assumption}[theorem]{Assumption}
\theoremstyle{remark}
\newtheorem{remark}[theorem]{Remark}
\newcommand{\dd}{\mathop{}\!\mathrm d}
\newcommand{\ii}{\mathrm i}
\DeclareRobustCommand{\bin}{\ensuremath{\Diamond}}
\newcommand{\bb}{\bin} % Compatibility alias; the sector-plane glyph is a diamond.
\newcommand{\RMink}{\mathbb R^{1,3}_{\mathrm{Mink}}}
\newcommand{\RGTL}{\mathbb R^{1,3}_{\mathrm{GTL}}}
\newcommand{\RGT}{\RGTL} % Locked representation, as in the monograph.
\newcommand{\BR}{\mathsf s}
\newcommand{\Ph}{\Phi}
\newcommand{\Tr}{\operatorname{tr}}
\newcommand{\PV}{\operatorname{PV}}
\newcommand{\Disc}{\operatorname{Disc}}
\newcommand{\Cut}{\operatorname{Cut}}
\newcommand{\supp}{\operatorname{supp}}
\newcommand{\diag}{\operatorname{diag}}
\newcommand{\rank}{\operatorname{rank}}
\newcommand{\sgn}{\operatorname{sgn}}
\newcommand{\erfc}{\operatorname{erfc}}
\newcommand{\erfcx}{\operatorname{erfcx}}
\newcommand{\RR}{\mathbb R}
\newcommand{\CC}{\mathbb C}
\newcommand{\one}{\mathbf 1}
\newcommand{\vct}[1]{\boldsymbol{#1}}
\newcommand{\norm}[1]{\left\lVert #1\right\rVert}
\newcommand{\abs}[1]{\left\lvert #1\right\rvert}
\newcommand{\cGT}{c_{\mathrm{GT}}}
\newcommand{\ee}{e_{\bb}}
\newcommand{\epsvec}{\hat\epsilon_{\bb}}
\newcommand{\je}{j_{\bb}}
\newcommand{\embed}{\varphi_{\ee}}
\newcommand{\ExpDir}{\operatorname{Exp}_{\ee}}
\newcommand{\SqrtDir}{\operatorname{Sqrt}_{\ee}}
\newcommand{\LogDir}{\operatorname{Log}_{\ee}}
\newcommand{\EoneDir}{(E_1)_{\ee}}
\newcommand{\Mzero}{M_0}
\newcommand{\bmu}{\boldsymbol{\mu}}
\newcommand{\sMand}{\mathsf{s}}
\newcommand{\tMand}{\mathsf{t}}
\newcommand{\uMand}{\mathsf{u}}
% Preserve terminal punctuation in run-in headings. The class adds a period
% even when the supplied heading already ends in a period or a question mark.
\makeatletter
\ExplSyntaxOn
\cs_new_protected:Npn \UVHeadingPunct #1#2
  {
    \regex_match:nnTF { [.!?] \s* \Z } {#2}
      { #1 {#2} } { #1 {#2.} }
  }
\cs_set_protected:Npn \els@aparagraph [#1]#2
  { \UVHeadingPunct {\elsparagraph[#1]} {#2} }
\cs_set_protected:Npn \els@bparagraph #1
  { \UVHeadingPunct {\elsparagraph*} {#1} }
\ExplSyntaxOff
\makeatother
% Two-column typography shared with the geometric-time companion articles.
\setlist{nosep,leftmargin=*}
\setlength{\emergencystretch}{2em}
\linespread{0.936}
\setlength{\dbltextfloatsep}{6pt plus 1pt minus 1pt}
\setlength{\textfloatsep}{7pt plus 1pt minus 1pt}
\setlength{\abovecaptionskip}{4pt}
\allowdisplaybreaks[1]
\interfootnotelinepenalty=10000
\raggedbottom
\AtBeginDocument{%
 \setlength{\abovedisplayskip}{7pt plus 2pt minus 2pt}%
 \setlength{\belowdisplayskip}{7pt plus 2pt minus 2pt}%
 \setlength{\abovedisplayshortskip}{4pt plus 2pt minus 2pt}%
 \setlength{\belowdisplayshortskip}{4pt plus 2pt minus 2pt}%
}
% Numeric references are embedded; no BibTeX or external figures are needed.
% Published articles below cite their versions of record. For books,
% original publication, corrected reissues, and later electronic formats
% are distinguished rather than treated as interchangeable editions.
% DLMF entries pin the publicly identified release and direct equation links.
% Versioned arXiv links identify cited versions; they need not be latest.
% Author-supplied Zenodo identifiers are retained without inferring
% a deposited-file or version match that has not been independently checked.
% Keep each bibliography entry together across columns and pages.
\AddToHook{env/thebibliography/begin}{\interlinepenalty=10000}
% elsarticle already includes Appendix in appendix reference labels.
\journal{}
\widowpenalty=10000
\clubpenalty=10000
\displaywidowpenalty=10000

\begin{document}
\begin{frontmatter}
\title{Geo-Phase Suppression and Ultraviolet Finiteness in PE-CGG-HOUQFT\\
\large The Curvature Action, Raw O--QED, and the Induced Electromagnetic Sector\tnoteref{preprintdoi}}
\tnotetext[preprintdoi]{Zenodo identifier for this version:
\href{https://doi.org/10.5281/zenodo.24407735}{\nolinkurl{10.5281/zenodo.24407735}}.}
\author[aff1]{Giovanni Joseph Chiappone\corref{cor1}\fnref{orcid}}
\cortext[cor1]{Corresponding author.}
\ead{giovannichiappone252@gmail.com}
\address[aff1]{Independent Researcher, New York, United States}
\fntext[orcid]{ORCID: \href{https://orcid.org/0009-0005-1689-463X}{0009-0005-1689-463X}.}
\begin{abstract}
Two charged insertions may share an endpoint while retaining a relative
displacement along their histories. Averaging the associated phase can
suppress hard loop momenta. We ask when this displacement survives the
curvature action, the charged determinant, and the boundary integration
on a locked sector plane of PE-CGG-HOUQFT. The regular groupoid trace
cancels a common displacement phase. A curvature pairing that retains
two distinct insertions gives an explicit positive result. On the
rotation-singlet A4 shell, five consecutive geometric diffusion
transits suffice for absolute ultraviolet convergence of every
interaction graph of the neutral curvature determinant, including
contacts and simultaneous short durations. A harmonic exit from a
radial segment gives the same bound for the full angular trace.

The charged sector imposes an independent test. Ordinary minimal
coupling of the neutral construction restores a magnetic logarithm.
A momentum-groupoid model has finite-volume graphs that converge at
each finite order, yet its conserved transverse current loses the
required charged-pair weight at large boost. In a scalar comparison, a
tree-eliminated shadow gives a Euclidean bound while its square-root
geo-phase removes positive weight from a physical cut. Under the
stated spectral assumptions a positive Born mass band also obstructs
absolute convergence of the ordinary coincident Maxwell closure.
On a quartic normal surface selected by the curvature action, we give
a necessary and sufficient short-duration test for every finite scalar
loop with specified angular shifts, inverse radial powers and ordinary
duration measure. It includes loops that pass through angular momentum
zero. For curvature insertions with square-integrable, zero-mean
angular profiles uniformly bounded in angular norm over radius, the
absolute sums converge at every positive order. For bounded real
profiles the relative scalar determinant is analytic near zero
source amplitude. We then couple a local charged scalar to a
horizontal magnetic field. In four carrier dimensions its quadratic
coefficient survives every bounded positive change of diffusion
rate along the normal path. On the shrinking normal circle even a
fixed nonzero winding leaves a divergence. A radial mark gives a
sharp criterion for suppressing it. The unrestricted claim still
requires the physical charged source to yield the needed relative
displacements, photon contacts and measured charged-pair spectrum.
\end{abstract}

\begin{keyword}
PE-CGG-HOUQFT \sep ultraviolet finiteness \sep geo-phase suppression
\sep effective action \sep O-QED \sep geometric-time locking
\sep path groupoid \sep BRST symmetry
\end{keyword}
\end{frontmatter}

\section{Introduction}
\label{sec:introduction}

Two virtual histories can close at the same point and have different
ultraviolet behavior. Consider normal diffusion generated by
$\partial_y^2$. A return from $0$ through a specified intermediate point
$d$ and back to $0$ in two intervals of length $t/2$ contributes
\[
 K_{t/2}(0,d)K_{t/2}(d,0)
   =\frac{1}{2\pi t}\exp\!\left(-\frac{d^2}{t}\right).
\]
The net normal displacement is zero. The route through $d\ne0$
nevertheless suppresses short proper times, which correspond to hard
loop momenta. Integrating over the intermediate point with ordinary
Lebesgue measure instead gives $K_t(0,0)=(4\pi t)^{-1/2}$. With a
quantum weight on $d$, the ultraviolet behavior depends on its mass near
the origin. In the positive midpoint trace below the magnetic coefficient
is finite exactly when the zero-visit contribution to this coefficient
vanishes and the midpoint law has a finite inverse-square moment,
counting an atom at zero as infinite cost. A charged insertion can
vanish on an unvisited channel. A nonzero positive boundary state
$\omega$ on a unital $*$-algebra cannot set $\omega(1)=0$, since
$|\omega(a)|^2\le\omega(a^*a)\omega(1)$ for every $a$.
For a return of zero net winding, the remaining question is how much
weight the sewn source gives to a nearly erased visit.

We use this distinction to analyze the Path-Enriched Chirally Graded
Groupoid Extension of the Hyper-Omniverse Unified Quantum Field Theory
(PE-CGG-HOUQFT). Its proposed microscopic action pairs curvatures of a
graded groupoid Cartan--matrix (GCM) superconnection. After coordinate
geometric-time locking (GTL) and the separately declared label-diagonal
restriction, calculations on one $1+3$ dimensional Lorentzian plane use
$\bb=(s,b_s)$, with $s$ the sector and $b_s$ the affine displacement
on its carrier. The source records the histories and states that survive
these restrictions.

With $\mathbb F_{\bin}=\mathbb D_{\bin}^2$, its curvature functional has scalar coefficient
\begin{equation}
\begin{aligned}
 S_{*,\mathrm{sc}}[q]={}&\mathcal N(G_N)\sum_{\bin}
  \int d^K X_{{\rm GT},\bin}\,d^3\mathbf x\
              \Delta_{{\rm GTL},\bin}[q](X_{{\rm GT},\bin})\\
 &{}\times\int_{\mathbb R^K}d^K v\
  \delta(\|v\|^2-R^2)\,
  \frac{2\sigma_\epsilon v_\Gamma\,
              \theta(\sigma_\epsilon v_\Gamma)}
       {\operatorname{vol}(B_R^{K-1})}\\
 &{}\times\int_{\mathcal I_r}dr\
              \mathcal C_{\bin}[q;X,\mathbf x,v,r],\\
 \mathcal C_{\bin}={}&\operatorname{Re}\tau_{\mathcal P,\bin}[q]\,
     \operatorname{Str}\!\left[
       \mathbb F_{\bin}[q]\mathbb P_{*,\bin}[q]\mathbb F_{\bin}[q]
                         \right]_{\rm top},
 \\[-2pt]
 \mathbb F_{\bin}={}&\mathbb D_{\bin}^{\,2}.
\end{aligned}
\label{eq:uvpeA-curvature}
\end{equation}
Throughout the curvature calculations $S_*$ abbreviates the scalar
coefficient $S_{*,\mathrm{sc}}$. The physical directed action takes
values on each carrier line. On an aligned support $\bin=(s,b_s)$,
write $u_{g,\bin}=\sigma_\epsilon\hat p_{0,\bin}$ and
$u_{m,\bin}=\sigma_s u_{g,\bin}$. The pure reductions give
$\boldsymbol S_{j,\bin}=u_{j,\bin}S_{j,\bin,\mathrm{sc}}$ for
$j=g,m$. The unit multiplies the real coefficient on its own support;
a mixed term requires a prescribed pullback before taking source
derivatives. Section~\ref{sec:path-action-type-uv} gives the
coefficient extraction and the quantum phase.
For the selected label-diagonal representatives, let
$R=\|\boldsymbol c\|>0$, $v_\Gamma=v\cdot e_\Gamma$, with
$e_\Gamma$ the chosen shell axis and $\sigma_\epsilon$ the
ordering-branch sign. The argument $X$ of $\mathcal C_{\bin}$ abbreviates
$X_{{\rm GT},\bin}$. The ambient base point is still integrated.
In an equal-component GT chart the invariant localization is
\begin{equation}
\begin{aligned}
 \Delta_{{\rm GTL},\bin}
 &=\delta_{\hat G_{\bin}^{\perp},\bin}
        (Q_{\bin}X_{{\rm GT},\bin})\\
 &=\sqrt K\prod_{\Lambda\in\mathcal T_K\setminus\{0\}}
          \delta^{(1)}_{\bin}(x^\Lambda_{\bin}-x^0_{\bin}).
\end{aligned}
\label{eq:uvfinal-normal-delta}
\end{equation}
With $Q_{\bin}$ the orthogonal projector normal to the carrier,
this normalization yields its induced length measure. The
equal-component constraint matrix has Gram determinant $K$.
For $K=1$ the product is empty. Each bare vertex uses its own
ambient base point with a single GTL pullback. The optional
dimensional normalization of the GTL insertion belongs in the
density once. The separate A2/GTP momentum binder restricts the
allowed geo-momentum. Neither delta changes the raw axis labels
$\{\bin_\Lambda\}$ into the representative $\bin=(s,b_s)$.
The displayed sum uses the declared copy measure when the family
is continuous.

The velocity integral is the full positive oriented scalar A4
shell measure $d\nu_{\sigma_\epsilon}^{[\Gamma]}(v)$. Its radius
has velocity units. The normal history variable $r$ and the
ambient GT coordinate have their own integration domains. The
directed lifts on an aligned chart are
$\sigma_\epsilon\hat p_{0,\bin}$ for the geometric reduction and
$\sigma_s\sigma_\epsilon\hat p_{0,\bin}$ for the matter reduction.
The pole offset on the latter uses
$i\sigma_{s\epsilon}\hat\epsilon_{\bin}$,
$\hat\epsilon_{\bin}=\varepsilon\hat p_{0,\bin}$ with
$\varepsilon\downarrow0$ and $\sigma_{s\epsilon}=\sigma_s\sigma_\epsilon$.
The two pole displacements have opposite signs along this same
prescribed vector. The shell sign has already selected its own
velocity hemisphere. The following Euclidean estimates are made on
that plane before interpreting a continued amplitude; the embedded
locked plane $\RGT$ retains its carrier and affine origin, while
$\RMink$ supplies the isometric Minkowski readout in a matched frame.

We distinguish three electromagnetic inputs. The raw O--QED source
starts with a local Maxwell--Dirac density directly on the locked
carrier. The fixed-depth O--QED prescription supplies an additional
multiplier on uncut virtual lines and requires compatible interaction
vertices. GCM-induced O--QED comes from integrating out ambient fields
while retaining the photon and charged probes. In that construction,
the effective charged operator, its closed determinant and the photon
contacts must all follow from the same source. Table~\ref{tab:g5R-scope}
gives the scope of the corresponding convergence statements.

At each finite perturbative order, the action and its boundary state
specify which geometric variables must be integrated to form the
complete graph coefficient. Ultraviolet finiteness here means absolute
convergence of its remaining loop integral or sum, including regions
where only a subset of the independent momenta grows. The order of
operations matters. A real displacement phase has unit modulus for
each history. The modulus of its average can decay because different
histories have different phases. Moving the absolute value inside that
average would lose this mechanism. A stronger estimate is available
for the heat-link model: the absolute graph sum at each fixed order
remains integrable over its common radial depth. The infrared
conditions are specified with each model. Continuation uses the
vectorial boundary prescription on the selected sector plane.

For the line-kernel estimates the comparison multiplier is
\begin{equation}
\begin{aligned}
 \Ph(z)&=\int_{[0,\infty)}e^{-\tau z}\rho_{\bb}(\dd\tau),\qquad \rho_{\bb}\geq0,\\
 0&<\Mzero:=\rho_{\bb}([0,\infty))<\infty .
\end{aligned}
 \label{eq:bernstein}
\end{equation}
It acts on the nonnegative Euclidean kinetic operator of an uncut
virtual line. A fixed $a>0$ sets the depth scale in $\Ph(a^2q^2)$.
Equation~\eqref{eq:virtual-edge-depths} gives the edgewise convention.
The physical kernel convention of Ref.~\cite{ChiapponeGHOUQFT2026}
uses $0<\Ph(0)\le1$. The comparison estimates allow any finite
positive $M_0$.

The suppression has two geometrical origins. For a heat displacement
$z$ with lapse $\kappa r$, the modulus of the phase $e^{\mathrm i k\cdot z}$
is one. Its average against the displacement law is $e^{-\kappa r|k|^2}$.
The independently derived harmonic boundary value at distance $r$
has multiplier $e^{-r|p|}$. Both distances are read at the same
position in the normal state. Corollary~\ref{cor:geo-free-slope-uv}
works out the three inputs together on a surface selected by its
curvature action.

The angular coordinate gives a second way to see why the probe
matters. A charged interval on a shrinking circle suppresses changes
in a suitable radial coordinate. Several junctions can still move
toward the tip together. Theorem~\ref{thm:uv-normal-word} gives a
finite test for every such group in a scalar loop. A second diagnostic
asks what happens when a probe changes the phase around the normal
circle. A pure gauge leaves a closed trace invariant. A change of
holonomy alters it even at arbitrarily short duration, as quantified
in Proposition~\ref{prop:quartic-holonomy-uv}. The physical electromagnetic
source determines whether it varies this connection at all.
Proposition~\ref{prop:horizontal-magnetic-word-uv} takes the
complementary route: it holds the normal connection fixed and
computes the response of an ordinary horizontal charged kinetic
term. This tests whether a fluctuating diffusion time or a retained
winding can protect that response.

The heat-link source permits two distinct tests. At fixed depth a marked
vertex suppresses the independent momentum of each charged loop, and
each end of a dynamical photon carries the harmonic readout.
Theorem~\ref{thm:s6-heatlink-uv} gives the corresponding graph bound.
When the depth fluctuates, all these factors share a single radial
coordinate. The bare charged determinant favors short visits, while the
radial ground state suppresses them. Proposition~\ref{prop:dualrev-uvheat-depth}
decides their competition for the free weight. Equation~\eqref{eq:dualrev-uvheat-tail}
bounds a hard momentum tail while any chosen photon durations shrink.
Theorem~\ref{thm:dualrev-uvheat-real} treats the unexpanded determinant
in the dynamical Maxwell field. Its source and normalization are the
same ones used to obtain the electromagnetic vertices. At matched
coefficients Proposition~\ref{prop:uvheat-critical-charge} uses the
Gaussian Maxwell law to test differentiability in the electric charge.

First read the displacements where they occur in the charged
operator, including closed cycles of any variables integrated out.
If two independent loop momenta enter a surviving phase only as
$k_1+k_2$, the direction $k_1=-k_2$ escapes that phase.
Proposition~\ref{prop:physical-uv-retained-minors} tests all such
directions on the actual boundary constraint and gives the cost when
several durations collapse. The same charged source must also supply the
photon response and the open charged line.
Proposition~\ref{prop:phys-uv-relative-chord} exhibits a finite
relative Wilson interaction with the ordinary free Dirac pole and
derives its one-loop soft matching condition. The summed physical cut
in Proposition~\ref{prop:physical-inclusive-opening-uv} explains why
one also needs the spectrum of the boundary states that current can reach.

The distinction between a finite graph and a physical particle is
sharpest in the photon and charged two-point functions. A massless
Maxwell susceptibility grows as the inverse square of the Euclidean
momentum. Theorem~\ref{thm:softward-allmass-uv} controls that
growth in the unexpanded fixed-depth source. For the relative Wilson
source, first averaging the one-loop response misses the fluctuations
between depths. Proposition~\ref{prop:spec-uv-soft-variance} calculates
their fourth-order contribution. Proposition~\ref{prop:uv-relative-spectrum}
tests positivity independently using measured frequency moments of the
interacting charged line. The spatial-transfer examples in
Section~\ref{sec:spex-photon-uv} show how a positive energy measure
changes those moments.

The physical question can first be tested on three source derivatives.
Proposition~\ref{prop:closure-native-uv} reads the frequency
dependence of the measured field directly from the normal curvature
form. Theorem~\ref{thm:bulk-magnetic-uv} tests that field in the
thermodynamic limit at fixed charge. Equation~\eqref{eq:closing-uv-density-hard}
tests a lone temporal bubble on a hard frequency region. In a two-band
hopping action the occupied onsite vacuum determines a second term.
Their common-routing difference is absolutely integrable before the
charged frequency is integrated
(Proposition~\ref{prop:native-relative-density-uv}).
Proposition~\ref{prop:closing-uv-ordered-Maxwell} follows the complete
charged paths through Coulomb exchange.
Corollary~\ref{cor:closing-uv-pair-bulk} constructs its gapless photon states.
The same action admits convergent bare integrals with open charged
lines (Theorem~\ref{thm:ct-uv-gauged-graphs}). The spinor application
in Corollary~\ref{cor:reference-spin-uv} shows which part of the
Dirac dispersion this construction recovers. Proposition~\ref{prop:reference-test-uv}
tests whether a proposed chiral reference vacuum can be removed in
a scalar source without changing the measured charge response.

The neutral return can be tested with photons alone.
Proposition~\ref{prop:seam-closed-uv} follows a virtual pair
between two distinct times and retains both spatial displacements.
The resulting obstruction depends on how finely the photon resolves
the shorter visit. Theorem~\ref{thm:seam-calibrated-joint-uv}
gives the corresponding all-orders estimate when the normal
Maxwell form fixes that resolution in proportion to the heat width.
Proposition~\ref{prop:seam-selection-uv} explains why measurements on
one charged leg cannot decide which bound applies to a sewn pair.
Proposition~\ref{prop:seam-uv-two-pair} tests the next independent
displacement when two pairs exchange partners.
Corollary~\ref{cor:seam-optical-shifts-uv} controls this exchange at
the matched optical length. The amendment fixes these lengths
before quantization. Its finite-region result on the physical Gauss
quotient is Theorem~\ref{thm:omega-quotient-uv}. The separated current
in Proposition~\ref{prop:omega-current-cone-uv} tests the physical
spectrum of that same action.

Section~\ref{sec:carrier-dispersive-uv} constructs a dispersive
charged comparison with a finite-region graph bound.
Section~\ref{sec:carrier-sewing-uv} tests the physical current and
photon sewing. We return to the ambient action in
Section~\ref{sec:monomial-selection-uv}, where its velocity
measure is evaluated before the charged trace. The two-arc
calculation in Proposition~\ref{prop:nativefull-shell-mark-uv}
compares a selected internal transition with the full electric sum. Proposition~\ref{prop:nativefull-alllegs-uv} includes all charged
circulations and exhibits a finite photon-mass obstruction.
Section~\ref{sec:ward-hop-uv} shows how gauging a complete shell
displacement removes that obstruction. Section~\ref{sec:ward-normal-uv}
tests a representation with a positive free charged spectrum and
compares two independent predictions for the production of a charged pair.
The charge and magnetic form factors in Section~\ref{sec:nov-transverse-uv}
test whether changing the open vertex is compatible with varying the
same determinant. Section~\ref{sec:nov-positive-uv} gives a positive
Hamiltonian comparison with explicit charged-pair production, and
Section~\ref{sec:nov-ring-uv} tests whether its physical response
could be sewn by ordinary covariant Maxwell lines.
The interval of exact optical agreement is determined in
\eqref{eq:fix-threshold-cuts-uv}. A charged Fock-space construction
tests equal-time locality in Section~\ref{sec:fix-fock-uv}. For the
ultraviolet question, Section~\ref{sec:fix-record-rank-uv} gives a
sharp record-rank criterion and Section~\ref{sec:nativefull-subgraphs-uv}
tracks the displacement left in each independent hard subgraph.
The measure is tested explicitly in Section~\ref{sec:next-physical-join-uv}.
Restoring its charged Jacobian recovers the missing pair weights and
revives a bare Maxwell logarithm. In Section~\ref{sec:nextsoft-uv}
a change to the charged vertex gives analytic soft response and
positive two-current spectral density in a common action. Its
remaining inclusive deficit is measurable. Section~\ref{sec:next-physical-photon-uv}
shows why the positive spectrum of the photon used to close a loop
is an additional restriction on ultraviolet convergence.

Section~\ref{sec:joint2-map-uv} follows the same measured charged
pair through its open transition, the spacelike current and the photon
used to join loops. The joint normal heat preparation in
Section~\ref{sec:joint2-charge-uv} gives an exact one-loop cut test
and an explicit negative reflection norm. The inverse construction in
Section~\ref{sec:joint2-inverse-uv} recovers a positive photon in the
neutral pair channel and permits its resolved determinant to be tested.
A fixed massless direction then permits the same photon to mediate a
nonlinear interaction with the neutral pair modes, as in
\eqref{eq:joint2-fixed-null-uv} and its observed-photon version.

The next test opens the pair into individual fermions and follows the same
current through its Gauss and transverse photon contacts
(Sections~\ref{sec:next3-map-uv}--\ref{sec:next3-gauss-car-uv}).
Equation~\eqref{eq:next3-dipole-energy-uv} computes the remaining bare cost
of a finite-mass continuum under boosts. Section~\ref{sec:next3-null-uv}
gives the first nonlinear reflection test for the earlier massless model.
The first higher-energy test is the sign of the charged determinant's
lowest pair cut (Section~\ref{sec:fourcone-uv}). Transporting spin
along the velocity history changes that sign and predicts an independent
magnetic moment (Section~\ref{sec:fourspin-berry-uv}).
Section~\ref{sec:four-angular-superconvergence-uv} tests the complete
current at all Euclidean frequencies. Section~\ref{sec:four-excursion-uv}
computes the contact that a rough closed excursion contributes and
compares it with the transition measured on an open charged line.
A static magnetic background distinguishes excursion laws with the
same vacuum pole structure. Section~\ref{sec:fifth-mag-uv} derives
the necessary spectral half-moment and constructs a law at the limiting
power. Section~\ref{sec:fifth-joint-completeness-uv} tests whether
unobserved normal records can suppress the photon contraction when all
physical outcomes are included. Section~\ref{sec:fifth-elastic-uv}
restores the single-charge vertex in the Hamiltonian comparison and
checks its forward energy cone.

The local curvature test in Section~\ref{sec:six-physical-uv} applies
directly to the measured photon. Equation~\eqref{eq:six-relative-ranks-uv}
then tests several relative displacements drawn from the same geometry.
Section~\ref{sec:six-recoil-uv} follows an ordered transit through the
kinetic and Coulomb contacts of a charged Hamiltonian.
Section~\ref{sec:six-mag-uv} repeats the magnetic test with the other
spatial polarization. Section~\ref{sec:six-joint-uv} identifies which
mixed source measurement tests a proposed joint law for the photon
and the charged path.

The groupoid and shadow calculations start in
Section~\ref{sec:seven-groupoid-uv}. First compare the phase of an
individual arrow with the positive norm that appears when it is
closed. Section~\ref{sec:seven-radial-uv} then computes a full
curvature determinant with relative displacements. The charged
pair and shadow tests in Sections~\ref{sec:seven-pair-uv}
and~\ref{sec:seven-shadow-uv} track the intermediate states when
the same source is opened.



\begin{table*}[t]
\centering\small
\renewcommand{\arraystretch}{1.13}
\caption{Electromagnetic sources and the scope of their convergence
results. Each entry uses its own charged determinant and photon measure.}
\label{tab:g5R-scope}
\begin{tabularx}{\textwidth}{@{}>{\raggedright\arraybackslash}p{.19\textwidth}
 >{\raggedright\arraybackslash}X>{\raggedright\arraybackslash}X@{}}
\toprule
Electromagnetic source & Data that define it & Result established here\\
\midrule
Raw local O--QED &
A locked Maxwell--Dirac density with the ordinary charged measure. &
The bare coefficient of $F_{\mu\nu}F_{\mu\nu}$ diverges
logarithmically (\eqref{eq:uvpeA-raw-log}).\\
Fixed-depth prescription &
A specified multiplier on uncut lines and the chosen vertices for each
test. &
The explicit locked kernel has all inverse moments
(\eqref{eq:OQED-inverse-moments}). The tested
Ward-compatible vertices yield divergent contractions.\\
First-order GCM representation &
Independently integrated charged fields, fixed curvature pairing,
normal constraint and its measure, visit readout, and dynamical Maxwell field. &
Every interacting bare graph is finite at each electric order. A
charged radial state also permits arbitrarily short visits
(Corollaries~\ref{cor:g5R-induced-Dirac} and~\ref{cor:g5R-fluctuating-induced}).
The soft transverse response in \eqref{eq:g5M-soft-mass} is nonzero.\\
GCM heat-link representation &
A graded charged curvature pairing, heat displacement, interval
transport and harmonic transverse photon readout. &
Every bare graph is UV finite on the curvature-minimizing surface
(Corollary~\ref{cor:geo-free-slope-uv}). The free-vacuum threshold for
power profiles and a sufficient real-charge partition bound are given in
Proposition~\ref{prop:dualrev-uvheat-depth} and
Theorem~\ref{thm:dualrev-uvheat-real}.
At fixed depth the exact weak-charge susceptibility has Maxwell
scaling under Theorem~\ref{thm:softward-allmass-uv}.
At each positive mass its bulk magnetic response fails the positive
spectral test (Theorem~\ref{thm:bulk-magnetic-uv}).\\
Relative Wilson chord source &
Massive free Dirac operator plus the difference of the Wilson chord
and its value at zero field, with the same harmonic readout. &
Every bare graph is UV finite. The continued free symbol has the
Dirac pole. A global fluctuating depth fails fourth-order soft
matching, and the interacting open field fails spectral positivity
(Propositions~\ref{prop:spec-uv-soft-variance}
and~\ref{prop:uv-relative-spectrum}).\\
Time-local magnetic pair source &
A spatially distributed magnetic transition creates a charged particle
and hole. The photon has its canonical time evolution. &
The bare magnetic graphs converge absolutely. The interacting transfer
has positive spectral weights
(Section~\ref{sec:pairferm-uv}).\\
Time-local electric pair source &
A two-band Hamiltonian with spatial Wilson hopping, a common electric
readout and the Gauss energy. &
At fixed widths the weak-charge transverse covariance obeys two-sided
Maxwell bounds (Theorem~\ref{thm:closing-uv-positive-susceptibility}).
The vacuum-normalized density pair is jointly integrable
(Proposition~\ref{prop:native-relative-density-uv}).\\
Spinor with geometric spatial hopping &
Dirac matrices contract the derivative of a heat displacement.
The onsite occupied slots fix the vacuum reference. &
At fixed widths, charged Euclidean coefficients are finite
at nonexceptional transfers and on joint momentum packets
(Theorem~\ref{thm:seam-uv-open}).\\
Moving neutral pair &
The charged lapse is proportional to $r$. A regular normal Maxwell
form samples the photon over a distance of order $r$. &
A connected transverse two-photon coefficient diverges
at $g^2 e^{26}$ on contractible returns
(Proposition~\ref{prop:seam-closed-uv}).\\
Pair with an angular transition at each hop &
An occupied pair raises the normal angular index by a nonzero
integer. Annihilation lowers it by the same amount. &
Every electric order of the closed two-hop return is finite.
An exchange gives a divergent connected two-photon coefficient
at $g^4e^{120}$
(Proposition~\ref{prop:seam-uv-two-pair}).\\
Spinor with matched optical length &
The charged lapse is $\kappa r^{2/3}$ and the photon readout
is its square root. The normal state evolves on the slab. &
All angular sectors, including transitions at pair vertices, admit
joint absolute bare graph bounds at each finite electric order
(Corollary~\ref{cor:seam-optical-shifts-uv}).\\
Electromagnetic $\Omega_\alpha$ amendment &
The equal-volume law fixes the readout. Gauss reduction and a paired
chiral mass fix the physical electric sources. &
All finite electric orders on the finite-region quotient
(Theorem~\ref{thm:omega-quotient-uv}). Its current tests forward-cone
support (Proposition~\ref{prop:omega-current-cone-uv}).\\
Dispersive charged replacement &
The diagonal carrier has relativistic energy. A spatial jump samples
the connection along its complete normal path. &
All finite electric orders on its slab. The zero-electric neutral
matter threshold obeys the forward-cone bound for every pair
hopping strength (Section~\ref{sec:carrier-dispersive-uv}).\\
\bottomrule
\end{tabularx}
\end{table*}

The monograph and companion articles supply the geometric conventions
and the fixed-depth benchmarks
\cite{ChiapponeGHOUQFT2026,ChiapponePartI2026,ChiapponePartII2026,ChiapponeIgnition}.
Postulate~$\Omega_\alpha$ specifies the proposed common source for
A4, A5, Q1, fixed-plane $\Omega$ and groupoid $\Omega_{\mathfrak G}$.
The resulting construction and recovery conditions are given in
Section~\ref{sec:pe-principle}. The Lean~4 inventory in
\ref{sec:lean-verification} refers to the 45 results in the
archived version~\cite{UVLean}.

\subsection{Convergence, gauge symmetry, and physical reconstruction}

\paragraph{Absolute convergence of complete graphs}
On the Euclidean carrier the scalar comparison covariance is
\begin{equation}
 C_{\Ph,E}(q)=\frac{\Ph(a^2q^2)}{q^2+m^2}.
 \label{eq:covariance-intro}
\end{equation}
The Laplace law in \eqref{eq:bernstein} makes its numerator a positive
average of Gaussians. Small $\tau$ gives broad support in momentum, so
the weight near zero controls how much damping is available. The
momentum routing determines which independent loop directions that
damping reaches. Section~\ref{sec:uv} combines these two facts with bounds
on the complete vertices.

\paragraph{Gauge compatibility of lines, vertices, and contacts}
A nonconstant kinetic modification can change the vertices required by
the master equation. A multiplier with
$\BR\Ph(a^2\mathscr L_{\bb})=0$ is inert under the quantum BRST operator
$\BR$. The Slavnov--Taylor argument also needs a completed action and an
invariant measure. Section~\ref{sec:brst} gives the identity for that
system. A factor that equals one at zero photon source can have a
nonzero second source derivative, which contributes a two-photon
contact in the same test.

\paragraph{Reconstruction of physical states and cuts}
The physical observable algebra must satisfy reflection positivity and
the other reconstruction hypotheses. A prescribed on-shell replacement
also has to reproduce the analytic discontinuity of the same graph.
The counterexamples and comparison calculations in
Sections~\ref{sec:analytic} and~\ref{sec:examples} test these conditions
for positive heat kernels and for explicitly evaluated bubbles.

\subsection{Assumptions and scope}
\label{sec:assumptions}

The following hypotheses govern the general fixed-plane estimates and
their physical interpretation. They are used in the combinations stated
below. When convergence and symmetry are claimed together, both sets
of hypotheses must hold for the same completed action.

\begin{description}[leftmargin=1.6cm,style=nextline]
\item[Geometry.]
The fixed-plane estimates use one locked support \(\bb=(s,b_s)\),
its label-diagonal restriction and the independent A2/GTP momentum binder.
A path travelling between several supports additionally requires an
explicit source and target for each arrow and a rule for the internal support
sum. On each chart the Euclidean comparison operator is nonnegative on
its specified domain.

\item[Kernel.]
For the line-kernel theorems, each selected uncut internal
covariance has the finite positive heat representation \eqref{eq:bernstein}.
The gap and inverse-moment assumptions apply to the particular estimates
that use them. Theorem~\ref{thm:update-joint-heat} gives a sufficient
budget for a correlated law. In a path-enriched action the relevant
positive majorant may instead arise after integrating a complete
history or its fluctuation determinant. Its derivation must include
the vertices produced by that action. The convenient normalization
\(\Ph(0)=1\) is optional for the Euclidean comparison.

\item[Vertices and IR.]
The same completed vertices used in the symmetry identities satisfy
the growth bounds of the ultraviolet theorem. Positive masses or a
separately specified integrable infrared envelope control the long
proper-time region. External data lie in a compact domain on which
these bounds are uniform.

\item[Quantum symmetry.]
An off-shell nilpotent BRST differential, a completed BV action,
compatible source couplings and covariance exchange symmetry, and an
invariant anomaly-free measure
exist on the stated perturbative domain in both carrier characters.
The identities concern this complete system with simultaneous
projection of all fields, antifields, sources, and coefficients.

\item[Continuation.]
A declared continuation of the kernels, the plane-resolved
vectorial boundary datum, and a common integrable majorant are
available on the complex domain actually used.
Contour deformations, when used, have their own boundary and
no-pinch estimates.

\item[Physical completion.]
Reflection positivity on the observable algebra, the other
reconstruction axioms, isolated asymptotic poles, and the matching of
cuts to discontinuities are additional requirements. Any optical
identity also uses scattering unitarity and a complete physical state
resolution, or the explicitly stated equivalent cut-factorization
hypotheses. They are invoked
only in the results that explicitly assume them.
\end{description}

These hypotheses enter in the combinations stated by the individual
theorems. The absolute bounds are coefficientwise. A sum over
perturbative orders, a sum over support-changing histories, or an
average over backgrounds adds another convergence problem.
Section~\ref{sec:pe-principle} supplies a sufficient bound for a
countable history sum. For a fluctuating positive depth, the proof
must integrate the bounds on the complete vertices with the actual
small-depth law. An abstract bound applies once these inputs are known.
In the explicit models they are calculated from their declared sources.

\paragraph{Reading the proofs}
Sections~\ref{sec:pe-actions} and~\ref{sec:pe-graphs} trace the
magnetic divergence to the charged determinant and show how a measured
return can change the graph estimate.
Sections~\ref{sec:pe-curvature}--\ref{sec:g5-induction} then test
specific representations and dynamical photon sewings. The convention
table below and Section~\ref{sec:framework} fix the geometric and
carrier notation used in the general graph, BRST and reconstruction
results. Section~\ref{sec:OQED-application} evaluates the original
fixed-depth O--QED prescription. The appendix on logical dependence
(\ref{sec:dependence}) records which assumptions each conclusion uses.

\begin{table*}[t]
\centering
\small
\renewcommand{\arraystretch}{1.15}
\caption{Principal notation. Each support label is fixed unless a sum
or integral is explicitly written. On that fixed copy, $X_0$ and
$k_\parallel$ abbreviate $X_{0,\bb}$ and $k_{\parallel,\bb}$. The
three ordinary spatial coordinates carry no sector-plane label.}
\label{tab:notation}
\begin{tabularx}{\textwidth}{@{}>{\raggedright\arraybackslash}p{0.24\textwidth}X@{}}
\toprule
Symbol & Meaning\\
\midrule
\(K\) & Number of ambient geometric-time directions.\\
\(\bin_\Gamma\) & Raw component label
\(\bin_\Gamma=(s_\Gamma,b_{s_\Gamma}^{\rm comp})\), before the
label-diagonal restriction.\\
\(\bb=(s,b_s)\) & Label-diagonal support with sector \(s\) and
normalized carrier-axis offset \(b_s\).\\
\(X_0,\ k_\parallel\) & Centred normalized carrier coordinate (the monograph's
\(G_{0\bb}\)) and its conjugate momentum coefficient.\\
\(a,\ \rho_{\bb},\ \Ph\) & Fixed inverse-momentum scale, finite positive
Laplace measure, and its dimensionless transform.\\
\(a_{\rm ker}^{(p)},\ M_{\rm ker},\ \mu\) & O--QED depth parameter,
auxiliary kernel scale, and their product \(\mu=a_{\rm ker}^{(p)}M_{\rm ker}\);
not the electron pole mass or laboratory response length.\\
\(M_r[\rho_{\bb}]\) & Inverse moment of order \(r>0\); absence of an atom
at zero is stated separately.\\
\(P_{r,M,\bb},\ \mathsf w_{\bb},\ \boldsymbol\varrho_{G,\bb}\) & Forward
normal transfer, a declared depth probability, and the joint heat law
conditional on that depth. The depth \(r\) is distinct from a moment order.\\
\(\dd\mu_{E,\bb},\delta_{\mu_{E,\bb}},\mathfrak j_{\bb}\) & Positive
continued support density, its identity kernel, and the interaction source;
\(\je=\ii\ee\) remains the pole generator.\\
\(\mathbb D,\mathbb F,\mathbb P_*\) & GCM superconnection, curvature, and curvature pairing.\\
\(S_*^{\rm BV},\ \Psi_\partial\) & Quantum bulk action on the chosen gauge slice (BV data) and boundary state.\\
\(R,\ K\) in \eqref{eq:uvpeP-groupoid} & Boundary precision and heat generator on the represented support space. Here the operator \(K\) is distinct from the number of GT directions.\\
\(Q,\ K_S\) in Section~\ref{sec:pe-graphs} & Quadratic form after the normal average, and the loop subspace left undamped and unseen by the lines outside \(S\).\\
\(\mathcal K,\ D\) & Compact external-momentum set and total
numerator-growth degree in the graph estimate.\\
\(\alpha_i,\tau_i,w_i\) & Schwinger parameter, dimensionless Laplace
parameter, and combined width \(w_i=\alpha_i+a^2\tau_i\).\\
\(d,\ j,\ \sigma^2=r^\alpha\) & Exponent of the radial metric,
angular number of its state, and charged heat lapse in the
fluctuating-readout model. The radial exponent \(d\) is distinct
from the carrier dimension, which is four.\\
\(\sigma_s,\sigma_\epsilon,\eta\) & Sector orientation, ordering
branch, and their product \(\eta=\sigma_s\sigma_\epsilon\).\\
\(\ee,\hat u_{\bb},\widehat u_{k,\bb}\) & Fixed pole carrier, branch-action unit
\(\hat u_{\bb}=u_{\epsilon,\bb}\), and running line unit.\\
\(\BR,\Sigma,\Gamma\) & BRST differential, completed action with
antifield sources, and one-particle-irreducible effective action.\\
\bottomrule
\end{tabularx}
\end{table*}

\section{The quantum source required by PE-CGG-HOUQFT}
\label{sec:pe-principle}

An internal history begins on one locked support and ends on a support
to which the field can propagate. A path records the entire ordered
list of such visits. Its endpoint arrow retains only the composite
in the groupoid. For example, a visit followed by its reverse gives
the identity arrow in the groupoid and a two-step return in the path
description. The two descriptions have the same endpoints. The charge transport,
return probability and momentum suppression must be computed from the
path and its state before replacing the path by its composite arrow.

\begin{assumption}[Postulate \texorpdfstring{$\Omega_\alpha$}{Omega-alpha}.
Unified Path-Enriched Quantum Source Principle]
\label{ass:uvpe-source}
A single geometric rule assigns a bulk and boundary quantum source to
admitted raw fields, the kinematical category of histories, and a physical
preparation. It fixes the represented superconnection and its domains,
the curvature pairing, the cyclic trace on histories, the BV measure
and the BFV boundary state. Gluing two regions uses the boundary
pairing of this assignment. Probes couple to this source before the
prescribed integrations and differentiations. In the single-parameter
proposal the rule fixes the constitutive couplings up to the
normalization in Newton's constant.
For the time-local electromagnetic branch, the equal-volume sampling
law, harmonic connection and charged boundary representation in the
constitutive amendment below are part of this assignment.
This specifies the onsite realization. The unrestricted physical
source of the full ambient monomial must also pass the translation-spectrum
test and the direct photon-sewing test of
Section~\ref{sec:carrier-sewing-uv} with its own current and
boundary measure.
\end{assumption}

\paragraph{Constitutive amendment to the $\Omega_\alpha$ postulate.}
We adopt the following electromagnetic specialization of the Unified
Path-Enriched Quantum Source Principle in this revision. Work on a
fixed locked sector plane with its given carrier and ordering branch.
In the time-local Abelian realization, let $s_\perp$ be proper normal
distance and let $L_*$ be a length supplied by the microscopic branch.
The three-dimensional sampling ball has volume proportional to
$s_\perp$ and radius $L_*$ at $s_\perp=L_*$. In that ball's length
coordinate $a$, the charged displacement accumulates covariance
$2a^2 I_3$. The normal connection is the harmonic extension of the
boundary connection and its conditional law is concentrated on that
extension. These provisions fix both the measured current along the
entire unreduced path and its variation at coincident insertions.

Choose the quartic normal transfer with its metric volume,
the complete Abelian Gauss constraint with covariant normal
derivatives, and the occupied onsite fermion reference. On a massive
paired multiplet the physical left and right internal fibers have a
charge-preserving unitary pairing compatible with their support
arrows. The four-arrow curvature product with the derivative of the
charged heat kernel generates the charged bilinear. All probes act
on the resulting quantum source before the constraint and boundary
integrations. The reference and the common trace-class boundary
preparation fix the vacuum denominator. In the basic case the normal electric connection is the horizontal
pullback of the measured boundary gauge, with no independent normal
one-form. The basic weight uses the ordinary boundary Maxwell action on the trivial bundle with the
harmonic mode prescribed. Distinct bulk or topological sectors
require their own measure in the postulate.

These clauses select one charged action by geometry, representation
and measure. None refers to a loop integral or its convergence.
They supply a physical Abelian Gauss quotient and a concrete chiral
mass pairing for this source. They also make quantitative predictions
against which the proposed recovery of Q1 can be tested. The spectral
test below is a reason to keep this constitutive choice distinct
from an identification of the entire relativistic theory.

Here $\alpha$ labels the proposed principle. It is independent of the
radial heat-lapse exponent, the Schwinger parameters $\alpha_i$, and any channel label on a scattering
wave operator. State choices and external preparations must be specified.
Postulate~$\Omega_\alpha$ proposes one origin for the five operations
below. To replace the five postulates, either the independent rule must
select the physical realization or the recovery statements must hold
for all admitted realizations on the claimed domain.

\begin{description}[leftmargin=1.4em,style=nextline]
\item[A4. Classical image.] Restrict the intrinsic action and its
variations to the composite affine metric and two-form image with the
GTL and shell measures. Equality of actions on stationary values alone
would miss the variations which generate vertices.
\item[A5. Gauge functional.] Normalize and differentiate the
gauge-fixed quantum source with its BRST differential, physical probes
and directed boundary values. Integrating out a field carries its
determinant into the surviving functional.
\item[Q1. Physical representation.] Reconstruct the retained correlation
functions on the physical quotient with the required positivity,
domains and translation generators. A5 and Q1 are the functional and
reconstruction clauses of the earlier quantization package.
\item[Fixed-plane $\Omega$.] Recover the positive heat-lapse law of a
covered uncut internal line on a fixed support, including its
geometric depth assignment. External attachments and cuts use the
corresponding on-shell prescription.
\item[Groupoid $\Omega_{\mathfrak G}$.] In a represented family of
supports recover the ordered internal covariance
$R^{-1/2}\Phi_a(K)R^{-1/2}$. Here $K\ge0$ is the heat generator
and $R>0$ is the boundary precision on the indicated domains.
The entries are composed with their support labels. A matched
reducing fibre gives the fixed-plane rule.
\end{description}

Section~\ref{sec:pe-actions} explains what must be calculated
even after one has recovered their classical low-energy forms. The
UV assertion concerns each complete graph of a specified bare
expansion after its stipulated geometric integrations. Two sources
which agree on every free propagator can give different
electromagnetic determinants. Adding a pair of opposite-statistics
auxiliaries with exactly reciprocal determinants also illustrates why
the physical expansion and the order of exact eliminations belong in
the assertion. Their product is one at finite regulator. Treating the
two determinants as separate physical loops would define a different
graphwise test.

The single-parameter claim concerns physical inputs. Once the geometry,
$G_N$ and the external preparation are specified, the proposed rule
must tell us which charged operator and which boundary pairing to
differentiate. A choice of representation, source-dependent density
or inner product can change a measured response without introducing
a tunable number. The return test below distinguishes two such
responses before the question of higher orders arises. A successful
derivation must retain the geometry that selects the same response in
the open and closed calculations.

\subsection{A calculable groupoid covariance and its domain of validity}

For a finite family of identified Euclidean supports, let $K,R$ be
strictly positive self-adjoint matrices on the same real mode space.
Use dimensionless coordinates and let $Z(r)$ have positive quadratic
action
\[
 \tfrac12\langle Z(0),K Z(0)\rangle+
 \tfrac12\int_0^\infty
  (\|Z'(r)\|^2+\|K Z(r)\|^2)\,\dd r .
\]
The harmonic extension of a given $z=Z(0)$ is $e^{-rK}z$.
Its boundary energy adds precision $K$ to the bare precision $K$,
so the endpoint covariance is $(2K)^{-1}$. The Dirichlet fluctuation
vanishes at the endpoint and has covariance
$(2K)^{-1}(e^{-|r-s|K}-e^{-(r+s)K})$. This is the reflected
half-line Green kernel, obtained in each eigenspace of $K$ by
the derivative jump and the zero boundary value. The independent
harmonic extension contributes $(2K)^{-1}e^{-(r+s)K}$, which cancels
the reflected term. The total covariance is $(2K)^{-1}e^{-|r-s|K}$. Set
$U(r)=R^{-1/2}(2K)^{1/2}Z(r)$ and average $U$ uniformly over a
visited normal interval of length $a>0$. Its endpoint and visit
covariances are
\begin{align}
 C_{\rm end}&=R^{-1},\nonumber\\
 C_{\rm visit}&= R^{-1/2}\Phi_a(K)R^{-1/2},\nonumber\\
 \Phi_a(z)&=\int_0^a e^{-tz}\frac{2}{a}(1-t/a)\,\dd t .
 \label{eq:uvpeP-groupoid}
\end{align}
The average of the two-point kernel samples the distance between
two independent uniform positions in the same interval. That distance
has the triangular density in \eqref{eq:uvpeP-groupoid}. The factors
$(2K)^{1/2}$ cancel $(2K)^{-1}$ inside the covariance, between the two
$R^{-1/2}$ factors. This explains their displayed order even when
$R$ and $K$ fail to commute. A common reducing scalar fibre with
eigenvalues $r_0,k_0$ gives $r_0^{-1}\Phi_a(k_0)$. The action, physical boundary
pairing and choice of visit observable have all entered this result.
An infinite support space additionally requires domains, trace or
summability estimates. Matching this finite covariance neither fixes
higher charged source derivatives nor proves Q1.

\subsection{When summing over internal paths preserves a UV estimate}

An infinite set of geometric histories can occur already at a fixed
electromagnetic order if a propagator has been defined by an exact
sum over its paths. Here is one criterion for keeping that sum in a
fixed-graph estimate. Take a finite or countable support set. Let
$T_{ab}(\xi)\ge0$ majorize the sum of norms of one-step weights for
all allowed directed steps $a\to b$, at fixed shared geometric data
$\xi$. Assume there are $w_a(\xi)>0$ and $0\le\theta(\xi)<1$ such
that
\begin{equation}
  \sum_b T_{ab}(\xi)w_b(\xi)\le\theta(\xi)w_a(\xi).
  \label{eq:uvpeP-word}
\end{equation}
By positivity and induction,
$(T^n)_{ab}\le (w_a/w_b)\theta^n$. Thus the absolute sum over
paths of arbitrary length joining fixed $a,b$, with the identity
allowed, is bounded by
$(w_a/w_b)/(1-\theta)$. The same bound for a closed line with
specified basepoint uses a ratio of one. If the basepoint itself is
traced, its initial measure or multiplicity must also be summable.

More precisely, suppose the integrand of each fully differentiated
graph contribution is bounded by one common momentum envelope
$E_G(k,\xi)$ times products of such step weights. Fix its junction
labels first and let $M$ be the number of path sums. Summing the
absolute bounds and dropping compatibility constraints only increases
them. The cost is at most $(1-\theta)^{-M}$ times the product of
endpoint weight ratios. Any remaining junction sums, multiplicities
and measures are taken on that bound. If the result integrated over
$k,\xi$ is finite and uniform in the regulator under a common
integrable majorant, Tonelli and dominated convergence give absolute
finiteness and cutoff removal, provided the regulated terms converge
pointwise. A uniform absolute tail together with local $L^1$
convergence gives the same conclusion. The step weights, including
those at marked vertices, must majorize the actual source expansions.
The empty identity path pays no travel cost. A positive-length return
with identity composite retains its prescribed path weight.

For example, with one support, $b$ distinguishable return steps
and norm $u$ per step, the sum of absolute weights of all words is
$\sum_{n\ge0}(bu)^n$. It converges exactly for $bu<1$, even if every
word at fixed length has its own finite momentum integral with the
same normalized envelope. Finite multiplicities of parallel histories
therefore matter. For a shared small depth, the factor
$(1-\theta(\xi))^{-M}$ must be integrated together with the
momentum cost of that depth. This criterion addresses a second
issue beyond four-dimensional power counting and requires no
independence of the depths of different segments.

\section{Raw O--QED and the electromagnetic sector induced by the GCM action}
\label{sec:pe-actions}

The raw O--QED action is a specified Maxwell--Dirac reference action
on locked supports. The raw fields of the graded groupoid
Cartan--matrix (GCM) action are the microscopic variables from which
a four-dimensional electromagnetic theory is to be obtained. The
presence of a Maxwell term and an electrically charged spinor in a
low-momentum expansion leaves open whether the two quantum sources
agree. Their vacuum loops, contact vertices and measures can differ
already at second order in an external photon.

\subsection{The reference action and the prescribed fixed-depth test}
\label{subsec:uvpeA-raw}
On a fixed post-GTL support $\bb=(s,b_s)$ the reference density
$\mathcal L_{{\rm ref},\bb}$ is given in
\eqref{eq:OQED-source-matter}--\eqref{eq:OQED-source-action}.
It contains minimal electromagnetic coupling, one Maxwell
term, and the Abelian gauge-fixing and ghost terms.
The starred spinor slots are componentwise conjugates with the
independent variations specified there. They count one physical
Dirac species in the one-species measure. The directed assembly is
\begin{equation}
 \vec S^{\rm lock}_{{\rm raw},\epsilon}=
 \sum_s\int\dd\mu_{\bb_s}\,
    \sigma_s\sigma_\epsilon\hat p_{0,\bb_s}
    Z_{\rm GT}\mathcal L_{{\rm ref},\bb_s}.
 \label{eq:uvpeA-raw}
\end{equation}
Here $\bb_s=(s,b_s)$ and $Z_{\rm GT}=\ell_{\rm GT}^{K-1}$ occurs
once. The shell measure and the cut-state measure keep their own
normalizations. The orientation sign is carried by the action. On
each copy the boundary displacement is $\ii\eta\epsvec$, with
$\eta=\sigma_s\sigma_\epsilon$ and
$\epsvec=\varepsilon\hat p_{0,\bb}$ as in \eqref{eq:signs}.
The following Euclidean coefficients can be evaluated on either
copy separately.

After the canonical rescaling \eqref{eq:fw-oqed-rescaling}, one
vectorlike Dirac species in a fixed Abelian test potential $A$
has operator $D_A=\gamma_E^\mu(\partial_\mu+\ii e A_\mu)+m_e$.
In this canonical Euclidean normalization the local action reads
\begin{equation}
 S_{\rm raw}^{E}=\int\dd^4x\,
  \bigl[\tfrac14 F_{\mu\nu}F_{\mu\nu}
       +\bar\psi D_A\psi+\mathcal L_{\rm gf+gh}^{E}\bigr].
 \label{eq:uvpeP-raw-euclidean}
\end{equation}
The directed support conventions precede this Euclidean comparison.
With the ordinary fermion measure the charged vacuum factor is $\det D_A$.
Quantizing this raw action with its local inverse introduces no
positive normal length. The standalone fixed-depth O--QED
prescription of Ref.~\cite{ChiapponeIgnition} uses the same local
density as a reference and supplies an additional internal-line
rule. On a covered free coefficient it prescribes
\begin{equation}
 C^{\rm test}_{a,M}(k)
   =e^{-a(\sqrt{k_E^2+M^2}-M)}\,C_0(k),\qquad a>0 .
 \label{eq:uvpeA-test}
\end{equation}
Its moment estimates are proved in
Proposition~\ref{prop:OQED-moment-certificate}. The value of
\eqref{eq:uvpeA-test} at zero background leaves the higher photon
derivatives undetermined. The covariant anticommutator in
\eqref{eq:polish-covariant-dirac} is one explicit choice of
those derivatives. Its Ward and contact tests belong to that trial.

\subsection{A bare vacuum response that a transferred probe cannot change}
\label{subsec:uvpeA-log}
On a flat Euclidean copy take $F_{12}=B$ constant and $m_e>0$.
Write $\Gamma_f=-\log|\det D_A|$ and compare its density with
$B=0$ using a common proper-time interval $\delta<t<T$.
The upper endpoint is fixed. The lower endpoint $\delta$ is a
diagnostic cutoff of units inverse mass squared. For one species,
\begin{equation}
 \frac{\Delta\Gamma_f^{\delta,T}}{\mathrm{Vol}}
 =\frac{2}{(4\pi)^2}\int_\delta^T
 \frac{e^{-t m_e^2}}{t^3}\,[x\coth x-1]\,\dd t,
 \qquad x=eBt .
 \label{eq:uvpeA-dirac-heat}
\end{equation}
Indeed, the scalar magnetic oscillator in the $1,2$ plane
has eigenvalues $(2n+1)|eB|$, $n\ge0$, and degeneracy
$|eB|/(2\pi)$ per unit area. Integrating the other two free
momenta gives $(4\pi t)^{-1}$. The two spin shifts contribute
a trace $4\cosh x$. Together these give the diagonal heat density
$4(4\pi t)^{-2}x\coth x$ before including the mass. Since
$-\log|\det D_A|=-\tfrac12\log\det(D_A^\dagger D_A)$, proper-time
integration has the positive prefactor in
\eqref{eq:uvpeA-dirac-heat}~\cite{Schwinger1951}.
The comparison removes only the field-independent vacuum density.

Using $x\coth x=1+x^2/3+O(x^4)$ gives
\begin{equation}
 [B^2]\,\frac{\Delta\Gamma_f^{\delta,T}}{\mathrm{Vol}}
   =\frac{e^2}{24\pi^2}\int_\delta^T e^{-t m_e^2}\frac{\dd t}{t}.
 \label{eq:uvpeA-raw-log}
\end{equation}
Thus the local $F_{\mu\nu}F_{\mu\nu}=2B^2$ coefficient grows as
$e^2\log(T/\delta)/(48\pi^2)$ for $e\ne0$.
It is the unrenormalized one-spinor response at fixed support.
The Abelian gauge and ghost Hessians of this raw action are
independent of this background. Adjusting a bare Maxwell
coefficient to absorb the logarithm would change the UV objective
of this paper. A positive length assigned elsewhere in the
calculation has to enter this response to affect its UV limit.

For example, couple external Grassmann probes $\bar\xi,\xi$ to
a normal readout $T_A\psi$, with the dual map on the barred
slot. On a finite regulator with invertible $D_A$, Berezin
integration of the charged Gaussian gives
\begin{equation}
 \frac{Z[A,\bar\xi,\xi]}{Z[0,0,0]}=
 \frac{\det D_A}{\det D_0}
 \exp\!\left(\bar\xi\,T_A D_A^{-1}T_A^\dagger\xi\right).
 \label{eq:uvpeA-probe}
\end{equation}
Here $D_0$ denotes $D_A$ at vanishing $A$. A transfer such as
$T_A=e^{-rH_A}$ may have arbitrarily rapid decay and may itself
depend on $A$. Differentiating it generates genuine mixed
probe--photon contacts. Setting $\xi=\bar\xi=0$ removes all those
contacts and leaves the same determinant as in
\eqref{eq:uvpeA-raw-log}. The closed charged sector retains its
logarithm even when the open two-point readout has the desired
forward geo-phase.

A sharper test allows changes of the charged operator.
Suppose at each regulator that $K_A=D_A+W_A$ and
$X_A=D_A^{-1}W_A$. Suppose $I+X_A$ has a uniformly bounded inverse
$R_A$ at the test background. For two fixed smooth photon variations
$i,j$, assume $X_i,X_j,X_{ij}$ have regulator-uniform trace
norms there. The finite matrix identity for the relative response is
\begin{equation}
 \partial_i\partial_j\log\frac{\det K_A}{\det D_A}
   =\operatorname{Tr}\bigl(R_A X_{ij}
                          -R_A X_j R_A X_i\bigr).
 \label{eq:uvpeA-relative}
\end{equation}
It follows by differentiating $\det K_A=\det D_A\det(I+X_A)$
on any local nonvanishing chart. Trace-ideal bounds give
$\|R_A\|\|X_{ij}\|_1+
\|R_A\|^2\|X_i\|_1\|X_j\|_1$ as a uniform upper bound for its
absolute value. For variations having a nonzero bare logarithm
this relative response is too small to remove it under those
hypotheses. This excludes a class of smoothing modifications to
the closed operator. A changed principal operator or a
source-dependent quantum measure requires a new calculation.

\subsection{Testing a positive charged spectrum in a magnetic field}
\label{subsec:physical-spectral-magnetic-uv}

A physical pole constrains the free charged two-point function. To see
what it implies for closed histories one must also specify how photons
enter the charged operator and its measure. On a flat, locked Euclidean
four-plane, consider the following two covariantisations with fixed
charged measure and nonzero charge \(e\). A constant magnetic field
\(F_{12}=B\) probes their determinants per unit four-volume. This test
includes all contacts of the chosen operator and compares only the
field-dependent part. No Lorentzian sector boundary value is changed.

\paragraph{Scalar charged Gaussian.}
Let \(L_A=-(\partial+ieA)^2\), and choose the inverse propagator
\(K_A=C(L_A)^{-1}\) for a complex scalar, with
\[
 C(x)=c+\int_{[0,\infty)}\frac{d\rho(s)}{x+s},\qquad
 c\ge0,\quad 0<C(0)<\infty .
\]
Assume \(\rho\) is positive and nonzero. This includes a positive
massive pole. Set \(a(x)=\log[C(0)/C(x)]\). The reciprocal of a
nonzero Stieltjes function is a complete Bernstein function.
Consequently \(C(0)/C(x)-1\) is Bernstein, as is its composition with
\(y\mapsto\log(1+y)\)~\cite{SchillingBernstein2012}. Some bounded
mass interval \([0,s_*]\) has weight \(Z>0\), so
\(C(x)\ge Z/(x+s_*)\). Thus \(a(0)=0\) and \(a(x)/x\to0\),
which remove the constant and drift from its Bernstein representation:
\begin{equation}
 \begin{aligned}
 a(x)&=\int_{(0,\infty)}(1-e^{-tx})\,d\nu(t),\\
 \nu((0,\infty))&=\log[C(0)/c],
 \end{aligned}
 \label{eq:physical-scalar-bernstein-uv}
\end{equation}
where the right-hand side is infinite when \(c=0\). The second identity
uses monotone convergence and \(C(x)\to c\).
For the complex scalar the effective action has sign \(+\log\det K_A\).
The magnetic scalar heat density is
\((4\pi t)^{-2}y/\sinh y\), with \(y=eBt\). The positive relative
integral of the Bernstein heat densities is therefore
\begin{equation}
 \frac{\Gamma_s(B)-\Gamma_s(0)}{\mathrm{Vol}}
 =\int_{(0,\infty)}\frac{1-y/\sinh y}{(4\pi t)^2}\,d\nu(t).
 \label{eq:physical-scalar-magnetic-uv}
\end{equation}
This equality defines the relative trace by exhausted positive
heat integrals, allowing the value \(+\infty\). At any fixed
\(eB\ne0\) the integrand tends to \((eB)^2/(96\pi^2)\) at the
short-time endpoint. Since \(\nu([\epsilon,\infty))<\infty\) for
each \(\epsilon>0\), the density diverges for \(c=0\). For \(c>0\)
it is finite and its exact \(B^2\) coefficient is
\(e^2\log[C(0)/c]/(96\pi^2)\), by dominated convergence using
\(0\le1-y/\sinh y\le C_1 y^2\).
Even in this latter case, a flat connection translates \(k\) to
\(k+eA\). At \(A=0\) the Hessian of the scalar log integrand in
any unit direction perpendicular to \(k\) equals
\(2e^2a'(|k|^2)\). Since
\(a'(x)\ge Z/[C(0)(x+s_*)^2]\), its tensor norm has a
logarithmically nonintegrable tail in four dimensions. This last
test concerns absolute integrability at zero photon transfer in the
specified routing. The magnetic response probes a field with nonzero
curvature.

\paragraph{Spinor charged Gaussian.}
Choose real spectral data on \([0,\infty)\) with
\(d\rho_V\ge0\), \(0<M=\int d\rho_V<\infty\), and
\(d|\rho_S|(s)\le\sqrt{s}\,d\rho_V(s)\). Define
\[
 V(x)=\int\frac{d\rho_V(s)}{x+s},\qquad
 U(x)=\int\frac{d\rho_S(s)}{x+s}.
\]
These unsubtracted propagators allow an elementary pole and a
continuum. Finite vector spectral mass is an additional hypothesis.
For the anti-self-adjoint Euclidean Dirac operator
\(D_A=\gamma^\mu(\partial_\mu+ieA_\mu)\) set
\[
 H_A=-D_A^2,\qquad
 S_A=U(H_A)-D_A V(H_A),\qquad K_A=S_A^{-1}.
\]
Use an invertible infrared sector; all estimates below concern
\(x\to\infty\). The identity \(S_A^\dagger S_A=g(H_A)\),
where \(g(x)=xV(x)^2+U(x)^2\), fixes the field-dependent real
fermionic log weight. With \(x_*>0\) and
\(f(x)=-\log[g(x)/g(x_*)]\), it is
\(-\tfrac12\operatorname{Tr}f(H_A)\) on finite spectral blocks,
up to the common reference constant.
For \(0\le j\le2\), differentiation under the spectral integrals gives
\[
 \begin{aligned}
 x^{j+1}V^{(j)}(x)&\longrightarrow(-1)^j j!M,\\
 x^{j+1/2}U^{(j)}(x)&\longrightarrow0.
 \end{aligned}
\]
For the second limit use the bounded kernel
\(\sqrt{s/x}(1+s/x)^{-j-1}\) against the finite measure
\(d\rho_V\). It tends pointwise to zero. Substitution gives
\(g\sim M^2/x\), \(g'\sim-M^2/x^2\), \(g''\sim2M^2/x^3\).
In particular \(f'\sim1/x\) and \(f''\sim-1/x^2\).

At zero photon transfer the action seagull and the two
one-photon insertions together differentiate the translated symbol
\(\log|\det K_0(k+eA)|=2f(|k+eA|^2)+\mathrm{const}\).
The tangential Hessian is \(4e^2 f'(|k|^2)\). On
\(r_0<|k|<R\), the unnormalised integral of the momentum
Laplacian of \(2f(|k|^2)\) is
\(8\pi^2[R^4 f'(R^2)-r_0^4 f'(r_0^2)]\sim8\pi^2R^2\).
Here \(r_0>0\) avoids an infrared assumption at the origin.
This translation test by itself says nothing about the integrated
physical response to a pure gauge.

The same calculation gives a necessary size for any compensating
boundary or operator term. Write $J=e^2\nabla^2 h$, where
$h(r)=2f(r^2)$ and $r=|k|$, and let $M$ collect the other two-mark
terms in that routing. On the shell $\mathcal A_R=\{R<|k|<2R\}$,
the operator norm of the angular Hessian is at least
$e^2|h'(r)|/r$. Thus
\[
 \int_{\mathcal A_R}\|J\|\,\mathrm d^4 k
 \ge e^2 |S^3| R^2 |h(2R)-h(R)|.
\]
Here $h(2R)-h(R)\to4\log 2$ by the free spectral asymptotic.
Absolute exterior integrability of $J+M$ would consequently require
\[
 \liminf_{R\to\infty} R^{-2}
   \int_{\mathcal A_R}\|M\|\,\mathrm d^4 k
 \ \ge\ 4 e^2 |S^3|\log 2.
\]
This follows by the triangle inequality before momentum integration.
The normalization here omits $(2\pi)^{-4}$. A separately integrable
measure variation cannot supply that cancellation in this continuum
constant-shift test.

One can also test the same momentum routing at fixed nonzero photon
transfer. This step uses gauge covariance alone once the free symbol
above has been fixed. Put \(A=\lambda\,d\chi\), with
\(\chi(y)=\cos(q\cdot y)\) and fixed \(q\ne0\). For any covariant
gauging of this free symbol, the operator \(L_A=\log|K_A|\) obeys
\(L_A=e^{-ie\lambda\chi}L_0e^{ie\lambda\chi}\); here
\(|K|=(K^\dagger K)^{1/2}\). Let
\(F(k)=\operatorname{tr}_{\rm spin} L_0(k)=2f(|k|^2)+\mathrm{const}\).
The diagonal in momentum of the second derivative, traced over spin
with the common volume factor removed, equals
\[
 \frac{e^2}{2}\,[F(k+q)+F(k-q)-2F(k)].
\]
Indeed, this is the diagonal of \(-e^2[\chi,[\chi,L_0]]\);
each of the two Fourier modes of \(\chi\) has coefficient \(1/2\).
Writing the second difference as the weighted integral of the
directional Hessian along \(k+tq\), \(-1\le t\le1\), gives, with
\(r=|k|\) and \(\theta=\angle(k,q)\),
\[
 \begin{aligned}
 &F(k+q)+F(k-q)-2F(k)\\
 &\qquad=\frac{4|q|^2}{r^2}(1-2\cos^2\theta)+o(r^{-2}).
 \end{aligned}
\]
The error is uniform in angle by the bounds on \(f'\) and \(f''\).
The cone \(|\cos\theta|\le1/2\) therefore already forces divergence
of the absolute momentum integral. This is a statement about the
diagonal of the covariant log modulus in this routing. A finite trace
is invariant under the pure-gauge conjugation, and the result places
no lower bound on a transverse observable.

A field with \(eB\ne0\) gives a second test. Write \(h=2|eB|\)
and \(u=p_3^2+p_4^2\). The Dirac Landau levels \(u+nh\) have
spin multiplicity two for \(n=0\) and four for \(n\ge1\).
Each transverse orbital level has degeneracy \(|eB|/(2\pi)\)
per area. Pair the endpoints of each Landau interval with the
zero-field continuum on that same interval. Its trace difference
uses the trapezoidal remainder
\[
 \begin{aligned}
 T_h(z)&=\frac h2[f(z)+f(z+h)]-\int_0^h f(z+t)\,dt\\
       &=\frac12\int_0^h t(h-t) f''(z+t)\,dt .
 \end{aligned}
\]
The common-cell magnetic contribution to the fermionic density per
\(du\) is
\begin{equation}
 -\frac{1}{8\pi^2}\sum_{n=0}^{\infty} T_h(u+nh)
 \ \sim\ \frac{(eB)^2}{24\pi^2 u}
 \qquad (u\to\infty).
 \label{eq:physical-spinor-magnetic-uv}
\end{equation}
Indeed, the transverse free spin density is \(1/\pi\) per unit
squared momentum, and the parallel measure is \(du/(4\pi)\).
For fixed \(u>0\) the cell sum converges absolutely. At large
\(u\), \(f''(z)\sim-z^{-2}\) uniformly on the tail,
\(\int_0^h t(h-t)\,dt/2=h^3/12\), and
\(\sum_{n\ge0}(u+nh)^{-2}\sim1/(hu)\). The result has the
ordinary positive logarithmic divergence already for one charged
loop. The cell comparison cancels the field-independent density
before the large-momentum limit and uses no Maxwell subtraction.

These computations test two specific actions and their own
photon contacts. A different transverse coupling, a source-dependent
boundary pairing, or a density obtained by retaining intermediate
displacements changes the calculation. To establish finiteness for
the physical GCM source, that change must be derived on the same
charged histories that contribute to the closed determinant.

\subsection{How the microscopic action defines an induced O--QED sector}
\label{subsec:uvpeA-native}
In PE-CGG-HOUQFT, O--QED denotes an electromagnetic sector of the
quantum GCM source once its charged and photon observables
have been identified. Write $q$ for the raw fields on a stationary
branch. Their superconnection has curvature $\mathbb F=\mathbb D^2$.
Equation~\eqref{eq:uvpeA-curvature} integrates the curvature
pairing at fixed raw fields, including its oriented velocity
argument and normal history. Coarea reduces each ambient event
integral to a locked event, as in \eqref{eq:update-labelled-coarea}.
The remaining quantum integration determines which intermediate
normal and charged states are paired.

A marked excursion may connect
locked endpoints through off-support normal positions.
The curvature pairing $\mathbb P_*$, graded
representation, domains and trace over composable histories enter
the definition of the integrand. A history traversed out and back
can return to the same endpoint while retaining a positive
unreduced length. Whether it contributes a decaying geo-phase
factor depends on the history insertion and the state with which
the itinerary is closed.

Fix a gauge-admissible chart and let $u$ collect the retained photon
and charged fields, and $z$ the integrated fields. Let $J$ denote
physical source probes. Specify the physical branch, the BV integration
cycle and the BFV boundary state. Transport source-dependent domains
to common reference domains. Write $\dd\nu_u$ and $\dd\nu_z$ for the
declared source-independent reference measures and pairings; put all
source dependence of their densities, Jacobians, constraints and
boundary states into the complete weight $w(u,z;J)$. The induced
functional is the pushforward followed by one vacuum normalization:
\begin{equation}
\begin{aligned}
 \mathcal W(u;J)&=\int\dd\nu_z(z)\,w(u,z;J),\\
 Z^{\rm ind}[J]
 &=\frac{\int\dd\nu_u(u)\,\mathcal W(u;J)}
        {\int\dd\nu_u(u)\,\mathcal W(u;0)} .
\end{aligned}
\label{eq:uvpeA-induced}
\end{equation}
The fixed Berezin and boundary pairings are included in this notation;
ordinary integrations and estimates for odd probes are interpreted
coefficientwise in each finite Grassmann expansion. With several
support charts, the required intermediate labels are summed or
integrated. The displayed formula defines a continuum functional only
when its stated integrals and pairings exist and the even vacuum
normalization has invertible body. An order-by-order construction
instead requires each specified coefficient of the complete bare source
to exist; it does not by itself sum the perturbation series. A common
finite-dimensional approximation can check the algebra, but its
removal is not the definition of UV finiteness used here.

Where $\mathcal W$ is even with invertible body, its negative logarithm
relative to the declared retained reference measure defines an induced
Euclidean action after a local logarithm of the body has been chosen.
The logarithm of its nilpotent part is a finite algebraic expansion.
One may differentiate $\mathcal W$ directly without taking a logarithm
of the integrand $w$. The 1PI functional comes from the connected
generating functional through the declared Legendre transform.
Matching the Maxwell normalization, pole mass, physical charge and
directed source maps to the raw sector is a separate calculation on
the chosen branch.

Even at tree level, holding the eliminated fields fixed can give the
wrong photon kernel. On a fixed gauge-reduced slice, if those
fields solve $S_z(u,z_*(u))=0$ with invertible $S_{zz}$, then
\begin{equation}
 \frac{\dd^2}{\dd u^2}S(u,z_*(u))
     =S_{uu}-S_{uz}S_{zz}^{-1}S_{zu}.
 \label{eq:uvpeA-schur}
\end{equation}
All blocks here are evaluated on the stationary branch. This formula
is written for the commuting photon probes and eliminated bosonic
coordinates. It follows by differentiating stationarity before
differentiating the reduced action. The quantum pushforward in
\eqref{eq:uvpeA-induced} also includes fluctuations about the branch.
For commuting source directions $i,j$, differentiate the complete
weight on the common reference domain. If its first and second
derivatives admit the coefficientwise integrable majorants needed to
differentiate under the integral, then
\[
\begin{aligned}
 \mathcal W_i&=\int\dd\nu_z\,w_i,\\
 \mathcal W_{ij}&=\int\dd\nu_z\,w_{ij}.
\end{aligned}
\]
Consequently,
\[
 \partial_{ij}[-\log\mathcal W]
   =-\mathcal W^{-1}\mathcal W_{ij}
     +\mathcal W^{-2}\mathcal W_i\mathcal W_j .
\]
Here $\mathcal W$ is even with invertible body on the source chart.
This grouped derivative formula does not require $w$ to have a
pointwise logarithm or require its individual score moments to exist.

If, in addition, $w$ is even with invertible body almost everywhere on
the integration chart, choose a local logarithm $w=e^{-Q}$ and put
$\langle F\rangle=\mathcal W^{-1}\int\dd\nu_z\,wF$.
When the separate terms on the right are integrable coefficientwise,
the same identity becomes
\begin{equation}
\begin{aligned}
 \partial_{ij}[-\log\mathcal W]
    ={}&\langle Q_{ij}\rangle-\langle Q_iQ_j\rangle\\
       &+\langle Q_i\rangle\langle Q_j\rangle .
\end{aligned}
 \label{eq:uvpeA-hessian}
\end{equation}
The first term is the mean direct source variation of order two.
The last two terms form the negative connected pair of first
variations. These brackets may be complex and need not define a
positive probability law. Both $Q_i$ and $Q_{ij}$ include variations
of source-dependent densities, transported constraints and boundary
states. Integrability of the grouped expression
$w_{ij}=w(Q_iQ_j-Q_{ij})$ alone does not justify splitting it into
two separately integrated score terms. Odd source derivatives require
their declared Berezin ordering and signs; the displayed covariance
identity is for commuting directions.

A source-dependent division by the zero-visit determinant would
change the normalized functional. For one invertible finite Gaussian
fermion block, the corresponding second variation of $\log\det D$ is
\[
 \operatorname{Tr}(D^{-1}D_{ij}-D^{-1}D_iD^{-1}D_j).
\]
The same identity applies to a differentiable ordinary Fredholm
determinant on an invertible determinant-class chart when the displayed
products are trace class and differentiation under the trace is
justified. Invertibility alone does not supply these analytic premises.
The current pair and the direct contact are derivatives of the same
source and enter its complete response before its absolute
integrability is tested. Higher derivatives of that source determine
its multi-photon contacts.

\paragraph{Closed and open charged coefficients at determinant zeros}
The original Berezin integral also gives a bound where a charged
inverse does not exist. On a fixed finite-dimensional chart, write
the complete closed weight as $b(A,z)\det K(A,z)$. The scalar $b$
contains the paired boundary half-densities and all remaining action,
constraint and measure factors. With $q$ prescribed open Grassmann
pairs in distinct fixed basis slots, the corresponding weight is,
up to its Berezin sign, $b\det K'$, where $K'$ is the complementary
square minor of $K$. This follows by saturating the indicated slots
in the original integral. The empty minor has determinant one.
The same formula covers the closed case by taking $K'=K$.

Let $N$ be the size of $K'$, $s(K')$ its singular values, and $e_j$
their elementary symmetric polynomial with $e_0=1$. Label the
commuting source differentiations by a finite set $I$. For
$J\subseteq I$ let $\Pi(J)$ be its partitions, with one empty
partition at $J=\varnothing$. Then
\[
\begin{aligned}
 |\partial_I(b\det K')|&\le
  \sum_{J\subseteq I}|\partial_{I\setminus J}b|\,M_J,\\
 M_J&=\sum_{\substack{\pi\in\Pi(J)\\h=|\pi|\le N}}
 h!\,e_{N-h}(s(K'))\prod_{B\in\pi}\|\partial_B K'\|.
\end{aligned}
\]
To see the determinant bound, use singular-vector bases at the
evaluation point and keep them fixed during differentiation. For each
partition choose the $h$ varied columns. The unchanged diagonal
entries multiply a complementary $h$-dimensional determinant.
Hadamard's bound controls it by the product of the insertion norms
for each of the $h!$ assignments. Summing over columns gives
$e_{N-h}$. The one-block partition includes the direct contact of
every order. Differentiating the whole weight also retains mixed
matrix and boundary derivatives. Absolute integration still
requires an integrable bound uniform in the relevant hard regions.
In particular, this polynomial estimate through zeros gives no
dimension-independent constant by itself.

For an ordinary trace-class charged perturbation there is also a bound
independent of the dimension of the one-particle space. Write
\(K=I+T\), and fix \(q\) incoming and outgoing wavefunctions
\(f_i,g_i\) in its complex Hilbert space. Put
\(R_i=f_i\langle g_i,\cdot\rangle\), with the inner product linear
in the second slot. The paired open coefficient, with the corresponding
ordering of external Grassmann variables, is
\[
 M_q(T)=[w_1\cdots w_q]\det\!\left(I+T+\sum_i w_i R_i\right).
\]
When \(K\) is invertible, the finite-rank determinant identity gives
\(M_q=\det K\,\det[\langle g_i,K^{-1} f_j\rangle]_{i,j=1}^q\).
The coefficient formula extends continuously through zeros of \(\det K\).
Indeed it agrees with the Berezin minors on finite matrices, and
finite-rank compressions of \(T\) converge in trace norm along
increasing orthogonal projections tending strongly to the identity.
The projected wavefunctions converge in Hilbert norm.

Suppose the commuting source derivatives \(T_B=\partial_B T\) with
\(\varnothing\ne B\subseteq J\) exist continuously in trace norm
on the source chart. For the partitions \(\pi\) of the finite
label set \(J\), with a single empty partition if \(J=\varnothing\),
\[
 \begin{aligned}
 |\partial_J M_q(T)| &\le
 e^{\|T\|_1+q}\prod_{i=1}^q\|f_i\|\,\|g_i\|\\
 &\quad\times\sum_{\pi}e^{|\pi|}
                  \prod_{B\in\pi}\|T_B\|_1 .
 \end{aligned}
\]
For each partition apply the Cauchy coefficient bound to
the Fredholm determinant with one complex variable per rank-one
insertion and per block \(T_B\). Choose each radius as the reciprocal
trace norm of its insertion. The determinant on that product of
circles is bounded by \(\exp(\|T\|_1+q+|\pi|)\), using the
ordinary trace-class determinant inequality~\cite{SimonUVpeDet}.
A zero insertion contributes zero. The mixed chain rule sums the
partitions, including the one in which all probes strike the same
vertex. For a scalar boundary and measure factor \(b\), apply the
product rule to \(bM_q\). Integrability of the resulting bound
against the actual joint law suffices for differentiation through
determinant zeros. Its trace norms must be estimated from the
chosen action over the full graph domain.

\paragraph{A charged-sector reduction with its determinant retained}
The two operations in \eqref{eq:uvpeA-induced} can be carried out
exactly on a quadratic graded-representation block. At a common
finite regulator fix a photon background $A$. In a decomposition
into retained and eliminated charged Grassmann slots let the
quadratic matrix be
\[
 \mathcal D_A=
 \begin{pmatrix} B_A&E_A\\ F_A&H_A\end{pmatrix},
 \qquad D_A^{\rm ret}=B_A-E_AH_A^{-1}F_A .
\]
Suppose $H_A$ and $D_A^{\rm ret}$ are invertible.
Integrating only the eliminated charged slots produces
$\det H_A\,\exp(-\bar\psi D_A^{\rm ret}\psi)$.
Indeed independent triangular shifts of the two eliminated
Grassmann variables complete the square and have unit Berezinian.
The retained open line and the closed functional consequently
satisfy
\begin{equation}
\begin{aligned}
 (\mathcal D_A^{-1})_{\rm ret,ret}&=(D_A^{\rm ret})^{-1},\\
 \det\mathcal D_A&=(\det H_A)\det D_A^{\rm ret}.
\end{aligned}
\label{eq:uvpeA-fermionic-schur}
\end{equation}
This elementary reduction supplies a direct test of a proposed
GCM derivation of O--QED. For invertible target $K_A$, choose any
compatible matrices $U_A,V_A$ and set
\[
 B_A=K_A+U_AH_A V_A,\qquad
 E_A=U_AH_A,\qquad F_A=H_A V_A .
\]
The retained inverse equals $K_A^{-1}$ for every background.
Nevertheless the closed photon response also contains the source
derivatives of $\log\det H_A$. In one realization $H_A$ can be a
fixed matrix. In another it can be the regulated raw charged
Dirac operator, in which case it contributes
\eqref{eq:uvpeA-raw-log}. Even matching the entire open two-point
kernel as a function of external background fails to determine the
closed loops. Which eliminated slots and statistics actually occur
must follow from the common geometric representation. This example
makes no assertion that the displayed extra charged species is
present in the physical branch.

\paragraph{Which charged loops survive elimination of the visit slots?}
On one locked Euclidean four-torus, separate the finitely many charged
slots in a proposed first-order action. Factor the actual free inverse
on each slot into its interaction vertices. This gives a dimensionless
block operator $I-U$ on their Hilbert direct sum, with $U_{ij}$ the sum
of the arrows $j\to i$. Use the sector and the common radial position
of the original source throughout this calculation. A closed word in
the arrows is a charged determinant loop. Some such words can stay
entirely within slots eliminated from the open probe.

\begin{proposition}[A cycle test for the full charged determinant]
\label{prop:visit-cycle-uv}
At a fixed background, give each arrow $A_b$ a weight $0\le w_b\le1$.
For $w_b>0$ require $A_b$ to belong to the Schatten class
$\mathfrak S_{1/w_b}$ between its two slots. Weight zero requires a
bounded operator. Suppose that the sum of weights on each directed
simple cycle is at least one, counting a self-arrow as a cycle.
On the chart $\sum_b\|A_b\|<1$, including a neighborhood of zero for
interactions vanishing there, the relative charged Gaussian has an
ordinary Fredholm determinant on each successive diagonal Schur
pivot. Their product is independent of the order of elimination.
The open coefficient uses the corresponding block of $(I-U)^{-1}$
and the free source factors.
\end{proposition}
\begin{proof}
Eliminating a slot $v$ uses $P_v=I-U_{vv}$ and the update
\[
 U'_{ij}=U_{ij}+U_{iv}P_v^{-1}U_{vj}.
\]
Every principal compression of the original $I-U$ is invertible
on the stated chart by its Neumann series. Each successive
$P_v^{-1}$ is a block of such a principal inverse and is bounded.
Treat these inverses as bounded insertions. Induction on the updates
writes every diagonal term $U_{vv}$ at its elimination as a finite
sum of nonempty closed arrow paths with such insertions. Each path
contains a simple cycle. One can lower some positive weights on
its factors until their sum is exactly one. The Schatten inclusions
and product inequality then make the entire path trace class~\cite{SimonUVpeDet}.
This proves that $\det_{\rm F}P_v$ exists, including for a slot not
seen directly by any open probe.

For verification, take finite-rank orthogonal projections tending
strongly to the identity separately in each slot. Finite Berezin
elimination gives the product of these pivot determinants exactly.
The finite principal inverses, extended by the identity off the
projections, are uniformly bounded and converge strongly together
with their adjoints, by the same Neumann estimate. Every positively
weighted arrow converges in its Schatten norm under compression.
Approximating these latter factors by finite-rank operators shows
trace-norm convergence on every closed path, also across the intervening
strongly convergent bounded operators and their adjoints. Continuity
of the ordinary Fredholm determinant gives the claimed product and
its independence of the projections and elimination order. The
product itself defines the relative weight without a spectral
cutoff or a subtracted determinant. Any free determinant, constraint
factor or boundary density that varies with the geometry still
multiplies this weight before the common integration.
\end{proof}

In particular, give a trace-class displacement arrow weight one
and all remaining arrows weight zero. The test then asks exactly
whether deleting the displacement arrows leaves an acyclic directed
graph. In that case the remaining block $N$ can be ordered strictly
triangular. If there are $h$ slots, $(I-N)^{-1}=\sum_{j=0}^{h-1}N^j$.
Writing $T=U-N$, the complete relative determinant reduces to
$\det_{\rm F}(I-(I-N)^{-1}T)$. In finite sections the triangular
factor has determinant one for every source. Its finite path sum
remains inside the second factor, so that excursions into eliminated
slots and derivatives striking their vertices are retained.

There is a converse for a reciprocal pair on a positive charged
block. After normalizing its two free positive diagonal factors,
suppose the actual real test field enters only in opposite arrows
$tB_\omega$ and $tB_\omega^*$, with no direct quadratic contact.
Here $\omega$ denotes all common geometric variables; the test
is applied to the arrows after the physical source projection.
Let $\nu_0$ be the positive normalized measure including the entire
free vacuum, and suppose the boundary weight has no variation in this
test. The first derivative of the conditional determinant at zero
vanishes because the variation is off-diagonal. Schur elimination in
each finite-dimensional pair gives
$\det(I-t^2B_\omega^* B_\omega)$. Hence the bare quadratic
coefficient of the logarithm of the integrated partition is exactly
\[
              -\int \|B_\omega\|_2^2\,\mathrm d\nu_0(\omega),
\]
with $\|\cdot\|_2$ the Hilbert--Schmidt norm, allowing an infinite
absolute cost. Increasing coordinate projections exhaust the sum of
squared matrix elements monotonically, so Tonelli proves both
necessity and sufficiency for this coefficient. There is no covariance
of first derivatives to cancel it. If a native action instead has a
quadratic contact or a varying density, its derivatives must be restored
before applying a lower bound.

For instance, in one scalar slot let $B_r$ be a fixed Fourier shift
on the torus times $e^{-r|k|}/\sqrt{m^2+|k|^2}$, with $m>0$.
Then $\|B_r\|_2^2\sim V/(32\pi^2r^2)$ at small $r$ for torus
volume $V$. This follows by scaling $k$ to $rk$ in the squared
lattice sum; the integrable continuum limit has radial factor
$2\pi^2\int_0^\infty u e^{-2u}\,\mathrm du$. The zero mode
contributes a bounded term. Thus the actual vacuum-weighted
boundary law must have a finite inverse second moment on this
test. An atom at $r=0$ already makes it divergent.

\paragraph{From the cycle test to absolute graph sums.}
Trace-class membership at a fixed background alone does not bound
the sum of the absolute momentum integrands of its source derivatives.
Here is a sufficient test that can be applied directly to those
derivatives. Suppose the visit expansion has finitely many types and
contacts through each fixed electric order and all interaction
vertices have positive electric order. At a vertex of a charged loop
write $k$ for its circulating momentum. For each actual contact
(including its adjoining free inverse) suppose the absolute momentum
bound before the photon readouts is of one of the following forms.
A displacement vertex has a translated envelope with uniformly
bounded lattice $\ell^1$ norm. An ordinary vertex has an envelope
$(1+|k+c_b|)^{-s_b}$ with $s_b\ge0$. The translation $c_b$ may depend
on the other momenta and on compact interval parameters, independently
of $k$. Allow each bound to be multiplied by a fixed inverse power
of $r\in(0,1]$ and a polynomial in the photon and fixed probe momenta.
Uniform averages of these envelopes obey the same estimates.
The estimates apply also to open-chain vertices in supremum norm.
Suppose every directed simple cycle either contains a displacement
vertex or has total ordinary order strictly greater than four.
The matrix of allowed arrows and the orders used in this test must
cover all types of source contacts that can occur.

Let $\Lambda^*$ denote the reciprocal momentum lattice.
Fix all photon transfers first. On a closed word containing a
displacement vertex use its translated $\ell^1$ bound and supremum
bounds elsewhere. For an ordinary word, its total order
$S=\sum_b s_b$ exceeds four. Discrete H\"older with
exponents $S/s_b$ on positive-order factors gives
\[
 \sum_{k\in\Lambda^*}\prod_b(1+|k+c_b|)^{-s_b}
 \le \sup_c\sum_{k\in\Lambda^*}(1+|k+c|)^{-S}<\infty .
\]
Factors of order zero use their supremum. Lattice shell counting
makes the bound uniform in all translations, even when the marks on
different arrows are correlated. This argument uses one independent
momentum for every closed charged line. An open chain is fixed by
its incoming momentum and transfers. Suppose now that all photon
ends come from the same harmonic readout at $r$. Each internal
nonzero transverse photon then contributes
$e^{-2r|p|}/|p|^2$. The translated charged estimates leave only a
polynomial in these photon momenta. Dropping remaining lattice
constraints in the positive majorant and summing by shells costs
a further finite inverse power of $r$. For any chosen photon durations
restricted to $0<t_i<\delta_i$ one can instead use
$\int_0^{\delta_i} e^{-t_i|p|^2}\,\mathrm dt_i\le\delta_i$,
at the cost of two further powers in the corresponding shell bound.
These estimates hold for simultaneous short durations. Consequently,
if the full free vacuum times the sewn boundary weight has all the
required inverse radial moments in absolute value, every such bare
graph is UV finite at each finite order, including those within
eliminated slots. This conclusion allows a common law on all visits.
Any other moving source or measure factor must contribute its own
derivatives with the same bounds.

For a shorter ordinary cycle there is also a direct failure test.
Hold its photon transfers fixed and form the actual spin trace and all
contacts belonging to the coefficient under examination. Suppose this
integrand has a finite large-$|k|$ expansion through degree $-4$ in
homogeneous powers of $|k|$, with continuous angular coefficients and
a uniform remainder $O(|k|^{-4-\epsilon})$, $\epsilon>0$.
Its absolute charged sum converges exactly when every coefficient
of degree at least $-4$ vanishes identically. Indeed the remainder
is summable in four dimensions. If the first nonzero term has degree
$-s$, its absolute angular coefficient is bounded below on a small
open cone. In each sufficiently large dyadic shell that cone contains
at least $cR^4$ reciprocal-lattice points, by counting lattice cells in
a smaller cone. Their absolute contribution is bounded below by a
constant times $R^{4-s}$, so even the borderline $s=4$ fails.
A vanishing angular integral is weaker than this pointwise condition.
Bounds uniform in growing photon transfers are still needed for sewing.

For example, in a directed ring of five massive first-order Dirac
slots coupled by smooth bounded local multiplication operators, each
free inverse is in $\mathfrak S_5$ on the four-torus. A traversal
has total weight one, and at fixed Fourier contacts of local
multiplication type it has total ordinary order five. Elimination
gives the ordinary determinant of identity minus the ordered product
around the ring. This is a test on a specified directed operator.
Adjoint, gauge and chiral conditions on a physical representation
must be imposed on the actual arrows; adding reverse arrows can
create shorter cycles.

A hidden two-slot ring shows what a test restricted to visible
arrows can miss. On a fixed transverse chart take equal massive
Dirac kinetic operators, with an off-diagonal scalar multiplication
$e^2 M_u$ in both directions. In fixed units use
$u=F_{\mu\nu}F_{\mu\nu}$ for a real finite Fourier photon probe
of nonzero curvature; harmonic readout of the probe is allowed
before forming $u$. Adjoin the ring as a direct summand to any
retained charged operator. The retained inverse stays fixed.
The formal order-$e^4$ coefficient of the hidden log determinant is
$-\operatorname{Tr}(D_0^{-1}M_uD_0^{-1}M_u)$, as is seen already
in finite sections. For a fixed
transfer $q$ its spin trace before the overall minus sign is
\[
 |\widehat u(q)|^2\,
 \frac{4[m^2-k\cdot(k+q)]}
 {(m^2+|k|^2)(m^2+|k+q|^2)} .
\]
For each of the finitely many nonzero Fourier coefficients the
leading term is $-4|\widehat u(q)|^2/|k|^2$. The signs agree and the
four-dimensional sum diverges quadratically. The example concerns
a transverse comparison source. A protected visible charged line
therefore cannot determine the closed response without an account
of the directed cycles of the operator and the state actually sewn.

To identify the induced O--QED sector, one must derive its retained
fields, physical charge and current, propagators, constraints and
sewings from \eqref{eq:uvpeA-induced}, with its support, BRST and
physical source identities. To prove its bare UV finiteness, one
must bound every complete graph of that prescribed expansion after
the exact normal integrations, including vacuum and zero-visit
loops and the photon contractions of the resulting vertices.
Integrability at fixed external photon momenta need not survive the
additional loop momenta created by sewing. For coupled supports
the source must supply the complete joint law. Extending an Abelian
result to all of PE-CGG-HOUQFT further requires the interacting
non-Abelian, chiral, constraint and geometric sectors of the same
physical branch.

\subsection{A normal return that makes the charged magnetic coefficient finite}
\label{subsec:pe-dirac-mark}
The raw logarithm in \eqref{eq:uvpeA-raw-log} also gives a quantitative
test for a proposed GCM charged history. Consider the following
specified one-loop Euclidean normal extension, in dimensionless
reference units. Tensor the four-dimensional squared Dirac heat
operator with the free normal heat flow generated by $\partial_y^2$
on $\RR$. Use a photon background independent of $y$. Follow the
normal path from $0$ to a recorded midpoint $y$ and back to $0$,
with two intervals of heat time $t/2$. The actual midpoint insertion
has positive finite measure $\mu(\dd y)$. For a density $w(y)$
this is precisely the kernel at $(0,0)$ of
$e^{(t/2)\partial_y^2} M_w e^{(t/2)\partial_y^2}$.
Keep the probabilities of both legs of the return. The normal factor is
\begin{equation}
 N_\mu(t)=\frac{1}{2\pi t}\int_{\RR}
                        e^{-y^2/t}\,\mu(\dd y).
 \label{eq:uvpeP-Dirac-mark}
\end{equation}
At order $B^2$ the Dirac trace and this scalar kernel factorize.
A direct-sum four-dimensional zero-visit component of nonnegative
weight $c_0$ contributes the raw trace separately. The absolute
UV cost of the combined magnetic coefficient, apart from its
positive factor $e^2/(24\pi^2)$ at $e\ne0$, is
\[
 \int_0^T e^{-m_e^2t}\,[c_0+N_\mu(t)]\,\frac{\dd t}{t}.
\]
For $T>0$ fixed this is finite precisely when $c_0=0$,
$\mu(\{0\})=0$ and
$\int_{0<|y|\le1}|y|^{-2}\mu(\dd y)<\infty$. Indeed the mass
factor is bounded above and below by positive constants on this
interval. For $y\ne0$ Tonelli's theorem gives the exact inner integral
\[
 \int_0^T e^{-y^2/t}\,\frac{\dd t}{t^2}
                  =\frac{e^{-y^2/T}}{y^2}.
\]
A zero-displacement atom costs the divergent integral at $y=0$.
Thus a compulsory return can make this coefficient finite even when
its positive separations accumulate at zero. An identity component
of this specified positive trace still carries its entire logarithm.

\begin{proposition}[Finite normal kinetic energy suffices on a nodal midpoint state]
\label{prop:uvpeP-Dirac-Hardy}
In the marked sector take $\mu(\dd y)=|\phi(y)|^2\dd y$ with a
normalized $\phi\in H^1(\RR)$ having trace $\phi(0)=0$. Then the
magnetic UV cost from the return is bounded by
$2\|\phi'\|_2^2/\pi$ in the above normalization.
\end{proposition}
\begin{proof}
On either half-line the zero trace permits the one-dimensional Hardy
bound
$\int |\phi(y)|^2/y^2\,\dd y\le4\int|\phi'(y)|^2\dd y$.
For a smooth test function compactly supported in the open positive
half-line, integration by parts gives
$\int|\phi|^2/y^2=2\operatorname{Re}\int\bar\phi\phi'/y$.
Cauchy--Schwarz proves the inequality. Reflection gives the other
half-line and density in $H^1_0$ gives the zero-trace case. Insert
the bound in \eqref{eq:uvpeP-Dirac-mark} after the heat integral.
\end{proof}
For example, a normal Hamiltonian $-\partial_y^2+a^2 y^2$ with
$a>0$ and its first odd eigenstate has density
$(2a^{3/2}/\sqrt\pi)y^2e^{-ay^2}$. In that preparation
$N_\mu(t)=a^{3/2}\sqrt{t}/[2\pi(1+at)^{3/2}]$. Its contribution
to the ultraviolet coefficient is integrable all the way down to
$t=0$. A continuous density positive at the origin instead yields
$N_\mu(t)\sim w(0)/(2\sqrt{\pi t})$. The parity here is reflection
of the normal coordinate. An assignment of internal fermion
chirality in PE-CGG-HOUQFT does not establish this normal boundary
condition. A derivation for the physical branch must show which
midpoint observable occurs in the charged determinant, whether its
normal form has this domain, and whether there is a zero-visit
component. Source-dependent measures, other charged histories and
sewings with additional photon lines have their own derivatives and
bounds. This calculation removes a specified bare one-loop
coefficient's UV divergence by a property of its normal state.

\subsection{Free normal data and the bare electromagnetic response}
\label{subsec:g5I-boundary}

The master integrand \eqref{eq:uvpeA-curvature} and the quantum
pushforward \eqref{eq:uvpeA-induced} take different inputs. The latter
requires a represented charged operator and a boundary state. For example,
a coefficient algebra generated solely by commuting raw fields, their
derivatives and products is purely even. Splitting a matrix representation
into left and right blocks does not supply an independent odd integration
variable. A Berezin determinant requires a graded raw algebra and measure,
or a quantum construction deriving its effects. Even after fixing the charged
Hessian, one can ask whether the state data that fix all free heat profiles
also fix the source derivative on a closed path.

Work on one Euclidean sector-plane and retain one noncompact normal
coordinate $y$ in reference units. In \eqref{eq:uvpeP-Dirac-mark} take
$c_0=0$ and retain the midpoint position as an observable with its
Born measure $n(y)\dd y$. The magnetic response is proportional to the
following positive cost:
\begin{align}
 J_\delta(n)&=\int_\delta^T N_n(t)\,\frac{\dd t}{t},
 &N_n(t)&=\frac{1}{2\pi t}\int e^{-y^2/t}n(y)\dd y .
 \label{eq:g5I-cost}
\end{align}
The possible mass weight $e^{-t m_e^2}$ lies between two positive
constants on this UV interval and equals $1+O(t)$ at zero. The same
$J_\delta$ occurs for the magnetic observable on the actual even GCM
charged gauge-root block of Section~\ref{sec:pe-curvature}: use two
half-time heat kernels, retain the intermediate normal position, and
close at zero. The constant normal parallel transports on the two legs
are inverse to each other. The finite, nonzero vector--ghost magnetic
coefficient in \eqref{eq:uvpeC-root-response} then multiplies
\eqref{eq:g5I-cost}, up to its sign and fixed normalization. This
is a common test on two specified charged Hessians. Only the second
is the even gauge-root restriction of the curvature action.

\begin{proposition}[The boundary node required of a positive preparation]
\label{prop:g5I-mixed}
Let the normalized normal density matrix be
$\rho=\sum_j p_j|\phi_j\rangle\langle\phi_j|$ with $p_j>0$ and
$\sum_j p_j\|\phi_j'\|_2^2<\infty$. Put
$n_\rho(y)=\sum_jp_j|\phi_j(y)|^2$. Then
$J_0(n_\rho)<\infty$ exactly when $\phi_j(0)=0$ for every occupied
state. In that case $J_0(n_\rho)\le
(2/\pi)\sum_j p_j\|\phi_j'\|_2^2$.
\end{proposition}
\begin{proof}
The finite-energy assumption puts each $\phi_j$ in $H^1(\RR)$, which
has a continuous representative. Any nonzero value at the origin
gives an infinite inverse-square integral for that summand.
For the converse use the half-line Hardy inequality of
Proposition~\ref{prop:uvpeP-Dirac-Hardy} on each state and then
sum by positivity. The condition can also be applied to finitely
many normal components with a unitary normal connection: parallel
transport to a common frame makes the covariant derivative ordinary
and preserves both the pointwise norm and the node.
\end{proof}

\begin{proposition}[Identical free normal spectral laws with opposite UV outcomes]
\label{prop:g5I-isospectral}
For the same free normal kinetic action, and the same midpoint
observable and electric charge, there are normalized even pure
preparations which give identical expectations of every bounded
Fourier multiplier of the normal momentum but respectively finite
and divergent costs \eqref{eq:g5I-cost}. They may both be real
Schwartz wavefunctions with finite normal kinetic energy.
\end{proposition}
\begin{proof}
Use the unitary Fourier transform and take two nonnegative, even
smooth functions $g_1,g_2$ of compact, disjoint momentum support
with equal nonzero integrals. Scale them so that
$\|g_1\|_2^2+\|g_2\|_2^2=1$. Define the two states by
$\widehat\phi_-=g_1-g_2$ and $\widehat\phi_+=g_1+g_2$.
All expectations of a multiplier $m(\xi)$ are integrals against
the same density $g_1^2+g_2^2$. This includes translations,
resolvents and the entire heat family. Both states have even
reflection parity. Nevertheless $\phi_-(0)=0$ and
$\phi_+(0)>0$. Proposition~\ref{prop:g5I-mixed} now applies.
More precisely put $n_\pm=|\phi_\pm|^2$. Changing variables
$y=\sqrt{t}z$ gives
\[
 N_{n_+}(t)=\frac{n_+(0)}{2\sqrt{\pi t}}+O(\sqrt{t}),
 \qquad N_{n_-}(t)=O(t^{3/2}).
\]
The second estimate uses the even smooth node
$\phi_-(y)=O(y^2)$. Thus
$J_\delta(n_+)=n_+(0)/(\sqrt{\pi\delta})+O(1)$ and
$J_\delta(n_-)\to J_0(n_-)$. Multiplication by $e^{-tm_e^2}$
leaves the leading divergence unchanged.
\end{proof}

This ambiguity has a cutoff-uniform test. Give the second momentum
band a relative phase by taking
$\widehat\phi_\theta=g_1-e^{\ii\theta}g_2$. The free spectral
laws are still identical for every $\theta$. In position space the
two inverse transforms of $g_1,g_2$ are real, hence exactly
\[
 n_\theta=\cos^2(\theta/2)n_-+\sin^2(\theta/2)n_+ .
\]
Even if $\theta_\delta\to0$ and the regulated states converge to
$\phi_-$ in every Schwartz seminorm, convergence of the magnetic
response to that of the limiting state requires and suffices that
\begin{equation}
           \sin^2(\theta_\delta/2)=o(\sqrt\delta).
 \label{eq:g5I-rate}
\end{equation}
For example $\theta_\delta=a\delta^{1/4}$ leaves a finite extra
cost $a^2 n_+(0)/(4\sqrt\pi)$. If
$\delta^{1/4}\ll|\theta_\delta|\ll1$, the cost diverges.
These statements follow from the positive convex combination and
the two preceding estimates. A source fixing only the limiting
normal vector and its free propagators therefore does not certify
uniform removal of the bare cutoff. A source whose entire cutoff
state remains in the nodal domain with uniformly bounded kinetic
energy satisfies the bound in Proposition~\ref{prop:g5I-mixed};
identifying its actual limit also requires control of the
inverse-square tails. For example, convergence in $H^1$ within
that domain suffices for a pure preparation by Hardy's inequality.

\paragraph{Can a node survive sewing a second return?}
The free normal action permits an exact test of stability under
gluing. Attach an ordinary heat interval of duration $u>0$ to a
nonzero boundary vector $\psi\in L^2(\mathbb R)$. In this paragraph
$H=-\partial_y^2$, and $f_u=e^{-uH}\psi$. Normalizing $f_u$ would
multiply each cost below by a finite positive constant. Keep the
seam at $y=0$ and set the unvisited weight $c_0$ to zero. For $r$
returns with independent durations $0<t_j<T$ but one shared
midpoint $y$, take the positive scalar term with weight
$|f_u(y)|^2\,\dd y$ and one factor
$(2\pi t_j)^{-1}e^{-y^2/t_j}\,\dd t_j/t_j$ for each return.
Integrating the durations by Tonelli gives its exact cost
\[
 J_r(u)=\frac1{(2\pi)^r}\int_{\mathbb R}
        |f_u(y)|^2\,\frac{e^{-r y^2/T}}{|y|^{2r}}\,\dd y .
\]
The factors $e^{-m_e^2t_j}$ preserve the test, since each lies between $e^{-m_e^2T}$ and one.
These are separate return durations coupled through a single measured
position. Two independently drawn midpoint positions give a different
integral.

Heat evolution makes $f_u$ smooth. For $r=1$ a finite cost requires
$f_u(0)=0$; for $r=2$ it also requires $f'_u(0)=0$. Taylor's formula
makes the same conditions sufficient, since $f_u\in L^2$ controls
the tails. A finite cost after one specified preparation therefore
asks for a node of the appropriate order. Requiring the cost to
stay finite as the duration of an attached interval varies has a
stronger consequence.

To see it, write $g(k)=\widehat\psi(k)$ in the unitary Fourier
convention. The boundary value and its derivative are, up to the
common Fourier normalization,
\[
 \begin{aligned}
 f_u(0)&\ \propto\ \int_0^\infty e^{-uk^2}[g(k)+g(-k)]\,\dd k,\\
 f'_u(0)&\ \propto\ \int_0^\infty e^{-uk^2} i k[g(k)-g(-k)]\,\dd k .
 \end{aligned}
\]
After $v=k^2$ these are Laplace transforms in $u$. For
$\operatorname{Re}u>0$, the Gaussian times every polynomial in $k$
is square integrable; Cauchy--Schwarz gives absolute convergence and
holomorphy on that half-plane. Values zero on any set with an
accumulation point in $(0,\infty)$ consequently vanish throughout
the half-plane. On any vertical line $\operatorname{Re}u=u_0>0$
these are Fourier transforms of integrable functions of $v$.
Uniqueness first forces $g(k)+g(-k)=0$ when the value vanishes, and
then $g(k)-g(-k)=0$ when the derivative vanishes.

Thus the states with finite one-return cost for such a set of
attached durations are exactly the odd states. For the specified
two-return term only the zero vector works. In particular, two
positive visits of zero net winding can retain their record and
still fail after sewing. This conclusion concerns the displayed free
normal transfer, shared position and positive contraction. The
physical action might change the transfer, introduce an independent
relative position, or cancel the term in the complete source.
Its determinant and photon derivatives determine which possibility
occurs. The monograph treats the related test in which a momentum
kick acts before the attached interval~\cite{ChiapponeGHOUQFT2026}.

\paragraph{What a coherent boundary pairing retains.}
A coordinate can occur on every history and yet disappear from the
closed coefficient. Consider a positive physical boundary channel
with incoming half-density map $W$ of Hilbert--Schmidt class, an
outgoing transfer $L$ bounded, and a commuting displacement observable
with joint spectral resolution $P$ on $\mathbb R^D$. Here $D$ counts
independent real momentum components on the selected Euclidean plane.
Write $U(k)=\int e^{i k\cdot x}P(\mathrm dx)$.
The positive paired coefficient is
$B(k)=\|L U(k)W\|_{\rm HS}^2$. Both boundary maps include the
propagations and probes on their respective sides of the pairing.

\begin{proposition}[A boundary test before the momentum integration]
\label{prop:phys-uv-boundary-overlap}
Choose orthonormal input and output bases $u_a,v_b$. Define the
complex overlap measures
\[
 \nu_{ba}(S)=\langle L^*v_b,P(S)Wu_a\rangle
\]
with the inner product linear in its second entry. For any
nonnegative integer $N$, the positive coefficient has finite
absolute moment
\[
 \int (1+|k|^2)^N B(k)\,
              \frac{\mathrm d^D k}{(2\pi)^D}
\]
exactly when every $\nu_{ba}$ has a density $h_{ba}$ in
$H^N(\mathbb R^D)$ and
$\sum_{a,b}\|(1-\Delta)^{N/2}h_{ba}\|_2^2<\infty$.
The integral equals this sum.
\end{proposition}
\begin{proof}
Each overlap measure has finite total variation by Cauchy--Schwarz
for spectral projections on disjoint sets. Its Fourier transform
with positive exponent is the $ba$ matrix entry of $L U(k)W$.
Expand the Hilbert--Schmidt square and apply Tonelli. For each entry
Plancherel's identity gives the stated norm. In the reverse
direction a transform with the finite weighted norm has an inverse
$H^N$ transform. Uniqueness as tempered distributions identifies
the measure with that density.
\end{proof}

For example, if $L^*L$ commutes with every $U(tv)$ in a nonzero
direction $v$, then $B(k+tv)=B(k)$ for all $t$. If $B$ is
nonzero anywhere, continuity gives a tube of positive transverse
width on which even the zeroth moment diverges. Taking $L=I$
loses every kick in this positive pairing. Noncommutation alone is
insufficient as well: on two slots with displacements $+1,-1$,
take both boundary vectors to be $(1,1)/\sqrt 2$. Their paired
coefficient is $\cos^2 k$ and the overlap measure is atomic.
More generally, if this displacement observable acts on a
finite-dimensional space, every overlap measure is atomic.
Any nonzero paired coefficient then fails the unit-weight test.
This conclusion concerns a hard direction for which the graph has
no other decaying factor.

\paragraph{How many independent coordinates does the detector retain?}
The preceding overlap test also quantifies the role of the residual
propagators. On a fixed Euclidean hard subspace of dimension $D$, suppose
a nonzero positive finite displacement law $\mu$ gives a squared
amplitude $|\widehat\mu(k)|^2$, with the characteristic-function
convention $\widehat\mu(k)=\int e^{ik\cdot x}\,\dd\mu(x)$. Use fixed
units and let $s>0$ specify the residual power of a radial comparison.
The Gamma integral and the Gaussian Fourier transform give
\[
\begin{aligned}
 &\int_{\mathbb R^D}\frac{|\widehat\mu(k)|^2}{(1+|k|^2)^s}
                   \frac{\dd^D k}{(2\pi)^D}\\
 &\quad=\frac{1}{(4\pi)^{D/2}\Gamma(s)}
   \int_0^\infty t^{s-1-D/2}e^{-t}\,
      \mathcal H_\mu(t)\,\dd t,\\
 &\mathcal H_\mu(t)=\iint
       e^{-|x-y|^2/(4t)}\,\dd\mu(x)\,\dd\mu(y).
\end{aligned}
\]
The heat variable here represents an exact denominator identity.
Tonelli applies first to the nonnegative momentum integrand, and at
each $t>0$ the Gaussian is integrable against the finite pair measure.

Suppose a positive amount of the sewn law lies in the Lipschitz
image of a bounded set of $m$ real geometric coordinates. Cover that
subset of displacement space by at most $C r^{-m}$ sets of diameter
at most $r$, for $0<r\le1$, and make them disjoint by taking successive
differences. The sum of their squared masses is at least the square of
the covered mass divided by their number. Pairs in each set therefore
give $\mathcal H_\mu(t)\ge c\,t^{m/2}$ at $r=\sqrt t$.
The comparison integral diverges whenever $m+2s\le D$, with at least
a logarithm at equality. If there are also constants with
$\mu(B(x,r))\le C' r^m$ uniformly on the whole support for
$0<r\le1$, then $\mathcal H_\mu(t)\le C''t^{m/2}$ for $0<t\le1$.
Indeed, use
$e^{-|x-y|^2/(4t)}=\int_{|x-y|^2/(4t)}^\infty e^{-v}\,\dd v$
and bound the ball of radius $2\sqrt{tv}$ inside the integral.
Under both estimates $m+2s>D$ is the exact threshold.

A long history need not have a high-dimensional image after sewing.
One shared smooth parameter can feed arbitrarily many line phases
while their joint image remains a curve. The count uses the actual
displacements dual to the independent hard momenta. To transfer a
divergence to a physical graph requires a lower bound on its complete
absolute integrand over the hard domain used in this comparison.
A convergence argument instead requires an upper bound throughout
all its hard regions. The result here applies to the positive squared
channel and radial weight just specified.

An ordered sequence of visits has a related test. Insert each
phase operator at its place between the actual transfer operators.
Whenever the scalar ordered spectral matrix elements define finite
complex measures on the visit coordinates, push each measure forward
by the linear combination of displacements paired with the independent
momenta. Its Fourier transform is the corresponding ordered matrix
entry, so the same proof tests these pushed-forward overlaps.
This step does not assume that successive observables have a joint
probability law. For source derivatives, a term pairing two
different marked amplitudes is controlled by Cauchy--Schwarz in
weighted momentum space, using the two separate overlap norms.
One must take every term from differentiation of the original maps
and measure, including direct contacts. If a positive residual
geometric measure remains, the corresponding norm sums must be
integrable against its actual weight, including the charged vacuum.

\paragraph{A charged return through successive coherent visits.}
An explicit source exhibits both retention and loss. Use $n$
displacement components on one fixed boundary channel, with the same
proper-time Hamiltonian on every segment,
\[
 H=\tfrac12(-\Delta_x+\omega^2|x|^2),\qquad
 E_0=n\omega/2,\qquad \omega>0 .
\]
Use its normalized ground state at the two ends. For $h\ge1$
marks at times $t_1=0<t_2<\cdots<t_h=T$, write
$\beta_j=t_{j+1}-t_j$ and $U_j=e^{i k_j\cdot X}$, where $X$ is
position. The ordered ground-to-ground amplitude is
\[
 A=\langle0|U_h e^{-\beta_{h-1}H}
               \cdots U_2 e^{-\beta_1H}U_1|0\rangle ,
\]
with just $U_1$ for a single mark. For instance, couple two
Grassmann slots through this visit amplitude and its adjoint
with diagonal entries one and off-diagonal entries $-g\bar A,gA$,
$g$ a real probe amplitude. Their full Berezin determinant is
$1+g^2|A|^2$. Hence the marked coefficient is the coherent square
of one amplitude. Its intermediate positions are integrated before
the square is taken.

Put $R_{ij}=e^{-\omega|t_i-t_j|}$. The ground wavefunction is
proportional to $e^{-\omega |x|^2/2}$. Conjugating $-(H-E_0)$ by
this function gives $\tfrac12\Delta-\omega x\cdot\nabla$.
Its stationary Gaussian has covariance $I_n/(2\omega)$ and successive
positions satisfy
\[
 X_{j+1}=a_j X_j+\sqrt{1-a_j^2}\,Z_j,\qquad
 a_j=e^{-\omega\beta_j},
\]
where the independent $Z_j$ have the stationary Gaussian law and
are independent of $X_1$. This transition follows by integrating
the linear drift and the unit-variance Brownian increment. The
ground-state transform thus gives the ordered characteristic
\[
 A=e^{-E_0T}
   \exp\!\left[-\frac1{4\omega}
              \sum_{i,j}R_{ij}\,k_i\cdot k_j\right].
\]
The triangular change from the independent innovations to the
visits has diagonal entries $1,\sqrt{1-a_1^2},\ldots,
\sqrt{1-a_{h-1}^2}$ in each component. Consequently
$\det R=\prod_{j<h}(1-a_j^2)$, and the complete uncut integral is
\begin{equation}
 \begin{aligned}
 \int |A|^2\prod_{i=1}^h\frac{\mathrm d^n k_i}{(2\pi)^n}
 &=e^{-2E_0 T}\left(\frac{\omega}{2\pi}\right)^{nh/2}\\[-1mm]
 &\quad\times\prod_{j<h}(1-e^{-2\omega\beta_j})^{-n/2}.
 \end{aligned}
 \label{eq:phys-uv-coherent-visit}
\end{equation}
Every positive interval supplies a new independent fluctuation.
If a junction has exactly zero duration, the two kicks at that
junction combine to $U(k_i+k_{i+1})$ and their difference is an
undamped hard momentum.

The formula also tests a common law of durations, including
correlations. Let $b$ be the minimum of one and the numbers
$1-e^{-2\omega\beta_j}$, taking $b=1$ for $h=1$. With all
$\beta_j>0$, multiplication of the integrand by
$(1+\sum_i|k_i|^2)^N$ multiplies its integral by a quantity bounded
above and below by positive constants times $b^{-N}$ at fixed
$h,n,\omega,N$. Indeed $R$ factors as $M\operatorname{diag}
(1,1-a_1^2,\ldots,1-a_{h-1}^2)M^{\mathsf T}$, where $M$ is
unit lower triangular and $M^{-1}$ has subdiagonal $-a_j$ and all
other off-diagonal entries zero. Both norms are bounded in terms
of $h$. The least eigenvalue of $R$ is therefore comparable to
$b$, and rescaling the Gaussian gives the weighted estimate.
Finiteness for a positive geometric law is equivalent to integrability
of this explicit product times $b^{-N}$, including the factor
$e^{-2E_0T}$. This is a test on the output of the charged and
boundary operators. In a physical GCM reduction the operators,
the prepared state and the inter-visit evolution must be derived
from that reduction before this test can be applied.

\paragraph{When the charged field selects the normal angular sector.}
The distinction between a charged visit and a neutral return can be
implemented by the same circle connection that supplies the Wilson
transport. On the fixed Euclidean carrier, consider the following
associated-line representation in reference units. The circle is
internal to the normal surface and its angle is $2\pi$-periodic.
For a real one-form $b_\mu(x,r)\,\mathrm dx^\mu$, independent of
$\theta$, the metric and angular coframe are
\[
 \mathrm ds^2=|\mathrm dx|^2+\mathrm dr^2+f(r)^2\Theta^2,
 \qquad \Theta=\mathrm d\theta+e\,b_\mu\,\mathrm dx^\mu.
\]
Here $e$ is the electric coupling in the representation. The volume
density is $f(r)$, and the horizontal lift of $\partial_\mu$ is
$\partial_\mu-e b_\mu\partial_\theta$. On a field
$e^{ij\theta}\phi(x,r)$, $j\in\mathbb Z$, the scalar Dirichlet form
per unit angular integral is exactly
\begin{equation}
 \begin{aligned}
 \int f(r)\Big[&
    |(\partial_\mu-i j e b_\mu)\phi|^2+|\partial_r\phi|^2\\
 &\hspace{13mm}+j^2 f(r)^{-2}|\phi|^2\Big]\,
             \mathrm dr\,\mathrm d^4x ,
 \end{aligned}
 \label{eq:phys-uv-angular-selection}
\end{equation}
with the horizontal index summed. This identity follows by using
the dual frame, including $f^{-1}\partial_\theta$. It applies to a
scalar normal degree of freedom with a spectator carrier spin index.
The change $b\mapsto b+\mathrm d_x\lambda$, $\theta\mapsto
\theta-e\lambda(x)$ preserves $\Theta$. Parallel transport of the
associated section along an oriented chord at fixed $r$ has factor
$\exp(i j e\int b)$, reversed on the opposite chord. Thus the same
integer gives the charge $je$ in the horizontal Wilson operator and
the repulsive angular potential at the tip. This identification uses
the charged representation on the normal circle. An electromagnetic
representation acting on an unrelated internal space supplies no
such relation.

There is an operator test of where that charge propagates. In the
source-free normal Hamiltonian write $P_j$ for the angular projection,
$H_j$ for its radial restriction on $L^2(f\,\mathrm dr)$, and $S_q$
for multiplication by $e^{iq\theta}$ at an oriented endpoint.
Then $S_qP_j=P_{j+q}S_q$. Identifying the neutral summand with its
radial space, two placements with the same total charge give
\[
 \begin{aligned}
 P_0 S_{-q}e^{-TH}S_qP_0 &\longleftrightarrow e^{-T H_q},\\
 P_0 e^{-TH}S_{-q}S_qP_0 &\longleftrightarrow e^{-T H_0}.
 \end{aligned}
\]
In the first placement each intervening radial state feels
$q^2/f^2$. In the second placement that potential is absent. More
generally a sequence starting in $j_0$ propagates in the sector
$j_0+\sum_{i\le l}q_i$ between the $l$th and $(l+1)$st insertions.
A single angle shared by two opposite charges can therefore carry
zero angular flux on an interval. Total charge neutrality of the
closed coefficient does not identify which cuts retain a charged
representation.

At zero background a charged collar provides one preparation of
the ground state used in Corollary~\ref{cor:geo-free-slope-uv}. Propagate a radial
vector of nonzero ground-state overlap by $e^{-T H_q}$ and normalize
it. For an isolated simple ground eigenvalue with a gap, its
excited component then vanishes exponentially as $T\to\infty$.
The resulting ground-state projector fixes the radial probability
in that corollary for $q\ne0$. For a finite collar one must use
the state it actually prepares. Convergence in Hilbert norm by
itself does not permit insertion of an unbounded charged vacuum
weight in this limit. The graph theorem uses the specified charged
endpoint diffusion, the harmonic photon readout and this common
boundary preparation. The local Dirichlet form identifies the
charge and the parallel transport; it does not alone select the
heat-displacement rule in the charged bilinear or the BFV state of
the full chiral theory.

\paragraph{What a finite charged collar prepares.}
The curvature-minimizing quartic chart gives a test of the boundary
state even when no ground-state projection is made. Use the normal
metric $\mathrm dr^2+r^8\mathrm d\theta^2$, with the outer slope
fixed at four, and the same electromagnetic heat link as in the
preceding radial calculations. Work in reference units on a fixed
four-torus, with $m,\kappa>0$, charged heat lapse $\kappa r$,
harmonic photon readout at $r$, and the harmonic photon mode fixed.
Let $Z_0(r)$ be this charged operator's full zero-field determinant.
At small $r$ its logarithm is bounded above and below by positive
multiples of $r^{-2}(1+\log(1/r))^3$. At fixed depth the preceding
charged-graph estimates, with all chord contacts and photon sewings,
cost at most $C_G r^{-N_G}Z_0(r)$ for each finite graph. Both
properties refer to this same represented source.

A collar is a region of normal evolution of duration $t>0$.
Retain an angular Fourier number $j\ne0$ on the entire collar.
In the unitary radial coordinate $u=r^2\psi$ the Friedrichs
generator with Dirichlet outer endpoint is
\[
 H_j=-\partial_r^2+2r^{-2}+j^2r^{-8}
       \quad\text{in }L^2((0,1),\mathrm dr).
\]
Sew an incoming positive trace-class boundary operator $D$ through
the two half-densities $e^{-tH_j}D e^{-tH_j}$. We use $t$ for the
duration of each half. The normalization is performed once after
including the determinant and all prescribed geometric integrations.

\begin{proposition}[Sharp short-collar cost of the bare charged determinant]
\label{prop:final-collar-uv}
Put $F_N(r)=r^{-N}Z_0(r)$ for $N\ge0$,
and $L_t=1+|\log t|$. For sufficiently small $t>0$,
\begin{align*}
 c\,t^{-1/3}L_t^4
 &\le \log\|F_0^{1/2}e^{-tH_j}\|_{\rm op}^2,\\
 \log\|F_N^{1/2}e^{-tH_j}\|_{\rm HS}^2
 &\le C_N\,t^{-1/3}L_t^4 .
\end{align*}
The constants may depend on the fixed volume and microscopic data.
Every nonzero positive trace-class input at fixed $t>0$ therefore
has a finite nonzero vacuum and absolutely finite bare graphs at
every finite electromagnetic order. The radial Gibbs preparation
at every inverse temperature $\beta>0$ has the same property.
\end{proposition}
\begin{proof}
Let $u$ be a normalized eigenfunction of energy $E$. Choose
$0<3c_0<|j|$, $a=j^2-9c_0^2$, and
$h=\min(1/2,(a/[2(E+1)])^{1/8})$. Use the bounded truncations of
$W=\exp[c_0(r^{-3}-h^{-3})_+]$, where $x_+=\max(x,0)$.
Testing the eigenvalue equation against $W^2u$ gives
\[
 \|(Wu)'\|_2^2+
 \int_0^1 (2r^{-2}+j^2r^{-8}-E)|Wu|^2\,\mathrm dr
 =\int_0^1 |W'u|^2\,\mathrm dr .
\]
For each truncation $|W'|\le3c_0r^{-4}W$, and $W=1$ for
$r\ge h$. Moving the derivative cost to the left yields
$(a/2)\int_0^h r^{-8}|Wu|^2\,\mathrm dr\le E$.
Monotone convergence removes the truncation. In particular each
eigenfunction retains exponential suppression, with the energy
dependence carried by $h^{-3}$.

For the actual vacuum weight one obtains a sharper estimate than
using an $r^{-3}$ exponential throughout. A decreasing upper envelope
of $\log F_N$ is
$M_N(r)=C[1+r^{-2}(1+\log(1/r))^3]+N\log(1/r)$.
The function $M_N(r)-2c_0r^{-3}$ is increasing sufficiently near
zero, by differentiation. It follows, first for small $h$ and then
with a constant on the remaining compact range, that
$F_N\le C_N\exp[M_N(h)]W^2$ on the whole interval. Thus
\[
 \langle u,F_Nu\rangle
 \le C_N(1+E)\exp\!\left[
 C_N(1+E)^{1/4}(1+\log(2+E))^3\right].
\]
The positive potential gives $E_n\ge\pi^2 n^2$ by comparison
with the Dirichlet Laplacian. Summing this bound against
$e^{-2tE_n}$ gives the Hilbert--Schmidt estimate. To see the
power of the logarithm, put $y=t^{-4/3}L_t^4$. Below $E=y$ the
positive exponent is $O(t^{-1/3}L_t^4)$. For $E=vy$, $v\ge1$,
its ratio to $t^{-1/3}L_t^4$ is bounded by
$C v^{1/4}(1+\log v)^3$, whereas $tE/(t^{-1/3}L_t^4)=v$.
The supremum after subtracting $tE$ is therefore bounded on this
range. A remaining heat factor absorbs the polynomial and has a
trace bounded by a power of $t^{-1}$.

For the reverse estimate, choose a normalized smooth bump $v_h$
supported in $[h,2h]$, obtained by rescaling one fixed test.
Its mean energy is at most $C h^{-8}$. Cauchy--Schwarz on its
support and Jensen's inequality for its spectral measure give
\[
 \|F_0^{1/2}e^{-tH_j}v_h\|^2
 \ge \inf_{[h,2h]}Z_0\,
     \exp[-2t\langle v_h,H_jv_h\rangle].
\]
Take $h=K t^{1/6}L_t^{-1/2}$. The two terms in the lower bound
for the logarithm have respective coefficients bounded below by
$c'K^{-2}$ and $-C'K^{-8}$ times $t^{-1/3}L_t^4$. A sufficiently
large fixed $K$ gives a positive coefficient. This proves sharpness
in order for the worst boundary input. At fixed $t$, the trace
pairing with $D$ is bounded by
$\|D\|_1\|F_N^{1/2}e^{-tH_j}\|_{\rm op}^2$.
The vacuum is strictly positive by $Z_0>0$ and injectivity of
the collar. The thermal trace is the squared Hilbert--Schmidt norm
at $t=\beta/2$. The fixed-depth graph bound completes the proof.
\end{proof}

The bound also justifies the long-collar ground-state limit inside
these graph integrals. Let $P_0$ be the ground projection and
$E_1-E_0>0$ the radial gap. For fixed $t_0>0$ and $t>t_0$,
factor $e^{tE_0}e^{-tH_j}(I-P_0)$ at $t_0$. Its norm after left
multiplication by $F_N^{1/2}$ is at most
\[
 e^{t_0E_0}\|F_N^{1/2}e^{-t_0H_j}\|_{\rm op}
           e^{-(t-t_0)(E_1-E_0)} .
\]
This proves convergence with the determinant and each graph weight
for preparations of nonzero ground overlap.

The collar can itself participate in a return with arbitrarily
small duration. For example, sewing two one-dimensional heat legs
of duration $t/2$ between points at separation $d>0$ supplies
$(2\pi t)^{-1}e^{-d^2/t}$. This integrable small-$t$ weight
dominates the worst collar cost above at every fixed graph order.
The conclusion also permits $d\downarrow0$ under a common mark law
$\mu$ with $\int e^{b/d^3}\mu(\mathrm dd)<\infty$ for some $b>0$.
Indeed $d^2/t+b/d^3\ge c_b t^{-3/5}$ by minimization, so the
integrated bridge is bounded by $C t^{-1}e^{-c_b t^{-3/5}}$.
One may fix an upper duration or use an integrable large-duration
preparation. This comparison concerns the stated sewing with an
external separation mark and a conditional collar input of uniformly
bounded trace. It does not identify that mark with the integrated
readout $r$. Both the return and collar must be present in the
source to use this joint estimate.

The finite preparation time has mathematical content. With the same
bulk action and $j\ne0$, a normalized unpropagated vector
$u(r)=C r^4\chi(r)$, with $\chi$ smooth, equal to one near zero
and vanishing near one, belongs to the Friedrichs form domain.
Both its derivative and $r^{-4}u$ are square integrable. Its full
vacuum integral $\int |u|^2 Z_0\,\mathrm dr$ nevertheless diverges.
Thus the angular label, positivity and finite mean energy alone
leave open even the zero-field normalization. For the unrestricted
chiral branch the bulk curvature rule must determine which physical
angular channel crosses a seam, its evolution and the state actually
sewn there. A neutral cut, a source-dependent collar, or a change in
the charged determinant introduces different terms in this test.

\paragraph{Closing through the same normal coordinate.}
The return can instead read the very coordinate $r$ used by the
charged determinant. On the preceding quartic chart, couple its two
ordinary one-dimensional heat legs of duration $t/2$ to this measured
separation at the seam, and propagate its angular boundary channel for
duration $t$ on either side. The common proper-time integration then
contains $e^{-r^2/t}/(2\pi t)$ inside the radial trace. This is a
specified joint sewing rule. It preserves the intermediate coordinate
of a return with zero net displacement.

\begin{corollary}[Jointly short visits with one shared normal readout]
\label{cor:final-collar-joint-uv}
For every fixed $T>0$, $N\ge0$ and $j\ne0$ the following trace is
finite, with $F_N$ and $H_j$ from Proposition~\ref{prop:final-collar-uv}:
\[
 \int_0^T\!
   \| F_N^{1/2}e^{-r^2/(2t)}e^{-tH_j}\|_{\rm HS}^2
            \frac{\mathrm dt}{2\pi t}<\infty .
\]
Hence the same-source heat-link determinant and every fixed-order bare
graph remain integrable for this sewing, including at $t=r=0$ in the
limiting integral. A bounded positive incoming density, in particular
a normalized state, is allowed before the collars. Its overall vacuum
normalization is positive for a nonzero density.
\end{corollary}
\begin{proof}
Write $K_{2t}(r,r)$ for the radial heat kernel with respect to
$\mathrm dr$. The squared Hilbert--Schmidt norm is
$\int_0^1 F_N(r)e^{-r^2/t}K_{2t}(r,r)\,\mathrm dr$ by the semigroup
law and positivity, first for bounded multipliers. The Dirichlet
kernel with potential $2r^{-2}+j^2r^{-8}$ is dominated by the
whole-line free kernel, giving $K_{2t}(r,r)\le C t^{-1/2}$.

There is a stronger bound when $0<r<1/4$ and $r^5\le t$. Use the
first $r^5$ units of heat evolution and bound the remaining kernel
by $C t^{-1/2}$. In the positive-potential Brownian representation,
a path started at $r$ and surviving below $2r$ over this first interval
pays at least $j^2 r^5/(2r)^8$ in the exponent. The probability
of reaching $2r$ during the interval is at most
$2\exp[-r^2/(4r^5)]$ for Brownian motion of generator $\partial_r^2$,
by the reflection principle. Killed paths contribute zero. The same
bounds follow by monotone approximation of the nonnegative singular
potential. Thus $K_{2t}(r,r)\le C t^{-1/2}e^{-c_j/r^3}$ in this
region, for some $c_j>0$.

For $0<t<r^5$ keep the ordinary free bound. Half the bridge exponent
then gives $e^{-r^2/(2t)}\le e^{-1/(2r^3)}$, and the other half
integrates against $t^{-3/2}\,\mathrm dt$ at a cost at most $C/r$.
For $r^5\le t\le T$, integration costs at most
$2r^{-5/2} e^{-c_j/r^3}$ times a constant. The resulting sum is
integrable against every $F_N$ because
$\log F_N=O_N(r^{-2}(1+|\log r|)^3)=o(r^{-3})$.
Away from the tip the bridge gives an integrable small-time bound.
Tonelli's theorem justifies the positive trace calculation. Each signed
complete graph is then controlled by its fixed-depth absolute majorant.
\end{proof}

Here the geometric factors control complementary ranges of the same
integral. The displacement retained by the return controls $t<r^5$.
For longer visits the charged angular evolution controls the small
readout. Applying this mechanism to a physical chiral graph requires
deriving the shared source insertion, its proper-time law and its
nonzero intervening angular channel after BV--BFV sewing.

\paragraph{Does a normal spin structure force a protected angular sector?}
A spinor on a locked Lorentzian plane and a spinor on its normal
surface require different bundle data. The distinction can be
tested by writing out the normal Dirac operator on the quartic
chart. Take the polar orthonormal frame
$(\partial_r,f^{-1}\partial_\theta)$ with $f=r^4$ and
$\theta\in\mathbb R/(2\pi\mathbb Z)$. Write the spin transition
in this frame as $e^{2\pi i\eta_\perp}$, where
$\eta_\perp=0$ or $1/2$. This $\eta_\perp$ labels the normal
spin structure only. The bounding spin structure has antiperiodic
components in the polar frame: lifting one full frame rotation
around a disk gives the central sign $-1$. A punctured collar
admits either transition. The quartic metric itself has a singular
tip and supplies no smooth filling condition.

Use the radial gauge $a(r)\,\mathrm d\theta$ for an Abelian normal
connection of charge $q$. For the source-free stationary normal
Maxwell action the radial functional is proportional to
$\int |a'|^2/f\,\mathrm dr$. Its Euler equation gives
$a=a_0+B r^5/5$ with constant $B$. The action leaves the flat
holonomy $e^{2\pi i q a_0}$ undetermined on the puncture.
A connection smoothly extending over a disk, with shrinking
contractible circles, has trivial limiting holonomy.

\begin{proposition}[Heat suppression from normal spin transport]
\label{prop:normal-spin-gap-uv}
Take the positive heat generator $H$ given by the Friedrichs
closure of the squared normal Dirac form on compactly supported
spinors, with Dirichlet outer endpoint. Suppose
$\delta=\operatorname{dist}(\eta_\perp-q a_0,\mathbb Z)>0$.
For each $F_N$ of Proposition~\ref{prop:final-collar-uv},
summing over the entire angular spectrum and both normal
chiralities gives constants $c,C>0$ such that at small $t>0$,
\[
 \big\|F_N^{1/2}e^{-r^2/(2t)} e^{-tH}\big\|_{\rm HS}^2
 \le C t^{-1}\exp[-c t^{-3/5}].
\]
In particular the joint short-return trace with measure
$\mathrm dt/t$ is finite for this spinorial propagation. A
bounding spin structure and trivial limiting Abelian holonomy
give $\delta=1/2$ without selecting an angular Fourier mode.
\end{proposition}
\begin{proof}
Use the Pauli matrices for the normal Clifford generators.
Metric compatibility gives
$\nabla_\theta=\partial_\theta-i f'\sigma_3/2-iqa$.
The radial spin connection in this frame vanishes. Conjugating
the Dirac operator by multiplication by $f^{1/2}$ takes the
radial measure to $\mathrm dr$ and removes $f'/(2f)$ from its
radial derivative. On the mode $e^{i(n+\eta_\perp)\theta}$
the operators are
\begin{gather*}
 D_n=-i\sigma_1\partial_r+\sigma_2 W_n,
 \qquad W_n=(n+\eta_\perp-qa)/f,\\
 D_n^2=-\partial_r^2+ W_n^2+\sigma_3 W_n'.
\end{gather*}
These differential expressions are used on the test core
before closing the positive form. Put $\nu_n=n+\eta_\perp-qa_0$.
Then $W_n=\nu_n r^{-4}-q B r/5$. Uniformly in $n$, either scalar
potential in $D_n^2$ is bounded below by
$c_1 \nu_n^2 r^{-8}+c_2 \nu_n^2-C$ on $(0,1)$. Indeed
$W_n^2\ge \nu_n^2 r^{-8}/2-C$ and
$|W_n'|\le4|\nu_n|r^{-5}+C$. The latter term is absorbed near
the tip using $|\nu_n|\ge\delta$ and away from the tip using the
quadratic bound in $\nu_n$.

Positive scalar heat-kernel comparison now extracts a factor
$e^{-c_2 t \nu_n^2+Ct}$ and leaves a Dirichlet heat kernel with
potential at least $c_1\delta^2 r^{-8}$. The Brownian estimate
of Corollary~\ref{cor:final-collar-joint-uv} applies with this
positive coefficient. In the radial integral on $r\le t^{1/5}$ it
gives $Ct^{-1/2}e^{-c/r^3}$; $\log F_N=o(r^{-3})$ absorbs half
the exponent. For $r>t^{1/5}$ the two halves of the retained
bridge exponent $-r^2/t$ absorb $F_N$ and give $-c t^{-3/5}$,
respectively. Thus the estimate per mode costs
$C t^{-1/2}e^{-c_2 t \nu_n^2}e^{-c t^{-3/5}}$ with adjusted
constants. The Gaussian sum in $n$ costs $O(t^{-1/2})$.
\end{proof}

The holonomy condition is sharp as an angular-gap test. If
$\nu_n=0$ for one mode, its $W_n$ is linear in $r$, and both
potentials $W_n^2\pm W_n'$ are bounded at zero. In this spinor
channel the Dirichlet ground density in the unitary radial measure
is comparable to $r^2\,\mathrm dr$ at the tip, by the regular
second-order equation and uniqueness for the initial data. The
choice of connection or spin structure must therefore be checked
on the actual charged slot.

\paragraph{What happens when the two spin slots are sewn.}
An electric-neutral bilinear can propagate in a different
bundle from its two legs. Suppose it propagates on one common
circle as a section of $\operatorname{Hom}(S_b,S_a)$, where
$a,b\in\{+1,-1\}$ are normal chiralities, the electric charges
of the two spin slots agree, and the kinetic form is the
connection Laplacian induced on this Hom bundle. These are
explicit assumptions on the represented action. The two spin
transitions and the common Abelian holonomy cancel in the
bilinear. Its polar-frame component is periodic. Write
$w_{ab}=(a-b)/2$ for its rotation weight. The angular derivative
is $\partial_\theta-i w_{ab} f'$, so its constant Fourier
component, after the same unitary radial change, has generator
\[
 H_{ab,0}=-\partial_r^2+(2+16w_{ab}^2)r^{-2}.
\]
Here $2r^{-2}$ comes from conjugating the scalar radial
Laplacian and $16w_{ab}^2r^{-2}$ from the induced connection.
For $a=b$ the identity channel is parallel. For opposite normal
chiralities the connection produces an inverse-square barrier.
In either case the regular ground solution at zero is comparable
to $r^{\gamma_{ab}}$, where
$\gamma_{ab}=(1+\sqrt{9+64w_{ab}^2})/2$, by the indicial
equation and the Friedrichs condition.

For example, use this single normal carrier for the one-hop
Wilson probe of Proposition~\ref{prop:physical-neutral-tip-uv},
keeping its hopping width, Gauss energy and endpoint preparation
rule. The hop acts as the identity on this normal slot and both
endpoints use its ground state. The same calculation reduces
finiteness of the coefficient of $g e^{2n}$ to integrability
of $r^{2\gamma_{ab}-n/2}\,\mathrm dr$ at zero. The Coulomb
moments remain finite because $1/r$ has arbitrarily small form
bound relative to this Dirichlet operator. The exact condition
is $n<4\gamma_{ab}+2$. The first divergent coefficient is
$g e^{20}$ for the scalar channel and $g e^{44}$ for either
opposite-chirality channel. These thresholds concern the term
linear in hopping. They test the stated common-circle kinetics;
separate normal propagation of the original spin slots retains
their individual squared angular costs until contraction. The raw
physical source must specify the normal spin lift, limiting
holonomy, kinetic pairing and placement of that contraction.

\paragraph{Integrating the flat normal holonomy.}
A nonzero gap at almost every boundary datum can still be
insufficient for the bare vacuum. On the punctured chart take
flat normal connections, with $B=0$, whose holonomy is integrated
once with positive law $\lambda$. It is a static variable shared
by the entire collar in this example. Near a resonance write
$0<\nu=n+\eta_\perp-qa_0<\nu_0$ for one angular mode. Choose
the chirality with repulsive derivative term, whose radial
operator in $L^2(\mathrm dr)$ is
\[
 H_\nu=-\partial_r^2+\nu^2r^{-8}+4\nu r^{-5}.
\]
Use $v(r)=C r^4(1-r)$ as a normalized nonnegative incoming
vector. Its form energy is uniformly finite for
$0<\nu<\nu_0<1$. The normal Maxwell action has the same value on
all these flat connections. Specify a joint representation in
which the normal holonomy twists this collar and the charged
heat-link determinant at its readout is the $Z_0(r)$ used above.

\Needspace{4\baselineskip}
\begin{proposition}[A finite collar with averaged normal holonomy]
\label{prop:normal-spin-average-uv}
For every fixed $t>0$ there is $c_t>0$ such that at sufficiently
small $\nu>0$ the prepared vacuum satisfies
\[
 \big\|Z_0^{1/2} e^{-tH_\nu}v\big\|_2^2
 \ge \exp\!\left[c_t\nu^{-2/3}(1+|\log\nu|)^4\right].
\]
It is finite at each $\nu>0$. A boundary law with a density bounded
below by a positive multiple of $\nu^a$ near zero, $a>-1$, gives
an infinite joint vacuum. The normalized Haar law on a nontrivially
charged circle of flat holonomies has this property.
\end{proposition}
\begin{proof}
Take a smooth nonnegative $L^2$-normalized bump $v_h$ supported in
$[h,2h]$, obtained by rescaling a fixed bump in $(1,2)$. Then
$v\ge c h^{9/2}v_h$. Scalar heat evolution preserves this
inequality. Localizing the norm to that interval and testing against
$v_h$ gives
\[
 \|Z_0^{1/2}e^{-tH_\nu}v\|_2^2
 \ge c h^9\inf_{[h,2h]}Z_0\,
       e^{-2t\langle v_h,H_\nu v_h\rangle}
\]
by Jensen's inequality for the spectral law of $v_h$.
The mean energy is bounded above by
$C(h^{-2}+\nu^2 h^{-8}+\nu h^{-5})$.
Write $L=1+|\log\nu|$ and choose
$h=K\nu^{1/3}L^{-1/2}$. The lower bound
$\log Z_0(r)\ge c r^{-2}(1+|\log r|)^3$ implies a positive term
at least $c'K^{-2}\nu^{-2/3}L^4$ in the logarithm. The angular
potential costs at most $Ct K^{-8}\nu^{-2/3}L^4$. The two other
energy costs grow at most as $C_{K,t}\nu^{-2/3}L^{5/2}$.
Choose $K$ sufficiently large for this fixed $t$ and then take
$\nu$ small. This proves the lower bound. Finiteness at fixed
$\nu$ follows by heat comparison with
$-\partial_r^2+2r^{-2}+\nu^2 r^{-8}$, up to a finite additive
constant since $4\nu r^{-5}-2r^{-2}$ is bounded below at that
$\nu$. The fixed-channel proof of Proposition~\ref{prop:final-collar-uv}
uses only a nonzero squared angular coefficient and applies to
$\nu^2$. Integrating the lower bound proves the assertion.
\end{proof}

This example uses the determinant's own asymptotic growth.
A uniform gap as in Proposition~\ref{prop:normal-spin-gap-uv}
avoids the resonance. If the geometric measure integrates flat
holonomies down to resonance, its rate of approach also enters the
vacuum test. A dynamical connection with a kinetic term for the
holonomy requires its corresponding joint evolution.

\paragraph{Several short visits in one physical coefficient.}
The single-return proof has a pointwise form which is useful when
several visits are constrained to have equal or overlapping durations.
Let $L_j(t)=F_N^{1/2}e^{-r^2/(2t)}e^{-tH_j}$ in the preceding
quartic chart. For every fixed $N$ there are $c,C,t_0>0$ such that
\begin{equation}
 \|L_j(t)\|_{\rm HS}^2
   \le Ct^{-1/2}\exp[-c t^{-3/5}]
 \quad (0<t<t_0,\ |j|\ge1).
 \label{eq:physical-seam-uv}
\end{equation}
Indeed on $r\le t^{1/5}$ the charged heat kernel from the proof of
Corollary~\ref{cor:final-collar-joint-uv} costs
$Ct^{-1/2}e^{-c_1/r^3}$ uniformly for integer $|j|\ge1$.
The bound $\log F_N=o(r^{-3})$ absorbs half of this exponent.
On $r>t^{1/5}$ split the bridge exponent in half: one half is at
most $-t^{-3/5}/2$, the other is at most $-1/(2r^3)$ since
$t<r^5$. This absorbs $F_N$ with the free heat bound. Integration
over the unit radial interval gives~\eqref{eq:physical-seam-uv}.

To apply the estimate, suppose the actual sewn coefficient has a
finite product of visit operators $L_j(t)^* B L_j(t)$ joined by
bounded propagation or insertion maps. A bounded map $B$ can
retain signs, spin data and correlations with the other visits.
Each such operator has trace norm at most
$\|B\|\,\|L_j(t)\|_{\rm HS}^2$. The trace ideal inequality bounds
the trace of the product by the product of these norms and of the
remaining operator norms, including any bounded incoming density.
In particular, take finitely many integration variables $t_a\in(0,T]$
and allow several visits to share each $t_a$. If each approaching
zero occurs in at least one of the estimated visits and all
remaining factors and the joint measure density have at most
product-power growth in the $t_a^{-1}$, the complete integral is
absolutely finite. The stretched exponential for each distinct
shrinking time absorbs every such power. This handles simultaneous
limits without multiplying marginal state probabilities. An
angular transition, an extra determinant, or a contact in the
native action can change the factorization. Its operator must
be included in the trace bound on the sector in which it acts.

\paragraph{Several measured normal coordinates}
If the actual Gaussian normal block has $b\ge1$ independent
coordinates, the same midpoint calculation gives
\begin{align}
 N_n^{(b)}(t)&=(2\pi t)^{-b}\int_{\RR^b}
                    e^{-|y|^2/t}n(y)\,\dd^b y,\nonumber\\
 J_0^{(b)}(n)&=(2\pi)^{-b}\int_{\RR^b}n(y)
   \frac{\Gamma(b,|y|^2/T)}{|y|^{2b}}\,\dd^b y .
 \label{eq:g5I-codim}
\end{align}
The second identity follows by Tonelli and $u=|y|^2/t$. For a
finite positive measure it requires the inverse $2b$ moment near
the coincident point, with a coincident atom giving infinite cost.
For a smooth state whose first nonzero Taylor term at zero is
homogeneous of degree $j$, angular integration of its squared
modulus has positive leading coefficient. The radial integrability
test is $2j>b$, with a logarithm at equality. The number $b$ here
counts the coordinates actually read by the normal heat kernel;
it must be calculated from the represented Hessian and its constraints.

The examples keep fixed the reduced normal action, free spectral probes
and chosen representation. Their boundary states give different mixed
source correlations. They leave open which state and which midpoint
operator the full BV--BFV source selects. The distinction persists
without adding hidden charged species, changing the sector orientation
or fitting a counterterm. To apply the finiteness calculation to
induced O--QED, the state selection and its cutoff-domain estimate
must follow from the same assignment that supplies the graded charged
measure and the vertices to be sewn.

\paragraph{A cyclic normal trace tests coincidence at every basepoint}
The midpoint estimate above holds a loop's starting position fixed.
If the cyclic trace also integrates its starting point $r$, the two
legs give $e^{-|y-r|^2/t}/(2\pi t)$. For a finite positive joint law
of $r,y$ the exact condition is integrability of $|y-r|^{-2}$
near the entire diagonal, with zero mass on it. This follows from
the same $t$ integral used in \eqref{eq:g5I-cost}. In particular,
on the full free line an ordinary cyclic trace with a single
position multiplier $M_n$ and a normalized $n\ge0$ gives exactly
\begin{align}
 &\int_{\RR^2}\frac{e^{-|y-r|^2/t}}{2\pi t}
                  n(y)\,\dd r\,\dd y
   =\frac{1}{\sqrt{4\pi t}} .
 \label{eq:g5I-cyclic-mark}
\end{align}
This is the full trace of the two normal half-time kernels with the
midpoint insertion, including the integration over the initial
normal position. It retains a divergent magnetic cost for every
such $n$, including a state nodal at $y=0$. A mark can retain the relative position through the joint weight
$\nu(\dd r)|\phi_r(y)|^2\dd y$, with a finite positive measure on
starting positions. If $\phi_r(r)=0$ for $\nu$-almost every $r$
and the normal kinetic energy is integrable over $r$, the translated
Hardy inequality at each basepoint gives a finite integrated cost.
To realize this observable, the source must retain the starting
position as part of the closed-path data.

\paragraph{The first connected test of a shared return}
Let $d=y-r$ be shared by all marked closed loops in a specified
comparison on the same support. Assume its actual normalized
zero-source law $\mu$ is fixed, including any zero-field
determinants in the depth weight. Use the massless normal and
tangential magnetic calculation on $\delta<t<T$; the finite upper
endpoint controls the infrared. At fixed nonzero $d$ the magnetic
cost per loop is
\[
 I_\delta(d)=
  \frac{e^{-d^2/T}-e^{-d^2/\delta}}{2\pi d^2}.
\]
At $d=0$ take the continuous cutoff value. In this comparison the
conditional source-dependent log determinant starts as
$c I_\delta(d)B^2+O(B^4)$ with $c\ne0$. Tag each such closed-loop
factor by a formal counting variable $z$ and integrate over $d$
before taking the connected logarithm. Then
\begin{equation}
 [z^2 B^4]\log Z_\delta
      =\frac{c^2}{2}\operatorname{Var}_{\mu}(I_\delta).
 \label{eq:g5I-two-loop}
\end{equation}
Expanding the exponential at fixed $d$ gives its second moment;
the outer logarithm subtracts the square of its mean. The $B^4$
jet of one determinant carries only one power of $z$. This is a
connected coefficient from the common geometric integration; it
does not identify an individual photon-sewn graph. When the first
moment converges, finiteness of this coefficient requires exactly
the inverse fourth moment near zero with no coincident atom, by
monotone convergence of $I_\delta^2$. For the odd oscillator
density $n(d)=C d^2e^{-a d^2}$ the one-loop mean is finite, while
\[
 \operatorname{Var}(I_\delta)\sim
       \frac{C(1-2^{-1/2})}{\pi^{3/2}\sqrt\delta}.
\]
Indeed $d=\sqrt\delta\,u$ reduces the leading second moment to
$C/(4\pi^2\sqrt\delta)$ times
$\int_{\RR}(1-e^{-u^2})^2u^{-2}\dd u
=4\sqrt\pi(1-2^{-1/2})$. Dominated convergence uses the
integrable bound $(1-e^{-u^2})^2/u^2$, and integration by parts
gives the displayed constant. The square of the mean is bounded.
Thus the one-loop nodal-state criterion and a positive traveled
length alone do not establish even this connected higher-order test.

\paragraph{Sewing the complete scalar response with a dynamical photon.}
Consider a massive charged complex scalar with kinetic operator
\((-i\nabla-eA)^2+m^2\) on \(\mathbb R^4\). Tensor it with the
normal generator \(-(2\kappa)^{-1}\partial_z^2\), \(\kappa>0\),
and pair the endpoints by a fixed finite positive law of their
displacement \(z\). All traces below are per unit carrier volume.
Put \(k_0(t)=(4\pi t)^{-2}e^{-m^2t}\) and
\[
 g(x)=\int_0^1 (1-2u)^2 e^{-xu(1-u)}\,\mathrm du .
\]
In the representation by the negative proper-time integral of the
scalar heat trace, differentiate its integrand at opposite external
Fourier momenta. The resulting Hessian per unit proper time is
\((|q|^2\delta_{ij}-q_iq_j)t k_0(t)g(t|q|^2)\), times \(e^2\)
and the normal kernel. Indeed the ordered current insertions have
Gaussian numerator \(2\delta_{ij}/t+(1-2u)^2q_iq_j\); the direct
\(A^2\) variation contributes the contact that completes the
transverse tensor. Integrating the derivative of
\((1-2u)e^{-xu(1-u)}\) on \([0,1]\) gives
\(xg(x)=2[1-\int_0^1 e^{-xu(1-u)}\mathrm du]\).
With the normal kernel \(R_z\) given below and \(z\ne0\), the bare vacuum costs
\(\int k_0 R_z\,\mathrm dt/t\), while a nonzero transverse
fixed-momentum coefficient costs a constant times
\(\int t k_0 R_z g\,\mathrm dt\). Substitution \(t=z^2v\)
gives short-displacement orders \(|z|^{-5}\) and \(|z|^{-1}\),
respectively. The integrands after rescaling have finite positive
limits by domination.
For each normal displacement write
\(R_z(t)=(\kappa/(2\pi t))^{1/2}e^{-\kappa z^2/(2t)}\).
For an ordinary Euclidean Maxwell field the transverse propagator
is \(\Pi^T(q)/|q|^2\) at nonzero momentum. Its contraction with
the scalar Hessian has positive density, up to the graph's
fixed symmetry factor and choice of log-determinant sign,
\[
  S_z(q)=3e^2\int_0^\infty t k_0(t)R_z(t)g(t|q|^2)\,\mathrm dt .
\]
The factor three is the rank of the transverse projection in four
dimensions. The direct two-photon contact is included in \(g\).
The identity for \(xg(x)\) gives \(0\le xg(x)\le2\) and
\(xg(x)\to2\) by dominated convergence on the unit interval.
Since \(m>0\) and \(z\ne0\), a second application yields
\[
  |q|^2 S_z(q)\longrightarrow
  6e^2\int_0^\infty k_0(t)R_z(t)\,\mathrm dt>0
\]
at nonzero charge. Thus the bare sewing with measure
\(\mathrm d^4q/(2\pi)^4\) diverges quadratically, including for a
boundary state supported at a single nonzero displacement.

A different specified source evaluates the transverse Maxwell field
at normal distance \(r=|z|\) by harmonic extension before inserting
it in the scalar action. Each photon end then supplies \(e^{-r|q|}\);
the scalar current pair and direct contact use that same field.
With this assignment the positive sewing over \(q\) at fixed
\(r>0\) is
\[
   C\,r^{-5}\, I(mr),\qquad C>0 ,
\]
where \(C\) contains the charge and Fourier normalization. To see
the full short-distance cost, set \(q=u/r\), \(t=r^2v\). The
dimensionless integral becomes
\[
 I(a)=\int_0^\infty v^{-3/2}e^{-\kappa/(2v)-a^2v}
   \int_{\mathbb R^4} e^{-2|u|}g(v|u|^2)\,\mathrm d^4u\,\mathrm dv .
\]
It obeys \(0<I(a)\le I(0)<\infty\), since \(0<g\le1/3\)
and the remaining positive factors are integrable. Moreover
\(I(a)\to I(0)\) as \(a\downarrow0\). For a finite positive
displacement law the complete bare one-photon sewing therefore
requires exactly a finite inverse fifth moment near zero, counting
a zero atom as infinite. This is the same moment as for the
unsubtracted scalar vacuum. The fixed-momentum transverse
response alone costs one inverse power. The comparison establishes
which additional decay actually survives the shared integral when
the photon is sewn.

\paragraph{When the charged vacuum selects the boundary length}
\phantomsection\label{dualrev-uvb-boundary}
A finite determinant at every positive length can change whether
the joint source exists. For an explicit test on one fixed Euclidean
sector plane take a square four-torus of side $h$, volume $V=h^4$,
and a single complex Grassmann species with a scalar kinetic
operator. Write $L_\theta$ for the positive covariant Laplacian
plus $m^2$, $m>0$, in a constant connection of Wilson angle
$\theta$ along the first cycle. Its eigenvalues are
$m^2+|(2\pi n+\theta e_1)/h|^2$, $n\in\mathbb Z^4$.
Take the bare charged bilinear with matrix $I+e^{-\ell\sqrt{L_\theta}}$
and let both boundary half-densities share the same length $\ell>0$.
The log of the ordinary Fredholm determinant is
$g_\ell(\theta)=\operatorname{Tr}\log(I+e^{-\ell\sqrt{L_\theta}})$.
With fixed $\lambda>0$, $s>0$ and a smooth positive boundary
coefficient $b(\theta)$, sew the half-densities to the probability
law
\[
 \dd\nu_b(\ell)\propto\ell^{-1/2}
          e^{-\lambda\ell-b/\ell^s}\dd\ell .
\]
Thus $Z(\theta)=\int e^{g_\ell(\theta)}\dd\nu_{b(\theta)}(\ell)$
includes the entire charged vacuum. All parameters here specify
the comparison source in fixed units.

Put $\eta(z)=\sum_{j\ge1}(-1)^{j+1}j^{-z}$ for $z>1$.
Uniformly in the Wilson angle at short length,
\[
\begin{aligned}
 g_\ell(\theta)&=c_{4}\ell^{-4}-c_{2}\ell^{-2}+O(1),\\
 c_{4}&=\frac{3V\eta(5)}{4\pi^2},\qquad
 c_{2}=\frac{Vm^2\eta(3)}{8\pi^2}>0 .
\end{aligned}
\]
The zero-winding term follows by radial integration of the
continuum symbol with measure $V\dd^4p/(2\pi)^4$. Setting
$x=\ell\sqrt{m^2+|p|^2}$ and $a=m\ell$ reduces its dimensionless
integral to
\[
 \int_a^\infty(x^3-a^2x)\log(1+e^{-x})\dd x
 =6\eta(5)-a^2\eta(3)+O(a^4).
\]
The omitted integral over $[0,a]$ has modulus at most
$a^4\log(2)/4$. For the winding remainder, expand the logarithm
in its alternating cyclic series. Subordination of the $j$th
Poisson transfer contributes
$j\ell\,e^{-(j\ell)^2/(4t)}/(2\sqrt\pi\,t^{3/2})$ at heat time
$t$, cancelling the series factor $1/j$. The alternating Gaussian
sum lies between zero and one. Poisson summation of the remaining
torus heat kernel shows that a nonzero winding $w$ contributes at
most a constant times
\[
 \ell\int_0^\infty t^{-7/2}
    e^{-m^2t-h^2|w|^2/(4t)}\dd t.
\]
The sum over $w\ne0$ is finite even with the extra factor
$1+w_1^2$ for two angle derivatives. Indeed it decays
exponentially at small $t$ after summing the windings, while
$m>0$ controls their polynomial growth at large $t$. This gives
an $O(\ell)$ winding remainder and the same bound on its first
two derivatives.

Consequently the full vacuum and all inverse length moments are
finite for $s>4$, and for $s=4$ exactly when $b\ge c_{4}$.
For $0<s<4$ the vacuum diverges. At the critical coefficient
$b=c_{4}$ the mass term $-c_{2}/\ell^2$ controls every fixed inverse
power. To differentiate the critical source, the boundary law
must remain on the normalizable side throughout the source
neighborhood. For instance, set $s=4$ and
$b(\theta)=c_{4}(2-\cos\theta)$. Let $\nu$ be the normalized
boundary law at zero angle and $P$ its weighting by $e^{g_\ell(0)}$
followed by one overall normalization. Evenness and the product
rule give the complete uncut response
\[
 (\log Z)''(0)=\langle g_\ell''(0)\rangle_P
   +c_{4}\big(\langle\ell^{-4}\rangle_\nu
                    -\langle\ell^{-4}\rangle_P\big).
\]
The winding estimate bounds the determinant derivative, which
includes its direct and sequential contacts. Both boundary
moments exist by the critical estimate. Their difference comes
from differentiating the normalized boundary law itself. Conversely,
if $b(0)=c_{4}$ and $b(\theta)<c_{4}$ for arbitrarily small angles,
the same expansion gives an infinite partition there. This is a
flat-holonomy response of the stated finite-volume action. The
physical GCM branch must determine its own charged operator and
paired boundary density.

The mass also controls what happens at the threshold. Set \(m=0\),
\(s=4\) and hold the boundary coefficient fixed at \(b=c_4\).
The zero winding now gives exactly \(c_4\,\ell^{-4}\). For the
undifferentiated nonzero windings, the sum in the heat estimate grows
at most as \(O(t^2)\) at infinity. Its product with \(t^{-7/2}\)
is integrable there even at zero mass. The zero of the Laplacian
at zero angle contributes \(\log 2\) to the mode sum, as follows
also by letting the positive mass decrease to zero. Consequently
\(g_\ell(\theta)=c_4\,\ell^{-4}+O(\ell)\) at short length,
uniformly in the angle. The combined density is comparable to
\(\ell^{-1/2}\) near zero. The massless critical vacuum exists,
but its moment of \(\ell^{-u}\), for \(u\ge0\), exists exactly
when \(u<1/2\).

There is a separate infrared effect in the angle response. The
zero Fourier mode contributes
\(\log(1+e^{-\ell v|\theta|})\), with \(v=1/h\).
Remove just this summand to compute the derivative of the other
modes near zero angle. In the subordinated heat trace for \(t\le1\),
the zero winding has zero angle derivative, the other windings
are exponentially suppressed, and the derivative of the removed
heat mode is \(O(t)\). For \(t\ge1\) the remaining Fourier
modes have a uniform positive gap near zero angle, so their
differentiated heat trace decays exponentially up to a power of
\(t\). The alternating subordination weight is bounded by
\(\ell/(2\sqrt\pi\,t^{3/2})\). These bounds give an
\(O(\ell)\) first derivative for the remainder, which is even
and has derivative zero at the origin. The complete one-sided
derivatives therefore satisfy
\[
  (\log Z)'(0^\pm)=\mp\tfrac12 v\,\mathbb E_{P_0}\ell ,
\]
where \(P_0\) is the normalized determinant-weighted law at zero
angle. Its first positive length moment is finite. The derivative
bound is integrable there and at infinity, where the boundary
law has exponential decay and the determinant is uniformly
bounded for \(\ell\ge1\). The positive expectation fixes a nonzero jump in slope. Thus the
finite critical vacuum has a cusp in its flat-holonomy response.

\section{From a history phase to an absolute graph bound}
\label{sec:pe-graphs}

Fix a post-GTL sector-plane $\bin=(s,b_s)$ and a complete graph on
its Euclidean carrier, with fixed external momenta. In reference units
write $k\in\RR^d$ for its independent loop coordinates, $d=4L$.
The following question arises even when the normal Hessian can be
inverted exactly. Which of the $d$ hard directions are still detected
by the phase after the history has been closed, the boundary state
paired and the external source differentiated?

\subsection{The response of a sewn Gaussian normal block}

At finite regulator let $Y$ be the real normal coordinates of a
positive quadratic block on the physical slice. For example, start
with mean $m_0$, covariance $G>0$, and a linear boundary condition
$DY=a$ with independent rows. Gaussian conditioning gives, with
$M=DGD^{\mathsf T}$,
\begin{equation}
\begin{aligned}
 m&=m_0+GD^{\mathsf T}M^{-1}(a-Dm_0),\\
 \Sigma&=G-GD^{\mathsf T}M^{-1}DG .
\end{aligned}
\label{eq:uvpeN-seam}
\end{equation}
This follows by projecting $G^{-1/2}(Y-m_0)$ onto the span of the
rows of $DG^{1/2}$ and its orthogonal complement. If the original
integral uses $\delta(DY-a)$, it also contains the Gaussian density
of $DY$ evaluated at $a$, with determinant factor $(\det M)^{-1/2}$.
For a seam with $r_D$ rows the full scalar density factor is
$(2\pi)^{-r_D/2}(\det M)^{-1/2}
\exp[-(a-Dm_0)^{\mathsf T}M^{-1}(a-Dm_0)/2]$.
A conditionally normalized state divides by this factor. A source
defined with the delta insertion retains the factor and all of
its source derivatives. In the absence of a seam take $m=m_0$ and
$\Sigma=G$.

Suppose one term of the complete source has affine scores
$u(k)=U k+u_0$ and $v(k)=V k+v_0$ on this block. The first is a
real source and the second enters the geo-phase. The matrices
include the actual routing, the oriented visits and the effects of
composition. They need not factor into independent edge maps.
The elementary Gaussian identity is
\begin{equation}
 \mathbb E e^{z^{\mathsf T}Y}
   =e^{z^{\mathsf T}m+z^{\mathsf T}\Sigma z/2},
 \qquad z=u(k)+\ii v(k).
 \label{eq:uvpeN-full-Gaussian}
\end{equation}
The transpose is bilinear in $z$. The imaginary score contributes
$-v^{\mathsf T}\Sigma v/2$ to the real exponent, while the real score
contributes $u^{\mathsf T}\Sigma u/2$. The modulus of the full expression
is $\exp(c+h^{\mathsf T}k-k^{\mathsf T}Qk)$, with $c$ independent
of $k$ and
\begin{equation}
\begin{aligned}
 2Q&=V^{\mathsf T}\Sigma V-U^{\mathsf T}\Sigma U,\\
 h&=U^{\mathsf T}(m+\Sigma u_0)-V^{\mathsf T}\Sigma v_0 .
\end{aligned}
\label{eq:uvpeN-signed}
\end{equation}
A polynomial insertion in $Y,k$ gives this same complex exponential
times a polynomial in $k$, by differentiation in $z$. This includes
the polynomial marks from differentiating an affine source on a fixed
chart. A source-dependent quadratic action can also generate polynomial
normal insertions at a fixed derivative order. Its differentiated
determinant, boundary condition, domains and density must first be
accounted for. Distinct terms with the same exponential and denominator
can be combined before identifying their surviving polynomial.

Consider an out-and-back word $e^{-1}e$ of positive length. If
the two legs carry opposite scores on the same normal variable,
their product is $e^{\ii w(k)Y}e^{-\ii w(k)Y}=1$. If the source instead
measures two centered Gaussian variables $Y_1,Y_2$ of variance
$2t$ and correlation $\rho$, the same pair of signs yields
\[
 \mathbb E e^{\ii w(k)(Y_1-Y_2)}
     =\exp[-2t(1-\rho) w(k)^2].
\]
This example makes the distinction quantitative. Both words return
to the starting point. Damping along $w(k)$ follows from the variance
of the score difference, which vanishes at perfect correlation.
For fluctuating $t,\rho$ the graph must also control the region
where $t(1-\rho)$ is small. The represented source fixes which
variables are measured. The ordered word and the zero-visit
determinant must both be differentiated when they depend on the
photon probe. Taking the
absolute value before the prescribed normal integration would also
discard the Fourier mechanism in \eqref{eq:uvpeN-full-Gaussian}.
All absolute bounds below apply after that physical integration and
before integration in the remaining loop momenta. These real
Euclidean normal scores do not specify a Lorentzian pole displacement.

\subsection{A sharp criterion for polynomial source insertions}

Let $C_e:\RR^d\to\RR^{b_e}$ route the momentum in a massive
Euclidean factor, with exponent $\nu_e>0$ and real shift $p_e$.
The unit in its denominator fixes a reference mass; any fixed positive
mass changes only comparison constants. Suppose that the modulus of
a complete term, up to a finite positive constant, is
\begin{equation}
\begin{aligned}
 I(k)={}&|P(k)|e^{-k^{\mathsf T}Qk+h^{\mathsf T}k}\\[-2pt]
 &\times\prod_{e\in E}(1+\|C_e k+p_e\|^2)^{-\nu_e}.
\end{aligned}
\label{eq:uvpeN-absolute-model}
\end{equation}
Here $P$ is a nonzero real or complex polynomial. A fixed signed
tensor component can be tested separately. If this expression is
only a majorant, its integrability suffices for that term.

For a subspace $K$, define $n_K(P)$ to be the largest degree in $t$
of the nonzero polynomials $P(x+tw)$ over $x\in\RR^d$, $w\in K$.
Thus $n_{\{0\}}(P)=0$, and translation of the entire polynomial by
a fixed vector preserves $n_K(P)$. This degree counts specifically
those momentum powers in the source insertion which can grow along
$K$.

For each set $S$ in the theorem below, $K_S$ consists of the
momentum directions left undamped by the Gaussian and unseen by
the denominator factors outside $S$. The factors indexed by $S$
must therefore supply more decay than the dimension of $K_S$ plus
the polynomial growth along it. Testing every $S$ also covers
regions with several different momentum scales.

\begin{theorem}[Absolute integrability after the normal source integral]
\label{thm:uvpeN-directional-power}
For $d>0$ and real symmetric $Q$,
$\int_{\RR^d}I(k)\dd^dk<\infty$ precisely under the following
conditions. The quadratic form satisfies $Q\ge0$ and $h\perp\ker Q$.
The combined map with rows from $Q$ and all $C_e$ is injective.
For each $S\subseteq E$ put
\[
 K_S=\ker Q\cap\bigcap_{e\notin S}\ker C_e .
\]
For each $K_S$ of positive dimension one must have
\begin{equation}
 2\sum_{e\in S}\nu_e>
       \dim K_S+n_{K_S}(P).
 \label{eq:uvpeN-directional-count}
\end{equation}
Equality on any such subspace fails the absolute integrability test.
\end{theorem}
\begin{proof}
A negative eigenvalue of $Q$ gives exponential growth on an open
cone, on a smaller cone of which the leading nonzero homogeneous
part of $P$ is bounded below on unit directions. Each massive
denominator there costs at most a power. For $Q\ge0$, if $h$ is
nonzero on $\ker Q$, keep the orthogonal components bounded and
take a cone in the kernel where $h^{\mathsf T}k$ grows linearly.
The leading polynomial in the kernel coordinates is nonzero for an
open set of transverse coordinates and directions. This again forces
divergence. Under the two required sign and drift conditions one can
complete the square by a fixed translation. For every fixed
$p$ and $q$ one has $1+|q+p|^2\asymp 1+|q|^2$, with comparison
constants depending only on $|p|$. The translated shifts in
the massive factors can therefore be removed for this test. A common unseen
direction of $Q$ and every $C_e$ leaves all weights constant along
it, and the integral of the absolute polynomial on generic parallel
lines is infinite.

It remains to treat a centered Gaussian with $Q=A^{\mathsf T}A$
and an injective combined map. For $z\ge0$ and $\nu>0$,
\[
 (1+z)^{-\nu}\asymp\int_0^1t^{\nu-1}e^{-tz}\dd t ;
\]
this follows by rescaling at $z\ge1$ and comparing constants
at $z\le1$. Tonelli's theorem gives a Gaussian integral with
precision in the exponent $H=Q+\sum_e t_e C_e^{\mathsf T}C_e$.
Order $0<t_1\le\cdots\le t_m\le1$ on one of the finitely many
sectors, with $m=|E|$. Select linearly independent rows greedily
from $A,C_m,\ldots,C_1$ until they form a basis. Let $x_i$ be the
resulting linear coordinates and $w_i$ their weights ($1$ or $t_e$).
Uniformly on the sector,
\[
 k^{\mathsf T}Hk\asymp\sum_{i=1}^d w_i x_i^2 .
\]
Indeed each discarded row is a linear combination of earlier
selected rows, all with at least its weight. The converse inequality
uses just the selected rows.

In these fixed coordinates write $P=\sum_\gamma a_\gamma x^\gamma$.
For either constant rescaling of the diagonal Gaussian, change variables
to $z_i=\sqrt{w_i}x_i$. The integral of $|P|$ against that Gaussian
is comparable, uniformly in the weights, to
\[
 \Big(\prod_i w_i^{-1/2}\Big)
    \sum_{\gamma:a_\gamma\ne0}|a_\gamma|
                  \prod_i w_i^{-\gamma_i/2}.
\]
The Gaussian $L^1$ norm and the sum of absolute coefficients are two
norms on the same finite-dimensional polynomial space. This observation
retains possible cancellations in $P$ before its absolute value is
taken.

Put $t_{m+1}=1$, $r_j=t_j/t_{j+1}$, and $S_j=\{1,\ldots,j\}$.
In these coordinates $K_{S_j}$ consists of the weak coordinate
directions associated to rows from $C_1,\ldots,C_j$; all the other
coordinate values vanish on it. Denote this set of coordinate indices
by $W_j$. It has size $\dim K_{S_j}$. The measure factor becomes
$\prod_j r_j^{\sum_{e\in S_j}\nu_e}\dd r_j/r_j$.
For a surviving monomial $\gamma$ the diagonal Gaussian factor
subtracts $(|W_j|+\sum_{i\in W_j}\gamma_i)/2$ from that power.
The maximal such monomial degree is $n_{K_{S_j}}(P)$, since the
coefficients of a highest-degree term on this subspace are polynomials
in the complementary coordinates and can be tested at generic values.
Each resulting one-dimensional integral on $(0,1)$ requires a
strictly positive exponent; for $W_j=\varnothing$ this is automatic.
Every nonempty $S$ occurs as an initial segment in some ordering.
If there are no line factors and $Q>0$, ordinary Gaussian moments
give the result directly.
\end{proof}

\subsection{An exact cost for source marks at collapsing depth}

Consider a single constrained or unconstrained centered normal mode
of variance $2t$, $t>0$. Give it real score $c u$ and imaginary score
$u$, where $k=(u,v)\in\RR\times\RR^3$. The whole source in
\eqref{eq:uvpeN-full-Gaussian} has modulus $e^{-\tau u^2}$ with
$\tau=t(1-c^2)$. At $|c|=1$ its normalized phase relative to the
real-source partition function still has Gaussian modulus, but the
absolute complete source has unit modulus. Combined with the
marginal denominator $(1+|k|^2)^{-2}$ it gives a logarithmic UV
divergence. If $|c|>1$ it grows exponentially.

For $|c|<1$ the blind subspace is three-dimensional. Any insertion
$P=u^m$, $m\in\mathbb N_0$, has blind degree zero, so every such
integral converges at fixed $\tau>0$. In contrast $P=v_1$ saturates
\eqref{eq:uvpeN-directional-count} and gives a logarithm already
at fixed $\tau$. The degree of the completed source vertex along the
undetected directions changes this conclusion even for equal total
polynomial degrees.

For the detected marks integrate $v$ first:
\begin{equation}
 J_m(\tau)=\pi^2\int_{\RR}
   \frac{|u|^m e^{-\tau u^2}}{\sqrt{1+u^2}}\dd u .
 \label{eq:uvpeN-mark-cost}
\end{equation}
Radial integration in $\RR^3$ gives the factor
$\pi^2/\sqrt{1+u^2}$. As $\tau\downarrow0$,
\begin{equation}
\begin{aligned}
 J_0(\tau)&=\pi^2\log(1/\tau)+O(1),\\
 J_m(\tau)&\sim \pi^2\Gamma(m/2)\tau^{-m/2},
                    \qquad m>0 .
\end{aligned}
\label{eq:uvpeN-mark-threshold}
\end{equation}
The first formula follows by replacing $(1+u^2)^{-1/2}$ with $1/|u|$
on $|u|\ge1$ at bounded integrated error, then substituting
$z=\tau u^2$. For the second rescale $u$ by $\tau^{-1/2}$
and dominate by $|z|^{m-1}e^{-z^2}$ in the rescaled variable $z$.
For a finite positive common law on $0<\tau\le1$ this gives the exact
necessary and sufficient budgets for the respective positive integrals,
namely the logarithmic moment at $m=0$ and the inverse $m/2$ moment
at $m>0$. A zero atom yields an infinite contribution in both cases.
If differentiating the actual source also supplies powers of $\tau$,
those powers multiply \eqref{eq:uvpeN-mark-cost} before integrating
the law.

To use Theorem~\ref{thm:uvpeN-directional-power} in the physical
PE-CGG-HOUQFT expansion, one must compute the scores, the polynomial
and all residual line powers from the specified complete GCM source.
This includes the empty identity history, positive-length returns,
repeated visits, ghost and boundary contributions, and the allowed
contacts between supports. The theorem gives an exact decision within
its quadratic, affine-score class. Non-Gaussian histories require a
bound for their own Fourier response. If the matrices or weights
depend on an integrated normal modulus, pointwise finiteness must be
supplemented by a joint integrable majorant. If finite-dimensional
approximations are used to verify a directly defined coefficient,
their agreement with it requires dominated convergence or local
convergence with uniformly vanishing absolute tails. Neither step
permits a divergent subtraction. The costs in
\eqref{eq:uvpeN-mark-threshold} exhibit the extra normal integration
in a case with one common fluctuating mode.

\paragraph{Two marks on a returned path.}
The marked-bridge calculation in Part~XXVI~\cite{ChiapponePartXXVI2026}
shows which logarithmic moment to test when the source chooses both
the duration and the mark positions. Take a normalized four-dimensional
Euclidean heat bridge of diffusivity $D>0$ and duration $T>0$.
If the marks are separated by a fraction $v$ of the duration, their
relative displacement has covariance $2\sigma I_4$, where
$\sigma=DTv(1-v)$. Conditional on $T,v$, the two opposite momentum
insertions therefore contribute $e^{-\sigma|q|^2}$. For the massive
scalar bubble at zero external momentum, put $b=m^2\sigma$, $m>0$.
Equation~\eqref{eq:gaussian-zero}, with $2t=\sigma$, gives
\begin{equation}
 \begin{aligned}
 \int_{\mathbb R^4}\frac{\dd^4q}{(2\pi)^4}\,
     \frac{e^{-\sigma|q|^2}}{(|q|^2+m^2)^2}
   &=\frac{F(b)}{16\pi^2},\\
 F(b)&=(1+b)e^bE_1(b)-1\quad(b>0).
 \end{aligned}
 \label{eq:uv-joint-bridge-bubble}
\end{equation}
Here $E_1(b)=\int_b^\infty e^{-x}\dd x/x$. The radial integral makes
$F$ positive and decreasing, and the expansion at zero gives
$F(b)=\log(1/b)-1-\gamma_{\rm E}+o(1)$, with $\gamma_{\rm E}$ Euler's
constant. Thus for any joint probability law of $T,v$ with this
conditional increment, Tonelli's theorem gives the sharp criterion
\[
 \mathbb E[F(b)]<\infty
 \quad\Longleftrightarrow\quad
 \mathbb E[\log^+(1/b)]<\infty .
\]
Here $\log^+x=\max(0,\log x)$. Both costs are infinite at $b=0$,
including a collapsed bridge. An
atom there makes the average diverge, but a continuous distribution
of strictly positive variances can do so as well.

For example, take a uniform $u\in(0,1)$, a fixed $T_0>0$,
$T=T_0u$ and $v=e^{-1/u}$. The marks are distinct for every draw
and $\mathbb E[|\log(T/T_0)|]=1$. Yet $b\le m^2DT_0 e^{-1/u}$,
so $\log^+(1/b)\ge 1/u-\log(m^2DT_0)$ and its mean diverges.
This example concerns the positive scalar comparison after the
coherent Gaussian average. For the charged coefficient, the sewn
source must supply its own joint law, including the determinant,
contact and boundary weights. The next calculation identifies which
relative displacements that sewing preserves.

\subsection{Which charged displacements survive the boundary pairing?}
\label{sec:physical-uv-retained-minors}
A positive variance on every separate line does not determine the
variance seen by a closed charged graph. The following calculation
starts with the list of displacements in its charged matrix products
and with the equations actually imposed by the boundary state. It
tests all independent loop momenta and all simultaneous collapses of
the displacement variances. Work on one Euclidean sector-plane
$\bin=(s,b_s)$ and use fixed reference units.

Let $Y_e\in\mathbb R^4$, $1\le e\le n$, have independent centered
Gaussian increments of covariance $t_e I_4$, $t_e>0$, before sewing.
The boundary half-densities together impose $CY=0$ by an unnormalized
delta function in each Cartesian component, where $C$ has $c$
independent rows. Read the $L$ loop scores from a term of the charged
source: a factor $e^{i q\cdot Y_e}$ adds the row coefficients of the
actual momentum $q$ to the corresponding column of $J$; the reverse
phase changes the sign. Repeated uses of the same displacement add in
that same column. At zero external shift the combined phase is
$\exp(i\sum_{a,e}J_{ae} k_a\cdot Y_e)$.
Set $T=\operatorname{diag}(t_1,\ldots,t_n)$ and write
\[
 M=\begin{pmatrix} C\\J\end{pmatrix},\qquad
 B_A(t)=\det(ATA^{\mathsf T}),\qquad B_{\varnothing}=1 .
\]
Here $B_{\varnothing}$ refers to a matrix with no rows.
Let $K_h(t)$ be the integral of the absolute value of the sewn
Fourier coefficient against $(1+|k|^2)^h\,\mathrm d^{4L}k/(2\pi)^{4L}$,
where $h\ge0$ and $L\ge1$. Other action and measure factors will be
included below.

\begin{proposition}[Retained loop scores and simultaneous width limits]
\label{prop:physical-uv-retained-minors}
All loop directions retain Gaussian suppression precisely when
$\operatorname{rank}M=c+L$. At that rank, writing $m=c+L$, one has
\begin{equation}
 K_0(t)=(2\pi)^{-2m}B_M(t)^{-2}.
 \label{eq:physical-uv-retained-delta}
\end{equation}
If the rank is smaller, $K_0(t)=\infty$. For full row rank suppose
the remaining positive weight on $(0,1]^n$ is bounded above and
below by constant multiples of
$\prod_e t_e^{\beta_e-1}\,\mathrm dt_e$. This weight includes the
remaining vacuum and boundary densities of the coefficient under
test. For $S\subseteq\{1,\ldots,n\}$ put
\[
 \begin{aligned}
 d_S&=m-\operatorname{rank}M_{S^c},\\
 a_S&=\operatorname{rank}M_{S^c}-\operatorname{rank}C_{S^c},
 \end{aligned}
\]
with the subscripts retaining columns. The integrated $K_h$ is
finite if and only if every nonempty $S$ satisfies
\begin{equation}
 \sum_{e\in S}\beta_e>2d_S+h\,\mathbf 1_{\{a_S<L\}} .
 \label{eq:physical-uv-retained-faces}
\end{equation}
An upper bound by the stated weight suffices for convergence under
these inequalities.
\end{proposition}
\begin{proof}
For each Cartesian component, the conditional loop covariance is
the Schur complement $Q$ of $CTC^{\mathsf T}$ in $MTM^{\mathsf T}$.
Equivalently $v^{\mathsf T}Qv$ is the minimum over $u\in\mathbb R^c$
of $\|T^{1/2}(J^{\mathsf T}v-C^{\mathsf T}u)\|^2$. A hidden
direction is therefore a nonzero loop vector whose charged score
lies in the span of the seam rows. In particular, replacing $J$ by
$J+UC$ changes no phase on the sewn space. The delta insertion
contributes $(2\pi)^{-2c}B_C^{-2}$. Integrating the centered
characteristic over the $L$ four-momenta gives
$(2\pi)^{-2L}(\det Q)^{-2}$ when $Q>0$ and an infinite factor
otherwise. The identity $B_C\det Q=B_M$ gives
\eqref{eq:physical-uv-retained-delta}. In particular one cannot omit the
seam density when estimating the unnormalized source.

Let $M^{[-j]}$ delete the $j$th score row, leaving every seam row.
The lower diagonal block of $(MTM^{\mathsf T})^{-1}$ is $Q^{-1}$.
The Gaussian integral with the momentum weight consequently gives
\[
 K_h(t)\asymp K_0(t)
  \left(1+\sum_{j=1}^L\frac{B_{M^{[-j]}}(t)}{B_M(t)}\right)^h .
\]
Cofactors give the trace in this formula. The Gaussian moment is
bounded above by scaling with the largest eigenvalue of $Q^{-1}$
and below by retaining its one-dimensional eigendirection. The
constants depend on $L,h$.

Cauchy--Binet expresses each full-row-rank $B_A$ as a sum of basis
monomials $\prod_{e\in I}t_e$ with coefficients $(\det A_I)^2>0$.
Order the widths from smallest to largest. On this sector a maximal
monomial is obtained by selecting independent columns from largest to
smallest. The finite positive sum is comparable to that monomial.
For a prefix $S$ of the sorted columns, the number selected in $S$
is $\operatorname{rank}A-\operatorname{rank}A_{S^c}$.
Writing each width as a product of successive ratios (with final
denominator one) turns the integral into integrals of powers on
$(0,1]$. For the ratio corresponding to $S$, the weight gives
power $\sum_{e\in S}\beta_e$ with logarithmic differential. The
factor $B_M^{-2}$ costs $2d_S$ powers. For a given $j$, the cofactor
ratio costs one further power precisely when deleting score row $j$
leaves the rank on $S^c$ unchanged, and zero powers otherwise.
Such a row exists precisely when the score rows modulo the seam rows
on $S^c$ are dependent, that is, $a_S<L$. The $h$th power of the
finite sum is comparable to one plus the sum of the individual
$h$th powers. This gives \eqref{eq:physical-uv-retained-faces}. Every
nonempty subset occurs as a prefix, and equality for one prefix leaves
a divergent logarithmic integral with the other ratios bounded
away from zero.
\end{proof}

To see how the score rows arise, consider two finite charged
Berezin blocks, one per tested carrier momentum. Each has matrix
$\left(\begin{smallmatrix}1&W\\-\overline W&1\end{smallmatrix}\right)$
and positive determinant $1+|W|^2$. For the first take
$W=a e^{ik\cdot Y_1}+b e^{ik\cdot Y_2}$ with real probes $a,b$.
For the second take independent probes and use $Y_2,Y_3$ with
momentum $l$. The full mixed coefficient, differentiated once in
each of the four probes at zero, is
$4\cos(k\cdot(Y_1-Y_2))\cos(l\cdot(Y_2-Y_3))$ before sewing.
After sewing it is a positive sum of Gaussian characteristics. Their
score rows are $(1,-1,0)$ and $\pm(0,1,-1)$.
A return constraint $C=(1,1,1)$ gives
$B_M=9t_1t_2t_3$ for either sign. Thus the unweighted two-loop
Fourier test needs each $\beta_e>2$. With a single shared width
$t$ and residual measure $t^{\beta-1}\,\mathrm dt$, the same
coefficient instead requires $\beta>6$. Conditional normalization
at each width would remove the delta density and change the latter
threshold to $\beta>4$. If the boundary state identifies the outer
observations by $C=(1,0,-1)$, one score is the negative of the
other on the seam. Along $k=l$ one orientation then retains no
suppression, although each individual difference still fluctuates.

A boundary modulus can cause rank loss in the limit even when all
positive-width configurations pass the rank test. In the same
two-block source use $t_e=t$ and a normalized seam row
\[
 C_\theta=\cos\theta\,(1,0,-1)/\sqrt 2
             +\sin\theta\,(1,1,1)/\sqrt 3 .
\]
Its variance per Cartesian component is $t$. For the plus orientation
$JJ^{\mathsf T}$ has eigenvalues $1,3$, with eigenvectors proportional
to $(1,1)$ and $(1,-1)$. Conditioning subtracts
$t\cos^2\theta$ on the first eigenspace, giving loop covariances
$t\sin^2\theta$ and $3t$. Reversing the second score conjugates the
matrix by a sign change and preserves these eigenvalues. Consequently,
for $0<t\le1$ and $0<\theta\le\pi/2$, the cost including the
unnormalized seam obeys
\begin{equation}
 K_h(t,\theta)\asymp
          t^{-6-h}(\sin\theta)^{-4-2h}.
 \label{eq:physical-uv-moving-seam}
\end{equation}
For example, if the actual boundary preparation has
$\sin\theta\asymp t^\delta$, $\delta\ge0$, near zero and its remaining
density is comparable to $t^{\beta-1}\,\mathrm dt$, this endpoint
is integrable exactly when $\beta>6+h+(4+2h)\delta$.
There is full retention at each $t>0$ in this family. Its rate of
loss, together with the boundary weight, decides the integrated result.

This finite charged test concerns a specified coefficient. For a
physical graph, the rows must be read from its determinant or open
minor after all contact terms and boundary identifications are made.
Polynomial momentum marks dominated by $(1+|k|^2)^h$ use the
corresponding $h$. Extra inverse powers of the widths in a contact
reduce its measure exponents. A rank defect leaves ordinary
propagator powers, other genuine source factors, and cancellations to
be tested on the blind directions. It disproves full Gaussian
retention on this chart, without deciding that every such physical
graph diverges.

\subsection{An exact geo-phase integral for a network of normal visits}
\label{sec:uv-network-phase}

The common phase covariance and its normalization can both be computed
from a local geometric action on a whole network. Work on one Euclidean
sector-plane in fixed units. In this comparison the normal displacement
has four real components $x$ to which the carrier Fourier probes couple,
and one transverse component $y$ whose position is sampled at the visits.
These five auxiliary coordinates specify the representation of the
normal transfer. The Dirichlet form $\int (|\nabla f|^2+m_e^2|f|^2)$
on $\mathbb R^5$, with $m_e>0$, gives a transfer over duration $s_e>0$
with kernel
\[
 (4\pi s_e)^{-5/2}\exp[-m_e^2s_e-(|\Delta x_e|^2+|\Delta y_e|^2)/(4s_e)] .
\]
Integration in $s_e$ uses the entire massive resolvent.

Take a finite connected graph of visits with one rooted vertex $x_0=0$
and $n\ge1$ other vertices. At each vertex shared by several edges
use one displacement. For prescribed transverse positions $y_v$, put
$z_e=|\Delta y_e|$. Give the remaining vertices the Fourier source
$\exp(i(Bk+b)\cdot x)$, where $x$ stacks the $n$ displacements, $k$
stacks the independent loop momenta, and $B,b$ specify their actual
routing and external shift. The variables $x$ and $k$ are paired in a
fixed Euclidean frame of the representation. Denote by $F_s(k)$ the integral of this
phase against the product of all edge kernels in the $n$ shared $x_v$.
The edge powers $(4\pi s_e)^{-5/2}$ belong to this integral, including
on a cycle.

\begin{proposition}[Integrated cost of the marked normal network]
\label{prop:uv-network-phase}
For independent visit probes $Bk=k\in\mathbb R^{4n}$ and $b=0$, the
uncut absolute integral over all loop momenta and positive durations is
\[
 \int\!\prod_e\mathrm ds_e
 \int\frac{\mathrm d^{4n}k}{(2\pi)^{4n}}\,|F_s(k)|
 =\prod_e \frac{(1+m_e z_e)e^{-m_e z_e}}{8\pi^2 z_e^3}.
\]
A zero $z_e$ gives an infinite factor. Thus for any finite positive
common law of the positions, the expectation of the displayed
product is the exact finiteness test for these probes. In particular,
a closed even polygon with $N$ edges alternating between $y=0$ and
$y=r$, with $0<r\le1$, costs precisely an inverse $3N$-th moment
up to positive graph-dependent constants. The total displacement
of each such return is zero.
\end{proposition}

\begin{proof}
Delete the root column from the oriented edge--vertex incidence matrix
to obtain $D$. Then $H=D^{\mathsf T}\operatorname{diag}(s_e^{-1})D$
is positive definite. Direct Gaussian integration in all $x_v$ gives
\[
 \frac{F_s(k)}{F_s(0)\big|_{b=0}}
 =\exp[-(Bk+b)^{\mathsf T}(H^{-1}\otimes I_4)(Bk+b)] .
\]
In particular $F_s$ is nonnegative. Its prefactor is the product
of the edge kernels evaluated at $\Delta x_e=0$, multiplied by
$(4\pi)^{2n}(\det H)^{-2}$. For independent visit probes the
momentum integral is Fourier inversion at $x=0$ and cancels this last
factor exactly. Tonelli then separates the duration integrals. Each
one equals the displayed five-dimensional Green function by differentiating
$2\sqrt\pi\,e^{-mz}/z=\int_0^\infty s^{-3/2}
 e^{-m^2s-z^2/(4s)}\mathrm ds$ in $z^2$ for $z>0$.
At $z=0$ the $s^{-5/2}$ integral diverges at zero. For the alternating
polygon all massive exponential factors are bounded above and below
by positive constants on $0<r\le1$.
\end{proof}

The same action gives a sufficient estimate for routed momenta and
source contacts. Suppose $B$ is fixed with full column rank, $b$ stays
bounded, and every edge satisfies $z_e\ge c_e r$ with $c_e>0$.
Allow fixed-order insertions polynomial in $x,k,s_e,s_e^{-1}$ with
coefficients bounded by a finite inverse power of $r$. Put
$s_{\min}=\min_e s_e$, $s_{\max}=\max_e s_e$.
Since $D$ has full column rank, $H^{-1}$ lies between
$c s_{\min}I$ and $C s_{\max}I$ as quadratic forms. A displacement
mark comes from differentiating the common characteristic in its
Fourier argument. This costs powers of $1+s_{\max}$ and $1+|k|$.
The Gaussian in $Bk+b$ bounds the momentum integral of any such
polynomial by a fixed power of $1+s_{\min}^{-1}$, uniformly in
the shift: $|Bk+b|\ge s_{\mathrm{low}}(B)|k|-|b|$, with
$s_{\mathrm{low}}(B)>0$ the smallest singular value.
The determinant prefactor costs at most $C s_{\max}^{2n}$ apart
from the displayed edge factors. Every resulting proper-time power
is controlled at infinity by $e^{-m_e^2s_e}$ and at zero by
$e^{-c_e^2r^2/(4s_e)}$. Indeed, for $0<s_e\le1$, the latter is at
most $C_M s_e^M r^{-2M}$ for any fixed $M>0$. The joint absolute
cost is therefore at most $C r^{-P}$ for some finite $P$ depending
on the graph and the insertion.

Consequently a common boundary measure with all inverse moments
controls every finite graph and polynomial contact of this type.
For a signed or complex residual measure it suffices to use its total
variation. For example, a normal ground-state density bounded by
$Cr^d e^{-a/r^p}$ with $a,p>0$ has all these moments. The bound covers simultaneous hard and
short-duration regions: inserting $(1+|k|)^a\prod_{e\in E_0}s_e^{-b_e}$
in the absolute integral costs only another fixed inverse moment, for
arbitrary fixed $a,b_e\ge0$ and any edge subset $E_0$.
This calculation integrates the physical geo-phase before taking
the absolute value in momentum. It also identifies two ways a
proposed application can fail. A common unseen momentum direction
in $\ker B$ has no Gaussian suppression, and a contact that adds
an independently probed edge at $z_e=0$ encounters the singular
Green function above. Applying the test to a charged GCM source
requires its actual visit routing, contact rules, and the complete
determinant-weighted boundary law.

\subsection{One shared normal score for several marginal loops}
\label{subsec:uvpeN-many-returns}

A sharper distinction appears when several histories use the same normal
variable. Suppose the source gives one real centered Gaussian of variance
$2\tau$ and phase score $s=\sum_{j=1}^L c_j u_j$, with $c_j\ne0$,
$u_j=(k_j)_1$, $k_j\in\RR^4$. Let the positive envelope have $L$
marginal factors $(1+|k_j|^2)^{-2}$ and a mark $|s|^m$, $m\ge0$.
There is only one normal integral for the whole configuration.
At each $\tau>0$ this envelope is integrable. Indeed
$\dim K_S=4|S|-1$ for nonempty $S$ and $s$ vanishes on $K_S$.
For integer $m$ this is an immediate application of
Theorem~\ref{thm:uvpeN-directional-power}; the direct argument below
also treats real $m\ge0$.

Write $J_{L,m}$ for the integral of this envelope with $e^{-\tau s^2}$.
For $0<\tau\le1$ put $\ell_\tau=1+\log(1/\tau)$. Then
\begin{equation}
 J_{L,m}(\tau)\asymp
 \begin{cases}
  \ell_\tau^L,&m=0,\\
  \tau^{-m/2}\ell_\tau^{L-1},&m>0 .
 \end{cases}
 \label{eq:uvpeN-many-mark-budget}
\end{equation}
To prove this, integrate the three orthogonal components of each loop
as in \eqref{eq:uvpeN-mark-cost}. In the variable $c_j u_j$ the
resulting density is proportional to $(c_j^2+x^2)^{-1/2}$ and hence
comparable to $f(x)=(1+|x|)^{-1}$. The density of the sum is therefore
bounded above and below by constants times the $L$-fold convolution
$f^{*L}(s)$. This convolution exists and satisfies
\[
 f^{*L}(s)\asymp
       \frac{[\log(2+|s|)]^{L-1}}{1+|s|}.
\]
For the induction step let $R=|s|\ge2$. The region $|x|\le R/2$
in the convolution of $f(x)$ with $f^{*L}(s-x)$ supplies a term of
order $(\log(2+R))^L/R$. The region $|s-x|\le R/2$ has the same
upper bound by integrating the induction estimate in $s-x$.
On the remainder with $|x|\le2R$ both denominators are of order $R$;
on $|x|>2R$ integrate a constant times
$[\log(2+|x|)]^{L-1}/|x|^2$. Bounded $s$ gives a positive finite
density by the same tail comparison. This proves the convolution
formula. Integrating it against $|s|^m e^{-\tau s^2}$ gives
\eqref{eq:uvpeN-many-mark-budget}, by splitting at
$|s|=\tau^{-1/2}$ for $m=0$ and rescaling there for $m>0$.
Integration of its right-hand side is the exact moment condition on
a finite positive depth law supported in $(0,1]$.

Source differentiation can instead mark an individual visit. For
$L\ge2$, replacing $|s|$ by $|u_1|$ already gives divergence at
every fixed $\tau>0$. Take $S=\{1,2\}$ in
\eqref{eq:uvpeN-directional-count}: the constraint
$c_1u_1+c_2u_2=0$ leaves seven undetected coordinates, and $u_1$
has degree one on them. The two marginal factors give exactly
$2(2+2)=8$. The absolute integral therefore diverges at equality.
The damped unmarked words, and even their detected source derivatives,
cannot decide this contact without determining which individual
momentum the differentiated native history measures.

\paragraph{A connected charged-source test on opposite hard momenta}
\phantomsection\label{dualrev-uvb-shared}
A cancellation in the routed frequency can survive even after all
direct and sequential source derivatives have been grouped. Consider
one scalar routing component on a fixed locked support. On two
pinned planar magnetic returns with independent unit-uniform durations
$t_1,t_2$, choose units in which the conditional area characteristic
is $f(t_j x)$, $f(u)=u/\sinh u$, $f(0)=1$. This is the ratio of the
constant-field diagonal heat kernel to its free value, from the
scalar Landau sum used in Section~\ref{sec:pe-curvature}.
Take the two areas conditionally independent and let $X$ be their sum.
Its characteristic is $\chi(x)=F(x)^2$ with
$F(x)=\int_0^1 f(tx)\dd t$.

Now let two charged blocks of momenta $k,l\in\RR$ see this
\emph{same} $X$. Each has two Grassmann slots, diagonal entries one,
and off-diagonal entries whose product is
$\rho e^{\ii(k+a)X}$ or $\rho e^{\ii(l+a)X}$ respectively, with
$\rho=1/4$. The real probe $a$ couples to the area. At fixed
history the two determinants multiply. Integrating them against the
common boundary state gives
\[
\begin{aligned}
 Z_{k,l}(a)&=1-\rho[\chi(k+a)+\chi(l+a)]\\
           &\hspace{6mm}+\rho^2\chi(k+l+2a),\\
 Z_k(a)&=1-\rho\chi(k+a).
\end{aligned}
\]
The two-block connected log response is the second derivative of
$C_{k,l}=\log Z_{k,l}-\log Z_k-\log Z_l$. All logarithms are
regular at real $a$ because $0\le\chi\le1$ and
$Z_{k,l}\ge1-2\rho$. For $n=0,1,2$, changing variables $u=t|x|$
bounds $|F^{(n)}(x)|$ by
$|x|^{-n-1}\int_0^\infty u^n |f^{(n)}(u)|\dd u$ away from zero.
Thus, uniformly on each compact interval of $y$,
\[
 \lim_{|k|\to\infty} C_{k,-k+y}''(0)
   =4\frac{\dd^2}{\dd y^2}\log(1+\rho^2\chi(y)).
\]
At $y=0$, $f''(0)=-1/3$ and $\mathbb E t_j^2=1/3$ give
$\chi''(0)=-2/9$, hence the limit is $-8/153$. Continuity supplies
a strip $|k+l|<\varepsilon$ of fixed positive width on which the
absolute contact remains bounded away from zero at large $|k|$.
Its integral with unit remaining momentum weight diverges. On an
integer lattice the line $l=-k$ suffices. This calculation includes
the integration over both durations down to zero. A different graph
can carry additional propagator decay on that strip. Identifying the
remaining weight and the routing of the complete charged source is
therefore part of a graphwise UV proof.

\section{Two tests of the GCM curvature action}
\label{sec:pe-curvature}

The curvature square permits a calculation of the normal weight as
well as the transfer on an internal line. We give two choices of
representation, trace and state on one locked support
$\bin=(s,b_s)$. One gives an interacting theory in which the UV cutoff
can be removed. The other retains a divergent magnetic response after
combining the gauge and ghost fluctuations. Both tests are Euclidean.
The normal coordinate measures a separation and carries neither a
clock orientation nor a pole sign.

\subsection{An interacting adjoint-scalar sector}
\label{sec:uvpeC-adjoint}

Use a four-torus of period $2\pi$ in dimensionless reference units,
with normalized volume and Fourier modes $k\in\mathbb Z^4$.
Let $X,Y$ be real fields valued in a compact non-Abelian Lie
algebra of dimension $n$, with positive invariant inner product
from a faithful unitary representation. Fix their zero modes to
zero. At normal separation $\ell\ge0$, the bounded harmonic
extension of a boundary mode solves
$u''=|k|^2u$, $u(0)=u_k$. Thus its trace at $\ell$ is $P_\ell u$,
where $P_\ell=\exp(-\ell\sqrt{-\Delta})$. For $\ell>0$ the
underlying geo-phase average is
\begin{equation}
 e^{-\ell |k|}
 =\int_0^\infty
 \frac{\ell e^{-\ell^2/(4t)}}{2\sqrt\pi\,t^{3/2}}
 e^{-t|k|^2}\,\dd t .
 \label{eq:uvpeC-phase}
\end{equation}
The last exponential is the Fourier transform of the periodized
tangential Gaussian displacement of covariance $2tI$. The density
in front is the first-passage law at distance $\ell$ of a normal
diffusion with generator $\partial_r^2$. Equation~\eqref{eq:uvpeC-phase}
also follows by the elementary Gaussian integral.

Use two anticommuting internal one-forms $\theta^1,\theta^2$ with
$\dd\theta^i=0$. The connection is block diagonal, with blocks
$\dd+X\theta^1$, $\dd+Y\theta^2$ and
$\dd+(P_\ell X)\theta^1+(P_\ell Y)\theta^2$.
Take the ordinary trace on the even internal blocks and orthonormal
coefficient pairings on two-forms. With $U=P_\ell X$, $W=P_\ell Y$,
the integrated curvature square is $S_\ell=Q_\ell+V_\ell$, where
\begin{align}
 Q_\ell={}&\tfrac12\!\int\!
  \left(|\dd X|^2+|\dd Y|^2+|\dd U|^2+|\dd W|^2\right)\dd x ,
 \nonumber\\
 V_\ell={}&\tfrac12\int\!\|[U,W]\|^2\,\dd x .
 \label{eq:uvpeC-sector-action}
\end{align}
Only the third block has a $\theta^1\wedge\theta^2$ component.
Its commutator gives the interacting vertex, while each copied
kinetic term changes the Gaussian normalization at depth $\ell$.
The fluctuating variables are $X,Y$ and one common $\ell$. The
harmonic extension is determined by the boundary values. The
finite positive input measure $\mu(\dd\ell)$ on $[0,\infty)$
satisfies $\mu((0,\infty))>0$. Expand in a marker $g$ multiplying
$V_\ell$ and keep the complete quadratic form exact. Its curvature
normalization is $g=1$. Use one normalization by the Gaussian
integral of the boundary-only kinetic terms. This constant is
independent of $\ell$ and of any external probes.

\begin{theorem}[Graphwise convergence from the curvature determinant]
\label{thm:uvpeC-curvature}
For \eqref{eq:uvpeC-sector-action}, every interacting Wick graph
at each fixed order in $g$ is absolutely UV finite on the stated
torus, including vacuum graphs and insertions of elementary fields
at finitely many fixed Fourier modes. The integrated absolute
momentum tail beyond $M$ is $O_s(M^{-s})$ for each fixed $s>0$.
The same estimate holds uniformly for the spherical Fourier
restrictions \(0<|k|\le\Lambda\) used in the proof.
The law $\mu$ may have an atom at zero, whose complete relative
Gaussian weight vanishes in the unrestricted source.
\end{theorem}

\begin{proof}
At each $\ell>0$, $e^{-2\ell\sqrt{-\Delta}}$ is trace class on the
mean-zero torus modes. The complete relative quadratic weight is
therefore the ordinary inverse Fredholm determinant power
\[
 H_\infty(\ell)=\prod_{k\ne0}(1+e^{-2\ell|k|})^{-n}>0.
\]
It is retained once for the common depth. The full-mode covariance
and this determinant define the graph directly. To verify its absolute
tail and the contribution at $\ell=0$, also restrict symmetrically to
$0<|k|\le\Lambda$. Integrating the same $Q_\ell$ then gives
the determinant and elementary covariance
\begin{align}
 H_\Lambda(\ell)&=\prod_{0<|k|\le\Lambda}
         (1+e^{-2\ell|k|})^{-n},\nonumber\\
 C_\ell(k)&=\bigl[|k|^2(1+e^{-2\ell|k|})\bigr]^{-1}.
 \label{eq:uvpeC-determinant}
\end{align}
The signed lattice counts one real sine and one real cosine
coordinate per opposite pair. The two independent adjoint fields
give the power $-n$. Each internal contraction at the interaction
carries $K_\ell(k)=e^{-2\ell|k|}C_\ell(k)$, including both ends of
a tadpole. There is one determinant per common depth.

Set $R=\min(\Lambda,\ell^{-1})$, taking $R=\Lambda$ at zero.
For $R\ge2$ the ball of radius $R$ contains at least $cR^4$
nonzero modes and $e^{-2\ell|k|}\ge e^{-2}$ there. Hence
$H_\Lambda(\ell)\le\exp(-c'R^4)$. On dyadic shells in four
dimensions, $C_\ell(k)\le|k|^{-2}$ and
\begin{equation}
 \sum_{0<|k|\le\Lambda}(1+|k|)^s K_\ell(k)
 \le C_s(1+R)^{2+s}.
 \label{eq:uvpeC-line-estimate}
\end{equation}
Below $R$ sum the power bound. Above $R$, when such shells occur,
$R=\ell^{-1}$ and their exponential tail is summable. The bound
also covers $R<2$.

For a fixed graph with $I$ internal interaction lines, discard the
momentum-conservation indicators in the absolute majorant. Its
weighted sum, with weight $(1+\max_e|k_e|)^s$, is at most
$C_{G,s}H_\Lambda(\ell)(1+R)^{(2+s)I}$. The color sums are finite,
and external elementary covariances at fixed nonzero modes are bounded.
The displayed product is bounded uniformly in $\ell,\Lambda$ and
tends to zero at $\ell=0$ as $\Lambda\to\infty$. Integration against
$\mu$ and restriction to $\max_e |k_e|>M$ give the stated tail.
Graphs without an internal line require no tail estimate. For
$\ell>0$, the same weighted estimate holds directly with
$R=\ell^{-1}$ for the full mode sum and the already defined
$H_\infty(\ell)$. At zero separation the additional quadratic
form relative to the boundary Gaussian is a sum of infinitely many
independent positive squares and is infinite almost surely; its
relative weight is therefore zero. Set $H_\infty(0)=0$.
The finite-mode estimate shows that this zero weight persists
with each fixed graph inserted. It suppresses the zero-depth atom
through the complete vacuum weight, without deleting
positive-length return histories. At positive depth, convergence on
finite mode sets and the uniform absolute tail identify the direct
full-mode graph with the limit of the Fourier restrictions.
The single zero-marker normalization is $\int H_\infty\dd\mu>0$,
as $\mu$ has positive mass at $\ell>0$. Division by this number and
the finite products forming connected coefficients are therefore
allowed.
\end{proof}

The mode count in the determinant is essential for the uniform
approximation estimate. For instance, retain only the four modes
\(\{\pm e_1,\pm N e_2\}\) and give the incoming law positive mass
both at zero and at a positive depth. At zero the restricted
determinant is \(2^{-4n}\), independent of \(N\). In the
first-order quartic vacuum graph the contraction with one low pair
and one hard pair has positive size proportional to \(N^{-2}\) for
a non-Abelian algebra. At threshold \(M=N/2\) it prevents a bound
uniform in such sparse projections for any tail power greater than
two. The unrestricted determinant and its spherical approximants
retain all modes in the low-frequency ball.

The strength of the short-visit suppression is calculable:
\begin{align}
 \lim_{\ell\downarrow0}\ell^4[-\log H_\infty(\ell)]
 &=2\pi^2 n\!\int_0^\infty\!r^3\log(1+e^{-2r})\,\dd r
 \nonumber\\
 &=\tfrac{45\pi^2}{64}\,n\zeta(5).
 \label{eq:uvpeC-small-depth}
\end{align}
The lattice sum converges to the radial integral after rescaling.
Expanding the logarithm gives an absolutely summable sequence of
integrals. The same rescaling in the edge sum gives
\begin{equation}
 \lim_{\ell\downarrow0}\ell^2\sum_{k\ne0}K_\ell(k)
 =2\pi^2\int_0^\infty\frac{r\,\dd r}{e^{2r}+1}
 =\frac{\pi^4}{24}.
 \label{eq:uvpeC-tadpole}
\end{equation}
The singularity $|u|^{-2}$ at the origin of the rescaled
four-dimensional integrand is controlled uniformly by summing
dyadic shells up to $|u|=\delta$, at cost $O(\delta^2)$.

These constants expose a normalization error that would change the
answer. At fixed $\ell>0$, choose an orthonormal basis $T_a$ in the
coefficient algebra. The Gaussian expectation of $V_\ell$ per
normalized torus volume is exactly
\begin{equation}
 \mathbb E_\ell V_\ell=
 \tfrac12\sum_{a,b}\|[T_a,T_b]\|^2
 \left(\sum_{k\ne0}K_\ell(k)\right)^2 .
 \label{eq:uvpeC-vacuum}
\end{equation}
Independence of $X,Y$ gives the two equal contractions. The color
factor is strictly positive and \eqref{eq:uvpeC-tadpole} gives a
cost of order $\ell^{-4}$. Normalizing separately at each depth
would remove $H_\Lambda$ and leave this graph divergent for
$\mu(\dd\ell)=\one_{[0,1]}\dd\ell$. If additional quantum fields
cancel the determinant, the full action must supply another bound
for this term.

\begin{corollary}[Finite coupling on the same compact volume]
\label{cor:uvpeC-coupling}
Let $\mu$ be the finite positive depth measure specified above, with
$\mu((0,\infty))>0$. For $g\ge0$, the same source defines a finite positive continuum
partition function and normalized distributions of the elementary
fields on the stated torus. Their polynomial moments at finitely many
fixed nonzero Fourier modes are finite. The corresponding finite
Fourier restrictions converge to these distributions and moments.
\end{corollary}
\begin{proof}
For $\ell>0$, let $\mathbb P_\ell$ be the product Gaussian law of
the two real adjoint fields with covariance $C_\ell(k)$ from
\eqref{eq:uvpeC-determinant}. This is a law on distributional fields.
Their harmonic transfers $U,W$ are almost surely smooth: every
derivative has an absolutely convergent Fourier series because
its expected absolute sum is bounded by a constant times
$\sum_{k\ne0}(1+|k|)^q e^{-\ell|k|}/|k|<\infty$ for each finite
derivative order $q$. In particular $V_\ell$ is finite and
nonnegative. Define directly
\[
\begin{aligned}
 Z_g={}&\int_{(0,\infty)}
      H_\infty(\ell)\,
      \mathbb E_\ell[e^{-gV_\ell}]\,\mu(\dd\ell),\\
 \dd\mathbb P_g={}&Z_g^{-1}H_\infty(\ell)e^{-gV_\ell}\\
 &\quad{}\times\mu(\dd\ell)\,\dd\mathbb P_\ell .
\end{aligned}
\]
The second formula is on $\ell>0$; the atom at zero has zero
weight. Since $\mu((0,\infty))>0$, all positive-depth factors
are strictly positive and $0<Z_g\le\mu((0,\infty))<\infty$.
The covariance bound $C_\ell(k)\le |k|^{-2}$ gives an integrable
Gaussian majorant for every polynomial in finitely many fixed modes.

For the approximation check, couple all Fourier restrictions by the
same standard real Gaussian coordinates. At $\ell>0$ the transferred
fields converge uniformly almost surely, so
$V_{\ell,\Lambda}\to V_\ell$. The density
$H_\Lambda e^{-gV_{\ell,\Lambda}}$ is at most one; at $\ell=0$
it tends to zero since $H_\Lambda(0)\to0$. Dominated convergence
therefore gives an $L^1$ limit with positive mass and convergence of
the normalized finite-mode distributions. The Gaussian polynomial
majorant gives convergence of their moments. The same argument gives
continuity in real $g\ge0$; it does not assert convergence of the
perturbation series in $g$.
\end{proof}

This construction keeps a global adjoint symmetry and fixes the
zero modes. Deriving a fluctuating physical gauge sector, a chiral
determinant, reflection positivity and an infinite-volume limit
requires further input. Neither the perturbative estimate nor
Corollary~\ref{cor:uvpeC-coupling} addresses summability of the
series in $g$.

\subsection{Magnetic response of a charged gauge-root pair}
\label{sec:uvpeC-magnetic}

For a different even GCM specialization take compact
$\mathfrak{su}(2)$ with positive invariant pairing on
$\mathbb R^4\times\mathbb R_r$. Set $\mathbb D=\dd+A$ and freeze
the shell and coarea factors at their common constant. The
curvature gives $S=\tfrac14\int(F_{MN},F_{MN})\dd^4x\,\dd r$.
Choose $[T_3,T_\pm]=\pm\ii T_\pm$ in the complexified coefficient
algebra and the stationary Cartan background
$\overline F_{12}=BT_3$, $\overline A_r=vT_3$. Consider first
the model of $r$-independent fields and gauge transformations, per
unit normal length. It retains four vector components and the
fluctuation of the normal component. Restricting in this way
defines a theory on the carrier; a constant field on the whole
normal line would have infinite squared norm.

\begin{proposition}[An unsuppressed transverse gauge response]
\label{prop:uvpeC-root}
In background Feynman gauge the complete root-pair contribution to
the bare one-loop background action, per unit four-volume on the
UV proper-time interval $\varepsilon<t<T$, satisfies
\begin{equation}
 [B^2]\bigl(\Gamma^{(1)}(B)-\Gamma^{(1)}(0)\bigr)
 =-\frac{7}{2(4\pi)^2}
 \int_\varepsilon^T e^{-tv^2}\frac{\dd t}{t}.
 \label{eq:uvpeC-root-response}
\end{equation}
The divergence persists for any fixed $v$ even after the normal
polarization and Faddeev--Popov ghost have been included.
\end{proposition}

\begin{proof}
Expansion of $F(\overline A+a)=\overline F+\overline D a+a\wedge a$
and addition of $\tfrac12\|\overline D_Ma_M\|^2$ give
\[
 (H_1)_{MN}=-\overline D^2\delta_{MN}
       -2\operatorname{ad}(\overline F_{MN}),\qquad
 H_0=-\overline D^2 .
\]
The weights of the vector and ghost traces are $1/2$ and $-1$,
respectively~\cite{Abbott1981}. On a root the normal kinetic
term contributes $v^2$. The Landau levels $(2j+1)|B|$, of
degeneracy $|B|/(2\pi)$ per unit area, give the scalar heat
density $(4\pi t)^{-2}x/\sinh x$ in four dimensions, $x=Bt$.
The vector spin matrix has two shifts $\pm2B$ and three zero
shifts. Summing the two roots and both statistics, the
$B$-dependent heat-trace difference is
\begin{equation}
 \frac{e^{-tv^2}}{(4\pi t)^2}
 \left[\frac{x}{\sinh x}(1+2\cosh(2x))-3\right].
 \label{eq:uvpeC-gauge-heat}
\end{equation}
The bracket equals $7x^2/2+O(x^4)$. Integrating with $-\dd t/t$
proves \eqref{eq:uvpeC-root-response}. Taking $T$ fixed isolates
the UV coefficient independently of large-time magnetic stability.
\end{proof}

A transverse pair of probes at finite wavelength gives the same
logarithm. Let $A_\pm=e_2T_3e^{\pm\ii p x^1}$, $p\ne0$, with
real source derivatives obtained by polarization.
Put $I(z)=\int_0^1e^{-z h(1-h)}\dd h$, $z=tp^2$. Duhamel
differentiation at zero magnetic field gives the mixed heat response
\begin{equation}
 \frac{e^{-tv^2}}{(4\pi t)^2}
 \left[6t\{I(z)-1\}+8t^2p^2 I(z)\right].
 \label{eq:uvpeC-mixed-heat}
\end{equation}
For one scalar component, the two-source contact contributes $-2t$.
The linear pair contributes $2tI(z)$ by Gaussian translation of
the loop momentum and its perpendicular variance $1/(2t)$. Their
orbital multiplicity after the ghost trace is three. The spin
insertions have vanishing individual polarization trace and
product trace $8p^2$. Their two-root statistical multiplicity is
one. With $I(z)=1-z/6+O(z^2)$ the bracket is
$7t^2p^2+O(t^3p^4)$. Thus field-independent normalization cannot
absorb this divergence. Even averaging $v$ over one
source-independent probability law on finite values leaves the
leading logarithm, since $\mathbb E(e^{-tv^2})\to1$ as $t\downarrow0$.

There is a useful distinction on the unrestricted normal line.
On a root the normal heat kernel evaluated at displacement $d$ is
\begin{align}
 \int_{\mathbb R}\frac{\dd\xi}{2\pi}
   e^{-t(\xi\pm v)^2+\ii\xi d}
 &=\frac{e^{\mp\ii vd}}{\sqrt{4\pi t}}e^{-d^2/(4t)} .
 \label{eq:uvpeC-offdiagonal}
\end{align}
Parallel transport removes the open phase of the flat connection.
At $B=0$ the transported resolvent at nonzero $|k|$ is
$e^{-|d||k|}/(2|k|)$. The geometric attenuation therefore arises
in an off-diagonal matrix element. The ordinary heat trace uses
$d=0$. For a transported magnetic insertion the remaining
small-$t$ cost obeys, at $d\ne0$,
\begin{equation}
 \int_0^T\frac{e^{-d^2/(4t)}}{\sqrt{4\pi t}}
                      \frac{\dd t}{t}
 =\frac{\erfc(|d|/(2\sqrt T))}{|d|}.
 \label{eq:uvpeC-displacement-cost}
\end{equation}
A finite positive distribution of displacements has finite absolute
cost for this observable precisely when its inverse first moment
near zero is finite, counting a zero atom as infinite. A
contractible return with an enforced intermediate visit and a
diagonal trace without that insertion are different source
operations. A claim for GCM-induced O--QED must determine from
the actual charged trace, measure and boundary state which
operation occurs. The root-pair calculation is an obstruction to
inferring that answer from the curvature square alone. Its field
content differs from the native physical chiral multiplet.

\paragraph{A geometric phase that survives on each leg and cancels on every return.}
The full root trace can be checked with any number $n\ge1$ of
noncompact flat normal directions. On the Euclidean carrier let
$y\in\mathbb R^n$ denote the normal coordinate. Replace
$\mathrm dx^3$ in the product orthonormal coframe by
$\mathrm dx^3-\lambda\cdot\mathrm dy$, with a fixed
$\lambda\in\mathbb R^n$. The coframe has unit Jacobian, and
$B\,\mathrm dx^1\wedge\mathrm dx^2$ is unchanged. The free normal
derivatives in the curvature Hessian become
$\partial_{y_a}+\lambda_a\partial_{x^3}$ in the magnetic gauge
$A_3=0$. At fixed tangential momentum $k_3$, the normal heat kernel
for a segment of time $s$ and displacement $z$ is exactly
\[
 (4\pi s)^{-n/2}\exp[-|z|^2/(4s)]
                 \exp[-i k_3\lambda\cdot z].
\]
Fourier translation of the symbol $|\xi+\lambda k_3|^2$ proves
this formula. On any unreduced polygon returning to its starting
point the last factors multiply to one, however large $k_3$ is.
Integrating all intermediate positions gives the diagonal density
$(4\pi t)^{-n/2}$ at total time $t$. For any fixed subdivision into
at least two positive time intervals, almost every polygon so
integrated has strictly positive total traveled distance.

The weighted sum over the two charged roots of the complete vector
and ghost heat traces, per unit carrier and normal volume and after
taking the field-dependent difference, is
\[
 \frac{1}{(4\pi t)^{2+n/2}}
 \left[\frac{x}{\sinh x}\{n+2\cosh(2x)\}-(n+2)\right],
 \qquad x=Bt .
\]
Two polarizations acquire the spin shifts, the other $n+2$ are
unshifted, and the ghost pair removes two of the latter in the
weighted trace. The coefficient of $x^2$ is $(22-n)/6$.
For $n\ne22$ the magnetic action coefficient therefore contains a
nonzero multiple of $\int_0^T t^{-1-n/2}\,\mathrm dt$. At $n=22$
the quadratic coefficient vanishes, but the coefficient of $x^4$
is $17/15$ and the quartic term costs
$\int_0^T t^{-10}\,\mathrm dt$. Here $T$ only isolates the
short-time behavior. This is a divergence of an even represented
curvature block with its actual metric phase, after combining the
statistics of the fluctuating fields. The physical chiral source has
different field content.

The geometric distinction can be quantified without a momentum
regulator. In the two-leg return the natural midpoint $Y$ has
conditional density $(\pi t)^{-n/2}e^{-|Y|^2/t}$. The probability
of $|Y|\ge h$ is $\Gamma(n/2,h^2/t)/\Gamma(n/2)$.
For example, a boundary observable which tests
$h(t)^2=t\beta\log(T/t)$ near $t=0$, $\beta>0$, gives the
nonzero quadratic coefficient a short-time absolute cost proportional
up to fixed factors to
\[
 \int_0^{T/e} t^{\beta-1-n/2}
       [\log(T/t)]^{n/2-1}\,\mathrm dt .
\]
It converges exactly when $\beta>n/2$. With
$\beta=n/2$ and the further term $t\gamma\log\log(T/t)$ in
$h(t)^2$, convergence requires exactly $\gamma>n/2$.
These tests follow from $\Gamma(a,z)\sim z^{a-1}e^{-z}$.
They describe how strongly a source would have to weight excursions
in this example. The ordinary trace uses the full midpoint law.
Merely recording all its positive-length returns changes no coefficient.

\paragraph{A closed root history with resolved intermediate visits}
The same action gives an exact test for an unreduced word. Evaluate the
normal heat kernel with point visits at $r_0,r_1,\ldots,r_m=r_0$,
at fractions $\lambda_j>0$ of the proper time,
$\sum_{j=1}^m\lambda_j=1$. These insertions specify a different
observable from the ordinary trace. For $d_j=r_j-r_{j-1}$ set
$D^2=\sum_j d_j^2/\lambda_j$. In the Cartan background the normal
kernels commute with the tangential operator. Multiplying
\eqref{eq:uvpeC-offdiagonal} along the word cancels the flat phases
since $\sum_j d_j=0$. It leaves the positive factor
\[
 (4\pi t)^{-m/2}
 \left(\prod_{j=1}^m\lambda_j\right)^{-1/2}e^{-D^2/(4t)} .
\]
The $B^2$ coefficient has the same tangential prefactor as in
\eqref{eq:uvpeC-root-response} and the following integrated normal
cost at $D>0$:
\begin{equation}
 \mathcal C_m(D,\lambda)
 = \frac{\Gamma(m/2,D^2/(4T))}
        {\pi^{m/2}D^m\sqrt{\prod_j\lambda_j}} .
 \label{eq:uvpeC-word-cost}
\end{equation}
Here $\Gamma(a,z)=\int_z^\infty u^{a-1}e^{-u}\dd u$ comes from
$u=D^2/(4t)$ in the $\dd t/t$ integral. For an out-and-return word
through $d\ne0$, $m=2$ and $a=\lambda_1(1-\lambda_1)$, this gives
$\sqrt{a}\,e^{-d^2/(4aT)}/(\pi d^2)$. For a fixed partition and
finite positive distribution of the positions the exact condition
is integrability of $D^{-m}$ near zero, counting an atom there as
infinite. With fluctuating partitions it is the joint cost
\eqref{eq:uvpeC-word-cost} that must be integrated. The actual
traveled distance satisfies $\sum_j|d_j|\le D$ even when the net
displacement vanishes. This calculation establishes the gain for the
specified magnetic insertion. Whether the physical source enforces
such visits on each dangerous charged, gauge or ghost loop is an
additional dynamical question.

\paragraph{Arbitrarily many external photons on one marked root loop}
One can dispense with the constant magnetic field when testing
convergence for a fixed word. Take the same flat normal connection
and finitely many tangential Cartan probes of fixed Fourier momenta,
independent of $r$. Expand the marked proper-time trace about zero
tangential source, at finite $T$. The normal and tangential operators
commute for these data, so the normal factor preceding
\eqref{eq:uvpeC-word-cost} is exact at every source order. Each
linear insertion in the carrier Hessian has a Fourier coefficient
at most linear in the loop momentum, including the bounded spin
insertion at fixed external momenta. A quadratic insertion is a
bounded contact. The ghost Hessian obeys the same bounds.
For an order-$b$ jet with $u$ linear and $v_2$ quadratic insertions,
$b=u+2v_2$. The Duhamel simplex contributes $t^{u+v_2}$. After
taking absolute values, complete the square in the product of heat
kernels. The Gaussian center is a convex combination of finitely
many routing shifts, hence bounded uniformly on the simplex. Its
degree-$u$ Gaussian moment costs at most $C t^{-2-u/2}$ on
$0<t\le T$. Each carrier term is thus $O(t^{b/2-2})$, without
using cancellations in the gauge trace.

For every fixed $D>0$, the product with the normal density is
integrable against $\dd t/t$ for arbitrary $b$. This proves bare
absolute UV convergence of each of these one-loop terms with
arbitrarily many external insertions. It also gives a depth budget:
for a fixed partition, put $\nu=(m+4-b)/2$. Up to constants the
sufficient cost near zero is
\[
 \int_0^T t^{-1-\nu} e^{-D^2/(4t)}\dd t
 \lesssim \begin{cases}
  D^{-2\nu},&\nu>0,\\
  1+|\log D|,&\nu=0,\\
  1,&\nu<0 .
 \end{cases}
\]
The sharper magnetic coefficient in \eqref{eq:uvpeC-word-cost}
uses the gauge contact and spin contributions at $b=2$. In this
estimate the external momenta are held fixed. Sewing those momenta
into loops requires a new uniform estimate for the sewn graph.

\paragraph{Sewing arbitrarily many marked root loops with one radial state.}
There is a positive result for a specified choice of trace and
boundary readout. Return here to the unsheared product metric
($\lambda=0$), with one flat normal line and vanishing normal
connection. The normal heat transfer then commutes with the
tangential Hessian for the background insertions used below. For each loop insert the midpoint at the signed normal position $+r$
relative to its basepoint and return. An equal average over $+r$
and $-r$ with total weight one gives the same normal kernel. On a Euclidean carrier four-torus
of volume $V$ take a flat Cartan holonomy whose shift
$a$ lies outside the reciprocal lattice $\Lambda^*$. This gives
the free charged roots a spectral gap
$\delta^2=\operatorname{dist}(a,\Lambda^*)^2>0$, an infrared
condition. Hold this reference holonomy fixed while $e$ multiplies
the fluctuating field. Fix the transverse Maxwell zero mode and use its harmonic
readout $P_r A=e^{-r\sqrt{-\Delta}}A$ in the whole tangential
Hessian. Thus the connection seen at this visit is the same in
each tangential insertion. There is one sampled $r$ for all loops.
In fixed units define the marked log weight by its coefficientwise
heat expansion
\[
 W_r[eA]=\int_0^\infty\frac{\mathrm dt}{t}\,
 \frac{e^{-r^2/t}}{2\pi t}\,
       \mathcal H_{\rm tan}(t;a,eP_r A).
\]
Here $\mathcal H_{\rm tan}$ is the full tangential root-pair trace
with vector weight $1/2$ and ghost weight $-1$, including its
zero-field value. At any fixed order it is obtained directly by
differentiating the two curvature Hessians. The normal factor is
the product of their two half-time Gaussian transition densities.
Exponentiation forms collections of these marked loops. This
specifies a boundary observable and its loop expansion. The ordinary
curvature determinant traced over all normal basepoints has the
different diagonal kernel derived above.

\begin{proposition}[Bare graphs of the marked curvature trace]
\label{prop:geometry-marked-root-sewing}
Give $r\in(0,1)$ the nonzero-angular-mode ground-state law
$|u_j(r)|^2\mathrm dr$ of the radial metric
$\mathrm dr^2+r^{2d}\mathrm d\theta^2$, $d>3$. Put
$p=d-1$, $A_j=2|j|/p$, and $b_V=3V/(16\pi^3)$. If $p>6$, or
$p=6$ and $A_j>b_V$, the free normalization
$\int e^{W_r[0]}|u_j(r)|^2\mathrm dr$ and every bare graph in
the formal electromagnetic expansion of this source are finite.
The graph bound is absolute over the full charged and photon
momenta and all proper times, with fixed external Fourier probes.
It includes multiphoton contacts and vacuum graphs.
\end{proposition}
\begin{proof}
To see the role of the free weight, the three effective
polarizations at zero field give
$\mathcal H_{\rm tan}(t;a,0)=3\sum_{k\in\Lambda^*}
e^{-t|k+a|^2}$. Poisson summation gives
$3V(4\pi t)^{-2}(1+O(e^{-c/t}))$ at small $t$. The holonomy gap
controls the large-time part. Integration gives
\[
 W_r[0]=b_V r^{-6}+O(1),\qquad r\downarrow0 ,
\]
using $\int_0^\infty t^{-4}e^{-r^2/t}\mathrm dt=2r^{-6}$.
The comparison for the radial ground state gives
$|u_j(r)|^2\asymp r^d e^{-A_j r^{-p}}$ near zero
(\eqref{eq:dualrev-uvheat-ground}). This absorbs the full vacuum and every fixed inverse
power under the strict conditions above. The free normalization
diverges for $p<6$ or $p=6$, $A_j<b_V$, and is finite at
$p=6$, $A_j=b_V$ by its $r^d$ prefactor. The quartic radial
metric already fails the free normalization for this marked
root-loop ensemble. Its inverse radial moments alone cannot
control a different charged vacuum.

For an individual graph, let $P$ denote the sum of the norms
of all photon and fixed external momenta. At a root loop with $u$
linear and $v$ quadratic insertions, the Duhamel simplex supplies
$t^{u+v}$ times a bounded simplex integral. The complete minimal
coupling has linear vertices of degree at most one in the loop and
incident momenta, and bounded quadratic contacts. The spin term is
linear in the photon momentum. For successive routing shifts
$c_i\in\Lambda^*$, $|c_i|\le P$, the exponent of its free heat
product is $-\sum_i t_i |k\pm a+c_i|^2$, $\sum_i t_i=t$.
Reserve half this exponent for $e^{-\delta^2t/2}$ and complete
the square in the rest. A translated lattice Gaussian moment then
bounds the charged momentum sum, including all these vertices, by
\[
 C(1+P)^u e^{-\delta^2t/2}
            (1+t^{-2})(1+t^{-u/2}) .
\]
The estimate is uniform on the simplex, including its collapsing
faces, since its center is shifted only by a convex combination of
the $c_i$. Multiplying by the midpoint density and integrating in
$t$ costs at most $C'(1+P)^u r^{-N}$ for some finite $N$.
Indeed $t\ge1$ is exponentially controlled, while any inverse
power at zero is absorbed after substituting $v'=r^2/t$.

Each sewn transverse photon has nonzero momentum $q$ and modulus
bounded by $C e^{-2r|q|}/|q|^2$, since both ends see $P_rA$.
After the charged sums one may drop the remaining routing
constraints in the positive majorant. Lattice shell sums of any
fixed polynomial in the photon momenta against these factors cost
only further inverse powers of $r$. Multiplying the bounds for all
loops at their common $r$ and integrating the full vacuum weight
proves the claim by Tonelli. In particular no hierarchy among
the vanishing proper times is imposed. At every finite order the
connected and normalized coefficients are finite algebraic
combinations divided by powers of the one free normalization.
\end{proof}
The proof uses no damping inserted after sewing and makes no
claim of convergence of the charge series. Applying it to the
native physical sector requires deriving these marked traces and
the two endpoint readouts from that sector's quantum source.

\section{A charged Wilson action with a Dirac limit and dynamical photons}
\label{sec:pe-charged}

A comparison close to the O--QED target should couple a spinor to a
transverse photon, give a nonzero electromagnetic loop, and include the
contraction of the two photon legs. The following construction meets
these tests on one Euclidean carrier $\bin=(s,b_s)$. Its long-wavelength
tree action has a minimally coupled Dirac term. The charged Hessian at
large momentum, the preparation of the history, and the normal photon
action are specified below. Derivation of these choices from the
physical GCM source is a separate question.

\subsection{A charged operator whose zero-visit determinant is fixed}

Use a four-torus of period $2\pi$ in dimensionless reference units,
Hermitian four-component Euclidean Dirac matrices $\gamma_j$, and
Grassmann fields $\bar\psi,\psi$. Remove the harmonic photon mode for
infrared control. Choose four normal-history readouts $h_j$ once
for the entire functional. They are independent with the same
$C^\infty$ probability density $u$ supported in $[1,2]$ and positive
on its interior, for example the normalized extension by zero of
$\exp[-1/((h-1)(2-h))]$. The $h_j$ specify lengths of directed
tangential translations read from the history. They are independent of
the normal interval length $L$ introduced below. Neither type of
length denotes a sector clock. Let $b$ be the transverse photon field
seen on the visit. Write $e_j$ for a Euclidean coordinate unit in
this calculation and $e$ for the electric coupling. Define
\begin{align}
 (W_{j,\sigma}f)(x)
  &=\exp\!\left[-\ii e\int_0^{\sigma h_j}
                  b_j(x+t e_j)\,\dd t\right]
       f(x+\sigma h_j e_j),\nonumber\\[-2pt]
 P&=(m+4w)(1-\Delta)^2,\qquad m,w>0,\nonumber\\
 T_{h,b}
  &=\sum_{j=1}^4\left[
     \gamma_j\frac{W_{j,+}-W_{j,-}}{2h_j}
       -w\frac{W_{j,+}+W_{j,-}}2\right].
 \label{eq:uvpeW-hops}
\end{align}
Here $\sigma=\pm1$ labels the two directions in a hopping term.
It is independent of the signs of the directed pole prescription.

In the transverse variables take the charged action
$\langle\bar\psi,(P+gT_{h,b})\psi\rangle$; $g$ counts visits.
For the unreduced small Abelian gauge symmetry let $a$ be the
boundary connection, $\Delta\chi_a=\partial\cdot a$ with
$\int\chi_a=0$, and $M_a=\exp(\ii e\chi_a)$. The original charged
block is $M_a(P+gT_{h,b})M_a^{-1}$, with $b$ read in transverse
variables. Under $a\mapsto a+\dd\lambda$ on the trivial bundle,
$\chi_a\mapsto\chi_a+\lambda-\overline\lambda$. The block transforms
by the charged conjugation. At finite Fourier cutoff in the
transverse variables, transport the charged cutoff domain by $M_a$
and compress the whole operator to it. The field and dual Berezin
Jacobians cancel. The full determinant at $g=0$ is $\det P$ for
every photon configuration. Hence division by this one free number
preserves every electromagnetic source derivative. For $g\ne0$ all
derivatives of the Wilson exponentials in~\eqref{eq:uvpeW-hops},
including their same-visit contacts, belong to the interaction.

For a smooth slowly varying test field the Wilson shift is the
covariant translation with generator $D_j=\partial_j-\ii e b_j$.
Its Taylor expansion gives
$(W_{j,+}-W_{j,-})/(2h_j)=D_j+O(D_j^3)$ and
$(W_{j,+}+W_{j,-})/2=I+O(D_j^2)$, counting $b$ and a derivative
as first order. This is a local derivative expansion on the fixed
torus. At $g=1$ the constant $4w$ cancels, so the zeroth
and first orders of the charged block are
$m+\sum_j\gamma_jD_j$ in the transverse variables for every choice
of $h$. This matches the minimal tree vertex as well as the free
Dirac symbol to that order. In identifying $b$ with the transverse
potential in the raw convention $\partial+\ii e_{\rm raw} A$, use
$e=-e_{\rm raw}$ here and the corresponding potential gauge
parameter convention. The fixed axes and higher derivative
terms are part of this finite-volume Euclidean comparison.

\subsection{An exact hard-photon test}

Suppose first that the visited field $b$ is a translation-invariant
centered transverse Gaussian with covariance $D(p)$ for $p\ne0$,
independent of $h$. Each $D(p)$ is positive semidefinite, with the
reality relation at $-p$. Put $d_j(p)=D_{jj}(p)$. The traces, open
charged lines and photon Wick contractions come from the same
charged source. Integrate the $h_j$ before taking absolute
values in the remaining carrier momentum sums.
\begin{proposition}[Exact photon condition for the Wilson comparison]
\label{thm:uvpeW-sewing}
Every individual interacting graph at every fixed order in $g,e$
has an absolutely convergent carrier momentum sum after the stated
normal integrations, with finitely many external Fourier probes,
precisely when
\begin{equation}
 B(D)=\sum_{j=1}^4\sum_{p\ne0}\frac{d_j(p)}{1+p_j^2}<\infty .
 \label{eq:uvpeW-budget}
\end{equation}
Failure is already detected by the order-$ge^2$ vacuum graph with
two photon marks contracted on a single charged visit. Vacuum
normalization here removes the free source-independent factor only.
\end{proposition}
\begin{proof}
At fixed photon momenta a mark on a segment supplies
$(e^{\ii p_j\sigma h_j}-1)/(\ii p_j)$ when $p_j\ne0$, and
$\sigma h_j$ at $p_j=0$. Expand each nonzero-frequency numerator into
two characters. The remaining $h_j$ weights are the smooth compact
density multiplied by fixed powers of $h_j$ and $h_j^{-1}$. Its
Fourier transform decreases faster than any inverse power, uniformly
in the resulting frequency offsets after translation of the argument.

For clarity, a trace with $r\ge2$ visits is absolutely summable over
its charged loop momentum at fixed offsets: the factors
$(1+|k+a_i|^2)^{-2}$ satisfy the generalized H\"older bound by the
$r$th power of the $\ell^r(\mathbb Z^4)$ norm of the unshifted
factor. In a one-visit trace the three components perpendicular to
the hop instead contribute at most $C/(1+|k_j|)$, by summing the
three-dimensional shells at fixed $k_j$. Such traces can share the
same $h_j$. For $t$ one-visit traces in direction $j$, integration
in $h_j$ gives a rapidly decreasing function of
$\sum_{i=1}^t\sigma_i k_{i,j}$ plus an offset. The convolution
of $t\ge2$ copies of $(1+|n|)^{-1}$ is bounded by Young's inequality,
since each copy is in $\ell^{t/(t-1)}(\mathbb Z)$. A single copy
is already bounded. Summing against the translated rapidly
decreasing function thus costs a constant independent of the other
momenta. Applying this bound to the one-visit traces first leaves
the H\"older sums for the longer traces. An open charged line has
fixed endpoint momentum, and its intermediate momenta are determined
by the photon transfers. Its free factors are bounded. Spin traces
and the number of character terms are finite at fixed order.

For a photon joining segments in directions $i,j$, positivity gives
$|D_{ij}(p)|\le\sqrt{d_i(p)d_j(p)}$. The interval factors bound its
sum by a fixed constant times the geometric mean of the two
directional sums in~\eqref{eq:uvpeW-budget}, by Cauchy--Schwarz.
Dropping residual momentum constraints gives an upper bound
$C_{G} B(D)^n$ for a graph with $n$ internal photons. The same
bounds remove bounded Fourier cutoffs tending to one.

For necessity write $K(h)$ for the periodic Fourier kernel of
$(1-\Delta)^{-2}$ at $h e_j$, with the torus trace normalization.
It is strictly positive and smooth for $1\le h\le2$. Indeed,
the heat representation integrates the positive periodized
four-dimensional Gaussian against $t e^{-t}\dd t$, and this compact
interval is separated from the diagonal. In the trace of one visit
the Dirac term has zero spin trace. The scalar hopping term has the
same sign for both orientations. Up to one fixed positive factor
and the choice of sign for the logarithm, the summed contraction is
$\sum_{p,j}d_j(p)J(p_j)$, where
\begin{align*}
 J(m')&=\int_1^2 u(h)K(h)\,|I_{m'}(h)|^2\,\dd h,\\
 I_{m'}(h)&=(e^{\ii m'h}-1)/(\ii m')\quad(m'\ne0),
 \qquad I_0(h)=h .
\end{align*}
The trace may first be evaluated with positive heat time bounded
away from zero, then by the heat bound on this interval. Equivalently
the integrated Fourier series converges by the one-visit estimate
above at every fixed $p$. Each $J(m')$ is positive, and the
Riemann--Lebesgue lemma gives $m'^2J(m')\to2\int uK>0$.
Thus $J(m')\asymp(1+m'^2)^{-1}$ and the condition is necessary.
\end{proof}

For the ordinary four-dimensional transverse Maxwell covariance
$D(p)=\Pi(p)/|p|^2$, the slice $p_j=0$ contributes
$\sum_{0<|p_\perp|\le\Lambda}|p_\perp|^{-2}$, of order $\Lambda$
in three dimensions. The suppression that sums the charged loop at fixed photon
momenta consequently leaves
a divergent complete photon sewing.

\subsection{Computing the visited photon law and the return weight}

The required extra gain can itself come from a quantum normal
action. Fix the transverse boundary value $A_p=U_p(0)$ and then
integrate it and the bulk field with the joint quadratic form per
independent real transverse component
\[
 \frac{z}{2}|A_p|^2+
 \frac12\int_0^\infty
       (|U_p'|^2+z^2|U_p|^2)\,\dd r,\qquad z=|p|^2 .
\]
The minimum at fixed $A_p$ is $e^{-zr}A_p$ and adds precision $z$
to the boundary term. The independent Dirichlet fluctuation has
kernel $(e^{-z|r-s|}-e^{-z(r+s)})/(2z)$. Adding the covariance
of the harmonic mean leaves the full kernel
$\Pi(p)e^{-z|r-s|}/(2z)$. Thus the boundary has an ordinary
massless transverse pole of residue $1/2$ in these units. The
visit uses $b=L^{-1}\int_0^L U(r)\dd r$. Its conditional mean at
fixed $A_p$ is $(1-e^{-Lz})A_p/(Lz)$, which has unit soft limit.
That multiplier controls the external boundary-to-visit transfer.
A sewn photon instead uses the full covariance between two visit
insertions, including the Dirichlet fluctuations. Integrating this
covariance over the two sample positions gives
\begin{align}
 D_L(p)&=\Pi(p)\,\frac{\Phi_L(z)}{2z},\nonumber\\
 \Phi_L(z)&=\int_0^L e^{-tz}\frac{2}{L}(1-t/L)\,\dd t
    =\frac{2(Lz-1+e^{-Lz})}{(Lz)^2}.
 \label{eq:uvpeW-depth}
\end{align}
The triangular law is the distribution of the distance between two
uniform locations on the shared visited interval. It gives
$\Phi_L(z)\asymp(1+Lz)^{-1}$. For $0<L\le1$,
\[
 B(D_L)\asymp
 \sum_{k\in\mathbb Z^3\setminus\{0\}}
          \frac{1}{|k|^2(1+L|k|^2)}
 \asymp L^{-1/2}.
\]
Use $p_j=0$ for the lower bound. For the upper bound only
$p_\perp\ne0$ contributes since $\Pi_{jj}=0$ on the remaining
axis; replace $|p|$ by $|p_\perp|$ and sum $(1+p_j^2)^{-1}$.
Splitting the remaining shells at $L^{-1/2}$ gives the result.
For $L\ge1$ the budget is bounded.

\begin{corollary}[Sharp inverse moment from a connected current response]
\label{cor:uvpeW-sharp-depth}
Replace the depth preparation by a finite positive law $\mu$ on
$[0,\infty)$, keeping the $h_j$ independent and using
$D_0(p)=\Pi(p)/(2|p|^2)$ at zero interval. For $n\ge1$, the absolute UV test holds for every graph with exactly
$n$ internal photons, over all finite visit orders, precisely when
\[
 \mu(\{0\})=0,\qquad
 \int_{(0,1]}L^{-n/2}\,\mu(\dd L)<\infty .
\]
\end{corollary}
\begin{proof}
Sufficiency follows from Proposition~\ref{thm:uvpeW-sewing}. For a
lower bound take one charged visit with $2n+2$ photon marks, with
$n$ Wick pairs attached to the same visit and a pair of external
transverse plane waves $e_j e^{\pm\ii p_*\cdot x}$ with
$p_*\ne0$, $(p_*)_j=0$. Choose $e_j$ polarizations. Only the hop
in direction $j$ survives these external marks. The spin trace again
retains just the scalar hopping term. Up to its fixed overall sign,
the integrand after the charged trace is positive. Restrict the $n$
independent sewn momenta to $(p_a)_j=0$, so every pair contributes
$h_j^2d_j(p_a)$ and the external pair contributes $h_j^2$ for
probes of $b$. For boundary probes it also carries the squared
fixed-mode boundary-to-visit factor, which has a strictly positive
limit as $L\downarrow0$. The common charged weight $u(h_j)K(h_j)$
has strictly positive integral on $[1,2]$. The restricted sewn
sums therefore give a lower bound $c_n L^{-n/2}$ at small positive
$L$, with divergence already at $L=0$. Positivity permits Tonelli's
theorem after this charged trace. The absolute joint sum dominates
the sum of the resulting absolute reduced coefficients. This is a
connected two-external-photon graph: every photon pair meets the
same Wilson contact on one closed charged line. Free vacuum
normalization leaves the test intact.
\end{proof}

Finally take a normal diffusion with generator $\partial_x^2$,
starting at $0$, marked at $x(L/2)=d>0$ and returning to $0$ at
$L$. Its free length measure is $e^{-L}\dd L$. Multiplication of
the two transition densities gives the actual relative weight
\begin{equation}
 (2\pi L)^{-1}\exp(-L-d^2/L)\,\dd L .
 \label{eq:uvpeW-return}
\end{equation}
This is one shared zero-net-displacement return of positive length.
Its normalization and every inverse moment are finite. The same
interval $L$ appears in every charged visit. Conditional on $L$,
the bulk and the $h_j$ are independent before interaction. The
bound $C_G B(D_L)^n$ now integrates at every fixed order, including
all mixed Wilson contacts, graphs with open charged lines, and the
first dynamical photon sewing. Fixed boundary photon probes only add
bounded cross covariances at their fixed nonzero momenta. The
absolute sums concern the bare graphs after the prescribed
history integration and require no ultraviolet counterterms.

\subsection{Finite coupling in the same finite-volume source}

At each $L>0$ the Wilson interval defines a real Gaussian integral in
the mode Hilbert space of the normal photon law. Its variance is
bounded by $C B(D_L)$ from \eqref{eq:uvpeW-budget}, uniformly in the
starting point and the allowed segment lengths. The resulting phase
is a measurable unitary multiplier for real $e$ and reversing the
segment gives the adjoint shift.

The relative perturbation $A=P^{-1}T_{h,b}$ is
Hilbert--Schmidt because $P^{-1}$ is Hilbert--Schmidt and $T_{h,b}$ is
bounded. Its single-insertion trace uses the displaced kernel of
$P^{-1}$. Write $T_{h,b}=\sum_{j,\sigma} C_{j,\sigma}
 M_{F_{j,\sigma}}U_{j,\sigma}$, where $C_{j,\sigma}$ is the
constant spin coefficient in \eqref{eq:uvpeW-hops},
$F_{j,\sigma}$ is the interval phase and
$U_{j,\sigma}f(x)=f(x+\sigma h_j e_j)$. On the original lift,
the shift evaluates the kernel at $\sigma h_j e_j$ when the
charged loop closes. Define
\begin{equation}
 a_1(h,b)=\sum_{j,\sigma}
 \frac{K(h_j)}{m+4w}\,\operatorname{tr}_{\rm spin}(C_{j,\sigma})
             \int F_{j,\sigma}(x)\,\dd x ,
 \label{eq:polish3-wilson-first}
\end{equation}
where $K$ and the normalized torus volume are as in
Proposition~\ref{thm:uvpeW-sewing}. The heat integral defining
$K(h_j)$ is bounded on the entire displacement interval. For an
individual shift the distributional kernel of $P^{-1}M_FU_{j,\sigma}$
has its singularity displaced from the diagonal. The displayed
trace uses that displaced kernel value and the integral of the
bounded phase. The pair of orientations makes $a_1$ real.

For later photon derivatives, write $K_P(x)$ for the convolution
kernel of $P^{-1}$ on the same normalized torus. If $f\in L^2$ and
$U_a u(x)=u(x+a)$, the kernel of $P^{-1}M_fU_a$ is
$K_P(x-y+a)f(y-a)$. Thus for scalar $f$ and a single spin component,
\begin{equation}
 \big\|P^{-1}M_fU_a\big\|_2^2
    =\|K_P\|_{L^2}^2\|f\|_{L^2}^2 .
 \label{eq:finaluv-wilson-HS}
\end{equation}
Tonelli and translation invariance give this identity directly.
It defines the Hilbert--Schmidt extension even when the multiplier
$f$ itself is unbounded. An additional fixed spin matrix only
changes the spin factor.

\begin{proposition}[Positive finite relative partition at physical visit coefficient]
\label{prop:uvpeW-finite-coupling}
With the normal law \eqref{eq:uvpeW-return}, the relative Wilson
partition is finite and strictly positive for real $e$ and
$0\le g\le1$. Its weight retains the complete finite first trace
\eqref{eq:polish3-wilson-first}.
\end{proposition}
\begin{proof}
The complete relative determinant prescribed by the displaced trace
and the ordinary higher trace powers is
\begin{equation}
 \det_{\rm rel}(P+gT_{h,b},P)
     =e^{g a_1(h,b)}\det{}_2(I+gA).
 \label{eq:uvpeW-relative-det}
\end{equation}
Here $\det_2(I+X)$ is the ordinary Fredholm determinant of
$(I+X)e^{-X}$, whose difference from the identity is trace class for
Hilbert--Schmidt $X$~\cite{SimonUVpeDet}. For small $g$, expansion
of its logarithm gives the absolutely defined operator traces
$\operatorname{Tr}(A^n)$ for $n\ge2$, with coefficients
$(-1)^{n+1}g^n/n$. The first insertion in
\eqref{eq:uvpeW-relative-det} is evaluated directly by the
off-diagonal heat kernel in \eqref{eq:polish3-wilson-first}.

For positivity, the bounded balanced operator
$C=P^{-1/2}T_{h,b}P^{-1/2}$ also belongs to the Hilbert--Schmidt
class, by the Schatten product estimate using $P^{-1/2}$ in the
fourth Schatten class. Its nonzero eigenvalues with multiplicities
agree with those of $A$, since $A$ and $C$ are the two products in
opposite orders of $P^{-1/2}$ and $P^{-1/2}T_{h,b}$.
The bound $\operatorname{Re} T_{h,b}\ge -4wI$ gives
$\operatorname{Re}(I+gC)\ge m/(m+4w)$, and chirality conjugation
maps $C$ to $C^*$. Thus $\det_2(I+gC)$ is real, nonzero and
continuous along $0\le g\le1$. It starts at one and stays positive.
The finite first trace is bounded independently of the history
and sample because $|F_{j,\sigma}|=1$ and $1\le h_j\le2$.
The deterministic bound on $\|T_{h,b}\|$ and
$|\det_2(I+X)|\le e^{\|X\|_2^2/2}$ give an integrable upper
bound on the full relative determinant. This proves the positive
finite normalization.

One can compare \eqref{eq:uvpeW-relative-det} with a
Fourier heat approximation of the charged kernel. Its first trace
becomes the bounded heat-approximated $K(h_j)$, which converges
uniformly on $[1,2]$. Its higher trace powers converge in the
Hilbert--Schmidt pairing. The direct displaced kernel formula
already defines the first trace without this approximation.
\end{proof}

The positive relative weight defines the real-coupling source on
this compact Euclidean chart. Q1 instead asks for positivity of
physical state pairings computed from its correlations.

This is an existence result for the specified charged action. Its
geometry selects a transverse readout of every visit, a fourth-order
charged reference Hessian, and a resolved normal return. The local
raw O--QED Dirac action has a transverse response already at zero
visits. Whether the GCM-induced charged determinant has that response,
which of its vertices retain the travel factors, and what contracts
two such vertices are computations on its own source. The
long-wavelength Dirac match above supplies a check on the tree
current. Q1 positivity, the physical chiral content, and a continuum
or infinite-volume limit require their own analysis.

\subsection{All electromagnetic orders after exact summation of unmarked visits}
\label{subsec:g5W-fixed}
The preceding fixed-order graph estimate counts both charged visits
and photon insertions. At the physical visit coefficient $g=1$, a
given order in the electric charge $e$ can contain arbitrarily many
visits carrying no photon. The finite-coupling existence of the
relative partition, by itself, would not justify differentiating it
to every order. For the action~\eqref{eq:uvpeW-hops} those visits
can be summed into the quadratic block before generating graphs.

Write $p_k=(m+4w)(1+|k|^2)^2$ and let $T_0=T_{h,0}$. Its
zero-photon value and the exact free charged inverse at $g=1$ are
\begin{align}
 t_k(h)&=\sum_{j=1}^4\left[
       \ii\gamma_j\frac{\sin(k_j h_j)}{h_j}
                      -w\cos(k_j h_j)I\right],\nonumber\\
 S_h(k)&=[p_k I+t_k(h)]^{-1}.
 \label{eq:g5W-exact-free}
\end{align}
Keep the positive relative determinant at zero photon coupling
$W_0(h)=\det_{\rm rel}(P+T_0,P)$ from
\eqref{eq:uvpeW-relative-det}. It depends on the four common
lengths and weights the full $h$ integration once. Its integral
enters the overall free normalization. External charged probes
couple to open lines of this inverse.

\begin{theorem}[Bare graph convergence at fixed physical visit coefficient]
\label{thm:g5W-electromagnetic-orders}
For $m,w>0$ in \eqref{eq:uvpeW-hops}, set $g=1$ and expand in the
electric coupling alone. Use the normal preparation
\eqref{eq:uvpeW-return}, covariance~\eqref{eq:uvpeW-depth}, the
shared smooth $h$ density and the displaced trace
\eqref{eq:polish3-wilson-first}.
After the common normal integrations, each graph at every finite
electromagnetic order has an absolutely convergent sum in all its
carrier loop momenta, including all dynamical photon contractions
and finitely many fixed external Fourier probes. The sums
define the bare coefficients without ultraviolet counterterms. At each order the only charged vertices required
are the derivatives of $T_{h,b}-T_0$ of positive photon order.
\end{theorem}
\begin{proof}
Let $A_0=P^{-1}T_0$, $X=P^{-1}(T_{h,b}-T_0)$ and
$Y=(I+A_0)^{-1}X$. Then $I+A_0+X=(I+A_0)(I+Y)$ and
$Y=S_h(T_{h,b}-T_0)$. The Hilbert--Schmidt determinant product
formula~\cite{SimonUVpeDet} gives the exact relative factor
\begin{equation}
 \frac{\det_{\rm rel}(P+T_{h,b},P)}{W_0(h)}
     =\exp\bigl(a_1(X)-\operatorname{Tr}(A_0Y)\bigr)
            \det_2(I+Y).
 \label{eq:polish3-wilson-factor}
\end{equation}
Here $a_1(X)$ is the linear displaced trace of the difference
in \eqref{eq:polish3-wilson-first}, and $A_0Y$ is trace class.
The linear term in the logarithm is the full trace of the
one-vertex loop, using $Y=X-A_0Y$. The higher traces of $Y$ and
the expansion of $(I+Y)^{-1}S_h$ give the closed and open charged
lines. Every vertex carries at least one photon mark, so the
number of vertices at a fixed electric order is finite. The
zero-photon determinant $W_0$ remains in the common weight.

Let $\mathcal A$ be the algebra of functions of $h$ on the
four-torus of period $2\pi$ with absolutely summable Fourier
coefficients. Its norm is the sum of their moduli; matrix
coefficients can use the operator norm. Multiplication is bounded
by convolution. For the proof extend $1/h_j$ smoothly and
periodically from $[1,2]$, where the preparation is supported.
The extended coefficients of $t_k$ have uniformly bounded
$\mathcal A$ norms, since multiplication by a character merely
shifts the frequency. For $p_k$ sufficiently large the Neumann
series in this algebra gives
\begin{equation}
 \|S_h(k)-p_k^{-1} I\|_{\mathcal A}\le C p_k^{-2},
 \qquad \|S_h(k)\|_{\mathcal A}\le C p_k^{-1}.
 \label{eq:g5W-Wiener}
\end{equation}
There are only
finitely many smaller $k$. Use a real smooth extension of $1/h_j$:
the real part of the full zero-photon operator is at least $m I$
for all $h$. These finitely many inverses have uniform smooth
bounds and hence absolutely summable Fourier coefficients. This
proves~\eqref{eq:g5W-Wiener} for every $k$.

The zero-photon determinant also has to satisfy this Fourier
estimate on the support of the preparation. At high $k$ write the
mode logarithm as
\[
 \log\det(I+t_k/p_k)
       =\Tr(t_k)/p_k+ r_k(h),\qquad
       \|r_k\|_{\mathcal A}\le C p_k^{-2}.
\]
This follows by the matrix-log series for a small algebra norm.
The remainders sum absolutely in $\mathcal A$. The finitely many low mode determinants are smooth
and strictly positive by the explicit Dirac symbol or the accretivity
and chirality argument above. The linear term in the sum is the
spin trace of the scalar hopping part, a sum of off-diagonal kernels
$K(h_j)/(m+4w)$ with fixed coefficients. At $1\le h_j\le2$ these
kernels are smooth because the periodized heat representation
can be differentiated away from the diagonal.
Localize each of these linear terms with a smooth periodic
cutoff equal to one on $[1,2]$ and supported away from zero.
Exponentiation in $\mathcal A$ and multiplication by the smooth
preparation density give an absolutely summable Fourier series for
the complete zero-photon weight.

At fixed photon momenta, each mark in a vertex of direction $j$
is either a bounded power of $h_j$ at $p_j=0$, or a difference
of two characters divided by $\ii p_j$. Its absolute Fourier norm,
with the smooth extensions on the preparation support, costs
$C(1+p_j^2)^{-1/2}$. The remaining shift by the charged momentum
is a character in $h_j$. A charged loop with $r\ge2$ vertices now
contributes at most a constant times the sum of
$\prod_{i=1}^r p_{k+a_i}^{-1}$ in its loop momentum $k$. The
shifts $a_i$ are partial external and photon transfers. This sum
is bounded uniformly in the shifts by H\"older since
$(p_k^{-1})\in\ell^r(\mathbb Z^4)$. For a one-vertex loop, the
remainder in \eqref{eq:g5W-Wiener} is absolutely summable on its
own. Open charged line momenta are fixed by the transfers and
cost uniform constants.

Only the leading $p_k^{-1}I$ on a one-vertex closed loop needs the
shared history average. Its dependence on its own loop momentum
in the $h$ variables is one character $e^{\ii\sigma k_jh_j}$.
After summing the three perpendicular components of $k$, the
remaining weight is at most $C f(k_j)$ with $f(n)=(1+|n|)^{-1}$.
Hold all photon momenta and the other charged loop momenta fixed.
The product of their bounded Fourier factors, the smooth densities
and the full $W_0$ is independent of these leading loop momenta.
For $t$ such loops in one direction the coefficient of any total
frequency is majorized by $C^t f^{*t}$. This convolution is
bounded: for $t=1$ use $f\in\ell^\infty$, and for $t\ge2$ use
Young's inequality with $f\in\ell^{t/(t-1)}$. The four directions
may share the same zero-photon determinant; summing against its
joint absolutely summable Fourier coefficients and those of the
other bounded factors gives the product bound. This uses the
Fourier decay of the surviving normal integration exactly where
the individual undamped one-vertex series is marginal.

The photon index estimate following \eqref{eq:uvpeW-budget} now gives
$C_G B(D_L)^n$ for $n$ internal photons and fixed external Fourier
probes. The return law integrates this bound. The absolute Fourier
estimate on the shared weight covers the leading one-vertex
sums after averaging the $h_j$. The charged remainders and all
photon sums are bounded by nonnegative summable coefficients.
Integrating the common $L$ and using the positive integrated
zero-photon normalization gives the stated bare coefficients.

\end{proof}

This calculation uses the actual inverse at the physical visit
coefficient and leaves the source-dependent Wilson contacts in the
electromagnetic expansion. Its stronger ultraviolet decay at large
$|k|$ and its off-diagonal first trace both use the specified
fourth-order principal charged block. The local minimally coupled
Dirac principal operator in raw O--QED has neither estimate.
Applying this proof to a GCM-induced electromagnetic sector requires
deriving that sector's zero-photon determinant, the corresponding
principal inverse and all its photon derivatives from the same
GCM measure and boundary state.

\begin{corollary}[Agreement with derivatives of the finite-coupling source]
\label{cor:g5W-smooth}
For this finite-volume model at $g=1$, the relative partition and
the normalized correlations with finitely many fixed charged Fourier
probes and boundary photon Fourier modes are smooth in the real
electric coupling. Their derivatives at zero are the
bare coefficients of Theorem~\ref{thm:g5W-electromagnetic-orders}
multiplied by the corresponding Taylor factorials.
\end{corollary}
\begin{proof}
Conditional on $L>0,h$, write $X(x)$ for any of the real photon
interval integrals in the hops. It is a centered Gaussian at
almost every $x$ with variance at most $C B(D_L)$, uniformly in
$x,h$. Its phase derivatives at real $e$ are
$(-\ii X)^j e^{-\ii eX}$. Gaussian moments bound their $L^p$ norms
by $C_{j,p}B(D_L)^{j/2}$ on the product of photon probability and
torus volume. These moment bounds and the mean-value estimate for
the real phase justify each derivative in $L^p$. Boundary
photon Fourier probes have finite Gaussian moments on their fixed
nonzero modes.

Equation~\eqref{eq:finaluv-wilson-HS} turns the phase estimates
into bounds on the Hilbert--Schmidt derivatives of
$A(e)=P^{-1}T_{h,b}(e)$. Apply the bound first in $L^p$ with
$p\ge2$, where the conditional moment estimates control the
spatial $L^2$ norm. Lower probability moments follow by
H\"older's inequality. The restored first trace $a_1$ is an off-diagonal
kernel times spatial averages of these phases; its derivatives
have the same bounds. At order zero both $a_1$ and $\|A\|_2$ are
uniformly bounded on the real axis. The map
$A\mapsto\det{}_2(I+A)$ is entire in the Hilbert--Schmidt class.
Its bound on balls, together with the Cauchy estimate and
polarization, bounds its derivatives of each fixed order on every
bounded ball. Consequently derivatives of
$e^{a_1}\det{}_2(I+A)$ are controlled in expectation by fixed
polynomials in $1+B(D_L)$, uniformly in the real coupling.

For open charged lines use
$S(e)=(I+A(e))^{-1}P^{-1}$. The accretivity bound gives
$\|S(e)\|\le m^{-1}$; the identity
$(I+A(e))^{-1}=I-S(e)T_{h,b}(e)$ bounds its inverse factor
uniformly as well. Differentiating the inverse uses finite products
of these factors and the derivatives of $A$. The return law
integrates every power of $1+B(D_L)$, so H\"older's inequality and
the fundamental theorem of calculus justify differentiation under
the joint expectation on compact real coupling intervals. Division
by the strictly positive partition preserves smoothness.
Equation~\eqref{eq:polish3-wilson-factor} gives the graph expansion
about $P+T_0$ for these derivatives.

\end{proof}

\subsection{A transverse soft test for the Wilson action at the physical visit coefficient}
\label{subsec:s6Wilson-soft}
The two-photon contact of \eqref{eq:uvpeW-hops} can be calculated
before making any assumption about a Ward identity at zero momentum.
To test arbitrarily small nonzero momenta, give the torus period
$2\pi R$, $R\ge1$, and retain the readout lengths in the original
units. Write $V_R=(2\pi R)^4$ and $M=m+4w$. Take
$b=e_j(z_+e^{\ii qx^1}+z_-e^{-\ii qx^1})$, $j\ne1$, with
$0\ne q\in R^{-1}\mathbb Z$. The two complex amplitudes simply
polarize the real quadratic form. Let $H_{R,h}(q)$ be the mixed
derivative in $z_+,z_-$ at zero of the negative log relative charged
determinant divided by $V_R$, for $g=1$.

\begin{proposition}[A finite nonvanishing transverse term in the Wilson comparison]
\label{prop:s6Wilson-soft}
Fix $w>0$. Uniformly in $R\ge1$, $q$ as above and all allowed $h$,
as $m\to\infty$,
\begin{equation}
 \begin{aligned}
 H_{R,h}(q)&=-\frac{4e^2 w h_j^2}{M}K_R(h_j e_j)+O_w(e^2 M^{-2}),\\
 K_R(x)&=\int_0^\infty t e^{-t}
    \sum_{n\in\mathbb Z^4}
    \frac{e^{-|x+2\pi R n|^2/(4t)}}{(4\pi t)^2}\,\dd t .
 \end{aligned}
 \label{eq:s6Wilson-soft}
\end{equation}
For sufficiently large $m$ and $e\ne0$, the order-$e^2$ transverse
boundary correction per unit volume stays strictly negative and
bounded away from zero along $q=1/R\to0$, including the stipulated
common history integration.
\end{proposition}
\begin{proof}
A $j$-link sees the same plane wave at every point of its interval.
Writing $T_\pm=\partial_{z_\pm}T|_0$ and
$T_{+-}=\partial_{z_+}\partial_{z_-}T|_0$, the latter is
$-e^2h_j^2$ times the zero-field $j$-hop. Its spin trace at $k$
is $4e^2 h_j^2 w\cos(k_jh_j)$. Differentiating
\eqref{eq:polish3-wilson-factor} gives
\[
 V_R H_{R,h}(q)=
   -\operatorname{Tr}(S_h T_{+-})+
     \operatorname{Tr}(S_h T_- S_h T_+),
\]
where the first trace uses the off-diagonal prescription already
present in \eqref{eq:uvpeW-relative-det}. Set
$a(k)=(1+|k|^2)^{-2}$. Since $\|t_k(h)\|\le4(1+w)$,
the Neumann estimate for $M>8(1+w)$ gives
$\|S_h(k)\|\le2a(k)/M$ and
$\|S_h(k)-a(k)I/M\|\le8(1+w)a(k)^2/M^2$.
The symbols of $T_\pm$ are shifts by $\pm qe_1$ with norm
at most $C_w |e|$, and $\|T_{+-}\|\le C_w e^2$.
The normalized lattice sums of $a^2$ are bounded uniformly for
$R\ge1$ by cube comparison with its integrable continuum majorant.
Cauchy--Schwarz gives the same bound for the sums of
$a(k)a(k+qe_1)$. Thus both the exchange trace and the error from
replacing $S_h$ in the contact trace are $O_w(V_R e^2 M^{-2})$.
The leading normalized contact sum is the Fourier kernel of $a$
at a nonzero displacement. It is the positive heat integral in
\eqref{eq:s6Wilson-soft}. The $n=0$
term has a strictly positive lower bound for $1\le h_j\le2$,
independently of $R$. This proves the uniform sign for large $m$.

At $q\ne0$ the one-photon trace vanishes for each $h$ by translation
invariance. The mixed derivative of the negative log of the
history-integrated determinant therefore averages the displayed
Hessian with its normalized positive zero-field weight $W_0(h)$.
At electromagnetic order two the conditional bulk photon
fluctuations supply only a boundary-independent term. Each external
boundary mode supplies its mean transfer from \eqref{eq:uvpeW-depth},
so the Hessian is multiplied by the independent return average of
$[(1-e^{-Lq^2})/(Lq^2)]^2$, which tends to one by dominated
convergence. The uniform lower bound on $-H_{R,h}(q)$ survives
both operations.
\end{proof}
In this action a flat background shifts the symbol $t_k(h)$ to
$t_{k-eb}(h)$ while the principal symbol $M(1+|k|^2)^2$ is held
fixed. In this comparison the exact first-derivative defect from a
common momentum shift of the full charged symbol $D_0(k)=p_kI+t_k$
is $(\partial_{b_j}D+e\partial_{k_j}D_0)|_{b=0}
=e(\partial_{k_j}p_k)I$. At second order the corresponding missing
flat-source contact is $e^2(\partial_{k_j}^2p_k)I$. These are
identities of symbols; small-gauge covariance alone does not
implement a common shift by a constant transverse field.
The calculation uses nonzero transverse modes throughout and
retains both derivatives of the Wilson links. It locates a finite
order-$e^2$ obstruction to a massless inverse kernel of order $q^2$
in this parameter range, consistent with the graphwise ultraviolet
theorem. An extra cancellation for the physical GCM source would
have to follow from that source's own charged and photon derivatives.

\section{Inducing a charged action from a represented GCM curvature}
\label{sec:g5-induction}
We now construct a curvature pairing that produces the charged Wilson
action of Section~\ref{sec:pe-charged}, including its closed
determinant, on one specified graded path chart.
Its matter coordinates are independent Grassmann variables. Its
geometric inputs include a finite representation of endpoint arrows,
a chosen set of jets of the fields, and a boundary state. This
construction tests existence within the indicated class of
superconnections and curvature pairings. The physical selection
of this class remains a separate question.

\subsection{A curvature product for an endpoint hopping term}

Fix one $\bin=(s,b_s)$, the four-torus, infrared convention and
normal fields of Section~\ref{sec:pe-charged}. Work first with smooth
fields at finite regulator. Use four auxiliary closed
anticommuting internal one-forms $\theta^1,\ldots,\theta^4$.
Their labels are unrelated to Lorentz or normal indices. Take an
even neutral line $o$ and three odd copies $u,v,w_0$ of the complex
four-dimensional spin representation. The following are the only
entries of the connection coefficients $B_a$ in one copy:
\[
 \begin{array}{c|c|c}
 a&\text{map}&\text{coefficient}\\ \hline
 1&u\longrightarrow v&U\\
 2&o\longrightarrow u&\xi\\
 3&v\longrightarrow w_0&V\\
 4&w_0\longrightarrow o&\bar\zeta .
 \end{array}
\]
Here $U,V$ are even matrices. The column $\xi$ and row $\bar\zeta$
are Grassmann odd: their odd internal homomorphism parity makes
each complete $B_a$ even. Use the usual composition of even
supermodule maps. In $\mathbb D=d+\sum_a\theta^a B_a$ the curvature
coefficient for two internal indices is $F_{ab}=[B_a,B_b]$,
since the coefficients are independent of the internal form
coordinates. In particular $F_{12}=U\xi:o\to v$ and
$F_{43}=\bar\zeta V:v\to o$. Give this copy the real curvature
density
\begin{equation}
 \tfrac12\operatorname{Re}\operatorname{Str}
   (F_{43}F_{12}+F_{12}F_{43})
       =\operatorname{Re}(\bar\zeta VU\xi).
 \label{eq:g5-arrow-product}
\end{equation}
This follows either by graded cyclicity or by moving the Grassmann
row past the column in the negative odd-block trace. The indicated
component pairing is constant and independent of every field and
photon probe. It can be written as the top-form pairing of
\eqref{eq:uvpeA-curvature}: contract the two indicated two-form
coefficients across one another and multiply by the common
volume form. All other coefficient pairings in this internal
copy are zero. Consequently a mixed curvature with an $x$ index,
though generally present in $\mathbb D^2$, makes no contribution on
this chart. For example, the additional bosonic coefficient $[B_3,B_1]$
contains the path $VU:u\to w_0$. It is present and unpaired.
The mixed contraction is symmetric on the represented coefficient
pair, with opposite signs after diagonalizing its off-diagonal
coefficient matrix. It is a graded matter pairing. The overall
positive transverse photon form is obtained separately below.
Nondegeneracy or positivity of a Hodge metric on every represented
curvature component, including the unused ones, would require new
terms and a new proof.

There is also an adjoint-closed realization. On the real Grassmann
form choose the product-reversing matrix adjoint, and replace each
oriented coefficient by $B_a-B_a^\dagger$ before complexifying.
All four nodes are distinct. The first selected commutator is then
$X-X^\dagger$ with $X=U\xi$, and the second is $Y-Y^\dagger$
with $Y=\bar\zeta V$. Products $YX^\dagger$ and $Y^\dagger X$
vanish by their endpoints. The two trace orders in
\eqref{eq:g5-arrow-product} give the original terms plus their
adjoints, so the normalization $1/4$ in place of $1/2$ restores
the same real density. This adjoint construction adds neither a
second spinor integration nor a new spinor species.

Take the direct sum of the following copies with the \emph{same}
component pairing. On each raw history the cyclic functional is evaluation of the
composable fibre maps, their graded matrix trace and integration
over $x$ with the torus volume. A change of basepoint along a
translated link preserves this volume. Interval holonomies and source
insertions are evaluated on their unreduced representative before
taking the trace. The histories are then integrated with the specified
common state. The fixed GTL restriction and a normalized,
velocity-independent full shell contribute one overall factor in this
Euclidean example.
The row/column sources are images of just one
independent spinor pair $\bar\chi,\chi$; the copies have no
independent Berezin integrations. Let $\gamma_5$ be the Hermitian
chirality involution anticommuting with the four Euclidean Dirac
matrices. Use the dressed small-Abelian-gauge variables of
Section~\ref{sec:pe-charged}.

\begin{enumerate}
\item For a local zero-visit block set $D_0=I$, $D_j=\partial_j$,
$j=1,\ldots,4$. For each ordered pair $a,b=0,\ldots,4$ take
$\xi=D_aD_b\chi(x)$ and
$\bar\zeta=D_aD_b\bar\chi(x)$, with
$U=I$, $V=(m+4w)\gamma_5$. The symbols $D_a$ here only index jets
of the raw coefficient field. Integrating over $x$ on the periodic
carrier gives the operator
\[
 (m+4w)\gamma_5
 \sum_{a,b}(D_aD_b)^\dagger D_aD_b
  =(m+4w)\gamma_5(1-\Delta)^2 .
\]
Indeed $\sum_a D_a^\dagger D_a=1-\Delta$, and the flat derivatives
commute. No fourth-order operator was inserted in the curvature
pairing.
\item For each directed segment $x\to x+\sigma h_j e_j$,
$\sigma=\pm1$, use
$\xi=\chi(x+\sigma h_j e_j)$ and $\bar\zeta=\bar\chi(x)$.
The map $U$ is parallel transport from the displaced endpoint
back to $x$ with the interval phase
of \eqref{eq:uvpeW-hops}. The spin endpoint map is
\[
 V=\tfrac g2\,\gamma_5
       (\sigma\gamma_j/h_j-wI).
\]
The copy $u$ is pulled back from the displaced endpoint, whereas
$v$ and $w_0$ are based at $x$. Thus the spin map $V$ acts at one
point and $U$ connects the two charged fibres. Multiplying the four coefficients gives one
ordinary covariant hopping bilinear. The endpoint label and the
interval that produced the holonomy are retained when the word is
traced. The prescribed spin map implements the Clifford difference
and the scalar average in each direction. It contains no photon
field.
\end{enumerate}

Let $P,T$ be as in \eqref{eq:uvpeW-hops}. The sum of the component
pairings is $\operatorname{Re}\langle\bar\chi,\gamma_5(P+gT)\chi\rangle$.
This equals the indicated bilinear on the real Grassmann form:
give rows and columns the standard adjoint involution with product
order reversed. For real gauge fields opposite Wilson translations
are adjoints, and $\gamma_5(P+gT)\gamma_5=(P+gT)^\dagger$.
Thus $\gamma_5(P+gT)$ is Hermitian. Complexification supplies
independent Berezin row and column variables for the Euclidean
source. The constant row substitution
$\bar\psi=\bar\chi\gamma_5$, $\psi=\chi$, with the same
substitution in its source, has constant unit determinant per
four-spinor. This recovers the entire charged bilinear
$\langle\bar\psi,(P+gT)\psi\rangle$ including the sign of its
oriented interval. After undoing the dressing, the endpoint maps
and dressed coefficient jets transform in their prescribed charged
frames. The resulting action has the small Abelian transformation
law given in Section~\ref{sec:pe-charged}.

\subsection{An even curvature chart, its measure and induced graphs}

The accompanying photon law can also be obtained with elementary
curvature components. In radial gauge use an Abelian connection
with transverse tangential components $U_i(x,r)$ on $r\ge0$,
where $U$ in this paragraph is the photon field. Include its
normal--tangential curvatures $F_{ri}=\partial_r U_i$.
Use four further Abelian connection images with tangential
coefficients $\partial_j U_i$. Pair their tangential curvatures
$\partial_j(\partial_i U_k-\partial_k U_i)$ with themselves.
The $j$ index is summed with the Euclidean metric. With the common
quadratic normalization $1/2$ and a boundary Maxwell form at
$r=0$, the transverse quadratic action per real mode has precisely
the terms
\[
 \frac{|p|^2}{2}|U_p(0)|^2+
 \frac12\int_0^\infty
    (|U_p'|^2+|p|^4|U_p|^2)\,\dd r .
\]
This follows from
$\sum_{i<k}|p_i U_k-p_k U_i|^2=|p|^2|U|^2$ on the transverse
subspace. Other form components of these copies are unpaired.
The copies are derivatives of this same field in the transverse
chart. Their elimination creates no determinant.
For the prepared one-dimensional normal path $X(r)$ the Abelian
connection image $X(r)\theta$ in a closed internal direction has
$F_{r\theta}=X'(r)$. Use the quadratic coefficient $1/4$,
the length weight $e^{-L}\dd L$ and the endpoint/midpoint
state in \eqref{eq:uvpeW-return}. The boundary state also specifies
the interval readout and the common $h_j$ law of
\eqref{eq:uvpeW-hops}. These distributions and the permitted jets
are part of the construction.

\begin{theorem}[Induced charged graph finiteness on the specified graded chart]
\label{thm:g5-curv-induced}
Take the above represented curvature blocks with the Berezin measure
on the one raw spinor pair, the normalized Gaussian normal and
transverse-photon measures just derived, and the stated history
preparation. Use one source-independent zero-interaction normalization.
The induced charged and electromagnetic source has absolutely
UV-convergent complete bare graphs at every fixed order in visits
and photon insertions, including dynamical photon contractions.
It also has absolute convergence for every graph at each finite
electromagnetic order with $g=1$ already fixed, and the corresponding
coefficients are derivatives of the same normalized source.
At $g=1$ its finite-volume relative partition is finite and strictly
positive at real $e$ with the displaced trace
of Proposition~\ref{prop:uvpeW-finite-coupling}.
\end{theorem}
\begin{proof}
The displayed curvature contraction equals the charged bilinear
of \eqref{eq:uvpeW-hops} with its same endpoint data. There is
one Berezin pair and the auxiliary copies are maps of that pair.
The charged relative weight is therefore the one defined by
\eqref{eq:uvpeW-relative-det}, with the finite displaced first
trace, and the charged probes use the inverse of this bilinear.
The only field eliminated in restoring the small Abelian longitudinal
gauge variable has the source-independent Abelian gauge Jacobian
on this zero-mode-fixed chart. The linear normal photon Dirichlet
fluctuations are integrated with their full covariance. The induced
visit covariance is \eqref{eq:uvpeW-depth}, while the two pieces of
the constrained normal path contribute both heat-kernel factors in
\eqref{eq:uvpeW-return}. The conditional photon sum with $n$
internal photons costs at most a constant times $1+L^{-n/2}$ by
Proposition~\ref{thm:uvpeW-sewing}. The common return weight
integrates this bound at every $n$. At fixed $g=1$,
Theorem~\ref{thm:g5W-electromagnetic-orders} and
Corollary~\ref{cor:g5W-smooth} apply with the same zero-photon
determinant in the history integration.
Proposition~\ref{prop:uvpeW-finite-coupling} gives the finite
positive relative partition through the displaced-kernel formula.
\end{proof}

This proves a source-to-graph result for an explicit graded choice
of the GCM data. The construction specifies the independent Grassmann
fields and the maps that place them in the auxiliary curvature copies.
Its inner form is equivariant under the retained small Abelian sector
and depends on the specified multiplicity and evaluation maps.
The propagating even quadratic form is positive on the zero-mode-fixed
transverse slice. An unbroken larger simple representation or a
uniquely selected full-curvature Hodge pairing need not admit this
projection. Channel parity, representation multiplicity and
Grassmann statistics have distinct roles. If the raw field algebra
is generated solely by commuting tensors, ordinary parity-preserving
algebraic operations and derivatives generate only even elements.
A choice of supermatrix indices cannot by itself supply the
independent Berezin variables used in this construction. Likewise,
the physical unrestricted action would have to select the indicated
jet domain, endpoint maps, curvature pairings, cyclic word trace and
boundary state, and control the remaining quantum fields on the same
branch. Agreement of the tree Dirac term or the boundary pole with
raw O--QED leaves those issues open. The physical parameter count,
Q1 representation and Lorentzian continuation with the directed
sector-plane boundary prescription require their own derivations.

\subsection{Exact charged elimination with an ordinary dynamical Maxwell field}
\label{subsec:g5M-parent}

A charged propagator can have the local Dirac principal symbol while its
interaction vertices sample a positive normal distance. The following parent action realizes this possibility and permits
the closed determinant and the open charged line to be tested within
the same electromagnetic source.

Work in reference units on a Euclidean four-torus of period $2\pi$
on one $\bb=(s,b_s)$. Use one massive vectorlike spinor,
$D_0=m+\gamma^\mu\partial_\mu$, $m>0$, and the ordinary
zero-mode-fixed transverse Maxwell covariance
$\Pi_{\mu\nu}(p)/|p|^2$ for $p\ne0$. Let $L>0$ be the recorded
normal separation. For a given boundary charged field $\psi$, require
a bulk charged field $\chi$ to solve
$(-\partial_y^2-\Delta)\chi=0$ on $y>0$ with
$\chi(0)=\psi$ and bounded extension. Its evaluation at $L$ is
$P_L\psi$ with $P_L(k)=e^{-L|k|}$. Indeed each nonzero Fourier mode
satisfies $f''=|k|^2f$ and the zero mode extends constantly.
The heat-lapse representation of this transfer is
\eqref{eq:uvpeC-phase}.

Here is a quantum implementation in which the determinant of the
constraint is calculable. Work first with finite carrier modes. Write
$H\chi=J\psi$ for a linear regulator of the normal Dirichlet
problem, with $H$ invertible on the unknowns, and let $C$ evaluate at
$L$. Eliminate the normal variables and take the normal limit at fixed
carrier cutoff so that $CH^{-1}J\to P_L$. For example on a finite
interval use the exact decaying boundary condition at its far end for
each nonzero mode (Neumann for the constant mode), followed by
convergent one-dimensional Dirichlet discretization.
With independent Berezin variables $\chi,\lambda$ and their duals,
the parent charged action on the transverse chart is
\begin{align}
 S_{\rm ch}={}&\bar\psi D_0\psi+
      \bar\lambda(H\chi-J\psi)
      +(\bar\chi H^\dagger-\bar\psi J^\dagger)\lambda\nonumber\\
   &+\ii e\,\bar\chi C^\dagger (\gamma\cdot A^T) C\chi .
 \label{eq:g5M-parent}
\end{align}
The local vertex at the visited layer is an ordinary vector coupling,
linear in the unsmoothed transverse photon. A small-gauge version on
the trivial bundle writes $A=A^T+\dd\xi_A$ and uses the dressed
spinor $e^{\ii e\xi_A}\psi$ and its dual for the transverse variables.
Transport the regulator and source maps with this change of variables.
The residual constant gauge rotation cancels in the action and the
two Berezin Jacobians cancel. The following graphs use transverse
photons and dressed or fixed-gauge charged probes.

Translate the bulk integration variables by $H^{-1}J\psi$ and its
dual, or take the Schur complement. The hidden matrix has the form
\[
 \begin{pmatrix}
 B_A&H^\dagger\\ H&0
 \end{pmatrix},\qquad
 B_A=\ii e C^\dagger(\gamma\cdot A^T)C.
\]
Its determinant is $(-1)^{\dim H}|\det H|^2$ at every $A^T$, as
seen by exchanging its two block columns. The upper left block of its
inverse is zero. Consequently the single free normal normalization
removes this factor for every photon source. The reduced charged
operator, including its determinant, is exactly
\begin{equation}
 D_A^{(L)}=D_0+\ii e P_L(\gamma\cdot A^T)P_L .
 \label{eq:g5M-induced}
\end{equation}
This also follows by applying an external charged source before the
elimination. The reduced Gaussian then gives
$\det D_A^{(L)}$ times its usual exponential in $(D_A^{(L)})^{-1}$.
At a photon vertex with incoming and outgoing charged momenta $k,k'$,
the vector coefficient is $\ii e\gamma_\mu
e^{-L(|k|+|k'|)}$. Thus the soft symbol has the canonically
normalized minimal vector coupling and the free inverse is precisely
Dirac. The two normal factors are fixed by the original charged
readout and its dual. This comparison has no interacting pure-photon
vertices or photon-dependent ghost determinant.

\begin{theorem}[All electromagnetic orders after the charged elimination]
\label{thm:g5M-Dirac}
For \eqref{eq:g5M-parent}, at each fixed order in $e$, every bare
interacting vacuum or fixed-external-mode graph with internal charged
and Maxwell lines is absolutely UV finite after the specified normal
pushforward. With one fluctuating $L>0$ common to the entire source,
the same holds if its normalized free law has all inverse moments
near zero and the constraint normalization above is retained.
For fixed $L>0$ the relative zero-source partition has an entire
extension in $e$ and is strictly positive for real $e$.
\end{theorem}
\begin{proof}
An internal fermion line between interacting vertices carries
$e^{-2L|q|}D_0(q)^{-1}$, also for a same-vertex contraction. Its
Dirac factor has norm $(m^2+|q|^2)^{-1/2}$. Every photon momentum is
the difference of the two ordered matter momenta at its vertex. With
fixed external momenta, the map from a consistent loop routing to
the list of internal matter momenta is injective: their differences
fix all internal photon momenta. Each nonzero Maxwell line has norm
at most one. Dropping momentum constraints after taking absolute
values bounds the graph with $I$ internal matter lines by a fixed
constant times
\[
 \left(\sum_{q\in\mathbb Z^4}
     \frac{e^{-2L|q|}}{\sqrt{m^2+|q|^2}}\right)^I
 \le C_G(1+L^{-3})^I .
\]
The shell count in four dimensions gives the last inequality. With
no internal matter lines there is no free loop direction. A positive
power of $1+|q|$ can be inserted on each line at the cost of the
corresponding additional inverse powers of $L$. This yields uniform
absolute tails, also under the stated common depth law. The proof is
at fixed electromagnetic order with no separate visit parameter.

For the last assertion put $X=D_0^{-1}P_L(\gamma\cdot A^T)P_L$.
At photon mode $p$ the coefficient of $X$ is a weighted shift of
the charged lattice. Its trace norm is bounded, up to a spin
constant, by
\[
 \sum_k\frac{e^{-L(|k|+|k-p|)}}{\sqrt{m^2+|k|^2}}
 \le C(1+L^{-3})e^{-L|p|/2}.
\]
A Maxwell Fourier coefficient has standard deviation at most
$|p|^{-1}$ per component. The sum of these standard deviations times
the last exponential is finite for each fixed $L>0$. Therefore the
series for $X$ converges absolutely in trace norm almost surely, and
its cutoff sums are dominated by a random variable $Y_L$ with
$\mathbb E e^{aY_L}<\infty$ for every finite $a>0$. Indeed normalize
the summable deterministic coefficients to a probability sequence
and apply Jensen's inequality to the exponential; the absolute
standardized Gaussian components have a common exponential moment
at the resulting fixed argument. This needs no independence between
$p$ and $-p$. The same estimates without the Dirac denominator show
that $P_L(\gamma\cdot A^T)P_L$ is a bounded self-adjoint operator
almost surely for the real field.

The trace-class operator $X$ defines the full Fredholm determinant
$\det(I+\ii eX)$ directly. Its modulus is bounded by $e^{|e|Y_L}$,
which is integrable uniformly on compact sets of complex $e$.
Holomorphy for trace-class determinants and the common bound give
an entire relative partition~\cite{SimonUVpeDet}. Increasing
Fourier compressions and compatible photon truncations converge to
this determinant in expectation by the same bound.

For real $e$, write $B_L=P_L(\gamma\cdot A^T)P_L$ and
$X=D_0^{-1}B_L$. The operator $D_A^{(L)}$ is $m$ plus a
skew-adjoint operator. A kernel vector of $I+\ii eX$ lies in the
domain of $D_0$ and would contradict the real-part bound $mI$.
The Fredholm determinant is therefore nonzero. It is real because
$\gamma_5$ conjugates $\ii e D_0^{-1}B_L$ to
$-\ii e(D_0^\dagger)^{-1}B_L$. Reversing the order of the two
bounded factors $-\ii e(D_0^\dagger)^{-1}$ and $B_L$ gives the
adjoint of $\ii eX$ and preserves the nonzero eigenvalues with
their algebraic multiplicities. The determinant is continuous in
real $e$ and equals one at zero. It follows that it is strictly
positive for every real $e$, and so is its expectation.
\end{proof}

An unconstrained bulk Dirac field at the visit would produce its
own photon determinant and change this proof. A source-dependent
normal operator $H_A$ would also require the complete contributions
of $\det H_A$ and its dual. The example derives an electromagnetic
theory from the displayed constrained parent with ordinary dynamical
photons. Identifying it as the electromagnetic pushforward of
\eqref{eq:uvpeA-curvature} requires deriving this constraint, the
minimal visited-layer coupling, its longitudinal source dressing and
the boundary preparation from the physical GCM action. The mild
nonlocality of the charged vertices and the fixed-gauge Euclidean
proof also leave Q1 reconstruction open.
\paragraph{A mode test on general charged blocks before photon sewing}
There is a sufficient test which uses the actual matrices of the
parent, including hidden ones. On a charged Gaussian sector of a
fixed finite regulator write
$B_\Lambda(A)=B_\Lambda(0)(I+X_\Lambda(A))$ with invertible
$B_\Lambda(0)$. This matrix must describe an actual Gaussian
factor in the finite-regulator source, possibly conditional on
shared geometry. If it follows a normal integration, that elimination
must preserve Gaussianity in the charged variables and retain its
determinant and measure. An average of charged matrices over geometry
generally gives neither the averaged determinant nor the open-line
source. Take variations in a common orthonormal mode basis. For $j\ge1$ denote by $X_{\boldsymbol p}$
the derivative with $j$ labelled transverse photon modes; the
polarizations are suppressed. A finite number of charged or
statistics-weighted determinant blocks can be treated separately.
For the physical photon expansion suppose each marked photon
counts one perturbative power and only finitely many terms enter
the specified coefficient.

Let $h(p)=h(-p)\ge0$ be a scalar envelope. Suppose for each $j$
that the following two bounds hold with constants independent of
the regulator:
\begin{align}
 \|X_{\boldsymbol p}\|_2
     &\le C_j\prod_{a=1}^j h(p_a),\nonumber\\
 \sum_v |(X_{\boldsymbol p})_{vv}|
     &\le C_j'\prod_{a=1}^j h(p_a).
 \label{eq:g5M-HSmarks}
\end{align}
Here $\|\cdot\|_2$ is the Hilbert--Schmidt norm. Require
Hilbert--Schmidt convergence and $\ell^1$ convergence of the
diagonal at fixed photon labels under the regulator embeddings.
With a finite number of spin and polarization components, suppose
the transverse photon covariance obeys
\begin{equation}
 \sum_{p\ne0}\|D(p)\|_{\mathrm{op}}h(p)^2<\infty .
 \label{eq:g5M-sewing}
\end{equation}
Then every electromagnetic graph generated by these blocks and a
free photon measure has an absolutely convergent bare mode sum,
including all contact, open-line and determinant terms. For open
lines assume the end probes, including their free inverse where
it occurs, converge in $\ell^2$. Apply the estimates to any
additional scalar contact vertices from the measure or boundary
state as absolute scalar bounds with the same photon envelope.
Any shared fluctuating data require joint integrability of the
product of the costs for the specified graph. Integrate a conditional
graph with its full zero-field weight. If a connected logarithm is
taken after that integration, its derivatives also contain the
corresponding cumulants, as in \eqref{eq:g5I-two-loop}. The same
joint bounds must apply to the products in these cumulants.

Indeed expand $\log\det(I+X)$ at $X(0)=0$. A derivative of order
$j$ is a finite sum of cyclic traces of products of the true
derivative matrices, including the order-$j$ single factor. For a
product of at least two factors replace each matrix entry by its
absolute value. This leaves its Hilbert--Schmidt norm unchanged.
The trace norm of a product of two Hilbert--Schmidt operators is
bounded by the product of their norms. Each remaining factor has
operator norm at most its Hilbert--Schmidt norm. Apply this bound
first to a common finite set of indices with each entry replaced by
its modulus. The finite traces then increase to the absolute sum
of the cyclic word, so the same bound holds on the full mode set. The single-factor term needs precisely
the additional diagonal estimate in \eqref{eq:g5M-HSmarks}.
The inverse expansion for open chains has the same entrywise
argument with $\ell^2$ endpoint vectors. Wick pairing two photon
slots costs at most \eqref{eq:g5M-sewing} after dropping routing
constraints. For fixed photon labels a telescoping difference
bounds the cutoff error of each product by the stated norm
convergence. The summable photon majorants then give convergence
after sewing. A uniformly bounded value of a signed photon trace
alone would not give these entrywise estimates.

For ordinary Maxwell $D(p)=\Pi(p)/|p|^2$ on this four-torus,
a two-sided isotropic envelope $h(p)\asymp(1+|p|)^{-s}$ satisfies
\eqref{eq:g5M-sewing} precisely for $s>1$; the endpoint gives a
logarithm. This threshold tests the sufficient envelope criterion,
and individual graphs can have more decay. In
\eqref{eq:g5M-induced} the one-photon trace-norm estimate in
Theorem~\ref{thm:g5M-Dirac} supplies both tests with exponential
$h$; there are no higher derivatives on the transverse chart.
In a native GCM reduction the matrices in \eqref{eq:g5M-HSmarks}
must be computed from its own determinant and inverse. In
particular a dependent constraint Jacobian or BV--BFV contact
must enter before invoking the estimate.

\paragraph{What smoothing of the charged operator implies for a sewn photon}
Consider an exactly reduced Gaussian block on a fixed transverse torus
chart. Suppose its relative perturbation and its first two variations
in each fixed smooth photon direction converge in trace norm on the
charged space, and that the reference inverse factor in
\eqref{eq:uvpeA-relative} has a bounded limit. With strong inverse
convergence and a uniform norm bound, the two traces in that formula
have finite limits at each fixed photon momentum. The trace is over
charged indices. A photon contraction sums a different index. Also,
if the reference determinant depends on the photon field, its own
Hessian must be restored when forming the full response.

There is an exact test for the quadratic part of this second sum.
Take a real translation-invariant reduced source on the fixed support.
Let $H(p)$ be the Hermitian matrix on the transverse polarization
space in its quadratic action, after any specified combinations
of the contact and the current pair. Use the convention
$\frac12\sum_{p\ne0} A(p)^* H(p) A(p)$ with $A(-p)=\overline{A(p)}$.
Let $D(p)$ be the positive semidefinite free photon covariance and put
$M(p)=D(p)^{1/2}H(p)D(p)^{1/2}$. At fixed nonzero momenta these
are finite matrices. Consider connected vacuum rings made exclusively
from copies of this quadratic vertex, with the finite polarization
trace taken before the absolute sum over the ring momentum. Every
such ring converges absolutely if and only if
\begin{equation}
 \sum_{p\ne0}|\operatorname{Tr} M(p)|<\infty,\qquad
 \sum_{p\ne0}\|M(p)\|_2^2<\infty .
 \label{eq:s6A-photon-ring}
\end{equation}
Here $\|\cdot\|_2$ is the matrix Hilbert--Schmidt norm. Indeed, a ring
with $n$ insertions has momentum coefficient proportional to
$\operatorname{Tr}M(p)^n$, by Wick contraction and cyclicity. The
finite nonzero symmetry factor is independent of $p$. The first two
rings give exactly the two tests because $M(p)$ is Hermitian. For
$n\ge2$ the absolute trace is bounded by
$\|M(p)\|_{\rm op}^{n-2}\|M(p)\|_2^2$, with the factor at $n=2$
understood as one. The second sum bounds $\sup_p\|M(p)\|_{\rm op}$.
In a positive channel $M(p)\ge0$ the first sum already implies
the second and is exactly the trace-class condition for the direct
sum on the photon space. In a signed channel polarization
cancellations in the one-insertion ring can leave a divergent
positive two-insertion ring. For regulator-dependent matrices,
convergence on each finite set of modes together with uniformly
vanishing tails in both absolute sums gives the stated
cutoff limits at every ring order.

This criterion concerns the grouped quadratic vertex and vacuum
graphs when these are included in the bare expansion. Entrywise
absolute convergence before the charged trace, higher photon
contacts, and graphs mixing those contacts with open charged lines
require their own estimates. Equations~\eqref{eq:g5M-HSmarks}
and~\eqref{eq:g5M-sewing} give one sufficient joint test. If
the reference is the ordinary minimally coupled Dirac sea, the
source-dependent bare logarithm in \eqref{eq:uvpeA-raw-log}
must also be accounted for before that test can be applied.

\subsection{Inducing the first-order charged block by the same curvature pairing}
\label{subsec:g5R-direct}
The preceding constrained parent can also be induced by the represented
curvature of Section~\ref{sec:g5-induction}. This gives a second example
with a first-order retained Dirac operator, a Maxwell photon and a single
expansion parameter $e$. We describe all charged fields, since omitting
one charged determinant could reverse a finiteness conclusion.

At finite regulator set $\Xi=(\psi,\chi,\lambda)$ with independent
Berezin pairs for all three entries, and let $\mathcal D_A$ be the
matrix of the parent bilinear in \eqref{eq:g5M-parent}. Use the fixed
spin involution $G=\operatorname{diag}(\gamma_5,\gamma_5,\gamma_5)$.
On a real transverse field $\mathcal D_A^\dagger=G\mathcal D_A G$:
the massive Dirac block and the skew interaction have this symmetry,
and $H,J,C$ are spin scalars with the indicated duals. We now produce
$\bar\Xi G\mathcal D_A\Xi$ from copies of \eqref{eq:g5-arrow-product}.
The common curvature pairing is fixed before any source is varied.

\begin{enumerate}
\item For the retained field choose row $\bar\psi$, with column $\psi$
and endpoint spin matrix $m\gamma_5$ for the mass. For each $\mu$
choose column $\partial_\mu\psi$ and endpoint matrix
$\gamma_5\gamma^\mu$ for the kinetic term. The connecting even
coefficient is the identity in these copies.
\item Use $\bar\lambda,\chi$ for the constraint row and column and
$\bar\chi,\lambda$ for its adjoint. At a normal lattice point each
second difference is the sum of two nearest-neighbor arrows and
one on-site term with their finite-difference weights. The spatial
Laplacian can equivalently be represented by first spatial derivatives
on both entries and integration by parts. The boundary neighbor is
the retained spinor $\psi$; moving it to the right in $H\chi=J\psi$
produces exactly the two $-J,-J^\dagger$ terms in the parent.
Use the spin matrix $\gamma_5$ multiplying these scalar coefficients.
In particular the normal matrix is assembled from geometric
differences without a dependence on the photon field.
\item At the visited layer use row $\bar\chi(L)$, column $\chi(L)$,
even coefficient $A_\mu^T(x)$ times the identity and endpoint spin
matrix $\ii e\gamma_5\gamma^\mu$ for each $\mu$. At finite normal
mesh both evaluations use the same sampling map $C$ and its dual.
The ordinary Abelian curvature pairing for the independent photon
supplies $\frac14\int F_{\mu\nu}F_{\mu\nu}$. No normal factor is
placed on its Maxwell covariance.
\end{enumerate}
Summing the real pairings gives the claimed bilinear on the real
Grassmann form because $G\mathcal D_A$ is Hermitian. The constant
linear change of every dual spinor by $\gamma_5$ recovers the
independent Euclidean sources and the full matrix $\mathcal D_A$.
The repeated rows and columns in the curvature copies are images of
these fields. There is one Berezin integration per owner. The
prescribed cycle, trace and transverse small-gauge chart are those
of the construction. The positive Euclidean Maxwell measure and
the Grassmann integral play different roles.

The normal regulator can be chosen without a momentum-dependent
term at its far boundary. Impose Neumann data at $y=R>L$, Dirichlet
data at $0$, and first take the normal mesh limit at fixed carrier
cutoff. At mode $k\ne0$ the transfer is
\[
 \frac{\cosh((R-L)|k|)}{\cosh(R|k|)}
   \le 2e^{-L|k|}.
\]
It tends to $e^{-L|k|}$ as $R\to\infty$; the zero mode extends
constantly. For example, a mesh of spacing $a$ solves the centered
interior difference equation with exponent $\kappa_a$ given by
$\cosh(a\kappa_a)=1+a^2|k|^2/2$. Reflection at the far boundary
gives the same cosh ratio with $|k|$ replaced by $\kappa_a$, and
$\kappa_a\to|k|$ at each fixed mode. One may interpolate at $L$
or take neighboring mesh readouts converging to it. At every such
regulator the hidden determinant is the same function of $H$ as at
zero source. This is why removal of the free normalization does
not remove a photon contact.

\begin{corollary}[Bare all-orders finiteness of an explicitly induced electromagnetic sector]
\label{cor:g5R-induced-Dirac}
For the represented curvature, independent field measure, relative
free ($e=0$) normalization and normal pushforward just specified, at
fixed $L>0$ every interacting graph at every finite electromagnetic
order is absolutely UV finite, including closed charged loops, open
charged lines with fixed Fourier probes and dynamical Maxwell sewings.
No UV counterterm is needed. The conclusion also holds graphwise
for a shared normalized positive-depth law having all inverse
moments near zero. At fixed $L$ the relative partition is finite
and strictly positive for real $e$.
\end{corollary}
\begin{proof}
The curvature computation gives the finite-regulator parent in
\eqref{eq:g5M-parent} before integrating any hidden charged owner.
Its exact determinant reduction induces \eqref{eq:g5M-induced}.
The two visits attached to every internal charged segment supply
$e^{-2L|k|}$. Charge conservation at the three-point vertex fixes
each internal photon momentum as a difference of charged momenta.
Theorem~\ref{thm:g5M-Dirac} then gives summability, uniform UV
tails, and the fixed-$L$ determinant conclusion. A normalized
law on $L$ is integrated once for the entire source. The upper
bound is integrable with that same law.
\end{proof}

The proof uses a field-independent mixed pairing on selected curvature
components, a finite jet representation, the harmonic charged
constraint and its multiplier measure, a rule for taking the normal
limit, and a specified visit observable. These are constitutive
choices of the chart. The physical GCM formulation asks for their
selection from its raw geometry, its chiral field content and its
BV--BFV state. The corollary is an existence theorem in the stated
class of curvature representations. It is not an assertion that a
physical positive Hodge pairing on every coefficient permits all
of these choices.

\paragraph{The longitudinal vertex and the soft photon test}
The difference between the induced current of this example and the raw
Dirac current can be computed explicitly. On the small Abelian chart
let $A=A^T+\dd\xi_A$, with the constant mode fixed, and restore the
bare charged variables by conjugating
\eqref{eq:g5M-induced} with $M=\exp(-\ii e\xi_A)$. The full kernel is
$M D_{A^T}^{(L)}M^{-1}$. At zero field its first variation consists
of the commutator $[-\ii e\,\delta\xi_A,D_0]$ and the variation of
the sampled transverse field. Write $p=k'-k\ne0$,
$\Pi^\parallel_{\mu\nu}=p_\mu p_\nu/|p|^2$ and
$\Pi^T=I-\Pi^\parallel$. With the factor $\ii e$ removed the vertex is
\begin{equation}
 \mathcal V_\mu(k',k)=\gamma^\nu\bigl[
 \Pi^\parallel_{\nu\mu}
  +e^{-L(|k'|+|k|)}\Pi^T_{\nu\mu}\bigr].
 \label{eq:g5R-Ward-vertex}
\end{equation}
Thus $\ii p^\mu\mathcal V_\mu=D_0(k')-D_0(k)$ exactly. In
transverse gauge only the suppressed part is sewn. The unsuppressed
longitudinal term acts on the external gauge variation. The commuting
Abelian change of charged coordinates with its transported cutoff
also makes the entire determinant independent of $\xi_A$. In a
calculation with unreduced probes the source map must be conjugated
as well.

On the continuous symbol, the difference between this vertex and
the minimal Dirac vertex tends to zero when \emph{both} charged
momenta tend to zero. If only the photon is made soft at fixed
$k\ne0$, the transverse coefficient tends to $e^{-2L|k|}$ and
the longitudinal projector retains its direction dependence.
This comparison of symbols makes no infinite-volume assertion about
the fixed torus. The distinction matters to the physical interpretation
of the finite induced loops.

\paragraph{A soft-photon test of the same finite charged determinant}
Absolute ultraviolet convergence and the physical massless pole can
be tested separately in \eqref{eq:g5M-induced}. Consider the continuum
momentum integral per unit four-volume for its quadratic photon
response, with flat physical transverse probes $e_2 e^{\pm\ii p x^1}$,
$p\ne0$. This compares arbitrarily small nonzero external momenta; on
one fixed torus the allowed nonzero momenta are discrete. With the
sign convention $-\log\det D_A^{(L)}$, the mixed derivative at $A=0$
is, with $k'=k+pe_1$ and $S(k)=D_0(k)^{-1}$,
\[
 -e^2\int\frac{\dd^4k}{(2\pi)^4}\,
 e^{-2L(|k|+|k'|)}
 \operatorname{Tr}_{\rm spin}
 [S(k)\gamma_2 S(k')\gamma_2].
\]
The normalized determinant has no transverse two-photon contact in
this chart. For $m,L>0$ dominated convergence applies as $p\to0$.
Using $D_0(k)^{-1}=(m-\ii\gamma\cdot k)/(m^2+|k|^2)$, the spin
numerator at $p=0$ is $4(m^2+|k|^2-2k_2^2)$. Angular averaging
on the unit three-sphere replaces $k_2^2$ by $|k|^2/4$.
The limit of the transverse inverse-kernel correction is therefore
\begin{equation}
 -\frac{e^2}{2\pi^2}\int_0^\infty
  r^3 e^{-4Lr}\frac{m^2+r^2/2}{(m^2+r^2)^2}\,\dd r .
 \label{eq:g5M-soft-mass}
\end{equation}
It is finite and strictly negative for $e\ne0$. At fixed $m$,
multiplication by $L^2$ and rescaling $r$ by $L^{-1}$ gives the
limit $-e^2/(64\pi^2)$ as $L\downarrow0$. The massless Maxwell
quadratic form itself contributes a term of order $p^2$.
Thus this completion generates a finite masslike transverse term
already at electromagnetic order two. Small-gauge covariance of the
nonlocal transverse dressing permits it. Equation~\eqref{eq:g5M-soft-mass}
is a two-point coefficient, with no inference here of an exact pole
mass or a reconstructed spectrum. Any derivation that also claims
the physical massless photon from the same GCM branch must exhibit
the additional source response cancelling this constant, or a
different charged vertex for which it vanishes.

\paragraph{A necessary and sufficient common-depth condition at all orders}
The nonzero coefficient also sharpens Theorem~\ref{thm:g5M-Dirac}
on the torus itself. For the model \eqref{eq:g5M-induced} let
$c_L(p)$ be its one-bubble transverse kernel per unit four-volume
at a fixed nonzero axial lattice momentum $p e_1$. Use
$(2\pi)^{-4}\sum_k$ in the preceding formula. Then
\[
 \lim_{L\downarrow0}L^2c_L(p)=-e^2/(64\pi^2).
\]
Indeed put $u=Lk$. On any compact annulus the Riemann sums converge
to the same integral with $Lp,Lm\to0$. The two inverse Dirac factors
are bounded by $C_{p,m}(1+|k|)^{-2}$ in product. Hence the
contribution from $|u|\le\eta$ to $L^2 c_L$ is at most
$C_{p,m}(\eta^2+L^2)$. Exponential bounds control the exterior
annuli uniformly. The signs and permutations of the last three
axes make the bubble a scalar on the transverse space at that $p$.

Now take the connected two-external-photon graph consisting of $n$
such bubbles linked in series. The joining photons all have this
same fixed nonzero momentum. Their transverse propagators give
$|p|^{-2(n-1)}$. All source vertices share the same $L$. At fixed
routed charged momenta, the complete scalar transfer is a positive
exponential in $L$ and the remaining spin and photon contraction is
independent of $L$. Consequently the absolute sum of this graph
after the common-depth average dominates, up to fixed graph and
volume factors,
\[
 |p|^{-2(n-1)}\int_{(0,1]} |c_L(p)|^n\,\mu(\dd L).
\]
To see this, take the absolute value of the $L$-independent
contraction, use Tonelli on its positive transfer, then use the
triangle inequality at each $L>0$. At $L=0$ a growing cutoff ball
already gives a quadratic divergence of the single bubble by the
same angular average. Thus every connected graph in this family
can be finite only if $\mu(\{0\})=0$ and every inverse moment of
$L$ near zero is finite: the order-$2n$ graph forces the inverse
$2n$ moment and $n$ is arbitrary. Conversely Theorem~\ref{thm:g5M-Dirac}
proves sufficiency for all graphs of this model. This necessity
uses connected graphs including one-particle-reducible chains, as
required by the bare graphwise objective. It does not follow by
charging a separate depth probability at each bubble.

The constant admits a wider test. Suppose the full order-$e$ transverse
vertex at zero photon momentum is $\ii e\gamma_\mu f(|k|)^2$ on
the massive Dirac block. Suppose the finite-momentum vertex tends to
it almost everywhere in the transverse component being tested and
has a common integrable two-vertex majorant near $p=0$.
For a bounded real radial $f$ let
\[
 B_f=\frac{1}{2\pi^2}\int_0^\infty
 r^3 |f(r)|^4\frac{m^2+r^2/2}{(m^2+r^2)^2}\,\dd r .
\]
When $f$ is nonzero on a set of positive radial measure and this
integral converges, $0<B_f<\infty$. The same angular trace gives
a soft two-vertex term $-e^2 B_f$. For a source that claims a
massless transverse photon order by order, the remaining bare
order-$e^2$ terms must therefore tend to $+e^2B_f$ in this
normalization, assuming those limits exist. This includes the direct
two-photon contact from the charged Hessian, field-dependent
Jacobians and any other fluctuating sectors in the same source.
The equation tests a massless pole's necessary zero-momentum
condition; its residue and analytic structure require further
information. In particular changing the radial attenuation on the
two-vertex diagram alone cannot set this positive integral to zero.

\paragraph{Restoring a positive mixed curvature can change the UV result}
There is a direct test of the unused mixed curvatures in this representation.
At one oriented link of fixed length $h\in[1,2]$ in direction $j$,
equip its two charged endpoint fibres with the same type of Abelian
connection $\partial_i-\ii e b_i$, evaluated at the respective
endpoints. This defines another chart, with the same visited field
and covariance $D_L$ as in \eqref{eq:uvpeW-depth} at $e=0$.
For its transporter $U=\exp(-\ii e\int_0^h b_j(x+t e_j)\dd t)$
the mixed curvature in directions $x^i,\theta^1$ contains
the covariant derivative
\[
 \nabla_i U=-\ii e\,U\int_0^h F_{ij}(x+t e_j)\,\dd t .
\]
Indeed the two endpoint connection terms subtract the integral of
$\partial_j b_i$ from the differentiated integral of $\partial_i b_j$.
Give this component its ordinary positive squared norm, in addition
to the terms of Theorem~\ref{thm:g5-curv-induced}. Since $|U|=1$,
its two-photon vertex is exactly $e^2\sum_i|\int F_{ij}|^2$ per
unit endpoint volume. Contract its two visited photons with the
same full covariance at $e=0$. Already on the slice $p_j=0$,
$0<|p|\le\Lambda$, the positive cost at fixed $L>0$ is comparable
at large $\Lambda$ to
\[
 h^2\!\sum_{\substack{p_j=0\\0<|p|\le\Lambda}}
  |p|^2(D_L(p))_{jj}
 \ \asymp_L\ \Lambda .
\]
Here $\Pi_{jj}=1$ on that slice, and the summand is comparable to
$h^2/(1+L|p|^2)$ on a three-dimensional lattice. This is an
individual divergent vacuum graph in the bare electric-coupling
expansion of the enlarged pairing. It persists at every fixed
positive separation. Its leading cost depends on the common $L$
and $h$ before normalization. A reformulation that integrates this
quadratic term exactly must retain its relative determinant in the
history weight. The calculation tests this specific extra component;
whether a raw GCM Hodge law includes it or couples its source to
different normal readouts is a further calculation.

\subsection{A normal kinetic action that supplies every required inverse moment}
\label{subsec:g5N-normal}
One marked Dirac loop tests an inverse-square moment of a midpoint
distance in \eqref{eq:uvpeP-Dirac-mark}. Photon sewing in
Corollary~\ref{cor:uvpeW-sharp-depth} tests increasing inverse
moments of the duration of the whole return. Distance and duration
have different laws. The following calculation derives both from one
preparation and allows visits arbitrarily close to zero.

In the marked bridge of \eqref{eq:uvpeW-return}, allow its positive
midpoint distance $r$ to fluctuate according to a probability
$\mu(\dd r)$ on $(0,1)$. With a common return in the entire
functional, the joint unnormalized measure in reference units is
\begin{equation}
 \frac{1}{2\pi L}e^{-L-r^2/L}\,\dd L\,\mu(\dd r),
 \qquad L>0 .
 \label{eq:g5N-bridge}
\end{equation}
Here the free normal line generates the return. A second radial
system provides the midpoint distribution $\mu$; coupling their
readouts at the same distance is part of the boundary state.
This expression multiplies both transition densities; fixing a
probability separately at every $L$ would change the source. For a
fixed integer $n>0$ the integral of $L^{-n/2}$ over $0<L<1$ is
finite precisely when $\int r^{-n}\mu(\dd r)<\infty$. In fact
$e^{-L}$ is bounded above and below there and
\begin{equation}
 \int_0^1 L^{-1-n/2}e^{-r^2/L}\,\dd L
    =r^{-n}\Gamma(n/2,r^2).
 \label{eq:g5N-moment}
\end{equation}
Here the incomplete gamma factor is bounded between positive
constants for $0<r<1$. At $n=0$ the normalization instead tests
$\int\log(1/r)\,\mu(\dd r)<\infty$. These equivalences concern the
stated path measure; a source-dependent midpoint operator must
also be differentiated.

\paragraph{Selecting a finite-action normal metric}
There is a restricted Euclidean Cartan curvature block for which the
metric of the preparation follows from a variational problem. Take a
rotational normal chart $(0,1)\times U(1)$ with unit proper-length
radial coordinate, metric $\dd r^2+f(r)^2\dd\theta^2$ and
angular Haar measure $\dd\theta/(2\pi)$. Use its torsion-free
$\mathfrak{so}(2)$ connection and the positive Hodge curvature pairing.
The connection form is $\pm f'(r)\dd\theta$ in an orthonormal frame,
with curvature $\pm f''(r)\dd r\wedge\dd\theta$. This even
curvature-square block therefore has radial action
\begin{equation}
 {\cal E}[f]=\int_0^1\frac{|f''(r)|^2}{f(r)}\,\dd r .
 \label{eq:g5N-normal-geometry}
\end{equation}
Hold the proper interval length and boundary preparation fixed.
Impose $f(0)=f'(0)=0$, $f(1)=1$, $f'(1)=4$ in the chosen reference
units. Take positive interior profiles with finite action, $C^1$ up
to the boundary and locally absolutely continuous first derivative.
The fixed endpoint slope is part of this preparation. The exact
minimum is $144$, uniquely at $f(r)=r^4$.

Indeed put $A(x)=\int_0^x |f''|^2/f$ and
$M(x)=\max_{[0,x]}f$. Integrating from zero and using
Cauchy--Schwarz gives $\sup_{[0,x]}|f'|\le\sqrt{xM(x)A(x)}$.
Hence $M(x)\le x^3 A(x)$ and
$\sup_{[0,x]}|f'|\le x^2 A(x)$, with $A(x)\to0$.
On $[\delta,1]$ complete the square and integrate by parts:
\begin{align*}
 \int_\delta^1 \frac{|f''|^2}{f}\,\dd r
    &=\int_\delta^1\frac{|f''-12f/r^2|^2}{f}\,\dd r\vphantom{\Big|}\\
    &\quad+24\left[\frac{f'}{r^2}+\frac{2f}{r^3}
                       \right]_{\delta}^{1}.
\end{align*}
The tip estimates make the lower boundary term vanish as $\delta\downarrow0$.
Equality in the resulting inequality requires $f''=12f/r^2$.
The solutions are linear combinations of $r^4,r^{-3}$, and the
boundary data select $r^4$. Its Gaussian curvature $-12/r^2$ is
unbounded at the excluded tip, but its integrated curvature
square is finite. The variational statement uses this singular
finite-action chart and the rotational, torsion-free class.

\begin{proposition}[A charged normal preparation at all photon orders]
\label{prop:g5N-fibre}
On this minimizing chart use the positive scalar kinetic form of the
metric, initially on smooth compactly supported radial data with its
Friedrichs closure. Let $j\ne0$ be an integer Fourier charge of the
angle. Take the radial probability in the lowest state of this
charge as $\mu$ in \eqref{eq:g5N-bridge}. The resulting probability
on $L$ has inverse moments of every finite order. The same conclusion
holds with the normalized heat state at any fixed $b>0$, in the
sector $j$ or the direct sum of all $j\ne0$ sectors. With either
preparation all fixed-order Wilson graphs in
Corollary~\ref{cor:uvpeW-sharp-depth} are absolutely UV finite.
The marked Dirac magnetic test is finite when the zero-visit
coefficient vanishes.
\end{proposition}
\begin{proof}
The radial form of the angular mode $e^{\ii j\theta}\psi(r)$ and
its norm read
\[
 \int_0^1\bigl(r^4|\psi'|^2+j^2r^{-4}|\psi|^2\bigr)\dd r,
 \qquad \int_0^1 r^4|\psi|^2\dd r .
\]
Conjugating to $u=r^2\psi$ and integrating by parts on the form
core yields the operator on $L^2(\dd r)$
\[
 H_j=-\partial_r^2+2r^{-2}+j^2r^{-8}.
\]
The charged term is the squared angular frequency divided by the
squared radius $f(r)^2$. It becomes large as the orbit shrinks,
pushing a charged state away from the tip. The closed form has
compact resolvent by the compact
embedding of $H^1_0(0,1)$ into $L^2(0,1)$. The lowest radial
eigenfunction $u_j$ can be chosen positive on $(0,1)$, by
minimizing the form on $|u|$ and uniqueness of the interior ODE
initial-value problem at any interior zero. Its squared modulus
is the normalized density of $\mu$.

Write $E$ for the ground energy and
$v_b(r)=r^2\exp[-b/(3r^3)]$ with $b>0$. It satisfies
\[
 (H_j-E)v_b=\big[(j^2-b^2)r^{-8}-E\big]v_b .
\]
For $0<b_-<|j|<b_+$ choose $r_0>0$ small. On $(0,r_0)$ the
potential minus $E$ is positive, and $v_{b_-},v_{b_+}$ are a
super- and subsolution, respectively. Match their values to
$u_j(r_0)$. The positive-part energy estimate for the difference,
whose traces at both endpoints vanish, gives
\[
 c\,r^2 e^{-b_+/(3r^3)}\le u_j(r)\le C\,r^2 e^{-b_-/(3r^3)},
 \qquad 0<r<r_0 .
\]
In particular every inverse radial moment is finite. More precisely,
put $N_j(L)=(2\pi L)^{-1}\int_0^1e^{-r^2/L}|u_j(r)|^2\dd r$.
Then
\begin{equation}
 \lim_{L\downarrow0}L^{3/5}\log[L N_j(L)]
           =-\tfrac53 |j|^{2/5}.
 \label{eq:g5N-rate}
\end{equation}
Indeed the minimum of $r^2/L+2b/(3r^3)$ is $(5/3)b^{2/5}L^{-3/5}$ at
$r=(bL)^{1/5}$. The upper barrier and this minimum give the
upper estimate, with an exponentially smaller contribution from
$r\ge r_0$. For the lower estimate integrate on a fixed small
interval about $b_+^{1/5}$ in $r/L^{1/5}$. Its length and the
factor $r^4$ contribute only logarithms. Shrink this interval,
then let $b_-,b_+$ tend to $|j|$.

For the heat-state assertion write $K_b^j(r,r)$ for the diagonal
heat kernel of $H_j$. A direct bound avoids summing individual
eigenfunction constants. For small $r$ put $s=r^5/|j|<b/2$.
In the killed Brownian motion with generator $\partial_r^2$,
the remaining evolution of duration $b-s$ has kernel at most
$(4\pi(b-s))^{-1/2}$. During $[0,s]$ a surviving path which stays
below $2r$ has heat weight at most $e^{-j^2s/(256r^8)}$.
The probability to reach $2r$ in that time from $r$ is at most
$2e^{-r^2/(4s)}$ by reflection. The heat-kernel representation
for a nonnegative potential (first truncate it, then pass by
monotone convergence) consequently gives
\[
 K_b^j(r,r)\le C_b e^{-|j|/(256r^3)} .
\]
Equivalently the bound follows by splitting the positive path
integral at time $s$~\cite{SimonG5Normal}. It is summable over
$j\ne0$ near zero. On the rest of the interval the free
heat bound times $e^{-b j^2}$ is summable since $r^{-8}\ge1$.
Normalizing the strictly positive finite heat trace therefore
gives all inverse radial moments, even for the full charged sum.
Use \eqref{eq:g5N-moment} to conclude.
\end{proof}

\paragraph{Suppression along a single shrinking interval}
The potential $j^2r^{-8}$ also controls arbitrary boundary vectors
when a normal heat interval is sewn onto them. This matters if the
preceding charged probe changes the radial state. In the same
Friedrichs realization on $(0,1)$ let $S_u=e^{-uH_j}$ for $u>0$
and integer $|j|\ge1$. For any inverse power $p>0$ there is a
constant $C_p$, independent of $u,j$, such that
\[
 \|r^{-p} S_u\|_{2\to2}
 \le C_p\big[(u j^2)^{-p/8}+|j|^{-p/3}\big].
\]
For each fixed nonzero $j$ the short-interval exponent $p/8$ is
sharp. The exponent $p$ specifies the inverse moment being tested;
it is independent of the couplings in the action.

To prove the bound, apply the heat path formula above to the survival
weight $S_u1(r)$. Follow the path for
$s=\min(u,r^5/|j|)$. A path that stays below $2r$ pays at least
$j^2s/(256r^8)$ in its integrated potential; the probability of
reaching $2r$ before $s$ is at most $2e^{-r^2/(4s)}$.
The remaining heat evolution has survival weight at most one.
For a fixed numerical $c>0$, this gives
\[
 S_u1(r)\le e^{-c u j^2/r^8}+3e^{-c|j|/r^3}.
\]
With the positive heat kernel, Cauchy--Schwarz gives
$|S_u h|^2\le(S_u1)S_u(|h|^2)$. The kernel is symmetric and its
total mass from each starting point is at most one. Integrating
therefore bounds $\|r^{-p}S_uh\|_2^2$ by
$\sup_r[r^{-2p}S_u1(r)]\,\|h\|_2^2$. Maximizing each exponential
times $r^{-2p}$ gives the stated powers.

For the lower bound set $\ell=(u j^2)^{1/8}$. Choose a unit smooth
test vector $h_\ell(r)=\ell^{-1/2}h(r/\ell)$ supported in
$(\ell,2\ell)$, for $u$ sufficiently small that $2\ell<1$.
Then $u\langle H_j\rangle_{h_\ell}\le C(u/\ell^2+1)$ stays
bounded. Jensen's inequality for its spectral measure gives
$\langle h_\ell,S_u h_\ell\rangle\ge
e^{-u\langle H_j\rangle_{h_\ell}}$. Since the test vector is
supported below $2\ell$, Cauchy--Schwarz gives a lower bound
$c_p\ell^{-p}$ on $\|r^{-p}S_u\|_{2\to2}$.

This estimate comes from the charged kinetic operator of the
minimizing metric. It holds uniformly over the incoming unit vector.
In a finite chain of such intervals joined by bounded operators
preserving the specified sector, submultiplicativity bounds the
chain by the product of these interval bounds and the joining
operator norms. Thus one insertion costing $r^{-p}$ uses at most
$u^{-p/8}$ of the short-duration budget for that interval, at fixed
$j$. In particular, if its absolute duration density is bounded by
$C u^{\beta-1}$ near zero, $\beta>p/8$ suffices for this factor.
A shared duration incurs the sum of the allocated powers. Equality
gives a logarithm in this norm bound; the actual paired matrix
element can be smaller. Applying the result to bare physical graphs
requires identifying the joining operators, their sector changes
and the joint duration density in the charged functional. No
choice of a momentum cutoff enters the estimate.

\paragraph{What a complete angular trace costs}
An operator bound on a single angular sector does not count how many
sectors a closed loop can occupy. On the same quartic surface this
count can be made before the duration integration. Use the periodic
circle and the Friedrichs radial kernels $K_u^j$ from the preceding
proof. The index $j$ is an angular Fourier number. An electric charge
or a net geometric-time winding does not by itself specify $j$.

Consider a cyclic word of $n\ge1$ positive intervals of lengths
$u_i$, with $U=\sum_i u_i$. At their junctions insert
$r^{-p_i}e^{i\ell_i\theta}$, with $p_i\ge0$ and integer shifts
$\ell_i$ of sum zero. This describes exact scalar radial factors; an
absolute upper envelope for an actual source gives only the upper
bounds below. Write $s_i$ for the cumulative shifts on the successive
intervals. Starting in $j$, the radial path evolves there with
$H_{j+s_i}$. Let $W_*(\boldsymbol u)$ be the sum over starting
sectors for which all $j+s_i$ are nonzero. Each summand is the
nonnegative cyclic integral of the radial kernels and the weights
$r_i^{-p_i}$. Put $P=\sum_i p_i$. For each fixed word, as
$U\downarrow0$,
\[
 W_*(\boldsymbol u)\asymp F_P(U),\qquad
 F_P(U)=
 \begin{cases}
 U^{-1}, &P<5,\\
 U^{-1}(1+|\log U|), &P=5,\\
 U^{-(P+3)/8}, &P>5 .
 \end{cases}
\]
Both comparison constants are independent of the fractions $u_i/U$.
In particular this estimate covers arbitrarily unequal shrinking
intervals.

For the upper bound, fix an allowed $j$ and let
$h=\min_i|j+s_i|$. Positivity of the heat path representation bounds
each interval kernel by $K^{h}_{u_i}$, since reducing the angular
potential increases its weight. At most $2n$ initial sectors can have
a given positive minimum $h$: one interval must attain $+h$ or $-h$.
For $P>0$ the weighted arithmetic--geometric mean inequality gives
$\prod_i r_i^{-p_i}\le\sum_i(p_i/P)r_i^{-P}$; when $P=0$ the product
is one. Integrating the other positions composes their heat kernels.
Hence it suffices to bound a constant times
\[
 \sum_{h\ge1}\int_0^1 r^{-P}K_U^h(r,r)\,\dd r .
\]
Split the heat evolution at $\min(U/2,r^5/h)$. The remaining kernel
is at most $C U^{-1/2}$, and the survival estimate above yields
\[
 K_U^h(r,r)\le C U^{-1/2}
    \left(e^{-cUh^2/r^8}+e^{-ch/r^3}\right).
\]
For $x>0$, comparison with the Gaussian integral gives
$\sum_{h\ge1}e^{-c x^2 h^2}\le C x^{-1}e^{-c x^2/2}$.
The $e^{-ch/r^3}$ terms are summable and integrable at the tip
against every fixed inverse power. The remaining bound is
\[
 C U^{-1}\int_0^1 r^{4-P}e^{-c' U/r^8}\,\dd r + C_P U^{-1/2}.
\]
Splitting at $r=U^{1/8}$ gives the three cases displayed above.

For the reverse inequality start a free radial Brownian bridge of
duration $U$ at $r$ and return to $r$. Restrict it to $(r/2,3r/2)$
throughout. For $A U^{1/8}\le r\le1/2$ its probability of remaining
there is bounded below uniformly for small $U$, by Brownian scaling;
the free diagonal density is $(4\pi U)^{-1/2}$. For $A$ large enough
depending on the finite list of shifts, there are at least
$c r^4/\sqrt U$ integers $j$ with
$0<|j+s_i|\le C r^4/\sqrt U$ for every $i$. On these paths the
integral of the full potential, including $2/r(t)^2$, is bounded
above by a constant. The product of the insertions is at least
$c_P r^{-P}$. Integrating over the displayed starting radii gives
$W_*\ge c U^{-1}\int_{A U^{1/8}}^{1/2}r^{4-P}\,\dd r$, as required.
The boundary between the three cases is where the available number
of angular modes, proportional to $r^4/\sqrt U$, just compensates
the inverse powers of radius.

For example, suppose the joint duration weight of this scalar word
is comparable to $\prod_i u_i^{\beta_i-1}\dd u_i$ on a small cube,
with every $\beta_i>0$. Changing variables to $U$ and the duration
fractions gives a finite simplex integral. This corner is integrable
exactly when
\[
 \sum_i\beta_i>\max\{1,(P+3)/8\}.
\]
At equality the duration integral diverges; when $P=5$ the angular
sum supplies an additional logarithm. With only an upper bound on
the duration density, the strict inequality remains sufficient.
Neither the heat intervals nor the Fourier sum in this calculation
are regulators. Both occur in the trace of the normal kinetic action.

\paragraph{When a word reaches zero angular momentum}
For the finitely many initial $j$ with some $j+s_i=0$, the preceding
count does not apply. The same kernels give an exact test at
\emph{fixed positive} interval lengths. An insertion between two
angular-zero intervals is integrable precisely for $p_i<5$.
An insertion adjoining a nonzero interval admits every finite
$p_i$. The whole positive scalar word is finite if and only if all
its junctions pass this local test.

Indeed, for zero angular momentum, multiplication by $r^2$
identifies radial $L^2(r^4\dd r)$ on the five-dimensional unit ball
with $L^2(\dd r)$, up to angular normalization. The Dirichlet
heat kernel of the ball is smooth and positive at its interior
center. Angular averaging and conjugation multiply that kernel
by $r^2x^2$ in the radial variables $r,x$. Thus
$K_t^0(r,r)$ is of order $r^4$ there, and
$K_t^0(r,x)$ is of order $r^2$ uniformly for $x$ in any compact
interior radial interval, at each fixed $t>0$. By the semigroup
identity and Cauchy--Schwarz,
\[
 K_t^k(r,x)\le [K_t^k(r,r)K_t^k(x,x)]^{1/2}.
\]
Each neutral side of a junction therefore contributes $O(r^2)$;
a charged side has exponential suppression at the tip. This proves
sufficiency, with independent position integrations. Necessity follows
by keeping the other junctions in a compact interior interval. Two
neutral sides then give a positive lower bound proportional to
$r_i^{4-p_i}\dd r_i$. For a single neutral interval the same conclusion
uses its diagonal. A photon vertex with zero angular shift between
neutral intervals is one place where this threshold must be checked.

A quantitative upper bound also covers simultaneous collapses in
such a word. For $0<t\le1$ and $0\le a<5/2$ the neutral evolution
satisfies $\|r^{-a}e^{-tH_0}\|_{2\to2}\le C_a t^{-a/2}$.
In the five-dimensional representation the free Gaussian
$L^2$-to-$L^\infty$ bound is $Ct^{-5/4}$, and killing at the outer
boundary preserves it. On $r<\sqrt t$ the squared radial norm
is bounded by $Ct^{-5/2}\int_0^{\sqrt t}r^{4-2a}\,\dd r$
times the input norm squared. On $r\ge\sqrt t$ use contraction.
The nonzero ground-state coefficient at the tip excludes $a\ge5/2$
even at fixed $t$.

Split the exponent at each junction between its two incident
intervals. Assign any power to a charged side if one is present;
at a neutral--neutral junction split $p_i<5$ equally. Let $a_i,b_i$
denote the powers assigned to the two ends of interval $i$, and set
$\gamma_i=(a_i+b_i)/8$ for charged intervals and
$\gamma_i=(a_i+b_i)/2$ for neutral ones. Splitting each semigroup
into intervals of lengths $u_i/4,u_i/2,u_i/4$ controls its weighted
ends by the operator norms above, using adjoints at the right end.
The middle factor is a contraction. On a longest interval use its
trace norm, at most $C u_i^{-1/2}$ by comparison with the free
one-dimensional heat kernel. The trace inequality bounds the word by
$C U^{-1/2}\prod_i u_i^{-\gamma_i}$ near zero. This argument can be
applied first to bounded radial weights and passed to the positive
kernel integral by monotone convergence. For a density bounded by
the product above, $\beta_i>\gamma_i$ for every $i$ and
$\sum_i(\beta_i-\gamma_i)>1/2$ suffice. This last estimate is an
upper bound, so equality alone does not decide a mixed word.

These results distinguish three questions about a physical contact:
which angular shifts its source derivative produces, which inverse
powers of radius accompany them, and which durations share a
boundary weight after sewing. The groupoid composition and velocity
shell specify how a word can close. The represented matrix elements
and boundary measure decide how much it contributes. An angular
trace bound can be carried over to the electromagnetic graph once
the same source also controls its remaining momentum factors and
its direct and intercycle photon contractions.

The neutral state on the same surface shows which part of the
estimate depends on angular charge. For $j=0$ the Hamiltonian is
$-\partial_r^2+2/r^2$. Its regular ground
function is proportional to $\sin(kr)/(kr)-\cos(kr)$, with $k$ the
first positive solution of $\tan k=k$. Its density is of order $r^4$.
It passes the Dirac inverse-square test. Its bridge has finite
normalization, but the $n=5$ photon cost in \eqref{eq:g5N-moment}
diverges logarithmically. This is realized by the connected
one-visit current graph of Corollary~\ref{cor:uvpeW-sharp-depth}.
Even at $j\ne0$, finite form energy gives an eighth inverse moment
but need not give the ninth: choose $u$ proportional to $r^4$
near zero and cut it off smoothly before $r=1$.
The stationary or heat-state condition in the proposition supplies
the higher-order control.

\paragraph{Two opposite shifts give a sharp mixed-duration test}
One closed word illustrates how a charged and a neutral interval
can collapse at very different rates. Use opposite shifts $+m,-m$
with fixed integer $m\ne0$, and the positive radial insertions
$r^{-p},s^{-q}$ at their two junctions. Write $P=p+q$ with
$p,q\ge0$. For an exceptional starting mode let $u$ be the charged
duration, $v$ the neutral duration and $U=u+v$. Its trace is
\[
 W_{m0}(u,v)=\int_0^1\!\int_0^1
 r^{-p}s^{-q} K_u^m(r,s)K_v^0(s,r)\,\dd s\,\dd r .
\]
Uniformly in the ratio $u/v$ for sufficiently small $U$,
\[
 W_{m0}(u,v)\asymp
 U^{-1/2}\int_{u^{1/8}}^{1/4}r^{-P}
                         \min\{1,r^4/U^2\}\,\dd r .
\]
The lower endpoint summarizes the charged exponential in a two-sided
estimate. The exact trace integrates all the way to the tip. The
neutral kernel resolves the tip on the scale $\sqrt U$. The two
length scales exchange order at $u=U^4$.

For the upper bound, weighted arithmetic--geometric mean and the
symmetry in $r,s$ reduce the two weights to a single $r^{-P}$.
When $u\ge U/2$, split the charged interval at most halfway as in
the survival proof, and bound the remaining return density by
$CU^{-1/2}$. This gives the exponential
$e^{-cu/r^8}+C e^{-c/r^3}$ at the weighted endpoint, with constants
depending on $m$. For a lower bound, a free bridge starting at
$r\in[u^{1/8},1/4]$ can stay within $(r/2,3r/2)$ throughout
with probability uniformly bounded below and bounded potential
cost. Here $u^{1/8}$ is larger than $\sqrt U$ for small $U$.

When $v\ge U/2$, use the neutral kernel estimate
\[
 K_v^0(r,s)\le C v^{-1/2}e^{-(r-s)^2/(4v)}
                      \min\{1,(rs/v)^2\}.
\]
For $r,s\le1/2$ the reverse inequality up to a constant also holds
at small $v$. Indeed the free five-dimensional Gaussian after radial
conjugation has angular integral proportional to
$\int_{-1}^1e^{-x(1-z)}(1-z^2)\,\dd z\asymp\min(1,x^{-2})$
with $x=rs/(2v)$. Killing at the outer sphere gives an upper
bound by the free kernel. Between two points in the half-radius ball
the straight chord stays in that ball; the bridge fluctuation stays
within the remaining distance with uniformly positive probability.
This proves the lower comparison there.

Put $S=(e^{-uH_m}1)(r)$. The free Brownian fourth moment and
Cauchy--Schwarz give
$\int s^2K_u^m(r,s)\,\dd s\le C S^{1/2}(r^2+u)$.
Use the neutral kernel estimate both with and without the factor
$(rs/v)^2$. The integrated return density at $r$ is bounded by
\[
 C U^{-1/2}\min(1,r^4/U^2)
             (e^{-c'u/r^8}+e^{-c'/r^3}).
\]
For $r\ge\sqrt u$ one uses $r^2+u\le2r^2$; for $r<\sqrt u$,
$u/r^8\ge r^{-6}$ and each exponential absorbs $1+u/r^2$
on reducing its constant. The integral below $R=u^{1/8}$ of the
first exponential term is bounded by the integral over $[R,2R]$
without that exponential, by rescaling $r=Rx$. The second
exponential has a bounded radial integral. For the lower bound on
$R\le r\le1/4$, restrict the charged Brownian path to a
$\sqrt u$-neighborhood of $r$. This event has uniformly positive
probability, the charged potential cost is bounded, and its endpoint
$s$ satisfies $s\asymp r$, $|s-r|\le\sqrt v$. The neutral lower
bound now retains both original insertion weights.

For a joint duration density comparable to
$u^{\alpha-1}v^{\beta-1}\dd u\,\dd v$ on a small square, take
$\alpha,\beta>0$. At fixed $U,r$, the integral over
$0<u<\min(U,r^8)$, including the $v=U-u$ weight, is comparable
to $U^{\alpha+\beta-1}\min(1,(r^8/U)^\alpha)$ by the beta
integral. Split the remaining radial integral at $\sqrt U$ and
$U^{1/8}$. Below $\sqrt U$ its radial power is $4+8\alpha-P$,
which must exceed $-1$. The contributions from an interior
annulus, from $r\asymp U^{1/8}$ and from $r\asymp\sqrt U$
respectively require
$\alpha+\beta>1/2$, $\alpha+\beta>(P+3)/8$ and
$4\alpha+\beta>P/2$. These strict inequalities, together with
$P<5+8\alpha$, also suffice: on each intervening radial interval
the integral is a single power, or a logarithm when that power is
$-1$. The strict duration margins absorb the latter logarithm.

The other exceptional starting mode interchanges the durations.
Adding both to the all-nonzero angular sum gives the exact test for
the \emph{complete positive scalar two-shift trace}:
\[
 \begin{gathered}
 \alpha+\beta>\max\{1,(P+3)/8\},\qquad
 P<5+8\min(\alpha,\beta),\\
 \min(4\alpha+\beta,\alpha+4\beta)>P/2 .
 \end{gathered}
\]
For $P=10$ and $\alpha=\beta=1$, every fixed-duration junction
and the all-nonzero sum pass. An exceptional sector nevertheless
has a logarithmic divergence where $u\asymp v^4$. A test limited
to $u\asymp v$ would miss it. The current, its powers, and the
duration law of an unrestricted graph must still be calculated
from the physical source before applying this criterion.

\begin{theorem}[All finite scalar words with ordinary duration measure]
\label{thm:uv-normal-word}
On the quartic chart of Proposition~\ref{prop:g5N-fibre}, consider a
cyclic word of $n$ insertions $r_i^{-p_i}e^{\ii\ell_i\theta}$
as above, with $p_i\ge0$ and $\sum_i\ell_i=0$. Integrate its
positive radial trace over $0<u_i<\delta$ with independent measures
$\dd u_i$, summing all starting angular momenta. For sufficiently
small $\delta>0$ the integral is finite exactly under the following
two conditions.
\begin{enumerate}
\item The sectors with no zero interval satisfy
$n>1$ and $\sum_i p_i<8n-3$.
\item For each of the finitely many remaining starting modes, draw
the cycle with one vertex per insertion and one edge per interval.
Join into one block all vertices connected by nonzero-mode edges.
Keep every zero-mode edge, including loops and repeated edges.
For a block $B$ put
\[
 d_B=1+\tfrac13\sum_{i\in B}(p_i-8).
\]
For every nonempty collection $S$ of blocks, let $e_0(S)$ count
the zero-mode edges with both ends in $S$, each counted once.
Then
\[
                 \sum_{B\in S}d_B+e_0(S)<0 .
\]
\end{enumerate}
The second test allows a group of insertions to approach the tip
while the others remain at ordinary radii.
\end{theorem}

\begin{proof}
The first test follows by integrating $F_P(U)$ above on the duration
simplex. For the remaining modes retain the actual order of the
intervals. Write $z=(r^{-3}-1)/3$, $w=(s^{-3}-1)/3$. The exact
zero-energy Green kernel on $L^2(\dd r)$ is
\[
 G_j(r,s)=r^2s^2 g_{|j|}(z,w),\qquad g_0(z,w)=\min(z,w),
\]
where, for $b>0$,
\[
 g_b(z,w)=\frac{e^{-b|z-w|}}{2b}
                     (1-e^{-2b\min(z,w)}).
\]
Indeed $u(r)=r^2 v(z)$ transforms the quadratic form into
$\int_0^\infty(|v'|^2+j^2|v|^2)\dd z$. These half-line kernels
vanish at $z=0$, have unit negative derivative jump at $z=w$,
and satisfy the finite-energy tip condition. The factors $r^2s^2$
convert the Green kernel back to the stated radial measure.
The spectrum on the original finite radial interval has a positive
lower bound. Positivity and Tonelli's theorem therefore allow each
heat kernel integrated over all positive durations to be replaced by
$G_j$. This bounds the integral over the small duration cube from above.

At vertex $i$ the two adjacent Green kernels and the radial measure
give $r_i^{4-p_i}|\dd r_i|=(1+3z_i)^{(p_i-8)/3}\dd z_i$.
A nonzero edge confines the difference $z_i-z_{i+1}$ exponentially.
Its common value can still grow without limit. Select a representative
$y_B$ in each connected block. Along a spanning tree in the block,
the ratios $(1+z_i)/(1+y_B)$ and their inverses are bounded by a
product of $1+|z_k-z_l|$ over tree edges. The same bounds replace
$\min(1+z_i,1+z_j)$ by $\min(1+y_B,1+y_C)$ at the ends of a
zero edge. Integration of the resulting polynomial in the tree
differences against their exponentials gives a finite constant.
The upper bound thus reduces to one integration in $y_B$ per block.
Its weights are $(1+y_B)^{d_B-1}$ and one
$\min(1+y_B,1+y_C)$ per zero edge.

Partition the $1+y_B$ into dyadic intervals with indices $k_B\ge0$.
For $h\ge1$ set $S_h=\{B:k_B\ge h\}$. Including the volumes,
the weight of a dyadic box is bounded by a constant times
\[
 2^{\sum_{h\ge1}\left(\sum_{B\in S_h}d_B+e_0(S_h)\right)} .
\]
If all nonempty subset sums are negative, each is at most
$-\eta$ for some $\eta>0$ because the word is finite. The number
of boxes with $\max k_B=k$ grows at most polynomially in $k$.
The sum converges.

For necessity we must also see the same costs at short durations.
For a charged edge take $|z-w|\le C$ with $r,s$ sufficiently
small. Then $|s-r|=O(r^4)$. Integrating its heat kernel over
$[r^8,2r^8]\subset(0,\delta)$ gives a lower bound $c r^4$:
the free radial bridge density is at least $c r^{-4}$, and the
bridge can stay in $(r/2,2r)$ with uniformly positive probability
and bounded integrated potential. For a neutral edge let
$R=\max(r,s)$. Integrate the five-dimensional radial kernel
estimate above over $[R^2,2R^2]\subset(0,\delta)$.
It gives at least $c r^2s^2/R^3$, comparable to
$r^2s^2\min(z,w)$ near the tip. These two estimates are uniform
on the indicated sets for the finitely many exceptional modes.

Keep every member of a block within a fixed bounded difference in
$z$ of its representative. Give the other members intervals of fixed
positive width there. If $S$ fails its test, let its representatives
range from $2^k$ to $2^{k+1}$ and keep all others in a fixed,
sufficiently deep annulus. The preceding short-duration bounds give
at least a constant times $2^{k(\sum_{B\in S}d_B+e_0(S))}$ on
each such disjoint box. Their sum diverges, including at equality.
\end{proof}

For two opposite shifts, an exceptional mode has one block with two
vertices and one neutral loop. Its exponent is $(p+q-10)/3$,
recovering $p+q<10$. With four alternating shifts, an exceptional
mode has two pairs of vertices connected by charged edges.
Two neutral edges join the pairs. If their respective sums of powers
are $P_1,P_2$, one needs $P_1<13$, $P_2<13$ and $P_1+P_2<20$
for this mode. Holding the second pair in an interior annulus shows
why $P_1=13$ already gives a logarithm even at $P_2=0$. The other
exceptional starting mode pairs the vertices in the other order
and must be tested in its own right.

\paragraph{Curvature readouts at the critical radial power}
The Gaussian curvature on this chart is $-12/r^2$. As an example
of a geometric observable, multiply it by a real trigonometric
polynomial in $\theta$ of zero angular mean. Every finite scalar
loop of these marked curvature insertions has a finite
short-duration coefficient with ordinary duration measure. A
one-insertion trace vanishes by angular orthogonality at every
positive duration. For $n\ge2$, expand the finite set of Fourier
shifts. Only closed words contribute, each with $p_i=2$. The
nonzero-sector test passes. In an exceptional mode there must be
a charged interval, since the shifts are nonzero. If $S$ comprises
$b$ blocks with $a$ original vertices, its exponent is
$b-2a+e_0(S)$. After contraction the neutral edges form a cycle
(with a single loop allowed), so $e_0(S)\le b$, strictly for
a nonempty proper $S$. For the whole cycle at least one charged
edge was contracted, hence $a>b$. All tests are strict.

A constant angular readout has a different answer at this same
radial power. The all-neutral mode has $a=b=n$ and $e_0(S)=n$
for the full set. Its contribution already diverges at equality.
Thus a change in what the source measures can change the UV result
while the kinetic action and the metric stay fixed. The
curvature readout is a specified scalar probe. Applying this result
to photons requires deriving their angular matrix elements, radial
factors and duration measure from the electromagnetic source.

\begin{corollary}[Curvature sampled by a general angular profile]
\label{cor:rough-curvature-uv}
At vertex $i$ insert $r^{-2}b_i(r,\theta)$, where at almost every
radius $b_i$ has zero angular mean and
$M_i=\mathop{\rm ess\,sup}_r\|b_i(r,\cdot)\|_{L^2(\dd\theta/2\pi)}
<\infty$. For $n\ge2$ the sum of the absolute angular contributions
to the scalar cyclic word, integrated over all positive independent
durations with measure $\prod_i\dd u_i$, is bounded by
\[
 \left(\tfrac23\sqrt{1+\pi^2/12}\right)^n\prod_i M_i .
\]
The one-insertion diagonal vanishes at positive duration. A fixed
multiple of Gaussian curvature is included by absorbing $-12$ into
$b_i$. The formula controls the coefficients individually for the
stated readout.
\end{corollary}
\begin{proof}
Use the radial Green kernels of Theorem~\ref{thm:uv-normal-word}.
Each vertex now has weight $(1+3z_i)^{-2}\dd z_i$ before its angular
matrix element. Put $h(0)=1$, $h(k)=(2|k|)^{-1}$ otherwise.
Each nonzero-mode Green factor is at most $h(k)$, after its radial
endpoint factors are extracted. A neutral factor is $\min(z,w)$,
bounded by $\sqrt{(1+3z)(1+3w)}$. Two neutral intervals cannot
meet in a contributing term because the intervening Fourier
coefficient $\widehat b_i(r,0)$ vanishes. Thus each vertex pays at
most one extra power $(1+3z_i)^{1/2}$.

At fixed radii distribute the $h$ factors equally between their two
ends. The nonnegative matrix for each angular sum is then bounded
using the Hilbert--Schmidt norm of
\[
 A_i(j,k)=\sqrt{h(j)h(k)}\,
                 |\widehat b_i(r_i,j-k)|.
\]
Indeed $\sum_k h(k+\ell)h(k)\le\|h\|_2^2=1+\pi^2/12$.
Parseval gives $\|A_i\|_{\rm HS}\le\|h\|_2 M_i$. The trace of a
product of $n\ge2$ finite matrix compressions is bounded by the
product of their Hilbert--Schmidt norms. All entries are
nonnegative, so the absolute routing sums increase to the full
answer. Finally $\int_0^\infty(1+3z)^{-3/2}\dd z=2/3$ at each
vertex. Tonelli's theorem justifies the positive majorant including
all durations. For $L^2$ profiles the word is defined by these
Fourier matrix elements; the estimate gives their absolute
integrability even when pointwise multiplication on the full
Hilbert space would be unbounded.
\end{proof}

The radial power can be increased if only convergence at each fixed
order is needed. For a common power $0\le p<5$ the same
zero-mean class still gives finite loops of every order $n\ge2$.
To see this, group the angular routings according to which edges
have mode zero. Keep the factor $e^{-|z-w|}$ from every charged
edge in addition to its $h(k)$. The angular sum in each pattern
obeys the same matrix bound. In a mixed pattern every contracted
charged block has at least two vertices. Its exponent in
Theorem~\ref{thm:uv-normal-word} is
$1+|B|(p-8)/3<-1$, and $e_0(S)\le |S|$. Thus its radial
integral converges. In an entirely charged pattern the tree
differences integrate exponentially, leaving a single common depth
with weight bounded by $(1+z)^{n(p-8)/3}$, which is integrable.
The all-neutral pattern vanishes. The two-insertion word with
$b_1=b_2=\cos\theta$ attains the logarithmic boundary at $p=5$,
since its two opposite shifts give $P=10$. This locates a sharp
threshold for uniform claims about this class of scalar probes.

\paragraph{A bare scalar determinant for the curvature probe}
The coefficient estimate also determines an actual relative
determinant. Let $H=\bigoplus_j H_j$ and let $V=r^{-2}b(r,\theta)$
with $b$ real, bounded and of zero angular mean at almost every
radius. Use its multiplication form in the same scalar space.
Write $M=\mathop{\rm ess\,sup}_r\|b(r,\cdot)\|_2$ and
$C=\tfrac23\sqrt{1+\pi^2/12}$. The source amplitude $\lambda$
specifies a measurement. For real $\lambda$ near zero with
$|\lambda|\|b\|_\infty<2$ and $|\lambda|CM<1$, the bare integral
\[
 \int_0^\infty
 \frac{\Tr e^{-tH}-\Tr e^{-t(H+\lambda V)}}{t}\,\dd t
\]
converges absolutely and is real analytic in $\lambda$. Each of its
order-$n$ coefficients for $n\ge2$ has absolute value at most
$(CM)^n/n$. The first coefficient is identically zero. In this
statement the reference is the same unperturbed scalar source. No
source-dependent subtraction is required.

Here are details of the relation between the loop coefficients and
the determinant. Approximate the radial factor in $V$ by
$\min(r^{-2},N)$. For this bounded insertion the ordered heat
expansion is valid at fixed $t$. In its trace the unmarked lengths
before the first and after the last insertion join into a single
gap. Integrating their possible split supplies that gap's length.
Cyclic symmetry makes its average over the $n$ identical marks
$t/n$. The measure $\dd t/t$ now leaves $1/n$ times the product
of the $n$ independent duration measures. The one-mark trace
vanishes at every $t$ by the zero Fourier coefficient. In every
higher term the corollary bounds the integral of the absolute
angular and radial contributions uniformly in $N$. For each $n$
the integrated terms converge in $L^1(\dd t/t)$ by domination with
that same positive, integrable majorant. The geometric series
bound for $|\lambda|CM<1$ permits passage to the limit in the sum.

To identify this limit with the stated heat trace, use the
Dirichlet radial Brownian bridge and, conditional on it, circle
diffusion of lapse $\int R_s^{-8}\dd s$. A surviving bridge over
a fixed positive duration has strictly positive minimum radius.
The approximating potential converges there to the true one.
The sum of the conjugation potential $2/r^2$ and the approximating
source potential stays nonnegative under the stated smallness
bound. The free radial diagonal bound times
$\sum_j e^{-t j^2}$ dominates the trace at each fixed $t$, since
the circle lapse is at least $t$. Bounded convergence gives the
Friedrichs heat trace of the stated potential. Equivalently one may
construct its form by bounded potential approximation on the
interior. The radius bound $N$ used in the proof is absent from
both the integral and its coefficients. This determinant belongs
to the scalar curvature probe; the physical charged determinant
also differentiates its connection, constraints and sewn measure.

\paragraph{A boundary response fixed by the same normal kinetic form}
The role of the zero harmonic can also be decided without a heat
estimate. In the scalar circle representation, allow a flat normal
connection with dimensionless holonomy parameter $a$, so the angular
derivative is $\partial_\theta-i a$. Hold the outer Dirichlet data
$g(\theta)=\sum_j g_j e^{ij\theta}$ fixed and minimize the static
normal kinetic action over finite-energy continuations. This
calculation varies a holonomy probe in the representation that
admits it; the basic horizontal electromagnetic source keeps the
normal connection fixed.

For orbit radius $f(r)=r^4$ set
$z=\int_r^1\dd s/f(s)=(r^{-3}-1)/3$. The tip is at $z=\infty$.
The kinetic energy of one Fourier amplitude $v(z)$ becomes exactly
\[
 \int_0^\infty \bigl(|v'(z)|^2+(j-a)^2|v(z)|^2\bigr)\,\dd z,
 \qquad v(0)=g_j.
\]
For $j\ne a$ let $b=|j-a|$. Completing the square gives
$b|g_j|^2+\int_0^\infty|v'+bv|^2\dd z$; the boundary term at
infinity vanishes for $v\in H^1(0,\infty)$. The minimizer is
$g_j e^{-bz}$. For $j=a$ the constant continuation has zero
gradient and finite norm in the original metric volume. It lies in
the finite-energy tip closure: interpolating it to zero over
$R<z<2R$ costs $O(R^{-1})$ in energy and changes the metric norm
by a quantity tending to zero. With angular Haar measure of total
mass one, the exact boundary energy is consequently
\[
 E_a(g)=\sum_{j\in\mathbb Z}|j-a|\,|g_j|^2,
 \qquad g\in H^{1/2}(S^1).
\]
Near an integer $a=k$, every other term is affine in $a$. The two
one-sided first derivatives differ by $2|g_k|^2$. In particular,
for fixed outer data the classical second derivative at that
holonomy exists precisely when $g_k=0$. The result tests a
boundary domain of the action with all modes included. It neither
asserts a divergence of a four-dimensional photon graph nor
identifies which normal boundary variation a measured photon
actually generates. That identification belongs to the source
map in Postulate~$\Omega_\alpha$.

\paragraph{Does the holonomy contact survive a quantum trace?}
There is a direct quantum test of that same flat normal connection.
At heat duration $t>0$ let
$T_t(a)=\sum_{j\in\mathbb Z}\Tr e^{-tH_{j-a}}$, with
$H_{j-a}=-\partial_r^2+2/r^2+(j-a)^2/r^8$ in its Friedrichs
realization. The trace is finite for every real $a$:
kinetic comparison bounds each summand by
$C_t e^{-t(j-a)^2}$. Nevertheless at the trivial holonomy
\[
 \lim_{a\to0}\frac{T_t(a)+T_t(-a)-2T_t(0)}{a^2}=-\infty .
\]
To see which contact produces the singularity, use the positive
neutral radial bridge integral at time $t$, including the diagonal
endpoint integration and the $2/r^2$ potential. Denote its finite
measure by $\mu_t$ and let $I=\int_0^t r(s)^{-8}\dd s$. The
$j=0$ summand is $\int e^{-a^2I}\dd\mu_t$. As $a^2$ decreases
to zero, $(1-e^{-a^2 I})/a^2$ increases to $I$. Positivity and
cyclic composition of the heat kernels give
\[
 \int I\dd\mu_t=t\int_0^1 r^{-8}K_t^0(r,r)\dd r=\infty .
\]
One may first use bounded approximations to $r^{-8}$ in this
identity and take their monotone limit. It is the derived
two-source contact: the linear current of the zero mode vanishes
at $a=0$, while its second variation of the angular kinetic
potential is $2r^{-8}$. The neutral diagonal costs $r^4$ at the tip.

No omitted angular sector can cancel this singularity. For
$|a|\le1/2$ and $j\ne0$, twice differentiating
$e^{-(j-a)^2I}$ costs at most
$C|j-a|^{-2}e^{-t(j-a)^2/2}$, since $I\ge t$.
The corresponding first derivatives are likewise summably
bounded. Thus their summed trace is twice continuously
differentiable near zero. This calculation uses the quantum
measure of the normal operator and its complete angular trace at
a fixed positive duration. Its source varies the flat angular
connection. A conclusion about the physical photon would require
the electromagnetic source derivative on this connection as well
as its other contributions to the sewn graph.

\paragraph{Which flat changes does a closed normal loop see?}
On this scalar chart a flat Abelian connection of the form
$\dd\chi+a\,\dd\theta$, with smooth single-valued real $\chi$ in
the punctured chart, is unitarily equivalent to
$a\,\dd\theta$. Multiplication by $e^{\ii\chi}$ intertwines the
covariant kinetic forms on compactly supported data and hence their
Friedrichs closures. A closed trace depends only on $a$ modulo
integers. For prescribed nonzero boundary data the same comparison
also transports the boundary state. In particular, a flat variation
with zero period cancels from the complete closed trace. Separate
formal current and contact integrals need not be absolutely
convergent before such a gauge identity is used.

\begin{proposition}[Relative heat trace of a flat normal connection]
\label{prop:quartic-holonomy-uv}
In the preceding periodic scalar realization let
$D_t(a)=T_t(0)-T_t(a)$. For each fixed $a\notin\mathbb Z$,
\[
 \lim_{t\downarrow0}t^{3/8}D_t(a)=
 \frac{\Gamma(5/8)}{8\pi^{5/4}}
 \sum_{w=1}^\infty\frac{1-\cos(2\pi wa)}{w^{5/4}}>0 .
\]
Consequently the bare relative proper-time integral for the
logarithm of the determinant, $\int_0^{t_0}D_t(a)\dd t/t$,
diverges for every $t_0>0$. At integer $a$ the two traces are equal.
\end{proposition}
\begin{proof}
For a surviving radial Brownian bridge $R_s$ generated by
$\partial_r^2$ put $I=\int_0^t R_s^{-8}\dd s$ as above.
Set $\Theta_v(a)=\sum_j e^{-(j-a)^2v}$ and
$L_a(v)=\Theta_v(0)-\Theta_v(a)$. Gaussian Poisson summation gives
\[
 L_a(v)=2\sqrt{\pi/v}\sum_{w\ge1}(1-\cos(2\pi wa))
                                     e^{-\pi^2w^2/v}.
\]
It is nonnegative, bounded for $v>0$, and at most $C e^{-c/v}$
for $0<v\le1$. The bounds at large $v$ follow from the original
Gaussian sum. The bridge integrand in $D_t(a)$ is exactly
$L_a(I)e^{-\int_0^t2/R_s^2\dd s}$ times the survival indicator,
with free diagonal density $(4\pi t)^{-1/2}\dd r$.

Set the starting radius to $r=t^{1/8}x$. At fixed $x>0$ the
bridge fluctuations have size $\sqrt t$, so in the rescaled
radius they tend uniformly to zero almost surely under Brownian
bridge scaling. The survival indicator tends to one,
$I\to x^{-8}$ and $\int_0^t2/R_s^2\dd s\to0$.
For large $x$ split according to whether the bridge ever falls
below $r/2$. The probability of that event is bounded by
$C e^{-c r^2/t}\le C e^{-c x^2}$ for $t\le1$. On its complement
$I\le 2^8 x^{-8}$, giving the integrable bound $Ce^{-c' x^8}$
for large $x$. Paths killed at the outer boundary contribute zero.
Together with boundedness on bounded $x$, these estimates permit
dominated convergence after extending the integrand by zero outside
the radial interval. The limit is
$(4\pi)^{-1/2}\int_0^\infty L_a(x^{-8})\dd x$. Integrating
the positive Poisson series term by term gives the displayed
coefficient. For nonintegral $a$ its first term is strictly positive.
\end{proof}

The paths that give the excess trace are concentrated on the scale
$r\asymp t^{1/8}$, where the angular orbit radius $r^4$ is of order
$\sqrt t$. The integer $w$ is dual to the normal circle momentum.
It labels normal-circle winding in this calculation. Identifying
it with geometric-time winding would require an additional map
between the two variables. The logarithm of the determinant includes
the short-duration weight $\dd t/t$; suppressing a momentum at a
fixed positive duration does not control this integral.

\paragraph{An exact comparison for the inverse operator}
The same geometry gives a finite relative inverse trace. Define
$F(b)=\Tr H_b^{-1}$ for one radial operator with real $b\ge0$.
The Green diagonal in the variable $z$ gives, by positive integration,
\[
 F(b)=\tfrac15\int_0^\infty e^{-2bv}(1+3v)^{-5/3}\dd v .
\]
Indeed $g_b(z,z)=\int_0^z e^{-2bv}\dd v$ and
$r^4|\dd r|=(1+3z)^{-8/3}\dd z$. Hence $F(0)=1/10$ and
$F(b)\sim1/(10b)$ as $b\to\infty$. The absolute inverse trace of
the entire angular family diverges. On the common angular Hilbert
space its difference at two holonomies is trace class near $a=0$.
Each radial difference is ordered by the value of $(j-a)^2$, so
its trace norm equals $|F(|j|)-F(|j-a|)|$. For $b>0$ the integral
implies $|F'(b)|\le1/(10b^2)$ and $|F''(b)|\le1/(5b^3)$.
The first bound sums the nonzero-sector trace norms, and the second
bounds their paired $j,-j$ trace contribution by $O(a^2)$.
Changing variables $x=2|a|v$ in $F(0)-F(|a|)$ now yields
\[
 \begin{aligned}
 \Tr\!\left(\bigoplus_j H_j^{-1}-
              \bigoplus_j H_{j-a}^{-1}\right)
 &\sim \tfrac{\Gamma(1/3)}{10}
             (2|a|/3)^{2/3}\qquad (a\to0).
 \end{aligned}
\]
The integrable comparison function is
$(1-e^{-x})x^{-5/3}$, of integral $3\Gamma(1/3)/2$.
This exponent describes the inverse integrated over all heat
durations. Proposition~\ref{prop:quartic-holonomy-uv} concerns the
short-duration logarithm of the determinant at fixed holonomy.
The two limits answer different questions about the same operator.

A horizontal pullback that keeps the normal holonomy fixed does
not turn on this relative determinant. For the unrestricted
electromagnetic source the relevant inputs are its actual normal
holonomy variation, the other fields in its determinant and the
boundary state used for gluing. The scalar calculation supplies a
test of those choices without inserting a damping function or a
counterterm.

\paragraph{What a horizontal magnetic field measures}
Does a normal return also suppress the momentum of a charged field
propagating on the four-dimensional carrier? Holding the normal
holonomy fixed makes this question accessible to a direct magnetic
test. Use one charged complex scalar on the Euclidean continuation
of a locked plane. On its product with the preceding normal surface
take the local kinetic form with generator
\[
 {\cal H}_b=H+\kappa(r)L_b,\qquad
 L_b=-(\partial_{x^\mu}+\ii e A_\mu)^2 .
\]
Here $H$ is the full normal operator of
Corollary~\ref{cor:rough-curvature-uv}, $F_{12}=B$ is constant,
the other components vanish, and $b=eB$. The horizontal connection
is pulled back without a normal component. For definiteness take
$0<\kappa_-\le\kappa(r)\le\kappa_+<\infty$ in the chosen
reference units. This allows the local horizontal diffusion rate to
vary along the normal path. The ordinary product metric gives the
special case $\kappa=1$. All ratios below are relative to zero
field and all horizontal traces are per unit four-volume.

\begin{proposition}[The quadratic magnetic coefficient before normal sewing]
\label{prop:horizontal-magnetic-word-uv}
Integrate out the horizontal coordinates in the heat difference
$e^{-t{\cal H}_0}-e^{-t{\cal H}_b}$ and keep both normal endpoints.
For every $t>0$ the coefficient of $b^2$ in the resulting normal
operator equals
\[
                         \frac{e^{-tH}}{6(4\pi)^2}.
\]
Consequently every nonzero, positive trace-class normal preparation
held fixed in this magnetic variation gives a logarithmically
divergent quadratic coefficient with bare measure $\dd t/t$ at zero.
With the full unweighted normal trace the heat difference has leading
term $b^2/[60(4\pi)^2t]$ at fixed $b\ne0$. Every higher fixed
magnetic coefficient of this scalar trace is integrable at zero.
\end{proposition}
\begin{proof}
Condition on a normal heat path with radial component $R_s$.
Its integrated horizontal diffusion time is
$S=\int_0^t\kappa(R_s)\dd s$. The horizontal evolution on that
path is $e^{-SL_b}$: every increment multiplies the same operator
$L_b$ by a scalar. The scalar magnetic heat diagonal in four
dimensions is $(4\pi S)^{-2}bS/\sinh(bS)$. For instance, the two
magnetic directions have levels $(2k+1)|b|$ of degeneracy
$|b|/(2\pi)$ per area, and the other two supply $(4\pi S)^{-1}$.
This also gives the local magnetic coefficient of the scalar
covariant Laplacian~\cite{WardHeatManual}. The loss on the path is
\[
 Q_b(S)=\frac{1-bS/\sinh(bS)}{(4\pi S)^2}
       =\frac{b^2}{6(4\pi)^2}+O(b^4 S^2).
\]
The quadratic term is exactly independent of $S$ before either
endpoint is integrated. The bound $S\le\kappa_+t$ makes the
remainder uniform at short duration. In absolute kernel value it is
at most $C b^4t^2$ times the positive kernel of $e^{-tH}$.
This also bounds its operator norm by $C b^4t^2$, by positivity
and contraction. Thus an off-diagonal trace-class preparation
$\rho\ge0$ gives a leading heat loss
$b^2\Tr\rho/[6(4\pi)^2]$, since
$\Tr(\rho e^{-tH})\to\Tr\rho$.

For the unweighted trace the exact short-time normal multiplicity is
$\lim_{t\downarrow0}t T_t(0)=1/10$. Indeed at each interior starting
$r$ a short radial bridge has $I/t\to r^{-8}$, with $I$ as in
Proposition~\ref{prop:quartic-holonomy-uv}. The angular factor
$\sqrt t\sum_j e^{-j^2 I}$ tends to $\sqrt\pi r^4$ and is bounded
uniformly for $t\le1$ by $I\ge t$ and Gaussian comparison.
The killing and the $2/r^2$ weight tend to one at each such $r$.
The remaining free radial factor gives
$\tfrac12\int_0^1r^4\dd r=1/10$. Combining this limit with the
uniform remainder gives the asserted asymptote. For order $b^{2n}$
with $n\ge2$ the path coefficient is bounded by $C_n t^{2n-2}$.
Multiplying by $T_t(0)=O(t^{-1})$ and by $\dd t/t$ proves
short-duration integrability.
\end{proof}

\paragraph{Can a nonzero winding carry the magnetic coefficient?}
On a cylinder of fixed circumference, a loop of nonzero winding
must travel a definite distance. On this surface the circumference
is $2\pi r^4$. To test the effect of its collapse, lift normal
paths to the angular cover and keep winding $w\ne0$. At the
trace basepoint insert the radial observable $r^p$, $p\ge0$.
This mark is in addition to the metric volume already present in
$H$. Call the resulting normal heat trace $W_{p,w}(t)$, with the
angular normalization of Proposition~\ref{prop:quartic-holonomy-uv}.
For each nonzero integer $w$ one has
\[
 \lim_{t\downarrow0}t^{(3-p)/8}W_{p,w}(t)
 =\frac{\Gamma((p+5)/8)}
        {16\pi^{(p+5)/4}|w|^{(p+5)/4}} .
\]
For the sum over all nonzero windings the limit is the sum of
these constants. In particular the magnetic proper-time integral
for that marked contribution is ultraviolet finite \emph{exactly}
for $p>3$. At $p=3$ it diverges logarithmically. With no mark
($p=0$), even winding one has too much short-duration weight.

To prove the formula, the contribution of one lift to the
conditional angular diagonal is
$\sqrt{\pi/I}\exp(-\pi^2w^2/I)$. Repeat the radial bridge scaling
$r=t^{1/8}x$. The limiting integrand for the rescaled trace is
$\tfrac12 x^{p+4}\exp(-\pi^2w^2x^8)$, whose integral is the
stated gamma factor. Domination holds even for the sum. The
conditional angular factor summed over $w\ne0$ is bounded by one
by comparison of the Gaussian series with its integral. At small
$I$ it is at most $C e^{-c/I}$. On paths staying above $r/2$,
$I\le 2^8 x^{-8}$. The probability of a larger inward excursion
is at most $Ce^{-c x^2}$ for $t\le1$, as in the proof for a flat
holonomy. Thus at large $x$ the rescaled integrand is bounded by
$C x^p(e^{-c x^8}+e^{-c x^2})$, and on bounded intervals by
$C x^p$. This justifies the limit for the complete positive sum.

The magnetic coefficient on each such path is still the constant
in Proposition~\ref{prop:horizontal-magnetic-word-uv}. Its
proper-time test is therefore precisely
$\int_0^{t_0}W_{p,w}(t)\dd t/t$, or its sum over windings.
The uniform remainder estimate gives the same threshold for the
unexpanded magnetic loss at any fixed $b\ne0$. All higher fixed
magnetic orders pass the short-time test here. For example,
$r^4=\|\partial_\theta\|$ is the length of the angular Killing
vector and yields a finite marked nonzero-winding contribution.
This is a specified geometric observable. It enters the physical
determinant only if the represented current or boundary pairing
actually inserts it. The winding tested here belongs to the
normal circle. Net winding in geometric time and the plane's
vectorial pole prescription have their own definitions.

These results locate two separate inputs to a claim about the
horizontal photon. The normal paths of the quartic kinetic action
can suppress an angular curvature insertion to every order.
Ordinary local minimal coupling on the carrier gives the quadratic
magnetic coefficient above, including its contact. For a different
coupling to evade that coefficient one must derive the changed
horizontal operator, the intervening displacement, or the required
cancellation with other fields from the same GCM source. The scalar
product realization alone settles none of those features of the
unrestricted charged theory.

The internal angle in this example labels an auxiliary charged
representation. It compactifies no GT direction. The proof holds
for a source which enforces the return, uses the indicated Cartan sector
and boundary preparation, and samples the same normal preparation at every charged visit.
The metric follows within the stated variational class. The absence
of an identity-history source and the physical assignment of this
normal charge and state remain to be derived. To apply it to induced O--QED one must
identify these data and check their behavior under electromagnetic
source variations in the GCM functional that generates the graphs.

\paragraph{The same charged suppression with a freely varying endpoint slope}
The power $r^4$ at the tip is also selected without fixing $f'(1)$.
Keep $f(0)=f'(0)=0$, $f(1)=1$, the unit proper length and the finite
curvature action \eqref{eq:g5N-normal-geometry}. Define $C>0$ by
\[
 1=\int_0^\infty (z^3/3+C)^{-1/2}\,\dd z .
\]
The integral is a positive finite constant times $C^{-1/6}$.
Define $w>0$ on $(0,1)$ by
$r=\int_{w(r)}^\infty(z^3/3+C)^{-1/2}\dd z$.
Then $w'^2=w^3/3+C$, $w'<0$, $2w''=w^2$,
$w(1)=0$ and $w'(1)=-\sqrt C$. At the tip
$w\sim12/r^2$ and $w'\sim-24/r^3$, by comparison in the integral.
For every admitted profile the square completion now gives
\[
 {\cal E}[f]=2\sqrt C+\int_0^1\frac{|f''-wf|^2}{f}\,\dd r .
\]
Indeed its integrated boundary term is $2[wf'-w'f]_0^1$ and the
interior residual $(2w''-w^2)f$ vanishes. The tip term vanishes by
the previous finite-action bounds; $w(1)=0$ removes the endpoint
slope. Put $g=-w'$. Then $g''=wg$ and the positive solution
\[
 f_*(r)=c\,g(r)\int_0^r g(s)^{-2}\,\dd s
\]
has $f_*''=wf_*$, with $c>0$ chosen by $f_*(1)=1$. It satisfies
$f_*\sim c_0 r^4$, $f_*'\sim4c_0r^3$ with $c_0>0$ by
$g\sim24/r^3$, $g'=-w^2/2$. In particular it has finite action.
The other independent solution $g$ diverges at the tip. This gives
a unique curvature minimizer for a free endpoint slope.

For this minimizer the flat radial Hamiltonian in charge $j$ is
$-\partial_r^2+(\sqrt{f_*})''/\sqrt{f_*}+j^2/f_*^2$.
The exact comparison functions
\[
 v_a(r)=\sqrt{f_*(r)}\,
   \exp\!\left[-a\int_r^1 f_*^{-1}(s)\,\dd s\right]
\]
satisfy $(H_j-E)v_a=[(j^2-a^2)/f_*^2-E]v_a$.
The integral in the exponent is asymptotic to $1/(3c_0r^3)$.
The same barriers give \eqref{eq:g5N-rate} with
$|j|$ replaced by $|j|/c_0$. The heat-state estimate also persists:
$(\sqrt{f_*})''/\sqrt{f_*}\sim2/r^2$ is bounded below on
$(0,1)$, and $j^2/f_*^2\ge c_1j^2/r^8$ near the tip.
In the preceding Brownian estimate a lower-bounded potential
only multiplies the upper bound by a constant at fixed heat time.
Within this positive rotational class, minimizing the curvature
action with the stated boundary data selects quartic collapse even
when the endpoint slope is free.

\paragraph{Integrating the normal geometry}
The bound can also be obtained before selecting a stationary metric.
Let $\pi$ be any nonzero finite positive measure on the admissible
profiles with $f(1)=1$, supported on finite actions
$A={\cal E}[f]$. In each metric use a normalized radial ground
state of a fixed $j\ne0$. Use a single joint normalization for
$e^{-A}\pi(\dd f)$ times the return measure \eqref{eq:g5N-bridge}
with that ground state's radial probability. Assume the indicated
state and kernel assignments are measurable. The resulting duration
law has every inverse moment, uniformly with respect to Fourier
cutoffs in the four-dimensional graphs. In particular the individual
Wilson graphs remain absolutely finite after integrating these
fluctuating normal geometries. The reference measure is held fixed.

To see why, the curvature estimate already proved gives
$f(r)\le A r^3$ and $|f'|\le A$ throughout the interval. In
the original radial inner product $f(r)\dd r$, let $\psi$ be a
normalized eigenfunction of energy $\lambda$ in charge $j$.
Write $M_t=\int r^{-t}|\psi|^2 f\,\dd r$ and $B=|j|/A>0$.
For the bounded weight $W=(r+\epsilon)^{-t/2}$ the eigenfunction
equation tested against $W^2\psi$ gives
\begin{align*}
 \int f|(W\psi)'|^2+\int \frac{j^2}{f^2}|W\psi|^2 f
 =\lambda\int|W\psi|^2 f+\int |W'|^2|\psi|^2 f .
\end{align*}
All integrals here use $\dd r$. Bounded Lipschitz multiplication
preserves the form domain for each $\epsilon>0$. Use
$j^2/f^2\ge B^2r^{-6}$ and
\[
 \tfrac14 t^2 r^{-2}\le \tfrac12 B^2r^{-6}+t^3/B ,
\]
the latter by maximizing a cubic in $r^{-2}$. Discard the
nonnegative derivative term and send $\epsilon\downarrow0$. Whenever
$M_t<\infty$, monotone convergence gives
\begin{equation}
 M_{t+6}\le 2B^{-2}(\lambda+t^3/B)M_t .
 \label{eq:g5N-boost}
\end{equation}
The charged radial form has compact resolvent for each such $f$.
Indeed bounded form norm controls the mass of $(0,\delta)$ by
$B^{-2}\delta^6$ times the form bound; on $[\delta,1]$ the
weight $f$ is bounded above and below and the one-dimensional
Sobolev embedding in $L^2$ is compact. In particular its ground
state exists.

For a uniform ground energy bound take
$h=[2(1+A)]^{-1}$. On $[1-h,1]$, $1/2\le f\le3/2$ by its
Lipschitz bound. A Dirichlet sine on this interval extended by zero
is an admissible trial function. Its Rayleigh quotient bounds
$\lambda_0\le 3\pi^2/h^2+4j^2$. Starting at $M_0=1$,
\eqref{eq:g5N-boost} thus bounds $M_{6k}$ by
$C_{k,j}(1+A)^{4k}$ at every fixed $k$, and $M_s\le M_{6k}$
for $0\le s\le6k$. Every such polynomial is integrable against
$e^{-A}\pi(\dd f)$ for a finite reference measure. This estimate
includes the return normalization, since the logarithmic cost at
$n=0$ is bounded by a constant times $1+M_1$.
Each profile has positive mass at $0<r<1$, so the global
normalization is strictly positive. Equation~\eqref{eq:g5N-moment}
and Corollary~\ref{cor:uvpeW-sharp-depth} then give the asserted
graph bounds and permit cutoff removal by the existing conditional
absolute estimates and dominated convergence. If the measure on
geometries itself varies with a cutoff, its normalized local limit
and tightness must also be justified.

\begin{corollary}[Charged curvature preparation with no fixed minimum separation]
\label{cor:g5R-fluctuating-induced}
In the first-order curvature representation of
Corollary~\ref{cor:g5R-induced-Dirac}, take the distance of the
charged readout to be the radial position $r$ of one normalized
ground state with $j\ne0$ on the finite-action surface just described.
Let the surface itself have the common normalized law proportional to
$e^{-{\cal E}[f]}\pi(\dd f)$ for the fixed finite reference measure
above. Integrate this surface and state once for the entire
electromagnetic source. Every fixed-order interacting bare graph,
including the vacuum and photon sewings, is absolutely UV finite
without counterterms.
\end{corollary}
\begin{proof}
In this joint preparation the harmonic charged field lives on the
flat normal half-line $y>0$. Its evaluation map samples $y=r$.
The radial operator selecting $r$ lives on the surface and has the
charged centrifugal term. Identifying the readout positions is the
specified boundary coupling. The two normal kinetic operators remain
distinct. No random duration is used as a distance. Elimination still
gives $P_r=\exp(-r\sqrt{-\Delta})$ with a field- and $r$-independent
hidden constraint determinant. At fixed $f,r$ a graph with $I$
internal matter lines has cost at most $C_G(1+r^{-3})^I$ by
Theorem~\ref{thm:g5M-Dirac}. A positive momentum moment on every
internal line costs at most a further fixed power of $r^{-1}$. Pick
$k$ with $6k$ larger than the resulting exponent. After averaging over
the charged state, \eqref{eq:g5N-boost} bounds the result by a fixed
polynomial in $\mathcal E[f]$ at that graph order. Its product with
$e^{-\mathcal E[f]}$ is bounded and integrable against $\pi$. The
normalization is finite and positive. This gives a uniformly vanishing
absolute UV tail after the joint preparation, and local convergence
of the regulated terms gives cutoff removal.
\end{proof}
The construction has arbitrarily short visits. The needed distribution
at short distance follows from the charged normal action and its
chosen boundary state. This corollary is coefficientwise. The entire function of $e$ obtained for the
fixed-distance partition of Theorem~\ref{thm:g5M-Dirac} would require
a stronger domination to integrate its finite-coupling bound over
these fluctuating geometries.

\paragraph{The charged ground state in conformal radial distance.}
For a fixed positive profile $f$ in the preceding curvature class,
the variable
\[
 s(r)=\int_r^{r_0}\frac{\mathrm du}{f(u)}
\]
gives a direct test of the short-distance geometry. Here $r_0$ is
a fixed interior point. Since $f(r)\le {\cal E}[f]r^3$, the
tip is at $s=+\infty$. In the angular sector $q=|j|>0$ write the
radial state in the original measure $f\,\mathrm dr$ as $\psi$.
The charged form on the tip interval becomes
$\int (|\psi_s|^2+q^2|\psi|^2)\,\mathrm ds$, while its Hilbert
space norm uses $f(r(s))^2\mathrm ds$. The ground-state equation of
energy $E$ therefore reads $\psi_{ss}=(q^2-E f^2)\psi$. In particular,
for some $C>0$ its normalized radial density satisfies
\begin{equation}
 \frac{\mathrm d\mu_f}{\mathrm dr}
       \sim C f(r)\exp[-2q s(r)]\qquad(r\downarrow0).
 \label{eq:geom-new-conformal-uv}
\end{equation}
Indeed $\int^{\infty} f(r(s))^2\mathrm ds$ is a finite radial
area integral. On a sufficiently long tail the Volterra equation
\[
 h(s)=1-\frac{E}{2q}\int_s^\infty
      (1-e^{-2q(t-s)})f(r(t))^2 h(t)\,\mathrm dt
\]
is a contraction in the supremum norm and its bounded solution tends
to one. Differentiation gives the decaying solution $e^{-qs}h(s)$ of
the radial equation. Reduction of order supplies the independent
solution growing as a nonzero multiple of $e^{qs}$. The charged form
requires $\psi\in L^2(\mathrm ds)$ at the tip and thus selects the
decaying solution. Normalization and positivity give the stated
constant. It may depend on the entire profile.

\paragraph{A finite-curvature test with the homogeneous heat lapse.}
To compare this geometric rate with the complete charged vacuum,
use the straight Wilson heat-link representation below on a fixed
four-torus of volume $V$, with positive mass $m$, zero photon mode
fixed, harmonic photon readout at $r$, and charged lapse
$\sigma^2=\kappa r$, $\kappa>0$ in the fixed reference units.
Use a profile equal to $c r^3\ell^{-\beta}$ near the tip,
$\ell=\log(1/r)$, $c>0$, $\beta>1$. It can be extended smoothly
and positively to the prescribed outer boundary. Its curvature
integrand is asymptotic to $36c/(r\ell^\beta)$, so its action is
finite. These are admissible configurations of the curvature
integral, whether or not they minimize the action.
Put $D_V=V/(192\pi^2\kappa^2)$. For this family the exact
classification of the free normalization is
\[
 \begin{array}{ll}
 \beta>3,\ \text{or }\beta=3,\ q/c>D_V &: \text{finite},\\
 1<\beta<3,\ \text{or }\beta=3,\ q/c\le D_V
       &: \text{infinite}.
 \end{array}
\]
In the finite cases the full free weight has all inverse radial
moments. Thus every bare graph at each finite electromagnetic
order is absolutely summable with fixed external Fourier probes.

In fact $2q s(r)=(q/c)r^{-2}\ell^\beta[1+O(1/\ell)]$ by
integration at the tip. Write $L=\log[1/(2m^2\kappa r)]$.
The full free logarithm in Proposition~\ref{prop:dualrev-uvheat-depth} is asymptotic to
$D_V r^{-2}L^3$. The comparison decides every case except equality
of both the power and the leading coefficient. At equality an
additional positive term in the logarithm is decisive. Put
$u_k=2\kappa r|k|^2$ and let $N(y)$ count $u_k\le y$ on the
reciprocal lattice. Comparing fundamental cells with balls gives
\[
 N(y)=\frac{V y^2}{128\pi^2(\kappa r)^2}
           +O\big(1+(y/(\kappa r))^{3/2}\big).
\]
The exact spin logarithm $2\log(1+u_k e^{L-u_k})$ bounds
$2(L-u_k+\log u_k)$ from below on $1\le u_k\le L$.
The sum of the linear terms is $2\int_0^L N(y)\,\mathrm dy$ up
to the omission of $u_k<1$, which costs $O(r^{-2}L)$ for small $r$.
The annulus $L/4\le u_k\le L/2$ alone contributes at least
$c_1r^{-2}L^2\log L$ from the nonnegative logarithms for large $L$.
Since $L=\ell+O(1)$, the resulting lower bound is
\[
 \log Z_0(r)\ \ge\ D_V r^{-2}\ell^3
        +c_2r^{-2}\ell^2\log\ell-O(r^{-2}\ell^2)
\]
with $c_2>0$. This exceeds the matched radial suppression;
the factor $f(r)$ in \eqref{eq:geom-new-conformal-uv} cannot restore integrability.
In the strict finite cases an exponential margin absorbs all
polynomial graph costs. This example keeps the same curvature
functional, charged kinetic form and photon source while varying
the allowed normal geometry. A reference measure with an atom on
one of the divergent profiles already gives infinite free mass,
since both the determinant and the geometric weight are positive.

\paragraph{Averaging over all finite-action profiles with the charged vacuum retained.}
For the original finite reference measure $\pi$ on profiles, let the
two radial half-densities sew to the normalized charged ground-state
law $\mu_f$. Keep their joint weight $e^{-A}\pi(\mathrm df)\mu_f(\mathrm dr)$,
where $A={\cal E}[f]$. This is the preparation before inserting the
charged determinant. In this paragraph the charged heat lapse is the
declared readout $\sigma^2=\kappa r^\alpha$. The geometry integral and
all Maxwell and charged factors still belong to one partition.
To see the available exponential margin, use $B=q/A$, $q=|j|$, and the
ground-energy bound $\lambda_0\le C_jA^2$ ($A\ge1$ here). In the
weighted eigenfunction identity preceding \eqref{eq:g5N-boost}, put
$W=\exp[(\delta B/2)\min(r^{-2},R)]$ for $0<\delta<1$.
Then $|W'|\le\delta B r^{-3}W$ almost everywhere. Discard the
nonnegative radial derivative form and use $j^2/f^2\ge B^2r^{-6}$ to
bound $(1-\delta^2)B^2\int r^{-6}W^2\,\mathrm d\mu_f$ by
$\lambda_0\int W^2\,\mathrm d\mu_f$. Set
$r_*=((1-\delta^2) B^2/(2\lambda_0))^{1/6}<1$.
The weighted mass below $r_*$ is at most the weighted mass above it.
Monotone convergence in $R$ then gives
\[
 \int e^{\delta q/(Ar^2)}\,\mathrm d\mu_f
 \le 2e^{\delta B/r_*^2}\le 2e^{C_{j,\delta} A^{1/3}} .
\]
For $\eta<2\sqrt q$, choose $0<\delta,\epsilon<1$ with
$\eta^2<4\delta\epsilon q$. Use
$\eta/r\le\epsilon A+\delta q/(Ar^2)$ and multiply by $e^{-A}$.
The resulting bound is integrable against every finite $\pi$.
In particular the joint preparation has an exponential inverse-distance
moment for every $0<\eta<2\sqrt{|j|}$. At fixed $V,m,\kappa$ the free
determinant obeys $\log Z_0=O(r^{-2\alpha}(1+|\log r|)^3)$.
Every fixed bare electromagnetic graph with fixed Fourier probes
costs at most $Z_0$ times a finite power of $r^{-1}$, including the
simultaneous momentum and duration estimates in \eqref{eq:dualrev-uvheat-tail}. For
$0<\alpha<1/2$ these costs are integrable under the same fluctuating
geometry law. Indeed they are bounded by $C e^{\eta/r}$ for any
positive $\eta$ at small $r$. There is one overall free normalization.
The result applies to an arbitrary fixed finite $\pi$ subject to the
measurability and boundary assumptions above.

For the real Maxwell field the uniform bound on the full determinant
in Theorem~\ref{thm:dualrev-uvheat-real} gives
\[
 \log {\cal D}_m\ \le\
 \frac{V}{32\pi^2m^2\kappa^3}r^{-3\alpha}
            (1+O(e^{-c_3/r^\alpha})),\qquad c_3>0.
\]
Consequently the unexpanded real-coupling partition and its derivatives
at each finite charge order exist for $0<\alpha<1/3$. They also
exist at $\alpha=1/3$ if $V/(32\pi^2m^2\kappa^3)<2\sqrt{|j|}$.
Choose $\eta$ strictly above the leading coefficient in this last case.
The marked-chord, inverse-mass and Gaussian readout bounds in the same
proof cost at most fixed powers of $r^{-1}$ after Maxwell expectation.
The strict exponential margin absorbs these powers, including for
smooth transverse probes and fixed open charged wavefunctions.
The comparison introduces no lower bound on $r$ or upper momentum
limit. The powers stated here are sufficient for integration over
this class of profiles. The choice of the heat-lapse law still
belongs to the representation.

\paragraph{Sharpness for a completely unspecified metric reference measure.}
The condition $\alpha<1/2$ for the free normalization cannot be relaxed
if the conclusion must hold for \emph{every} fixed finite positive
reference measure on the admitted profiles. Choose $h=2^{-n}$
small. Near the tip take
$f_h(r)=r^2\chi(r/h)$ with a fixed smooth positive $\chi(u)$ equal to
$u^2$ for $u\le1$ and to one for $u\ge2$. Beyond a fixed interior
radius extend $r^2$ by the same positive smooth outer profile for all
$h$, with the required boundary values. There is a constant $C$
independent of $h$ such that
\[
 A_h\le C/h,\qquad
 \mu_{f_h}([h,2h])\ge C^{-1}h^3 e^{-C/h}.
\]
For the first inequality, on $(0,h)$ the curvature integrand is
$144/h^2$, on $(h,2h)$ it is $O(h^{-2})$, and on the subsequent
$r^2$ region it is $4/r^2$.

For the second inequality a fixed trial function in the outer
region bounds all ground energies $E_h$ by $E_*<\infty$. Near
the tip $f_h\le C r^2$, so the angular part of the energy gives
$\mu_{f_h}((0,b))\le C E_* b^4/q^2$ for small fixed $b$. Choose
$b$ so this mass is below $1/2$ and $q^2-E_h f_h^2>0$ on $(0,b)$.
The positive ground wavefunctions, normalized in $L^2(f_h\,\mathrm dr)$,
then satisfy $\inf_h\psi_h(b)>0$ for all sufficiently small $h$.
Indeed the coefficients on $[b/2,1]$ are fixed, and the norm and
energy give a uniform $H^1$ bound there. A sequence approaching zero
at $b$ would have a uniformly convergent subsequence with a nonzero
nonnegative limit (its mass to the right of $b$ is at least $1/2$).
Passing to a subsequence of the bounded energies gives the same
regular radial equation at the limiting energy. An interior zero
would give zero value and derivative and hence the zero solution.

In the conformal coordinate measured from $b$, the function
$\psi_h(b)e^{-qs}$ is a subsolution of the ground equation.
Both solutions vanish at the tip. Positivity of $q^2-E_h f_h^2$
gives $\psi_h(r)\ge\psi_h(b)e^{-qs(r)}$ for $r<b$. On $[h,2h]$
we have $f_h\ge c h^2$ and $s(r)\le C/h$ by integration of
$1/f_h$, proving the mass bound. Now give these profiles reference
masses $2^{-n}$, independent of the carrier momenta. The free
vacuum is decreasing in $r$, and Proposition~\ref{prop:dualrev-uvheat-depth} gives
$\log Z_0(2h)\ge c h^{-2\alpha}[\log(1/h)]^3$ for small $h$.
For $\alpha\ge1/2$ this outruns the combined $O(1/h)$ cost of
the curvature action and charged ground state. The sum of the
positive contributions from these intervals diverges. For
$1/2\le\alpha<3/2$, each member separately still has a finite vacuum
and every finite graph: its own tip is quartic, with a fixed
$h$-dependent coefficient in \eqref{eq:geom-new-conformal-uv}.
Thus a metric-by-metric proof can fail after the geometry is
integrated, even though its curvature action and all the individual
bare graph sums are finite.

\subsection{The photon cost of a local curvature contact}
\label{subsec:g5F-local}
The tree-level normalization of the photon fixes only the quadratic
part of the reduced action. Consider a real transverse photon on the
four-torus of period $2\pi$, with its harmonic mode fixed. Use normalized
volume, a translation-invariant positive Gaussian covariance $D(p)$
on $p^\perp$ at $p\ne0$, and the Euclidean exterior norm
$|F|^2=\sum_{i<j}|F_{ij}|^2$. The curvature in this paragraph is the
Abelian $F=\dd A$. Its one-point Gaussian contraction is governed by
\begin{equation}
   B_F(D)=\sum_{p\ne0}|p|^2\operatorname{Tr}_{p^\perp}D(p).
   \label{eq:g5F-budget}
\end{equation}
The following statement tests local curvature contacts in the chosen
physical expansion. It does not require power counting in a preferred
momentum direction.

\begin{proposition}[Exact criterion for local curvature interactions]
\label{prop:g5F-contacts}
For a fixed integer $N\ge2$ introduce the bare interaction
$V_N=\int |F(x)|^{2N}\dd x$. The following are equivalent:
$B_F(D)<\infty$, and the entire family of complete graphs generated
by $V_N$ is graphwise absolutely UV finite at every finite order,
including vacuum graphs and all fixed external Fourier tests.
The sufficiency statement also holds for any finite list of local
polynomials in $F$ with bounded constant coefficients.
\end{proposition}
\begin{proof}
For an internal line at momentum $p$, each endpoint is
differentiated once by $p\wedge$. The absolute value of any
polarized curvature covariance is bounded by
$|p|^2\operatorname{Tr}_{p^\perp}D(p)$. Each polynomial vertex at
a fixed order has finitely many tensor components. Drop the
momentum-conservation restrictions in the absolute majorant and sum
independently over the internal line momenta. A graph with $I$ such
lines costs at most a constant times $B_F(D)^I$. External momenta
are fixed and cutoff removal follows by domination.

Conversely let $t f(x)$ be a nonzero real external curvature of a
fixed finite Fourier potential. Denote the internal Gaussian curvature
by $G=\dd A$. In $|tf+G|^{2N}$ the coefficient of $t^{2N-2}$ with
one Wick contraction is the integral of
\[
 N|f|^{2N-2}\mathbb E|G|^2+
 2N(N-1)|f|^{2N-4}\mathbb E[(f,G)^2].
\]
At symmetric finite cutoff $\Lambda$,
$\mathbb E|G(x)|^2=B_F(D_\Lambda)$ independently of $x$ because
$|p\wedge u|^2=|p|^2|u|^2$ on $p^\perp$. The second term is
nonnegative and bounded by $2N(N-1)|f|^{2N-2}B_F(D_\Lambda)$.
Consequently this complete one-vertex contact is bounded above and
below by positive fixed multiples of $B_F(D_\Lambda)$. It has
external legs and survives division by the zero-source partition.
Since the test is a finite linear combination of fixed external
Fourier configurations, divergence implies failure of the universal
graphwise assertion on this support. Individual graphs, including
trees, can be finite for larger covariances.
\end{proof}

For example, take $D(p)=\Pi(p)\Phi(|p|^2)/(2|p|^2)$ with
$\Pi$ the transverse projector and finite positive Laplace law
$\rho$ for $\Phi$. Tonelli's theorem gives
\[
 B_F(D)=\tfrac32\int_{[0,\infty)}\rho(\dd t)
                   \sum_{p\ne0}e^{-t|p|^2}.
\]
The sum is comparable to $t^{-2}$ for $0<t\le1$, by summing
four-dimensional shells above and below $t^{-1/2}$; for $t\ge1$ it
is bounded. Hence the local curvature criterion requires exactly
zero mass at $t=0$ and a finite inverse second moment there. The
triangular law of \eqref{eq:uvpeW-depth} has a nonzero density at
zero for each $L>0$. Its Wilson sewing is finite at fixed $L$,
whereas its local curvature contact has growth of order
$\Lambda^2/L$ as $\Lambda\to\infty$. The difference comes from
the actual interval factors at the Wilson vertices. A shared
fluctuating law requires integrating the whole graph cost at its
conditional covariance.

This test probes a genuine ambiguity when the curvature pairing is
specified only by its symmetry and value at the vacuum. On a flat
Abelian block with unit pairing, the positive local choices
$1$ and $1+|F|^{2N-2}$ in dimensionless units give the same
stationary Maxwell Hessian and all classical Taylor coefficients
through order $2N-1$. Their curvature-square actions first differ by
a nonzero multiple of $V_N$. Gauge covariance of this scalar
pairing follows from Abelian curvature invariance. On a comparison
branch with an ordinary Maxwell covariance the new contact costs
order $\Lambda^4$. On a branch where the coupling is to the visited
Gaussian photon with covariance \eqref{eq:uvpeW-depth}, it instead
costs order $\Lambda^2/L$. These two sewings are distinct.
The ambiguity applies if the allowed constitutive rule admits this
curvature dependence. A native representation can exclude the
deformation or resolve its vertex through a different geometric
history. Establishing which alternative holds requires its
functional pairing and the complete induced source jets. Matching
an arbitrarily long finite list of classical jets does not by
itself settle that question.

\paragraph{A normal action that controls the curvature and Wilson vertices together}
A complete positive example tests both requirements on the same
photon. On the transverse four-torus keep $p\ne0$, $z=|p|^2$, and
introduce $d$ comparison normal coordinates in $\RR^d$. Start with a
real Gaussian transverse field $Z$, with translation-invariant normal
precision symbol
$\prod_{j=1}^d(z^2+\xi_j^2)$, where $\xi$ is the normal Fourier
variable. For $d=3$ its positive local quadratic form is
\begin{equation}
 \tfrac12\sum_{S\subseteq\{1,2,3\}}\int
  \big|\partial_{r_S}(-\Delta_x)^{3-|S|}Z\big|^2\,
                                         \dd^4x\,\dd^3 r .
 \label{eq:g5F-normal-action}
\end{equation}
Here $\partial_{r_S}$ differentiates once in each coordinate in
$S$. This quadratic form has a representation by the derivative
connection images in Section~\ref{sec:g5-induction}. In radial
gauge, for each nonempty $S$ choose $j\in S$ and the tangential
connection image $\partial_{r_{S\setminus\{j\}}}
(-\Delta_x)^{3-|S|}Z_i$. Pair its $r_j,i$ curvatures. For the
empty set pair the tangential curls of the four images
$\partial_a(-\Delta_x)^2 Z_i$ and sum over $a$. Transversality
gives the missing two powers of $z$ in the resulting squared
symbol $z^6$. These are copies of the same transverse field with
constant component pairings. For a general $d$ choose the observed mode $U_p=(2z)^{(d-1)/2}Z_p$.
Fourier inversion of each normal factor gives
\[
 \mathbb E[U_p(r)U_p(s)^*]
 =\frac{\Pi(p)}{2z}\exp\!\left(-z\sum_{j=1}^d|r_j-s_j|\right).
\]
Thus $U(0)$ has an ordinary transverse Maxwell pole for every $d$.
Use the normalized average of $U$ over the cube $[0,L]^d$ as
the same photon $b$ in every Wilson and curvature vertex. Its exact
covariance is
\[
 D^{(d)}_L(p)=\frac{\Pi(p)}{2z}\Phi_L(z)^d ,
\]
where $\Phi_L$ is the triangular transform in
\eqref{eq:uvpeW-depth}. The product follows by integrating the
single full covariance over two independent cube coordinates at
the two insertions. It retains the quantum normal fluctuations.

For $L>0$ fixed, the curvature budget is finite exactly for $d>2$.
Indeed it is comparable to
$\sum_{p\ne0}(1+L|p|^2)^{-d}$. Its four-dimensional dyadic shell
at radius $R\gg L^{-1/2}$ contributes order
$L^{-d} R^{4-2d}$. In particular $d=2$ gives a logarithm. For
$d=3$ one obtains $B_F(D^{(3)}_L)\asymp L^{-2}$ as $L\downarrow0$
by adding the shells on both sides of $R=L^{-1/2}$; the budget is
bounded for $L\ge1$.

Now use the charged Wilson block of \eqref{eq:uvpeW-hops}, the same
common marked depth law \eqref{eq:uvpeW-return}, and the field $b$
also in a finite collection of local Abelian curvature polynomials.
Every complete graph at fixed visit, electric and curvature
interaction orders is absolutely UV finite after normal averaging
on this chart. In fact the charged loop sums after averaging over
$h_j$ obey the argument of Proposition~\ref{thm:uvpeW-sewing}.
A Wilson--Wilson line has the Wilson budget, which is at most a
constant times $B_F$ since $|p|\ge1$. Each curvature--curvature
line costs $B_F$. A mixed line has a product of the Wilson interval
and one curvature derivative, bounded in its absolute sum by the
geometric mean of the two budgets using positivity of the photon
covariance and Cauchy--Schwarz. The joint bound per graph with $I$
sewn photon lines is $C_G(1+B_F)^I$ before the common $L$ integral.
Every required inverse moment exists under the one chosen return law.
The same argument applies at physical visit coefficient $g=1$
using Theorem~\ref{thm:g5W-electromagnetic-orders}. The curvature
vertices introduce no charged line or additional dependence on
the $h_j$. The proof's charged Fourier sums and restored zero-photon
determinant are therefore unchanged, and each new photon line
obeys the curvature or mixed budget above. This yields absolute
convergence at each fixed order in the electric and nonlinear
curvature couplings. In the enlarged constitutive class of
Section~\ref{sec:g5-induction}, a term $|\dd b|^{2N}$ can be
implemented by the positive scalar factor $|\dd b|^{2N-2}$
pairing the curvatures of the specified averaged photon image.
This permits a zero pairing at zero field in that interaction
block. It is a choice about the microscopic representation and
its state-dependent readout.

For each positive $L$, take $V=\sum_N a_N\int|\dd b|^{2N}\dd x$,
a finite sum with $a_N\ge0$. The pointwise curvature has finite
Gaussian variance by \eqref{eq:g5F-budget}. Hence
$\mathbb E\int |\dd b|^{2N}\dd x<\infty$ by the Gaussian moments and
Tonelli, and $V$ is finite almost surely. Multiplying the relative
determinant in Proposition~\ref{prop:uvpeW-finite-coupling} by
$e^{-V}$ defines a strictly positive integrable weight. The
determinant bound holds for all allowed samples and displacements,
so the common history average has a finite, positive relative
partition at real $e$ and $g=1$. Its choice of a
three-coordinate normal precision, the cube readout, and which
photon enters the curvature contacts are independent inputs in
this comparison. Identifying them on a physical GCM branch would
be a further dynamical result.

\subsection{A covariant heat-link source with a soft Maxwell response}
\label{subsec:s6-heatlink}

A charged loop has an independent momentum even after all photon
momenta are fixed. The Gaussian charged displacement controls that
momentum. Harmonic extension supplies a separate factor for each
photon attached to the charged line. Both factors come from one
source. Its interval transport generates the electromagnetic current
and the multiphoton contacts by differentiation~\cite{TerningS6}.
The construction uses one locked sector plane $\bin=(s,b_s)$.

Use Hermitian Euclidean Dirac matrices and $a,m>0$ in reference units.
On the $2\pi$-periodic four-torus, fix the harmonic photon mode in the
dynamical integral. In the trivial small-gauge chart decompose the
nonzero modes as $A=A^\parallel+A^T$. The sampled connection is
\[
 B=A^\parallel+P_a A^T,\qquad P_a(p)=e^{-a|p|}.
\]
The bounded solution of $\partial_y^2 u_p=|p|^2u_p$ with boundary
value $A^T_p$ gives $P_a A^T$ at normal distance $a$. The boundary
field has the Maxwell Gaussian law with covariance
$\Pi^T(p)/|p|^2$ at $p\ne0$, where $\Pi^T(p)$ projects onto
$p^\perp$. The harmonic readout is deterministic conditional on
the boundary field. Equation~\eqref{eq:uvpeC-phase} expresses its
Fourier multiplier as a positive average of Gaussian geo-phases.

Choose a displacement endpoint at positive heat lapse $a^2$ with
density $g(z)=(4\pi a^2)^{-2}\exp(-|z|^2/(4a^2))$ on $\RR^4$.
The displacement is tangential to the four-dimensional Euclidean
carrier and includes its three spatial directions. This full tangential
law, together with the specified harmonic photon operator, controls
all four momentum directions in the graph estimate. Use the straight
segment between the endpoints on the universal cover, keeping its
winding on projection to the torus. Its charged
transport and spin-weighted density are
\begin{align}
 W_B(x,z)&=\exp\!\left(\ii e\int_0^1 z\cdot B(x+tz)\,\dd t\right),
 \nonumber\\
 K(z)&=-\gamma^\mu\partial_\mu g(z)
       =g(z)\,\gamma\cdot z/(2a^2). \label{eq:s6-heatK}
\end{align}
In the four-arrow curvature contraction \eqref{eq:g5-arrow-product},
take a local block with $U=I$, $V=m\gamma_5$ and the coincident
spinor row and column. Take a direct integral of segment blocks with
the same spinor pair, row at $x$, column at $x+z$, $U=W_B(x,z)$,
$V=\gamma_5\gamma\cdot z/(2a^2)$ and pairing measure $g(z)\,\dd z$.
The Gaussian is the normalized fundamental solution of the heat
equation at the chosen lapse. The displacement insertion is its first
derivative as in \eqref{eq:s6-heatK}. The cyclic functional integrates
the contracted density over $x,z$ before the single spinor Gaussian.
The Hermitian real form and constant $\gamma_5$ row substitution of
Section~\ref{sec:g5-induction} therefore give the exact operator
\begin{align}
 D_B&=mI+T_B, \nonumber\\
 (T_B f)(x)&=\int_{\RR^4} K(z)W_B(x,z)f(x+z)\,\dd z .
 \label{eq:s6-D}
\end{align}
Indeed opposite paths are adjoints and $K$ is odd, so $T_B^*=-T_B$
for real $B$. Also $\gamma_5T_B\gamma_5=-T_B$. The shift
$A\mapsto A+\dd\lambda$ changes $D_B$ to
$e^{-\ii e\lambda}D_B e^{\ii e\lambda}$, including the contact
vertices. For soft fields $\int K=0$ and $\int z_\mu K(z)\,\dd z=\gamma_\mu$.
Thus the leading charged operator is
$m+\gamma^\mu(\partial_\mu+\ii e B_\mu)$, and $P_a(0)=1$.
The exact free symbol and its inverse bound are
\begin{equation}
 d(k)=mI+\ii\gamma\cdot k\,e^{-a^2|k|^2},
 \qquad \|d(k)^{-1}\|\le m^{-1}.
 \label{eq:s6-d}
\end{equation}
For the unsampled photon the one-mark vertex is
$e\int_0^1\partial_{k_\mu}d(k+tp)\,\dd t$. Contracting it with
$p_\mu$ gives $e[d(k+p)-d(k)]$. The transverse attenuation leaves
this Ward identity intact. All higher derivatives are fixed by
the same transport. For $a(|k|+|p|)\ll1$, the unsampled
one-photon vertex is $\ii e\gamma^\mu$ up to
$O(|e|a^2(|k|+|p|)^2)$. The transverse boundary vertex has the
additional factor $e^{-a|p|}=1+O(a|p|)$. At large $|k|$,
$d(k)$ in \eqref{eq:s6-d} tends to $mI$, whereas the charged
symbol in \eqref{eq:g5M-induced} grows linearly.

On the transverse slice, write $\mathbb P_M$ for the normalized
boundary Maxwell law. The charged vacuum factor relative to the
mass operator $mI$ is the ordinary Fredholm determinant
$\mathcal D_m(B)=\det_{\mathrm F}(I+m^{-1}T_B)$ on
$L^2(\mathbb T^4,\CC^4)$. For independent Grassmann sources
$\bar\xi,\xi$, the normalized induced functional is
\begin{equation}
 Z_e[J,\bar\xi,\xi]=
 \frac{\displaystyle\int\mathcal D_m(B)
       e^{\langle J,A^T\rangle+\bar\xi D_B^{-1}\xi}
           \,\dd\mathbb P_M(A^T)}
      {\displaystyle\int\mathcal D_m(B)\,\dd\mathbb P_M(A^T)} .
 \label{eq:s6-heatlink-source}
\end{equation}
Here $J$ and the charged sources test finitely many fixed Fourier
modes. At fixed $a,m$, the vacuum value $\mathcal D_m(0)$ is a
common scalar factor in the numerator and denominator. If the
readout fluctuates, that value depends on the geometry and weights
its law before the single overall normalization. This weighted law
is displayed in Corollary~\ref{cor:s6-depth}.

\begin{theorem}[All electromagnetic orders of the heat-link source]
\label{thm:s6-heatlink-uv}
For every $a,m>0$, each interacting bare graph of
\eqref{eq:s6-heatlink-source} at every finite order in $e$ has an
absolutely convergent momentum sum on the zero-mode-fixed torus,
without UV regularization or counterterms. The finite-volume
partition is finite and strictly positive for real $e$. For the
stated Fourier probes the normalized source is smooth in real $e$.
Its Taylor coefficients are given by the graph expansion.
\end{theorem}
\begin{proof}
\emph{Existence of the source.}
In transverse gauge $B=P_a A^T$ is almost surely smooth. For any
integer $d\ge0$, the random series
$\sum_{p\ne0}(1+|p|)^d e^{-a|p|}|A^T_p|$ has moments of all
finite orders by Minkowski's inequality and the Gaussian mode
bounds. The periodization of $K(z)W_B(x,z)$ gives a smooth matrix
kernel on the torus. Indeed, for real $B$ each spatial derivative
is dominated by the periodization of a polynomial in $|z|$
times the Gaussian, with coefficients depending on finitely many
derivatives of $B$. In dimension four $(1-\Delta)^{-2}$ is
Hilbert--Schmidt. The factorization
\[
 T_B=(1-\Delta)^{-2}\big[(1-\Delta)^2 T_B\big]
\]
writes $T_B$ as the product of two Hilbert--Schmidt operators
because the second factor has a square-integrable kernel. Hence
$T_B$ is trace class and $\mathcal D_m(B)$ is an ordinary
Fredholm determinant~\cite{SimonUVpeDet}. The open inverse is a
bounded operator with $\|D_B^{-1}\|\le m^{-1}$ by skew adjointness.

There is a bound independent of the photon sample and the real
coupling. Use $|W_B|=1$ in the periodized kernel to get
$\|T_B\|_2\le H_a$, where $H_a$ is the Hilbert--Schmidt norm
bound from the periodization of $\|K(z)\|$ and $\|\cdot\|_2$
denotes the Hilbert--Schmidt norm. Skew adjointness and
$\gamma_5 T_B\gamma_5=-T_B$ pair its nonzero eigenvalues
as $\ii t_j,-\ii t_j$, with $t_j>0$ and multiplicities included.
The determinant then obeys
\begin{equation}
 1\le \mathcal D_m(B)=\prod_j(1+t_j^2/m^2)
    \le\exp\!\left(H_a^2/(2m^2)\right).
 \label{eq:s6-ordinary-paired-det}
\end{equation}
For real $e$ these bounds prove that \eqref{eq:s6-heatlink-source}
exists. For any compact real $e$ interval, differentiation of the
periodized kernel and the same trace-class factorization bound each
fixed $\|\partial_e^j T_B\|_1$ by a polynomial in $\|B\|_{C^4}$.
Here $\|\cdot\|_1$ is the trace norm. Derivatives of the logarithm
come from $\partial_e\log\mathcal D_m(B)
=\operatorname{Tr}(D_B^{-1}\partial_e T_B)$ and the inverse
differentiation formula. The bounds on the inverse and the
determinant dominate every fixed derivative of the integrand by a
polynomial in Gaussian norms, including the moments of the fixed
Fourier probes. Thus one can differentiate the normalized integrals
at every finite order. The Grassmann-source exponential is taken
coefficientwise and the fixed photon source also has Gaussian
exponential moments.

\emph{Absolute graph sums.}
At a vertex with $r\ge1$ photon marks, the Fourier coefficient is
an average over $t_j\in[0,1]$ of the Fourier transform with positive
exponent of $K(z)z_{\mu_1}\cdots z_{\mu_r}$ evaluated at
$k+\sum_j t_jp_j$, where $k$ is the incoming charged momentum.
Each transverse mark also supplies $P_a(p_j)$. This Fourier
transform is Schwartz. For $N>4$,
$\sup_{b\in\RR^4}\sum_{k\in\mathbb Z^4}(1+|k+b|)^{-N}<\infty$
by reducing $b$ modulo the lattice and counting shells.
Fix all the photon momenta first. The charged loop momentum flows
through every vertex on its loop, so the translated lattice bound
at one marked vertex controls that whole sum. The other vertices
use their supremum bounds and each free charged inverse costs at
most $m^{-1}$. The incoming mode and photon transfers already fix
all momenta on an open charged line. Each internal photon has two
marks, giving $e^{-a|p|}$ at each end, including contacts for which
the two ends are on the same vertex. Taking absolute values and
dropping the remaining lattice constraints now costs at most
\[
 \sum_{p\in\mathbb Z^4\setminus\{0\}}|p|^{-2}e^{-2a|p|}<\infty
\]
per photon. The uniformity of the charged bound in the photon
momenta is what controls overlapping hard regions. The traced products and open
chains obtained by differentiating the determinant and the
inverse are therefore the usual sums over the resulting bare
graphs with their full matrix and source coefficients.
Normalization or passage to connected graphs uses finite products
of these finite coefficients at each order.

The continuum UV estimate at fixed external data uses the
translated $L^1$ norm on one vertex per closed charged loop.
With the stated infrared bounds on photon covariances, choose
the independent photon momenta on chords of the graph with its
charged lines collapsed. Apply the $L^1$ bounds on chord
covariances and the uniform bounds on the remaining edges. The
massless photon-only soft estimates below allow these photon
contractions down to zero momentum.
\end{proof}

\paragraph{A graph with fixed winding}
In the torus proof the winding of a closed charged loop can be kept
fixed, including winding zero. On the universal cover let $z_i$
be the lifted displacements around that loop. The closing
condition is $\sum_i z_i=2\pi n$ with $n\in\mathbb Z^4$.
The free inverse on the cover is an identity delta plus a
Schwartz kernel, by \eqref{eq:s6-d}. For a marker
$\vartheta\in[0,1]^4$, multiply each lifted charged kernel, including
the free inverse, by $e^{\ii\vartheta\cdot z_i}$ prior to
periodization. The resulting factor is
$e^{2\pi\ii\vartheta\cdot n}$ on the loop. Projection by its
Fourier character $e^{-2\pi\ii\vartheta\cdot n}$ selects winding
$n$. The marker shifts the charged Fourier arguments by
$\vartheta$, preserving the uniform lattice and inverse bounds.
Projection therefore commutes with the absolutely convergent
momentum sums at each fixed winding. In particular, the Gaussian
at a marked vertex still controls the charged loop in the
zero-winding sector even though the displacements sum to zero.
That projection retains positive-length returns; it does not restrict
the loop to the empty identity path.

One can also sum the winding-resolved coefficients absolutely.
Give each of the $h$ closed charged loops its own marker
$\vartheta_i$. After its charged momentum sum, the coefficient is
periodic in each marker by reindexing that loop momentum. Every
marker derivative either increases the number of momentum
derivatives on a Schwartz link multiplier or differentiates the
bounded inverse in \eqref{eq:s6-d}. At each finite derivative order
the resulting loop still admits the translated lattice bound on
one marked link and supremum bounds on all other factors. These
bounds justify differentiation of the summed coefficient and
are uniform in the remaining photon momenta after extracting
their readout factors. Applying $(1-\Delta_{\boldsymbol\vartheta})^M$
under the Fourier projection bounds a winding coefficient at
$\mathbf n=(n_1,\ldots,n_h)$ by a constant times
$(1+4\pi^2|\mathbf n|^2)^{-M}$. Taking an integer $M>2h$ makes the winding
sum finite even after summing the photon momenta.
Unmarked charged loops belong to the free determinant. Their
geometric weight becomes explicit in Corollary~\ref{cor:s6-depth}.

\paragraph{Continuum soft test for every massive spinor}
Use the translation-invariant Euclidean integral per unit volume on
$\RR^4$ to probe $p\to0$. A constant test connection $b$ gives the
symbol $d(k+eb)$ before the transverse readout. Differentiating the
induced negative log determinant twice gives the combined seagull
and bubble at $p=0$:
\begin{align}
 \Pi^{B}_{\mu\nu}(0)
 &=-e^2\int_{\RR^4}\partial_\mu\partial_\nu
      \log\det[d(k)/m]\,\frac{\dd^4 k}{(2\pi)^4}=0 ,
 \label{eq:s6-zero}\\[-.4ex]
 \det[d(k)/m]
 &=\left(1+m^{-2}|k|^2e^{-2a^2|k|^2}\right)^2 .\nonumber
\end{align}
Indeed the two differentiated terms are
$\Tr(d^{-1}\partial_{\mu\nu}d)$ and
$-\Tr(d^{-1}\partial_\mu d\,d^{-1}\partial_\nu d)$.
Each is integrable on its own, and the logarithm is Schwartz.
For small external $p$, two derivatives of the unsampled loop
integrand have integrable majorants: derivatives of the segment
vertices add Gaussian moments, all translates are bounded on
compact $p$-sets, and derivatives of $d^{-1}$ are bounded.
Simultaneous inversion of the Euclidean coordinates and both vector
indices leaves the tensor unchanged, by
$\gamma_5d(k)\gamma_5=d(-k)$. Accordingly the unsampled loop is
even and $C^2$, and its correction is $O(|p|^2)$. Restoring the two
external transverse readouts multiplies by $e^{-2a|p|}\le1$.
A flat connection can carry nontrivial torus holonomy; the integral
identity \eqref{eq:s6-zero} uses continuum translations.

\begin{proposition}[A soft factor on every photon leg for every positive mass]
\label{prop:s6-heatlink-soft}
On the translation-invariant continuum chart, fix $a,m>0$. Let
$V^{(r)}$ be the complete $r$th derivative of the relative negative
log charged determinant at zero connection, with its overall momentum
delta removed. Use transverse boundary photons of unit polarization.
For $r\ge1$,
\begin{equation}
 |V^{(r)}(p_1,\ldots,p_r)|
 \le C_{r,a,m}|e|^r
       \prod_{j=1}^r e^{-a|p_j|}\min(1,|p_j|).
 \label{eq:s6-eachleg}
\end{equation}
Every photon-only graph obtained by sewing these complete vertices is
absolutely integrable at both ends of each photon loop. At each finite
electromagnetic order its contribution to the amputated proper
transverse two-point function is $O(|p|^2)$ as $p\to0$.
\end{proposition}
\begin{proof}
First vary the sampled connection $B$ before restoring the transverse
readouts. For an auxiliary positive mass $v\ge m$ write
$R_{v,B}=(vI+T_B)^{-1}$. The positive-order coefficients of the
relative negative log determinant are those of
\begin{equation}
 \int_m^\infty \operatorname{Tr}(R_{v,B}-R_{v,0})\,\dd v .
 \label{eq:s6-massres}
\end{equation}
For real bounded smooth probes the difference has a well-defined
diagonal density. Indeed
$R_{v,B}-v^{-1}I=-v^{-2}T_B+v^{-2}T_B R_{v,B} T_B$;
the two outer kernels have continuous $L^2$ rows and columns and
$\|R_{v,B}\|\le v^{-1}$. Equation~\eqref{eq:s6-massres}
is the logarithmic derivative in the mass with its value at
infinite mass fixed by free normalization. In infinite volume the
trace means the integral of the diagonal kernel, with the overall
translation delta taken out for plane-wave probes. We justify this
formula directly at each marked coefficient. Finite matrices satisfy
it by differentiation. For the continuum coefficient the same
identity follows by cyclically differentiating each finite-order
trace of marked vertices and free inverses in $-\log\det$; the
charged Fourier integral converges absolutely by the argument of
Theorem~\ref{thm:s6-heatlink-uv}, using an integral for the one
Schwartz factor. Each positive-order trace vanishes at infinite mass.

In position space $R_{v,0}=v^{-1}I+L_v$, with $I$ the identity
delta kernel. Writing $t(k)=\ii\gamma\cdot k\,e^{-a^2|k|^2}$,
the multiplier of $L_v$ is $-v^{-1}(vI+t(k))^{-1}t(k)$.
All Schwartz seminorms of $v^2 L_v$ are bounded uniformly for
$v\ge m$. Indeed every derivative of the inverse is a finite sum
of products of inverses of norm at most $v^{-1}$ and derivatives
of the fixed Schwartz multiplier. Fourier inversion gives the
corresponding bounds for the kernel. In particular the absolute
kernel measure of $R_{v,0}$ has weighted $L^1$ norm bounded by
$C_{d,a,m}/v$ with each fixed weight $(1+|z|)^d$.

Anchor the diagonal of the resolvent at $x$. For a soft plane-wave
probe $b_j(y)=\epsilon_j e^{\ii p_j\cdot y}$ replace it by
$b_j(y)-b_j(x)$. The removed constant one-form is an exact gradient
on $\RR^4$. Gauge covariance conjugates the full resolvent by the unitary
associated with that gradient. The two phases cancel on the
diagonal at $x$. Holding the other probes fixed, this equality
allows each soft probe to be replaced separately. The calculation
uses real sine and cosine probes and extends to plane waves by
multilinearity. On a link from $y$ to $z$ its insertion has magnitude
at most
\[
 |p_j|\,|z-y|\bigl(|y-x|+|z-x|\bigr).
\]
For $|p_j|>1$ use the original bounded plane wave, whose insertion
costs at most $|z-y|$. These operations take place at the anchored
density, before imposing momentum conservation. A constant
connection on a torus would require a separate holonomy test.

At order $r$ the inverse derivative is a finite sum of strings
$R_{v,0}T_{J_1}R_{v,0}\cdots T_{J_h}R_{v,0}$, where the $h\ge1$
nonempty ordered sets partition the marks and $T_J$ differentiates
one link. Every $T_J$ supplies a Gaussian derivative kernel $K$
times its interval insertions. In each closed diagonal convolution,
the distance of any vertex from $x$ is at most the sum of the step
lengths. The polynomial in those lengths from the preceding bounds
is dominated by a product of fixed polynomial weights on the steps.
Use the weighted supremum of one marked $K$ to impose closure and
the weighted $L^1$ norms on the other steps. This bounds each
anchored term by
$C_{r,a,m}|e|^r v^{-h-1}\prod_j\min(1,|p_j|)$ for all $p_j$.
The bound is integrable in $v$. It justifies the mass integral and
the marked convolutions, including coincident-link contacts.
Restoring the transverse factors proves \eqref{eq:s6-eachleg}.
For the unsampled vertex, a longitudinal insertion is an exact
gradient and its full Ward contraction vanishes by the same endpoint
identity.

Each sewn transverse photon now costs at most
$C e^{-2a|q|}\min(1,|q|)^2/|q|^2$, bounded by $Ce^{-2a|q|}$
and integrable on $\RR^4$. In a connected photon graph choose
the independent chord momenta relative to a spanning tree,
including self-loops. Use the integrable bounds on those chords and
the uniform bounds on the other edges. This controls every overlap
and every soft region uniformly in the external momenta. At exceptional external momenta a bridge can be forced to carry
exactly zero momentum. Such a term uses the separately declared
zero-mode prescription. A proper photon graph has no photon bridge.
At each finite order there are only finitely many marked vertices
and pairings. The two external leg factors give the quadratic
small-momentum bound.
\end{proof}

\begin{corollary}[Continuum photon sewing with open charged wave packets]
\label{cor:geom-uv-commonsink}
At fixed positive charged width and harmonic readout, take a connected
Euclidean graph on $\mathbb R^4$ with at least one open charged chain.
Smear both ends of every open chain with $L^2$ spinor wavefunctions.
External photon momenta are fixed and their legs are amputated.
Collect all contacts of each closed charged loop into its determinant
vertex. In the graph of photon lines, contract each charged chain or
loop to a vertex, then identify all open vertices with a single root.
Assume this rooted multigraph has no bridge, meaning that removing
any single photon edge leaves it connected. The remaining momentum integral is
absolutely convergent at zero and at infinity at every finite order,
also at exceptional external photon momenta. This includes graphs
with no closed charged loops. The conclusion concerns fixed-depth
continuum coefficients.
\end{corollary}
\begin{proof}
For fixed photon transfers, the absolute open-chain coefficient is
bounded uniformly after its endpoint momenta are integrated. The
free charged inverses cost at most $m^{-1}$, and each marked Wilson
vertex has a bounded Fourier multiplier before harmonic readout.
If $Q$ is the total transfer, the remaining integral of the two
wavefunction norms is bounded by
\[
 \int |\widehat f(k)|\,|\widehat g(k+Q)|\,\mathrm d^4 k
       \le \|\widehat f\|_2\|\widehat g\|_2.
\]
No constraint on the photon transfers remains from an open chain.
At each end on a closed vertex the complete Ward estimate
\eqref{eq:s6-eachleg} supplies $\min(1,|q|)$. Thus a photon of
momentum $q$ with two, one, or zero ends on a closed vertex is
bounded, respectively, by a constant times
\[
 e^{-2a|q|},\qquad g(q)=e^{-2a|q|}/|q|,\qquad
 h(q)=e^{-2a|q|}/|q|^2 ,
\]
where $a>0$ is the harmonic height and units are fixed as in the
preceding estimate. All three functions belong to $L^1(\mathbb R^4)$
and $g$ also belongs to $L^2(\mathbb R^4)$.

In each connected component of closed vertices choose a spanning tree.
Momentum conservation fixes its tree momenta with unit routing
Jacobian. Their photon bounds are uniform. Each remaining internal
chord in that component has an integrable exponential bound.
One constraint per component survives, fixing the signed sum of its
$n$ photon transfers to open chains. Connectedness and absence of
bridges give $n\ge2$. Integration of these transfers
is bounded by $\|g*g\|_\infty\|g\|_1^{n-2}$, uniformly in the
prescribed total, since $n\ge2$ and the functions are even.
Cauchy--Schwarz gives $\|g*g\|_\infty\le\|g\|_2^2$.
Transfers between open chains, including self-contractions, are free
and cost $\|h\|_1$ each. This bounds all overlapping soft and hard
regions with the full massless covariance. It also explains the
bridge condition. Every edge of the rooted graph lies on a cycle,
so its transfer varies with at least one integration variable. A
zero transfer then occupies a null set even for exceptional external
data. A bridge could instead be forced to carry exactly zero, where
the transverse tensor needs a separate prescription. Soft cancellation here uses
complete charged-loop derivatives. At fixed volume the earlier
ultraviolet estimate controls individual contact diagrams as well.
\end{proof}

The ultraviolet assertion of Theorem~\ref{thm:s6-heatlink-uv}
holds for individual bare graphs including distinct link contacts.
The stronger infrared bound in \eqref{eq:s6-eachleg} uses the
complete derivative of each charged determinant. At a fixed
perturbative truncation the bare Maxwell coefficient is positive
and the uniformly bounded relative correction starts at $e^2$.
For sufficiently small real $|e|$, the eigenvalues of the
truncated Euclidean inverse response on the transverse space
lie between positive multiples of $|p|^2$ for nonzero soft $p$.

\paragraph{The coefficient of the soft Maxwell response}
In the continuum rotationally invariant chart, fix a unit vector $v_0$
and put $p=tv_0$. In a proper photon graph after collapsing its
charged vertices, every internal photon belongs to a cycle. Its
momentum thus depends on an independent loop variable and is nonzero
at $t=0$ except on a set of measure zero. At either external mark
use the constant subtraction in the proof of
Proposition~\ref{prop:s6-heatlink-soft}. After division by $t$
its plane wave tends to the linear insertion
$\ii\epsilon_j v_0\cdot(y-x)$, with the sign reversed at the
opposite momentum. The same weighted-kernel estimate dominates
both divided insertions. The internal vertices are continuous at
the remaining momenta by their marked integral formulas. The chord
bounds in that proof justify the resulting limit of
$t^{-2}\Pi^{\rm int}(tv_0)$ at each fixed order. Choose each
independent integration variable to be the actual momentum on
one chord. Then the product of the chord envelopes is independent
of $t$. The remaining photon factors are uniformly bounded, and
at $t=0$ the vanishing of one of their momenta occurs on a null set.
A single charged vertex with no photon chords needs only the
weighted charged-kernel estimate. Rotation covariance,
the Ward contraction, and symmetry on the transverse three-space
give $c_j\Pi^T(v_0)$, with one finite scalar $c_j$ independent
of direction at electric order $j$. The transverse inverse response divided by $|p|^2$ therefore has
the formal expansion $1+\sum_{j\ge1}c_je^j$. Its reciprocal exists
order by order and is the corresponding Euclidean soft residue.


\begin{proposition}[The soft Maxwell coefficient along the lowest torus modes]
\label{prop:softward-torus-uv}
Use the fixed-distance Wilson heat-link operator just defined with
\(a,m>0\) and the boundary Maxwell covariance. Take the torus period
to be \(2\pi R\) in each direction, keep the dynamical harmonic mode
fixed at zero, and use the same lengths in physical units for all
\(R\ge1\). Let \(\Pi_{R,j}(p)\) be the order-\(e^j\) coefficient
of the proper photon two-point correction per unit volume with the
overall momentum delta removed. Normalize each independent momentum
sum by \((2\pi R)^{-4}\). For any nonzero lattice momenta
\(p_R\in R^{-1}\mathbb Z^4\) such that \(p_R\to0\) and
\(p_R/|p_R|\to v\), one has, on the transverse space,
\begin{equation}
 |p_R|^{-2}\Pi_{R,j}(p_R)\longrightarrow c_j\Pi^T(v).
 \label{eq:softward-torus-uv}
\end{equation}
Here \(c_j\) is the continuum Maxwell coefficient obtained by first
using continuous momenta, and \(\Pi^T(v)=I-vv^{\mathsf T}\).
In particular, the limit is valid for \(p_R=v/R\) with \(v\) a
coordinate unit vector.
\end{proposition}
\begin{proof}
First compare a complete charged vertex on the torus with the continuum
vertex at the same real photon transfers. After extracting the
transverse readout at each end, every term is a trace of free
inverses and marked links with one charged sum over \(k\). The free
symbol \(d(k)=mI+i\gamma\cdot k\,e^{-a^2 k^2}\) extends
holomorphically using \(k^2=\sum_\mu k_\mu^2\). Its derivative is
bounded on a tube of fixed, sufficiently small radius about real
momentum space. Since \(\|d(k)^{-1}\|\le m^{-1}\) for real \(k\),
one can choose a radius \(2\eta>0\) on which the inverse has norm at
most \(2/m\). Real translations by any photon transfers preserve this
tube. The Fourier multiplier of each marked link is an average of
translated polynomial Gaussians with bounded supremum on the tube.
Use one of them in \(L^1(d^4 k)\) on each shifted contour and the
others in supremum norm. This bounds the charged trace integrand
there in \(L^1\), uniformly in the photon momenta. The same argument
applies to coincident marks on a link.

For the Fourier transform
\(\widehat f(w)=\int e^{-ik\cdot w}f(k)\,d^4 k/(2\pi)^4\),
shifting \(k\) by \(-i\eta w/|w|\) gives
\(|\widehat f(w)|\le C_j e^{-\eta|w|}\).
Poisson summation therefore bounds the difference between the
normalized charged sum and its continuum integral by
\(C_j e^{-2\pi\eta R}\) for \(R\ge1\), after a change of constant.
This bound is uniform even when the photon momenta being sewn depend
on \(R\). The suppressed external and internal readouts multiply both
sides.

Now fix a proper photon graph of this order and telescope the
replacement of its torus charged vertices by their continuum versions.
The vertex suprema are uniform, including on the torus. Every
nonzero sewn momentum has \(|q|^{-2}\le R^2\), and a missing zero
mode is omitted before this comparison. Choose the loop variables
to be the photon chord momenta of a spanning tree. Their exponential
readouts have uniformly bounded normalized lattice sums, while the
remaining exponentials are at most one. For a graph with \(I\)
internal photons the error is at most
\(C_j R^{2I}e^{-2\pi\eta R}\). Division by \(|p_R|^2\) costs at
most \(R^2\). This error tends to zero also at the first Fourier
mode.

With continuum charged vertices use their two external soft factors
before dividing by \(|p_R|^2\). Both ends of every internal photon
supply \(\min(1,|q|)\), leaving an integrable exponential envelope
in the independent chord momenta. On compact sets where no internal
momentum vanishes, the divided integrands converge uniformly to the
continuum soft integrand. Indeed each subtracted soft plane wave on
a charged loop, divided by \(|p_R|\), tends to its linear insertion
along \(v\), and the weighted charged kernels dominate these insertions
and their products. The transverse projectors are continuous on
these sets. To check the omitted sets, every internal edge of a
proper graph lies on a cycle. Its momentum contains a chord variable
with coefficient \(1\) or \(-1\). With all other variables in a fixed
compact set, the normalized lattice measure where \(|q|<\epsilon\)
is at most \(C(\epsilon+R^{-1})^4\), uniformly in the shift by the
other momenta. The finite union of these regions thus contributes
arbitrarily little. The exponential chord bounds control the complement
of the compact set uniformly in \(R\). Ordinary Riemann sums on
the remaining region give the continuum integral. When there are
no internal photons only the external divided-vertex limit is
needed. Summing the finitely many proper graphs proves the claim.
\end{proof}

A separate flat-connection calculation explains the role of the
removed zero mode. Couple an external constant \(b\) to the Wilson
links and put
\(f(k)=2\log(1+m^{-2}|k|^2e^{-2a^2|k|^2})\).
The negative log charged determinant density is
\(-(2\pi R)^{-4}\sum_{k\in R^{-1}\mathbb Z^4} f(k+eb)\),
up to a constant. Its Hessian at zero is exactly
\[
 e^2(2\pi R)^2\sum_{n\ne0}
      n_\mu n_\nu\,\widehat f(2\pi R n).
\]
The zero winding has no holonomy derivative. On a sufficiently narrow
tube the logarithm has the branch tending to zero at real infinity
and has bounded translated \(L^1\) norms for imaginary shifts of
size at most \(\eta\), reducing \(\eta\) if necessary. The same
contour shift bounds the Hessian by
\(C e^2 R^2e^{-2\pi\eta R}\). Thus torus holonomy can contribute
a finite-size response even though the propagating transverse
modes have the common soft limit above. The conclusion is
order by order for the specified action and concerns its Euclidean
inverse photon response.

\begin{theorem}[Two-sided massless susceptibility at fixed real charge]
\label{thm:softward-real-charge-uv}
For the same heat-link source let
\(A_0=\int_{\mathbb R^4}\|K(z)\|\,d^4z\), where \(K\) is its
spin-weighted displacement kernel. Its Gaussian gives \(A_0=3\sqrt{\pi}/(4a)\). If \(m>A_0\), there is a
finite \(C=C(a,m)>0\), independent of \(R\ge1\), such that for
real charge with \(\delta=Ce^2<1\) the exact transverse boundary
photon covariance satisfies the quadratic-form bounds
\begin{equation}
 \frac{Q_R^{-1}}{1+\delta}
 \ \preceq\ \operatorname{Cov}_{e,R}(A^T)
 \ \preceq\ \frac{Q_R^{-1}}{1-\delta}.
 \label{eq:softward-real-charge-uv}
\end{equation}
Here \(Q_R=-\Delta\) acts on real transverse fields of zero mean with
the torus \(L^2\) metric. In particular, the susceptibility of each
unit-normalized transverse Fourier polarization grows as
\(|p|^{-2}\) between the displayed constants, even at the first
nonzero mode of arbitrarily large tori. This is a bound for the
unexpanded source.
\end{theorem}
\begin{proof}
Young's inequality and \(|W_B|=1\) give \(\|T_B\|\le A_0<m\).
The ordinary log determinant therefore has its convergent trace
series in the powers \(m^{-n}T_B^n\), with coefficients
\((-1)^{n+1}/n\). A trace of length \(n\) is a closed polygon on
the torus. For fixed lifted steps \(z_1,\ldots,z_n\), write
\(\ell=\sum_i|z_i|\) and \(w=\sum_i z_i\in2\pi R\mathbb Z^4\).
At base point \(x\), twice varying the Wilson phase in the same
real sampled direction \(b\) multiplies it by
\(-e^2 J_b(x)^2\), where \(J_b(x)=\int_{x+C}b\cdot dy\) and
\(C\) is the lifted polygon. This includes all pairs of marks on
the same or different links.

For smooth periodic \(b\) of zero mean, subtract \(b(x)\) along the
polygon. The torus translation estimate
\(\|b(\cdot+y)-b\|_2\le |y|\,\|\nabla b\|_2\) bounds the
resulting integral in \(L^2(dx)\) by \(\ell^2\|\nabla b\|_2\).
The subtracted constant contributes \(w\cdot b(x)\). If
\(w\ne0\), then \(2\pi R\le|w|\le\ell\), while
\(\|b\|_2\le R\|\nabla b\|_2\) by the first nonzero torus
eigenvalue. Consequently, for every winding,
\[
 \|J_b\|_2 \le (1+(2\pi)^{-1})\ell^2\|\nabla b\|_2 .
\]
Choose \(s>0\) so small that
\(A_s=\int e^{s|z|}\|K(z)\|\,d^4z<m\).
The cost \(\ell^4\) is bounded by \(C_s e^{s\ell}\).
For the last step the Gaussian makes the periodization of
\(e^{s|z|}\|K(z)\|\) bounded uniformly in \(R\ge1\), and the other
\(n-1\) steps integrate to \(A_s^{n-1}\).
The spin trace costs at most four times the product of the
kernel norms. The series of absolute Hessian bounds thus converges
uniformly in the real background and in the volume after
factoring out \(\|\nabla b\|_2^2\). With \(b=P_a h\), the
harmonic readout is a contraction in this norm. We have proved
\[
 |\partial_h^2\log\mathcal D_m(P_a A^T)|
 \le C e^2\langle h,Q_R h\rangle.
\]

To obtain the covariance, first evaluate the same charged determinant
on a finite Fourier projection of the Maxwell field. The Hessian
of the real negative log density \(S\) of these modes lies between
\((1-\delta) Q_R\) and \((1+\delta)Q_R\) on the projected space.
In coordinates in which the free Maxwell covariance is identity the
Brascamp--Lieb variance bound~\cite{SpectrumBrascampLieb} for a linear functional with coefficient
\(u\) gives \(\operatorname{Var}(u\cdot X)\le |u|^2/(1-\delta)\).
Integration by parts in these coordinates gives
\[
 \begin{aligned}
 \operatorname{Cov}(u\cdot X,u\cdot\nabla S)&=|u|^2,\\
 \mathbb E[(u\cdot\nabla S)^2]
 &=\mathbb E[u^{\mathsf T}S''u]\le(1+\delta)|u|^2 .
 \end{aligned}
\]
Cauchy--Schwarz gives the other side of the bound. Increasing the
Fourier projections at fixed \(R\) gives almost sure convergence in
every smooth norm of the sampled field and hence trace-norm convergence
of its charged kernel. The positive determinant is bounded above
independently of the real photon sample at fixed \(R\). Dominated
convergence for the tested Maxwell moments proves the inequalities
for the exact source. The finite projections serve only to prove
the estimate for that source.
\end{proof}

The bounds persist for limiting smeared covariances whenever a
thermodynamic limit is taken. They also specify the additional input
for a particle interpretation. If an invariant limit has a physical
transverse scalar coefficient of the form
\(G(z)=\int_{[0,\infty)}(z+t)^{-1}d\rho(t)\) with
\(\rho\ge0\) and \(G(1)<\infty\), then
\(\lim_{z\downarrow0} zG(z)=\rho(\{0\})\). Indeed
\(0\le z/(z+t)\le1/(1+t)\) for \(0<z\le1\). In that case the
mass-zero atom lies between \((1+\delta)^{-1}\) and
\((1-\delta)^{-1}\). Establishing a positive spectral representation
for the physical projection is a separate requirement. The Euclidean
susceptibility bound itself uses only the displayed heat-link action
and its positive finite-volume weight.

\begin{theorem}[Two-sided photon susceptibility for every positive charged mass]
\label{thm:softward-allmass-uv}
For each fixed \(a,m>0\), the bounds of
\eqref{eq:softward-real-charge-uv} hold for some finite
\(C(a,m)\) independent of \(R\ge1\) and all real \(e\) with
\(C(a,m)e^2<1\). The lower covariance bound
\((1+C(a,m)e^2)^{-1}Q_R^{-1}\) holds at every finite real charge.
\end{theorem}
\begin{proof}
It suffices to establish the same uniform Hessian estimate without
expanding the determinant in \(1/m\). Lift the actual smooth real
periodic connection \(B\) to the cover. There
\(R_{v,B}=(vI+T_B)^{-1}\) exists with norm at most \(v^{-1}\)
for \(v\ge m\) by skew adjointness. Write \(k_0(z)=\|K(z)\|\).
Choose \(\eta>0\) so small that
\(\int k_0(z)(e^{\eta|z|}-1)\,d^4z<m/2\). For any bounded real
\(\eta\)-Lipschitz weight \(\phi\),
\(T_\phi=e^\phi T_B e^{-\phi}\) differs from \(T_B\) by at
most this integral in operator norm, by convolution domination.
Hence \(\|(vI+T_\phi)^{-1}\|\le2/v\). The rows and columns of
\(T_\phi\) have uniformly bounded \(L^2\) norms since their
kernels are dominated by \(e^{\eta|z|}k_0(z)\). Use the exact identity
\[
 R_{v,B}-v^{-1}I=-v^{-2}T_B+v^{-2}T_B R_{v,B}T_B .
\]
After conjugation, Cauchy--Schwarz applied to the two outer kernels
bounds the continuous kernel of the right-hand side by
\(C_m v^{-2}\), uniformly in \(\phi\) and real \(B\).
For the kernel at \((x,y)\) take
\(\phi(t)=\eta\min(|t-y|,M)\), with \(M>|x-y|\). Undoing the
conjugation gives
\begin{equation}
 R_{v,B}=v^{-1}I+L_{v,B},\qquad
 \|L_{v,B}(x,y)\|\le C_m v^{-2}e^{-\eta|x-y|}.
 \label{eq:softward-magnetic-resolvent-uv}
\end{equation}
The constant may depend on \(a\). This estimate also justifies using
the cover inverse on the periodic source. The kernel is invariant
under simultaneous period translations. Sum it in its second
argument over all period translates, interpreting the identity
as the periodic delta. The exponential remainder is absolutely
summable. Applying \(vI+T_B\) on the torus unfolds its periodized
convolution to the cover identity. The interchange follows from
the Gaussian and exponential bounds. Thus the periodization is a
right inverse, and skew adjointness on the torus identifies it with
the torus inverse.

At fixed \(R\), the mass integral of the trace of
\(R_{v,B}^{\rm tor}-v^{-1}I\) from \(m\) to infinity equals
\(-\log\mathcal D_m(B)\). This follows by differentiating the
ordinary Fredholm determinant in \(v\) and using its value one
at infinite mass. To differentiate the diagonal trace twice
along a real periodic direction \(b=P_a h\), consider one lifted
winding \(w\in 2\pi R\mathbb Z^4\) based at \(x\). Gauge
covariance on the cover gives the exact identity
\[
 \begin{aligned}
 R_{v,B+\lambda b}(x,x+w)
 &=e^{ie\lambda w\cdot b(x)}\\[-1mm]
 &\quad\times R_{v,B+\lambda(b-b(x))}(x,x+w).
 \end{aligned}
\]
It is an identity of the continuous remainders for \(w=0\),
since the unchanged identity delta has zero variation.
For each fixed \(x\), write \(T_1,T_2\) for the first and second
Wilson derivatives in the residual direction \(b(\cdot)-b(x)\)
and \(R_v=R_{v,B}\). The first inverse derivative is
\(-R_vT_1 R_v\), and the second is
\(2R_v T_1 R_v T_1 R_v-R_v T_2 R_v\). Differentiating the
displayed phase also gives a mixed term with one winding factor
and the first inverse derivative, and a term with two winding
factors and the undifferentiated inverse. This restores the
dependence on every periodic holonomy.

For a fixed sequence of displacements in any such kernel product,
let \(\ell\) be their total length, including the inverse-kernel
steps, so that \(|w|\le\ell\). A residual marked line integral
has \(L^2(dx)\) norm at most \(\ell^2\|\nabla b\|_2\) by the
translation estimate in the preceding proof. A winding factor
has norm at most \(\ell^2\|\nabla b\|_2/(2\pi)\) by the same
Poincar\'{e} bound. The two marks, including the two on one link,
thus cost at most \(C e^2\ell^4\|\nabla b\|_2^2\) in
\(L^1(dx)\). All background phases have unit modulus.
Each inverse in a marked string is dominated as a kernel measure
by \(C_m v^{-1}(\delta_0+e^{-\eta|z|}\,d^4 z)\). Each marked
link supplies \(k_0(z)\,d^4 z\) with the marks just estimated.
Absorb \(\ell^4\) into \(C_\eta e^{\eta\ell/2}\). Periodize one
Gaussian marked kernel to impose closure, using its bounded
periodized supremum, and integrate the other steps in total
variation. The finite strings have at least two inverses and so
cost at most \(C_m v^{-2}\) after the common
\(e^2\|\nabla b\|_2^2\) is removed. For the term with two winding
marks there is no marked link. Only \(w\ne0\) contributes, and
\eqref{eq:softward-magnetic-resolvent-uv} bounds the remaining sum by
\(C_m v^{-2}\sum_{w\ne0}|w|^4 e^{-\eta|w|}\), uniformly in \(R\ge1\).

All these estimates hold at every smooth periodic real background,
also along the real line of variation. For instance first
differentiate the mass integral on a compact interval in \(v\),
then apply the integrable bound to the centered second difference
by the one-variable integral remainder formula. The same bound
permits its limit at zero increment. This proves
\( |\partial_h^2\log\mathcal D_m(P_a A^T)|
 \le C(a,m)e^2\|\nabla h\|_2^2\).
The finite-dimensional variance inequalities and convergence at
fixed \(R\) in the preceding proof now give the assertion. At
arbitrary real charge the upper Hessian bound
\(S''\preceq(1+C(a,m)e^2)Q_R\) still holds. The fixed-volume
determinant bound gives Gaussian tails and permits the same score
integration by parts. Its resulting lower covariance estimate
requires no convexity assumption.
\end{proof}

\paragraph{Momentum dependence of the exact susceptibility bound.}
The harmonic readout also fixes the large-momentum behavior of the
unsampled boundary covariance. In the all-mass proof the estimated
variation is \(b=P_a h\). Before bounding \(P_a\) by its operator
norm, the Hessian estimate reads
\[
 |\partial_h^2\log\mathcal D_m(P_a A^T)|
 \le C(a,m)e^2\langle h,Q_R P_a^2h\rangle .
\]
Set \(M_\pm=Q_R(I\pm C(a,m)e^2 P_a^2)\).
On every finite set of transverse Maxwell modes, integration by
parts gives
\(\operatorname{Cov}(h\cdot X,u\cdot\nabla S)=h\cdot u\) and
\(\mathbb E[(u\cdot\nabla S)^2]\le u^{\mathsf T}M_+u\).
Choose \(u=M_+^{-1}h\) in Cauchy--Schwarz. For the reverse comparison
when \(C(a,m)e^2<1\), use \(S''\succeq M_->0\) in the
Brascamp--Lieb inequality. This proof works even though the
background-dependent Hessian need not commute with \(Q_R\) or
\(P_a\). Passing to the exact finite-volume measure gives
\begin{equation}
 M_+^{-1}\preceq\operatorname{Cov}_{e,R}(A^T)\preceq M_-^{-1},
 \label{eq:softward-anisotropic-uv}
\end{equation}
where the lower bound holds at every real charge.

In particular, if a rotation-invariant continuum limit at weak
charge has a transverse scalar coefficient \(G(z)\), \(z=|p|^2\),
its large-\(z\) behavior is
\(G(z)=z^{-1}+O(z^{-1}e^{-2a\sqrt z})\), in the sense of
the covariance bounds and pointwise for a continuous representative.
An unsubtracted positive representation
\(G(z)=\int (z+s)^{-1}d\rho(s)\) would then require
\(\rho([0,\infty))=\lim_{z\to\infty}zG(z)=1\).
By monotone convergence again,
\[
 \int s\,d\rho(s)
 =\lim_{z\to\infty}\int\frac{zs}{z+s}\,d\rho(s)
 =\lim_{z\to\infty}z^2[z^{-1}-G(z)]=0 .
\]
Consequently such a positive limit must have
\(\rho=\delta_0\) and exactly the free transverse coefficient.
This implication follows directly from the exact covariance of
the unsampled photon in the fixed-distance model.

\paragraph{A nonzero finite Maxwell correction from a four-segment loop}
For the same heat-link action \eqref{eq:s6-D}, one can compute the
first nonzero term of the one-loop soft coefficient at large $m$.
Write $\sigma^2=2a^2$, and denote by $g_{4a^2}(0)=(16\pi a^2)^{-2}$
the density at the origin of the sum of four independent hops.
For $p=t e_1$ and polarization $e_2$, the order-$e^2$ part obeys
\begin{equation}
 \lim_{t\to0}\frac{\Pi^{(2)}_{22}(t e_1)}{t^2}
 =e^2\left(\frac{3g_{4a^2}(0)}{4m^4}+O_a(m^{-6})\right)
 \qquad (m\to\infty).
 \label{eq:s6-square-response}
\end{equation}
Here $\Pi^{(2)}$ includes both source derivatives of the negative
logarithm of the charged determinant. The coefficient is strictly
positive for sufficiently large $m$.

For a closed polygon $C$, the photon couples to
$J_\mu(p)=\oint_C e^{\ii p\cdot x}\,\dd x_\mu$. The integral
of $\dd x_2$ around $C$ vanishes, so $J_2(t e_1)/t$ tends
to $\ii A_{12}$, where $A_{12}=\oint_C x_1\,\dd x_2$ is the
signed projected area. Odd spin traces vanish by chirality.
On the cover a two-segment loop of zero net displacement retraces its
first segment, so its photon derivatives vanish even though its unmarked
vacuum trace can survive. At four
segments condition independent Gaussian steps $z_1,\ldots,z_4$ on
$\sum z_i=0$. Set
\[
 U=(z_1+z_2)/\sigma,\quad V=(z_2+z_3)/\sigma,\quad
 W=(z_1-z_2+z_3-z_4)/(2\sigma).
\]
These are independent standard Gaussian four-vectors: their coefficient
vectors on the closure subspace are orthonormal. The signed area is
$A_{12}=\sigma^2(U_1 V_2-U_2 V_1)/2$.
With $r=|U|^2$, $s=|V|^2$, $c=U\cdot V$ and
$T=\operatorname{tr}\prod_{i=1}^4(\gamma\cdot z_i)$, the
four-gamma trace followed by the $W$ average gives
\[
 4\,\mathbb E[T\mid U,V]/\sigma^4
   = r^2+s^2-6rs+4r+4s+4c^2+24=:F .
\]
Indeed the dot-product trace uses the three ordered pairings with
signs $+,-,+$, and
$z_1/\sigma=(U-V+W)/2$,
$z_2/\sigma=(U+V-W)/2$,
$z_3/\sigma=(-U+V+W)/2$, with $z_4$ fixed by closure.
Rotational symmetry gives each of the six coordinate planes an
equal share of $rs-c^2=\sum_{i<j}(U_iV_j-U_jV_i)^2$. Hence
\[
 \frac{\mathbb E[T A_{12}^2]}{\sigma^8}
     =\frac{\mathbb E[F(rs-c^2)]}{96}=-3 .
\]
For the last equality put $b=c^2/(rs)$.
The radii $r,s$ and $b$ are independent,
$\mathbb E r^j=\mathbb E s^j=2^j(2)_j$ with $(2)_j$ the rising
factorial, $\mathbb E b=1/4$, and $\mathbb E b^2=1/8$.

The four kernels supply $\sigma^{-8}$ and the closure integral
supplies $g_{4a^2}(0)$. The trace-log factor is $+1/(4m^4)$,
and the two Wilson derivatives supply $-e^2|J_2(p)|^2$.
The external normal readouts approach one. To bound the remainder, put $M_a=\int\|K(z)\|\dd z$ and expand
the large-mass logarithm for $m>M_a$. For a closed $n$-step trace
let $s=\sum_i|z_i|$. Subtracting the constant from the plane wave
round the closed polygon gives $|J_2(t e_1)|\le |t|s^2$. Bound one
closing kernel by its supremum; closure gives
$s\le2\sum_{i<n}|z_i|$. The inequality
$(\sum_{i<n}|z_i|)^4\le(n-1)^3\sum_{i<n}|z_i|^4$
bounds the absolute length moment by
$C_a n^4 M_a^{n-2}$. After the factor $1/(nm^n)$ the sum of the
even lengths $n\ge6$ is $O_a(m^{-6})$ for $m\ge2M_a$. This
also justifies differentiation under the loop integrals and the
directional soft limit. The large-mass expansion is used only to
prove the nonzero coefficient in \eqref{eq:s6-square-response}.

\begin{proposition}[Strict sign of the soft Maxwell coefficient at every mass]
\label{prop:closure-heat-maxwell-sign-uv}
For the fixed-width Wilson heat link the coefficient $c_2$ of $e^2$
in the continuum transverse inverse response divided by $|p|^2$ is
strictly positive for every $m,a>0$. In this convention the charged
action contributes the negative logarithm of its determinant.
\end{proposition}
\begin{proof}
In the straight Wilson chord, set $x=a x'$ and $B'(x')=aB(ax')$.
The two-mark kernel per unit volume scales as $a^{-2}$ times the
unit-width kernel at momentum $ap$ and mass $\mu=ma$. Thus $c_2$
depends only on $\mu$. Work at unit width, put $h(k)=k e^{-|k|^2}$
and $D=\mu^2+|h|^2$. Choose $p=t e_1$ and photon polarization $e_2$.
Write $h_\pm=h(k\pm p/2)$ and $D_\pm=\mu^2+|h_\pm|^2$. Differentiation
of the chord phase gives the first and second mark vectors
\[
 \begin{aligned}
 V&=\int_{-1/2}^{1/2}\partial_2h(k+s p)\,\mathrm ds,\\
 W&=\int_0^1\!\int_0^1
       \partial_2^2h(k+(s-v)p)\,\mathrm ds\,\mathrm dv .
 \end{aligned}
\]
The first formula uses the midpoint of the two fermion momenta.
The two marks in the direct contact act on a single chord. The
Euclidean Clifford trace identities give the bubble plus its contact,
with the factor $4e^2$ and the two external readouts removed, as
\[
 F(k,t)=
 \frac{2(h_-\cdot V)(h_+\cdot V)
       -(\mu^2+h_-\cdot h_+)|V|^2}{D_-D_+}
       -\frac{h\cdot W}{D}.
\]
At $t=0$ its integral is zero by the constant-shift identity. The
external harmonic factors approach one, so $c_2=4\int[t^2]F(k,t)
\,\mathrm d^4k/(2\pi)^4$. Here $[t^2]$ denotes one half of the
second derivative at zero. The mass bound and the Gaussian marks
dominate the derivatives in a neighborhood of zero.

Set $u=|k|^2$, $y=u e^{-2u}$, $z=y/(\mu^2+y)$ and
$P(u)=2u^3-7u^2+2u+2$. Normalized angular averaging on $S^3$ gives
\[
 u\langle[t^2]F\rangle
 =\frac{z^2P(u)}{6u}
     +\frac{\mathrm d}{\mathrm du}
           \big(f_1z+f_2z^2+f_3z^3\big),
\]
where
\[
 \begin{aligned}
 f_1&=u(2u^2-9u+3)/36,\qquad f_3=(2u-1)^3/36,\\
 f_2&=-(2u-1)(2u^2-6u+5)/24 .
 \end{aligned}
\]
This reduction can be checked by expanding $h_\pm=h\pm t\partial_1h/2+
t^2\partial_1^2h/8+O(t^3)$ in $F$. The quadratic terms in $V,W$ are
$\partial_1^2\partial_2 h/24$ and
$\partial_1^2\partial_2^2h/12$, respectively. Use
$\langle k_1^{2i}k_2^{2j}\rangle
=u^{i+j}(1/2)_i(1/2)_j/(2)_{i+j}$ and
$z'=z(1-z)(1/u-2)$; $(b)_j$ is the rising factorial.
The boundary terms vanish since $z=O(u)$ at zero and decays
exponentially at infinity for every $\mu>0$. With the radial
four-dimensional measure this leaves
\begin{equation}
 c_2=\frac{1}{24\pi^2}\int_0^\infty z^2 P(u)\,
               \frac{\mathrm du}{u}.
 \label{eq:closure-heat-maxwell-sign-uv}
\end{equation}

The signed polynomial in this integral changes sign, but the same
value of $y$ occurs on two radii. For $0<y<\exp(-1)/2$ let
$0<l<1/2<b$ solve $l e^{-2l}=b e^{-2b}=y$.
After changing variables on both branches, the remaining signed
factor multiplying the positive $z^2\,\mathrm dy/y$ is
\[
 \frac{P(l)}{1-2l}+\frac{P(b)}{2b-1}.
\]
The first term exceeds $2+6l$, since
$P(l)-(2+6l)(1-2l)=l^2(2l+5)>0$.
The second is at least $-2$ except possibly on $3/2<b<2$:
$P(b)+2(2b-1)=b(2b-3)(b-2)$.
On that interval its deficit below $-2$ is at most $3/32$,
using $b/(2b-1)\le3/4$ and $(2b-3)(2-b)\le1/8$.
There $l>y\ge2\exp(-4)$ and $12\exp(-4)>3/32$. The sum is therefore
strictly positive at every interior $y$, which proves the claim.
\end{proof}

The branch calculation also determines the dependence on the mass.
Put $x=(ma)^2$. The coefficient is an integral of $(x+y)^{-2}$
against a positive measure on $0<y<\exp(-1)/2$, with
density equal to $y/(24\pi^2)$ times the
positive two-branch factor in the proof. Differentiation on compact
subsets of $x>0$ shows that $(-1)^n\partial_x^n c_2>0$ for every
integer $n\ge0$. Varying the massive representation changes the size
of the magnetic response but cannot tune it to zero.

\paragraph{A physical reconstruction test for this same charged field}
There is a further test available directly from the free inverse in
\eqref{eq:s6-d}. Write $z=k_E^2$ and keep the fixed positive heat
lapse $a^2$. The scalar coefficient of $-\ii\gamma\cdot k$ in the
open propagator is
\[
 F_V(z)=\frac{e^{-a^2 z}}{m^2+z e^{-2a^2z}} .
\]
The denominator vanishes at
\begin{equation}
 z_n=-\frac{W_n(2a^2m^2)}{2a^2},\qquad n\in\mathbb Z,
 \label{eq:s6-freepole}
\end{equation}
where $W_n$ are the branches of the Lambert function
\cite{DLMF4W}. Indeed $w=-2a^2z$ solves $we^w=2a^2m^2$.
The roots are simple. The only critical value of $we^w$ is
$-\exp(-1)$, whereas $2a^2m^2>0$. There is one negative real root $-s_*$; its vector
residue is $e^{-a^2s_*}/(1+2a^2s_*)>0$. Every other root is
nonreal, and the numerator has no zero that could remove it.
A standard positive spinor spectral representation would give a
vector coefficient analytic off the negative real $z$ axis, even
if local polynomial contact terms were added. There is also a real-axis check. The coefficient obeys
$z F_V(z)\to0$ as $z\to+\infty$. For an unsubtracted nonzero
positive vector measure $\rho_1$, monotone convergence would instead give
$\lim_{z\to+\infty}\int z\,(z+s)^{-1}\rho_1(\dd s)
=\rho_1([0,\infty))>0$. Consequently the
unmodified retained open field of this chart fails that spectral
test already at $e=0$. Gauge-fixed or auxiliary variables need
not have such a representation; an identification with a physical
positive spinor would require deriving its retained operator and
source anew. The nonreal singularities belong to the analytic continuation
of this Euclidean form factor. The vectorial datum
$\ii\sigma_{s\epsilon}\varepsilon\hat p_{0,\bb}$ specifies the
approach to the real rim on the fixed sector plane. For Q1 the
physical charged source must yield its own spectral coefficient.

\subsubsection{Homogeneous transfer along the proper radial distance}
The preceding graph estimates ask how the charged smoothing scale varies
with the photon readout distance. There is a geometric way to identify the
two parameters on the curvature-minimizing normal surface. In its radial
metric the coordinate $r$ is proper distance from the tip. Suppose that the
\emph{free endpoint displacement} has the same local tangential kinetic
form on each successive unit of this distance. In fixed reference units
write that form as $\kappa\int_{\mathbb R^4}|\nabla_z h|^2\,\mathrm d^4z$,
with $\kappa>0$ the heat calibration. Its nonnegative operator is $-\kappa\Delta_z$.
Equivalently, on a subdivision with positive proper lengths $s_l$,
the normalized transition density of a displacement $z_l$ is
$(4\pi \kappa s_l)^{-2}\exp[-|z_l|^2/(4\kappa s_l)]$.
The positive quantity $|z_l|^2/(4\kappa\,s_l)$ is the kinetic action of the
straight segment joining its endpoints over that length. Integrating all increments
with the endpoint phase $e^{ik\cdot\sum_l z_l}$ multiplies their
characteristic functions to give $e^{-\kappa(\sum_l s_l)|k|^2}$ exactly at
every subdivision. This calculation uses all displacement momenta
and requires no subdivision limit. Transfer along length $r$ starting
at zero displacement thus gives
\[
 \int e^{i k\cdot z}g_{\kappa\,r}(z)\,\mathrm d^4z=e^{-\kappa\,r|k|^2}.
\]
Indeed the Fourier modes satisfy $\partial_r\widehat g=-\kappa|k|^2\widehat g$
with $\widehat g(0,k)=1$. Equivalently, for an endpoint law already known
to have a heat lapse $\tau(r)$, homogeneous composition of successive
displacements requires $\tau(r+s)=\tau(r)+\tau(s)$ for $r+s\le1$.
Positivity and the calibration $\tau(1)=\kappa$ give $\tau(r)=\kappa\,r$.
Thus the charged heat exponent is one for this kinetic form. In dimensional
variables $\kappa$ has units of displacement squared per unit proper radial
length. The unit normalization used below corresponds to $\kappa=1$
in the chosen reference units.

This composition takes place before the Wilson chord is attached to the
endpoint. The photon readout still follows from its own harmonic boundary
problem and gives $e^{-r|p|}$ at each transverse end. Curvature generally
makes the holonomy of two consecutive chords differ from the holonomy of
the single chord between their outer endpoints. No charged semigroup
identity is needed. All photon contacts still arise by differentiating
the transported endpoint in the specified curvature contraction.

\begin{corollary}[Bare electromagnetic graphs on the free-slope minimizing surface]
\label{cor:geo-free-slope-uv}
Use the positive rotational curvature block in \eqref{eq:g5N-normal-geometry} with the
proper interval and outer circumference fixed as there, and leave the
outer slope free. Prepare its minimizing profile $f_*$ in the sense
of that variational problem, then use the lowest radial state in an
angular sector $j\ne0$. Couple one draw of $r$ from this state to the
charged heat transfer of lapse $\kappa\,r$ and to the harmonic photon readout
at distance $r$ throughout the represented curvature source. On a fixed
flat Euclidean four-torus with positive charged mass and the Maxwell
zero mode fixed, the determinant-weighted free normalization is finite.
Every fixed bare electromagnetic graph with fixed Fourier probes is
absolutely ultraviolet finite at every finite order in the electric
charge. These sums use the entire reciprocal lattice.
\end{corollary}
\begin{proof}
Write $C$ for the positive constant in the free-slope minimizer and let
$q=|j|$. The exact tip coefficient follows from its quadrature. Set
$I(C)=\int_0^\infty(z^3/3+C)^{-1/2}\,\mathrm dz=1$ at the selected
$C$. Scaling gives $I(C)\propto C^{-1/6}$ as a function of $C$.
In the minimizing formula $g=-w'=\sqrt{w^3/3+C}$ and
$f_*(r)=c\,g(r)\int_0^r g(s)^{-2}\,\mathrm ds$.
Changing variable from $s$ to $w$ gives
$\int_0^1g^{-2}\,\mathrm ds=-2I'(C)=1/(3C)$.
Since $g(1)=\sqrt C$ and $f_*(1)=1$, one has $c=3\sqrt C$.
For large $w$, the quadrature for $r$ gives
$r=2\sqrt3\,w^{-1/2}(1+O(w^{-3}))$. Substitution in $g$ and then
in $f_*$ yields
\[
 \begin{aligned}
 g(r)&=24r^{-3}(1+O(r^6)),\\
 f_*(r)&=c_0 r^4(1+O(r^6)),\qquad c_0=\sqrt C/56 .
 \end{aligned}
\]

The minimum curvature action is $\mathcal E[f_*]=2\sqrt C$, so the
tip coefficient is also $c_0=\mathcal E[f_*]/112$.

Put $f=f_*$ and denote by $H_j$ the radial
Hamiltonian in $L^2(\mathrm dr)$ and by $E_0$ its ground energy.
Use the positive comparison functions
\[
 v_b=\sqrt f\exp\!\left[-q\int_r^1\frac{\mathrm ds}{f(s)}
                         -b\int_0^r f(s)\,\mathrm ds\right].
\]
Direct differentiation gives
$(H_j-E_0)v_b/v_b=2bq-E_0+2bf'-b^2 f^2$.
Choosing $2bq$ on opposite sides of $E_0$ gives a subsolution and a
supersolution near the tip. Both have zero Friedrichs trace.
The potential exceeds $E_0$ on a sufficiently short interval because
$j^2/f^2$ dominates the radial term $(\sqrt f)''/\sqrt f=O(r^{-2})$.
Matching positive multiples at the outer end of that interval
and testing against the positive part compares the ground state with
the two barriers. Since $1/f-c_0^{-1}r^{-4}=O(r^2)$, its normalized
density $|u_j|^2$ is bounded above and below by positive multiples of
\[
 r^4\exp[-A_*r^{-3}],\qquad A_*=2q/(3c_0).
\]
At $\sigma^2=\kappa\,r$ the free log determinant is of order
$r^{-2}(1+|\log r|)^3$ at the tip by Proposition~\ref{prop:dualrev-uvheat-depth}.
It is dominated by the radial exponent, even after multiplication of
the full density by any fixed inverse power of $r$. At fixed $r$ one
Fourier derivative of the displacement characteristic, with further
derivatives for multiphoton contacts, gives a summable translated
vertex on each marked charged loop. Each sewn transverse photon
uses the harmonic factor twice. The graph bounds in Theorem~\ref{thm:s6-heatlink-uv} and Proposition~\ref{prop:dualrev-uvheat-depth} then
cost only fixed inverse powers of $r$. Integration against the
same determinant-weighted radial state proves the assertion, including
simultaneous hard regions and arbitrarily short readouts.
\end{proof}

The unexpanded real-charge source on this minimizing surface also
has a sufficient finite-volume test. Its determinant upper exponent
has coefficient $B_*=V/(32\pi^2m^2\kappa^3)$ on $r^{-3}$.
If $B_*<A_*$ the strict radial margin integrates every finite
charge derivative. At $B_*=A_*$ the common density bound is $Cr^4$.
The Maxwell estimate in Proposition~\ref{prop:uvheat-critical-charge} uses the normal geometry only
through this density bound and the lapse proportional to $r$.
It therefore gives all finite charge derivatives at equality as well.
All coefficients here come from the positive normal curvature action,
its boundary preparation and the specified charged and photon
representations. A different displacement generator or representation
requires its own transfer and determinant calculation.

\subsection{A common geometric readout including its charged vacuum weight}
\label{subsec:s6-depth}
The two smoothing scales can shrink together. Sample a single
radial position $r$ for the whole source, so all charged and
photon lines share the same geometry. The charged displacement has heat
lapse $r^\alpha$ and smoothing scale $r^{\alpha/2}$. The harmonic
photon readout has height $r$. Their correlation matters when
several independently hard loops probe the same short depth.

Use fixed reference units on one locked sector plane with a
finite Euclidean carrier torus and its photon zero mode fixed.
The geometry law comes from the normal surface
$\dd r^2+f(r)^2\dd\theta^2$, $0<r<1$. The angle labels an internal
$U(1)$ representation, with noncompact geometric-time directions.

Take $f(r)=r^d$. Two instructive choices of geometry and readout are
$(d,\alpha)=(4,1)$ or $(6,2)$. At the sampled position solve the
following source constraints on the lift of the carrier and on the
harmonic half-space, respectively:
\begin{align*}
 \partial_\tau g=\Delta_z g,\quad g(0,\cdot)=\delta_0,
 &\qquad K_r(z)=-\gamma\cdot\partial_z g(r^\alpha,z),\\
 (\partial_y^2+\Delta_x)b^T=0,\quad b^T(0)=A^T,
 &\qquad b=A^\parallel+b^T(r).
\end{align*}
Use the bounded harmonic solution. Write $\sigma=r^{\alpha/2}$ for the heat smoothing scale. The
endpoint density is
$(4\pi r^\alpha)^{-2}\exp[-|z|^2/(4r^\alpha)]$.
The oriented charged arrow in the Gaussian model uses the straight
Wilson chord to that endpoint in the connection $b$. This chord
rule specifies how the normal scalar diffusion couples to charge.
It is distinct from inserting a Wilson functional on the entire
Brownian path. On a torus, integrate over lifts $z$ and transport
on the projected chord, retaining all windings. The heat kernel
and its derivative are fixed coefficients of the endpoint representation. The harmonic
constraint specifies a deterministic readout, with conditional law
a point mass at the bounded solution. Longitudinal gauge
transformations act through the endpoint phases of the Wilson line.
The dynamical Gaussian is evaluated on the transverse section.

\begin{corollary}[All-orders UV bound with a single bare vacuum normalization]
\label{cor:s6-depth}
For either choice of $(d,\alpha)$, let the unweighted radial measure be the
normalized ground-state probability of a fixed angular mode $j\ne0$ of
the scalar kinetic operator on this surface, with Friedrichs boundary
conditions. For the massive Gaussian charged-hop action and transverse
Maxwell measure, every coefficient at every finite electromagnetic order
with fixed Fourier probes has graphwise absolutely convergent UV momentum
sums after the common geometric integration, with the charged
vacuum weight included in the geometry law and one overall free
normalization.
\end{corollary}
\begin{proof}
In the unitary radial variable $u=r^{d/2}\psi$ the kinetic operator is
\[
 H_{d,j}=-\partial_r^2+
       \tfrac d2(\tfrac d2-1)r^{-2}+j^2r^{-2d}.
\]
Its positive ground state exists by the form argument in
Proposition~\ref{prop:g5N-fibre}. For $0<\beta<|j|$ put
$v=r^{d/2}\exp[-\beta r^{1-d}/(d-1)]$.
Direct differentiation gives
$(H_{d,j}-E)v=[(j^2-\beta^2)r^{-2d}-E]v$.
For sufficiently small $r$ this is positive. Comparison with a
constant multiple of $v$ on that interval as in the cited proposition
bounds the ground density by
$C r^d\exp[-2\beta r^{1-d}/(d-1)]$.
The $d=4$ chart minimizes the curvature action under the stated
endpoint conditions of Section~\ref{subsec:g5N-normal}. The $d=6$
chart has finite action $\int_0^1 900 r^2\,\dd r=300$ and has its
own chosen boundary profile. Its role is to allow the Gaussian
width and the harmonic height to coincide.

At zero gauge field the free Dirac matrix at carrier mode $k$ is
$m+\ii\gamma\cdot k\,e^{-\sigma^2|k|^2}$, $m>0$.
The ordinary determinant relative to the mass operator in
\eqref{eq:s6-heatlink-source} gives the positive free factor
\[
 Z_0(r)=\prod_{k\in\mathcal L^*}
    \big(1+m^{-2}|k|^2 e^{-2\sigma^2|k|^2}\big)^2 .
\]
Here $\mathcal L^*$ denotes the reciprocal momentum lattice with
$e^{\mathrm i k\cdot w}=1$ on every period $w$. For each $r>0$ the
logarithm converges absolutely. The free factor weights the common geometry. For $0<\delta<1$, use
$\log(1+x)\le C_\delta x^{\delta/2}$ and a Gaussian lattice
sum to get
\[
 0\le\log Z_0(r)\le C_{V,m,\delta}
            \sigma^{-4-\delta}
        = C_{V,m,\delta}r^{-\alpha(4+\delta)/2}.
\]
For both geometries the exponent $\alpha(4+\delta)/2$ is strictly
less than $d-1$. Denote the ground-state law in the nonzero angular sector $j$ by $\mu$.
For each fixed $N\ge0$, $Z_0(r)r^{-N}$ is integrable against $\mu$.
In particular $\int Z_0(r)\mu(\dd r)$ is finite and positive.
If $G_r$ is a graph coefficient computed at fixed depth after
factoring out its conditional free determinant, the shared-depth
coefficient uses
\begin{equation}
 \int G_r\,\nu_0(\dd r),\qquad
 \nu_0(\dd r)=
 \frac{Z_0(r)\mu(\dd r)}{\int Z_0(u)\mu(\dd u)} .
 \label{eq:s6-tilted-depth-law}
\end{equation}
The product of separate depth averages would assign independent
geometry to each factor. Equation~\eqref{eq:s6-tilted-depth-law}
uses the same $r$ in every factor of $G_r$.

For completeness let $h$ count closed charged loops, $I$ count
sewn transverse photon edges and $s\ge1$ count photons at a
charged vertex. Differentiating the chord in Fourier space gives
an average over $[0,1]^s$ of $s$ momentum derivatives of
$\gamma\cdot k\,e^{-\sigma^2|k|^2}$ at shifted momenta.
Scaling the Schwartz function gives an operator norm at most
$C_s\sigma^{s-1}$. Its summed norm over one charged-loop anchor
is at most $C_{s,V}\sigma^{s-5}$, uniformly in every real
translation, by lattice cube comparison. The other charged
propagators cost at most $m^{-1}$. A photon edge costs at most
$C_V r^{-2}$ from the sum of
$e^{-2r|p|}/|p|^2$ over $p\ne0$. Dropping routing constraints on
the photon edges after bounding the charged sums gives a bound
$C_G\sigma^{-4h}r^{-2I}$ at each fixed graph. Free open
propagators tested between fixed normalized Fourier modes obey the
same mass bound, as do pairings with fixed $L^2$ test vectors, with
their norms included. This is an operator estimate; the identity
part of the inverse has a delta-function kernel and does not give
a pointwise bound at configuration-space endpoints.
Multiplying the fixed-graph bound by $Z_0(r)$ and integrating
over the shared position gives a finite absolute majorant for the
joint sum and geometric integral by Tonelli's theorem. The same
bound controls the sum of winding-resolved coefficients. A
winding-marker derivative on a marked link adds another momentum
derivative, hence another factor of $\sigma$ in its bounds. All
positive-order derivatives of the free inverse are uniformly
bounded for $0<\sigma\le1$, by differentiating the inverse and
using $\|d^{-1}\|\le m^{-1}$. Thus the constants for the
winding sum are bounded by a multiple of
$\sigma^{-4h}r^{-2I}$ as before.
Taking the formal logarithm or dividing normalized perturbative
sources at a fixed order then uses finitely many products of
finite coefficients.
\end{proof}

The two examples illustrate a sharp condition for this family. In what
follows \(V\) is the Lebesgue volume of a fixed flat four-torus,
\(\mathcal L\) is its period lattice and \(\mathcal L^*\) its momentum
lattice, of covolume \((2\pi)^4/V\). The estimates are on one locked
Euclidean copy in fixed reference units. Keep \(j\ne0\), \(m>0\),
\(d>1\), \(\alpha>0\), and put
\[
 p=d-1,\qquad A_j=2|j|/p,\qquad
 B_V=V/(32\pi^2m^2).
\]
The radial spectral estimate is valid for \(d>1\). If a finite
normal curvature-square action is also required, the same family
has \(\int_0^1|(r^d)''|^2/r^d\,\dd r<\infty\) for \(d>3\).

\begin{proposition}[Free vacuum and simultaneous short durations]
\label{prop:dualrev-uvheat-depth}
The common free normalization in \eqref{eq:s6-tilted-depth-law}
is finite exactly when \(p>2\alpha\). In this range every fixed
bare graph of the specified electromagnetic action at fixed external
Fourier modes is absolutely
summable on the full momentum lattice. More quantitatively, restrict
the common charged lapse by \(r^\alpha<\eta\le1\), the Schwinger
times of any subset \(E\) of the sewn photons by
\(0<s_i<\delta_i\le1\), and the largest internal momentum norm by
\(\max|k_{\rm int}|>R\), \(R\ge0\). The unnormalized absolute
contribution on this region satisfies
\begin{equation}
 \mathcal A^{\rm abs}_G(\eta,\boldsymbol\delta;R)
 \le C(1+R)^{-M}e^{-c\eta^{-p/\alpha}}
       \prod_{i\in E}\delta_i
 \label{eq:dualrev-uvheat-tail}
\end{equation}
for each fixed \(M>0\) and \(0<c<A_j\). The constant may depend on
the graph, volume, probes, \(c,M\) and the action data. The same
estimate with the momentum restriction omitted holds without the
factor \((1+R)^{-M}\).
\end{proposition}
\begin{proof}
First compare the boundary weight with the actual charged
determinant. Write \(q=|j|\), \(E_0\) for the radial ground
energy, and
\[
 w_\kappa(r)=r^{d/2}
     \exp[-q r^{-p}/p-\kappa r^{d+1}].
\]
Direct differentiation gives
\[
 \begin{aligned}
 (H_{d,j}-E_0)w_\kappa/w_\kappa
   ={}&2q\kappa(d+1)-E_0\\
      &+2d\kappa(d+1)r^{d-1}
                 -\kappa^2(d+1)^2r^{2d}.
 \end{aligned}
\]
Choose \(\kappa\) on either side of \(E_0/[2q(d+1)]\).
The smaller value gives a subsolution and the larger value a
supersolution near zero. Both have zero Friedrichs trace there.
On a sufficiently short interval the potential minus \(E_0\)
is positive. Match positive multiples of the comparison functions
to the ground state at the other endpoint, and test the difference
with its positive part. This bounds the ground state on both sides
and gives the exact exponential coefficient, including its prefactor
up to constants:
\begin{equation}
 c_0 r^d e^{-A_j r^{-p}}\le |u(r)|^2
              \le C_0 r^d e^{-A_j r^{-p}}
 \qquad (0<r<r_0).
 \label{eq:dualrev-uvheat-ground}
\end{equation}

Set \(L=\log[1/(2m^2\sigma^2)]\), \(\sigma=r^{\alpha/2}\).
Then the full product over \(\mathcal L^*\) obeys
\[
 \log Z_0(r)\sim\frac{V}{192\pi^2}\sigma^{-4}L^3
       \qquad (r\downarrow0).
\]
Indeed, under \(y=\sqrt2\sigma k/\sqrt L\) the one-mode
logarithm before spin multiplicity, divided by \(L\), is
\(L^{-1}\log(1+L|y|^2 e^{L(1-|y|^2)})\).
It tends uniformly to \((1-|y|^2)_+\) on compact annuli
away from zero. It is uniformly bounded on \(|y|\le2\) and is
at most \(|y|^2 e^{1-|y|^2}\) outside that ball for \(L\ge1\).
Lattice-cell comparison controls the tail and the omitted ball at
zero. The scaled Riemann sum thus has integral
\(\int (1-|y|^2)_+\,\dd^4 y=\pi^2/6\), times
\(V/(2\pi)^4\). Restoring the mesh and the spin multiplicity
two gives the asserted coefficient. Combining with
\eqref{eq:dualrev-uvheat-ground} proves the criterion. At
\(p=2\alpha\) the cubic logarithm makes the free normalization
diverge. For \(p>2\alpha\) the combined density still obeys
a bound \(C e^{-c r^{-p}}\) after any fixed inverse power of
\(r\) is included, for every \(c<A_j\).

For each restricted photon use
\( |k|^{-2}=\int_0^\infty e^{-s |k|^2}\,\dd s\) at nonzero
momentum. Its marked sum costs at most
\(C_V\min(\delta_i r^{-4}, r^{-2})\): integrate up to
\(\delta_i\) and sum the harmonic readout by shells, or use the
whole proper-time integral. Thus the estimate in the preceding
corollary gains \(\prod_{i\in E}\delta_i\) at a cost
\(r^{-2|E|}\). To test the hard tail, insert the \(M\)-th power
of one plus the largest internal momentum norm. Each charged
momentum is an anchor or a fixed external mode plus photon
transfers. The weighted translated Schwartz sum at an anchor
costs at most a further \(C\sigma^{-M}\) for its translated
momentum. Its translation has norm bounded by a fixed multiple
of the sum of the photon and probe norms. A weight on any one
sewn photon costs \(Cr^{-M}\) in either shell estimate. The
fixed-graph density therefore costs at most another fixed inverse
power of \(r\). On the stated region use
\(r^{-p}\ge\eta^{-p/\alpha}\) and divide by \((1+R)^M\).
Tonelli applies directly to the full weighted absolute sum.
All direct contacts and open charged chains use the same scales
as their graph. No step requires a hierarchy among the \(s_i\).
\end{proof}

\Needspace{4\baselineskip}
\begin{theorem}[The bare interacting source at a fluctuating readout]
\label{thm:dualrev-uvheat-real}
For the same representation, consider the full free Maxwell
probability space and integrate the real-background determinant
in \eqref{eq:s6-heatlink-source} against \(|u(r)|^2\dd r\)
before the single normalization. The resulting vacuum partition
is finite, continuous in real \(e\), and at least one if
\begin{equation}
 p>3\alpha,\qquad\text{or}\qquad
 p=3\alpha\quad\text{and}\quad B_V\le A_j .
 \label{eq:dualrev-uvheat-real-cond}
\end{equation}
It admits finite open charged coefficients tested in \(L^2\),
smooth real transverse zero-mean background shifts applied before
the harmonic readout, and smooth linear Maxwell probes.
If \(p>3\alpha\), or \(p=3\alpha\) with \(B_V<A_j\), the source
is also smooth in the real electric coupling under the joint
integral. Proposition~\ref{prop:uvheat-critical-charge} extends
it to $p=3\alpha$, $B_V=A_j$ if $\alpha>6/7$, and proves two
continuous charge derivatives at matching coefficients for any $\alpha>0$.
In these smooth cases the Taylor coefficients at zero for fixed
Fourier probes coincide with the absolutely convergent bare graph
expansion at each order.
\end{theorem}
\begin{proof}
Use the transverse slice and write \(b=P_r A^T\).
The kernel \(K_\sigma\) in \eqref{eq:s6-heatK} uses the charged
width \(\sigma\) in place of \(a\). Write \(T=T_b\) for its Wilson
operator in \eqref{eq:s6-D}; the phase already includes \(e\).
For each \(r>0\) the fixed-depth proof supplies a smooth readout
and a trace-class kernel. A single Maxwell sample has polynomially
bounded Fourier coefficients almost surely, so the construction
holds at every positive \(r\). Pairing the spin eigenvalues
gives \(1\le\mathcal D_m(b)\le\exp(\|T\|_2^2/(2m^2))\).
It remains to track the cost as the common depth shrinks.

The spin Hilbert--Schmidt norm of the covering kernel equals
\(F_\sigma(z)=|z|g_\sigma(z)/\sigma^2\). Thus
\(\int F_\sigma^2\,\dd^4z=(16\pi^2\sigma^6)^{-1}\).
Two distinct lifts separated by \(w\in\mathcal L\) have overlap
at most this same factor times
\((1+|w|^2/(16\sigma^2))e^{-|w|^2/(8\sigma^2)}\).
This follows by shifting by \(w/2\) and bounding
\(|z-w/2|\,|z+w/2|\le|z|^2+|w|^2/4\). Since the Wilson phases
have modulus one, integrating over the other endpoint gives
\[
 \log\mathcal D_m(b)\le B_V\sigma^{-6}(1+\varepsilon_\sigma),
 \qquad
 0\le\varepsilon_\sigma\le C_{\mathcal L}e^{-c_{\mathcal L}/\sigma^2}.
\]
The last estimate follows by summing over the nonzero periods
and reducing the Gaussian exponent to absorb the polynomial.
These bounds are uniform in the Maxwell realization and real charge.
In the strict cases of \eqref{eq:dualrev-uvheat-real-cond},
\eqref{eq:dualrev-uvheat-ground} absorbs the determinant and
every fixed inverse power of \(r\). At \(p=3\alpha\) and
\(B_V=A_j\) the leading exponents cancel, while
\(r^{-p}\varepsilon_{r^{\alpha/2}}\to0\). The resulting
density is bounded by \(C r^d\) near zero, which is integrable.
Away from zero the bounds are uniform. For fixed \(r,A^T\) the
determinant is continuous in real \(e\). Dominated convergence gives
continuity of the integrated source. The ground state and Maxwell
law are normalized, so its vacuum has mass at least one.

Open charged coefficients multiply the determinant by finite minors
of \(D_b^{-1}\) paired with the external wavefunctions. They are
bounded by a constant depending on the probes and \(m\), by
\(\|D_b^{-1}\|\le m^{-1}\). For a smooth mean-zero transverse shift
\(h\), change the integration variable to \(A^T+h\) in all functions
of the shifted readout. The Maxwell density then acquires
\[
 \exp\!\left(\langle (-\Delta)h,A^T\rangle
             -\tfrac12\|(-\Delta)^{1/2}h\|_2^2\right),
\]
with norms taken in the Maxwell covariance convention. This is the
Gaussian translation formula~\cite{DualRevBogachev1998}, valid since smooth transverse \(h\)
has finite Cameron--Martin norm. Its derivatives are integrable
Gaussian polynomials times a Gaussian exponential on every compact
set of shift amplitudes. A simultaneous linear photon source acquires
the corresponding scalar translation and also has finite exponential
moments. At equality of the coefficients, each such integrated derivative
on the region \(r^\alpha<\eta\) costs at most \(C\eta^{(d+1)/\alpha}\),
by integration of \(r^d\). This proves the asserted probe regularity. It uses the prescribed
source-independent boundary law and shifts of the physical transverse
Maxwell variable before readout.

To verify differentiation of the charged kernel itself in the
strict cases, let \(J\) contain \(n\ge1\) labelled transverse probe
directions \(a_i\) and write \(T_J=\partial_J T\),
\(L_J=\prod_{i\in J}|e|\|a_i\|_\infty\).
Each derivative of one Wilson chord brings its line integral,
bounded in modulus by \(|z|\|a_i\|_\infty\) times the charge.
Schur's bound on the integral kernel and the preceding squared
kernel calculation, now with the extra powers of \(|z|\), give
\[
 \|T_J\|\le C_n L_J\sigma^{n-1},\qquad
 \|T_J\|_2\le C_{n,V}L_J\sigma^{n-3}.
\]
A one-block derivative term in \(\log\mathcal D_m\) is
\(\operatorname{Tr}(D_b^{-1}T_J)\). The diagonal spin trace
of \(T_J\) is zero on every lift. Substituting
\(D_b^{-1}-m^{-1}I=-m^{-1}D_b^{-1}T\) gives a bound
\(C L_J\sigma^{n-6}\). A cyclic term with \(k\ge2\) contact
blocks uses two Hilbert--Schmidt norms and \(k-2\) operator
norms; its bound is \(C L_J\sigma^{n-k-4}\).
The finite partition formula for derivatives of an exponential
therefore gives
\[
 |\partial_J\mathcal D_m(b)|
     \le C_{n,V,m}\,L_J\,\mathcal D_m(b)\,\sigma^{-5n}
       \qquad(0<\sigma\le1).
\]
The open inverse obeys the corresponding product rule. These
estimates include a contact with every probe on a single chord.

A charge derivative instead inserts the random \(P_r A^T\) in
the chord. For finite \(h\ge1\), the \(L^h\) norm of
\(\|P_r A^T\|_\infty\) under the Maxwell law is at most
\(C_{h,V}\sum_{k\ne0}e^{-r|k|}/|k|\le C'_{h,V}r^{-3}\).
The absolute Fourier sum gives the same random majorant for every
real coupling. Its fixed moments, the inverse mass bounds, and
the radial exponential margin dominate each derivative locally
under the integral. At zero charge Wick contraction gives the
sewn transverse photons and all chord contacts in
Proposition~\ref{prop:dualrev-uvheat-depth}. Absolute convergence
justifies the resulting momentum sums. This proves coefficient
identification without assuming convergence of the perturbation series.
At exact coefficient equality the pointwise bound $C r^d$ alone need
not absorb these inverse powers. The following estimate uses the
distribution of the Maxwell readout.
\end{proof}

\begin{proposition}[Charge derivatives at matching vacuum and boundary coefficients]
\label{prop:uvheat-critical-charge}
Assume that $d-1=3\alpha$ and
$V/(32\pi^2m^2)=2|j|/(d-1)$ for the straight Wilson chords and
zero-mean Gaussian Maxwell field of the preceding theorem. The
integrated charged source is twice continuously differentiable in real
$e$ for every $\alpha>0$ and is smooth in $e$ for $\alpha>6/7$.
Both conclusions include fixed open charged probes and mixed derivatives
in finitely many smooth transverse shifts and linear Maxwell probes.
For the fixed Fourier probes at $e=0$, the derivatives give the
bare graph coefficients with their Taylor factorials. The exponent $6/7$ is a sufficient condition for
the estimate below.
\end{proposition}
\begin{proof}
At matching coefficients the direct bound on the squared kernel
leaves only $C r^d$ at the tip. We estimate how much the exact
determinant falls below that bound under Maxwell averaging. Write
$X=T^*T/m^2$, $H=\operatorname{Tr}X$ and
$f(x)=x-\log(1+x)$ for $x\ge0$. Spectral pairing gives the identity
\[
 \log\det(I+T/m)=\tfrac12 H-\tfrac12\operatorname{Tr}f(X).
\]
The diagonal lift terms in $H$ give $V/(16\pi^2m^2\sigma^6)$
exactly. The sum of the absolute values of all cross-lift integrals
is bounded by this quantity times $C e^{-c/\sigma^2}$, by the
overlap estimate above. In particular $H\ge c\sigma^{-6}$ uniformly
in the field near the tip. Here and below positive constants may
depend on the torus, the action and a compact set of real sources.

\emph{Two charge derivatives.} The following estimates hold for every $\alpha>0$ at this
endpoint. Set $M=T/m$ and $Q=\operatorname{Tr}f(X)$, $X=M^*M$.
Primes denote real charge derivatives. A smooth transverse shift can first
be transferred to the Gaussian density as in the preceding proof. On each individual lift the
squared modulus of the Wilson phase is one, so its contribution to $H$
is independent of the charge. Only cross terms between distinct
lifts remain in $H'$ and $H''$. If $I(x,z)$ denotes the uncharged
chord integral of the readout, Gaussian moments give
$\|I(x,z)\|_{L^h}\le C_h |z|r^{-1}$ for fixed $h\ge2$. Indeed
the Maxwell variance at a point of the readout is at most
$C\sum_{k\ne0}e^{-2r|k|}|k|^{-2}\le C' r^{-2}$.
Minkowski's inequality integrates this bound along the chord.
For either cross derivative insert at most two factors of the
difference of the two chord integrals. Gaussian overlap on distinct
lifts then bounds its $L^h$ norm by a fixed power of $r^{-1}$ times
$e^{-c/\sigma^2}$. For the differentiated operators the same
Gaussian scaling gives
\[
 \big\|\|M^{(k)}\|_2\big\|_{L^h}
       \le C_{k,h}\,\sigma^{k-3}r^{-k},\qquad k=1,2.
\]
Here the diagonal squared kernel costs $C\sigma^{-6}$ and each
power of $|z|$ supplies a factor $\sigma$. Weighted Cauchy--Schwarz
over the lifts preserves these powers. The absolute envelopes are independent of the real charge because
the phases have modulus one, so they also bound compact suprema.

For a skew-adjoint trace-class variation $B$, differentiating
$Q=\operatorname{Tr}X-\operatorname{Tr}\log(1+X)$ gives
\[
 |\mathrm d Q[M](B)|
 =2\left|\operatorname{Re}\operatorname{Tr}
   \big[(MX(1+X)^{-1})^* B\big]\right|
 \le C\sqrt Q\,\|B\|_2 .
\]
In fact $x^3/(1+x)^2\le 2 f(x)$ for $x\ge0$.
Along $M+tB$ the second derivative at zero equals
$2\|B\|_2^2+2\operatorname{Re}\operatorname{Tr}(RBRB)$, with
$R=(1+M)^{-1}$. This follows equally from
$Q=H-2\log|\det(I+M)|$. Since $\|R\|\le1$, its modulus is at
most $4\|B\|_2^2$. Applying the chain rule to $M(e)$ therefore
bounds $|Q''|$ by
$4\|M'\|_2^2+C\sqrt Q\,\|M''\|_2$.
The full weighted determinant is at most $Cr^d e^{-Q/2}$.
In its first two derivatives the factors $\sqrt Q$ and $Q$
are absorbed by this exponential. The inverse factors of open
coefficients have mass-bounded norm and their first two
derivatives cost at most fixed multiples of $\|M'\|_2$ and
$\|M'\|_2^2+\|M''\|_2$, respectively. H\"older's inequality
includes the Gaussian densities and polynomials for any finite
collection of smooth transverse shifts or linear probes.
The expected absolute second derivative is thus bounded at the
tip by $C r^d\sigma^{-4}r^{-2}=C r^{\alpha-1}$, using
$d=1+3\alpha$. The first derivative has an integrable bound as well.
The estimates give compact-uniform integrable majorants and hence
two continuous charge derivatives of the full source. At zero charge
they agree with the corresponding bare coefficients, with their
Taylor factorials.

\emph{Higher charge derivatives.} To track the momenta that carry $H$, differentiate the kernel in
its input endpoint $y$ with $x$ fixed. On each lift put $z=y-x$
and $b_r=e^{-r\sqrt{-\Delta}}A^T$. The derivative of the real
chord integral $I(x,z)=\int_0^1 z\cdot b_r(x+tz)\,\mathrm dt$ is
\[
 \partial_{y_\mu} I
  =\int_0^1 \big[b_{r,\mu}(x+tz)
           +t z\cdot\partial_\mu b_r(x+tz)\big]\,\mathrm dt .
\]
The first term has variance at most $C r^{-2}$ by summing the
Maxwell mode variances. For the second term use
$|\int_0^1 t e^{itu}\,\mathrm dt|\le C/(1+|u|)$.
On the unit three-sphere the component parallel to a fixed
direction has a bounded density. Its angular average therefore
satisfies
\[
 \int_{S^3}\frac{\mathrm d\omega}
                   {(1+|k\cdot\omega|s)^2}
       \le \frac{C}{1+s|k|},\qquad s\ge0 .
\]
For each $k\ne0$ the differentiated covariance is bounded by
$C e^{-2r|k|}$. The angularly averaged variance of the second
term at $|z|=s$ is consequently bounded by
$C s\sum_{k\ne0}e^{-2r|k|}/|k|\le C s r^{-3}$.
The estimate uses the integration over displacement directions
in the action.

Let $S_r^2=\sigma^{-2}+r^{-2}+\sigma r^{-3}$.
For $h\ge2$, with $\mathbb E$ denoting Maxwell expectation, the kernel
derivative has a common random upper bound $Y_r$ over the compact
set of sources, with
\[
  \|\nabla_y T\|_2\le Y_r,\qquad
  (\mathbb E Y_r^h)^{1/h}
       \le C\sqrt h\,\sigma^{-3}S_r .
\]
Here $\|\nabla_y T\|_2$ is the $L^2$ norm of the differentiated
matrix kernel. To obtain $Y_r$, bound the absolute derivative on
each lift by $|\nabla K_\sigma|+
C|K_\sigma|(1+|z|+|\nabla_y I|)$, using matrix norms.
A smooth deterministic shift adds only the displayed $1+|z|$
cost. In the sum over lifts use Cauchy--Schwarz with weights
$e^{-|z|^2/(8\sigma^2)}$, whose periodization is uniformly bounded.
The square integral then uses Gaussian weights broadened by a
constant factor. Their mass for $|K_\sigma|^2$ is $O(\sigma^{-6})$,
and insertion of $|z|$ costs an additional factor $O(\sigma)$.
The derivative of $K_\sigma$ costs $O(\sigma^{-1})$ in norm.
Minkowski's integral inequality and the $C\sqrt h$ Gaussian
moment bound now give the estimate.

For small $r$ project onto input Fourier modes $|k|\le L$,
including spin, where $L=r^{-\epsilon}S_r$ and $\epsilon>0$.
Write $P$ for this projection, of rank $N\le C L^4$. Parseval gives
\[
 \operatorname{Tr}((1-P)X)
       \le m^{-2} L^{-2}Y_r^2 .
\]
Outside an event of probability at most $C e^{-c r^{-2\epsilon}}$
the right-hand side is less than $c_0\sigma^{-6}/2$, where
$c_0\sigma^{-6}$ is the uniform lower bound on $H$. This follows
from Markov's inequality with an integer moment of order
$c' r^{-2\epsilon}$ and $c'>0$ sufficiently small.
Hence $a=\operatorname{Tr}(PX)\ge c_0\sigma^{-6}/2$ on this event's
complement. Applying scalar Jensen in an orthonormal basis of
$P$ and then averaging the $N$ diagonal entries gives
\[
 \operatorname{Tr} f(X)\ge N f(a/N)
           \ge \frac{a^2}{2(N+a)} ,
\]
where $f(x)\ge x^2/[2(1+x)]$ follows by integration.
The comparison acts on the complete trace. It introduces no
restriction in the defining action or in its momentum sums.

For $0<\alpha\le2$ we have $S_r^4\le C r^{\alpha-6}$.
For $\alpha\ge2$ the bound is $S_r^4\le C r^{-2\alpha}$.
On the event where $a$ is large the last trace is therefore
bounded below by $c r^{-\beta}$ for some $\beta>0$ if
$\alpha>6/7$. Indeed it suffices that both $\sigma^{-6}$ and
$\sigma^{-12}/(r^{-4\epsilon}S_r^4)$ grow as inverse powers.
The second power is $7\alpha-6-4\epsilon$ for $\alpha\le2$ and
$4\alpha-4\epsilon$ for $\alpha\ge2$ in these lower bounds.
Choose $\epsilon$ small enough that the appropriate power is positive.

Multiplying the ground-state density by the exact determinant now gives a bound of
at most $Cr^d e^{-c r^{-\beta}}$ on the event with this deficit
and $Cr^d$ everywhere. Each fixed derivative and open minor obeys
the contact and inverse estimates above. After factoring out the determinant,
its compact-uniform majorant has every fixed Maxwell moment bounded
by a power of $r^{-1}$. Smooth linear Maxwell sources preserve these
bounds by H\"older's inequality. On the exceptional event apply
Cauchy--Schwarz to this majorant and the event indicator. The
expected absolute derivative density is then bounded near the tip by
$C e^{-c r^{-\gamma}}$ for some $\gamma>0$, uniformly on the compact
source set. On the quartic chart $\alpha=1$, the choice
$\epsilon=1/6$ allows $\gamma=1/3$. This justifies differentiation of the joint integral at
each finite order. At zero charge the same Wick expansion as in the
strict cases identifies each derivative for fixed Fourier probes with
the corresponding absolute bare sums, with its Taylor factorial.
\end{proof}

\begin{corollary}[Weak-charge expansion at the matched radial endpoint]
\label{cor:uv-geom-charge-origin}
Keep the matched coefficients, straight Wilson chords and Gaussian
Maxwell state of Proposition~\ref{prop:uvheat-critical-charge}. At $\alpha=6/7$ the
full source has a real asymptotic expansion at $e=0$ through every
finite order. For $0<\alpha<6/7$ the same conclusion holds through
every nonnegative integer order $N$ satisfying
\[
             N<\frac{4(d+1)}{6-7\alpha},\qquad d=1+3\alpha .
\]
Each coefficient is the corresponding bare coefficient of the
original action. The statement applies to the vacuum, fixed open
charged probes, and mixed variations of finitely many smooth
transverse shifts and linear Maxwell probes. The normalized source
uses the quotient of the two expansions. For $6/11<\alpha\le6/7$
the third ordinary charge derivative also exists at zero.
\end{corollary}
\begin{proof}
Write $F(r,e)$ for the Maxwell average of any of these unnormalized
source densities, including the radial probability density. At fixed
$r>0$ it has charge derivatives of all finite orders. At matching
coefficients the unmarked bound is $|F(r,e)|\le C r^d$ near the tip.
For external shift derivatives, translate the smooth shift into the
Gaussian density first. The resulting Gaussian polynomials and
linear-source exponentials have finite moments, uniformly on compact
sets of source amplitudes. Open inverse factors obey the mass bound.

The dependence on the electric charge in the previous proof matters
near $e=0$. Differentiating a Wilson phase in its spatial endpoint
inserts $e\,\partial_y I$. The derivative of its Gaussian edge
kernel is independent of $e$. Thus on $|e|\le r^\lambda$, $\lambda>0$,
the same random Sobolev estimate and Fourier projection use
\[
 S_{\lambda,r}^2=\sigma^{-2}
       +r^{2\lambda}(r^{-2}+\sigma r^{-3}),
 \qquad L=r^{-\epsilon}S_{\lambda,r}.
\]
The exceptional-event bound remains $C\exp(-c r^{-2\epsilon})$:
the Gaussian part of the absolute kernel derivative carries the
same factor $r^\lambda$, uniformly on this interval of charges.
The lower bound on the full Hilbert--Schmidt trace is uniform
there as well. For $0<\alpha\le6/7$,
\[
 S_{\lambda,r}^4\le
 C\big(r^{-2\alpha}+r^{4\lambda+\alpha-6}\big).
\]
The compressed trace estimate of the preceding proof now yields
a determinant deficit growing at least as a positive inverse power
of $r$ on the complementary event whenever
$7\alpha-6+4\lambda>0$. Indeed the two powers in the lower
bound for $\sigma^{-12}/L^4$ are $4\alpha-4\epsilon$ and
$7\alpha-6+4\lambda-4\epsilon$. Both are positive for
sufficiently small $\epsilon>0$. Each finite charge derivative
costs polynomial moments of the marked kernels. The same
Cauchy--Schwarz estimate on the exceptional event therefore gives,
for every fixed $k$, some $c,\gamma>0$ and $C_k<\infty$ with
\[
   \sup_{|e|\le r^\lambda}|\partial_e^k F(r,e)|
        \le C_k\exp(-c r^{-\gamma})
\]
near the tip. This estimate averages the original determinant and
covers every charged contact assigned by the Wilson action.

For $|e|<1$ split the radial integral at $r_e=|e|^{1/\lambda}$.
Throughout $r\ge r_e$ the segment joining zero to $e$ lies in the
charge interval just estimated. Taylor's formula there has integrated
remainder $O(|e|^{N+1})$. For $r<r_e$ the original source contributes
at most $C r_e^{d+1}$. Each coefficient density at zero charge is
bounded by the preceding exponential, so its integral over this
smaller interval decays faster than any power of $|e|$. Thus, in
the fixed reference units,
\[
 \begin{aligned}
 \int_0^1 F(r,e)\,\mathrm dr
 &=\sum_{k=0}^N c_k e^k+
 O\!\left(|e|^{N+1}+|e|^{(d+1)/\lambda}\right),\\
 c_k&=\frac1{k!}\int_0^1\partial_e^k F(r,0)\,\mathrm dr .
 \end{aligned}
\]
The radial split estimates a remainder within one full source.
Wick contraction at zero identifies $c_k$ with the absolutely
summable bare graphs at that order for the specified Fourier probes.
A positive zeroth vacuum coefficient permits ordinary finite-order
series division for the normalized source.
When $\alpha=6/7$ take $\lambda>0$ sufficiently small to make
$(d+1)/\lambda>N+1$. Below $6/7$ choose
$\lambda>(6-7\alpha)/4$ with $(d+1)/\lambda>N$. This gives the
claimed asymptotic remainder $o(|e|^N)$.

For the last assertion, apply the same split to
$\partial_e^2F(r,e)$. The global second-derivative bound
$C r^{\alpha-1}$ in the preceding proof gives a tip contribution
$O(|e|^{\alpha/\lambda})$. If $\alpha>6/11$, one can choose
$\max(0,(6-7\alpha)/4)<\lambda<\alpha$. The result is a
first-order expansion of the second derivative at zero, with linear
coefficient $3!c_3$. Thus the third iterated derivative at zero is
defined. Higher-order asymptotic coefficients at the origin do not
by themselves imply iterated differentiability on a neighborhood or
convergence of the infinite series.
\end{proof}

For the quartic chart \((d,\alpha)=(4,1)\) the sufficient real-source
condition is \(|j|\ge3V/(64\pi^2m^2)\) in the fixed reference units.
The proven charge derivatives extend to equality because $1>6/7$.
The chart \((6,2)\) satisfies the graph criterion, while the above
uniform-field determinant estimate does not decide its interacting
partition. Both statements keep the original angular preparation and
all lifted Wilson chords, including returns of zero net winding.
Identifying the physical GCM-induced sector would additionally require
its chiral charged operator, its boundary half-densities and their
common variations. The spectral test above still applies to the
unmodified open spinor of this represented source.

\paragraph{Which angular states protect the free vacuum?}
The nonzero angular number in the graph theorem is a condition on the
normal preparation. To see what happens if a neutral component is
occupied, keep the metric, both Friedrichs endpoints, and the same heat
lapse. For $j=0$ put $\nu=(d-1)/2>0$. In the radial variable normalized
with respect to $\mathrm dr$, the regular ground state is proportional to
$\sqrt r\,J_\nu(\sqrt{E_0}r)$, where $E_0>0$ is its ground energy.
Indeed, substitution reduces the radial equation to Bessel's equation.
The independent singular solution has leading physical radial wavefunction
proportional to $r^{1-d}$ and has infinite kinetic energy at the tip.
The regular series starts with a nonzero multiple of $r^{d/2}$ in
the normalized variable, so the neutral probability density is
asymptotic to $c r^d$ with $c>0$. The free determinant grows faster
than every inverse power. It follows that any positive weight on this
neutral ground state makes the common free normalization infinite.
For a finite diagonal mixture supported entirely on nonzero angular
ground states, integration is a finite positive sum. The same sharp
criterion $d-1>2\alpha$ and the graph estimates therefore apply.
Here $j$ labels the normal preparation independently of the electric
coupling multiplying a Wilson chord.

\paragraph{What changes when the carrier volume grows}
The vacuum tilt in \eqref{eq:s6-tilted-depth-law} can concentrate
the shared depth near zero in the large-volume limit. Dilate a
periodic box of fixed shape to volume $V\to\infty$ while holding
the radial preparation $\mu$ fixed. Denote the vacuum factor in
that box by $Z_{0,V}$. For each fixed $r>0$ the Riemann sums give
\[
 \begin{aligned}
 V^{-1}\log Z_{0,V}(r)&\longrightarrow F(r),\\
 F(r)&=2\int_{\RR^4}\log\!\big(1+m^{-2}|k|^2e^{-2r^\alpha |k|^2}\big)
                  \,\frac{\dd^4 k}{(2\pi)^4}.
 \end{aligned}
\]
Both $F(r)$ and each finite-volume determinant are strictly decreasing
in $r$. For this fixed radial law even the logarithm of the free normalization
per unit volume has no finite limit. Put $\mathcal Z_{0,V}=\int
Z_{0,V}(r)\,\mu(\mathrm dr)$ and suppose $d-1>2\alpha$, so every
$\mathcal Z_{0,V}$ exists. For fixed $0<\eta<1$, monotonicity gives
\[
 \liminf_{V\to\infty} V^{-1}\log\mathcal Z_{0,V}\ \ge F(\eta),
\]
since $\mu([\eta/2,\eta])>0$ is independent of $V$. As $\eta$ decreases
to zero, the integrand defining $F(\eta)$ increases to
$2\log(1+|k|^2/m^2)$ with the same Fourier measure. Its integral
is infinite. Thus $V^{-1}\log\mathcal Z_{0,V}\to+\infty$. For each
fixed $0<\varepsilon<1$ the tilted probability of $r\ge\varepsilon$
is at most $Z_{0,V}(\varepsilon)/\mathcal Z_{0,V}$ and consequently
decays faster than $e^{-cV}$ for every fixed $c>0$. In particular
each positive inverse moment diverges as $V\to\infty$. This conclusion
uses the full depth-dependent free determinant relative to the fixed
mass reference and one radial state shared by the expanding carrier.
The finite-volume graph bounds remain valid at every fixed $V$.

\paragraph{Exactly retaining one chirality of a massive pair}
Let $P_\pm=(I\pm\gamma_5)/2$, write $\mathcal H_\pm$ for their ranges,
and set $W=P_-T_BP_+$ on the fixed
Euclidean sector plane. For a real connection, the two identities
$T_B^*=-T_B$ and $\gamma_5 T_B\gamma_5=-T_B$ give
\[
 D_B=\begin{pmatrix}mI&-W^*\\W&mI\end{pmatrix},\qquad
 Q_+=mI+m^{-1}W^*W .
\]
Integrating the minus Berezin pair, with probes only in the plus
row and column slots, gives the Schur complement $Q_+$. In particular
\begin{equation}
 \det(I+T_B/m)=\det(I+W^*W/m^2),\qquad
 \left.P_+D_B^{-1}P_+\right|_{\mathcal H_+}=Q_+^{-1}.
 \label{eq:geom-uv-paired-schur}
\end{equation}
Indeed, factor $D_B/m$ into upper and lower block triangular operators
using the lower diagonal identity. The off-diagonal triangular factors
differ from the identity by trace-class operators of square zero and
have ordinary Fredholm determinant one. The remaining diagonal block
is $Q_+/m$. Block inversion gives the second identity and
$\|Q_+^{-1}\|\le m^{-1}$. This calculation retains the full
depth-dependent vacuum before the one common normalization. Gauge
conjugation preserves both chiral summands on this Abelian chart.

There is also a direct graph check in the resulting variables.
For a nonempty set $J$ of labelled photon marks,
\[
 \partial_J Q_+=m^{-1}\sum_{I\subseteq J}
       (\partial_I W)^*\,\partial_{J\setminus I}W
\]
for real variations. Thus every contact is a sum of products of
two Wilson factors with the original Gaussian displacement. At zero
background, each factor at $s\ge0$ marks is an average of a translated
Schwartz momentum symbol, of size $O_s(\sigma^{s-1})$ in supremum norm
and $O_{s,V}(\sigma^{s-5})$ in the sum over a translated loop momentum,
where $\sigma^2=r^\alpha$. The case $s=0$ is the free single-chord symbol.
In each product the routed momentum of the intermediate charged slot
is fixed by the incoming momentum and the marked photons.
For every closed loop with a vertex, sum one such factor over its loop
momentum and bound the other factors by their suprema. Free open
lines at prescribed Fourier modes have the mass bound above. The
transverse readout at the ends of each internal photon is unchanged.
The total cost at any fixed graph is at most a fixed inverse power
of $r$, which the same free determinant-weighted law integrates for
$d-1>2\alpha$. Vacuum factors without marks are already included in
\eqref{eq:geom-uv-paired-schur}. This proves absolute finiteness of the bare graphs of
this chiral marginal. Its quadratic operator is the positive Schur
complement of a massive paired action. An independent first-order
Weyl determinant would also require its own domain, measure and phase.

\paragraph{A paired charged source with a free Dirac pole and finite bare graphs.}
There is a direct test of whether the pole and the geometric suppression can
coexist. Keep the same displacement kernel, readout and common radial
state. On the transverse slice use the relative Wilson transport
$W_b(x,z)-1$ for the interaction on each chord, and use the ordinary massive
Dirac kinetic form for the free charged field. In the same sign convention for
the chord as above, write
\begin{equation}
 \begin{aligned}
 D_{\rm D}&=m+\gamma\cdot\partial,\qquad R_0=D_{\rm D}^{-1},\\
 Y_b&=T_b-T_0,\qquad D^{\rm rel}_b=D_{\rm D}+Y_b .
 \end{aligned}
 \label{eq:phys-uv-relative-chord}
\end{equation}
Here $b$ is the harmonically read transverse field, $T_b$ already includes
the real electric charge, and $\sigma^2$ is the charged lapse. Relative
transport compares the marked chord with its value in the flat reference
connection. In the original gauge variables conjugate the entire block,
including its domain and external charged slots, by the longitudinal
phase of the source on the trivial bundle. The dual Berezin factors cancel
under this unitary change of variables. The transverse action
\eqref{eq:phys-uv-relative-chord} is a specified comparison coupling. Covariance under
the small Abelian gauge group in these variables leaves its physical
transverse response to be checked.

At zero photon field the full charged matrix is $D_{\rm D}$ at every
$r$. The continued free Dirac Fourier symbol has its usual pole and positive
vector residue; the finite torus samples this symbol at discrete modes. An absolute
vacuum normalization relative to the single field- and depth-independent
Dirac reference is the ordinary Fredholm determinant
\[
  \mathcal D^{\rm rel}_m(b)=\det(I+Y_b R_0).
\]
This definition only takes a determinant of an identity plus a
trace-class operator. The massive free pole and all positive-order Wilson
contacts now belong to one functional.

\begin{proposition}[Bounds for the relative Wilson chord source]
\label{prop:phys-uv-relative-chord}
On the fixed four-torus with $m>0$ and the Maxwell zero mode fixed,
all finite-order bare electromagnetic graphs with fixed Fourier probes
of \eqref{eq:phys-uv-relative-chord} are absolutely summable when the one radial law has
every inverse moment and the widths are $r$ and $\sigma^2=\kappa r^\alpha$,
$\kappa,\alpha>0$. On the minimizing surface of Corollary~\ref{cor:geo-free-slope-uv}
with $\alpha=1$ and $j\ne0$, the unexpanded interacting partition also
exists at every real charge and is smooth there at fixed finite volume.
\end{proposition}
\begin{proof}
The free determinant equals one at each depth by the single reference
above. Every positive-order derivative of $Y_b$ is the corresponding
Wilson derivative of $T_b$, including the contact with all marks on one
chord. Thus one charged vertex per closed loop still has a translated
summable Fourier factor; the other charged vertices have uniform bounds.
The free inverses have norm at most $m^{-1}$. The two harmonic photon
endpoints give $e^{-2r|p|}/|p|^2$ on a nonzero sewn momentum. The absolute
cost is again at most a fixed inverse power of $r$ at each graph. This
proves the perturbative assertion, including simultaneous short visits.

For the unexpanded source, $Y_b$ is skew-adjoint and anticommutes with
$\gamma_5$. For each $r>0$ its kernel on a smooth harmonic readout is
smooth, hence trace class. The relative determinant is strictly positive.
Indeed, on finite Fourier sections commuting with $D_{\rm D}$ and
$\gamma_5$, the numerator and denominator have conjugate mass-shifted
eigenvalue pairs and positive real determinants. Trace-norm convergence
gives a nonnegative limit. The massive skew-adjoint inverse bound for
$D^{\rm rel}_b$ gives invertibility, excluding zero. These finite
sections test the sign of the already defined determinant.

A useful field-uniform upper bound is
\begin{equation}
     \log \mathcal D^{\rm rel}_m(b)\le C_{V,m}\sigma^{-4},
                 \qquad 0<\sigma\le1 .
 \label{eq:phys-uv-relative-chord-bound}
\end{equation}
To verify it, put $X=Y_bR_0$. The ordinary determinant inequality
$\log|\det(I+X)|\le\operatorname{Re}\operatorname{Tr}X+
\|X\|_2^2/2$ follows from
$\log|1+z|\le\operatorname{Re}z+|z|^2/2$ and finite-rank approximation
in trace norm. The modulus of $W_b-1$ is at most two. Scaling the
periodized Gaussian derivative gives $\|Y_b\|\le C\sigma^{-1}$ and
$\|Y_b\|_2^2\le C_V\sigma^{-6}$. In the free Fourier basis
$R_0(k)R_0(k)^*=I/(m^2+|k|^2)$. Modes with $|k|\le\sigma^{-1}$
therefore contribute at most $C\sigma^{-2}
\sum_{|k|\le\sigma^{-1}}(m^2+|k|^2)^{-1}=O(\sigma^{-4})$ to
$\|X\|_2^2$. The complement costs at most
$\sigma^2\|Y_b\|_2^2=O(\sigma^{-4})$.
Both estimates concern the full operator.

For the ordinary trace, the massive Dirac inverse kernel on the torus is
bounded by $C(1+\operatorname{dist}(x,0)^{-3})$ off the diagonal.
This follows by applying $m-\gamma\cdot\partial$ to the massive
scalar heat integral; its singular lift has heat-time integrands bounded
by constants times $(t^{-2}+|x|t^{-3})e^{-|x|^2/(4t)-m^2t}$.
Other lifts are smooth near zero. Pairing this locally integrable
kernel with the periodization of
$C|z|\sigma^{-6}e^{-|z|^2/(4\sigma^2)}$ gives
$|\operatorname{Tr}X|=O(\sigma^{-4})$, including the integral over
the free endpoint. This proves \eqref{eq:phys-uv-relative-chord-bound}.

On the minimizing surface the upper bound grows in the exponent only
as $C r^{-2}$, while the radial law decays as $r^4e^{-A_* r^{-3}}$ up to
positive constants. The bound is uniform in the real coupling and the
Maxwell realization. For differentiability, each marked-chord derivative
inserts line integrals of the harmonic field, with finite Gaussian
moments bounded by powers of $r^{-1}$ as before. The single-block trace
$\operatorname{Tr}[(D^{\rm rel}_b)^{-1}Y_b^{(h)}]$ can be split using
$(D^{\rm rel}_b)^{-1}=R_0-(D^{\rm rel}_b)^{-1}Y_bR_0$.
The first trace uses the same locally integrable free kernel with the
chord insertions; the second uses two Hilbert--Schmidt factors.
Multiple-block traces and derivatives of open inverses use the mass
bound as in the heat-link proof. Each fixed derivative consequently
costs only further inverse powers after Maxwell expectation, locally
uniformly in real charge. The radial margin justifies differentiation
and the usual finite-order identifications at zero charge.
\end{proof}

The transverse soft response detects a remaining difference from a
physical electromagnetic source. To calculate it, use the fixed-width
translation-invariant continuum coefficient per unit volume. Put
$h(k)=k e^{-\sigma^2|k|^2}$ and $L(k)=2\log(m^2+|k|^2)$.
For a constant transverse test $B$ the logarithmic determinant density
has integrand
$L(k+h(k+eB)-h(k))-L(k)$, up to the immaterial choice of sign of $eB$.
This is the continuous zero-momentum limit of the two marked vertices
and the direct contact: the Gaussian factors dominate their charged
Fourier integrals uniformly for soft external momenta. It does not
assign a value to a removed torus zero mode. Let $H_m$ be one quarter
of the trace of its Hessian in $B$ at zero. An integration by parts,
with no boundary term because $h$ is Schwartz, gives
\[
 H_m=\frac{e^2}{4}\int \operatorname{tr}
 \big[L''(k)((Dh(k))^2-Dh(k))\big]\,
                 \frac{\mathrm d^4 k}{(2\pi)^4}.
\]
Here $Dh$ is a real symmetric matrix. Since
$m^2L''(k)\to4I$ as $m\to\infty$ and $\|m^2L''(k)\|\le4$,
dominated convergence, $\int\operatorname{tr}Dh=0$, and a radial
Gaussian integral give
\[
  \lim_{m\to\infty}m^2 H_m
   =e^2\int\operatorname{tr}\!\big[(Dh)^2\big]\frac{\mathrm d^4 k}{(2\pi)^4}
   =\frac{3e^2}{32\pi^2\sigma^4}.
\]
The trace here is the trace of the matrix square. At the other end of
the mass range, set $u=|k|^2$ and $y=\sigma^2u$. The tangential and
radial eigenvalues of $Dh$ are $e^{-y}$ and $(1-2y)e^{-y}$.
Those of $L''$ are $4/(m^2+u)$ and $4(m^2-u)/(m^2+u)^2$.
The four-dimensional radial measure supplies a factor $u$, so each
weighted eigenvalue of $L''$ is bounded in absolute value by $4$.
Dominated convergence gives
\[
 \begin{split}
 \lim_{m\downarrow0}H_m
  &=\frac{e^2}{16\pi^2\sigma^2}\int_0^\infty
   \big[e^{-2y}(2+4y-4y^2)-e^{-y}(2+2y)\big]\,\mathrm dy\\
  &=-\frac{3e^2}{16\pi^2\sigma^2}.
 \end{split}
\]
For $e\ne0$, continuity therefore gives at least one positive value of
$m\sigma$ at which the full one-loop transverse constant vanishes.
It is nonzero for sufficiently large and sufficiently small $m\sigma$;
changing to the negative logarithm reverses both signs.

The same argument covers one common fluctuating depth. Write
$H_m(\sigma)=e^2\sigma^{-2}F(m\sigma)$ by rescaling $k$.
The two limits and continuity give
$\sup_{s>0}(1+s^2)|F(s)|<\infty$. For a mass-independent normalized
depth law $\nu_0$ with $\langle\sigma^{-4}\rangle_0<\infty$, the
continuum one-loop coefficient is
$\overline H_m=\langle H_m(\sigma(r))\rangle_0$.
Indeed, the free relative determinant is identically one and its first
variation in the constant test field vanishes by inversion at every
depth. Differentiating the logarithm of the depth average consequently
adds no covariance of first variations. A harmonic readout on each
external leg tends to one in this soft limit. For the radial ground state,
the marked graph bounds dominate the integrands uniformly for bounded
soft momenta at fixed $m>0$. Dominated convergence also gives
\[
 \lim_{m\downarrow0}\overline H_m
   =-\frac{3e^2}{16\pi^2}\langle\sigma^{-2}\rangle_0,\qquad
 \lim_{m\to\infty}m^2\overline H_m
   =\frac{3e^2}{32\pi^2}\langle\sigma^{-4}\rangle_0 .
\]
Thus at least one positive mass solves the averaged one-loop matching
condition for each such fixed geometric law and $e\ne0$. The physical
action must select that relation if it is to use this source. The fourth-order soft response and the interacting open covariance
are tested below on this same source.

\paragraph{Measuring the electric charge of the same Wilson source.}
The free pole and the transverse constant test different derivatives of
the action. There is a third quantity one can calculate without a loop:
the current measured by a transverse temporal probe on the free mass
shell. Use the continuum Fourier symbol on the selected sector plane,
with its previously specified analytic continuation. Take a spatial
photon momentum with zero carrier component and let its spatial
momentum go to zero in a Breit frame. A probe polarized along the
carrier is transverse along this sequence. The harmonic readout tends
to one. The Gaussian vertex is entire in the charged momentum, and the
two external massive momenta tend to the rest shell. In the Euclidean
carrier basis let $k_\parallel=i m$ there and
$P=(1+\gamma_\parallel)/2$. In the convention where a minimally coupled
Dirac vertex has unit charge ratio, the rest-spin matrix of the
one-photon chord vertex is
\[
 \begin{aligned}
 h(k)&=k e^{-\sigma^2 k^2},\\
 P\gamma_\nu (\partial_\parallel h_\nu)(k)P
    &=g(x)P,\\
 g(x)&=e^x(1+2x),\qquad x=m^2\sigma^2 .
 \end{aligned}
\]
Indeed $k^2=-m^2$ and $P\gamma_\parallel P=P$ on the pole.
The factors common to the two Euclidean vertex conventions cancel
in the ratio. The rest residue is one on this spin space. The measured
tree charge is therefore $e g(x)$ for the original transverse phase.
The overall sign follows the orientation of the charged chord.
This calculation concerns a static transverse limit on a specified
analytic symbol. It makes no identification of a torus zero mode with
an asymptotic particle.

It matters how a proposed physical normalization enters the action.
First multiply the \emph{entire interaction}
$Y_b=T_b-T_0$ by $c=1/g(x)$ at fixed depth. This gives unit tree
charge ratio at every width. The contact with two photon marks is
multiplied by $c$, whereas the two one-photon vertices are multiplied
by $c^2$. In the sign convention of the logarithmic determinant their
combined one-loop transverse constant is strictly negative for every
$m,\sigma>0$ and nonzero $e$. To see this, rescale $y=\sigma^2|k|^2$
in the Hessian calculation. It gives
\[
 H(c)=\frac{e^2 c}{16\pi^2\sigma^2}(c A_x-B_x),
\]
where
\[
 \begin{aligned}
 A_x=\int_0^\infty &y e^{-2y}
  \left[\frac{3}{x+y}
    +\frac{(x-y)(1-2y)^2}{(x+y)^2}\right]\,\mathrm dy,\\
 B_x={}&4\int_0^\infty e^{-y}
      \frac{y^2(3x+y)}{(x+y)^3}\,\mathrm dy .
 \end{aligned}
\]
For $B_x$ integrate the term linear in $Dh$ by parts and use
$\nabla\Delta L=-16k(3m^2+|k|^2)/(m^2+|k|^2)^3$.
The first integral is bounded above by
$\int_0^\infty e^{-2y}[3+(1-2y)^2]\,\mathrm dy=2$.
In the second integral, $v=y/(x+y)\ge1/2$ for $y\ge x$, and
$v^2(3-2v)\ge1/2$ on that interval. Hence
$B_x\ge2e^{-x}>2c$ for $x>0$. The strict sign survives any positive
average of these fixed-depth normalized coefficients whenever it exists.
This rules out that particular way of fixing the charge and keeping a
vanishing one-loop transverse constant.

The failure also holds if one multiplies the entire interaction by
a \emph{single} constant $c_* = 1/\langle g(x)\rangle_0$ to fix
only the averaged tree charge. Assume bounded widths and the inverse
moments above. To compare the two depth averages, write
$B_x=\int_0^\infty e^{-xt}a(t)\,\mathrm dt$ and
$A_x=\int_0^\infty e^{-xt}b(t)\,\mathrm dt$, where
\[
 a(t)=\frac{24t}{(1+t)^4},\qquad
 b(t)=\frac{24(2-3t+2t^2)}{(2+t)^5}>0 .
\]
Indeed the integral over $y$ for the respective Laplace densities
is $\int_0^\infty y e^{-(t+j)y}
[3+(1-2y)^j(1-2ty)]\,\mathrm dy$, with $j=1,2$.
The density for $B_x-e^{-x}A_x$ is $a(t)$ below $1$ and
$24(t-1)(7-t)/(1+t)^5$ above $1$. It changes sign once, at $7$,
and has integral $4-1=3$. Weighting by $e^{-xt}$ therefore gives
$B_x-e^{-x}A_x\ge 3e^{-7x}>0$. Now $w(x)=\sigma^{-2} A_x$ and
$e^{-x}$ are both decreasing in $x$ at fixed $m>0$. Their covariance
is nonnegative, as follows by integrating the product of their
differences on two independent copies of the depth law. Hence
\[
 \langle\sigma^{-2} B_x\rangle_0
 >\langle w e^{-x}\rangle_0
 \ge\langle w\rangle_0\langle e^{-x}\rangle_0
 \ge c_*\langle w\rangle_0 .
\]
The last step uses
$\langle e^{-x}\rangle_0\langle e^x\rangle_0\ge1$ and
$g(x)\ge e^x$. This proves the strict sign for the single outside
normalization as well.

Alternatively, the transverse coupling can be defined by a measured
charge \emph{inside the Wilson phase}. For a common normalized free
depth law with bounded $\sigma$ define
\[
 z_m=\left(\int g(m^2\sigma(r)^2)\,\mathrm d\nu_0(r)\right)^{-1}.
\]
Keep the longitudinal gauge representation fixed and put $e z_m$ in
the entire transverse chord phase. The free propagator is independent
of $r$, so its tree temporal vertex averages linearly. The resulting
averaged rest charge is $e$. Both terms of the quadratic response now
carry $z_m^2$. The one-loop matching masses are therefore precisely
those of $\overline H_m=0$ derived above. At each fixed $m>0$ the
factor is a finite scalar independent of all graph momenta. The
determinant bound and every graph estimate still apply. This gives
one comparison action satisfying tree charge normalization, one-loop
soft matching and all finite-order ultraviolet bounds together. It
uses the freedom to choose the coupling to the transverse field
separately from the longitudinal representation on the trivial chart.
If the native circle connection fixes both couplings, their relative
normalization and the mass must be obtained from that connection and
its physical state. The fourth-order and spectral tests below determine what this
normalization can achieve in the interacting source.

\paragraph{The fourth-order photon response of the relative Wilson source.}
A common fluctuating depth leaves a measurable term after one-loop
soft matching. It is important here to invert the photon two-point
function of the whole source. Inverting a separate two-point function
at each depth gives a different response.

Use a torus of period $2\pi R$, $R\ge1$, volume $V=(2\pi R)^4$,
and the Maxwell covariance
$C(q)=P_T(q)/|q|^2$ on nonzero momenta. Keep a fixed mass $m>0$
and a volume-independent depth law $\nu_0$. Write $A(x)=\sum_p A_p e^{ipx}$,
so that $\mathbb E[A_p A_{-p}^{\mathsf T}]=C(p)/V$.
At each $r$, let $K_{n,r}$ be the $n$th derivative per unit volume
of the log relative determinant at zero field, with its momentum delta
removed, one harmonic readout included at each mark,
and the power $e^n$ stripped. If the phase uses $e z_m$, retain
$z_m^n$ in $K_{n,r}$. Write $K_r(q)=K_{2,r}(q,-q)$ on transverse
modes. For $q=t e_1$, $t>0$, reflection and permutation of the
other three axes make this matrix scalar on the transverse space;
denote its eigenvalue by $k_r(t)$. The formulas below are for the
paired Dirac comparison with these symmetries.

No individual contact can be dropped from $K_{4,r}$. Explicitly, if
$Y_I$ is the derivative of the Wilson chord with the labelled marks
in $I$ (and the corresponding powers of $e$ stripped), the
four-linear kernel is the following finite sum on conserved external
momenta:
\begin{equation}
 K_{4,r}=\frac1V\sum_{a=1}^4\frac{(-1)^{a+1}}a
 \sum_{(I_1,\ldots,I_a)}
 \operatorname{Tr}(Y_{I_1}R_0\cdots Y_{I_a}R_0).
 \label{eq:spec-uv-four-chord}
\end{equation}
The inner sum runs over ordered partitions of the four labelled marks
into nonempty blocks. This is the fourth derivative of
$\operatorname{Tr}\log(1+Y R_0)$ at $Y=0$; in particular $a=1$
is the single four-mark contact. Define the total conditional vacuum
coefficient $v_r$ and the one-charged-loop contraction $u_r(t)$ by
\[
 \begin{aligned}
 v_r&=\tfrac12\sum_{p\ne0}\operatorname{tr}[C(p)K_r(p)],\\
 u_r(t)&=\frac1{2V}\sum_{p\ne0}
 C_{ab}(p) K_{4,r}^{22ab}(q,-q,p,-p),
 \qquad q=t e_1 .
 \end{aligned}
\]
Repeated polarization indices are summed. Each $K_r(p)$ in the
first formula is the transverse quadratic kernel at the indicated
momentum. The common vacuum coefficient $v_r$ is extensive.

For the momentum covariance $G$ normalized so that
$G(q)=V\langle A_q A_{-q}^{\mathsf T}\rangle_{\rm conn}$, put
$\Pi=C^{-1}-G^{-1}$ on the nonzero transverse subspace. At each
fixed $R$ its order-$e^4$ coefficient in polarization $e_2$ is
\begin{equation}
 [e^4]\Pi_{22}(t e_1)
 =\langle u_r(t)\rangle_0+
   \operatorname{Cov}_0(v_r,k_r(t))+
   \frac{\operatorname{Var}_0(k_r(t))}{t^2}.
 \label{eq:spec-uv-four-photon}
\end{equation}
Here and below coefficient extraction concerns the ordinary finite-order
expansion at zero charge. To verify the identity, shift the Maxwell
field by its mean $B=CJ$ in the source with current $J$. Take two
derivatives in $B$ of the logarithm of the shifted determinant average,
per unit volume, and call the resulting matrix $W$. There is no linear
term in the conditional log determinant: translation invariance
excludes every nonzero mode and inversion kills the constant test.
Wick contraction of its quartic term gives $\langle u_r\rangle_0$.
The variance of its quadratic term supplies
$\langle k_r(t)^2\rangle_0/t^2$ and
$\operatorname{Cov}_0(v_r,k_r(t))$; the latter also accounts for
the single overall vacuum normalization. Finally
$G=C+CWC$, and inversion through fourth order subtracts
$\langle k_r(t)\rangle_0^2/t^2$. This proves
\eqref{eq:spec-uv-four-photon}. The sign convention is fixed by the
definition of $\Pi$.

\begin{proposition}[A common-depth obstruction after one-loop matching]
\label{prop:spec-uv-soft-variance}
Suppose $0<r<1$, $\sigma^2=\kappa r^\alpha$ with $\kappa,\alpha>0$,
and $\nu_0$ has all inverse moments and charges every interval
$(0,\delta)$, $0<\delta<1$. Let $k_r^0$ be the continuum soft
quadratic coefficient at this depth, namely the $H_m(\sigma)$ above
divided by $e^2$, with the additional factor $z_m^2$ if used.
With $w=(1,-3,3,-1)$ the exact fourth-order coefficients satisfy
\begin{equation}
 \lim_{R\to\infty}\frac1{R^2}
 \sum_{j=1}^4 w_j [e^4]\Pi_{22}(j e_1/R)
  =\frac{25}{48}\operatorname{Var}_0(k_r^0)>0 .
 \label{eq:spec-uv-soft-variance}
\end{equation}
In particular tuning the mass to $\langle k_r^0\rangle_0=0$ still
leaves an unbounded sequence of proper fourth-order soft responses.
\end{proposition}
\begin{proof}
For fixed $m$, the marked-kernel estimates above give constants $C,N$,
independent of $R\ge1$, such that
\[
 |u_r(j/R)|+V^{-1}|v_r|+
  \max_{0\le a\le3}\sup_{0\le t\le4}
           |\partial_t^a k_r(t)|\le C(1+r^{-N}).
\]
For the ray derivatives, interpolate the two-mark Fourier formula
from lattice $t$ to real $t\ge0$: use the same loop grid, the shifted
massive inverse and the shifted Gaussian chord symbols. Each derivative
adds momentum derivatives of these symbols or a factor from
$e^{-2rt}$. At least one differentiated or undifferentiated Gaussian
on each charged loop still gives the translated lattice bound per
unit volume. On this ray the transverse projection is constant.
The photon sum per unit volume is bounded by
$C(1+r^{-N})$ by shell counting with
$|p|^{-2}e^{-2r|p|}$; the four-dimensional shell sum is uniformly
integrable at the origin. These estimates prove the displayed bounds
for all blocks in \eqref{eq:spec-uv-four-chord}.

The divided volume $V^{-1}$ in the vacuum bound matters. The
third finite difference annihilates the constant, linear and quadratic
Taylor terms of $k_r(t)$ at $0^+$. Its absolute covariance with
$v_r$ is therefore $O(V R^{-3})=O(R)$ after the depth integral.
The four single-loop terms are bounded independently of $R$.
For each fixed $r>0$, the normalized charged-momentum Riemann sums
at $t=j/R$ converge to their fixed-depth soft integral $k_r^0$.
The displayed bounds give convergence in $L^2(\nu_0)$ as well.
The last term of \eqref{eq:spec-uv-four-photon} now gives
\eqref{eq:spec-uv-soft-variance} because
$\sum_j w_j/j^2=25/48$. Finally $k_r^0$ is continuous for $r>0$
and is asymptotic to $-3/(16\pi^2\sigma^2)$ at the tip, times
$z_m^2$ if used. It is finite at each positive depth and unbounded
on sets of positive law, so its variance is strictly positive.
\end{proof}

This soft obstruction concerns the single depth shared by the full
source and its ordinary Maxwell normalization. Every coefficient in
\eqref{eq:spec-uv-four-photon} remains UV finite at fixed volume
and nonzero external mode. A deterministic depth has zero variance;
its fourth-order coefficient reduces to the explicit contraction
$u_r(t)$. The result concerns fixed Taylor coefficients and by itself
does not determine a resummed photon pole.

\paragraph{What changes when the normal degree of freedom propagates?}
A common depth can correlate two charged insertions arbitrarily far
apart. Local motion of the depth changes the momentum carried by
that correlation. The following calculation isolates the part of
a conditional charged determinant quadratic in the transverse field.
It gives an exact fourth-order test for this part, including the
vacuum weight. Other fourth-order marks of a native determinant
must be added with their signs.

Use a Euclidean torus of volume $V$, and write
$h(p)=e^{-a|p|}$ for a fixed harmonic readout, $a>0$. Let $\xi(x)$
be a real centered stationary coefficient of the quadratic response.
To derive the finite-box formula take $\xi$ bounded and the photon
zero mode fixed. Add $e^2\int \xi\,|hA^T|^2/2$ to the logarithm
of the Maxwell weight before normalizing. Put
$\xi_l=V^{-1}\int e^{-ilx}\xi(x)\,\mathrm d^4 x$,
$s_l=\mathbb E|\xi_l|^2$. For a unit transverse polarization
$\epsilon$ at $q\ne0$ set $P_p=I-pp^{\mathsf T}/|p|^2$ for $p\ne0$.
In the convention $\Pi=C^{-1}-G^{-1}$ used above,
\begin{equation}
\begin{aligned}
 \frac{[e^4]\Pi_{\epsilon}(q)}{h(q)^2}
 &=\sum_{p\ne0}s_{q-p}\frac{h(p)^2|P_p\epsilon|^2}{|p|^2}\\
 &\quad+\frac32\,s_0\sum_{p\ne0}\frac{h(p)^2}{|p|^2}.
\end{aligned}
\label{eq:closure-uv-depth-spectrum}
\end{equation}
Indeed the conditional precision is $C^{-1}-e^2K_\xi$ with
$K_{\xi,pq}=h(p)h(q)\xi_{p-q}P_pP_q$ on the transverse space.
For sufficiently small charge at fixed box it is positive. The
relative Gaussian normalization is the ordinary determinant of
$I-e^2C^{1/2}K_\xi C^{1/2}$ to the power $-1/2$; the perturbation
is trace class by the two readouts. Its quadratic log coefficient
is $\tfrac32\xi_0\sum_{p\ne0}h(p)^2/|p|^2$. Expanding the conditional
inverse and averaging with this very determinant yields
\eqref{eq:closure-uv-depth-spectrum}. The mean first response is
zero. The second line is the covariance with the conditional vacuum.
In particular both displayed terms are nonnegative.

For a single centered random constant one recovers the
$\operatorname{Var}(\xi)/|q|^2$ soft behavior of the first line.
For motion in the carrier direction only, consider a normal radial
Hamiltonian with normalized ground state $u_0$, energy $E_0$ and an
excited-state gap $g>0$. The centered conditional response
$\xi(r)=k_r^0-\langle k_r^0\rangle$ has $\xi u_0\in L^2$ under
the inverse-moment bounds above. Define
$\nu(I)=\|P_{H-E_0}(I)\xi u_0\|^2$ on Borel sets $I$. Its support is in
$[g,\infty)$ for a simple ground state. The two ordered insertions
in this state have connected correlation
$\int e^{-\lambda|t|}\nu(\mathrm d\lambda)$. Its Fourier density is
$\chi(\omega)=\int 2\lambda(\lambda^2+\omega^2)^{-1}\nu(\mathrm d\lambda)$.
This last calculation also holds for unbounded $\xi$ with the stated
second moment, by spectral truncation of its matrix element.

If all spatial points still share this degree of freedom, no spatial
momentum is exchanged between the response marks. In continuous
carrier frequency, the first line of
\eqref{eq:closure-uv-depth-spectrum} at a purely spatial transfer
of magnitude $Q$ becomes
\[
 B(Q)=h(Q)^2\int_{\mathbb R}
  \frac{h\!\left(\sqrt{Q^2+\omega^2}\right)^2\chi(\omega)}
       {Q^2+\omega^2}\frac{\mathrm d\omega}{2\pi}.
\]
Here $h(t)=e^{-at}$ for a scalar argument, and the measured spatial
polarization is perpendicular to the transfer. With $\omega=Qz$,
dominated convergence gives
$\lim_{Q\downarrow0} Q B(Q)=\chi(0)/2
=\int\lambda^{-1}\nu(\mathrm d\lambda)$.
Thus exponential decorrelation along the carrier still leaves an
inverse first power of the spatial momentum whenever $\xi u_0\ne0$.
This is a bridge coefficient for the specified placement of the
response marks; promoting a depth to a normal operator requires
deriving those marks anew in the full action.

For comparison, take an isotropic, four-dimensionally varying response
with spectra $s_l=V^{-1}S(l)$ on the expanding tori, where $S$ is
continuous, bounded, nonnegative and integrable. Ordinary momentum
sums away from zero and shell counting at zero then give the
following soft limit along a coordinate ray:
\[
 \int_{\mathbb R^4}\frac{h(p)^2}{|p|^2}
       \left[\frac34 S(p)+\frac32 S(0)\right]
                \frac{\mathrm d^4p}{(2\pi)^4}.
\]
The angular average of $|P_p\epsilon|^2$ is $3/4$. The result is
finite and strictly positive for nonzero $S$. For example, a bounded
stationary coefficient $\cos(l\cdot x+\theta)$ with uniform phase
and a random reciprocal wave vector sampled from a normalized
isotropic Gaussian has precisely this continuum limit, with the
corresponding Gaussian $S$; the finite-box normalization tends to
its continuum value. Spreading the response over four directions
can remove the inverse-power divergence and still leave a constant
in the inverse photon response. A Ward identity, derivative coupling,
or compensating fourth-order mark of the actual action must account
for that constant before masslessness can follow.

\paragraph{An interacting spectral test on the relative Wilson chord.}
The full open covariance can be tested on the Euclidean frequencies of
the torus. Fix a sector plane, $m>0$, and zero spatial charged
momentum. Write $\omega$ for the real frequency along the Euclidean
carrier, and $\gamma_\parallel$ for its Hermitian Dirac matrix.
Use periodic spatial spinors. In the carrier direction one may
take antiperiodic spinors for the thermal interpretation below.
The resulting signs on winding lifts have unit modulus and preserve the
kernel and massive-inverse estimates. Keep the same transverse Maxwell
field, including its fixed zero mode.

Write $Y_b=T_b-T_0$, $R_b=(D_{\rm D}+Y_b)^{-1}$ and
$R_0=D_{\rm D}^{-1}$. Let $\overline R$ be the expectation of $R_b$
under the normalized, positive relative determinant and the original
Maxwell and depth law. This is the two-point function of the retained
charged source in the transverse variables, at zero external field.
The law is translation invariant, so $\overline R$ is diagonal in
momentum. Use either fixed positive widths or the strict radial regime
$\sigma^2=\kappa r$, bounded upper widths, and density near zero
bounded by $Cr^4\exp(-A/r^3)$, $A,\kappa>0$.

\begin{proposition}[Positive spectral weights would force a free open covariance]
\label{prop:uv-relative-spectrum}
At every nonzero real transverse coupling, the rest-fibre matrices
$\overline R(\omega,\mathbf0)\gamma_\parallel$ of this relative chord
cannot have the unsubtracted representation
\begin{equation}
 \overline R(\omega,\mathbf0)\gamma_\parallel
    =\int_{\mathbb R}\frac{M(\mathrm dE)}{i\omega+E},
 \qquad M\geq0,
 \label{eq:uv-relative-spectral-law}
\end{equation}
on all nonzero carrier frequencies. Here $M$ is a positive
matrix measure, including both signs of $E$. This test also applies
to a fixed nonzero rescaling of the coupling inside the phase.
\end{proposition}
\begin{proof}
First, the interacting covariance has the free asymptotics to
every inverse power on the real frequency sequence. For a unit
normalized spinor plane wave $v_\omega$, integration by parts in
the incoming variable of the smooth Wilson kernel gives, for each
integer $N\geq0$,
\[
 \|Y_b v_\omega\|\leq C_N(b,r)(1+|\omega|)^{-N}.
\]
The constants can be taken polynomial in inverse widths and finitely
many uniform derivatives of the harmonic readout. Indeed derivatives
of a real Wilson phase produce powers of the chord displacement
and of these derivatives. Gaussian integration of each displacement
power and each differentiated kernel costs at most a power of
$\sigma^{-1}$ for bounded widths. Maxwell moments of the readout
derivatives cost powers of $r^{-1}$. The field-uniform determinant
bound $\exp(C\sigma^{-4})$ leaves every such moment integrable
in the stated strict radial regime. The same moments are finite at
fixed widths. The resolvent identity
$R_b-R_0=-R_bY_bR_0$, the mass bound $\|R_b\|\le m^{-1}$, and
the free Fourier multiplier now give, for every $N$,
\[
 \overline R(\omega,\mathbf0)\gamma_\parallel
  =(i\omega+m\gamma_\parallel)^{-1}+O_N(|\omega|^{-N}).
\]

To see what positivity would imply, take a unit eigenvector $v$ of
$m\gamma_\parallel$ of eigenvalue $\lambda=\pm m$. Set
$\mu=v^*Mv$ and call the corresponding diagonal coefficient $g$.
Along positive frequencies tending to infinity,
$-\omega\operatorname{Im}g=\int\omega^2/(\omega^2+E^2)\,\mathrm d\mu$
tends to one, so $\mu$ has mass one by monotone convergence. Next,
\[
 \omega^2(1+\omega\operatorname{Im}g)
   =\int\frac{\omega^2 E^2}{\omega^2+E^2}\,\mathrm d\mu
   \longrightarrow\lambda^2 .
\]
This gives a finite second moment. Dominated convergence applied to
$\omega^2\operatorname{Re}g$ then gives first moment $\lambda$.
Thus $\int(E-\lambda)^2\,\mathrm d\mu=0$.
Every vector in the same eigenspace has its measure concentrated at
that eigenvalue. Polarization there and Cauchy--Schwarz for $M$ on
each Borel set show that the two opposite eigenspaces have zero cross
measure. Consequently \eqref{eq:uv-relative-spectral-law} would equal
the free resolvent on the entire rest fibre.

The charged operator supplies a strict test of that conclusion.
For each realization $D_{\rm D}+Y_b=m+K_b$ with $K_b$ skew-adjoint,
so $\operatorname{Re}R_b=mR_b^*R_b$. The free inverse obeys the
same identity. For any incoming Fourier spinor $f$, equality
$\overline R f=R_0 f$ would imply
\[
 \mathbb E\|R_b f-R_0 f\|^2
   =m^{-1}\operatorname{Re}\langle f,\overline R f\rangle
          -\|R_0 f\|^2=0 .
\]
In particular $Y_b R_0 f=0$ almost surely. Choose a nonzero rest-fibre
spinor at any nonzero incoming $\omega$. Vary one real Maxwell cosine
mode of frequency $p\ne0$ along the carrier with spatial polarization
$a$. In the Fourier component at $\omega+p$, the first variation of
$Y_b$ at zero Maxwell field is a nonzero multiple of
\[
 i e\, e^{-r|p|}\gamma_a
     \int_0^1\exp[-\sigma^2(\omega+up)^2]\,\mathrm du .
\]
The cosine contributes a factor $1/2$. The integral is strictly
positive and $\gamma_a$ and $R_0(\omega,\mathbf0)$ are invertible.
A sufficiently small nonzero amplitude therefore gives
$Y_bR_0 f\ne0$. This persists on an open set of readouts in the
uniform norm by continuity of the Wilson kernel. The Gaussian
law of the harmonic readout gives this set positive mass at each
positive depth. Its support contains every transverse trigonometric
polynomial. Finitely many independent modes can approximate the
prescribed coefficients, and the remaining tail can be made uniformly
small with positive probability by its absolute Fourier moment bound.
The strictly positive determinant preserves this positive probability
under normalization. This contradicts the variance identity.
\end{proof}

The obstruction can also be detected with two finite carrier
frequencies. For a unit rest eigenspinor of eigenvalue $\lambda=\pm m$,
write $g(\omega)=v^*\overline R(\omega,\mathbf0)\gamma_\parallel v$.
After the mass-one normalization derived above, any positive measure
in \eqref{eq:uv-relative-spectral-law} would obey, for $\Omega\geq\omega>0$,
\[
 \big|(i\omega+\lambda)g(\omega)-1\big|^2
     \leq \omega^{-2} V_\Omega ,
\]
where
\[
 \begin{aligned}
 V_\Omega={}&\Omega^2(1+\Omega\operatorname{Im} g(\Omega))
     -2\lambda\Omega^2\operatorname{Re}g(\Omega)\\
     &-\lambda^2\Omega\operatorname{Im}g(\Omega).
 \end{aligned}
\]
Under that representation $V_\Omega$ is the integral of
$\Omega^2(E-\lambda)^2/(\Omega^2+E^2)$. Apply Cauchy--Schwarz
to the integral of $(\lambda-E)/(i\omega+E)$ and use
$(\Omega^2+E^2)/(\Omega^2(\omega^2+E^2))\leq\omega^{-2}$.
For the actual chord, $V_\Omega\to0$ by the high-frequency estimate.
The left side is strictly positive at each fixed $\omega>0$.
Indeed $R_0 f$ is a scalar multiple of $f$ on this incoming spinor,
so equality of its single diagonal coefficient to the free value
already gives $\mathbb E\|R_bf-R_0 f\|^2=0$ by expanding the square.
The transverse mark above excludes this equality. Thus sufficiently
large finite $\Omega$ violate the inequality.

For a thermal positive-Hilbert-space charged field paired with its
adjoint, the Lehmann sum gives the form in
\eqref{eq:uv-relative-spectral-law}. Each transition contributes
the outer product of its matrix elements with the positive sum
of the two Boltzmann weights, divided by the partition function.
The choice of sign for $E$ above agrees with the free Euclidean
inverse. This obstruction concerns the interacting charged readout
specified by the relative chord. It does not require analytic
interpolation of discrete frequencies or an assumption about remote
complex poles. A different physical endpoint observable requires
testing its own covariance. Positivity of the gauge-invariant even
algebra and the infinite-volume limit are separate questions.

\paragraph{What a positive charged boundary channel contributes.}
The failure of the relative chord suggests testing the dynamics of an
eliminated boundary channel directly. Fix a spatial momentum on the
selected sector plane and let $z=i\omega$ denote the Euclidean carrier
frequency. The following result concerns the charged two-point
function with its physical adjoint. Its matrix spectral measure can
have support at both signs of energy.

\begin{proposition}[Positive boundary realization and the charge vertex]
\label{prop:closing-uv-positive-boundary}
Let $G(z)$ on the imaginary axis be the charged matrix of a finite
set of source slots. A representation
\[
 \begin{gathered}
 G(z)=\int_{\mathbb R}(z+E)^{-1}M(\mathrm dE),\\
 M\ge0,\quad \int M=I,\quad \int E^2\,\operatorname{tr}M<\infty ,
 \end{gathered}
\]
is equivalent to the existence of a Hermitian matrix $A$, a
self-adjoint boundary operator $B$, and a bounded map $T$ from the
slots to the boundary space such that, for nonreal $z$,
\[
 D(z)=z+A-T^*(z+B)^{-1}T,\qquad G(z)=D(z)^{-1}.
\]
For data given only at Matsubara frequencies, equivalence means that
the right-hand side agrees at those frequencies. The first moment
is $A$, and the second moment is $A^2+T^*T$. If both sets of slots
carry the same conserved charge $q$, the uniform temporal insertion
of their common first-order kinetic action gives
\begin{equation}
 \Gamma(z',z)=q\left[I+T^*(z'+B)^{-1}(z+B)^{-1}T\right].
 \label{eq:closing-uv-temporal-mark}
\end{equation}
Here the scalar source is in energy units, with $q$ factored out
when defining a unit temporal variation.
\end{proposition}
\begin{proof}
For the reverse construction from $M$, complete vector-valued simple
functions in the norm $\int f^*M(\mathrm dE)f$, quotienting null
vectors. Multiplication by $E$ gives a self-adjoint operator $H$,
and the map $J$ sending a slot vector to its constant function is an
isometry. The second moment puts the range of $J$ in $D(H)$.
Writing $P=JJ^*$, take $A=J^*HJ$ and $T=(1-P)HJ$, which is
bounded since the slot space has finite dimension. Subtract the
bounded self-adjoint off-diagonal blocks $T,T^*$ from $H$.
The resulting operator reduces $P$ since $P D(H)\subset D(H)$ and
the off-diagonal blocks have been removed. Its restriction to
$(1-P)$ is self-adjoint; call it $B$. Block elimination now gives
$G=J^*(z+H)^{-1}J=D^{-1}$. Conversely take the spectral resolution
of $H=\left(\begin{smallmatrix}A&T^*\\T&B\end{smallmatrix}\right)$
on the direct sum. Its compression gives a positive normalized
measure, with second moment $(HJ)^*HJ=A^2+T^*T$ in the form sense.
To obtain the current, insert $qI$ in both diagonal kinetic blocks
before eliminating $B$. Differentiating its inverse gives
\eqref{eq:closing-uv-temporal-mark}. In particular
$(z'-z)\Gamma=q[D(z')-D(z)]$.
\end{proof}

For a simple pole at $z=-E_0$ away from $-\operatorname{spec}B$,
let $v$ be a normalized null vector of $D(-E_0)$. Its slot residue
along $v$ is
$Z=[1+\|(B-E_0)^{-1}Tv\|^2]^{-1}$. Thus
$Z\langle v,\Gamma(-E_0,-E_0)v\rangle=q$. The boundary term
carries exactly the charge missing from the slot residue. Likewise a
frequency-independent normalized slot readout with $T\ne0$ has a
strictly nonzero second-moment excess on some vector. Agreement with
all the free high-frequency moments would force $T=0$ on that
subspace. The two-point construction is an existence statement
about a positive transfer channel. It does not establish that its
boundary determinant, the photon current, or higher-point functions
are those of the native theory.

\paragraph{Two charged bands with a spatial Wilson link.}
There is an explicit source in which both the temporal charge vertex
and a positive pair response can be checked. At zero temperature
give each of two fermion slots the same electric charge $q$.
Use a local carrier-time kinetic term, fill the negative band, and
choose the one-body spatial operator at zero field to be
\[
 \begin{gathered}
 h(k)=m\sigma_3+t(k)\sigma_1,\\
 t(k)=g\exp(-\sigma^2|k|^2),\qquad E_k=\sqrt{m^2+t(k)^2},
 \end{gathered}
\]
where $m,\sigma>0$, $g$ is real, $k\in\mathbb R^3$, and the
$\sigma_i$ are Pauli matrices on the two slots. The off-diagonal
kernel is the spatial heat displacement of lapse $\sigma^2$.
Multiply it by the straight Wilson phase with charge $q$ and
couple the temporal potential on both slots. For a smooth external
connection one may smear \emph{all} components at fixed time by
the same unit-mass spatial kernel. Gauge transformation then acts on
both endpoint slots by the phase of the spatially smeared gauge
parameter. Its zero-spatial-transfer charge stays $q$.
Eliminating the lower slot of this source gives $A=m$, $B=-m$,
$T=t(k)$. Its charged spectral weights at $+E_k,-E_k$ are
$(1+m/E_k)/2,(1-m/E_k)/2$. The temporal vertex of
\eqref{eq:closing-uv-temporal-mark} has the correct residue on
both poles when $t(k)\ne0$.

A spatially uniform external vector potential shifts the symbol
by $k\mapsto k+q\mathbf A$ in the chosen orientation.
Write $t_i=\partial_{k_i}t$ and $t_{ii}=\partial_{k_i}^2 t$.
The direct and two-insertion responses follow from the same filled
Hamiltonian. At each $k$ the squared current matrix element from
the occupied to the empty band is
\[
 W_i(k)=q^2t_i(k)^2\,m^2/E_k^2 .
\]
Indeed the diagonal matrix element of $\sigma_1$ in the lower band
is $-t/E_k$ and $\sigma_1^2=I$. At nonzero $qg$ the positive
pair measure, with squared excitation energy $s$, is
\[
 \mathrm d\rho_i(s)=
 \int\frac{\mathrm d^3k}{(2\pi)^3}\,4E_kW_i(k)
                   \,\delta_{4E_k^2}(\mathrm ds).
\]
Its Stieltjes transform is the paramagnetic current correlation:
the transition energy is $2E_k$ and the two time orderings supply
$4E_k/(x+4E_k^2)$ at squared frequency $x$. The seagull is
$C_i=-q^2\int t\,t_{ii}/E_k\,\mathrm d^3k/(2\pi)^3$.
Integration by parts gives exactly $C_i=\int\mathrm d\rho_i(s)/s$.
It is finite, since all derivatives of $t$ have Gaussian decay.
Thus the complete uniform proper response is
\begin{equation}
 K_i(x)=C_i-\int\frac{\mathrm d\rho_i(s)}{x+s}
       = x\int\frac{\mathrm d\rho_i(s)}{s(x+s)}.
 \label{eq:closing-uv-optical-ward}
\end{equation}
It vanishes in the static limit and has nonzero positive optical
spectral weight. This is the zero-temperature continuum matter
response to an external photon. These bounded bands do not have
relativistic Dirac dispersion. Even in a finite spatial box they
have infinitely many modes of bounded energy, so the proof does not
use an absolute thermal Gibbs trace for the two-band matter.

\paragraph{A hard-frequency test for the electric density.}
Positive time transfer and optical matching give no automatic bound
on an absolute four-momentum integral. At zero external frequency
and fixed nonzero spatial transfer $p$, consider a finite-band
time-local operator, with massive invertible Hermitian free symbol
$h(k)$, linear temporal coupling, charge vertex $Q=Q(k+p,k)$ and
reverse vertex $Q^*$. Suppose that no quadratic temporal contact or
boundary term accompanies this two-mark coefficient. Put
$R_k=(i\omega+h(k))^{-1}$ and
$\Lambda_k=\mu+\|h(k)\|+\|h(k+p)\|$ for a fixed energy unit $\mu>0$.
On $|\omega|\ge2\Lambda_k$, the matrices $i\omega R_k$ are
contractions that differ from $I$ by at most $\|h(k)\|/|\omega|$.
In the trace of their product with $Q,Q^*$ use
$|\operatorname{tr}(UQVQ^*)|\le\|U\|\|V\|\|Q\|_2^2$, where
$\|\cdot\|_2$ is the Hilbert--Schmidt norm. Subtracting the identity
in one resolvent factor at a time proves
\[
 \operatorname{Re}\{-\omega^2\operatorname{tr}
      (R_{k+p}Q R_kQ^*)\}\ \ge\ \tfrac12\|Q\|_2^2 .
\]
Integration on this hard-frequency region gives the necessary condition
\begin{equation}
 \int_{\mathbb R^3}
 \frac{\|Q(k+p,k)\|_2^2}{\mu+\|h(k)\|+\|h(k+p)\|}\,\mathrm d^3k
 <\infty
 \label{eq:closing-uv-density-hard}
\end{equation}
for joint absolute convergence of this single trace. The same
argument gives a spatial sum in a box. For the two-band example
$Q=q\widehat w(p) I$, so the individual trace fails the test
at every transfer where the readout is nonzero. The filled vacuum also
fixes the onsite term in the complete density coefficient, derived in
Proposition~\ref{prop:native-relative-density-uv}. More generally,
a nondecaying charge on an exterior cone fails when both operator
norms grow there at most as $(1+|k|)^3$, including the borderline
logarithm. The cone condition refers to a positive lower bound
on $\|Q\|_2$ throughout that cone.

The finite answer in ordered time can be made explicit. Put
$n(k)=(t(k),0,m)/E_k$ and $D_{k,p}=E_k+E_{k+p}$.
The occupied-to-empty density matrix element has squared norm
$q^2|\widehat w(p)|^2|n(k+p)-n(k)|^2/4$ by the Pauli projectors.
At squared carrier frequency $x$ its contribution to the continuum
density correlation is this norm times
$2D_{k,p}/(x+D_{k,p}^2)$, integrated with
$\mathrm d^3 k/(2\pi)^3$. Inversion symmetry pairs the two time
orderings. Since $2m\le D_{k,p}\le2\sqrt{m^2+g^2}$ and
$\nabla n\in L^2$, the translation bound
$\int |n(k+p)-n(k)|^2\,\mathrm d^3k
 \le |p|^2\int|\nabla n|^2\,\mathrm d^3k$
gives $0\le G_{00}(x,p)\le C|\widehat w(p)|^2|p|^2/(x+4m^2)$.
It is strictly positive for $qg\widehat w(p)\ne0$ and $p\ne0$.
This spectral formula determines the finite physical density response.
Its representation with the charged frequency still present depends on
the vacuum determinant used to generate both insertions.

In the bounded-band example the limiting flat-band trace at fixed
$\omega$ is proportional to
$2(m^2-\omega^2)/(m^2+\omega^2)^2$, whose frequency integral is
zero but whose absolute integral is positive. The normalized
two-band source supplies this flat-band trace as its onsite reference
at the same frequency and momentum. Equation~\eqref{eq:native-relative-bubble-uv}
exhibits the resulting integrable density coefficient. In a GCM
representation the charged action and boundary state must determine
whether a reference derivative, temporal contact or other term
accompanies the two density marks.

\subsection{Local charge, spatial form factors and instantaneous Gauss exchange}
\label{sec:final-uv-gauss}
Work on one locked sector plane. There are two ways a spatial
form factor can enter the density seen by a photon. A change of basis
can move charge among the slots that a detector couples to. A
polarization term in the action can change the density that couples
to an electric field. These operations have different implications
for a complete charged graph.

\begin{proposition}[Momentum dependence of a commuting local charge]
\label{prop:final-uv-charge-rigidity}
In a fixed nonzero charge sector let the one-particle slots be
$L^2(\mathbb R^3;\mathbb C^d)$, $1\le d<\infty$. Suppose that a
field-independent Abelian gauge action has commuting bounded charge
modes $T_p$ for spatial plane waves, with $T_{-p}=T_p^*$. Translation
covariance gives the form
\[
 (T_p\psi)(k+p)=K_p(k)\psi(k).
\]
Assume the matrix representatives are continuous in $(p,k)$ and
$K_0(k)=qI$, $q\ne0$. At each fixed $p$ the eigenvalues of
$K_p(k)^*K_p(k)$, counted with multiplicity, are independent of
$k$. In particular the squared norm of the inclusive temporal
vertex, summed over all incoming and outgoing slots, is a strictly
positive constant in $k$ for all $p$ sufficiently close to zero.
\end{proposition}
\begin{proof}
$T_p^*T_p$ acts by the positive multiplication matrix
$M_p(k)=K_p(k)^*K_p(k)$. Since $T_h$ commutes with $T_p$ and $T_p^*$,
\[
 M_p(k+h)K_h(k)=K_h(k)M_p(k).
\]
The continuous representatives turn this operator equality into a
pointwise one for each pair of transfers. At any fixed $k$,
$K_h(k)$ is invertible for every sufficiently small $h$, by its
value $qI$ at zero. Thus the two positive matrices are similar,
and their ordered eigenvalues are locally constant in $k$.
Connectedness gives the assertion. At $k=0$, continuity in $p$
gives $\operatorname{tr}M_p(0)\to d q^2$ as $p\to0$.
\end{proof}

For a time-local, linear minimal coupling to these generators the
hypotheses of \eqref{eq:closing-uv-density-hard} therefore exclude
joint absolute convergence whenever the two symbol norms in its
denominator grow at most cubically on an exterior cone. This
covers a momentum-dependent unitary choice of charged basis and
any finite enlargement of the slots obeying the same hypotheses.
Compressing the vertex to a selected slot omits its transitions to
the complementary slots. For instance, for a sharp position
observable $X$ on a larger Hilbert space, $e^{ip\cdot X}$ is
unitary. For a unit slot vector $v$, the sum of squared amplitudes
$|\langle e_n,e^{ip\cdot X}v\rangle|^2$ over any full orthonormal
basis is one. Dispersion and the energies of these intermediate
states still enter the loop test.

A polarization current illustrates the scope of the proposition.
Use the Fourier convention $F_{i0}=ip_i A_0-i\nu A_i$ at transfer
$(\nu,p)$. At carrier momentum $k$ set
\[
 W(k,p)=\exp\!\left[-a^2|p|^2
               \bigl(1+\ell^2|k+p/2|^2\bigr)\right],
 \qquad a,\ell>0.
\]
The candidate density $qW(k,p)$ satisfies the adjoint rule, but its
squared modulus at $k$ and $k-p$ differs in general. The two shifts
$T_p,T_p^*$ already fail to commute. One can instead add to an
action a charged bilinear times $F_{i0}$, with matrix element
$P_i=iq p_i(1-W)/|p|^2$ at zero spatial gauge field. The limit at
$p=0$ is zero for fixed $k$. The relation
$P_i(k+p,-p)=\overline{P_i(k,p)}$ gives the adjoint rule.
For smooth external fields and test wavefunctions the spatial
bilinear can be covariantized by transporting between its two
charged endpoints at equal time. Its product with $F_{i0}$ is
then invariant, and its linear density contribution is
$i\sum_i p_iP_i=q(W-1)$. The same term supplies the spatial vertex
$-i\nu P_i$. The effective density becomes $qW$ while the original
gauge generators continue to commute. With a uniformly gapped
free symbol the fixed-$p\ne0$ temporal bubble with these two marks
is absolutely integrable at zero external frequency: the two resolvents are bounded by
$(\omega^2+m^2)^{-1}$ in product norm and $W(\cdot,p)^2$ is
integrable in $k$. This example specifies an additional interaction.
Deriving such a polarization, its higher variations and its scales
from the chiral curvature pairing is a separate physical question.

\paragraph{A carrier-frequency obstruction on an open charged line.}
Solving Gauss's law on each carrier-time slice gives an
instantaneous Coulomb covariance $C_{00}(p)>0$ independent of
carrier frequency at nonzero $p$. Consider the one-exchange
self-energy between charged endpoints. At a fixed external
momentum $k$ its temporal factor in the ordinary spacetime routing
is
\[
 C_{00}(p)\,Q_p^*G_{k+p}(iu)Q_p .
\]
Here $u$ is the internal charged frequency, shifted from the
photon frequency by the external one. $Q_p=Q(k+p,k)$ is a
frequency-independent vertex. Take this expression when it is the
complete temporal contribution in this routing, with the reverse
vertex its physical adjoint and no coincident-time or boundary
term accompanying the exchange.

\begin{proposition}[Instantaneous electric exchange against a positive charged measure]
\label{prop:final-uv-gauss-frequency}
Suppose the free charged covariance at the intermediate momentum
has a positive matrix or operator measure,
\[
 G_{k+p}(iu)=\int_{\mathbb R}\frac{M_{k+p}(\mathrm dE)}{iu+E}.
\]
For a fixed external slot $v$ suppose $w_p=Q_p v$ has finite,
strictly positive total spectral weight on a positive-measure set
of transfers $p$ with $C_{00}(p)>0$. The displayed exchange is not
absolutely integrable even in this matrix element. Arbitrarily
strong suppression at large spatial momentum leaves this
conclusion valid.
\end{proposition}
\begin{proof}
Put $\mu_p(\mathrm dE)=\langle w_p,M_{k+p}(\mathrm dE)w_p\rangle$.
For $u>0$,
\[
 -u\operatorname{Im}\langle w_p,G_{k+p}(iu)w_p\rangle
   =\int_{\mathbb R}\frac{u^2}{u^2+E^2}\,\mu_p(\mathrm dE)
   \ \uparrow\ \mu_p(\mathbb R).
\]
At each such $p$ the modulus is at least
$\mu_p(\mathbb R)/(2u)$ for sufficiently large $u$. Tonelli's
theorem gives divergence of the absolute exchange on any bounded
subset of this set of positive measure. No uniform threshold in
$p$ or finite-dimensional bound on the intermediate energy is
needed. The divergence is present with external charged endpoints.
\end{proof}

The order of operations matters here. For a frequency-independent
vertex the symmetric combination at fixed $p$ is exactly
\[
 \frac12\langle w_p,[G_{k+p}(iu)+G_{k+p}(-iu)]w_p\rangle
   =\int_{\mathbb R}\frac{E}{u^2+E^2}\,\mu_p(\mathrm dE).
\]
Its absolute integral over $u>0$ is at most
$\pi\mu_p(\mathbb R)/2$, since
$\int_0^\infty |E|/(u^2+E^2)\,\mathrm du=\pi/2$ for $E\ne0$ and
the integrand vanishes for $E=0$. This bound even dispenses with
a finite first energy moment. It explains how an integration rule
or an endpoint time-ordering prescribed by an action can give a
finite electric response. The oriented absolute integral and the
absolute integral of the symmetric combination are different
questions; spatial integrability of the latter is still required.

There is also a two-frequency test of the closed vacuum
contraction. For finite bands with self-adjoint symbols, use
$R_k(iu)=(iu+h(k))^{-1}$. The independent original carrier
energies can be taken as $u,v=u+\nu$, with unit Jacobian.
On a bounded spatial region and for
$u,v\ge4(1+\|h(k)\|+\|h(k+p)\|)$, in fixed energy units,
\[
 \operatorname{Re}\bigl[-uv\operatorname{tr}
       (R_{k+p}(iv)Q_p R_k(iu)Q_p^*)\bigr]
       \ge \tfrac12\|Q_p\|_{\rm HS}^2.
\]
Indeed $iuR_k(iu)$ is a contraction and its difference from the
identity has norm at most $\|h(k)\|/u$; use the two bounds in
turn under the trace. Instantaneous sewing then leaves
$\int^\infty \mathrm du/u\int^\infty \mathrm dv/v$ wherever the
vertex has positive squared norm. This estimate concerns the oriented trace before any source-derived
vacuum normalization, symmetrization or equal-time assignment.

Neither oriented example tests a gauge-invariant observable until all
terms required by its source have been assembled. Both concern absolute
convergence in the specified instantaneous four-momentum routing. Exact ordered carrier propagation
or a contact derived from the same source can change the complete
integrand; the ordered Wilson calculation in
Proposition~\ref{prop:closing-uv-ordered-Maxwell} illustrates the
first possibility. A formulation with a propagating temporal gauge
field also has a different frequency kernel and must combine all
its gauge and boundary contributions in its own routing. For a
frequency-dependent adjoint pair $Q_p(u)^*,Q_p(u)$ the open-line
proof itself supplies a useful necessary test: writing
$w_p(u)=Q_p(u)v$, absolute convergence requires
\begin{equation}
 \int\! \mathrm d^3p\, C_{00}(p)\int_1^\infty\!
 \frac{\mathrm du}{u}\,
 \langle w_p(u),M_{k+p}([-u,u])w_p(u)\rangle<\infty .
 \label{eq:final-uv-gauss-window}
\end{equation}
On this energy interval the factor $u^2/(u^2+E^2)$ is at
least $1/2$. Thus the joint bound tests the spectral weight
actually reached by the temporal insertion, including every
accessible boundary channel. A spatial form factor measured at
one nonzero transfer cannot certify it.

\paragraph{Which temporal contacts can remove a hard density loop?}
A local contact can cancel a loop when both arise from an exact
change of variables. To decide whether it cancels an electric
density insertion, one can inspect the poles of the complete
coefficient before integrating over the carrier frequency.
Work on the selected Euclidean sector plane and use energy units
for a scalar source. Two factors of $i$ in the unitary connection
amount to a common sign convention for the mixed variation.

At zero external frequency and spatial transfer $p$, let $Q_k$
map the finite-dimensional charged slot at $k$ to the one at
$k+p$. Its reverse mark is $Q_k^*$. Write $C_k$ for the mixed
second derivative of the same first-order charged operator at $k$.
Suppose this contact is independent of the internal carrier frequency.
For $R_k(z)=(z+h(k))^{-1}$ the mixed derivative of the ordinary
log determinant has the common routing
\begin{equation}
 \begin{split}
 K_k(z)={}&\operatorname{tr}(R_k(z)C_k)\\
 &-\operatorname{tr}\!
       \left(R_k(z)Q_k^*R_{k+p}(z)Q_k\right).
 \end{split}
 \label{eq:ct-uv-complete-density}
\end{equation}
This follows by differentiating the determinant twice. A
source-dependent boundary factor would add its own mixed logarithmic
derivative. A source-independent measure has zero derivative here.

\begin{proposition}[An equal-energy obstruction including temporal contacts]
\label{prop:ct-uv-temporal-residue}
Consider $k=t n$ along a set of unit directions of positive solid
angle. Suppose that for each such $n$ there is an energy scale
$\lambda=\lambda(t,n)>0$ for which both Hermitian matrices
$h(k)/\lambda$ and $h(k+p)/\lambda$ converge to the same
$H=H(n)$. In fixed orthonormal slot coordinates suppose
$Q_k\to Q$ and $\lambda\|C_k\|$ stays bounded. If some spectral
projection $P_E$ of $H$ satisfies $P_EQP_E\ne0$ on each such
ray, then
\[
 \liminf_{t\to\infty}\lambda
   \int_{\lambda}^{2\lambda}|K_{tn}(i\omega)|\,\mathrm d\omega>0
\]
on that ray. Hence the original four-momentum integral fails
absolute convergence whenever
$\int^{\infty} t^2\mathrm dt/\lambda(t,n)=\infty$
on that set. This includes growth at most cubic in the spatial
momentum. The conclusion also holds for a finite sum of ordinary
fermion log determinants satisfying these scaling assumptions with
a common $\lambda$, if at least one equal-energy charge survives.
\end{proposition}
\begin{proof}
From any sequence escaping along a fixed ray extract a subsequence
on which $\lambda C_k\to C$. The rescaled
$\lambda^2 K_k(\lambda z)$ converges uniformly for
$z=iu$, $1\le u\le2$, to the rational function
\[
 F(z)=\operatorname{tr}(R C)
          -\operatorname{tr}(R Q^* R Q),
 \qquad R=(z+H)^{-1}.
\]
At $z=-E$ its double-pole coefficient is
$-\|P_EQP_E\|_{\mathrm{HS}}^2$. Terms involving different
eigenvalues give simple poles, as does every term of the contact.
Thus $F$ cannot vanish throughout the indicated segment. If the
claimed lower limit were zero, a minimizing sequence and the
subsequence just constructed would give
$\int_1^2|F(iu)|\,\mathrm du=0$, a contradiction.
Tonelli's theorem in polar coordinates gives the divergence.
For several ordinary fermion determinants the squared norms at
each common pole add, including for opposite electric charges.
The argument also allows a sum of bounded rescaled contacts routed
at $k$ and at $k+p$.
\end{proof}

One can also decide exactly when the entire limiting coefficient
vanishes. Write $a_{EF}=\|P_EQP_F\|_{\mathrm{HS}}^2$.
In addition to $a_{EE}=0$ for every $E$, one needs
\begin{equation}
 \operatorname{tr}(P_E C)=
       \sum_{F\ne E}\frac{a_{EF}+a_{FE}}{F-E}
       \qquad\hbox{for every }E .
 \label{eq:ct-uv-simple-residues}
\end{equation}
Partial fractions give the right side as the simple-pole residue
of the bubble. These conditions are sufficient for $F(z)=0$ as
well, since the rational function tends to zero at infinity. They
test the leading hard coefficient; its vanishing alone gives no
integrable rate for the remainder.

This test includes the bounded two-band symbol ($\lambda=1$ in
reference units) and the massive canonical Dirac symbol
($\lambda=|k|$). In the second case the limiting eigenvalues are
$+1,-1$ in units of carrier speed one. A scalar charge
$q\widehat w(p)I$ has a nonzero diagonal block at each energy
whenever $q\widehat w(p)\ne0$. This lower bound comes from a
finite range of frequency-to-energy ratios. Canceling only the
first coefficient of a large-frequency expansion does not suffice.
For the bounded bands and a fixed nonzero external frequency
$\nu$, the second resolvent in
\eqref{eq:ct-uv-complete-density} is instead $R_{k+p}(z+i\nu)$.
In the limit with $\lambda=1$ it has poles at $z=-E-i\nu$.
At such a pole the bubble residue, up to its overall sign, is
\[
  \sum_{E'}\frac{\|P_EQP_{E'}\|_{\mathrm{HS}}^2}{E'-E-i\nu}.
\]
All nonzero summands have imaginary parts of the same sign,
and the contact has only real poles. At least one such residue
survives for $Q\ne0$. The spatial divergence therefore persists
also at nonzero external frequencies for this bounded-band source.

Schur reduction can produce a frequency-dependent contact in a
retained block. Its effect must be added to the eliminated
determinant. On finite marked blocks the Schur identity gives the
same second variation before the common trace is taken. Ordinary
positive-norm fermion slots with canonical time derivative obey
the preceding residue test in the full block. A source-dependent
linear change of Grassmann variables likewise has its cancelling
Jacobian, by the determinant product identity. If an actual
boundary normalization instead contributes a reciprocal determinant,
the necessary cancellation condition is more specific than a count
of slots. For each limiting energy $E$, the signed sum of the
squared equal-energy charge blocks must vanish, or its double pole
must be supplied by another term of the common source. This
condition concerns a complete coefficient with the indicated
scaling; a genuinely nonlocal temporal symbol requires its own test.

The role of the equal-energy hypothesis can be verified without
taking a limit. For a Hermitian generator $S$ and finite Hermitian
$H$ put $h_a=e^{iaS}He^{-iaS}$. Differentiating gives
$Q=i[S,H]$ and $C=-[S,[S,H]]$. In this case $P_EQP_E=0$ for
every $E$, and the two terms in
\eqref{eq:ct-uv-complete-density} at equal spatial slots cancel
identically for every nonreal $z$: $\det(z+h_a)$ is constant in
$a$. This is an exact source rotation. In a physical density
calculation the equal-energy insertion and any compensating boundary
determinant have to be derived from the gauged action itself.

\paragraph{Carrier-frequency integration after eliminating the Gauss constraint.}
For the two-band charge on a spatial torus of volume $V$, the Coulomb
form in \eqref{eq:cl-gauss-charge-uv} can instead be used directly
as an instantaneous Hamiltonian vertex. At fixed width $a>0$ set
$w(p)=e^{-a|p|}$. With the onsite particle--hole vacuum as reference,
$\rho_p$ is the sum of the two number densities with signs $+1,-1$.
On a configuration of $N$ occupied particle or hole positions, the
diagonal $j=\ell$ in $\rho_p^*\rho_p$ is exactly $N$. The other
ordered pairs are the normally ordered quartic operator. Equivalently
this splitting is the single contraction in the canonical
anticommutation relation. On finite particle vectors it yields
\[
 V_{\rm C}={:V_{\rm C}:}+
      \frac{q^2}{2V}\sum_{p\ne0}\frac{|w(p)|^2}{|p|^2}\,N .
\]
The coefficient is finite even after replacing the normalized sum
by the continuum integral. This one-body energy is fixed by the
charge square, including its sign and factor of one half. Both
terms must be retained. For a different vacuum polarization or
boundary determinant its contractions must be derived anew.

\begin{proposition}[Absolute carrier-frequency bounds for time-local graphs]
\label{prop:final-uv-frequency-trees}
Let a connected finite graph have time-local vertices with
frequency-independent bounds. Remove the overall conservation
delta and fix external frequencies satisfying conservation. Suppose
each internal edge has an envelope
$D_e\langle u_e\rangle^{-a_e}$ at its routed carrier frequency,
with $D_e$ independent of frequencies, $a_e\ge1$, and
$\langle u\rangle=(1+u^2)^{1/2}$ in fixed units. If every
self-loop (an edge with both ends at one vertex) has $a_e>1$, the
absolute integral over all independent internal frequencies is at
most $C_G\prod_e D_e$ times the vertex bounds, uniformly in the
external frequencies.
\end{proposition}
\begin{proof}
Use a uniform spanning tree of the undirected graph, counting parallel
edges separately. Let $b_e$ be the probability that edge $e$ is a
chord, meaning outside this tree. Bridges have $b_e=0$ and may be
bounded by their supremum. A self-loop has $b_e=1$. Every remaining
edge belongs to some spanning tree and can be omitted from another,
so $0<b_e<1$. For each tree $T$ define
\[
 H_T=\prod_{e\notin T}\langle u_e\rangle^{-a_e/b_e}.
\]
The tree probabilities $\pi_T$ give
$\prod_T H_T^{\pi_T}=\prod_{e:b_e>0}\langle u_e\rangle^{-a_e}$.
H\"older's inequality bounds its integral by
$\prod_T(\int H_T)^{\pi_T}$. After imposing the vertex deltas the
chord frequencies of any tree are independent coordinates:
elimination from the leaves fixes each tree frequency with
coefficient $+1$ or $-1$. This gives the same integration measure
in every tree basis, including in the presence of the affine
external shifts. Thus $\int H_T$ is a product of one-dimensional
integrals of $\langle u\rangle^{-a_e/b_e}$ over the real line.
Every exponent is strictly greater than one. The empty product
treats a tree graph with no independent frequency integrals.
\end{proof}

For a gapped self-adjoint free charged symbol,
$\|(iu+h(k))^{-1}\|\le (u^2+m^2)^{-1/2}$ is uniform in $k$.
A propagating transverse Maxwell line at $p\ne0$ has the envelope
$(1+|p|^{-2})\langle u\rangle^{-2}$. In the onsite Hamiltonian
expansion just described, Wick ordering excludes a charged line
with both ends at the same normally ordered vertex, including the
one-body vertex. A transverse photon returning to the same vertex
has exponent two. For bounded remaining time-local vertices the
proposition consequently controls all carrier-frequency regions
of the resulting graphs, at every finite order. A proof of full
convergence still needs the independent spatial integrals and
short visits, and any additional charge, boundary or measure
vertices of the chosen action. The comparison with the oriented
exchange in Proposition~\ref{prop:final-uv-gauss-frequency}
depends on how its coincident charges are defined in the source.

\paragraph{How large must an endpoint map be to restore a physical pole?}
The rate proved below in \eqref{eq:geom-uv-open-vector-rate} tests physical source maps as well
as the unmodified charged field. Suppose the two linear, translation-invariant
endpoint maps preserve chirality and, at fixed depth, the product of their
operator norms is bounded by $C(1+|k|)^M r^{-N}$ for finite $M,N$.
The norm of the opposite-chirality block after applying the maps and
integrating $r$ is at most
\[
 C(1+|k|)^M |k|\int r^{-N}
    \frac{e^{-r^\alpha z}}{m^2+z e^{-2r^\alpha z}}\,\nu_0(\mathrm dr),
 \qquad z=|k|^2 .
\]
The extra powers leave the upper logarithmic rate in \eqref{eq:geom-uv-open-vector-rate}
unchanged: $N|\log r|=o(r^{-p})$ at the tip. This block cannot supply
a nonzero unsubtracted positive vector spectral measure. In particular a
pole of weight $Z>0$ at mass $M_0$ alone requires norm at least
$Z|k|/(z+M_0^2)$ on these chiral components under that representation.
This conclusion uses preservation of the chiral blocks, since a
chirality-changing map can mix the scalar coefficient into the vector slot.

For a depth-independent scalar endpoint map with its adjoint at the other
end, the required multiplier $|U(z)|^2$ must grow at least as
$\exp[(c_*-o(1))z^{p/(p+\alpha)}]$ if it restores such a pole;
$c_*$ is the positive constant on the right of \eqref{eq:geom-uv-open-vector-rate} with its
minus sign removed. But the scalar coefficient before this operation is
$\int m/(m^2+ze^{-2r^\alpha z})\,\nu_0(\mathrm dr)\to1/m$ by
dominated convergence. The same multiplier makes that scalar coefficient
grow faster than any polynomial along the full momentum lattice.
It fails to define a torus distribution, even after addition of any finite
polynomial contact term. A prescription to remove the resulting nonlocal
scalar term or change the measure must therefore be part of a new source.


\paragraph{Spectrum of the same open line after the radial average}
Fix $j\ne0$, $m>0$, and $p=d-1>2\alpha$. Write $A=2|j|/p$ and let
$\nu_0$ be the probability obtained by weighting the radial ground
state with the full zero-photon determinant. At zero electric charge,
the normalized free open line averages $D_0^{-1}$ with precisely
this measure. Its vector coefficient multiplying $-\mathrm i\gamma\cdot k$ is
\begin{equation}
 \overline F_V(z)=\int_0^1
 \frac{e^{-r^\alpha z}}{m^2+z e^{-2 r^\alpha z}}\,\nu_0(\mathrm dr),
 \qquad z=|k|^2 .
 \label{eq:geom-uv-open-vector}
\end{equation}
In particular the soft open line and the ultraviolet graphs use the same
boundary weight. The opposite-chirality probe block $P_+D_0^{-1}P_-$
retains this coefficient, since the scalar term vanishes between the two
projectors. For any $N>0$ one has
$\overline F_V(z)\le C_{N,m} z^{-N}\int_0^1 r^{-\alpha N}\nu_0(\mathrm dr)$
for $z>0$. Every moment on the right is finite.

The rate can be read off from the geometry itself. For the positive
interpolation in \eqref{eq:geom-uv-open-vector}, and hence along the torus momenta,
\begin{equation}
 \lim_{z\to\infty}\frac{\log \overline F_V(z)}{z^{p/(p+\alpha)}}
   =-\left(1+\frac{\alpha}{p}\right)
      \left(\frac{p A}{\alpha}\right)^{\alpha/(p+\alpha)} .
 \label{eq:geom-uv-open-vector-rate}
\end{equation}
To prove it, write $n_0=\mathrm d\nu_0/\mathrm dr$. The ground-state
bounds and the free determinant asymptotics give
$\log n_0(r)=-A r^{-p}+o(r^{-p})$ at the tip: the vacuum costs
$O(r^{-2\alpha}(1+|\log r|)^3)$ there. The denominator in
\eqref{eq:geom-uv-open-vector} lies between $m^2$ and $m^2+z$, so its logarithm
changes the scaled limit by zero. On $0<r<r_0$ bound $n_0$ above and
below by $\exp[-(A\mp\epsilon)r^{-p}]$, respectively, for fixed
$0<\epsilon<A$ and sufficiently small $r_0$. The complementary
interval contributes at most $e^{-r_0^\alpha z}$ to the Laplace
average. Set $r=z^{-1/(p+\alpha)}y$. For $C=A\mp\epsilon$ the exponent
is $-z^{p/(p+\alpha)}(Cy^{-p}+y^\alpha)$. Its global minimum cost
is $(1+\alpha/p)(pC/\alpha)^{\alpha/(p+\alpha)}$. The interval length
in $y$ grows only by a power of $z$, giving the upper estimate.
A fixed neighborhood of the positive minimizing $y$ gives the lower
estimate. Shrink that neighborhood and send $\epsilon$ to zero.

This also tests a proposed physical source identification directly on
the momentum sequence. An unsubtracted positive vector spectral
coefficient $F(z)=\int (z+s)^{-1}\rho_1(\mathrm ds)$ with
$\rho_1\ne0$ has $zF(z)\to\rho_1([0,\infty))>0$, allowing an
infinite total mass. A positive Dirac pole alone would give a strictly
positive lower limit under this hypothesis. Equation~\eqref{eq:geom-uv-open-vector}
instead gives $z\overline F_V(z)\to0$ on every unbounded momentum
sequence. Thus this averaged open coefficient cannot be the restriction
of such a positive spectral coefficient. The result concerns this
specific retained source and this spectral representation. A finite-torus
reflection-positivity test and a reconstruction using a different
physical source map are separate questions.

\paragraph{A spectral test needing only a finite inverse half-moment}
The stretched-exponential estimate above uses the detailed radial ground state. The spectral obstruction has a weaker sufficient hypothesis, and it can be detected on the torus itself~\cite{ChiapponePartXXIX2026}. Keep the free vector symbol in \eqref{eq:geom-uv-open-vector}, but let its heat lapse $\tau$ have any positive, momentum-independent probability law on $(0,\infty)$. Write $F_V(z;\tau)$ for the integrand with $r^\alpha=\tau$, and $\overline F_V=\mathbb E[F_V]$. It suffices that $\mathbb E[(m^2\tau)^{-1/2}]<\infty$.

To see why, use dimensionless variables $x=z/m^2$, $a=m^2\tau$. At fixed lapse,
\[
 g_a(x):=zF_V(z;\tau)
 =\frac{x e^{-ax}}{1+x e^{-2ax}}
 =\frac{\sqrt{x}}{2\cosh(ax-\tfrac12\log x)} .
\]
For $ax\le1$, $g_a(x)\le\sqrt{x}/2$, so its integral against $dx/x$ on that interval is at most $a^{-1/2}$. For $y=ax\ge1$ put $v=y-\tfrac12\log(y/a)$. Since $dv/dy\ge1/2$ and $y^{-1/2}\le1$, the remainder is bounded by
\[
 \frac{1}{2\sqrt a}\int_1^\infty
 \frac{dy}{\sqrt y\cosh v}
 \le\frac{1}{\sqrt a}\int_{\mathbb R}\frac{dv}{\cosh v}
 =\frac{\pi}{\sqrt a}.
\]
Tonelli's theorem now yields $\int_0^\infty\overline F_V(z)\,dz\le(1+\pi)\mathbb E[a^{-1/2}]$. Every nonzero positive Stieltjes spectral measure makes this integral infinite, since $\int_1^\infty dz/(z+s)=\infty$ at each finite spectral mass squared $s\ge0$.

For a periodic torus ray let $z_n=m^2c n^2$ with fixed $c>0$. Its samples satisfy
\[
 \sum_{n=1}^\infty\frac{z_n\overline F_V(z_n)}{n}
 \le 2\sqrt c+\tfrac12(1+\pi)\mathbb E[a^{-1/2}].
\]
Indeed, at fixed $a$, $g_a(ct^2)/t$ equals $\sqrt c/2$ times the hyperbolic secant of $act^2-\log(\sqrt c\,t)$. This argument decreases to a unique minimum and then increases. On each branch the secant rises and falls at most once, so its total variation on both branches together is at most four. Comparing the right endpoint of each interval $[n-1,n]$ with the integral there costs at most the variation of $g_a(ct^2)/t$ over that interval. Summing costs at most $2\sqrt c$, and the integral is $\tfrac12\int_0^\infty g_a(x)\,dx/x$. Averaging proves the bound. Any nonzero positive spectral measure would contribute a positive lower bound to $z_n\overline F_V(z_n)$ from some finite spectral band at all large $n$, forcing harmonic divergence. This tests a common unsubtracted vector spectral measure using the actual sampled values. It requires neither interpolation between torus momenta nor pointwise decay of the rescaled coefficient.

\paragraph{Does a restored free pole carry the electric charge?}
A physical charged readout must pass a test at its mass shell as well
as the estimates at large Euclidean momentum. Consider the analytic
coefficient on a single locked plane, at zero spatial momentum. Write
$\omega$ for the Euclidean momentum component along its chosen carrier,
$\gamma_\parallel^2=1$, and put
\[
 t=m+\mathrm i\omega,qquad
 P=(1+\gamma_\parallel)/2,\qquad
 D_{\rm D}(t)=tP+(2m-t)(1-P).
\]
The positive-energy rest pole on the selected sheet is $t=0$. Its
inverse has residue $P$ in this local variable. Analytic coefficients
here are taken before the declared sector-plane vectorial boundary
value. The computation requires a specified analytic continuation;
a discrete torus sequence alone does not determine one.

For example, let $K(B)$ be the fixed-width Wilson heat-link operator
and $K_0=K(0)$. At the rest pole
$K_0(0)=m(1-e^{a^2m^2}\gamma_\parallel)$ is invertible for $a,m>0$;
the argument $0$ in $K_0(0)$ denotes $t=0$. On real Euclidean momenta
one can form the operator $D'(B)=K(B)K_0^{-1}D_{\rm D}$, with
the free Dirac domain. It has exactly the free Dirac inverse at zero
background. On the finite torus its ordinary relative determinant
is the heat-link one because $D'(B)D_{\rm D}^{-1}=K(B)K_0^{-1}$.
To first order in a photon probe its connected open response is
\[
 -D_{\rm D}^{-1}(k')\,\delta K(k',k)\,K_0^{-1}(k).
\]
The right end is regular at the Dirac shell wherever $K_0^{-1}$ and
the mark are analytic. Its double Dirac pole residue is zero. This
example preserves the closed determinant but changes the open source.
It illustrates why the two tests have to be made on the same operator.

\begin{proposition}[Forward charge of an analytic two-sided pole assignment]
\label{prop:phys-uv-pole-pencil}
On this rest chart let $K(t,\lambda)$ be a holomorphic finite spin
matrix invertible near $(0,0)$, where $\lambda$ is a
translation-invariant external probe. Let $L(t),R(t)$ be holomorphic
there, independent of $\lambda$ and invertible for $0<|t|$ sufficiently
small. Suppose
\[
 D(t,\lambda)=L(t)K(t,\lambda)R(t),\qquad
 D(t,0)=D_{\rm D}(t).
\]
For every forward core mark $V=(\partial_\lambda K)(0,0)$ the
amputated on-shell spin matrix $C(V)=P L(0)V R(0)P$ obeys
$C(V)C(V')=0$ for any two such marks. A Hermitian forward
current on the positive pole space therefore vanishes. In particular
it cannot equal $qP$ for $q\ne0$, in the normalization in which a
uniform charge probe shifts $t$ to $t+q\lambda$ on that space.
\end{proposition}
\begin{proof}
For $t\ne0$ the factorization implies
$R(t)D_{\rm D}^{-1}(t)L(t)=K(t,0)^{-1}$. The residue on the right
vanishes, giving $R(0)PL(0)=0$. This proves the product identity.
With the usual positive adjoint on the pole space a Hermitian
square-zero matrix vanishes. There is also an exact finite-probe
check. At every small $\lambda$,
\[
 \det D(t,\lambda)=\det D_{\rm D}(t)
             \frac{\det K(t,\lambda)}{\det K(t,0)} .
\]
The ratio is a nonvanishing holomorphic factor. The location and
multiplicity of the determinant zero at the mass shell stay fixed.
\end{proof}

If the outer maps depend on $\lambda$, their own determinant marks
can carry the missing charge. In the same holomorphic factorization
suppose $P(\partial_\lambda D)(0,0)P=qP$. For $t\ne0$ differentiate
the ordinary matrix determinants at $\lambda=0$. The simple-pole
coefficient is
\[
 \mathop{\rm Res}_{t=0}\partial_\lambda
       \log(\det L\,\det R)
 =\operatorname{tr}\!\left[P(\partial_\lambda D)(0,0)P\right]
 =2q .
\]
Indeed $\partial_\lambda\log\det D
=\operatorname{tr}(D_{\rm D}^{-1}\partial_\lambda D)$ and the
corresponding derivative of $\log\det K$ is regular.
The factor two counts the two rest spin states of this Dirac channel.
The constant probe tests a background or a conserved charge coupled
as specified; it does not add a fluctuating photon zero mode. This
is a fixed-momentum statement. Whether the newly required marks
retain intermediate displacements on every hard cycle, or cancel
against a Jacobian in a specified physical functional, requires their
full source variation. Their shell residue alone gives no bound on
the ultraviolet momentum integral.

\section{Framework on a fixed sector plane}
\label{sec:framework}

The explicit sources above use one fixed Euclidean carrier. We now
state the geometric, source and carrier-algebra conventions needed
to connect these constructions with the general fixed-plane estimates
and their Lorentzian interfaces.

\subsection{Locked geometry and canonical normalization}

Let \(\mathcal G\simeq\RR^K\), \(K\geq1\), be the geometric-time space
with its positive inner product, and let \(\hat G_{\bb}\in\mathcal G\)
be a unit carrier. On the ambient product with three spatial directions,
use the quadratic form
\[
 \dd s^2=\norm{\dd X_{\rm GT}}_{\mathcal G}^2-\norm{\dd\mathbf x}^2.
\]
The raw GT indices are $\mathcal T_1=\{0\}$ and
$\mathcal T_K=\{0,5,\ldots,K+3\}$ for $K\ge2$. The three spatial
coordinates retain the indices $1,2,3$. Before locking, each
$x^\Gamma_{\bin_\Gamma}$ has its own component label
$\bin_\Gamma=(s_\Gamma,b_{s_\Gamma}^{\rm comp})$ with
$s_\Gamma\in\{N,A\}$. In the
label-diagonal restriction, write
$g_\Gamma=\hat G_{\bb}\cdot\hat x_\Gamma$ for the component of the
unit carrier in the raw GT basis. The label and offset assignments are
\begin{equation}
 \begin{gathered}
 s_\Gamma=s,\qquad
 b_{s_\Gamma}^{\rm comp}=g_\Gamma b_s,\qquad
 \bb=(s,b_s).
 \end{gathered}
 \label{eq:uvfinal-offset-labels}
\end{equation}
Thus $b_s$ is a normalized carrier-axis offset and the total ambient
offset is $B_{{\rm GT},\bb}=b_s\hat G_{\bb}$.
The coordinate pullback uses the centred carrier coordinate
$X_{0,\bb}=\hat G_{\bb}\cdot X_{{\rm GT},\bb}-b_s$, denoted
$G_{0\bb}$ in Ref.~\cite{ChiapponeGHOUQFT2026}. The fibrewise embedding is
\begin{equation}
 \begin{gathered}
 \iota_{\bb}(X_0,\mathbf x)
   =\bigl((X_0+b_s)\hat G_{\bb},\mathbf x\bigr),\\
 Q_{\bb}X_{\rm GT}=0,\qquad
 Q_{\bb}=\one-\hat G_{\bb}\hat G_{\bb}^{\,T}.
 \end{gathered}
 \label{eq:locking}
% Monograph source: eq:GTL-master-affine-embedding; eq:pre-post-GTL-affine-offset-pushforward.
\end{equation}
The event on the labelled space is $(\bb,X_0,\mathbf x)$.
The map to ambient coordinates can send different labelled events to
the same point. For example, changing the offset on a complete carrier
line relabels its clock-zero section while leaving the set of points
on that line fixed. A copy identifier distinguishes independently
populated copies at the same $(s,b_s)$ and carrier. Throughout the
fixed-copy calculations, the diamond $\bin$ labels the field and its GT
slot. The carrier fixes the pole unit $e_{\bb}=\hat p_{0,\bb}=\hat G_{\bb}$,
and the sector sign sets the direction of drift $\sigma_s\hat G_{\bb}$.

On a fixed fibre the pullback metric and the positive intrinsic measure are
\begin{equation}
 \iota_{\bb}^{*}(\dd s^2)=\dd X_0^2-\dd\mathbf x^2,
 \qquad
 \dd^4x_{\bb}=|\dd X_{0,\bb}|\,\dd^3\mathbf x.
 \label{eq:pullback-metric}
\end{equation}
The subscript in $X_{0,\bb}$ records the GT-coordinate slot; its value
is $X_0=G_{0\bb}$ on this fixed copy and has dimensions of length.
The spatial triple $\mathbf x=(x^1,x^2,x^3)$ carries no sector-plane
label. The subscript on $\dd^4x_{\bb}$ labels the complete measure,
not its spatial components; likewise $\dd\Sigma_{\bb}=\dd^3\mathbf x$
on a flat fixed-GT section. Fields and their GT arguments retain $\bb$.
Independent spatial slots use $\mathbf x,\mathbf y$ without identifying
their field copies. Let $t$ be the calibrated clock readout in the
chosen carrier-adapted inertial chart. A comparison between copies uses
a common frame, matched clock origins and calibration, and matching
local geometry. With \(\sigma_N=+1\) and \(\sigma_A=-1\),
\[
 X_0=\sigma_s\cGT t,
 \qquad t=\frac{\sigma_s X_0}{\cGT},
 \qquad
 \vct c_{\bb}=\sigma_s\cGT\hat G_{\bb}.
\]
Thus \(\dd X_0=\sigma_s\cGT\dd t\), whereas
\(\dd^4x_{\bb}=\cGT|\dd t|\,\dd^3\mathbf x\) is positive
and contains no orientation sign. The common clock readout is distinct
from the directed carrier coordinate; the antiverse orientation does not
mean a backward-running physical clock.

\paragraph{Locked and Minkowski representations}
Write $\RGT$ for this assigned locked representation, retaining its
carrier, orientation, affine placement and copy data; write $\RMink$
for the ordinary intrinsic Minkowski reference representation. Neither
symbol denotes the full $(K+3)$-dimensional ambient geometry. On one
fixed copy the comparison map is
\begin{equation}
 \begin{aligned}
 \mathcal F_{\bb}:\RGTL&\longrightarrow\RMink,\\
 \bigl((X_0+b_s)\hat G_{\bb},\mathbf x\bigr)
    &\longmapsto(\gamma_{\bb},\mathbf x),\qquad
       \gamma_{\bb}=\sigma_sX_0,\\
 \mathcal F_{\bb}^*(\dd\gamma_{\bb}^2-\dd\mathbf x^2)
    &=\dd X_0^2-\dd\mathbf x^2.
 \end{aligned}
 \label{eq:polish-representation-map}
\end{equation}
Its inverse requires those fixed geometric data. Forgetting them across
an entire family is not an identification of independent fields or states.
The constant $\cGT>0$ denotes the norm of the calibrated vector
$\vct c_{\bb}$, not a replacement for that vector. The action-derived,
branch-normalized geo-momentum vector $\overline{\vct Y}_{\bb}$ supplies
$Y_{\parallel,\bb}=\ee\cdot\overline{\vct Y}_{\bb}$ and
$Y_{\gamma,\bb}=\sigma_sY_{\parallel,\bb}$. On an admitted self-adjoint
domain its carrier translation is
\begin{equation}
 U_{\bb}(\Delta X_0)
  =e^{-\ii\Delta X_0Y_{\parallel,\bb}/\hbar}
  =e^{-\ii\Delta\gamma_{\bb}Y_{\gamma,\bb}/\hbar}.
 \label{eq:polish-geomomentum-generator}
\end{equation}
This paragraph temporarily displays $\hbar$ explicitly. Only after clock
calibration is $\vct c_{\bb}\cdot\overline{\vct Y}_{\bb}
=\cGT Y_{\gamma,\bb}$ a scalar energy readout. GT order is order along
the declared oriented carrier, not an additional primitive clock-time flow.
All subsequent Euclidean operators and Fourier integrals are coefficient
constructions on the selected continuation of $\RGTL$. Agreement with
$\RMink$ requires the corresponding action, source and state comparison.

In the equal-component chart $g_\Gamma=K^{-1/2}$.
On the coordinate-locked support, after the independent label-diagonal
assignment, the raw GT components share the value $x^0:=x^0_{\bb}$.
Each component offset equals $b_s/\sqrt K$, and the coordinate
dictionary is
\[
\begin{aligned}
 X_0&=\sqrt K\,x^0-b_s,\\
 \dd^4x_{\bb}&=\sqrt K\,|\dd x^0_{\bb}|\,\dd^3\mathbf x,\\
 p_\parallel&=K^{-1/2}\sum_\Gamma p_\Gamma.
\end{aligned}
\]
The measure, density, Fourier phase, and external normalization must be
transformed together. In particular, one cannot insert \(Kp_0\)
into a Dirac slash while keeping the unrescaled denominator and measure.
Throughout the loop examples,
\begin{equation}
 k=(k_\parallel,\mathbf k),\qquad
 k^2=k_\parallel^2-\mathbf k^2,\qquad
 \slashed k=\gamma^\mu k_\mu
 \label{eq:intrinsic-momentum}
\end{equation}
are canonically normalized intrinsic quantities on the fixed copy. A
change to the oriented coordinate $\gamma_{\bb}=\sigma_sX_0$ transports
its conjugate coefficient as $k_\gamma=\sigma_sk_\parallel$; the shell
square is unchanged. Coordinates, Fourier covectors and the spin frame
must be transformed together when changing charts. The Clifford algebra
is \(\mathrm{Cl}(1,3)\); no additional gamma matrices or multiplicity
factors arise from this coordinate choice.

We use \(\hbar=1\) and momentum-dimension mass parameter \(m\). With
explicit units the shell parameter is \(m_{\rm phys}\cGT\), and the
length corresponding to the inverse-momentum kernel parameter \(a\)
is \(\ell_{\Ph}=\hbar a\). Proper-time variables
\(\alpha\) and \(a^2\tau\) have inverse-momentum-squared dimensions.
The variable \(z=a^2q_E^2\) in \(\Ph\) is dimensionless.

\subsection{GT order, physical endpoints, and positive state measures}
\label{sec:fw-endpoints}

\paragraph{Clock order and carrier order}
Physical preparation and detection use increasing calibrated clock readings.
On a fixed copy, however, $\dd X_{0,\bb}/\dd t=\sigma_s\cGT$.
Consequently, after matching the clock origins, a large carrier interval
has physical endpoints
\begin{equation}
 X_{0,\bb}^{\rm in}=-\sigma_s R,\qquad
 X_{0,\bb}^{\rm out}=+\sigma_s R,\qquad R\longrightarrow\infty.
 \label{eq:fw-physical-endpoints}
\end{equation}
The fixed affine offset cancels from differences on one copy. Comparing
different offsets requires their absolute carrier coordinates or a
specified synchronization transport; equal numerical $X_0$ is not by
itself an identification of events on different copies.

For bosons on this copy define the signed-carrier orderers by
\begin{equation}
\begin{aligned}
 X_0^+\{A_{\bb}(x)B_{\bb}(y)\}
 &=\theta(X_{0,x}-X_{0,y})A_{\bb}(x)B_{\bb}(y)\\
 &\quad+\theta(X_{0,y}-X_{0,x})B_{\bb}(y)A_{\bb}(x),\\
 X_0^-\{A_{\bb}(x)B_{\bb}(y)\}
 &=\theta(X_{0,y}-X_{0,x})A_{\bb}(x)B_{\bb}(y)\\
 &\quad+\theta(X_{0,x}-X_{0,y})B_{\bb}(y)A_{\bb}(x).
\end{aligned}
\label{eq:fw-signed-orderers}
\end{equation}
For homogeneous odd fields the exchanged term has the appropriate Koszul
sign. At coincident GT coordinate, regular kernels use $\theta(0)=1/2$;
singular products require a common distributional extension with its
contact terms, not arbitrary multiplication of step functions and deltas.

The orderer with sign $\sigma_\epsilon$ places its late carrier boundary
in the bra and its early boundary in the ket. Those boundaries are
$X_0=\sigma_\epsilon R$ and $X_0=-\sigma_\epsilon R$, respectively.
Comparison with \eqref{eq:fw-physical-endpoints} gives Table~\ref{tab:fw-endpoints}.
In particular, the product $\sigma_s\sigma_\epsilon$, rather than either
sign alone, determines whether a physical final state occupies the bra.

\begin{table}[t]
\centering
\begin{tabular}{@{}llll@{}}
\toprule
GT order & Sector & Bra slot & Ket slot\\
\midrule
$X_0^+$ & N & physical out & physical in\\
$X_0^+$ & A & physical in & physical out\\
$X_0^-$ & N & physical in & physical out\\
$X_0^-$ & A & physical out & physical in\\
\bottomrule
\end{tabular}
\caption{External slots on a common oriented carrier. ``In'' and ``out''
refer to physical preparation and detection; neither is redefined by its
position in the bracket.}
\label{tab:fw-endpoints}
\end{table}

If the common branch-normalized interaction geo-momentum is
self-adjoint on an admitted domain, its two ordered coefficient
evolutions on the interval with endpoints $-R,R$ are given by
\begin{equation}
 U^\epsilon_{\bb}(R,-R)
 =X_0^\epsilon\exp\!\left[-\ii\sigma_\epsilon
        \int_{-R}^{R}\!\dd X_{0,\bb}\,
                Y_{{\rm I},\parallel,\bb}(X_{0,\bb})\right].
 \label{eq:fw-dyson-direction}
\end{equation}
Here $\hbar=1$ as elsewhere in the loop calculations; otherwise divide
the exponent by $\hbar$. The integral in this display uses the positive
coordinate interval, with its reversal sign outside once. It does not
multiply a generator already carrying the raw action's branch sign by
that sign a second time. Under these assumptions $U^-=(U^+)^\dagger$.
For an isolated sector, clock-forward transport uses
$\sigma_\epsilon=\sigma_s$. In particular,
\begin{equation}
 \langle f_A|U^-_A|i_A\rangle
   =\overline{\langle i_A|U^+_A|f_A\rangle}.
 \label{eq:fw-antiverse-adjoint}
\end{equation}
This identity is conditional on the admitted physical adjoint and the
same generator; the vectorial pole prescription alone proves neither.

\paragraph{Mixed boundary kernels}
Let $\bb_N=(N,b_N)$ and $\bb_A=(A,b_A)$ have a declared common carrier
comparison. Their carrier-ordered boundary expressions are
\begin{equation}
\begin{aligned}
 \mathcal A_+[f;i]
 &=\langle f_{N,\bb_N},i_{A,\bb_A}|\mathscr S_+
                     |i_{N,\bb_N},f_{A,\bb_A}\rangle_{\rm bdry},\\
 \mathcal A_-[f;i]
 &=\langle i_{N,\bb_N},f_{A,\bb_A}|\mathscr S_-
                     |f_{N,\bb_N},i_{A,\bb_A}\rangle_{\rm bdry}.
\end{aligned}
\label{eq:fw-mixed-boundary-amplitudes}
\end{equation}
Reversing an A boundary uses its conjugate Hilbert space: the Riesz map
$\mathcal J_A:v\mapsto\langle v,\cdot\rangle$ is antilinear. Thus its
wave packet is dualized and conjugated, not merely moved typographically.
A fixed ordering of fermionic external slots supplies the corresponding
permutation signs. ``bdry'' records these endpoint maps. It is not an
assertion that an arbitrary rearrangement of a unitary matrix is unitary.
For example, the swap on $\mathbb C^d\otimes\mathbb C^d$ is unitary,
but its partial transpose is the rank-one operator
$|\Omega\rangle\langle\Omega|$, $|\Omega\rangle=\sum_j|j,j\rangle$,
for $d>1$. A mixed boundary kernel must therefore establish its physical
composition and adjoint law rather than importing them from relabeling.

For physical creation with no incoming particles, the same assignments give
\begin{equation}
\begin{aligned}
 \mathcal C_+(f_N,f_A)
 &=\langle f_{N,\bb_N},0_{A,\bb_A}|\mathscr S_+
                      |0_{N,\bb_N},f_{A,\bb_A}\rangle_{\rm bdry},\\
 \mathcal C_-(f_N,f_A)
 &=\langle0_{N,\bb_N},f_{A,\bb_A}|\mathscr S_-
                      |f_{N,\bb_N},0_{A,\bb_A}\rangle_{\rm bdry}.
\end{aligned}
\label{eq:fw-creation-slots}
\end{equation}
The final A particles in the ascending-GT ket are not physically incoming.
An excited Fock ket or a vacuum-to-one-particle field matrix element is
not such a creation amplitude and keeps its ordinary meaning.

The necessity of mixed orientations in a closed creation process has a
separate conservation premise. Suppose the initial total directed
geo-momentum is zero, the complete action conserves it, no source or
background recoil has been omitted, and every final mode has positive
nonzero oriented magnitude $\kappa_j$. On a common carrier,
\begin{equation}
 \Delta\overline{\vct Y}
  =\hat p_0\left(\sum_{f_N}\kappa_j-\sum_{f_A}\kappa_j\right)=0.
 \label{eq:fw-creation-balance}
\end{equation}
A nonempty all-N or all-A collection fails this condition. Mixed sectors
are necessary, not sufficient: spatial momentum, charge, species rules,
and an actual coupling still matter. With different carriers the
transported vector sum replaces the scalar difference. Driven processes
and zero modes are not covered by this argument. In particular, the
support-diagonal O--QED reference action below contains no N--A creation
vertex. Neither the slot convention nor this charge balance supplies one.

\paragraph{Positive states and branch-routed cut momenta}
For a positive one-particle channel on $\bb$, let
$p^{\rm ph}_\parallel=\sigma_s\kappa_{\mathbf p}$ with
$\kappa_{\mathbf p}=\sqrt{\mathbf p^2+m^2}>0$. Covariant normalization is
\begin{equation}
\begin{aligned}
 \langle p,h;\bb|p',h';\bb'\rangle
  &=(2\pi)^3 2\kappa_{\mathbf p}\,
    \mathcal P_{\bb\bb'}\delta_{hh'}
           \delta^{(3)}_{\bb}(\mathbf p-\mathbf p'),\\
 \dd\Pi^+_{\bb}(p^{\rm ph})
  &=\frac{\dd^4p^{\rm ph}_{\bb}}{(2\pi)^4}
       2\pi\delta^{\rm sh}_{\bb}((p^{\rm ph})^2-m^2)
                    \theta(\sigma_s p^{\rm ph}_\parallel)\\
  &=\frac{\dd^3\mathbf p}{(2\pi)^3 2\kappa_{\mathbf p}}.
\end{aligned}
\label{eq:fw-positive-state-measure}
\end{equation}
The support label belongs to $\dd\Pi^+_{\bb}$, not to the spatial
integration variable $\mathbf p$. In multiparticle measures, spatial
momenta such as $\mathbf q_i$ retain the particle index $i$ only.
The angular/spin labels $h$ and the support identity are independent.
This coefficient normalization is standard~\cite{TongQFTFramework}; the
assignment of its signed carrier sheet is the framework convention above.
No negative state measure follows from an A final state occupying a ket.

The branch-selected cut of \eqref{eq:physical-cut-map} instead uses
$\theta(\eta k_\parallel)$, with $\eta=\sigma_s\sigma_\epsilon$.
For a cut arrow represented by $k=\sigma_\epsilon p^{\rm ph}$,
\begin{equation}
 \eta k_\parallel=\sigma_s p^{\rm ph}_\parallel
                   =\kappa_{\mathbf p}>0,
 \qquad |\det(\partial k/\partial p^{\rm ph})|=1.
 \label{eq:fw-cut-routing}
\end{equation}
Thus that branch-selected expression is a positive physical completeness
measure only with the matching endpoint/arrow convention. Reversing the
arrow reverses the complete covector and all its attachments, not just
the step function. At fixed physical $p^{\rm ph}$, changing
$\sigma_\epsilon$ does not change which states have positive norm. This
routing is a real change of Fourier convention, not a shift by a
vectorial infinitesimal or a proof that a virtual discontinuity equals a cut.

\subsection{Coordinate restriction and momentum alignment}
\label{sec:fourier-locking}

Coordinate locking constrains the coordinates of one event. To see
its effect by itself, fix the raw component labels $\bin_\Gamma$, put
$u_\Gamma=x^\Gamma_{\bin_\Gamma}$, and let $\kappa_\Gamma$ be arbitrary
ambient Fourier variables. In the equal-component chart,
\begin{equation}
 \begin{aligned}
 &\int_{\RR^K}\dd^K u\,
  e^{-\ii\sum_{\Gamma\in\mathcal T_K}\kappa_\Gamma u_\Gamma}
  \prod_{\Gamma\in\mathcal T_K\setminus\{0\}}\delta(u_\Gamma-u_0)\\
 &\qquad=2\pi\delta\!\left(\sum_{\Gamma\in\mathcal T_K}\kappa_\Gamma\right).
 \end{aligned}
 \label{eq:fourier-locking}
\end{equation}
Evaluation of the transverse deltas leaves just the Fourier integral in
$u_0$. A function $f(u_0)$ on the diagonal gives
$\widetilde f(\sum_\Gamma\kappa_\Gamma)$ in place of the delta. These
expressions see the sum of the Fourier components and leave their
differences arbitrary. The component labels are fixed throughout this
calculation. Equation~\eqref{eq:uvfinal-offset-labels} is the independent
label-diagonal assignment.

In the selected ambient chart, the invariant normal delta has the
coarea normalization given in \eqref{eq:uvfinal-normal-delta}.
It converts the common-component measure to the normalized carrier
measure, including the empty product when $K=1$.

The physical post-GTL theory also imposes the independent A2/GTP
momentum binder
\begin{equation}
 Q_{\bb}p_{\rm GT}=0,\qquad
 p_{\rm GT}=p_\parallel\hat G_{\bb},\qquad
 p_\parallel=\hat G_{\bb}\!\cdot p_{\rm GT}.
 \label{eq:independent-momentum-binder}
% Monograph source: eq:A2-carrier-covector-pullback; eq:GTL-covector-pullback.
\end{equation}
The pullback of any ambient covector in the Fourier phase gives
$p_{\rm GT}\cdot[(X_0+b_s)\hat G_{\bb}]=p_\parallel(X_0+b_s)$.
Equation~\eqref{eq:independent-momentum-binder} selects the physical
ambient covectors that lie entirely along that carrier. Even after
alignment, $p_\parallel$ may have either sign in an unrestricted
Fourier integral. The positive physical ray used for external states
is selected in \eqref{eq:fw-positive-state-measure}.

\subsection{Plane binding and functional derivatives}
\label{sec:fw-plane-binding}

For a discrete support ledger equipped with counting measure, let
$\bb=(s,b_s)$ and $\bb'=(s',b'_{s'})$ denote independent labels. Define
\begin{equation}
 \mathcal P_{\bb\bb'}=\delta_{ss'}\delta_{b_s b'_{s'}},
 \qquad
 C_{\bb\bb'}(x,y)=\mathcal P_{\bb\bb'}C_{\bb}(x,y).
 \label{eq:plane-projector}
\end{equation}
Here the Kronecker symbols are the identity kernel for counting
measure on the declared labels. If instead the label integral uses an
atomic measure $\mathfrak m(\dd\bb')$ with nonunit atom weights, its
identity kernel is $\mathcal P_{\bb\bb'}/\mathfrak m(\{\bb'\})$ on those
atoms. The action and interactions on this chart are support diagonal by
assumption. For a continuous family, fix a carrier and choose a
post-GTL measure $\nu_s(\dd b_s)$. Fibrewise identities then hold for
$\nu_s$-almost every $b_s$. The corresponding diagonal component
measure is the pushforward of $\nu_s$ by
$b_s\mapsto(g_\Gamma b_s)_{\Gamma\in\mathcal T_K}$ as in
\eqref{eq:uvfinal-offset-labels}. For example, with
$\nu_s(\dd b_s)=w_s(b_s)\dd b_s$ in the equal-component chart, the
marginal for any one raw offset $u=b_s/\sqrt K$ is
$\sqrt K\,w_s(\sqrt K\,u)\dd u$. This choice supplies the measure on
the label diagonal even when that diagonal has zero mass in an
independent product of component measures. The diagonal covariance on
a continuous family acts as a multiplication operator on the direct
integral over these fibres.

Fix also the ordering branch \(\epsilon\). With
\(\ee:=\hat p_{0,\bb}\), the post--GTL action unit is
\begin{equation}
 u_{\epsilon,\bb}:=\sigma_\epsilon\ee,\qquad
 u_{\epsilon,\bb}^{\,2}=1,\qquad
 \BR u_{\epsilon,\bb}=0.
 \label{eq:action-branch-unit}
\end{equation}
We use \(\hat u_{\bb}\equiv u_{\epsilon,\bb}\) only after this
branch has been fixed. The regulator-free physical action and density
factor into a carrier direction and a scalar coefficient:
\begin{equation}
 \begin{aligned}
 \vct S_{\epsilon,\bb}
   &=u_{\epsilon,\bb}S_{{\rm sc},\epsilon,\bb},\\
 S_{{\rm sc},\epsilon,\bb}
   &=u_{\epsilon,\bb}\vct S_{\epsilon,\bb},\\
 \vct{\mathcal L}_{\epsilon,\bb}
   &=u_{\epsilon,\bb}\mathcal L_{{\rm sc},\epsilon,\bb}.
 \end{aligned}
 \label{eq:directed-action}
% Monograph source: eq:A5-postGTL-scalar-directed-action-ledger.
\end{equation}
For the standalone O--QFT representative, the scalar action
contains the sector-fold sign \(\sigma_s\) exactly once. The branch
unit \(u_{\epsilon,\bb}\), the fixed pole carrier \(\ee\), and the
running line unit \(\widehat u_{k,\bb}\) have different roles and are not
interchanged. In particular, \(u_{\epsilon,\bb}\) is not the
numerator convention \(\sigma_s\ee\) of a fixed virtual-carrier
benchmark.

The scalar gauge-fixing and ghost coefficients obey
\[
 S_{{\rm gf},{\rm sc}}+S_{{\rm gh},{\rm sc}}
       =\BR\Psi_{\rm sc},\qquad
 \vct\Psi=u_{\epsilon,\bb}\Psi_{\rm sc},\qquad
 \vct S_{\rm gf}+\vct S_{\rm gh}=\BR\vct\Psi.
\]
The coefficient representation is useful for Euclidean estimates.
The factorization in \eqref{eq:directed-action} applies to the
regulator-free physical action terms; it is not a rule for attaching
one carrier factor to every regulated functional, determinant,
connected graph, or loop-corrected effective action. Those objects
are reconstructed from their carrier characters as described below.

The variational prescription is pullback first. At fixed carrier and
fixed support, the action, connections, inverse metrics, and other
derived geometric objects are formed intrinsically from the pulled-back
fields before Euler--Lagrange variation or integration by parts.
Holding the support fixed means
\(\delta\hat G_{\bb}=0\) and
\(\delta\Delta_{{\rm GTL},\bb}=0\); it does not mean
\(\partial_M\Delta_{{\rm GTL},\bb}=0\).
An ambient integration by parts would generally differentiate the
locking distribution. Ambient source and contraction bookkeeping may
precede restriction, but it does not change this intrinsic
variational prescription.

On a fixed carrier, whenever the corresponding
canonical or variational objects exist on the admitted domain,
they inherit the same carrier factor:
\[
 \vct\pi_{\bb}=\hat u_{\bb}\pi_{{\rm sc},\bb},
 \qquad \vct j_{\bb}^{\,\mu}
          =\hat u_{\bb}j_{{\rm sc},\bb}^{\,\mu}.
\]
The carrier-directed translation charge is geo-momentum. Its scalar
coefficient is useful in the analytic estimates, while the charge
itself retains the direction and sector sign inherited from the action.

On an admitted canonically normalized bosonic free-field domain,
at equal \(X_0\),
\begin{align}
 [\phi_{\bb}(X_0,\mathbf x),
     \pi_{{\rm sc},\bb'}(X_0,\mathbf y)]
   &=\ii\mathcal P_{\bb\bb'}\delta^{(3)}_{\bb}(\mathbf x-\mathbf y),
 \label{eq:ccr}\\
 [\phi_{\bb},\vct\pi_{\bb'}]
   &=\ii\hat u_{\bb}\mathcal P_{\bb\bb'}
       \delta^{(3)}_{\bb}(\mathbf x-\mathbf y).
 \label{eq:vector-ccr}
\end{align}
Fermions have the corresponding canonically normalized CAR with
the same label identity. Positivity and stability of that
quantization remain properties of the admitted physical completion.

Independent intrinsic sources obey the usual identity
\begin{equation}
 \frac{\delta J^A_{\bb}(x)}{\delta J^B_{\bb'}(y)}
   =\delta^A{}_B\,\mathcal P_{\bb\bb'}\delta^{(4)}_{\bb}(x-y).
 \label{eq:source-identity}
\end{equation}
The Dirac distribution in \eqref{eq:source-identity} integrates against
$\dd^4y_{\bb'}$ in the source slot. On a nonconstant Euclidean density it
is replaced by $\delta_{\mu_{E,\bb}}$, as defined below. Ambient transport
phases belong to a transport kernel; they do not change this identity
operator or identify two independent sources.

\paragraph{Measures, copy identities, and geometric variation}
The sector-plane pair is sufficient here because one field copy on the
selected carrier is fixed. If independently populated copies share that
pair, retain a further identifier $\mathfrak p$, with assigned
$\bb_{\mathfrak p}$ and carrier. Their source identity has an additional
$\delta_{\mathfrak p\mathfrak q}$; equal ambient coordinates do not
identify their integration slots. The numerical estimates below hold
on each fixed copy; summing over the additional copy index requires
a separate bound.

To make the integration convention explicit, put
\begin{equation}
\begin{aligned}
 \dd\mu_{\bb}&:=|\dd G_{0\bb}|\,\dd^3\mathbf x,
       &G_{0\bb}&=X_0,\\
 \dd\lambda_{\bb}&:=\dd^KX_{{\rm GT},\bb}\,\dd^3\mathbf x,\\
 \int\dd\lambda_{\bb}\,\Delta_{{\rm GTL},\bb}F_{\bb}
   &=\int\dd\mu_{\bb}\,\iota_{\bb}^{*}F_{\bb}.
\end{aligned}
\label{eq:update-labelled-coarea}
\end{equation}
Here $\Delta_{{\rm GTL},\bb}$ is the invariant normal delta, with the
$\sqrt K$ Jacobian already specified above. The optional
$\ell_{\rm GT}^{K-1}$ making that insertion dimensionless is a different
normalization, absorbed in the density only once. Coordinate deltas,
normal functional constraints, spectral atoms, and mass-shell deltas
are different distributions even when their labels agree.

\paragraph{The ambient action and its exact intrinsic reduction}
Choose an orthonormal normal frame $E_{\perp,\bb}$ and write
$X_{{\rm GT},\bb}=(X_{0,\bb}+b_s)\hat G_{\bb}+E_{\perp,\bb}y_{\bb}$.
Then $\Delta_{{\rm GTL},\bb}=\delta^{(K-1)}(y_{\bb})$ relative to the
normal Lebesgue measure. Let $A$ be the $(K-1)\times K$ matrix
whose $\Gamma$th row, for $\Gamma\ne0$, is
$(\hat x_\Gamma-\hat x_0)^{\mathsf T}$. Then the components of
$AX_{\rm GT}$ are the differences $x^\Gamma-x^0$ in the diagonal
chart. Since $AA^{\mathsf T}=I+\mathbf1\mathbf1^{\mathsf T}$,
$\det(AA^{\mathsf T})=K$. The coarea formula gives the $\sqrt K$
in \eqref{eq:uvfinal-normal-delta}.

Let $Z_{\rm GT}=\ell_{\rm GT}^{K-1}$ and let
$\widetilde{\mathcal L}^{\rm pull}_{\bb}$ be an intrinsic matter
coefficient before this normalization. Extend that coefficient off the
support only to interpret the distributional integral. With the labels already restricted by
\eqref{eq:uvfinal-offset-labels}, display the standalone matter
action first as an ambient integral with one coordinate delta and
then as its intrinsic integral.
\begin{equation}
\begin{aligned}
 \vct S^{\rm amb}_{\epsilon}
  &=\sum_s\sigma_s\sigma_\epsilon\hat p_{0,\bb_s}
       \int\dd\lambda_{\bb_s}\,Z_{\rm GT}
       \Delta_{{\rm GTL},\bb_s}
                   \widetilde{\mathcal L}^{\rm pull}_{\bb_s}\\
  &=\sum_s\int\dd\mu_{\bb_s}\,
                   \vct{\mathcal L}^{(4)}_{\epsilon,\bb_s},\\
 \vct{\mathcal L}^{(4)}_{\epsilon,\bb}
  &=\sigma_s\sigma_\epsilon\hat p_{0,\bb}
                    Z_{\rm GT}\widetilde{\mathcal L}^{\rm pull}_{\bb}.
\end{aligned}
\label{eq:fw-ambient-directed-action}
\end{equation}
Only the density's restriction enters. Its derivatives, connections and
variation retain the intrinsic pullback-first prescription. The common
unit may be factored out only after the declared carrier transport. In
\eqref{eq:directed-action}, $S_{{\rm sc},\epsilon,\bb}$ already includes
$\sigma_s$ for this representative, so that sign is not multiplied twice.
The metric-origin A4 representative has its different lift
$\sigma_\epsilon\hat p_{0,\bb}$ and neutral metric-copy fold; it does not
acquire the standalone sector fold by changing notation.

The right-hand action in \eqref{eq:fw-ambient-directed-action} already
has the normal delta integrated out. Inserting it again would not be the
same action. Likewise $Z_{\rm GT}$ belongs in the density or measure once,
not both. Different interaction vertices have independent event variables
and independent support deltas; coincidence after contraction does not
prescribe a square of a delta at one vertex.

A normal-transfer model has a different extra integration variable. A
marked layer can be represented by
$\int_0^\infty\dd r_{\bb}\,\delta(r_{\bb}-r_\star)V_{\bb}[U_{\bb}(r_{\bb})]$,
where $r_\star>0$. Multiplying its entire normal kinetic action by
$\delta(r_{\bb})$ would collapse that propagation problem. Coordinate
locking, the label projector, the A2/GTP momentum binder, and the
pre-GTL shell-export unit $\hat v_{\mathcal G}^{[\Gamma]}$ therefore
remain distinct operations or data. The shell unit is identified with a
post-GTL carrier only through its declared binder, not before it.

Write $\dd\mu_{E,\bb}(x)=\varrho_{\bb}(x)\dd^4x_{E,\bb}$ for a
continued positive density, and use its actual adjoint in all source
pairings. The invariant identity is
\begin{equation}
\begin{aligned}
 \delta_{\mu_{E,\bb}}(x,y)
   &=\frac{\delta^{(4)}_{\bb}(x-y)}{\varrho_{\bb}(y)},\\
 \int\dd\mu_{E,\bb}(y)\,\delta_{\mu_{E,\bb}}(x,y)f(y)&=f(x).
\end{aligned}
\label{eq:update-measure-delta}
\end{equation}
Thus \eqref{eq:source-identity} is its flat-density specialization.
In the displayed fixed-plane actions $\dd^4x_{\bb}$ denotes the
appropriate positive intrinsic density (continued to
$\dd\mu_{E,\bb}$ in the Euclidean coefficient calculation); it never
contains $\sigma_s$. Similarly, $\delta^{(3)}_{\bb}$ is the equal-section
identity and $\delta^{\rm sh}_{\bb}(D)$ acts in the declared Fourier
momentum measure. For readability, fixed-copy fields and scalar graph functions may still
suppress their $\bb$ subscript, but their integration measures and identity
kernels retain it. In a flat Euclidean chart define
\begin{equation}
\begin{aligned}
 \dd\mu_{E,\bb}(x)&=\dd^4x_{E,\bb},
 &\dd\Sigma_{\bb}&=\dd^3\mathbf x,\\
 \dd^{4L}k_{\bb}&:=\prod_{\ell=1}^L\dd^4k_{\ell,\bb},
 &\langle f,g\rangle_{E,\bb}
   &:=\int\dd\mu_{E,\bb}(x)\,f(x)^*g(x).
\end{aligned}
\label{eq:final-integration-ledger}
\end{equation}
The Fourier factors $(2\pi)^{-4L}$ remain explicit. In scalar graph
estimates $k,q,\ell$ are Euclidean variables unless a Lorentzian
boundary is stated; adding their copy subscript does not change the
routing, Jacobian or Fourier convention. Radial and Feynman-parameter
integrals obtained from them are ordinary scalar calculations.

The same rule applies to functional integration: $\mathcal D_{\bb}\phi$
refers to all declared field species on this copy, not an integration
over the family of planes. It remains a formal notation unless a measure,
finite-mode realization or perturbative construction has been specified.
An actual boundary Gaussian law is written $\dd\mathbb P_{0,\bb}$.
The left source derivative is relative to the displayed density:
\[
 \delta F[J]=\sum_A\int\dd\mu_{E,\bb}(x)\,
       \delta J_A(x)\frac{\delta_lF}{\delta J_A(x)}.
\]
For odd sources this source-first order is part of the convention.
Kernel composition and adjoints therefore integrate over
$\dd\mu_{E,\bb}$; no coordinate delta may silently replace its
measure-dependent identity kernel on a nonconstant density.

Heat and depth variables are not additional event coordinates.
The dimensionless heat measure is $\rho_{\bb}(\dd\tau)$; edge-specific
laws retain the edge index as well. A heat atom
$\delta^{\rm heat}_{\bb,c}$ means evaluation at $\tau=c$ in that law,
not localization of a physical point. Ordinary auxiliary differentials
such as $\dd\tau$, $\dd\alpha$ and a contour $\dd z$ do not carry an
independent field-copy degree of freedom. Their enclosing measure or
fixed-copy calculation specifies their role. Kronecker identities on
field species, characters and copy labels likewise remain distinct from
coordinate, mass-shell, heat and depth Dirac distributions.

Background response must transport the density as well as the kernel.
For $\dd\mu_{E,\bb;y}=r_y\dd\mu_{E,\bb;*}$, with smooth positive
$r_y$, identify the coefficient Hilbert spaces by $W_y=r_y^{1/2}$.
On a common differentiable operator domain, let
$A_y=\tfrac12\partial_y\log r_y$. Then
\begin{equation}
\begin{aligned}
 \nabla_y B_y&:=\partial_yB_y+[A_y,B_y],\\
 \partial_y(W_yB_yW_y^{-1})&=W_y(\nabla_yB_y)W_y^{-1},\\
 \nabla_y C_y&=-C_y(\nabla_yB_y)C_y,\qquad C_y=B_y^{-1}.
\end{aligned}
\label{eq:update-measure-response}
\end{equation}
These follow by differentiating the unitary identification and
$B_yC_y=\one$. At kernel level $W_yC_yW_y^{-1}$ has a factor
$\sqrt{r_y}$ at both endpoints relative to the reference measure.
Sources held fixed in this reference representation acquire the same
half-density transport; a different source convention adds its own
variation. A fixed-background convergence theorem therefore does not provide a
metric-insertion bound by itself. The moving-support coarea and
raw-action variation needed for such an extension are developed in
Refs.~\cite{ChiapponeGHOUQFT2026,ChiapponePartI2026,ChiapponePartII2026}.


The identity kernel in \eqref{eq:source-identity} does not absorb
the coefficient of the physical source coupling. On this fixed
support the latter is
\begin{equation}
\begin{aligned}
 a_{{\rm src},\epsilon,\bb}
   &:=\varsigma_{\rm rep}(s)u_{\epsilon,\bb},\\
 \varsigma_{\rm OQFT}(s)&=\sigma_s,\\
 a_{{\rm src},\epsilon,\bb}&=\eta\ee\quad\text{in standalone O--QFT},\\
 \eta&:=\sigma_s\sigma_\epsilon.
\end{aligned}
 \label{eq:physical-source-carrier}
% Monograph source: eq:A5-postGTL-connected-cumulant-rule; def:W-Gamma-corrected.
\end{equation}
The metric-origin representative of HOUQFT uses
\(\varsigma_{\rm HOU}(s)=1\); that different ledger is not obtained
by inserting another O--QFT sector fold. We denote physical source
coordinates by \(\mathcal J_A\), reserving \(J_A\) in the
Euclidean BRST calculation for canonically normalized coefficient
sources. Any map from auxiliary Wick fields to physical source
fields is part of the completed source pairing. Neither the label
identity nor the carrier factor determines that map.

Repeated support labels never imply a sum. For an
\(n\)-leg amplitude in an atomic ledger, the external binder is
\begin{equation}
 \mathcal P^{(n)}_{\bb\mid{\rm ext}}
   =\prod_{\ell=1}^n\delta_{s s_\ell}
                         \delta_{b_s b_{s_\ell}}.
 \label{eq:external-binder}
\end{equation}
This binder restricts all attached legs to one support in the diagonal
model. It is not the binder for a mixed N--A creation channel, which
requires support-changing vertices absent from that model. Support
diagonality gives a selection rule, not a proof of unitarity.

\subsection{Virtual-line depths and placement}
In the monograph's Postulate~\(\Omega\), the positive virtual-lapse
law and its depth assignment are additional model data. Coordinate
locking, a unit-modulus endpoint phase, or an angular or Gaussian
support average does not uniquely determine that law.
For a graph \(G\), the nondegenerate assignment is
\begin{equation}
 e\in E_{\rm u}(G)\longmapsto a_{e,\bb_e}>0,
 \qquad \ell_{\Omega,e,\bb_e}=\hbar a_{e,\bb_e},
 \label{eq:virtual-edge-depths}
% Monograph source: post:Omega-subordination; eq:Omega-edge-depth-assignment.
\end{equation}
where \(E_{\rm u}(G)\) consists of uncut virtual edges. If
\(\alpha:G\to G'\) is a label-preserving graph isomorphism, and
\(\mathcal I_{A\omega\to Nz}\) is a declared pairing of normal
and antiverse support representatives, covariance of the assignment
requires
\begin{align}
 a_{\alpha(e),\bb'_{\alpha(e)}}&=a_{e,\bb_e},
 \label{eq:depth-isomorphism-covariance}\\
 a_{\mathcal I_{A\omega\to Nz}(e),Nz}&=a_{e,A\omega}.
 \label{eq:depth-pairing-covariance}
\end{align}
These conditions concern the declared edge data. Coincident projected
endpoints, including a tadpole, do not set its depth to zero.
Depth additivity is imposed only for a separately declared concatenation
of parameterized virtual segments; inserting a bivalent Feynman vertex
does not by itself impose such a composition law.

% Monograph source: post:Omega-subordination; eq:Omega-lapse-and-depth-dimensions; cor:fixed-order-uniform-UV-package.
For the fixed-plane graphs below every edge has $\bb_e=\bb$.
Retaining the edge slot as well, write its dimensionful heat measure
as \(\rho^{\rm heat}_{e,\bb}(\dd\lambda)\). At its fixed depth
\(a_e>0\), the change of variable \(\tau=\lambda/a_e^2\)
defines a dimensionless measure \(\rho_{e,\bb}\), with
\[
\begin{aligned}
 \Phi^{\rm edge}_e(q_E^2)
 &=\int e^{-\lambda q_E^2}\rho^{\rm heat}_{e,\bb}(\dd\lambda)\\
 &=\int e^{-a_e^2\tau q_E^2}\rho_{e,\bb}(\dd\tau),\\
 \int\lambda^{-r}\rho^{\rm heat}_{e,\bb}(\dd\lambda)
 &=a_e^{-2r}M_r[\rho_{e,\bb}].
\end{aligned}
\]
The inverse-moment equality is over the positive half-lines, with
absence of an atom at zero imposed separately when required.
The common-depth formulas below use the allowed benchmark
\(a_e=a>0\), \(\rho_{e,\bb}=\rho_{\bb}\), hence
\(\Phi^{\rm edge}_e(q_E^2)=\Ph(a^2q_E^2)\).
One common rescaled measure is a specialization, not a consequence
of the general depth assignment.

For a fixed finite graph, unequal positive depths and edge measures
are handled by retaining \(a_e\) and the corresponding linewise
gap or moments in the estimates. This gives no uniform theorem for
an infinite graph family. Such a theorem additionally needs control
of the effective damping data, routing and graph constants, vertex
growth, and any graph or label sums. With a fixed rescaled kernel
family, a positive uniform lower bound on the controlling depths is
one sufficient ingredient in those estimates.

The canonical branch of Postulate~\(\Omega\) selects
\(\Ph(z)=e^{-\sqrt z}\), hence \(e^{-a|q_E|}\), as an
additional stable--\(1/2\) specialization. The other kernel families
and the Gaussian bubble below are comparison models within the stated
mathematical class. Once a law has been selected, changing it is a
change of theory data, not a finite mass, residue, or coupling matching
condition.

\subsection{Background operator and internal covariance}

Write the total gauge connection as
\(\mathcal A_\mu=\overline{\mathcal A}_\mu+a_\mu\).
In Euclidean signature choose
\begin{equation}
 \mathscr L_{\bb}=-\overline D_E^{\,\mu}\overline D_{E,\mu}
    \succeq0,\qquad
 \Ph(a^2\mathscr L_{\bb})
     =\int e^{-a^2\tau\mathscr L_{\bb}}\rho_{\bb}(\dd\tau)
 \label{eq:positive-operator}
\end{equation}
on a specified self-adjoint domain. The minus sign makes the flat
eigenvalue \(q_E^2\) nonnegative. Writing
\(\overline D^2\) without this sign convention would obscure which
heat semigroup is being used.

In a common flat spectral representation, the covered internal
covariance is \eqref{eq:covariance-intro}. More general background
Hessians, spin terms, and gauge projectors need not commute with
\(\Ph(\mathscr L_{\bb})\); their ordering and domain must come from
the completed action. No general commutation with a gauge projector
is assumed here. The momentum-space theorem below concerns the
displayed flat representation, or any specified background reduction
for which its absolute envelope has been established.

External LSZ attachments receive no direct \(\Ph\) multiplier.
Operational physical cuts use their declared on-shell projectors.
Both are interfaces of the completed theory. Neither interface is
derived by putting an external momentum into
\(\Ph(0)\).

\begin{proposition}[Free source normalization and its interacting limit]
\label{prop:free-source-normalization}
% Monograph source: prop:source-normalized-realization; rem:source-normalized-why.
On a fixed commuting Euclidean field block, let
\(\mathscr K_0\) be a positive self-adjoint injective free kinetic
operator and \(\Delta=\mathscr K_0^{-1}\). Put
\[
 F=\Ph(a^2\mathscr L_{\bb}),\qquad \Delta_\Ph=F\Delta.
\]
Assume that \(F\) is positive and injective and strongly commutes
with \(\mathscr K_0\): their spectral projections commute.
Work on the common reducing spectral core consisting of vectors
whose joint spectral support is bounded above and separated from
zero in each spectral coordinate, taking the union over such
bounds. Every displayed product
and inverse is then defined there. In particular,
\(T:=F^{-1/2}\) is self-adjoint on its spectral domain.
Introduce an auxiliary Wick field \(\varphi_{\bb}\) with covariance
\(\Delta_\Ph\). The notation \(\varphi_{\bb}^{\rm phys}=T\varphi_{\bb}\)
means the smeared-field identity
\(\varphi_{\bb}^{\rm phys}(f)=\varphi_{\bb}(T^\dagger f)\) on a common
admitted test space; it does not apply an unbounded operator to a
pointwise random field. The covariance forms must extend to any
larger test space subsequently used for local observables.
We use the Hermitian covariance pairing for a complex commuting
field; on a compatible real field block the adjoint below is the
transpose with respect to the real source pairing.
Then its free source--source covariance is
\begin{equation}
 \langle\varphi_{\bb}^{\rm phys}
             (\varphi_{\bb}^{\rm phys})^\dagger\rangle_0
 =T\Delta_\Ph T^\dagger
 =T\Delta_\Ph T=\Delta.
 \label{eq:undressed-source-covariance}
\end{equation}
Contractions between auxiliary interaction insertions still use
\(\Delta_\Ph\). However, a mixed source--vertex contraction is
\begin{equation}
 T\Delta_\Ph=F^{1/2}\Delta,
 \qquad
 (T\Delta_\Ph)\Delta^{-1}=F^{1/2}.
 \label{eq:mixed-source-half-factor}
\end{equation}
Thus free source normalization alone does not prove an undressed
interacting external attachment.
\end{proposition}

\begin{proof}
In the common reducing spectral representation,
\(F\), \(\Delta\), and \(T\) are simultaneous spectral
multipliers, and \(T^\dagger=T\). Thus on the stated core,
\[
 T\Delta_\Ph T^\dagger
 =F^{-1/2}F\Delta F^{-1/2}=\Delta,
 \qquad
 T\Delta_\Ph=F^{-1/2}F\Delta=F^{1/2}\Delta.
\]
These are the two asserted covariance identities. The mixed factor
survives ordinary amputation unless the completed endpoint map or
compatible interactions compensate it, or an appropriate on-shell
normalization is imposed separately.
\end{proof}

For a damping kernel with \(\Ph(q)\to0\), \(T\) is generally
unbounded; the domain clause is essential. The proposition supplies
a free Gaussian realization of the source convention. It does not
construct the completed interacting endpoint map, prove the BRST
master equation, or derive the physical cut rule. Spectral commutation
does not say that \(T^\dagger\) preserves compact or positive-GT
support, so it is not a locality or reflection-positivity argument. The source couplings
in the completed functional must distinguish these physical source
coordinates from auxiliary Wick variables, or display the corresponding
source map explicitly. For other components of the BRST multiplet the corresponding identity
is algebraic: the right endpoint operator is the transpose or graded
dual prescribed by the declared source pairing. Its domain and
quadratic-form signs must be checked blockwise. No positive bosonic
Gaussian measure is thereby assigned to fermions or ghosts.

\begin{proposition}[A pure field redefinition cancels the damping]
\label{prop:pure-redefinition-cancellation}
In the real commuting block of
Proposition~\ref{prop:free-source-normalization}, suppose that the
entire interaction is chosen as \(V[T\varphi_{\bb}]\), with sources
coupled to \(T\varphi_{\bb}\), and that the quadratic coefficient is
\(\Delta_\Ph^{-1}\). At fixed background, take an invertible
finite-mode realization preserving the quadratic identities below
and include its change-of-variables Jacobian. Its normalized source
functional equals that of the undamped theory with quadratic
coefficient \(\mathscr K_0\) and interaction \(V\).
In particular this construction is not, by itself, an ultraviolet
improvement of that theory.
\end{proposition}

\begin{proof}
Put \(\chi_{\bb}=T\varphi_{\bb}\). Strong commutation gives
\[
 \Delta_\Ph^{-1}=T\mathscr K_0T,
 \qquad
 \tfrac12\langle\varphi_{\bb},\Delta_\Ph^{-1}\varphi_{\bb}\rangle
   =\tfrac12\langle\chi_{\bb},\mathscr K_0\chi_{\bb}\rangle.
\]
The interaction and source term become \(V[\chi_{\bb}]\) and
\(J\cdot\chi_{\bb}\). The finite-mode Jacobian is independent of
fields and sources and cancels in the normalized functional.
Equivalently, every internal line of the transformed vertices has
one \(T\) factor at each endpoint, hence
\begin{equation}
 T(q)\,[F(q)\Delta(q)]\,T(q)=\Delta(q).
 \label{eq:internal-damping-cancellation}
\end{equation}
The cancellation holds graph by graph before any ultraviolet limit.
If the corresponding undamped coefficient diverges when the
finite-mode restriction is removed, this reparametrization leaves
that divergence unchanged.
\end{proof}

For example, \(V[\chi_{\bb}]=\lambda\int\dd\mu_{E,\bb}\,\chi_{\bb}^4/4!\) produces the
ordinary undamped tadpole after its two internal endpoint factors
are included. The argument is an exact test of this particular
choice of interaction, not a theorem excluding other completions.
It sharpens the distinction between auxiliary Wick contractions and
the full vertex-weighted graph. The monograph's free source identity
is retained; it does not select \(V[T\varphi_{\bb}]\), and the extra
interacting test here is not attributed to that identity.

\paragraph{The charged determinant under a source-dependent change of fields}
There is a corresponding test for a proposed charged lift. On a finite
regulated space let $D[A]$ be the invertible local Dirac matrix, and let
$T[A]>0$ be a differentiable invertible endpoint map on the same space.
Include every matrix determinant factor in the chosen quantum measure.
For a candidate auxiliary operator $D_T[A]=T[A]D[A]T[A]$, its normalized
fermionic Gaussian determinant obeys the exact identity
\begin{equation}
 \frac{\det D_T[A]}{\det D_T[0]}=
 \frac{\det D[A]}{\det D[0]}
 \left(\frac{\det T[A]}{\det T[0]}\right)^2 .
 \label{eq:uvpeB-covariant-lift}
\end{equation}
Multiplicativity of finite determinants proves the formula also when
$D[A]$ and $T[A]$ do not commute. If $D_T$ is obtained by changing
variables in the original charged functional, set $\psi=T\chi$ and
$\bar\psi=\bar\chi T$ and transform the charged sources with these
fields. The Berezin Jacobian is $(\det T[A])^{-2}$. It then cancels
the two endpoint determinants before differentiating with respect to
the photon field, including at contact order. With an independently
declared flat $\chi$ measure the extra factor survives, and its source
variations must instead be evaluated. For example the second variation
of its logarithm in one direction is
\[
 2\Tr\!\left(T^{-1}T''-T^{-1}T'T^{-1}T'\right).
\]
The equality concerns a common regulated product of the matrices.
Compressing an infinite-dimensional product instead of multiplying its
compressed factors need not implement this change of variables.
The GCM reduction must supply its own measure and ordering before
either conclusion can be transferred to its O--QED sector.

\subsection{Geometric transfer and the complete mixed-source functional}
\label{sec:update-geometric-source}

The normal-transfer construction of
Refs.~\cite{ChiapponeGHOUQFT2026,ChiapponePartII2026} makes the forward
source map explicit. Work in one real bosonic coefficient block on
$L^2(\dd\mu_{E,\bb})$ with nonnegative self-adjoint $L_{\bb}$.
The symbol $x\ge0$ below is a kinetic \emph{spectral value}, such as
$q_E^2$; it is not an event coordinate or the vector $q_E$.
For $M\ge0$ define
\begin{equation}
\begin{aligned}
 H_{M,\bb}&=\sqrt{L_{\bb}+M^2\one}-M\one,\\
 P_{r,M,\bb}&=e^{-rH_{M,\bb}},\qquad r>0,\\
 h_M(x)&=\sqrt{x+M^2}-M.
\end{aligned}
\label{eq:update-forward-transfer}
\end{equation}
In the natural units used here $r$ is an inverse-momentum depth,
not $G_{0\bb}$, a clock reading, or a heat-lapse integration variable.
For boundary datum $\chi_{\bb}$ on the admitted spectral domain, the
bounded/decaying solution of the specified half-line problem is
\begin{equation}
\begin{gathered}
 (-\partial_r^2+L_{\bb}+M^2)U_{\bb}=0,\qquad U_{\bb}(0)=\chi_{\bb},\\
 U_{\bb}(r)=e^{-r\sqrt{L_{\bb}+M^2}}\chi_{\bb},
 \qquad \psi_{r,\bb}=e^{rM}U_{\bb}(r)=P_{r,M,\bb}\chi_{\bb}.
\label{eq:update-normal-equation}
\end{gathered}
\end{equation}
The square-root transfer is derived from this local propagation problem,
rather than chosen afterward as a removable cutoff. The underlying
extension/Dirichlet-to-Neumann construction is standard
\cite{StingaTorrea2010}; its identification with a specific ambient
A3/A4 channel is additional action and domain information. In particular,
GTL does not select the normal operator or its positive Euclidean sign.
For a Gaussian boundary field, the last relation in
\eqref{eq:update-normal-equation} is a smeared identity:
$\psi_{r,\bb}(f)=\chi_{\bb}(P_{r,M,\bb}^*f)$ on the admitted covariance
form domain. It does not require a typical Gaussian sample to have finite
classical action or to lie in the operator domain of $L_{\bb}$.
For $K>1$, its proposed ambient ray has
$X_{\rm GT}=(X_0+b_s)\hat G_{\bb}+r\hat n_\perp$ in the local
$\hbar=1$ units, with $\hat n_\perp\cdot\hat G_{\bb}=0$ and
$\|\hat n_\perp\|=1$. The ray coordinate is not $X_0$ or clock time.
The ray measure and operator are additional choices; they are not the
full radial Laplacian in $K-1$ dimensions. At $K=1$ no transverse GT
ray exists. For a mode of $L$ with value $x$, normal propagation gives
$e^{-r\sqrt{x+M^2}}$: larger tangential wave numbers penetrate a
shorter depth. Two endpoint traces then multiply to $e^{-2r h_M(x)}$
on a commuting block. Coincident retained coordinates do not set
$r$ to zero. Passing from $\RGTL$ to $\RMink$ transports this same
operator; the isometry alone neither creates nor removes the damping.

A deterministic normal extension and an independently fluctuating bulk
are different theories. At finite mode number the former has the
conditionally normalized measure
\begin{equation}
\begin{aligned}
 \dd\Pi_{\bb}(z\mid\chi_{\bb})
   ={}&|\det D_{D,\bb}|\,\dd^{N_{\bb}}z_{\bb}\\
 &\times\prod_{A=1}^{N_{\bb}}
 \delta^{\rm con}_{\bb,A}
       ((D_{D,\bb}z_{\bb}+D_{B,\bb}\chi_{\bb})_A),
\end{aligned}
\label{eq:update-normal-constraint}
\end{equation}
where $D_{D,\bb}$ is invertible. For a test functional $F$, the ordinary finite-dimensional change
$y=D_{D,\bb}z+D_{B,\bb}\chi_{\bb}$ gives
\[
 \int\dd\Pi_{\bb}(z\mid\chi_{\bb})F(z)
   =F(-D_{D,\bb}^{-1}D_{B,\bb}\chi_{\bb}).
\]
In particular, $\int\dd\Pi_{\bb}(z\mid\chi_{\bb})=1$; there is no residual
Gaussian noise. This is a
\emph{functional} constraint, not a second coordinate GTL delta.
In the continuum use the deterministic pushforward on an existing
Gaussian/test-function space, not a formal infinite determinant.
A background-dependent fiber Jacobian must be varied together with
its normalization; this does not remove the retained boundary
partition determinant or selector stress from a coupled A4 variation.

Let $C_{0,\bb}$ be a retained covariance, and let
$\mathsf w_{\bb}(\dd r_{\bb})$ be the specified probability law of one
shared depth. With real sources $J_{\bb}$ and $\mathfrak j_{\bb}$,
the full interacting source functional is
\begin{equation}
\begin{aligned}
 Z_V[J,\mathfrak j]
  =\frac1{\mathcal N_V}\int\mathsf w_{\bb}(\dd r_{\bb})\,
     \mathbb E_{C_{0,\bb}}\exp\bigl[&-V[P_r\chi_{\bb}]\\
      &+\langle J+P_r^*\mathfrak j,\chi_{\bb}\rangle_{E,\bb}\bigr],\\
 \mathcal N_V&=Z_V^{\rm unnorm}[0,0].
\end{aligned}
\label{eq:update-mixed-functional}
\end{equation}
All operators and fields in this fixed-copy formula retain $\bb$;
on the real block, the pairing of \eqref{eq:final-integration-ledger}
reduces to $\int\dd\mu_{E,\bb}\,fg$. Complex covariance formulas use its
Hermitian polarization; an independently complexified source functional
uses the corresponding bilinear or graded dual pairing, not an undeclared
Hermitian replacement of its analytic source variables.
The letter $\mathfrak j_{\bb}$ is an interaction source, distinct from
$j_{\bb}=\ii e_{\bb}$ in the pole algebra. The displayed integrals must exist on the declared source domain. For
example, if $C_0$ is bounded, $P_r$ is a real contraction, and $V\ge0$
is finite almost surely under the joint law, then $0<\mathcal N_V\le1$.
For real $L^2$ sources the numerator is bounded by
\[
 \exp\!\left[\tfrac12\|C_0\|
       (\|J\|_{E,\bb}+\|\mathfrak j\|_{E,\bb})^2\right],
\]
independently of $r$. The Gaussian moment-generating identity gives this
bound after dropping $e^{-V}\le1$. Unbounded covariances and interactions
not bounded below need their own domain and integrability estimates.
This is one global normalization.
Normalizing separately at every $r$ removes its matter reweighting
$\mathsf w_{\bb}(\dd r_{\bb})Z_{\chi,\bb}(r)/\mathcal N_V$ and changes the theory.
Here $Z_{\chi,\bb}(r):=\mathbb E_{C_{0,\bb}}e^{-V[P_r\chi_{\bb}]}$
is the conditional matter partition factor, not a new depth prior.

At fixed $r$ and $V=0$, completing the Gaussian square gives
\begin{equation}
\begin{aligned}
 Z_r[J,\mathfrak j]
  &=\exp\left[\tfrac12\langle J+P_r^*\mathfrak j,
               C_0(J+P_r^*\mathfrak j)\rangle_{E,\bb}\right],\\
 \mathbb C_{r,\bb}
  &=\begin{pmatrix}
      C_0&C_0P_r^*\\ P_rC_0&P_rC_0P_r^*
     \end{pmatrix}_{\bb}.
\end{aligned}
\label{eq:update-three-source-blocks}
\end{equation}
On a commuting block this reproduces
Proposition~\ref{prop:free-source-normalization} with
$P_r=F^{1/2}$ and $T=P_r^{-1}$ on its reducing spectral core.
The forward formula is meaningful without applying the unbounded
inverse $T$ to arbitrary sources. The two source slots represent
one boundary field and its transferred observable, not two
independent Gaussian fields.

\begin{proposition}[Exact condition for the one-sided virtual covariance]
\label{prop:update-commuting-transfer}
Let $L\ge0$ be self-adjoint, let $C_0$ be bounded and positive, and
let $P_r=e^{-r h_M(L)}$, $r>0$. Then
\begin{equation}
 P_rC_0P_r=P_r^2C_0
 \quad\Longleftrightarrow\quad [C_0,P_r]=0.
\label{eq:update-commuting-transfer}
\end{equation}
This is equivalent to commutation of $C_0$ with the spectral
projections of $L$. Without it the exact remainder is
$P_r[C_0,P_r]$.
\end{proposition}
\begin{proof}
Subtract the two expressions. The bounded operator $P_r$ has zero
kernel, so $P_r[C_0,P_r]=0$ implies $[C_0,P_r]=0$ without any bounded
inverse assumption. Commutation with a bounded self-adjoint operator
implies commutation with its spectral projections. The scalar
function $e^{-r h_M(x)}$ is injective, so those projections recover
those of $L$. The converse follows by functional calculus.
\end{proof}
For $C_0=(L+m^2)^{-1}$, $m>0$, the criterion is automatic and
\begin{equation}
 \widetilde C_{\psi_{\bb}\psi_{\bb},r}(q_E)
  =\frac{e^{-2r h_M(q_E^2)}}{q_E^2+m^2},\qquad a=2r.
\label{eq:update-flat-geometric-covariance}
\end{equation}
A general spin or reduced geometric Hessian need not commute with
$L$. Unbounded massless covariances require common covariance-form
domains, not just the bounded-operator proof. The article's scalar
momentum estimates apply after an admitted flat reduction or a
separately proved complete envelope.

\begin{proposition}[Conditional determinism and non-Gaussian depth mixtures]
\label{prop:update-depth-source-law}
For centered fields on a common real smeared-field domain with finite
second moments, the three blocks in \eqref{eq:update-three-source-blocks} imply
$\psi_{\bb}(f)=\chi_{\bb}(P_r^*f)$ almost surely for each test at fixed $r$.
For a shared random depth, put
$v_f(r)=\langle f,P_rC_0P_r^*f\rangle_{E,\bb}$.
If $v_f^2$ is integrable, the centered conditional Gaussian model has
\begin{equation}
 \langle\psi_{\bb}(f)^4\rangle_{0,c}
      =3\operatorname{Var}_{\mathsf w}(v_f).
\label{eq:update-depth-fourth-cumulant}
\end{equation}
Under the same square-integrability hypothesis, if its conditional
translation-invariant correlations $c_r(R)=\mathbb E_r[\psi_{\bb}(f)\psi_{\bb}(f_R)]$ tend to zero, then
\begin{equation}
 \lim_{\|R\|\to\infty}
 \operatorname{Cov}_{\mathsf w}(\psi_{\bb}(f)^2,\psi_{\bb}(f_R)^2)
      =\operatorname{Var}_{\mathsf w}(v_f).
\label{eq:update-depth-cluster}
\end{equation}
\end{proposition}
\begin{proof}
The variance of $\psi_{\bb}(f)-\chi_{\bb}(P_r^*f)$ is zero by all three blocks,
and its mean is zero by centering. For noncentered fields the additional
condition $\mathbb E\psi_{\bb}(f)=\mathbb E\chi_{\bb}(P_r^*f)$ is necessary;
covariance matching alone permits an arbitrary deterministic shift.
A countable separating test core gives simultaneous equality whenever
the random distributions are defined; Gaussianity is not needed for
this implication. Conditional Gaussian moments give
$\mathbb E\psi_{\bb}(f)^4=3\mathbb E v_f^2$ and
$\operatorname{Cov}(\psi_{\bb}(f)^2,\psi_{\bb}(f_R)^2)
=2\mathbb E c_r(R)^2+\operatorname{Var}(v_f)$.
Since $|c_r(R)|\le v_f(r)$, dominated convergence proves the limit.
\end{proof}
For a nonzero flat massive test packet, $v_f$ decreases strictly with
$r$ unless its spectrum is wholly in a transfer zero mode. A
nondegenerate common-depth state can therefore fail clustering on
this two-source algebra even when its boundary-only free marginal
clusters. This is a correlation statement, not a signaling claim.
Its full generating functional, not its marginal two-point function,
defines higher contractions and connected functions. In particular
$\log\int\mathsf w_{\bb}(\dd r_{\bb})Z_r$ is not generally
$\int\mathsf w_{\bb}(\dd r_{\bb})\log Z_r$.

The underlying A3/A4 support functional does not acquire
\eqref{eq:update-normal-equation} merely by postponing GTL evaluation.
Normal variations invisible to all retained fields remain invisible
to a functional depending only on those fields. The construction
above is an identified \emph{enlargement}: it specifies the normal
operator, boundary domain, deterministic constraint, interaction map,
and depth measure. Once specified, it derives the transfer and its
three covariance blocks. It does not prove that unchanged A3/A4
uniquely supplies those inputs. The applicable physical reconstruction
tests are stated in Section~\ref{sec:analytic}.


\subsection{Carrier algebra and directed boundary values}
\label{sec:carrier}

The fixed geometric pole carrier, its algebraic representative, and
the intrinsic GT axis are identified by
\begin{equation}
\begin{aligned}
 \ee&:=\hat p_{0,\bb}=\hat G_{\bb},\qquad \ee\!\cdot\!\ee=1,\\
 \mathfrak A_{\bb}&=\CC[\ee]/(\ee^2-1),\\
 \je&=\ii\ee,\qquad\je^2=-1.
\end{aligned}
 \label{eq:carrier-algebra}
\end{equation}
The geometric norm and the Clifford product \(\ee^2=1\) are separate
structures. The generator is odd in the \emph{carrier} grading and
Grassmann-even; it commutes with field, spinor, and BRST coefficients.
Thus \(\ee^{2n}=1\) and \(\ee^{2n+1}=\ee\), while all ordinary
Grassmann signs are retained independently.

The directed complex line and its real-algebra isomorphism are
\begin{equation}
 \CC_{\ee}=\operatorname{span}_{\RR}\{1,\je\},\qquad
 \embed(x+\ii y)=x+\ii y\ee,\qquad x,y\in\RR.
 \label{eq:directed-embedding}
\end{equation}
The central complex line \(\operatorname{span}_{\RR}\{1,\ii\}\)
is a different subalgebra. An untransported complex phase times a
directed coefficient need not lie in \(\CC_{\ee}\).
For \(A=A_0+A_1\ee\), with \(A_0,A_1\in\CC\), define
\begin{align}
 \chi_\kappa(A)&=A_0+\kappa A_1,&
 \mathsf P_\kappa&=\frac{1+\kappa\ee}{2},&\kappa&=\pm1,
 \nonumber\\
 A&=\sum_{\kappa=\pm1}\mathsf P_\kappa\chi_\kappa(A),&
 \mathsf P_\kappa\mathsf P_{\kappa'}
   &=\delta_{\kappa\kappa'}\mathsf P_\kappa.
 \label{eq:carrier-characters}
\end{align}
The character label \(\kappa\) is independent of the sector and
evolution labels. Carrier-even and carrier-odd parts are respectively
\(A_0\) and \(A_1\ee\). The inverse \(\embed^{-1}\) is defined
only on \(\CC_{\ee}\), where
\(\chi_+\embed(z)=z\) and \(\chi_-\embed(z)=\overline z\).
In particular, the tag map \(z\mapsto z\ee\) is neither
\(\embed\) nor an algebra homomorphism.

With independent complex and Grassmann sources the coefficient algebra is
\begin{equation}
 \mathfrak B_{\bb}
   =\mathfrak A_{\bb}\widehat\otimes_{\CC}\mathfrak G,
 \label{eq:graded-carrier-source-algebra}
\end{equation}
where \(\mathfrak G\) is the declared Grassmann algebra and
\(\mathfrak A_{\bb}\) lies in its graded centre. The characters
extend by the identity on \(\mathfrak G\). The carrier reality
involution is
\[
 (z_+\mathsf P_++z_-\mathsf P_-)^{\ddagger}
    =\overline{z_-}\mathsf P_++\overline{z_+}\mathsf P_-;
 \qquad \operatorname{Fix}(\ddagger)\cap\mathfrak A_{\bb}
    =\CC_{\ee}.
\]
% Monograph source: def:A5-directed-exp-log; eq:A5-postGTL-reality-graded-OQFT-pairings.
Its extension uses the declared antilinear graded field/source
involution \(\star\):
\[
 (A\widehat\otimes\gamma)^\ddagger
     =A^\ddagger\widehat\otimes\gamma^\star,
 \qquad
 (XY)^\ddagger=(-1)^{|X||Y|}Y^\ddagger X^\ddagger .
\]
Here \(X,Y\) are homogeneous and \(|\cdot|\) denotes Grassmann
parity. On the physical conjugate-source contour, the complete
unweighted pairing \(\int\dd^4x_{\bb}\,\mathcal J_A\phi^A\) is even
and \(\ddagger\)-invariant. Its product with
\(a_{{\rm src},\epsilon,\bb}\) is anti-fixed, whereas the full
source term \(\ii a_{{\rm src},\epsilon,\bb}
\int\dd^4x_{\bb}\,\mathcal J_A\phi^A\) in the Lorentzian exponent
is fixed, since \(\ee^\ddagger=-\ee\) and
\(\ii^\ddagger=-\ii\).
Independent source differentiation precedes restriction to a physical
conjugate-source contour; arbitrary complex sources need not give
conjugate character components.

\begin{proposition}[Carrier reality and positive adjoints]
\label{prop:carrier-reality-adjoint}
On the scalar carrier algebra, the reality operation \(\ddagger\)
interchanges the two character representations:
\begin{equation}
 \chi_\kappa(A^\ddagger)
       =\overline{\chi_{-\kappa}(A)}.
 \label{eq:reality-character-interchange}
\end{equation}
There is no normalized complex-linear functional \(\omega\) on the
full unital algebra satisfying
\(\omega(A^\ddagger A)\geq0\) for every \(A\).
In particular, \(\ddagger\) cannot be the adjoint of a unital
representation of this algebra on a nonzero positive Hilbert space.
\end{proposition}

\begin{proof}
The character formula follows from the definition of \(\ddagger\).
Taking \(A=\ee\) gives
\[
 \ee^\ddagger\ee=-\one,\qquad
 \omega(\ee^\ddagger\ee)=-\omega(\one)=-1.
\]
This contradicts the proposed positivity. Likewise a representation
preserving that involution would give
\(\pi(\ee)^\dagger\pi(\ee)=-I\), impossible on a nonzero
positive Hilbert space.
\end{proof}

For comparison, the same scalar algebra admits a different involution,
which we denote by \(\sharp\):
\begin{equation}
 \left(\sum_\kappa z_\kappa\mathsf P_\kappa\right)^\sharp
      :=\sum_\kappa\overline{z_\kappa}\mathsf P_\kappa.
 \label{eq:coefficient-positive-involution}
\end{equation}
It has \(\ee^\sharp=\ee\), makes each character a
\(\sharp\)-preserving map to ordinary complex numbers, and obeys
\(\|A^\sharp A\|_\chi=\|A\|_\chi^2\) for the norm defined
below. If \(\alpha(A_0+A_1\ee)=A_0-A_1\ee\), then
\(A^\ddagger=\alpha(A^\sharp)\). Thus the two involutions
are not interchangeable. Here \(\sharp\) is only an algebraic
comparison; it neither replaces the declared carrier reality nor
constructs a physical inner product.

The physical adjoint \(\dagger\) belongs to the separately admitted
observable representation, with its own involution on physical fields. No positive \(\ddagger\)-preserving
representation of the entire carrier algebra is assumed. In particular,
carrier reality of a generating functional, positivity of the
character norm, and positivity of physical state sums are three
separate statements. This distinction leaves the vectorial pole
prescription and the two-character reconstruction unchanged.

For a specified scalar function and branch, distinguish
\begin{equation}
 F_{\mathfrak A}(A)
   =\sum_\kappa\mathsf P_\kappa F(\chi_\kappa(A)),
 \qquad
 F_{\ee}=\embed\circ F\circ\embed^{-1}.
 \label{eq:carrier-functional-calculus}
% Monograph source: eq:general-fixed-carrier-functional-calculus; eq:general-character-versus-directed-functional-calculus.
\end{equation}
These agree on the directed line only on a conjugation-invariant
domain with \(F(\overline z)=\overline{F(z)}\) for the chosen
branch. Inverses, logarithms, square roots, and kernel functions are
always interpreted in the indicated calculus. Determinants and
functional logarithms are defined on the two character operators and
reconstructed afterward. The coefficient maps extend entrywise to
matrix-valued factors, preserving the order of their matrix products.

Keep the sector, evolution branch, and positive boundary magnitude
separate:
\begin{equation}
\begin{aligned}
 \sigma_s,\sigma_\epsilon&\in\{+1,-1\},\\
 \eta&:=\sigma_{s\epsilon}=\sigma_s\sigma_\epsilon,\\
 \epsvec&:=\varepsilon\hat p_{0,\bb},\qquad\varepsilon\downarrow0^+.
\end{aligned}
 \label{eq:signs}
\end{equation}
The branch label \(\epsilon\) and the continuous height \(\varepsilon\)
occupy different slots. The height has the units of the denominator to
which it is added. The sign \(\eta\) is outside \(\epsvec\).
On a paired N/A carrier class, use the same fixed pole carrier
\(\hat p_{0,Nb_N}=\hat p_{0,Ab_A}=\hat G\); the opposite
canonical mode orientations are recorded by \(\sigma_s\).
They do not reverse an observer's physical clock.

For a real invariant denominator \(D\), the prescribed boundary is
\begin{equation}
 D\longmapsto D+\ii\eta\epsvec,
 \qquad
 (D+\ii\eta\epsvec)^{-1}
   =\frac{D-\ii\eta\varepsilon\ee}{D^2+\varepsilon^2}.
 \label{eq:directed-inverse}
\end{equation}
This boundary height multiplies the fixed carrier in
$\CC_{\ee}$ and leaves the real routing variable $k_\parallel$
intact. In the $\chi_+$ coefficient the constant denominator
signs $\eta=+1$ and $\eta=-1$ give the Feynman and
anti-Feynman sides, respectively. The $\chi_-$ coefficient
has the opposite denominator sign. Retarded and advanced
prescriptions use frequency-dependent sides. The positive
branch-neutral A5 quadratic convergence regulator belongs
to the source functional.

Keep finite \(\varepsilon>0\) through the required source
differentiations and graph integrations, and take the common boundary
limit afterward. Exchanging that limit with an integral or derivative
requires a separate convergence argument. In particular, products
containing running units denote this common regulated expression;
they are not arbitrary products of already singular boundary
distributions. A formula for one line does not prove existence of a
full Lorentzian graph boundary.

The pre-GTL exported radial unit $\hat v_{\mathcal G}^{[\Gamma]}$, the
transverse unit $\hat n_{\perp,\bb}$, the branch action unit
$u_{\epsilon,\bb}$ and the fixed pole carrier $\ee$ are not synonyms.
The first requires the declared binder, the second specifies a normal
channel, and the last two remain in action assembly and inversion.
Only the running unit below depends on the integration momentum.

For a nonzero aligned GT momentum set
\begin{equation}
 k_{0k,\bb}=k_\parallel\ee,\qquad
 \widehat u_{k,\bb}
  =\frac{k_{0k,\bb}}{\norm{k_{0k,\bb}}}
  =\sgn(k_\parallel)\ee.
 \label{eq:running-carrier}
\end{equation}
The quotient is undefined at \(k_\parallel=0\). In an absolutely
convergent Lebesgue integral, its value on that null set can be chosen
measurably without changing the integral. This does not define a
normalized vector or a continuous limit there. Restricting a boundary
distribution or evaluating an external transfer on that hyperplane needs
a completed source pairing; see Section~\ref{sec:fw-elastic-benchmark}.

The running sign also distinguishes a mass shell from the full loop
domain. On a nonzero causal momentum, the sign of its oriented
frequency is invariant under proper orthochronous changes of inertial
frame. A spacelike momentum can have positive frequency in one frame
and negative frequency in another. In units with intrinsic boost speed
one, take \(0<k_\parallel<k_1\) and boost with
\(k_\parallel/k_1<v<1\). Its new frequency component is
\((k_\parallel-vk_1)/\sqrt{1-v^2}<0\). Thus an off-shell
running sign must retain the chosen readout covector. The pairing
itself is covariant when momentum and readout are transported together.

The monograph's oriented scalar-line prescription is the following
carrier-valued expression. Calling it a covariance of a particular
source field additionally requires the exchange test below:
\begin{equation}
 \widetilde G^\Ph_{\bb,\eta}(k)
 =\lim_{\varepsilon\downarrow0^+}
 \eta\ii\widehat u_{k,\bb}\,
 \frac{\Ph_{\ee}\!\left(-a^2[k^2+\ii\eta\epsvec]\right)}
      {k^2-m^2+\ii\eta\epsvec}.
 \label{eq:directed-propagator}
% Monograph source: prop:vectorial-poles; eq:general-directed-scalar-coefficient.
\end{equation}
Its continuation and existence as a boundary distribution require
the analytic hypotheses of Section~\ref{sec:analytic}. Its identification
with an inverse Hessian or source covariance requires the completed
source. On test domains where the running-unit multiplication
commutes with the boundary limit, the same expression can be written
using its extracted directed resolvent coefficient.
\begin{align}
 R_\eta^\Ph(k)
  &:=\lim_{\varepsilon\downarrow0^+}\embed^{-1}\!\left[
  \frac{\Ph_{\ee}(-a^2[k^2+\ii\eta\epsvec])}
       {k^2-m^2+\ii\eta\epsvec}\right],\nonumber\\
 \widetilde G^\Ph_{\bb,\eta}(k)
  &=\eta\ii\widehat u_{k,\bb}\embed(R_\eta^\Ph(k)).
 \label{eq:scalar-line-coefficient}
\end{align}
For example, the multiplication is smooth near the massive
shell when the numerator admits regular boundary values there.
At a singular intersection with $k_\parallel=0$, use the common
finite-height product or establish a multiplier domain before
taking the boundary limit. In scalar coefficients below,
$+\ii\eta0$ denotes the extracted form of $\ii\eta\epsvec$.

Only on a fixed-orientation ray is
\(\widehat u_{k,\bb}=\sigma_s\ee\). An explicitly declared
\emph{fixed virtual-carrier benchmark} may assign that unit to
every virtual line independently of its Fourier coefficient. In
that specialization,
\begin{equation}
 \widetilde G^\Ph_{\bb,\eta}
    =\embed(\sigma_\epsilon\ii R_\eta^\Ph).
 \label{eq:fixed-carrier-line}
\end{equation}
For the undamped fixed-carrier line the finite-magnitude decomposition
and its boundary are
\begin{align}
 \frac{\eta\ii\sigma_s\ee}{D+\ii\eta\epsvec}
 &=\sigma_s\frac{\varepsilon}{D^2+\varepsilon^2}
   +\sigma_\epsilon\ii\ee\frac{D}{D^2+\varepsilon^2},\nonumber\\
 \widetilde G_{\bb,\eta}
 &\longrightarrow\sigma_s\pi\delta^{\rm sh}_{\bb}(D)
     +\sigma_\epsilon\ii\ee\,\PV\frac1D.
 \label{eq:carrier-pv-delta}
\end{align}
Already one line has both carrier parities. Counting its explicit
unit numerators therefore does not determine the parity of a graph.
The Euclidean estimates below concern ordinary coefficient
representations; an unrestricted Lorentzian loop also retains every
running carrier sign and the complete branch data.

For finite \(\varepsilon>0\), let
\(R_\eta^\Ph(k;\varepsilon)\) denote the extracted expression
in \eqref{eq:scalar-line-coefficient} before its limit, and put
\(r_k=\sgn(k_\parallel)\) for a nonzero mode. On a
reflection-compatible kernel branch the full line has characters
\begin{equation}
 \chi_\kappa\!\left(\widetilde G^\Ph_{\bb,\eta}
                  (k;\varepsilon)\right)
 =\kappa\eta r_k\ii\,
  \frac{\Ph(-a^2[k^2+\ii\kappa\eta\varepsilon])}
       {k^2-m^2+\ii\kappa\eta\varepsilon}.
 \label{eq:line-character-values}
\end{equation}
Without that reflection hypothesis, compute the characters of the
specified directed function itself. The factors \(r_{q_i}\)
remain inside every unrestricted running-unit loop integral.

\subsection{Source exchange symmetry and oriented lines}
\label{sec:source-exchange}

\begin{proposition}[Exchange test for a source covariance]
\label{prop:source-exchange}
Fix the support, ordering branch, carrier, and source pairing. Let
\(C^{AB}(x,y)\) be a two-point coefficient obtained by two
independent source derivatives of the same even generating
functional, with the source-first convention of
Section~\ref{sec:bv}. Whenever the derivatives exist and graded
commute, it obeys
\begin{equation}
 C^{AB}(x,y)=(-1)^{\epsilon_A\epsilon_B}C^{BA}(y,x).
 \label{eq:source-covariance-reciprocity}
\end{equation}
For a translation-invariant coefficient this becomes
\(C^{AB}(k)=(-1)^{\epsilon_A\epsilon_B}C^{BA}(-k)\).
In particular, a single identical commuting real scalar has
\(C(k)=C(-k)\), in each carrier character and in the reconstructed
carrier algebra.
\end{proposition}

\begin{proof}
Interchanging the two left source derivatives gives precisely the
Koszul sign in \eqref{eq:source-covariance-reciprocity}.
The common even carrier and source-normalization factors do not
change it. Translation invariance turns endpoint interchange into
momentum reversal. The character projections preserve the identity.
\end{proof}

This necessary test distinguishes an ordered virtual datum from a
source-generated covariance. At fixed finite \(\varepsilon>0\),
the denominator and kernel in \eqref{eq:directed-propagator} are
even under \(k\mapsto-k\), whereas
\(\widehat u_{-k,\bb}=-\widehat u_{k,\bb}\) away from
\(k_\parallel=0\). Consequently its unrestricted running-unit
expression satisfies
\begin{equation}
 \widetilde G^\Ph_{\bb,\eta}(-k;\varepsilon)
       =-\widetilde G^\Ph_{\bb,\eta}(k;\varepsilon).
 \label{eq:running-line-reversal}
\end{equation}
The endpoint value on a measure-zero hyperplane does not repair
this mismatch. A nonzero expression of this type cannot itself be
the covariance of the single identical commuting scalar in the
proposition. This obstruction is already present at finite boundary
magnitude, so it is not cured by a choice of distributional limit
or by combining the two carrier characters.

An oriented two-field or endpoint-labelled kernel is a different
possibility: it must specify the reversed component and satisfy
\eqref{eq:source-covariance-reciprocity} for the complete source
matrix. No such completion is constructed by merely writing the
running unit. Nor does source differentiation at a fixed branch
reverse that branch automatically. The fixed virtual-carrier scalar
benchmark used in Section~\ref{sec:literal-bubble} is even and
passes this necessary test, but that fact alone proves neither
positivity nor unitarity.

Accordingly the running-unit formulas are retained as the declared
oriented line data of Ref.~\cite{ChiapponeGHOUQFT2026}, not promoted
here to a proved Gaussian realization for an identical real scalar.
The unrestricted running-unit source/endpoint realization remains
an explicit completion requirement. The test is an additional
consistency analysis, rather than a replacement of the monograph's
vectorial pole prescription. It also prevents importing the
fixed-carrier scalar benchmark into that unconstructed realization.

\subsection{Complete graphs and their character components}
\label{sec:graph-characters}

For every finite product on one fixed carrier,
\begin{equation}
 \prod_{r=1}^{N}A_r
  =\sum_{\kappa=\pm1}\mathsf P_{\kappa}
           \prod_{r=1}^{N}\chi_\kappa(A_r).
 \label{eq:common-character-product}
\end{equation}
The same character acts on every line, vertex, phase and endpoint
factor. There is no independent character sum for each internal
line: the mixed products vanish because
\(\mathsf P_+\mathsf P_-=0\).

For the complete carrier graph \(\vec{\mathcal G}_{s,\epsilon}\), define
\(\mathcal G^{\rm rep}_{s,\epsilon}
:=\chi_+(\vec{\mathcal G}_{s,\epsilon})\) and
\(\mathcal G^{\rm rep}_{s,\bar\epsilon}
:=\chi_-(\vec{\mathcal G}_{s,\epsilon})\), with
\(\sigma_{\bar\epsilon}=-\sigma_\epsilon\).
The opposite-branch notation refers to the negative-character
projection of the same complete functional. Its calculation uses
the projected orderer, Dyson exponent, vertices and adjoints, pole
prescriptions, external-state conventions and all running line signs.
The carrier-valued graph is therefore
reconstructed from these two projections, not from two unrelated
scalar theories. The bar on the branch label is bookkeeping for the
projected data; it is not complex conjugation. In particular,
\(\chi_-(\ii)=\ii\), so central phases are not conjugated merely
by changing characters. Explicitly,
\begin{equation}
 \vec{\mathcal G}_{s,\epsilon}
  =\mathsf P_+\mathcal G^{\rm rep}_{s,\epsilon}
   +\mathsf P_-\mathcal G^{\rm rep}_{s,\bar\epsilon}.
 \label{eq:complete-graph-character-reconstruction}
% Monograph source: eq:general-directed-graph-character-reconstruction; eq:general-directed-graph-parity-reconstruction.
\end{equation}
Its carrier-even and carrier-odd parts are respectively
\begin{equation}
 \frac{\mathcal G^{\rm rep}_{s,\epsilon}
             +\mathcal G^{\rm rep}_{s,\bar\epsilon}}{2},
 \qquad
 \frac{\mathcal G^{\rm rep}_{s,\epsilon}
             -\mathcal G^{\rm rep}_{s,\bar\epsilon}}{2}\,\ee.
 \label{eq:complete-graph-parity}
\end{equation}
Only when the complete reflection relation
\(\mathcal G^{\rm rep}_{s,\bar\epsilon}
=\overline{\mathcal G^{\rm rep}_{s,\epsilon}}\) has been
established may this be compressed to
\(\vec{\mathcal G}_{s,\epsilon}
=\embed(\mathcal G^{\rm rep}_{s,\epsilon})\).
For example, a central factor \(\ii\) has both character values
equal to \(\ii\); reconstructing it as
\(\embed(\ii)=\ii\ee\) would change the element.
The inverse directed map therefore cannot scalarize an arbitrary
mixed carrier expression.

\begin{proposition}[Character estimates for complete graphs]
\label{prop:character-majorants}
For scalar carrier coefficients define
\begin{equation}
 \norm{A}_{\chi}:=\max_{\kappa=\pm1}|\chi_\kappa(A)|.
 \label{eq:character-norm}
\end{equation}
This is a submultiplicative norm. Let \(\nu_{\bb}\) be a positive
integration measure. If a measurable complete graph integrand is
\(g(\xi)=\sum_\kappa\mathsf P_\kappa g_\kappa(\xi)\)
and \(|g_\kappa(\xi)|\leq M_\kappa(\xi)\), with
\(M_\pm\in L^1(\dd\nu_{\bb})\), then
\begin{equation}
\begin{aligned}
 \int\norm{g(\xi)}_\chi\dd\nu_{\bb}
 &\leq\int(M_++M_-)\dd\nu_{\bb}<\infty,\\
 \int g\,\dd\nu_{\bb}
 &=\sum_\kappa\mathsf P_\kappa\int g_\kappa\,\dd\nu_{\bb}.
\end{aligned}
 \label{eq:character-integral-majorant}
\end{equation}
The same statement holds uniformly in external parameters whenever
the two majorants are uniform. For matrix coefficients replace the
absolute value by a fixed submultiplicative matrix norm.
\end{proposition}

\begin{proof}
The character map identifies \(\mathfrak A_{\bb}\) with
\(\CC\oplus\CC\). Multiplicativity of each character proves
the norm inequality for products. The integral estimate follows from
\(\max(|g_+|,|g_-|)\leq M_++M_-\); the two finite-dimensional
projections commute with the absolutely convergent integral.
\end{proof}

This norm is an absolute-convergence device. Its compatibility with
\(\sharp\) in \eqref{eq:coefficient-positive-involution} does not
turn \(\ddagger\) into a positive adjoint or supply physical
state norms; Proposition~\ref{prop:carrier-reality-adjoint} prevents
that identification on the full carrier algebra.

For formal source or Grassmann expansions, apply this estimate after
forming the complete graph and extracting a specified coefficient.
Coefficient extraction alone is not multiplicative; products retain
the convolution over source degrees and their Koszul signs. Equivalently,
a finite Grassmann truncation admits a submultiplicative coefficient
norm before integration.

This proposition transfers the UV bounds below to the full carrier
graph when both representatives satisfy them. On real integration
domains the running units contribute signs of modulus at most one,
so they preserve an established absolute envelope. Holomorphy still
requires the hypotheses of Theorem~\ref{thm:holomorphy} in each
character's declared complex variables. In particular, \(\embed\)
is a real-algebra map: holomorphy in a directed argument is not an
unstated identification of \(\ii\) with \(\je\).


\section{Ultraviolet finiteness without divergent subtraction}
\label{sec:uv}

To test a bare graph on a fixed Euclidean plane, first identify
every independent loop direction. Then ask whether the virtual
covariances control those directions after the complete vertices
and contacts have been included. For independent heat variables the
Gaussian proof below gives a bound in terms of one finite inverse
moment. Selected line momenta give a more precise test when different
lines carry different numerator powers. If the lines share one
random depth, the corresponding test uses a joint moment. Sector
bounds are available when different ultraviolet regions require
different estimates.

These tests bound the absolute value of the complete graph integrand,
with its action-defined signs, tensor contractions and contact terms
already included. Any geometric integration that defines that integrand
is performed first; a remaining positive heat or depth law is controlled
by the stated joint majorant. Physical
application requires the same vertices as the BRST completion in
Section~\ref{sec:brst}. The analytic hypotheses for the Lorentzian
boundary values are given in Section~\ref{sec:analytic}.

\subsection{Graph data and infrared control}

Fix \(\bb\), a selected evolution branch, and a finite connected
Euclidean virtual graph \(G\) with
\(E_{\rm u}(G)=\{1,\ldots,I\}\) and \(L\geq1\)
independent loop momenta. Every line counted by \(I\) is an
\emph{uncut virtual covariance}; physical external attachments and
on-shell cut projectors are not factors in this comparison integral.
Write
\begin{equation}
\begin{gathered}
 q_i(k,p)=\sum_{\ell=1}^L C_{i\ell}k_\ell+R_i(p),\\
 k\in\RR^{4L},\qquad p\in\mathcal K,\qquad\rank C=L.
\end{gathered}
 \label{eq:routing}
\end{equation}
where \(\mathcal K\) is a compact set of real Euclidean external data.
The routing shifts \(R_i(p)\) are affine-linear and bounded on \(\mathcal K\).
Each row of $C$ tells which loop directions an internal line sees.
For example, take $q_1=k_1$, $q_2=k_1+p$, and $q_3=k_2$.
The first two rows have rank one, so their kernels have the
same values however large $k_2$ becomes at fixed $k_1$.
The first and third rows span the two independent loop directions.
This is the reason that coverage depends on the routing rank.

Hold the fields and orientations on their declared support. Assume
measurable graph factors and bound the full product of vertex and
numerator factors by
\begin{equation}
 \abs{N_G(k,p)}\leq C_{G,\mathcal K}\bigl(1+\norm{k}^D\bigr),
 \qquad D\geq0.
 \label{eq:numerator-bound}
\end{equation}
Norms in polynomial inequalities are expressed in a fixed reference
momentum unit; the constants restore the dimensions. Nonlocal
completion vertices are included in \eqref{eq:numerator-bound}.
A formal power series for a vertex does not establish this bound for
its untruncated sum.

Use nonnegative measurable Schwinger envelopes
\begin{align}
 C_i(q)&=\int_0^\infty E_i(\alpha)e^{-\alpha q^2}\dd\alpha,
 \label{eq:IR-envelope}\\
 \mathcal I_i&:=\int_0^\infty E_i(\alpha)\dd\alpha<\infty,
 &\mathcal E_i&:=\sup_{0<\alpha\leq\alpha_0}E_i(\alpha)<\infty ,
 \label{eq:IR-conditions}
\end{align}
for a fixed \(\alpha_0>0\). A positive mass gives
\(E_i(\alpha)=e^{-m_i^2\alpha}\) and
\(C_i(q)=(q^2+m_i^2)^{-1}\).
If a massless theory is considered, \eqref{eq:IR-envelope} is an
explicit infrared comparison prescription; an infrared-finite graph
can instead be treated with its own established IR estimate.
The theorem does not assert that all massless graphs are finite,
or that an auxiliary IR prescription can be removed. In particular,
selected massless denominators in
Theorem~\ref{thm:optimized-selected-lines} use a separate direct
line-integral argument and do not retroactively satisfy
\eqref{eq:IR-conditions}.

For clarity, first give every line the same kernel
\eqref{eq:bernstein}. A \emph{majorant} is a nonnegative integrand whose
integral bounds the absolute size of the graph. Here it is
\begin{equation}
 \mathcal A_G(p)=\int_{\RR^{4L}}
 \frac{\dd^{4L}k_{\bb}}{(2\pi)^{4L}}\,
 \abs{N_G(k,p)}
 \prod_{i=1}^I \Ph(a^2q_i^2)C_i(q_i).
 \label{eq:graph-majorant}
\end{equation}
If this positive integral is finite, it bounds each matrix component
and any fixed finite-dimensional tensor norm of the graph with the stated
denominators and infrared prescriptions. Thus the argument needs no
cancellation between components. This is a \emph{graphwise}
claim at fixed support and finite loop order, not a bound on the number
of graphs or their sum. An IR-modified covariance is not presumed to
majorize an unregulated massless covariance.

\subsection{Global fixed-graph theorem}

For \(r>0\), define the inverse moment
\begin{equation}
 M_r[\rho_{\bb}]=\int_{(0,\infty)}\tau^{-r}\rho_{\bb}(\dd\tau).
 \label{eq:inverse-moment}
\end{equation}
The condition that no atom is present at zero is imposed
\emph{separately}: the integral defining $M_r$ excludes that point.
For example, $M_r[\delta^{\rm heat}_{\bb,0}]=0$ for every $r>0$, even though the
associated multiplier is $\Ph(z)\equiv1$. For $d\geq0$, the
corresponding undamped massless comparison has
\[
 \int_{|q|<R}|q|^{d-2}\dd^4q_{\bb}
    =\frac{2\pi^2}{d+2}R^{d+2}\longrightarrow\infty.
\]
Thus finiteness of the inverse moment as defined here does not imply
the no-atom hypothesis. The latter is a separate premise of the
convergence statements below. Equivalently, for every $r>0$ the two
requirements can be recorded as one extended-valued integral,
\begin{equation}
 \begin{gathered}
 \rho_{\bb}(\{0\})=0,\qquad M_r[\rho_{\bb}]<\infty\\
 \Longleftrightarrow\quad
 \int_{[0,\infty)}\tau^{-r}\rho_{\bb}(\dd\tau)<\infty.
 \end{gathered}
 \label{eq:extended-inverse-moment}
\end{equation}
Here the integrand on the right takes the extended value
$0^{-r}:=+\infty$. That integral is \emph{not} the definition
of $M_r$ in \eqref{eq:inverse-moment}. This formulation makes the independent
atom condition explicit in a machine-checkable theorem statement.

\begin{theorem}[Global Euclidean convergence]
\label{thm:global-uv}
Assume \eqref{eq:routing}--\eqref{eq:IR-conditions}, finite
\(\Mzero\), and fixed \(a>0\). Either of the following is sufficient
for \(\sup_{p\in\mathcal K}\mathcal A_G(p)<\infty\):
\begin{description}[leftmargin=1cm]
\item[(G)] There is \(\tau_*>0\) with
\(\supp\rho_{\bb}\subset[\tau_*,\infty)\).
\item[(M)] \(\rho_{\bb}(\{0\})=0\) and
\(M_{1+D/2}[\rho_{\bb}]<\infty\).
\end{description}
The joint absolute majorant in loop momenta, Schwinger parameters,
and Laplace--Stieltjes parameters is integrable. All Euclidean
integration-order exchanges used in the proof are therefore valid.
No UV-divergent subtraction is needed to define this graph.
This conclusion concerns the stated Euclidean comparison integral;
it does not establish a physical spectral representation or the
validity of an undressed operational cut.
\end{theorem}

A more conservative sufficient condition is
\begin{equation}
 \int_{(0,1]}\tau^{-\beta}\rho_{\bb}(\dd\tau)<\infty,\qquad
 \beta>1+\frac D2,\qquad\rho_{\bb}(\{0\})=0.
 \label{eq:conservative-moment}
\end{equation}
It implies (M), but is not necessary for this proof. The endpoint
formulation requires the \emph{actual} inverse moment
$M_{1+D/2}[\rho_{\bb}]$ to be finite; borderline power-law behavior
alone does not establish that finiteness.

\begin{proof}
Introduce independent parameters \(\tau_i,\alpha_i\) and set
\begin{equation}
 w_i=\alpha_i+a^2\tau_i,\qquad W=\diag(w_1,\ldots,w_I),
 \qquad A=C^TWC .
 \label{eq:width-matrix}
\end{equation}
The two parameters enter the same Gaussian. The denominator supplies
$\alpha_i$ and the geo-phase supplies $a^2\tau_i$, so their sum is the
width $w_i$. Tonelli's theorem permits the loop integrals of the
nonnegative majorant at fixed widths before their convergence has
been proved. Since $C$ has full column rank, $A\succ0$ for every
strictly positive width vector. Completing the square in all four
Euclidean components gives
\begin{equation}
 \sum_iw_iq_i^2
  =(k+\bmu)^T(A\otimes\one_4)(k+\bmu)+\Delta(w,p),
 \qquad \Delta(w,p)\geq0 .
 \label{eq:complete-square}
\end{equation}

The Gaussian center depends on the external data. A bound on the
inverse of $A$ alone would become large when some widths shrink.
The exact center is instead an average of finitely many centers
computed from independent row sets. For one Euclidean component put
\[
 \mu=(C^TWC)^{-1}C^TWR.
\]
If \(\mathfrak B\) is the finite family of invertible \(L\)-row
submatrices \(C_T\), then
\begin{equation}
 \mu=\sum_{T\in\mathfrak B}\pi_T(w)C_T^{-1}R_T,\qquad
 \pi_T(w)=
 \frac{(\det C_T)^2\prod_{i\in T}w_i}
      {\sum_{U\in\mathfrak B}(\det C_U)^2\prod_{i\in U}w_i}.
 \label{eq:shift-mixture}
\end{equation}
The weights are nonnegative and sum to one.
The identity is proved in \ref{sec:appendix-gaussian}.
Since $R$ is bounded on $\mathcal K$, this gives
$\norm{\bmu}\leq C_{G,\mathcal K}$ even when different widths
approach zero at different rates.

After the change of variables
\(y=(A^{1/2}\otimes\one_4)(k+\bmu)\), the Jacobian is
\((\det A)^{-2}\). The Gaussian moments of orders \(0\) and \(D\), together with the
bounded shift, give the following estimate. For \(D=0\), only the
zeroth moment is needed:
\begin{equation}
 \int_{\RR^{4L}}\dd^{4L}k_{\bb}\,\abs{N_G(k,p)}
                 e^{-\sum_iw_iq_i^2}
 \leq C_{G,\mathcal K}\frac{1+\norm{A^{-1}}^{D/2}}{(\det A)^2}.
 \label{eq:gaussian-bound}
\end{equation}
The constant term on the right is necessary. Neither lower-degree
numerator terms nor a nonzero bounded affine shift can be bounded
solely by \(\norm{A^{-1}}^{D/2}\) when this quantity tends to zero.

Choose any invertible \(L\)-row submatrix \(C_S\) and let
\(s_S>0\) be its smallest singular value. Cauchy--Binet and the
Rayleigh quotient give the explicit routing-dependent bounds
\begin{equation}
\begin{aligned}
 \det A&\geq |\det C_S|^2\prod_{i\in S}w_i,\\
 \norm{A^{-1}}&\leq\frac{s_S^{-2}}{\min_{i\in S}w_i}.
\end{aligned}
 \label{eq:chord-bounds}
\end{equation}
Indeed, the determinant contains the indicated positive minor, and
\(v^TAv\geq(\min_{i\in S}w_i)\norm{C_Sv}^2\).
In a fundamental chord basis, \(C_S\) is a signed identity and
both routing constants are one. For a general fixed routing they
are absorbed into the graph constant below.
Using
\((\min_iw_i)^{-r}\leq\sum_iw_i^{-r}\), \(r\geq0\),
the loop-integrated majorant is at most
\begin{equation}
 C_{G,\mathcal K}\prod_{i=1}^IE_i(\alpha_i)
  \prod_{i\in S}(\alpha_i+a^2\tau_i)^{-2}
 \left[1+\sum_{j\in S}
             (\alpha_j+a^2\tau_j)^{-D/2}\right].
 \label{eq:parameter-majorant}
\end{equation}
This inequality holds throughout the positive parameter orthant;
it is not restricted to simultaneous radial scaling.

Every nonchord Schwinger integral contributes
\(\mathcal I_i<\infty\). For a chord, split its integral at
\(\alpha_0\). For \(\gamma>1\) and \(\kappa>0\),
\begin{align}
 \int_0^\infty
  \frac{E_i(\alpha)}{(\alpha+\kappa)^\gamma}\dd\alpha
 &\leq\frac{\mathcal E_i}{\gamma-1}
           \kappa^{1-\gamma}
       +\alpha_0^{-\gamma}\mathcal I_i .
 \label{eq:alpha-envelope-bound}
\end{align}
The first term bounds \(0<\alpha\leq\alpha_0\); the second bounds
the remaining tail. In particular, one must not discard
\(\int_0^\infty E_i(\alpha)\dd\alpha\) for an undamped nonchord
parameter and then claim uniformity at \(m_i=0\).

Let \(h_r(\tau)=1+\tau^{-r}\). Applying
\eqref{eq:alpha-envelope-bound} with
\(\gamma=2\) or \(2+D/2\) gives the remaining majorant
\begin{equation}
 C_{G,\mathcal K,a,E}\left[
   \prod_{i\in S}h_1(\tau_i)
  +\sum_{j\in S}h_{1+D/2}(\tau_j)
          \prod_{i\in S\setminus\{j\}}h_1(\tau_i)
 \right].
 \label{eq:final-tau-majorant}
\end{equation}
Each nonchord \(\tau\)-integration contributes only \(\Mzero\).
Under (G), every factor in \eqref{eq:final-tau-majorant} is bounded
on the support of \(\rho^{\otimes I}_{\bb}\). Under (M), both
\(\int h_1\dd\rho_{\bb}\) and
\(\int h_{1+D/2}\dd\rho_{\bb}\) are finite. Independence of the line
measures factorizes these finite integrals. The positive majorant is
therefore integrable, uniformly on \(\mathcal K\); Fubini now applies to the
signed or tensor graph.
\end{proof}

\begin{remark}[Strength and limits of the estimate]
The global bound already controls all partial, nested, and overlapping
UV regions for the graph class in Theorem~\ref{thm:global-uv}.
It deliberately discards some useful nonchord decay and is generally
not a necessary convergence criterion. Its constants can diverge as
\(a\downarrow0\), as an unregulated positive mass tends to zero, or
as the graph family becomes infinite. No such uniform limit is asserted.
\end{remark}

\begin{remark}[Endpoint moments need not imply stronger moments]
\label{rem:strict-endpoint-example}
For a fixed graph degree put \(r=1+D/2\geq1\) and consider
\begin{equation}
 \rho_{r,\bb}(\dd\tau)
 =\mathbf1_{(0,\mathrm e^{-1})}(\tau)
   \frac{\tau^{r-1}}{\log^2(1/\tau)}\dd\tau.
 \label{eq:endpoint-measure-example}
\end{equation}
This is a finite positive measure without an atom at zero. The
substitution \(u=\log(1/\tau)\) gives
\begin{equation}
 \begin{aligned}
  M_r[\rho_{r,\bb}]&=\int_1^\infty\frac{\dd u}{u^2}=1,\\
  M_{r+\delta}[\rho_{r,\bb}]
     &=\int_1^\infty\frac{e^{\delta u}}{u^2}\dd u=\infty
       \quad(\delta>0).
 \end{aligned}
 \label{eq:endpoint-measure-moments}
\end{equation}
Normalizing its total mass to one changes neither finiteness statement.
The support reaches zero, so there is no positive gap. Nevertheless,
its bounded support gives, for every \(\lambda>0\),
\[
 \int e^{\lambda\tau}\rho_{r,\bb}(\dd\tau)
       \leq e^{\lambda/\mathrm e}\rho_{r,\bb}((0,\mathrm e^{-1}))<\infty.
\]
Its Laplace transform is therefore entire. This example separates
near-zero inverse-moment control from the large-depth condition used
for local holomorphy; neither should be confused with a positive gap.
Thus the endpoint hypothesis admits measures excluded by the stronger
sufficient condition \eqref{eq:conservative-moment}. It is not claimed to be
necessary for a particular graph, nor does this example provide one
law covering unbounded numerator degrees.
\end{remark}

For example, consider the scalar four-dimensional one-loop bubble
with numerator one, two positive masses, fixed real Euclidean
external momentum, and two power kernels
\(\Ph(z)=(1+z)^{-\nu}\). Its radial UV tail is proportional to
\(r^{-1-4\nu}\), so it is UV convergent for every \(\nu>0\).
The sufficient global criterion with \(D=0\) asks for \(\nu>1\).
Thus the moment condition is sufficient and need not be a sharp
convergence threshold for an individual graph.

\begin{corollary}[Damping a full-rank set of loop lines]
\label{cor:damped-basis}
Retain the routing, numerator, and infrared hypotheses of
Theorem~\ref{thm:global-uv}, and let \(S\subset E_{\rm u}(G)\) be a set of \(L\) uncut
virtual lines for which \(C_S\) is invertible. For each \(i\in S\), suppose
\(\Ph_i\) has a finite positive representing measure \(\rho_{i,\bb}\)
with either a positive gap or
\[
 \rho_{i,\bb}(\{0\})=0,\qquad M_{1+D/2}[\rho_{i,\bb}]<\infty.
\]
On every remaining line, it is enough that the multiplier be measurable
and bounded on the nonnegative real axis; the identity multiplier is allowed.
The selected multipliers may have distinct fixed scales $a_i>0$.
The graph with these line-dependent multipliers is absolutely
convergent, uniformly on \(\mathcal K\).
\end{corollary}

\begin{proof}
Bound the multipliers outside \(S\) by their finite suprema.
In the Schwinger proof use
\(w_i=\alpha_i+a_i^2\tau_i\) for \(i\in S\) and
\(w_i=\alpha_i\) otherwise. The invertible minor \(C_S\) supplies
the determinant and inverse-matrix bounds, with fixed routing
constants. Equation~\eqref{eq:final-tau-majorant} then contains
only the measures on \(S\). Its factors are integrable by the
individual gap or moment assumptions, while each remaining
Schwinger integral contributes \(\mathcal I_i<\infty\).
A direct proof with explicit one-line integral bounds is given in
\ref{sec:direct-loop-proof}.
\end{proof}

The bounded multipliers outside \(S\) are a freedom in this mathematical
comparison criterion. In applying it to a selected model, all uncut
virtual lines still carry the kernels declared by that model; one may
bound some of those kernels by constants. The allowed identity multiplier
does not turn a physical cut projector into a Euclidean covariance,
or authorize changing the placement rule. A cut removes the cut-line
covariance and its kernel together. The UV estimates concern the
remaining virtual integrations on their established kinematic domain;
they do not prove convergence of the on-shell phase-space integral or
remove infrared singularities.

An invertible selected block lets the proof estimate each loop
direction through one line. A lower-rank block leaves some directions
to the other factors. The forest criterion below describes how to
test those directions with their actual denominators and vertices.

\paragraph{Real momentum limits}
A uniform bound on the integrated absolute value does not by itself
justify dominated convergence. To obtain a common pointwise majorant
for real external limits, use the selected line momenta as integration
variables. Put
\[
 \ell=q_S,\qquad
 k=(C_S^{-1}\otimes\one_4)(\ell-R_S(p)).
\]
The Jacobian is the constant \(|\det C_S|^{-4}\). Bound each
unselected line by its multiplier supremum times \(\mathcal I_i\).
The complete transformed integrand is then bounded independently of
\(p\in\mathcal K\) by
\begin{equation}
 \mathcal M_S(\ell)=C\Bigl(1+\sum_{j\in S}|\ell_j|^D\Bigr)
       \prod_{j\in S}\Ph_j(a_j^2|\ell_j|^2)C_j(\ell_j).
 \label{eq:real-parameter-majorant}
\end{equation}
For \(D=0\), the parenthesis is simply a constant. The one-line
moments in \ref{sec:direct-loop-proof}, with the corresponding fixed
\(a_j>0\), show that \(\mathcal M_S\in L^1(\RR^{4L})\).
Thus, if the transformed numerator and unselected line factors are
continuous in the real external data for almost every \(\ell\),
the graph is continuous and its real external limits pass through
this representation integral. This continuity assumption is not
automatic for arbitrary measurable comparison multipliers. Derivatives
require their own dominating functions; complex continuation and
Lorentzian boundary limits still use Section~\ref{sec:analytic}.
No common pre-Gaussian complex majorant is inferred from a uniform
bound on the integrated absolute values.


\subsection{Optimized selected-line coverage, including massless denominators}
\label{sec:optimized-coverage}

Selected line momenta can improve the estimate in two ways.
First, each independent coordinate can carry its own numerator degree.
Second, the selected massless pole $|q|^{-2}$ is locally integrable
in four dimensions. The theorem uses that exact pole and requires
an explicit bound on the complete product of the unselected factors.

\begin{theorem}[Line-specific inverse moments on a full-rank routing]
\label{thm:optimized-selected-lines}
Fix the finite routing and compact real external set of
\eqref{eq:routing}. Choose $S$ with $|S|=L$ and invertible $C_S$, and
write $\ell_j=q_j$ for $j\in S$. On selected line $j$, take the actual
nonnegative comparison factor
\begin{equation}
 C_j^{\Ph}(\ell_j)
  :=\frac{\Ph_j(a_j^2|\ell_j|^2)}{|\ell_j|^2+m_j^2},
 \qquad a_j>0,\quad m_j\geq0,
 \label{eq:optimized-selected-factor}
\end{equation}
with a finite positive representing measure $\rho_{j,\bb}$.
Assume the \emph{complete product} of the numerator and all
unselected virtual-line factors, expressed in these variables, obeys
\begin{equation}
 \left|N_G(k(\ell,p),p)\prod_{i\notin S} C_i^{\rm full}(q_i)\right|
 \leq\sum_{\alpha\in\mathcal A}c_\alpha
       \prod_{j\in S}(1+|\ell_j|^{d_{j\alpha}}),
 \label{eq:anisotropic-complete-envelope}
\end{equation}
uniformly in $p\in\mathcal K$, with $\mathcal A$ finite,
$c_\alpha\geq0$ and $d_{j\alpha}\geq0$. Thus an unselected
massless denominator must be controlled as part of the \emph{full}
envelope; it is not assumed bounded merely because its multiplier is.
Set $d_j=\max_\alpha d_{j\alpha}$ and suppose
\begin{equation}
 \rho_{j,\bb}(\{0\})=0,\qquad
 M_{1+d_j/2}[\rho_{j,\bb}]<\infty\quad(j\in S).
 \label{eq:optimized-moment-budget}
\end{equation}
Then the specified complete graph is absolutely convergent,
uniformly over $\mathcal K$, including mixed soft and UV regions
in the selected variables. The following one-line integral does
\emph{not} include the Fourier factor $(2\pi)^{-4}$:
\begin{equation}
 \begin{aligned}
 Q_{d,j}(m_j)
  &:=\int_{\RR^4}\frac{|q|^d\Ph_j(a_j^2|q|^2)}{q^2+m_j^2}
         \dd^4q_{\bb}\\
  &\leq\pi^2\Gamma(1+d/2)a_j^{-2-d}
          M_{1+d/2}[\rho_{j,\bb}].
 \end{aligned}
 \label{eq:optimized-one-line-moment}
\end{equation}
For $m_j=0$ this is an equality, possibly infinite. For each fixed
$m_j\ge0$ the one-line integral is finite exactly when the heat
measure has no atom at zero and the indicated inverse moment is
finite. This exact one-line threshold gives a sufficient condition
for the complete graph under \eqref{eq:anisotropic-complete-envelope}.
\end{theorem}
\begin{proof}
For $q\ne0$ and $m_j\geq0$,
$(q^2+m_j^2)^{-1}\leq|q|^{-2}$. By Tonelli and four-dimensional
polar integration,
\[
 \begin{aligned}
 Q_{d,j}(m_j)&\leq2\pi^2\int\rho_{j,\bb}(\dd\tau)
   \int_0^\infty r^{d+1}e^{-a_j^2\tau r^2}\dd r\\
 &=\pi^2\Gamma(1+d/2)a_j^{-2-d}
                   M_{1+d/2}[\rho_{j,\bb}].
 \end{aligned}
\]
Equality holds for a massless selected denominator. Near $q=0$ the
radial integrand is proportional to $r^{d+1}$ at fixed positive
Laplace depth, so no additional selected-line infrared singularity
is created for $d\geq0$. An atom at $\tau=0$ would contribute an undamped constant and
make the one-line integral diverge at infinity for any fixed $m_j$.
Conversely, for $|q|\ge m_j>0$ the massive denominator supplies at
least half of $|q|^{-2}$. On $|q|<m_j$ the massless weighted
integral is finite by the boundedness of $\Ph_j$ and $d\ge0$.
Finiteness at positive mass therefore implies finiteness of the
massless integral and hence the moment condition.

The real affine change
$k=(C_S^{-1}\otimes\one_4)(\ell-R_S(p))$ has fixed Jacobian
$|\det C_S|^{-4}$. Insert \eqref{eq:anisotropic-complete-envelope} and integrate
the selected factors independently. With the graph's stated
$(2\pi)^{-4L}$ Fourier convention, the explicit bound is
\begin{equation}
 \mathcal A_G(p)\leq
 \frac{|\det C_S|^{-4}}{(2\pi)^{4L}}
 \sum_{\alpha\in\mathcal A}c_\alpha
 \prod_{j\in S}\bigl[Q_{0,j}(m_j)+Q_{d_{j\alpha},j}(m_j)\bigr].
 \label{eq:anisotropic-integrated-bound}
\end{equation}
Here $\mathcal A_G$ denotes the absolute comparison graph with
the selected factors \eqref{eq:optimized-selected-factor} and
the complete remaining product bounded in
\eqref{eq:anisotropic-complete-envelope}.
Finite total measure and finiteness of the moment of order $1+d_j/2$
imply every smaller required positive inverse moment. This is an integrable pointwise
majorant in independent selected coordinates, controlling their
partial and simultaneous ultraviolet limits. The theorem assumes
rather than proves the requisite estimate for unselected factors and
completed vertices.
\end{proof}

\begin{corollary}[An explicit selected-line numerical tail]
\label{cor:optimized-selected-tail}
Under Theorem~\ref{thm:optimized-selected-lines}, if for some
$\delta>0$ the moments
$M_{1+(d_j+\delta)/2}[\rho_{j,\bb}]$ are finite, then
\begin{equation}
 \begin{aligned}
 &\int_{|q|>R}\frac{|q|^d\Ph_j(a_j^2|q|^2)}{q^2+m_j^2}
  \dd^4q_{\bb}\\
 &\quad\leq \frac{\pi^2\Gamma(1+(d+\delta)/2)}
                       {R^\delta a_j^{2+d+\delta}}
       M_{1+(d+\delta)/2}[\rho_{j,\bb}],
 \quad 0\leq d\leq d_j.
 \end{aligned}
 \label{eq:optimized-tail-bound}
\end{equation}
Replacing one factor at a time in
\eqref{eq:anisotropic-integrated-bound} and taking the union over
$\{|\ell_j|>R\}$ gives an $O(R^{-\delta})$ tail for the complete
specified graph in the selected variables.
\end{corollary}
\begin{proof}
On $|q|>R$, $|q|^d\leq R^{-\delta}|q|^{d+\delta}$.
Apply \eqref{eq:optimized-one-line-moment} and the finite-product
bound. The union estimate follows from positivity. For each summand in
\eqref{eq:anisotropic-integrated-bound}, apply the tail estimate
separately to its constant and $d_{j\alpha}$ terms; the remaining
line integrals retain their own moment budgets. The bound concerns
\emph{selected-coordinate} truncations, not arbitrary independently
routed cutoffs on divergent constituent integrals.
\end{proof}

For illustration, a massless selected line with
$\Ph_j(z)=(1+z/b_j)^{-\nu_j}$ has the exact moment
\[
 Q_{d,j}(0)=\frac{\pi^2\Gamma(1+d/2)b_j^{1+d/2}}
    {a_j^{2+d}}
 \frac{\Gamma(\nu_j-1-d/2)}{\Gamma(\nu_j)},
 \quad \nu_j>1+d/2.
\]
This particular integral diverges at or below the stated threshold.
If two selected rows require degrees $0$ and $4$, they need
$M_1[\rho_{1,\bb}]$ and $M_3[\rho_{2,\bb}]$, respectively, rather than
$M_3$ on both rows. The improvement relies on the anisotropic
\emph{complete} envelope; decomposing a numerator after discarding
unbounded nonselected factors is not valid.

\paragraph{A two-loop certificate with massless selected lines}
\label{par:two-loop-certificate}
As a checkable example, let \(k_1,k_2\in\RR^4\), take the routing
\(q_1=k_1\), \(q_2=k_2\), \(q_3=k_1-k_2+p\). The third line has
denominator \(q_3^2+1\) and the identity multiplier. The selected
routing \(C_{\{1,2\}}=I_2\) has full rank. In fixed reference momentum
units the comparison graph is
\begin{equation}
 \begin{aligned}
 \mathcal A_{\rm ex}(p)=\int
 \frac{\dd^4k_{1,\bb}\dd^4k_{2,\bb}}{(2\pi)^8}
 &\frac{1+|k_1|^2|k_2|^4}{k_1^2k_2^2(q_3^2+1)}\,\\[-2pt]
 &\quad\times(1+k_1^2)^{-3}(1+k_2^2)^{-4}.
 \end{aligned}
 \label{eq:two-loop-selected-example}
\end{equation}
The denominators are interpreted almost everywhere, and \(p\)
can range over any real compact set. Bounding the third line by one
and applying the exact one-line identity yields
\begin{equation}
 \begin{aligned}
 \mathcal A_{\rm ex}(p)
 &\leq\frac{Q_{0,1}(0)Q_{0,2}(0)
              +Q_{2,1}(0)Q_{4,2}(0)}{(2\pi)^8}\\
 &=\frac{1}{768\pi^4},
 \end{aligned}
 \label{eq:two-loop-selected-certificate}
\end{equation}
using \(Q_{0,1}=Q_{2,1}=\pi^2/2\) and
\(Q_{0,2}=Q_{4,2}=\pi^2/3\). To verify the kernel hypotheses
independently, the representing Gamma laws are
\[
 \rho_{\nu,\bb}(\dd\tau)
 =\frac{\tau^{\nu-1}e^{-\tau}}{\Gamma(\nu)}\dd\tau,
 \qquad
 M_r[\rho_{\nu,\bb}]=\frac{\Gamma(\nu-r)}{\Gamma(\nu)}
 \quad(0<r<\nu).
\]
They have no atom at zero, and the highest selected-line moments are
\(M_2[\rho_{3,\bb}]=1/2\) and \(M_3[\rho_{4,\bb}]=1/6\).
The four displayed $Q$ values follow directly from
\eqref{eq:optimized-one-line-moment} with unit depth: the two
selected propagators contribute their exact radial integrals,
whereas the third line contributes only the bound
$(q_3^2+1)^{-1}\leq1$.
Thus one concrete two-loop comparison is uniformly bounded despite
two massless selected propagators. This example tests the
mathematical certificate and its Fourier normalization; it is not
claimed to be a graph generated by a completed O--QED action.

\begin{proposition}[Complex and real external-momentum certificates for the two-loop example]
\label{prop:two-loop-complex-certificate}
In \eqref{eq:two-loop-selected-example}, interpret the third
factor as the analytic denominator
$[(k_1-k_2+p)^2+1]^{-1}$ for complex $p$, with
$q^2=\sum_{\mu=1}^4q_\mu^2$. Denote this complex integral by
$\mathcal B_{\rm ex}(p)$. It is holomorphic on the tube
$\mathcal T_1=\{p=x+\ii y\in\CC^4:|y|<1\}$, and
\begin{equation}
 |\mathcal B_{\rm ex}(x+\ii y)|
 \leq\frac{1}{768\pi^4(1-|y|^2)}.
 \label{eq:two-loop-complex-bound}
\end{equation}
For every real $p_0$, $0<r<1/2$, and multi-index
$\alpha\in\mathbb N_0^4$, its external derivatives satisfy
\begin{equation}
 |\partial_p^\alpha\mathcal B_{\rm ex}(p_0)|
 \leq\frac{\alpha!\,r^{-|\alpha|}}
 {768\pi^4(1-4r^2)}.
 \label{eq:two-loop-complex-derivative}
\end{equation}
For real external momenta $p,p'\in\RR^4$ one also has the
global stability estimate
\begin{equation}
 \begin{aligned}
 \norm{\nabla_p\mathcal A_{\rm ex}(p)}
 &\leq\frac{9}{8\sqrt3}\,\frac1{768\pi^4},\\
 |\mathcal A_{\rm ex}(p)-\mathcal A_{\rm ex}(p')|
 &\leq\frac{9}{6144\sqrt3\,\pi^4}\norm{p-p'}.
 \end{aligned}
 \label{eq:two-loop-real-Lipschitz}
\end{equation}
For real $p$, $\mathcal B_{\rm ex}(p)=\mathcal A_{\rm ex}(p)$.
The gradient is also globally Lipschitz. For $p,p',h\in\RR^4$,
\begin{equation}
 \begin{aligned}
 &\norm{\nabla\mathcal A_{\rm ex}(p)
             -\nabla\mathcal A_{\rm ex}(p')}\leq
             \frac{\norm{p-p'}}{384\pi^4},\\
 &|\mathcal A_{\rm ex}(p+h)-\mathcal A_{\rm ex}(p)
        -\nabla\mathcal A_{\rm ex}(p)\cdot h|\leq
             \frac{\norm{h}^2}{768\pi^4}.
 \end{aligned}
 \label{eq:two-loop-real-second-order}
\end{equation}
These are coefficient-level real and complex stability bounds, not
assertions about Lorentzian cuts or scattering amplitudes. The tube
can be checked explicitly here because all external dependence
resides in the massive third denominator under the selected routing.
\end{proposition}
\begin{proof}
Write $p=x+\ii y$ and $u=k_1-k_2+x\in\RR^4$. Then
\[
 \Re[(u+\ii y)^2+1]=|u|^2+1-|y|^2
 \geq1-|y|^2>0.
\]
The absolute integrand is therefore bounded by
$(1-|y|^2)^{-1}$ times the same product of the two selected
massless-line factors and its declared numerator. The already
computed one-line integrals yield
\eqref{eq:two-loop-complex-bound} and give a common integrable
majorant on every compact subset of $\mathcal T_1$.
Theorem~\ref{thm:holomorphy} proves joint holomorphy. On a
polydisc of coordinate radius $r$ about real $p_0$ one has
$|\Im p|^2\leq4r^2$; the multivariable Cauchy estimate
then proves \eqref{eq:two-loop-complex-derivative}.
For real $u$, the denominator has the explicit derivative bound
\[
 \norm{\nabla_u(1+|u|^2)^{-1}}
 =\frac{2|u|}{(1+|u|^2)^2}
 \leq\frac9{8\sqrt3},
\]
with equality for $|u|=1/\sqrt3$. The remaining, nonnegative
selected-line weight is independent of $p$ and has integral bounded
by \eqref{eq:two-loop-selected-certificate}. Dominated
differentiation proves the gradient estimate in
\eqref{eq:two-loop-real-Lipschitz}; integration along the real
line segment from $p$ to $p'$ proves its Lipschitz bound.
For $g(u)=(1+|u|^2)^{-1}$, moreover,
\[
 \nabla^2g(u)=-\frac{2I}{(1+|u|^2)^2}
             +\frac{8uu^T}{(1+|u|^2)^3}.
\]
Its tangential and radial eigenvalues are respectively
$-2/(1+t)^2$ and $(6t-2)/(1+t)^3$, with $t=|u|^2$.
The tangential eigenvalue has absolute value at most $2$; the radial
eigenvalue ranges from $-2$ at $t=0$ to a positive maximum $1/2$
at $t=1$. Thus $\norm{\nabla^2g(u)}_{\rm op}\leq2$.
The same integrable selected-line majorant controls the integrated
Hessian, giving
$\norm{\nabla^2\mathcal A_{\rm ex}(p)}_{\rm op}
\leq1/(384\pi^4)$. The fundamental theorem of calculus along
real segments and its integral Taylor remainder prove
\eqref{eq:two-loop-real-second-order}. For each fixed pair of loop momenta, $g(k_1-k_2+p)$ and its
first two derivatives tend to zero as $|p|\to\infty$.
Their globally bounded derivatives are dominated by the same
integrable selected-line weight; dominated convergence therefore
implies that $\mathcal A_{\rm ex}(p)$ and its first two
external derivatives vanish in that limit. This argument uses the
special external-momentum placement in this example and is not a
general consequence of graphwise UV finiteness.
\end{proof}

\subsection{Correlated heat measures and shared geometric depths}
\label{sec:update-joint-depth}

A shared depth changes how a graph is averaged. At a fixed depth
$r$ the transferred Gaussian example has a product of line factors
with independent heat variables. Averaging over $r$ can make small
heat variables occur together. For instance, two selected line
integrals can both increase as $r^{-2}$, so their product costs
$r^{-4}$ before averaging. The joint measure captures this cost.
If the source couples the histories even after $r$ is fixed, use
its actual conditional law. This is the input from the mixed-source
construction of Section~\ref{sec:update-geometric-source} and
Refs.~\cite{ChiapponeGHOUQFT2026,ChiapponePartII2026}.

\begin{theorem}[Selected-line convergence for a joint heat measure]
\label{thm:update-joint-heat}
Fix a finite routing with $L\ge1$ and a set $S$ of $L$ independent
rows, as in Theorem~\ref{thm:optimized-selected-lines}. Let
$\boldsymbol\varrho_{G,\bb}(\dd\boldsymbol\lambda)$ be a finite positive
measure on $[0,\infty)^I$, not necessarily a product. The parameters
$\lambda_i$ have inverse-momentum-squared units and include the
line-depth scale. Assume $\lambda_j>0$ almost everywhere for $j\in S$.
Conditional on these parameters, write the graph with its selected
factors
\[
 \frac{e^{-\lambda_j|\ell_j|^2}}{\ell_j^2+m_j^2},
 \qquad \ell_j=q_j,\quad m_j\ge0.
\]
Let $\mathcal R_G(k,p;\boldsymbol\lambda)$ denote the entire remaining
coefficient, including unselected lines, vertices, sources, and contacts.
The joint absolute comparison integral is
\begin{equation}
\begin{aligned}
 \mathcal A_G^{\rm joint}(p)
   &:=\int\boldsymbol\varrho_{G,\bb}(\dd\boldsymbol\lambda)
       \int\frac{\dd^{4L}k_{\bb}}{(2\pi)^{4L}}\,
          |\mathcal R_G(k,p;\boldsymbol\lambda)|\\
   &\hspace{10mm}\times\prod_{j\in S}
       \frac{e^{-\lambda_j|q_j|^2}}{q_j^2+m_j^2}.
\end{aligned}
\label{eq:update-joint-graph-integral}
\end{equation}
Assume $\mathcal R_G$ obeys the complete envelope
\eqref{eq:anisotropic-complete-envelope} uniformly in the joint
parameters and external data. Unselected heat factors may be discarded
only through $e^{-\lambda_i q_i^2}\le1$; unselected massless
singularities still need that envelope. Define
\begin{equation}
\begin{aligned}
 b_d(\lambda)&:=\pi^2\Gamma(1+d/2)\lambda^{-1-d/2},\\
 \mathcal M_{G,S}
  &:=\sum_{\alpha\in\mathcal A}c_\alpha
    \int\boldsymbol\varrho_{G,\bb}(\dd\boldsymbol\lambda)
       \prod_{j\in S}
       [b_0(\lambda_j)+b_{d_{j\alpha}}(\lambda_j)].
\end{aligned}
\label{eq:update-joint-moment-budget}
\end{equation}
If $\mathcal M_{G,S}<\infty$, the joint absolute graph integral is
uniformly finite and is bounded by
\begin{equation}
 \sup_{p\in\mathcal K}\mathcal A_G^{\rm joint}(p)
 \le\frac{|\det C_S|^{-4}}{(2\pi)^{4L}}\mathcal M_{G,S}.
\label{eq:update-joint-graph-bound}
\end{equation}
This justifies the exchanges of loop and joint heat/depth
integrations for that graph.
\end{theorem}
\begin{proof}
Apply the same fixed affine change from loop variables to the selected
$\ell_j$. For each positive $\lambda_j$, polar integration gives
\[
 \int_{\RR^4}\frac{|\ell|^d e^{-\lambda_j\ell^2}}
                      {\ell^2+m_j^2}\dd^4\ell_{\bb}
       \le b_d(\lambda_j),
\]
with equality for $m_j=0$. Insert the complete envelope, perform the
selected integrations, and then integrate the resulting product
against the \emph{joint} measure. Tonelli applies to this nonnegative
majorant without an independence assumption. Its finiteness gives
Fubini for the original signed/tensor coefficient.
\end{proof}
Stronger marginal moments can certify this joint bound even when
the heat variables are dependent. For a term $\prod_{j\in S}
\lambda_j^{-u_j}$ in the expanded budget choose finite
$p_j\ge1$ with $\sum_{j\in S}1/p_j=1$. If $\varrho_j$ is the
$j$th marginal of the same finite joint measure, H\"older's inequality gives
\begin{equation}
 \int \prod_{j\in S}\lambda_j^{-u_j}\,
                \dd\boldsymbol\varrho_{G,\bb}
 \le \prod_{j\in S}\left(\int
         \lambda^{-p_ju_j}\,\dd\varrho_j(\lambda)\right)^{1/p_j}.
 \label{eq:uvpeB-joint-holder}
\end{equation}
All parameters here are on their stated measurable spaces.
Consequently finite inverse moments of every positive order for
each selected marginal control each such fixed finite product.
Knowing only its individual factors at the original orders $u_j$
generally does not suffice. The argument assumes the stipulated
complete uniform envelope for all remaining vertex and measure factors.

More generally $c_\alpha$ can depend on common geometric data, provided
that the corresponding weighted right side is finite. In that case retain
those data in the joint integration space, for example $(r,\boldsymbol\lambda)$,
rather than integrating $r$ out of the measure while still using it in
the vertices. Bounds for a
fixed background or depth cannot be averaged when their constants
are not integrable. The same qualification applies to a common
complex majorant in Theorem~\ref{thm:holomorphy}, to insertion
bounds, and to the simultaneous forest envelopes.

For the canonical two-ended transfer, $a=2r$ and $M=0$, the
conditional dimensionful heat law on an edge is
\begin{equation}
 \rho^{\rm heat}_{r,\bb}(\dd\lambda)
 =\frac{r}{\sqrt\pi}\lambda^{-3/2}
           e^{-r^2/\lambda}\dd\lambda_{\bb},\qquad\lambda>0,
\label{eq:update-shared-half-law}
\end{equation}
whose transform is $e^{-2r\sqrt x}$. The graph law is
\begin{equation}
 \boldsymbol\varrho_{G,\bb}(\dd\boldsymbol\lambda)
   =\int\mathsf w_{\bb}(\dd r_{\bb})
          \bigotimes_{e\in E_{\rm u}(G)}
                    \rho^{\rm heat}_{r,\bb}(\dd\lambda_e).
\label{eq:update-joint-lapse-law}
\end{equation}
The tensor product is conditional on $r$. It is not the tensor
product of the depth-averaged marginals. The subscript on each heat
differential records its copy; its edge index records its independent
conditional integration slot.

\begin{proposition}[The exact joint inverse-lapse budget]
\label{prop:update-joint-inverse-lapse}
For positive exponents $u_j$ on any finite selected set,
\begin{equation}
 \int\prod_{j\in S}\lambda_j^{-u_j}
           \boldsymbol\varrho_{G,\bb}(\dd\boldsymbol\lambda)
  =\prod_{j\in S}\frac{\Gamma(u_j+\tfrac12)}{\sqrt\pi}
       \int\mathsf w_{\bb}(\dd r_{\bb})\,r^{-2\sum_{j\in S}u_j}.
\label{eq:update-joint-inverse-lapse}
\end{equation}
Both sides are understood as nonnegative extended integrals for a
probability law supported on strictly positive depths. Under a uniform
complete envelope, a sufficient shared-depth condition for
Theorem~\ref{thm:update-joint-heat} is
\begin{equation}
 \begin{gathered}
 \mathsf w_{\bb}(\{0\})=0,\qquad
 \int_{(0,r_{\rm ref}]}r^{-R_G}\mathsf w_{\bb}(\dd r_{\bb})<\infty,\\
 R_G=2L+\max_{\alpha\in\mathcal A}\sum_{j\in S}d_{j\alpha}.
 \end{gathered}
\label{eq:update-depth-budget}
\end{equation}
with arbitrary fixed $r_{\rm ref}>0$. This reference only separates
the small-depth region; it is not a new physical scale.
\end{proposition}
\begin{proof}
For fixed $r$, substitute $v=r^2/\lambda$ in
\eqref{eq:update-shared-half-law}. The moment is
$\Gamma(u+\tfrac12)/(\sqrt\pi r^{2u})$. Condition on $r$, multiply
these independent conditional moments, and use Tonelli.
Expanding the finite product in \eqref{eq:update-joint-moment-budget}
requires exponents at most $R_G$ in $r^{-1}$. At $r>r_{\rm ref}$
all these terms are bounded, so a probability law controls the tail.
\end{proof}
The atom condition cannot be dropped: the displayed truncated moment
excludes $r=0$, whereas a zero-depth atom contributes an undamped
conditional line. Equivalently one may use $0^{-R_G}=+\infty$ and
integrate the extended power over $[0,r_{\rm ref}]$ for $R_G>0$.
This is a sufficient graph budget, not a sharp necessary threshold
for arbitrary graphs. It is stronger than a collection of marginal
conditions. For example, take
\[
 \mathsf w_{\bb}(\dd r_{\bb})
    =3r^2R_0^{-3}\mathbf1_{(0,R_0)}(r)\dd r_{\bb},\qquad R_0>0.
\]
Each marginal first inverse-lapse moment is finite because
$\mathbb E r^{-2}=3/R_0^2$, but
$\mathbb E(\lambda_1^{-1}\lambda_2^{-1})=\infty$.
Indeed the two independent massless selected-loop integrations
are proportional to $\mathbb E r^{-4}$. This supplies a two-loop
comparison counterexample to replacing a shared-depth law by its marginals;
it is not asserted to be a graph of a completed gauge action.

\paragraph{Shifted transfer does not require an artificial large-depth cutoff}
For fixed $M>0$ the conditional law for $e^{-2r h_M(x)}$ is
\begin{equation}
 \rho^{\rm heat}_{r,M,\bb}(\dd\lambda)
 =\frac{r e^{2rM}}{\sqrt\pi}\lambda^{-3/2}
       e^{-r^2/\lambda-M^2\lambda}\dd\lambda_{\bb}.
\label{eq:update-shifted-conditional-law}
\end{equation}
Its exact inverse moment is
\begin{equation}
 m_u(r,M)=\frac{2r e^{2rM}}{\sqrt\pi}
       \left(\frac Mr\right)^{u+1/2}K_{u+1/2}(2rM).
\label{eq:update-shifted-conditional-moment}
\end{equation}
The small- and large-argument Bessel bounds give, for fixed $u>0$,
\begin{equation}
 m_u(r,M)\le C_u\bigl[r^{-2u}+(M/r)^u\bigr].
\label{eq:update-shifted-moment-bound}
\end{equation}
For example, $e^zK_{u+1/2}(z)$ is bounded by a constant times
$z^{-u-1/2}+z^{-1/2}$ on $z>0$, proving the inequality by rescaling.
This is also consistent with the exact $M=0$ limit above.
For a finite family of fixed $M_e\ge0$, joint conditional moments
are therefore bounded by finitely many depth moments with maximal
small-depth order $2\sum u_e$. No condition such as
$A>2\sum_e M_e$ is needed for the selector below; discarding
$e^{-M_e^2\lambda_e}$ too early would produce that unnecessarily
strong exponential-moment requirement.

\subsection{Selector measure, exact ultraviolet tail, and concentration}
\label{sec:update-selector}

A preferred classical depth and a probability law for that depth answer
different questions. The law includes the integration measure as well as
the selector potential. On the fixed copy consider the monograph's declared family
\begin{equation}
 \mathsf w_{\vartheta,\bb}(\dd r_{\bb})
  =\frac{r^{\vartheta-1}e^{-A_*r-B_*/r}}{Z_\vartheta}
          \dd r_{\bb},\qquad A_*,B_*>0,
\label{eq:update-selector-law}
\end{equation}
with $[A_*]={\rm length}^{-1}$, $[B_*]={\rm length}$ in
$\hbar=1$ units. The exponent $\vartheta$ denotes the monograph's
collective-measure exponent; it is not the reflection $\theta$, a
sector sign, or a contour branch. The homogeneous nonzero-winding
selector uses $\vartheta=1/2$ with its stated coefficients. A normal
ray measure, an isotropic radial measure, and this collective measure
must not be substituted for one another without rederiving the action
and its Jacobian.

\begin{proposition}[Depth moments and the exact marginal multiplier]
\label{prop:update-selector-moments}
For \eqref{eq:update-selector-law}, every real depth moment exists, and
\begin{equation}
\begin{aligned}
 Z_\vartheta&=2(B_*/A_*)^{\vartheta/2}
                  K_\vartheta(2\sqrt{A_*B_*}),\\
 \mathbb E r^u&=(B_*/A_*)^{u/2}
       \frac{K_{\vartheta+u}(2\sqrt{A_*B_*})}
            {K_\vartheta(2\sqrt{A_*B_*})},\\
 \mathcal F_\vartheta(\omega)&:=\mathbb E e^{-\omega r}\\
  &=\left(\frac{A_*}{A_*+\omega}\right)^{\vartheta/2}
       \frac{K_\vartheta(2\sqrt{B_*(A_*+\omega)})}
            {K_\vartheta(2\sqrt{A_*B_*})},\quad\omega\ge0.
\end{aligned}
\label{eq:update-selector-transform}
\end{equation}
If $B_*=0$ instead, normalization requires $\vartheta>0$, and
$\mathbb E r^{-u}<\infty$ holds exactly for $u<\vartheta$.
\end{proposition}
\begin{proof}
The substitution $r=\sqrt{B_*/A_*}\,t$ in the modified-Bessel
integral \cite{DLMF10} gives the normalization. Replacing
$\vartheta$ by $\vartheta+u$ gives the moments; replacing $A_*$ by
$A_*+\omega$ gives the transform. The exponentials at both endpoints
control all powers. For $B_*=0$ the law is gamma-normalized and
its small-$r$ integral is $r^{\vartheta-1-u}\dd r$, giving the threshold.
\end{proof}
Thus $B_*>0$ supplies every \emph{finite-order} joint moment budget
in \eqref{eq:update-depth-budget}, provided the complete vertex/IR
bounds are uniform or have integrable constants. It does not prove
those bounds for a gauge or gravitational action or convergence of
an infinite perturbation series. At $B_*=0$ this sufficient budget
requires $\vartheta>R_G$. In particular the earlier
$\vartheta=1/2$ zero-winding law has a divergent massless coincidence
average, since
\[
 C_{2r}(0)=\frac1{32\pi^2r^2},\qquad
 \mathbb E_{\mathsf w_{1/2}}C_{2r}(0)=\infty.
\]
Almost-sure positivity of the depth alone does not control its
probability near zero.

\begin{proposition}[The marginal ultraviolet rate and fixed-depth limit]
\label{prop:update-selector-tail}
For fixed positive selector coefficients and real $\vartheta$, put
\[
 C_\vartheta=\frac{\sqrt\pi B_*^{\vartheta/2-1/4}}{Z_\vartheta}.
\]
Then
\begin{equation}
 \mathcal F_\vartheta(\omega)
 \sim C_\vartheta(A_*+\omega)^{-\vartheta/2-1/4}
       e^{-2\sqrt{B_*(A_*+\omega)}}.
\label{eq:update-depth-transform-tail}
\end{equation}
At fixed $M\ge0$ the marginal internal multiplier consequently obeys
\begin{equation}
\begin{aligned}
 \overline\Phi_\vartheta(x)
   &:=\mathcal F_\vartheta(2h_M(x))\\
   &\sim C_\vartheta2^{-\vartheta/2-1/4}
       x^{-\vartheta/4-1/8}e^{-2\sqrt{2B_*}\,x^{1/4}}.
\end{aligned}
\label{eq:update-selector-uv-tail}
\end{equation}
An exact equality $\int\mathsf w_{\bb}(\dd r_{\bb})e^{-2rh_M(x)}
=e^{-a h_M(x)}$ for every $x\ge0$ instead forces
$\mathsf w_{\bb}=\delta^{\rm depth}_{\bb,a/2}$.
\end{proposition}
\begin{proof}
Insert $K_\vartheta(z)\sim\sqrt{\pi/(2z)}e^{-z}$
\cite{DLMF10} in \eqref{eq:update-selector-transform}. For the second
equivalent use $A_*+2h_M(x)=2\sqrt{x}+A_*-2M+O(x^{-1/2})$ at
fixed $M$. Its square root differs from $\sqrt{2}\,x^{1/4}$ by
$O(x^{-1/4})$, so the ratio of the exponential factors tends to one
along with the ratio of the stated power prefactors.
For the last assertion, $h_M$ maps $[0,\infty)$ onto itself. Uniqueness of the Laplace transform of finite positive measures
\cite{Widder1941} identifies the depth measure with the indicated atom.
That atom is not a coordinate GTL delta.
\end{proof}
For one common depth, a complete product of transferred internal
lines is averaged as
\begin{equation}
 \mathbb E_r\prod_{e=1}^I e^{-2r h_{M_e}(q_e^2)}
  =\mathcal F_\vartheta\!\left(2\sum_{e=1}^I h_{M_e}(q_e^2)\right),
\label{eq:update-joint-line-transform}
\end{equation}
not as $\prod_e\overline\Phi_\vartheta(q_e^2)$. External half-transfers
and induced vertices must remain in the same conditional expression.
This distinction persists even though every finite inverse moment exists.

If $A_*=\Lambda A_0$, $B_*=\Lambda B_0$, put $V(r)=A_0r+B_0/r$.
It tends to infinity at both endpoints and has the unique minimum
$r_0=\sqrt{B_0/A_0}$. Outside any neighborhood of $r_0$,
$V\geq V(r_0)+2\epsilon$ for some $\epsilon>0$. There the
unnormalized mass is bounded by
$e^{-(\Lambda-1)(V(r_0)+2\epsilon)}\int_0^\infty r^{\vartheta-1}e^{-V(r)}\dd r$
for $\Lambda\geq1$. An interval where $V\leq V(r_0)+\epsilon$
gives a lower normalization bound proportional to
$e^{-\Lambda(V(r_0)+\epsilon)}$. Hence the exterior probability
tends to zero. Applying this concentration to the bounded continuous
function $e^{-r\omega}$ gives the transform limit $e^{-r_0\omega}$
at each fixed $\omega\geq0$. At every finite $\Lambda$, however,
\eqref{eq:update-depth-transform-tail} gives
$-\log\mathcal F_\vartheta(\omega)/\omega\to0$ as $\omega\to\infty$,
whereas the point-mass limit gives $r_0$. A fixed-spectral-value
concentration argument is not uniform control of the exact ultraviolet
exponential rate. No interchange of this limit with a graph integral
is made without the required common majorant.


\subsection{Residual damping after nonpolynomial vertex growth}
\label{sec:residual-damping}

The polynomial bound in \eqref{eq:numerator-bound} is sufficient,
not necessary. The following test compares the \emph{complete}
vertex product with the damping in the same graph. It does not
construct the vertices required by a Ward identity.

\begin{proposition}[Residual-coercivity criterion]
\label{prop:residual-coercivity}
Retain the fixed routing and compact external set in
\eqref{eq:routing}, and assume \(0\leq C_i(q)\leq\mathcal I_i<\infty\).
Either of the following conditions gives uniform absolute convergence
of the graph.

\textit{Quadratic case.} Suppose
\(0\leq\Ph_i(a_i^2q_i^2)\leq B_i^{\rm env}e^{-d_i|q_i|^2}\), with
\(d_i>0\), and the complete numerator satisfies
\begin{equation}
 |N_G(k,p)|\leq C(1+\norm{k}^D)
            \exp\!\left(\sum_i v_i|q_i|^2\right),
 \label{eq:vertex-quadratic-budget}
\end{equation}
where \(v_i\geq0\) are fixed and \(D\geq0\). Put
\begin{equation}
 H_{\rm res}:=C^T\diag(d_i-v_i)C.
 \label{eq:residual-Gram-matrix}
\end{equation}
It is sufficient that \(H_{\rm res}\succ0\); individual
\(d_i-v_i\) need not all be positive. A Laplace gap supplies
\(d_i=a_i^2\tau_{i,*}\).

\textit{Linear case.} Suppose instead
\(0\leq\Ph_i(a_i^2q_i^2)\leq B_i^{\rm env}e^{-a_i|q_i|}\) and replace
\(v_i|q_i|^2\) in \eqref{eq:vertex-quadratic-budget} by
\(v_i|q_i|\). Put \(h_i=a_i-v_i\). The residual coefficients may
have either sign, provided
\begin{equation}
 \gamma_h:=\min_{\substack{\omega\in(\RR^4)^L\\\norm{\omega}=1}}
     \sum_i h_i\left|\sum_{\ell=1}^L C_{i\ell}\omega_\ell\right|>0.
 \label{eq:signed-linear-coercivity}
\end{equation}
In particular, nonnegative \(h_i\) and full loop rank of the rows
with \(h_i>0\) suffice. Here \(B_i^{\rm env}>0\) are finite
upper-envelope constants, not necessarily the masses \(\Ph_i(0)\)
of the representing measures. The growth coefficients have
inverse-momentum-squared units in the quadratic case and
inverse-momentum units in the linear case.
All constants and the positive coercivity margin are uniform on the
external set. The asserted conclusion is Euclidean and does not
supply a Lorentzian continuation.
\end{proposition}

\begin{proof}
The bounded covariances and finite \(B_i^{\rm env}\) can be absorbed
into the constant. In the quadratic case, the remaining exponent is
\(-\sum_i(d_i-v_i)|(Ck)_i+R_i(p)|^2\). Its quadratic part is
\(-k^T(H_{\rm res}\otimes\one_4)k\); its linear and constant
coefficients are bounded on \(\mathcal K\). If
\(\lambda=\lambda_{\min}(H_{\rm res})>0\), Young's inequality gives
\begin{equation}
 \sum_i(d_i-v_i)|q_i|^2
       \geq\tfrac12\lambda\norm{k}^2-C_{\mathcal K}.
 \label{eq:residual-coercive-bound}
\end{equation}
An integrable polynomial times a Gaussian therefore bounds the
complete graph.

For either sign of \(h_i\), the triangle inequalities imply
\[
 h_i|(Ck)_i+R_i|\geq h_i|(Ck)_i|-|h_i|\,|R_i|.
\]
Homogeneity and \eqref{eq:signed-linear-coercivity} therefore give
\begin{equation}
 \sum_i h_i|q_i|\geq\gamma_h\norm{k}
       -\sum_i|h_i|\sup_{p\in\mathcal K}|R_i(p)|.
 \label{eq:signed-linear-affine-bound}
\end{equation}
The remaining polynomial times \(e^{-\gamma_h\norm{k}}\) is
integrable. For a directly computable sufficient test, choose an
invertible \(C_S\) with \(h_*:=\min_{i\in S}h_i>0\). Then
\begin{equation}
 \gamma_h\geq h_*s_{\min}(C_S)
                  -\sum_{i:h_i<0}|h_i|\,\norm{C_i}_2,
 \label{eq:signed-linear-basis-certificate}
\end{equation}
where \(C_i\) is a routing row. The positive rows bound the selected
block norm, and each negative row has operator norm \(\norm{C_i}_2\)
on \((\RR^4)^L\). If all \(h_i\geq0\), the negative sum vanishes
and the previous full-rank sufficient test follows. Both cases give
common pointwise bounds, uniform on \(\mathcal K\).
For example, one loop with two equal routing rows and residuals
\((h_1,h_2)=(2,-1)\), in fixed inverse-momentum units, has
\(\gamma_h=1\). Thus a negative linewise residual need not destroy
convergence of the complete graph. The actual O--QED vertex product
must still establish the stated envelope; this example supplies no
new vertex or contact completion.
\end{proof}


\paragraph{An exact finite test for the linear damping margin}
The loop momenta in \eqref{eq:signed-linear-coercivity} are
four-vectors. Nevertheless, collinear configurations suffice to
decide strict positivity, even when some residuals $h_i$ are negative.
Let $g(x)=\sum_i h_i|C_i x|$ on the scalar routing space $\RR^L$.
For $L>1$, form the finite set $\mathcal U_C$ of unit vectors along
the null lines of all rank-$(L-1)$ subsets of routing rows, taking
both signs. For $L=1$ take $\mathcal U_C=\{-1,1\}$. Put
\[
 \lambda_C=\min_{u\in\mathcal U_C} g(u).
\]
Then $\gamma_h>0$ exactly when $\lambda_C>0$, and in that case
$\gamma_h=\lambda_C$.

To prove the test, divide $\RR^L$ into the closed cones on which
every nonzero row $C_i x$ has fixed sign. Full column rank makes
each cone pointed. Its extreme rays lie in intersections of $L-1$
independent row hyperplanes. Each $x$ in the cone is a sum of its
unit ray generators with nonnegative coefficients. On that cone
$g$ is linear, so $\lambda_C>0$ and the triangle inequality imply
$g(x)\ge\lambda_C|x|$. To lift this scalar inequality, let
$T:\RR^4\to\RR^L$ have rows $k_\ell^{\mathsf T}$, and average
over a uniform unit vector $\theta\in S^3$. Write
$c_4=\mathbb E|\theta_1|>0$. Rotational invariance gives
\[
 c_4\sum_i h_i|(Ck)_i|=\mathbb E\,g(T\theta)
       \ge\lambda_C\mathbb E|T\theta|.
\]
For $T\ne0$, let $w_j$ be its squared singular values divided
by $\|T\|_{\rm F}^2$, extended by zeros to four entries.
Concavity of the square root gives
$\mathbb E(\sum_j w_j\theta_j^2)^{1/2}
 \ge\sum_j w_j\mathbb E|\theta_j|=c_4$.
Thus $\gamma_h\ge\lambda_C$ for a positive ray minimum. Conversely,
the collinear tuple $k_\ell=u_\ell e$ for a unit four-vector $e$
attains $g(u)$ when $u$ minimizes the ray values. It also witnesses
failure of strict coercivity if $\lambda_C\le0$.
For example, with $L=2$, rows $(1,0),(0,1),(1,1)$ and residuals
$(1,1,-3/4)$, the ray test gives $\gamma_h=1/4$.
The positive-basis lower bound in \eqref{eq:signed-linear-basis-certificate}
is $1-3\sqrt2/4<0$ for these same data.

A vanishing coercivity margin does not imply divergence: denominator
power counting or sector estimates may still suffice. Conversely,
rapid line decay alone does not establish a positive residual margin.
For example, cancelling both endpoint factors in
\eqref{eq:internal-damping-cancellation} leaves no such gain.
The criterion allows some growing vertices without assuming that the
Ward-compatible completion actually meets its hypotheses.

\subsection{A forest criterion for weaker linewise hypotheses}
\label{sec:forests}

When the available decay does not meet
Corollary~\ref{cor:damped-basis}, retaining the individual UV regions
can yield weaker sufficient conditions. Failure of that full-rank
criterion does not imply divergence. The relevant principle is the
Euclidean power-counting principle that convergence requires control
of every hard loop subspace, including directions in which some routed
momenta stay bounded \cite{Weinberg1960}. For a history-dependent source
the needed estimate is on its complete absolute envelope after sewing.
Separate scaling checks supply such an estimate only if their constants
remain controlled in the other scales.

\begin{definition}[Admissible sector envelope]
\label{def:sector-envelope}
An admissible finite sector decomposition for \(G\) covers the
ultraviolet region of loop-momentum space up to a null set, with
finite overlap; the remaining bounded region is covered by the
stated infrared estimate. Each sector parametrization includes its
Jacobian, and the transformed integrand has the stated bound.
The decomposition consists of finitely many sectors
\(\mathfrak h\) with scale variables
\(R_j\geq1\), \(j=1,\ldots,r_{\mathfrak h}\), and an absolute
ordinary envelope, including the measure, of the form
\begin{equation}
 C_{G,\mathcal K,\mathfrak h}
 \prod_{j=1}^{r_{\mathfrak h}}
     R_j^{\omega_0(\mathfrak h,j)-1}\dd R_j\,
 \dd\nu_{\mathfrak h,\bb},
 \qquad \int\dd\nu_{\mathfrak h,\bb}<\infty .
 \label{eq:sector-envelope}
\end{equation}
Angular, lower-scale, and infrared variables are controlled in
\(\dd\nu_{\mathfrak h,\bb}\). The degrees \(\omega_0\) include every
completed-action vertex and insertion. These are simultaneous
bounds on the whole sector, not pointwise bounds with constants
depending uncontrollably on other \(R_j\).
\end{definition}

Let a covered line \(e\) have a finite positive representing
measure \(\rho_{e,\bb}\), with \(\rho_{e,\bb}(\{0\})=0\), and suppose
\(M_{\beta_e,e}:=M_{\beta_e}[\rho_{e,\bb}]<\infty\) for an exponent
\(\beta_e>0\).
For \(z>0\),
\begin{equation}
 \Ph_e(z)\leq
 \left(\frac{\beta_e}{\mathrm e}\right)^{\beta_e}
 M_{\beta_e,e}\,z^{-\beta_e},
 \label{eq:moment-power-bound}
\end{equation}
because \(u^{\beta_e}e^{-u}\leq
(\beta_e/\mathrm e)^{\beta_e}\).
A sufficient simultaneous routing bound in sector \(\mathfrak h\)
is
\begin{equation}
 \abs{q_e(k,p)}\geq c_{e,\mathfrak h}
       \prod_jR_j^{\chi_{ej}},\qquad
 c_{e,\mathfrak h}>0,\quad \chi_{ej}\geq0.
 \label{eq:simultaneous-routing}
\end{equation}
It gives total decay
\begin{equation}
 \delta_\Ph(\mathfrak h,j)=2\sum_e\beta_e\chi_{ej},
 \qquad
 \omega_\Ph(\mathfrak h,j)=
       \omega_0(\mathfrak h,j)-\delta_\Ph(\mathfrak h,j).
 \label{eq:forest-degree}
\end{equation}
Only covered uncut internal lines appear in this sum. The simultaneous
bound or an equivalent explicit sector factorization prevents the
same gain from being counted twice for independent scales.

With a positive heat gap $\tau_{e,*}$, line $e$ supplies the
Gaussian width $a_e^2\tau_{e,*}$. If $J$ is the set of sector scales
to be controlled by gaps, a sufficient estimate in the actual sector
coordinates is
\begin{equation}
 \begin{aligned}
 \sum_{e\in E_{\rm ctrl}}a_e^2\tau_{e,*}|q_e(k,p)|^2
 &\geq \sum_{j\in J} c_j R_j^{\eta_j}-C_{\mathcal K,\mathfrak h},\\
 c_j&>0,\qquad \eta_j>0 .
 \end{aligned}
 \label{eq:gap-coercivity}
\end{equation}
The constants must be uniform in the other scales and in the
variables integrated by $\nu_{\mathfrak h,\bb}$. A quadratic bound
in scaled loop momenta gives this estimate when their norm controls
the indicated powers of $R_j$. For instance, a bound on $|k_1+k_2|$
cannot control the direction $k_1=-k_2$. The routing of the gapped
lines must detect every direction for which their decay is used.

\begin{theorem}[Conditional forest convergence]
\label{thm:forest}
Assume an admissible finite sector decomposition and the stated
infrared control. On each sector suppose the complete integrand,
including its Jacobian, has a uniform absolute envelope of the form
\[
 C\prod_j R_j^{\omega(\mathfrak h,j)-1}\dd R_j\,
    \prod_{j\in J}e^{-c_j R_j^{\eta_j}}\,
    \dd\nu_{\mathfrak h,\bb},
 \qquad c_j,\eta_j>0 .
\]
Require $\omega(\mathfrak h,j)<0$ for every $j\notin J$.
Equations~\eqref{eq:forest-degree} and~\eqref{eq:gap-coercivity}
give such a bound when their kernel estimates hold simultaneously.
Using disjoint line sets for the power and gap bounds is sufficient
for multiplying the bounds. Then the graph is uniformly absolutely
UV convergent on \(\mathcal K\).
\end{theorem}

\begin{proof}
By positivity extend the radial integrations to the full
product of $[1,\infty)$ intervals. For $j\notin J$ the integral is
$-1/\omega(\mathfrak h,j)$, which is finite. For $j\in J$ the
exponential factor is integrable against every fixed radial power,
by the change of variables $u=c_jR_j^{\eta_j}$. The remaining
measure has finite mass. Summing the finite sector bounds and
adding the assumed infrared bound proves the result.
\end{proof}

The criterion applies once a finite simultaneous sector envelope has
been established; it does not construct that envelope for an arbitrary
completed action. Proper subgraphs and overlapping regions are covered
through their sector scales. Checking only the overall degree of the
graph is insufficient. On a fixed flat graph with all line kernels gapped,
a full-rank chord set automatically supplies global coercivity, as
already captured by Theorem~\ref{thm:global-uv}.

A commonly tempting determinant shortcut is false:
\[
 \mathcal U_G(w)\not\geq c\left(\sum_iw_i\right)^L
 \quad\hbox{uniformly on the whole positive cone}.
\]
For the sunset graph,
\(\mathcal U=w_1w_2+w_1w_3+w_2w_3\), while
\((w_1,w_2,w_3)=(1,\delta,\delta)\) makes
\(\mathcal U\to0\) as \(\delta\downarrow0\).
The chord minor \eqref{eq:chord-bounds} and sector envelopes avoid
this error.

\paragraph{Joint-measure use of the preceding criteria}
The independent-line forest proof is unchanged. Applying it conditionally
on a common geometric variable requires a final integrable average of
its simultaneous absolute majorant. Theorem~\ref{thm:update-joint-heat}
provides a direct sufficient alternative when a selected routing set has
the stated joint inverse moments. Marginal moments at the unsplit orders need not suffice.
Equation~\eqref{eq:uvpeB-joint-holder} states which stronger marginal
bounds do control joint products. An averaged propagator alone does
not supply the forest, insertion or complex-continuation estimates
for the completed source.

A useful variant chooses the controlling lines after the widths have
been sampled. For fixed routing there are only finitely many linearly
independent choices of $L$ loop rows. If $B_S$ is the conditional
absolute majorant obtained with choice $S$, then their pointwise
minimum still bounds the same conditional integral and is measurable.
Only this minimum needs an integrable joint average. For example, take
a four-dimensional scalar bubble of mass $m>0$, with momentum measure
$\mathrm d^4q/(2\pi)^4$. In fixed units give the width pairs
$(1,2^{-n})$ and $(2^{-n},1)$ probability $2^{-n-1}$ each for
$n\ge1$. Both marginal inverse first moments are infinite.
The line of width one gives a translated Gaussian at every draw, and
its full line integral is at most
$\int e^{-|q|^2}|q|^{-2}\mathrm d^4q/(2\pi)^4=\pi^2/(2\pi)^4$.
The other line is bounded by $m^{-2}$. Their product therefore has an
integrable joint majorant at every fixed external momentum.

\subsection{All orders, insertions, and finite matching}

\begin{corollary}[Finite-order meaning of all orders]
\label{cor:all-orders}
Fix a declared completed action and its kernel data on the selected
support and evolution branch, a scale \(a>0\), and an infrared
prescription. Suppose every finite loop graph \(G\) with
\(L(G)\geq1\), generated by this fixed set of data, satisfies
Theorem~\ref{thm:global-uv}, Corollary~\ref{cor:damped-basis},
Theorem~\ref{thm:optimized-selected-lines}, the joint-law
Theorem~\ref{thm:update-joint-heat}, Proposition~\ref{prop:residual-coercivity},
or Theorem~\ref{thm:forest}, with the corresponding complete-numerator,
joint-measure and infrared bounds. Then every such loop graph is UV finite.
Tree coefficients contain no loop integration; they are included only
on external domains where their vertices and internal exchange factors
are defined, with any poles treated by the declared boundary prescription.
The loop estimates do not remove tree-level kinematic singularities.

For a finite collection of these graphs and a common compact external
set, their bounds can be chosen uniformly over the collection.
If residual field or support labels are summed, the summed amplitudes
additionally require finite support or an absolutely summable family
of graph majorants. This gives no uniform bound in perturbative
order and no convergence theorem for the infinite graph sum.
\end{corollary}

A gap, or inverse moments of every positive order in the fixed kernel
family, can accommodate every finite polynomial numerator degree in
the globally covered graph class. A single finite inverse moment does not automatically
cover an unbounded tower of derivative vertices. Infinite-valence
or nonlocal vertices also require control of graph multiplicity at a
fixed perturbative order.

External-momentum derivatives and composite insertions are covered
only when their \emph{complete differentiated integrands} have an
appropriate common majorant on the claimed domain. Derivatives of a
kernel can change both numerator growth and infrared behaviour;
replacing them by an assumed polynomial-degree cost without an
actual bound is insufficient. BRST and Nielsen insertions must pass
the same test.

For the covered graphs, mass, residue, and coupling conditions are
finite matching conditions on already finite amplitudes. They remain
physical choices of normalization and parameter determination.
UV finiteness does not make amplitudes independent of \(a\), of the
kernel family, or of the completed vertices. It removes the need
for UV-divergent subtraction on the stated domain.
\section{Exact BRST and Slavnov--Taylor identities of the completed action}
\label{sec:brst}

The quadratic inverse determines the free lines, while derivatives of
the completed action and its source map determine the vertices and
contact terms. The quantum BRST change of variables must use these same
objects. The Slavnov--Taylor argument therefore assumes an invariant
functional measure and a common domain for the required insertions;
the separate estimates in Section~\ref{sec:uv} establish absolute
convergence when its graph coefficients satisfy their hypotheses.

\subsection{Quantum BRST and background covariance}

Let \(G_{\rm gauge}\) be a compact gauge group with structure
constants \(f^{abc}\), and let \(\BR\) be a left odd derivation.
For homogeneous Lie-algebra-valued fields we use
\[
 [X,Y]^a=f^{abc}X^bY^c,\qquad
 (D_\mu[\mathcal A]X)^a
       =\partial_\mu X^a+g f^{abc}\mathcal A_\mu^bX^c.
\]
Grassmann parities are carried by the coefficients, so that
\([X,Y]=-(-1)^{\epsilon_X\epsilon_Y}[Y,X]\).
For \(\mathcal A=\overline{\mathcal A}+a\), take
\begin{align}
 \BR a_\mu^a&=(D_\mu[\mathcal A]c)^a,
 &\BR c^a&=-\frac g2 f^{abc}c^bc^c,
 \label{eq:brst-gauge}\\
 \BR\bar c^a&=B^a,&
 \BR B^a&=0,&
 \BR\overline{\mathcal A}_\mu^a&=0 .
 \label{eq:brst-doublet}
\end{align}
Matter fields transform in their declared gauge representations.
All quantities carry the same suppressed fixed label \(\bb\).
The field \(B\) is the Nakanishi--Lautrup auxiliary field, distinct
from the support label \(b_s\).

The graded Jacobi identity gives \(\BR^2c=0\); on the connection,
\begin{equation}
 \BR^2a_\mu
 =D_\mu(\BR c)+g[\BR a_\mu,c]
 =-\frac g2D_\mu[c,c]+g[D_\mu c,c]=0 .
 \label{eq:nilpotence}
\end{equation}
The same nilpotence follows on the auxiliary doublet from
\eqref{eq:brst-doublet} and on matter from the representation identity.
The quadratic background operator is fixed in this algebraic
calculation. Since \(\BR\overline{\mathcal A}=0\),
\([\BR,\overline D_E]=0\). On the common functional-calculus
domain this gives
\begin{equation}
 \BR\Ph(a^2\mathscr L_{\bb})=0.
 \label{eq:kernel-inert}
\end{equation}
Under a background gauge transformation the situation is different:
\begin{equation}
 \delta_\omega\overline{\mathcal A}_\mu=\overline D_\mu\omega,
 \qquad
 \delta_\omega a_\mu=g[a_\mu,\omega].
 \label{eq:background-transform}
\end{equation}
Ghosts, \(B\), and matter transform homogeneously in their respective
representations. The operator in \eqref{eq:positive-operator}
transforms by conjugation. Its functional calculus is therefore
background covariant. Thus the background transformation conjugates the kernel,
while the quantum BRST variation holds the background kernel fixed
\cite{Abbott1981,BRST1976,Tyutin1975}.

\subsection{Why a completed interaction is required}
\label{sec:completion-required}

Suppose, in a fixed momentum representation,
\[
 S_\Ph(p)=F(p^2)S_0(p),\qquad
 S_\Ph^{-1}(p)=F(p^2)^{-1}S_0^{-1}(p).
\]
A Ward-compatible photon vertex must satisfy
\begin{equation}
 q_\mu\Gamma_\Ph^\mu(p+q,p)
   =S_\Ph^{-1}(p+q)-S_\Ph^{-1}(p).
 \label{eq:dressed-WTI}
\end{equation}
The bare vertex satisfies the same equation with \(S_0^{-1}\).
For nonconstant \(F\), the two right sides are generally different.
Thus \eqref{eq:kernel-inert} cannot prove the interacting Ward
identity for a propagator-only modification.

We assume a completed gauge-fixed action
\begin{equation}
 S_{\rm gf}^{\rm comp}
   =S_{\rm inv}^{\rm comp}+\BR\Psi,\qquad
 \BR S_{\rm inv}^{\rm comp}=0,
 \label{eq:completed-action}
\end{equation}
with all vertices and source interfaces generated consistently.
The existence of such an action with the desired propagators,
physical interfaces, and sufficient UV bounds is substantive input.

The first consistency conditions can be seen in a formal pure-gauge
expansion around the trivial background. Consider
\[
 S_{\rm inv}^{\rm comp}=S_2+gS_3+g^2S_4+\cdots,\qquad
 \BR=\BR_0+g\BR_1,
\]
where \(S_n\) is homogeneous of degree \(n\) in the gauge field,
\(\BR_0a_\mu=\partial_\mu c\), and
\(\BR_1a_\mu=[a_\mu,c]\). BRST invariance requires
\begin{equation}
 \BR_0S_2=0,\qquad
 \BR_0S_3=-\BR_1S_2,\qquad
 \BR_0S_4=-\BR_1S_3.
 \label{eq:quadratic-vertex-hierarchy}
\end{equation}
For two actions with this same BRST transformation,
\(\BR_0\Delta S_3=-\BR_1\Delta S_2\).
A changed quadratic kernel can retain the old cubic vertex only if
\(\BR_1\Delta S_2=0\). This hierarchy constrains the interaction
vertices; their momentum growth must additionally satisfy the UV
estimates for that same completed action.

\subsection{A covariant action and its interaction vertices}

An inverse form factor can enter an invariant action through the
full connection. Let \(H(z)=\Ph(z)^{-1}\) wherever \(\Ph(z)>0\).
On a gauge-stable common form domain, consider the quadratic form
in the curvature
\begin{equation}
 S_{\rm inv}^{H}
 =\frac14\int\dd^4x_{\bb}\,
 \bigl\langle F_{\mu\nu},
 H(a^2[-D[\mathcal A]^2])F_{\mu\nu}\bigr\rangle_{\rm color}.
 \label{eq:covariant-H-action}
\end{equation}
Here \(\langle\ ,\,\rangle_{\rm color}\) is an invariant positive
color pairing. The form pairing is interpreted as the inner product of two
$H^{1/2}$-transformed curvatures on its form domain.
Gauge transformations act unitarily on the curvature and
transport the spectral resolution of the full covariant Laplacian.
Thus the form is invariant on its stated domain. At zero
background the linear curvature gives the transverse quadratic
operator \(p^2H(a^2p^2)\). On the invertible transverse
subspace its covariance is \(\Ph(a^2p^2)P_T/p^2\).

In the invariant action, \(H\) is evaluated on the full covariant
Laplacian and transforms by conjugation under quantum gauge
transformations. In the gauge-fixing sector below, the same scalar
function is evaluated on the background operator
\(\mathscr L_{\bb}\), which is quantum-BRST inert.

For the gauge-fixing function
\(\mathcal F=\overline D^\mu a_\mu\), consider
\begin{equation}
 \Psi_H=\int\dd^4x_{\bb}\,
 \bar c\,H(a^2\mathscr L_{\bb})
        \left(\mathcal F-\frac{\zeta}{2}B\right),
 \qquad \zeta>0.
 \label{eq:hyper-gauge-fermion}
\end{equation}
Since \(\BR H=0\),
\begin{equation}
 \BR\Psi_H
 =\int\dd^4x_{\bb}\left[
 B H\mathcal F-\frac{\zeta}{2}BHB
             -\bar c\,H\,\BR\mathcal F\right].
 \label{eq:hyper-gauge-expansion}
\end{equation}
With the convention \(\BR\bar c=B\), a literal Euclidean Gaussian
uses the usual imaginary \(B\)-contour, or its equivalent real
auxiliary variable. Eliminating \(B\) yields
\[
 S_{\rm gauge\ fixing}^{H}
       =\frac1{2\zeta}\langle\mathcal F,H\mathcal F\rangle .
\]
If \(\BR\mathcal F=-M c\), the ghost action is
\(\int\bar c\,HM c\). On the zero-mode-reduced domain assume
the ordered composition admits the indicated inverse, including
the domains of both factors. Then
\begin{equation}
 (HM)^{-1}=M^{-1}\Ph .
 \label{eq:ghost-ordering}
\end{equation}
It can be written \(\Ph M^{-1}\) when the two inverse factors commute as operators on the relevant
domain. The Gaussian in \(B\) contributes
\(\det(\zeta H)^{-1/2}\) at finite resolution. Its continuum
interpretation belongs to the declared measure or compatible
normalization, including its possible background dependence.

This is the Gaussian determinant in the chosen coefficient
representation. Its contribution to a carrier-valued logarithm is
formed from the two branch determinants, with their declared
logarithm branches, and the central idempotents. Multiplying the
scalar determinant or its logarithm by \(\ee\) is not that
reconstruction.

Expanding \(H(-D[\mathcal A]^2)\) in
\eqref{eq:covariant-H-action} generates the interaction terms. If
\(\Ph\) decreases exponentially, its inverse can grow exponentially,
so the induced vertices can compensate the propagator suppression.
The construction establishes algebraic gauge invariance on the stated
domain. A successful completion must additionally verify the UV bounds
for the full vertex expansion, construct the special external source
interface, and establish unitarity.

\subsection{Geometric transfer vertices and background response}
\label{sec:update-induced-vertices}

The inert multiplier in \eqref{eq:kernel-inert} is a function of the
\emph{fixed background} operator. If a normal transfer is instead built
from the full field-dependent $L[\mathcal A]$, its operator variation is
part of the interaction, not an optional correction to an otherwise
unchanged vertex. At fixed $M$ and $r$, write
$A=\sqrt{L[\mathcal A]+M^2}$ and $P_r=e^{-r(A-M)}$.
At finite spectral resolution these variations are matrix identities.
For a continuum operator, use a common differentiable domain on
which the integral and product formulas hold. One sufficient
setting for the norm estimates below is $A\ge a_0\one>0$ and
bounded differentiable perturbations of $L$ on a fixed operator
domain. Differentiation gives
\begin{equation}
\begin{aligned}
 A\,\delta A+\delta A\,A&=\delta L,\\
 \delta A&=\int_0^\infty e^{-tA}(\delta L)e^{-tA}\dd t,\\
 \delta P_r&=-\int_0^r
       P_{r-v}(\delta A)P_v\dd v.
\end{aligned}
\label{eq:update-transfer-variation}
\end{equation}
Here $t$ and $v$ are semigroup parameters. The first equation
comes from $A^2=L+M^2$. For the second, differentiate
$e^{-tA}(\delta L)e^{-tA}$ in $t$ between vectors in the domain
of $A$ and use the endpoint limits at zero and infinity.
The norm integrand has majorant $e^{-2ta_0}\|\delta L\|$, so
$\|\delta A\|\le \|\delta L\|/(2a_0)$ for a bounded variation.
The third equation follows by differentiating the exponential at
fixed $r$ and $M$. Gapless or unbounded variations require their
own domain estimates.

In a discrete spectral representation $L|a\rangle=x_a|a\rangle$,
these identities give the explicit vertex coefficient
\begin{equation}
 (\delta P_r)_{ab}
  =\frac{e^{-rh_M(x_a)}-e^{-rh_M(x_b)}}{x_a-x_b}
                  (\delta L)_{ab},
\label{eq:update-transfer-divided-difference}
\end{equation}
with the derivative limit on coincident eigenvalues. For the BRST
convention $\BR L=[L,c]$, $\BR\chi=-c\chi$, they imply
\begin{equation}
 \BR P_r=[P_r,c],\qquad \BR(P_r\chi)=-cP_r\chi.
\label{eq:update-full-field-brst-transfer}
\end{equation}
Thus the transferred field is covariant even though the full-field
transfer is not literally BRST-inert. Background inertness and
full-field covariance remain compatible, distinct statements.

Second variations also belong in the contact coefficients:
\begin{equation}
\begin{aligned}
 A\,\delta_1\delta_2A+(\delta_1\delta_2A)A
    ={}&\delta_1\delta_2L\\
      &-(\delta_1A)(\delta_2A)-(\delta_2A)(\delta_1A).
\end{aligned}
\label{eq:update-transfer-second-variation}
\end{equation}
Differentiating Duhamel's formula gives one mixed variation of
$A$ and two terms with ordered single variations. In the bounded
variation setting, using $L\ge0$ and $M\ge0$ so that
$\|P_v\|\le1$, these three terms give
\begin{equation}
 \|\delta_1\delta_2 P_r\|
 \le r\|\delta_1\delta_2 A\|
       +r^2\|\delta_1 A\|\|\delta_2 A\|.
 \label{eq:final-brst-transfer-contact-bound}
\end{equation}
Each ordered insertion has simplex volume $r^2/2$. For the first
term, the same Sylvester estimate bounds $\|\delta_1\delta_2 A\|$
by $(\|\delta_1\delta_2 L\|+2\|\delta_1 A\|\|\delta_2 A\|)/(2a_0)$.
These estimates concern the specified operator domain. A graph
additionally needs a bound uniform in its routed momenta.
For an interaction $V[P_r\chi]$, its second background
variation contains both $V''[P_r\chi](\delta_1P_r\chi,\delta_2P_r\chi)$
and $V'[P_r\chi](\delta_1\delta_2P_r\chi)$, in addition to variations
of any other fields, measures, or coefficients. Omitting these terms
would test a different action. Their full routed growth must enter
\eqref{eq:anisotropic-complete-envelope} or the residual/forest estimates.

The positive Gaussian Schur calculation supplies a separate matching
test. For a finite quadratic block
$\left(\begin{smallmatrix}B&E^*\\E&D\end{smallmatrix}\right)>0$,
let $\psi=F\chi+Oz$ and $D>0$. Completing the square replaces
$z$ by the independent Gaussian fluctuation $z+D^{-1}E\chi$.
Thus $\psi$ is the sum of a linear readout of $\chi$ and a
readout of that fluctuation. With
\begin{equation}
 \mathscr K=B-E^*D^{-1}E,\qquad Q=F-OD^{-1}E,
\label{eq:update-schur-data}
\end{equation}
its covariance blocks are
\begin{equation}
\begin{aligned}
 C_{\chi\chi}&=\mathscr K^{-1},\qquad
 C_{\psi\chi}=Q\mathscr K^{-1},\\
 C_{\psi\psi}&=Q\mathscr K^{-1}Q^*+OD^{-1}O^*.
\end{aligned}
\label{eq:update-schur-covariances}
\end{equation}
Matching \eqref{eq:update-three-source-blocks} requires
$\mathscr K=C_0^{-1}$, $Q=P_r$, and $OD^{-1}O^*=0$.
For positive invertible $D$ the last condition implies $O=0$, since
its quadratic form is $\|D^{-1/2}O^*f\|^2$. A sampled independent
positive bulk fluctuation therefore cannot be removed while keeping
the same complete three-block law.

At finite resolution the unconstrained Gaussian integration leaves
$\tfrac12\operatorname{Tr}\log\mathscr K+
\tfrac12\operatorname{Tr}\log D$ in the negative logarithm of its
unnormalized partition function. These terms and
\eqref{eq:update-measure-response} contribute to a varying-background
A4 equation. A normalized deterministic constraint is a different
measure. Infinite-dimensional determinant traces, a dynamical carrier,
and the gauge/ghost contributions need their own definitions; the
fixed-plane UV theorems do not certify them simply by the notation
$\operatorname{Tr}\log$.


\subsection{BV fields, source signs, and the master identity}
\label{sec:bv}

Use a complete field set \(\phi^A\), including \(a,c,\bar c,B\)
and matter, with Grassmann parity \(\epsilon_A\).
Introduce antifields \(\phi_A^*\) of parity \(\epsilon_A+1\)
and ghost number \(-1-\operatorname{gh}\phi^A\).
The nonminimal pair \((\bar c,\bar c^*)\) is retained.
Our antibracket convention is
\begin{equation}
 (F,G)=\int\dd^4x_{\bb}\sum_A
 \left[
 \frac{\delta_rF}{\delta\phi^A}
 \frac{\delta_lG}{\delta\phi_A^*}
 -
 \frac{\delta_rF}{\delta\phi_A^*}
 \frac{\delta_lG}{\delta\phi^A}
 \right].
 \label{eq:antibracket}
\end{equation}
For a homogeneous functional $F$ and coordinate $z$, our
variation convention is
$\delta F=\delta z\,\delta_lF/\delta z
             =(\delta_rF/\delta z)\,\delta z$.
The two derivatives differ by $(-1)^{|z|(|F|+|z|)}$.
For a closed off-shell BRST algebra, choose
\begin{equation}
 \Sigma=S_{\rm gf}^{\rm comp}
       +\int\dd^4x_{\bb}\sum_A(-1)^{\epsilon_A}
                   \phi_A^*\,\BR\phi^A.
 \label{eq:bv-action}
\end{equation}
With this convention, \(\BR=(\Sigma,\,\cdot\,)\) on fields.
The displayed gauge and nonminimal terms are
\[
\begin{aligned}
 \Sigma={}&S_{\rm gf}^{\rm comp}
  +\int\dd^4x_{\bb}\,\Bigl[
       a^{*a\mu}(D_\mu c)^a
       -c^{*a}\BR c^a\\
   &\hspace{3.8em}-\bar c^{*a}B^a
       +\hbox{matter antifield terms}\Bigr].
\end{aligned}
\]
The antifield-independent part of $\tfrac12(\Sigma,\Sigma)$
is $\BR S_{\rm gf}^{\rm comp}$, which vanishes by
\eqref{eq:completed-action} and nilpotence. The terms linear in
antifields vanish because $\BR^2\phi^A=0$. Thus
\((\Sigma,\Sigma)=0\).
For more general gauge systems the assumed completed master action
includes the additional antifield terms needed by its algebra.

The physical Lorentzian functional and the Euclidean coefficient
functional have distinct source conventions. To state the dictionary,
first choose the completed independent physical source coordinates;
any auxiliary-to-physical field map is included in the source pairing
and the transformed action. The BRST transformations are pushed
forward with that field map, and the antifield map preserves the BV
pairing. A nonlocal source normalization does not automatically
preserve an undeformed nonlinear gauge transformation. A real field has one source coordinate,
a genuinely complex pair has independent conjugate source coordinates,
and Dirac and ghost pairs retain their declared odd ordering. The
physical conjugation contour is imposed only after differentiation.

% Source conditions: lem:ST-stability-internal-only-early; thm:BRST-Nielsen.
% The following linear cotangent-lift identities are derived in the stated BV convention.
For example, let a Grassmann-even, parity- and ghost-number-preserving
linear map \(T\) be invertible on the admitted reduced domain,
with \(\BR T=0\). Write \(\psi^A\) for auxiliary coordinates
and \(\phi^A=T^A{}_B\psi^B\) for the physical source coordinates.
The nonlocal field map and its source cotangent map use the same
sector-plane density. On the Euclidean coefficient space,
\begin{equation}
\begin{aligned}
 \phi^A(x)&=\sum_B\int\dd\mu_{E,\bb}(y)\,
          T^A{}_{B;\bb}(x,y)\psi^B(y),\\
 \mathcal J^{\rm aux}_B(y)
   &=\sum_A\int\dd\mu_{E,\bb}(x)\,
             \mathcal J_A(x)T^A{}_{B;\bb}(x,y),\\
 \psi_B^*(y)&=\sum_A\int\dd\mu_{E,\bb}(x)\,
             \phi_A^*(x)T^A{}_{B;\bb}(x,y).
\end{aligned}
\label{eq:final-labelled-source-map}
\end{equation}
Every field and source in this display belongs to $\bb$. The analogous
Lorentzian formula uses $\dd^4x_{\bb}$ and the kernel on its own admitted
source domain; its existence is not inferred from a Euclidean Gaussian
identity. In condensed notation $\BR\phi^A=T^A{}_B\BR\psi^B$; repeated
field labels then include precisely these measure-weighted integrations. The nonzero
entries of \(T\) are Grassmann even, have ghost number zero,
and connect fields of equal parity and ghost number.
At fixed background and other data determining
\(T\), the antifield transformation is the cotangent lift:
\[
 \sum_A\phi_A^*\,\delta\phi^A
  =\sum_B\psi_B^*\,\delta\psi^B,\qquad \delta T=0.
\]
It preserves the canonical field--antifield pairing. The same
transformation preserves the source pairing and the signed
antifield term:
\[
 \sum_B(-1)^{\epsilon_B}\psi_B^*\BR\psi^B
    =\sum_A(-1)^{\epsilon_A}\phi_A^*\BR\phi^A .
\]
Thus the canonical antibracket and its declared signs are
preserved character by character. For the corresponding
finite-dimensional change of variables, the measure factor is
\(\mathcal D_{\bb}\psi=\operatorname{Ber}(T^{-1})\mathcal D_{\bb}\phi\);
its continuum counterpart requires the declared determinant and
measure prescription. It may be absorbed only into a compatible
normalization.
For a nonlinear map one needs its graded canonical lift and
Jacobian on the reduced domain; the free covariance identity alone
does not provide them.

On the fixed carrier algebra, with antifields suppressed in this
display, the Lorentzian source template is
\begin{equation}
\begin{split}
 \vct Z^L_{\epsilon,\bb}[\mathcal J]
  =\int\mathcal D_{\bb}\phi\;
  \operatorname{Exp}_{\mathfrak A,\bb}\Bigg[
   &\ii u_{\epsilon,\bb}S_{{\rm sc},\epsilon,\bb}[\phi]\\
   &+\ii a_{{\rm src},\epsilon,\bb}
       \int\dd^4x_{\bb}\,\mathcal J_A\phi^A\\
   &-\frac{\varepsilon}{2}\mathcal Q_{\rm B}[\phi]
  \Bigg].
 \label{eq:physical-L-source-functional}
% Monograph source: eq:A5-postGTL; def:A5-directed-exp-log.
\end{split}
\end{equation}
Here \(S_{{\rm sc},\epsilon,\bb}\) denotes the complete
gauge-fixed coefficient
\(S_{{\rm inv},{\rm sc}}^{\rm comp}+\BR\Psi_{\rm sc}\),
including the ghost sector; only antifield couplings are suppressed
in this display. The exponential is evaluated componentwise in the
central idempotents. The optional convergence form \(\mathcal Q_{\rm B}\)
is nonnegative on its declared commuting field block; fermionic and
ghost quadratic forms remain in the BRST action. The positive
convergence term is branch neutral. It is distinct from the ordered
Green-function datum \(\ii\eta\epsvec\); an exact regulated
matrix inverse requires specification of this datum with the
complete Lorentzian Hessian before inversion. If a convergence
regulator is retained in a BRST
identity, it must satisfy the stated BRST-invariance or BRST-exactness
hypothesis; positivity of \(\mathcal Q_{\rm B}\) alone does not
establish that hypothesis.

% Monograph source: def:A5-directed-exp-log; eq:A5-physical-carrier-involution.
Choose a normalization \(\vct{\mathcal N}\) independent of
sources and, when present, antifields, with nonzero character bodies.
It must also respect each background symmetry and extended BRST
transformation used in the corresponding functional identity.
On a source neighbourhood where both normalized character bodies
are nonzero, choose local logarithm branches. For an even character
\(z_\kappa=b_\kappa+n_\kappa\), with nonzero scalar body
\(b_\kappa\) and positive Grassmann degree \(n_\kappa\), this means
\[
 \log z_\kappa=\log b_\kappa+
  \sum_{r\geq1}\frac{(-1)^{r+1}}{r}
       \left(\frac{n_\kappa}{b_\kappa}\right)^r .
\]
The nilpotent sum terminates in a finite Grassmann algebra and is
otherwise a formal source-degree series. When a compatible vacuum
normalization gives \(\vct Z^L[0]\vct{\mathcal N}^{-1}=1\), use
the local logarithms through \(1\). On a physical conjugate-source
contour, the normalization and branches are chosen so that
\[
 \vct{\mathcal N}^{\ddagger}=\vct{\mathcal N},\qquad
 \operatorname{Log}_{\mathfrak A,\bb}(X^\ddagger)
  =\bigl(\operatorname{Log}_{\mathfrak A,\bb}X\bigr)^\ddagger .
\]
The physical connected generator and its mean fields are
\begin{equation}
 \vct{\mathcal W}^L_{\epsilon,\bb}
   =-\ii\operatorname{Log}_{\mathfrak A,\bb}
       (\vct Z^L_{\epsilon,\bb}\vct{\mathcal N}^{-1}),
 \qquad
 \bar\phi^A
  =a_{{\rm src},\epsilon,\bb}^{-1}
       \frac{\delta_l\vct{\mathcal W}^L_{\epsilon,\bb}}
            {\delta\mathcal J_A}.
 \label{eq:physical-connected-source-rule}
\end{equation}
In this functional subsection, compact symbols $\phi^A$, $J_A$ and
$\phi_A^*$ mean $\phi^A_{\bb}$, $J_{A,\bb}$ and $\phi^*_{A,\bb}$ at
one fixed branch and copy; their variational pairings use the displayed
$\dd^4x_{\bb}$. This local abbreviation is not a sum or an equality of
fields on different supports. The source carrier is still
$a_{{\rm src},\epsilon,\bb}$, including its sector factor only once.

For ordered source insertions on this same support,
\begin{equation}
\begin{aligned}
 \langle\phi^{A_1}\cdots\phi^{A_n}\rangle_c
  &=\ii^{1-n}a_{{\rm src},\epsilon,\bb}^{-n}
    \frac{\delta^n\vct{\mathcal W}^L_{\epsilon,\bb}}
     {\delta\mathcal J_{A_1}\cdots\delta\mathcal J_{A_n}},\\
 a_{{\rm src},\epsilon,\bb}^{-n}
   &=\eta^{-n}\ee^{\,n}\quad({\rm OQFT}).
\end{aligned}
 \label{eq:physical-source-derivative-parity}
\end{equation}
The rightmost derivative acts first, and the fields on the left
are in the displayed order. All source terms use source-first
coordinates: rewriting a field-first odd source term includes its
Koszul minus sign in the definition of that source coordinate.
For example, \(\bar\eta_\psi\psi+\bar\psi\eta_\psi\) becomes
\(\mathcal J_\psi\psi+\mathcal J_{\bar\psi}\bar\psi\) with
\(\mathcal J_\psi=\bar\eta_\psi\) and
\(\mathcal J_{\bar\psi}=-\eta_\psi\), and likewise for ghosts.
Their carrier power is even for even \(n\) and odd for odd \(n\),
independently of Grassmann parity. On a reduced source domain with
locally invertible graded source-to-mean-field map, the Lorentzian
Legendre transform is
\begin{equation}
 \vct\Gamma^L_{\epsilon,\bb}
  =\vct{\mathcal W}^L_{\epsilon,\bb}
   -a_{{\rm src},\epsilon,\bb}
        \int\dd^4x_{\bb}\,\mathcal J_A\bar\phi^A .
 \label{eq:physical-L-Legendre}
\end{equation}
All products and logarithms here belong to the graded,
source-complexified carrier algebra.

% Monograph source: eq:A5-total-physical-real-spaces; eq:A5-directed-connected-generator-algebra.
If the action, measure, integration contour and source pairing
respect the total \(\ddagger\) involution, then on the physical
conjugate-source contour
\[
 (\vct Z^L\vct{\mathcal N}^{-1})^\ddagger
       =\vct Z^L\vct{\mathcal N}^{-1},\qquad
 (\vct{\mathcal W}^L)^\ddagger=-\vct{\mathcal W}^L,\qquad
 (\vct\Gamma^L)^\ddagger=-\vct\Gamma^L .
\]
Thus the physical connected and effective functionals belong to
the Grassmann-even anti-fixed subspace of \(\mathfrak B_{\bb}\).
Their bosonic scalar bodies lie in
\(\operatorname{span}_{\RR}\{\ii,\ee\}\); this is a reality
condition, not a statement of homogeneous carrier parity.
Independent complex-source calculations still take place in the
full graded carrier algebra.

For the algebraic BRST proof below we work in one admitted
Euclidean character representation, with its source and antifield
coordinates normalized consistently. Thus the unbold symbols
\(\Sigma,Z,W,\Gamma\) in this subsection denote coefficient
functionals. Before any admitted continuation, the Lorentzian source
coefficient has character value
\(\chi_\kappa(a_{{\rm src},\epsilon,\bb})
 =\kappa\eta\) in O--QFT; this sign is part of the map from
\(\mathcal J\) to the coefficient source, not a second sector
fold. The corresponding Euclidean source functional and Legendre
transform are
\begin{align}
 Z[J,\phi^*]
 &=\frac1{\mathcal N}
     \int\mathcal D_{\bb}\phi\,
        \exp\left[-\Sigma[\phi,\phi^*]
                    +\int\dd^4x_{\bb}\,J_A\phi^A\right],
 \label{eq:Z-functional}\\
 W&=\log Z,\qquad
 \varphi^A=\frac{\delta_lW}{\delta J_A},\nonumber\\
 \Gamma[\varphi,\phi^*]&=\int\dd^4x_{\bb}\,J_A\varphi^A-W.
 \label{eq:Legendre}
\end{align}
The graded source pairing is written source first and used
consistently; \(J\) is eliminated at fixed antifields.
The sources have \(\epsilon(J_A)=\epsilon_A\) and
\(\operatorname{gh}(J_A)=-\operatorname{gh}(\phi^A)\), making
\(J_A\phi^A\) even and of ghost number zero.
The normalization \(\mathcal N\) is source and antifield
independent, background-gauge invariant, and chosen invariant
under the gauge-parameter BRST extension used below. A vacuum normalization is allowed
when it has these properties. The field- and antifield-dependent
tree action is \(\Sigma\); its irrelevant vacuum constant is
fixed by \(\mathcal N\).
Equation~\eqref{eq:directed-action} restores the regulator-free
physical action terms that have already been shown to be purely
carrier odd. A complete regulated or loop-corrected effective
functional is instead reconstructed from both characters, using one
signature and one Legendre convention for both representatives:
\begin{equation}
 \vct\Gamma_{\epsilon,\bb}
   =\mathsf P_+\Gamma^{\rm rep}_{\epsilon,\bb}
    +\mathsf P_-
                  \Gamma^{\rm rep}_{\bar\epsilon,\bb},
 \qquad \sigma_{\bar\epsilon}=-\sigma_\epsilon.
 \label{eq:effective-character-reconstruction}
% Monograph source: eq:Omni-connected-effective-carrier-restoration.
\end{equation}
The two representatives are the simultaneous character projections of
one completed functional, including its field and antifield coordinates.
The opposite-branch notation retains the projected source signs,
ordering, adjoint and endpoint rules, and logarithm branches. It is
not obtained by changing only the sign of an isolated denominator,
nor by choosing an unrelated solution of a scalar master equation.
Its even and odd carrier parts are respectively the half-sum and
the half-difference times \(\ee\). The Euclidean proof below
applies to the Euclidean coefficient convention. A physical
Lorentzian identity additionally requires an admitted continuation
that transports the source and antifield normalization and the
Legendre convention consistently. In general neither
\(u_{\epsilon,\bb}\Gamma^{\rm rep}_{\epsilon,\bb}\) nor one
transported complex coefficient reconstructs this functional.

\begin{theorem}[Conditional exact Slavnov--Taylor identity]
\label{thm:ST}
For the closed off-shell system and master action
\eqref{eq:bv-action}, assume an invariant anomaly-free measure
and a common domain for the BRST change of variables, the indicated
source and antifield derivatives (with the differentiations admitted
under the functional integral), the logarithm, and the locally
invertible source-to-mean-field map. Assume the variation has no
field-space boundary term. Then the fixed-plane
effective action satisfies
\begin{equation}
 \mathcal S(\Gamma_{\bb})
      :=\frac12(\Gamma_{\bb},\Gamma_{\bb})=0
 \label{eq:ST-master}
\end{equation}
as a formal perturbative identity at every admitted loop order.
If the same completed graphs and required insertions satisfy the
absolute estimates of Section~\ref{sec:uv} on this domain, the
identity holds for their finite coefficients. When the coefficients
are obtained by a comparison limit, the contractions in the
antibracket must admit that limit on the same domain.
\end{theorem}

\begin{proof}
Perform the infinitesimal change of variables
\(\phi^A\mapsto\phi^A+\theta\BR\phi^A\), with constant odd
\(\theta\), in \eqref{eq:Z-functional}.
The measure and the completed action contribute no breaking.
Since the sources are inert, the left-BRST variation of the source
term is
\[
 \BR(J_A\phi^A)=(-1)^{\epsilon_A}J_A\BR\phi^A .
\]
From \eqref{eq:bv-action},
\(\delta_l\Sigma/\delta\phi_A^*
   =(-1)^{\epsilon_A}\BR\phi^A\).
The change of variables therefore gives the functional identity
\begin{equation}
 \int\dd^4x_{\bb}\sum_A J_A
       \frac{\delta_l Z}{\delta\phi_A^*}=0,
 \qquad
 \int\dd^4x_{\bb}\sum_A J_A
       \frac{\delta_l W}{\delta\phi_A^*}=0.
 \label{eq:Z-W-identity}
\end{equation}
In the Legendre variation, the change of $J$ cancels between
$J_A\varphi^A$ and $W$. The remaining field variation has
$J_A$ on the left. Hence
\(\delta_r\Gamma/\delta\varphi^A=J_A\) and
\(\delta_l\Gamma/\delta\phi_A^*
      =-\delta_lW/\delta\phi_A^*\).
For even \(\Gamma\), the two terms of
\eqref{eq:antibracket} combine into twice the first term.
Substitution of \eqref{eq:Z-W-identity} gives
\eqref{eq:ST-master}.
Expanding that identity in the loop parameter produces the
Slavnov--Taylor hierarchy. The UV conclusion uses the independent
absolute bounds for these very graphs, rather than a cancellation
of divergent terms.
\end{proof}

% Monograph source: thm:BRST-Nielsen; prop:STexact.
% The regulator-insertion signs follow from eq:bv-action and eq:Legendre.
For completeness, an antifield-independent Euclidean regulator
\(R_\varepsilon[\phi]\) inserted as
\(\exp[-\Sigma-R_\varepsilon+J\cdot\phi]\) gives, with the
same invariant measure and ordinary Legendre transform,
\begin{equation}
 \int\dd^4x_{\bb}\,J_A
       \frac{\delta_l W_\varepsilon}{\delta\phi_A^*}
       =-\langle\BR R_\varepsilon\rangle,\qquad
 \frac12(\Gamma_\varepsilon,\Gamma_\varepsilon)
       =\left.\langle\BR R_\varepsilon\rangle
          \right|_{J=J(\varphi)} .
 \label{eq:regulated-ST-breaking}
\end{equation}
Here no regulator term is subtracted from the Legendre transform.
Indeed, the change of variables gives
\(\langle\BR(J\cdot\phi)\rangle
  =\langle\BR R_\varepsilon\rangle\), and the antifield and
Legendre signs then give the displayed equations. A BRST-inert
regulator, or an exact regulator \(R_\varepsilon=\BR\Upsilon_\varepsilon\),
has zero insertion. Otherwise its removal requires a limit in
which this insertion vanishes; an explicit factor of
\(\varepsilon\) alone does not prove that limit. The same
statement holds separately in both carrier characters.

The full carrier-valued identity requires the same hypotheses in
both character components. Project every field, antifield, source,
and coefficient simultaneously, and write
\(\Gamma_{\kappa}
  :=\chi_\kappa(\vct\Gamma)\).
Because the carrier coefficients are Grassmann even and central,
the character maps commute with the declared graded derivatives.
The carrier-valued Slavnov functional is
\begin{equation}
 \mathcal S_{\mathfrak A,\bb}(\vct\Gamma)
  :=\sum_{\kappa=\pm1}\mathsf P_\kappa
        \frac12(\Gamma_{\kappa},\Gamma_{\kappa})
  =0.
 \label{eq:carrier-ST}
% Monograph source: thm:BRST-Nielsen; eq:A5-algebra-valued-Slavnov-functional.
\end{equation}
Each bracket uses its projected canonical field and antifield
coordinates, with exactly the ordering of \eqref{eq:antibracket}.
At a retained regulator, apply \eqref{eq:regulated-ST-breaking}
in each character. The unbroken form holds when both breaking
insertions vanish on the common source domain. In the removal
limit this requires control of those insertions together with
the measure and any field-space boundary terms. The two
character identities enter the displayed reconstruction independently.

When $\Gamma=\Gamma_{\rm red}-\int\bar c^*B$ and the reduced
functional is independent of $\bar c^*$ and $B^*$, the only
nonminimal term in the antibracket is the one coming from this
linear dependence. Since $\bar c$ is odd,
$\delta_r\Gamma_{\rm red}/\delta\bar c
=-\delta_l\Gamma_{\rm red}/\delta\bar c$. This sign combines
with the coefficient $-B$ to give
\begin{equation}
 \frac12(\Gamma_{\rm red},\Gamma_{\rm red})_{\rm min}
       +\int\dd^4x_{\bb}\,B^a
             \frac{\delta_l\Gamma_{\rm red}}{\delta\bar c^a}=0 .
 \label{eq:reduced-ST}
\end{equation}
The minimal bracket includes gauge, ghost, and matter antifields.
A bracket containing only the minimal pairs cannot by itself
represent the displayed \(B\,\delta\Gamma/\delta\bar c\) term.

Restoring \(\hbar\) as a formal loop-counting parameter, write
\(\Gamma=\Sigma+\sum_{\ell\geq1}\hbar^\ell\Gamma_\ell\).
The linearized operator
\(\mathcal B_\Sigma F=(\Sigma,F)\) obeys
\(\mathcal B_\Sigma^2=0\), and
\begin{equation}
 \mathcal B_\Sigma\Gamma_\ell
 =-\frac12\sum_{r=1}^{\ell-1}
                 (\Gamma_r,\Gamma_{\ell-r}).
 \label{eq:loop-ST}
\end{equation}
The effective action satisfies the nonlinear Slavnov--Taylor identity.
BRST-closed physical representatives need not be BRST-exact.

\subsection{Anomalies and the quantum action principle}

The anomaly-free measure hypothesis is essential. UV convergence
does not cancel an anomalous Jacobian or prove that a formal
field-space integration has no boundary term. Algebraic
renormalization organizes possible symmetry breakings cohomologically
\cite{PiguetSorella,Slavnov1972,Taylor1971,ZinnJustin1975}.
Its locality and quantum-action-principle hypotheses must nevertheless
be established for the admitted nonlocal completion and insertions;
their validity in the undeformed local theory does not establish them
for that completion.

Theorem~\ref{thm:ST} uses the explicitly assumed invariant
completed functional. It is not an existence proof of a
nonperturbative measure, a proof that a chiral representation is
anomaly free, or a proof of global gauge fixing. These tasks are
part of constructing the completion to which the theorem applies.

\subsection{Background Ward identity}

Assume the complete action, measure, domains, and antifield/source
transformations are covariant under
\eqref{eq:background-transform}. A background change of variables
then yields
\begin{equation}
 \mathcal W_\omega\Gamma_{\bb}=0.
 \label{eq:background-Ward}
\end{equation}
For a functional \(F\), the generator is most unambiguously
specified by its variations:
\begin{align}
 \mathcal W_\omega F
 =\int\dd^4x_{\bb}\bigg\{&
  (\overline D_\mu\omega)^a
      \frac{\delta F}{\delta\overline{\mathcal A}_\mu^a}
      \nonumber\\
 &+\sum_A(\delta_\omega\phi^A)
               \frac{\delta_lF}{\delta\phi^A}
  +\sum_A(\delta_\omega\phi_A^*)
               \frac{\delta_lF}{\delta\phi_A^*}\bigg\}.
 \label{eq:Ward-generator}
\end{align}
Here \(\phi^A\) is the complete quantum field set, and the
antifields transform in the corresponding dual representations.
When \(\mathcal W_\omega\) acts on \(\Gamma\), its quantum-field
arguments are the mean fields \(\varphi^A\).
Cyclicity of an invariant trace shows covariance of a kernel
insertion, but invariance of the whole functional is required for
\eqref{eq:background-Ward}.

\subsection{Nielsen identity and the meaning of gauge independence}
\label{sec:Nielsen}

Let the gauge parameter \(\zeta\) enter only through the
gauge-fixing fermion. The derivative below is taken at fixed
canonical mean fields and antifields, fixed support and carrier,
and fixed kernel and physical-source normalization maps. A
\(\zeta\)-dependent source or field redefinition would contribute
its own variation and is not included in this fixed-map formula.
Introduce the BRST doublet
\[
 \BR\zeta=\chi,\qquad \BR\chi=0,\qquad\chi^2=0.
\]
Assume the extended completed functional and its insertions obey
the extended identity, in the convention of
\eqref{eq:antibracket},
\begin{equation}
 \frac12(\Gamma,\Gamma)+\chi\,\partial_\zeta\Gamma=0.
 \label{eq:extended-ST}
\end{equation}
Expand $\Gamma=\Gamma_0+\chi\Xi_\zeta$, with $\chi$ placed
on the left, $\Gamma_0=\Gamma|_{\chi=0}$, and
$\Xi_\zeta=(\partial_l\Gamma/\partial\chi)|_{\chi=0}$.
The bracket is odd, so
$(\Gamma_0,\chi\Xi_\zeta)=-\chi(\Gamma_0,\Xi_\zeta)$.
The coefficient of \(\chi\) in \eqref{eq:extended-ST} gives
\begin{equation}
 \partial_\zeta\Gamma\big|_{\chi=0}
       =(\Gamma,\Xi_\zeta)\big|_{\chi=0}
       =\mathcal B_\Gamma\Xi_\zeta\big|_{\chi=0}.
 \label{eq:Nielsen}
\end{equation}
For example, at the classical level
\(\Xi_\zeta=\partial_\zeta\Psi\) in a compatible gauge-fixing
construction, so the variation is BRST exact.
This is the Nielsen identity in the declared convention
\cite{Nielsen1975}.

The corresponding carrier-valued Nielsen identity is read in the
same componentwise sense:
\begin{equation}
 \partial_\zeta\vct\Gamma\big|_{\chi=0}
  =\sum_{\kappa=\pm1}\mathsf P_\kappa
       (\Gamma_{\kappa},\Xi_{\zeta,\kappa})\big|_{\chi=0},
 \qquad
 \Xi_{\zeta,\kappa}
  :=\chi_\kappa(\vct\Xi_\zeta).
 \label{eq:carrier-Nielsen}
% Monograph source: thm:BRST-Nielsen.
\end{equation}
The background Ward identity is lifted with the same simultaneous
projection of all variables. These lifts preserve the coefficient
sign conventions; they do not assign a fixed carrier parity to a
nonlinear functional or to its loop corrections.

At zero external sources and antifields, consider an admissible
$\zeta$-independent BRST-closed observable $O$. When the measure
and field domain are fixed along the gauge-fixing variation,
differentiation of the normalized expectation gives
\begin{equation}
 \partial_\zeta\langle O\rangle
 =-\langle O\,\BR(\partial_\zeta\Psi)\rangle
   +\langle O\rangle\langle\BR(\partial_\zeta\Psi)\rangle=0.
 \label{eq:final-brst-gauge-expectation}
\end{equation}
The change-of-variables identity makes the two derivative
insertions vanish if $\partial_\zeta\Psi$ and
$O\partial_\zeta\Psi$ lie in its common domain. At the
effective-action level \eqref{eq:Nielsen} generates the
canonical transport $\partial_\zeta F=-(\Xi_\zeta,F)$.
The graded Jacobi identity shows that it transports
$(\Gamma_0,F)$ by the same equation. This preserves closed
and exact
representatives whenever the transport exists on the domain.

Gauge independence of physical LSZ quantities also requires the
same physical state and normalization along the gauge slice,
together with the necessary spectral and insertion limits.
Off-shell Green functions can depend on $\zeta$.
The physical kernel scale $a$ and measure $\rho_{\bb}$
enter the dynamics and are held fixed in \eqref{eq:Nielsen}.
\section{Analyticity, positivity, and the physical cut interface}
\label{sec:analytic}

The preceding symmetry identities constrain a completed functional on
its admitted domain. To extract physical amplitudes, one must also
control analytic continuation, positive physical states, and the relation
between cuts and boundary discontinuities. This section states these
requirements separately and tests them on the specified covariances.

\subsection{Complete monotonicity and analytic continuation}

Equation~\eqref{eq:bernstein} gives a holomorphic function on
\(\Re z>0\). For a compact subset of this half-plane, derivatives
may be taken under the integral because
\(\tau^n e^{-c\tau}\) is bounded for every \(c>0\).
If
\begin{equation}
 \int e^{\epsilon_{\rho,\bb}\tau}\rho_{\bb}(\dd\tau)<\infty
 \quad\hbox{for some }\epsilon_{\rho,\bb}>0,
 \label{eq:exponential-moment}
\end{equation}
the same argument extends holomorphy to
\(\Re z>-\epsilon_{\rho,\bb}\).
Compact support of \(\rho_{\bb}\), or positive exponential moments of
every order, gives an entire Laplace transform. A general
completely monotone function need not extend to a slit plane
\cite{Bernstein1929,Widder1941}.

\begin{proposition}[A gapped counterexample to automatic slit analyticity]
\label{prop:meromorphic-counterexample}
The normalized measure and kernel
\begin{equation}
 \rho_{\bb}=\sum_{n=1}^\infty2^{-n}\delta^{\rm heat}_{\bb,n},\qquad
 \Ph(z)=\frac1{2e^z-1}
 \label{eq:meromorphic-counterexample}
\end{equation}
satisfy the gap condition with \(\tau_*=1\), but their complex
continuation has poles away from the negative real axis.
\end{proposition}

\begin{proof}
For \(z\geq0\), the positive geometric series
\[
 \Ph(z)=\sum_{n\geq1}2^{-n}e^{-nz}
\]
proves complete monotonicity and \(\Ph(0)=1\). The gap bound is
\(\Ph(z)\leq e^{-z}\). The denominator vanishes at
\[
 z=-\log2+2\pi\ii n,\qquad n\in\mathbb Z.
\]
For \(n\neq0\), these singularities lie off the negative real axis.
\end{proof}

Even the elementary power kernel
\(\Ph(z)=(1+z/b)^{-\nu}\) has a singularity at \(z=-b\):
a pole for positive integer \(\nu\), and generally a branch point.
Strong Euclidean decay therefore does not imply a pole-free
Lorentzian continuation.

A second distinction concerns growth. Define the directed kinetic
invariant on the fixed plane by
\begin{equation}
 z_{\eta,\bb}(k):=k^2+\ii\eta\epsvec,
 \qquad \epsvec=\varepsilon\ee,
 \qquad \eta=\sigma_s\sigma_\epsilon.
 \label{eq:directed-kinetic-invariant}
\end{equation}
The continuation of \(e^{-tq_E^2}\) is the directed function
\begin{equation}
 \ExpDir\!\left(tz_{\eta,\bb}(k)\right)
 =\embed\!\left[e^{t(k^2+\ii\eta\varepsilon)}\right].
 \label{eq:heat-continuation}
\end{equation}
Its extracted coefficient grows as \(e^{tk^2}\) on real
timelike directions. It is not an oscillatory semigroup factor.
Entireness alone does not justify discarding a large contour arc.

\subsection{Local holomorphy from a common majorant}

Choose the independent external complex momenta \(z\) after imposing
momentum conservation. In the all-incoming convention,
\(p_n=-\sum_{\ell<n}p_\ell\), so their imaginary parts must sum to
zero. A common condition \(\Im p_\ell^0>0\) on all \(n\) legs
would violate that constraint.

\begin{theorem}[Holomorphy from a common graph majorant]
\label{thm:holomorphy}
Let \((\Xi,\nu_{\bb})\) be the positive, sigma-finite measure space of a graph
representation on a connected open set \(\Omega\subset\CC^N\).
For each \(z\), assume \(J(\cdot;z)\) is measurable, and outside
one \(\nu_{\bb}\)-null set assume \(J(\xi;\cdot)\) is holomorphic on
\(\Omega\). Suppose that for every \(U\Subset\Omega\) there is
\(M_U\in L^1(\dd\nu_{\bb})\) with
\begin{equation}
 |J(\xi;z)|\leq M_U(\xi),\qquad z\in U.
 \label{eq:complex-majorant}
\end{equation}
The bound holds outside a common null set for all \(z\in U\).
Then \(\mathcal M(z)=\int J(\xi;z)\dd\nu_{\bb}(\xi)\) is jointly
holomorphic on \(\Omega\). Local contour integrals and
derivatives may be exchanged with the representation integrals.
\end{theorem}

\begin{proof}
Fubini permits integration over a small contour in any one complex
variable before integration in \(\xi\). Cauchy's theorem makes
that contour integral vanish. Dominated convergence supplies
continuity and local boundedness. Morera's theorem gives separate
holomorphy, and Hartogs' theorem gives joint holomorphy.
Cauchy's formula on a slightly larger compact polydisc supplies an
integrable derivative bound.
\end{proof}


\begin{corollary}[Cauchy budgets for external derivatives]
\label{cor:complex-derivative-budget}
Under Theorem~\ref{thm:holomorphy}, suppose a closed polydisc
$\overline{\mathbb D}(z_0;r_1,\ldots,r_N)$ is inside the admitted
complex domain and one $M\in L^1(\dd\nu_{\bb})$ majorizes the complete
integrand there. For a multi-index $\alpha\in\mathbb N_0^N$,
\begin{equation}
 |\partial_z^\alpha\mathcal M(z_0)|
  \leq\alpha!\prod_{j=1}^N r_j^{-\alpha_j}
                  \int_\Xi M(\xi)\dd\nu_{\bb}(\xi).
 \label{eq:complex-Cauchy-budget}
\end{equation}
The estimate is for the stated complex-domain representation; a
uniform bound on an \emph{integrated real} Euclidean amplitude is not
a substitute for its pointwise complex majorant. No LSZ, retarded, or
physical boundary conclusion is obtained from this inequality alone.
\end{corollary}
\begin{proof}
The common majorant permits the graph integration and the
multivariable Cauchy contour integrals to be interchanged. Estimate
the resulting contour integrals by $\int M$ and the product of the
inverse radii.
\end{proof}

% Monograph source: lem:wick; eq:external-energy-majorant-bound; rem:two-forest-complex-majorant.
\begin{lemma}[A sufficient massive complex neighbourhood]
\label{lem:massive-neighbourhood}
Fix a compact set \(\mathcal K\) of real Euclidean external data
and use the routing of \eqref{eq:routing}, with momentum conservation
already imposed. Let every Schwinger factor be
\(E_i(\alpha)=e^{-m_i^2\alpha}\), with \(m_i>0\).
Suppose that on an open complex neighbourhood \(\mathcal U_0\)
of \(\mathcal K\), the full numerator has the finite expansion
\[
 N_G(k,z)=\sum_{|\gamma|\leq D}n_\gamma(z)k^\gamma,
 \qquad n_\gamma\in\mathcal O(\mathcal U_0),
\]
with this same loop-momentum degree \(D\) throughout
\(\mathcal U_0\). Assume the gap or inverse-moment hypothesis of
Theorem~\ref{thm:global-uv} for that \(D\), and
\eqref{eq:exponential-moment}. Then there is an open complex
neighbourhood \(\mathcal U\) of \(\mathcal K\), with
\(\overline{\mathcal U}\Subset\mathcal U_0\), on which the
pre-Gaussian Schwinger--Laplace graph has one common pointwise
\(L^1\) majorant. Its integral is holomorphic on each connected
component of \(\mathcal U\).
\end{lemma}

\begin{proof}
Choose finitely many real centres \(p^{(r)}\in\mathcal K\)
and sufficiently small complex balls around them whose closures lie
in \(\mathcal U_0\). The finite holomorphic coefficient family is
bounded on their union, so
\( |N_G(k,z)|\leq C(1+\norm{k}^{D})\) there with the same \(D\).
For a ball of radius \(\delta\), write
\[
 q_i(k,z)=q_i(k,p^{(r)})+h_i(z),\qquad
 |h_i(z)|\leq C_G\delta.
\]
The first term is real. With \(q^2=\sum_\mu q_\mu^2\),
the complex bilinear square, one has
\[
 \Re\bigl(q_i(k,z)^2\bigr)
 \geq\tfrac12|q_i(k,p^{(r)})|^2-C'_G\delta^2.
\]
Consequently a fixed-centre pointwise envelope is
\[
 C(1+\norm{k}^{D})\prod_i
 e^{-\frac12(\alpha_i+a^2\tau_i)|q_i(k,p^{(r)})|^2}
 e^{-\alpha_i(m_i^2-C'_G\delta^2)}
 e^{a^2C'_G\delta^2\tau_i}.
\]
Shrink \(\delta\) so that
\(C'_G\delta^2<\frac12\min_i m_i^2\) and
\(\lambda:=a^2C'_G\delta^2<\epsilon_{\rho,\bb}\).
The reduced masses remain positive, and the tilted measure
\(\rho_{\lambda,\bb}(\dd\tau):=e^{\lambda\tau}\rho_{\bb}(\dd\tau)\)
has finite mass. It has the same support and no new atom at zero.
For every inverse moment \(M_r[\rho_{\bb}]\) used in the UV estimate,
\[
 M_r[\rho_{\lambda,\bb}]
 \leq e^\lambda M_r[\rho_{\bb}]
       +\int_{(1,\infty)}e^{\lambda\tau}\rho_{\bb}(\dd\tau)
 <\infty.
\]
Thus the global Gaussian proof, with the widths multiplied by
\(1/2\), makes this envelope integrable in \(k,\alpha,\tau\).
The finite sum over the covering centres is a common majorant on
their union. The integrand is holomorphic there for every finite
parameter point, so Theorem~\ref{thm:holomorphy} applies on each
connected component.
\end{proof}

This is a local statement near Euclidean data. It does not produce
a full physical sheet or a scattering crossing path.
For stable kernels such as \(e^{-\sqrt z}\), the representing
measure has no positive exponential moment, so
Lemma~\ref{lem:massive-neighbourhood} does not apply. Their explicit
square-root continuation can still be studied on a declared
domain, but its branch geometry and graph majorant require an
independent argument. The missing hypothesis can correspond to an actual loss of regularity:
Proposition~\ref{prop:stable-bubble-nonanalytic} proves that the literal
massive stable-half scalar bubble is not \(C^6\) at zero Euclidean
external momentum.

The complex variables in Theorem~\ref{thm:holomorphy} and
Lemma~\ref{lem:massive-neighbourhood} are ordinary coefficient
variables. The Lorentzian boundary value is fixed separately by
\(\ii\eta\epsvec\); it is not derived by this Euclidean
holomorphy argument. For an invariant coefficient function
\(F^{\rm sheet}\), its directed continuation and extracted
boundary value are
\begin{align}
 F^{\rm dir}_{\bb}
   &:=\embed\circ F^{\rm sheet}\circ\embed^{-1},
       \nonumber\\
 F^{\rm coeff}_{\eta,\bb}(s)
   &:=\lim_{\varepsilon\downarrow0}
       \embed^{-1}\!\left[
       F^{\rm dir}_{\bb}(s+\ii\eta\epsvec)\right].
 \label{eq:analytic-function-directed-interface}
\end{align}
The entire physical graph is then reconstructed by
\eqref{eq:complete-graph-character-reconstruction}.
The same distinction applies to multiple invariant variables.

\subsection{Candidate poles, branches, and contour deformations}

Suppose a continued line numerator \(f\) is holomorphic near
\(D(k)=k^2-m^2=0\). If \(f\) is finite and nonzero there,
multiplication by \(f\) preserves the simple pole and multiplies
its residue by \(f|_{D=0}\). A zero can soften or cancel the pole.
A kernel singularity can add a pole, increase an existing order,
or introduce a new branch surface. Consequently, the precise
inertness statement is
\begin{quote}
The ordinary denominator family is unchanged. Pole preservation
holds locally where the numerator is holomorphic, finite, and
nonzero; the singularities and growth of that numerator require
their own analysis.
\end{quote}

The ordinary denominator Landau equations remain candidate
conditions associated with those denominators
\cite{Landau1959}. They are not the complete singularity
analysis when the kernel or completed vertices have additional
singular surfaces.

A permissible deformation of a loop contour requires an analytic
integrand between the contours, no crossed singularities or
pinches, and sufficient bounds on boundary arcs. A holomorphic
Euclidean parameter integral can instead define a continuation
without such a deformation, as in
Section~\ref{sec:gaussian-bubble}.
These two constructions should not be conflated.
Compact real-momentum bounds do not, on their own, establish
tempered physical boundary values.

\subsection{Reflection positivity is an independent requirement}
\label{sec:OS}

Here $x_E^0$ denotes the continued oriented GT length
$\gamma_{E,\bb}$, not an additional physical clock coordinate.
Let \(\theta(x_E^0,\mathbf x)=(-x_E^0,\mathbf x)\), and let
\(\Theta\) include reflection and conjugation. The ungraded formulas
in this subsection concern the Grassmann-even, gauge-invariant
observable algebra. Reflection positivity on that algebra requires
\begin{equation}
 \sum_{i,j}\overline{c_i}c_j
    \langle\Theta F_i\,F_j\rangle_E\geq0
 \label{eq:OS}
\end{equation}
for all finite families of observables smeared at
\(x_E^0>0\). Ordinary positivity of a Euclidean convolution
kernel is different. In particular, \(e^{t\Delta_E}\) spreads
positive-GT support across \(x_E^0=0\). Moving this heat
operator onto test functions does not preserve the OS test space.

\begin{proposition}[Heat damping does not preserve OS positivity]
\label{prop:OS-counterexample}
For \(t>0\) and \(\omega>0\), the following Euclidean GT-coordinate covariance
fails reflection positivity:
\begin{equation}
 C_t(u)=\int_{\RR}\frac{\dd p_{\bb}}{2\pi}\,
       \frac{e^{\ii pu}e^{-tp^2}}{p^2+\omega^2}
 \label{eq:heat-time-covariance}
\end{equation}
\end{proposition}

\begin{proof}
All carrier-frequency moments exist. Set
\[
 c_0=\int\frac{\dd p_{\bb}}{2\pi}
           \frac{e^{-tp^2}}{p^2+\omega^2}>0,\qquad
 c_2=\int\frac{\dd p_{\bb}}{2\pi}
           \frac{p^2e^{-tp^2}}{p^2+\omega^2}>0 .
\]
Then \(C_t(u)=c_0-\tfrac12c_2u^2+O(u^4)\).
Positive-GT probes at \(u=\epsilon,2\epsilon\) would require
the Hankel matrix
\[
 \begin{pmatrix}
 C_t(2\epsilon)&C_t(3\epsilon)\\
 C_t(3\epsilon)&C_t(4\epsilon)
 \end{pmatrix}
\]
to be positive semidefinite. Its determinant is
\begin{equation}
 -c_0c_2\epsilon^2+O(\epsilon^4)<0
 \label{eq:OS-negative-determinant}
\end{equation}
for sufficiently small \(\epsilon>0\).
Smooth positive-GT approximations to these point probes preserve
the violation by continuity.
\end{proof}

At fixed spatial momentum, a four-dimensional heat-damped scalar
covariance reduces to \eqref{eq:heat-time-covariance} times a
positive constant, with
\(\omega^2=\mathbf p^2+m^2\).
Spatial wave packets localized near that momentum give the same
obstruction. Thus the failure is not an artifact of a nonnormalizable
single spatial mode.


\paragraph{A two-probe test for arbitrarily slow damping}
The preceding Taylor expansion uses a second frequency moment. The
same obstruction survives when that moment diverges. On a fixed
spatial fibre let $\omega>0$ and replace $e^{-tp^2}$ in
\eqref{eq:heat-time-covariance} by a bounded, measurable $F(p^2)\ge0$
which is positive on a set of positive measure and tends to zero
at infinity. Call the covariance $C_F$. Its Fourier coefficient is
integrable, so $C_F$ is real, even and continuous. Since
$1-\cos(pu)>0$ outside a discrete set for $u>0$, one has
$C_F(u)<C_F(0)$. Yet
\[
 \frac{C_F(0)-C_F(u)}{u}
 =\int_{\RR}\frac{\dd v}{2\pi}
    \frac{1-\cos v}{v^2+\omega^2u^2}F(v^2/u^2)
 \longrightarrow0\qquad (u\downarrow0).
\]
Dominated convergence uses the integrable majorant
$\|F\|_\infty(1-\cos v)/v^2$.

This determines a pair of reflected probes as close to the origin
as desired. Choose $u_0>0$ so small that $C_F$ is positive on
$[0,2u_0]$ and set $f(u)=\log[C_F(u)/C_F(0)]$. Then $f(u_0)<0$
and $f(u)/u\to0$ as $u\downarrow0$. If
$C_F(0)C_F(2u)\ge C_F(u)^2$ at each $u=2^{-n}u_0$, $n\ge1$,
then $f(u_0)/u_0\ge f(2^{-n}u_0)/(2^{-n}u_0)$ for every $n$,
a contradiction. At one of these scales the determinant is strictly
negative. Continuity permits the two positions $0,u$ to be moved
to $\delta,\delta+u$ with $\delta>0$ and smeared inside positive
GT. In particular this argument covers every nonzero finite positive
Laplace damping law with no atom at zero on a massive spatial fibre.
For these laws the smeared fibre test depends continuously on the
spatial momentum; a wave packet near the selected momentum gives a
four-dimensional violation. This test concerns the specified scalar
covariance. Its application to an observable in the reconstructed
theory requires the covariance of that observable.

There is a related spectral restriction. If a nonzero physical
scalar two-point coefficient has an unsubtracted positive spectral
representation
\begin{equation}
 C_{\rm phys}(x)=\int_{[0,\infty)}
          \frac{\rho_{\rm phys,\bb}(\dd\mu^2)}{x+\mu^2},
 \qquad \rho_{\rm phys,\bb}\geq0,\quad x>0,
 \label{eq:physical-spectral}
\end{equation}
then monotone convergence implies
\begin{equation}
 \lim_{x\to\infty}xC_{\rm phys}(x)
       =\rho_{\rm phys,\bb}([0,\infty))\in(0,\infty].
 \label{eq:spectral-obstruction}
\end{equation}
A covariance \(\Ph(a^2x)/(x+m^2)\) with
\(\Ph(a^2x)\to0\) cannot itself be that nonzero positive physical
spectral covariance. Auxiliary virtual-line data must therefore
be distinguished from the full positive physical two-point
function. The broader role of spectral positivity in Euclidean
reconstruction is established in
Refs.~\cite{OS1973,OS1975}; rational-propagator obstructions are
also discussed in Ref.~\cite{Arici2018}.

\begin{assumption}[Physical OS completion]
\label{ass:OS-completion}
The completed Schwinger functions of the admitted Grassmann-even,
gauge-invariant observable fields obey reflection positivity
\eqref{eq:OS}, Euclidean covariance, permutation symmetry, and
\(S_0=1\). We assume the sufficient regularity condition \(E0'\)
of Ref.~\cite{OS1975}, Section~IV.1, Eq.~(4.1): for one integer
\(r_{\rm OS}\geq1\) and constants \(C_{\rm OS},\beta_{\rm OS}>0\),
\begin{equation}
 \begin{aligned}
  |S_n(f)|&\leq\sigma_n\,\mathfrak p^{\rm OS}_{n r_{\rm OS}}(f),\\
  0<\sigma_n&\leq C_{\rm OS}(n!)^{\beta_{\rm OS}},\qquad n\geq1.
 \end{aligned}
 \label{eq:OS-hierarchy-growth}
\end{equation}
Here \(S_n\) is the full \(n\)-point Schwinger distribution,
not a connected cumulant or a coefficient of the coupling expansion.
The space \(\mathscr S_0(\RR^{4n})\) consists of Schwartz functions
whose derivatives of every order vanish on coincident-point diagonals;
\(\mathfrak p^{\rm OS}_j\) is the seminorm in Ref.~\cite{OS1975},
Section~II, Eq.~(2.1). Use self-adjoint observable components, or
include their adjoints and the corresponding conjugation in reflection.
For a finite multiplet, impose the conditions jointly on all mixed
correlations, not only on each component separately. Clustering is
additionally assumed when vacuum clustering or uniqueness is invoked.
\end{assumption}

These are assumptions on the completed correlation functions,
not consequences of convergence at each perturbative order.
The integer \(n\) in \eqref{eq:OS-hierarchy-growth} counts
observable insertions, not loops, and the displayed constants are
independent of \(n\). A finite bound for each graph does not supply
this hierarchy estimate or reflection positivity of a full correlation
functional. Nor is positivity a coefficientwise requirement: for
real \(g\), the matrix
\[
 \begin{pmatrix}1&g\\g&g^2\end{pmatrix}
 =\begin{pmatrix}1\\g\end{pmatrix}\begin{pmatrix}1&g\end{pmatrix}
 \succeq0
\]
has a first-order truncation with determinant \(-g^2\).
This elementary example illustrates the truncation issue; it is not
itself a Schwinger hierarchy or an OS counterexample.
The applicable bosonic OS reconstruction then provides a positive
physical Hilbert space and tempered Wightman distributions with
the spectrum condition and local commutativity on the admitted
observable algebra \cite{OS1973,OS1975,StreaterWightman}.
In particular, reconstructed observables obey
\[
 [\mathcal O_{\bb}(x),\mathcal O'_{\bb}(y)]=0
 \quad\hbox{for }(x-y)^2<0.
\]
If physical odd fields are included, their reconstruction needs the
corresponding spin, graded reflection and graded permutation rules.
The locality statement then uses the graded commutator, not an
ordinary commutator for two odd fields. The ungraded assumption above
does not supply that fermionic reconstruction.
No such reconstruction is obtained from positive Laplace weights
alone. Clustering, global distributional bounds, and positivity
must concern the completed Schwinger functions; the absence of a
direct external multiplier does not prove them.
\subsection{Joint-source positivity and the transferred scalar interaction}
\label{sec:update-joint-positivity}

A positive Fourier multiplier defines a nonnegative covariance form,
but a physical reflected state is a stronger requirement. The full source
matrix tests this distinction without discarding the mixed correlations. The following results specialize the monograph
and Part~II to the article's flat scalar block
\cite{ChiapponeGHOUQFT2026,ChiapponePartII2026}. The reflection reverses
the continued oriented GT length. Its transport from N to A changes
the chart orientation, not the sign of an OS norm.

\begin{proposition}[A free joint-source GT-reflection obstruction]
\label{prop:update-joint-OS}
Let $C_0=(-\Delta_E+m^2)^{-1}$ on $\RR^4$, $m>0$, and
$P_r=e^{-r\sqrt{-\Delta_E}}$, $r>0$. The positive covariance
matrix \eqref{eq:update-three-source-blocks} is not OS-positive on
the joint algebra in which both fields have independent positive-GT
smearings with the same reflection.
\end{proposition}
\begin{proof}
Choose nonzero nonnegative $u,g\in C_c^\infty$ in the positive
oriented-GT half-space, away from its boundary, and set
$f=(-\Delta_E+m^2)u$. Locality and disjoint reflected supports give
\begin{equation}
 \langle\Theta f,C_0f\rangle_{E,\bb}=0,
 \qquad
 \langle\Theta f,C_0P_rg\rangle_{E,\bb}
       =\langle\Theta u,P_rg\rangle_{E,\bb}>0.
\label{eq:update-joint-OS-witness}
\end{equation}
The strict inequality uses the Poisson kernel
$P_r(x)=3r/[4\pi^2(r^2+|x|^2)^{5/2}]>0$.
In a positive semidefinite sesquilinear form a null vector has zero
pairing with every vector. The displayed nonzero mixed pairing
violates this consequence of Cauchy--Schwarz.
\end{proof}
This excludes the specified simultaneous physical localization of
both source slots; it does not exclude $P_r\chi$ as a Euclidean
nonlocal random variable. A positive depth average preserves the
null first block and the strictly positive mixed pairing, so the same
joint test fails whenever that average is finite. Neither sector
reversal nor a carrier character converts it into a positive norm.

\begin{proposition}[Boundary-only weak-coupling reflection test]
\label{prop:update-boundary-quartic-OS}
Keep the same massive boundary Gaussian and fixed $r>0$. Let
$\zeta\ge0$ be a nonzero smooth compactly supported reflection-even
switch, and define the nonnegative, unsubtracted interaction
\begin{equation}
 V[\psi]=\frac1{4!}\int\dd\mu_{E,\bb}(x)\,
                           \zeta(x)\psi(x)^4,
 \qquad \psi=P_r\chi.
\label{eq:update-bounded-quartic}
\end{equation}
In the normalized measure proportional to $e^{-gV}\dd\mathbb P_{0,\bb}$,
$g\ge0$, a boundary-field OS test is negative for all sufficiently
small $g>0$.
\end{proposition}
\begin{proof}
The transferred field has finite coincidence covariance
$v_r=C_{\psi\psi,r}(0)>0$ and all finite Gaussian moments on the
support of $\zeta$. Hence $V$ is almost surely finite, the measure
normalizes, and the required right derivative at $g=0$ exists by
dominated Gaussian moment bounds. With $u\ge0$ as above, put
$f=(-\Delta_E+m^2)u$ and $Q(g)=\langle\chi(\Theta f)\chi(f)\rangle_g$.
Then $Q(0)=0$. Wick expansion of the differentiated normalized
expectation gives
\begin{equation}
 Q'(0^+)=-\frac{v_r}{2}\int\dd\mu_{E,\bb}(x)\,
       \zeta(x)(P_r\Theta u)(x)(P_ru)(x)<0.
\label{eq:update-quartic-OS-slope}
\end{equation}
There are twelve pairings connecting the two boundary fields to two
of the four interaction fields; the remaining pair gives $v_r$,
so the coefficient is $12/4!=1/2$. The depth transfer has a strictly
positive kernel. Thus the integrand is positive on a nonzero region,
and $Q(g)=gQ'(0^+)+o(g)<0$.
\end{proof}
For a common-depth law with sufficient inverse moments, including
\eqref{eq:update-selector-law} with $B_*>0$, the same differentiated
\emph{globally normalized} expectation averages this strictly negative
slope. For instance $v_r\le1/(32\pi^2r^2)$ and
$\|P_ru\|_\infty\le\|u\|_\infty$ control that average; the higher
moments of $V$ justify differentiation. This is a specific
unsubtracted scalar model, not a no-go theorem for every gauge/BRST
observable completion. Its covered graphs can be UV finite while
its physical boundary reconstruction fails.

\paragraph{Wick ordering changes the witness, not automatically the conclusion}
The tadpole pairing in \eqref{eq:update-quartic-OS-slope} disappears
if the transferred interaction is Wick ordered. A quadratic boundary observable gives the following stronger
fixed-depth test, consistent with the detailed construction in
Ref.~\cite{ChiapponePartII2026}. The argument is established analytically
below and lies outside the pinned Lean inventory.

\begin{proposition}[Fixed-depth Wick-ordered boundary reflection test]
\label{prop:polish-wick-OS}
Keep $m,r>0$, $C_0$, $P_r$, and the switch $\zeta$ of
Proposition~\ref{prop:update-boundary-quartic-OS}. Put
$I_r=(P_rC_0P_r)(0)$, where zero denotes position-space coincidence, and
\begin{equation}
 W_r=\frac1{4!}\int\dd\mu_{E,\bb}(x)\,\zeta(x)
       [\psi(x)^4-6I_r\psi(x)^2+3I_r^2],\quad\psi=P_r\chi.
 \label{eq:polish-wick-potential}
\end{equation}
For every fixed $g\ge0$, the probability measure
$Z_{g,r}^{-1}e^{-gW_r}\dd\mathbb P_{0,\bb}$ exists and has finite polynomial
moments of boundary smearings in the free covariance-form domain. Nevertheless its boundary-field GT-reflection
form is negative for some quadratic observable whenever $g>0$ is
sufficiently small.
\end{proposition}
\begin{proof}
The identity
$y^4-6I_ry^2+3I_r^2=(y^2-3I_r)^2-6I_r^2$ gives
\begin{equation}
 W_r\ge-\frac{I_r^2}{4}\int\zeta\,\dd\mu_{E,\bb},\qquad
 1\le Z_{g,r}\le
 \exp\left[\frac{gI_r^2}{4}\int\zeta\,\dd\mu_{E,\bb}\right].
 \label{eq:polish-wick-measure-bound}
\end{equation}
Here $\mathbb E_0W_r=0$, so Jensen's inequality supplies the lower
partition bound. Smoothing at fixed positive depth gives every Gaussian
moment required below. The lower bound on $W_r$ and Gaussian moment
integrability justify the right derivative at $g=0$ of the normalized
expectation, without a convergent perturbation series.

Choose nonzero real $u,h\ge0$ smooth with compact support strictly in
the positive oriented-GT half-space. Write $B_m=-\Delta_E+m^2$,
$f=B_mu$, and
$\mathcal A=\chi(f)\chi(h)-\langle f,C_0h\rangle_{E,\bb}$.
Locality of $B_m$ and separation of the reflected supports give
\[
 \langle\Theta f,C_0f\rangle_{E,\bb}
 =\langle\Theta f,C_0h\rangle_{E,\bb}
 =\langle\Theta h,C_0f\rangle_{E,\bb}=0.
\]
Wick's rule therefore gives $Q(0):=\mathbb E_0[(\Theta\mathcal A)\mathcal A]=0$.
In the insertion of the Wick-ordered quartic, all four interaction
fields must pair with the four boundary fields. There are $4!$ such
pairings; contractions within the interaction are removed by Wick
ordering. The centering constants in $\mathcal A$ leave fewer than four boundary
fields, so their contractions with this fourth Wick polynomial vanish. Consequently
\begin{equation}
 \begin{aligned}
 Q'(0^+)=-\int\dd\mu_{E,\bb}(x)\,\zeta(x)
   &(P_r\Theta u)(x)(P_ru)(x)\\
   &\times(P_rC_0\Theta h)(x)(P_rC_0h)(x)<0.
 \end{aligned}
 \label{eq:polish-wick-negative-slope}
\end{equation}
Both the Poisson kernel and the massive scalar Green function are
strictly positive away from their harmless integration singularities.
Thus the four smooth factors are positive on the positive-measure
support of $\zeta$. Since $Q(g)=gQ'(0^+)+o(g)$, the assertion follows.
\end{proof}
This is a result for the specified fixed-depth scalar action and its
ordinary boundary observable algebra, not a universal obstruction to
every nonlocal or BRST-reduced theory. It does not define a shared-depth
Wick measure: the upper bound in \eqref{eq:polish-wick-measure-bound}
contains $e^{O(g r^{-4})}$ as $r\downarrow0$, whereas the winding law
only supplies an $e^{-B_*/r}$ small-depth factor. All inverse moments
alone therefore do not justify this additional exponential average.
Failure of that sufficient bound is not by itself a proof that the
actual shared-depth partition function diverges.

\subsection{Positive spectral completions and their ultraviolet cost}
\label{sec:update-spectral-cost}

The scalar obstruction \eqref{eq:spectral-obstruction} is not removed
by putting the damper in a cross entry of an otherwise unspecified
positive matrix. Fix a copy $\bb$. In this subsection $s_{\rm sp}$
is a nonnegative spectral mass squared, distinct from the sector
$s$ and every heat parameter. The atoms
$\delta^{\rm sp}_{\bb,c}(\dd s_{\rm sp})$ are unit measures at $c$,
not coordinate or shell deltas. Suppose
\begin{equation}
 C_{ij}(x)=\int_{[0,\infty)}
              \frac{\nu_{ij,\bb}(\dd s_{\rm sp})}{x+s_{\rm sp}},
 \qquad \boldsymbol\nu_{\bb}\succeq0,
\label{eq:update-matrix-spectral}
\end{equation}
with locally finite measures and
$\int(x_{\rm ref}+s_{\rm sp})^{-1}\nu_{ii,\bb}(\dd s_{\rm sp})<\infty$
for one arbitrary reference $x_{\rm ref}>0$.

\begin{proposition}[Pure-pole cross restriction and diagonal budget]
\label{prop:update-spectral-budget}
If $\nu_{11,\bb}=Z\delta^{\rm sp}_{\bb,m^2}$, $Z>0$, then every
$\nu_{1j,\bb}$ is supported at $m^2$. Thus
$C_{1j}(x)=c_{1j}/(x+m^2)$.
More generally, put $\nu_{c,\bb}=\nu_{12,\bb}$ and
$D_c(x)=\int(x+s_{\rm sp})^{-1}|\nu_{c,\bb}|(\dd s_{\rm sp})$.
Every positive two-channel completion obeys
\begin{equation}
 C_{11}(x)C_{22}(x)\ge D_c(x)^2,\qquad
 C_{11}(x)+C_{22}(x)\ge2D_c(x).
\label{eq:update-spectral-product-budget}
\end{equation}
These bounds do not assume equal diagonals.
\end{proposition}
\begin{proof}
On a bounded Borel set $B$, positivity gives
$|\nu_{1j,\bb}(B)|^2\le\nu_{11,\bb}(B)\nu_{jj,\bb}(B)$, so the cross measure
vanishes off the pole; exhaust the complement by bounded sets.
For the general statement use a dominating trace measure
$\nu_{11,\bb}+\nu_{22,\bb}$, with densities $a,b,c$.
Pointwise positivity gives $|c|^2\le ab$. Cauchy--Schwarz in the
weighted measure $(x+s_{\rm sp})^{-1}(\nu_{11,\bb}+\nu_{22,\bb})$ yields
$D_c\le\sqrt{C_{11}C_{22}}$. The arithmetic--geometric mean
inequality proves the sum bound.
\end{proof}
In particular the undressed pure boundary pole cannot coexist with
a nonconstant damped cross numerator in the same ordinary positive
spectral multiplet. Ghosts and separately declared diagrammatic
kernels are not asserted to have representation
\eqref{eq:update-matrix-spectral}.

The shifted square-root family gives an explicit test. The
discontinuity across its square-root cut determines a signed cross
measure, and the \emph{total variation} of that measure sets the lower
bound on the two positive diagonal covariances. Define
\begin{equation}
 C_c(x)=\frac{e^{-a(\sqrt{x+M^2}-M)}}{x+m^2},
 \qquad a>0,\quad m,M\ge0.
\label{eq:update-cross-family}
\end{equation}

\begin{proposition}[Cross-spectrum classification and its unavoidable tail]
\label{prop:update-cross-classification}
The function \eqref{eq:update-cross-family} has an ordinary real signed
spectral measure with finite weighted total variation precisely when
$m\le M$, or when $m>M$ and
$\sin(a\sqrt{m^2-M^2})=0$. In these cases
\begin{equation}
\begin{aligned}
 \nu_{c,\bb}(\dd s_{\rm sp})
   &=c_*\delta^{\rm sp}_{\bb,m^2}(\dd s_{\rm sp})
                      +h(s_{\rm sp})\dd s_{{\rm sp},\bb},\\
 c_*&=\begin{cases}
 e^{a(M-\sqrt{M^2-m^2})},&m\le M,\\
 e^{aM}\cos(a\sqrt{m^2-M^2}),&m>M,
 \end{cases}\\
 h(s_{\rm sp})&=
 \frac{e^{aM}\sin(a\sqrt{s_{\rm sp}-M^2})}{\pi(m^2-s_{\rm sp})}
                    \mathbf1_{s_{\rm sp}>M^2}.
\end{aligned}
\label{eq:update-cross-measure}
\end{equation}
At an admitted embedded pole the density has its removable value.
At $m=M$ it behaves as
$-ae^{aM}/(\pi\sqrt{s_{\rm sp}-M^2})$, with a separate threshold
atom of weight $e^{aM}$. The total-variation transform satisfies
\begin{equation}
 D_c(x)=\frac{2e^{aM}}{\pi^2}\frac{\log(a^2x)}{x}+O(x^{-1})
       \quad(x\to\infty).
\label{eq:update-cross-cost-tail}
\end{equation}
\end{proposition}
\begin{proof}
Use the principal square root. At $x=-s_{\rm sp}+\ii0$ above threshold,
its value is $\ii\sqrt{s_{\rm sp}-M^2}$. Stieltjes inversion gives
\eqref{eq:update-cross-measure}, except that a generic embedded pole
has the nonzero distribution
$\sin(a\sqrt{m^2-M^2})\PV(m^2-s_{\rm sp})^{-1}$.
That distribution is not a locally finite signed measure.
A vanishing sine removes it. At threshold the expansion of
$e^{-a\sqrt z}/z$ is $z^{-1}-az^{-1/2}+O(1)$, giving the atom and
integrable continuum stated. Below threshold the isolated real residue
is $e^{a(M-\sqrt{M^2-m^2})}$.
To verify sufficiency, subtract $c_* /(z+m^2)$ from the analytic
continuation $C_c(z)$. On the two sides of the cut the remainder
is bounded near an admitted embedded pole because both
one-sided residues equal $c_*$. At the threshold the remainder
is at worst $O(|z+M^2|^{-1/2})$. On a large circle it is
$O(|z|^{-1})$ since the principal square root has nonnegative
real part. Cauchy's formula for this remainder on the slit domain
therefore receives no contribution from the small arcs at the
threshold or embedded pole, or from the large circle. The two lips
give the Stieltjes transform of $h$. It is absolutely convergent
by its local integrability and the bound $|h(s_{\rm sp})|=O(s_{\rm sp}^{-1})$
at infinity. Restoring the subtracted pole gives the stated
measure. Stieltjes inversion uniquely determines the boundary
distribution. In the generic embedded case that distribution
contains the nonzero principal value found above, so it cannot
come from an ordinary signed measure.

At large $k$, its continuous total variation satisfies
\[
 |h(s_{\rm sp})|\dd s_{{\rm sp},\bb}
   =\frac{2e^{aM}|\sin(ak)|}{\pi k}\dd k_{\bb}
      +O(k^{-3})\dd k_{\bb}.
\]
The mean of $|\sin(ak)|$ is $2/\pi$ and its mean-zero part has a
bounded primitive. Hence
$|\nu_c|([0,Q])=2e^{aM}\log(a^2Q)/\pi^2+O(1)$.
To get the transform, multiply it by $x$. For $s_{\rm sp}\le x$,
replacing $x/(x+s_{\rm sp})$ by one has bounded error. The fixed
low-mass interval has finite variation, and on the rest use
$s_{\rm sp}/(x+s_{\rm sp})\le s_{\rm sp}/x$ with
$|h(s_{\rm sp})|\le C/s_{\rm sp}$. For $s_{\rm sp}>x$ the
bound $(x/s_{\rm sp})|h(s_{\rm sp})|\le Cx/s_{\rm sp}^{2}$
also integrates to $O(1)$. Division by $x$ proves
\eqref{eq:update-cross-cost-tail}.
\end{proof}
The signed measure has a positive matrix realization with
$|\nu_c|$ on both diagonals and $\nu_c$ on the cross entries, but its
diagonals are not the undressed reference covariance. This is a
conditional spectral construction, not a particle-mass quantization
law or an A3/A4 necessity result.

\begin{corollary}[Redistributing the spectral weight does not cure this bubble]
\label{cor:update-completed-bubble}
For any positive two-channel spectral completion of an admitted
\eqref{eq:update-cross-family}, the four-dimensional integral
\begin{equation}
 \int_{\RR^4}\dd^4Q_{E,\bb}\,
       C_{11}(|Q_E|^2)C_{22}(|Q_E+P_E|^2)
\label{eq:update-completed-bubble}
\end{equation}
is ultraviolet divergent for fixed $P_E$.
\end{corollary}
\begin{proof}
For sufficiently large $|Q_E|$, $|Q_E+P_E|^2\le2|Q_E|^2$.
A positive Stieltjes covariance obeys
$C_{22}(2x)\ge C_{22}(x)/2$. The integrand is therefore at least
$D_c(|Q_E|^2)^2/2$ by
\eqref{eq:update-spectral-product-budget}. Its radial lower bound
is a positive multiple of $\int^\infty(\log k)^2\dd k/k$.
\end{proof}
This statement concerns precisely the displayed diagonal contractions,
for example a corresponding scalar bubble with nonzero constant
vertices. It does not assume that every proposed gauge interaction
contains this graph. It does show why a damped cross entry cannot
certify all contractions of its unchanged positive Gaussian completion.

\paragraph{A local positive bulk illustrates the same action-level distinction}
On one specified positive normal channel, the coefficient action
\begin{equation}
\begin{aligned}
 S^{\rm bulk}_{E,\bb}
   ={}&\tfrac12\langle\chi,(L+m_b^2)\chi\rangle_{E,\bb}\\
    &+\tfrac\beta2\int_0^\infty\dd r_{\bb}
       [\|\partial_rU\|_{E,\bb}^2+
                       \langle U,(L+M^2)U\rangle_{E,\bb}],\\
 &U(0)=\chi,\qquad\beta>0,
\end{aligned}
\label{eq:update-positive-bulk}
\end{equation}
with $L=-\Delta_E$ and $m_b^2,M^2\ge0$, has a positive boundary/bulk
spectral realization on its admitted trace domains
\cite{ChiapponeGHOUQFT2026,ChiapponePartII2026}. Related Wentzell
boundary dynamics are studied in Ref.~\cite{DappiaggiEtAl2022};
the flat comparison here has its own stated action and parameters.
Every scalar coefficient retains the declared A3 or A4 directed lift
when returned to the corresponding action; $\beta$ is additional
normal-sector data. With $\kappa=\sqrt{x+M^2}$ and
$d(x)=x+m_b^2+\beta\kappa$, completing the quadratic form gives
\begin{equation}
\begin{aligned}
 C_{bb}(x)&=d(x)^{-1},\qquad C_{b,r}(x)=e^{-r\kappa}d(x)^{-1},\\
 C_{r,r'}(x)&=\frac{e^{-(r+r')\kappa}}{d(x)}
       +\frac{e^{-|r-r'|\kappa}-e^{-(r+r')\kappa}}{2\beta\kappa}.
\end{aligned}
\label{eq:update-positive-bulk-blocks}
\end{equation}
The conditional Dirichlet noise in the second term is nonzero for
a positive sampled depth. In particular
$C_{rr}(x)\sim(2\beta\sqrt x)^{-1}$ at every fixed $r>0$,
and $C_{bb}(x)\sim x^{-1}$ on the boundary. Sending $\beta\to\infty$ suppresses noise
but also sends $C_{bb}\to0$ at fixed $x+M^2>0$; maintaining
$C_0=(L+m^2)^{-1}$ would require the different boundary operator
$L+m^2-\beta\sqrt{L+M^2}$. Its normalized deterministic alternative
cannot borrow the positive bulk's physical spectrum or cuts.
The bound mode, when present, and the normal continuum contribute to
the same GT-ordered source contour and physical state sum. Omitting
one while keeping that action's retarded response changes the state.
The full spectral and causal constructions are in
Refs.~\cite{ChiapponeGHOUQFT2026,ChiapponePartII2026}; they are comparisons,
not witnesses for all the UV hypotheses of this article.

\paragraph{The positive comparison uses one common spectral state sum}
The following gap branch exhibits the additional physical weight
explicitly. It supplies a free comparison contour, not a proof that the
exactly constrained source model has the same state space.

\begin{proposition}[Boundary/layer spectrum of the positive normal comparison]
\label{prop:update-positive-bulk-spectrum}
Assume $0<m_b^2<M^2$ in \eqref{eq:update-positive-bulk}, and let
$r_i\ge0$ label finitely many sampled layers, including the boundary
$r_b=0$. Define
\begin{equation}
\begin{aligned}
 \kappa_*&=\frac{\sqrt{\beta^2+4(M^2-m_b^2)}-\beta}{2},\\
 s_*&=M^2-\kappa_*^2,\qquad
 Z_*=(1+\beta/(2\kappa_*))^{-1},\qquad h_i=e^{-r_i\kappa_*},\\
 k_n&=\sqrt{s_{\rm sp}-M^2},\qquad t_s=m_b^2-s_{\rm sp},\\
 f_i(s_{\rm sp})&=\cos(k_n r_i)+\frac{t_s}{\beta k_n}\sin(k_n r_i),\\
 \rho_{b,\bb}(s_{\rm sp})&=\frac{\beta k_n}{\pi(t_s^2+\beta^2 k_n^2)}.
\end{aligned}
\label{eq:update-positive-spectrum-data}
\end{equation}
On the admitted layer-source forms the positive matrix measure
\begin{equation}
\begin{aligned}
 \nu_{ij,\bb}(\dd s_{\rm sp})
   ={}&Z_*h_i h_j\delta^{\rm sp}_{\bb,s_*}(\dd s_{\rm sp})\\
     &+\rho_{b,\bb}(s_{\rm sp})f_i(s_{\rm sp})f_j(s_{\rm sp})
         \mathbf1_{s_{\rm sp}>M^2}\dd s_{{\rm sp},\bb}
\end{aligned}
\label{eq:update-positive-spectrum}
\end{equation}
represents every block in \eqref{eq:update-positive-bulk-blocks}.
The boundary weights satisfy $Z_*+\int_{M^2}^\infty\rho_{b,\bb}(s_{\rm sp})\dd s_{{\rm sp},\bb}=1$.
\end{proposition}
\begin{proof}
The equation for a pole below threshold is
$m_b^2-s+\beta\sqrt{M^2-s}=0$; it decreases strictly and has the
unique root $s_*\in(m_b^2,M^2)$ and residue $Z_*>0$.
At $x=-s_{\rm sp}+\ii0$ the square root is $\ii k_n$, so the
boundary spectral density is
$-\pi^{-1}\operatorname{Im}(t_s+\ii\beta k_n)^{-1}=\rho_{b,\bb}$.
The normal eigenfunction obeys
$-\beta u'(0)+(m_b^2-s_{\rm sp})u(0)=0$, giving $u(r_i)=f_i u(0)$.
Taking the boundary values of the complete layer Green function,
including its Dirichlet term, gives the outer product $\rho_{b,\bb} f_i f_j$;
the pole gives $Z_*h_i h_j$. Both matrices are positive.
The resolvent of the nonnegative boundary/bulk quadratic form therefore
has this measure. Equivalently its branch jump, isolated residue, and
resolvent decay determine the displayed Stieltjes transforms.
Finally $xC_{bb}(x)\to1$; monotone convergence of its positive
spectral integral gives the boundary normalization.
\end{proof}
Point-layer evaluation is a trace Green form on admitted sources,
not a bounded evaluation operator on an unrestricted bulk $L^2$ space.
The positive normal kinetic sign and its domain remain part of this
comparison action, not consequences of the vectorial pole choice.

For example, the two canonical directed orderings built from this
same measure are
\begin{equation}
 \mathcal G_{ij,\alpha}(w)
    =\int\nu_{ij,\bb}(\dd s_{\rm sp})
        \frac{\alpha\je}{w-s_{\rm sp}+\alpha0\je}.
\label{eq:update-positive-spectral-orderers}
\end{equation}
For a smooth compactly supported spectral test $f$, their sum gives
\[
\begin{aligned}
 &\int\dd w_{\bb}\,f(w)[\mathcal G_{ij,+}(w)+\mathcal G_{ij,-}(w)]\\
 &\hspace{15mm}=2\pi\int\nu_{ij,\bb}(\dd s_{\rm sp})f(s_{\rm sp}).
\end{aligned}
\]
Together with the two oriented-GT frequency selectors this gives
\eqref{eq:update-GT-contour-identity}, including the atom and continuum.
It is not an independently imposed undressed cut rule.

At fixed spatial momentum the boundary retarded coefficient is
\begin{equation}
 D_{R,bb}(z,\mathbf k)
   =\frac1{\mathbf k^2+m_b^2-z^2
                   +\beta\sqrt{\mathbf k^2+M^2-z^2}},
 \qquad\operatorname{Im}z>0,
\label{eq:update-positive-retarded}
\end{equation}
with the decaying normal branch. In the vacuum its symmetrized
boundary spectral distribution is, for $w=p_\gamma^2-\mathbf k^2$,
\begin{equation}
\begin{aligned}
 N_{bb}(p_\gamma,\mathbf k)
   ={}&\pi Z_*\delta^{\rm sh}_{\bb}(w-s_*)\\
     &+|D_{R,bb}(p_\gamma+\ii0,\mathbf k)|^2
            \beta\sqrt{w-M^2}\,\mathbf1_{w>M^2}.
\end{aligned}
\label{eq:update-positive-noise}
\end{equation}
The atom cannot be obtained by retaining only continuum forcing.
Shell Jacobians are understood in the declared Fourier measure.
A retarded approach $z=p_\gamma+\delta\je$ has a quadratic
infinitesimal proportional to $p_\gamma\delta\je$, unlike the
constant directed Feynman boundary height. Both follow their own
GT-ordering/response definitions. This model realizes free physical
positivity and a common contour while its diagonal loops retain the
power-law ultraviolet behavior already displayed; its properties
cannot be assigned to the constrained transfer by omission of terms.


\subsection{External attachments and interacting LSZ residues}
\label{sec:LSZ}

Under the external-leg placement convention, the LSZ attachment and
amputation operators carry no additional factor of \(\Ph\). The
interacting two-point function still depends on the completed dynamics,
so equality with its undeformed counterpart requires a separate argument.

The following single-sheet transport describes a resolved scalar
two-point channel whose completed carrier coefficient belongs to
the directed subalgebra. Equivalently, its two complete character
coefficients obey the reflection relation specified below. For a
general mixed carrier two-point function, pole extraction is performed
in both characters before reconstruction; choosing one sheet does
not impose that reflection condition on the other character.

\paragraph{LSZ limits follow the physical endpoint, not the printed slot}
Assume a scalar physical channel with positive residue
$Z^{\rm phys}_{\bb}>0$ and the requisite scattering limits. On that
channel, a positive-oriented wave packet $f_{\bb}$ has phase
$e^{-\ii\kappa_{\mathbf p}\gamma_{\bb}+\ii\mathbf p\cdot\mathbf x}$.
Its Klein--Gordon pairing uses the clock-forward oriented normal:
\begin{equation}
 (f,g)_{\gamma,\bb}
 =\ii\int\dd\Sigma_{\bb}\,
      f^*\overleftrightarrow{\partial}_{\gamma_{\bb}}g,
 \qquad \partial_{\gamma_{\bb}}=\sigma_s\partial_{X_{0,\bb}}.
 \label{eq:fw-lsz-surface}
\end{equation}
For the chosen covariant packet normalization this gives the annihilation
limits
\begin{equation}
\begin{aligned}
 a_{{\rm in},\bb}(f)
  &=(Z^{\rm phys}_{\bb})^{-1/2}
      \lim_{X_{0,\bb}\to-\sigma_s\infty}
                 (f_{\bb},\varphi^{\rm phys}_{\bb})_{\gamma,\bb},\\
 a_{{\rm out},\bb}(f)
  &=(Z^{\rm phys}_{\bb})^{-1/2}
      \lim_{X_{0,\bb}\to+\sigma_s\infty}
                 (f_{\bb},\varphi^{\rm phys}_{\bb})_{\gamma,\bb}.
\end{aligned}
\label{eq:fw-lsz-limits}
\end{equation}
Creation limits are the admitted physical adjoints. These formulas are
not obtained from the tagged virtual numerator alone. Spinor channels use
their Dirac surface pairing and the correspondingly transported spin
frame. The endpoint assignment is always Table~\ref{tab:fw-endpoints}:
physical final A particles take outgoing limits and outgoing phase space
even when written in an ascending-GT ket. An unqualified rule ``all bras
are outgoing'' is not valid for the mixed boundary kernel.

Vacuum-to-one-particle overlaps used to form a residue Gram matrix below
are legitimate matrix elements of a field operator. They are not
closed vacuum-creation S-matrix elements, so the mixed creation rule
\eqref{eq:fw-creation-slots} does not delete those overlaps.

For a chosen completed-theory sheet, let
\(\mathscr R^{(2),\rm sheet}_{\eta,\bb}(z)\) be the
ordinary meromorphic two-point coefficient with its propagator
numerator removed. Here \(m_{{\rm phys},\bb}\) is the
momentum-dimension shell parameter of Section~\ref{sec:framework};
with explicit units it is the physical mass times \(\cGT\). Assume that a neighbourhood of the physical
mass contains precisely one simple pole:
\begin{equation}
 \mathscr R^{(2),\rm sheet}_{\eta,\bb}(z)
  =\frac{Z_{\bb}}{z-m_{\rm phys,\bb}^2}
            +H^{\rm sheet}_{\eta,\bb}(z),
 \qquad 0<|Z_{\bb}|<\infty,
 \label{eq:full-pole}
\end{equation}
where \(H^{\rm sheet}_{\eta,\bb}\) is holomorphic across the
pole. This is a sheet statement before a boundary distribution
is taken. Set
\begin{align}
 \mathscr R^{(2),\rm dir}_{\eta,\bb}
   &:=\embed\circ
          \mathscr R^{(2),\rm sheet}_{\eta,\bb}
                   \circ\embed^{-1},\nonumber\\
 \mathcal R^{(2),\rm coeff}_{\eta,\bb}(p)
   &:=\lim_{\varepsilon\downarrow0}
       \embed^{-1}\!\left[
       \mathscr R^{(2),\rm dir}_{\eta,\bb}
              (p^2+\ii\eta\epsvec)\right].
 \label{eq:LSZ-explicit-coefficient-extraction}
% Monograph source: prop:LSZ; eq:LSZ-directed-two-point-coefficient.
\end{align}
Under this directed two-point hypothesis, and only after a compatible
source/endpoint channel has been specified, the oriented boundary
template is
\begin{align}
 \vec G^{(2)}_{\eta,\bb}(p)
 &=\eta\ii\widehat u_{p,\bb}\,
       \embed\!\left(
           \mathcal R^{(2),\rm coeff}_{\eta,\bb}(p)\right)
       \nonumber\\
 &=\embed\!\left(\sigma_\epsilon\ii
           \mathcal R^{(2),\rm coeff}_{\eta,\bb}(p)\right)
       \quad\text{when }\widehat u_{p,\bb}=\sigma_s\ee.
 \label{eq:LSZ-physical-directed-kernel}
\end{align}
The first line is only an oriented component or raywise pole template.
It is not, by itself, the covariance of an identical real scalar:
Proposition~\ref{prop:source-exchange} still applies. For that field,
the even fixed-carrier specialization passes this necessary test but
requires its own completed source pairing; it must not be substituted
silently for the running-unit model. An endpoint-labelled realization of the
first line must supply its reversed component and verify reciprocity
for the complete source matrix. A null-set convention for the undefined running unit does not
supply that missing matrix or an external zero-carrier limit.
In particular, its directed coefficient pole part is
\(\embed(Z_{\bb})[p^2-m_{\rm phys,\bb}^2
+\ii\eta\epsvec]^{-1}\). The complete coefficient is
embedded before multiplication by \(\eta\ii\widehat u_{p,\bb}\).

For completeness, write \(r_p=\sgn(p_\parallel)\) at a
nonzero mode and let
\(\mathscr R^{(2),\rm sheet}_{\kappa,\eta,\bb}\) be the
two independently determined, numerator-stripped character
coefficients of a resolved scalar channel. At finite
\(\varepsilon>0\), the reconstruction within this oriented,
numerator-factorized ansatz is
\begin{equation}
 \vec G^{(2)}_{\eta,\bb}(p;\varepsilon)
  =\sum_{\kappa=\pm1}\mathsf P_\kappa\,
     \kappa\eta\ii r_p\,
     \mathscr R^{(2),\rm sheet}_{\kappa,\eta,\bb}
                       (p^2+\ii\kappa\eta\varepsilon).
 \label{eq:LSZ-two-character-reconstruction}
% Monograph source: eq:general-directed-graph-character-reconstruction; prop:LSZ.
\end{equation}
Each character has its own Laurent data on its declared pole
neighbourhood. Equality of their pole locations and the required
relation between their residues are properties of the physical
completion. If
\begin{equation}
 \mathscr R^{(2),\rm sheet}_{-,\eta,\bb}(\overline z)
   =\overline{
      \mathscr R^{(2),\rm sheet}_{+,\eta,\bb}(z)},
 \label{eq:LSZ-character-reflection}
\end{equation}
then \eqref{eq:LSZ-two-character-reconstruction} reduces to the
single-sheet directed formula above. Otherwise both character
coefficients are retained. The common physical boundary limit is
taken only after the admitted reconstruction and integrations.

For an arbitrary completed two-point source matrix, the genuinely
general formula is instead
\begin{equation}
 \vec C^{AB}_{\bb}(p)
   =\sum_{\kappa=\pm1}\mathsf P_\kappa C^{AB}_{\kappa,\bb}(p),
 \label{eq:LSZ-general-source-matrix}
\end{equation}
where each coefficient includes its own full numerator and obeys
\(C^{AB}_{\kappa,\bb}(p)
=(-1)^{\epsilon_A\epsilon_B}C^{BA}_{\kappa,\bb}(-p)\).
Neither scalar dependence on \(p^2\) nor extraction of a common
running sign is automatic. LSZ reduction additionally requires that
the pole belong to the specified physical source channel and admit
the stated spectral and asymptotic limits. A scalar meromorphic
coefficient alone does not construct those limits.

In the selected scalar source channel, put
\(M^2=m_{\rm phys,\bb}^2\) and write
\(h_{\bb}=H^{\rm sheet}_{\eta,\bb}\). The Laurent form
\eqref{eq:full-pole} has the local reciprocal
\begin{equation}
 \begin{aligned}
  \mathfrak F_{\bb}(z)
    &=\frac{z-M^2}{Z_{\bb}+(z-M^2)h_{\bb}(z)},\\
  \mathfrak F_{\bb}(M^2)&=0,\qquad
  \mathfrak F_{\bb}'(M^2)=Z_{\bb}^{-1}.
 \end{aligned}
 \label{eq:local-Laurent-inverse}
\end{equation}
Because \(Z_{\bb}\ne0\), the denominator stays nonzero in a
sufficiently small neighbourhood. Thus this reciprocal extends
holomorphically through the pole with a simple zero; no global
invertibility is asserted. Define its correction relative to the
reference inverse \(z-m_{0,\bb}^2\) by
\begin{equation}
 \mathfrak F_{\bb}(z)
       :=z-m_{0,\bb}^2-\Sigma^{\rm sheet}_{\Ph,\bb}(z),
 \qquad
 \mathfrak F^{\rm dir}_{\bb}
       :=\embed\circ\mathfrak F_{\bb}\circ\embed^{-1}.
 \label{eq:Dyson-directed-inverse-function}
% Monograph source: eq:Dyson-completed-inverse-coefficient; eq:Dyson-local-pole-expansion.
\end{equation}
Here \(\Sigma^{\rm sheet}_{\Ph,\bb}\) includes any tree-level
kinetic modification as well as loop corrections. It is a loop-only
self-energy when the full tree-level inverse in that source channel
is \(z-m_{0,\bb}^2\).

For a general admitted tree inverse \(K_{{\rm tree},\bb}(z)\),
write instead
\begin{equation}
 \begin{aligned}
  \mathfrak F_{\bb}(z)
     &=K_{{\rm tree},\bb}(z)-\Sigma_{{\rm loop},\bb}(z),\\
  \Sigma^{\rm sheet}_{\Ph,\bb}(z)
     &=z-m_{0,\bb}^2-K_{{\rm tree},\bb}(z)
         +\Sigma_{{\rm loop},\bb}(z).
 \end{aligned}
 \label{eq:tree-versus-loop-inverse}
\end{equation}
Both decompositions refer to the same inverse, not to different
choices of the pole prescription.
The inverse in the full Dyson covariance is
\(\mathfrak F^{\rm dir}_{\bb}
(p^2+\ii\eta\epsvec)^{-1}\), taken in the directed
subalgebra. On the local pole neighbourhood established in
\eqref{eq:local-Laurent-inverse}, the finite matching equations are
\begin{equation}
 \begin{gathered}
 m_{\rm phys,\bb}^2-m_{0,\bb}^2
     -\Sigma^{\rm sheet}_{\Ph,\bb}(m_{\rm phys,\bb}^2)=0,\\
 Z_{\bb}^{-1}
     =1-(\Sigma^{\rm sheet}_{\Ph,\bb})'
                               (m_{\rm phys,\bb}^2).
 \end{gathered}
 \label{eq:finite-pole-matching}
\end{equation}
Equivalently, for the same pole \(M^2\),
\begin{equation}
 \begin{aligned}
  K_{{\rm tree},\bb}(M^2)-\Sigma_{{\rm loop},\bb}(M^2)&=0,\\
  Z_{\bb}^{-1}
    &=K_{{\rm tree},\bb}'(M^2)-\Sigma_{{\rm loop},\bb}'(M^2).
 \end{aligned}
 \label{eq:general-tree-pole-matching}
\end{equation}
For example, a specified scalar tree coefficient
\(f(z)/(z-m_0^2)\), with \(f\) holomorphic and nonzero near
\(m_0^2\), has inverse \((z-m_0^2)/f(z)\) and tree residue
\(f(m_0^2)\). Calling the entire correction in
\eqref{eq:Dyson-directed-inverse-function} ``loop only'' would
miss this contribution. This algebraic example does not identify
that coefficient with an independently completed physical source field.
In a general mixed carrier channel the inverse coefficient and
these pole-matching equations are defined for each character
separately. Their one-sheet form here uses the preceding directed
two-point hypothesis.
The pole-matching quantities in
\eqref{eq:finite-pole-matching} and \eqref{eq:general-tree-pole-matching}
are coefficient Laurent data on the selected sheet. Within an admitted
channel they determine the pole part of
\eqref{eq:LSZ-physical-directed-kernel}; they neither define an independent
scalar replacement nor prove source reciprocity.

Internally damped loops and compatible vertices can change
\(\Sigma_\Ph\). A unit-residue convention requires an admissible
finite field normalization that preserves the physical source
pairing and adjoint. In a positive Hermitian scalar channel the
isolated spectral weight is \(Z_{\bb}>0\), and rescaling by
\(Z_{\bb}^{-1/2}\) supplies it. The weaker algebraic condition
\(0<|Z_{\bb}|<\infty\) alone does not establish that physical
normalization \cite{LSZ1955}; neither does \(\Ph(0)=1\).

\begin{proposition}[Positive residues and field--source duality]
\label{prop:positive-residue-matrix}
Consider a finite multiplet \(\varphi_A\) of admitted physical scalar
fields at an isolated stable one-particle mass. Use the adjoint-paired
two-point matrix, with endpoints \(\varphi_A\) and
\(\varphi_B^\dagger\). After removing the conventional propagator
and state-normalization factors, suppose its residue comes from the
positive physical one-particle resolution. Then
\begin{equation}
 Z^{\rm phys}_{AB}
       =\sum_\alpha v_{A\alpha}\overline{v_{B\alpha}},
 \qquad Z^{\rm phys}\succeq0,
 \label{eq:positive-pole-Gram-matrix}
\end{equation}
with finite or convergent overlap sums. A constant field-basis map
\(\varphi'=T\varphi\), with
\((\varphi')^\dagger=\varphi^\dagger T^\dagger\), gives
\(Z'=T Z^{\rm phys}T^\dagger\). An invertible such map can
normalize the residue to the identity on its range, but cannot make
a singular residue the identity on the whole field space.
\end{proposition}

\begin{proof}
The overlaps \(v_{A\alpha}\) are the vacuum-to-one-particle
matrix elements in the chosen normalization; wave packets may be
used before removing common momentum distributions. For every finite
vector \(c\),
\[
 c^\dagger Z^{\rm phys}c
       =\sum_\alpha\left|\sum_A\overline{c_A}v_{A\alpha}\right|^2
       \geq0.
\]
The field map acts on the two endpoints and gives the asserted
congruence. Write the finite-dimensional spectral decomposition as
\(Z^{\rm phys}=\sum_{\lambda>0}\lambda P_\lambda\), and set
\(P_0=P_{\ker Z^{\rm phys}}\) and
\(P_{\mathcal R}=I-P_0\). Define
\begin{equation}
 \begin{aligned}
  T_{\mathcal R}&=\sum_{\lambda>0}\lambda^{-1/2}P_\lambda,\\
  T_{\rm inv}&=T_{\mathcal R}+P_0,\\
  T_{\rm inv}Z^{\rm phys}T_{\rm inv}^\dagger&=P_{\mathcal R}.
 \end{aligned}
 \label{eq:invertible-residue-normalization}
\end{equation}
The last identity also holds with \(T_{\mathcal R}\), but only
\(T_{\rm inv}\) retains an invertible map on the entire field space.
Invertible congruence preserves rank, so the right side can be the
full identity only when \(Z^{\rm phys}\) is nonsingular.
\end{proof}

For this commuting field block, the dual source change is distinct.
With an invertible \(T\) and the Hermitian pairing it is
\begin{equation}
 \begin{aligned}
  J'&=(T^{-1})^\dagger J,\\
  J'^\dagger\varphi'+\varphi'^\dagger J'
    &=J^\dagger\varphi+\varphi^\dagger J.
 \end{aligned}
 \label{eq:physical-field-source-duality}
\end{equation}
A real bilinear pairing uses the inverse transpose instead. Thus
\(Z\mapsto TZT^\dagger\) describes the field change with its
compatible dual sources, not an unspecified reparametrization of \(J\).
A null residue direction has no overlap with this one-particle mass;
it need not be a zero field or have vanishing correlations elsewhere.
For a mixed field matrix, the scalar inverse formulas above apply only
after an admitted channel reduction; entrywise reciprocals are not a
matrix inverse.

This test presupposes a positive physical channel; it does not construct
one from a virtual kernel or impose positivity on raw carrier characters
with their ordering phases. Its adjoint is the physical \(\dagger\),
not the carrier \(\ddagger\) of
Proposition~\ref{prop:carrier-reality-adjoint}.

The massive internal shell samples the directed numerator boundary
\(\lim_{\varepsilon\downarrow0}
\Ph_{\ee}(-a^2[m^2+\ii\eta\epsvec])\).
Its ordinary coefficient is
\(\Ph(-a^2[m^2+\ii\eta0])\) only after extraction.
It need not equal \(\Ph(0)\), and its two sides need not agree.

Even tree exchange graphs can depend on the kernel through their
internal lines and completed vertices. No blanket equality of
tree amplitudes or interacting residues with the undeformed theory
is asserted.

\subsection{GT ordering, nonlocal transfer, and the closed contour}
\label{sec:update-GT-contour}

GT ordering is primitive. On the fixed support set
$\gamma_{\bb}=\sigma_sG_{0\bb}=\sigma_sX_0$, an oriented GT
\emph{length}. Its increasing order is $X_0^+$ on N and $X_0^-$ on A.
Clock calibration may identify $\gamma_{\bb}=\cGT t$, but does not
supply a second ordering law. If an admitted self-adjoint
geo-momentum coefficient $Y_{\gamma,\bb}$ generates this flow,
$U(\Delta\gamma)=\exp[-\ii\Delta\gamma Y_{\gamma,\bb}]$ in the
chosen coefficient representation. Existence and its domain remain
physical assumptions, not consequences of the Fourier notation.

For bosonic fields in one physical state, let
$W^>_{\bb}(x,y)=\langle\chi_{\bb}(x)\chi_{\bb}(y)\rangle$ and
$W^<_{\bb}(x,y)=\langle\chi_{\bb}(y)\chi_{\bb}(x)\rangle$. With no extra factor of
$\ii$ in a two-point function, write $W^\gtrless=W^\gtrless_{\bb}$
only within the following fixed-copy display. Its signs refer to the two
$\gamma_{\bb}$ contour branches, not uniformly to $X_0^+$ and $X_0^-$
in both sectors. The two orderers obey
\begin{equation}
\begin{aligned}
 G^{++}&=\theta(\gamma_x-\gamma_y)W^>
              +\theta(\gamma_y-\gamma_x)W^<,\\
 G^{--}&=\theta(\gamma_y-\gamma_x)W^>
              +\theta(\gamma_x-\gamma_y)W^<,\\
 G^{-+}&=W^>,\qquad G^{+-}=W^<,\\
 G^{++}+G^{--}&=G^{-+}+G^{+-}.
\end{aligned}
\label{eq:update-GT-contour-identity}
\end{equation}
Use the equal-GT convention of Section~\ref{sec:fw-endpoints} and
common smearing in all four terms. Both endpoints carry $\bb$ (or the full copy identifier when
needed). The contour signs are not new sector labels. A trace-one positive state operator $\widehat\varrho_{{\rm st},\bb}$
and unitary GT evolution give
\[
 Z_{\bb}[J_+,J_-]
  =\operatorname{Tr}(U_{\bb}[J_+]\widehat\varrho_{{\rm st},\bb}U_{\bb}[J_-]^\dagger),
 \qquad Z_{\bb}[J,J]=1.
\]
The state operator $\widehat\varrho_{{\rm st},\bb}$ is not the heat
measure $\rho_{\bb}$. The familiar closed-contour normalization \cite{GaoGloriosoLiu2020}
is used here with the declared GT generator, not a silently substituted
clock-time Hamiltonian.

\begin{proposition}[Transfer does not in general commute with GT ordering]
\label{prop:update-GT-ordering-transfer}
Suppose an actual Lorentzian smeared operator is defined by
$\psi_{\bb}(x)=\int\dd\mu_{\bb}(u)K_{\bb}(x,u)\chi_{\bb}(u)$ on an admitted
common domain. For a real scalar kernel, the difference between its
ordered covariance and transferring the already ordered covariance is
\begin{equation}
\begin{aligned}
 &\langle X_\gamma^+\{\psi_{\bb}(x)\psi_{\bb}(y)\}\rangle
          -(K\otimes K)G^{++}_{\chi_{\bb}}(x,y)\\
 &=\int\dd\mu_{\bb}(u)\dd\mu_{\bb}(v)\,
       K_{\bb}(x,u)K_{\bb}(y,v)\\
 &\quad\times[\theta(\gamma_x-\gamma_y)-\theta(\gamma_u-\gamma_v)]
                       \langle[\chi_{\bb}(u),\chi_{\bb}(v)]\rangle.
\end{aligned}
\label{eq:update-GT-ordering-correction}
\end{equation}
\end{proposition}
\begin{proof}
Expand each ordered product into $W^>$ and $W^<$, pass the admitted
smearing through those distributions, and subtract. Their difference
is the commutator shown. No field equation or change of pole side
is used.
\end{proof}
A transfer acting within each GT slice without GT derivatives commutes
with the step weights. Carrier-derivative transfers can instead produce
contact terms even though their kernels are supported at equal GT.
An isotropic Euclidean multiplier does not by itself define an
admissible Lorentzian transfer. The distinction between diagrammatic and
operator-ordered nonlocal propagators is also discussed in
Ref.~\cite{BrisceseEtAl2024}; results there for entire form factors
are not applied as though the square-root multiplier were entire.

For an undressed free scalar positive channel with unit residue, the
signed-carrier Fourier convention gives
\begin{equation}
 \widetilde W^>_{\bb}(p)=2\pi\delta^{\rm sh}_{\bb}(p^2-m^2)
                             \theta(\sigma_s p_\parallel),\qquad
 \widetilde W^<_{\bb}(p)=\widetilde W^>_{\bb}(-p).
 \label{eq:fw-wightman-sheet}
\end{equation}
The physical spectral condition is positivity of the oriented coefficient
$p_\gamma=\sigma_s p_\parallel$, not positivity of $p_\parallel$ in both
sectors. An interacting spectral measure and residue replace this free
example when needed. The sum of the two free sheets is the usual complete
shell distribution, consistently with
\eqref{eq:update-GT-contour-identity}. A branch-routed cut uses
\eqref{eq:fw-cut-routing}; it must not silently change this state spectrum.

\paragraph{The vectorial boundary value in the contour test}
Retain $\je=\ii\ee$ and $\eta=\sigma_s\sigma_\epsilon$ before inversion:
\begin{equation}
\begin{aligned}
 (D+\eta0\je)^{-1}
    &=\PV D^{-1}-\pi\eta\je\delta^{\rm sh}_{\bb}(D),\\
 \eta\je(D+\eta0\je)^{-1}
    &=\pi\delta^{\rm sh}_{\bb}(D)+\eta\je\PV D^{-1}.
\end{aligned}
\label{eq:update-vectorial-canonical-split}
\end{equation}
The second expression is the positive-CCR canonical scalar convention
used in this contour comparison. It does not replace the tagged
running numerator or either carrier character of
Section~\ref{sec:carrier}. If $\alpha=\pm1$ denotes the two oriented
$\gamma$ orderings, their fixed-sector branch choice is
$\sigma_\epsilon=\alpha\sigma_s$, so $\eta=\alpha$.

For the canonical stable-half profile, $a,m>0$, set
$w=k_\parallel^2-\mathbf k^2$ and consider on $w>0$
\begin{equation}
 \mathcal D_{a,\alpha}(w)
  =\lim_{\varepsilon\downarrow0}
   \frac{\alpha\je\,e^{-a\sqrt{-(w+\alpha\varepsilon\je)}}}
        {w-m^2+\alpha\varepsilon\je}.
\label{eq:update-directed-virtual-orderers}
\end{equation}
The square root is taken in the same directed line as the inverse.
Since $\sqrt{-(w+\alpha0\je)}=-\alpha\je\sqrt w$, multiplication
of the smooth numerator near $w=m^2>0$ gives
\begin{equation}
\begin{aligned}
 \mathcal D_{a,+}+\mathcal D_{a,-}
   ={}&-2\sin(a\sqrt w)\PV\frac1{w-m^2}\\
     &+2\pi\cos(am)\delta^{\rm sh}_{\bb}(w-m^2).
\end{aligned}
\label{eq:update-directed-contour-sum}
\end{equation}
If both were identified with physical orderings while the Wightman
sum remained the undressed $2\pi\delta^{\rm sh}_{\bb}(w-m^2)$,
the required sum of corrections would be
\begin{equation}
 \mathcal E_a(w)=2\sin(a\sqrt w)\PV\frac1{w-m^2}
       +2\pi[1-\cos(am)]\delta^{\rm sh}_{\bb}(w-m^2).
\label{eq:update-GT-contour-defect}
\end{equation}
This is a regular-shell distributional comparison; the branch point
and critical massless shells need separate domains. The identity fixes
a sum, not a unique correction or a positive state. Corrections
vanishing with the interaction coupling cannot repair a defect already
present at zero coupling. A virtual diagrammatic line need not be a
physical ordered covariance, but its relation to the physical contour
must then follow from the same completed action and source maps.

\paragraph{Geometric variation before the boundary limit}
Let the regular scalar shell $D_y$ and unit
$j_{\bb,y}=\ii e_{\bb,y}$ vary differentiably, with the discrete
signs fixed. Compare the varying units in a common algebra
using the declared carrier transport. At finite height put
$R_\varepsilon=(D_y+\eta\varepsilon j_{\bb,y})^{-1}$.
Differentiation of the inverse gives
\begin{equation}
 \delta R_\varepsilon
    =-R_\varepsilon(\delta D_y+\eta\varepsilon\delta j_{\bb,y})
                  R_\varepsilon.
\label{eq:update-vectorial-inverse-variation}
\end{equation}
For the canonical covariance $\mathcal C_\varepsilon=\eta jR_\varepsilon$,
on compact regular-shell charts with fixed momentum density,
\begin{equation}
\begin{aligned}
 \dot{\mathcal C}_{\rm bv}
  ={}&\eta\dot j\,\PV D^{-1}
      -\eta j\dot D\,\operatorname{Fp}D^{-2}\\
     &+\pi\dot D\,(\delta^{\rm sh}_{\bb})'(D).
\end{aligned}
\label{eq:update-vectorial-boundary-response}
\end{equation}
The differentiation uses a regular local shell coordinate.
With this convention
$\operatorname{Fp}D^{-2}=-\partial_D(\PV D^{-1})$ is the
derivative distribution determined by
\eqref{eq:update-vectorial-canonical-split}.
A running numerator and a varying integration density have their own
variations. Noncommuting second variations are taken before the limit,
not by multiplying separated deltas and principal values.
These qualifications matter for gauge/metric insertions and do not
alter the fixed-background Euclidean convergence estimate.


\paragraph{Coincident shells at a common boundary height}
Work on a compact chart of a regular shell, set \(y=D(k)\), and
include the smooth coarea Jacobian in each test. After the directed
coefficient extraction in \eqref{eq:scalar-line-coefficient}, the
two algebraically distinct quadratic products satisfy
\[
\begin{aligned}
 (y+\ii\eta\varepsilon)^{-2}
 &\longrightarrow \operatorname{Fp}y^{-2}
                   +\ii\eta\pi\delta'(y),\\
 [(y+\ii\eta\varepsilon)(y-\ii\eta\varepsilon)]^{-1}
 &=\frac{\pi}{\varepsilon}\delta(y)
                  +\operatorname{Fp}y^{-2}+o(1)
\end{aligned}
\]
in distributions. Here \(\operatorname{Fp}y^{-2}
=-\partial_y\PV(1/y)\). Differentiating the simple boundary gives
the first identity. For the second, test against \(\phi\), take
its even part \(\phi_{\rm ev}\) in \(y\), and remove the exactly
integrable constant term \(\pi\phi(0)/\varepsilon\).
The remaining integral tends by domination to
\(\int[\phi_{\rm ev}(y)-\phi(0)]y^{-2}\dd y\), which is the
stated finite part.

This also tests contacts whose coefficients depend on the height.
Suppose their complete ordinary numerator is
\(n_\varepsilon=n_0+\varepsilon n_1+o(\varepsilon)\) locally
uniformly, with smooth \(n_0,n_1\). The opposite-side product has
a distributional limit exactly when \(n_0\) vanishes on the shell.
In that case it tends to
\[
 n_0\operatorname{Fp}y^{-2}+\pi n_1|_{y=0}\,\delta(y).
\]
The remainder vanishes since the positive kernel has mass
\(\pi/\varepsilon\). A single zero of \(n_0\) generally leaves a
principal value. Where its normal derivative is nonzero, the local
absolute integral grows logarithmically. A double zero gives a locally
uniform absolute bound for this pairing. These are shell tests for a
completed numerator on one routed support, independent of its behavior
at large momentum.
\subsection{Operational cuts and analytic discontinuities}
\label{sec:cut-disc}

The operational cut is prescribed for the entire oriented virtual
line. Its interpretation as a physical completeness projector requires
the admitted source/endpoint channel of Section~\ref{sec:LSZ}. On that
normalized cut-state interface, with the branch-routed covector
$k=\sigma_\epsilon p^{\rm ph}$ of \eqref{eq:fw-cut-routing},
\begin{equation}
\begin{aligned}
 &\Cut_{{\rm phys},\bb}\!\left[
  \eta\ii\widehat u_{k,\bb}\,
  \frac{\Ph_{\ee}(-a^2[k^2+\ii\eta\epsvec])}
       {D(k)+\ii\eta\epsvec}\right]\\
 &\qquad:=2\pi\delta^{\rm sh}_{\bb}(D(k))\,\theta(\eta\ee\!\cdot k),\\
 &D(k)=k^2-m^2.
\end{aligned}
 \label{eq:physical-cut-map}
% Monograph source: eq:stress-cutkosky; eq:operational-directed-line-cut-replacement.
\end{equation}
The step function here is not a claim that the physical A spectrum has
$p_\parallel>0$ when $\eta=+1$. In physical signed-carrier momentum its
argument is $\sigma_s p^{\rm ph}_\parallel$ and its measure is
\eqref{eq:fw-positive-state-measure}. Endpoint dualization and the
spin/adjoint state sum belong to that same map. Changing the selector
without transporting the line routing and subamplitudes would change
the cut rule.

This operation removes the virtual line, its ordering
coefficient, its mode unit and its kernel together, replacing them
by the normalized locked completeness projector. It is not a
distributional identity for the uncut propagator. On a canonical
aligned ray the extracted uncut coefficient is
\(\sigma_\epsilon\ii R^\Ph_\eta\), as in
\eqref{eq:scalar-line-coefficient}; it is not obtained by
discarding the sector factor in that equation.

For comparison with analytic boundary values, define
\begin{equation}
 R^{\rm coeff}_{\pm,\bb}(k)
 :=\lim_{\varepsilon\downarrow0}
       \embed^{-1}\!\left[
       f^{\rm dir}_{\pm,\bb}(k;\varepsilon)
        (D(k)\mathbin{\pm}\ii\epsvec)^{-1}\right].
 \label{eq:cut-directed-coefficient-extraction}
\end{equation}
Define the corresponding directed boundary coefficients by
\(\mathcal R^{\rm dir}_{\pm,\bb}
:=\embed(R^{\rm coeff}_{\pm,\bb})\).
Only after this extraction write
\(R^{\rm coeff}_{\pm,\bb}=f_\pm/(D\pm\ii0)\).
Their scalar jump is
\(\Disc^{\rm coeff}R=R^{\rm coeff}_+-R^{\rm coeff}_-\).
The corresponding jump of the directed coefficient is
\begin{equation}
 \Disc^{\rm dir}R
 :=\mathcal R^{\rm dir}_{+,\bb}
       -\mathcal R^{\rm dir}_{-,\bb}
 =\embed(\Disc^{\rm coeff}R).
 \label{eq:directed-discontinuity-reconstruction}
\end{equation}
It is distinct from reversing a complete graph's evolution
branch, which also reverses its ordering and Dyson data.
The auxiliary coefficient identity
\((D+\ii0)^{-1}=\PV(1/D)-\ii\pi\delta^{\rm sh}_{\bb}(D)\)
does not supply the on-shell orientation selector.

Let the continued numerator have boundary values \(f_+\) and
\(f_-\), smooth near a regular shell \(D=0\) with \(\dd D\ne0\).
Assume also that the regulated numerators converge as multipliers of
the resolvent boundary distributions; local \(C^1\) convergence on
compact shell charts is sufficient. The identities below are local
there (with test functions of compact support); a global identity
requires control of other singularities and infinity. Pointwise
numerator limits alone do not justify multiplying two distributional
limits. The ordinary
coefficient jump is
\begin{align}
 \frac{f_+}{D+\ii0}-\frac{f_-}{D-\ii0}
  &=(f_+-f_-)\PV\frac1D
     -\ii\pi(f_++f_-)\big|_{D=0}\delta^{\rm sh}_{\bb}(D).
 \label{eq:numerator-disc}
% Monograph source: crit:cut-vs-disc; eq:cut-disc-general-boundary-numerator.
\end{align}
To verify the stated sufficient multiplier condition, use a regular
shell chart \((y,\xi)\) with \(y=D\), and absorb its smooth
Jacobian into a compactly supported test function. If the numerator
error tends to zero in \(C^1\), its tested product can be written
\(g_\varepsilon(y,\xi)=g_\varepsilon(0,\xi)\chi(y)
+y r_\varepsilon(y,\xi)\), where \(\chi=1\) near zero and
both coefficients tend uniformly to zero on a fixed compact set.
The pairing of \(\chi\) with \((y\pm\ii\varepsilon)^{-1}\)
is bounded, and \(|y/(y\pm\ii\varepsilon)|\leq1\).
Both error terms therefore vanish. Applying the ordinary resolvent
boundary formula to the remaining fixed smooth numerators proves
\eqref{eq:numerator-disc} on the chart.
If \(f_+=f_-=f\), this reduces to
\begin{equation}
 \Disc^{\rm coeff}\frac{f}{D+\ii0}
       =-2\pi\ii\,f|_{D=0}\delta^{\rm sh}_{\bb}(D).
 \label{eq:weighted-pole-disc}
\end{equation}
For the directed coefficient, the same formula is transported in
its entirety:
\begin{equation}
 \Disc^{\rm dir}R
 =\embed\!\left[
 (f_+-f_-)\PV\frac1D
   -\ii\pi(f_++f_-)|_{D=0}\delta^{\rm sh}_{\bb}(D)\right].
 \label{eq:directed-numerator-discontinuity}
\end{equation}
If the common numerator is real, the continuous-numerator
specialization is
\(-2\pi\je f|_{D=0}\delta^{\rm sh}_{\bb}(D)\).
For complex numerators one must retain the complete embedding.
Neither expression discards the mode numerator or the remaining
graph factors; those enter the full reconstruction
\eqref{eq:complete-graph-character-reconstruction}.

Thus even an entire numerator can change the pole weight.
A numerator branch jump adds the principal-value contribution in
\eqref{eq:numerator-disc}. Removing the damper by a prescribed
\(\Cut_{\rm phys}\) operation does not remove these terms from
\(\Disc_{\rm an}\).

For the stable-half kernel, form the full directed argument and
transported functions before taking a boundary value:
\begin{equation}
\begin{aligned}
 W_{\eta,\bb}(k)&:=-(k^2+\ii\eta\epsvec),\\
 f^{\rm dir}_{a,\eta,\bb}(k;\varepsilon)
   &:=\ExpDir\!\left[-a\SqrtDir(W_{\eta,\bb}(k))\right].
\end{aligned}
 \label{eq:stable-full-directed-kernel}
% Monograph source: eq:stable-half-full-directed-argument; eq:stable-half-directed-boundary-values.
\end{equation}
Here \(\SqrtDir:=\embed\circ\sqrt{\,\cdot\,}
\circ\embed^{-1}\) uses the principal scalar square-root
sheet, and \(\ExpDir:=\embed\circ\exp\circ\embed^{-1}\).
For \(k^2>0\), their directed boundary is
\begin{equation}
\begin{aligned}
 &\lim_{\varepsilon\downarrow0}
 f^{\rm dir}_{a,\eta,\bb}(k;\varepsilon)\\
 &\qquad=\cos(a\sqrt{k^2})+\eta\je\sin(a\sqrt{k^2}),\\
 &\lim_{\varepsilon\downarrow0}
 \embed^{-1}\!\left(f^{\rm dir}_{a,\eta,\bb}(k;\varepsilon)\right)
       =e^{\ii\eta a\sqrt{k^2}}.
\end{aligned}
 \label{eq:stable-directed-boundary}
\end{equation}
The square root of the positive real number \(k^2\) in this
boundary formula is an ordinary nonnegative real root.
For \(m>0\), the shell is regular and lies away from the branch
point. The ordinary coefficient jump and its directed reconstruction
are
\begin{equation}
\begin{aligned}
 \Disc^{\rm coeff}R_a
  ={}&2\ii\theta(k^2)\sin(a\sqrt{k^2})\PV\frac1D\\
     &-2\pi\ii\cos(am)\delta^{\rm sh}_{\bb}(D),\\
 \Disc^{\rm dir}R_a
  ={}&2\je\theta(k^2)\sin(a\sqrt{k^2})\PV\frac1D\\
     &-2\pi\je\cos(am)\delta^{\rm sh}_{\bb}(D).
\end{aligned}
 \label{eq:stable-disc}
\end{equation}
The full line remains
\(\widetilde G^\Ph_{F,\bb}
=\eta\ii\widehat u_{k,\bb}\mathcal R^{\rm dir}_{\eta,\bb}\).
Thus the displayed jump concerns its directed denominator/kernel
coefficient, not a graph obtained by omitting that numerator.
At \(m=0\), put
\(z_{\eta,\bb}=k^2+\ii\eta\epsvec\). On the selected
directed sheet,
\begin{equation}
 \frac{\ExpDir[-a\SqrtDir(-z_{\eta,\bb})]}
              {z_{\eta,\bb}}
 =\frac1{z_{\eta,\bb}}
      +\frac{a}{\SqrtDir(-z_{\eta,\bb})}+O(1).
 \label{eq:massless-branch}
\end{equation}
The branch point therefore prevents an isolated meromorphic pole
from following merely from \(\Ph(0)=1\).
For this particular coefficient its single-line jump extends through
the cone tip in four momentum dimensions. It equals
\[
 2\ii\,\mathbf1_{k^2>0}\frac{\sin(a\sqrt{k^2})}{k^2}
                  -2\pi\ii\,\delta^{\rm sh}_{\bb}(k^2).
\]
The first term is locally integrable and the cone measure has finite
mass near the origin. To see why the exponent of a denominator matters,
write \(k=r\omega\) in Euclidean polar coordinates on the real
momentum chart. Then \(k^2=r^2 U\), with \(U=2\omega_0^2-1\).
For normalized spherical measure, \(U=u\) has density
\((1-u)/(\pi\sqrt{1-u^2})\) on \((-1,1)\).
On the unit sphere the cone \(U=0\) is regular. Angular testing of
\((U+\ii\eta t)^{-1}\) stays bounded uniformly for \(t>0\) by
the regular-shell formula and subtraction of the test value. The radial
majorant for the simple pole is therefore \(Cr\,\dd r\).
The remainder after this pole in \eqref{eq:massless-branch} has at
worst an angular singularity \(|U|^{-1/2}\) and radial degree \(-1\).
These bounds justify the jump on tests including the tip.

The uncontracted scalar square behaves differently. Its exact
normalized angular average is
\[
 \int_{-1}^1\frac{1-u}{\pi\sqrt{1-u^2}}
       \frac{\dd u}{(u+\ii\eta t)^2}
     =\frac{\ii\eta-t}{(1+t^2)^{3/2}},\qquad t>0 .
\]
Indeed the simple resolvent average at \(z=\ii\eta t\) is
\((1+z)/\sqrt{z^2-1}-1\), on the branch asymptotic to \(z\)
for the square root. Its negative derivative gives the square.
With \(t=\varepsilon/r^2\), a compactly supported radial test
equal to one near the tip gives imaginary part
\(\eta\pi^2\log(\varepsilon_*/\varepsilon)+O(1)\)
for the square paired against \(\dd^4 k_{\bb}\), for any fixed
positive reference height \(\varepsilon_*\). The coefficient uses
\(|S^3|=2\pi^2\). In the longitudinal gauge term of
\eqref{eq:physical-directed-gauge-ghost-lines}, the actual numerator
\(k_\mu k_\nu\) supplies two powers of \(r\). The angular
double-pole bound follows by one derivative of the simple-pole
distribution at the regular angular shell. For the stable-half
multiplier the next angular power has order \(-3/2\) and is one
derivative of a locally integrable power of order \(-1/2\).
The worst radial majorant for the complete tensor is then
\(Cr\,\dd r\). The running sign is smooth near either angular
shell component. This argument treats a single tensor line.

\Needspace{4\baselineskip}
\begin{assumption}[Cut--discontinuity compatibility]
\label{ass:cut-matching}
For the completed physical amplitudes, the prescribed on-shell
cut operation agrees with the analytic discontinuity, including
the actual pole residues, completed subamplitudes, kernel branch
terms, and contour contributions. Unphysical states cancel under
the admitted BRST construction, and the remaining state sum has
positive physical norm.
\end{assumption}

Cut--discontinuity compatibility can be established by an action-level
unitarity argument with the requisite analytic conditions, or by an
explicit matching calculation. Neither \eqref{eq:kernel-inert} nor a
positive Laplace measure establishes this compatibility.
The standard Cutkosky theorem identifies discontinuities under its
analytic and Feynman-integral hypotheses \cite{Cutkosky1960}.
Section~\ref{sec:gaussian-bubble} provides an exact example in which
the naive unweighted cut fails this test.

% Monograph source: prop:LSZ; eq:LSZ-plane-wise-operational-cut-sum; crit:cut-vs-disc.
For the optical relation, assume in addition that the physical
scattering operator is unitary on the chosen positive-norm asymptotic
state space and that the intermediate-state resolution is complete,
including every open physical channel. Here $S^{\rm phys}$ acts on physical incoming/outgoing state spaces;
when represented by mixed carrier boundaries, their endpoint dualities
are those of \eqref{eq:fw-mixed-boundary-amplitudes}. It is not an
arbitrary partially transposed unitary. Writing
\(S^{\rm phys}=\one+\ii T^{\rm phys}\), one obtains
\begin{equation}
 (S^{\rm phys})^\dagger S^{\rm phys}=\one
 \quad\Longrightarrow\quad
 -\ii\bigl[T^{\rm phys}-(T^{\rm phys})^\dagger\bigr]
       =(T^{\rm phys})^\dagger T^{\rm phys}.
 \label{eq:physical-unitarity-identity}
\end{equation}
Here \(\dagger\) is the adjoint on that physical Hilbert space,
not the coefficient reality operation \(\ddagger\).
Proposition~\ref{prop:carrier-reality-adjoint} rules out identifying
these on the full carrier algebra. Neither a character norm nor a
reality relation can replace the positive physical state resolution.
Take a forward matrix element and insert the complete physical
state resolution between the two operators on the right. This
argument first uses normalizable wave packets. Passing to
momentum-space matrix-element coefficients and factoring out the
common momentum-conservation distribution then gives
\begin{equation}
 2\Im\mathcal M^{\rm phys,coeff}_{\bb}(i\to i)
  =\sum_n\int\dd\Pi_{n,\bb}\,
       |\mathcal M^{\rm phys,coeff}_{\bb}(i\to n)|^2.
 \label{eq:optical}
\end{equation}
Here \(\dd\Pi_{n,\bb}\) includes the physical state normalization,
momentum conservation and identical-particle counting factors.
In perturbation theory both sides, including the squared amplitudes,
are expanded to the same order in all couplings. Identifying the state sum with an analytic
discontinuity
additionally requires Assumption~\ref{ass:cut-matching} and the
appropriate analytic boundary relations.

Alternatively, the same forward optical identity follows from a
normalized analytic matching calculation if the physical amplitude
has the Hermitian-analytic boundary relation
\[
 \mathcal M_-^{\rm phys,coeff}(i\to i)
   =\overline{\mathcal M_+^{\rm phys,coeff}(i\to i)}
\]
and its matched discontinuity factorizes as
\[
\begin{aligned}
 &\mathcal M_+^{\rm phys,coeff}(i\to i)
  -\mathcal M_-^{\rm phys,coeff}(i\to i)\\
 &\qquad=\ii\sum_n\int\dd\Pi_{n,\bb}\,
       |\mathcal M_+^{\rm phys,coeff}(i\to n)|^2.
\end{aligned}
\]
The subscripts \(+\) and \(-\) here denote the two boundary
approaches to the selected external-invariant cut after physical
representation and the admitted coefficient extraction. They do
not denote the carrier characters \(\chi_\pm\). A relation
between carrier components alone does not establish this
Hermitian analyticity. Without it, an imaginary entire addition
can change the imaginary part of a one-loop amplitude while
leaving its discontinuity and its tree-level cut unchanged.

A carrier-valued graph, including its cut graph, is reconstructed
from both complete character calculations by
\eqref{eq:complete-graph-character-reconstruction}; its carrier-odd
part is not defined to be an absorptive probability.

Here the amplitudes are those of the completed theory, with
physical residues and state normalization. Plane diagonality
restricts the intermediate states to the same support in the
absence of support-changing interactions. It does not supply
unitarity by itself, and it does not replace completed on-shell
subamplitudes by undeformed ones.

\subsection{What the physical sum over intermediate boundary states retains}
\label{subsec:physical-inclusive-uv}
A geometric phase can suppress every fixed final-state overlap and
still leave an unsuppressed inclusive current. To test this issue,
fix one reconstructed Lorentzian sector plane $\bb=(s,b_s)$ with
positive-norm asymptotic states and a Hermitian conserved current.
Write $S=P^2$ for the timelike invariant on that plane. In the
following calculation the physical cut and its orientation are those
of that reconstruction. The normal phase is a matrix element on the
cut. Its argument must be determined there by the same source that
produces the Euclidean vertices.

Take distinguishable spinor--antispinor pairs in mutually orthogonal
channels $n$, with equal mass $m_n$ in each pair. Consider a vector
current with pair matrix element
$e\,c_n(S)\bar u(p)\gamma_\mu v(p')$ and covariant state
normalization in the sense $\langle\mathbf p|\mathbf p'\rangle
=2E(2\pi)^3\delta^3(\mathbf p-\mathbf p')$ on the selected plane.
Use the transverse scalar convention
$\rho=(1/\pi)\operatorname{Im}\Pi$ on the positive cut. The pair
contribution is
\begin{equation}
 \begin{aligned}
 \rho_{\rm pair}(S)&=\frac{e^2}{12\pi^2}
  \sum_{4m_n^2\le S}|c_n(S)|^2\,
       \beta_n(S)\left(1+\frac{2m_n^2}{S}\right),\\
 \beta_n(S)&=\sqrt{1-4m_n^2/S}.
 \end{aligned}
 \label{eq:physical-inclusive-pair-uv}
\end{equation}
Indeed the spin trace summed over three spatial components in the
pair rest frame is $4(S+2m_n^2)$. Multiplication by the integrated
two-body phase space $\beta_n/(8\pi)$ and division by $6\pi S$
gives the formula. This is the pure Dirac-vector case of the
timelike form-factor decomposition~\cite{PhysicalDenigSalme}.
Further channels enter the inclusive spectral measure positively.
If the physical transverse scalar is represented by an unsubtracted
positive dispersion integral at Euclidean $X>0$, its pair tail
therefore requires $\int^{\infty}\rho_{\rm pair}(S)\,\mathrm dS/S<\infty$.
This extra dispersive hypothesis has to be proved on the physical
source. An absolute Euclidean graph estimate alone does not impose it.

For a unitary displacement $U$ on the boundary space with incoming
state $v$ of norm one, the channel weights are
$|c_n|^2=|\langle n|U v\rangle|^2$. Their sum over a complete
orthonormal basis is one. The relevant object on the cut is therefore
$\langle v,U^*W_S Uv\rangle$, where $W_S$ is the nonnegative
diagonal operator of threshold and phase-space weights in
\eqref{eq:physical-inclusive-pair-uv}. This formula evaluates the accessible part
of the displaced state. An elastic matrix element alone omits it.

\begin{proposition}[An explicit test at the opening of infinitely many channels]
\label{prop:physical-inclusive-opening-uv}
As a comparison, take a boundary harmonic oscillator of frequency
$\omega>0$ in units of oscillator mass and $\hbar$ equal to one.
Let $X_b=(a+a^\dagger)/\sqrt{2\omega}$. Specify the squared masses
$m_n^2=m^2+n\Delta$, $m,\Delta>0$, and the on-cut vector-current
overlaps $c_n(S)=\langle n|e^{i\xi(S)X_b}|0\rangle$ with
$\xi(S)^2/(2\omega)=bS$, $b>0$. These are hypotheses about this
comparison's boundary mass operator and timelike current. Put
$M=S/(4\Delta)$ and $\ell=4\Delta b$. For nonzero $e$, the pair
spectral density has a positive limit if $\ell<1$, falls as a positive
constant times $M^{-1/4}$ if $\ell=1$, and is exponentially bounded
at large $S$ if $\ell>1$. Its unsubtracted ultraviolet tail is
finite exactly for $\ell\ge1$.
\end{proposition}
\begin{proof}
The oscillator commutator gives
$|c_n(S)|^2=e^{-bS}(bS)^n/n!$. Thus, with
$N$ Poisson of mean $\ell M$ and $a_0=m^2/\Delta$, the sum in
\eqref{eq:physical-inclusive-pair-uv} without its constant factor is the expectation
of $w((N+a_0)/M)$ for
\[
  w(x)=\sqrt{(1-x)_+}(1+x/2),\qquad x\ge0 .
\]
If $\ell<1$, the law of large numbers and bounded continuity of $w$
give the positive limit $w(\ell)$. For $\ell>1$, the expectation is
at most $(3/2)\Pr(N\le M)$, bounded by
$(3/2)\exp[-M(\ell-1-\log\ell)]$ by the Poisson exponential
moment. At $\ell=1$, $(N-M)/\sqrt M$ tends in law to a standard
normal variable $Z$. On the open channels, $M^{1/4}$ times the
threshold factor equals
$[ -(N-M)/\sqrt M-a_0/\sqrt M]_+^{1/2}$.
Its fourth moment is uniformly bounded by
$\mathbb E[(N-M)^2/M]=1$. The second factor may be written as $1+\min\{(N+a_0)/M,1\}/2$
on the entire sample space, since the first factor vanishes on closed
channels. It is bounded and tends to $3/2$ in probability. Consequently
the rescaled expectation tends to
$(3/2)\mathbb E[(-Z)_+^{1/2}]>0$. Integration against
$\mathrm dS/S$ gives the stated test.
\end{proof}

At every fixed $n$ the overlap in this example has an exponential
tail. Whether that decay survives the physical sum depends on the
excitation spectrum and on the current that the incoming state actually
creates. For the chiral GCM branch, both must follow from the same
charged operator and boundary state. The calculation makes no
continuation of a Euclidean heat multiplier to this timelike overlap.

\paragraph{The pair current produced by one relative Wilson chord.}
The transverse coupling in \eqref{eq:phys-uv-relative-chord} specifies a
test in which the timelike matrix element can be computed from the vertex.
Use the continuum symbols on one reconstructed sector plane, take $m>0$,
and write $a=\sigma^2=\kappa r^\alpha$, $0<r\le1$. Take the principal
square root of the scalar Euclidean transfer $X$. Its two boundary values
give $h_{\pm}(S)=\exp(\mp i r\sqrt S)$ at $X=-S\pm i0$, $S>0$.
These signs label approaches in the $X$-plane. Exchanging these two scalar boundary values conjugates $h_\pm$.
In a physical matrix element the continuation follows the carrier
and the ordering branch on $\bin=(s,b_s)$, whose pole displacement is
$i\sigma_s\sigma_\epsilon\varepsilon\hat p_{0,\bin}$,
$\varepsilon\downarrow0$. The signs on $h_\pm$ label the two
sides of the scalar $X$ cut.

For an outgoing particle and antiparticle write $P=p+p'$, $d=p-p'$,
$p^2=p'^2=m^2$ and $S=P^2\ge4m^2$, in Lorentz signature $(+---)$.
The straight chord assigns incoming charged momentum $-p'$ and outgoing
momentum $p$. Its first Wilson derivative inserts a uniform chord
parameter $t\in[0,1]$. The charged symbol $k\exp(-a k_E^2)$ is
differentiated at $k=-p'+tP$ before continuation. Thus the continued
vertex, up to a common current orientation, is $e h_\pm$ times
\[
 \int_0^1 e^{a k^2}
       \big(\gamma^\mu+2a k^\mu\slashed k\big)\,\mathrm dt .
\]
Here $k^2=m^2-S t(1-t)$. The free on-shell equations give
$\bar u\slashed k v=m\bar u v$ at every $t$, and symmetry about
$t=1/2$ removes the component of $k=d/2+(t-1/2)P$ parallel to $P$.
The resulting matrix element is
\begin{equation}
 \begin{aligned}
  J^\mu_\pm&=e h_\pm A_a(S)\,
        \bar u(p)\big(\gamma^\mu+am d^\mu I\big)v(p'),\qquad
  P_\mu J^\mu_\pm=0,\\
  A_a(S)&=e^{am^2}I(aS),\qquad
  I(x)=\int_0^1 e^{-xt(1-t)}\,\mathrm dt .
 \end{aligned}
 \label{eq:physical-relative-pair-uv}
\end{equation}
In particular the timelike photon factor has unit modulus at fixed depth.
The charged interpolation retains the two ends of the chord at large $S$.

With the convention
$\Gamma^\mu=F_1\gamma^\mu+F_2 i\sigma^{\mu\nu}P_\nu/(2m)$,
the pair Gordon identity gives $F_1=h_\pm A_a(1+2am^2)$ and
$F_2=-2am^2h_\pm A_a$. The magnetic and electric form factors and
the corresponding positive two-particle density are therefore
\begin{equation}
 \begin{aligned}
  G_M&=h_\pm A_a,\qquad
  G_E=h_\pm A_a(1+2am^2-aS/2),\\
  \rho_{r,\mathrm{pair}}(S)
   &=\frac{e^2\beta}{12\pi^2}
       \left(|G_M|^2+\frac{2m^2}{S}|G_E|^2\right),
       \qquad \beta=\sqrt{1-4m^2/S}.
 \end{aligned}
 \label{eq:physical-relative-pair-density-uv}
\end{equation}
At fixed depth the zero-transfer limit is
$F_1(0)=e^{am^2}(1+2am^2)$. For the normalized boundary vector with
common free depth law $\nu_0$, the elastic value is its $\nu_0$ average.
Using $ez_m$ in the transverse Wilson phase normalizes the averaged
elastic charge as above and multiplies the pair densities by $z_m^2$.
The normalization follows from the spin sum in the pair rest frame.
For momentum along axis three, each of the first two current components
has spin sum $2e^2 S|G_M|^2$ and the third has
$8e^2m^2|G_E|^2$. Multiplication by $\beta/(8\pi)$ and division
by $6\pi S$ yields the displayed normalization. This is the positive
state-sum contribution if the indicated charged states and current
exist in the reconstruction. At threshold $G_E=G_M$. At fixed
$a>0$, changing variable to $v=t(1-t)$ on the half interval gives
$I(x)=2\int_0^{1/4}e^{-xv}(1-4v)^{-1/2}\,\mathrm dv
=2/x+O(x^{-2})$ as $x\to\infty$. In particular
\[
 \lim_{S\to\infty} S\rho_{r,\mathrm{pair}}(S)
       =\frac{e^2 m^2}{6\pi^2}e^{2am^2}.
\]
The unsubtracted pair tail $\int^\infty \rho_{r,\mathrm{pair}}(S)
\,\mathrm dS/S$ is finite. The electric term determines its leading
power even though the magnetic amplitude decreases as $S^{-1}$.
In fact the unsubtracted Stieltjes transform of this pair contribution
alone obeys
\[
 \lim_{X\to\infty}
 \frac{X}{\log(X/m^2)}
 \int_{4m^2}^\infty\frac{\rho_{r,\mathrm{pair}}(S)}{S+X}
         \,\mathrm dS
   =\frac{e^2m^2}{6\pi^2}e^{2am^2}.
\]
To see this, split at a fixed large $S_0$, bound $S\rho$ above and
below on the tail by its limit plus or minus any fixed positive
tolerance, and use
$X\int_{S_0}^\infty [S(S+X)]^{-1}\,\mathrm dS=\log(1+X/S_0)$.
The bounded interval contributes $O(1)$ after multiplication by $X$.

For one common normalized free depth law $\nu_0$, two distinct state
sums are meaningful. If the depth acts by multiplication in
$L^2(\nu_0)$, the external boundary vector is $1$, and all boundary
channels in this pair sector have the same threshold and are accessible,
closure gives $\int\rho_{r,\mathrm{pair}}\,\mathrm d\nu_0$.
Its limit after multiplication by $S$ is the depth average of
the fixed-width coefficient $e^2m^2e^{2am^2}/(6\pi^2)$ when $\langle a^{-1}\rangle_0<\infty$. Indeed
$I(x)\le C/(1+x)$, so $S|G_M|^2\le C/a$ and $|G_E|\le C$
uniformly for $0<a\le\kappa$. The elastic projection onto that same
boundary vector instead replaces $G_M,G_E$ individually by their
averages before taking squared moduli. It is bounded above by the
closure result. It alone already gives a nonzero pair density on some
bounded open interval above threshold for $e\ne0$. In fact its magnetic
amplitude is entire as a function of $u=\sqrt S$, since $r,a$ are
bounded and the chord integrand is entire and locally uniformly bounded
there. At $u=0$ its value is $\langle e^{am^2}\rangle_0>0$.
It cannot vanish on any real interval. A different mass operator on
the boundary space requires the actual accessible-state projection.

\paragraph{Testing the Euclidean response against the same pair current.}
There is a stronger consistency check than convergence of the positive
pair tail. Let $K_T(X)$ be the average over transverse unit
polarizations of the continuum, translation-invariant, order-$e^2$
quadratic response of the relative determinant per unit volume.
Its overall sign can be fixed to the convention for the current
correlator. Both the bubble and the two-mark contact contain
$e^{-2r\sqrt X}$. Each one-mark charged symbol is an average of
translated Gaussian derivatives, bounded by $C$ in supremum norm and
by $C\sigma^{-4}$ in momentum integral. The two-mark symbol has
integral bound $C\sigma^{-3}$. These estimates are uniform in the
external Euclidean momentum, since the chord parameters only translate
the argument of a rescaled Schwartz function. The massive free
inverses have norm at most $m^{-1}$. Translation invariance excludes
one-mark means at nonzero momentum, so the common-depth average adds
no cross term. With the fixed free radial law this gives
\[
 |K_T(X)|\le C e^2
  \int e^{-2r\sqrt X}
  \left(m^{-2}\sigma^{-4}+m^{-1}\sigma^{-3}\right)
  \,\mathrm d\nu_0(r).
\]
This statement concerns the continuum graph coefficient at fixed law.
If every inverse depth moment is finite, $K_T(X)=O_L(X^{-L})$ for
every $L>0$. The scalar response $\Pi_E(X)=K_T(X)/X$ has the same
rapid decrease at infinity. Both terms required by differentiation of
the charged determinant have been included.

Suppose a positive dispersive representation of this same leading
scalar response has the pair current \eqref{eq:physical-relative-pair-uv}
in an accessible channel and admits at most finitely many local
subtractions. This is incompatible with the preceding bound. Indeed
write $\Delta_h f(X)=f(X+h)-f(X)$. For $N$ large enough to kill the
subtraction polynomials, the exact kernel identity is
\[
 (-1)^N\Delta_h^N\frac{1}{s+X}
   =\frac{N! h^N}{\prod_{j=0}^N(s+X+jh)} .
\]
Apply it to $K_T=X\Pi_E$, using
$X/(s+X)=1-s/(s+X)$. Taking $N\ge1$ also kills the constant and
allows a possible local constant contact in $K_T$. With the dispersive
sign chosen positive, $(-1)^{N+1}\Delta_X^N K_T(X)$ is an integral
of the displayed positive kernel with weight $s\,\mathrm d\rho(s)$.
Here the step is held fixed during the differencing and then set to
$h=X$. On any occupied compact mass interval inside $(4m^2,\infty)$
the integral is bounded below by $c/X$ for large $X$. In contrast,
every term $K_T((j+1)X)$ decreases faster than every inverse power.
An equal-mass closure is unnecessary for this contradiction if the
elastic pair is accessible. Additional positive channels only increase
the lower bound.

Thus the conditional pair matrix element has a finite dispersive
ultraviolet tail, yet its positive spectral contribution cannot
reproduce the full Euclidean response of this coupling under the stated
dispersion hypotheses. A two-mark Wilson contact contributes to the
second derivative of the source and carries the same nonlocal readout;
it does not create a pair at first order in the current. Its analytic
continuation and the timelike square-root boundary values can affect
the full discontinuity. An optical theorem for that discontinuity,
or a completion supplying further physical current terms, requires
an independent construction.

\subsection{Which curvature pairing transports the electromagnetic probe?}
\label{sec:closure-native-uv}
The spectral tests distinguish normal transport of all four Euclidean
components from transport confined to one spatial slice. This distinction
can be read directly from the quadratic curvature form before introducing
a charged path. On one fixed sector plane, write $\tau$ for the
continued carrier coordinate, $r$ for proper normal distance and
$\mathbf k$ for a spatial momentum. Choose one spatially transverse Abelian
polarization, with $k=|\mathbf k|>0$. On a diagonal, spatially isotropic
curvature block, the terms $|F_{ri}|^2$, $|F_{\tau i}|^2$ and
$\sum_{i<j}|F_{ij}|^2$ give the mode energy below in reference units. Put $y=\omega^2$.
\[
 \mathcal E_y[u]=\tfrac12\int_0^\infty
 \left(a(r)|u'|^2+
                  [b_{\rm sp}(r)k^2+b_{\parallel}(r)y]|u|^2\right)\,\mathrm dr .
\]
Indeed the longitudinal and scalar potentials drop out of this
polarization. Volume Jacobians and inverse metric coefficients are
included in $a,b_{\rm sp},b_{\parallel}$. For this test assume that they are real
bounded continuous functions, $a,b_{\rm sp}$ bounded below by positive
constants and $b_{\parallel}\ge0$. Fix the boundary amplitude to one and use
the finite-energy solution at infinity. The same deterministic
extension may be imposed by its normalized constraint. A separately
integrated normal photon would also contribute its Dirichlet noise.

\begin{proposition}[Temporal response of the normal curvature block]
\label{prop:closure-native-uv}
Let $u_y(r)$ be the minimizing transfer and let
$\Lambda(y)=-a(0)u_y'(0)$ be its boundary flux. With $G_y^D$ the
Dirichlet Green function of the mode operator, one has
\begin{equation}
 \begin{aligned}
 \partial_y\Lambda(y)&=\int_0^\infty b_{\parallel}(s)u_y(s)^2\,\mathrm ds,\\
 -\partial_yu_y(r)&=\int_0^\infty
      G_y^D(r,s)b_{\parallel}(s)u_y(s)\,\mathrm ds .
 \end{aligned}
 \label{eq:closure-native-flux-uv}
\end{equation}
An interior readout at any $r>0$ is independent of carrier frequency
exactly when $b_{\parallel}$ vanishes throughout the connected ray.
\end{proposition}
\begin{proof}
The coercive Dirichlet problem has a unique minimizer. Its equation
is $-(a u_y')'+(b_{\rm sp} k^2+b_{\parallel} y)u_y=0$. The positive potential and the
one-dimensional maximum principle give $0<u_y\le1$ and $G_y^D(r,s)>0$
for positive $r,s$. Integration by parts gives
$\Lambda(y)=2\mathcal E_y[u_y]$. Differentiating this identity uses
stationarity on variations with zero boundary value, giving the first
formula. For the second, differentiate the mode equation with the
boundary amplitude held fixed and invert the Dirichlet operator.
These derivatives follow also by the resolvent identity, since
multiplication by $b_{\parallel}$ is bounded on $L^2$. Both integrands are
strictly positive wherever $b_{\parallel}>0$. If $b_{\parallel}=0$, both the energy and
its boundary condition are independent of $y$, proving the converse.
\end{proof}

A positive temporal coefficient on a strip already constrains the
ultraviolet readout beyond it. Write $w(r)=\int_0^r a(s)^{-1}\mathrm ds$.
In this coordinate $u_{ww}=a(b_{\rm sp}k^2+b_{\parallel}y)u$. Its finite-energy
positive solution is decreasing, since $u_{ww}>0$ and a positive
eventual slope would violate square integrability. If $a b_{\parallel}\ge d>0$
on $[r_1,r_2]$, comparison backwards from $r_2$ with the constant
coefficient equation gives
\[
 u_y(r)\le
 \frac{1}{\cosh\!\left(\sqrt{d}\,|\omega|
                [w(r_2)-w(r_1)]\right)},\qquad r\ge r_2 .
\]
The terminal derivative in that comparison has the favorable sign.
A finite-order local time derivative at the retained endpoint costs
only a polynomial in $|\omega|$. Thus it cannot undo exponential decay
across this strip. A mixture reaching $r=0$ must integrate the actual
attenuation length: for example, if $a b_{\parallel}$ has a uniform positive
lower bound up to the sampled distances, all inverse moments of $r$
suffice for decay faster than every power. The flux identity also
shows why a temporal kinetic term solely \emph{beyond} a measured
layer can still change its finite-frequency transfer.

The two fixed pairings $a=b_{\rm sp}=b_{\parallel}=1$ and $a=b_{\rm sp}=1$, $b_{\parallel}=0$, have
identical static extension, static boundary flux, and the same Wilson
charge on a fixed spatial probe. They give respectively
$e^{-r\sqrt{k^2+\omega^2}}$ and $e^{-rk}$ at nonzero frequency.
Both are obtained by contracting the indicated curvatures on this
transverse channel with the same fixed boundary value. The second
pairing is degenerate on $F_{\tau i}$ and uses a distinguished
carrier splitting. A nondegenerate $O(4)$-invariant Euclidean
tangential two-form norm on this diagonal block gives equal weights
to $F_{\tau i}$ and to the spatial pairs in an orthonormal frame.
If the full action requires that norm, the spatial-only option
is unavailable on this channel. If the geometric boundary state
supplies the splitting, its covariance under changes of sector-plane
frame and its action on the charged current must be derived.
The vectorial pole specifies the boundary value on the reconstructed
plane; it does not determine the real curvature coefficient $b_{\parallel}$.
Equations~\eqref{eq:closure-native-flux-uv} give a test on the
stationary quadratic pairing once that pairing has been selected.

\paragraph{A perimeter test on the photon actually sampled by a closed history.}
The positive spectrum of a measured photon constrains its coupling
to a thin charged path independently of the normal mode calculation.
Take a translation- and rotation-invariant Euclidean four-plane.
Suppose that the photon tested on conserved currents has transverse
covariance $P_T(p)G(p^2)$ with an unsubtracted spectral law
$G(x)=\int_{[0,\infty)}(x+s)^{-1}\rho(\mathrm ds)$,
$\rho\ge0$, $\rho\ne0$, finite for $x>0$. A bounded interval then
has weight $Z>0$, so $G(p^2)\ge Z/(p^2+M^2)$ for some finite $M$.
Couple this very observable by the ordinary Abelian line integral
along a closed polygon $C$. In its quadratic response to the test
charge, the two marks on a single edge and all inter-edge terms
combine into the current norm
\[
 \int_{\mathbb R^4}G(p^2)|J_C(p)|^2
                       \frac{\mathrm d^4p}{(2\pi)^4},\qquad
 J_C(p)=\oint_C e^{ip\cdot x}\,\mathrm dx .
\]
The longitudinal projection drops out because $p\cdot J_C=0$.
This expression, up to the sign and factorial of the Taylor
coefficient, follows by differentiating the path exponential twice.
It uses only the photon two-point function. The perimeter divergence
of thin Wilson contours is familiar~\cite{ClosureDotsenkoVergeles};
the following estimate also handles retracing and a positive distribution
of paths.

\begin{proposition}[The ultraviolet cost of a thin closed current]
\label{prop:closure-native-perimeter-uv}
Split collinear overlapping edges into disjoint subsegments and add
their signed multiplicities. Write $m_j$ and $L_j$ for the resulting
multiplicities and lengths. For a nonnegative nonzero smooth radial
window $\varphi$ supported in $1<|p|<2$,
\[
 \begin{aligned}
 D_R&=\int\varphi(p/R)|J_C(p)|^2\frac{\mathrm d^4p}{(2\pi)^4},\\
 \lim_{R\to\infty} R^{-3}D_R&=c_\varphi\sum_j m_j^2 L_j,\\
 c_\varphi&=\int_{\mathbb R^3}
              \varphi(0,\mathbf q)\frac{\mathrm d^3q}{(2\pi)^3}>0 .
 \end{aligned}
\]
If the net current is nonzero, the quadratic response is infinite.
\end{proposition}
\begin{proof}
Write $\check\varphi$ for the inverse Fourier transform with measure
$\mathrm d^4p/(2\pi)^4$. After expansion of the squared current,
the spatial kernel is $R^4\check\varphi(R(x-x'))$. For an edge
with unit tangent $v$ the diagonal integral divided by $R^3$ is
\[
 m_j^2\int_{-RL_j}^{RL_j}(L_j-|u|/R)\check\varphi(uv)\,\mathrm du
 \ \longrightarrow\ m_j^2L_j c_\varphi .
\]
Fourier inversion on the line gives $c_\varphi$, independently of
$v$. For two nonparallel edges, change their two arclengths to
coordinates on the plane spanned by the two tangents. The absolute
integral is $O(R^2)$ by Schwartz decay, including edges that
intersect. Disjoint collinear intervals meeting at an endpoint also
give $O(R^2)$, since their distance near that endpoint is the sum
of their distances to it. Separated parallel edges give rapid decay.
Only the diagonal terms survive the limit. On the window
$G(p^2)\ge Z/(4R^2+M^2)$, so each sufficiently high shell in the
full current norm costs at least a positive constant times $R$
when $\sum_j m_j^2L_j>0$.
\end{proof}

For a positive ensemble of such histories the same lower bound on
the averaged shell follows by Fatou's lemma if a nonzero net current
has positive probability. A contractible polygon can have such a
current even though its endpoint displacement and its winding are
zero. Exact backtracking instead cancels from this Abelian line
current. It can still be detected by other ordered probes in the
path-enriched source. The estimate concerns averaging the squared
current of a thin path; a coherent average of different displacements
\emph{before} the two marks are paired produces a different current
kernel. Source-dependent boundary contacts or signed spin traces
also require their own calculation. Thus changing the measured
photon to restore a positive spectrum does not by itself preserve
the earlier absolute graph bounds. The native action has to specify
how that photon enters the charged Wilson marks and their boundary
pairing.

There is a constructive converse for a spatially extended charged
history. Suppose the same positive photon has finite total spectral
weight $M_0$. A spatial loop on a single carrier slice already couples
only to the spatially transverse channel computed above. For a general
closed spacetime path, specify that the same spatial Dirichlet readout
acts on every component of its conserved current. This further choice
of charged source is realized, for example, by translating the whole
curve by $\boldsymbol\xi$ on the slice and integrating its
\emph{linear} coupling to the gauge field against the spatial
Poisson probability of $\boldsymbol\xi$ before exponentiation.
Fourier transformation gives
$J_{C,r}(\omega,\mathbf k)=e^{-r|\mathbf k|}J_C(\omega,\mathbf k)$.
Convolution commutes with the divergence, so the closed source is
still gauge invariant. An open line under a gauge change acquires
the difference of the same spatially convolved gauge function at its
two endpoints. If $L$ is the geometric length of the original path,
$|J_C|\le L$. The full self-contraction is bounded by
\[
 \begin{aligned}
 Q(C,r)
 &\le M_0 L^2\int
 \frac{e^{-2r|\mathbf k|}}{\omega^2+|\mathbf k|^2}
 \frac{\mathrm d\omega\,\mathrm d^3\mathbf k}{(2\pi)^4}\\
 &=\frac{M_0 L^2}{16\pi^2 r^2}.
 \end{aligned}
\]
Here the carrier frequency was integrated on the entire real line.
The absolute integral for two distinct currents is bounded by the
geometric mean of their respective bounds. These same bounds with
individual edge lengths apply before the edge contacts are summed.
For a Gaussian photon with this covariance the Wilson expectation
for prescribed closed histories is exactly the exponential of minus
one half their total current norm squared, with each charge included
in that current. The bound controls every finite Wick contraction on
those paths, including when endpoints of two marks approach each
other. With a common positive law of paths and depths it suffices
at order $2n$ that the moment of
$(\sum_j |e_j| L_j/r_j)^{2n}$ be finite. This computes the photon
contractions of specified histories from one positive spectrum and
one spatial curvature readout. It leaves the charged path amplitudes,
their spin measure and the convergence of the determinant sum to be
derived from the physical charged action.

\subsection{What a positive photon spectrum requires of an interacting readout}
\label{sec:spex-photon-uv}
On a fixed Euclidean sector plane, choose one nonzero spatial wave
vector $\mathbf k$ and a spatial polarization perpendicular to it.
This is an observable transverse component. It can also be read from
the magnetic field after division by $|\mathbf k|$. Use real spatial standing waves and Fourier modes
of period $\beta$ in the continued carrier coordinate. When $\beta$
has length units, the energy variable in the transfer formulas has
inverse-length units. The physical clock calibration multiplies that
variable by $\hbar\|\mathbf c_{\bin}\|$.
Normalize the Maxwell coordinates so that their free covariance is
$C_0(x)=(x+a)^{-1}$, where $a=|\mathbf k|^2$ and
$x=(2\pi n/\beta)^2$. We only use $n\ne0$.

\begin{proposition}[A spectral moment test for the measured photon]
\label{prop:spex-photon-uv}
For the relative Wilson source at fixed widths or on the minimizing
surface with $\sigma^2=\kappa r$, the exact transverse covariance at
each real charge satisfies, for every $N>0$,
\[
 G(x)=C_0(x)+O_{N,\mathbf k,\beta,V}(x^{-N})
 \quad (n\to\infty).
\]
If these nonzero-frequency data have an unsubtracted positive thermal
spectral representation, its measure in squared energy must be the
single atom $\delta_a$. In particular this photon readout cannot both
change its nonzero-frequency two-point function and retain that
positive representation.
\end{proposition}
\begin{proof}
For a single normalized real Fourier coordinate $A_f$, let $D$ be
the full ordinary charged determinant, and let $\partial_f$ vary
the Maxwell field along its corresponding wave. Gaussian integration
by parts twice gives
\[
 G-C_0=C_0^2
 \left(\frac{\int\partial_f^2D}{\int D}
       -\left[\frac{\int\partial_fD}{\int D}\right]^2\right).
\]
The integrals include the Maxwell law and the same radial state.
Translation invariance makes the linear term zero for these modes.
Equivalently one may average the sine and cosine coordinates with
the same normalization. At each marked Wilson endpoint this
variation brings the factor $\exp(-r\sqrt{x+a})$. After taking out
both factors, the twice varied determinant is bounded in absolute
value by $C(1+r^{-L})\exp(Cr^{-2})$ on the minimizing surface.
Indeed the line integral of a unit-amplitude Fourier wave is bounded
by its amplitude times the chord length, independently of its
frequency. The one-block trace of the second mark uses the locally
integrable free Dirac kernel and a product of Hilbert--Schmidt
factors, as in the relative-determinant proof. The two first marks
use the mass bound on both inverses and the Hilbert--Schmidt bound
on each marked kernel. The product of the two first traces in
$\partial_f^2 D$ obeys the same type of bound. The fixed-volume mode
normalization changes only the constant. All these estimates hold for
real Maxwell backgrounds. At fixed widths they are finite constants.
At the radial endpoint the density supplies
$r^4\exp(-A_*r^{-3})$ up to constant factors. Consequently one can
use $e^{-2r\sqrt{x+a}}\le C_j(r\sqrt{x+a})^{-j}$ for every $j$.
This proves the claimed decay for the full normalized observable.

For a positive transfer and a Hermitian spatial mode, the thermal
Lehmann representation at nonzero frequency is
\[
 G(x)=\int_{(0,\infty)}\frac{\rho(\mathrm ds)}{x+s},
 \qquad \rho\ge0 .
\]
Pair transitions of energies $E_l<E_h$. Each pair has
$s=(E_h-E_l)^2$ and contributes its squared matrix element times
$2(E_h-E_l)(e^{-\beta E_l}-e^{-\beta E_h})/Z$ to the measure.
Equal-energy transitions have only zero frequency. This derivation
also identifies the spectral hypothesis when passing from discrete
levels to an energy measure~\cite{SpectrumLaineVuorinen}.
A finite time-local polynomial contact in this scalar convention
is fixed to zero by decay, provided the displayed transform exists.

Along the unbounded sequence $x_n$, monotone convergence applied to
$x/(x+s)$ gives $\int\rho=1$. Next $x^2(1/x-G)$ gives
$\int s\,\rho(\mathrm ds)=a$ by the same argument. The exact identity
\[
 (x+a)^2[G(x)-C_0(x)]
   =\int\frac{(s-a)^2}{x+s}\,\rho(\mathrm ds)
\]
follows by expanding the numerator and using these two moments.
Multiplication by $x$ and monotone convergence give zero on the right
in the limit only if $\rho=\delta_a$. Only the sampled carrier frequencies enter the proof.
\end{proof}

There is a quantitative implication for a proposed physical readout.
If a positive measure with these first two moments has any weight on a
compact interval $I$ separated from $a$ by $d>0$, then
\[
 G(x)-C_0(x)\ge
  \frac{d^2\rho(I)}{(x+a)^2(x+\sup I)} .
\]
Any positive spectral weight separated from the free energy must
therefore change at least one of the first three high-frequency
coefficients, or invalidate the corresponding finite-moment expansion.
The same test applies to a magnetic probe at fixed spatial momentum. Positivity
for a different endpoint or an enlarged set of physical states
requires computing the response of that endpoint with its own source.
The argument tests the specified nonzero frequencies and leaves any
static thermal contribution undetermined.

\begin{corollary}[A magnetic spectral obstruction in the fixed-depth heat-link source]
\label{cor:spex-magnetic-uv}
For every $a,m>0$, each sufficiently large four-torus has a nonzero
spatial magnetic mode whose exact interacting covariance fails the positive spectral
test of Proposition~\ref{prop:spex-photon-uv} for all sufficiently
small nonzero real $e$. The weak-coupling interval in this statement
may depend on the torus.
\end{corollary}
\begin{proof}
The two marked external readouts also give rapid frequency decay in
the fixed-depth determinant $\det(I+T_b/m)$. To verify the needed
bound without differentiating a Fourier wave, set $R_b=(mI+T_b)^{-1}$.
In the single trace with the twice marked kernel, use
$R_b=m^{-1}I-m^{-1}R_bT_b$. The identity term integrates the diagonal
of the twice marked kernel. Its absolute value is bounded by the
Gaussian sum over winding displacements with two additional length
factors. The other term is bounded by two Hilbert--Schmidt norms.
A product of first marks, and the two separate first traces, obey the
same type of estimates. All bounds after extracting the two geometric
readouts are independent of the Fourier frequency. The same Gaussian
integration by parts and moment calculation thus apply at fixed
spatial momentum to this interacting observable.

It remains to demonstrate a change in an observable at a
\emph{nonzero carrier frequency}. Take external Euclidean momentum
with carrier and first spatial components equal to $1/R$, set its
other components to zero, and choose the second spatial polarization.
Its spatial wave number is $1/R$. At every positive mass the continuum one-loop coefficient
of the soft inverse response is strictly positive by
Proposition~\ref{prop:closure-heat-maxwell-sign-uv}.
Rotation covariance makes it independent of the chosen direction.
Proposition~\ref{prop:softward-torus-uv} implies that the quadratic
coefficient of the inverse transverse covariance at this mode differs
from the free one for every sufficiently large $R$. Inversion preserves
nonvanishing of the quadratic correction to the covariance. There is
no linear term. Fixed-volume differentiability in $e$ then gives a
nonzero difference throughout a punctured interval about zero.
The same proof permits antiperiodic charged carrier frequencies:
the translated loop grid introduces unit-modulus winding signs in
Poisson summation and leaves its error bound intact.

If the thermal spectral test held for this Hermitian spatial mode, the
frequency moments would force the free covariance at every nonzero
carrier frequency for this fixed $R$ and spatial momentum. This
contradicts the result at the first carrier frequency. The magnetic
mode is a spatial derivative of the transverse potential and is
invariant under small Abelian gauge transformations. Thus the
obstruction is visible with a gauge-invariant photon observable on
these boxes. The conclusion concerns the magnetic two-point function at the
specified finite volumes and couplings.
\end{proof}

\paragraph{The magnetic test at fixed charge in the bulk.}
The finite-box obstruction above allows the charge interval to shrink
as the box grows. One can remove that qualification for the
fixed-width heat link by estimating the first interaction correction
directly under the exact measure. Write \(C_R=Q_R^{-1}\) on mean-zero
transverse fields of period \(2\pi R\), \(P_a=e^{-a\sqrt{Q_R}}\), and
\(W_e(A)=\log\mathcal D_m(P_a A)\) for its real log determinant. All
norms of test fields on a box use its ordinary volume measure.

\begin{theorem}[Bulk photon laws and a magnetic obstruction at fixed weak charge]
\label{thm:bulk-magnetic-uv}
At fixed \(a,m>0\) and sufficiently small real \(|e|\), any sequence
\(R\to\infty\) has a subsequence on which the exact transverse
photon laws converge as random tempered distributions. Their limits
are stationary and obey the two-sided bounds
\eqref{eq:softward-anisotropic-uv} in infinite volume.
For each fixed \(m>0\) there are \(p_*,c,e_*>0\)
such that, for \(0<|e|<e_*\), every such limit satisfies
\[
 G(p)\preceq (1-c e^2)|p|^{-2}I
 \qquad (0<|p|<p_*)
\]
almost everywhere on the transverse space. A spatially smeared
magnetic observable of each limit fails the unsubtracted positive
energy spectral representation. Both \(e_*\) and the momentum
neighborhood are independent of the chosen subsequence.
\end{theorem}
\begin{proof}
Periodize each Schwartz test field on the box and project out the
longitudinal and constant modes. The upper Hessian-based covariance
bound of Theorem~\ref{thm:softward-allmass-uv} also holds after a
real linear tilt of the density, since the tilt leaves its Hessian
unchanged. Its proof by finite Maxwell projections passes to the
tilted moments at fixed volume by Gaussian domination. Integrating
the bound on the second derivative of the log moment generating
function gives a centered sub-Gaussian bound with variance proxy
\[
 \frac{C}{(2\pi R)^4}
 \sum_{p\in R^{-1}\mathbb Z^4,\ p\ne0}
          |\widehat f(p)|^2 |p|^{-2}.
\]
The untilted law has mean zero by translation invariance and removal
of the constant mode. Lattice shell counting bounds this proxy
uniformly by a continuous Schwartz seminorm squared. The
characteristic functionals are thus equicontinuous at the origin.
A diagonal extraction on a countable dense set, extended by
equicontinuity, gives a continuous positive-definite limit with
value one at zero. The characteristic-functional continuity theorem
on tempered distributions gives the asserted limit law
\cite{BulkLevyContinuity}. The exponential bound implies uniform
integrability of every fixed smeared moment. Translation invariance
and the covariance inequalities pass to the limit. For the free
comparison multipliers, Riemann sums near \(p=0\) are controlled by
the integrable four-dimensional singularity \(|p|^{-2}\).
Stationarity and those inequalities give a transverse Fourier
covariance density bounded by the same multipliers almost everywhere.

For the strict low-momentum bound use the all-mass one-loop Maxwell
coefficient of the straight Wilson heat link in
Proposition~\ref{prop:closure-heat-maxwell-sign-uv}. Define the deterministic quadratic form
\(B_R=e^{-2}W_e''(0)\), where the derivative is in the field at
zero field. It is independent of \(e\ne0\). If \(G_R\) denotes the
exact covariance, then for real linear probes one has
\[
 G_R=C_R+C_R\big(\langle W_e''\rangle
                     +\operatorname{Cov}(\nabla W_e)\big) C_R .
\]
This identity follows by translating the Maxwell source and
differentiating its logarithm twice. It includes the ordinary
determinant and the single vacuum normalization. The Hessian bound
\(-\delta Q_R\preceq W_e''\preceq\delta Q_R\), \(\delta=C_m e^2\),
and the Brascamp--Lieb variance inequality give
\[
 \operatorname{Var}(\partial_h W_e)
 \le\frac{\delta^2}{1-\delta}\langle h,Q_R h\rangle .
\]
Indeed the gradient of this observable is \(W_e''h\), and
\(Q_R^{-1/2}W_e''Q_R^{-1/2}\) has norm at most \(\delta\).
Both identities may first be used in the finite projections in the
preceding proof.

It remains to bound the change in the mean Hessian at nonzero
coupling, uniformly in volume. This estimate also works at arbitrarily
small positive charged mass. Set \(B=P_a A\), \(b=P_a h\) and regard
\(\partial_h^2 W_e(tA)\) as a function of \(0\le t\le1\) while
keeping the averaging law fixed at its full coupling \(e\). Its
first derivative at zero is a deterministic linear functional of
\(A\), with mean zero. We bound its second derivative using the
lifted mass-resolvent proof of Theorem~\ref{thm:softward-allmass-uv}.

In each lifted trace keep the two \(h\)-marks and all holonomy
terms exactly as in that proof. Differentiation in \(t\) adds marks
along the actual field \(B\). A first or second derivative of a
Wilson link inserts one or two line integrals of \(ieB\). A first
resolvent derivative is \(-R_v T' R_v\); the second is
\(2R_v T'R_v T'R_v-R_v T''R_v\), with derivatives here in \(t\).
Thus there are finitely many strings. Each has two \(h\)-marks
(possibly winding factors) and two \(B\)-marks. Its inverse kernels
obey the same deterministic, exponential bounds at every real \(tB\).
Every string now contains at least one Gaussian link and two inverse
factors, including the derivatives of an otherwise unmarked winding
resolvent. The identity part of an inverse is its delta kernel.
For \(v\ge m\) the inverse factors together cost at most \(C_m v^{-2}\).

Independently of the base point, the second moment
\(\langle |B_i(y)|^2\rangle\) of each sampled component
is bounded uniformly on the boxes for \(\delta\le1/2\) by the Maxwell
sum with weight \(e^{-2a|p|}|p|^{-2}\). For a fixed sequence of
displacements of total length \(\ell\), Cauchy--Schwarz therefore
bounds the mean absolute product of the two \(B\)-marks by
\(C_a e^2\ell^2\). All other background-dependent kernels in the
string have deterministic absolute bounds, so no statistical
independence is required. The two test marks cost at most
\(C e^2\ell^4\|\nabla h\|_2^2\) after the base-point integration,
including nonzero windings. Absorb the resulting length moment of
order six into the Gaussian and exponential kernels. Periodize one
Gaussian link to impose closure, and use total variation on the other
steps. The resulting bound
\(C_{a,m} e^4 v^{-2}\langle h,Q_Rh\rangle\) is integrable over
the auxiliary mass \(v\). This argument justifies Taylor's integral
remainder and the interchange with expectation, for example first on
a compact mass interval and then by the same domination. Hence
\[
 \big|\langle \partial_h^2W_e\rangle-e^2 B_R[h,h]\big|
       \le C_{a,m}e^4\langle h,Q_R h\rangle .
\]
The estimates are uniform in the finite Maxwell projections. Their
limit gives the full-box expansion, in operator norm on transverse
modes,
\begin{equation}
 \left\|Q_R^{1/2}G_R Q_R^{1/2}-I
          -e^2 Q_R^{-1/2}B_R Q_R^{-1/2}\right\|
       \le C_{a,m}e^4 .
 \label{eq:bulk-uniform-four-uv}
\end{equation}

The soft coefficient of the \emph{negative} log determinant is
strictly positive at each fixed \(m>0\). Rotation covariance
and the soft limit therefore bound the continuum symbol of \(B_R\)
above by a strictly negative multiple of \(|p|^2 I\) throughout
a sufficiently small punctured ball. The Poisson estimate in
Proposition~\ref{prop:softward-torus-uv} bounds the unnormalized
symbol error by \(C e^{-\eta R}\) for some \(\eta>0\), uniformly in
\(p\). Division by \(|p|^2\) costs at most \(R^2\) at nonzero
lattice momenta. Taking \(R\) large and then \(|e|<e_*\) in
\eqref{eq:bulk-uniform-four-uv} yields the strict inequality on
that entire ball. It passes to every subsequential limit against
probes supported there.

Finally, choose a normalized real spatial test mode orthogonal to
its spatial momentum \(\mathbf k\), with smooth transform supported
in a thin annulus \(b\le|\mathbf k|^2\le b+\epsilon\), \(b>0\).
Spatial normalization uses the Fourier \(L^2\) measure. Since its
support avoids zero, this observable is a spatially smeared magnetic
field: the curl admits a smooth inverse on these transverse tests.
Choose \(b\) small enough that \(x+|\mathbf k|^2<p_*^2\) on an
open interval of positive squared carrier frequencies \(x\).
For the smeared scalar covariance \(G_f(x)\), the large-frequency
bounds \eqref{eq:softward-anisotropic-uv} fix the first two terms
of its asymptotic to \(x^{-1}-\bar b\,x^{-2}\), where
\(\bar b=\int|\mathbf k|^2|\widehat f|^2\) in this normalization.
An unsubtracted positive spectral representation in \(x\) would
therefore have total weight one and first moment \(\bar b\), by
monotone convergence. Possible static weight affects only zero
frequency. Jensen's inequality would require
\(G_f(x)\ge (x+\bar b)^{-1}\). On the chosen interval the strict
upper estimate gives
\[
 G_f(x)\le\frac{1-ce^2}{x+b}<\frac{1}{x+\bar b}
 \qquad\hbox{if }0<\epsilon<ce^2 b .
\]
The inequalities hold almost everywhere, which suffices on an open
interval. This contradiction tests positive transfer along the chosen
carrier without assuming a rotation-invariant limiting state.
\end{proof}

This construction gives subsequential real-charge photon laws of
the fixed-width heat link, without asserting thermodynamic uniqueness
or an infinite-volume charged determinant. The spectral obstruction
concerns an actual magnetic observable in the stated mass and charge
range. The same laws retain their massless Euclidean susceptibility.

\paragraph{The relative determinant under a local bulk probe.}
The limiting law also remembers how the action responds to a local
change of the photon field. At fixed \(a,m>0\), let \(h\) be a smooth
real compactly supported transverse four-vector on the Euclidean
carrier. Its integral vanishes componentwise by integration of
\(\partial_j(x_i h_j)=h_i\), so its periodization is an allowed
mean-zero shift. Write \(B=P_aA\), \(b=P_a h\). The Poisson kernel
has unit integral and \(\|b\|_1\le\|h\|_1\).

The field-dependent part of the bulk determinant under this shift has
a direct definition. Write \(T'_U[b]\) for the first derivative of
the Wilson operator along \(b\) at background \(U\).
With the inverse \(R_{m,U}=(mI+T_U)^{-1}\) on the cover, define
\[
 \Delta_h W(B)=
 \int_0^1 \mathrm dt\int_{\mathbb R^4}
       \operatorname{tr}\big(R_{m,B+tb}T'_{B+tb}[b]\big)(x,x)
                \,\mathrm d^4x .
\]
This is a diagonal kernel integral, independent of an absolute
infinite-volume vacuum determinant. For each marked segment of
displacement \(z\), the \(L^1\) norm over its base point of the
line integral of \(b\) is at most \(|z|\|b\|_1\).
The marked link has a Gaussian envelope. The inverse has the
exponential bound \eqref{eq:softward-magnetic-resolvent-uv}
and its stated identity part. These facts give absolute convergence
and the bound \(C_{a,m}|e|\|h\|_1\) uniformly in the real
background. Periodizing the same kernels gives the same uniform
estimate on the tori. There the integral is the derivative formula
for the ordinary log determinant integrated from \(A\) to \(A+h\).

For clarity obtain the readout jointly with the photon law. The
expectation of each squared finite-order Sobolev norm of \(P_a A_R\)
on a fixed compact set is bounded by summing polynomials in \(|p|\)
times \(e^{-2a|p|}|p|^{-2}\). Local compact embedding gives tightness
in the topology of smooth convergence on compact sets. The limiting
readout is \(P_a A\) in mean square against tests; approximate its
multiplier by smooth Fourier functions supported away from zero and
infinity, then remove these restrictions by the covariance bounds.

The relative functional passes to this limit. First approximate \(b\)
in \(L^1\) by a compactly supported smooth field. The error in the
relative action is bounded by \(C_{a,m}|e|\) times the \(L^1\)
error, by integrating the derivative along the difference of the
shifts. This argument is valid on each box. With a compact shift,
long displacement tails in the marked trace are uniformly controlled
by the exponential kernels. Windings then vanish in the bounded-length
part for sufficiently large boxes. Local continuity of the inverse
kernel under changes of the background follows from
\(R_U-R_V=-R_U(T_U-T_V)R_V\). In this product the difference of the
two link phases tends uniformly to zero on chords of bounded length
within each compact set when \(U-V\) tends locally uniformly to zero.
The remaining kernels have the same deterministic exponential tail.
This also proves the limiting \(L^1\) approximation bound by passage
from compact shifts, or directly by truncation in the diagonal
trace. Thus the distributional limits retain the relative determinant
with its full dependence on the sampled field.

For each limiting photon law \(\mu\) and bounded continuous cylinder
observable \(F\),
\[
 \begin{split}
 \int F(A+h)\,\mathrm d\mu(A)
 &=\int F(A)\exp\!\left[
   \langle A,-\Delta h\rangle-\tfrac12\|\nabla h\|_2^2\right.\\[-1mm]
 &\hspace{28mm}\left.{}+\Delta_{-h}W(P_aA)\right]\,\mathrm d\mu(A).
 \end{split}
\]
On the torus this follows by translating the Gaussian field in
its exact determinant-weighted law. The relative action is uniformly
bounded, and the Gaussian linear term has a uniform exponential
moment by the sub-Gaussian estimate. The identity therefore survives
the limit at fixed weak charge for every positive \(m\). It fixes the
response to these compact transverse probes and leaves uniqueness of
the infinite-volume state open.

\paragraph{A positive photon response with a geometric spatial readout.}
A solvable polarization model shows where a positive spectral band
can enter while the photon remains massless. Work in one rest frame
on the selected plane with two transverse photon polarizations. Let
$k=|\mathbf k|>0$. A harmonic extension of spatial boundary data
to normal distance $r$ obeys $(\partial_r^2-k^2)u=0$ with bounded
solution $u(r)=h_r(k)u(0)$, $h_r=e^{-rk}$. Couple a neutral transverse
polarization oscillator $X$ of frequency $\Omega_0>0$ to this
magnetic readout, at the same carrier time. In a real transverse
mode basis, with the curl phase absorbed in $X$, its Hamiltonian is
\[
 H_k=\tfrac12\left(\Pi_A^2+k^2A^2+\Pi_X^2
                   +[\Omega_0 X-g k h_r(k) A]^2\right).
\]
Here $\Omega_0$ denotes a material frequency, independently of the
path postulates. The square fixes the direct quadratic photon
contact and the strength of the mixing. Oscillator theories of
polarization provide the physical comparison~\cite{SpectrumHopfield}.
This example has neutral polarization, two propagating mixed modes,
and no charged Dirac determinant. We compare each mode's vacuum
energy with the corresponding uncoupled vacuum. The zero spatial
photon mode is fixed in a periodic box. No absolute thermal trace
over an infinitely degenerate material band is required.

\begin{proposition}[A massless photon with positive interacting spectral weights]
\label{prop:spex-polarization-uv}
At every real $g$ and every fixed $r>0$, the above action has a
transverse massless branch, a positive spectral measure for the
interacting photon, and finite bare vacuum coefficients per unit
volume at every positive perturbative order.
\end{proposition}
\begin{proof}
For a Euclidean squared carrier frequency $x$, integration of $X$
gives the exact inverse response
\[
 G_r(x,k)^{-1}=x+k^2+g^2 k^2h_r^2\frac{x}{x+\Omega_0^2}.
\]
Let $\lambda_\pm$ be the squared normal frequencies, in increasing
order. Their sum is $\Omega_0^2+k^2+g^2 k^2h_r^2$, and their product
is $\Omega_0^2 k^2$. Thus $\lambda_-/k^2\to1$ and
$\lambda_+\to\Omega_0^2$ as $k\to0$. For $gk h_r\ne0$, the
characteristic polynomial is strictly negative at $\Omega_0^2$ and
positive at zero. Its roots straddle $\Omega_0^2$. The partial fractions
\[
 \begin{aligned}
 G_r(x,k)&=\frac{Z_-}{x+\lambda_-}+\frac{Z_+}{x+\lambda_+},\\
 Z_-&=\frac{\Omega_0^2-\lambda_-}{\lambda_+-\lambda_-},
 \qquad Z_+=1-Z_-
 \end{aligned}
\]
give two strictly positive weights and $Z_-\to1$ at small $k$. At
$g=0$ take the decoupled limit. The coefficient of $-x^{-2}$ at
large $x$ is $k^2+g^2 k^2h_r^2$. This is precisely where the physical
readout differs from the free frequency moments in
Proposition~\ref{prop:spex-photon-uv}.

The bare shift of the vacuum energy per mode relative to the same
two free oscillators is
\[
 \tfrac12\!\left[\sqrt{(k+\Omega_0)^2+g^2 k^2h_r^2}-(k+\Omega_0)\right]
 \le \frac{g^2 k^2h_r^2}{4(k+\Omega_0)} .
\]
It is nonnegative, and its integral in $\mathrm d^3\mathbf k$ is
finite at both endpoints. One can also examine individual contractions
before combining the quadratic contact with two mixing vertices.
Apart from the fixed powers of $g$, the two possible insertions into
a connected Gaussian vacuum ring have absolute size bounded by
\[
 U=\frac{k^2e^{-2rk}}{\omega^2+k^2},\qquad
 U\,\frac{\Omega_0^2}{\omega^2+\Omega_0^2}\le U .
\]
Every interacting vacuum ring contains an insertion, $0\le U\le1$,
and $\int U\,\mathrm d\omega\,\mathrm d^3\mathbf k=C r^{-4}$
with a finite numerical constant. This establishes absolute
convergence at all finite orders, including the frequency integration,
for the individual rings. Products of disconnected rings require the
corresponding higher inverse moments if $r$ fluctuates. For a
prescribed common probability mixture of fixed-depth states with all
inverse moments these bounds integrate, and the resulting two-point
measure is a positive mixture. The two discrete poles belong to each fixed depth. A volume limit with a dynamical
vacuum-dependent depth distribution requires its own measure estimate.
\end{proof}

Here spatial harmonic propagation supplies ultraviolet decay, while
ordinary time evolution fixes the positive spectral weights. Coupling
the normal transfer to all four Euclidean components, as in the
relative Wilson source, changes the temporal moments even when its
bare graphs remain finite. Deriving which propagation and contact
the GCM boundary state induces is therefore part of the physical
masslessness and positivity problem.

\paragraph{A mobile charge with positive transfer and retained spatial displacements.}
A time-local comparison separates two measurements on the same open line.
The positive Hamiltonian determines its energy distribution. The spatial
charge profile determines which hard photon momenta its interactions sample.
Work at one locked support with Euclidean transfer length $t\geq0$ along
the continued carrier. Use the three spatial coordinates on that support.
At a common depth $0<r\leq1$, displace a spatial endpoint by the
heat law with measurable width parameter $\sigma(r)^2>0$ and use the harmonic
readout at distance $r$. Their convolution is a probability density
$w_r$ with characteristic function
\[
 F_r(k)=\exp[-\sigma(r)^2|k|^2-r|k|],\qquad k\in\mathbb R^3 .
\]
Both kernels act within one spatial slice. The same $w_r$ samples the
charge density, the transverse potential and external probes.

Consider one spinless mobile charge with mass $m>0$, rest energy $M>0$
and real coupling $e$. The neutral sector has free transverse photons.
Write $H_f=\mathrm d\Gamma(|k|)$ and $P_f=\mathrm d\Gamma(k)$ on their
bosonic Fock space, with vacuum $v$. Take real measurable orthonormal
transverse polarizations $\epsilon_\lambda(k)$, $\lambda=1,2$, and set
\[
 \begin{aligned}
 g_i(k,\lambda)&=
 \frac{F_r(k)\epsilon_{\lambda,i}(k)}
 {(2\pi)^{3/2}\sqrt{2|k|}},\qquad
 A_i=a(g_i)+a^*(g_i),\\
 c_r&=\frac{1}{2(2\pi)^3}\int_{\mathbb R^3}
        \frac{|F_r(k)|^2}{|k|^2}\,\mathrm d^3k .
 \end{aligned}
\]
The Coulomb energy $e^2c_r$ is half the squared longitudinal
electric-field norm obtained by solving Gauss's law for $e w_r$.
In the fibre of conserved total spatial momentum $P$, the
Pauli--Fierz kinetic square~\cite{HiroshimaPF2007} gives
\begin{equation}
 H_{e,r}(P)= M+H_f+e^2c_r+
     \frac{1}{2m}\sum_{i=1}^3(P_i-P_{f,i}-eA_i)^2 .
 \label{eq:sp-uv-mobile-transfer}
\end{equation}
The square denotes the sum of squared norms in the quadratic form.
In particular its $A_i^2$ contacts retain the vacuum contraction.
In the two-photon kernel their two independent marks have factors
$F_r(k_1)F_r(k_2)$ even at equal transfer lengths. Replacing the two
marks by one displacement of their sum would change this Hamiltonian.

\begin{proposition}[Interacting energy measure of the mobile charge]
\label{prop:sp-uv-mobile-transfer}
At every $P$ and real $e$, the vacuum matrix element of
$\exp[-tH_{e,r}(P)]$ is the Laplace transform of a probability
$\rho_{P,e,r}$ on $[M,\infty)$. Set
$\gamma_r=(6\pi^2)^{-1}\int_0^\infty s|F_r(s)|^2\,\mathrm ds$.
Its mean and variance are
\[
 \begin{aligned}
 \langle E\rangle&=M+\frac{|P|^2}{2m}
             +e^2\left(c_r+\frac{3\gamma_r}{2m}\right),\\
 \operatorname{Var}(E)&=\frac{e^2\gamma_r |P|^2}{m^2}
                   +\frac{3e^4\gamma_r^2}{2m^2}.
 \end{aligned}
\]
At nonzero charge the measure has more than one energy in its
support, including at $P=0$. If $r$ is a conserved boundary label
with probability $\nu$ having all inverse moments, the same
statements of positivity and nontrivial support hold after integration
against $\nu$. Every finite-order coefficient of this mixed open
transfer is absolutely UV finite, including its contacts.
\end{proposition}
\begin{proof}
Each $g_i$ and $g_i/\sqrt{|k|}$ belongs to the one-photon space.
The annihilation bound
$\|a(g_i)u\|\leq\|g_i/\sqrt{|k|}\|\,\|H_f^{1/2}u\|$ and the
creation norm identity bound $A_i$ on the field-energy form domain.
By the triangle inequality in both directions, the form norm in
\eqref{eq:sp-uv-mobile-transfer} is equivalent, for fixed $e,r,P$,
to the intersection norm of $D(P_f)$ and $D(H_f^{1/2})$, where
$D(P_f)$ is the common domain of its three components.
This gives a closed densely defined form with lower bound $M$.
The spectral theorem applied to its unit vacuum vector gives the
stated probability. The vector is in the operator domain since the
form applied to it pairs against a finite-particle vector with all
momentum moments. Transversality gives
$\sum_i P_{f,i}a^*(g_i)v=0$, and angular integration gives
$\langle g_i,g_j\rangle=\gamma_r\delta_{ij}$. After subtracting the
vacuum component, $H_{e,r}(P)v$ has one-photon part
$-(e/m)\sum_iP_i a^*(g_i)v$ and two-photon part
$(e^2/2m)\sum_i a^*(g_i)^2v$. Their squared norms add and give
the variance. For the direct integral in $r$ the source vector is the
same unit vector $v$ in each fibre. Its spectral law is the mixture
of the fibre laws; its variance includes the variance of their means.

For the coefficient statement expand $W=H_{e,r}-H_{0,r}=eV_1+e^2V_2$
on finite-photon vectors with all polynomial momentum moments. Each
free transfer is multiplication by a number between zero and one
on these wave functions. The $P_f$ factors cost a polynomial in
the photon momenta. Each creation or annihilation carries its own
$g_i$. Thus for $j$ marks the norms of all words of $V_1,V_2$ and
interleaved free transfers acting on $v$ are bounded by $C(1+r^{-K})$
for some finite $C,K$ depending on $j,P,m$. Indeed, after using the
creation and annihilation formulas on at most $2j$ photons, the
required weighted one-photon squared norms are bounded by constants
times $\int_0^\infty(1+s)^L s e^{-2rs}\,\mathrm ds$ for a finite $L$;
$c_r$ costs at most a constant times $r^{-1}$. These absolute bounds
are uniform over the closed ordered time simplex. Iterating Duhamel's
identity with the exact contraction leftmost gives the free words
through $j$ factors of $W$. The remainder contains $j+1$ factors and
is bounded for $|e|\leq1$ by
$C_{j,t}(1+r^{-K_j})|e|^{j+1}$. The identity follows by differentiating
the product of exact and free semigroups on the indicated vectors,
which remain in both operator domains after each free word.
Integration over $\nu$ therefore gives the claimed finite-order
expansion at zero charge even on simultaneous short intervals.
\end{proof}

The external gauge test uses this same spatial source. In position
space couple a smooth external potential by substituting
$p-eA_r(x)-e(w_r* b)(t,x)$ in the kinetic square and adding
$e(w_r*b_0)(t,x)$; the convolution is spatial. In real transfer
time the change $b\mapsto b+\nabla\lambda$,
$b_0\mapsto b_0-\partial_t\lambda$ conjugates the evolution by the
endpoint operators $\exp[i e(w_r*\lambda)(t,x)]$. This follows by
differentiating the phase in the kinetic and time derivatives.
A constant scalar probe thus shifts the entire charged spectral
measure by $e b_0$, including at nonzero coupling. The neutral
photon sector is unchanged by that probe.

Extending the charged autocorrelation evenly in Euclidean length
gives the energy-frequency coefficient
$\int 2E/(\omega^2+E^2)\,\rho_{P,e,r}(\mathrm dE)$. Its product with
$\omega^2$ tends to $2\langle E\rangle$ by monotone convergence:
the geometry damps momenta entering spatial marks, while this
positive transfer fixes the frequency tail of the open line.
The measure proved here is a fixed-carrier energy measure. An
invariant mass spectrum, a charged mass-shell atom and a chiral
pair-producing determinant require additional dynamics. In this
one-particle sector there are no closed matter vacuum loops to
which the spatial coefficient proof could be applied.

\subsection{A positive photon coupled to a propagating geometric depth}
\label{subsec:closure-uv-radial-light}
There is a way to keep the radial degree of freedom dynamical and to
test the photon with the same Hamiltonian that weights the radial
histories. At a fixed sector plane let $0<r\le r_*$ be the proper
depth in the charged normal radial ground state used above. Assume
its Hamiltonian has compact resolvent, a simple positive ground
state $u_0$ and all inverse moments in that state. Consider the
magnetic field harmonically extended from a spatial slice at two
points of the same normal segment, its midpoint and its endpoint.
The difference has symbol
\[
 d_r(k)=e^{-rk/2}-e^{-rk},\qquad k=|\mathbf k|.
\]
It is fixed by solving the bounded spatial harmonic extension at
the two depths. At long wavelengths the measured fields agree.
At high spatial momentum both readouts decay. Using the squared
difference as a magnetic energy is a specific choice of the
constitutive pairing. The normal metric by itself does not select it.

In a spatial three-torus remove the constant photon mode. On the
product of the radial space and the transverse photon space use the
free photon Hamiltonian $H_f$, referred to its free vacuum, and the
nonnegative quadratic form
\begin{equation}
 H=H_{\rm rad}-E_0+H_f+
       \frac{e^2}{2}\int\!
       \big|d_r(|\nabla|)\,\nabla\times A(\mathbf x)\big|^2
            \,\mathrm d^3\mathbf x .
 \label{eq:closure-uv-radial-light}
\end{equation}
Keep the vacuum contraction of the magnetic square. The
interaction is multiplication by a nonnegative function in the
joint configuration representation. Its sum with the closed free
form is closed on the intersection of form domains. Radial vectors
supported away from zero and smooth finite-photon vectors give a
dense domain. Since the free radial and photon operators have compact
resolvent at this spatial volume, form domination implies the same
for $H$. The radial and free photon Euclidean transfers preserve
positivity. Their Feynman--Kac weight is multiplied by the real
exponential of minus the time integral of this magnetic energy.
Consequently there is a nonnegative ground vector with nonzero
overlap with the product of the two free ground vectors.

\begin{proposition}[Unit-speed photon spectral weight with fluctuating depth]
\label{prop:closure-uv-radial-light}
For every finite real $e$, let $G(\omega,k)$ denote the
ground-state covariance of a normalized real transverse spatial mode
in the state obtained by long Euclidean evolution of that product
vector. In units of the limiting speed on the selected carrier, set
$\delta_k=\min(1,r_*^2 k^2/4)$. At nonzero frequency,
\begin{equation}
 \frac1{\omega^2+(1+e^2\delta_k)k^2}
   \ \le\ G(\omega,k)\ \le\ \frac1{\omega^2+k^2}.
 \label{eq:closure-uv-radial-sandwich}
\end{equation}
Its positive measure in squared energy has total weight one.
If $Y$ denotes squared energy divided by $k^2$ in this measure,
then for $0<\varepsilon<1$,
\[
 \Pr(|Y-1|\ge\varepsilon)
       \le \frac{(2+\varepsilon)e^2\delta_k}{\varepsilon^2}.
\]
The estimate is uniform in the spatial volume. It implies
concentration at unit speed along any sequence of nonzero momenta
tending to zero. At $e\ne0$ and finite $k>0$ the measure has more than one point
in its support for a radial Schr\"odinger coordinate with absolutely
continuous position measure.
\end{proposition}
\begin{proof}
Use a symmetric time slab of duration $2T$, with free ground vectors
at its ends. The radial path before weighting is its stationary
ground-state process. Conditionally on that entire path the photon
action is Gaussian. For a mode $k$ its precision on the time line
is $-\partial_t^2+k^2+e^2k^2 d_{r(t)}(k)^2
\mathbf 1_{[-T,T]}(t)$. Conditional normalization contributes its
ordinary relative determinant to the power $-1/2$. Since
$0\le d_r(k)\le\min(1,r_*k/2)$, inverse-operator order yields the
two free comparisons in \eqref{eq:closure-uv-radial-sandwich}
as quadratic forms on time probes, even after averaging with
this path-dependent determinant. Long-slab evolution converges to
the normalized ground projection of the product vector. For
unbounded linear probes this limit follows first for bounded
truncations: the conditional Gaussian variance is uniformly
bounded by the free one, hence fourth moments of each smeared probe
are uniformly bounded. The limiting covariance has the same form
bounds and is stationary. The spectral theorem gives its
nonzero-frequency transform in the form $\int(\omega^2+s)^{-1}
\rho_k(\mathrm ds)$ with positive measure (a transition with energy
$E>0$ contributes $2E$ times its squared matrix element at $s=E^2$).
The time-probe upper bound also rules out a static delta component.
The frequency inequalities hold everywhere away from zero by continuity.

Put $x=\omega^2$. Monotone convergence of $xG$ gives
$\int\rho_k=1$. That of $x^2(1/x-G)$ gives
$1\le\mathbb E Y\le 1+e^2\delta_k$ after division by $k^2$.
At $\omega=k$ the upper bound gives
$\mathbb E[(1+Y)^{-1}]\le1/2$. Consequently
\[
 \mathbb E\frac{(Y-1)^2}{Y+1}
 =\mathbb E Y-3+4\mathbb E[(1+Y)^{-1}]
 \le e^2\delta_k .
\]
The function on the left inside the expectation is at least
$\varepsilon^2/(2+\varepsilon)$ on $|Y-1|\ge\varepsilon$.
This proves the concentration estimate, without interchanging a
thermal or zero-frequency mode with a particle pole.
\end{proof}

There is also a direct test that the interacting spectrum at fixed
momentum has changed. Select one real transverse mode with coordinate
$y$, in the ordinary $L^2(\mathrm dy)$ oscillator representation.
The Hamiltonian takes the form
\[
 H=-\tfrac12\partial_y^2+\tfrac12 b(r)y^2+H_\circ,\qquad
 b(r)=k^2(1+e^2d_r(k)^2),
\]
where $H_\circ$ acts only on the remaining variables and includes
the fixed vacuum reference. Let $\Psi$ be the chosen ground vector,
of energy $E_{\rm g}$. If the measure had only one positive
frequency $\lambda$, absence of a static component would make
$y\Psi$ an eigenvector with energy $E_{\rm g}+\lambda$. Subtract
$y$ times the ground equation from this eigenvector equation, in
the weak sense on compactly supported smooth tests in $y$.
It gives $-\partial_y\Psi=\lambda y\Psi$ and hence
$\Psi=e^{-\lambda y^2/2}\psi$ with $\psi$ independent of $y$.
Substitution back into the ground equation forces
$(b(r)-\lambda^2)\psi=0$ by the coefficient of $y^2$.
For $e\ne0$ and $k>0$ this analytic function of $r$ is nonconstant.
Its level sets have zero radial measure, contradicting unit norm.
Thus the finite-momentum spectral spread and the sharp limiting
speed are properties of the same interacting state.

The complete determinant is finite on each slab without restricting
the photon momenta. In the free continuous-time covariance a
normalized transverse mode has diagonal $1/(2k)$. The trace of the
positive relative perturbation with its $e^2$ removed is bounded,
by
\[
 C_V\int_{-T}^T(1+r(t)^{-4})\,\mathrm dt ,
\]
since $\sum_{\mathbf k\ne0} k e^{-r k}\le C_V(1+r^{-4})$.
The free expectation of this trace is finite. The ordinary relative
determinant therefore exists almost surely, and its inverse square
root is positive and bounded by one. For fixed order, the free
Gaussian contractions of the time-integrated magnetic square have
absolute bounds by constants times moments of the same trace.
For example, one can diagonalize the positive trace-class covariance
of the measured field on the slab; each ring is bounded by the
corresponding power of its trace. In the momentum and time basis
every two-mode contraction for real normal polarizations has kernel
$k d_{r(t)}(k)d_{r(u)}(k)e^{-k|t-u|}/2\ge0$.
H\"older's inequality and stationarity bound every time-ordered
product of the inverse radial powers by one higher inverse moment
of the initial state, uniformly when insertion times coincide.
Taylor remainders of the nonpositive exponential are bounded by
the next moment at each finite order. This establishes the bare
finite-order slab amplitudes of this magnetic sector and their
full vacuum normalization. Its positive spectra use genuine
propagating normal histories. Introducing a charged Dirac
determinant with this pairing and deriving the same coupling from
the chiral GCM action remain additional physical requirements.

\subsection{A positive transfer with a neutral magnetic current that creates charged pairs}
\label{sec:pairferm-uv}
A neutral polarization oscillator does not test a closed fermion loop.
A spatial geometric readout can also be tested on a source that
creates a particle and its antiparticle, permits their annihilation,
and includes both operations in the same thermal trace.

Fix a spatial torus of volume $V$ on one selected plane. On its full
reciprocal lattice take two independent species of spin-$1/2$ fermions
$b,d$ with energies $E_p=(m^2+|\mathbf p|^2)^{1/2}$, $m>0$, and
global charge $Q=q_{\rm c}(N_b-N_d)$. The free Hamiltonian $H_0$
is their positive second quantization plus that of the transverse
Maxwell field with its spatial zero mode fixed. All spatial scales
below are positive. A heat displacement of each charged endpoint
has characteristic function $h(p)=e^{-\sigma^2|\mathbf p|^2}$.
For example $b_h(\mathbf x)=V^{-1/2}\sum_p h(p)b_p
 e^{i\mathbf p\cdot\mathbf x}$, with the spin index suppressed.
The independent particle and antiparticle endpoint displacements
are retained at the same vertex. The Maxwell readout $P_r$ is
harmonic extension through a spatial half-space to depth $r$, with
multiplier $e^{-r|\mathbf k|}$. Both operations act at one carrier
time. With the Pauli matrices $\sigma_i$ set $S_i=\sigma_i i\sigma_2$
and define the spin-one pair current and Hamiltonian by
\begin{equation}
 \begin{aligned}
 J_i(\mathbf x)&=\sum_{s,t} b^\dagger_{h,s}(\mathbf x)(S_i)_{st}
                 d^\dagger_{h,t}(\mathbf x)+\mathrm{h.c.},\\
 H_g&=H_0+g\int_V\mathbf J(\mathbf x)\cdot
           P_r(\nabla\times \mathbf A^T)(\mathbf x)\,\mathrm d^3x .
 \end{aligned}
 \label{eq:pairferm-H-uv}
\end{equation}
The vertex preserves $Q$ and is invariant under shifts of the
potential by a spatial gradient. The photon coupling here is a
magnetic transition moment. The global charge label $q_{\rm c}$
does not assert a minimal electric coupling or a local Gauss-law
realization for the two species.

\begin{proposition}[Positive photon weight down to the soft energy scale]
\label{prop:pairferm-uv}
At every finite volume, inverse temperature $\beta>0$ and real $g$,
\eqref{eq:pairferm-H-uv} defines a semibounded Hamiltonian with a
finite Gibbs trace. Let $A$ be a real, unit-normalized, transverse
photon standing mode with nonzero wave number $k$. At the nonzero
carrier frequencies $\nu=2\pi n/\beta$ its connected covariance
$G(\nu,k)$ has a positive measure $\rho_k$ in squared transition energy
$s$. This measure satisfies
\begin{equation}
 \begin{gathered}
 \int \rho_k(\mathrm ds)=1,\qquad \int s\,\rho_k(\mathrm ds)=k^2,\\
 \rho_k((0,\Lambda^2 k^2])\ge 1-\Lambda^{-2}
 \qquad(\Lambda>1).
 \end{gathered}
 \label{eq:pairferm-weight-uv}
\end{equation}
For $g\ne0$ one has $G(\nu,k)>(\nu^2+k^2)^{-1}$. In particular,
nonzero magnetic transition weight occurs at energies of order $k$
or below, uniformly down to the lowest nonzero mode on growing tori.
Every finite-order bare vacuum and proper photon graph is absolutely
ultraviolet convergent. The interaction part of the proper photon
two-point coefficient is $O(k^2)$ at each finite order in the joint
soft limit.
\end{proposition}
\begin{proof}
The canonical anticommutation relations bound each $b_{h,s}(\mathbf x)$
and $d_{h,t}(\mathbf x)$ by
$(V^{-1}\sum_p|h(p)|^2)^{1/2}$. The smeared magnetic field has
one-photon wave functions bounded by a constant times
$k^{1/2}e^{-rk}$ in magnitude in continuum normalization.
Their squared norms and their squared norms divided by the photon
energy are summable on the torus. The creation and annihilation bounds
therefore bound the interaction operator from the free field-energy
form domain to the Hilbert space. Its form has arbitrarily small
relative bound with respect to $H_0$. In particular
$H_g\ge H_0/2-C_{g,V}$. The common form domain embeds compactly,
and the eigenvalue comparison with $H_0/2-C_{g,V}$ gives a finite
Gibbs trace. This positive transfer supplies the usual positive
thermal spectral weights~\cite{SpectrumLaineVuorinen} for charged operators paired with their
adjoints as well as for the Hermitian current.

Choose the matching real standing mode $J$ of the current, normalized
by the spatial $L^2$ metric; the curl rotates the polarization and
interchanges sine and cosine. The chosen photon oscillator of coordinate
$A$ and canonical momentum $\pi$ then couples by $c A J$ with
$c=g k e^{-rk}$. Neither $A$ nor $\pi$ has an equal-time commutator
with $J$. Write $C_0=(\nu^2+k^2)^{-1}$ and let $\chi_{JJ}(\nu,k)$
be the fully interacting connected current covariance. The two
oscillator equations give the exact identity
\begin{equation}
 G=C_0+c^2 C_0^2\chi_{JJ}.
 \label{eq:pairferm-dyson-uv}
\end{equation}
Indeed $[H_g,A]=-i\pi$ and $[H_g,[H_g,A]]=k^2A+cJ$.
Applying the Euclidean oscillator operator to the time-ordered
$AA$ correlation gives the contact $1$ in Fourier space and the
term $-cG_{JA}$. Applying it at the photon insertion in $G_{JA}$
gives $-c\chi_{JJ}$ with zero contact. These identities also follow
by integration on the thermal circle against a smooth periodic
test function. The commutators can first be taken as forms on
finite-particle vectors. The bounded current and the oscillator
energy bound justify the resulting Gibbs two-point identities.

At nonzero frequency $\chi_{JJ}$ is nonnegative by its thermal
spectral representation. It is strictly positive: the matrix
element of $[H_g,J]$ between the free empty state and a two-fermion
state of suitable momenta and spin, with zero photons at both ends,
is the nonzero pair matrix element of $J$ times $E_p+E_{k-p}$.
The interaction commutator contains a photon coordinate and has
zero matrix element there. Thus $J$ does not commute with $H_g$.
For a finite-volume Gibbs state some transition of nonzero energy has
a nonzero matrix element, and its positive weight contributes at every
nonzero $\nu$. At zero coupling and zero temperature, for example,
the free pair contribution in a unit polarization is explicitly
\[
 \chi^{(0)}_{JJ}(\nu,k)=
 \frac{2}{V}\sum_p |h(p)h(k-p)|^2
       \frac{2(E_p+E_{k-p})}{\nu^2+(E_p+E_{k-p})^2}.
\]
The leading factor $2$ is the spin trace.
This is the same current that appears in the first interacting
photon coefficient of \eqref{eq:pairferm-dyson-uv}. In the continuum
its zero-spatial-transfer spectral density can be read out exactly:
\[
 \frac{\mathrm d\rho^{(0)}_{J,0}}{\mathrm ds}=
 \frac{\sqrt{s}\sqrt{s-4m^2}}{8\pi^2}
 e^{-\sigma^2(s-4m^2)}\,\mathbf 1_{\{s>4m^2\}} .
\]
Here $s=4E_p^2$ and $\mathrm ds=8|p|\,\mathrm d|p|$.
The formula concerns the current after factoring out the external
magnetic curl; a spatially constant photon has zero curl.

Since $J$ is bounded at fixed volume, its thermal correlation is
integrable on the time circle and $\chi_{JJ}\to0$ as $|\nu|\to\infty$.
For the Hermitian photon coordinate write
$G=\int(\nu^2+s)^{-1}\rho_k(\mathrm ds)$ at nonzero frequencies.
The measure sums positive transitions of nonzero energy. In fact the
oscillator commutator gives a second exact identity for the measure
$\rho_{J,k}$ of this same current. Between energy eigenstates,
$(s-k^2) A_{lh}=c J_{lh}$, where $s=(E_h-E_l)^2$.
Multiplying by the common thermal transition factor gives
\[
 (s-k^2)^2\,\rho_k(\mathrm ds)=c^2\,\rho_{J,k}(\mathrm ds).
\]
Thus the nonzero current spectrum forces interacting magnetic
transition weight away from the free shell.
Equal-energy matrix elements contribute only to the static mode.
Equation~\eqref{eq:pairferm-dyson-uv} gives
$\nu^2G\to1$ and $\nu^4(\nu^{-2}-G)\to k^2$. Monotone convergence
gives the first two identities in \eqref{eq:pairferm-weight-uv};
Markov's inequality gives the third. Multiplication by $k^2$
converts the photon weight into magnetic-field weight, with the
same relative fraction on the stated low-energy interval.
Since the current response is strict and the mean squared energy is
$k^2$, there is positive photon weight on both sides of $s=k^2$
at nonzero coupling. The argument does not fix an individual pole.

For the graph bound, expand the thermal trace around $H_0$ with
its ordinary time ordering. At fixed volume the Duhamel coefficients
are justified by the preceding form bounds. More explicitly, if
$W=\partial_g H_g$, the norm of $W e^{-tH_{g'}}$ is at most
$C(1+t^{-1/2})$ at bounded $t$ and bounded real $g'$. In the trace
remainder at a fixed order, split the longest heat interval to bound
the trace norm by a finite free-energy trace using eigenvalue comparison.
The other short-interval powers $t^{-1/2}$ are integrable on the
ordered simplex. This gives the finite-order expansion at $g=0$.

Each closed fermion cycle in a contraction has at least two vertices:
a single creation or annihilation pair has zero free expectation.
At fixed photon transfers each free charged line is bounded by
$(\omega^2+m^2)^{-1/2}$ with a translated or reversed fermionic
frequency. Any two such factors have uniformly bounded summed product
by Cauchy--Schwarz; use the mass bound on the others. In the spatial
charged sum one factor $|h(p+p_0)|^2$ suffices, uniformly in the
photon transfers. Each remaining photon line with momentum
$(\eta,\mathbf l)$ has an absolute factor bounded by
\[
   \frac{|\mathbf l|^2 e^{-2r|\mathbf l|}}{\eta^2+|\mathbf l|^2}.
\]
After contracting the charged cycles to vertices, take the photon
momenta outside a spanning tree as the independent variables.
For each of these lines the sum over $\eta$, divided by $\beta$,
is $e^{-2r|\mathbf l|}|\mathbf l|
\coth(\beta|\mathbf l|/2)/2$; its normalized spatial sum is finite.
The tree factors are bounded by one. The same bounds hold for
frequency integrals at zero temperature. The normalized loop
estimates are uniform over growing tori and $\beta$ bounded away
from zero, by lattice shell counting. They cover all simultaneous
short-time and hard-momentum regions of these graphs. Each external
photon retains its own curl factor, proving the proper two-point
bound. The spatial and frequency sums run over their full prescribed
lattices, with the Maxwell zero mode fixed. No ultraviolet counterterms
are needed.
\end{proof}

\begin{corollary}[The thermal weight at a common fluctuating depth]
\label{cor:pairferm-cusp-uv}
Use $\sigma^2=\kappa r$, $\kappa>0$, $0<r\le1$, and one normalized
boundary law with density bounded near zero by
$C r^4\exp(-A_*/r^3)$, $A_*>0$, as on the previously derived normal
surface. At fixed $\beta,V$ there is an interval of nonzero real
couplings about zero on which the complete common-depth partition
\[
 Z_{\beta,g}=\int \operatorname{Tr} e^{-\beta H_g(r)}\,\mathrm d\nu_0(r)
\]
is finite. The resulting positive state obeys
\eqref{eq:pairferm-weight-uv} with the same constants.
\end{corollary}
\begin{proof}
In the real normalized photon mode basis let $J_j$ be the
corresponding spatial projections of the pair current, including
the polarization rotation. Write $k_j$ for the wave number and
$a_j$ for the annihilation operator. For any unit form vector
$\psi$, commuting the photon and fermion factors and applying
Cauchy--Schwarz yields
\[
 \begin{aligned}
 |\langle H_g-H_0\rangle_\psi|
 &\le \sqrt{2}|g|\langle H_\gamma\rangle_\psi^{1/2}
             \left(\sum_j e^{-2rk_j}\|J_j\psi\|^2\right)^{1/2},\\
 \sum_j\|J_j\psi\|^2
 &\le\int_V\sum_i\|J_i(\mathbf x)\psi\|^2\,\mathrm d^3x
 \le C V\sigma^{-6}.
 \end{aligned}
\]
Here $H_\gamma=\sum_j k_j a_j^\dagger a_j$, and the second line
uses spatial Parseval and the anticommutation bounds. The normalized
heat sum $V^{-1}\sum_p e^{-2\sigma^2|p|^2}$ is $O(\sigma^{-3})$
for bounded widths, uniformly on cubic tori with a fixed lower
bound on side length. Vacuum and all fermion number sectors are
included. Consequently, if $H_{\rm fermi}$ is the free fermion
Hamiltonian,
\[
 H_g\ge H_{\rm fermi}+\tfrac12 H_\gamma-C g^2 V\kappa^{-3}r^{-3}.
\]
Eigenvalue comparison bounds the full thermal trace by the trace
of the displayed free comparison times
$\exp(C\beta V g^2\kappa^{-3}r^{-3})$. The strict margin
$C\beta V g^2\kappa^{-3}<A_*$ suffices for absolute depth integration.
The remaining margin also controls every fixed inverse power of the
width generated by a finite set of marked vertices. In particular the
finite-order bare coefficients at zero coupling are integrable.

The normalized thermal depth law is proportional to
$\operatorname{Tr}e^{-\beta H_g(r)}\,\mathrm d\nu_0(r)$. Momentum
translation symmetry makes the mean of every nonzero photon mode
zero at each depth. Its complete two-point spectral measure is
therefore the positive mixture of the fibre measures with this
one vacuum-weighted law. Both moment identities hold at every depth
before integration, so the stated concentration bound follows exactly.
This constructs the integrated state at fixed volume. Uniform
selection of a limiting depth law as $V$ or $\beta$ increases
is a separate question.
\end{proof}

This comparison places pair creation, absolute bare graph bounds, and
positive interacting photon and charged spectral weights in the same
source. The low-energy fraction in \eqref{eq:pairferm-weight-uv}
establishes gaplessness along finite-volume modes. An isolated
Lorentz-invariant photon pole, minimally coupled electric charges and
the stationary PE-CGG-HOUQFT boundary state each require further
information from the physical action.

The same calculation gives an action-level test when the prospective
physical charged operator has a photon contact. Suppose the Hamiltonian
in a real, unit-normalized transverse mode can be written
$\pi^2/2+k^2 A^2/2+H_{\rm m}(A)$ with canonical momentum $\pi$,
where the last term may act on all other modes and matter variables,
but is independent of $\pi$. Assume its twice differentiated forms
and the following thermal commutators are integrable. The positive
nonstatic spectral measure of $A$ then obeys
\[
 \begin{aligned}
 \int\rho_A(\mathrm ds)
   &=\langle[A,[H,A]]\rangle=1,\\
 \int s\,\rho_A(\mathrm ds)
   &=\langle[\pi,[H,\pi]]\rangle
     =k^2+\langle\partial_A^2 H_{\rm m}\rangle .
 \end{aligned}
\]
Indeed $\pi_{lh}=i(E_l-E_h)A_{lh}$, and expansion of either double
commutator in an energy basis gives the same positive thermal
transition weights as the spectral sum. Hence an upper bound
$Ck^2$ on the right of the second line implies at least
$1-C/\Lambda^2$ of the normalized transition weight below
energy $\Lambda k$. This is a sufficient test on the complete
Hamiltonian contact. A static Ward identity instead compares a
contact with a susceptibility integrated over time and need not
give this moment bound. The linear magnetic coupling above has
$\partial_A^2H_{\rm m}=0$ exactly.

\subsection{A common charged source for Wilson hopping, the Gauss field and a moving normal coordinate}
\label{sec:cl-gauss-uv}
The two-band symbol with hopping amplitude
$t_r(k)=g e^{-\kappa r|k|^2}$ admits a canonical completion in which
the scalar and spatial electric probes couple to the same extended
charge. This construction tests more than the positivity of a single
free hopping matrix. Its normal variable moves even in the interacting
vacuum, its charged particle and hole can be created together, and
Gauss's law acts on this pair. The completion specifies a representation
and a radial preparation on one sector plane.

Take a spatial torus of volume $V$, in reference units, and retain all
its Fourier modes. Work in Coulomb gauge, fix the harmonic photon mode
and use the real transverse Maxwell field in its equal-time Gaussian
representation. Let $N=N_++N_-$ on the fermion Fock space of particles
and holes, of opposite charges $e,-e$ and reference energy $m>0$ per
excitation. The physical vacuum sector has zero net charge. For charged endpoint
sectors we choose a uniform compensating background on the torus.
This is the convention implemented by omitting the Gauss zero mode. Use $0<r<1$ and the
Friedrichs realization of the radial form
\[
 \int_0^1 f(r)\!
       \left(\|\partial_r u\|^2+\frac{j^2}{f(r)^2}\|u\|^2\right)
             \mathrm dr .
\]
Here $f(r)\sim f_*r^4$ with $f_*>0$. The inner product uses
$f\,\mathrm dr$. Take $f$ smooth and positive
away from the tip, including at the outer boundary. The integer $j$ is the Fourier index on the normal circle in the
prepared state. Electric endpoints here act on a separate Fock factor
and preserve $j$. Identifying it
with the total electric charge of a neutral closed history would give
$j=0$, an important distinction below.

Let $P_r$ be the harmonic readout in the three spatial directions,
with multiplier $e^{-r|p|}$. In the position representation
$\rho_p=\sum_\ell s_\ell e^{-ip\cdot x_\ell}$, with $s_\ell=+1$ for a
particle and $-1$ for a hole. Gauss's law with charge density sampled
by $P_r$ gives the positive quadratic form
\[
 V_{\rm C}(r)=\frac{e^2}{2V}\sum_{p\ne0}
                  \frac{e^{-2r|p|}}{|p|^2}\rho_p^*\rho_p .
\]
Each particle's diagonal contribution is included here. On a fixed
particle sector this is a multiplication operator bounded at fixed
$r$; the full positive form has its natural multiplication domain.

The spatial hopping kernel $K_{r,A}(x,y)$ is the periodization, over
$z=y-x$ and its lifts, of the heat density $g_{\kappa r}(z)$ times the
straight Wilson phase of the readout. Here the heat kernel has Fourier
transform $e^{-\kappa r|p|^2}$. Orient the phase as transport from the
hole to the particle. Under a spatial gauge change it acquires the
factor $\exp[ie(\lambda_r(x)-\lambda_r(y))]$, where
$\lambda_r=P_r\lambda$. Let $B^+(K)$ create one particle at $x$ and
one hole at $y$ with this kernel, and let $B^-(K)$ be its adjoint.
The completed, time-local Hamiltonian on this chart is the form sum
\begin{equation}
 \begin{aligned}
 H={}&H_j+H_{\rm ph}+mN+V_{\rm C}(r)\\
    &+g\{B^+(K_{r,A})+B^-(K_{r,A})\}.
 \end{aligned}
 \label{eq:cl-gauss-charge-uv}
\end{equation}
The hopping changes $N$ by two and preserves $N_+-N_-$. Gauge
covariance of the external source matters when $r$ moves. In a sector
of positions $x_\ell$ the charge phase acts by multiplication with
$\exp[ie\sum_\ell s_\ell\lambda_r(x_\ell)]$. External radial parallel
transport therefore replaces $\partial_r$ by
$\partial_r-ie\sum_\ell s_\ell b_r(x_\ell)$, with
$b_r\mapsto b_r+\partial_r\lambda_r$; the temporal scalar couples to
$e\sum_\ell s_\ell(P_r b_0)(x_\ell)$. In real time
$b_0\mapsto b_0-\partial_t\lambda$. One explicit extension from the
transverse chart is to use $b_r=\partial_rP_r\chi$ for the longitudinal
part $\nabla\chi$ of an external spatial potential. Then $\chi\mapsto
\chi+\lambda$ modulo a spatial constant, whose radial derivative is
zero. Direct differentiation of the endpoint phase verifies the
covariance of both derivatives. In the unforced transverse chart
$b_r=0$. This specifies the external gauge completion on the trivial
bundle. A separate bulk dynamics for the radial gauge connection would
require deriving its own Gauss form.

\begin{proposition}[A positive charged transfer with its full Coulomb energy]
\label{prop:cl-gauss-charge-uv}
For each finite volume, $m,\kappa>0$, and every real $e,g$,
\eqref{eq:cl-gauss-charge-uv} defines a semibounded self-adjoint
Hamiltonian. For sufficiently small $|g|$ it has an isolated neutral
ground state, with a nonzero pair component when $g\ne0$ is sufficiently
small. In this ground-state regime its charged endpoint correlations
have positive energy measures.
If $j\ne0$, every eigenvector of finite eigenvalue, in particular this
interacting ground state, has finite inverse radial moments of all
orders.
\end{proposition}
\begin{proof}
For almost every Maxwell field the positive-distance readout is smooth.
The Wilson factor has unit modulus on every lifted chord. Its
periodized Gaussian majorant and Parseval on the torus give
\[
 \|K_{r,A}\|_2^2
 \le \sum_{p}e^{-2\kappa r|p|^2}
 \le C_{V,\kappa}(1+r^{-3/2}).
\]
Here $\|\cdot\|_2$ is the spatial Hilbert--Schmidt norm.
In an orthonormal spatial basis use the annihilators $a_i$ of the
particle and the hole creators $d^*(K_i)$ for the kernel rows.
Cauchy--Schwarz and the fermion creation bound give
\[
 |\langle u,B^+(K)u\rangle|
 \le\left(\sum_i\|a_i u\|^2\right)^{1/2}
     \|K\|_2\|u\| .
\]
These bounds also apply after integrating the photon and radial
variables. In particular, pointwise in $r$ the hopping form is bounded
in absolute value by $\epsilon\langle mN\rangle+
C_{\epsilon,g,m}(1+r^{-3/2})\|u\|^2$ for every $\epsilon>0$.
At $j=0$ one can still use the radial kinetic energy. Integration by
parts on its core gives the weighted Hardy inequality
$\int r^2\|u\|^2\,\mathrm dr\le(4/9)\int r^4\|u'\|^2\,\mathrm dr$
for compact support in $(0,1)$. Use a smooth localization near the
tip, the comparison $f\asymp r^4$, and
$r^{-3/2}\le\eta r^{-2}+C_\eta$. Thus the hopping has arbitrarily
small form bound with respect to the nonnegative diagonal form.
The sum of the diagonal forms is closed on their dense intersection.
The closed-form perturbation theorem now gives the asserted
Hamiltonian. This proof uses no field cutoff.

At $g=0$ the ground vector is the product of the lowest radial
state and the photon and fermion vacua. The radial ground state is
simple and gapped on the finite interval, the first transverse photon
has positive energy on the box, and every charged excitation costs
at least $m$. The positive Coulomb term annihilates the vacuum.
The isolated eigenprojection persists for small $|g|$: the forms are
analytic in $g$ on a common domain, or one can use the min--max bound
on the orthogonal complement and continuity of the resolvent. The
first variation of the ground vector includes the reduced resolvent
applied to $B^+(K_{r,A})$ on the product vacuum. This is a nonzero
pair vector. For example its overlap with constant spatial particle
and hole waves and photon vacuum contains the Gaussian expectation of
a Wilson phase, $\exp[-e^2\operatorname{Var}(\int P_r A)/2]>0$ at
each $r>0$, integrated against the positive heat kernel. The kernel
bound and the ground state's tip behavior put this vector in Hilbert
space (for $j=0$ use Hardy's inequality). A bounded smeared charged
endpoint applied to the interacting ground state has a positive energy
measure by spectral resolution of $H-\inf H$ in the corresponding
charge sector. This argument uses ground-state transfer; the flat
high-momentum reference bands supply no finite thermal trace.

Finally consider a normalized eigenvector $\Psi$ with eigenvalue
$E$ and $j\ne0$. Multiplication by a bounded Lipschitz radial weight
$W$ commutes with every term except the radial derivative, so that
\[
 \langle W\Psi,(H-E)W\Psi\rangle_{\rm form}
    =\int_0^1 f(r)|W'|^2\|\Psi(r)\|^2\,\mathrm dr.
\]
Use $W=\exp[c/(r+\epsilon)^3]$ and the preceding pointwise hopping
bound. Since $|W'|^2\le9c^2 r^{-8}W^2$, the angular potential
dominates the derivative cost as well as the $O(1+r^{-3/2})$
hopping cost near zero whenever $0<3c<|j|/f_*$. Off a fixed tip
neighborhood the weights are uniformly bounded. Splitting the integral
there and letting $\epsilon\downarrow0$ yields
\[
 \int_0^1 f(r)r^{-8}e^{2c/r^3}\|\Psi(r)\|^2\,\mathrm dr<\infty .
\]
This proves the moment claim for the full state with a moving radial
coordinate. Stability alone gives a weaker result: on the exact
quartic chart at $g=0,j=0$, the regular ground solution is constant
to leading order at zero. Its density is proportional to $r^4$ there
and inverse moments of order five or larger diverge.
\end{proof}

This construction exhibits a common Gauss source, charged pair
creation, a positive transfer and geometric control of the
interacting boundary state. Its bounded free bands and canonical
time ordering are additional choices. In particular, the finite
Coulomb quadratic form alone leaves open the joint ultraviolet estimate
for a moving normal coordinate. At fixed widths the occupied onsite
vacuum identifies the precise reference term in the density derivative.
The following calculations keep track of that datum as photons are sewn.

\paragraph{A neutral moving-depth test at twentieth electric order.}
Fixed widths cannot simply be integrated over an arbitrary stationary
normal state. There is a test in the Hamiltonian
\eqref{eq:cl-gauss-charge-uv} which includes the complete Gauss
energy. Work in its physical transverse chart on a spatial torus
of volume $V$, omit the harmonic photon mode, and choose the exact
quartic normal profile $f(r)=r^4$, $0<r<1$, with its Friedrichs
tip and Dirichlet outer endpoint. Take $j=0$, and let $\psi_0$ be
the positive normalized uncoupled radial ground state in
$L^2(r^4\mathrm dr)$. Its regular radial solution tends to a
strictly positive constant at zero. Indeed the regular solution of
$-(r^4\psi')'=E_0r^4\psi$ is constant to leading order, whereas
the independent $r^{-3}$ solution is outside this Hilbert space.

\begin{proposition}[A divergent complete electric coefficient in the neutral normal sector]
\label{prop:physical-neutral-tip-uv}
Prepare the photon vacuum and $\psi_0$ at both ends of a slab of
duration $T>0$. Take the onsite matter vacuum at the initial end
and one particle and one hole, each in its normalized constant
spatial wave, at the final end. Both states have zero total
electric charge. In this matrix element of the fixed Hamiltonian
the coefficient linear in hopping is finite for every real electric
charge $e$. Its coefficients through $e^{18}$ are finite; its
coefficient of $e^{20}$ has infinite absolute magnitude.
Division by the vacuum amplitude preserves this conclusion.
\end{proposition}
\begin{proof}
Choose the sign of the final pair state to make the pair creation
matrix element positive. Differentiation at $g=0$ inserts the one
spatial Wilson hop at some $0<s<T$. The charged particles have
fixed positions after creation in the onsite dynamics. For displacement
$z$ their full Coulomb energy at depth $r$ is $e^2 C_r(z)$ with
\[
 C_r(z)=\frac1V\sum_{p\ne0}
       \frac{e^{-2r|p|}}{|p|^2}(1-\cos(p\cdot z))\ge0 .
\]
This includes both self-energies. In this unforced Coulomb chart
there is no radial gauge source. The Maxwell field is free before
and after the one hop, so the photon vacuum expectation of its
phase is $\exp[-e^2 W_r(z)/2]$. Its exact equal-time transverse
variance is
\[
 W_r(z)=\frac1{2V}\sum_{p\ne0}\frac{e^{-2r|p|}}{|p|}
 \left(|z|^2-|\widehat p\cdot z|^2\right)
 \left|\int_0^1 e^{iu p\cdot z}\,\mathrm du\right|^2 .
\]
The momenta here are all spatial lattice momenta on the chosen
torus. Each lifted chord in the periodized hopping kernel is
included. Integration against the constant endpoint waves converts
the sum over lifts into the ordinary heat law
$G_{\kappa r}(z)\,\mathrm d^3z$ on $\mathbb R^3$.

Let $r_t$ have the stationary radial path law obtained by the
ground-state transform of $H_0$ in the normal coordinate, with
marginal $r^4|\psi_0(r)|^2\mathrm dr$. The radial semigroup with
the nonnegative multiplication potential $e^2 C_r(z)$ has its
Feynman--Kac weight on this same path. One may first bound the
potential and pass monotonically to its form realization.
Consequently, after dropping a fixed positive factor and the
Duhamel minus sign, the coefficient linear in $g$ is exactly
\begin{gather*}
 \int_0^T e^{-2m(T-s)}\,\mathrm ds\,
 \mathbb E\!\int G_{\kappa r_s}(z)e^{-e^2 X_s(z)}\,\mathrm d^3z ,
 \\
 X_s(z)=\tfrac12 W_{r_s}(z)+\int_s^T C_{r_t}(z)\,\mathrm dt .
\end{gather*}
The radial coordinate continues to evolve after the hop.
The nonnegative integral in $X_s$ is finite almost surely, using
$C_r(z)\le C(1+r^{-1})$ and the stationary one-time marginal.
The vacuum denominator at $g=0$ has no electric dependence and
the unperturbed vacuum-to-pair numerator vanishes.

It suffices to bound the transverse contribution from below. At
fixed $r$ average $W_r(z)$ in the centered Gaussian heat law of
covariance $2\kappa r I$. For each $p$, the perpendicular
components of $z$ are independent of its parallel component.
Their squared-length expectation is $4\kappa r$, while the mean
square of the integral in $u$ equals
\[
 2\int_0^1 (1-v)e^{-\kappa r |p|^2 v^2}\,\mathrm dv
 \ \ge\ \frac{c}{1+\sqrt{\kappa r}\,|p|}.
\]
Restrict the positive sum to $r^{-1}\le|p|\le2r^{-1}$.
For all sufficiently small $r$ this shell contains at least
$c_V r^{-3}$ spatial modes, by lattice counting. Each term in the
Gaussian average on the shell is at least $c_{V,\kappa}r^{5/2}$.
Thus $\int G_{\kappa r}W_r\,\mathrm d^3z\ge c r^{-1/2}$.
Jensen's inequality and $X_s\ge W_{r_s}/2$ give, for every
integer $n\ge1$, a lower bound proportional to
\[
   \int_0^{r_0} r^{4-n/2}\,\mathrm dr
\]
on the $n$th moment of $X_s$ integrated in this amplitude. At
$n=10$ the lower bound already diverges logarithmically.
We also need an upper bound to locate the first failure. Dropping
the transverse projector bounds $W_r$ above by the variance from
the scalar covariance with the same equal-time multiplier. Within
half the spatial injectivity radius its kernel obeys
\[
 |K_r(x)|\le C_V+C(4r^2+|x|^2)^{-1}.
\]
Indeed it equals one half of the integral from $2r$ to infinity
of the periodized Poisson kernel minus its constant mode.
For bounded Poisson times the zero image is proportional to
$t/(t^2+|x|^2)^2$. The other images are bounded by a constant
times $t\sum_{\ell\ne0}|\ell|^{-4}$ on that spatial ball.
The large-time integral converges by the nonzero spatial spectral
gap. For a chord inside this ball, integration along its two
parameters gives $W_r(z)\le C_V|z|^2+C|z|/r$. On longer lifts
use $W_r(z)\le C_V(1+r^{-2})|z|^2$ and the Gaussian heat tail.
Consequently $\int G_{\kappa r} W_r^n\,\mathrm d^3z\le
C_{n,V}(1+r^{-n/2})$. All moments of the Coulomb time integral
are finite. For any $b>0$ its majorant
$C\int_0^T(1+r_t^{-1})\,\mathrm dt$ has a finite exponential
moment by radial Feynman--Kac. In fact $r^{-1}$ is infinitesimally
form bounded for the radial kinetic energy by the weighted Hardy
inequality, so $H_0-b r^{-1}$ is bounded below. The resulting
semigroup has a finite matrix element in the radial ground state.
These two upper estimates guarantee finite moments through $n=9$.

The complete first-order hopping amplitude is a Laplace transform
of a finite positive measure in $u=e^2$. Its alternating $n$th
forward difference at zero divided by $h^n$ integrates
$[(1-e^{-hX_s})/h]^n$, which increases to $X_s^n$ as
$h\downarrow0$. Hence the divergence cannot be canceled by
Coulomb insertions at other times or by a rearrangement of
transverse contacts. The real-charge transform itself is bounded
by its value at zero.
\end{proof}

The endpoints here are globally neutral canonical states in the
chosen Gauss chart. They preserve the normal Fourier number
$j=0$. This example obstructs extending the fixed-width theorem
uniformly to the neutral moving normal sector of this action.
Identification of this preparation, or of an alternative
preparation that suppresses it, is a separate question for the
physical chiral branch.


There is an endpoint-independent way to state the one-hop test.
Take any two nonzero nonnegative radial boundary functions $v,w$ in
the same $j$ sector. Define their time-integrated occupation measure by
\begin{gather*}
 \Pi_T(\mathrm dr)=r^4\,\mathrm dr\int_0^T
       e^{-2m(T-s)}v_s(r)w_{T-s}(r)\,\mathrm ds,\\
 v_s=e^{-sH_j}v,\qquad w_s=e^{-sH_j}w .
\end{gather*}
The preceding two bounds on the transverse variance, Jensen's
inequality and the Coulomb exponential estimate show that the
coefficient of $e^{2n}$ in this transition is finite precisely when
$\int r^{-n/2}\Pi_T(\mathrm dr)<\infty$. The tilted radial
semigroup estimate for the Coulomb term holds for these endpoints
by their $L^2$ norms. At $j=0$ the radial semigroup is the radial
part of the Dirichlet heat semigroup of the five-dimensional unit
ball. For positive interior times each evolved nonnegative nonzero
function is continuous and strictly positive at the center, by the
positive heat kernel on the ball. Thus even the middle third of
the time integral forces the twentieth-order obstruction for every
such choice of endpoints. At $j\ne0$ the finite collar estimate
in Proposition~\ref{prop:final-collar-uv} bounds
$r^{-n/2}e^{-(T/2)H_j}$ on the radial Hilbert space for each $n$.
At least one side of the marked time has duration $T/2$ or longer.
Cauchy--Schwarz on the two sides therefore gives finite moments of
every order for this one-hop transition. This calculation uses the
radial generator common to the two constructions.

The neutral-sector obstruction also occurs in the ground-state
projection at first order in the hopping. Retain the radial
ground preparation of Proposition~\ref{prop:physical-neutral-tip-uv}.
At $g=0$ its vacuum vector $v_0$ has energy $E_0$ for every $e$.
The two-particle sector has energy at least $E_0+2m$, including
the positive Gauss energy. By the isolated-eigenprojection
calculation in Proposition~\ref{prop:cl-gauss-charge-uv},
the pair component of the derivative of the normalized ground
vector at $g=0$ is
\[
 -(H_{2,e}-E_0)^{-1} B^+(K_{r,A})v_0 .
\]
Here $H_{2,e}$ is the zero-hopping Hamiltonian restricted to the
one-particle, one-hole sector, including the photon and radial
operators. This is also the corresponding off-diagonal matrix
element of the derivative of the ground projection, tested on the
same constant pair and photon vacuum. The hopping kernel on $v_0$
is square integrable by the spatial Hilbert--Schmidt estimate.

In the vacuum-normalized slab expression put $\tau=T-s$.
Stationarity makes the path expectation depend only on $\tau$.
Extending its integral to $\tau=\infty$ with weight $e^{-2m\tau}$
gives the resolvent matrix element above. Any fixed interval of
positive $\tau$ already gives the infinite tenth moment in $e^2$.
The first nine remain finite in this limit. Indeed the $n$th
moment of $\int_0^\tau(1+r_t^{-1})\,\mathrm dt$ is bounded by
$C_n(1+\tau)^n$: split time into intervals of length at most one,
use convexity of the $n$th power and stationarity, and apply the
unit-interval radial exponential bound. The exponential pair weight
is integrable against that polynomial. Thus $g e^{20}$ is already
divergent for this stationary transition coefficient. This argument
takes the hopping derivative at zero and does not assert divergence
of the exact transition at fixed nonzero hopping.

\paragraph{The short-chord coefficient at fixed hopping.}
The instantaneous pair transition is an observable of the same
Hamiltonian at a prescribed hopping $g$. Let $v_0,w_0$ be the
normalized vacuum and constant pair preparations of
Proposition~\ref{prop:physical-neutral-tip-uv}, with the same neutral
radial ground state $\psi_0$ and photon vacuum. Orient the final
state so the hopping form at $e=0$ is positive. For each fixed real
$e,g$ the polarized semigroup derivative gives
\begin{equation}
 \begin{aligned}
 I_g(e)&=-\lim_{T\downarrow0}
                 \langle w_0,e^{-TH(e,g)}v_0\rangle/T\\
       &=g\int_0^1r^4|\psi_0(r)|^2\,\mathrm dr
            \int G_{\kappa r}(z)e^{-e^2 W_r(z)/2}\,\mathrm d^3z .
 \end{aligned}
 \label{eq:seam-uv-instant}
\end{equation}
Indeed the endpoint vectors belong to the form domain: the pair's
Coulomb energy costs at most $C(1+r^{-1})$ against the displayed
radial density. Orthogonality in particle number removes every
diagonal form, leaving exactly the single hopping matrix element.
The derivative follows by the spectral theorem for a lower-bounded
form and polarization. Dividing the numerator by the vacuum slab
amplitude has the same initial derivative because the denominator
tends to one. The electric variation in what follows is taken
\emph{after} defining this observable.

\begin{proposition}[Neutral-tip charge asymptotic]
\label{prop:seam-uv-tip-log}
Put $u=e^2$, let $Y$ be a three-dimensional Gaussian vector of
covariance $2\kappa I$, and set
\[
 D=\frac{|\psi_0(0)|^2\mathbb E|Y|^{10}}
          {5(16\pi)^{10}}.
\]
For $g\ne0$ the degree-nine Taylor polynomial $P_9$ at $u=0$ of
$I_g(\sqrt u)/g$ exists. The first nonpolynomial term is
\[
 I_g(\sqrt u)/g-P_9(u)
    =\frac{D}{9!}u^{10}\log(1/u)\,[1+o(1)]
       \qquad (u\downarrow0)
\]
in the reference units of the normal chart. The same coefficient,
multiplied by $\int_0^T e^{-2m(T-s)}\mathrm ds$, holds for the
positive one-hop slab integral in
Proposition~\ref{prop:physical-neutral-tip-uv}, including Coulomb
evolution. For its stationary reduced resolvent replace this factor
by $1/(2m)$.
\end{proposition}
\begin{proof}
For a short lifted chord of length $l$ begin with the scalar
equal-time covariance in place of its transverse part. Its kernel
on a ball less than half the injectivity radius equals
$[4\pi^2(4r^2+|x|^2)]^{-1}+O(1)$ uniformly for small $r$.
Indeed the zero image of the Poisson kernel integrates to this
explicit expression, and the remaining images and large-time
integral are bounded as in the preceding proof. Integrating along
the two segment parameters gives the scalar variance
\[
 \frac1{2\pi^2}
 \left[\frac l{2r}\arctan\frac l{2r}
        -\frac12\log\!\left(1+\frac{l^2}{4r^2}\right)\right]
       +O(l^2).
\]
The longitudinal variance removed by the transverse projection is
$V^{-1}\sum_{p\ne0}e^{-2r|p|}(1-\cos(p\cdot z))/|p|^3$.
Spatial shell counting bounds this term by $C(1+|\log r|)$.
It follows that
$\sqrt r\, W_r(\sqrt r Y)/2\to |Y|/(16\pi)$,
uniformly on compact sets of $Y$. The previously proved estimates
also give $\mathbb E_Y W_r(\sqrt r Y)^n\le C_n r^{-n/2}$ for
$0<r\le1$ and every fixed integer $n$.

Let $A=W_r(\sqrt rY)/2$ under the normalized radial and heat-step
law of~\eqref{eq:seam-uv-instant}. Changing variables to $r=v/x^2$
in the tail integral yields
\[
 \lim_{x\to\infty}x^{10}\Pr(A>x)
    =|\psi_0(0)|^2\int_0^\infty v^4
       \Pr\bigl(|Y|>16\pi\sqrt v\bigr)\,\mathrm dv=D .
\]
Dominated convergence uses the bounded ground function and
$v^4\min(1,C_n v^{-n/2})$ for any $n>10$. For the positive
one-hop time integral the extra Coulomb contribution to $A$ has
a finite moment of order greater than ten, as in the preceding
radial semigroup estimates. In the stationary integral the
corresponding moment grows at most polynomially in its time
duration, which is integrated against $e^{-2m\tau}$. For any
$0<\epsilon<1$, split the event that the sum exceeds $x$ according
as its Coulomb term exceeds $\epsilon x$. The upper and lower tail
limits, followed by $\epsilon\downarrow0$, give the same tail
constant times the stated time factor. No independence between
the two contributions is needed.

Finally the truncated tenth moment below $1/u$ is asymptotic to
$10D\log(1/u)$ for the normalized instantaneous law, by integration
of the tail. Taylor expansion to degree ten on $A\le1/u$ has
integrated remainder $O(u^{10})$, since the truncated eleventh
moment is $O(1/u)$. On $A>1/u$ the exponential and the first nine
polynomial terms together cost $O(u^{10})$. This proves the
formula and its time-integrated versions.
\end{proof}
Thus the boundary matrix element of the generator already fails at
twentieth electric order at every fixed nonzero hopping. A
positive-duration transition of this pair-creating Hamiltonian at
the same fixed $g$ sums other histories; the initial derivative
does not determine its electric regularity.

\paragraph{A closed charged return in a connected transverse measurement.}
The same neutral normal sector can be tested using the photon alone.
Keep the quartic two-band Hamiltonian, neutral radial ground state and
vacuum slab of Proposition~\ref{prop:physical-neutral-tip-uv}. The
hopping amplitude is $g$, and $e$ is the electric charge. The normal
coordinate moves throughout the slab. Create a pair at time $s$
and annihilate it at time $t>s$. This is the entire term quadratic
in $g$ in the vacuum matter logarithm. The linear term vanishes by
particle number. The two spatial chords close, but their measured
Wilson phases sample different times and different normal depths.

\begin{proposition}[A divergent connected two-photon coefficient on the neutral chart]
\label{prop:seam-closed-uv}
On a fixed spatial torus there is a smooth real transverse background
$b$ of nonzero spatial momentum for which the connected two-insertion
response has a divergent bare joint coefficient of order $g^2e^{26}$.
The bare vacuum logarithm before vacuum normalization already has
a divergent joint coefficient at $g^2e^{20}$. These statements
concern the indicated coefficients in a two-parameter expansion.
\end{proposition}
\begin{proof}
Use the neutral onsite and photon vacua at both ends. At order
$g^2$ there is only one intermediate pair. The two hopping operators
are adjoints; the vacuum contraction pairs the particle and hole
respectively with themselves. In position space the weights of both
heat kernels are positive and the sites do not move between the
two insertions. For a given radial trajectory and lifted chords,
write $Y$ for half the variance of their oriented transverse Wilson
sum plus $\int_s^t C_{r_u}(z)\,\mathrm du$, where $C$ is the full
Gauss energy per $e^2$ from the preceding proof. Thus $Y\ge0$.
The positive measure on paths and hop times contains
$e^{-2m(t-s)}$, the two heat kernels and the radial ground-state
path law. Its total mass is finite: bound one periodized heat
kernel by $C(1+r_s^{-3/2})$, integrate the other at mass one and
use the stationary $r^4|\psi_0|^2$ marginal. Fixed overall volume
and vacuum factors are harmless. In this measure the photon and
Coulomb integrations give $e^{-e^2Y}$. At this hopping order the
logarithm has no products of lower hopping derivatives.

Adding the transverse background $\lambda b$ multiplies this integrand
by the phase $\exp(ie\lambda L_b)$, where $L_b$ is the oriented sum
of the two measured classical chord integrals, with the respective
spatial readouts. Differentiation twice in $\lambda$ gives
$-e^2 L_b^2 e^{-e^2Y}$. Equivalently use the photon source
$J=(-\partial_t^2-\Delta)b$, with $b$ supported in the interior of
the slab. The shift of the free transverse Maxwell Gaussian is then
$b$; its free quadratic source term has no $g^2$ coefficient. The
normalization at zero source removes every background-independent
term. It leaves this connected response.

Choose $s,t$ in disjoint interior time windows with positive
separation, put $r=r_s$ and restrict $r_t$ to a fixed compact
subinterval of $(0,1)$. In the $j=0$ sector the radial heat operator
is the radial Dirichlet Laplacian of the five-dimensional unit ball.
For separated times its kernel and the ground state are continuous
and strictly positive at the center. On the indicated windows the
joint radial measure therefore bounds a positive constant times
$r^4\,\mathrm dr$ near $r=0$, after integrating $r_t$ in that
subinterval. For each early hop retain displacements $z=\sqrt r\,y$
with $y$ in a fixed compact cone segment away from zero, and use the
reverse lift $-z$ for the late hop. The first heat density has a
strictly positive mass on that scaled set, while the second is
bounded below there uniformly. These restrictions take a subset of
an integral with a common sign.

By the short-chord estimate for the equal-time Wilson variance,
$W_r(\sqrt r\,y)\ge c r^{-1/2}$ on this cone segment, uniformly
for sufficiently small $r$. The late chord variance is bounded
by $C|z|^2$ because its positive readout depth stays away from
zero. The reverse triangle inequality for the Gaussian
$L^2$ norm of the sum gives $Y\ge c' r^{-1/2}$ even after including
the exact photon cross term. The nonnegative Gauss term only
increases the bound. This proves the vacuum claim by integrating
$Y^{10}$ against $r^4\,\mathrm dr$.

For the connected measurement take
$b(u,x)=\eta(u)v\cos(p\cdot x)$, $p\ne0$, $v\cdot p=0$.
Let the smooth function $\eta$ equal one on the early window and
vanish on the late one. The harmonic factor $e^{-r|p|}$ tends to
one. On a positive-volume set of base points where the cosine
stays away from zero, choose the cone so $|v\cdot y|$ is bounded
below. The early chord then gives $L_b^2\ge c r$. On the same
subset the coefficient involving the $n$th moment of $Y$ costs
at least
\[
 c_n\int_0^{r_0}r^{5-n/2}\,\mathrm dr .
\]
It diverges for $n=12$, which gives $e^{26}$ with the two external
couplings included. All lifts and hop times have the same sign
in this moment. Equivalently, alternating finite differences
in $e^2$ of the positive transform with weight $L_b^2$ recover
the moment by monotone convergence, after factoring out the
external $e^2$. This argument supplies a lower bound and does
not specify the first divergent electric order of the closed
response.
The uniform short-chord lower bound also follows directly on the
torus. In the shell $1/r\le |k|\le 2/r$ keep modes with
$|k\cdot\widehat y|\le c/\sqrt r$ for a sufficiently small $c$.
The sinc factor of the chord is bounded below. The slab within
the shell contains at least $c_R r^{-5/2}$ reciprocal-lattice
points uniformly in the unit normal for small $r$: removing the
fixed cell diameter from its radial edges and slab faces leaves
a volume of that same order. Each retained transverse covariance
term is bounded below by $c' r^2$.

\end{proof}

This test has only external photons. The singular term occurs within
a contractible charged return and survives connected vacuum
normalization. Its relevance to physical PE-CGG-HOUQFT is decided
by the normal dynamics during the intermediate pair and the
transverse readout at the two hops. The same argument cannot be
applied to a branch with a different radial law or probe.

\paragraph{The first divergent orders of the closed coefficient.}
The joint distribution of the two hop times also gives matching
upper bounds. Keep the finite slab and stationary $j=0$ radial
preparation of Proposition~\ref{prop:seam-closed-uv}. Write
$r=r_s$, $u=r_t$ and $Y$ for its nonnegative full variance and
Gauss cost. For histories with zero net torus winding the lifted
displacements are $z,-z$. For bounded smooth transverse tests,
$L_b^2\le C_b|z|^2$ because the harmonic readout is a contraction
on bounded fields. For general windings use
$L_b^2\le C_b(|z_1|^2+|z_2|^2)$ on the two lifts.

Let $G_r$ denote the spatial heat density at lapse $\kappa r$ in
this paragraph. For a short lift,
$W_r(z)\le C(|z|^2+|z|/r)$ by the equal-time kernel bound in
Proposition~\ref{prop:physical-neutral-tip-uv}. On a long lift
$W_r(z)\le C(1+r^{-2})|z|^2$ suffices. Polynomial powers on
long lifts are absorbed by widening the Gaussian, since those
lifts are bounded away from zero. For short lifts use
$|z|^k G_r(z)\le C_k r^{k/2}G_{2r}(z)$. Multiplying the two
widened heat densities and integrating over the common endpoint
gives a factor bounded by $C[1+(r+u)^{-3/2}]$. For opposite
lifts this follows by ordinary convolution on $\mathbb R^3$.
Summing all lifts gives the periodization of the convolution at
zero, with the same bound on a fixed torus.

Consequently the spatially integrated term containing $W_r^n$
is bounded by
\[
 C_n[1+(r+u)^{-3/2}][1+r^{d-n/2}],
\]
including the two external insertions when present. Here $d=1$
for the two-insertion coefficient restricted to opposite lifts.
The extra $r$ comes from $|z|^2$ on the very same early heat
density. Take $d=0$ for vacuum graphs and for the estimate
including all windings. In the latter case a polynomial in either
lift can be absorbed in its heat kernel. The corresponding bound
holds with $r,u$ interchanged. For $n\ge1$ the $n$th power of $Y$
is bounded by a fixed constant times the sum of the separate
$n$th powers of $W_r$, $W_u$ and the Gauss time integral, by
Cauchy--Schwarz for the joint photon variance and convexity.

The small time lag costs no extra inverse power of $r$. Conditional
on $r_s=r$ near zero, the expectation of $r_{s+v}^{-3/2}$ is
bounded by $C_T v^{-3/4}$ for $0<v\le T$. Indeed, in the
five-dimensional ball the ground-state-transformed radial
transition kernel is bounded by the free heat kernel times
$e^{E_0 v}\psi_0(y)/\psi_0(r)$; the ratio is bounded for such
starting radii. Convolution of $|y|^{-3/2}$ with the free kernel
costs at most $C v^{-3/4}$, by rescaling and splitting the
rescaled integrand at $|y|=1$. For starting radii away from zero
$(r+u)^{-3/2}$ is bounded directly. Integrating the lag therefore
leaves the moment of $1+r^{d-n/2}$ in the stationary radial law.
Time reversal treats the late hop. For the Gauss term use
$C_{r_v}(z)\le C(1+r_v^{-1})$. Its time integral has moments
of all orders by the radial exponential estimate. Cauchy--Schwarz
bounds the remaining product since
$\mathbb E[1+(r+u)^{-3}]\le 1+\mathbb E[r^{-3}]<\infty$.
This controls the complete terms, including the domain $t-s\downarrow0$.

At hopping order two the vacuum coefficients thus exist for
$n<10$, and the connected response on opposite lifts exists for
$n<12$, where the respective powers are $e^{2n}$ and $e^{2n+2}$.
Together with the lower bounds above, $g^2e^{20}$ in the vacuum
and $g^2 e^{26}$ in the contractible two-insertion measurement
are the first failures under the stated preparation and the
indicated nonzero probe. The contractible estimates also apply
to the continuum coefficient per spatial volume: use the ordinary
Gaussian convolution and the corresponding Fourier integrals.

If all torus windings are admitted, an earlier connected failure
is visible with a probe supported on the late window. Keep the
same small early chord and use late lift $2\pi R\ell-z$, with
$\ell$ a nonzero integer vector. Take nonzero lattice momentum
$p$ perpendicular to $\ell$ and polarization parallel to $\ell$.
The late heat density at this lift stays strictly positive as
$r\downarrow0$ at the specified finite late depths. Its variance
is bounded, while the background line integral has a nonzero
limit on a set of base points of positive volume.
The lower bound for $Y$ survives and $L_b^2$ now has a positive
constant lower bound. This gives divergence at $n=10$.
The $d=0$ upper estimate proves that $g^2e^{22}$ is the first
divergent coefficient for this all-winding connected measurement.
The lower bound in this last statement uses a wrapped torus chord.
For this same late-window probe an opposite short lift gives
$|L_b|\ge c\sqrt r$ on a cone of early displacements and a
positive-volume set of base points, since its late depth stays in
a fixed compact interval. Thus one and the same transverse test
can detect both thresholds.

\paragraph{A test on the coefficients of the normal Maxwell action.}
The failure of the connected photon test is stable under smooth
changes of the normal transport law. On the neutral quartic chart,
suppose a charged displacement accumulates lapse
$\tau(r)=\int_0^r d(u)\,\mathrm du$, where $d$ is continuous with
$d(0)>0$. On each carrier slice let the normal Maxwell form for a
transverse Fourier mode of spatial length $k$ have radial coefficient
$P$ and tangential coefficient $Q$. Its measured harmonic extension
$h_k$ solves
\[
 -(P h_k')'+Q k^2 h_k=0,\qquad h_k(0)=1 .
\]
Take the positive finite-energy extension with $0\le h_k\le1$ on
the continuing collar. Suppose on an initial interval $[0,b]$
that $P,Q$ are bounded above and below by positive constants and
$P'$ is bounded. Keep the canonical onsite vacuum, the same readout
on the temporal charge, the full Gauss square, and the moving neutral
pair with radial law specified above. These conditions describe a
class of represented transverse curvature forms. The unrestricted
chiral pairing need not belong to it.

For large $k$ choose fixed $L>l>0$ so
$Q-PL^2+LP'/k<0$ and both $Q-Pl^2\pm lP'/k>0$ throughout the
interval. The one-dimensional comparison principle gives
\begin{equation}
 \frac{e^{-Lkr}-e^{-Lkb}}{1-e^{-Lkb}}
 \ \le h_k(r)\le e^{-lkr}+e^{-lk(b-r)} .
 \label{eq:seam-regular-optical-uv}
\end{equation}
The lower function is a subsolution with boundary values one and
zero. The two upper exponentials are supersolutions whose sum
dominates the boundary values at both ends. In particular on
$1/r\le k\le2/r$ the measured multiplier stays bounded below as
$r\downarrow0$. For depths in a compact subinterval of $(0,b)$
the high-frequency extension decays exponentially. At small $r$
the upper estimate also bounds the Coulomb self-energy by
$C(1+r^{-1})$ after summing $h_k(r)^2/k^2$ in three dimensions.
For the two-hop lower bound only nonnegativity of the Gauss
energy is needed; one can keep the radial trajectory below an
interior upper radius throughout. The killed radial kernel still
has the required positive small-radius lower bound on the two
separated time windows.

Since $\tau(r)\asymp r$, the early hop has typical displacement
of order $\sqrt r$. The shell-slab proof of
Proposition~\ref{prop:seam-closed-uv} now applies using
\eqref{eq:seam-regular-optical-uv} in place of $e^{-r k}$.
For a fixed nonzero external momentum the measured probe tends to
its boundary value as $r\downarrow0$. Thus the connected
contractible $g^2 e^{26}$ coefficient diverges throughout this
class, including in its infinite-space graph per unit spatial
volume. In the latter calculation the shell-slab lattice estimate
is the ordinary momentum-volume estimate, and averaging the
squared external sinusoid per unit volume gives its positive
mean. This statement is about a perturbative coefficient and
requires no assumption about thermodynamic uniqueness.

One way for the physical theory to evade this test is to change the
relative geometric lengths. If the constitutive readout is Poisson
and its symbol is invariant under the joint dilation of displacement
and heat lapse, $z\mapsto\lambda z$ and $\tau\mapsto\lambda^2\tau$,
then $a(\lambda^2\tau)=\lambda a(\tau)$ as a function of the
accumulated lapse. This identity fixes $a=C\sqrt\tau$ for $\tau>0$.
Dilation covariance here is an additional constitutive condition.
The units and a nonzero value of $G_N$ alone do not fix the relative
calibration or prove that the condition holds.

\paragraph{Which angular energy survives the physical Gauss constraint?}
On a fixed sector plane, a charged spinor and its physical dual admit
a neutral pair. If a compact symmetry acts unitarily on a nonzero
finite-dimensional retained spin space $W$, the identity in
$W\otimes W^*\simeq\operatorname{End}(W)$ is invariant under the
diagonal action. Its Hilbert--Schmidt norm squared is $\dim W$.
Restricting to a symmetry-invariant chiral subspace replaces the
identity by that of the subspace. In particular, when
$P_{\rm phys}=(1-\gamma_5\otimes\Gamma_F)/2$ acts on internal and
spin indices independently of the normal angles, it commutes with
the angular Fourier projections. A dual row in this statement acts
on the retained ket; an assumed compression of the Weyl kinetic
operator on both sides would be a different operation.

Even for an Abelian normal circle, a neutral pair can carry a large
angular energy. To calculate it one must know how the two
orientations at the endpoints enter the normal metric.
Consider two $2\pi$-periodic angles at a seam with metric
$\mathrm dr^2+\mathrm d\theta^{\mathsf T}M(r)\mathrm d\theta$,
$M(r)>0$ for $0<r<1$. Use the invariant subspace of the metric
Hilbert space with measure $\sqrt{\det M(r)}\,\mathrm dr$ times
Haar measure on the angles. In this realization projection onto
the diagonal Gauss constraint is normalized Haar averaging on
$\theta\mapsto\theta+\zeta (1,1)$. The invariant Fourier weights
are $q n$, $q\in\mathbb Z$, $n=(1,-1)^{\mathsf T}$. Put
$u=(1,1)^{\mathsf T}$. On this sector the angular contribution to
the radial Hamiltonian is exactly
\begin{equation}
 V_q(r)=q^2 n^{\mathsf T}M(r)^{-1}n
       =q^2\frac{u^{\mathsf T}M(r)u}{\det M(r)} .
 \label{eq:seam-selection-uv}
\end{equation}
The first equality is the angular Dirichlet form evaluated on its
Fourier character. The second follows from the two-by-two inverse
formula. It expresses the relative angular energy in terms of the
length of the diagonal orbit and the area of the entire two-angle
fiber. The invariant subspace inherits its inner product from the
original fields; an additional orbit-volume factor in a different
boundary measure would have to be carried through the calculation.

\begin{proposition}[Gauss-neutral pairs with indistinguishable single-leg barriers]
\label{prop:seam-selection-uv}
There are two positive angular metrics with the same radial measure,
the same diagonal Gauss action, and the same angular energy for an
isolated charge on either leg, for which the neutral pair energies
are $2q^2r^{-8}$ and $2q^2$, respectively. In the moving-depth
electric hopping model, coupling the hop to the corresponding pair
character gives finite one-hop coefficients at every electric order
in the first ground-state preparation. The second ground-state
preparation has its first divergent one-hop coefficient at electric
order twenty.
\end{proposition}
\begin{proof}
Write $P_d=uu^{\mathsf T}/2$ and $P_a=nn^{\mathsf T}/2$. Take
\[
 M_1=P_d+r^8P_a,\qquad M_2=r^8P_d+P_a .
\]
Both volume factors are $r^4$. For each coordinate vector $e_i$,
$e_i^{\mathsf T}M_b^{-1}e_i=(1+r^{-8})/2$ for $b=1,2$.
Equation~\eqref{eq:seam-selection-uv} gives the two stated
values on the pair. Reversing the sign of the second angle
interchanges the two metrics but also interchanges which circle
is gauged. They are distinguished once the gauge embedding is fixed.

For the transfer calculation use the charged action and zero-mode
convention of Proposition~\ref{prop:physical-neutral-tip-uv}.
Attach $e^{iq(\theta_1-\theta_2)}$, $q\ne0$, to its pair creation
operator, and the conjugate factor to pair annihilation. This
unitary angular shift respects the diagonal Gauss projection and
leaves the spatial Wilson contact and the Coulomb charge square
as in that proposition. The electromagnetic charge still belongs
to its original Fock representation. If spin slots are included one
can take the identity pairing on $W$ and test its normalized singlet.
Use initial angular weight zero and final weight $qn$, with the
positive radial ground vector in each sector, the photon vacuum
at both ends, and the same spatial pair probe. The transfer at first
hopping order has zero angular energy before the hop and potential
$V_q$ after it. With $H_0=-r^{-4}\partial_r(r^4\partial_r)$ and
the same Friedrichs and outer conditions, let $v_0$ and $v_q$ be
the ground vectors of $H_0$ and $H_0+V_q$. At the marked time $s$
the unperturbed radial weight is proportional to
$r^4 v_0(r)v_q(r)\,\mathrm dr$, with positive scalar factors from
the two ground energies and the pair mass.

For $M_1$ the charged collar estimate of
Proposition~\ref{prop:final-collar-uv} applies with angular
potential $2q^2r^{-8}$: its proof uses only a strictly positive
coefficient of $r^{-8}$. Thus $v_q$ has inverse radial moments of
all orders, and Cauchy--Schwarz gives the same integrability with
$v_0v_q$. The $n$th moment of the transverse phase variance after
the heat-step average costs at most $C_n(1+r^{-n/2})$. The
integrated Coulomb potential has all finite polynomial moments on
each finite slab. Indeed its exponential moment is bounded by the
radial semigroup with potential $V_q-C_b(1+r^{-1})$ after the
hop, since the full Coulomb energy at fixed displacement is
bounded by $C(1+r^{-1})$. Weighted Hardy control gives a lower
form bound for this semigroup for every finite $b$. Positivity
and the elementary power bound for the sum of the two phase costs
now give all finite electric orders. For $M_2$, $v_q=v_0$ and the
constant $V_q=2q^2$ merely multiplies the one-hop integrand by
$e^{-2q^2(T-s)}$. This is bounded above and below by positive
constants on a fixed slab. The lower moment and the finite lower
orders of Proposition~\ref{prop:physical-neutral-tip-uv}
therefore apply to the complete coefficient.
\end{proof}

This test concerns the diagonal angular constraint. At distinct
spatial endpoints an electrically neutral pair still requires its
Wilson or Gauss dressing under local gauge transformations, which
is retained in the hopping source here. A full BFV differential
can impose further constraints or make the singlet exact.
Postulate~$\Omega_\alpha$ calls for that differential, the
represented charged block and the sewn state from the same raw
datum. To transfer the radial proof to that branch one must evaluate
the angular form on the surviving cohomology and the matrix element
of the actual marked operator there. Both comparison metrics give
the same single-leg measurement, and the chiral projection described
above makes the same spin selection in both.

\paragraph{A test on every independent angular circulation.}
The two-angle computation extends to a finite connected network of
seam angles. Assign one circle to each oriented edge. Let $B$ be
the vertex--edge incidence matrix with one redundant vertex row
removed. A change of vertex phases acts on the angles by
$\theta\mapsto\theta+B^{\mathsf T}\lambda$. Haar projection
keeps precisely the Fourier characters $\ell\in\mathbb Z^E$
satisfying $B\ell=0$, where $E$ counts the edges. External
charges replace the zero on the right by their prescribed net
injection. For a neutral return choose a spanning tree and let
$C$ have as its columns the fundamental oriented cycles, with
coefficient one on their respective chords. Every invariant
character is uniquely $\ell=Cm$ with $m$ integral. If $M(r)>0$ is
the angular metric, its energy on that character is
$m^{\mathsf T}S(r)m$, where
\begin{equation}
 S(r)=C^{\mathsf T}M(r)^{-1}C,\qquad
 \det S(r)=\frac{\det\bigl(BM(r)B^{\mathsf T}\bigr)}{\det M(r)} .
 \label{eq:seam-cycle-uv}
\end{equation}
To verify the determinant, use the tree columns of $B$, an integer
matrix of determinant $\pm1$, to construct a right inverse $D$ of
$B$ supported on the tree. The square matrix $U=[D\ C]$ has
determinant $\pm1$ by its tree--chord block form.
The lower diagonal block of $U^{\mathsf T}M^{-1}U$ is $S$.
The complementary upper diagonal block of its inverse is
$BMB^{\mathsf T}$, since $BU=[I\ 0]$. The determinant identity
follows from the Schur complement. For a graph with more than one
independent circulation the determinant gives their combined
volume; the least eigenvalue governs a bound valid in every
circulation direction. For example $S=\operatorname{diag}
(r^{-8},r^{-8})$ and $S=\operatorname{diag}(1,r^{-16})$ have
the same determinant. The latter leaves a nonzero integral
circulation with constant energy.

On a quartic chart, $S(r)\succeq c r^{-8}I$ near the tip supplies
the charged barrier for each nonzero invariant character on this
network. At a charged visit the cut can carry a prescribed nonzero
vertex injection $b$. The actual characters then have the form
$\ell=\ell_0+Cm$ with $B\ell_0=b$. Their barriers are evaluated
on $\ell_0+Cm$, including its cross term. A cycle bound alone
cannot certify an affine cut on which the inserted source has
support. With finitely many actual characters at a given order
one can instead test their individual quadratic forms. A state
at $m=0$ needs its own estimate. If a graph sums over unbounded
internal charges, positivity on each individual charge alone gives
no uniform constant. For instance $(m_1-\sqrt{2}m_2)^2$ is
positive on each nonzero integer pair but approaches zero along
integer approximants to $\sqrt{2}$. A physical GCM graph would
have to supply the actual angular metric, the algebra of junction
phases, and the image of its boundary insertions in these
characters. The angular index $m$ describes normal transport at a
seam; it is separate from the independent four-momenta whose
convergence is tested after sewing.

\paragraph{A finite-duration test in a number-conserving representation.}
One can test whether summing the hopping series removes the logarithm
in a setting where the entire sum is accessible. Keep the neutral
quartic radial ground preparation, Maxwell field, torus, heat chord
and common harmonic electric readout just used. Choose two equal-charge
active slots, both empty in the reference vacuum, and one spectator
of opposite charge. Each species has conserved particle number, and
a hop exchanges the active slots. On the one-active-particle space
the hopping Hamiltonian is $-g\sigma_1 K_{r,A}$ with the same scalar
Wilson heat kernel as in~\eqref{eq:cl-gauss-charge-uv}. It is
Hermitian, with $\|K_{r,A}\|\le1$ by the periodized unit heat-density
majorant and the Schur test. Give the active slots the same mass
$M>0$ and the static spectator mass $M_s>0$. One may take
$0<|g|<M$ to have a positive matter Hamiltonian on all number
sectors. The full Coulomb square and Maxwell energy are included.
In the empty reference the matter density and every hopping operator
annihilate the vacuum, so its determinant and relative vacuum
normalization are one. These choices specify a distinct charged
representation: the active hop conserves number in this model, whereas
the earlier filled reference permits pair creation.

Consider propagation for time $T>0$ in the sector containing one
active charge and one spectator. Prepare the active slot in state 1
and detect it in state 2. Both particles have normalized constant
spatial waves at both ends; the photon and radial endpoint states
are as before. These globally neutral pair probes are specified in
the Gauss-reduced chart. Write $F_{g,T}(u)$ for the transition divided
by the free radial and photon vacuum amplitude and multiplied by
$\operatorname{sgn}g$, where $u=e^2$. It is a connected two-point
matrix element of the composite pair endpoints since their one-point
functions vanish.

\begin{proposition}[Resummation at fixed hopping in a neutral comparison]
\label{prop:seam-uv-resummed}
At each fixed $0<|g|<M$ the complete finite-duration amplitude is finite
at every $u\ge0$. Its Taylor polynomial $P_{9,g,T}$ through degree nine
exists at zero. With $D$ as in Proposition~\ref{prop:seam-uv-tip-log},
\[
 \begin{aligned}
 F_{g,T}(u)-P_{9,g,T}(u)
     &\sim \frac{D_{g,T}}{9!}\,u^{10}\log(1/u),\\
 D_{g,T}&=e^{-T(M+M_s)}D\,|g|T\cosh(|g|T).
 \end{aligned}
\]
This is the electric expansion after summing all hopping orders.
\end{proposition}
\begin{proof}
In the fixed two-particle sector the hopping is bounded, so its
time-ordered Duhamel series converges in norm. Only odd numbers
$h$ of hops connect the two slots. After multiplying by
$\operatorname{sgn}g$, all coefficients are positive weights
$|g|^h$. For fixed ordered times $t_i$, condition on the stationary
radial path. Integration between constant spatial endpoints gives a
uniform initial position and independent conditional increments
$z_i=\sqrt{r(t_i)}Y_i$, with the Gaussian vectors $Y_i$ as in the
preceding proof. The spectator stays at its initial position.
All lifts are included. Gaussian integration of the transverse
photons yields $e^{-uA}$ with
$A=\|\sum_i J_i\|^2/2$, where $J_i$ is the real smeared chord
current at time $t_i$ in the spacetime Maxwell covariance norm.
Its individual value $A_i=\|J_i\|^2/2$ is precisely
$W_{r(t_i)}(z_i)/2$. The time dependent active--spectator
separation is $d_t$. The Gauss term adds
$B=\int_0^T C_{r(t)}(d_t)\,\mathrm dt\ge0$. Its moments of every
fixed order are bounded uniformly in all hop times and paths after
radial averaging, by the separation-independent estimate
$C_r(d)\le C(1+r^{-1})$ and the radial exponential bound above.
The same time-sliced transfer or Feynman--Kac construction yields this
Gaussian formula with both self-energies. The total positive measure
before $e^{-u(A+B)}$ hence has mass
$e^{-T(M+M_s)}\sinh(|g|T)$.

At any fixed distinct hop times each $A_i$ has tail $Dx^{-10}(1+o(1))$.
For two distinct times, the joint radial law has density at most
$C_{t_i,t_j}r^4{r'}^4$ with respect to $\mathrm dr\,\mathrm dr'$.
Indeed its ground-state kernel is the ordinary Dirichlet heat kernel
on the five-dimensional ball restricted to radial functions, multiplied
by the bounded ground function at both ends and the ground-energy
factor. At a fixed positive separation in time this heat kernel
is bounded. Conditional independence of the two Gaussian steps and
the same change of variable $r=v/x^2$ as above show that
$\Pr(A_i>\epsilon x,A_j>\epsilon x)=O_{t_i,t_j,\epsilon}(x^{-20})$
for every $\epsilon>0$. On the event that all but the largest $A_i$
are at most $\epsilon x$, the triangle inequality and its reverse
give
\[
 |\sqrt A-\sqrt{\max_i A_i}|\le (h-1)\sqrt{\epsilon x}.
\]
Inclusion--exclusion for the individual exceedances, followed by
$\epsilon\downarrow0$, therefore gives
$x^{10}\Pr(A>x)\to hD$. Adding the nonnegative $B$ leaves this
limit unchanged because its eleventh moment is finite.

The time-dependent constants from the two-time density are needed
only for pointwise limits at distinct times. For domination,
$\sqrt A\le h\sqrt{\max_i A_i}$ and the individual tail bound
give $x^{10}\Pr(A>x)\le C h^{21}$ uniformly in those times.
Together with the bound on $B$ this is summable against
$|g|^hT^h/h!$. The resummed tail constant is consequently
$e^{-T(M+M_s)}D\sum_{h\ {\rm odd}}h(|g|T)^h/h!=D_{g,T}$.
The truncated-moment argument of Proposition~\ref{prop:seam-uv-tip-log}
now proves the asserted logarithm.
\end{proof}
Thus positivity of the transfer, total electric neutrality, and a
finite real-charge amplitude can coexist with a divergent electric
coefficient even at fixed hopping and positive duration. This example
does not determine the positive-duration amplitude in the filled-sea,
pair-creating representation or select a physical chiral BFV state.

\paragraph{A radial trajectory during electric pair creation.}
The same normal coordinate can move between charged insertions.
Consider the Hamiltonian of
Proposition~\ref{prop:cl-gauss-charge-uv} on a fixed spatial torus.
Use a finite vacuum slab of duration $T$, the neutral onsite and
photon vacua at its ends and the ground vector $u_j$ of the free
radial operator $H_j$, with angular index $j\ne0$. Factor out its
free ground energy. The ground-state transform of the radial
semigroup gives a stationary process $r(t)$ on $(0,1)$ with
one-time law $|u_j|^2 f\,\mathrm dr$. Denote expectation in this
process by $\mathbb E_j$. The tip estimate gives
$\mathbb E_j e^{c/r(t)^3}<\infty$ for some $c>0$.
For every $b>0$ Jensen's inequality in time then gives
\begin{equation}
 \mathbb E_j\exp\!\left[b\int_0^T r(t)^{-3/2}\,\mathrm dt\right]
 \le \mathbb E_j \exp[bT r(0)^{-3/2}]<\infty .
 \label{eq:physical-moving-radial-cost-uv}
\end{equation}
This estimate uses no independence of depths at different visits.

\begin{proposition}[Ordered electromagnetic graphs with a moving charged depth]
\label{prop:physical-moving-ordered-uv}
For $2|g|<m$ the bare vacuum normalization on this slab is
positive and finite at every real electric charge. At every finite
electric order the connected vacuum and externally tested photon
graphs, including their Coulomb self-contractions, admit an absolute
bound jointly in the radial process, charged durations, charged
spatial momenta and independent photon four-momenta. This result
uses the canonical ordered charged transfer, with the Gauss zero
mode fixed.
\end{proposition}
\begin{proof}
Condition on the full radial trajectory in the transverse chart. At
a hop time $t_i$ the spatial kernel has heat lapse $\kappa r(t_i)$.
On a vertical interval the temporal source couples to the density
sampled by $P_{r(t)}$. Integrating the auxiliary scalar on disjoint
time slices therefore gives exactly the positive Coulomb energy
with both self and cross terms. This path decomposition follows by
inserting radial positions in the time-sliced transfer: the radial
weights are the transition kernels of the ground-state transform,
whereas the electromagnetic readout depends on the instantaneous
$r$. To justify the identification first impose absorbing radial boundaries
on a compact subinterval and expand the hopping by Duhamel between
the diagonal transfers. There its Hilbert--Schmidt norm is bounded
uniformly in the photon field. On sectors of at most $N$ particles
its creation and annihilation operators have norm at most
$C\sqrt{N+2}$, by the row estimate used in the pair-form bound and
duality. Starting from the vacuum, the time-ordered expansion is
therefore absolutely convergent on the slab: its $n$th norm bound is
at most $(CT)^n\sqrt{2^n n!}/n!$ up to a fixed shift of the
factorial. The diagonal transfers preserve particle number. In the
position representation they give the radial and photon transition
kernels and the diagonal Coulomb Feynman--Kac factor; expansion
in the auxiliary scalar recovers the same joint law. The onsite
quasifree contraction of the complete hopping series is the closed
loop determinant. Finally the localized Friedrichs forms exhaust the
original form on its core, so their semigroups converge strongly.
The stopped radial paths carry the original ground-state weight
with their survival indicator. The absolute bounds below allow the
indicators to be removed in each integral.

At fixed trajectory the logarithm is again a sum of alternating
hopping loops. Pin one hop at $t$ and integrate the spatial closure
using the supremum of its periodized heat density, bounded by
$C(1+r(t)^{-3/2})$. The other heat steps have mass one. Each
one-sided time kernel is bounded by $e^{-m|u|}$, even in a real
auxiliary scalar field with its Wilson phase included. After $n-1$
duration integrations the cost is at most $(2/m)^{n-1}$.
Summing yields the pathwise bound
$C\int_0^T(1+r(t)^{-3/2})\,\mathrm dt$ on the absolute logarithm.
Reversal of each closed word and exchange of the slots preserve the
actual hop times and their heat widths and conjugate its phases.
Thus the conditional determinant is positive. For almost every
radial path, scalar phase integrals have finite Gaussian variance
since $\int_0^T(1+r(t)^{-1})\,\mathrm dt<\infty$. Equation
\eqref{eq:physical-moving-radial-cost-uv} controls the path
average of the unexpanded vacuum weight.

For the marked estimates, a scalar end on a directed interval $I$
at fixed spatial position contributes, apart from its spatial phase,
\[
 J_I(\nu,p)=\int_I e^{i\nu t-r(t)|p|}\,\mathrm dt .
\]
The radial path need not be differentiable. Plancherel and
$2xy\le x^2+y^2$ give
\[
 \int_{\mathbb R}|J_I(\nu,p)J_{I'}(-\nu,-p)|\,\mathrm d\nu
 \le \pi\left(\int_{|I|}+\int_{|I'|}\right)
                         e^{-2r(t)|p|}\,\mathrm dt ,
\]
where $|I|$ denotes the same interval without its orientation.
Several marks on one interval are products of these single integrals
when all orders on the interval are included. After assigning marks
to segments the remaining horizontal time phases have unit modulus.
Each internal photon has its own frequency before any hop-time
integral, so the estimate applies independently to every Coulomb
contraction, including when both ends lie on one interval. Summing
with $|p|^{-2}$ on the torus costs at most
$C\int_0^T(1+r(t)^{-1})\,\mathrm dt$ per such line. A transverse
photon with ends at $t_i,t_l$ costs at most $C(1+r(t_i)^{-2})$
after its frequency integral and spatial sum, by
$\int(\nu^2+|p|^2)^{-1}\,\mathrm d\nu=\pi/|p|$ and its two
harmonic readouts. Momentum conservation in space can here be
discarded in a positive sum over discrete internal momenta to obtain
an upper bound.

Keep one spatial Fourier momentum on each charged loop. At fixed
times its integral or sum uses a single translated heat factor
at cost $C(1+r(t_i)^{-3/2})$. Other heat factors are bounded;
a fixed number of horizontal marks produces Gaussian derivatives
bounded at every $0<r\le1$ with the same spatial summability.
All remaining depth costs form products of a fixed number of inverse
powers at specified times or integrated over intervals of length
at most $T$. Their expectations are bounded uniformly in all these
times by H\"older's inequality and stationarity of the one-time
law. The same conclusion holds with the zero-charge vacuum weight
by Cauchy--Schwarz and \eqref{eq:physical-moving-radial-cost-uv}.
The number of required one-time moments depends on the electric
order, while the number of possible mark positions grows only
polynomially with the hopping length. The unchanged geometric sum
in $2|g|/m$ finishes the estimate, also when the marked visits
approach one another. Connected coefficients are obtained after the
radial average: the zero-charge vacuum normalization is positive, and
every finite product of the averaged bare derivatives is finite.
\end{proof}

This estimate controls the photon frequencies before integrating the
charged durations. To keep a carrier Fourier variable on each
charged loop as well, use the time kernel on the slab and take its
diagonal Fourier trace. Conditioning on a moving trajectory
removes time-translation invariance, so frequencies need not be
conserved at each intermediate hop.

\paragraph{The charged carrier frequencies at a moving depth.}
A moving radial coordinate breaks time translation invariance
conditional on its trajectory. There is nevertheless a Fourier
variable for the temporal closure of each charged loop. The
reversibility of the prepared radial diffusion controls the electric
marks uniformly when this closure is opened.

\begin{proposition}[Joint bare graph bound for a moving charged normal sector]
\label{prop:seam-moving-frequencies-uv}
Use the finite slab, spatial torus, charged free radial ground
preparation $j\ne0$, scalar two-band hopping and $2|g|<m$ of
Proposition~\ref{prop:physical-moving-ordered-uv}. Every finite
electric order admits an absolute bound jointly over the radial
law, one charged spatial momentum and carrier frequency per closed
line, the photon four-momenta, and the remaining hop times.
This holds for normalized vacuum and photon coefficients and for
finitely many elementary charged endpoints in neutral combinations,
tested with smooth time packets inside the slab and fixed spatial
Fourier modes. The time packets can be integrated in absolute value.
This is the diagonal Fourier trace of the original charged integral
operator in its carrier variable, with the endpoints extended by zero
outside the slab. Its intermediate factors need not conserve carrier
frequency on any individual moving trajectory.
\end{proposition}
\begin{proof}
Let $X_t=r(t)$, $0\le t\le T$, have the stationary radial law of
the preceding proposition, with probability marginal $\pi$. For
$s\ge2$ and $b\ge0$ put $f_b(r)=e^{-br}$. We claim that for
every $N$ and all real $\nu$,
\begin{equation}
 \left\|\sup_{I\subset[0,T]}
   \left|\int_I e^{i\nu t}f_b(X_t)\,\mathrm dt\right|\right\|_{L^s}
 \le \frac{C_{s,N,T}}{(1+|\nu|)(1+b)^N},
 \label{eq:seam-moving-mark-uv}
\end{equation}
where $I$ ranges over intervals. At $\nu=0$ the same spatial
bound holds for $\|\sup_t f_b(X_t)\|_{L^s}$.

Here is the time estimate. For a smooth real radial function $f$
write its increments in forward and reversed time as the sum of a
martingale increment and the integral of its generator. Reversibility
gives the same drift integral in the two orientations. Subtracting
the two formulas gives twice the forward increment of $f(X)$ as the
forward martingale increment minus the reverse martingale increment.
In this metric each martingale has bracket
$2\int |f'(X_t)|^2\,\mathrm dt$ in its time orientation.
One may first use It\^o's formula in the interior; the conservative
ground-state process stays inside the interval on a finite time slab.
For functions with bounded derivative both stochastic integrals are
true martingales. Integrating the increment identity against a smooth
deterministic $h(t)$ of modulus at most one gives the corresponding
forward and reversed stochastic integrals, with $h(T-t)$ in the
reverse orientation. This follows also from deterministic Riemann
sums and integration by parts. The martingale maximal inequality
and stationarity therefore give
\[
 \left\|\sup_{I}\left|\int_I h(t)\,\mathrm d f(X_t)\right|\right\|_s
      \le C_s \sqrt T\,\|f'\|_{L^s(\pi)} .
\]
Taking $h=1$ bounds $\|\sup_t|f(X_t)|\|_s$ by the right side
plus $\|f\|_{L^s(\pi)}$. Integration by parts with $h=e^{i\nu t}$
now supplies $|\nu|^{-1}$ for $|\nu|\ge1$, with the two interval
boundary values bounded by $2\sup_t|f(X_t)|$. For $|\nu|<1$ use
the length of the slab. Both $\|f_b\|_{L^s(\pi)}$ and
$\|f_b'\|_{L^s(\pi)}$ decay faster than any inverse power of $1+b$,
by the inverse moments of the charged radial ground state.
This proves~\eqref{eq:seam-moving-mark-uv}. The same maximal
estimate applied to $f(r)=(r+\epsilon)^{-d}$, followed by monotone
convergence, gives all fixed moments of $\sup_t X_t^{-d}$.

At fixed spatial momenta and photon frequencies, open the time
closure at one hop of each charged loop. Keep its two adjacent
end times $y,z$ distinct and restore equality by Fourier inversion
in $\xi$, the frequency conjugate to $y-z$. Fix all other hop
times and the whole radial trajectory first. After assigning every
electric mark to a segment and using all temporal orders on it, the
dependence on these two times factors. A scalar mark contributes
an integral from~\eqref{eq:seam-moving-mark-uv} on its own
interval; a horizontal mark uses the field at its own hop.
Assign the cut hop to one end only. Its two neighboring onsite
intervals give a function $A(y)$ and a function $B(z)$, extended
by zero outside the slab. Thus the absolute integral in the
charged carrier frequency is bounded by
\[
 \int |\widehat A(-\xi)\widehat B(\xi)|
                  \frac{\mathrm d\xi}{2\pi}
   \le \|A\|_2\|B\|_2 .
\]
Parseval gives the original seam integral $\int A(y)B(y)\,\mathrm dy$
before the inequality. The two split end times have already been
Fourier transformed here. The joint estimate concerns $\xi$ and the
remaining hop times; the unresolved delta representation would have
no absolute $\xi$ bound. Each of the two $L^2$ norms costs at most
$(2m)^{-1/2}$ from its one-sided exponential, times the suprema
of its assigned marks. Each closed loop has at least two hops in
the vacuum-relative determinant. The remaining uncut intervals
form a chain with one base time in the slab and $n-2$ durations
for a loop of $n$ hops. Their absolute integrals cost at most
$T(2/m)^{n-2}$. The procedure retains a separate $\xi$ for every
closed line, including when photons connect different cuts.

Keep one spatial Fourier variable per charged loop as well. A
translated heat hop on the loop, with any fixed number of spatial
marks, is bounded by a constant times
$\exp[-c\kappa X_t|k+d|^2]$ or its average over chord parameters,
with $c>0$. The shifts $d$ depend on photon momenta and the
parameters, independently of $X$. An upper bound uniform in $t$
has summable translated spatial tail at the cost of a fixed
power of $\sup_t X_t^{-1}$. Other unmarked spatial heat factors
have supremum one. Open lines have two prescribed endpoint times
and external spatial momentum: there is no additional charged
closure. Their onsite inverse expansion follows by differentiating
the same fixed-trajectory fermion Gaussian. It has one more
one-sided interval than spatial hops and includes the zero-hop
line with its temporal marks. Between the two endpoints its
time integrals cost at most $(2/m)^n$ for $n$ hops, by integrating
from one end and bounding the last interval by one.

Finally apply H\"older's inequality over the radial process to
the finite number of marked factors, translated heat tails and
any zero-charge vacuum weight. The latter has every fixed
positive moment by~\eqref{eq:physical-moving-radial-cost-uv}
and the conditional word estimate. Every scalar photon end supplies
$(1+|\nu|)^{-1}(1+|p|)^{-N}$ in this bound. For horizontal ends
the spatial factor alone follows from the supremum estimate for
$f_{|p|}$. Integration in $\nu$ now uses the two scalar factors
on a Coulomb line or the transverse Maxwell propagator. The
spatial lattice sums converge for sufficiently large $N$, with
the usual Gauss zero-mode convention. At each finite electric
order there are finitely many marked and open lines. Position
choices for their marks grow polynomially with the number of
hops. The factors $|g|^n(2/m)^n$ are summable. Unmarked vacuum
components exponentiate at fixed trajectory before the radial
average. Normalizing and taking connected coefficients uses
finite products of the resulting derivatives. This proves the
joint absolute bound, including simultaneous short visits and
hard charged and photon frequencies.
\end{proof}

This calculation uses a stationary charged radial preparation on a
finite slab and the onsite-relative determinant. The estimate
of~\eqref{eq:seam-moving-mark-uv} makes explicit where a different
angular channel or physical boundary state enters the frequency proof.

\paragraph{Summing the hopping before the electric expansion.}
On a finite slab the moving-depth result also has a formulation at
arbitrary hopping strength. It applies to the spatial Dirac heat
derivative in~\eqref{eq:reference-spin-uv}. Both the determinant
and the open inverse must be expanded about the same moving-depth
operator for this conclusion.

\begin{theorem}[Moving-depth spinor graphs at arbitrary hopping strength]
\label{thm:seam-moving-allg-uv}
Retain the finite slab, torus, boundary preparation and external
tests of Proposition~\ref{prop:seam-moving-frequencies-uv}.
Allow either the two-band heat hopping with arbitrary real $g$,
or the four-slot derivative heat hopping with arbitrary real $v$ in
\eqref{eq:reference-spin-uv}. The electromagnetic coefficients at
every finite order have complete, jointly absolutely convergent bare
graph representations, with all independent charged closure and photon
four-momenta retained. The full vacuum matrix element is positive
and finite at every real electric charge in these preparations.
\end{theorem}
\begin{proof}
Condition on $r(t)=X_t$. Work first with smooth real spacetime
probes. Multiply the charged field by the unitary phase given by
the carrier-time primitive of the \emph{measured} scalar
potential, with the sign cancelling its term in the covariant
time derivative. At fixed radial path the primitive of its mode
$(\nu,p)$ is, apart from its spatial phase,
\[
 J_{0,t}(\nu,p)=\int_0^t e^{i\nu u-|p|X_u}\,\mathrm du .
\]
The transformation acts on the matter variables in the conditional
fermion functional. It fixes the onsite vacuum at the slab boundaries.
Each open endpoint of the original charged field keeps its
compensating phase. On a spatial hop a scalar mark gives this
primitive times the difference of the endpoint plane waves.
There is no differentiation of a Brownian path. In particular using
$e^{-|p|X_t}\int_0^t e^{i\nu u}\mathrm du$ here would omit part
of the charge source. Every scalar end is estimated by
\eqref{eq:seam-moving-mark-uv}; the difference of endpoint phases
on a hop gives at most two translated hopping symbols. Spatial
transverse marks have the usual Wilson derivatives and readouts.

At zero electric field keep the full hopping $V_r(t,k)$, switched
on for the slab and zero outside. It is Hermitian and anticommutes
with the constant onsite involution $\beta$. For each trajectory
put $K_* = \sup_{t,k}\|V_r(t,k)\|$. This is bounded by $|g|$ for
the scalar heat hop and by $C|v|\sup_t X_t^{-1/2}$ for the spinor.
The full-line inverse
\[
 G=(\partial_t+m\beta+V_r)^{-1}
\]
implements the onsite vacuum boundary conditions on its two slots,
with their free propagation outside the slab. Multiplication on
the left by $\beta$ gives $m+\beta\partial_t+\beta V_r$.
The two latter terms have a skew-adjoint sum on the free derivative
domain, since the second is bounded on each trajectory.
Thus $\|G\|_{2\to2}\le m^{-1}$, uniformly in spatial momentum.
Write $S=(\partial_t+m\beta)^{-1}$. Its time kernel is a
one-sided exponential in each onsite slot, bounded by one and
with uniformly bounded $L^2$ rows and columns. The resolvent identity
\[
 G=S-SV_r S+SV_r G V_r S
\]
bounds the dressed time kernel on either side of the diagonal by
$C_m(1+K_*^2)$, using Cauchy--Schwarz for the outer rows and columns.
These estimates allow an arbitrary number of unmarked hops between
source variations.

To identify the source derivatives, the off-diagonal $S V_r$ with
hopping supported on the slab is Hilbert--Schmidt at fixed path and
volume. The spatial square norm of one hop is bounded by
$C(1+X_t^{-3/2})$ in the two-slot case and
$C(1+X_t^{-5/2})$ in the derivative case. Abelian chords have
unit modulus in the corresponding position-kernel bounds. Eliminating
one onsite slot therefore gives an ordinary Fredholm determinant
of a trace-class product. Its logarithmic derivatives and its open
lines give the usual words of the dressed inverse and the marked
hopping vertices. One may first check the formulas at small hopping
by the convergent expansion. Along the real hopping interpolation
the slot operator remains invertible by the skew-adjoint argument,
also in smooth real spatial backgrounds after the scalar change of
variables. The Fredholm identities thus extend to every stated
hopping strength. Only finitely many marked dressed vertices occur
at a fixed positive electric order. The reference determinant and
end phases are the ones of the prepared charged action.

On an active closed line choose a marked spatial hop as the temporal
cut. One of its spatially translated heat symbols or Wilson
derivatives controls the spatial sum uniformly in the shift. For
the spinor it can cost an additional fixed power of $\sup_t X_t^{-1}$.
All other spatial marked symbols have polynomial bounds in that
same random variable. When the loop has at least two propagators,
the two adjacent dressed time kernels in the split closure have
bounded $L^2$ endpoint norms on $[0,T]$ by the kernel estimate
above. All the other hop times lie in this fixed slab. For a loop
consisting of one contact $W$, its onsite trace
$\operatorname{tr}_{\rm slots}(S W)$ vanishes by the off-diagonal
grading, regardless of the convention for the diagonal jump.
The identity $G=S-SV_rG$ then puts a second interval and one
unmarked hop in the remaining trace. The same two-endpoint Fourier
bound applies. An open chain has no closure variable and its
dressed time kernels have the same uniform bound, including the
case in which all scalar insertions occur on its endpoint phases.
This proves the charged-frequency and spatial bounds, with polynomial
random factors and the scalar or transverse mark estimates already
described.

The zero-electric-charge vacuum normalization at fixed path is also
needed. In the charged Fock space the instantaneous Hamiltonian
at zero electromagnetic field has lower form bound
$-C_m\|V_r(t)\|_{\rm HS}^2$, by the pair-form estimate. Time-ordered
product evolution on the slab consequently bounds its vacuum
amplitude in modulus by
\[
 \exp\!\left[C\int_0^T (1+X_t^{-5/2})\,\mathrm dt\right] .
\]
The estimate follows first for step approximations of a path staying
in a compact radial interval. There the heat kernels vary continuously
in Hilbert--Schmidt norm and their number-graded vacuum Duhamel
series converges, giving the limit. Almost every prepared path stays
in the interior on this slab. The displayed majorant has moments of
every fixed order: time Jensen and stationarity reduce them to
one-time exponential moments, and $r^{-5/2}=o(r^{-3})$.
Unmarked closed vacuum contributions at zero electric order are
kept in this determinant. For its logarithm, if a Fourier trace
is required, interpolate the hopping from zero and use the same
one-contact trace and its two-interval identity with one
undifferentiated heat insertion.

In the derivative-spinor model the Hamiltonian form exists on the
same charged radial channel. At any real electromagnetic field
the hopping form costs $\epsilon mN+C_\epsilon(1+r^{-5/2})$.
For $j\ne0$ its radial part is infinitesimally form bounded by
the angular potential $j^2/f(r)^2$, which grows as $r^{-8}$.
The positive Coulomb and photon forms can be kept in the diagonal
sum. Thus the vacuum matrix element of the semibounded Hamiltonian
is finite and strictly positive. Compact-interval radial
exhaustion, as in Proposition~\ref{prop:physical-moving-ordered-uv},
identifies the prepared transfer and its conditional matter
functional; the estimates above dominate its marked coefficients.
H\"older's inequality over the radial law absorbs the determinant
and the finitely many other random factors. The two scalar
frequency bounds on each instantaneous photon, or the transverse
Maxwell frequency integral, give convergent photon sums as in
the preceding proof. This proves absolute convergence of the
complete dressed graphs on the stated slab.
\end{proof}

\paragraph{Arbitrary radial boundary data in the retained angular sectors.}
The ground vector in the preceding proof can serve as a reference
path measure even when the experiment uses a different boundary state.
On the exact quartic chart it suffices to keep the normal propagation
in a nonzero angular sector on the slab. There is no need to choose
the charged radial ground vector as the incoming state.

\begin{corollary}[Boundary states supported in the nonzero angular sectors]
\label{cor:seam-moving-bfv-uv}
Take the constant-$j$ representation of
Theorem~\ref{thm:seam-moving-allg-uv} with the exact quartic
metric and Dirichlet outer endpoint. Its unnormalized joint graph
bounds at each fixed order extend to any two $L^2$ radial endpoint
vectors $v,w$. Their cost is at most $C_{T,\Gamma}\|v\|_2\|w\|_2$
for the graph $\Gamma$, uniformly over integer $|j|\ge1$.
Thus the result includes arbitrary nonzero positive trace-class
radial and angular preparations supported in the direct sum of these
sectors, with a single positive vacuum normalization. The charged
operators and probes here preserve $j$ and have the same radial
matrix elements in each such sector.
\end{corollary}
\begin{proof}
Use the unitary radial coordinate of
Proposition~\ref{prop:final-collar-uv}, where $H_j=-\partial_r^2+
2r^{-2}+j^2r^{-8}$ in $L^2(\mathrm dr)$. First take $j=1$.
Let $u>0$ be its normalized ground vector of energy $E_0$,
$\pi(\mathrm dr)=u^2\mathrm dr$, and $K_T$ its heat kernel.
We show that the transition density
$k_T(r,s)=e^{TE_0}K_T(r,s)/(u(r)u(s))$ belongs to
$L^{5/2}(\pi\otimes\pi)$ for each fixed $T>0$.

For $0<c<1/3$, the eigenfunction estimate in the finite-collar
proof uses $W=\exp[c(r^{-3}-h^{-3})_+]$ with
$h^{-3}\le C_c(1+E)^{3/8}$. It gives
$\|W u_E\|_2^2\le C_c(1+E)$ for each normalized eigenfunction.
The same weighted identity also gives
$\|(Wu_E)'\|_2^2\le E\|Wu_E\|_2^2$: the angular potential
dominates $|W'|^2/W^2$. These inequalities follow first with the
bounded truncations of $W$. The one-dimensional Sobolev bound then
gives
\[
 |u_E(r)|\le C_c(1+E)
   \exp[C_c(1+E)^{3/8}-c/r^3].
\]
Sum the heat-kernel expansion at fixed $T>0$ using
$E_n\ge\pi^2 n^2$. This bounds $K_T(r,s)$ by
$C_{T,c}\exp[-c(r^{-3}+s^{-3})]$ in absolute value.
For a lower bound on $u$ take $c'>1/3$. Near the tip
$(H_1-E_0)e^{-c'/r^3}<0$, since the leading multiplier is
$(1-9(c')^2)r^{-8}$. Match a small positive multiple to $u$ at
an interior point and apply the one-dimensional maximum principle
on the tip interval, using the common vanishing endpoint and positive
potential there. Thus $u(r)\ge c_1 e^{-c'/r^3}$ near zero.
At the regular outer Dirichlet endpoint $u(r)$ vanishes linearly
by uniqueness of the second-order initial value problem.
With $q=5/2$, choose $c,c'$ so that $qc>(q-2)c'$.
Each endpoint integral of $|K_T|^q u(r)^{2-q}u(s)^{2-q}$
then converges at zero and at one. This proves the claim.

In the stationary $j=1$ process, the unnormalized radial matrix
element from $v$ to $w$ is obtained by inserting
$e^{-TE_0}(v/u)(X_0)\overline{(w/u)(X_T)}$. Its random part has
$L^{6/5}$ norm bounded by $C_T\|v\|_2\|w\|_2$: integrate the
$6/5$ power against $k_T\,\mathrm d\pi\,\mathrm d\pi$ and use
H\"older with conjugate exponents $5/2,5/3$.
The marked estimates above allow a sixth moment of the random
majorants, with constants still integrable in the loop variables.
This follows by the same maximal estimates and time Jensen with
larger fixed exponents. H\"older now gives the claimed bound for
the actual endpoints.

For $|j|\ge1$, radial Feynman--Kac relative to this same process
adds the multiplier
\[
 \exp\!\left[-(j^2-1)\int_0^T X_t^{-8}\,\mathrm dt\right]
       \le e^{-(j^2-1)T}.
\]
The multiplier depends on the trajectory alone. In particular it
has no dependence on the two half-end times of a Fourier cut.
All absolute graph bounds for $j=1$ therefore apply uniformly to
these sectors before normalization. For a direct sum with
$j$-preserving probes, Cauchy--Schwarz over the sectors bounds
each boundary matrix element. A positive trace-class preparation
then pairs with the resulting bounded operator. Its real-coupling
vacuum is strictly positive for a nonzero preparation by positivity
and injectivity of the Hamiltonian transfer in each sector.
Free vacuum energies use the common unnormalized convention.
\end{proof}

This extension concerns changing the radial and angular boundary
state on a fixed slab. A representation with a charged insertion
that shifts the angular sector must track that shift in its operator
and can pass through a neutral interval.

\paragraph{Angular shifts at the charged hops.}
A conserved angular label is sufficient for the preceding arbitrary-hopping
bound. There is also a construction in which the pair itself changes
that label. Take the scalar heat-hop Hamiltonian on the exact quartic
normal circle, multiply its pair-creation kernel by $e^{i\ell\theta}$
and use its adjoint for annihilation, with $\ell\in\mathbb Z$.
The radial kinetic term is the full angular Laplacian on a flat
circle line bundle with Fourier lattice $\mathbb Z+\vartheta$. In a
word with successive creation or annihilation signs
$\epsilon_i=+1$ or $-1$, the circle multiplication gives
$j(t)=j_0+\ell\sum_{t_i\le t}\epsilon_i$.
Suppose the allowed orbit of every incoming label has distance
at least $\delta>0$ from zero. For example, a prescribed
nonintegral holonomy $\vartheta$ gives this property for every
integer shift. It can also hold on an invariant subspace of an
untwisted bundle. This is a condition on the represented charge
action and the actual prepared module.

\begin{proposition}[Moving angular sectors with a common gap]
\label{prop:seam-moving-shifts-uv}
On the same finite torus and slab, all connected normalized electric
coefficients of this scalar heat-hop action admit joint absolute
bounds in the moving-depth representation of
Proposition~\ref{prop:seam-moving-frequencies-uv} at $2|g|<m$.
Arbitrary $L^2$ radial and angular boundary vectors on the stated
orbits are allowed in matrix elements. The vacuum normalization
uses a nonzero positive preparation. The normal angular index
can change during a closed charged return.
\end{proposition}
\begin{proof}
Expand the time-ordered hops before collapsing the charged
contractions. Between hops the angular transfer is diagonal.
Relative to the radial reference with centrifugal potential
$\delta^2/r^8$, its full effect beyond the common reference is
the factor
\[
 W=\exp\!\left[-\int_0^T
        (j(t)^2-\delta^2)X_t^{-8}\,\mathrm dt\right],
 \qquad 0<W\le1.
\]
The positive ground process for $H_\delta$
has the same moment and endpoint estimates, with the Agmon exponent
chosen below $\delta/3$. Temporal density marks and spatial Wilson
marked phases keep the interval and hopping estimates of the moving
scalar model. They do not themselves change the angular label.
Summing the allowed orderings of scalar marks on an interval thus
continues to give the product of the measured interval primitives.

The one angular factor for a word can correlate all its hop times.
For $L$ closed lines open one hop time on each, use variables
$\mathbf y,\mathbf z\in[0,T]^L$ on the two sides, and attach the
actual cut hops, including their angular shifts, to the $\mathbf y$
side. Fix the trajectory, other times and momenta. On the seam
$\mathbf y=\mathbf z$, the factor $W(\mathbf y)$ computes exactly
the chronological sequence of all angular shifts. Its extension
away from that seam may use this same function of $\mathbf y$ alone.
Multidimensional Parseval restores the seam. In absolute value
the integral in the $L$ closure frequencies is at most
\[
 \left\|W(\mathbf y)\prod_{a=1}^L A_a(y_a)\right\|_2
             \left\|\prod_{a=1}^L B_a(z_a)\right\|_2
 \le \prod_{a=1}^L \|A_a\|_2\|B_a\|_2 .
\]
Open lines and their other times may be present in $W$ as fixed
parameters in this inequality. This keeps the angular factor at
the level of the actual complete word, including for overlapping
loops. No independence of angular histories is used.

There may be arbitrarily many unmarked loops coupled by this common
angular evolution. Write $R_* = \min_{0\le t\le T}X_t$. A single
unmarked spatial heat hop on each loop gives translated lattice
sum at most $C(1+R_*^{-3/2})$, by the three-dimensional Gaussian
sum. Other unmarked hops cost one in space. The interval and hopping
bounds give a convergent geometric series per unmarked loop at
$2|g|<m$. The usual $1/L_0!$ for $L_0$ unlabelled unmarked cycles
in the expansion of the fermion vacuum bounds their full sum in
absolute value by $\exp[C(1+R_*^{-3/2})]$. Indeed the Wick cycle
expansion is valid prior to the angular average, or by fixing the
angular multipliers at their hop times. Angular ordering leaves the
permutation symmetry of the cycles intact. For a given number of
marked and open lines the remaining factors have the estimates of
Proposition~\ref{prop:seam-moving-frequencies-uv}.
The exponential in $R_*$ has all finite moments. Apply its proof's
maximal inequality to
$f(r)=\exp[b/(r+\varepsilon)^{3/2}]$, then send
$\varepsilon\downarrow0$; $f$ and $f'$ have uniformly finite
one-time moments for every $b$ by the charged tip bound.
The radial endpoint change of measure from
Corollary~\ref{cor:seam-moving-bfv-uv} now applies on each angular
word. For its net shift $d$ the sum of the endpoint norm products
$\sum_j\|v_j\|_2\|w_{j+d}\|_2$ is bounded by
$\|v\|_2\|w\|_2$. This also permits positive trace-class preparation.
The common gap bounds the scalar hopping form and gives a
semibounded Hamiltonian, so its diagonal vacuum amplitude is finite
and strictly positive. The preceding absolute estimates justify
the marked expansion and its normalization.
\end{proof}

A circle twist used in this construction belongs to its chosen
angular representation. A physical chiral multiplet or a
particle--hole tensor product may contain zero angular weight even
when each constituent is shifted. The relevant gap is that of
the full orbit carried between the insertions.

\paragraph{The charged angular interval between creation and annihilation.}
A local angular transition at a pair vertex changes the ultraviolet
test even for a globally neutral measurement. Use the same quartic
scalar heat-hop Hamiltonian and slab as in the neutral test, but
multiply the pair creator by the normal-circle character
$e^{i\ell\theta}$, with integer $\ell\ne0$, and its adjoint by
the conjugate. The angular heat generator acts on the full circle.
Start and end in the angular-zero onsite and photon vacua with the
regular free radial ground vector. A pair created at $s$ and
annihilated at $t>s$ now propagates in sector $\ell$ for the actual
duration $t-s$. The heat widths $\kappa r$ and harmonic charge
readout $e^{-r|p|}$ are the same at the two types of vertex as
in the neutral calculation.

\begin{proposition}[All electric orders of the closed two-hop return]
\label{prop:seam-switched-pair-uv}
In this preparation the complete order-$g^2$ vacuum coefficient
and its connected derivatives in smooth transverse background
probes are ultraviolet finite at every finite electric order,
including the full photon and Coulomb sewings. The estimate includes
contractible and winding returns, simultaneously short dwell time
and depth, and both orderings of the two conjugate insertions when
present. It applies on the fixed spatial torus and finite slab.
\end{proposition}
\begin{proof}
Put $\Delta=t-s$, $r=X_s$, $r'=X_t$ and $d=\min(r,r')$.
Use the unitary radial coordinate in $L^2(\mathrm dr)$.
Between the two hops the radial kernel is that of
$H_\ell=-\partial_r^2+2r^{-2}+\ell^2r^{-8}$. The two outer radial
ground vectors are bounded on the interval. At quadratic hopping
order there is one intermediate pair; the vacuum logarithm has no
product of linear vacuum terms. Its weight in a given pair of
lifted spatial chords $z,z'$ includes
$G_{\kappa r}(z)G_{\kappa r'}(z')$ with $z+z'$ a period.
The electric variable $e^2$ multiplies the sum of half the
transverse Gaussian variance of these two oriented chords and
the full positive Coulomb integral. Every power of the latter
can be absorbed by half the angular killing factor. Indeed
$C_x(z)\le C(1+x^{-1})$ whereas the angular cost is
$\ell^2\int_s^t X_u^{-8}\,\mathrm du$. For fixed slab duration
and each fixed power this comparison follows from $x^{-1}\le
1+x^{-8}$ on the interval and boundedness of $y^n e^{-c y}$.
It suffices to bound all moments of the transverse variance with
the remaining angular strength $\ell^2/2$.

We may restrict to $d$ near zero. For $\Delta\ge d^5$ the
remaining kernel is bounded by
$C\Delta^{-1/2}\exp(-c/d^3)$ uniformly in the other endpoint.
By symmetry start at the smaller endpoint. Over the first
$d^5/2$ time units a surviving path staying below $2d$ pays
at least $c/d^3$ in angular action. Reaching $2d$ has probability
at most $2\exp(-c'/d^3)$ under the free Brownian majorant. The rest
of the kernel costs $C\Delta^{-1/2}$. For $\Delta<d^5$ and
$\max(r,r')>2d$, the free kernel bound
$C\Delta^{-1/2}\exp[-(r-r')^2/(4\Delta)]$ supplies the same
exponential. In either region it absorbs all inverse depth powers.
For example the transverse variance is everywhere bounded by
$C(1+d^{-2})(|z|+|z'|)^2$ from the equal-time Maxwell sum.
All polynomial chord moments in the heat product, summed over
lifts, cost at most an inverse power of $d$.

There remains $\Delta<d^5$ with $r,r'\in[d,2d]$. Nonzero
$z+z'$ costs $\exp(-c/d)$ from the two heat densities: reserve
half each Gaussian and use $|z|^2+|z'|^2\ge |z+z'|^2/2$.
The remaining lifts and chord moments are summable. For $z'=-z$
the two currents cancel at coincident time and depth. The exact
transverse covariance has equal-time weight $1/(2|p|)$ and time
factor $e^{-\Delta |p|}$. The variance of their sum is bounded by
\[
 C |z|^2\left[(r-r')^2 d^{-4}+\Delta d^{-3}\right].
\]
In each mode the difference of the readouts is bounded by
$|p||r-r'|e^{-d|p|}$ and the remaining time difference by
$1-e^{-\Delta|p|}\le\Delta|p|$. The transverse projection and
the modulus of the straight-chord average are at most one.
The spatial lattice sums of the resulting multipliers give the
displayed powers. This estimate uses the full covariance, including
the cross term between the two visits.

The integral of $|z|^{2n}$ against the two reverse heat densities
is at most $C_n d^{n-3/2}$. For the free radial heat upper bound,
integration in the larger endpoint of $|r-r'|^{2n}$ costs at most
$C_n\Delta^n$; without that power it costs a constant. Thus the
$n$th transverse moment after these integrations is bounded by
\[
 C_n d^{-3/2}\left[(\Delta d^{-3})^n+
                                    (\Delta d^{-2})^n\right].
\]
This is integrable in $\mathrm dd\,\mathrm d\Delta$ over
$0<\Delta<d^5$. The estimates also apply at $n=0$. Smooth spatial
and temporal external probes at the two Wilson links contribute
at most fixed polynomial chord factors. Finally the positive
transverse variance itself is the integral of the squared complete
Fourier current in the free Maxwell covariance; this retains the
relative signs of the two visits before taking the absolute bound.
The scalar covariance gives the full Coulomb square by its canonical
time sewing. This proves the assertion for the complete return.
\end{proof}

\paragraph{An open charged interval at a physical boundary.}
The preceding cancellation uses both ends of a return. On the same
angular-shifting action prepare the initial angular-zero radial ground
function $\psi_0$ in $L^2(r^4\mathrm dr)$ and, at the final end,
one particle and one hole in constant spatial waves and the photon
vacuum. Give their angular-$\ell$ radial state the normalized function
proportional to $r^2\psi_0(r)$. This is a finite-energy state:
its angular contribution at the tip is bounded by a constant times
$\int_0^{r_0} r^4 r^4 r^{-8}\,\mathrm dr$, and the derivative
and Coulomb form terms are finite as well. One creation at time $s$
now has no return chord. At the coefficient linear in $g$ it gives
the positive Laplace transform of the one-chord variance and the
full Coulomb energy, with the additional radial weight
\[
 X_T^2\exp\!\left[-\ell^2\int_s^T X_u^{-8}\,\mathrm du\right]
\]
relative to the angular-zero radial ground process. The other factors
are those in Proposition~\ref{prop:physical-neutral-tip-uv}.
For $r=X_s$ small, restrict to $0<T-s<c r^8$ and to radial paths
staying in $[r/2,2r]$ until $T$. The conditional probability of
this event is bounded below uniformly: in the five-dimensional ball
a Brownian displacement smaller than $r/2$ throughout this time has
probability tending to one, and the ground-state density change
is bounded above and below on the event. The angular weight is then
at least a constant times $r^2$. The photon variance has Gaussian
chord mean at least $c' r^{-1/2}$. Its $n$th moment is bounded below
by the $n$th power of that mean. The time window supplies $r^8$,
and the stationary incoming density is comparable to $r^4\mathrm dr$.
The result is a moment lower bound proportional to
\[
 \int_0^{r_0}r^{14-n/2}\,\mathrm dr .
\]
This diverges at $n=30$. The coefficient of $g e^{60}$ in this
prepared transition consequently diverges, with the Coulomb and
Wilson terms already included. For these endpoints the threshold is exact. The complete
$n$th moment is bounded by a constant times the sum of the pure
Coulomb moment and the pure chord moment. Absorb the Coulomb moment
into half the angular killing as above. The heat-averaged chord
moment costs at most $C_n(1+r^{-n/2})$ by
Proposition~\ref{prop:physical-neutral-tip-uv}.
Extending the remaining positive charged propagation time to infinity
bounds its radial factor, in the original $r^4\mathrm dr$ coordinate,
by $U=(H_{\ell/\sqrt2}+2m-E_0)^{-1}(r^2\psi_0)$ up to a
fixed constant. Here the subscript gives the reduced strength of the
potential; it need not be a Fourier label. The shift is invertible
since $H_{\ell/\sqrt2}\ge H_0$ and $m>0$.
Near zero $U(r)\le C r^{10}$. Indeed applying this shifted radial
operator to $r^{10}$ gives $(\ell^2/2)r^2+O(r^8)$ there.
Choose a large multiple to dominate the forcing and the value of
$U$ at a fixed interior point, then compare in the Friedrichs form
using the positive part of the difference. The resulting tip
upper bound is $C_n r^{14}(1+r^{-n/2})\,\mathrm dr$.
Every coefficient through $g e^{58}$ is therefore finite and
$g e^{60}$ is logarithmic. The same bounds apply in the stationary
ground-projection derivative: the exact charged time integral then
runs to infinity with pair weight $e^{-2m\tau}$, and the interval
$0<\tau<c r^8$ still yields the lower bound. For the Coulomb term
in this extension, $\int(1+X_u^{-1})\,\mathrm du$ is bounded by
a constant times $\int X_u^{-8}\,\mathrm du$ on the unit interval
in $r$, so absorption is uniform in the duration.
Vacuum normalization has no linear hopping term. Thus propagation in a genuinely charged angular sector for
a positive but shrinking time does not on its own control a
one-sided boundary insertion on every finite-energy radial state.
Whether such endpoint probes occur in the physical chiral reduction
must be decided by its boundary domain.

\paragraph{What happens when two charged pairs can exchange partners?}
A nonzero angular potential throughout every occupied interval
protects a retraced two-hop chord in the original $a(r)=r$,
$\tau(r)=\kappa r$ model. The four-hop test is different. Use the
scalar fermion hopping of~\eqref{eq:cl-gauss-charge-uv} with one
particle species and one hole species and attach $e^{i\ell\theta}$
to pair creation, with its adjoint for annihilation, for a fixed
integer $\ell\ne0$. Start with the neutral radial ground state and
angular index zero. On every interval the angular index is then
$j=\ell N_{\rm pair}$. Keep the full transverse and Gauss fields on
the finite spatial torus and vacuum slab. The following coefficients
refer to the joint expansion in hopping $g$ and electric charge $e$.

\begin{proposition}[A contractible four-hop exchange obstruction]
\label{prop:seam-uv-two-pair}
In this switched angular model the bare vacuum logarithm has a
divergent joint coefficient at $g^4e^{110}$. For some smooth real
transverse backgrounds of nonzero spatial momentum the normalized
connected two-insertion response has a divergent joint coefficient
at $g^4e^{120}$. Each divergence is present on a contractible
quadrilateral. The lower estimates below give at least a logarithmic
divergence.
\end{proposition}
\begin{proof}
At four ordered hop times a vacuum history has number signs
$+ - + -$ or $++--$. The first consists of two retraced pairs. In
the second there are two choices for the order of particle
annihilation and two for the holes. Their fermionic sign is the
product of the two permutation signs, with the convention that
retracing both pairs has positive weight. This also follows from
the product of the two $2$-by-$2$ antisymmetric position
contractions. The two equal permutations retrace the two original
pairs. The two different permutations give a quadrilateral, each
with negative weight. Once the positions are fixed the diagonal
evolution between hops depends only on the occupied sites, the
radial coordinate and the photon field. Thus the positive spatial
weights are products of four heat densities. The radial path gets
the factor
\[
 W_\theta=\exp\!\left[-\ell^2\int N_{\rm pair}(v)^2
                                      r_v^{-8}\,\mathrm dv\right].
\]
The Maxwell and Coulomb integrals multiply each signed term by
$e^{-e^2 Y}$, where $Y\ge0$ is half the transverse variance of the
sum of the four currents plus the integrated full Coulomb energy
per $e^2$. The latter is bounded by $C\int_0^T(1+r_v^{-1})\mathrm dv$
also with two pairs. A transverse test contributes the square
$L_b^2$ of the oriented line integral to the second variation, with
the common prefactor $-e^2$. Unmarked spatial masses are finite
after radial averaging: a retraced pair has the periodized
two-heat convolution at zero, and a quadrilateral has the
four-heat convolution at zero, up to fixed volume factors.

We first establish finiteness of the direct terms at every electric
order, including coincident times. For one retraced pair let $s,t$
be its hop times, $\Delta=t-s>0$, $r=\min(r_s,r_t)$. Its two-chord
mass is bounded by $C[1+(r_s+r_t)^{-3/2}]$. Conditional on their
net winding, the two lifted displacements have Gaussian bridge
laws. Their winding vector has weights proportional to
$\exp[-|w|^2/(4\kappa(r_s+r_t))]$ on the period lattice. Hence
all fixed displacement moments are uniformly bounded in the
normalized spatial law. If both depths are comparable to small
$r$, displacements with $w\ne0$ cost $O(e^{-c/r})$ even with
fixed polynomial moments. In general the transverse variance
of the retraced currents is bounded by
$C(1+r^{-2})(|z_1|^2+|z_2|^2)$ from the two equal-time covariances.
A bounded smooth classical test costs at most $C_b(|z_1|+|z_2|)$.

For small $r$ use as starting endpoint the one attaining the minimum;
the stationary unperturbed radial law is reversible. If
$\Delta>r^5$, a path remaining below $2r$ for the first $r^5$
units of the occupied interval pays $e^{-c/r^3}$. The probability
of exiting above $2r$ in that time is at most $Ce^{-c/r^3}$,
by the five-dimensional Brownian maximal bound. The bounded
ground-state transform at such an exit changes only the constant,
since the starting ground function is bounded below near zero.
If $\Delta\le r^5$, an endpoint change exceeding $r^2$ likewise
has conditional probability at most $Ce^{-c/r}$.
On the remaining event the depths differ by at most $r^2$.
For opposite lifts $z,-z$ the time covariance in each mode gives
the variance bound
\[
 C|z|^2\left[\frac{|r_s-r_t|^2}{r^4}
                    +\frac{\Delta}{r^3}\right].
\]
Indeed the squared difference of the two readouts costs
$|r_s-r_t|^2 |p|^2 e^{-2r|p|}$ and the remaining time difference
costs $2\Delta |p|e^{-2r|p|}$. Sum with the equal-time multiplier
$1/(2V|p|)$. The bound now costs at most $C|z|^2$.
This proves uniform radial moment bounds of all finite orders for
a direct pair's normalized spatial variance with its occupied
angular weight, after averaging. The same reasoning works for any
positive multiple of that angular potential. To restore the two
spatial masses observe that their product is uniformly in $L^p$
of the stationary radial law for $1<p<5/3$, by Cauchy--Schwarz
and the $r^4\,\mathrm dr$ marginal at the first endpoint of each
pair. The combined angular potential dominates the sum on the
two individual occupied intervals. H\"older's inequality, the
radial exponential bound for the integrated Coulomb potential and
the triangle inequality for Maxwell currents now give all direct
moments even when the two pairs overlap. The same bounds apply
to the two-hop terms and their products in the logarithm.

For a crossed term restrict the first time to an interior slab
interval. Given its depth $r\ll1$, take each of the three subsequent
time gaps between $r^8$ and $2r^8$. With conditional probability
bounded below, the radial path stays within $r^3$ of its starting
depth throughout these gaps, by the Brownian maximal estimate and
the ground transform. The angular weight then has a strictly
positive lower bound since $N_{\rm pair}\le2$. Restrict the four
sites to a small neighborhood, parameterizing their three relative
displacements in units of $\sqrt r$. In these units choose a
nondegenerate quadrilateral with no two edges parallel, nonzero
oriented area in a chosen coordinate plane, and a compact
neighborhood retaining these properties. Take the lifts in this
neighborhood. The spatial volume element contributes $r^{9/2}$
and the four positive heat densities contribute at least $c r^{-6}$.

The oriented currents on this quadrilateral have transverse
variance at least $c r^{-1/2}$. To see this, take a small patch
of directions perpendicular to the first edge and bounded away
from perpendicularity to each of the other three. In the shell
$1/r<|p|<2/r$ restrict $|p\cdot z_1|$ to a sufficiently small
constant. A rectangular box inside this region can be chosen with
side lengths comparable to $r^{-1},r^{-1},r^{-1/2}$, uniformly on
a smaller neighborhood of the quadrilateral. It contains at least
$c_R r^{-5/2}$ reciprocal-lattice modes by elementary lattice-cell
counting. In each mode the transverse projection, smeared first
chord and its sinc factor give norm at least $c\sqrt r$ before
the equal-time covariance multiplier. The other three chord
norms in that same mode are $O(r)$, since
$|p\cdot z_i|\ge c/\sqrt r$. Applying the reverse triangle
inequality in the Gaussian norm of this mode, at the actual four
times, leaves a squared contribution at least $c r^2$ including
the multiplier. The full Coulomb energy is nonnegative.

For a connected test choose a transverse sine mode with a time
window constant during all four hops, and restrict the base point
so that its curl pairs nontrivially with the oriented area.
Stokes' formula on these contractible straight segments gives
$|L_b|\ge c r$ for small $r$. The spatial readout for the fixed
test mode changes each chord by $O(r^{3/2})$; the lower bound
survives those four errors. This is a physical transverse
covariance insertion using the source shift described in
Proposition~\ref{prop:seam-closed-uv}.

The crossed $n$th moment of $Y$, with either no marks ($d=0$)
or the two background marks ($d=2$), is thus bounded below by a
positive multiple of
\[
 \int_0^{r_0} r^{4+24-3/2+d-n/2}\,\mathrm dr .
\]
For $d=0,n=55$, and for $d=2,n=59$, the exponent is $-1$.
The other crossed histories have the same sign in the respective
moment. Monotone electric finite differences on the positive
crossed measure give the same conclusion if an earlier moment
already diverges. The direct terms and the lower-order products
in the vacuum logarithm have finite moments. Normalization
therefore cannot cancel the marked divergence. The proof uses a
finite hopping order and does not determine the fixed-$g$,
finite-time answer of this switched Hamiltonian.
\end{proof}

The angular barrier controls the length of each occupied interval.
Four sufficiently rapid exchanges can still trace a non-retraced
spatial polygon whose independent transverse displacement is resolved
on wavelengths of order $r$. Testing only two opposite traversals
would miss this contribution.

The four-hop thresholds can in fact be located exactly in the
switched scalar model. The direct terms have already been bounded
to every order. For a crossed term set $t_1<t_2<t_3<t_4$ and choose
one index $i$ attaining the least of the four endpoint depths,
$r=r_{t_i}$. Sum over the four possible choices in the bound.
The pair is occupied on the entire interval $[t_1,t_4]$. At fixed
radii the four Gaussian step densities on the oriented cross cycle
give a mass at most $C(1+r^{-3/2})$ after spatial closure. Indeed
the convolution at zero is periodized with lapse
$\kappa\sum_j r_{t_j}$. In the resulting normalized Gaussian
bridge the moments of every fixed displacement power are bounded,
uniformly for $r_{t_j}\le1$. Each spatially integrated moment of
the Wilson cost and bounded smooth probe therefore costs at most
some inverse power of the least radius. The Coulomb time integral
can be treated separately using the power bound for a sum of
nonnegative costs. Its moments are finite to every order even with
the spatial mass and the probe: apply H\"older using the radial
exponential estimate and the fact that the spatial mass has a
moment strictly greater than one, uniformly in the four times.

At small $r$ write $h=r^5$. If $t_4-t_1>2h$, at least one
side of $t_i$ within the occupied interval extends for a duration
$h$. Start there at radius $r$, using reversibility on the past
side. A trajectory staying below $2r$ for the next $h$ pays
$\exp[-c h/r^8]$ in angular energy, and an exit above $2r$ on
that time scale costs at most $C\exp[-c r^2/h]$ in probability.
Both bounds absorb every fixed inverse power of $r$. If
$\Delta=t_4-t_1\le2h$, an excursion outside $[r/2,2r]$ within
this interval costs at most $C\exp[-c r^2/h]$, treating each side
of $t_i$ separately. All these estimates use only the stationary
ground process conditioned on $r_{t_i}=r$. Restricting to the
event that this radius is the least at the four marked times
gives the polynomial majorant above.

On the remaining paths the four depths are comparable and the
angular weight is at most $\exp[-c\Delta/r^8]$. At fixed $r,t_i$
the integral in the three other times against this exponential
is at most $C r^{24}$: rescale their three time differences from
$t_i$ by $r^8$ and use $\Delta\ge\max_j|t_j-t_i|$.
In the normalized four-step spatial bridge,
nonzero winding costs $O(e^{-c/r})$ including any fixed polynomial
in the displacements. At zero winding the conditional Gaussian
increments have absolute moments of order $k$ bounded by $C_k r^{k/2}$.
A single transverse variance then has $n$th moment at most
$C_n(1+r^{-n/2})$, by the short-chord covariance bound and the
Gaussian estimate on long lifts. The same bound holds for half
the total variance, by the triangle inequality. For $L_b$, subtract
the same value of $b$ at the base point from every line integral;
the four constant contributions cancel. Lipschitz continuity in
space and time and smoothness under the small harmonic readout
give
\[
 |L_b|\le C_b\big[ Z^2+(2h+2r)Z\big],\qquad
 Z=\sum_{j=1}^4 |z_j|,
\]
for the smooth mode used above (and for smooth tests with bounded
first time, spatial and Poisson derivatives).
Thus $\mathbb E_{\rm bridge}|L_b|^4\le C_b r^4$.
Cauchy--Schwarz against the variance moment gives the extra
factor $r^2$ for the two-insertion response. After integrating
$t_i$ in the slab and using its stationary one-time law, the
remaining upper bound at the tip is proportional to
\[
 \int_0^{r_0}r^{4+24-3/2+d-n/2}\,\mathrm dr ,
\]
with the same $d$ as in Proposition~\ref{prop:seam-uv-two-pair}.
It is finite for $n<55$ in the vacuum and $n<59$ in the
two-insertion measurement. Hence at hopping order four the
first failures are precisely $g^4e^{110}$ and $g^4e^{120}$,
respectively, for the stated preparation and probe. At each threshold
the divergent crossed moment on $\min_j r_{t_j}\ge\epsilon$ is
bounded above and below by positive multiples of $\log(1/\epsilon)$
as $\epsilon\downarrow0$.

\paragraph{Which geometric length is measured by a neutral charged pair?}
In the moving two-band action the spatial heat step and the electric
readout have distinct operational meanings. The heat step fixes
the distribution of the chord; the other length fixes the
resolution of the photon field sampled along it. We can test both
lengths in the full neutral radial sector. Use the quartic normal
metric, the scalar off-diagonal hopping and the finite spatial
torus of Proposition~\ref{prop:physical-neutral-tip-uv}. At depth
$r$ let the heat lapse be $\tau(r)=\kappa r^\alpha$ and the
electric readout length be $a(r)=a_*r^\gamma$, with positive
parameters in reference units. The original homogeneous
displacement has $\alpha=1$.

There is a local variational realization of these two functions.
Transport the charged displacement by tangential heat generation
$\tau'(r)\Delta_x$ per unit of normal distance. The independent
Gaussian increments from zero to $r$ then have total lapse
$\tau(r)$. On each carrier slice take the positive normal Maxwell
form for the harmonic readout, componentwise in the Abelian chart,
to be proportional to
\[
 \int_0^\infty\!\mathrm dr\,\mathrm d^3x
 \left(\frac{|\partial_r B|^2}{a'(r)}
              +a'(r)|\nabla_x B|^2\right).
\]
Continue the collar so $a(r)\to\infty$ at its far end. In the
coordinate $y=a(r)$ this is the ordinary half-space Dirichlet
form. With given data at zero, the finite-energy extension of each
nonconstant mode reaches the sampling point with multiplier
$\exp[-a(r)|p|]$. Use it for both electric potentials and for the
gauge parameter, with the induced radial connection for a moving
charge. The gauge-fixed transverse chart then has the same Gauss
square and Wilson contacts as before, at the new sampling length.
In particular one may keep $\tau(r)=\kappa r$ and take
$a(r)=\sqrt{\kappa r}$, the standard deviation of a single heat
coordinate up to a fixed factor. This choice changes only the
photon constitutive form, and introduces no additional adjustable
width. The metric of the radial state remains fixed. A coordinate
change applied to every block would not change their two physical
lengths relative to each other.

\begin{proposition}[All finite electric orders for a neutral moving depth]
\label{prop:optical-neutral-uv}
Suppose $0<\alpha<4/3$ and $0<\gamma\le\alpha/2$. With $2|g|<m$
on the finite vacuum slab, the scalar hopping model has a positive
finite bare normalization at every real electric charge. Its
complete ordered vacuum and tested photon graphs converge
absolutely at every finite electric order, jointly in charged
spatial loop momenta, photon four-momenta, hopping times and
radial paths. The radial preparation may include the entire
neutral angular channel $j=0$. Nonnegative nonzero radial
$L^2$ endpoint functions in this channel are allowed.
The same bounds apply to open fermion chains with finitely many
smooth spatial endpoint tests at the slab boundaries and the
canonical endpoint charges.
\end{proposition}
\begin{proof}
On a fixed radial path retain all temporal mark orderings of the
onsite vacuum determinant as in
Proposition~\ref{prop:physical-moving-ordered-uv}. A spatial
Wilson hop with $d$ electromagnetic marks is an average of
translated $d$th derivatives of $e^{-\tau(r)|k|^2}$. At the hop
at time $t_i$ it has the absolute majorant
$C_d\tau_i^{d/2}e^{-\tau_i|k+u|^2/2}$ under that average, for
the corresponding translate $u$, where $\tau_i=\tau(r(t_i))$.
Each transverse photon joining two such marks therefore costs at
most
\[
 C\sqrt{\tau_i\tau_l}
       \left(1+(a_i+a_l)^{-2}\right)\le C' .
\]
Here $a_i=a(r(t_i))$. Integrating the photon frequency gives
$\pi/|p|$, and spatial shell counting with
$e^{-(a_i+a_l)|p|}$ gives the displayed power. The last inequality
uses $\sqrt{\tau(r)}/a(r)\le C$ and
$a_i a_l/(a_i+a_l)^2\le1/4$. It applies also to two marks on the
same visit. A fixed external spatial mark only costs a bounded
power of $\sqrt\tau$.

For a temporal photon on two charged intervals, Plancherel in its
own frequency again bounds its integral and full spatial sum by
$C\int_0^T[1+r(t)^{-\gamma}]\,\mathrm dt$. Keep one charged
spatial momentum on each closed word. One marked or unmarked
Gaussian at an anchor hop sums over that momentum at a cost at
most $C[1+r(t_i)^{-3\alpha/2}]$, uniformly in all other transfers.
An open word needs only supremum bounds for its spatial factors.
Each distinct closed word has its own anchor time. At fixed
trajectory integrate all other durations using the one-sided
onsite exponential; the free time integral per additional step
is at most $2/m$. The integrations over the distinct anchor times
are then bounded by powers of
\[
 A_T=\int_0^T [1+r(t)^{-3\alpha/2}+r(t)^{-\gamma}]
                  \,\mathrm dt .
\]
At each fixed electric order the extra number of marks is fixed,
and the count of their placements is polynomial in each word
length. Thus summing the word lengths converges for $2|g|<m$.
The absolute log weight of all unmarked closed words is bounded
by $C A_T$. This gives an integrable majorant
$C_M(1+A_T)^{k_M}e^{C A_T}$ for the complete coefficient,
including the zero-charge determinant. At arbitrary real electric
charge the unexpanded phases have unit modulus and the real
conditional determinant is strictly positive, by the same
reversal of closed words. Smooth external charged endpoints give
summable spatial Fourier coefficients; the open words have no
charged spatial trace and no vacuum determinant of their own.

It remains to prove integrability with the actual radial law.
For $j=0$ the kinetic form in $L^2(r^4\mathrm dr)$ obeys
$\int r^2|u|^2\mathrm dr\le(4/9)\int r^4|u'|^2\mathrm dr$
by integration by parts on the Friedrichs core. For every
$0<s<2$ and every $b>0$, $b r^{-s}$ therefore has arbitrarily
small relative form bound. The potential in $A_T$ has only
such exponents. For any fixed $b$, the Feynman--Kac radial
matrix element weighted by $e^{b A_T}$ is consequently finite:
truncate the positive weight, apply the uniform lower semibound
of the tilted generator, and take the monotone limit. Arbitrary
nonnegative $L^2$ endpoints cost their two norms. Positivity of
the radial kernel and the conditional determinant gives a nonzero
vacuum denominator for nonzero endpoint functions. The argument
also controls signed or complex endpoint tests in absolute value
by using their moduli, whenever normalization uses a separate
nonzero vacuum. This completes the estimate for coincident visits
and for overlapping Coulomb intervals as well.
\end{proof}

\paragraph{Why shortening the charged step alone can destabilize the vacuum.}
For the same neutral pair transition at first hopping order,
when $d=\gamma-\alpha/2>0$, the first divergent electric power
in the radial ground preparation is $e^{2n}$ with
$n=\lceil5/d\rceil$, provided $\gamma<2$. To see the lower
bound, average the equal-time transverse chord variance over
its Gaussian displacement, and keep the shell
$a^{-1}\le |p|\le2a^{-1}$. As in
Proposition~\ref{prop:physical-neutral-tip-uv}, the mean is
bounded below by a constant times
\[
 \frac{\tau/a^2}{1+\sqrt\tau/a}\asymp r^{-d}.
\]
The positive trace includes all Wilson and Coulomb terms. Jensen's
inequality gives the lower moment cost $r^{-dn}$. For the upper
bound the short-chord kernel estimate in that same proof gives
$W_a(z)\le C(|z|^2+|z|/a)$ inside the injectivity radius; longer
lifts are bounded by a power times a Gaussian tail. The $n$th
Gaussian moment is consequently $O_n(1+r^{-dn})$.
The Coulomb time integral has all moments by the preceding radial
form estimate for $s=\gamma$. The ground marginal is comparable
to $r^4\mathrm dr$ at the tip, which gives exactly the threshold
claimed. If $d\le0$ the same calculation bounds every one-hop
coefficient.

Retain now the ordinary optical distance $a(r)=r$ and consider
increasing the heat-lapse exponent to $\alpha\ge2$. The one-hop
Wilson test then passes, but the zero-electric-charge Hamiltonian
itself is unbounded below at fixed nonzero hopping. In fact this
holds whenever $3\alpha/2>2$. Localize a normalized neutral
radial test in $\delta<r<2\delta$. A rescaled smooth test has
radial kinetic cost $O(\delta^{-2})$. For each spatial mode with
$|k|\le c\delta^{-\alpha/2}$ the hopping magnitude
$|g|e^{-\kappa r^\alpha|k|^2}$ is bounded below on the support
by a fixed $b>0$. There are at least $c_V\delta^{-3\alpha/2}$
such modes. In each particle--hole pair use the normalized
superposition of the empty and doubly occupied states with mixing
angle $\theta>0$ and relative sign opposite to $g$.
Its bare energy expectation, relative to the empty onsite sea,
is at most $2m\sin^2\theta-2b\sin\theta\cos\theta<0$ at any
sufficiently small fixed $\theta$. Take their product and the
onsite vacuum on every other mode. This finite-particle test is
independent of $r$ on the support. No derivative of its mixing
angle enters the radial energy. With the free photon vacuum the
total energy is bounded above by
$C\delta^{-2}-c_{V,g,m}\delta^{-3\alpha/2}$, which tends to
minus infinity. This argument requires no regularization of the
Hamiltonian. The critical stability case $\alpha=4/3$ is not
needed: at $a=r$ the positive one-hop test still diverges there.
Thus every positive power-law lapse with the original readout
$a=r$ either fails a neutral complete coefficient or exhibits
vacuum collapse in this family. The modified optical curvature
form in Proposition~\ref{prop:optical-neutral-uv} resolves this
specific conflict. A physical chiral identification must derive
the relative calibration from its own curvature pairing and
retain the same boundary and charge terms.

\paragraph{The measure of the optical field in the represented action.}
For these estimates the optical field on a fixed normal half-space
is the harmonic extension of the boundary connection. The depth
$r$ specifies where that extension is evaluated. In particular,
the optical field has a precise measure in the finite-slab
Hamiltonian. If $\mathcal E$ denotes the extension map on smooth
boundary data, its joint law with a boundary photon of law
$\mu_{\rm ph}$ is the pushforward of $\mu_{\rm ph}$ under
$A\mapsto(A,\mathcal E A)$. Hence evaluating at a radial
trajectory introduces no second photon determinant. At positive
sampling distances the covariance kernel defines the extension
also for Gaussian Maxwell data. The spatial heat law in the hop,
the ordinary charged determinant relative to the occupied onsite
slots, and the radial transfer between prepared endpoints remain
in the integral. This describes the complete sources used below.
If an unrestricted BFV integral independently quantizes bulk
fluctuations about the harmonic extension, their covariance,
determinant and source contacts have to be calculated in that
integral. The variational equation for the mean field alone cannot
replace that calculation.

\paragraph{Arbitrary hopping in the moving neutral and charged channels.}
One can dispense with the small-hopping expansion by gauging
the scalar source before differentiating the filled-vacuum
determinant. On a finite slab this admits a time primitive even
at zero carrier frequency. The argument works for a spinor hop
as well as for the scalar one.

\begin{theorem}[Dressed ordered graphs for all normal angular sectors]
\label{thm:optical-moving-dressed-uv}
Keep the spatial torus, the real hopping amplitude and the two
geometric transfer lengths of Proposition~\ref{prop:optical-neutral-uv}.
For the scalar two-band hop take $0<\alpha<4/3$. For the spinor
hop with Fourier symbol $v\boldsymbol\alpha\cdot k\,e^{-\tau(r)|k|^2}$
from~\eqref{eq:reference-spin-uv} take $0<\alpha<4/5$. In both
cases let $0<\gamma\le\alpha/2$. In every angular sector the
moving-depth Hamiltonian is semibounded and self-adjoint at every
finite real hopping strength and electric charge. The finite-slab
onsite vacuum normalization is finite and positive. At every finite
electric order the complete ordered vacuum, photon and elementary
open charged graphs admit absolute joint bounds in the charged
spatial momenta, photon four-momenta, marked times and radial paths.
Their radial preparations may be nonnegative, nonzero $L^2$ endpoint
functions. The proof permits smooth spatial charged endpoint probes
on the slab, with their charge phases retained.
\end{theorem}
\begin{proof}
Write $b=0$ for the scalar kernel and $b=1$ for the spatial derivative
of the heat kernel in the spinor model. Denote their gauged
off-diagonal one-particle hopping by $K_t$. The Schur test on its
Wilson kernel and its squared Hilbert--Schmidt spatial bound give,
uniformly in real gauge fields,
\[
 \begin{aligned}
 \|K_t\|&\le C\big(1+r(t)^{-b\alpha/2}\big),\\
 \|K_t\|_{\rm HS}^2&\le C_V
      \big(1+r(t)^{-(3+2b)\alpha/2}\big).
 \end{aligned}
\]
The constants include the chosen hopping strength. Indeed the
absolute kernel is bounded by $C\tau^{-b/2}$ times a heat
kernel at doubled lapse. Start the Hamiltonian on the intersection
of the radial, Maxwell, $mN$ and positive Coulomb form domains.
It is dense, for example on finite-particle test vectors localized
away from the radial endpoints with smooth spatial data. The
pair-creation form estimate in Proposition~\ref{prop:cl-gauss-charge-uv}
costs an arbitrarily small multiple of $mN$ plus a constant times
$\|K_t\|_{\rm HS}^2$. Its radial singularity is infinitesimally
form bounded because $(3+2b)\alpha/2<2$. Thus the closed perturbed
form defines the claimed Hamiltonian. The angular potential is
nonnegative and the neutral-channel Hardy bound suffices for every
$j$. In particular the construction does not require a gap in
the normal distance sampled by a state.

Now condition on the full radial trajectory. For smooth real scalar
test fields carry out the time-dependent change of the charged
variables with phase given, up to the temporal connection sign, by
\[
 \chi_t(x)=\int_0^t
    \big(e^{-a(r(s))|\nabla|}\phi(s)\big)(x)\,\mathrm ds .
\]
The scalar term in the carrier derivative disappears. The spatial
hop acquires $e^{iq(\chi_t(x)-\chi_t(y))}$ and the open endpoints
acquire their compensating phases. The onsite filled vacuum is
invariant. This conditional identity differentiates no radial
sample path. The phase is a time integral of the actually sampled
potential. With the instantaneous real Coulomb auxiliary field,
its interval integrals exist in Gaussian mean square for almost
every radial path by the bound on $\int_0^T(1+r^{-\gamma})$.
One may first use smooth sources and pass to their Wilson
increments and contact derivatives.

Extend the hopping by zero outside the slab on the full carrier
line. Write $R_0=(\partial_t+m\beta)^{-1}$ for the onsite vacuum
inverse and $D=R_0^{-1}+K$. The canonical matter vacuum factor
equals $\det_2(I+R_0K)$: the two diagonal onsite slots have
off-diagonal hopping, hence the single-hop trace vanishes
identically. Squaring $R_0K$ gives trace-class block products, and
Schur elimination gives the ordinary determinant on one slot, also
when the two slots each carry spin. This is the same relative
determinant as in Proposition~\ref{prop:native-relative-density-uv}.
The integral kernel of $R_0$ is bounded by a one-sided exponential
in each slot, giving
\[
 \|R_0K\|_{\rm HS}^2\le C_m\int_0^T\|K_t\|_{\rm HS}^2\,\mathrm dt .
\]
The field-independent determinant bound
$|\det_2(I+X)|\le e^{\|X\|_{\rm HS}^2/2}$ therefore has finite
radial expectation at each finite real coupling. At fixed
trajectory away from the tip the operator
$\beta D=m+\beta\partial_t+\beta K$ is a positive constant
plus a skew-adjoint operator on the onsite carrier Sobolev domain,
so $R=D^{-1}$ exists and $\|R\|\le1/m$. The matrix which
anticommutes with both $\beta$ and the scalar hopping matrix
(or with $\beta$ and the three spinor hopping matrices) takes
$D^*$ to $-D$. The corresponding relation for $R_0$ and cyclicity
of the nonzero spectrum of the two product orders make
$\det_2(I+R_0K)$ real. Invertibility excludes real zeros, so
continuation from zero hopping makes it positive. These statements
can first be used on a compact radial subinterval. There the
time-sliced transfer agrees with the quasifree vacuum matrix
element by Duhamel expansion, as in
Proposition~\ref{prop:physical-moving-ordered-uv}. Exhaustion by
the radial Friedrichs cores and the exponential majorant above
give the full vacuum integral. For nonnegative radial endpoints
one may equivalently use killed radial kernels directly with their
Feynman--Kac weights.

For the electric derivatives evaluate the resolvent at zero
electromagnetic source but at the given hopping. It then
preserves spatial momentum, while its time kernel need not be
translation invariant. At every spatial $k$ the identity
\[
 R=R_0-R_0 K R_0+R_0 K R K R_0
\]
and the $L^2$ operator bound $1/m$ give a uniform time-kernel
majorant $C_T(1+\int_0^T r(s)^{-b\alpha}\,\mathrm ds)$. At the two
outer ends of the last term use the $L^2$ time norms of
$e^{-m|t-s|}\|K_s\|$. The onsite equal-time jump is diagonal and
has zero trace against a single off-diagonal hop. Equivalently,
in the single-mark trace use $(R-R_0)K'$ by differentiation of
the canonical relative determinant. Products of inverse and vertex
kernels then have the stated time bound even at contact.

After differentiating the gauge phases, every nonempty closed
component has a marked hopping vertex. Its spatial Fourier symbol
with $d$ transverse marks is a sum of averages of translated
derivatives of the zero-field hop. These have envelope
$C_d\tau_i^{(d-b)/2}\exp[-c\tau_i|k+u|^2]$. Each temporal phase
only shifts the arguments at its two ends. At fixed hop times its
frequency factor is
$\int_0^{t_i}e^{i\nu s-a(r(s))|p|}\,\mathrm ds$.
The independent photon-frequency integration now applies on each
line before the hop times are integrated. The transverse section
of the preceding proof accounts for every $\sqrt{\tau_i}$ attached
to a transverse mark. For each Coulomb line, Plancherel bounds
both of its phase factors by the same additive cost in
$r^{-\gamma}$ as before. A closed spatial trace needs one anchor
Gaussian, so its anchor time has a residual cost
$C(1+r^{-(3+b)\alpha/2})$. Each other hop time costs at most
$C(1+r^{-b\alpha/2})$. A fixed number of closed and open chains
occurs at each positive electric order in this expansion in the
exact inverses. Products integrate over their distinct times,
with multiple contacts on a single hop already included in its
vertex. Together with the determinant, all costs are a polynomial
in integrals of subquadratic inverse powers times an exponential
of such an integral. The radial form estimate proves absolute
integrability of the whole product for arbitrary radial $L^2$
endpoints in each allowed channel. No expansion in the hopping
amplitude has been used in this graph bound.
\end{proof}

For example, take $\tau(r)=\kappa r^{2/3}$ and
$a(r)=\sqrt{\tau(r)}$ for the spinor. The represented heat
generator and optical form are fixed before the fields are
varied. The small-momentum velocity $v$ stays constant as $r$
moves. Both a neutral stationary radial preparation and its
open charge insertions satisfy the theorem. This is an explicit
way to remove the earlier neutral obstruction by changing the
curvature pairing. The theorem concerns the specified model's
ordered time representation. The unrestricted chiral reduction
has to identify its physical operator and the normal propagation
that its charged vertices actually sample.

\paragraph{Joint carrier-frequency estimates at the optical scale.}
At the matched optical length the neutral radial state can visit
the tip often enough to lose high inverse moments at a fixed time.
It still permits an absolute momentum representation. The argument
retains one carrier Fourier variable for each closed charged trace
and uses a maximal Fourier estimate for the Coulomb insertions.

\begin{theorem}[Joint bare bounds with all normal angular sectors]
\label{thm:seam-calibrated-joint-uv}
On the finite spatial torus and finite slab use the common Maxwell
and Gauss readout, the onsite-relative Fock vacuum and the
angular-sector-preserving hoppings of the optical model. Let $b=0$
for its scalar hopping and $b=1$ for the first spatial derivative
contracted with the Dirac matrices. Take positive exponents with
\[
 \frac{(3+2b)\alpha}{2}<2,\qquad 0<\gamma\le\frac{\alpha}{2}.
\]
For arbitrary finite hopping strength, every finite electric order
of the connected bare vacuum, photon and elementary charged functions
has a joint absolute bound in every independent charged carrier
frequency and spatial loop momentum, the photon four-momenta, the
remaining vertex times and the moving radial path. The statement
uses the diagonal Fourier representation of a closed trace on the
slab. Charged endpoints may be tested at the slab boundaries as in
the optical model. The absolute numerator bound includes the neutral angular channel
and is uniform in the conserved integer $j$. It persists for square
summable angular superpositions of $L^2$ radial endpoints.
Normalization uses the nonzero vacuum of the chosen preparation;
its inverse can depend on that preparation.
\end{theorem}
\begin{proof}
Condition first on the radial path. Remove the measured temporal
potential by its time primitive on that path, carrying its endpoint
phase on each charged source. In particular this step does not
differentiate a radial sample path. The exact dressed inverse in the
preceding arbitrary-hopping argument has a time kernel bounded
uniformly in spatial momentum by a polynomial in
\[
 B_T=\int_0^T\!
 \left[1+r_t^{-(3+2b)\alpha/2}+r_t^{-b\alpha}
                                  +r_t^{-\gamma}\right]\mathrm dt .
\]
Indeed the free time kernel from either end, multiplied by the
time-dependent hopping $V$, has $L^2$ operator-norm envelope
bounded by $(\int_0^T\|V(u)\|^2\,\mathrm du)^{1/2}$, and the
dressed inverse on time $L^2$ has norm at most $1/m$.
The pathwise matter determinant costs at most $\exp(CB_T)$.
All constants below may depend on the electric order, duration,
box and microscopic couplings.

At a closed loop with at least two hop vertices, split the time
at one hop, Fourier transform the two half-endpoint functions and
retain their common Fourier variable $\xi$. Parseval restores the
original sewing; Cauchy--Schwarz bounds the $L^1(\mathrm d\xi)$
norm by the two time $L^2$ norms. The time transforms are taken
before the absolute value. At a loop with just one marked hop the
onsite inverse has zero slot trace against that off-diagonal hop.
The resolvent identity supplies a second hop and an additional
propagator. Use another, unsplit hop for the spatial trace.
These operations are performed after all the scalar endpoint phases
on each marked hop have been included. Open lines have no charged
trace to split.

A transverse mark on a hop at time $u$ supplies a factor
$\sqrt{\tau(r_u)}e^{-a(r_u)|p|}$. For the derivative spinor the
vertex has in addition at most one factor
$C\tau(r_u)^{-1/2}$ per hop. These bounds follow by differentiating
the translated Gaussian symbol before taking the chord-parameter
average. Keep half the Gaussian on the unsplit spatial trace hop.
The remaining spatial momentum sum at that time costs at most
$C(1+r_u^{-3\alpha/2})$, uniformly in all shifts. With the spinor
factor the cost at this hop is $O(1+r_u^{-(3+b)\alpha/2})$.
Other unsplit hop times carry at most $1+r_u^{-b\alpha/2}$.
The square entering a split end's $L^2$ norm costs at most
$1+r_u^{-b\alpha}$. All these time integrals are bounded by
polynomials in $B_T$. The number of dressed vertices at fixed
electric order is finite.

It remains to bound the electromagnetic marks uniformly over the
split end times. Put $r_{\min}=\min_{0\le t\le T}r_t$. Since
$\sqrt\tau/a$ is bounded, each transverse end is bounded by
$C(1+|p|)^{-1}e^{-a(r_{\min})|p|/2}$. Integrating the transverse
propagator in frequency costs $\pi/|p|$. The spatial sum for a
complete line therefore costs at most $C(1+|\log r_{\min}|)$
by three-dimensional shell counting with the zero mode removed.
Only finitely many such logarithms occur.

For a scalar end let
$f_p(t)={\bf1}_{[0,T]}e^{-a(r_t)|p|}$, and define
\[
 M_p(\nu)=\sup_{I\subset[0,T]}
     \left|\int_I e^{i\nu u} f_p(u)\,\mathrm du\right|,
\]
with the supremum over intervals. The strong $L^2$ maximal Fourier
inequality on the real line gives
$\|M_p\|_2\le C\|f_p\|_2$. For example, apply the continuous
variation theorem with $L^2$ exponent two and variation exponent
three to the inverse Fourier transform of $f_p$; a single
interval difference is bounded by that variation
\cite{SeamMaximalFourier}. Each of the two ends of an instantaneous
Coulomb line is bounded by this same $M_p$ in absolute value,
including ends at a split time. Its frequency integral is thus
bounded by $C\int_0^T e^{-2a(r_u)|p|}\,\mathrm du$. Summing with
$|p|^{-2}$ gives $C\int_0^T(1+r_u^{-\gamma})\,\mathrm du$.
Before integrating the matter vertex times the Wick photons have
independent frequencies on the slab. Spatial constraints can be
dropped in these positive lattice majorants. This proves the
conditional estimate with a polynomial in $B_T$ and
$1+|\log r_{\min}|$ times $\exp(CB_T)$.

For completeness the minimum contributes integrably even for
unbounded $L^2$ radial endpoint functions $v,w$. At $j=0$ identify
the unnormalized radial law with Brownian propagation killed on
the boundary of the five-dimensional unit ball, restricted to
radial functions. On the first half of the slab, the subsequent
free heat evolution bounds the terminal function by $C_T\|w\|_2$.
The probability for free Brownian motion started at radius $r$ to
ever hit the ball of radius $\epsilon$ is at most
$\min(1,(\epsilon/r)^3)$, by the radial harmonic function $r^{-3}$.
Killing only decreases this probability. Cauchy--Schwarz against
the initial function in $r^4\,\mathrm dr$ therefore bounds the
mass of such paths by $C_T\|v\|_2\|w\|_2\epsilon^{5/2}$:
the integral of the squared hitting bound is $O(\epsilon^5)$.
Time reversal treats the other half. Every fixed power of
$|\log r_{\min}|$ is thus integrable. The exponential moments
of $B_T$ for the same endpoints follow from the radial Hardy
estimate because all its exponents are strictly below two. Apply
Cauchy--Schwarz with respect to the positive endpoint path measure
to combine these estimates. The resulting bound is proportional
to $\|v\|_2\|w\|_2$. At any $j$ the positive angular potential
can be dropped in the absolute radial majorant. Summing the endpoint
norm products uses Cauchy--Schwarz in $j$ because the assumed
hoppings preserve that angular label. The positive vacuum
denominator and finite products of differentiated coefficients give
the connected bounds.
\end{proof}

For the scalar step one can take $\alpha=1$ with
$a(r)=\sqrt{\kappa r}$. For the derivative spinor the choice
$\alpha=2/3$ and $a(r)=\sqrt{\kappa r^{2/3}}$ satisfies the
stricter determinant bound. The latter hopping is realized by a
heat generator of variable speed in the normal coordinate. Its
weight and the corresponding optical Maxwell coefficients can be
singular at the collapsed tip. Their domains are those of the
weighted forms defined here. The proof uses precisely this curvature pairing and heat transport. The theorem covers every angular sector of these
represented actions. It leaves the identification of the physical
chiral action, including its measure and any angular transitions,
as an independent question.

\paragraph{The optical bound when the charged vertices change angular sector.}
Keep the full periodic normal circle. Let a pair-creation hop
carry $e^{i\ell\theta}$ and its adjoint carry $e^{-i\ell\theta}$,
as in the switched two-hop return. Use the optical heat lapse and
electric readout of Theorem~\ref{thm:seam-calibrated-joint-uv} at
every visited depth. The phase acts in the geometric angle and
commutes with the measured electric gauge parameter.

\begin{corollary}[All electric orders with angular transitions through zero]
\label{cor:seam-optical-shifts-uv}
Under the scalar or spinor exponent hypotheses of
Theorem~\ref{thm:seam-calibrated-joint-uv}, its full joint graph
bound holds at arbitrary finite hopping strength on the finite
slab and torus for this angular-shifting action. The bound
includes the neutral angular sector, simultaneous transitions of
overlapping charged loops, vacuum and photon coefficients, and
the same elementary charged endpoint tests. Unnormalized radial
and angular matrix elements are bounded by
$C_{T,\Gamma}\|v\|_2\|w\|_2$ for arbitrary endpoint vectors.
A nonzero positive trace-class preparation has a strictly positive
finite vacuum normalization. Each numerator and the one vacuum
denominator use the same geometric transfer.
\end{corollary}
\begin{proof}
In the weighted coordinate $r^4\mathrm dr\,\mathrm d\theta/(2\pi)$
the quartic normal heat generator in the interior is
\[
 \partial_r^2+4r^{-1}\partial_r+r^{-8}\partial_\theta^2 .
\]
Start with its killed angular-zero radial diffusion $X$ and its
independent circle Brownian motion of generator $\partial_\theta^2$.
Conditional on a surviving radial path use the angular clock
$A_t=\int_0^t X_u^{-8}\,\mathrm du$. It is finite because a
surviving path on the fixed slab avoids zero and is continuous.
The outer boundary is absorbing. This construction gives the
full normal transfer: on mode $j$ a time interval contributes
$\exp[-j^2(A_t-A_s)]$, exactly its normal angular heat factor.
One can first work on a compact radial subinterval and pass by
Friedrichs exhaustion. The angular starting measure is normalized
Haar measure. These paths describe one geometric degree of freedom
with its actual endpoints; there is no second sum of angular
determinants.

Fix both paths, writing $\Theta_t$ for the time-changed angle.
The pair hop now has the numerical multiplier
$e^{i\ell\Theta_t}$ and its adjoint the conjugate. The full
one-particle hop is Hermitian and off diagonal for the onsite
grading. It still preserves spatial momentum at zero electric
field. Its spatial derivatives, Wilson contacts, Hilbert--Schmidt
norm and operator norm satisfy the same absolute bounds uniformly
in $\Theta$. The scalar time primitive and its charged endpoint
phases commute with this multiplier. In the diagonal carrier
Fourier trace, a cut hop simply retains the angular factor at
its own endpoint time. Conditional on the paths the two adjacent
dressed propagators still depend on the respective split time.
It follows that the resolvent, Coulomb maximal Fourier, and
spatial envelope bounds in the preceding theorem apply with
exactly the same radial costs. The conditional determinant can
be bounded in modulus by $\exp(C B_T)$ from its Hilbert--Schmidt
estimate. Positivity of this determinant for each angular
trajectory is unnecessary.

In the full Hamiltonian the angular kinetic energy is positive
and the hopping form has the same subquadratic radial singularity
as in the optical construction. The neutral radial Hardy bound
therefore defines the perturbed semibounded form with the angle
left dynamical. To identify the conditional functional, insert
both normal coordinates in the vacuum transfer on a compact radial
subinterval. Expansion between diagonal transfers gives the angular
phase at each hop and the ordinary filled-vacuum determinant at
fixed paths. The bounded spatial Hilbert--Schmidt kernels on the
subinterval justify this expansion on the finite slab, as in
Proposition~\ref{prop:physical-moving-ordered-uv}. The radial
majorants just obtained allow removal of the exhausting boundaries
in the vacuum and marked integrals.

Finally let $V(r)=\|v(r,\cdot)\|_{L^2(\mathrm d\theta/(2\pi))}$
and define $W$ from $w$ in the same way. After using a graph
majorant independent of $\Theta$, the conditional angular endpoint
integral costs at most $V(X_0)W(X_T)$. Indeed it is the pairing
of the absolute endpoint functions with the circle heat semigroup
at time $A_T$, an $L^2$ contraction for every clock value.
The radial Hardy, minimum-depth and exponential estimates thus
apply with $V,W$ as endpoints. They give the claimed bilinear
bound since $\|V\|_2=\|v\|_2$ and $\|W\|_2=\|w\|_2$ in their
respective measures. Pairing with a positive trace-class preparation
preserves finiteness. Strict positivity of its real-coupling
diagonal vacuum follows from injectivity and positivity of
$e^{-TH}$ for the full semibounded Hamiltonian. Normalization and
connected coefficients now use finite products of these derivatives.
\end{proof}

The same proof permits a finite internal multiplicity. In the upper
off-diagonal hopping block take a finite sum of the scalar or
first-derivative heat kernels above multiplied by bounded matrices
$U_a(r,\theta)$, and take its adjoint in the lower block.
All matrix elements connect slots with the same electric gauge
transformation in the one-particle description. Use the spinor
exponent bound if a derivative kernel occurs. These matrices act by
multiplication at the hop, with no derivative on the normal fields.
For a fixed normal trajectory they are bounded matrices of time
alone. The Schur and Hilbert--Schmidt estimates cost a finite sum of
their norm bounds, the spatial Fourier estimates hold componentwise,
and the off-diagonal self-adjoint hopping retains the inverse bound
$1/m$. Bounded measurable multipliers can be used by defining the transfer
through its time integrals. The diagonal normalization again follows from
the full Hamiltonian. This includes finite superpositions of angular
charge shifts without requiring a gap for the angles they reach.

In this representation the charge can propagate from zero angular
weight to a charged sector and back, even in a four-hop exchange.
The optical estimate controls the complete graph through that
history. Its use for an unrestricted physical action requires the
same heat lapse, electric readout, electric source and normal
transfer to be identified on its physical domain.

\subsection{A constitutive electromagnetic action and its physical tests}
\label{sec:omega-completion-uv}
How do the charged operator and the physical quotient follow from
the amended action? Do its bare graphs converge? Does its current
have the relativistic spectrum required on the sector plane?
The following results address each question using the same source.

\paragraph{Fixing the two sampling lengths in the raw electromagnetic action.}
One can prescribe the relative calibration before defining a Feynman
integral. The following constitutive clauses give an example on each
chosen locked plane. Let $s_\perp>0$ be proper distance from the
tip of the quartic collar, and let $L_*$ be a length already
assigned on the branch. Use normalized spatial coordinates in an
orthonormal carrier frame. Require that equal increments of
$s_\perp$ sweep equal three-dimensional sampling volumes, with
the ball of radius $L_*$ swept at $s_\perp=L_*$. Use the
Euclidean volume of the model ball in the tangent space even on
the torus. Its radius $a$ is thereby fixed by
\begin{equation}
 a(s_\perp)^3=L_*^2s_\perp .
 \label{eq:constitutive-volume-uv}
\end{equation}
The second clause transports displacements by normalized independent
isotropic Gaussian increments with accumulated variance $2a^2$ per
orthonormal coordinate. Equivalently its kernel obeys
$\partial_a K=2a\Delta_x K$, $K(0)=\delta_0$.
In Fourier space this gives $\widehat K_a(k)=e^{-a^2|k|^2}$.
Thus $r=s_\perp/L_*$, $a/L_*=r^{1/3}$ and
$\tau/L_*^2=r^{2/3}$. These exponents follow uniquely from the
two clauses and their normalization. The three-volume clause is
new geometric input. Spatial dimension alone leaves the relation
between sampling volume and proper depth undetermined. Taking
$L_*$ from the common microscopic assignment leaves its existing
normalization in place; the construction does not derive the
matter masses or the measured electric charge from $G_N$.

\paragraph{A constrained normal connection and its boundary term.}
On a carrier slice use optical distance $y=a(s_\perp)$ and extend
the boundary connection by the local equation
$(\partial_y^2+\Delta_x)B=0$ in $y>0$, $B(0)=A$, with the
finite-energy continuation and the chart's harmonic-mode convention.
Its unique extension on nonzero modes is
$B_y=\mathcal E_y A=e^{-y|\nabla|}A$. Carry the same extension
on the scalar potential and on the gauge parameter. In the physical
transverse chart the radial connection for a pure boundary gauge
can be reconstructed from its longitudinal primitive. In fact, for
$A_i=A_i^{\rm T}+\partial_i\chi$ set
$B_\perp=\partial_{s_\perp}(\mathcal E_y\chi)$. The pullback
of a pure gauge to a smooth charged path then gives its exact
endpoint difference. On a Brownian radial lift this statement
uses the Stratonovich chain rule, or the endpoint definition
obtained from smooth approximations on compact interior collars.
In the transverse chart $\chi=0$, and the radial gauge term vanishes.
The electromagnetic zero mode is treated by the same convention
as the charged transfer.

Take the conditional optical measure at fixed $A$ to be the
point mass at $\mathcal E A$. This constraint has a local
multiplier formulation. At finite dimension choose a lift $A^0$
of the boundary value and write $B=A^0+b$ with Dirichlet $b$.
The interior equation is $D b=f(A)$ with positive invertible
Dirichlet matrix $D$. The normalized measure
$|\det D|\,\delta(D b-f(A))\,\mathrm db$ integrates any
source functional by evaluation at $b=D^{-1}f(A)$. In the
optical coordinate $D$ is fixed independently of matter,
sources and the trajectory at which the extension is sampled.
The continuum measure is defined directly by the graph of
$\mathcal E$. Its evaluations on successive radial slices sew by
$\mathcal E_{y_2}\mathcal E_{y_1}=\mathcal E_{y_1+y_2}$.
This derivation keeps the full ordinary charged determinant.
Assigning an independent propagating normal photon would change
the conditional measure.

In the basic constitutive choice the Dirichlet principle defines
the map $\mathcal E$ and the constrained measure pushes forward
the ordinary boundary Maxwell law. Its energy is not a separate
term in the bare weight. An alternative bare action also weights
the constrained connection by a positive normal energy. Its value
on the harmonic graph then contributes to the boundary action. For one spatial mode $p\ne0$,
\[
 \int_0^\infty
 (|\partial_y B_p|^2+|p|^2|B_p|^2)\,\mathrm dy
 =|p|\,|A_p|^2 .
\]
It cannot be absorbed into source-independent normalization. One
gauge-invariant choice retaining the low-frequency Maxwell speed
is to use this normal form on the magnetic curvature
$F_{ij}(B)=\partial_iB_j-\partial_jB_i$. Sum the Dirichlet energies of the independent components
$i<j$ with coefficient $L_*/2$. For a transverse
mode the resulting action adds $L_*|p|^3|A^{\rm T}_p|^2/2$
per unit carrier time. The transverse Euclidean covariance is
therefore
\begin{equation}
 \frac{P^{\rm T}(p)}{\nu^2+|p|^2+L_*|p|^3}.
 \label{eq:constitutive-boundary-symbol-uv}
\end{equation}
The Coulomb charge square is unchanged because the added
magnetic energy is independent of the scalar potential. The
source for a measured connection $B_y$ still carries its
individual factor $e^{-y|p|}$.

\begin{proposition}[Two optical curvature sources with explicit boundary weights]
\label{prop:constitutive-source-uv}
Use the volume and displacement laws above, the constrained
normal measure, the quartic radial transfer and the onsite-relative
charged boundary state of
Corollary~\ref{cor:seam-optical-shifts-uv}.
Use the four-arrow curvature pairing of the charged construction
with $U$ the measured spatial Wilson link and
$V=-iv\sum_i\alpha_i\partial_i K_a$, together with the
onsite mass and time jet. Use either the ordinary boundary Maxwell weight or that same
weight with the positive magnetic normal energy giving
\eqref{eq:constitutive-boundary-symbol-uv}.
On the finite slab and spatial torus its complete bare charged
and electromagnetic graph coefficients satisfy the joint
all-electric-orders bound of that corollary, at arbitrary finite
hopping strength.
\end{proposition}
\begin{proof}
The arrow contraction is the real part of $\bar\zeta VU\xi$.
Reversing the link gives its adjoint. All copies of a charged jet
evaluate the same Berezin field and hence give its single Fock
determinant. Differentiating the Wilson link gives all coincident
photon insertions, and integrating the scalar potential gives the
same two-end Gauss square. The constrained normal integration
above gives no extra curvature weight in the basic choice and
gives the displayed magnetic curvature energy in the alternative.
The corresponding boundary oscillators have positive frequencies
$|p|$ and $(|p|^2+L_*|p|^3)^{1/2}$, respectively.
Their frequency integrals are bounded by $\pi/|p|$; all
transverse photon majorants in the optical proof therefore still
apply. At the free infrared pole the squared frequency divided
by $|p|^2$ tends to one. The spinor exponent in the charged
Hilbert--Schmidt estimate is $5(2/3)/2=5/3<2$ and the
charge-width ratio $\sqrt\tau/a$ equals one. Applying the
corollary with the actual conditional measure proves the result,
including its marked determinant, endpoint phases and angular
transitions. This construction leaves the chiral and
Lorentzian identification to the physical source comparison.
\end{proof}

\paragraph{Why independently fluctuating normal photons change the answer.}
This comparison can be tested in a transverse, gauge-invariant
observable before any charged expansion. Give an independently
integrated optical vector component the quadratic action
$\tfrac12\int\mathrm dt\,\mathrm dy\,\mathrm d^3x
(|\partial_yB|^2+|\nabla_x B|^2)$ at fixed boundary data.
For wavenumber $k=|p|>0$ its fluctuation about the harmonic
extension has Dirichlet Green function
\[
 G_D(k;y,y')=\frac{e^{-k|y-y'|}-e^{-k(y+y')}}{2k}.
\]
This follows by the derivative jump at $y=y'$ and the zero
boundary trace. The covariance is white in carrier time for
this particular normal action. Smear it against a real smooth
time test of unit $L^2$ norm and use a rectangular spatial
Wilson contour at any fixed $y>0$ in a cubic torus of period
$L$. Take one side of length $h\in(0,L)$ along axis one and
the adjacent side of length $L/2$ along axis two. For momenta
$p=(2\pi/L)(0,n_2,n_3)$ with $n_2$ odd, the two
axis-two contributions cancel and the axis-one current has
modulus $2h$. This conserved current is exactly transverse.
In the shell $K<|p|<2K$ there are at least $c_L K^2$ such
modes for large $K$, each with
$G_D(k;y,y)\ge c/K$. The variance of this time-smeared
circulation is thus infinite. This calculation distinguishes
the independent Gaussian measure of this example from the
deterministic extension in the proposition. It tests the full
closed current; a fluctuation determinant independent of the
external probe cannot remove the variance. Other dynamical
bulk actions or admissible probe classes require their own
normal covariance calculation.

\paragraph{Solving the electric constraint before taking the source derivatives.}
The charged readout transports the gauge parameter as well as the
connection. This matters when a charge moves in the normal coordinates
while its spatial position stays fixed. We give the constrained
realization on the spatial torus of~\eqref{eq:cl-gauss-charge-uv}.
Work on its trivial Abelian bundle with periodic, zero-mean gauge
parameters. The harmonic photon is held fixed. In a charged sector
the prescribed uniform compensating density supplies the missing
zero-mode charge. Write $y=(r,\theta)$ for the normal position and
$P_y=\exp[-a(y)|\nabla|]$, with $a>0$ smooth in the interior.
Hodge decomposition gives $A=A^\perp+\nabla\chi$, uniquely with
$\chi$ of zero mean. For a configuration of particles and holes set
$\rho(u)=\sum_j s_j u(x_j)$, with $s_j=+1,-1$ as in the Gauss
square. Let $\rho^\circ$ be the corresponding density with its
mean removed. Multiplication on charged wave functions by
\[
 U_{y,\chi}=\exp\!\big[ie\rho(P_y\chi)\big]
\]
implements the local charge transformation on every particle sector.

\begin{proposition}[The Abelian physical quotient at a moving readout]
\label{prop:constitutive-gauss-uv}
Use the Maxwell electric norm, the onsite charged representation,
and the normal Friedrichs form with its electric derivatives
\[
 D_i=\partial_{y_i}-ie\rho((\partial_{y_i}P_y)\chi),
 \qquad y_i=r,\theta .
\]
Use the complete hopping and its adjoint with the spatial Wilson
link of $P_y A$. Solving Gauss's law and taking the gauge quotient
gives exactly the transverse optical Hamiltonian, including its
full Coulomb energy, on the stated torus. Gauge-invariant endpoint
operators obtained by dressing the elementary charged probes have
the same matrix elements as their transverse representatives.
\end{proposition}
\begin{proof}
At fixed time the momentum conjugate to $\chi$ is
$\Pi_\chi=-\nabla\cdot\Pi_A$. In the Schr\"odinger representation
the Gauss constraint smeared against a smooth zero-mean $u$ reads
\[
 \big[-i\partial_{\chi,u}-e\rho(P_yu)\big]\Psi=0 .
\]
Its solutions on a gauge orbit are precisely
$\Psi(A^\perp,\chi,y)=U_{y,\chi}\psi(A^\perp,y)$.
Indeed the derivative of $U_{y,\chi}^{-1}\Psi$ in every gauge
direction vanishes. The physical inner product integrates $\psi$
on the transverse slice with the original normal and matter
measures. Restriction to $\chi=0$ is unitary for this inner
product, with inverse the displayed equivariant extension. This
defines the longitudinal quotient directly, without a Lebesgue
measure on an infinite-dimensional space of gauge functions.

This constraint also has an explicit minimal Abelian BFV complex.
Here BFV denotes the Hamiltonian ghost complex on the boundary.
In an orbit coordinate $\chi$ assign one exterior ghost to each
gauge direction. Conjugation by $U_{y,\chi}$ identifies the
constraint differential, up to the conventional factor $-i$,
with the exterior derivative on the vector space of zero-mean
gauge functions. On smooth cylindrical forms the ordinary
contraction of the orbit to $\chi=0$ is a chain homotopy. In
coordinates $q$ for the finitely many arguments of a $k$-form,
it is integration of $t^{k-1}\iota_q\omega(tq)$ over $0<t<1$.
Thus its degree-zero cohomology consists of precisely one copy of
the transverse data, with the Hilbert completion just described.
The homotopy is compatible with adjoining further cylindrical
arguments. Decoupled gauge-fixing doublets have their usual
two-term contraction. This chooses the boundary ghost domain
and physical norm explicitly on this chart.

The orthogonal splitting of the electric norm gives the longitudinal
energy $\frac12\langle\Pi_\chi,(-\Delta)^{-1}\Pi_\chi\rangle$.
The generators of all the smeared charge densities commute on each
position sector. Inserting the constraint therefore gives the
positive multiplication form
\[
 \frac{e^2}{2V}\sum_{p\ne0}
       \frac{e^{-2a(y)|p|}}{|p|^2}\rho_p^*\rho_p .
\]
It includes the diagonal and mixed contributions for every occupied
configuration. For $r>0$ the charge smearing makes each finite-particle
matrix element well defined. The full positive form has its natural
domain. The Wilson hop acquires exactly the phase of the created
particle and hole when $A^\perp$ is replaced by
$A^\perp+\nabla\chi$. Also
$D_i(U_{y,\chi}\psi)=U_{y,\chi}\partial_{y_i}\psi$. These are
identities of the hopping and normal kinetic forms on smooth
interior cores and hence of their closed realizations. In particular
the contact quadratic in $\partial_iP_y\chi$ is already in the
covariant normal kinetic energy. The longitudinal Hodge Jacobian
and the Abelian gauge-slice Jacobian are independent of $y$, the
charged fields and the probes. The unitary quotient establishes
this reduction without assigning a numerical value to their formal
infinite determinants.

For a changing external gauge one has $b_0\mapsto b_0-\partial_t u$
and $\chi\mapsto\chi+u$. The temporal charge is $e\rho(P_y b_0)$.
At fixed $y$, differentiating the phase in time gives the
corresponding transformed Hamiltonian. Along a smooth normal
trajectory there is in addition a longitudinal current
\[
 J_{\rm n}(t)=e\nabla(-\Delta)^{-1}
             \big[(\mathrm d P_y/\mathrm dt)\rho^\circ\big]
\]
on a segment with constant occupation and spatial positions.
Its divergence cancels $e(\mathrm d P_y/\mathrm dt)\rho^\circ$ in
the time derivative of the measured charge. The normal covariant
derivatives implement this identity at the operator level; their
definition requires no velocity of a diffusion path. At a jump
the associated change of charge is carried by the Wilson link.
For an endpoint operator $O$ on the transverse slice, the
equivariant representative acts as $U_{y,\chi} O U_{y,\chi}^{-1}$
with the appropriate incoming and outgoing charge sectors. Its
matrix elements, including charged boundary phases, are preserved
by the quotient. A single charged insertion changes the global
charge sector and uses the stated background convention. In a
strictly neutral torus only globally neutral combinations are
admitted.
\end{proof}

\paragraph{Angular dependence of the geometric widths.}
The quotient allows the heat lapse $\tau$ and optical distance
$a$ themselves to vary with the normal angle. On the same quartic
chart suppose both are bounded and strictly positive in the
interior, and for some $c_1,c_2>0$ and $\alpha_0>0$ satisfy
\[
 \tau(r,\theta)\ge c_1r^{\alpha_0},\qquad
 a(r,\theta)\ge c_2\sqrt{\tau(r,\theta)},\qquad
 (3+2b)\alpha_0/2<2 .
\]
Here $b=0$ for a scalar heat kernel and $b=1$ when a spatial
derivative occurs in the hop. For these widths the graph conclusion
of Corollary~\ref{cor:seam-optical-shifts-uv} holds throughout
the Abelian quotient of Proposition~\ref{prop:constitutive-gauss-uv},
with the same permitted boundary data and finite slab. Indeed,
conditioning on the full normal path gives a squared spatial
Hilbert--Schmidt cost $C(1+r^{-(3+2b)\alpha_0/2})$ per unit time.
The transverse mark has $\sqrt\tau/a\le1/c_2$, and
$\min_t a(y_t)\ge c_2\sqrt{c_1}\,r_{\min}^{\alpha_0/2}$.
Thus the transverse spatial sum has at most the same logarithmic
cost in $r_{\min}$. The maximal Fourier estimate for a Coulomb
mark applies to the actual function $t\mapsto e^{-a(y_t)|p|}$
on the slab. After summing the line it costs at most
$C\int_0^T(1+r_t^{-\alpha_0/2})\,\mathrm dt$. The radial Hardy
and hitting bounds integrate these quantities uniformly in the
angular history. Conditioning and then contracting the angular
endpoint heat kernel proves the claimed bound. All normal
connection terms have already been included in the quotient.
A dynamically independent bulk connection, or a change in the
boundary differential beyond this electric constraint, would supply
additional physical degrees of freedom and its own source terms.

\paragraph{The physical chiral rows and the filled charged reference.}
The electromagnetic sector can retain the physical Weyl constraint and
still give the complete optical spinor determinant. This requires a
specific choice of mass on the retained module. Let $S_L,S_R$ be the
negative and positive eigenspaces of $\gamma_5$, and let $V_L,V_R$
have internal grading $+1,-1$, respectively. On a fixed locked sector
plane the physical kets for $T_\chi=\gamma_5\otimes\Gamma_F$ with
eigenvalue $-1$ form
\[
 E=(S_L\otimes V_L)\oplus(S_R\otimes V_R).
\]
Take finite internal multiplicities. Fix a charge-equivariant unitary
$J:V_L\longrightarrow V_R$, so $Q_RJ=JQ_L$ for the Hermitian
electric charge generators. Write $J_o$ for the odd self-adjoint
involution on $V_L\oplus V_R$ taking $v_L$ to $Jv_L$ and $v_R$ to
$J^*v_R$.

\begin{proposition}[Physical chiral identification of the optical charged block]
\label{prop:omega-chiral-uv}
Use the onsite operator $m\beta\otimes J_o$, where $m>0$ and
$\beta=\gamma^0$ in the Hamiltonian frame. For each electric charge
take the optical spatial Dirac hopping of~\eqref{eq:reference-spin-uv}
with the chosen heat lapse, acting identically on paired internal
multiplicities. Use the common optical readout on the time connection
and spatial Wilson links. Restrict the kets and their canonical
adjoints to $E$, and fill the negative onsite eigenspace on $E$.
In this quantization the physical electromagnetic action, boundary
reference and every charged or photon source derivative are unitarily
equivalent to those of the optical spinor on $S\otimes V_L$, with the
corresponding electric charges. The graph bound of
Theorem~\ref{thm:seam-calibrated-joint-uv} and its angular extension
in Corollary~\ref{cor:seam-optical-shifts-uv} therefore apply to the
full retained chiral block on their stated domains.
\end{proposition}
\begin{proof}
Let $p_L,p_R$ be the spin chirality projections. The isometry
\[
 I_J u=(p_L\otimes 1)u+(p_R\otimes J)u
\]
from $S\otimes V_L$ has range exactly $E$. The Dirac algebra gives
$[\alpha_i,\gamma_5]=0$ and $\{\beta,\gamma_5\}=0$. The odd
internal map anticommutes with $\Gamma_F$. Consequently both
$\alpha_i\otimes 1$ and $\beta\otimes J_o$ preserve $E$.
Conjugation by $I_J$ on this range gives, respectively,
$\alpha_i\otimes 1$ and $\beta\otimes 1$ on $S\otimes V_L$.
Conjugating the charge generator by this isometry gives
$1\otimes Q_L$, so it intertwines charged endpoint phases, Wilson
links, the temporal readout and the full Gauss square. The charge
here is the one-particle generator; the holes in its filled
reference have opposite charge. The two mass slots and their filled
polarizations are intertwined before the Fock space is formed.
Angular factors multiplying forward and reverse hops act on these
physical mass slots, with conjugate factors on the two orientations.
Equivalently one may transport the allowed off-diagonal optical
kernels and their adjoints by $I_J$.

The four-arrow curvature product yielding the bilinear uses these
constant matrices on its onsite and spatial arrows. Its canonical
charged time derivative on $E$ is intertwined as well. In a covariant
row convention $\bar\Psi=\Psi^\dagger(\beta\otimes1)$ has the
opposite $T_\chi$ grading. Multiplication of the canonical operator
by $\beta\otimes 1$ accordingly maps physical kets into the correct
row space; its mass has internal map $mJ_o$. Compressing a covariant
Dirac operator between two identical ket projections would discard
the kinetic term. Equivalently one may carry out the entire
calculation with physical kets and their canonical adjoints.

The isometry is independent of the optical fields and moving depth on
the specified plane. On Fock space its particle and conjugate hole
maps carry one vacuum polarization into the other. They intertwine
the charged creation operators, complete transfer and each endpoint
insertion. The normal boundary measure is unchanged by this unitary
identification. Thus the normalized generating functionals agree with
transported probes before any source differentiation. The relative
determinant operators are unitarily equivalent slot by slot. There
is one charged determinant per retained Dirac multiplicity, with all
source contacts.
\end{proof}

\paragraph{A test of the mass assignment on support arrows.}
For a massive electromagnetic block, gauge invariance of a
nondegenerate Dirac scalar mass with no pairing of a field to its
charge conjugate requires $\dim V_{L,q}=\dim V_{R,q}$ for every
charge $q$. Indeed its chiral pairing has an internal map
$M:V_L\to V_R$ with $Q_R M=M Q_L$. Each charge block must be
bijective. Conversely, equality of the dimensions gives a unitary
$J$ charge by charge. The polar part of any invertible such $M$
is one such $J$. An unpaired Weyl mode has no gap from this mass
construction. Its charged determinant requires a separate analysis.

For a groupoid-equivariant choice there is a further, exact test.
Write $U_L(a),U_R(a)$ for the charge-preserving unitary transports
over an arrow $a:x\to y$. A coherent mass pairing obeys
\[
 J_y U_L(a)=U_R(a)J_x .
\]
In each connected component choose a base object $x_0$. A pairing
exists algebraically exactly when its two internal fibers at $x_0$
admit a charge-equivariant unitary $J_{x_0}$ intertwining the
transports of every isotropy arrow $h:x_0\to x_0$. Necessity follows
by restriction. For sufficiency choose a path $a:x_0\to x$ and set
$J_x=U_R(a)J_{x_0}U_L(a)^*$. Two such paths differ by an isotropy
arrow, so the formula is independent of the choice and intertwines
all transports by composition. Integration over a family of support
objects also requires a measurable pairing on its components. Even
on one object, distinct circle characters on two one-dimensional
fibers have equal electric multiplicity but fail the intertwining
test. The completed electromagnetic assignment may choose matched
fibers and pairings at its outset. Additional charged multiplets,
independent fluctuating connections and any measure over supports
must enter the source estimate on a branch that includes them.

\begin{theorem}[Bare finiteness on the constitutively selected electromagnetic quotient]
\label{thm:omega-quotient-uv}
For the amended $\Omega_\alpha$ electromagnetic action on a finite
slab and spatial torus with its prescribed harmonic mode, take a
massive paired multiplet and a nonzero positive trace-class normal
preparation, with the onsite matter and photon vacuum at the boundary.
At every finite electric order and arbitrary finite hopping strength,
the complete connected electromagnetic and elementary dressed charged
coefficients on the Abelian physical quotient are ultraviolet finite
without counterterms or momentum cutoffs. This includes every normal
occupation reached by the specified hopping, loops with zero net
winding, exchange of partners, instantaneous Coulomb sewings and
simultaneous short-duration limits. The joint absolute estimate uses
one carrier frequency and spatial loop momentum for each closed charged
trace and the four-momentum of each internal photon on the finite
slab, as in Theorem~\ref{thm:seam-calibrated-joint-uv}.
\end{theorem}
\begin{proof}
Equation~\eqref{eq:constitutive-volume-uv} and the independent
geometric increments give the spinor exponents $2/3$ and $1/3$.
The harmonic boundary constraint and the four-arrow curvature
product give the measured fields and their complete hopping
operator (Proposition~\ref{prop:constitutive-source-uv}).
Proposition~\ref{prop:constitutive-gauss-uv} solves the moving
Gauss constraint, including both normal components of a pure gauge.
The physical ket, its adjoint and the filled reference are identified
in Proposition~\ref{prop:omega-chiral-uv}. These equalities
of sourced matrix elements transport all derivatives, including
the closed determinant and contact terms. Corollary~\ref{cor:seam-optical-shifts-uv}
then bounds the coefficients on the full normal path space.
Positivity of the diagonal transfer gives the shared nonzero vacuum
denominator. Each coefficient of its normalized logarithm is a
finite sum of products of the bounded derivatives.
\end{proof}

Here the physical Hilbert space is obtained by solving the specified
Abelian constraints. Whether its current implements the spectrum
and locality required by Q1 is a separate assertion about this
same sourced theory. The tests below evaluate that assertion.

\paragraph{The dynamical flat Maxwell connection.}
The preceding torus proof can include the harmonic spatial photon.
Take a cubic torus of side $L$ and volume $V$ in the sector of zero
magnetic flux. Denote the spatially constant vector potential by
$A_{\rm h}$. The dimensionless holonomy $b=eL A_{\rm h}$ has each
component defined modulo $2\pi$ for charge unit $e$. Trivializing
on the cover of the holonomy torus, the identification
$b\mapsto b+2\pi n$ acts on a charged configuration by
$U_n=\exp[2\pi i\sum_j s_j n\cdot x_j/L]$. Wilson hops carry
the additional factor $\exp(ib\cdot z/L)$ on each oriented
lifted chord $z$. Under the identification it transforms by the
same endpoint phases. The onsite vacuum is invariant under every
$U_n$.

\begin{corollary}[All spatial Maxwell modes in a neutral vacuum response]
\label{cor:constitutive-holonomy-uv}
In the topologically trivial magnetic bundle take the constant
normalized holonomy wave function at the two ends with the onsite
vacuum. Keep the positive diagonal normal preparation of the optical
theorem. Quantize the flat connection with its Maxwell kinetic
energy and impose zero total electric charge. All finite electric
orders of the bare vacuum logarithm and the connected nonzero-mode
transverse photon responses obey the same joint ultraviolet bounds,
with every dynamical spatial photon mode included.
\end{corollary}
\begin{proof}
In the coordinate $b$ the harmonic Maxwell generator is
$-e^2L^2\Delta_b/(2V)$ with the equivariant identifications above.
Its coefficient follows from the real-time kinetic energy
$V|\dot A_{\rm h}|^2/2$. Thus expansion in $e$ here holds the
dimensionless holonomy and its Haar preparation fixed. At $e=0$
the generator vanishes and the constant twist $b$ still enters
each spatial hop. A smooth simultaneous representative at real
$e$ on the universal cover is
\[
 b_t=b_0+\frac{eL}{\sqrt V}W_t ,
\]
where $b_0$ is uniform in one fundamental cell and $W$ is
three-dimensional standard Brownian motion starting at zero,
independent of the nonzero Maxwell and normal variables. Conditional
on $b_0$, the lift is used once. In a representation folded back
into the cell the boundary crossings supply the same $U_n$ by
equivariance. Integrating this lift includes all windings. The
final vacuum wave function is constant, so the endpoint of the
Brownian path is unrestricted. The uniform $b_0$ average retains
the holonomy projection appropriate to these vacuum boundaries.

At $e=0$ each spatial symbol is translated by $b_0/L$ up to the
orientation sign. Gaussian spatial lattice bounds and the
off-diagonal inverse estimate hold uniformly under this shift.
At a fixed electric order a Brownian mark from a chord at $t$
is $i z\cdot W_t/\sqrt V$ in addition to the static twist.
Its spatial displacement has the same $\sqrt\tau$ bound as a
zero-spatial-momentum transverse mark. A pairing of Cartesian
components $i,j$ at $s,t$ has the exact frequency representation
\[
 \frac{\delta_{ij}}{V}\min(s,t)=
 \frac{\delta_{ij}}{2\pi V}\int_{\mathbb R}
   \left(\int_0^s e^{i\nu u}\,\mathrm du\right)
   \left(\int_0^t e^{-i\nu v}\,\mathrm dv\right)
   \mathrm d\nu .
\]
Each interval transform is at most $\min(T,2/|\nu|)$, including
at a split charged endpoint time. Thus the frequency integral is
absolutely convergent uniformly on the slab. There are finitely
many such lines per electric order. The other hard momenta and
the shared normal paths have their previous majorants.

The modulus of the harmonic Wilson phase is one before expansion.
The Maxwell rotor adds a nonnegative kinetic form to the original
semibounded physical Hamiltonian on equivariant sections. On a
compact interior normal collar, the bounded heat-hop expansion
with these phases is precisely its transfer in the lifted Brownian
representation. Differentiating the phases to fixed order there
adds finite Brownian moments and derivatives of spatial heat
symbols. The radial majorants allow Friedrichs exhaustion also
for the resulting marked integrals. The diagonal vacuum transfer
is strictly positive by spectral resolution. The constant scalar
potential, in contrast to the spatial rotor, is a Lagrange
multiplier: it enforces zero total charge and supplies no free
Coulomb oscillator. This proof describes the vacuum sector of the
specified magnetic bundle.
\end{proof}

\paragraph{Magnetic flux and charged momenta without translation invariance.}
On a flat cubic torus a magnetic flux sector is indexed by
$n\in\mathbb Z^3$. Choose a unitary reference connection
$\mathcal A_n$ with constant curvature of periods $2\pi n$ on
its charge-one line bundle. Extend this reference constantly along the
normal collar and apply the harmonic readout to differences of
connections, which are globally defined one-forms. Changing the
reference by a flat torus mode translates the integrated holonomy.
Write the full dimensionless
connection as this reference plus the flat holonomy and the
globally defined Maxwell fluctuations with charge factor $e$.
The reference parallel transport and the bundle transition maps
enter every lifted Wilson chord. Spatial displacement still follows
the heat law of the charged action. The physical Maxwell energy
of the background is $E_n=2\pi^2|n|^2/(e^2 L)$. Its cross term
with the curl of the periodic fluctuation vanishes by integration
over the torus.

\begin{proposition}[Absolute graph bounds in each magnetic bundle]
\label{prop:constitutive-magnetic-uv}
The optical graph estimates extend to the charged action on each
such line bundle, including its dynamical flat connection in the
neutral vacuum preparation of Corollary~\ref{cor:constitutive-holonomy-uv}.
At fixed $\mathcal A_n$, every finite order in the coupling to
the Maxwell fluctuations has an absolute carrier-mode and
photon-momentum bound uniform in $n$ for unnormalized marked
numerators with the classical background factor taken out.
The same carrier estimates allow bounded, source-independent
gauge-intertwining matrix multipliers on the spatial hopping kernels
provided the
lower block is the adjoint of the upper. External charged slots
may be tested by smooth bundle sections with the required
endpoint equivariance (smooth in holonomy for differentiated terminal
sections). A normalized estimate uses the nonzero
vacuum of its own sector.
\end{proposition}
\begin{proof}
Condition on the complete surviving normal path and the initial
holonomy. Let $\mathcal H_n$ be spatial $L^2$ of the charged
bundle with its finite spin slots. Put $R_0=(\partial_t+m\beta)^{-1}$
on $L^2(\mathbb R,\mathcal H_n)$, extending the hopping by zero
outside the slab. At zero fluctuating electromagnetic field write
$K_t$ for that Hermitian off-diagonal hop, including the static
holonomy and magnetic parallel transport. Define $X=R_0K$,
$D=R_0^{-1}+K$, and $R=D^{-1}$. Unitarity of line transport
on each lift leaves every spatial Schur or squared
Hilbert--Schmidt estimate unchanged: bound the magnitude of the
periodized kernel by the sum of its Gaussian envelopes. The
matrix-valued spatial case only changes the finite norm constants.
The two integrals $\int\|K_t\|_{\rm op}^2\,\mathrm dt$ and
$\int\|K_t\|_{\rm HS}^2\,\mathrm dt$ are bounded by $C B_T$,
where $B_T$ is the additive radial cost of the optical proof.
For $R_0$ use its one-sided exponential time kernels. They give
$\|X\|_{\rm HS}^2\le C_m B_T$. Also
\[
 \|KR_0K\|_{\rm HS}^2
 \le C\int\|K_t\|_{\rm op}^2\,\mathrm dt\,
           \int\|K_s\|_{\rm HS}^2\,\mathrm ds .
\]
On each surviving path the hopping is bounded. The skew-adjoint
part of $\beta D$ gives $\|R\|\le1/m$. Hence
$A=(1+X)^{-1}=1-RK=1-R_0K+RKR_0K$ has operator norm at most
$C_m(1+B_T)$. No spatial Fourier diagonalization enters this
argument.

For a nonempty set $I$ of electric marks on a hop let
$X_I=R_0K_I$. A transverse mark multiplies its chord kernel by
a line integral of a plane wave, in modulus at most
$|z|e^{-a(y_t)|p|}$. Absorbing a fixed power of $|z|$ into a
widened Gaussian gives one factor $C_I\sqrt{\tau(y_t)}$ per
mark. A scalar mark on the rephased hop costs at most twice the
maximal interval transform $M_p(\nu)$ of the optical proof.
Consequently $\|X_I\|_{\rm HS}$ is bounded by $C_I B_T^{1/2}$
times the same individual photon envelopes as in
Theorem~\ref{thm:seam-calibrated-joint-uv}. Indeed the squared
spatial heat kernel after its lift sum costs at most
$C_I(1+\tau^{-(3+2b)/2})$, with $b=0$ or $1$ the order of the
spatial derivative. This remains true with all magnetic transition
phases and requires no cancellation between different lifts.
Harmonic marks use the interval factors of the preceding corollary.

Each term in a derivative of the relative matter logarithm is a
trace of a product of $AX_I$. For just one differentiated factor
the trace is $\operatorname{Tr}[(A-1)X_I]=
-\operatorname{Tr}(AXX_I)$, which again has two Hilbert--Schmidt
factors. In any spatial orthonormal basis and the carrier time
Fourier representation, write $z=(j,\nu)$ and use counting
measure times $\mathrm d\nu/(2\pi)$. For $k\ge2$ Hilbert--Schmidt
operators $T_i$ have square-integrable kernels, and
\[
 \int\prod_{i=1}^k |T_i(z_i,z_{i+1})|\,\prod_i\mathrm dz_i
 \le\prod_{i=1}^k\|T_i\|_{\rm HS},\qquad z_{k+1}=z_1 .
\]
This follows by composing the absolute kernels and using the
Hilbert--Schmidt Cauchy--Schwarz inequality, first on finite
projections or simple kernels and then by approximation. The bound
retains every integrated charged carrier variable, even in a
magnetic eigenbasis with no ordinary conserved spatial momentum.
For open words, the two $L^2$ endpoint vectors replace the
cyclic trace. At a fixed boundary time the kernel of $R$ is
uniformly bounded in spatial operator norm by $C_m(1+B_T)$ from
$R=R_0-R_0KR_0+R_0KRKR_0$, using the $L^2$ time envelopes on
its two outer ends. Thus the first marked hop against a boundary
vector gives an $L^2$ time vector at cost
$C_m(1+B_T)(\int\|K_I(t)\|_{\rm op}^2\,\mathrm dt)^{1/2}$
times the endpoint norm. The other boundary gives an $L^2$
vector by $R=A R_0$ acting on a time delta with a spatial test.
The same kernel inequality applies to the resulting open chain.

The absolute determinant is bounded by $\exp(C_m B_T)$.
Pairing the transverse, scalar and harmonic photon envelopes
and averaging the normal path now use exactly the integrable
radial costs already proved. The vacuum denominator for a positive
diagonal preparation is strictly positive by the full Hamiltonian
transfer. The classical factor $e^{-T E_n}$ is common to this
denominator and every numerator in its flux sector; it cancels
in normalized responses. In the bare vacuum logarithm it is the
explicit finite background contribution $-T E_n$ for $e\ne0$.
This is a background-field expansion at fixed quantized
$\mathcal A_n$. Its classical $e^{-2}$ term and the winding
weights preclude identifying it with an ordinary Taylor expansion
of a flux sum at $e=0$.
\end{proof}

If the compact gauge group is $U(1)$, take integer matter charges
in units of its fundamental charge and tensor powers of the given
line bundle. At fixed nonzero topological normalization $e_*$ one
may also sum over all magnetic fluxes. Give the $n$th bundle its
classical weight
\[
 w_n=\exp[-2\pi^2 T|n|^2/(e_*^2L)]
\]
and the same normalized neutral vacuum preparation in its own
Hilbert space. Use the free electric fluctuation order of
Proposition~\ref{prop:constitutive-magnetic-uv}, keeping $e_*$
fixed when counting fluctuation vertices. If $N_n^{(k)}$ is an
unnormalized marked coefficient with the classical weight removed,
that proposition gives $|N_n^{(k)}|\le C_k$ uniformly, for the
vacuum and neutral probe family with uniformly bounded test norms.
Thus $\sum_n w_n|N_n^{(k)}|<\infty$. The common diagonal vacuum
sum is finite and strictly positive, since even its $n=0$ term
is positive. Taking its logarithm or differentiating the single
normalized sum preserves finiteness at every fluctuation order.
This sums prior to normalization, retaining the correlations caused
by sampling one bundle for a whole diagram. When $e_*$ is identified
with the expansion variable, the nontrivial sectors have weights
flat at zero; electric Taylor coefficients alone do not reconstruct
these contributions at nonzero charge.

\paragraph{An interacting neutral curvature block and a mixed charged vertex.}
Convergence on the electromagnetic quotient settles only the fields
which enter its normalized source. Consider an independent real
gauge singlet $\varphi$ of mass $\mu>0$ on the same slab, sharing
the actual normal trajectory at every interaction. Give its free
field the boundary vacuum covariance
$(\nu^2+|k|^2+\mu^2)^{-1}$. For $P_a=e^{-a|\nabla|}$ let
$\varphi_a=P_a\varphi$. In addition to its kinetic and mass terms
take the positive bare potential
\[
 \lambda\int \mathrm dt\,\mathrm d^3x\,
                   \varphi_{a(t)}(t,x)^4,\qquad \lambda\ge0 .
\]
The quartic is the square of the curvature term $\Phi^2$ of an odd
two-slot scalar endomorphism $\Phi=\varphi_a
\left(\begin{smallmatrix}0&1\\1&0\end{smallmatrix}\right)$ with a
positive pairing on this block. The ordinary free kinetic and mass
terms here are fixed separately. On a forward charged hopping arrow
from $z$ to $x$ include the additional unitary factor
\[
 \exp\!\left(i y\,a(t)
       \frac{\varphi_{a(t)}(t,x)+\varphi_{a(t)}(t,z)}2\right),
\]
and use its adjoint on the reversed arrow. The field has inverse
length dimension, so the measured scalar in this character is
dimensionless. This coupling preserves the electric Gauss law.
Unlike two independent systems conditioned on the same radius,
these fields can exchange a scalar quantum on a charged loop.

\begin{theorem}[Bare mixed scalar and electromagnetic coefficients]
\label{thm:omega-singlet-uv}
On the optical source selected above, with the stated finite slab,
torus and positive diagonal preparation, every finite joint order
about $e=y=\lambda=0$ has an absolute graph bound for connected
vacuum, photon and elementary dressed charged functions. The hopping
strength may take any fixed finite real value. Finite scalar
insertions smeared by smooth external tests are also allowed. The
statement includes the unaltered quartic self-contractions and every
mixed scalar and Wilson contact.
\end{theorem}
\begin{proof}
At a positive depth the free scalar averaged by $P_a$ has finite
equal-time variance. With a moving depth its unrenormalized quartic
is a positive multiplication form. Its sum with the free radial
and field forms is densely defined: use finite-particle smooth
tests supported in a compact subinterval of the collar. The sum
of the positive closed forms is closed. For real $e,y$ the hopping
kernel has unchanged modulus. Its small relative form bound and the
positive diagonal vacuum normalization follow as in the optical
construction. This defines the interacting transfer on the finite
region. The graph assertion concerns joint perturbation coefficients
of this action. Positive $\lambda$ permits the usual right
derivatives at zero. The free oscillator vacuum has its canonical
zero of energy; the interaction is the literal fourth power.

Fix the complete normal path and expand at a specified finite order.
A scalar attachment on a hop has multiplier $a_i e^{-a_i|p|}$.
Its two possible spatial endpoints give only momentum translations.
An attachment on a quartic vertex has multiplier $e^{-a_i|p|}$.
Put $a_{\min}=\min_t a(t)$ in reference units. At either endpoint,
retain $(1+|p|)^{-1}e^{-a_{\min}|p|/2}$; at a quartic endpoint
pay in addition $C(1+a_i^{-1})$. The inequalities follow from
boundedness of $(1+x)e^{-x/2}$. Summing one full scalar line after
integrating its own frequency therefore costs at most
\[
 C\sum_{p}\frac{e^{-a_{\min}|p|}}
 {(1+|p|)^2\sqrt{|p|^2+\mu^2}}
 \le C_V(1+|\log a_{\min}|).
\]
The mass controls the spatial zero mode. Momentum conservation can
be imposed after this positive lattice bound. The same estimate
applies to a contraction at one quartic vertex. Smooth external
tests have integrable Fourier weights and introduce only finite
constants. The charged trace, its independent closure frequency,
and the electric and Coulomb contractions have their previous optical
bounds. In particular no derivative falls on a sampled scalar when
the translated charged Gaussian is used.

Each quartic vertex has a single time and at most four incidence
costs. Its time integral is bounded by
$C\int_0^T (1+r_t^{-4/3})\,\mathrm dt$. Each scalar line contributes
one logarithm and there are only finitely many at the chosen order.
The five-dimensional radial Hardy and minimum-depth estimates of
Theorem~\ref{thm:seam-calibrated-joint-uv} control the products,
together with the actual matter determinant. This proof includes
simultaneous approaches to the tip, all angular transitions and
the common boundary average. The same absolute moment estimates
justify source differentiation, or the ordered Gaussian moment
expansion on compact interior collars followed by exhaustion.
Vacuum normalization takes finite products of these coefficients.
\end{proof}

\paragraph{A neutral sector which is invisible to electric matching.}
There is also a test of any proposal to infer the unrestricted
theorem solely from the photon and charged source. Form a block
direct sum with a neutral real scalar and choose its curvature
pairing to give the ordinary local $\lambda\varphi^4/4!$ and
massive scalar kinetic action. The odd endomorphism above with
unmeasured $\varphi$ has squared zero-form curvature
$\varphi^2 I$; tracing its squared norm gives a positive quartic
term. Different form degrees can have separate positive pairings.
The singlet is invariant under the electric constraint. Choose the
tensor-product preparation and no mixed arrows. Its normalized
zero-source factor then cancels from every electromagnetic
derivative. Its own four-point 1PI coefficient at order $\lambda^2$
contains the positive Euclidean bubble, proportional to
\[
 \int_{\mathbb R^4}
 \frac{\mathrm d^4 k}
 {(|k|^2+\mu^2)(|k+p|^2+\mu^2)} .
\]
For fixed external transfer $p$, on each sufficiently large
dyadic shell the integrand is bounded below by $c|k|^{-4}$,
and the shell contributes at least a positive constant. The three
channels of the four-point vertex have the same sign. A vacuum
normalization cannot remove this four-source derivative.
This is an explicit independent sector allowed by electromagnetic
matching alone. It is not an assertion that the physical GCM
already contains this local block. To exclude it one needs a
constitutive law for the neutral block, such as the optically
measured quartic just proved, or a calculation showing that the
physical field content omits it.

\paragraph{An independently quantized normal field with its own hard momentum.}
The graph measure above is sufficient, but the estimate can also
include dynamical bulk fluctuations. Here is an explicit extension
in which the complete charged determinant sees a bulk Gaussian
field. In addition to the charged, radial and boundary photon blocks,
quantize a tangentially transverse vector $C$ on the optical
half-space with Dirichlet boundary condition. Its raw quadratic action is
\[
 \tfrac12\int \mathrm dt\,\mathrm d^3x\,\mathrm dy\,
 \big(|\partial_t C|^2+|\nabla_xC|^2+|\partial_yC|^2\big).
\]
Use its free Fock vacuum on the slab and the same convention
fixing the spatial harmonic mode. This is an independent neutral
transverse block. Its boundary value vanishes. Its coupling to a
charge is specified by a matrix element across the normal fiber.
Take a fixed real $h\in C_c^\infty((0,\infty))$ of unit
$L^2(\mathrm dy)$ norm and at optical depth $a$ use
$h_a(y)=a^{-1/2}h(y/a)$. This normalization is fixed by the
isometric identification of the normal packet at different scales.
Use the following additional spatial potential inside the same
charged Wilson link:
\begin{equation}
 e^{-a|\nabla_x|}\int_0^\infty h_a(y) C(t,x,y)\,\mathrm dy .
 \label{eq:constitutive-bulk-test-uv}
\end{equation}
The potential is transverse and invariant under the boundary
electric gauge group. The scalar boundary potential and its
Gauss coupling remain those of the charged source. The profile
is part of the microscopic choice. The theorem below holds for every
such normalized smooth profile. It is possible to prescribe one
profile once for all charged hops.

\begin{corollary}[Bare graphs with a dynamical normal momentum]
\label{cor:constitutive-bulk-uv}
For either source of Proposition~\ref{prop:constitutive-source-uv},
adding the free transverse bulk action and coupling
\eqref{eq:constitutive-bulk-test-uv} preserves the all-orders bare
graph bound. It holds jointly over each bulk photon's independent
normal momentum, carrier frequency and spatial momentum, including
closed and open charged graphs in the preceding finite-slab
preparations.
\end{corollary}
\begin{proof}
Let $p$ be a tangential spatial momentum, $\eta>0$ the sine momentum
in $y$, and $R=(|p|^2+\eta^2)^{1/2}$. The free field has mixed
Fourier covariance $P^{\rm T}(p)/(\nu^2+R^2)$ in these variables.
With a unitary sine transform the packet is
$a^{1/2}\widehat h(a\eta)$, with rapid decay at large argument
since $h$ has smooth compact support in the interior. One bulk
mark on a hop carries the usual heat-width factor $\sqrt\tau=a$.
The full endpoint thus has the absolute envelope
\[
 C_N a^{3/2} e^{-a|p|}(1+a\eta)^{-N}.
\]
This accounts for a new momentum, which cannot be removed by
the boundary Maxwell constraint. Put $a_{\min}=\min_t a(r_t)$
in fixed reference units. Since $a$ is bounded above on the
quartic collar the endpoint is at most
$C(1+R)^{-3/2}(1+a_{\min}R)^{-2}$; use sufficiently large
$N$ and $|p|+\eta\ge R$.
The bound holds before opening a charged temporal trace, including
at its variable endpoint time. For a complete bulk line integrate
$\nu$ to obtain $\pi/R$. The measure of the lattice-continuum
shell $K\le R<2K$ under $\sum_p\int_0^\infty\mathrm d\eta$
is at most $C_V K^4$ for $K$ larger than the reciprocal box
scale. Low modes are integrable with the stated infrared convention.
The product of the two endpoint estimates therefore gives a total
cost $C(1+|\log a_{\min}|)$ per bulk contraction. With the
chosen depth law, every finite product of these logarithms is
integrable against the radial path and boundary estimates of
Theorem~\ref{thm:seam-calibrated-joint-uv}.

For almost every surviving radial path its minimum is positive,
so the Wilson increments of the measured bulk field have finite
Gaussian variance. At real electric coupling each such phase
has modulus one. The pathwise charged Schur bound and absolute
Hilbert--Schmidt determinant bound are uniform in $C$. The bulk
Maxwell Fock block is positive. Adding it to the free energy
preserves the closed semibounded form construction because the
charged pair bound only uses those same two spatial norms.
On a compact radial subinterval the time-sliced transfer of the
independent bulk Gaussian inserts exactly its Wilson phases in
the matter determinant. The uniform radial majorants permit
Friedrichs exhaustion as in the optical proof. Thus integrating
the raw bulk field retains all photon contacts of this determinant.
The one diagonal vacuum normalization is strictly positive by
the full Hamiltonian and its free bulk Fock vacuum.
\end{proof}

A normal packet read as point evaluation has a different
covariance. Here the normalization of its normal wavefunction
provides the factor $a^{1/2}$ at each end, and the charged
displacement provides the factor $a$. Together they control the
ordinary four-spatial-dimensional Maxwell integral without a
bulk momentum cutoff or a higher-derivative photon kinetic term.
The calculation treats the specified transverse bulk field.
Additional charged bulk matter, a normal electric Gauss field or
a different boundary pairing would require their own joint
source calculation.

\paragraph{Continuous loop momenta with a moving normal state.}
A finite-time, infinite-space version of the optical estimate can be
proved directly for complete charged components joined by photons.
This also identifies the spatial volume factor when several such
components communicate only through the normal state. Keep the
time slab, the exponents and the optical action of
Corollary~\ref{cor:seam-optical-shifts-uv}. Use continuous spatial
momenta on the selected Euclidean plane. The number of independent
carrier frequencies still equals the number of closed charged traces.
Temporal photon frequencies are independent on the slab.

\begin{proposition}[Continuous spatial graph densities at moving depth]
\label{prop:omega-continuum-uv}
At each finite joint hopping and electric order, remove the overall
spatial translation delta of each component connected by charged
propagation and photon edges. Its complete graph density has a joint
absolute bound in the continuous charged and photon loop variables,
the remaining times, and the full radial and angular measure.
Square-integrable normal endpoint states are allowed. The bound
applies to products of finitely many such densities on the same
normal path. For elementary charged endpoints use nonexceptional
external spatial transfers or integrate over joint Schwartz
momentum packets as in Theorem~\ref{thm:seam-uv-open}.
An internal photon bridge forced to zero spatial momentum is omitted
with the same harmonic-mode prescription as in that theorem.
\end{proposition}
\begin{proof}
First fix a surviving normal path. A closed trace is rooted at a
spatial vertex $X$. For a soft transverse test of momentum $(\nu,p)$,
subtract the constant spatial vector equal to its unmeasured value
at $(0,X)$. A constant time-independent connection couples through
the sum of oriented displacements, which vanishes on the spatial
continuum for each closed path. The subtraction can be made for
every selected test before taking any coefficient. At a point $x$
on a Wilson segment the difference of the phase factors costs at
most
\[
 C\big(T|\nu|+|p|[a(r_t)+|x-X|]\big).
\]
Indeed the phase is $e^{i\nu t-a(r_t)|p|+ip\cdot x}$ before
subtraction. Represent each exponential difference by integrating
its first derivative on the unit segment. The coordinates $x-X$
are sums of displacements on the rooted trace and differentiate
the translated spatial kernels in momentum space. Multiplication
by any fixed number of such coordinates costs only finitely many
further spatial derivatives. A differentiated heat hop gains a
power of $\sqrt{\tau}$ per derivative. The Gaussian that controls
the charged spatial trace remains integrable. For a scalar test
the time primitive at a hop multiplies the difference of the two
spatial endpoint phases. At $|p|\le1$ this yields a factor $|p|$
and a displacement by the same identity. These improvements apply
at a closed end. An elementary open endpoint retains its actual
phase and needs no soft improvement. This proof uses the moving
optical readout without differentiating the normal path.

Write $B_T$ for the additive radial majorant in
Theorem~\ref{thm:seam-calibrated-joint-uv} and $a_{\min}$ for the
minimum optical length on the path. That theorem's split time
trace and maximal interval estimate keep the integrals absolute
in charged carrier frequencies and each photon frequency.
The spatial trace now uses the ordinary translated Gaussian
integral. After the independent photon frequency integrations,
let $b=0,1,2$ count the open ends of a photon line. Its positive
spatial envelope at $0<|p|\le1$ is at most $C_T(1+|p|^{-b})$.
For a transverse line this follows by integrating
$\min(1,|\nu|+|p|)^{2-b}/(\nu^2+|p|^2)$ over $\nu$.
For a Coulomb line the maximal interval estimate costs $C T$
and the two closed-end gains, when present, multiply $|p|^{-2}$.
For $|p|>1$ one can use the sum of
\[
 C |p|^{-3}e^{-a_{\min}|p|},\qquad
 C |p|^{-2}\int_0^T e^{-2a(r_t)|p|}\,\mathrm dt .
\]
The first follows from the two transverse end estimates and the
frequency integral $\pi/|p|$. The second is the instantaneous
line estimate. Their combined spatial integral is bounded by
$C(1+|\log a_{\min}|+B_T)$, and their supremum here is bounded
by $C_T$. Thus every envelope is in $L^1(\mathbb R^3)$, the
two-closed-end envelope is bounded, and the one-open-end envelope
is in $L^2(\mathbb R^3)$. All needed norms are bounded by radial
costs with finite moments in the actual endpoint measure.

Spatial momentum is conserved at each charged component. For
closed components alone, use the suprema on a spanning tree and
integrate the remaining lines in their own independent momenta.
With open lines the spatial tree and convolution argument of
Theorem~\ref{thm:seam-uv-open} uses precisely the three norms
just obtained. Two root edges with one closed end cost two
$L^2$ norms, even when their soft momenta coincide. Joint packets
integrate the freely variable external transfers before this
estimate. The remaining radial majorant is a polynomial in $B_T$
and $1+|\log a_{\min}|$. Its normal integral is bounded by a
constant times the product of the endpoint norms, by the Hardy
and hitting estimates of Theorem~\ref{thm:seam-calibrated-joint-uv}. A finite product
of the conditional bounds only raises the polynomial degree.
The exact angular heat contraction handles both endpoints.
\end{proof}

There is also a fixed-hopping interpretation of this estimate.
At arbitrary finite hopping strength apply it to the conditional
connected matter vertices, with the normal trajectory fixed before
taking the charged logarithm. Use its rooted derivative in the
hopping multiplier as in Proposition~\ref{prop:ct-uv-rooted-soft}.
In continuous space this is a trace density at $X$. In zero
electromagnetic background the dressed inverse preserves spatial
momentum. Its differentiated kernels obey the same time bounds:
differentiate $R(k)$ by $\partial_i R=-R(\partial_i K)R$ and
iterate. Every positive spatial derivative of the scalar or spinor
hop $K$ has uniformly bounded operator norm for the given profiles.
The internal times are integrated on a finite slab and the endpoint
kernel bound is polynomial in $B_T$. At the one-vertex trace use
the off-diagonal cancellation against the onsite inverse to recover
a second hop. The root Gaussian and split time trace then bound
the charged momentum integrals as above. This proves the continuum
estimate for these conditional photon-connected densities at each
finite electric order, also after integrating them against the
original normal heat law with $L^2$ endpoints.

In a full spatially homogeneous source the normal path may be
shared by otherwise disjoint pieces. At fixed hopping order expand
the ordinary determinant and its normalization before integrating
the normal variables. If there are $c$ spatially disconnected
pieces, each conserves its own spatial momentum. A radial
cumulant can correlate their densities but does not exchange
any spatial momentum. For a vacuum term there are $c$ independent
spatial centers, each giving a volume factor. This is why the
proposition removes one translation delta for each such piece.
At fixed hopping strength the full normal law is also weighted by
the extensive zero-field matter determinant. The continuum
density estimate alone neither constructs that limiting law nor
proves convergence after division by a single spatial volume.

\paragraph{The volume dependence of a single moving normal state.}
A spatial torus in the optical proof fixes an infrared convention,
but its volume also sets the number of charged modes coupled to the
same radial path. This second dependence matters even with the
electromagnetic coupling set to zero. Use the neutral normal ground
preparation, the onsite matter vacuum, a cubic torus of volume $V$,
and a fixed slab duration $T>0$. All lengths are in the same
reference units. Keep the ordinary, angular-preserving optical hop.
Write $b=0$ for the scalar and $b=1$ for the Dirac derivative,
and take a nonzero hopping coefficient $v$. On paired spin slots
the singular value of the off-diagonal hop at fixed $r,k$ is
\[
 B_k(r)=|v|\,|k|^b e^{-\kappa r^\alpha|k|^2}.
\]
For $b=0$ the momentum factor is identically one. A fixed finite
spin multiplicity only changes constants. Let $Z_V(T)$ be the
diagonal vacuum matrix element divided by its value at zero
hopping and zero electric charge. This definition uses the
onsite vacuum energy fixed by the given action.

\begin{proposition}[Volume growth of the optical vacuum with one common radius]
\label{prop:upgrade-common-volume-uv}
Put $s=(3+b)\alpha/2$ and use the exponent hypotheses of
Theorem~\ref{thm:seam-calibrated-joint-uv}. For every fixed
$T>0$ there are positive constants independent of $V$ such that,
for all sufficiently large cubes,
\[
 c_T V^{2/(2-s)}\le\log Z_V(T)\le C_T V^{2/(2-s)} .
\]
Under the resulting positive zero-electric-charge vacuum path law,
$\Pr(\min_{[0,T]}r_t\ge\epsilon)\to0$ as $V\to\infty$ for
every $\epsilon>0$. In particular, for the derivative spinor
with $\alpha=2/3$ the logarithm grows as $V^3$ up to constant
factors. The statement concerns the actual vacuum normalization
before division into connected correlation functions.
\end{proposition}
\begin{proof}
On a frozen radial fiber the exact pair ground-energy shift in
one slot is $m-\sqrt{m^2+B_k^2}$, of magnitude at most $B_k$.
This follows by diagonalizing the vacuum and pair block with
diagonal entries $0,2m$ and off-diagonal entries of magnitude
$B_k$. The two singly occupied states have energy $m$.
Summing a Gaussian on the reciprocal lattice gives
$\sum_k B_k(r)\le C V(1+r^{-s})$ uniformly in $V\ge1$.
Thus on the full direct integral the Hamiltonian is bounded
below in form sense by $H_0-C V(1+r^{-s})$, where $H_0$ is the
radial Friedrichs kinetic operator. Radial Hardy control of
$r^{-2}$ and the pointwise inequality
$C V r^{-s}\le \eta r^{-2}+C_\eta V^{2/(2-s)}$ give the upper
bound on $\log Z_V(T)$. Here $s<2$ and $\eta$ is fixed sufficiently
small. The free radial normalization is independent of $V$.

For the lower bound restrict the radial path to an annulus
$\delta<r<2\delta$ throughout the slab. In the positive radial
ground preparation its unperturbed mass is bounded below by
$c\delta^5\exp(-C T/\delta^2)$. Indeed the incoming and outgoing
ground functions are bounded below near the tip. In the annulus
the first Dirichlet eigenvalue scales as $\delta^{-2}$ and the
squared overlap of its normalized positive eigenfunction with the
constant function scales as $\delta^5$ in $r^4\mathrm dr$.
For $|k|$ between fixed positive multiples of
$\delta^{-\alpha/2}$, each of at least $cV\delta^{-3\alpha/2}$
modes has $B_k(r)\ge c'\delta^{-b\alpha/2}$ over the whole path.
For every such $k$ a time-independent slot basis makes its pair
matrix element real and negative. At zero electric charge the
conditional time-ordered Fock vacuum factor then has an expansion
with positive even-hop terms, monotone in $B_k$ at every time.
All other modes give factors at least one. For the constant lower
value $B$ the individual factor is exactly
\[
 e^{-mT}\left[\cosh(TE)+\frac mE\sinh(TE)\right],
 \qquad E=\sqrt{m^2+B^2}.
\]
It exceeds one for $B>0$; its logarithm grows as $TB+O_T(1)$
at large $B$. The shell therefore contributes at least
$\exp(c_T V\delta^{-s})$. These bounds retain the determinant;
they can first be applied to any finite subset of the shell.
Set $\delta=D V^{-1/(2-s)}$ with a sufficiently large fixed $D$.
The gain exceeds the radial cost by a positive multiple of
$V^{2/(2-s)}$, proving the lower bound. On paths with minimum
at least $\epsilon$, the fiber lower bound, or the product norm
bound on the transfer along the path, yields at most
$\exp(C_{T,\epsilon}V)$. Division by the lower bound proves
the probability assertion.
\end{proof}

This loss of an extensive vacuum pressure is compatible with
absolute ultraviolet convergence on every fixed torus. It shows
why taking a volume limit inside the moving-depth proof requires
an independent argument. Connected local probes could still have
limits after division by the same vacuum. The proposition alone
does not decide those limits.

\paragraph{A geometric pressure test for a modified volume law.}
One can test an alternative microscopic assignment without changing
a Wilson or Gaussian spatial kernel. On the same two-dimensional
quartic normal metric $\mathrm dr^2+r^8\mathrm d\theta^2$,
the interior scalar curvature is
$\mathcal R=-2(r^4)''/r^4=-24/r^2$. Include an intensive positive
normal curvature energy $\lambda(-\mathcal R/24)$ throughout
space, with $\lambda>0$ in the chosen units. For a common radius
the extra bare Hamiltonian is $V\lambda/r^2$. Its sign and its
volume law are constitutive choices. It is a geometric term
present already at zero hopping and electric charge. Its
nonnegative Feynman--Kac weight preserves the finite-volume
absolute numerator estimates and the diagonal vacuum normalization
remains strictly positive. To claim a single fundamental
normalization the value of $\lambda$ must be fixed in that same
microscopic assignment.

In this modified branch the zero-electric-charge ground energy
$E_V$ per unit volume has the limit
\begin{equation}
 \begin{aligned}
 \lim_{V\to\infty}\frac{E_V}{V}
 &=\min_{0<r\le1}\left\{\frac{\lambda}{r^2}-I(r)\right\},\\
 I(r)&=\nu\int\frac{\mathrm d^3k}{(2\pi)^3}
          \left(\sqrt{m^2+B_k(r)^2}-m\right).
 \end{aligned}
 \label{eq:upgrade-pressure-min-uv}
\end{equation}
where $\nu$ counts the paired spin slots. To verify the lower
bound, replace $I$ by the corresponding sum $I_V$ with the factor
$1/V$. The exact fermionic ground-energy bound gives
$H\ge H_0+V(\lambda/r^2-I_V(r))$. The lattice sum converges to
$I$ uniformly on each closed subinterval away from zero by the
common Gaussian tails and uniform continuity. The same bound as
in the proposition gives $I_V(r)\le C(1+r^{-s})$. The effective
potential therefore tends to $+\infty$ at the tip uniformly in
the volume. For the upper bound localize a normalized radial test
near any fixed $r_*>0$ and use the fermionic ground state for the
frozen value $r_*$ at every radius in the test. Its lack of radial
dependence costs no fermionic gradient. On compact subintervals
$V^{-1}\sum_k\sup_r|\partial_r B_k(r)|$ is bounded independently
of $V$ by Gaussian summability. The change of the energy density
from a packet of width $h$ is therefore $O(h)$, including the
curvature potential, while its radial kinetic cost per volume is
$O(1/(Vh^2))$. First send $V$ to infinity, then $h$ to zero,
and approach the outer endpoint from the interior if necessary.
A state sequence attaining this energy density concentrates in
radius on the minimizing set, since off each of its neighborhoods
the limiting potential has a positive gap above its minimum.
This ground-state statement supplies a pressure criterion at zero
electric charge. Infinite-volume connected source derivatives and
their infrared and large-time limits require additional estimates.

\paragraph{What the zero-charge equilibrium measures.}
The same volume law also gives a statement at fixed nonzero
hopping strength. Let $E_V^0$ be the spectral lower bound of the
unmodified Hamiltonian at zero electric charge. The Hardy bound
and the diagonal transfer lower estimate in
Proposition~\ref{prop:upgrade-common-volume-uv} place $E_V^0$
between two negative constants times $V^{2/(2-s)}$ at large $V$.
Let $\Psi_V$ be a unit vector whose energy is at most $E_V^0+1$.
For fixed $0<\epsilon<1/2$ take a smooth radial partition
$\chi^2+\eta^2=1$, with $\eta=0$ below $\epsilon$ and $\eta=1$
above $2\epsilon$. The radial localization identity costs at most
$C_\epsilon$ from $|\chi'|^2+|\eta'|^2$. Every other part of this
Hamiltonian commutes with radial multipliers. On the support of
$\eta$ the fiber bound is $-C_\epsilon V$; the part localized
by $\chi$ has lower bound $E_V^0$. Hence
\[
 (-C_\epsilon V-E_V^0)\|\eta\Psi_V\|^2\le 1+C_\epsilon .
\]
Thus the mass above $2\epsilon$ tends to zero also for actual
approximate ground states at finite hopping. For any fixed
momenta $p,k$, the expectations of $e^{-a(r)|p|}$ and
$e^{-\kappa r^\alpha|k|^2}$ consequently tend to one. An
equilibrium readout on this common-radius branch loses its finite
optical width in the spatial volume limit. This pointwise result
does not interchange that limit with integration over hard momenta.

\paragraph{A volume test after vacuum normalization.}
There is a related test that keeps a charged endpoint response in
the numerator. Remain on the ordinary optical branch of
Proposition~\ref{prop:upgrade-common-volume-uv} at zero electric
charge. For this calculation write $g$ for the hopping coefficient
and $a_k(x)=|k|^b e^{-\kappa x^\alpha|k|^2}$ for its magnitude with
that coefficient removed. Prepare the positive radial ground
function at the two ends and let $\mathbb E_0$ denote its stationary
path law. Along this law let $L_{2,V}$ denote the coefficient of
$g^2$ in the conditional matter vacuum logarithm. A pair created at
$u$ and annihilated at $w$ gives exactly
\[
 L_{2,V}=\nu\int_{0<u<w<T}\!e^{-2m(w-u)}
       \sum_k a_k(X_u) a_k(X_w)\,\mathrm du\,\mathrm dw .
\]
As $V\to\infty$, $L_{2,V}/V$ converges in $L^2(\mathbb E_0)$ to
the positive functional $F$ defined by the same formula with the
sum replaced by $\int\mathrm d^3k/(2\pi)^3$. Indeed, each path
has a positive minimum on the surviving slab. The Gaussian
Riemann sums converge there. For domination the spatial sum per
unit volume is bounded by $C(1+X_u^{-(3+2b)\alpha/2})$.
The integrated right-hand side has all fixed moments by the
radial Hardy estimate. The continuum momentum integral is a
positive constant times
$(X_u^\alpha+X_w^\alpha)^{-(3+2b)/2}$.
Thus $F$ strictly decreases when the whole trajectory increases.
Its variance is positive, since the free ground process has
positive probability of remaining for time $T$ in either of two
disjoint annuli.

For comparison, the conditional coefficient $L_{4,V}$ in the
logarithm has $\mathbb E_0|L_{4,V}|\le C_T V$. Conditional on the
path the log is a sum of independent mode contributions. In one
mode the absolute fourth coefficient is bounded by a constant
times the integral of $\prod_{i=1}^4 a_k(X_{t_i})$ over the four
times, by the pair evolution and the quadratic log subtraction.
The reciprocal lattice sum divided by $V$ is bounded by
$C[1+(\sum_i X_{t_i}^\alpha)^{-(3+4b)/2}]$. The geometric mean
bound allocates exponent $(3+4b)\alpha/8<2$ to each integration
time. The radial Hardy estimate again integrates the fourth
power of that time integral. Consequently the complete, connected
vacuum satisfies
\[
 \lim_{V\to\infty} V^{-2}[g^4]\log Z_V(T)
     =\tfrac12\operatorname{Var}_0(F)>0 .
\]
Two charged returns that are independent at fixed geometry are
coupled by averaging their common normal trajectory.

For a normalized charged test, fix a nonzero momentum $k$ on a
sequence of increasing cubes whose reciprocal lattices contain
$k$. Put the particle and its paired hole in that normalized mode
at the final end only, with photon vacuum at both ends. A fixed
slot phase makes the first hop matrix element positive.
At fixed path the ratio of this mode's pair amplitude to its
vacuum amplitude is $g B_k+g^3D_k+O(g^5)$, where
\[
 B_k=\int_0^T e^{-2m(T-u)}a_k(X_u)\,\mathrm du ,
 \qquad |D_k|\le C_{T,k}.
\]
The bound follows by expanding the single two-level transfer to
third order with $a_k$ bounded uniformly in depth. Other modes
contribute their original vacuum factors. In the full ratio
$P_{k,V}$ the common-vacuum division gives exactly
\[
 [g^3]P_{k,V}
     =\mathbb E_0 D_k+\operatorname{Cov}_0(L_{2,V},B_k).
\]
This is a connected two-endpoint response because the separate
one-fermion expectations vanish by fermion parity.

In this preparation $\operatorname{Cov}_0(F,B_k)>0$. For details,
the positive radial ground transformation gives a one-dimensional
diffusion with constant noise coefficient and smooth drift
$4/r+2\psi'_0/\psi_0$ in the interior. Its paths from two ordered
starting points can be driven by the same Brownian motion. They
stay strictly ordered on the finite slab by uniqueness and the
linear difference equation for the drift difference, localized
on compact interior intervals. Neither boundary is reached in the
ground process. In particular its transition kernels preserve
monotonicity. Conditional on the starting point, any finite set
of ordered-time observations is associated: given the past, the
last coordinate has nonnegative covariance for two monotone
functions in that coordinate; the conditional means preserve
monotonicity in the past by the monotone transition. Induction
proves the assertion. Quadrature of bounded monotone truncations
and then the $L^2$ bounds extend it to $F$ and $B_k$, both
decreasing path functionals. Their conditional covariance is
therefore nonnegative. Their two conditional means strictly
decrease with the starting radius by the strictly ordered
coupling; the starting ground distribution is nondegenerate.
The covariance of those two means is strictly positive, proving
the claim. Thus $V^{-1}[g^3]P_{k,V}$ tends to this positive
covariance. The unbounded volume factor survives division by the
vacuum. It comes from a common geometric degree of freedom and a
joint hopping expansion. It does not diagnose a divergent
electric ultraviolet subgraph at fixed hopping strength.

The same conclusion holds with elementary spatial packets of fixed
continuum profile. For the scalar take a smooth Fourier packet
with compact support in a ball away from zero. Test the particle
and the conjugate hole packet at the same boundary. Total spatial
momentum conservation pairs only the corresponding modes, with
weights proportional to $|\widehat f(k)|^2/V$. Normalize the
two packets by their one-particle norms. For the Dirac derivative
choose a sufficiently small momentum ball and a nonzero constant
spin vector for the particle. On the hole test use its corresponding
unitary spin image under the off-diagonal block of
$\boldsymbol\alpha\cdot k/|k|$, with the adjoint orientation chosen
for the reverse map. This smooth normalized slot pairing on the
packet has a positive $a_k$ as above. In either case the inverse
Fourier transforms are Schwartz tests. In the normalized lattice
sum $B_k$ and its cubic remainder are uniformly bounded on the
chosen compact momentum support. Riemann convergence of the packet
averages gives a positive weighted average of the strict
covariances, using once more $L_{2,V}/V\to F$ in $L^2$.

\paragraph{Changing which normal states the charged determinant traces.}
Replacing a normal coordinate shared by the prepared source with a
separate normal label on every fermion changes the short-duration
test. Consider the corresponding scalar comparison action on the
same finite time slab. The one-particle spatial and normal space is
$L^2(\mathbb T^3)\otimes L^2(M,\mathrm d\mu)$, where
$M=(0,1)\times S^1$ has the quartic collar metric. Let
$H$ be the nonnegative normal Friedrichs Laplacian and let the
two onsite slots have energies $\pm(m+H)$. At zero electric charge
a scalar pair hop in spatial mode $k$ multiplies the normal fields
by $F_k(r)=e^{-\tau(r)|k|^2}$, with the adjoint on the reverse
hop. There is no angular projection or angular shift in this test.
The onsite vacuum fixes the relative determinant as before.

At second hopping order the coefficient in its bare vacuum logarithm
is a positive multiple of
\begin{equation}
 \sum_k\int_0^T (T-u)e^{-2mu}
       \operatorname{Tr}_M\!\left(e^{-uH}F_k
                                   e^{-uH}F_k\right)\,\mathrm du .
 \label{eq:omega-fiber-trace-uv}
\end{equation}
Indeed a hop at $s$ and its adjoint at $s+u$ propagate a particle
and a hole, each with its own normal heat kernel; there is no
vacuum term linear in hopping. The trace at each $u>0$ is the
integral of $F_k(x)F_k(y)|K_u(x,y)|^2$ in the two normal points.
For the zero spatial charged mode $k=0$ it is simply
$\operatorname{Tr}_M e^{-2uH}$. Even a compact annulus of normal
states gives a logarithmic divergence at $u=0$. To see this, take
a coordinate rectangle compactly inside the collar. A sine-product
trial subspace with both mode numbers at most $n$ has dimension
$n^2$. On that rectangle the normal metric and density are uniformly
comparable to the Euclidean ones, and the largest Rayleigh quotient
on the extended-by-zero subspace is at most $C n^2$. By the
min--max principle the Dirichlet rectangle heat trace at time $2u$
is bounded below by $n^2 e^{-2Cun^2}$. Take
$n=\lfloor u^{-1/2}\rfloor$ for small $u$. Positivity and killing
on the rectangle give the same $c/u$ lower bound for the full
normal trace. On the spatial continuum per unit spatial volume,
restrict $k$ to a fixed compact ball. There $F_k$ is uniformly
positive on the rectangle, so the positive two-kernel integral
is again bounded below by $c/u$ there. This is an ultraviolet
failure of the bare vacuum logarithm already at $g^2 e^0$, even
with matched photon and charged lengths.

The test changes the normal sewing and keeps all its physical
multiplicities in the ordinary fermion trace. In the common-clock
source a single prepared normal heat kernel connects the charged
events and the slab boundaries; it does not trace two independently
propagating normal carriers in a shrinking interval. Equation
\eqref{eq:omega-fiber-trace-uv} concerns the unsubtracted vacuum
coefficient for the alternative choice. Its leading divergence
does not by itself give a divergence in a normalized photon
response. It does show that a change of normal sewing made to
restore spatial clustering requires a new ultraviolet proof.

\paragraph{A transverse current tests the sector-plane spectrum.}
A positive Hamiltonian gives nonnegative excitation energies in its
own carrier frame. Exact relativistic reconstruction on a fixed
sector plane requires more. In intrinsic carrier units with light
speed one, the vacuum energy--momentum measure of each physical
current must be supported on $E\ge |P|$. This follows directly by
applying the joint translation spectral projection to the current
acting on the vacuum \cite{StreaterWightman}. The condition also applies to an external
spatially smeared current, if that current acts in the same
relativistic theory. Using the coordinates of its Minkowski
readout leaves the joint spectrum of the physical translations
unchanged.

\begin{proposition}[A spectral test of the optically calibrated spinor]
\label{prop:omega-current-cone-uv}
Use the degree-one optical spinor with nonzero velocity $v$,
the scalar phase $e^{i\ell\theta}$ at pair creation and its adjoint
at annihilation. The integer $\ell$ may vanish. Use the neutral
normal ground state and the onsite Fock vacuum for the free
stationary preparation. Write $g$ for the multiplier of the whole
spatial hopping term and $e$ for the electric coupling. In the
separated-time autocorrelation of the transverse matter current,
the complete coefficient of $g^2 e^2$ has a positive spectral atom at
\[
 E_* =2m+\lambda_\ell-\lambda_0
\]
at every nonzero spatial momentum $P$. Here $\lambda_j$ is the
ground energy of the radial generator in angular sector $j$.
In particular this coefficient has support outside the forward
cone for $|P|>E_*$. The statement holds for the continuum spatial
coefficient per unit volume as well as for the allowed box momenta.
\end{proposition}
\begin{proof}
Take a real unit polarization $u$ perpendicular to $P$ and the
current obtained by differentiating the actual spatial Wilson
link at a background Fourier mode $P$. At noncoincident insertion
times the order in question consists of two first derivatives,
each of order $g e$. A second derivative of a single hop is
supported at coincident times. Number conservation in the onsite
vacuum gives zero one-point function. At this order the expectation
of the two first derivatives propagates with the unperturbed
Hamiltonian. Neither a Coulomb insertion, which carries two more
electric powers, nor another hopping vertex enters. This computes
the separated-current coefficient in the filled-vacuum determinant.

Write $h(r,k)=v\boldsymbol\alpha\cdot k\,e^{-\tau(r)|k|^2}$
for the spatial symbol before multiplication by $g$. Up to a
common sign and the angular character, the spin matrix of the
current, with $ge$ removed, is
\[
 V_u(r;k,P)=e^{-a(r)|P|}\!
       \int_0^1 (u\cdot\nabla_k)h(r,k+sP)\,\mathrm ds .
\]
Indeed the point on the straight Wilson link has parameter
$s\in[0,1]$, and its linear insertion multiplies the heat
displacement by its component along $u$. Let $\phi_j$ be the
real positive normalized radial ground function for
$-r^{-4}\partial_r(r^4\partial_r)+j^2r^{-8}$ with the prescribed
outer boundary condition. The current makes a particle and a
hole from the onsite vacuum. Each costs $m$ at zeroth hopping
order, independently of its spatial momentum. Project the
intermediate normal evolution onto
$e^{i\ell\theta}\phi_\ell$. Its transition matrix, between the
opposite $\beta$ slots, is the corresponding block of
\[
 \int_0^1 r^4\phi_\ell(r)\phi_0(r)V_u(r;k,P)\,\mathrm dr .
\]
At $k=0$, transversality eliminates the derivative of the
exponential. The resulting matrix is $v\boldsymbol\alpha\cdot u$
times the strictly positive average of
$e^{-a(r)|P|}\int_0^1 e^{-\tau(r)s^2|P|^2}\,\mathrm ds$.
Its off-diagonal block is nonzero. Continuity in $k$ supplies a
positive-measure set of relative momenta in the continuum. In the
box the $k=0$ term already suffices. Spectral decomposition of the
two-current product squares this matrix element and assigns it
the factor $e^{-tE_*}$ at positive carrier separation $t$.
Other normal eigenstates contribute nonnegative squared norms.

The continuum coefficient is finite at fixed $P$. By the radial
Cauchy--Schwarz inequality and the $s$ average, the spatial
integral of the squared current vertex in the initial state is
bounded by a constant times
$\int_0^1 r^4|\phi_0(r)|^2[1+\tau(r)^{-3/2}]\,\mathrm dr$:
a translated derivative of $k e^{-\tau |k|^2}$ has squared
spatial integral of order $\tau^{-3/2}$. This is integrable at
the calibrated spinor exponents. Thus the nonzero atom is an
ordinary positive spectral contribution to the separated-time
response. A local contact term cannot remove its exponential at
arbitrarily large time separation.
\end{proof}

The distinction can also be tested without expanding the hopping
at a fixed positive depth. In the continuum filled ground state at
zero dynamical electric charge, the two bands of
\eqref{eq:reference-spin-uv} have energies $\pm E_r(k)$, with
$E_r(k)\le E_{\max}=(m^2+v^2/(2\exp(1)\tau(r)))^{1/2}$.
Absorb the chosen nonzero hopping strength into $v$. A derivative
with respect to a unit external density probe carries the optical
factor $e^{-a(r)|P|}$. Its positive particle--hole spectral
measure is supported at $E_r(k+P)+E_r(k)\le2E_{\max}$.
It has nonzero mass at arbitrarily large $|P|$. In fact the spin
sum is $\operatorname{tr}(P_+(k+P)P_-(k))$, with
$P_\pm(k)=(1\pm h_r(k)/E_r(k))/2$. For fixed $k\ne0$ this tends
to $1-m/E_r(k)>0$ as $|P|\to\infty$, and is positive on a
neighborhood of such $k$. Both the fixed-depth signal and the
moving-depth atom can become small at large momentum. The
spectral support test is sensitive to any strictly positive
weight there.

The moving-depth result is a statement about a joint perturbative
coefficient. It obstructs imposing forward-cone support order by
order on that same source. Equivalently, an exactly covariant
family with fixed translation generators, whose Wightman measures
are differentiable in the couplings as distributions, has
derivatives that vanish on every test supported strictly outside
the cone. The calculation alone makes no corresponding assertion
about an isolated nonzero coupling of the moving theory. The
fixed-depth calculation uses the exact free charged sea at its
chosen hopping strength. In both cases a physical reconstruction
with a different charged current or translation generator must
derive its own determinant and its own UV bounds. The oriented
vectorial pole prescription on the chosen plane selects boundary
values of those coefficients; it cannot displace a positive
real-transfer spectral mass of the source just evaluated.

\paragraph{Can a chiral projection remove this spectral support?}
The preceding test extends to the bounded internal matrix hoppings
of Corollary~\ref{cor:seam-optical-shifts-uv}, including destructive
interference between their contributions. Suppose a candidate
physical quotient is represented by closed orthogonal projections
on the normal and finite spin slots that commute with the free
normal transfer. Assume the same projections act at every spatial
momentum in a fixed set of canonical slots. Use positive inner
products on the retained states. If any transverse pair-current
matrix element survives at one momentum in these slots, the quotient
still fails the order-$g^2 e^2$ forward-cone test.

To see this, fix a direction $n$, polarization $u\perp n$, and
a relative momentum $k_0$ for which that matrix element is nonzero
at total momentum $P=t_0n$, $t_0\ge0$. Apply the spectral projections
of the outgoing free normal Hamiltonian on bounded energy intervals.
They tend strongly to the identity, so some bounded energy interval
and some vector $w$ in its range give a nonzero matrix element
there. In the onsite theory the particle--hole excitation adds
$2m$ to the normal energy difference at all $P$. Thus this
fixed outgoing interval lies below one common finite energy $C$.
Keeping the same spin and normal vectors, denote the matrix
element at $P=t n$ by $M(t)$. For $t\ge0$ the Poisson readout is
$e^{-a(r)t}$. Along this ray it extends to an entire function of
the one complex variable $t$. The differentiated Gaussian
$h(r,k_0+stn)$ extends using its polynomial momentum square,
and its first spatial derivatives are bounded uniformly in $r,s$
on each compact subset of the complex $t$ plane, since $a$ and
$\tau$ are bounded above on the collar. The finite sum of bounded
internal multipliers and the fixed normal projections therefore
give an entire $M$ by dominated integration. When $M(t_0)\ne0$,
the identity theorem precludes $M(t)=0$ for all $t>C$. There
are arbitrarily large real $t$ at which the projection of the
current into this same energy interval has strictly positive
squared norm. Continuity in the relative momentum supplies a
set of nonzero spatial measure there. This argument uses the
continuum momentum ray; in a finite box the explicit
positive matrix element in Proposition~\ref{prop:omega-current-cone-uv}
suffices for the scalar angular character.

Consequently a momentum-independent projection preserving a
nonzero low-energy transverse pair current cannot repair this
coefficient by cancellations among internal slots or angular
histories. Changing the free carrier dispersion, the momentum
dependence of the physical projection, or the physical current
changes a hypothesis of the argument. Each such change affects
the common source whose absolute graph estimates are required.

\paragraph{A spacelike locality test for the proposed physical observable map.}
The neutral, gauge-invariant occupation density and the electric
density measured through the normal Maxwell extension have different
localization properties. Consider the former with a smooth real test
$f$ on the spatial slice. In the physical particle--hole representation
write $Q(f)$ for the particle number weighted by $f$ minus the hole
number weighted by $f$. It acts diagonally and annihilates the onsite
vacuum. Its exponential $U_f=\exp[iQ(f)]$ is bounded and neutral.
These definitions use the canonical positions before harmonic
readout. If $Q(f)$ is assigned to the local observable algebra of
$\operatorname{supp} f$, as a pointlike neutral density would be, the
following calculation tests relativistic locality directly.

\begin{proposition}[Instantaneous spatial hopping and the local density test]
\label{prop:omega-locality-uv}
Suppose the reduced Hamiltonian has a nonzero instantaneous
particle--hole hop with kernel between two separated spatial open
sets. Assume the other terms commute with the occupation-density
phases. In particular one may take the optical Dirac heat hop at
zero electromagnetic and scalar coupling, retaining nonzero
hopping strength. If all $U_f$ belong to local relativistic
algebras with the indicated supports on the sector plane, spacelike
commutativity fails.
\end{proposition}
\begin{proof}
Take smooth real functions $f,h$ of disjoint compact supports at
positive spatial separation. The required real-time commutator
$[e^{itH}U_f e^{-itH},U_h]$ must vanish for an open interval about
$t=0$ by spacelike separation. Test its derivative at zero between
an onsite vacuum vector and a particle--hole vector, both with
smooth normal wave functions supported in a compact interior collar.
Use smooth spatial wave packets, the free photon vacuum and, for
the direct test, zero electric coupling. These are Hamiltonian
domain vectors: the differential terms act on smooth data and the
spatial hopping is Hilbert--Schmidt uniformly on the collar. The
bounded phases $U_f,U_h$ and their adjoints preserve this test
domain on the matrix elements in question. Differentiation in the
weak sense therefore gives $i[[H,U_f],U_h]$.

On a pair at $x,y$, $U_f$ acts by
$\exp[i(f(x)-f(y))]$. Since the initial vacuum has occupation
zero, the vacuum-to-pair kernel of this double commutator, apart
from the overall $i$, is
\[
 (1-e^{i(f(x)-f(y))})(1-e^{i(h(x)-h(y))})
                   K_r^{+-}(x,y) .
\]
Here $K_r^{+-}$ is the actual occupied-to-empty block of the
creation kernel. Both factors can be made nonzero on the separated
open sets. For the optical spatial derivative of a Gaussian there
are separated points where this block is nonzero for some spin
components: the spatial kernel is proportional to
$\boldsymbol\alpha\cdot\nabla G_{\tau(r)}$ with nonzero
off-diagonal mass slots. The same holds on the torus away from
the stationary points of the periodized Gaussian. Choose smaller
neighborhoods and normal support on which a matrix element remains
nonzero, and choose the final test vector to detect it. This gives
a nonzero derivative, contradicting the required interval of
vanishing. The argument takes place inside a spatial coordinate
patch and at arbitrarily short times.
\end{proof}

More generally, if all such canonical density phases are strictly
local observables, testing every pair of separated neighborhoods
forces the instantaneous charged creation kernel to be supported
on the spatial diagonal in this class of Hamiltonians. This
necessary condition is independent of an ultraviolet bound. The
electric density obtained with $e^{-a(r)|p|}$ is spatially
extended; a compact test on its readout does not construct the
pointlike operator used here. A recovery of historical Q1 locality
on a different set of physical observables must specify its
localization and prove the spacelike commutators for that algebra.

\subsection{Relativistic charged propagation on the locked plane}
\label{sec:carrier-dispersive-uv}

\paragraph{What an unrestricted action must pass.}
We retain unrestricted physical finiteness as a selection criterion for
PE-CGG-HOUQFT. A successful choice must determine its translation
generators, observable current, Gauss constraint, vacuum, and sewing
from one raw source. On a locked plane these data are assigned before
transport to a Minkowski readout. A positive proper depth may control
relative momentum in a graph even if it gives the wrong energy at a
fixed total momentum. Conversely a current with the correct invariant
pair spectrum can leave a divergent internal photon integral. The
following tests distinguish these outcomes without assuming the
criterion as an axiom.

\paragraph{A dispersive charged carrier and its normal Wilson path.}
Consider an alternative on a single locked plane. In the physical
particle and hole representation give both carriers the dispersion
$E(k)=\sqrt{m^2+|k|^2}$ at zero hopping and electric charge,
where $m>0$ and the carrier speed is one. Keep the marked optical
pair-creation heat hop and its adjoint, with constant strength $g$,
at lapse $\tau=a^2$, $a=r^{1/3}$ in the reference units. An angular
factor $e^{i\ell\theta}$ is allowed on creation. We specify the
electromagnetic coupling of the new diagonal kinetic operator as
well. On the cover of the spatial torus its jump density is
\[
 j_m(z)=\frac1{2\sqrt\pi}\int_0^\infty
 \frac{e^{-m^2h-|z|^2/(4h)}}{h^{3/2}(4\pi h)^{3/2}}\,\mathrm dh .
\]
Integration in $z$ gives
$\int(1-\cos k\cdot z)j_m(z)\,\mathrm dz=E(k)-m$ by the
Gaussian characteristic function and the elementary square-root
Laplace integral. The density is $O(|z|^{-4})$ near zero and
exponentially decreasing at infinity. Use all lifts of a periodic
spatial wave function. At fixed normal position the kinetic
form for charge sign $q=\pm1$ is
\[
 m\|u\|^2+\frac12\int\mathrm d^3x\,\mathrm d^3z\,
 j_m(z)|u(x)-U_{q,r}(x,x+z)u(x+z)|^2 .
\]
For this diagonal transport choose a path that goes from optical
height $a$ to $a+|z|$, traverses the straight chord $z$ there,
and returns to $a$ over the new endpoint. $U$ is its Wilson
parallel transport with the appropriate orientation. In the
transverse slice the readout on the horizontal segment is
$e^{-(a+|z|)|p|}$. The vertical segments carry the normal
connection of Proposition~\ref{prop:constitutive-gauss-uv}.
Consequently a pure gauge contributes precisely its two endpoint
values at height $a$, even though the horizontal segment is
measured farther out. The temporal charge density, its normal
derivatives, the full Gauss energy and the electrically dressed
endpoints then use the original $P_a$. This prescription changes
the action before differentiation. It introduces no tunable
momentum scale. The constrained harmonic collar is continued to
the largest optical height sampled by a lifted jump.

\begin{proposition}[Moving-depth bounds with relativistic free flights]
\label{prop:carrier-flight-uv}
On the finite slab and torus of Theorem~\ref{thm:omega-quotient-uv}
with the spatial harmonic photon held fixed at zero, the complete
vacuum, electromagnetic and tested charged coefficients at each
finite joint order in $g,e$ have an absolute bound in the normal
law, charged flight histories, hopping times, and photon
four-momenta. One can also retain one spatial trace momentum on
each closed charged word. Smooth spatial boundary probes and the
same square-integrable normal preparations are allowed. The result
uses the ordered transfer of the new action and the actual onsite
particle--hole vacuum. It makes no assertion of absolute summability
in $g$ or of a joint Fourier estimate for the charged time variables.
\end{proposition}
\begin{proof}
Between heat hops, the free positive or negative frequency slot
propagates in its oriented carrier time by $e^{-tE}$, $t>0$.
After the factor $e^{-mt}$ this is the symmetric jump process
with intensity $j_m(z)\,\mathrm dz\,\mathrm dt$. The magnetic
kinetic form inserts the indicated parallel transport on its jumps.
For smooth transverse fields this can first be seen for
$|z|>\epsilon$ by expanding the finite jump generator. Symmetric
small-jump exhaustion gives the original closed kinetic form.
Indeed the magnetic forms are increasing sums of squared
differences on inclusion of the missing jumps. The following
estimate also constructs their Wilson currents in the Gaussian
Maxwell field and controls the marked limits.

Fix the surviving normal path. In the transverse Maxwell covariance
space a jump at $t$ from $x$ contributes the Fourier current
\[
 F_{t,x,z}(\nu,p)=e^{i\nu t+ip\cdot x}
 e^{-(a_t+|z|)|p|}\,P^\perp_p z
       \int_0^1 e^{iu p\cdot z}\,\mathrm du .
\]
The Hilbert norm uses the full positive covariance
$(\nu^2+|p|^2)^{-1}$ with its standard box and Fourier measure.
Frequency integration and spatial shell counting give
\[
 \|F_{t,x,z}\|\le C\left(|z|+
             \frac{|z|}{a_t+|z|}\right).
\]
On every spherical annulus $\int F_{t,x,z}j_m(z)\,\mathrm dz=0$:
before transverse projection its vector integral in each mode is
parallel to $p$ by rotation symmetry. The integral here is in the
covering displacement $z$ even on the torus. The sum of the
currents along one flight is thus a compensated Hilbert-space
jump martingale. For every integer $n\ge1$,
\[
 \int \|F_{t,x,z}\|^{2n}j_m(z)\,\mathrm dz
       \le C_n(1+a_t^{-1}).
\]
In the unit ball split at $|z|=a_t$ and use the radial density
bound $C|z|^{-2}\,\mathrm d|z|$; outside use the exponential
tail. It follows that all its terminal $2n$th moments are
bounded by polynomials in $1+\int_0^T a_t^{-1}\,\mathrm dt$.
For example, apply the jump identity for $\|M\|^{2n}$ to the
compensated sum. Its linear term has zero conditional mean. Its
remainder is bounded by $C_n(\|M\|^{2n-2}\|F\|^2+\|F\|^{2n})$.
Induction gives the polynomial bound uniformly in the small-jump
exhaustion. The omitted currents vanish in mean square, and the
higher uniform moments give convergence at every lower finite order.
These bounds include the diagonal contacts of the new kinetic
term: its arbitrarily many marks on one jump are present when
one differentiates the complete exponential. Temporarily discarding
small jumps defines the stochastic integral; no jump cutoff enters
the coefficients. Absolute photon-momentum integration applies to the
completed compensated current. The small jumps of the individual
uncompensated vector current need not be absolutely summable.

At fixed hopping order the canonical determinant is a finite sum
of closed particle--hole words and open chains, with their
fermionic signs. The empty vacuum of the diagonal number-preserving
Hamiltonian has unit matter vacuum factor, including in an external
field. Thus every nonempty closed word has at least one charged
heat hop. Cut one such hop on each closed word. The remaining
segments can be sampled in cyclic order as independent free
increments started at the preceding endpoint, with the appropriate
time orientation. Reversal of a symmetric flight has the same
increment law when run from its cut end. A non-anchor heat hop of
derivative order $b=0$ or $1$ is bounded in absolute value by
$C\tau^{-b/2}$ times a normalized Gaussian of larger lapse.
Its transverse current has all fixed norm moments bounded under
this Gaussian, since $|z|/a$ has bounded moments at $\tau=a^2$.
A Coulomb end is the time transform of the measured occupation
on a flight segment. At each spatial momentum its phase has
modulus one before multiplication by $e^{-a_t|p|}$.
Plancherel and Cauchy--Schwarz for two ends, followed by the
$|p|^{-2}$ spatial sum, bound each complete Coulomb line by
$C\int_0^T(1+a_t^{-1})\,\mathrm dt$. This applies to two
charges on different segments and includes their self-energies.

For clarity one may Fourier invert the remaining cut heat kernel.
Each set of $h$ transverse marks on it differentiates its
translated symbol $h$ times and multiplies it by the measured
photon factors. The spatial trace integral or sum is bounded
uniformly in the translates by
$C_h(1+\tau^{-3/2})\tau^{-b/2}$, with a remaining
$\sqrt{\tau}\,e^{-a_t|p|}$ for each mark. The same bound follows
in position space by bounding the periodized Gaussian at the
unpaired endpoint. With $a_{\min}=\min_t a_t$, each remaining
mark has envelope $C(1+|p|)^{-1} e^{-a_{\min}|p|/2}$.
Its squared transverse covariance norm costs at most
$C(1+|\log a_{\min}|)$. After the spatial trace bound, products
of photon contractions converge absolutely by the Hilbert-space
Cauchy--Schwarz inequality and the flight moment estimate above.
This also treats coincident marks before the normal average.

The anchor costs at worst $1+r^{-4/3}$ and any other differentiated
heat hop at worst $1+r^{-1/3}$. Their distinct times are integrated.
All positive powers of these additive costs have finite matrix
elements by the radial Hardy estimate, and all fixed powers of the
minimum-depth logarithm are integrable by the hitting bound in
Theorem~\ref{thm:seam-calibrated-joint-uv}. Angular multipliers
have modulus one on the conditional normal path. Finitely many
endpoint and vacuum subtractions preserve the estimates. This
proves the claimed finite-order bounds using the same determinants,
charges and Wilson products that define the transfer.
\end{proof}

\begin{proposition}[Exact moving-depth pair threshold at zero electric coupling]
\label{prop:carrier-blocked-uv}
Use the same action at $e=0$ and the standard isotropic spinor heat
pair hop with nonzero real strength $g$, or its scalar version.
Take the trivial bundle and zero flat background. The coupled
charged and normal Hamiltonian has a neutral ground state with
zero spatial momentum. In its neutral momentum-$P$ fiber, $P\ne0$,
the energy above the neutral vacuum is at least
$\sqrt{|P|^2+4m^2}$. Thus the transverse matter-current spectral
measure on that vacuum obeys this lower bound, including normal
motion and angular transitions, at arbitrary finite $g$. Tensoring
with free transverse photons gives the forward-cone bound $E\ge|P|$
for the combined neutral spectrum after subtracting the vacuum energy.
\end{proposition}
\begin{proof}
At momentum $k$ the pair operator only couples a particle of
momentum $k$ to a hole of momentum $-k$. The off-diagonal spin
block is a fixed matrix at that $k$ multiplied by a scalar
function of $(r,\theta)$. A singular-vector basis independent of
the normal position therefore decomposes it into pairs of slots.
In each pair the empty and doubly occupied states mix. Either
singly occupied state is blocked by the Pauli principle. The
kinetic energy of a blocked slot is exactly $E(k)$, independent
of the geometry. Each unblocked slot contributes zero net
momentum and charge, in both its basis states. We use the full periodic angular Hilbert space and the
neutral vacuum at the bottom of its unrestricted neutral fiber. If an
additional superselection rule excludes the comparison empty-slot
vector, this threshold argument needs a new variational comparison.

Decompose a neutral state into its mutually invariant blocked
configurations. For a normalized vector in one such configuration
replace every blocked slot by its empty state. This is an isometry
on all remaining slots and on the entire normal wave function.
The normal kinetic form and all unblocked pair expectations are
unchanged. On a newly emptied slot the pair term has zero diagonal
expectation. The energy expectation thus decreases by exactly
$\sum_{i\in B}E(k_i)$, where $B$ is the blocked set. The resulting
vector lies in the sector with all slots unblocked. Its energy is
at least the ground infimum there. Finite-mode occupation tests
give the same inequality on the closed form by density. Compactness
of the free resolvent on the finite torus and bounded collar,
and the infinitesimal form bound of the heat pair hop relative
to the free kinetic, mass and normal forms, give existence of a
ground vector. The argument shows it can be taken in the all-unblocked
sector. For $P\ne0$ neutrality requires at least two blocked charges.
The triangle inequality in four Euclidean components applied to
$(k_i^{\rm phys},m)$ gives
$\sum E(k_i)\ge\sqrt{|P|^2+(m|B|)^2}$, where each physical
momentum is that of the particle or the hole in its slot. This
bound is independent of the sign of $g$ and the rate of the
normal angular clock. A current obtained by differentiation in an
external field takes the neutral vacuum into the corresponding
neutral momentum fiber. Free photons add their nonnegative
energies and their momenta obey the triangle inequality.
\end{proof}

The diagonal dispersion and its normal Wilson transport remove the
finite-energy, arbitrarily-large-momentum pair of
Proposition~\ref{prop:omega-current-cone-uv} in this neutral test.
They retain an equal-carrier-time spatial readout. The bound on
spectral support does not give boost covariance of the current,
spacelike commutativity of its assigned observables, or a
forward-cone theorem after dynamical photons couple to matter.

\paragraph{Summing the heat pair hopping before expanding the electric source.}
The charged flight provides a uniform estimate for each separate
photon mark. This also restores the Fourier variables for charged
carrier time and removes the small-hopping expansion.

\begin{theorem}[All electric orders for the dispersive charged action]
\label{thm:carrier-flight-dressed-uv}
Under the finite-slab, torus and preparation hypotheses of
Proposition~\ref{prop:carrier-flight-uv}, the same action has
absolutely UV-finite dressed graph coefficients at every finite
electric order and every finite real pair-hopping strength.
The bound includes charged carrier space and time modes, independent
photon four-momenta and the moving normal path. The real-coupling
diagonal vacuum transfer is finite and strictly positive.
\end{theorem}
\begin{proof}
Write $k=|p|$. Fix the complete normal path and let $I$ be a
deterministic carrier interval of a charged flight. There are
uniform single-mark estimates before conditioning a flight on its
final spatial position. For a transverse component with carrier
frequency $\nu$, let $J(\nu,p)$ be the sum over the actual jumps
in $I$. For every finite $h\ge2$,
\[
 \|J(\nu,p)\|_{L^h({\rm flight})}
 \le C_h(1+B_T)^{c_h}
       \frac{e^{-a_{\min}k/2}}{1+k},
\]
where $B_T$ can be taken to be $C\int_0^T(1+r_t^{-5/3})\,\mathrm dt$.
Indeed, for every integer $l\ge2$ the $l$th power of the jump
size after multiplying by $(1+k)e^{a_{\min}k/2}$ integrates
against its intensity to at most
\[
 \begin{aligned}
 &C_l(1+k)^l e^{-la_tk/2}
       \int |z|^l e^{-l|z|k}j_m(z)\,\mathrm dz\\
 &\hspace{1em}\le C_l(1+k)e^{-la_tk/2}
                   \le C'_l(1+a_t^{-1}).
 \end{aligned}
\]
For $k\ge1$ use the radial integral of
$s^{l-2} e^{-lsk}$; at small $k$ use the massive tail. The
conditional compensator still vanishes by transverse projection.
The jump moment induction of Proposition~\ref{prop:carrier-flight-uv}
now proves the bound at each Fourier mode. Signs or reversal of
the deterministic carrier interval leave it intact.

For a temporal photon define the deterministic maximal primitive
\[
 M_p(\nu)=\sup_{I'\subset[0,T]}
     \left|\int_{I'} e^{i\nu u-a_u k}\,\mathrm du\right|.
\]
For a flight starting at $x$ write $Z_u=e^{ip\cdot X_u}$.
Its jump decomposition has drift $-(E(p)-m)Z_u$ and a
martingale whose conditional jump intensities obey
$\int|\Delta Z|^l j_m\,\mathrm dz\le C_l(E(p)-m)$ for $l\ge2$;
this follows from $|e^{ip\cdot z}-1|\le2$ and the defining
characteristic exponent. Integrate $Z$ against the deterministic
primitive of the temporal mark by parts on $I$. The terminal
term is at most $M_p$, the drift term at most
$T(E(p)-m)M_p$, and the stochastic integral has integrand
bounded by $M_p$ times the martingale jump. Its $L^h$ norm is
at most $C_{h,T}(1+k)M_p$ by the same moment induction (or the
bounded-jump moment inequality). Therefore
\[
 \left\|\int_I e^{i\nu t-a_t k} Z_t\,\mathrm dt\right\|_{L^h}
             \le C_{h,T}(1+k) M_p(\nu).
\]
This estimate allows random spatial positions without differentiating
the radial path. Products of a fixed number of marks use H\"older
with a higher finite $h$ on this same flight. Two temporal ends
now cost an additional $(1+k)^2$. The maximal Fourier estimate
$\int M_p^2\,\mathrm d\nu\le C\int_0^T e^{-2a_t k}\,\mathrm dt$
from Theorem~\ref{thm:seam-calibrated-joint-uv} still yields
a summable Coulomb line:
\[
 C\sum_{p\ne0}\frac{(1+k)^2}{k^2}\int_0^T e^{-2a_t k}\,\mathrm dt
     \le C_V\int_0^T(1+a_t^{-3})\,\mathrm dt .
\]
Here $a_t^{-3}=r_t^{-1}$. The two transverse ends cost the same
finite minimum-depth logarithm as in the optical estimate.

For the charged determinant put $R_d(e)$ equal to the diagonal
vacuum inverse in an external electric field, including the two
one-sided oriented flights, and let $K(e)$ be the off-diagonal
heat pair hop. Use the canonical number-preserving diagonal
second quantization in each particle and hole sector. Its empty
vacuum has amplitude one for every external field. The ordinary
pair determinant relative to it is $\det_2(1+X(e))$,
$X(e)=R_d(e)K(e)$: on a compact interior collar the ordered vacuum
expansion gives exactly the cycles of these propagators, and its
one-hop slot trace vanishes. The same statement follows by the
two-slot Schur identity. This choice of reference is part of the
new Hamiltonian. Identifying it with a local Dirac sea would require
an additional equivalence of the two sourced vacua.

Denote by $X_I$ the derivative at zero field in a given set of
photon modes and polarizations, with temporal sources supported
on the given slab, and let $X=X(0)$. A flight crossing an outer
end of the slab sees these marks only on its interval of overlap. Its time kernel
is a sum over distributions of marks between a single diagonal
flight and its terminal heat hop. For fixed times $t,s$ start the
flight at $x$ and run it to its other time, where it meets the
heat kernel with free endpoint $y$. Use Jensen's inequality for
the squared kernel and integrate in $y$ first. Conditional on the
flight, every transverse heat-hop mark supplies $|z|e^{-a_s k}$
at the terminal displacement $z$. After the squared heat kernel
and its periodized lifts are integrated, this supplies a factor
$C\sqrt{\tau_s} e^{-a_s k}$ per mark and a common squared cost
$C(1+\tau_s^{-5/2})$ for the spin derivative hop. The endpoint
phase factors have modulus one. For the flight marks the two
preceding $L^h$ estimates are uniform in $x$. The remaining time
envelope is $e^{-m|t-s|}{\bf1}_{s\in[0,T]}$ with the appropriate
one-sided slot. Consequently $\|X_I\|_{\rm HS}$ on carrier
space-time is bounded by a polynomial in $1+B_T$ times the
product of the individual transverse and temporal envelopes above.
This estimate uses one realization of each marked flight and includes
all multiple insertions on a jump. For $I=\varnothing$ it gives
$\|X\|_{\rm HS}^2\le C B_T$. The same last bound holds with the
full real Wilson factors by their unit modulus and the conditional
heat-kernel majorant. For a real instantaneous Coulomb auxiliary
field use its time-integrated phases, defined by the positive charge
covariance on the flight, or first test with smooth temporal potentials.

At zero electric field the dressed inverse $R$ obeys $\|R\|\le1/m$.
Indeed the real part of $\beta(\partial_t+\beta E+K)$ is
$E\ge m$ while the off-diagonal $\beta K$ is skew adjoint.
At each fixed surviving path $K$ is bounded and the free Dirac
transfer has its ordinary time and space graph domain. The same
bound on the adjoint and the closed-range criterion give the inverse.
For $A=(1+X)^{-1}$ the resolvent identity gives
$A=1-X+R K R_d(0)K$. The free one-sided time kernels have
operator norm bounded by $e^{-m|t-s|}$, so
$\|K R_d(0)K\|_{\rm HS}$ is bounded by the product of
$(\int\|K_t\|_{\rm op}^2\,\mathrm dt)^{1/2}$ and
$(\int\|K_t\|_{\rm HS}^2\,\mathrm dt)^{1/2}$. Thus
$\|A\|\le C(1+B_T)$. Every mixed derivative of the conditional
matter logarithm is a finite sum of traces of products of $AX_I$.
With one differentiated factor the regularized determinant gives
$\operatorname{Tr}[(A-1)X_I]=-\operatorname{Tr}(AXX_I)$, again
with two Hilbert--Schmidt factors. The absolute cyclic
Hilbert--Schmidt kernel bound of Proposition~\ref{prop:constitutive-magnetic-uv}
keeps all spatial and temporal carrier modes in these traces.
It costs only their product of Hilbert--Schmidt norms. Wick-pair
the individual photon envelopes and integrate their independent
frequencies. The preceding Coulomb bound and the transverse
logarithm apply before the radial expectation.

For open charged lines the free flight without a heat hop is
bounded by the same photon moments when tested at its endpoints.
For terms with heat hops use
$R(e)=(1+X(e))^{-1}R_d(e)$. A terminal free flight against an
endpoint time and bounded smooth spatial probe gives an $L^2$
space-time vector by its exponential time envelope and the flight
mark bounds. At the measured initial time, a first marked $X_I$
or an unmarked $X$ followed by $A$ has the preceding spatial
Hilbert--Schmidt bound integrated only over $s$, hence defines
the other $L^2$ vector. This bounds the absolute open chains.
In the determinant prefactor
$|\det_2(1+X(e))|\le\exp(\|X(e)\|_{\rm HS}^2/2)$.
The radial Hardy and minimum-depth estimates now control all
these bounds for the stated normal endpoints. The Hamiltonian
uses the positive covariant kinetic form in each occupied slot,
the ordinary Maxwell form, the complete positive Gauss energy
and the pair heat hop of infinitesimal relative form bound.
They define a dense semibounded closed form, for example by
adding the positive jump forms to the free field and normal forms
before adding the pair term. On compact interior collars their
transfer is the ordered charged flight representation, by
symmetric small-jump approximation and time slicing. The
bounds above pass to the full collar. The full Hamiltonian
has strictly positive diagonal vacuum transfer on each nonzero
prepared vector. Its use as the common vacuum denominator proves
the claimed connected estimates.
\end{proof}

The integral in the diagonal kinetic term makes jumps at all
spatial lengths. Retaining their complete normal path gives each
jump of length $|z|$ a further depth $|z|$ at which it samples the
transverse photon. The resulting marked flight estimates and the
Pauli-blocked threshold in Proposition~\ref{prop:carrier-blocked-uv}
concern this same action.

\paragraph{A nonzero current inside the cone and a residual boost test.}
First freeze a positive depth and trivial angular character in the
dispersive action, with dynamical electric charge zero. Differentiate
with respect to a unit external transverse probe. In the continuum
take a momentum transfer $P\ne0$, $p=P/2$ and a real unit polarization
$u\perp P$. At the one-particle momenta $-p\to p$ the diagonal
L\'evy vertex in this polarization vanishes. Its displacement integral
contains $u\cdot z$ multiplied, apart from the radial weights, by
a function only of $P\cdot z$; the transverse angular integral
is zero. This is an evaluation of the actual diagonal Wilson term.
The heat-pair vertex at these momenta is
\[
 gv e^{-a|P|}(\boldsymbol\alpha\cdot u)
       \int_0^1 e^{-\tau(s-1/2)^2|P|^2}\,\mathrm ds .
\]
In the standard isotropic spin representation the zero-field symbol
is $h(p)=\beta E(p)+gv\boldsymbol\alpha\cdot p\,e^{-\tau|p|^2}$.
Let $b=gv|p|e^{-\tau|p|^2}$ and $D=(E(p)^2+b^2)^{1/2}$.
Between the exact band projectors $P_\pm=(1\pm h/D)/2$ the spin sum
of $\boldsymbol\alpha\cdot u$ is
\[
 {\rm tr}\!\left[P_+(p)(\boldsymbol\alpha\cdot u)
                    P_-(-p)(\boldsymbol\alpha\cdot u)\right]
        =\frac{2E(p)^2}{E(p)^2+b^2}.
\]
The equality follows by anticommuting $\boldsymbol\alpha\cdot u$
with $\beta$ and $\boldsymbol\alpha\cdot p$. For every nonzero $gv$
the current therefore has nonzero pair weight, with continuity in
the relative momentum giving a set of positive measure. Its
spectrum in this frozen action is the sum of the two positive band
energies and satisfies the preceding cone bound.

Now let the normal state move and give a pair the nonzero angular
shift $\ell$. For the leading joint $g^2 e^2$ transverse current
measure use the stationary vacuum about $g=e=0$, counting $e$ on
each probe and keeping the electromagnetic dynamics free at this
order. The unperturbed diagonal current annihilates the empty
vacuum. The first nonzero current acting on the vacuum is of
order $ge$. Its pair component consists of the differentiated
pair hop and the diagonal current acting on the first-order vacuum
correction. Both create a particle and hole in angular sector
$\ell$. At this order there are no other intermediate matter
sectors in the separated two-current measure. A contact at one
carrier time cannot supply separated-time spectral weight.
The normal excitation costs at least
$\delta=\lambda_\ell-\lambda_0>0$, since the angular potential
is $\ell^2/r^8$ on the unit collar. Thus this measure is supported
on
\[
 E\ge\sqrt{|P|^2+4m^2}+\delta .
\]
This is a statement about the spatial continuum coefficient of the
two-current density (or allowed discrete momenta before removing
the spatial volume), with the free normal ground preparation.
The first-order vacuum correction is understood at the coefficient
level: its single-pair amplitude is the free positive pair
resolvent on that source, equivalently its decaying ordered time
integral. It requires no interacting infinite-volume vacuum.

The displayed support contains a signal at its lower boundary in
the continuum sense of support. Fix $P_0\ne0$ and again take the
symmetric momenta and transverse polarization. The diagonal
vertex vanishes there for both the particle and the hole regardless
of the depth, so the first-order vacuum correction contributes
zero at those relative momenta. The explicit pair vertex has strictly
positive overlap in magnitude with the positive radial ground vector
in sector $\ell$, by its displayed integral and positivity of the
initial radial ground vector. Its off-diagonal spin block is nonzero.
Both contributions are continuous for nearby relative momenta.
For the vacuum correction this follows from the pair gap $2m$,
the smooth heat symbol on compact momentum sets and the bounded
radial readout. The diagonal vertex is continuous there by
symmetric small-jump cancellation. Consequently each neighborhood
of this boundary point has positive spectral mass. For sufficiently
large $|P|$, every bounded invariant mass interval $E^2-|P|^2$
lies below the stated support, because the squared lower edge minus
$|P|^2$ grows as $2\delta\sqrt{|P|^2+4m^2}$. A transverse boost
test can keep the same polarization by boosting along the ray of
$P_0$. The perturbative spectral support is therefore inside the
forward cone but fails boost invariance. The normal gap in this
source adds to carrier energy in its selected frame. An internal
excitation entered as part of an invariant mass would obey a
different threshold law.

\paragraph{Independent visits and a normal excitation carried as invariant mass.}
A second dispersive action addresses the common-radius volume
problem and the unboosted angular threshold together. On the
quartic collar let $N=H_{\rm n}-\lambda_0\ge0$ be its dimensionless
Friedrichs Laplacian shifted by its lowest eigenvalue. Its eigenstates
are internal one-particle states. In reference units assign them mass
$m_j=m+\lambda_j(N)$; in dimensional units multiply $N$ by the
inverse branch length. Both particle and hole slots have free energies
$E_j(k)=\sqrt{m_j^2+|k|^2}$. All internal levels are retained. The
operator $N$ now enters a mass in the energy dispersion. There is no
global fluctuating radial energy shared by all spatial modes.

At each interaction use the normalized collar volume
$5r^4\,\mathrm dr$ on $0<r<1$, independently at different visits
and with all marks of a single visit evaluated together.
Put $a=r^{1/3}$, $\tau=a^2$, and $d=a^{1/2}$ in reference units.
On a pair hop of spin derivative type multiply its spatial kernel by
\[
 D_d U_\ell D_d,\qquad D_d=e^{-dN},\qquad
                   U_\ell u=e^{i\ell\theta}u .
\]
These are two actual normal heat transfers with a character between
them. The hop in the action is the integral of this marked kernel
over the stated collar volume. Its adjoint is the reverse hop.
On a diagonal flight at level $j$ use the L\'evy form above with
mass $m_j$. Average its squared covariant difference over the same
single-visit law, with transverse segment at height $a+|z|$.
Thus its stochastic transfer assigns an independent mark $r$ to
each jump. A gauge transformation at a physical charged endpoint
uses the represented readout
\[
 \overline P=\int_0^1 5r^4 P_{r^{1/3}}\,\mathrm dr .
\]
This also measures the temporal potential. Explicitly, in
$A=A^\perp+\nabla\chi$, a link from $x$ to $y$ carries in addition
to the stated transverse integral the longitudinal difference
$\overline P\chi(y)-\overline P\chi(x)$, with the orientation
sign appropriate to its transport. This horizontal connection
specifies an equivariant link even when its sampled height differs
from the physical endpoint readout. Abelian Gauss reduction gives
the positive Coulomb square with $\overline P$ on both ends.
The bare Hamiltonian is the sum of these positive diagonal kinetic
forms, Maxwell and Gauss energies, and the pair hop. The visit law,
endpoint connection and internal transfer are constitutive choices
of this second action.

\begin{proposition}[Internal mass tower with independent optical visits]
\label{prop:carrier-mass-visits-uv}
On a finite slab and spatial torus, with the harmonic mode held
fixed, the normalized vacuum, photon and smooth elementary charged
coefficients of this action obey the complete all-electric-orders
bound at arbitrary finite pair hopping. The bound includes the
infinite internal state sums. At zero electric charge its finite-slab
vacuum logarithm per unit spatial volume has a finite thermodynamic
limit. At zero pair hopping a particle and hole in mass channels
$i,j$ have continuum threshold $\sqrt{|P|^2+(m_i+m_j)^2}$.
At any fixed pair hopping the stationary zero-electric-charge
matter response in a neutral nonzero momentum fiber still obeys
$E\ge\sqrt{|P|^2+4m^2}$ above its ground energy.
\end{proposition}
\begin{proof}
The unitary radial transformation by $r^2$ yields in angular
mode $j$ the one-dimensional operator
$-\partial_r^2+2/r^2+j^2/r^8$. Form comparison with the
Dirichlet interval Laplacian plus $j^2$, followed by summation
of the two Gaussian series, gives
$\operatorname{Tr}e^{-2dN}\le C(1+d^{-1})$ for $0<d\le1$.
In particular $\|D_dU_\ell D_d\|_{\rm HS}\le C(1+d^{-1/2})$.
The squared spatial Gaussian estimate and this internal norm
are integrable against the visit law. They give a spatial and
internal Hilbert--Schmidt pair kernel, as well as the usual Schur
bound. The mass gap controls the pair form as before. The
canonical empty vacuum has amplitude one for the number-preserving
diagonal Hamiltonian even when all normal levels are included.

Here are Fourier estimates uniform in the mass index. Fix electric
order $M$ and use the actual measure $15a^{14}\,\mathrm da$.
Its mean of $e^{-ak}$ is bounded by $C(1+k)^{-15}$ for spatial
frequency length $k$. The jump density satisfies $j_{m_j}\le j_m$.
For a transverse marked flight choose a fixed integer
$h\ge\max(2,2M)$ and then $\epsilon>0$ with $h\epsilon<1$.
After multiplying a Fourier mark by $(1+k)^{1+\epsilon}$, the
integrated $l$th jump moment for $2\le l\le h$ is at most
\[
 C_l(1+k)^{1+l\epsilon}(1+k)^{-15}.
\]
Indeed the jump length gives the integral of
$|z|^l e^{-lk|z|}j_{m_j}(z)$ and the independent visit gives
the mean of $e^{-lak}$. Transversality cancels its compensator.
The per-mode moment proof of Theorem~\ref{thm:carrier-flight-dressed-uv}
therefore supplies $(1+k)^{-1-\epsilon}$ for each transverse
flight end. At a heat hop with $b$ marked ends, extracting these
same envelopes costs at most $C_M a^{-b\epsilon}$ after allowing
one power of the displacement per mark. Its integrated spatial
and internal Hilbert--Schmidt estimate is bounded by
\[
 C_{M,V}\int_0^1 a^{14}a^{-5/2}(1+a^{-1/4})
                            a^{-b\epsilon}\,\mathrm da<\infty .
\]
This holds at coincident marks on one visit. A complete transverse
line has a convergent hard spatial sum after its frequency integral.
For a temporal mark the semimartingale drift
$E_j(p)-m_j\le |p|$ is uniform in $j$. The endpoint envelope is
$C_{M,T}(1+k)^{-14}\min(T,2/|\nu|)$ by the same interval
integration by parts, now using the deterministic $\overline P$.
Its Coulomb sewings converge as well. In the estimate of a marked
$R_d K$ sum the squared heat kernel over both internal labels
after integrating the final spatial position. The incoming flight
bound is uniform in its fixed mass label. Minkowski's inequality
for the heat visit and the last integral give the full
Hilbert--Schmidt bound. The dressed inverse, determinant and
absolute trace-kernel argument of
Theorem~\ref{thm:carrier-flight-dressed-uv} now apply, including
the actual infinite level sums. The constants are finite at each
fixed electric order.

At zero electric coupling the time-independent action separates in
momentum. The pair kernel has
$\|K(k)\|_{\rm HS}\le C(1+|k|)^{-13}$: for $|k|\ge1$ integrate
$C|k|\,a^{14}(1+a^{-1/4})e^{-a^2|k|^2}$ and rescale $a|k|$.
For the onsite vacuum amplitude $Z_k$ in each mode,
\[
 0\le\log Z_k\le C_{m,T}\|K(k)\|_{\rm HS}^2 .
\]
The lower bound is Jensen's inequality since the onsite vacuum
has zero expectation of the normal-ordered free energy and the
off-diagonal pair term. The upper bound is the relative determinant
estimate. These amplitudes are continuous in $k$ by the integrated
heat formula and continuity of the diagonal transfer kernel
($|E_j(k)-E_j(k')|\le |k-k'|$ uniformly). Hence the Riemann sums
$V^{-1}\sum_k\log Z_k$ converge to the ordinary finite integral
with measure $\mathrm d^3k/(2\pi)^3$. Independent visits and the
internal Hilbert trace have changed the volume dependence already
at the level of the bare action.

At zero hopping the threshold follows by minimizing two ordinary
positive relativistic energies at fixed total momentum; the
minimum occurs at equal velocities. At nonzero hopping the
one-particle matrix at $k$ has diagonal entries
$E_N(k),-E_N(k)$ and off-diagonal entries $K(k),K(k)^*$.
The Schur complement excludes every real spectral value in
$(-E_m(k),E_m(k))$, since $E_N(k)\ge E_m(k)$. Here
$E_m(k)=\sqrt{m^2+|k|^2}$. In finite volume the full off-diagonal
operator is Hilbert--Schmidt by the decay just proved. The
difference of the interacting and onsite sign projections is
Hilbert--Schmidt, as follows by integrating their resolvent
difference on the imaginary axis with the majorant
$C\|K\|_{\rm HS}/(m^2+t^2)$. Thus the two polarizations give
unitarily equivalent fermion Fock representations. The interpolating
gapped projections keep the vacuum in the neutral sector.
A particle and a hole in these bands cost at least their
corresponding $E_m$ energies. The triangle inequality gives the
asserted neutral threshold. The finite-volume construction allows
the continuum band and trace densities to be read without asserting
a unitary equivalence of global vacua in infinite volume.
\end{proof}

In this version opening a normal channel changes its mass before
the Lorentzian dispersion is formed. This repairs the additive
normal-gap threshold and the zero-charge volume test. Equal-time
optical sampling still selects a spatial slicing. Covariance of
the complete measured current, interacting cuts and the normalized
thermodynamic limit with electric interactions require separate
tests on this action.

\paragraph{The charge-density boost test on the full internal tower.}
The mass assignment repairs both the free threshold and the
zero-charge vacuum pressure. We can also test its actual density at
nonzero pair hopping and zero dynamical electric charge. In the
continuum denote the filled projection of the interacting band
operator at $k$ by $\Pi(k)$ and the constant onsite negative-slot
projection by $\Pi_0$. The sign-resolvent estimate in
Proposition~\ref{prop:carrier-mass-visits-uv} gives
\[
 L=\int_{\mathbb R^3}\|\Pi(k)-\Pi_0\|_{\rm HS}^2
                         \,\mathrm d^3k<\infty .
\]
This norm includes all the normal channels. The density has the
one-particle identity as its internal matrix, multiplied by the
optical factor $\overline P(P)$ at spatial transfer $P$. In the
filled sea its equal-time spectral weight per unit spatial volume,
up to the fixed charge and Fourier normalization, is
\[
 |\overline P(P)|^2\int
 \|(1-\Pi(k+P))\Pi(k)\|_{\rm HS}^2\,\mathrm d^3k .
\]
Indeed an occupied state can be promoted only into an unoccupied
state. The integrand before squaring is bounded in Hilbert--Schmidt
norm by $\|\Pi(k)-\Pi(k+P)\|_{\rm HS}$. Translation invariance of
Lebesgue measure bounds the full integral by $4L$, even if the
optical factor is removed.

For a nonzero spinor hop $L>0$. For example its normal angular
matrix element from the lowest radial state at angular label zero
to the lowest at label $\ell$ has strictly positive radial overlap
on every positive visit. A nonzero off-diagonal hopping matrix
cannot commute with $\Pi_0$ at those momenta. The sum of the
unmeasured excitation weights at $P$ and $-P$ is exactly the
integral of $\|\Pi(k+P)-\Pi(k)\|_{\rm HS}^2$. It tends to $2L$
as $|P|\to\infty$, since translates of an $L^2$ function tend
weakly to zero. The harmonic readout is strictly positive at each
finite $P$. Some nonzero spectral weight therefore survives the
readout, and the pair threshold puts it in a positive timelike
mass interval (choose a bounded subinterval of nonzero weight).

A Lorentz-covariant conserved current with positive invariant
spectral measure in such a band would contribute
$|P|^2\int_{\rm band}\mathrm d\rho(s)/(2\sqrt{|P|^2+s})$
to the equal-time density. This grows at least as $c|P|$.
Thus the internal tower retains a preferred slicing in the actual
charge response even after its masses and matter threshold have
been repaired. Increasing the number of gapped internal channels
with the stated Hilbert--Schmidt summability cannot supply the
missing boost weight.

\paragraph{Connected continuum coefficients with independent normal visits.}
For the mass-tower action the spatial continuum calculation extends
beyond the zero-charge vacuum. At every finite joint order in pair
hopping and electric coupling, each connected coefficient with only
closed charged lines (vacuum or external photon probes) has a finite
density after its single overall spatial translation factor is
removed. Here the slab duration remains finite. In the ordered
representation the relative charged flight integrals and normal
state sums converge, and the remaining photon four-momenta are
absolutely integrable. Use smooth external photon probes, or fixed
nonzero spatial transfers with smooth time tests. A bridge forced
to exactly zero spatial momentum is omitted in the prescribed
harmonic-mode convention. This statement concerns finite-order
continuum coefficients; it does not exchange a sum over hopping
orders with a thermodynamic limit.

For the soft estimate fix the root point $X$ of a closed charged
word and use displacements on $\mathbb R^3$. At a transverse
photon end of momentum $(\nu,p)$, $k=|p|\le1$, subtract the
constant spatial test equal to its transverse polarization at
the root, with its constant phase. The sum of oriented displacements
around the word is zero. This identity holds for the compensated
flights as well, since their zero-frequency current is their
net displacement. Thus each of several soft marks may use its
subtracted current before multiplying out the complete vertices.
On a diagonal jump at displacement $z$ starting at $x$, the
additional difference of phase and readout is bounded by
\[
 C\min\{1,\ T|\nu|+k(1+|z|+|x-X|)\},
\]
since $a\le1$. The resulting integrand of the compensated
transverse current still has zero predictable mean on spherical
annuli, by the same transverse projection for each of the two
terms separately. For every $l\ge2$ its $l$th jump moment is
bounded by
$C_{l,T}\min(1,|\nu|+k)^l(1+|x-X|)^l$.
Indeed $j_{m_j}\le j_m$ has every displacement moment of order at
least two. The maximal displacement of a flight on the finite
slab has moments of all finite orders uniformly in $j$, by the
same jump moment bound and the martingale maximal inequality.
Applying the jump moment inequality with these excursion moments
gives $C_{h,T}\min(1,|\nu|+k)$ in $L^h$ for each soft
transverse end. A finite number of preceding flights and heat
chords only add their excursion moments to $|x-X|$. A soft mark
on a pair hop obeys the same estimate by Gaussian displacement
moments. The Gaussian used to close the word may instead be
bounded in supremum norm, paying its spatial heat trace as before.
Powers of its own displacement are absorbed by widening that
Gaussian. All powers introduced this way are integrable against
the independent visit weights.

For a temporal mark subtract the constant spatial phase on each
time slice. A closed oriented word crosses a regular time slice
as many times in either direction. Its total temporal occupation
therefore has zero charge at $p=0$. At $k\le1$ this replaces the
flight factor by $Z_t-e^{ip\cdot X}$, up to the common scalar
readout $\overline P(p)$. The initial and terminal differences
have $L^h$ norm bounded by $Ck$ times the excursion moments.
The drift costs $E_j(p)-m_j\le k^2/(2m)$. The $l$th martingale
jump moments cost at most $C_l k^l$ by
$|e^{ip\cdot z}-1|\le k|z|$. Integration by parts on each
ordered interval now bounds a temporal end by
$C_{h,T} k\min(T,2/|\nu|)$ in $L^h$. H\"older's inequality on
the flight histories permits these soft estimates and the hard
estimates of Proposition~\ref{prop:carrier-mass-visits-uv} to be
used simultaneously on one word. One chooses the small hard
exponent $\epsilon$ after fixing the necessary finite moments.

After integrating each internal photon frequency, every line between
closed words has a bounded envelope at small $k$. For transverse
lines this follows from
\[
 \sup_{0<k\le1}\int_{\mathbb R}
 \frac{\min(1,|\nu|+k)^2}{\nu^2+k^2}\,\mathrm d\nu<\infty .
\]
For temporal lines the two soft gains cancel $k^{-2}$ and the
squared interval envelope integrates in $\nu$. At $k>1$ the
transverse line costs $Ck^{-3-2\epsilon}$ and the Coulomb line
costs $Ck^{-2}(1+k)^{-28}$. Thus each spatial envelope lies in
both $L^\infty$ and $L^1$. Root each charged word spatially
and integrate its relative positions and flights as above. Its
free root position enforces spatial momentum conservation for
that word. In the connected photon graph on these roots take
the bounded envelope on a spanning tree; every remaining line
integrates its $L^1$ envelope in an independent spatial momentum.
Self-sewings are among these independent lines. The Hilbert--Schmidt
internal heat transfers bound all internal mode sums. Independent
normal visit laws give a product over distinct words before photon
sewing. The normalized logarithm therefore removes the spatially
disconnected products without leaving a common-radius cumulant.
This proves the single-volume connected density assertion.

\subsection{Testing the physical density after correcting the mass shell}
\label{sec:carrier-density-uv}

\paragraph{An exact relativistic shell with an optically transferred spinor.}
On one fixed sector plane consider a positive physical sampling
length $a$, charged lapse $\tau>0$ and mass $m>0$.
Define a number-preserving kinetic term and a pair hopping term
separately before introducing the gauge connection. Set
\[
 \begin{aligned}
 E(k)&=(m^2+|k|^2)^{1/2},&
 t(k)&=\boldsymbol\alpha\cdot k\,e^{-\tau|k|^2},\\
 c(k)&=\big(E(k)^2-|k|^2e^{-2\tau|k|^2}\big)^{1/2},&
 h(k)&=\beta c(k)+t(k).
 \end{aligned}
\]
Here $\beta$ and $\alpha_i$ are the Dirac Hamiltonian matrices.
The kinetic energy $c(k)$ is at least $m$. At the selected pair
hopping strength one, $h(k)^2=E(k)^2 I$ exactly. At low momentum
$t(k)=\boldsymbol\alpha\cdot k+O(|k|^3)$ and
$c(k)=m+O(|k|^4)$. At large momentum the normalized eigenspinors
approach the onsite mass slots. Thus this choice preserves the
relativistic pole while changing its momentum-dependent spin
projections.

For a spatial symbol $b$ let $b_A$ be its kernel transported
along the straight lifted chord with the Wilson connection of
$P_a A$, $P_a=e^{-a|\nabla|}$. In particular take
\[
 l(k)=\sqrt{c(k)-m},\qquad C_A=m+l_A^*l_A .
\]
This square fixes all kinetic photon contacts and makes the
diagonal spatial operator positive. The definition on smooth
backgrounds is magnetic quantization of an order-$1/2$ symbol,
using the closure of its quadratic form. Indeed
$l(k)=|k|^2[(1-e^{-2\tau|k|^2})/(|k|^2(c(k)+m))]^{1/2}$
extends smoothly over zero. Its derivatives of order $j$ are
$O_j(\langle k\rangle^{1/2-j})$ at infinity, in fixed units
with $\langle k\rangle=(1+|k|^2)^{1/2}$.

Use the magnetic kinetic energy $C_A$ on each particle and its
charge conjugate on each hole, with the corresponding normal
ordered number operators in the empty particle--hole Fock
representation. Use the off-diagonal heat kernel $t_A$ and its
adjoint for pair creation and annihilation. This specifies a
different native charged action. The covariant time derivative
couples to $P_a A_0$. Both spatial blocks transform by the endpoint
phases of $P_a\chi$; the complete Gauss square follows from the
same temporal density. In particular the spatial response of the
number-preserving kinetic term has to be differentiated along
with the pair response. On the charged four-arrow chart the
composition $l_A^*l_A$ uses the two transported kernels with a
common intermediate spatial point.

\begin{proposition}[Fixed-width bare graphs with the exact free mass shell]
\label{prop:repair-dirac-uv}
On the finite slab and spatial torus, with the harmonic mode
prescribed and onsite particle--hole and photon vacuum endpoints,
the complete bare coefficients of this action have joint absolute
bounds at each finite electric order. The result includes neutral
vacuum coefficients, smooth photon tests and smooth charged
endpoint tests with the canonical Gauss dressing. The off-diagonal
hopping multiplier can be any finite real constant for the graph
estimate; take it equal to one for the mass-shell identity above.
The bound keeps the independent carrier and photon integration
variables and includes kinetic, hopping and Coulomb contacts.
\end{proposition}
\begin{proof}
The number-preserving diagonal dynamics takes the empty
particle--hole vacuum into itself with unit amplitude for every
smooth background. Conditional on that background the charged
vacuum factor is therefore
$\det_2(1+S_A K_A)$ with $K_A$ the pair hop and $S_A$ the
diagonal time inverse. For smooth temporal fields first remove
the sampled scalar potential by its time primitive and keep its
endpoint phases. This conjugates both spatial operators. The
positive and negative one-particle slots of the time inverse then
have the one-sided kernels of the forward particle and backward
hole transfers, bounded by $e^{-m|t-s|}$. The formula for the
relative determinant follows by expanding the bounded
Hilbert--Schmidt pair hops between these transfers, or by the
two-slot elimination used for the optical source. Its linear
trace vanishes by the off-diagonal grading. This vacuum
assignment eliminates pure number-preserving cycles at the
source, including all their photon derivatives.

Here are detailed bounds at the expansion point. Extend the pair
hopping by zero outside the slab and put
$S_0=(\partial_t+\beta c)^{-1}$ and $X=S_0K_0$.
One has $\|S_0\|\le 1/m$ and $\|X\|_{\rm HS}<\infty$ by the
one-sided time kernel and the spatial Gaussian. The operator
$\beta(\partial_t+\beta c+K_0)$ has real part $c\ge m$ and its
adjoint has the same lower bound. Fourier transformation in space,
or bounded perturbation of the diagonal generator, then gives
the inverse with norm at most $1/m$. Thus $A_*=(1+X)^{-1}$ is
bounded at arbitrary fixed hopping.

A spatial Wilson derivative of order $j\ge1$ of $l_A$ is an
average of translates of a $j$th spatial derivative of $l$, one
factor $e^{-a|p|}$ for each transverse photon. The same estimate
for a temporal primitive uses momentum differences generated by
endpoint phases: a difference with shift $p$ is an integral of
$p\cdot\nabla l$ over a unit segment. Each such insertion contributes
at most an extra $|p|$ and the factor
$e^{-a|p|}\min(T,2/|\nu|)$, with $\nu$ its own carrier frequency.
In a nonempty derivative of $l_A^*l_A$ an undifferentiated
order-$1/2$ factor can multiply a derivative of order at most
$-1/2$. The ratio of the two translated momentum weights costs a
polynomial in the external transfers, by
$\langle k+v\rangle\le C\langle k+u\rangle\langle v-u\rangle$.
All nonempty kinetic derivatives are therefore bounded spatial
operators, with those photon envelopes and a fixed polynomial.
For a differentiated pair hop the translated Gaussian instead
gives a Hilbert--Schmidt spatial bound with the same envelopes.
These claims follow first at finite spatial Fourier transfer,
then by integration of the bounds. They include multiple
derivatives on a single chord and the product contacts in $C_A$.

Differentiate $S_AK_A$ at zero field before tracing. Every term
is a product of $S_0$, bounded marked diagonal vertices and a
rightmost factor $S_0K_I$, where $K_I$ is a possibly unmarked
pair hop. This last factor is Hilbert--Schmidt in space and
carrier time. All other factors are bounded on the same space.
Derivatives of its logarithmic determinant are therefore cycles
of at least two Hilbert--Schmidt factors after multiplication
by $A_*$. For a single differentiated factor use
$A_*-1=-A_*X$. The absolute cyclic kernel estimate of
Proposition~\ref{prop:constitutive-magnetic-uv} retains the
carrier variables of these grouped factors in a Fourier trace.
For open lines, the derivatives of the full inverse also allow
pure diagonal insertions. At the expansion point its time kernel
preserves spatial momentum and is bounded in operator norm on
the slab, by the one-sided kernel and the resolvent identity.
Between smooth boundary tests a finite open chain therefore has
bounded time integrals and a convergent spatial sum, even without
a pair insertion. Any carrier-frequency Fourier closure is needed
only on closed lines.

Each transverse line costs $\pi/|p|$ after frequency integration.
A temporal line costs an integrable product of its two interval
envelopes in $\nu$, times the Coulomb factor $|p|^{-2}$.
All remaining polynomials in spatial transfers are integrable
against the fixed positive optical widths. This proves the
absolute graph bounds without a momentum cutoff.
In particular the argument uses the number-preserving kinetic
contacts derived from $C_A$, not a vacuum trace for a local
minimally coupled Dirac sea. The statement makes no uniform claim
as the widths shrink.
\end{proof}

\paragraph{The pole and the relativistic current are separately measurable.}
Now take the translation-invariant continuum free charged sea at
the selected hopping strength. Let $P_\pm(k)=(I\pm h(k)/E(k))/2$.
An insertion of either the density or full spatial current at
momentum $P$ can create a particle--hole pair of energy
$E(k+P)+E(k)\ge\sqrt{|P|^2+4m^2}$. This follows by convexity
of $E$ about the midpoint of $k+P$ and $-k$. Both spatial
variations of the kinetic square and the optical variation enter
the physical current. Their on-shell pair matrix elements at fixed
$P$ are square integrable: the pair vertex has a Gaussian tail
and the bounded diagonal vertex connects opposite energy slots
only through the rapidly decreasing difference between $P_\pm$
and the onsite projections. This establishes forward-cone support
for the separated free current at full hopping strength. There is
a further exact test of its matrix elements.

\begin{proposition}[Boost test of the filled reference]
\label{prop:repair-projector-uv}
For the density delivered by this action, even its unsmeared
canonical part cannot be the time component of a nonzero
conserved Lorentz-vector vacuum current with a positive
timelike spectral measure. The obstruction also applies to
translation-invariant finite-slot charged seas whose occupied
projector differs from a constant projector by a spatial
square-integrable matrix and whose local density acts by the
identity shift.
\end{proposition}
\begin{proof}
With the charge factored out, its equal-time density weight per
unit spatial volume before optical smearing is
\[
 D(P)=\int_{\mathbb R^3}
       \operatorname{tr}\big[P_+(k+P)P_-(k)\big]
                    \frac{\mathrm d^3k}{(2\pi)^3}.
\]
The actual measured weight is $e^{-2a|P|}D(P)$. With equal
constant ranks of the occupied projectors at each momentum,
the trace is half the squared Hilbert--Schmidt distance between
$P_-(k+P)$ and $P_-(k)$. Put $F=P_--(I-\beta)/2$ and take its
$L^2$ norm in this same momentum measure. Here $F$ is square
integrable and nonzero. Translation of an $L^2$ function to
infinity has vanishing overlap with the original function, as
follows by compact-support approximation. Thus
$0\le D(P)\le2\|F\|_2^2$ and $D(P)\to\|F\|_2^2$.
At every nonzero $P$ there is pair weight in a bounded
timelike energy interval: at $k=0$ the spin trace is positive,
and so is its integral on a small neighborhood.

By Lorentz covariance, conservation and positivity the timelike
part of a vacuum vector-current spectral tensor is proportional,
in signature $(+---)$, to
$(p^\mu p^\nu-\eta^{\mu\nu}p^2)
\,\theta(p^0)\delta(p^2-s)\,\mathrm d\rho(s)$
integrated over $s>0$ with positive spectral measure.
In the rest frame conservation kills the time component and
rotations make the three spatial components equal, giving this
tensor by a boost. Absorb common positive normalization into
$\rho$. Any finite timelike band $[s_1,s_2]$ of positive mass
contributes to the equal-time $00$ weight at least
\[
 \frac{|P|^2\,\rho([s_1,s_2])}{2\sqrt{|P|^2+s_2}}.
\]
This grows at least linearly for large $|P|$. Dividing out the
optical test factor on both sides, if the spatial smearing is an
external measurement, gives the same contradiction with the
bounded $D(P)$. The comparison involves positive intermediate
states at separated spectral energies. Time-ordering contacts
cannot cancel that spectral weight. Thus a relativistic energy
shell removes the earlier spacelike atom while the filled
reference and physical density still fail the boost test.
\end{proof}

\paragraph{Can a different rotation of the relativistic spinor close the gap?}
The boost test has a stronger consequence which does not require a
square-integrable change of the sea. On the continuum plane set
\[
 h(k)=E(k)\big[u(r)\beta+v(r)\boldsymbol\alpha\cdot\hat k\big],
 \quad r=|k|,\quad u\ge0,\quad u^2+v^2=1 .
\]
Here $u,v$ are real radial functions; the value at $k=0$ can be
fixed separately. The condition $u\ge0$ expresses nonnegativity
of the diagonal kinetic slot. Keep the canonical identity-shift
charge density with linear temporal source. The free vacuum fills
the entire negative-energy subspace. Its intermediate density
states are the free particle--hole pairs of this spinor. A bounded
scalar optical multiplier depending only on external spatial
momentum may be applied to the density. Take it nonzero at the
transfer used for the ultraviolet test.

\begin{proposition}[A boost obstruction for the full radial mixing law]
\label{prop:repair-boost-tail-uv}
Suppose that the free density and its conserved current are to
transform as a positive-vacuum-measure Lorentz vector with
nonzero spectral mass in a compact timelike band. The complete
unsubtracted Euclidean density correlation at nonzero fixed
spatial transfer cannot then be ultraviolet finite on this
class of exact-shell actions. This includes the case where the
optical factor is applied only to an external test of the
purported covariant current.
\end{proposition}
\begin{proof}
Write $n(k)=(u(|k|),v(|k|)\hat k)$ and $p\ne0$. The filled
spin trace for density pair creation is
$M(k,p)=1-n(k+p)\cdot n(k)=|n(k+p)-n(k)|^2/2$.
Its transverse projection in the three spatial components
perpendicular to $k+p$ gives
\[
 M(k,p)\ \ge\
 \frac{v(r)^2}{2}\,
    \frac{|p|^2-(p\cdot\hat k)^2}{|k+p|^2}.
\]
At Euclidean frequency $\omega$ the two temporal orderings
supply $2D/(\omega^2+D^2)$, $D=E(k)+E(k+p)$, after spatial
inversion of one orientation. This is the complete connected
density coefficient for the canonical field-independent charge:
the vacuum one-point function drops out and the source has no
second temporal vertex. Its integrand is nonnegative. On an
exterior cone with $|p\cdot\hat k|\le |p|/2$, the temporal
factor is bounded below by $c/r$. Finiteness at this single
transfer therefore requires
\[
 \int_{r_0}^\infty v(r)^2\,\frac{\mathrm dr}{r}<\infty .
\]

Now evaluate the equal-time density weight only on the intermediate
states with invariant squared mass $s=D^2-|P|^2\in[s_1,s_2]$,
where the assumed positive timelike measure has nonzero mass.
The pair momenta are $k+P$ and $-k$. In their common rest frame
their energies are $\sqrt{s}/2$ and their opposite three-momenta
have magnitude $\sqrt{s/4-m^2}$. For $s\le s_2<\infty$ this
magnitude is at most a fixed fraction strictly less than one
of the rest-frame constituent energy. A boost to large $|P|$
therefore gives both constituents positive longitudinal momentum
between $c|P|$ and $C|P|$, and transverse components of bounded
magnitude. The constants use $m>0$ and $s_2$. This change of
coordinates describes the support of a pair of mass shells;
it assumes no boost law for the optical spin matrix.

Since $u\ge0$, $|1-u|\le v^2$. Hence
$M(k,P)\le C[v(|k|)^2+v(|k+P|)^2]$. In each bounded transverse
disk the absolute longitudinal differential is bounded by a
constant times $\mathrm dr$ for the corresponding constituent
radius. Integrating over the restricted pair states gives
\[
 D_{[s_1,s_2]}(P)\le
 C\int_{c|P|}^{C|P|}v(r)^2\,\mathrm dr .
\]
Covariance of a conserved timelike spectral tensor gives the
lower bound $c_1|P|$ for the same band by the preceding
proposition. If the current itself carries a bounded optical
factor, its removal in this upper bound costs at most the
original uniform supremum constant; if it is external smearing,
compare the unsmeared quantities. It follows in either case that
\[
 \int_{c|P|}^{C|P|}v(r)^2\,\frac{\mathrm dr}{r}\ge c_2>0
\]
for every sufficiently large $|P|$. Choose the successive boost
momenta in a geometric progression of ratio greater than $C/c$.
The intervals are disjoint, contradicting the necessary
ultraviolet integral. This proof does not require monotonicity
or a limiting value of $v$.
\end{proof}

The incompatibility uses the canonical two-slot charge acting on the
complete free pair space, a nonnegative radial mass slot, and the
absence of a temporal source contact. A polarization contribution
to the density, a collective charged operator with a different
intermediate spectrum, or a different dependence on the temporal
probe changes the coefficient to be tested. A replacement requires
the complete measured current, its charge at zero transfer, and the
full determinant from the same action. Improving only the dispersion
or only the transverse spatial vertex leaves the density test intact.

\subsection{The measured current and its dynamical photon sewing}
\label{sec:carrier-sewing-uv}

\paragraph{A relativistic current from a propagating normal amplitude.}
Use the ordinary massive free Dirac translation generators on the
locked plane, in carrier units $c=\hbar=1$. Assign the current a
proper normal half-line with Euclidean boundary problem
$(-\partial_u^2+M^2-\Delta_4)B=0$, $B(0)=A$, decaying at
infinity. Set $M=2m$ for the particle mass $m>0$. A boundary mode
of squared Euclidean momentum $X$ reaches depth $u$ with amplitude
$e^{-u\sqrt{M^2+X}}$. Let $u$ be the sum of $N$ independent
exponentially distributed positive distances, each with mean $\ell$.
The measured amplitude is the average over these distances, divided
by its value at zero transfer. With the square root positive below
threshold the resulting analytic form factor is
\[
 F_N(z)=\left(\frac{1+\ell M}
                      {1+\ell\sqrt{M^2-z}}\right)^N .
\]
Here $\ell$ must be specified by the branch. This is an amplitude
readout of a boundary current, with no independent set of excited
normal channels. Promoting the distances to final-state quantum
labels would change the inclusive cut of the source. On the positive
timelike cut the normal exponential is a phase. The boundary value
on the negative-energy mass transfer is its conjugate; the two are
fixed by Hermiticity. Equal invariants on a given time-oriented
Lorentz orbit have equal multipliers. The scalar boundary values
do not change the sector-plane vectorial pole of the charged field.

Define a current on the free Dirac particle--hole Fock space by
multiplying each matrix element of its vacuum-relative vector current
with transfer $q$ by $F_N(q^2+i0\,\operatorname{sgn}q^0)$.
This prescription applies also to elastic scattering, where
$q^2\le0$. The spinor matrix elements satisfy the free vector
Ward identity before multiplication. The new current is Hermitian,
conserved and Lorentz covariant under the proper orthochronous group.
Its integrated charge equals the Dirac charge: same-band zero spatial
transfer has $F_N(0)=1$, and opposite-band zero spatial transfer of
the density vanishes by orthogonality. This construction uses
the physical free sea and its translations at the outset.

\begin{proposition}[Positive pair spectrum and all finite classical-source orders]
\label{prop:carrier-invariant-current-uv}
The positive two-current spectrum of this comparison lies in
$q^0>0$, $s=q^2\ge4m^2$. Its scalar pair measure has density
\[
 \rho_N(s)=\frac{e^2}{12\pi^2}\sqrt{1-4m^2/s}
       (1+2m^2/s)
       \frac{(1+\ell M)^{2N}}
            {[1+\ell^2(s-M^2)]^{N}} .
\]
For $N\ge2$, the spatial two-current has a finite unsubtracted
time-Fourier integral at fixed spatial momentum, including its
finite contact cost. On a finite time slab all finite coefficients
of the Dirac vacuum transfer coupled linearly to smooth classical
source packets are absolutely integrable in the charged spatial
momenta and ordered vertex times. These assertions concern the
current as defined on the mass shell and time-ordered after the
normal amplitude has been evaluated.
\end{proposition}
\begin{proof}
The vacuum current creates one free particle and one hole.
Their future mass shells place the total transfer in $s\ge4m^2$.
Multiplying the spin trace and phase space in
\eqref{eq:physical-inclusive-pair-uv} by the stated absolute
square gives the density. At zero spatial momentum its positive
separated-time spatial kernel has weight
$s\, e^{-\sqrt{s}|t|}/(2\sqrt{s})$ against this density, per
spatial component. With $C_N=\int_{4m^2}^\infty\rho_N(s)\,\mathrm ds$,
the time transform for $N\ge2$ is
\[
 \int_{4m^2}^{\infty}\frac{s\rho_N(s)}{s+\omega^2}\,\mathrm ds
   = C_N-\omega^2\int_{4m^2}^\infty
                      \frac{\rho_N(s)}{s+\omega^2}\,\mathrm ds .
\]
At other fixed spatial momenta the covariance numerator adds bounded
spatial momentum products and the denominator acquires $|\mathbf q|^2$.
All these integrals are finite for $N\ge2$. A single normal interval
has $\rho_1(s)\asymp s^{-1}$, so its contact cost $C_1$ diverges
although the scalar dispersion integral exists.

For the higher orders write $E_k=\sqrt{m^2+|k|^2}$ and
$P_\epsilon(k)$ for the two free Dirac energy projections,
$\epsilon=\pm1$. At a source insertion carrying spatial momentum
$p$, the transfer invariant in the two slots is
\[
 z_{\epsilon\eta}=
 (\epsilon E_{k+p}-\eta E_k)^2-|p|^2 .
\]
The spin factor between projections is bounded. The free time
propagator has modulus at most one between ordered Euclidean
vertices in its forward or backward band. A nonzero closed word
with at least two distinct times must contain both bands: a cycle
cannot have strictly increasing times along every segment or
strictly decreasing times along every segment. It consequently
has at least two band changes. At those vertices, for $|k|$
larger than a sufficiently large multiple of $m+\sum|p_i|$, the
invariant is bounded below by $c|k|^2$ and
$|F_N|\le C(1+|k|)^{-N}$. The same circulating momentum may be
used around the loop because the other momenta are fixed transfers.
The product for the complete cycle is bounded there by
$C(1+|k|)^{-2N}$, integrable in three dimensions for $N\ge2$.
On the remaining ball all multipliers are uniformly bounded and
the volume grows at most polynomially in the transfers.
The vacuum-relative charge convention fixes the one-point value.
Integration against spatial Schwartz packets uniformly on a
finite slab completes the estimate, including products of finitely
many cycles. No dynamical photon momentum has been integrated.
\end{proof}

\paragraph{Why the conserved density does not yet generate local electric gauge transformations.}
Let $K_p(k+p,k)$ be the one-particle Fourier density for the
unit charge, the sum of $F_N(z_{\epsilon\eta})P_\epsilon(k+p)P_\eta(k)$
over the two bands, with the boundary orientation above. Define
$T_p(k)=\operatorname{tr}K_p(k+p,k)^*K_p(k+p,k)$.
Orthogonality gives a sum of the four squared moduli weighted
by $\operatorname{tr}(P_\epsilon(k+p)P_\eta(k))$.
For $k=0$ and $|p|\to\infty$ the two distinct invariants are
$2m^2\mp2mE_p$. Both multipliers tend to zero, so $T_p(0)\to0$.
For fixed $p\ne0$ and $k=np$, $n\to+\infty$, the two same-band
invariants tend to zero and the corresponding projections become
identical along this ray. The opposite-band weights vanish.
Thus $T_p(np)\to4$ for the four-component spinor. Fix $p$ large
enough that $T_p(0)<4$. At some integer $n\ge1$ the difference
$T_p(np)-T_p((n-1)p)$ is nonzero. This is precisely the spin trace
of the momentum-diagonal block of $[K_{-p},K_p]$, since
$K_{-p}=K_p^*$. Continuity gives a nonzero matrix element on a positive-measure
neighborhood of this nonzero momentum. On a torus choose a large
allowed $p$; its real sine and cosine already give a noncommuting
pair. In the continuum one may also approximate the two Fourier
modes by smooth approximate identities in transfer momentum.
Their density operators converge strongly on spatial packets,
by continuity of translations and the uniform bound on the band
matrices. Products converge strongly by the uniform operator
bounds. Commutation for all real test densities would pass to
this limit.
Second quantization retains this operator part of the commutator;
a scalar vacuum term cannot make it disappear.
Thus the integrated charge and the on-shell vector Ward identity
do not give the field-independent commuting local density
generators of the Abelian Gauss construction.

A different prescription couples a Euclidean Dirac determinant
through the connection
$A_{\rm eff}=P_{\rm L}A+F_N(-Q_E^2)P_{\rm T}A$
on nonzero modes. Its longitudinal gauge variation is the usual
one. Each of its two transverse source derivatives multiplies the
ordinary Dirac bubble by the same fixed external-momentum factor.
For fixed $Q_E^2>0$ this factor is nonzero and gives no suppression
of a hard charged loop. This prescription therefore does not
derive the finite ordered coefficients above. A dynamical gauge
completion has to specify the constraint and the contacts that
implement it and then pass the photon sewing test.

\paragraph{A spectral test of the complete photon-sewn vacuum graph.}
A finite two-current response at each fixed momentum leaves another
integration when a photon closes its two ends. Suppose that on a
reconstructed sector plane its Lorentz-invariant transverse eigenvalue
is $K(X)$, with $X=Q_E^2>0$. Assume that the \emph{same} two-source
kernel, including its contacts, is sewn with a translation-invariant
transverse photon covariance $D_T(X)$. A change of the joint
normal measure when the ports are sewn would change this hypothesis.

\begin{proposition}[A positive physical band obstructs direct Maxwell sewing]
\label{prop:carrier-sewn-uv}
Suppose $D_T(X)\ge c/X$ for sufficiently large $X$.
Up to an overall sign and finitely many contact polynomials and
dispersion subtractions, let $K(X)$ have the spectral representation
with kernel $X/(s+X)$ and a positive measure $\mathrm d\rho(s)$.
Suppose $0<\rho([s_0,s_1])<\infty$ on some
$0<s_0\le s_1<\infty$. Then the complete bare vacuum contraction
of these two ports fails absolute UV convergence. The conclusion
allows any integrable high-mass tail of the current and any finite
number of polynomial contacts. It holds in particular when $D_T$
is the Maxwell propagator. A positive spectral photon covariance
with a nonzero massless pole also satisfies the lower bound.
\end{proposition}
\begin{proof}
Choose a finite difference order $N\ge1$ large enough to remove
all polynomial terms and justify differencing under the subtracted
dispersion integral. For a fixed step $h$, write
$\Delta_h f(X)=f(X+h)-f(X)$. After differencing set $h=X$. Using
$X/(s+X)=1-s/(s+X)$, the absolute difference is exactly
\[
 |\Delta_X^N K(X)| =
 \int \frac{N!\,sX^N}{\prod_{j=0}^N [s+(1+j)X]}\,\mathrm d\rho(s).
\]
Indeed the subtractions in the Cauchy kernel are polynomials of
bounded degree in the evaluation variable; sufficient differencing
eliminates them. The integrand has one sign. On the fixed mass
interval it gives a lower bound $c_N/X$ for $X\ge\max(1,s_1)$.
But the difference is a fixed linear combination of $K((1+j)X)$.
Integration and a change of variables therefore give, for large $R$,
\[
 \int_{1}^{(N+1)R}|K(X)|\,\mathrm dX
 \ \ge\ c'_N\log R-O(1).
\]
Four-dimensional polar integration supplies a constant times
$X\,\mathrm dX$. Thus the radial absolute integral of
$D_T(X)K(X)$ diverges. When $K=X\int\mathrm d\rho(s)/(s+X)$
with no further contact, $K$ itself is bounded below by a strictly
positive constant at large $X$. In that case the shell integral
grows at least linearly in $R$, or quadratically in $|Q_E|$.
\end{proof}

To identify the graph, write the quadratic term of the complete
matter functional at the background as
$e^2\langle A,\mathcal K A\rangle/2$, with zero one-current
expectation. Its Hessian $\mathcal K$ includes the two one-current
terms and the direct second variation of the source, including the
determinant and any measure contacts. When the remaining photon
integration is the centered Gaussian specified above, its order-$e^2$
contribution to the connected vacuum logarithm is, up to the
action sign, $\operatorname{Tr}(D\mathcal K)/2$ per volume with
the coupling removed. Transversality gives a nonzero fixed
polarization multiplicity times $D_TK$. The Abelian free gauge
determinant and the Gaussian normalization are independent of $e$.
Any further source-dependent field or boundary term must be included
in the complete Hessian before using its spectral representation.
Nonlocal time ordering can produce contact kernels beyond the finite
polynomials in the hypothesis.

This test is on a connected vacuum graph per unit carrier volume.
A division by the source-free vacuum can remove it from normalized
external responses; that division does not prove absolute
convergence of the bare graph. The occupied finite mass interval
already forces the lower bound. Reducing pair production at large
invariant mass alone cannot evade it. To use a different geometric
sewing, one must derive the two-port contraction from the microscopic
gluing measure, exhibit where the product $D_TK$ ceases to be the
complete graph, and recover the physical two-current and its positive
cut from the same construction.

If the source family also probes the Maxwell kinetic density, there
is a normalized response with the same obstruction. Use the uniform
gauge-invariant kinetic variation $S_M\mapsto(1+h)S_M$ for $h>-1$,
holding $e$, the matter functional and the measured current fixed.
The transverse covariance becomes $D_T/(1+h)$ for the Gaussian
Maxwell field. At order $e^2$ the logarithm of $Z(h)/Z(0)$ per unit
volume has the same complete contraction multiplied by
$(1+h)^{-1}-1$. Differentiating at zero leaves its negative.
Equivalently the derivative inserts the integrated Maxwell kinetic
density; its free vacuum normalization has no term of order $e^2$.
This statement uses a kinetic-source insertion at zero transfer.
It does not assert divergence of scattering amplitudes with
elementary external photons or of a current sourced by a metric
variation that also changes the matter action.

\paragraph{Which geometric port carries the physical photon?}
For finitely many Hermitian current ports the residue provides
a direct test. Suppose that after summing all diagonal and cross
contractions the photon covariance is the positive matrix
$D(X)\succeq Z/X$ and that $Z$ is its nonnegative massless residue.
Write the complete unsubtracted current matrix as
$K(X)=X\int (s+X)^{-1}\mathrm dR(s)$ with a positive matrix
measure. For $B=[s_0,s_1]\subset(0,\infty)$, positivity gives
\[
 \operatorname{tr}(D(X)K(X))
     \ \ge\ \frac{\operatorname{tr}(Z R(B))}{s_1+X}.
\]
Indeed $K(X)\succeq X R(B)/(s_1+X)$, and the trace of the product
of two nonnegative matrices is nonnegative even when they do not
commute. A single band with $\operatorname{tr}(Z R(B))>0$ forces
the quadratic bare divergence upon four-momentum sewing. For
the pure Maxwell covariance $D=Z/X$, the finite-contact argument
in Proposition~\ref{prop:carrier-sewn-uv} also applies with scalar
measure $\operatorname{tr}(Z\,\mathrm dR)$. In these tests an empty
overlap is equivalent to $R(B)^{1/2}Z^{1/2}=0$. An exponentially
small transfer between two distinct ports does not establish this
orthogonality for a physical photon coupled to a charged current.

\paragraph{Does a common normal trajectory change this conclusion?}
This question can be decided before selecting any distribution for
the geometric depths. Let the zero-charge measure factor as the
free Gaussian boundary photon of covariance $D$ and a joint
matter--geometry measure $\mu$. The latter may correlate the
normal positions on every charged segment in the diagram. Suppose
the external connection $b$ and the dynamical connection enter
the interaction and the boundary source by the same replacement
$A\mapsto A+b$. Suppose the quadratic Taylor coefficient in
$e(A+b)$ exists as a distribution on smooth test fields and all
its source terms are included. Denote the resulting two-port
kernel after the $\mu$ integration by $K$. With a vanishing
one-point current its contribution to the vacuum logarithm from
closing those two ports is exactly $\tfrac12\operatorname{Tr}(DK)$,
up to the overall action sign. This follows first for finitely many
smooth test directions from the second moment of a Gaussian. The
equality then holds for every finite test projection. If the
unprojected complete graph converged absolutely, its projected
integrals on expanding momentum domains would converge as well.
A correlated matter path by itself therefore leaves the argument
of Proposition~\ref{prop:carrier-sewn-uv} intact. Its effect has
already been included in the complete kernel $K$.

For a different proposed gluing write $S(X)$ for its traced radial
integrand per $\mathrm dX$, absorbing the common transverse
multiplicity and angular area into the normalization. For ordinary
Maxwell sewing this is $K(X)$ in the same normalization.
If the new gluing is absolutely integrable and the measured
$K$ has the positive band in the proposition, its difference from
ordinary sewing must satisfy, for some fixed $N$ and $c>0$,
\[
 \int_1^{(N+1)R}|S(X)-K(X)|\,\mathrm dX
           \ge c\log R-O(1) .
\]
This follows directly from the reverse triangle inequality and
the finite-difference lower bound. A correction whose radial
absolute integral is already finite cannot restore finiteness.
Changing a boundary measure at sewing must therefore have a
calculable effect on the same hard-momentum region as the measured
current. For a positive mixture of sources with the same Maxwell
sewing, a common polynomial subtraction degree, and a fixed
nonnegative residue weight, the averaged spectral measure remains
positive whenever the weighted current occupies a compact
positive-mass band. The same test then applies to the average,
provided the contact coefficients and the subtracted measures
admit the stated integration. These conditions make clear where
a genuinely different geometry could enter a future construction.

\paragraph{A sharp test for a correlated photon and normal preparation.}
Keeping the exact free Maxwell law at zero electric coupling
already restricts an attempted change of gluing. Suppose that,
conditional on the geometric history $\xi$, the transverse boundary
field is centered Gaussian with covariance $D_\xi$. If its marginal
law is the centered Gaussian with covariance $D$, then $D_\xi=D$
almost surely on the test space. For a real test $h$ its conditional
variance $v_\xi$ has $\mathbb E v_\xi=\langle h,Dh\rangle$.
The Gaussian fourth moment gives
$\mathbb E v_\xi^2=\langle h,Dh\rangle^2$, hence the variance
of $v_\xi$ vanishes. Polarize on a countable dense test family
and use continuity of the covariance forms. Under this conditional
Gaussian assumption the boundary photon and the geometry are
independent at zero charge. Randomizing only its centered covariance
cannot change the Maxwell contraction of the averaged Hessian in
the preceding calculation.

There is a related estimate even for non-Gaussian conditional
laws. Consider a real transverse mode of a finite test projection.
Divide its amplitude by the free standard deviation to obtain
$Y$ with standard normal marginal. Suppose the complete contribution
on this mode has a nonnegative weight $W$ with $0<w=\mathbb E W<\infty$.
The weight can depend on $Y$ and on all the normal variables.
In a model where the direct quadratic contraction on this mode is
weighted by $W$, its ratio to independent sewing is
$\eta=\mathbb E(WY^2)/w$. For any $q>1$ a uniform bound
$\|W\|_q/w\le C_q$ forces $\eta\ge c_q(C_q)>0$.
Indeed for $\delta>0$,
\[
 \mathbb E[W\,{\bf1}_{|Y|\le\delta}]
 \le\|W\|_q\,\Pr(|Y|\le\delta)^{1-1/q}.
\]
Choose $\delta$ so the right side is at most $w/2$.
Outside this interval $Y^2\ge\delta^2$.
More quantitatively, if $\eta$ tends to zero, at least half of the
$W$-weighted mass lies in $|Y|\le\sqrt{2\eta}$ by Markov's
inequality. The ordinary Gaussian measure of this set is
$O(\sqrt\eta)$. Consequently
\[
 \frac{\|W\|_q}{w}\ \ge\ c_q\,\eta^{-(1-1/q)/2}
 \qquad (0<\eta\le1).
\]
Conditioning with the weight proportional to $\mathbf 1_{|Y|\le\delta}$
gives $\eta$ of order $\delta^2$ and the ratio of norms of order
$\delta^{-(1-1/q)}$, so the exponent is sharp. Even when
normal data depend on the photon amplitude, a positive coupling
with uniformly bounded relative moments cannot suppress hard modes
by an arbitrarily large factor in this class. This estimate applies
to a positive mode contribution with an actual Gaussian marginal.
Signed seagulls, mode-mixing vertices and non-Gaussian marginal
fields have to be evaluated with their own source laws.

\paragraph{A relativistic connection with a geometric current transfer.}
A possible repair keeps exact Dirac propagation and mediates the
measured current through a second connection. This can be tested on a
single locked sector plane, in its intrinsic carrier units. Let $A$
be the probed Abelian connection and $B$ the one in the covariant
derivative of a Dirac field of mass $m>0$. Take $0<b<4m^2$ and
$D_{B,\mu}=\partial_\mu+ieB_\mu$. The Euclidean action density is
\[
 \begin{aligned}
 \mathcal L_E
 &=\tfrac14 F_B^{\mu\nu}F_{B,\mu\nu}
      +\tfrac{b}{2}(B-A)^\mu(B-A)_\mu\\
 &\quad+\bar\psi(\gamma^\mu D_{B,\mu}+m)\psi .
 \end{aligned}
\]
A simultaneous gauge shift of $A$ and $B$ preserves this source
with the usual matter phase. At zero external source the relative
vector has Proca mass $\sqrt b$. If both connections are dynamical
one may add a Maxwell term for $A$; here we differentiate at
prescribed $A$. The four-dimensional kinetic terms use the
relativistic metric on continuation. For Euclidean squared
momentum $X>0$ the transverse Gaussian conditional mean is
$B^{\rm T}=f(X)A^{\rm T}$ with
\[
 f(X)=\frac{b}{X+b}
      =\int_0^\infty b e^{-b u} e^{-uX}\,\mathrm du .
\]
The same multiplier is a geometric phase average of isotropic
Euclidean heat displacements at exponential normal lapse $u$.
The second connection provides a local Lorentzian realization of
this rational transfer. No assumption about a Lorentz-invariant
probability distribution of real spacetime displacements is needed.
The transfer is one at zero momentum.

\begin{proposition}[A finite pair cut with a divergent complete source response]
\label{prop:carrier-portal-uv}
The measured tree-level charged-pair current of this action has
forward-cone support and an integrable high-energy tail in the
unsubtracted dispersion weight. Its full bare two-source coefficient
at one charged loop, of order $e^2$, diverges at every fixed
nonzero Euclidean transverse momentum.
\end{proposition}
\begin{proof}
Write $s$ for squared timelike momentum. At the first order which
produces a charged pair, the ordinary Dirac matrix element is
multiplied by $b/(b-s)$. Hence the pair spectral density is
\[
 \begin{aligned}
 \rho_{\rm pair}(s)&=\frac{b^2}{(s-b)^2}\rho_D(s),\\
 \rho_D(s)&=\frac{e^2}{12\pi^2}
 \sqrt{1-4m^2/s}\,(1+2m^2/s)
 \end{aligned}
\]
for $s\ge4m^2$ and zero below threshold. The expression follows
from the two-body spin trace in
\eqref{eq:physical-inclusive-pair-uv}. The two particles obey
the exact relativistic dispersion, and the transfer pole lies
strictly below the pair threshold. In particular
$\int_{4m^2}^\infty \rho_{\rm pair}(s)\,\mathrm ds/s$ is finite.

For the complete source, split $B$ into its conditional mean and
a centered Gaussian fluctuation. The first-order term of the
charged determinant vanishes in the empty background. Its quadratic
term is bilinear in $B$. Averaging that term evaluates both external
legs at the conditional mean; contracting the fluctuation with
itself gives a source-independent vacuum contribution. The
quadratic contact $b A^2/2$ has already entered the free Gaussian
integration. Its combination with $B$ exchange gives the free
transverse Hessian $X b/(X+b)$. It has no $e^2$ term. Thus the
charged part of the Hessian is precisely the free Dirac polarization
tensor multiplied by $f(X)^2$.

For clarity, at proper time $t>0$ the transverse scalar of the
unsubtracted Euclidean Dirac loop has proper-time density, up to
the sign convention for the effective action,
\[
 \frac{e^2}{2\pi^2t}e^{-tm^2}
 \int_0^1 v(1-v)e^{-tXv(1-v)}\,\mathrm dv .
\]
One obtains this expression by taking the two connection variations
of the squared Dirac heat trace. In transverse polarization the
spin term contributes $F(z)$ and the orbital and contact terms
contribute $2(F(z)-1)/z$, where $F(z)=\int_0^1e^{-zv(1-v)}\mathrm dv$
and $z=tX$. Gaussian carrier integration gives the factor
$e^2 e^{-tm^2}/(8\pi^2t)$. Integration of
$\partial_v[(1-2v)e^{-zv(1-v)}]$ gives
$F+2(F-1)/z=4\int_0^1 v(1-v)e^{-zv(1-v)}\,\mathrm dv$.
All these operations take place at positive $t$.
The result is asymptotic to $e^2/(12\pi^2t)$ at zero.
The factor $f(X)^2$ is independent of this internal proper time
and strictly positive. The complete coefficient therefore diverges
even with the contact and vacuum terms included.
\end{proof}

The spectral bookkeeping locates the missing contribution if only
the pair cut is kept. On any finite interval of the Dirac continuum
the following algebraic identity holds:
\[
 \begin{aligned}
 \frac{b^2}{(X+b)^2(X+s)}
 &=\frac{b^2}{(s-b)^2}
       \left(\frac1{X+s}-\frac1{X+b}\right)\\
 &\quad+\frac{b^2}{s-b}\frac1{(X+b)^2}.
 \end{aligned}
\]
Integrate this identity against $\rho_D(s)\,\mathrm ds$. The
continuum and the simple-pole integrals converge at infinity.
The double-pole coefficient grows logarithmically, since
$\rho_D(s)$ tends to a nonzero constant. This is an exact diagnostic
of the unsubtracted spectral kernel. It does not declare a
double pole to be the spectrum of a completed interacting
particle. A finite source would have to control the vector's own
polarization as well as the on-shell pair matrix element.

More generally, any normalized nonnegative law of finite Euclidean
heat lapses gives a strictly positive $f(X)$ at finite $X\ge0$.
If it enters a local Dirac determinant only through this external
readout, its value at fixed momentum cannot suppress the independent
charged loop. This test suggests placing geometric displacement
in the relative charged propagation of the replacement action, and
repeating the contour, closed-determinant and pair-production tests
for that same source.

\paragraph{Giving pair creation its own curvature matrix element.}
The boost test suggests changing how the electromagnetic field creates
a pair. The following action has a finite two-source response and an
actual timelike pair current. Take the free massive relativistic
particle and antiparticle spin spaces with their usual Wigner
transformation and covariantly normalized modes, together with an
empty particle--hole vacuum. The free energy of both species is
$E(k)$. Use a positive number-preserving magnetic kinetic energy
$\sqrt{m^2+(-i\nabla\mp eA)^2}$ on each slot, with the corresponding
linear temporal charge. Its scalar square root is defined by the
positive covariant spatial form. At this stage choose smooth
background probes. Normal ordered second quantization of this
number-preserving kinetic energy fixes the vacuum energy at zero
for every such background. It supplies ordinary charge $\pm e$ on
one-particle states but cannot create a pair from this vacuum.

Add the Hermitian interaction
$-e\int F_{\nu\mu} B_A^{\nu\mu}\,\mathrm d^4x/2$.
The following specification of its zero-background kernel is an
equal-time spatial operator prescription. For on-shell future
particle and antiparticle momenta $p,p'$ let $P=p+p'$, $s=P^2$,
and use the usual spinor matrix element
$j^\mu=\bar u(p)\gamma^\mu v(p')$. Define the
pair-creating part by its matrix element at $x=0$,
\begin{equation}
 \langle p,p'|B_0^{\nu\mu}(0)|0\rangle
 =-i\,f(s)\frac{P^\nu j^\mu-P^\mu j^\nu}{s},
 \qquad f(s)=e^{-\ell^4s^2}.
 \label{eq:repair-pair-curvature-uv}
\end{equation}
Add the adjoint annihilator. The denominator has $s\ge4m^2$ on all
these modes. Equal-time spatial Fourier transformation at the two
charged endpoints defines the bilinear kernel, including the
relativistic one-particle normalization factors. Translation of
the insertion supplies the usual total momentum. The assigned
length $\ell>0$ can come from the same geometric branch. One
constitutive origin for $f$ is a common real Gaussian displacement
in a variable conjugate to squared invariant mass, with mean zero
and variance $2\ell^4$. Its characteristic function at $s$ gives
exactly the factor above. This is a new law for the pair curvature,
with invariant-mass dependence chosen before taking the current
derivative. It is independent of a continuation of the
previous spatial heat hop to timelike momenta.

On the trivial Abelian chart transport the two spatial charge
endpoints of this neutral bilinear with oppositely oriented
electric phases, so it is gauge invariant. For example use the
Coulomb dressing built from the longitudinal primitive $\chi$ at
the two endpoints. The two opposite gauge variations then cancel
the two charged transformations, and on the transverse slice
$B_A=B_0$. The Coulomb dressing preserves the number change
of each term. The kernel and its first two source variations can
be evaluated as quadratic forms on smooth charged packets and
smooth background tests. The calculation below concerns only
these background derivatives and their continuum densities.

\begin{proposition}[A covariant pair cut and a finite two-source coefficient]
\label{prop:repair-pair-curvature-uv}
This action produces a conserved Lorentz-vector vacuum-to-pair
matrix element with the exact massive threshold. Its timelike
spectral density has a finite unsubtracted pair tail. At
electromagnetic order two the complete bare matter vacuum
response to smooth external photons is ultraviolet finite at
each fixed Euclidean transfer and is $O(|Q|^2)$ as the Euclidean
four-momentum $Q$ tends to zero.
\end{proposition}
\begin{proof}
Integration by parts in the first background variation gives
the pair current $e\partial_\nu B_0^{\nu\mu}$. On the free
on-shell modes its derivative is $iP_\nu$ in the stated
matrix-element convention. Since $P_\mu j^\mu=0$, this gives
exactly $e f(s)j^\mu$. Both the vector transformation and the
positive intermediate particle--antiparticle spectrum follow
from the covariant free modes and the scalar character $f(s)$.
The number-preserving charge contributes zero to this matrix
element. By the ordinary spin sum and two-body phase space the
positive scalar pair density is
\[
 \rho_{\rm pair}(s)=\frac{e^2}{12\pi^2}
       \sqrt{1-4m^2/s}\,(1+2m^2/s)e^{-2\ell^4s^2}
          {\bf1}_{s\ge4m^2}.
\]
In particular its integral with weight $1/s$ is finite.

For the time-ordered background coefficient keep its two external
field strengths on the two sides of the free transfer of $B_0$.
This prescription includes the distinction between differentiation
of a time-ordered product and inserting a time derivative inside it.
All pure number-preserving insertions, including their second
variation, annihilate the vacuum. A mixed kinetic and pair
insertion cannot return it to the vacuum. The second derivative
of the curvature coupling differentiates a dressing on a single
pair creator or annihilator, whose vacuum expectation is zero.
The two curvature insertions therefore give the complete coefficient.
At fixed spatial transfer $\mathbf P$ the intermediate energy is
$D=E(k)+E(\mathbf P-k)$ and $s=D^2-|\mathbf P|^2$. The two
one-sided time kernels at Euclidean frequency $\omega$ have
denominators $D\pm i\omega$ with $D\ge2m$. For $\mathbf P$ in
a compact set, $s\ge c|k|^2$ at sufficiently large $|k|$,
uniformly in that set. The remaining spin and tensor factors
in the squared matrix element grow at most polynomially.
Thus every component of the $B_0$ transfer is absolutely
integrable in the relative spatial momentum, uniformly near
$Q=0$. Equivalently the free relative carrier-frequency
integral of two massive time propagators is absolutely bounded
at fixed spatial momenta by Cauchy--Schwarz; it adds only a
polynomial majorant for this estimate. The two external field
strengths each supply one power of $Q$. In particular this
coefficient introduces no constant photon mass term.
\end{proof}

This pair-curvature choice passes the separated-current spectral
test and the zero-momentum photon test at the order computed.
There is an important difference between suppressing a large
invariant relative mass and suppressing a simultaneous boost of
both charged endpoints. A spatially unbounded pair momentum can
carry any fixed $s>4m^2$. Internal photon sewings integrate over
such regions. The above fixed-transfer estimate does not bound
them. The full elastic and multipair currents, the dynamical
Gauss square with this polarization, and every complete
higher source derivative must still be controlled. The
number-preserving kinetic choice by itself establishes no boost
law for its elastic current. These qualifications matter when
using the pair test to choose a physical replacement action.

\paragraph{Sewing the dynamical Maxwell field of the pair-curvature action.}
The distinction between invariant pair suppression and large total
momentum is decisive in this example. Quantize the ordinary boundary
Maxwell field and use the Coulomb dressing just specified, so the
pair tensor on the transverse slice is $B_0$. Keep the canonical
vacuum expectation of its squares. Write $X_i=B_0^{0i}$ for the
electric polarization and $M_i=\epsilon_{ijk}B_0^{jk}/2$ for its
magnetic counterpart. In this paragraph $p=|\mathbf P|$ is the total
spatial pair momentum magnitude, not the relative momentum.
Let $\rho_0=\rho_{\rm pair}/e^2$ in
Proposition~\ref{prop:repair-pair-curvature-uv}.
It is positive and both $\int\rho_0(s)\,\mathrm ds$ and
$\int \rho_0(s)\,\mathrm ds/s$ are finite and nonzero.

\begin{proposition}[Complete Maxwell vacuum sewing of the curvature pair]
\label{prop:repair-pair-sewing-uv}
At electromagnetic order two the bare vacuum energy density of the
pair-curvature action has a divergent photon momentum integral.
With the usual spatial measure $\mathrm d^3\mathbf P/(2\pi)^3$,
the coefficient per unit of that measure is
\[
 \frac{e^2}{4}\int_{4m^2}^\infty
 \rho_0(s)\left[\frac{3}{\sqrt{p^2+s}}-\frac{2p}{s}\right]
 \,\mathrm ds .
\]
The integrand after this mass integration grows as a negative constant
times $p$. Completing the magnetic energy to a square in this
action changes the large-momentum power but leaves the integral
divergent.
\end{proposition}
\begin{proof}
The electric coupling contains one time derivative of the transverse
photon. Its Maxwell Legendre transform supplies the contact
$e^2\|X_{\rm T}\|^2/2$ from the shifted transverse canonical
momentum. Eliminating the scalar potential supplies the longitudinal
energy of the number charge together with the induced polarization
charge, the latter proportional to $i\mathbf P\cdot X$. In the
vacuum the number charge annihilates the ket and bra, so this
longitudinal square contributes $e^2\|X_{\rm L}\|^2/2$.
The sign of the shift does not affect these two contacts.
All terms of the number-preserving kinetic energy annihilate the
vacuum, including its Hessian. The chosen pair dressing has no
transverse-field dependence on this slice. These facts exhaust the
direct order-two contacts of the action.

For the second-order intermediate state take the pair momentum
$P=(D,0,0,p)$ with $D=\sqrt{p^2+s}$ and an emitted photon
$q=(p,0,0,-p)$. The state has energy $D+p$. The linear
Hamiltonian vertex, evaluated with free transverse canonical
momentum, has the same squared emission matrix element as
$F_{\nu\mu}B_0^{\nu\mu}/2$ evaluated on the on-shell photon.
For either transverse polarization $i=1,2$ one has
\[
 |q_\nu B_0^{\nu i}|^2
   =p^2(D+p)^2 |f(s)j^i|^2/s^2 .
\]
Here the opposite spatial momenta of the photon and pair give
$q\cdot P=p(D+p)$. The second-order vacuum energy contributes
the negative of this square times $e^2/[2p(D+p)]$.
It includes the electric--magnetic interference.

The timelike spectral tensor of $f j^\mu$ at fixed $s$
is $P^\mu P^\nu-\eta^{\mu\nu}s$ times the positive weight
$\rho_0(s)\,\mathrm ds/(2D)$. Thus $|f j^i|^2$ for each
transverse component supplies $s$ in this weight and
$|f j^3|^2$ supplies $D^2$. Since $Dj^0=pj^3$, the
longitudinal polarization $X_3$ has squared weight one.
Each transverse $X_i$ has squared weight $D^2/s$. In units
of $e^2\rho_0(s)\,\mathrm ds/(4D)$ the sum of all terms is
therefore
\[
 1+\frac{2D^2}{s}-\frac{2p(D+p)}s
 =3-\frac{2pD}{s}.
\]
This gives the claimed formula. Its positive term is
$O(p^{-1})$ after mass integration and its negative term is
exactly linear in $p$. The same ultraviolet failure occurs on a
fixed vacuum slab. In the energy extracted from its negative
logarithm divided by $T$, the exchange denominator $1/(D+p)$
is replaced by
$1/(D+p)-(1-e^{-T(D+p)})/[T(D+p)^2]$.
The resulting correction is bounded uniformly in $p$ by
$C_T e^2\int\rho_0(s)\,\mathrm ds/s$. It cannot cancel the
linear term. All comparisons can be made with finite shells
before testing the improper integral.

As a different bare action, completing the magnetic energy to
$\frac12\|\mathbf B_{\rm Maxwell}\mp eM\|^2$ with the sign of
the same linear interaction adds $e^2\|M\|^2/2$. Its vacuum
spin weight in the same units is $2p^2/s$. The integrand in
brackets before division by $D$ becomes
$3-2p/(D+p)\ge2$. Consequently this second action still gives
an ultraviolet vacuum energy integral of at least quadratic
degree in the radial upper scale. These are two tests of specified
Hamiltonians and their own contacts. The divergent coefficient is
a vacuum graph. It does not on its own establish divergence of
a normalized scattering response on either action.
\end{proof}

A source coupled to the curvature interaction itself provides a
connected probe of this divergence. Multiply only that invariant
interaction by a smooth real function $\eta(t)$ supported inside
the slab and keep the number charge convention fixed. At second
order in this source the two direct contacts carry
$\int\eta^2\,\mathrm dt$. For the exchange put $b=D+p$ and
extend $\eta$ by zero outside the slab. One has
\[
 \int_{t>u}\eta(t)\eta(u)e^{-b(t-u)}\,\mathrm du\,\mathrm dt
   =\frac{\|\eta\|_2^2}{b}+O_\eta(b^{-2}).
\]
Indeed
$|\int\eta(t)[\eta(t-v)-\eta(t)]\,\mathrm dt|
 \le v\|\eta\|_2\|\eta'\|_2$ for $v>0$.
The error in the energy integrand is uniformly bounded in $p$
as in the slab estimate. The linear large-$p$ term of the
ordinary action therefore remains, multiplied by $\|\eta\|_2^2$.
Taking the connected derivative of $\log[Z(\eta)/Z(0)]$ leaves
this term even when $\int\eta(t)\,\mathrm dt=0$.
This is a measurement of the specified interaction density.
It makes no identification of that composite probe with an
external photon.

\subsection{Local normal dynamics and an interacting gapless photon}
\label{sec:carrier-local-uv}

\paragraph{A spatially local normal state in place of the common radius.}
The volume obstruction in Proposition~\ref{prop:upgrade-common-volume-uv}
comes from letting spatially distant loops measure one normal degree of
freedom. A replacement can be defined in continuous space. Take a real
free normal scalar $\chi$ of mass $M>0$ in its Gaussian vacuum. In
reference length units fix a real radial, smooth spatial packet $h$ of
compact support. Normalize its equal-time variance so that
$X_t(c)=\chi_t(h(c-\cdot))$ has variance one. Set
\[
 r_t(c)=(1+X_t(c)^2)^{-\delta},\qquad
 a_t(c)=L_* r_t(c)^{1/3},\qquad \tau_t(c)=a_t(c)^2 .
\]
The symbol $c$ here is a spatial location. The field can be sampled at
arbitrarily small positive depth. The packet, the normal mass and the
choice of the exponent $\delta$ are constitutive inputs in this
comparison. They would have to be fixed by the microscopic theory to
obtain a single-parameter model. The normal propagator has energy
$\sqrt{|k|^2+M^2}$, including for excitations which carry momentum
between two charged loops.

First take the previous scalar or differentiated heat hop at zero
electric charge. Assign to each actual charged chord its width at
its midpoint and time. Use its adjoint on the reverse hop. Compress the
two endpoints to a spatial cube $D$ for the finite-region test.
The normal field lives on all of space in its own vacuum. The
particle and hole slots have their onsite mass $m>0$. Alternatively
use an exact free Dirac field of mass $m$ in its filled vacuum and
place the heat averages in the interaction. For a concrete interacting
test take another real free scalar $\phi$ of mass $\mu>0$, put
$\phi_a=P_a\phi$ for $P_a=e^{-a|\nabla|}$, and use
\[
 \begin{aligned}
 H_{\rm int}={}&g\int_D\mathrm d^3c\,
  \psi_{\tau,c}^\dagger B\psi_{\tau,c}
                         \sin\!\big(y a\phi_a(c)\big)\\
 &+\lambda\int_D\mathrm d^3c\,\phi_a(c)^4 ,
 \qquad \psi_{\tau,c}=\int G_\tau(c-x)\psi(x)\,\mathrm d^3x .
 \end{aligned}
\]
Here $G_\tau$ is the spatial heat kernel, $a,\tau$ are evaluated at
the vertex and $B$ is a fixed Hermitian spin matrix. The Dirac field
has its ordinary two mass-shell time orderings. The bilinear and the
quartic keep all of their equal-time self-contractions. In particular
the two spinor factors have separate heat displacements even when
their common center is integrated. This action has global electric
charge symmetry. Its interaction does not specify the current of a
locally gauged Dirac kinetic term.

\begin{proposition}[Extensive coefficients with moving local normal measurements]
\label{prop:replacement-local-normal-uv}
Take $0<\delta<3/5$ and a finite slab duration $T$. Both models
above define semibounded finite-region Hamiltonians. In cubes of
volume $V\ge1$ their diagonal vacuum transfers satisfy
$|\log Z_D(T)|\le C_T V$ at arbitrary fixed real couplings, with
$\lambda\ge0$ for the scalar potential. At each fixed perturbative
order, the connected vacuum coefficients divided by $V$ have finite
limits as the cubes exhaust space. In the first model this is the
complete hopping expansion at zero electric charge. In the second
it is the joint expansion in $g,y,\lambda$ about zero. Rooted
connected densities at that order are absolutely integrable in the
relative spatial centers, with the actual local normal field averaged.
\end{proposition}
\begin{proof}
Let $H_\chi$ be the free normal Hamiltonian. At fixed normal
configuration the squared Hilbert--Schmidt norm of the original hop
is bounded by a constant times $\int_D(1+a^{-5})\,\mathrm d^3c$
for a first derivative and by $\int_D(1+a^{-3})\,\mathrm d^3c$
for a scalar hop. Indeed one integrates the squared Gaussian in the
relative displacement at its own midpoint, which is in the convex
cube. The pair form pays an arbitrarily small fraction of $mN$ and
a multiple of this integral as in the optical proof. For the Dirac
alternative the equal-time canonical anticommutator gives
$\|\psi_{\tau,c}\|\le C\|G_\tau\|_2=C'\tau^{-3/4}$ componentwise.
The norm of its interaction density apart from the positive quartic is
therefore at most $C|g|a^{-3}$, uniformly in the real scalar phase.
All these bounds are dominated by a constant times
$Q_D=\int_D(1+X(c)^2)^{5\delta/3}\,\mathrm d^3c$.

For every $b>0$, $\log\mathbb E e^{bQ_D}\le C_bV$. In fact
$b(1+u^2)^{5\delta/3}\le C_{b,\epsilon}+\epsilon u^2$.
The covariance of $X$ on $L^2(D)$ has trace $V$ and operator norm
bounded uniformly by the supremum of the spatial spectral density
$|\widehat h(k)|^2/(2\sqrt{|k|^2+M^2})$ in the chosen Fourier
normalization. Diagonalize this positive covariance. For sufficiently
small fixed $\epsilon$ the logarithm of the Gaussian quadratic
moment is $-\tfrac12\sum_j\log(1-2\epsilon\lambda_j)\le C V$,
where $\lambda_j$ are its eigenvalues. No momentum truncation
enters the estimate.

For a unit vector $f$ in Gaussian coordinates the entropy inequality
bounds $\langle f,Q_D f\rangle$ by
$b^{-1}(\operatorname{Ent}(\|f\|^2)+C_b V)$.
The Gaussian entropy satisfies
$\operatorname{Ent}(\|f\|^2)\le 2\langle f,N_\chi f\rangle$.
For instance integrate entropy dissipation for the stationary unit-speed
Gaussian Ornstein--Uhlenbeck semigroup. Its gradient commutes with the
semigroup with factor $e^{-t}$; Cauchy--Schwarz bounds the Fisher
information of the evolved density by $e^{-2t}$ times its initial
value. Integration proves the inequality on cylindrical functions.
Apply it to the pointwise norm and take form closure. Since
$H_\chi\ge M N_\chi$, choosing $b$ large gives an arbitrarily small
relative form bound and a remainder $C V$. The positive quartic is
a multiplication form with a dense finite-energy domain. Smooth
finite-particle vectors belong to that domain by the finite Gaussian
moments and the inverse moments of $a$. The closed form sum gives
the finite-region Hamiltonian with lower bound $-C V$. Jensen's
inequality applied to its vacuum spectral measure gives
$\log Z_D(T)\ge-T\langle H\rangle_{\rm vac}\ge-C_T V$; the free
vacuum energies have their canonical zero. This proves the two-sided
transfer bound without assuming a pressure limit at fixed couplings.

For the fixed-order assertion first integrate matter and the measured
scalar at a fixed set of normal-field samples. In the onsite model,
spatial delta functions between adjacent hops identify the endpoints
on each charged cycle. Each chord kernel is bounded by a polynomial
in $1+|X|$ times $C e^{-c_1|z|^2}$ for its displacement $z$.
The one-sided time propagators are bounded on the slab.
For the second model each charged contraction has at its two ends
heat averaging with lapses $\tau_i,\tau_j$ and the bounded
positive- or negative-energy Dirac multiplier with exponential
$e^{-|t|\sqrt{|k|^2+m^2}}$. Its inverse Fourier kernel is bounded
by $C_N(1+|c_i-c_j|)^{-N}$ times a polynomial in the inverse
endpoint widths, uniformly for $|t|\le T$. This follows by integrating
by parts in spatial momentum: $m>0$ bounds the derivatives of the
mass-shell projections and the Gaussian integrates all polynomial
weights. It applies to an occupied equal-time tadpole with its
specified ordering. For a scalar line the unsmeared Euclidean
covariance is positive, has spatial integral at most $1/(2\mu)$
and at $|z|\ge1$ is bounded by $C e^{-c_\mu|z|}$ uniformly on the
slab, by the massive heat integral. Convolution with the Poisson
kernel of width $a_i+a_j$ gives a bound
$C(1+a_i^{-2}+a_j^{-2})(1+|z|)^{-4}$. Near zero use the
equal-time Fourier integral, and at large separation split the
convolution at $|z|/2$. The same type of integrable spatial
envelope holds after any fixed number of derivatives in the normal
samples. For Poisson factors, differentiation in the width is
bounded in absolute value by inverse powers of the width times
a constant multiple of the Poisson kernel; this follows from its
explicit three-dimensional density. All sample and inverse-width
powers have bounded Gaussian moments at each fixed order, uniformly
in their positions and times.

It remains to control distinct matter components joined only by the
normal state. Here is a direct clustering estimate. For two groups
of samples at spatial separation $d$, each cross covariance of $X$
is at most $C e^{-c_M d}$ uniformly over the slab. This follows by
convolving the massive scalar covariance with the two compact
packets; the equal-time packet variance controls bounded $d$.
To bound a joint cumulant of $l$ local factors, temporarily multiply
the covariance between two groups by $u\in[0,1]$, retaining the
covariances within the groups. This is a convex combination of
positive covariance matrices. At $u=0$ the mixed cumulant vanishes.
Differentiating its finite moment-partition formula differentiates a
factor on each side and supplies their cross covariance, by Gaussian
integration by parts. All differentiated moments have the bounds
just established. In applying this to scalar and spinor lines,
one may keep a separate position-space variable for each of their
two smearing kernels and perform the free contractions afterwards.
Equivalently, group the entire product on each connected matter
component as a smooth function of the samples at all its vertices.

For several components join two when a pair of their sampled centers
is within distance $d$, and use the smallest $d=d_*$ at which this
graph becomes connected. Below the last merger the components split
into two groups separated by at least $d_*$. The interpolation
therefore gives an additional factor $C_l e^{-c_M d_*}$ for their
mixed cumulant. A spanning tree of the components has $l-1$ links
with lengths at most $d_*$. Thus one may bound this exponential by
a finite sum over choices of linking vertices and trees, with an
integrable factor $e^{-c_M |c_i-c_j|/(l-1)}$ on each chosen link.
Inside each matter component the chord or field-line envelopes
already connect its vertices. At a midpoint use its offset from an
endpoint as half the corresponding displacement. Integrate all
relative coordinates on the combined tree, taking bounded envelopes
on the remaining links. This proves absolute integrability at every
fixed order, including coincident vertices. At one component the
matter envelopes alone suffice. The cumulant expansion follows by
first taking the connected matter logarithm at fixed geometry and
then taking ordinary cumulants of these conditional components.
Division by the full vacuum is thus accounted for. Translation
invariance leaves one spatial center; the fraction of its cube
translations allowed for any fixed relative configuration tends
to one. Dominated convergence gives the asserted extensive limits.
\end{proof}

This choice proves finiteness and spatial clustering for the stated
coefficients with unrestricted local depth excursions. It also gives
a diagnostic for attempts to localize the original radius: assigning
an independent normal oscillator of fixed energy to each momentum
would supply finite-energy excitations at arbitrarily large momentum.
Here all free normal, scalar, and Dirac excitations in the second
model have relativistic dispersions. Their positive energies
therefore place the intermediate states of the leading pair-production
test inside the forward cone. Equal-time spatial sampling still
selects a carrier frame in the interaction. An interacting cone,
the electric Gauss current of a relativistic charged action, and a
cutoff-free unified gauge theory do not follow from this construction.

There is also a momentum-space version of the fixed-order spatial
estimate. At a vertex of the Dirac--scalar model, before imposing
momentum conservation, each Dirac incidence gives a factor
$e^{-a^2|k|^2}$ and each scalar incidence a factor $e^{-a|p|}$.
Any finite set of inverse powers of $a$ has finite expectation.
Consequently the moment bounds for products of vertex multipliers
and their normal-sample derivatives decay to arbitrary polynomial
order in every individual incidence momentum. The local factors
are smooth in their spatial centers in mean to all orders, since
$h$ is smooth. Differentiating a center introduces spatial jets of
$X$ with the same finite moments and exponential cross-covariance
decay. Expanding the joint expectation into cumulants of local
vertices, the same interpolation proof bounds every center derivative
of a cumulant by an exponential in its spatial diameter times the
incidence-momentum envelopes. For a block of one vertex there are
only the incidence envelopes. Fourier transforming the relative
centers of each block gives rapid decay in every independent
spatial momentum transferred by the normal state, with a delta
function for total momentum in that block. This follows by repeated
integration by parts against the integrable derivative bounds.
The mass-shell multipliers in the time representation are uniformly
bounded, so every spatial loop momentum is integrable after the
linear conservation constraints are imposed. For the onsite
midpoint hop the heat symbol instead controls the mean of the two
oriented carrier momenta. Their difference is the momentum delivered
at that vertex by the normal state. The mean and difference together
control both spatial momenta. A single overall translation delta
belongs to a fully connected component. These are spatial Fourier
estimates with the canonical one-sided Dirac time orderings,
including the specified value of a one-vertex density contraction.

\paragraph{A transverse current with the relativistic free carrier.}
The local normal construction can pass the leading physical
current test with actual dynamical photons. Keep the second model of
Proposition~\ref{prop:replacement-local-normal-uv} and include the
free Maxwell field on its transverse physical space. For example,
take $B=1$ in the Dirac bilinear and replace its sine by the
matrix sine
\[
 \sin\!\left(y a\phi_a(c)\,1
       +u a^2\,\boldsymbol\Sigma\cdot{\bf B}_a(c)\right),
 \qquad {\bf B}_a=P_a(\nabla\times A).
\]
The fixed Hermitian Dirac spin matrices $\boldsymbol\Sigma$ make
this interaction invariant under spatial rotations. Its arguments
have zero length dimension. Here the spinor transforms trivially
under the local Maxwell gauge group. The Pauli interaction
therefore furnishes a magnetic response with zero total electric
charge. This is a choice of physical current, distinct from the
electric current of a charged Dirac field.

The matrix sine has norm at most one for real $y,u$. The same
Hamiltonian form and extensive transfer bounds apply upon adding
the nonnegative free Maxwell Hamiltonian. At every finite joint
order a magnetic contraction between two insertions has Fourier
multiplier
\[
 \frac{|p|^2\delta_{ij}-p_i p_j}{2|p|}
       e^{-\left(|t-s|+a_i+a_j\right)|p|}.
\]
The subscripts on $a_i,a_j$ here label the two vertices. With
$d=|t-s|+a_i+a_j$ its inverse Fourier transform equals
$(-\delta_{ij}\Delta+\partial_i\partial_j)
 [4\pi^2(|x|^2+d^2)]^{-1}$. Its magnitude is bounded by
$C(|x|^2+d^2)^{-2}$. Thus the massless field supplies an
integrable spatial tail even when the two times coincide.
A derivative with respect to a normal sample changes the estimate
only by fixed powers of the inverse widths. The Gaussian and
cumulant estimates of the proposition give the same thermodynamic
limits for these complete connected coefficients, including all
magnetic self-contractions. In the Fourier proof the field
strength also makes the soft spatial multiplier $O(|p|)$.
There is no auxiliary mass for the photon.

To test the current, vary a smooth transverse external connection
in this action. At first order in $gu$ the insertion is the curl
of the measured, heat-smeared spin density, with the derivative
acting on the external test through the prescribed spatial
readout. The current is transverse, and every pure-gauge test
gives zero. Its leading two-insertion, separated-time spectral
measure is the ordinary positive spectral measure of this
operator in the free vacuum. At zeroth order all translations
are generated by free relativistic fields, including the local
normal scalar. For a pair accompanied by any number of normal
quanta the energy at total momentum $P$ obeys
\[
 \sum_j\sqrt{|k_j|^2+m_j^2}
 \ \ge\sqrt{|P|^2+\left(\sum_j m_j\right)^2},
 \qquad \sum_j k_j=P .
\]
This is the energy bound for a sum of future-directed mass-shell
vectors. It persists on the closure of the resulting Fock spectral
subspaces. For the pair contribution it gives in particular
$E\ge\sqrt{|P|^2+4m^2}$. The leading current is nonzero:
at a nonzero external momentum its curl selects a transverse spin
matrix. The free positive and negative Dirac spin projectors at
the two momenta have a nonvanishing matrix element of that matrix
on an open set. For example at relative momentum zero take one
spinor at rest and the other at momentum $P$; the latter has
nonzero upper and lower two-spinor components. Contracting
$\boldsymbol\Sigma$ with a direction perpendicular to $P$
has an invertible two-spinor block. Projection onto the normal
vacuum supplies the expectation of $a^2 e^{-a|P|}$ times the
two positive heat multipliers, hence a strictly positive scalar
factor at these momenta. This gives a measured transverse
two-point response with positive weight on timelike pair states.

This example combines local normal fluctuations, extensive
coefficients, uncut bare graphs and the correct leading spectral
support for a physical transverse probe. Its charge form factor
at zero photon momentum vanishes. An interacting relativistic
electric theory must also produce a unit electric charge,
its Gauss constraint and the complete current of the same action.

\paragraph{An interacting cone from a positive conversion term.}
There is a further choice for which one can prove the forward cone
at nonzero coupling. Keep the massive local normal field and its
depth law above, in continuous space. Use positive-energy particle
and antiparticle spin slots $b,d$ of mass $m$, the free transverse
Maxwell field, and their common Fock vacuum with the free normal
vacuum. The free Hamiltonian and momentum $H_0,\mathbf P$ are the
sums of their ordinary relativistic second-quantized generators.
The spin-slot annihilators carry energy $\sqrt{m^2+|k|^2}$ before
they couple. They are positive-frequency one-particle fields on
the chosen spatial slice.

At a point $c$ let $C_i(c)$ annihilate a particle and an antiparticle,
each measured with the heat kernel $G_{\tau(c)}$ at that same point.
Contract the two spin indices in the three spin-triplet components
using the matrices $i\sigma_2\sigma_i$. The two measured displacements
are independent integrations. Let
\[
 Y_i(c)=L_*^{-1}\big(P_{a(c)} B_i^{(+)}\big)(c),\qquad
 O_i(c)=C_i(c)+Y_i(c) ,
\]
where $\mathbf B^{(+)}$ is the annihilation part of the transverse
magnetic curvature. Its plane-wave coefficient is proportional to
$i(p\times\epsilon_\lambda)_i/\sqrt{2|p|}$. Take the bare form
\[
 H=H_0+\kappa_{\rm c}\int_{\mathbb R^3}\!
          \mathrm d^3c\,\sum_{i=1}^3 O_i(c)^\dagger O_i(c),
 \qquad \kappa_{\rm c}\ge0 .
\]
This square includes the conversion of a photon into a pair and
the reverse conversion, pair scattering and a quadratic magnetic
term. The annihilation amplitudes are part of the definition of
the curvature block. Squaring a Hermitian field strength without
this splitting would generate a different Hamiltonian and vacuum.
The geometric smearing gives finite equal-time vacuum contractions
whenever they occur; the conversion square requires no subtraction.

\begin{proposition}[Exact interacting forward cone for a magnetic conversion action]
\label{prop:replacement-conversion-cone-uv}
The conversion form in full space defines a nonnegative
self-adjoint Hamiltonian with an exact translation-invariant
vacuum. At every $\kappa_{\rm c}\ge0$ its joint energy--momentum
spectrum lies in $E\ge|\mathbf P|$. For $\kappa_{\rm c}>0$
a variation by a transverse external magnetic field gives a
nonzero current with positive spectral weight and this same
support. The connected, spatially tested Euclidean source
coefficients at each finite conversion order are absolutely
integrable in continuous independent spatial momenta and the
remaining ordered slab times, with the actual normal measurements.
\end{proposition}
\begin{proof}
The operator taking a vector $\Psi$ to the triple of functions
$c\mapsto O_i(c)\Psi$ has a densely defined closure into spatial
$L^2$ with Hilbert-space values. For the domain one may take finite
products of smooth particle packets, transverse photon packets
with smooth compact momentum support away from zero, and
finite-particle normal test vectors. The heat-convolved matter
packets have rapid spatial decay for every $0<a\le L_*$,
uniformly in this interval. For a photon packet the same follows
by spatial Fourier integration by parts on its compact momentum
support. Thus the image of each test has finite spatial norm.
Its adjoint domain contains compact smooth functions of $c$ times
these test vectors. At fixed $c$ the two created matter packets
have product norm at most $C a^{-3}$, and the magnetic photon has
one-particle norm at most $C a^{-2}$. Integration on a compact
set is finite by Cauchy--Schwarz and the finite Gaussian moments
of every inverse width, also against finite normal test vectors.
These adjoint tests are dense. The squared norm of the closure
defines a positive closed form. Sum it with the free positive form
on their dense common domain. The resulting form is invariant under
simultaneous translations of all the fields and of the center of
integration. Its Hamiltonian commutes with spatial translations.

Every annihilation amplitude kills the free vacuum, so its energy
remains exactly zero. The one-particle relativistic dispersion and
the triangle inequality on each joint Fock momentum give
$H_0\ge|\mathbf P|$, including the photon with zero rest mass.
Adding the positive conversion form gives $H\ge|\mathbf P|$ on
its form domain. The joint spectral theorem gives the claimed
support at fixed interaction strength, including arbitrarily many
occupied slots.

Shift $Y_i$ in the external source by the real classical field
$L_*^{-1}P_{a(c)}(\nabla\times b)_i(c)$ and take the same shift
in its adjoint. One derivative measures $O_i+O_i^\dagger$
against this readout. Two derivatives also produce the contact
prescribed by the square. Pure gauge tests have zero curl.
At zero background, the current on the vacuum contains the pair
created by $C_i^\dagger$. Projecting the normal field onto its
vacuum gives the strictly positive scalar factor
\[
 \mathbb E\!\left[
 e^{-a|P|-a^2(|k_1|^2+|k_2|^2)}\right],\qquad k_1+k_2=P .
\]
For a nonzero transverse curl choose a matching triplet component
and packets on a small allowed momentum region. The matrix
element is nonzero. The full current vector has a positive
spectral measure whose support obeys the exact bound after the
photon and the pair mix in the interacting evolution. For
chronological vacuum products on their transfer domains the
half-slab norm square likewise gives reflection positivity in
carrier time.

At fixed order keep the sampled spatial kernels at every
occurrence of $C_i,Y_i$. The two massive matter contractions have
the bounded one-sided multiplier $e^{-|t-s|\sqrt{|k|^2+m^2}}$
on their allowed time orderings. A magnetic contraction is the
one-sided version of the field-strength kernel just bounded.
Each independent spatial incidence retains its heat or Poisson
factor. The normal-field moment, cumulant and spatial translation
proof of Proposition~\ref{prop:replacement-local-normal-uv}
applies to these contacts and contractions, including the source
derivatives of the square. The vacuum transfer with no external
marks equals one. The connected coefficients obey the same
spatial integrability and single-center volume bounds on finite
slabs.
\end{proof}

The exact cone concerns the interacting translations. The model
uses a positive-frequency decomposition and equal-time spatial
measurement. Lorentz transformation laws for this current and
spacelike local commutativity remain independent tests. The
magnetic current has zero total electric charge. A charge-one
Coulomb response would require additional interactions and its
own bounds.

\paragraph{Gapless photon poles at finite conversion strength.}
The same conversion action gives more than forward-cone support.
Both $2N_\gamma+N_b+N_d$ and $N_b-N_d$ commute with the
Hamiltonian: $C_i$ and $Y_i$ lower the first count by two and
preserve the second. Multiplication by the optical widths can
change the number of massive normal quanta. In the neutral
count-two subspace there is either one photon or a particle and an
antiparticle. Resolve the full-space Hamiltonian by total spatial
momentum $P\ne0$. These assertions about fibers hold at almost
every momentum of the direct integral. Let $\mathcal P_P$ project
onto the two transverse one-photon states with no normal quanta.
On its orthogonal complement the free Hamiltonian is at least
$d_0=\min(2m,M)>0$: every such state has a massive pair or a massive
normal quantum. Free soft photons of arbitrary multiplicity are
excluded in this sector by the conserved count. Acting on the
vacuum, the transverse vector potential creates precisely the
one-photon space with that count.

For either polarization the expectation of the conversion term
in a bare one-photon trial equals
\[
 \kappa_{\rm c}\frac{|P|}{2L_*^2}\,
               \mathbb E\,e^{-2a|P|} .
\]
Indeed $C_i$ annihilates this trial. Apply $Y_i(c)$ to a photon
packet, square, sum the three magnetic components and integrate
over $c$. Stationarity of the normal vacuum gives the spatial
momentum delta and the transverse norm
$\sum_i |(P\times\epsilon)_i|^2=|P|^2$. Polarization mixed terms
give the same multiple of the identity. Thus for sufficiently
small $|P|$ the two-dimensional trial space has maximum
Rayleigh quotient below $d_0$. On its orthogonal complement the
bound $H\ge H_0$ gives a Rayleigh quotient at least $d_0$.
The min--max principle gives exactly two eigenvalues below $d_0$
counting multiplicities, each isolated from the remaining
spectrum in this sector. They satisfy
\[
 |P|\le E_j(P)\le
 \left(1+\frac{\kappa_{\rm c}}{2L_*^2}\right)|P|,\qquad j=1,2 .
\]
For every unit vector $v$ in their common spectral subspace,
$d_0\|(1-\mathcal P_P)v\|^2\le\langle v,Hv\rangle\le C|P|$.
Since the two subspaces have equal finite dimension, their
orthogonal projections approach each other in norm. Thus the
transverse field sees these gapless poles with positive spectral
weight, and their combined weight relative to the free transverse
residue tends to one. Degenerate eigenvalues give a single energy
with a polarization weight matrix. This gives an exact zero-mass-gap statement at finite interaction
strength. The following estimate determines on which side of the
carrier light cone the low branches lie.

\paragraph{The photon speed in the positive magnetic square.}
The conversion Hamiltonian of
Proposition~\ref{prop:replacement-conversion-cone-uv}
has two gapless photon states. Their speed can be tested against the
same free carrier used to define the action. Work in carrier units
with light speed one. Write $p=|P|$, $L=L_*$,
$d_0=\min(2m,M)$ and $\kappa=\kappa_{\rm c}$. In the neutral
count-two fiber let $u$ be the projection of a unit vector $v$ onto
the two bare photon polarizations and write $w=v-u$. All fiber
bounds and limits below are understood up to sets of momentum
measure zero.

\begin{proposition}[Strict displacement of the low photon poles]
\label{prop:monomial-magnetic-speed-uv}
There is a finite constant $K$ determined by the free normal
preparation and the measured pair and magnetic rows such that, for
either of the two low eigenvalues and every $\kappa>0$,
\[
 \mathop{\rm ess\,liminf}_{|P|\downarrow0}
   \left(\frac{E_j(P)}{|P|}-1\right)
 \ \ge\ \frac{\kappa d_0}{2L^2(d_0+\kappa K^2)}>0 .
\]
The conclusion requires no differentiability of an eigenvalue branch.
\end{proposition}
\begin{proof}
Project the output of the annihilation row $v\mapsto (O_i(c)v)_i$
onto the common vacuum and Fourier transform the center $c$.
This gives a row $B_P$ to three magnetic components.
Since $C_i$ kills a bare photon, the two singular values of $B_P$
on the bare photon space are both
\[
 b(p)=\frac{\sqrt p}{\sqrt{2}L}\,\mathbb E[e^{-a p}] .
\]
In particular the projection retains a nonzero part of the
magnetic term at every finite $p>0$.

The matrix $B_P B_P^*$ is the Fourier transform, with the unitary
direct-integral normalization, of the translation kernel
$\langle 0|O_i(x) O_j(0)^*|0\rangle$. This kernel is spatially
integrable. At equal time the normal readouts at the two centers
commute. Conditional on their widths $a,b\le L$, the two
matter contractions give a product of Gaussian heat densities with
variance parameter $a^2+b^2$. Its absolute expectation is bounded by
$C\mathbb E[a^{-6}]\exp[-|x|^2/(4L^2)]$. For the magnetic
contraction, differentiating the Poisson-smeared photon kernel gives
a bound $C L^{-2}(|x|^2+(a+b)^2)^{-2}$. The latter has a bounded
expectation at the origin by the inverse-width moments and decays
as $|x|^{-4}$ at infinity. The mixed terms vanish in the free
vacuum. Thus $\|B_P\|\le K$ uniformly; in particular this estimate
does not require a bound on the full, unprojected annihilation row.

For $0<p<d_0$, positivity of the full square and the free gap off
the bare photon subspace give
\begin{align*}
 \langle v,(H-p)v\rangle
 &\ge (d_0-p)\|w\|^2
   +\kappa\big[b(p)\|u\|-K\|w\|\big]_+^2 \\
 &\ge \frac{\kappa(d_0-p)}{d_0-p+\kappa K^2}
     \,b(p)^2\|u\|^2 .
\end{align*}
Here $[x]_+=\max(x,0)$, and the second inequality minimizes the
first line over $\|w\|\ge0$ with $\|u\|$ fixed.
For a low eigenvector the bare photon trial bound gives
$E_j(P)\le (1+\kappa/(2L^2))p$ and $d_0\|w\|^2\le E_j(P)$.
Consequently $\|u\|^2\to1$. Since $b(p)^2/p\to1/(2L^2)$,
division by $p$ proves the bound.
\end{proof}
In four-dimensional natural units $\kappa$ has length dimension
two, since $C_i$ and $L^{-1}B_i^{(+)}$ both have mass dimension
three. Thus even the choice $\kappa=L^2$, $L=\sqrt{G_N}$ gives a
strict shift relative to the prescribed carrier cone. Assigning the
coupling its geometric dimensions leaves the speed test in force.

\paragraph{An invariant-mass square with an exact massless pole.}
One can fix the dispersion defect and separately evaluate the
cost of closing the source. Consider one transverse response
component with internal space $\mathbb C\oplus
L^2(\mathbb R^3,\mathrm d^3k)$. Set $m=L^{-1}$,
$D(k)=4(m^2+|k|^2)$ and $h=f/\|f\|_2$, where the normal
displacement preparation is
\[
 f(k)=\int_0^1 5r^4\exp[-L^2r^{2/3}|k|^2]\,\mathrm dr .
\]
The change of variable $s=L^2r^{2/3}|k|^2$ gives
$f(k)\le C(1+L|k|)^{-15}$, so $Dh\in L^2$. Choose $L=\sqrt{G_N}$
in natural units. Use a real marker $z$ to count the mixing and
fix $z=1$ for the model. The geometric measure, shape and mass
ratio in this comparison are specified choices. Its invariant-mass
form consists of one square,
\[
 \begin{gathered}
 \|B_z(u,\psi)\|^2=\|D^{1/2}(\psi+z h u)\|^2,\\
 M_z^2=B_z^* B_z,\qquad H_z(P)=(p^2+M_z^2)^{1/2}.
 \end{gathered}
\]
Here $B_z$ is closed on $\mathbb C\oplus\mathcal D(D^{1/2})$.
Its kernel is spanned by $(1,-zh)$, while
$B_zB_z^*=D+z^2|D^{1/2}h\rangle\langle D^{1/2}h|\ge4m^2$.
The positive spectra of the two products agree by polar
decomposition. Thus the transverse component sees an isolated
zero of invariant mass, has dispersion $E=p$ on that state, and
assigns it residue $(1+z^2)^{-1}$. The complementary spectral
weight $z^2/(1+z^2)$ is positive at nonzero mixing.

To test sewing, treat this quadratic block as a positive Gaussian
field with its own four-dimensional center. At $X=|q_E|^2>0$ the
difference $M_z^2-M_0^2$ is finite rank and the relative Fredholm
determinant is unambiguous. Taking the Schur complement of $X+D$
in the quadratic block gives
\begin{align}
 \det\big[(X+M_z^2)(X+M_0^2)^{-1}\big]
  &=1+z^2\langle h,D(X+D)^{-1}h\rangle .
 \label{eq:monomial-mass-sew-uv}
\end{align}
In fact the upper diagonal contact is $z^2\langle h,Dh\rangle$
and the exchange is $-z^2\langle Dh,(X+D)^{-1}Dh\rangle$.
Their sum is $z^2 X\langle h,D(X+D)^{-1}h\rangle$ before
division by the free entry $X$. These terms are fixed together by
the original square. The entire internal integral is finite.
By monotone convergence,
$X\langle h,D(X+D)^{-1}h\rangle\to\langle h,Dh\rangle>0$.
For $z\ne0$ the log determinant therefore has leading term
$z^2\langle h,Dh\rangle/X$. Its four-dimensional center integral
diverges quadratically, already at the coefficient of $z^2$.
This calculation uses the complete contact and exchange for that
coefficient and the ordinary Gaussian closure. The exact null
dispersion and positive spectral measure solve two tests on this
channel. Its unsuppressed center momentum and its absence of an
electric Gauss source preclude identifying it with the unrestricted
charged action.

The same failure survives normalization by the zero-source vacuum
when the mixing itself is probed. Replace the marker by a smooth
real source $t\eta(x)$ of compact support in the center and keep
both kinetic terms. Put $X=-\Delta_E$, let $M_\eta$ denote
multiplication by $\eta$, and set
$W(y)=y\langle h,D(y+D)^{-1}h\rangle$. Integrating the massive
block gives the exact scalar Hessian
$X+t^2 M_\eta W(X)M_\eta$. The order-$t^2$ coefficient of the
relative half log determinant has nonnegative momentum density
proportional to
\[
 \int\!\mathrm d^4k\,\mathrm d^4l\;
  \frac{|\widehat\eta(k-l)|^2\,W(|l|^2)}{|k|^2}.
\]
For each $\eta\ne0$ a bounded set of differences $k-l$ carries
positive Fourier weight. On that set $W(|l|^2)$ is uniformly
bounded below at large $|k|$ because its limit is
$\langle h,Dh\rangle>0$. Thus this connected mixing response
also diverges. It is a source test of the mass block itself.

\subsection{Invariant mass, physical positivity and normal transfer}
\label{sec:carrier-mass-tests-uv}

\paragraph{Moving the geometrical interaction into the invariant mass.}
There is a way to retain exact boosts in an interacting comparison.
Instead of applying a spatial heat hop in one laboratory frame, apply
it to relative coordinates in the rest frame of a massive channel.
The following test specifies the action and the sewing. It concerns
a scattering mass operator and does not identify its states with the
charged determinant of O--QED.

In each sector of fixed total charge take spinless channels with
$n\ge2$ equal-mass constituents of mass $m>0$. In their rest frame
use independent Jacobi momenta $k_1,\ldots,k_{n-1}$ with Lebesgue
measure. The physical constituent momenta $\pi_i$ are the corresponding
linear combinations with $\sum_i\pi_i=0$. The free squared mass $D_0$
is multiplication on each channel by
$(\sum_i\sqrt{m^2+|\pi_i|^2})^2$. Integrate one independent normal
displacement for each relative coordinate. For example take the metric
volume $5r^4\,\mathrm dr$ on $0<r<1$ with the amended heat lapse
$L_*^2r^{2/3}$. The characteristic function of one displacement is
\[
 f(k)=\int_0^1 5r^4
               e^{-L_*^2r^{2/3}|k|^2}\,\mathrm dr .
\]
It obeys $0<f(k)\le C(1+L_*|k|)^{-15}$, by the substitution
$u=L_*^2|k|^2r^{2/3}$. Thus $\widehat f=f/\|f\|_{L^2}$ is a
unit vector in the relative momentum space. Use
$\chi_n=c_n\widehat f^{\otimes(n-1)}$, with
$\sum_n|c_n|^2=1$, in the channel direct sum. Two nonzero coefficients
suffice for a transition changing the constituent count. For identical
bosons one may first take the finite permutation average and normalize
the nonzero result in each channel. The following argument needs only
rotation invariance and the stated norms.

Take $D=D_0+\lambda|\chi\rangle\langle\chi|$, $\lambda\ge0$ in
squared mass units. At total momentum $\mathbf P$ the bare quadratic
form in the internal wave function defines the mass-shell equation
$(P^2-D)\Psi=0$ and the positive-energy evolution
$H(\mathbf P)=(|\mathbf P|^2+D)^{1/2}$. The interaction contributes
$\lambda|\langle\chi,\Psi\rangle|^2$ to the internal mass form.
Its momentum kernel therefore contains two independently evaluated
normal preparations, one at each side of the interaction. Collapsing
the two preparations into the difference of their endpoint momenta
would change the operator.

\begin{proposition}[A covariant mass interaction with finite uncut scattering integrals]
\label{prop:carrier-mass-uv}
This action has a positive-energy unitary Poincar\'e representation
at every $\lambda\ge0$. Each finite Born coefficient of its internal
transition operator is absolutely UV finite, including all channel
and relative-momentum integrations. The subthreshold interaction and
its resolvent are given by a single finite integral. A separated
spectator has an explicit recoil-dependent embedding with the same
UV control.
\end{proposition}
\begin{proof}
The internal operator is self-adjoint on the domain of $D_0$ by
bounded perturbation, commutes with rotations, and $D\ge4m^2$.
Decompose into its spectral masses $\sqrt{\sigma}$ and spins, and
on each such fiber use the ordinary massive unitary representation
with invariant measure $\mathrm d^3P/(2\sqrt{|\mathbf P|^2+\sigma})$.
The boosts are unitary between the mass-shell fibers and the joint
spectrum is in the future cone. This construction puts the interaction
into the invariant mass before defining boosts.

For real $\zeta<4m^2$ let $R_0=(D_0-\zeta)^{-1}$ and
$B(\zeta)=\langle\chi,R_0\chi\rangle$. Every internal integral in a
Born chain is an independent copy of $B$. Its absolute integral,
summed over all channels, is bounded by
$\|\chi\|^2/(4m^2-\zeta)$. Multiplication verifies
\[
 (D-\zeta)^{-1}
 =R_0-\frac{\lambda R_0|\chi\rangle\langle\chi|R_0}
                   {1+\lambda B(\zeta)} .
\]
For a nonreal mass parameter the corresponding absolute bound uses
$1/|\operatorname{Im}\zeta|$. At a real scattering energy the
large-momentum integrals stay absolutely convergent. On its
energy shell the usual boundary values are distributions. For
example in the distinguishable two-body channel
$D_0=4(m^2+|k|^2)$, and
$\operatorname{Im}B(s+i0)=\pi\langle\chi,\delta(D_0-s)\chi\rangle$
at regular values. This is strictly positive above its threshold
when $c_2\ne0$. The transition factor
$\lambda/(1+\lambda B)$ has the corresponding nonzero absorptive
part for positive $\lambda$. No ultraviolet regulator or subtraction
is involved. The scalar internal resolvent variable $\zeta$ does not
replace the vectorial pole prescription of a charged carrier on its
locked sector plane.

For a free spectator of mass $m_3>0$ and recoil $q$, the tensor
product of the interacting pair and the spectator has rest energy
$\sqrt{D+|q|^2}+E_3(q)$ where $E_3(q)=\sqrt{m_3^2+|q|^2}$.
Its square differs from the naive free three-body squared mass plus
$\lambda|\chi\rangle\langle\chi|$ by
\[
 2E_3(q)\left(\sqrt{D+|q|^2}-\sqrt{D_0+|q|^2}\right).
\]
This term is fixed by cluster composition. Its magnitude can also
be controlled. Set $R_t=(D_0+|q|^2+t)^{-1}$. The positive
resolvent representation of the square root gives the difference
of roots as
\[
 \frac{\lambda}{\pi}\int_0^\infty
 \frac{t^{1/2}R_t|\chi\rangle\langle\chi|R_t}
 {1+\lambda\langle\chi,R_t\chi\rangle}\,\mathrm dt .
\]
It follows by subtracting the scalar square-root integral using
spectral calculus and the preceding rank-one identity. The integral
is positive and its trace is bounded by
$\lambda\|\chi\|^2/(2\sqrt{4m^2+|q|^2})$. Thus multiplication by
$2E_3(q)$ leaves a uniform bound in spectator momentum. The
spectator carries a diagonal momentum delta and its prepared packet
has its ordinary squared norm. This proves the exact embedding for
one separated partition. Combining several interacting partitions
requires a consistent set of embeddings and currents, a known issue
in relativistic few-body dynamics \cite{CarrierClusters}.
\end{proof}

The mass form produces nontrivial transitions, positive spectra and
cut weights with simultaneous integrations over independent geometrical
displacements. It supplies no local electric Gauss generator or crossed
vacuum source for the original constituents. Promoting its relative
channel amplitudes to elementary fields would introduce additional
closed contractions and a different theorem.

\paragraph{An exact reflection test for isotropic sampling.}
Can one restore plane covariance simply by sampling in all four
Euclidean directions with the same geometric displacement law?
Take a real scalar of mass $m>0$ on the Euclidean continuation of
one locked plane. Its boundary field $\phi$ has free covariance
$C=(-\Delta_E+m^2)^{-1}$. Give each incidence an independent
isotropic Brownian displacement of positive lapse $\tau$. Its
density $h(x)=(4\pi\tau)^{-2}e^{-|x|^2/(4\tau)}$ determines the
interaction field $Y=h*\phi$. Use the ordinary boundary-field
source, and put
\[
 V=\int\chi(x)Y(x)^4\,\mathrm d^4x,\qquad
 \mathrm d\mu_\lambda=Z_\lambda^{-1}e^{-\lambda V}\,\mathrm d\mu_0 .
\]
The nonzero switch $\chi$ is smooth, nonnegative, compactly
supported and even under the Euclidean carrier-time reflection
$\Theta$. It may be radial. The exact bare measure exists for
every $\lambda\ge0$ since $V$ is finite almost surely. At any
finite order each internal scalar contraction has the integrable
momentum density $e^{-2\tau p^2}/(p^2+m^2)$. The switch has a
bounded Fourier transform, so all internal momenta of switched
graphs can be integrated absolutely. External boundary fields may
be smeared against Schwartz tests. Thus the source passes the
Euclidean UV test without a subtraction.

\begin{proposition}[Failure of the interacting reflected norm at every positive coupling]
\label{prop:carrier-reflect-uv}
For each $\lambda>0$ there is a bounded observable of the boundary
field at positive times whose squared reflected norm in
$\mu_\lambda$ is negative. This holds when only $\phi$ is declared
a physical source. The test uses the complete measure at the chosen
coupling.
\end{proposition}
\begin{proof}
Let $\xi=\phi(0,\cdot)$ denote the distributional time-zero field.
Conditional on $\xi$, the positive half and the reflected negative
half are independent copies. Indeed, at spatial momentum $k$ the
free time kernel is $e^{-\omega|t-s|}/(2\omega)$ for
$\omega=(|k|^2+m^2)^{1/2}$. Subtracting the covariance of the
conditional means $e^{-t\omega}\xi$ at the two positive time
distances makes the cross covariance zero. The residual on each
half has the massive Dirichlet covariance. Time-zero smearings
are limits in mean square of positive-time smearings.

To construct the test, take a compact reflection-invariant set $D$
containing the switch support. The contribution of the positive
half to $Y(x)$, $x\in D$, is a continuous Gaussian random function.
It can be defined by smearing $\phi$ against
$\mathbf{1}_{z_0>0}h(x-z)$; the covariance is bounded on $L^2$, and
the same bounds for all $x$ derivatives give a continuous version
on $D$. Decompose this contribution into its conditional mean
$M_\xi$ and an independent centered Dirichlet part $U$. The
reflected half has an independent copy of $U$ given $\xi$ and
the reflected copy of $M_\xi$. The centered Gaussian $M_\xi$ has
positive probability of lying in any uniform neighborhood of zero.
These small-ball statements can be read in $C(D)$: the support of
a centered Gaussian contains zero and every Cameron--Martin shift.

Choose smooth nonnegative, nonzero $u$ supported strictly in the
positive half and put $v=h*u$ on $D$. Translating the Dirichlet
half field by $u$ has finite squared Cameron--Martin norm
$\int_{t>0}(|\nabla u|^2+m^2u^2)$. The uniform neighborhoods of
$v$ and $-v$ for $U$ therefore have positive and equal probabilities.
For opposite halves with respective Dirichlet shifts $\sigma u$ and
$\sigma' u$, $\sigma,\sigma'=\pm1$, the full sampled field near zero
boundary mean is close to $\sigma v(x)+\sigma' v(\Theta x)$.
The two ideal diagonal histories yield the same value $V_+$.
The off-diagonal histories have value $V_-$, where
\[
 V_+-V_-=
 8\int\chi(x)v(x)v(\Theta x)
        [v(x)^2+v(\Theta x)^2]\,\mathrm d^4x>0 .
\]
Strictness uses positivity of the heat density on both halves.

Now let the test be the indicator of $\|M_\xi\|_\infty<\epsilon$
times the difference of the two neighborhood indicators for $U$.
For sufficiently small $\epsilon$ and sufficiently small disjoint
neighborhoods, every diagonal weight is bounded above by
$e^{-\lambda(V_+-\delta)}$ and every off-diagonal weight below by
$e^{-\lambda(V_-+\delta)}$, with $2\delta<V_+-V_-$. The conditional
free probabilities of the four pairs of events are the same.
The two positive diagonal contributions to the reflected norm
are strictly smaller than the magnitudes of the two negative
off-diagonal contributions. Averaging over the positive-probability
boundary event and dividing by $Z_\lambda>0$ preserves the sign.
All variables defining the test are measurable on the positive half.
Bounded cylinder approximants from strictly positive times suffice
as well, by $L^2(\mu_0)$ approximation and $e^{-\lambda V}\le1$.
\end{proof}

This gives a design test before expanding a proposed Euclidean
action. For continuous cross-seam weight
$K(u,v)=e^{-\lambda V(u,v)}$ at a common boundary value, reflection
positivity requires its two-history minors to be positive
semidefinite on configurations in the conditional support, wherever
the conditional test can be localized as above. In particular
$V(u,u)+V(v,v)\le2V(u,v)$ for a real symmetric interaction.
Terms additive in the two halves cancel from this inequality.
A sufficient construction at a seam is a positive kernel of the
form $K(u,v)=e^{-A(u)-A(v)}\exp[\langle F(u),F(v)\rangle]$
with real $F$, provided the resulting functional is integrable:
the exponential is a sum of Gram kernels of tensor powers with
positive coefficients. A homogeneous covariant action has to pass
the reflection test for its full observable algebra.

\paragraph{Local relativistic propagation between separated normal profiles.}
There is a way to keep positivity and exponential suppression in
the same relativistic calculation. It also identifies which source
derivatives lose the suppression. Introduce one noncompact
spacelike normal coordinate $y$ and use the free positive scalar
action
\[
 S_0=\tfrac12\int\mathrm d^4x\,\mathrm dy\,
 \big(|\partial_x\Phi|^2+|\partial_y\Phi|^2+M^2\Phi^2\big),
 \qquad M>0.
\]
Read out two real smooth unit $L^2$ profiles $h_1,h_2$ with compact
supports separated by normal distance $d>0$, and take
$\Phi_i(x)=\int h_i(y)\Phi(x,y)\,\mathrm dy$. Timelike
translations act on the usual sector-plane carrier. At normal
momentum $\eta$, the positive-energy mass is
$\sqrt{M^2+\eta^2}$. The profile matrix of physical spectral
measures is positive: its integrand is the Gram matrix of
$\widehat h_i(\eta)$ with measure $\mathrm d\eta/(2\pi)$.
Consequently the joint free source has the relativistic spectrum
on the plane. Reflection positivity also follows mode by mode
from the massive scalar transfer. The spacelike normal integration
does not alter the sector-plane orientation or the directed pole.

\begin{proposition}[Which quadratic responses retain a normal crossing?]
\label{prop:carrier-bridge-uv}
Couple classical sources $J_i$ by the quadratic vertices
$\tfrac12 J_i(x)\Phi_i(x)^2$. Every connected Gaussian coefficient
containing both types of source is absolutely finite at any finite
number of insertions and fixed external Euclidean momenta, after
factoring the one overall momentum-conservation delta. At second
order in a single $J_i$ the connected response is logarithmically
UV divergent. In each statement the vacuum disconnected terms have
already been removed.
\end{proposition}
\begin{proof}
Writing $\omega=\sqrt{|p|^2+M^2}$, the free normal Green function
gives
\[
 C_{ij}(p)=\int h_i(y)h_j(z)
   \frac{e^{-\omega|y-z|}}{2\omega}\,\mathrm dy\,\mathrm dz .
\]
This expression follows by Fourier inversion of
$(\eta^2+\omega^2)^{-1}$. In particular
$|C_{12}|\le\|h_1\|_1\|h_2\|_1e^{-d\omega}/(2\omega)$.
The Fourier representation and Plancherel also give
$0<C_{ii}\le\omega^{-2}$ and
$\omega^2 C_{ii}\to1$ by dominated convergence. The latter limit
counts all of the normal spectral weight on the individual profile.

For two or more quadratic insertions a connected Wick contraction
is one closed cycle. Each internal plane momentum is $k+q_j$,
where the finitely many shifts are fixed by the external momenta.
A cycle visiting both supports changes its profile at least twice.
In the region $|k|>2\max_j|q_j|$ its integrand therefore has an
exponential upper bound; on the remaining compact set $M>0$
bounds every propagator. Inserting the Fourier representation also
causes no hidden divergent normal trace. One may first perform the
exact geometric normal propagation to obtain $C_{ij}$. Alternatively,
taking absolute values before normal integration gives
$\int|\widehat h_i\widehat h_j|/(\eta^2+\omega^2)\,
\mathrm d\eta/(2\pi)\le\omega^{-2}$, which suffices by itself
for cycles of length at least three. At length two the stated
exponential estimate uses the exact normal integrations first,
as is appropriate for a geo-phase graph.

For two equal labels the positive integrand is instead
$C_{ii}(k)C_{ii}(k+q)$. Its ratio to $|k|^{-4}$ tends to one,
so each sufficiently high dyadic shell contributes a fixed
positive amount. The coefficient is a derivative with respect
to the same physical source at both ends. Discarding it would
change the measured observable.
\end{proof}

This example has genuine forward-cone propagation and a geometric
high-momentum bound on all transmitted cycles. Smooth profiles in
the normal coordinate leave the short-distance self response on
each profile intact. A full interacting current would have to
control that response as well as the transmitted one, and to retain
the same control when the external fields are themselves sewn.

\paragraph{Even an intrinsically separated quadratic vertex has a self response.}
One might try to make the normal crossing unavoidable by choosing
a single source coupled to $\Phi_1\Phi_2$ instead. Its two-point
contraction contains both $C_{12}(k)C_{21}(k+q)$ and
$C_{11}(k)C_{22}(k+q)$. The second term has the same logarithmic
tail as above. This conclusion is independent of the choice of
normal coordinates or profiles. Let $H_\perp\ge0$ be any
self-adjoint normal Hamiltonian with a positive Hilbert space and
$C(k)=(|k|^2+M^2+H_\perp)^{-1}$. Couple a real quadratic
source through any nonzero bounded self-adjoint operator $T$ on
this space, with a real structure when the fields are real.
Apart from its positive Wick multiplicity, the complete
two-insertion loop has integrand
\[
 \left\|C(k)^{1/2}T C(k+q)^{1/2}\right\|_{\rm HS}^2 .
\]
If this norm is infinite the ordinary normal trace already fails.
Otherwise the spectral theorem gives
$|k|C(k)^{1/2}\to I$ and $|k|C(k+q)^{1/2}\to I$ strongly
with bounded norms for large $|k|$. For Hilbert--Schmidt $T$,
approximating it in that norm by finite-rank operators yields
\[
 \lim_{|k|\to\infty}|k|^4
 \left\|C(k)^{1/2}T C(k+q)^{1/2}\right\|_{\rm HS}^2
       =\|T\|_{\rm HS}^2>0 .
\]
For any other bounded nonzero $T$, a single matrix element
between unit vectors with $\langle f,Tg\rangle\ne0$ gives a
strictly positive lower limit for the rescaled squared norm.
Thus no nonzero plane-momentum-independent normal bilinear of
this type can make its own full second response UV finite.
The hypothesis is the stated positive quadratic carrier and a
source linear in $T$. A momentum-dependent current or additional
source contacts generated by a different action are separate cases.

\paragraph{A relativistic normal current with a finite complete holonomy response.}
The normal connection itself provides a source for which the contact
term and the oriented histories can be evaluated together. Use one
spacelike normal circle of circumference $L$ as an internal comparison
geometry, independent of the locked GT directions. On
$\mathbb R^4\times S^1_L$ give a complex scalar of mass $m>0$ the bare Euclidean
action. Set $D_y=\partial_y+i\vartheta(x)/L$. Then
\[
 S[\Phi,\vartheta]=\int\mathrm d^4x\,\mathrm dy\,
 \left(|\partial_x\Phi|^2+
       |D_y\Phi|^2+m^2|\Phi|^2\right).
\]
Here $\vartheta$ is the dimensionless normal holonomy in the chart
with no $y$ dependence and with the other components of the
background connection zero. A constant $\vartheta$ is flat and
periodic modulo $2\pi$. At this background every normal Fourier
mode has positive carrier energy
\[
 E_n(\mathbf p)=\sqrt{|\mathbf p|^2+m^2+
                         (2\pi n+\vartheta)^2/L^2}.
\]
Thus the same source has a relativistic free spectrum, including
its charged modes at arbitrarily large spatial momentum. The
finite potential of a flat Wilson line at one matter loop also
occurs in gauge--Higgs models
(\href{https://doi.org/10.1093/ptep/ptv153}{Hosotani and Yamatsu, 2015}).
The following calculation tests both the flat potential and the
first local variation of this particular Gaussian action.

\begin{proposition}[Flat holonomy and its first nonconstant perturbation]
\label{prop:carrier-holonomy-uv}
Per unit carrier four-volume the complete complex-scalar determinant
relative to its zero-source value is UV finite for every constant
real $\vartheta$. It has finite derivatives of all orders in
$\vartheta$. At any nonzero Euclidean external momentum $q$ the
bare two-insertion response for $\vartheta(x)$ instead diverges.
\end{proposition}
\begin{proof}
Write $\omega_p=\sqrt{|p|^2+m^2}$. At fixed carrier momentum the
normal heat trace at proper time $s>0$ is
\[
 \sum_n e^{-s(2\pi n+\vartheta)^2/L^2}
  =\frac{L}{\sqrt{4\pi s}}\sum_{w\in\mathbb Z}
            e^{-w^2L^2/(4s)}e^{iw\vartheta}.
\]
This is Poisson summation of the free normal heat law. The history
of winding $w$ has phase $e^{iw\vartheta}$. In the relative
logarithm the $w=0$ contribution vanishes identically since a
flat connection has trivial phase on every contractible return.
Pair opposite windings. The proper-time formula for
$\Gamma=\operatorname{Tr}\log(-D^2+m^2)$ gives
\[
 \frac{\Gamma(\vartheta)-\Gamma(0)}{\mathrm{Vol}_4}
  =2\int\frac{\mathrm d^4p}{(2\pi)^4}
       \sum_{w=1}^{\infty}\frac{1-\cos(w\vartheta)}{w}
                   e^{-wL\omega_p}.
\]
Indeed the heat-time integral of
$s^{-3/2}e^{-s\omega_p^2-w^2L^2/(4s)}$ is
$2\sqrt\pi\,e^{-wL\omega_p}/(wL)$. The paired integral is
nonnegative at real holonomy. Each fixed derivative costs at most
a fixed power of $w$, while
$\omega_p\ge(m+|p|)/2$. Thus winding and unbounded carrier
momentum are jointly absolutely integrable after every finite
differentiation. Only the field-independent vacuum was divided out.
This is the complete one-loop Gaussian determinant at all source
orders, including both derivatives of the covariant normal kinetic
term. It involves no expansion in a truncated set of winding sectors.

For the nonconstant test let $H(q)$ denote the two-source Hessian
at zero background with its overall momentum delta removed, and set
$k_n=2\pi n/L$, $M_n^2=m^2+k_n^2$. The two linear vertices on
mode $n$ each supply $2k_n/L$. The second derivative of the action
supplies the local contact $2/L^2$. Its contribution is independent
of $q$. Use the proper-time formula for both source derivatives
at each $s>0$, combine $H(q)-H(0)$ there, and only then integrate
in $s$. The contact cancels from this comparison. Completing the
Gaussian carrier square in the two-vertex term, put $b=u(1-u)|q|^2$
and $c_L=4/[L^2(4\pi)^2]$. This gives
\[
 \begin{aligned}
 H(q)-H(0)
  &=c_L\sum_n k_n^2
      \int_0^1\!\mathrm du\int_0^\infty
      \frac{\mathrm ds}{s}\,e^{-s M_n^2}
                       (1-e^{-s b})\\
  &=c_L
      \sum_n k_n^2\int_0^1\!\mathrm du\,
       \log(1+b/M_n^2).
 \end{aligned}
\]
All terms in this last comparison have the same sign. At fixed
$q\ne0$ the summand including the prefactor tends to
$|q|^2/(24\pi^2L^2)$ as $|n|\to\infty$, using
$\int_0^1u(1-u)\,\mathrm du=1/6$. The zero-momentum Hessian
itself was just proved finite using its actual contact. Consequently
the local response diverges. The comparison diagnoses a divergence;
it does not subtract a source-dependent term from the bare action.
\end{proof}

A varying $\vartheta$ has field strength $F_{\mu y}=\partial_\mu
\vartheta/L$. A contractible virtual path can now measure nonzero
flux even at zero net winding. Thus exact relativistic propagation
and the complete flat-source determinant can coexist with
geometric UV suppression. The local current tests a curvature
component to which flat-source measurements are insensitive.

\subsection{The ambient curvature action and its charged source}
\label{sec:monomial-selection-uv}

\paragraph{The order of integration in the ambient action.}
Take the full curvature monomial in \eqref{eq:uvpeA-curvature} and
keep a common raw configuration $q$. On fixed regular charts the
GTL coarea and axial shell integration give the identity
\begin{equation}
\begin{aligned}
 S_{*,\mathrm{sc}}[q]={}&\mathcal N(G_N)\sum_{\bin}
 \int |dG_{0\bin}|\,d^3\mathbf x
 \int_{B^{K-1}} \frac{dy}{|B^{K-1}|}\int_{\mathcal I_r}dr\\
 &{}\times
 \mathcal C_{\bin}[q;\iota_{\bin}(G_{0\bin}),\mathbf x,v(y),r],\\
 v(y)={}&R\big(y+\sigma_\epsilon\sqrt{1-|y|^2}\,e_\Gamma\big).
\end{aligned}
\label{eq:monomial-owner-uv}
\end{equation}
Here $y$ is a dimensionless transverse velocity and
$\iota_{\bin}(G_{0\bin})=(G_{0\bin}+b_s)\hat G_{\bin}$.
The equality is the explicit pushforward of the two delta
integrals. It uses no restriction on an internal charged
momentum. The auxiliary interval and circle treated below are
comparison problems for this source.

Even the coefficient of the Coulomb force tests the one-parameter
claim. Normalize the compact electromagnetic generator to act on a
primitive charge-one state. Write the projected Maxwell density as
$Z_A F^2/4$ and the covariant derivative as $\partial+iA$ in that
charge convention. Normalizing the photon gives $e_{\rm can}^2=1/Z_A$.
A change in the pairing on this curvature component changes $Z_A$
at fixed $G_N$. A common rescaling of $A$ also changes the unit
charge in the derivative. Hence both the gravitational normalization
and the trace in that charge convention have to be derived on the
same branch. With $\hbar=1$ and unit invariant speed magnitude, setting $L=\sqrt{G_N}$ supplies
a length. Dimensionless quantities such as $mL$, the shape of the
normal state and the relative trace weights still require a law.
In the examples their numerical choices specify the model, and
source markers count insertions at fixed model data.

There is a finite first test on an ordinary perturbative expansion.
Expand at a stationary background, put the full quadratic part in
the propagators and count vertices of valence $n_v\ge3$. A connected
graph with $E$ external legs and $b$ loops satisfies
$\sum_v(n_v-2)=2b+E-2$, by counting half-edges and using
$b=I-V+1$. For a one-loop two-photon response and for its two-loop
vacuum closure the right side is two. The candidates are therefore
one quartic vertex or two cubic vertices. If sources enter by
shifting the same connection, this tests the cubic and quartic
variations of the action, including the direct two-photon contact.
An explicitly loop-dependent Jacobian, determinant or boundary
vertex has to be counted at its actual order. Treating it as an
ordinary bare vertex after eliminating fields can change the count.
A positive band in an exact current does not by itself imply a
positive band in each perturbative coefficient.

\paragraph{At which stage is the velocity shell integrated?}
Suppose a charged variable is shared by every velocity in the bare
curvature integral. On a finite set of physical test modes with a
fixed Berezin measure, write $H_v(z)$ for its invertible quadratic
matrix at a source value $z$. The charged Gaussian integral then
gives $\det H(z)$ for $H(z)=\int H_v(z)\,d\nu_\sigma^{[\Gamma]}(v)$.
In particular the quadratic coefficient of its logarithm uses
\[
 \partial_z^2\log\det H
  =\operatorname{tr}\!\left(H^{-1}H''
                    -H^{-1}H'H^{-1}H'\right).
\]
This follows by doing the shell integration in the action at a fixed
charged field and differentiating the resulting determinant.
An independent field for each shell label would have a different
kinetic space. Even in a finite direct sum its trace logarithm
sums the logarithms of its blocks.

There is an exact sign test if $H_v(0)=H_0>0$ on the same finite
test space and all $H_v'(0)$ are Hermitian. Put
$T_v=H_0^{-1/2}H_v'(0)H_0^{-1/2}$ and let the bar denote that
shell average. Direct subtraction of the two second derivatives
gives
\[
 \left.\partial_z^2
   \big(\log\det\overline H-\overline{\log\det H_v}\big)
 \right|_{0}
   = \int \operatorname{tr}\big[(T_v-\overline T)^2\big]\,d\nu .
\]
The direct second variations cancel in this difference and the
remaining term is the sum of the variances of all matrix entries
in the Hilbert--Schmidt norm. This identity concerns the stated
positive test block, whereas a graded determinant and a moving
boundary density require their own variations.

A calculation on a positive even test matrix displays the size of
this distinction. For $K\ge2$, take a transverse unit $e_j$ and
$h(v)=v\cdot e_j/R_{\mathcal G}$. Let $H_v(z)=e^{z h(v)}H_0$ with
$H_0>0$ of dimension $m$. The transverse-ball measure gives
$\int h\,d\nu=0$ and $\int h^2\,d\nu=1/(K+1)$. The second
logarithmic derivative at zero of the common-field determinant is
$m/(K+1)$. By contrast $\int\log\det H_v(z)\,d\nu(v)$ is constant
in $z$. The mark here tests a moving pairing with a fixed shell
chart. This difference occurs before a continuum momentum integral
is considered. In the physical source, the raw representation and
boundary field space must say which variables are shared, and which
normal or shell labels propagate through a contraction.

\paragraph{A two-source contact from the oriented shell.}
The full velocity shell can be differentiated without choosing an
axial representative. Work on one aligned sector-plane and put
$R=R_{\mathcal G}>0$, $\omega=v/R$, $e=e_\Gamma$,
$\sigma=\sigma_\epsilon$, and $V_{K-1}=\operatorname{vol}(B^{K-1}_1)$,
with $K\ge2$. The radial Dirac delta in the scalar shell measure
gives $d\nu=(\sigma e\cdot\omega)_+\,dS(\omega)/V_{K-1}$,
where $x_+=\max(x,0)$. All other traces, densities and the
single coordinate GTL pullback can be kept in the marked test.
This calculation differentiates the scalar shell factor. Directed
action lifts and the Lorentzian boundary value still belong to their
own factors in the same source.

In a fixed reference chart consider a twice differentiable family
of unit axes $e(z)$ and smooth marked weights $f_z(\omega)$.
Define $h=\sigma e(z)\cdot\omega$. For two commuting source
derivatives, subscripts denote differentiation at the reference
background. The mixed derivative of $I=\int f_z\,d\nu_z$ obeys
\begin{equation}
\begin{aligned}
 V_{K-1}I_{ij}={}&\int_{S^{K-1}}
 \bigl[h_+ f_{ij}+\mathbf1_{h>0}
 (h_i f_j+h_j f_i+h_{ij} f)\bigr]\,dS\\
 &+\int_{h=0} h_i h_j f\,dS_{\rm eq}.
\end{aligned}
\label{eq:nativefull-shell-contact-uv}
\end{equation}
The second derivative of the positive part contains
$\delta(h)h_i h_j$. On its equator $|\nabla_{S^{K-1}}h|=1$,
giving the last integral by coarea. For $K=2$ that integral uses
counting measure on two points. With a fixed raw axis and a
probe that varies only the matter insertion, $h_i=h_j=h_{ij}=0$.
If the axis and every mark undergo the same rotation, changing
variables on the sphere makes the entire variation zero.

A relative rotation provides a test even when the raw axis stays
fixed. Let $n\perp e$ be a unit vector, $e(z)=e\cos z+n\sin z$
and take the marked insertion
$f_0(\omega)=\varphi(\sigma e\cdot\omega)$. Rotating the
insertion through the opposite angle at fixed axis gives the same
integral as rotating the axis at fixed insertion. In the latter
coordinates, $h''=-h$. Symmetry on the equator gives
$\int_{h=0}(n\cdot\omega)^2dS_{\rm eq}=V_{K-1}$. Thus
\[
 I''(0)=\varphi(0)-I(0).
\]
For example $\varphi(x)=\exp[-L^2Q^2x^2]$ is the actual
characteristic function in a trial normal measurement
$d=\sqrt{2}\,L x\,\xi$ with $\xi$ a standard real Gaussian.
Here $Q$ is the Fourier momentum conjugate to $d$, and one
can set $L=\sqrt{G_N}$ in natural units. Slicing the full
hemisphere gives
\[
 I(0)=(K-1)\int_0^1 x(1-x^2)^{(K-3)/2}
                            e^{-L^2Q^2x^2}\,dx
       \ \sim\ \frac{K-1}{2L^2Q^2}.
\]
The asymptotic follows by splitting at $x=1/2$ and rescaling
the interval adjacent to zero. The shell average of the unmarked
response decays, while its complete second relative-rotation
derivative tends to one. The contact reads the short-displacement
equator, even though that equator has zero unmarked shell measure.
This establishes a quantitative test of a proposed raw curvature
and boundary pairing. Its application to an electromagnetic
Hessian requires an actual relative rotation in that source,
and all other determinant and boundary variations must be
included before drawing a conclusion about a physical graph.

\paragraph{A local boundary transfer on the original velocity shell.}
The full shell makes it possible to test a proposed boundary state
without compactifying a geometric-time direction. For $K\ge2$ put
$n=K-1$ and $y=(v-v_\Gamma e_\Gamma)/R$. Eliminating the single
axial root in the shell delta sends its scalar measure to
$d\beta(y)=dy/|B^n|$ on the unit ball $B^n$. The inverse map has
$v_\Gamma=\sigma_\epsilon R\sqrt{1-|y|^2}$. The axial weight
cancels the radial-delta Jacobian exactly. At $K=1$ the oriented
shell is a point and offers no transverse displacement.

For a boundary dynamics test choose on this ball the flat Neumann
form $\int|\nabla_y g|^2d\beta$, with generator $H=-\Delta_N$.
Let $p_s(y,z)$ be its heat kernel relative to $\beta$, and let
$L=\sqrt{G_N}$ in units with $\hbar=1$ and unit invariant speed
magnitude. A proper time $t$ has units $L^2$. The A4 measure
determines $\beta$. The quadratic form and the internal readouts
below are additional, explicitly specified choices for testing
the history trace in the curvature pairing.

For integer $j\ge1$ take $j$ copies of a two-state internal space.
A single chosen transverse coordinate $f(y)=y_1$ suffices.
Conjugate $H\otimes I$ by the multiplication operator
\[
 U_j(y)=\operatorname{diag}(1,e^{if(y)})^{\otimes j}.
\]
This is a local flat normal connection with the transported
Neumann domain. Its heat kernel is
$p_s(y,z)U_j(y)U_j(z)^*$. Write
$\xi_\pm=((1,\pm1)/\sqrt{2})^{\otimes j}$. A transition from
$\xi_+$ to $\xi_-$ and back, with the same intermediate velocity
in the two kernels, has the positive trace weight
\begin{equation}
\begin{aligned}
 M_j(s_1,s_2)={}&
 \int d\beta(y)d\beta(z)\,p_{s_1}(y,z)p_{s_2}(z,y)\\
 &{}\times\left|\sin\frac{f(y)-f(z)}2\right|^{2j}.
\end{aligned}
\label{eq:nativefull-shell-mark-uv}
\end{equation}
Indeed each copy contributes the matrix element
$(1-e^{i(f(y)-f(z))})/2$ on one arc and its conjugate on the
other. The variable $z$ is integrated even for a contractible
return. Set $M_0=\operatorname{Tr}e^{-(s_1+s_2)H}$ for an
unmarked return. These are unnormalized loop traces. Dividing by
the heat trace would remove part of the determinant cost.

\begin{proposition}[Short-time and short-segment test on the A4 shell]
\label{prop:nativefull-shell-mark-uv}
Fix $m=1/L$ and put $h=u(1-u)$,
$s_1=tu/L^2$, $s_2=t(1-u)/L^2$.
Consider a transverse two-curvature form whose scalar polarization
coefficient at $X=|Q_E|^2$ is
\[
 P_j(X)=\int_0^\infty \frac{dt}{t}e^{-m^2t}
          \int_0^1du\,h\, e^{-thX}M_j(s_1,s_2).
\]
The curvature insertions give the transverse Hessian eigenvalue
$XP_j(X)$, up to its overall positive normalization. Sewing a
Maxwell port of covariance $1/X$ gives an absolutely finite
four-momentum integral exactly when $j>2+n/2$. The statement
includes coincident marks and simultaneous short durations.
It concerns this specified transverse block of a pairing.
\end{proposition}
\begin{proof}
After the shell integrations, all terms are nonnegative. Tonelli's
theorem and $\int_0^\infty X e^{-thX}dX=(th)^{-2}$ reduce the
test to
\begin{equation}
 \int_0^\infty\frac{dt}{t^3}e^{-m^2t}
       \int_0^1\frac{du}{h}\,M_j(s_1,s_2).
 \label{eq:nativefull-shell-power-uv}
\end{equation}
Let $s=s_1+s_2$. For $s\le1$ the Neumann heat kernel on the
unit ball obeys a Gaussian upper bound of the form
$C s_i^{-n/2}\exp[-|y-z|^2/(C s_i)]$
at time $s_i$~\cite{NativeFullNeumann}. Use
$|\sin[(f(y)-f(z))/2]|\le |y-z|/2$. Integration of the
product of the two Gaussian bounds against $|y-z|^{2j}$
gives
\[
 M_j(s_1,s_2)\le C_j s^{-n/2}(s h)^j.
\]
For $s\ge1$ the kernel on the longer arc is bounded by a constant
since its time is at least $1/2$. On the shorter arc its
$2j$-th displacement moment is at most
$C_j\min(1,\min(s_1,s_2)^j)$ by the short-time bound and the
bounded diameter. Thus this region is integrable for $j>0$
against the mass and $du/h$, using $\min(s_1,s_2)\le2s h$.

The small-time power also gives a lower bound on
$1/3\le u\le2/3$. Restrict $y$ to a smaller ball in the interior
and $z-y$ to $\sqrt{s}$ times a fixed bounded open set on which
the first component stays away from zero. Each short heat kernel
is bounded below there by $c s^{-n/2}$. For instance, restrict
paths to an interior box containing both points away from its
boundary. Its Dirichlet kernel has this bound by the elementary
one-dimensional reflection formula in each coordinate, and the
Neumann kernel includes these paths. The displacement factor is
at least $c_j s^j$, with the same conclusion for $j=0$ using
the unmarked trace. The volume of the allowed $z$ is a positive
constant times $s^{n/2}$. Hence the total-time integral has lower
power $t^{j-n/2-3}dt$ and the upper bound has this same power
times $h^{j-1}du$. Equality at $j=2+n/2$ gives a logarithm.
\end{proof}

\paragraph{What the inclusive electric trace sees.}
Keep the same local connection and the same initial state $\xi_+$.
Sum the return through a complete orthonormal family $\{\eta_b\}$
at the intermediate mark. Unitarity gives exactly
\[
 \sum_b |\langle\eta_b,U_j(y)U_j(z)^*\xi_+\rangle|^2=1 .
\]
The full trace with that initial state consequently recovers
$M_0$ for all $t,u$. In particular a local electric generator
acting as the identity on these internal copies cannot inherit the
finite off-diagonal answer by adding the missing positive channels.
A unitary change of the normal connection and its domain also
leaves the unmarked determinant invariant. This is a direct test
on the physical pairing required by $\Omega_\alpha$. A viable
charged representation would have to derive different matrix
elements in the complete source or a different joint sewing and
retain the physical electric charge. The finite off-diagonal form
is a curvature response defined above. A minimally coupled Dirac
determinant supplies additional orbital, spin and direct terms.
Their weights and Ward identity must come from that determinant
before the form can be used for induced O--QED.

For the tensor readout the symmetric tensor power already contains
both states and is invariant under $U_j$. Its dimension $j+1$ has
an algebraic interpretation. For a finite Hermitian generator $T$
the transition amplitude of $e^{i\delta T}$ is a sum of exponentials
over the distinct eigenvalues. A zero of order $j$ at $\delta=0$
sets its first $j$ moment sums (orders zero through $j-1$) to zero.
With at most $j$ distinct eigenvalues the Vandermonde matrix has
full column rank and all coefficients vanish. A nontrivial
order-$j$ transition therefore needs at least $j+1$ contributing
eigenvalues. The symmetric tensor construction attains this count.
It supplies a representation test on a proposed chiral block.

There is also a direct determinant check. Take a vectorlike Dirac
field whose primitive connection is independent of the velocity
and acts equally on these states. Suppose its squared operator in
a four-dimensional external field separates as
\[
 H_A^{(4)}\otimes I+L^{-2}U_j(H\otimes I)U_j^*+m^2 .
\]
The two terms on the two factors commute. The parity-even heat
trace then has the shell factor
$(j+1)\operatorname{Tr} e^{-tH/L^2}$ on the symmetric power,
at every background. The orbital insertions, spin term and direct
quadratic contact of the four-dimensional trace give the same
$h\,e^{-thX}$ as in \eqref{eq:monomial-wound-bubble-uv}, with
the circle trace replaced by this shell trace. Thus, in this
specific determinant, the complete scalar polarization is
$(j+1)P_0(X)/(2\pi^2)$. Even the constant Neumann eigenfunction
gives a divergent bare Maxwell coefficient. In this separable
class a selected transition can pass the short-duration test
while the compulsory electric contraction fails it.

\paragraph{Comparison on an auxiliary interval: the determinant of normal crossings.}
There is a case where the normal crossing can be followed through
the entire determinant. Let $y\in[0,\ell]$ denote an internal
spacelike normal coordinate, independent of the locked geometric-time
coordinates. Take $\ell=\sqrt{G_N}$ in natural units and the
Abelian action
\begin{equation}
 S_{\mathrm{int}}[A]
  =\frac{1}{4\ell}\int_0^\ell \mathrm dy\int\mathrm d^4x\,
                F_{MN}F_{MN}.
 \label{eq:monomial-interval-uv}
\end{equation}
The factor $1/\ell$ gives a $y$-independent field the customary
four-dimensional Maxwell kinetic normalization. At each nonzero
Euclidean momentum $k=|Q_E|$ a transverse component has normal
energy $(2\ell)^{-1}\int(|u'|^2+k^2|u|^2)\,\mathrm dy$.
Fix both endpoint values, solve $u''=k^2u$, and evaluate the
boundary flux $[\bar u u']_0^\ell/(2\ell)$. This is the cylinder
Dirichlet-to-Neumann construction~\cite{MonomialSteklov}. The resulting
two-endpoint Hessian is
\[
 \begin{gathered}
 a(k)\begin{pmatrix}1&-h(k)\\-h(k)&1\end{pmatrix},\\
 a(k)=\frac{k}{\ell}\coth(\ell k),\qquad
 h(k)=\frac{1}{\cosh(\ell k)}.
 \end{gathered}
\]
The Gaussian integral of the zero-endpoint fluctuations gives the
same Dirichlet determinant for all endpoint values. Mark each
off-diagonal entry by $z$, taking $z=1$ in the physical slab.
This marker counts crossings in the exact Hessian. The diagonal
part gives the unmarked reference. The relative determinant per
transverse mode is $1-z^2h(k)^2$, including both endpoint fields.
Every connected ring of $2n$ crossings has absolute log-determinant
coefficient $h(k)^{2n}/n$; the real Gaussian effective action adds
its universal factor $1/2$. Odd rings do not close. Because
$h(k)\le 2e^{-\ell k}$ for $k>0$, every such ring is absolutely
integrable in $\mathrm d^4Q_E$. At $z=1$ the summed absolute series
is $-\log(1-h^2)$. It grows only as $-2\log(\ell k)+O(1)$ at
$k\downarrow0$ and decays exponentially at infinity. This proves
absolute convergence of the relative log determinant itself per
unit four-volume. The unmarked bulk vacuum factors are independent
of the crossing count; the calculation does not assign a finite
absolute density to a source-independent free vacuum.

At $z=1$ integration of both endpoints gives the free Neumann
normal problem. Its eigenfunctions are constant or proportional to
$\cos(n\pi y/\ell)$. An endpoint observable with a nonzero component
along the constant mode therefore has a massless pole. The massive
levels have squared masses $(n\pi/\ell)^2$ and nonnegative spectral
weights equal to squared boundary matrix elements. The cross
covariance between distinct ends has alternating signs from
$\cos(n\pi)$, fully consistent with this positive matrix spectrum.
In particular the inter-end attenuation follows from propagation
through the geometric length in the bare action. This one-scale
result establishes all orders in the number of actual crossings,
within a Gaussian photon sector. The action has no charged
determinant. Coupling charged fields at a single endpoint allows
same-end contractions, for which there is no such crossing factor.

The value of this series also fixes the cost when the same geometric
length is integrated over a boundary state. With Fourier measure
$d^4 Q_E/(2\pi)^4$, the relative half log determinant per real
transverse component and unit four-volume at $z=1$ equals
\[
 -\frac{93\,\zeta(5)}{1024\pi^2\ell^4}.
\]
Indeed $\log\tanh^2 u=-4\sum_{j\ge1,\ j\ {\rm odd}}e^{-2ju}/j$,
and $\int_0^\infty u^3e^{-2ju}\,du=3/(8j^4)$. For a positive
law on lengths shared by an entire ring, Tonelli's theorem thus
makes finiteness equivalent to a finite fourth inverse length
moment. This statement allows arbitrarily small positive lengths.
A law on lengths that vary along the same ring requires its own
joint integral.

\paragraph{The massless photon in the interval comparison.}
Invert the two-endpoint Hessian of
\eqref{eq:monomial-interval-uv} at the physical value of the marker.
A real linear combination with endpoint coefficients $v_0,v_1$ has
covariance
\[
 D_v(k)=\frac{L}{k}\left[
 (v_0^2+v_1^2)\coth(Lk)+2v_0v_1\operatorname{csch}(Lk)\right],
 \qquad L=\ell .
\]
Its massless residue is $(v_0+v_1)^2$. Every nonzero fixed endpoint
combination has leading large-$k$ covariance $L(v_0^2+v_1^2)/k$.
For a bulk profile $w\in L^2([0,L])$ normalized by
$\int w(y)\,dy=1$, the spectral expansion of the Neumann resolvent
gives $L\langle w,(k^2-\partial_y^2)^{-1}_{N}w\rangle\ge1/k^2$.
The constant eigenfunction alone supplies this term. Exponential
propagation between two different endpoints is compatible with
these positive self-contractions.

\paragraph{A charged determinant with all proper-time faces retained.}
To see the distinction in the fermion determinant itself, close an
auxiliary flat normal coordinate into a circle of circumference $L=\sqrt{G_N}$.
Use periodic spinors in this direction, a five-dimensional Euclidean
Dirac operator with four spin components, mass $m=1/L$, and the
primitive charge-one connection. Take a smooth real external
four-potential $A(x)$, independent of the normal coordinate and of
Schwartz class on the four-plane. Let $e^{i\theta}$ be a flat normal
holonomy supplied as an external probe. Its gauge-invariant Fourier
coefficients label windings. This specifies a vectorlike comparison
for the positive determinant modulus; the phase of a chiral
determinant would require an additional calculation. The massless
photon profile is constant, so its coupling in an orthonormal basis
of charged normal modes $u_n$ is $\int\bar u_i u_j=\delta_{ij}$.
Thus every charged normal mode has the same coupling to the
constant photon, including the constant spinor mode at trivial holonomy.

For the squared Dirac operator in this background the normal and
four-dimensional heat factors commute. The proper-time integrand
of the parity-even effective action is, up to the common sign,
$e^{-m^2t}\Theta_\theta(t)\operatorname{Tr}_4 e^{-tH_A}/(2t)$,
where $H_A=(\slashed D_4)^*\slashed D_4$ and
\begin{equation}
\begin{aligned}
 \Theta_\theta(t)
 &=\sum_{j\in\mathbb Z}e^{-t(2\pi j+\theta)^2/L^2}\\
 &=\frac{L}{\sqrt{4\pi t}}\sum_{w\in\mathbb Z}
       e^{-L^2w^2/(4t)}e^{iw\theta}.
\end{aligned}
\label{eq:monomial-circle-uv}
\end{equation}
The second equality follows by Poisson summation. The spin trace
is included in $\operatorname{Tr}_4$. At zero external field its
extensive part is understood per unit four-volume.

\begin{proposition}[Charged winding coefficients and the contractible sector]
\label{prop:monomial-charged-winding-uv}
At every finite order in the external four-potential, the sum of
the nonzero winding coefficients of this one-fermion-loop source
converges absolutely in winding, the continuous loop momentum and
the ordered proper times. The same holds after any fixed number of
normal holonomy derivatives. The complete parity-even determinant
has a divergent bare magnetic quadratic coefficient.
\end{proposition}
\begin{proof}
For a coefficient of positive order $n$, expand $H_A$ about the
flat Laplacian. It has insertions at most linear in the loop
momentum (from the first-order covariant derivative and the spin
term) and a quadratic contact in $A$. For $r$ insertions the
ordered simplex has volume $t^r/r!$ and the free periods $t_j$
between insertions sum to $t$. For cumulative transfers $p_j$,
$\bar p=\sum_j t_j p_j/t$, one has
\[
 \sum_j t_j|k+p_j|^2
 =t|k+\bar p|^2+\sum_j t_j|p_j-\bar p|^2.
\]
The absolute Gaussian momentum integral is bounded by
$C_n t^{-2}(1+t^{-n/2})$ times a polynomial in the external
momenta. Schwartz tests integrate that polynomial after overall
momentum conservation. The bound is uniform even on faces where
several $t_j$ vanish together. The free term uses its heat
density per volume. For $0<t\le L^2$, the sum over $w\ne0$,
including a fixed power of $|w|$ from differentiation, is at most
$C e^{-L^2/(8t)}$ before the prefactor in
\eqref{eq:monomial-circle-uv}. At large $t$ the same sum grows
at most polynomially and $e^{-m^2t}$ is integrable against each
such power. These estimates prove the first claims at every fixed
insertion order. They do not require a convergent expansion in the
strength of the external field.

A uniform magnetic probe $B$ gives a direct check of the remaining
local coefficient per unit four-volume. In a constant field the
four-dimensional heat trace relative to its free density has
factor $Bt\coth(Bt)-1$. Indeed the two plane directions give the
Landau oscillator factor $Bt/\sinh(Bt)$, the spin trace relative
to four components gives $\cosh(Bt)$, and the other two momenta
are free~\cite{Schwinger1951}. The absolute bare coefficient of
$B^2$ is therefore
\[
 \frac{1}{24\pi^2}\int_0^\infty
        \frac{dt}{t}e^{-m^2t}\Theta_\theta(t).
\]
Here the field-independent term has no magnetic derivative. The
zero winding in \eqref{eq:monomial-circle-uv} gives a positive
multiple of $L\int_0 dt\,t^{-3/2}$ at the lower limit. All
nonzero windings are integrable there. This establishes the
divergence of a field-dependent coefficient at every holonomy
without fitting or subtracting a Maxwell term.
\end{proof}

\paragraph{A winding on the whole graph need not protect a subgraph.}
The compactly supported electromagnetic probe can also be sewn.
For clarity, write $X=|Q_E|^2>0$, $a=tX$, and
$I(a)=\int_0^1 e^{-a u(1-u)}\,du$. For a transverse Fourier
component of a real Schwartz source $A$, the quadratic coefficient
of the four-dimensional spin heat trace per
$|\widehat A_T(Q_E)|^2\,d^4Q_E/(2\pi)^4$ is
\[
 (4\pi t)^{-2}\{4t[I(a)-1]+2t^2X I(a)\}.
\]
To obtain this expression, the two orbital derivative insertions
give $4tI(a)$ after the Gaussian momentum shift. Their direct
quadratic contact gives $-4t$. The Pauli insertions have spin trace
$4X$ in a unit transverse polarization, hence give $2t^2X I(a)$;
the mixed spin trace vanishes. Transversality follows from
covariance of the heat trace. Integrating the derivative of
$(1-2u)e^{-a u(1-u)}$ from zero to one gives
$2[1-I(a)]=a\int_0^1(1-2u)^2 e^{-a u(1-u)}\,du$.
This combines the direct and iterated terms before the short-time
integral is taken.

Write the winding-$w$ quadratic effective action as the form with
Hessian $X\Pi_w(X)$ on transverse tests, with the usual factor
$1/2$ in the action. The preceding identity gives
\begin{equation}
\begin{aligned}
 \Pi_w(X)
   &=\frac{L}{2\pi^2\sqrt{4\pi}}\int_0^1du\,u(1-u)\\
   &\quad\times\int_0^\infty\frac{dt}{t^{3/2}}\,
      e^{-t[m^2+u(1-u)X]-L^2w^2/(4t)} .
\end{aligned}
\label{eq:monomial-wound-bubble-uv}
\end{equation}
Here $\Pi_0$ denotes \emph{zero winding}, distinct from evaluation
at $\theta=0$. For $w\ne0$ the integral is finite and strictly
positive. The $w=0$ proper-time density is positive and diverges
upon integration at zero, uniformly for $X$ in any compact
positive interval. Consequently the divergent quadratic response
is already visible for smooth wave packets away from zero momentum.

Already a single wound loop gives a sharper photon test. For
$w\ne0$ the proper-time integral in
\eqref{eq:monomial-wound-bubble-uv} can be done explicitly:
\[
 \Pi_w(X)=\frac{1}{2\pi^2|w|}\int_0^1 du\,u(1-u)
    e^{-|w|L\sqrt{m^2+u(1-u)X}} .
\]
Use $\int_0^\infty t^{-3/2}e^{-at-b/t}\,dt
=\sqrt{\pi/b}\,e^{-2\sqrt{ab}}$ for $a,b>0$. At fixed
$0<u<1$ the integrand decays exponentially with
$\sqrt X$. The short segment at either endpoint changes the
large-$X$ law of the complete coefficient. By symmetry restrict
to $0<u<1/2$ and put $y=Xu$. Dominated convergence then gives
\begin{equation}
 \lim_{X\to\infty}X^2\Pi_w(X)
   =\frac{1}{\pi^2 |w|}\int_0^\infty
              y\,e^{-|w|L\sqrt{m^2+y}}\,dy>0.
 \label{eq:monomial-wound-tail-uv}
\end{equation}
Indeed the rescaled integrand is bounded by
$y e^{-|w|L\sqrt{m^2+y/2}}$, an integrable function independent
of $X$. Quantize the constant normal Maxwell mode as a
four-dimensional photon, retaining the entire charged tower.
Its contraction with this same winding Hessian has radial absolute
density proportional to $X\Pi_w(X)\,dX$. Thus the two-loop
bare vacuum graph for \emph{each fixed nonzero winding} diverges
logarithmically. This assertion concerns the specified dynamical
zero-mode photon. A construction that glues other normal photon
modes coherently would require its own combined calculation.
The slow tail here is derived from the charged source and its
contact, with no change to the propagator or the boundary pole.

Now propagate a Maxwell zero-mode photon between two complete
fermion loops, with a transverse external mark on each loop. The
connected diagram is reducible upon cutting that photon. Normal
momentum conservation allows the constant internal photon in this
graph. Closing the two Hessians with its covariance $1/X$ gives
$X\Pi_{w_1}(X)\Pi_{w_2}(X)$, up to the common graph sign and
symmetry factor. In the Fourier coefficient of total winding
$w\ne0$ the term $(w_1,w_2)=(0,w)$ has a divergent absolute
integral, including at fixed nonzero external momentum. This
statement follows by restricting the proper times to positive
finite intervals first and then using monotonicity of the displayed
densities. Every term with $w_1+w_2=w$ has the same sign at this
stage. Summing the normal labels of this graph preserves the
divergence. It persists for smooth external transverse packets
with nonzero weight in a compact annulus. This example has one
hard loop with zero winding while the other loop retains a finite
geometric excursion. Demanding nonzero total winding of an
amplitude would miss precisely this region.

The local obstruction extends to a class of normal mass operators.
Suppose at a constant magnetic probe the complete squared charged
operator separates as $H_B\otimes 1+1\otimes M^2$, with the same
charge-one spinor $H_B$ as above, $M^2\ge0$, and a positive
normal trace and boundary preparation independent of $B$.
Assume these give a nonnegative mass measure $\nu$ for the heat
factor, with finite heat trace at positive times. Any interval
$[0,\Lambda_m^2]$ with $0<\nu([0,\Lambda_m^2])<\infty$ supplies
the lower bound
\[
 \int e^{-t\lambda}\,d\nu(\lambda)
 \ \ge\ e^{-t\Lambda_m^2}\nu([0,\Lambda_m^2]).
\]
The bare $B^2$ coefficient therefore contains at least a logarithmic
short-time divergence even for a single occupied finite mass band.
The interval here selects part of a positive measure in a lower
bound; all masses are retained in the action. An additional
source-dependent normal operator or boundary contact would change
this factorization. Its size and sign must be calculated in the
same magnetic coefficient. Changing the positive mass distribution
alone cannot cure the divergence in this class.

The nonzero-winding projection is an observable of this specified charged
source. It leaves out contractible loops, including positive-length
unreduced returns. The theorem covers infinitely many classical
source orders of one determinant loop. A dynamical photon can
still close at the same normal location. A physical GCM
specialization must derive how the intermediate displacement of
that return enters the charge-one determinant, the seagull and the
actual photon gluing measure simultaneously. The circle comparison
does not identify this extra spacelike coordinate with any ambient
geometric-time coordinate.

\paragraph{A charged operator that retains the intermediate shell average.}
There is a different placement of the same A4 shell law for which
the complete charged operator can be tested at every order. Take
$K\ge4$, let $n=K-1$ and choose a fixed linear map
$T:\mathbb R^4\to e_\Gamma^\perp$ of rank at least three in
Euclidean orthonormal charts. Set $L=\sqrt{G_N}$ and $m=1/L$
in the carrier units used above. The map and mass ratio specify
this comparison. The normalized ball measure of the A4 shell has
characteristic function
\[
 \chi_n(p)=\Gamma(1+n/2)(2/|p|)^{n/2}J_{n/2}(|p|),
 \qquad \chi_n(0)=1.
\]
This follows by integrating the slices perpendicular to $p$, or
by the power series of the Bessel function. Two shell averages
at opposite marks give $|\chi_n|^2$ when their measures factor.
The same velocity at both marks gives one. In the following
charged source the placement follows from the chosen bilinear.

Let $P_0 f=\langle1,f\rangle_\beta 1$ be the projection onto the
normalized constant shell state. Compress the connection to this
state \emph{inside the affine charged bilinear}. In a nonzero
four-momentum chart decompose $A=\partial\phi+A_T$ and put
\begin{equation}
 \begin{aligned}
 \mathcal A_\mu^{(0)}(x)
    &=\partial_\mu\phi(x)+
        \int_{B^n}(A_T)_\mu(x+LT^*y)\,d\beta(y),\\
 D_{\theta,A}
    &=m+\gamma^\mu(\partial_\mu+i\mathcal A_\mu^{(0)})
                 +\gamma^5(\partial_z+i\theta/L).
 \end{aligned}
 \label{eq:nativefull-compressed-operator-uv}
\end{equation}
Use the ordinary quadratic Maxwell action for the
four-dimensional, normal-constant photon, periodic spinors on
the normal circle of circumference $L$, and Hermitian Euclidean
Clifford matrices. At zero four-transfer use the continuous
factor $\chi_n(0)=1$ in the transverse coupling. The primitive
charge transforms at $x$ since $A\mapsto A+\partial\lambda$
adds $\partial\lambda(x)$ to the connection. In the transverse
gauge every ordinary Dirac insertion has multiplier
$\chi_n(LTQ)$, where $n=K-1$. The Abelian gauge determinant in
this gauge is independent of the electromagnetic coupling.
There is a single charged determinant. Squaring the operator
to form its modulus gives the product of the two compressed
connections in its direct quadratic contact. At opposite
transverse transfers this gives $|\chi_n(LTQ)|^2$. Averaging
the already squared connection over a shared shell state gives
a different contact.

\begin{proposition}[All finite orders for protected charged loops]
\label{prop:nativefull-compressed-all-uv}
For \eqref{eq:nativefull-compressed-operator-uv}, consider a
connected graph at positive electromagnetic order, with closed
charged loops and optional transverse external photon packets.
Resolve each charged loop in normal winding. The sum of the
absolute integrals converges in all ultraviolet regions at every
finite order if each closed charged loop either has nonzero winding
or has at least six photon insertions. The latter loops can have
zero winding. This includes internal photon contractions and simultaneous
short normal segments. Hold the infrared treatment fixed. One
explicit version uses a four-dimensional Euclidean torus with
fixed harmonic photon sector and the massless Maxwell covariance
on all nonzero Fourier modes. The torus can have any finite size;
no ultraviolet mode is omitted. In the continuum the same hard
estimate applies with soft photon singularities controlled
separately.
\end{proposition}
\begin{proof}
On the real cover of the normal circle, Fourier transform the
free circle momentum $\xi$ to a displacement $d$. For four-momentum
$p$ the inverse Dirac multiplier is
$(m-i\gamma\cdot p-i\gamma^5\xi)/(\xi^2+E_p^2)$ with
$E_p=\sqrt{m^2+|p|^2}$. The oriented displacement kernel is
\[
 S(p;d)=\left[\frac{m-i\gamma\cdot p}{2E_p}
             +\frac{\gamma^5}{2}\operatorname{sgn}d\right]
                  e^{-E_p|d|}.
\]
This follows by the one-dimensional Fourier transform with phase
$e^{i\xi d}$; its jump at zero is bounded. Multiplication of circle
propagators followed by the normal trace becomes convolution on the
cover. The winding-$w$ term restricts the real increments to
$\sum_{j=1}^r d_j=wL$ for a loop with $r\ge1$ insertions.
The holonomy gives a phase of modulus one; the common center
supplies an overall factor $L$ per loop. Equivalently, the
periodizations of the cover kernels can first be taken as
distributions on smooth tests. The following estimate gives
absolute integrability on the nonzero windings.

Route the carrier momenta around the loop as $k+p_j$, where the
$p_j$ are cumulative photon transfers. On the stated hyperplane
some $|d_i|\ge |w|L/r$. Partition or bound by the union of these
$r$ regions. The norm bound on the product of Dirac propagators
and unit polarization vertices is then at most a fixed
$r$-dependent constant times
\[
 e^{-|w|L E_{k+p_i}/(2r)}
 e^{-m\sum_j|d_j|/4}\,e^{-m|w|L/4}.
\]
Integrating in $k$ is uniform in every transferred momentum by
a translation. On a finite torus the corresponding lattice sum
is bounded uniformly in the shift, e.g.\ by comparing each cell
with a translated radial exponential. Eliminate one displacement
with the hyperplane constraint. The other $r-1$ integrations cost
at most $(8/m)^{r-1}$. The last exponential sums over all nonzero
integers $w$. This estimate imposes no positive lower bound on
any individual segment and is uniform in all photon transfers,
including momenta shared between charged loops.

A closed loop with zero winding has the unrestricted Fourier
integral over its continuous momentum $\xi$ on the cover. For
$r>5$ one can estimate it directly in this representation.
Each Dirac propagator has norm at most
$(m^2+|k+p_j|^2+\xi^2)^{-1/2}$ by the Euclidean Clifford
relations. H\"older with $r$ equal exponents bounds the integral
of their product by
\[
 \int_{\mathbb R^5}(m^2+|\kappa|^2)^{-r/2}\,d^5\kappa<\infty .
\]
This is uniform in all four-dimensional transfers. The corresponding
four-dimensional lattice sum with continuous $\xi$ on the torus
has the same bound up to a volume-dependent constant.
No statement about absolute convergence of four- or five-insertion
loops is needed for this criterion.

It remains to integrate those transfers. Shrink each charged loop
to a vertex, retaining self-edges and multiple photon edges. After
overall momentum conservation choose a spanning tree. The photon
momenta on the omitted edges are independent variables; the tree
momenta are determined by these and the external momenta. Every
internal transverse photon has a factor
$|\chi_n(LTQ)|^2/|Q|^2$ from its two endpoints and its Maxwell
covariance. The large-argument Bessel bound gives
$|\chi_n(p)|^2\le C_n(1+|p|)^{-(n+1)}$. At rank four
the photon factor is integrable at large $Q$. At rank three
write $Q=(q,\omega)$ in the three directions seen by $T$ and
the remaining direction. Its hard integral is bounded by
a constant times
\[
 \int_{\mathbb R^3}\!
    \frac{\pi\,d^3q}{|q|(1+cL|q|)^{n+1}},
 \qquad
 \int_{\mathbb R}\frac{d\omega}{|q|^2+\omega^2}
       =\frac{\pi}{|q|}.
\]
Here $c>0$ is the smallest nonzero singular value of $T$.
Since $n\ge3$, the integral is finite. For the torus, away from
the fixed soft region the same estimate follows by comparison
of unit lattice cells with the integrable envelope
$(1+|Q|^2)^{-1}(1+c'L|TQ|)^{-n-1}$. Its values at points a
bounded distance apart have bounded ratio. Tree factors are bounded
under the stated infrared convention, and smooth packets integrate
the external transfers. Rank three is sufficient even when the
readout sees only the three spatial components of the locked
carrier. For rank at most two, keeping $TQ$ small leaves at least
two unsuppressed directions with a nonintegrable $|Q|^{-2}$ tail.
This last test concerns the single photon majorant; further
fermionic decay could matter in an individual graph.
These estimates cover each graph of the first-order determinant
expansion; taking its parity-even part preserves the bound.
\end{proof}

\begin{corollary}[Open charged paths]
\label{cor:nativefull-open-uv}
The same estimates allow any finite number of open charged
lines with smooth tested fermion endpoints. No winding condition
is required on an open line. In particular all graphs with no
closed fermion loop obey the ultraviolet bound at every finite
order in this comparison, with the same infrared prescription.
\end{corollary}
\begin{proof}
For a chain of $r+1$ propagators let $z,z'$ be its external normal
coordinates in a chosen interval of length $L$. On the cover
the displacements sum to $D=z'-z+wL$. The carrier momentum of
each propagator is fixed by an external fermion momentum and the
photon transfers. There is no independent momentum circulating
in the chain. Its norm is bounded uniformly in those transfers
by $C^{r+1}\exp[-m\sum_j|d_j|]$. Reserve $e^{-m|D|/2}$,
eliminate one increment, and integrate the others at a cost at
most $(4/m)^r$. The sum of $e^{-m|z'-z+wL|/2}$ over all
integers is uniformly bounded for endpoints on the circle.
This also treats coincident normal endpoints on an open
line, with the cover kernel at a jump understood on smooth
tests. Contract each open chain as a vertex supplying its net
external momentum in the same photon spanning-tree argument.
Schwartz carrier packets integrate the remaining independent
external momenta.
\end{proof}

In particular the previously logarithmically divergent
constant-mode photon closure at fixed $w\ne0$ now has radial
density proportional to $X\,|\chi_n(LTQ)|^2\Pi_w(X)\,dX$.
The factor is present in the entire two-source variation.
This improvement comes from retaining two shell integrals in
the prescribed bare charged coupling. The zero-winding charged
polarization still has its divergent proper-time integral at
small nonzero external $Q$, where $\chi_n(LTQ)$ remains nonzero.
A nonzero \emph{total} winding also permits some charged loops
to have zero winding. A
rank-three spatial readout preserves the carrier time of the
transverse connection at each measurement. The projection onto
four-dimensional transverse tests and the physical BRST quotient
still require a reconstruction argument before a positive
Lorentzian current or a covariant cut can be claimed. The
comparison derives a counterterm-free restricted graph theorem.
Its hypotheses specify a particular physical pairing and boundary
preparation beyond the displayed native curvature monomial.


\paragraph{Protecting the circulating charged momentum.}
The same oriented shell also permits a test of the remaining
zero-winding ultraviolet problem. Choose $K\ge6$ and an
isometric five-dimensional readout into $e_\Gamma^\perp$,
whose restriction to the four-plane is the rank-four photon readout. The fifth
momentum is conjugate to the auxiliary spacelike circle.
If $P=(k,\xi_j)$, $\xi_j=(2\pi j+\theta)/L$, let $F_\theta$
multiply a charged mode by
\[
 f(P)=\int_{B^{K-1}}e^{iL(T_5P)\cdot y}\,d\beta(y)
       =\chi_{K-1}(LT_5P).
\]
This is the average of translations through $LT_5^*y$ in
the four carrier coordinates and the lifted normal circle,
including its holonomy in the fifth translation. Its norm is
at most one. Keep the compressed transverse photon
$A_T^{(0)}$ of \eqref{eq:nativefull-compressed-operator-uv}.
For the longitudinal potential $\phi$ defined there, specify the
charged operator
\begin{equation}
 \begin{aligned}
 \mathscr D_{\theta,A}
 &=e^{-i\phi}\big[D_{\theta,0}
       +i F_\theta\,\gamma^\mu(A_T^{(0)})_\mu
                              F_\theta\big]e^{i\phi}.
 \end{aligned}
 \label{eq:nativefull-three-compressions-uv}
\end{equation}
The two charged factors and the transverse connection are
separate compressions onto the constant shell state, before the
determinant is taken. For example the matrix element of a
product $U_1P_0U_2P_0U_3$ between unit constant states is the
product of the three one-mark matrix elements. This specifies
additional history and pairing data for a comparison using the
native A4 law. A single occurrence of the velocity-shell delta
in the curvature action would not imply this composition rule.
All profiles and $m=1/L$ are fixed once this model is chosen.
Rotational invariance of the shell ball makes the scalar readout
identical for all such isometric embeddings. The representation
and the intervening projections remain definite model choices.
The conjugation implements charge-one gauge covariance under
$A\mapsto A+\partial\lambda$ within the declared harmonic
sector. The electromagnetic matrix elements and the physical
photon pole still require calculation.

\begin{proposition}[Unrestricted charged circulations and a finite photon obstruction]
\label{prop:nativefull-alllegs-uv}
In the Euclidean comparison \eqref{eq:nativefull-three-compressions-uv},
every complete transverse graph at a positive finite
electromagnetic order is absolutely ultraviolet convergent with
the preceding infrared convention, including interacting vacuum
graphs and arbitrary charged normal modes. No winding projection,
momentum cutoff or counterterm enters this estimate. Its one-loop
inverse transverse photon propagator nevertheless acquires a
nonzero constant at vanishing momentum transfer.
\end{proposition}
\begin{proof}
Use the transverse gauge, in which $\phi=0$ and the first-order
operator is affine in the photon. Every internal charged line
between two insertions now has norm bounded by
\[
 C\,|f(P)|^2/E_P\le C_{m,L}(1+L|P|)^{-(K+1)},
 \quad E_P=\sqrt{m^2+|P|^2}.
\]
This majorant is integrable on $\mathbb R^4\times\mathbb Z$
with the indicated circle momenta. On a finite carrier torus the
corresponding five-dimensional lattice sum converges. On each
closed fermion loop select one propagator, assigning its four
carrier components and its circle mode as the independent loop
variables. A translation of these four components by the photon
transfers leaves a uniform integral or lattice bound. All other
charged propagators are uniformly bounded since $m>0$. This
also bounds open charged chains with tested endpoints. The
remaining independent variables are the photon momenta on edges
omitted from a spanning tree of the contracted photon graph.
Each such edge has its two compressed transverse factors and a
Maxwell covariance, as in the preceding proof. Consequently all
hard integrations converge simultaneously. This proof works in
the original normal mode sum, including the contributions that
caused the zero-winding divergence in the earlier calculation.
The free Gaussian normalization has electromagnetic order zero.

For the physical test take infinite carrier volume and approach
$Q=0$ in a fixed transverse polarization $i$. At zero transfer
the two vertices supply $f(P)^4$ and the photon readout has
$\chi_{K-1}(0)=1$. There is no direct transverse contact in this
affine first-order operator. In the Euclidean convention
$\Gamma_f=-\log|\det\mathscr D|$, the second derivative is
$\operatorname{Re}\operatorname{Tr}(S\,iV_i S\,iV_i)$ with
$S=D_{\theta,0}^{-1}$ and $V_i$ the vertex with its factor $i$
removed. The spin trace before the sign $i^2$ is
$4(m^2+|P|^2-2k_i^2)/(m^2+|P|^2)^2$ times $f(P)^4$.
Rotational averaging in the four carrier variables replaces
$k_i^2$ by $|k|^2/4$. Thus the transverse Hessian in this
convention has a strictly negative limit, whose magnitude per
unit four-volume and unit polarization is
\[
 4\sum_{j\in\mathbb Z}\int\frac{d^4 k}{(2\pi)^4}
 \frac{f(P)^4\,[m^2+\xi_j^2+|k|^2/2]}
      {(m^2+\xi_j^2+|k|^2)^2}.
\]
The number is finite by the same majorant and strictly positive
because the real analytic form factor is nonzero on a set of
positive measure. Dominated convergence applies also at nearby
transfers: bounded translations preserve the hard bound at
large $|P|$ and $m>0$ controls the small denominator.
The longitudinal dressing permits a constant multiplying the
nonlocal transverse projector. The bare Maxwell eigenvalue $X$
cannot cancel that constant at the same electromagnetic order.
\end{proof}

This calculation removes the zero-winding ultraviolet obstruction
for an explicit charged bilinear with shell-derived profiles.
The finite shift fails the perturbative massless-photon condition.
A cure in the preferred GCM action must generate the compensating
physical source terms through its own pairing, determinant or
boundary dynamics, while retaining their hard-momentum bounds.
The gauge transformation law by itself supplies neither the
missing transverse response nor Lorentzian spectral positivity.

\paragraph{What would have to change to recover the electric vertex?}
The failure can be seen even before doing the photon loop. For
$Q\ne0$ let $P_T$ and $P_L$ be the transverse and longitudinal
projectors of the Euclidean four-plane, and write
$P+Q=(k+Q,\xi_j)$. Differentiating the same operator at zero
background gives the one-photon vertex
\[
 \Gamma_\mu(P+Q,P)=i\gamma^\nu
  \left[f(P+Q)f(P)\chi_{K-1}(LTQ)(P_T)_{\nu\mu}
                         +(P_L)_{\nu\mu}\right].
\]
Thus $Q^\mu\Gamma_\mu=i\gamma\cdot Q$ is the difference of the
free inverse propagators on this fixed normal mode. In contrast,
for $Q$ parallel to the $\mu$ axis the limiting vertex is
$i\gamma^\mu$, whereas for $Q$ perpendicular to it the limit is
$i f(P)^2\gamma^\mu$. At momenta where $f(P)^2\ne1$ the soft
vertex depends on the direction of approach. Gauge covariance at
nonzero transfer survives this defect.

Suppose instead that a differentiable charge-one vertex has the
same soft limit $i g(k)\gamma^\mu$ in every direction and obeys
the Ward identity with a $C^2$ inverse charged propagator $D(k)$.
Then $\partial_\mu D=i g\gamma^\mu$. Equality of mixed derivatives
requires $(\partial_\nu g)\gamma^\mu=(\partial_\mu g)\gamma^\nu$.
The independent Clifford matrices force $\nabla g=0$.
Consequently changing only the free symbol cannot accommodate a
nonconstant scalar leg weight of this type while giving the
vertex a regular soft limit. A physical curvature pairing that
retains charged ultraviolet suppression would have to supply
additional vertex structures or change the specified sewing,
and its second variation would have to pass the masslessness test.

\subsection{A massless photon from complete shell displacements}
\label{sec:ward-hop-uv}

Work on a fixed post-locking sector-plane $\bin=(s,b_s)$.
The momentum norms in the estimates are Euclidean. In the cut
tests the continuation uses the matter boundary value on the
same plane, with opposite infinitesimal shifts along its prescribed
geo-momentum unit vector. The intermediate Wilson paths and normal
marks below are choices for the charged representation of the full
ambient curvature integral. They specify what survives before the
trace over the charged field.

\paragraph{Gauging the entire shell hopping operator.}
The finite photon constant in the three-compression example
suggests changing the charged inverse and its photon vertices
together. On a reconstructed Euclidean four-plane take $K\ge5$,
$n=K-1$, $L=\sqrt{G_N}$ in the carrier units used above, and an
isometry $T:\mathbb R^4\to e_\Gamma^\perp$. Use the actual A4
scalar shell chart $y\in B^n$ with probability $d\beta$.
Let $h(z)\,d^4z$ be the law of $LT^*(y+y')$ for two independently
integrated shell marks. Its characteristic function is
$g(p)=\chi_n(LTp)^2$. Thus $h$ is a bounded probability density
supported on $|z|\le2L$. For every fixed multi-index $\alpha$,
\[
 |\partial_p^\alpha g(p)|\le C_{\alpha,L}(1+L|p|)^{-K}.
\]
Indeed the radial Bessel formula and its derivative recurrences
give the envelope $(1+L|p|)^{-K/2}$ for each fixed derivative of
one shell factor; the compact measure controls derivatives near
the origin. This also proves boundedness of the density by
Fourier inversion. The profile is fixed by the shell and its
placement, without a momentum cutoff.

Write $S$ for the photon multiplier $\chi_n(LTQ)$ and use the
connection $B=SA+(1-S)\partial\phi$ of
\eqref{eq:nativefull-compressed-operator-uv}, where
$\partial\phi=P_L A$ at nonzero Fourier momentum. It obeys
$B[A+\partial\alpha]=B[A]+\partial\alpha$. Its multiplier
$S I+(1-S)P_L$ extends smoothly through $Q=0$: radial symmetry
gives $1-S=Q^2 b(Q^2)$ with $b$ regular there.
For Hermitian Euclidean Dirac matrices set
\begin{equation}
\begin{aligned}
 D_B&=L^{-1}(I-\lambda H_B),\qquad \lambda=1/4,\\
 (H_B\psi)(x)&=\int d^4z\,h(z)W(z)\,
          e^{\,i\int_0^1 z^\mu B_\mu(x+s z)\,ds}\psi(x+z),\\
 W(z)&=I-\gamma\cdot z/(2L).
\end{aligned}
\label{eq:covhop-source-uv}
\end{equation}
The dimensionless weight, spin factor, shell composition and
straight path are specified choices of the charged representation.
They are not uniquely selected by the master curvature pairing.
For a real connection, the link transforms by its endpoint
phases, so $D_{B+\partial\alpha}=e^{-i\alpha}D_B e^{i\alpha}$.
Every order of the Wilson exponential supplies its own contact
vertices. Here $\|W(z)\|\le2$, and Schur's bound gives
$\|\lambda H_B\|\le1/2$. The charged inverse is its norm-convergent
geometric series. The source-dependent part of the fermionic
negative trace logarithm uses
$\sum_{N\ge1}\lambda^N\operatorname{Tr}H_B^N/N$ with the same
weights in its marked derivatives. For the determinant modulus
take the real part. In infinite volume use the connected
coefficient per unit volume or external packets. The bounds
below define these coefficients directly. The ultralocal constant
in the free normalization has electromagnetic order zero.

\begin{proposition}[Ultraviolet convergence and the perturbative photon pole]
\label{prop:covhop-photon-uv}
For \eqref{eq:covhop-source-uv} with the ordinary transverse
Maxwell field, the complete bare graphs at every fixed
electromagnetic order, including all closed charged circulations,
are absolutely ultraviolet finite. The sum over hopping lengths
converges at that order. On the infinite Euclidean plane each
finite-order contribution to the transverse one-particle-irreducible
photon inverse propagator, divided by $Q^2$, has a finite limit
as $Q\to0$. The formal inverse propagator thus has zero constant
term and a finite perturbative coefficient of $Q^2$ with tree
value one. This assertion concerns the perturbative pole, before
any Lorentzian spectral reconstruction.
\end{proposition}
\begin{proof}
In transverse gauge $B=SA$. A hop bearing $a$ photon insertions
has the Fourier transform of $h(z)W(z)$ times $a$ components of
$z$, evaluated at the circulating momentum shifted by photon
transfers and path fractions $s\in[0,1]$. Each photon also has
its factor $\chi_n(LTQ)$. The shifted Fourier transform is
uniformly integrable in the four circulating components by the
derivative envelope above. Its supremum is at most $2(2L)^a$.
Use one such integrable hop on each closed charged loop and the
supremum bounds on the others. No assumption on the other
transfers is needed. A marked $N$-hop loop at fixed total order
$a$ then costs at most a constant times
$(2\lambda)^N(1+N)^a$, after counting the mark placements.
An open charged chain has no independent carrier circulation
and satisfies the same bound with tested endpoints; the local
identity part of the free inverse creates no additional loop.
At fixed positive electromagnetic order only finitely many
marked charged components occur in a connected graph.

After contracting each charged component to a vertex, independent
photon transfers can be assigned to the edges omitted from a
spanning tree. Each has hard majorant
$|\chi_n(LTQ)|^2/|Q|^2=O(|Q|^{-K-2})$. The tree factors are
bounded with the previously stated infrared convention. For fixed
soft treatment this proves absolute convergence of the full hard
integrals in the charged and photon momenta together, including
partially hard regions, with an absolutely summable bound in $N$.

To examine the actual soft behavior on the infinite plane,
perform the absolutely convergent charged-loop integrals first.
The trace of $N$ hops closes a polygon with $\sum_i z_i=0$.
Group the placements of each photon on this polygon, including
multiple photons on one link. Its geometric current, before the
single photon shell factor, is
\[
 J^\mu(Q)=\oint dx^\mu e^{iQ\cdot x}.
\]
It obeys $Q_\mu J^\mu=0$ and $J^\mu(0)=0$. Taking the initial
vertex as origin gives
$|J(Q)|\le\min(2LN,\,|Q|(2LN)^2)$. The estimate uses
$|e^{iQ\cdot x}-1|\le |Q||x|$ and the total perimeter bound.
The remaining relative displacement integral on a closed polygon
is bounded by $\|h\|_\infty$, since convolution with a probability
density contracts the supremum norm. Spin traces cost at most
$4\,2^N$. These bounds give a polynomial in $N$ for every finite
number of currents, absorbed by $(2\lambda)^N$. Consequently
each incident photon on a closed loop contributes
$\min(1,L|Q|)$ to a bound on its completed kernel, as well as
its shell factor. The two incidences of each internal photon
control the soft pole even when several momenta vanish together.
The resulting photon majorant is bounded at the origin and
integrable in every independent momentum at infinity.

For a two-point one-particle-irreducible photon graph both
external marks lie on closed charged components. Dividing each
external current by $|Q|$ leaves a bounded expression with a
limit along each direction, obtained by expanding its finite
polygon phase. For almost every internal momentum the other
photon kernels have limits as well. In a one-particle-irreducible
graph each internal photon lies on a cycle after the charged
loops are contracted, so its momentum can vary with a chord
variable. Its zero set has measure zero. Tree bounds and the
integrable chord majorants above permit dominated convergence,
including the hop sums. Euclidean rotation invariance and
transversality make the limiting transverse scalar independent
of direction. This proves the pole assertion order by order,
without assuming convergence of the electromagnetic series.
\end{proof}

The one-photon Ward identity is regular in this model. A constant
connection gives $D_B(p)=D_0(p+B)$, so the soft insertion is
$\partial_\mu D_0(p)$, including the differentiated spin and
shell terms. For $Q\ne0$ contracting a link insertion with $Q$
gives its endpoint difference; together the links give
$D_0(p+Q)-D_0(p)$. Higher variations of that same link supply
the seagulls used in the polygon argument.

\paragraph{The charged poles of the same Wilson operator.}
At zero photon background the source \eqref{eq:covhop-source-uv}
has an exact charged-spectrum test independent of any continuation
of the interacting photon. Work at zero spatial momentum in
the continuum, in the selected Euclidean time coordinate. Continue
its frequency to $\omega=i z$, where $z$ is a complex energy variable
on that sector-plane. This continuation tests the intrinsic singularities
of the free resolvent; choosing the directed boundary value on the real
energy axis occurs separately. Let $\mu$ be the marginal law of the
time component of the hop displacement. It is symmetric, supported on
$[-2L,2L]$, and has a density positive on the interior. Define
\[
 G(z)=\int e^{z s}\,d\mu(s),\qquad
 H_\pm(z)=G(z)\pm\frac{G'(z)}{2L}.
\]
Indeed the hop at $\omega=i z$ carries $e^{-zs}$, and its time-spin
weight on the $\gamma_\parallel=\pm1$ subspace is
$1\mp s/(2L)$. Spatial spin terms vanish by spatial reflection.
Changing $s$ to $-s$ gives $H_\pm$.

\begin{proposition}[A positive and a negative charged residue]
\label{prop:upgrade-charged-poles-uv}
At positive rest energy the time-ordered endpoint
$D_0^{-1}\gamma_\parallel$ has two distinct simple poles
$M_+<M_-$ with weights of opposite signs. Each time-spin block
also has infinitely many nonreal poles.
\end{proposition}
\begin{proof}
Each $H_\pm$ is the Laplace transform of the positive probability
measure $(1\pm s/(2L))\,d\mu(s)$. Its second derivative is strictly
positive on the real axis. Its support contains subintervals of
both half axes, so $H_\pm$ tends to infinity at either real
infinity. Since $\lambda H_\pm(0)=\lambda<1$, the equation
$\lambda H_\pm(z)=1$ has exactly one positive and one negative real
root, both simple. On the positive half-axis $G'(z)>0$ and
$H_+>H_-$. This places the positive $+$ root first.
For $\sigma=\pm1$ the corresponding endpoint eigenvalue is
$\sigma L/(1-\lambda H_\sigma(z))$. In the convention
$Z/(M_\sigma-z)$ its spectral weight is
\[
 Z_\sigma=\frac{\sigma L}{\lambda H'_\sigma(M_\sigma)} .
\]
Both derivatives are strictly positive. The pole in the negative
time-spin subspace therefore carries negative weight at
\emph{positive} energy. A positive charged transfer for this same
endpoint would have a positive matrix measure there. An instantaneous
contact cannot change either residue.

The scalar denominator $1-\lambda H_\sigma$ is entire of exponential
type. The Fourier transform of its absolutely continuous hop
measure tends to zero along the real Euclidean frequency axis, so
$1-\lambda H_\sigma(i\omega)\to1$ as $\omega\to\pm\infty$.
If it had only finitely many zeros, the finite-order entire-function
factorization theorem would give $e^{az+b}P(z)$ with $P$ a polynomial.
The two nonzero finite limits force $\operatorname{Im}a=0$ and
$P$ constant. Convergence of the phase then forces $a=0$. The
denominator is nonconstant and has just two real zeros; it must
have infinitely many nonreal zeros. The numerator on each block
is the nonzero constant $\sigma L$.
\end{proof}

\paragraph{Retention of the full charged circle tower.}
The proof also keeps the auxiliary normal circle and its unbounded
mode sum from \eqref{eq:nativefull-three-compressions-uv}.
For $K\ge6$ replace $T$ on the charged hop by an isometry
$T_5:\mathbb R^5\to e_\Gamma^\perp$. The probability law of the
five-component displacement $v$ is the pushforward of the same
two shell draws. Use $W(v)=I-(\gamma\cdot v_{1:4}+
\gamma^5v_5)/(2L)$, so its norm is again at most two.
Wrap the fifth translation on a circle of length $L$ with
the flat holonomy $e^{i\theta v_5/L}$ in the periodic chart.
The Wilson line for the dynamical zero-mode photon uses the four
carrier components $v_{1:4}$. Its gauge transformation is the
same endpoint phase on this circle. The symbol estimate on a
chosen charged hop is now $C(1+L|(p,\xi_j)|)^{-K}$,
$\xi_j=(2\pi j+\theta)/L$. It is summable jointly in
$d^4p$ and $j\in\mathbb Z$ since $K>5$, uniformly in all
carrier momentum shifts. Thus the preceding charged-loop
estimate includes every normal mode, even when the normal
history returns with zero winding. For the soft estimate,
periodize the bounded compactly supported five-dimensional hop
density in the fifth coordinate. Its supremum is finite, and
convolution on $\mathbb R^4\times S^1$ by this probability density
still contracts the supremum norm. Holonomy phases have unit
modulus. The projected four-dimensional polygon closes for
every normal winding, so the photon-current estimates and the
order-by-order massless pole persist for the whole tower.
At trivial holonomy the untwisted charged normal mode has the
same rest-frame poles and the same negative residue as the
four-dimensional Wilson source.

\subsection{The normal Green function and the physical pair cut}
\label{sec:ward-normal-uv}

A charged pole can carry positive weight even when the photon
probes of the same action violate unitarity. To test the two
predictions against one another, we use the separation of two
points drawn from the velocity measure. It fixes a normal heat
factor and thereby determines both the open propagator and the
closed charged trace.

\paragraph{A normal Green function between two integrated shell records.}
Fix a locked sector-plane and the full oriented A4 shell integral
in the ambient action. Write $n=K-1\ge4$; the transverse velocity
coordinate has normalized uniform law on $B^n$.
For this charged representation use two independent shell records
$y,y'$ with that law. Project their difference onto an orthonormal
two-plane by $P_2$, and set $a=L|P_2(y-y')|$, $L=\sqrt{G_N}$ in
the stated carrier units, $m=1/L$. The distribution is rotation
invariant and includes arbitrarily short returns. Independence is a
choice of the boundary pairing before the determinant. The normal
readout is a map to an auxiliary two-plane. If its two axes must
also be orthogonal to all four readout axes in the same shell,
that realization requires $n\ge6$.

At positive normal mass parameter $x$, the two-dimensional
Green function between records a distance $a$ apart, with the
normalization $2\pi(-\Delta+x)^{-1}$, is $K_0(a\sqrt x)$. Define
$\Phi(x)=\mathbb E K_0(a\sqrt x)$ and $W(t)=\mathbb E e^{-a^2/(4t)}$.
If $\chi_n$ is the real radial characteristic function of the unit
ball, put $b(v)=|\chi_n(L\sqrt v)|^2$. The two-dimensional
Gaussian Fourier transform and radial angular integration give
\begin{equation}
\begin{aligned}
 \mathbb E J_0(a\sqrt v)&=b(v),\qquad
 W(t)=t\int_0^\infty e^{-tv}b(v)\,dv,\\
 \Phi(x)&=\frac12\int_0^\infty e^{-tx}W(t)\,\frac{dt}{t}
         =\frac12\int_0^\infty\frac{b(v)}{x+v}\,dv .
\end{aligned}
\label{eq:upgrade-shell-pair-log-uv}
\end{equation}
Here $0\le b\le1$ and $b(0)=1$. In fact
\[
 \lim_{t\downarrow0}\frac{W(t)}t
 = C_b=\int_0^\infty b(v)\,dv=\frac{n^2}{(n-1)L^2}.
\]
Each projected record has density
\[
 \frac{n}{2\pi L^2}(1-|Y|^2/L^2)^{(n-2)/2},
 \qquad |Y|<L ,
\]
and density zero outside the disk.
The integral defining $C_b$ equals $4\pi$ times its squared
$L^2$ norm by Plancherel. These identities retain both individual
records and the full interval of proper times.

Use the four-dimensional transverse shell multiplier
$c_L(Q)=\chi_n(L|Q|)$ and the connection
$B(Q)=[c_L(Q)P_T(Q)+P_L(Q)]A(Q)$. This matrix extends smoothly
to the identity at $Q=0$ since $1-c_L(Q)=O(|Q|^2)$ with a smooth
radial quotient. Gauge shifts of $A$ give the same shifts of $B$.
Let $d_B=\gamma^\mu(\partial_\mu+iB_\mu)$ be the anti-self-adjoint
Euclidean Dirac operator and $H_B=m^2-d_B^2$. Specify the charged
bilinear with kinetic operator
\begin{equation}
 {\cal D}_B=L^{-1}(m+d_B)H_B^{-1/2}e^{-\Phi(H_B)} .
 \label{eq:upgrade-shell-pair-dirac-uv}
\end{equation}
Choosing the exponential of this normal Green function is a
constitutive choice. The prefactor $(m+d_B)H_B^{-1/2}$ is unitary.
For a vectorlike charged field the nonzero eigenvalues of $d_B$
are paired by $\gamma^5$, and its zero modes give polar factor one.
Consequently $\log|{\cal D}_B|=-\log L-\Phi(H_B)$ as an operator
identity in fixed units. The field-dependent part of the parity-even
fermionic effective action comes from varying $\operatorname{Tr}\Phi(H_B)$.
All terms in a photon jet, including the direct quadratic contact
and the spin term, follow from $H_B$. At zero background the four-dimensional heat trace per volume
is proportional to $e^{-m^2t}t^{-2}$. Its integral against
$W(t)dt/t$ diverges as $\int_0^\epsilon t^{-2}\,dt$ at the lower endpoint.
Only this field-independent normalization is dropped.
At constant $B$ the symbol shifts in momentum, so the regular
soft vertex is the momentum derivative of the complete free inverse.
The subsequent ultraviolet test distinguishes absolute convergence of
a charged momentum integral from convergence after its gauge-covariant
heat trace has been computed.

\paragraph{Positive free spectrum and the open charged source.}
The free propagator is $L(m-i\gamma\cdot p)F(m^2+p^2)$, where
$F(z)=e^{\Phi(z)}/\sqrt z$. To see that $F$ has a positive
Stieltjes measure, take $\operatorname{Im}z>0$ and
$\theta=\arg z$. Then
\[
 0\le-\operatorname{Im}\Phi(z)
 \le \frac12\int_0^\infty\frac{\operatorname{Im}z}{|z+v|^2}\,dv
 =\frac{\theta}{2}.
\]
Thus $\arg F$ lies in $[-\theta,-\theta/2]$. The Stieltjes
characterization~\cite{SchillingBernstein2012} applies because $F$ is analytic off the
negative real axis, positive on $(0,\infty)$, and tends to zero
along the positive axis. Write $F(x)=\int d\nu(u)/(x+u)$. Its atom
at $u=0$ is explicit. Set
\[
 c_b=\frac12\int_0^\infty
   [b(v)-\mathbf1_{v<L^{-2}}]\,\frac{dv}{v}.
\]
The integral converges since $b(v)=1+O(L^2v)$ at zero and is
integrable at infinity. The indicator evaluates the logarithmic
constant analytically; it does not truncate any momentum integral. Separating the elementary integral in
\eqref{eq:upgrade-shell-pair-log-uv} gives
$\lim_{x\downarrow0}xF(x)=L^{-1}e^{c_b}$.
The charge-one mass atom has residue $Z_\psi=e^{c_b}>0$.
On $u>0$ the continuous density is
\[
 \frac{1}{\pi\sqrt u}
  \exp\!\left[\frac12\operatorname{PV}\!
                \int_0^\infty\frac{b(v)}{v-u}\,dv\right]
         \cos\!\left(\frac{\pi b(u)}2\right).
\]
This follows from the upper cut boundary of $F$. For squared mass
$m^2+u$ the scalar Dirac weight is $m$ times its positive vector
weight, so the free Dirac spectral positivity inequality holds.

The same functional calculus gives
$\mathcal D_B^{-1}=L(m-d_B)F(H_B)$. With $M=\sqrt{m^2+u}$,
\[
 \frac{m-d_B}{M^2-d_B^2}
 =\frac{1+m/M}{2}(d_B+M)^{-1}
  +\frac{1-m/M}{2}(d_B-M)^{-1}.
\]
This identity holds at nonzero background before forming an open
chain. A chain of $r\ge1$ transverse insertions at fixed momenta
costs at most $M^{-r-1}\prod_a|c_L(Q_a)|$ up to a constant.
Its mass integral converges since
$\int d\nu(u)/(m^2+u)=F(m^2)<\infty$. Thus the pole component
and all the marked continuum components have assigned source vertices.

\paragraph{The open pole and the inclusive pair cut.}
The ordinary spin, orbital and direct terms of the four-dimensional
heat trace in \eqref{eq:monomial-wound-bubble-uv} give the closed scalar polarization before the two external
shell factors as
$\Pi(X)=\pi^{-2}\int_0^1 h\Phi(m^2+hX)\,du$, $h=u(1-u)$,
$X=|Q_E|^2$. This normalization reproduces the unit Dirac cut
when the heat weight equals one. From
\eqref{eq:upgrade-shell-pair-log-uv} the absorptive density at
timelike $S>4m^2$, in that sign convention, is exactly
\begin{equation}
 \rho_{\rm pair}(S)=\frac{1}{2\pi^2}
   \int_{(1-\beta)/2}^{(1+\beta)/2}
   h\, b(S h-m^2)\,du,\qquad \beta=\sqrt{1-4m^2/S}.
 \label{eq:upgrade-shell-pair-cut-uv}
\end{equation}
This cut is positive. It is strictly smaller than the unit Dirac
density $\rho_{\rm D}(S)=\beta(1+2m^2/S)/(12\pi^2)$ at every
such $S$. Indeed $b(v)<1$ for $v>0$ because the projected
single-record measure has volume support. Replacing $b$ by one
in the integral gives $\rho_{\rm D}$. The analytic transverse
shell readout multiplies both densities in this test by the same
strictly positive square at the timelike momentum.

The normalized open source gives an independent measurement.
Its mass atom contributes exactly $Z_\psi(d_B+m)^{-1}$.
At one insertion the two external pole factors therefore carry
the usual Dirac pair vertex and the same shell multiplier.
The continuum has no double mass atom; the external factors
$x/(x+u)$ and $y/(y+u)$ select $u=0$ when the two mass
parameters tend to their poles from below. The tail is controlled
by the open-chain bound. Normalizing both spinors cancels
$Z_\psi$. If this pair is a positive physical channel in the usual crossed interpretation and the
closed current cut sums the squared norms of all such channels, the closed
spectral weight must be at least the pole-pair weight. Equation
\eqref{eq:upgrade-shell-pair-cut-uv} violates that requirement
already at first nontrivial electromagnetic order. The conclusion
uses the pole component in the specified free source; an isolated
interacting charged pole is not assumed.

For comparison a fixed separation $a=L$ gives the stronger
short-time factor $e^{-L^2/(4t)}$ but replaces $b(v)$ in the cut
by $J_0(L\sqrt v)$. Angular integration yields
$\beta[\,\sin Y/Y-\beta^2(\sin Y-Y\cos Y)/Y^3\,]/(8\pi^2)$,
where $Y=L\sqrt{S-4m^2}/2$. This expression is negative at
$Y=3\pi/2+2\pi j$, $j=0,1,\ldots$.
The two independently integrated records remove these negative
bands while leaving the strict deficit relative to the open
pole pair. A separate change in the physical source is needed
if its open amplitude and its closed cut are to agree. More generally,
for any nonnegative law of separations with a positive-residue
mass atom and the same minimally gauged mass decomposition, the
Dirac cut integrand is multiplied by $\mathbb E J_0(a\sqrt v)$.
The angular representation gives $J_0(x)<1$ for $x>0$. Every
such law with $\Pr(a>0)>0$ produces a strict deficit at every
energy above threshold.

\paragraph{A joint momentum bound for the paired-shell determinant.}
Use the paired normal weight $W(t)$ above, the four-dimensional
sector-plane pullback, and $c_L(q)=\chi_{K-1}(L|q|)$ on a transverse
photon. The following proof keeps the circulating charged momentum
as an integration variable. It applies to the source functional
normalized at zero background. 
\begin{proposition}[Convergence of complete normal heat vertices]
\label{prop:ward-normal-graphs-uv}
For the representation \eqref{eq:upgrade-shell-pair-dirac-uv},
each source-dependent Euclidean graph assembled from complete
charged heat vertices admits an absolutely integrable rooted
representation at every fixed electromagnetic order. The bound
covers simultaneous short durations and independent hard photon
momenta. The transverse photon has a massless perturbative pole.
\end{proposition}

Fix a root $x$ in the diagonal spin heat kernel of $H_B=m^2-d_B^2$.
In its multilinear variation a transverse source
$\epsilon_a e^{iq_a\cdot(x+y)}$ may be replaced by
\[
 \epsilon_a e^{iq_a\cdot x}(e^{iq_a\cdot y}-1).
\]
The removed constant is a pure gauge on the Euclidean plane; the
two phases at a diagonal root cancel for every smooth background.
This identity holds at each fixed positive proper time for the diagonal kernel and
extends to complex source directions by multilinearity. The
integration over $x$ still gives conservation of the original
momenta. In a spin insertion the curl supplies its own $q_a$
and differentiates away the constant. Every remaining potential
provides a factor $e^{iq_a\cdot y}-1$. Its divergence vanishes
because $q_a\cdot\epsilon_a=0$.

In an ordered heat-kernel term let $o,f,c$ count the orbital, spin
and direct two-source insertions, respectively. Write
$j=o+f+c$, $N=o+f+2c$, and let $u_0,\ldots,u_j$ be the $j+1$
successive free heat intervals with sum $t$. Represent each of
the $d=o+2c$ potential differences by
\[
 e^{iq\cdot y}-1
 =\int_0^1 (q\cdot\partial_r)e^{ir\cdot y}
              \big|_{r=\theta q}\,d\theta .
\]
The $r$'s are taken independent even for the two potentials at a
direct contact. The mixed derivatives therefore introduce one
$|q_a|$ per potential. The rooted Fourier kernel has one independent
charged momentum $k$; its intervals have momenta $k+P_l$, where
$P_l$ is a partial sum of insertion transfers, including the
interpolated $r$'s. No momentum-conserving delta is differentiated
at the root. Each orbital insertion contributes at most one linear
numerator on its outgoing free interval. These intervals are
distinct; the spin and direct insertions have no loop numerator.
In particular a mark can differentiate an orbital numerator as
well as a Gaussian.

For a unit vector $\epsilon$, $e\in\{0,1\}$, and a multi-index
$\beta$, Gaussian scaling gives
\[
 \left|\partial_p^\beta
    \big[(\epsilon\cdot p)^e e^{-u|p|^2}\big]\right|
 \le C_\beta u^{(|\beta|-e)/2}e^{-u|p|^2/2}.
\]
Apply this inequality by Leibniz's rule before integrating the
absolute kernel. For each placement of the $d$ marks, the factors
$u_l^{|\beta_l|/2}$ have product at most $t^{d/2}$. Completing the
square gives uniformly in the interpolation parameters and all
original momenta
\[
 \int_{\mathbb R^4}\!
 e^{-\frac12\sum_lu_l|k+P_l|^2}\,d^4k\le (2\pi)^2t^{-2}.
\]
Thus the joint integral over $k$ and the interpolation parameters
is bounded by a constant times
$t^{-2+d/2}\prod_{l\,\mathrm{orbital}}u_l^{-1/2}
\prod_{a=1}^N |q_a c_L(q_a)|$ before the time-simplex integration.
The simplex integral of the inverse square roots scales as
$t^{j-o/2}$ and is finite on every face, including intersections.
Since $j+(d-o)/2=N$, the full rooted density satisfies
\begin{equation}
 \int |\mathcal I_N(t;k,u,\theta;q)|\,d^4k\,du\,d\theta
 \ \le\ C_N e^{-m^2t}t^{N-2}\prod_{a=1}^N |q_a c_L(q_a)| .
 \label{eq:upgrade-paired-root-bound-uv}
\end{equation}
Here $du$ denotes the simplex measure at total time $t$ and the
finite spin trace is included in $C_N$. Multiplication by
$W(t)dt/(2t)$ is integrable for $N\ge2$, since $W(t)=O(t)$ at
zero and $0\le W(t)\le1$. The one-point term vanishes by charge
conjugation. The zero-source vacuum term has the short-time cost
given above. This construction proves joint absolute convergence
in a gauge-completed representation obtained by an exact identity
at each positive $t$. Individual terms in the usual unrooted
orbital and contact expansion need not have absolutely integrable
charged momentum and time near their common ultraviolet boundary.

In a connected graph of closed charged loops, each photon edge
now carries the scalar majorant $C|c_L(q)|^2$ after the two factors
$|q|$ absorb its $|q|^{-2}$ propagator. This majorant is bounded
and belongs to $L^1(\mathbb R^4)$ for $K>4$ and an isometric
rank-four readout, since its tail is $O(|q|^{-K})$. Choose a
spanning tree after contracting each charged loop to a vertex.
The photon momenta on edges outside the tree, including self-loops,
are independent integration variables; bound the tree factors
uniformly and integrate the other factors in $L^1$. This bound
integrates all charged and photon loop variables, including
regions in which only a subset become large.

Each external transverse photon on a closed charged component
retains its own $|q|$. The proper two-point function consequently
vanishes to order $|q|^2$ at every finite order. There is a finite
directional limit of its quotient by $|q|^2$. In the rooted
formula divide each of the two external marks or curls by $|q|$,
then take $q$ to zero along a ray. At fixed positive times the
charged Gaussian integrals are continuous in their translations;
on compact sets of transfers they admit integrable Gaussian
envelopes. Bound \eqref{eq:upgrade-paired-root-bound-uv} then
permits dominated convergence over the simplex and proper time
after the charged momentum integration. For the remaining photon
integral use the preceding spanning-tree majorant, choosing the
chord momenta independent of the external transfer. An internal
photon in a proper graph belongs to a cycle, so the directions of
its transverse projector are defined almost everywhere at the
limit. Rotation covariance and the Ward identity give a finite
coefficient of the transverse Maxwell eigenvalue. Its tree value
is one, hence its formal perturbative inverse retains the massless
pole.

The positive mass decomposition of the open source above gives
a uniform finite bound at fixed transfers for every marked
open chain. With tested external endpoints an open chain has no
independent charged circulation. Each end of a hard photon still
carries $|c_L(q)|$ and at most one power of $|q|$. The same tree
estimate controls the hard integrals for these graphs with the
chosen infrared prescription for charged external states.

The unshifted loop routing gives a useful check on the scope of
this result. At a constant connection its two-source integrand
contains $\partial_{k_i}\partial_{k_j}\Phi(m^2+|k|^2)$.
Since $b$ is integrable, differentiation under its Stieltjes
integral gives $\Phi^{(j)}(x)\sim(-1)^j j!C_b/(2x^{j+1})$
for $j=0,1,2$. The leading Hessian is therefore
\[
 C_b |k|^{-4}(-\delta_{ij}+4\widehat k_i\widehat k_j).
\]
For $i=j$ its absolute integral on an angular region of positive
measure diverges logarithmically. At each positive heat time the
constant connection is exactly pure gauge. The rooted expression
uses that identity before the time integration and supplies an
absolute bound for a complete gauge-covariant loop. An arbitrary
order of integration of the ordinary unrooted Taylor diagrams
would give a stronger claim than this proof supports.

\paragraph{Summing the charged circle before the short-time limit.}
For a photon independent of the normal circle coordinate let
$\xi_j=(2\pi j+\theta)/L$ and
$d_{5,B,j}=d_B+i\gamma^5\xi_j$. Since $\{d_B,\gamma^5\}=0$,
the positive block is $H_{5,B,j}=m^2+\xi_j^2-d_B^2$. Each
source derivative of its heat kernel therefore has precisely the
same four-dimensional integrand, multiplied by $e^{-t\xi_j^2}$.
Comparison of the shifted Gaussian sum with its integral gives
$\sum_{j\in\mathbb Z}e^{-t\xi_j^2}\le C(1+L/\sqrt t)$,
uniformly in the holonomy angle. The worst new power near zero
in the absolute joint bound is $t^{N-5/2}dt$; already the
two-source coefficient is integrable. The photon spanning-tree
estimate and the soft two-point limit thus include all charged
modes of this parity-even determinant. The source-free trace
density still diverges. At trivial holonomy
$\gamma^5d_{5,B,j}\gamma^5=-d_{5,B,-j}$, which pairs the polar
phases in the full determinant; the zero mode is the vectorlike
four-dimensional block. Moreover, on
$4m^2<S<4(m^2+(2\pi/L)^2)$ only that zero mode reaches its pair
threshold. Its optical deficit remains throughout this interval
even after the entire circle tower is restored. The separate
high-energy four-photon test below is for the four-dimensional
representation. The circle bound is sharp for this family of
flat charged normal modes. With $d$ independent periodic
directions the small-time trace gains a factor of order
$t^{-d/2}$. A nonzero transverse magnetic test has a nonzero
quadratic coefficient of order $t^0$ in the four-dimensional
Dirac heat trace, as in \eqref{eq:monomial-wound-bubble-uv}.
Since $W(t)/t\to C_b>0$, its complete field-dependent coefficient
then contains $\int_0^1 t^{-d/2}\,dt$, divergent for $d\ge2$.
The normal Green function between two marks does not itself add
a tower of charged dynamical fields.

\paragraph{A limitation on improving the normal small-time power.}
Consider more generally a positive probability law of the
separation $a>0$ in the same logarithm. Suppose $W(t)=o(t)$
as $t\downarrow0$ and $\mathbb E a^{-1/2}<\infty$.
The second hypothesis follows, for example, if $W(t)=O(t)$:
$\Pr(a<r)\le e^{1/4}W(r^2)$ at small $r$.
The scalar polarization before the common shell factor then cannot
have an everywhere nonnegative spectral density. Indeed Laplace
integration gives $x\Phi(x)\to0$. The bounded quantities
$Xh\Phi(m^2+hX)$ tend to zero for almost every $0<u<1$,
so $X\Pi(X)\to0$ as well.
The exact unsubtracted dispersion formula is
\[
\begin{aligned}
 \Pi(X)&=\int_{4m^2}^{\infty}\frac{\rho(S)\,dS}{X+S},\\
 \rho(S)&=\frac{1}{2\pi^2}
   \int_{(1-\beta)/2}^{(1+\beta)/2}
     h\,\mathbb E J_0(a\sqrt{S h-m^2})\,du .
\end{aligned}
\]
To justify it, use
$2K_0(a\sqrt x)=\int_0^\infty J_0(a\sqrt v)\,dv/(v+x)$
and then $v=S h-m^2$. The bound
$|J_0(a\sqrt v)|\le C\min(1,a^{-1/2}v^{-1/4})$
makes the average absolutely integrable against $dv/(m^2+v)$.
If $\rho$ were nonnegative, $X\Pi(X)$ would increase to
$\int\rho(S)\,dS>0$ since $\Pi(X)>0$. The assumed faster
short-time decay gives the opposite limit. Thus simply
strengthening the small-displacement suppression of this
spectral ansatz cannot recover both that stronger ultraviolet
power and a positive current cut. The result does not assume
a physical pole-pair interpretation.

\paragraph{Rigidity of a single covariantly gauged scalar spectrum.}
\label{par:ward-scalar-rigidity-uv}
The open--closed comparison extends beyond the normal Green function.
Suppose the complete open covariance is
$(m-d_B)F(m^2-d_B^2)$ up to a positive constant, $F$ is a nonzero
Stieltjes function, and $F$ has a positive atom at zero. Allow any nonnegative continuum measure and any local constant
term in the Stieltjes representation. Suppose also that this same Gaussian charged field generates the closed current
by its parity-even determinant. Define
\[
 \Phi(z)=\log(\sqrt z\,F(z)),\qquad
 b_F(v)=-1-\frac{2}{\pi}\arg F(-v+i0).
\]
Use the analytic logarithm real on the positive axis, with constants
expressed in fixed carrier units. In the upper
half-plane $-\arg z\le\arg F(z)\le0$ by the Stieltjes formula.
Boundary values exist almost everywhere and $-1\le b_F\le1$.
More explicitly, for $x,x_*>0$,
\[
 \Phi(x)-\Phi(x_*)=
 \frac12\int_0^\infty b_F(v)
       \left(\frac{1}{x+v}-\frac{1}{x_*+v}\right)dv .
\]
Indeed the imaginary part of $\Phi$ is bounded harmonic, is zero
on the positive axis and has boundary values $-\pi b_F/2$ on the
negative axis. The Poisson formula gives the imaginary part of
the integral, and evaluation at $x_*$ fixes its real constant.
This subtraction compares values of the analytic symbol; its
absorptive discontinuity is independent of every constant.

The constant in the symbol contributes only a constant trace.
For the resolvent in the integral use the mass derivative of the
Dirac heat trace, $-\partial_v e^{-t(H_B+v)}=t e^{-t(H_B+v)}$.
The two-mark trace therefore has exactly the cut
in \eqref{eq:upgrade-shell-pair-cut-uv} with $b_F$ in place
of $b$. One may use its difference at two Euclidean momenta
to derive the formula even when the unvaried polarization diverges:
the scalar difference is
$\pi^{-2}\int_0^1 h[\Phi(m^2+hX)-\Phi(m^2+hX_*)]du$.
For real positive $X,X_*$, the resolvent difference has an
integrable $O((m^2+v)^{-2})$ majorant in $v$. The bounded
imaginary part permits taking the cut under the
$u$ integral. The common transverse shell factor cancels from
the comparison. In the open source the mass atom alone gives the double-pole
vertex of the unit-charge Dirac current after endpoint normalization;
a local constant term has no double pole. A positive inclusive
cut containing this pair would require $\rho_F(S)\ge\rho_{\rm D}(S)$
for every $S>4m^2$. But $b_F\le1$, so equality is necessary.
For each such $S$ the averaging weight on
$0<v<S/4-m^2$ is strictly positive after the change
$v=S u(1-u)-m^2$. Hence $b_F=1$ almost everywhere on
$(0,\infty)$. Substitution in the difference formula gives
$\Phi(x)=C-\tfrac12\log x$ in fixed units, or $F(x)=e^C/x$.

Thus the two-source optical requirement uniquely selects the
ordinary gauged Dirac resolvent in this single-scalar functional
class. Its bare polarization coefficient contains
$\int_0^1 e^{-m^2t}dt/t$ for nonzero external shell readout.
This excludes a simultaneous solution by changing only the
spectral function of this covariantly gauged operator. Independent
transverse current structure or additional charged dynamics would
require its own open--closed comparison.

\paragraph{Can a scalar electric form factor supply the missing weight?}
Let $s_\star=4m^2$ and $r(S)$ be the ratio of the shell-difference
cut to the Dirac pole-pair cut after removing the common external
readout. The preceding calculation gives $0<r(S)<1$ on the
entire open cut. Suppose the only change in the open source is a
scalar form factor $f(S)$ multiplying its on-shell Dirac matrix
element. Impose unit charge $f(0)=1$, real analyticity and
bounded analyticity on the physical sheet
$\mathbb C\setminus[s_\star,\infty)$, with boundary values almost
everywhere. Keeping the closed response forces
$|f(S+i0)|^2\le r(S)$ by positivity of the inclusive channel sum.
Map the cut plane to the disk by sending $S=0$ to its center:
\[
 z=\frac{\sqrt{s_\star}-\sqrt{s_\star-S}}
         {\sqrt{s_\star}+\sqrt{s_\star-S}} .
\]
On either rim $S=s_\star\sec^2(\arg z/2)$. The disk mean-value
formula at the center, followed by dominated convergence and
real analyticity, implies
\[
 |f(0)|\le
 \frac{\sqrt{s_\star}}{\pi}\int_{s_\star}^{\infty}
 \frac{\sqrt{r(S)}\,dS}{S\sqrt{S-s_\star}} < 1 .
\]
The integration weight has total mass one and $r<1$ on every
compact interval inside the cut. Thus this particular change
in the electric vertex cannot recover unit charge. More general
vertices, new physical channels and changes to the closed trace
require their own open and closed amplitudes.

\paragraph{An equal-helicity test of the shell-difference determinant.}
Let the scalar weight of a four-dimensional spinor heat trace be
\(W(t)=t\int_0^\infty e^{-t\lambda}b(\lambda)\,d\lambda\), with
\(b(\lambda)=|\chi_n(L\sqrt{\lambda})|^2\), \(L,m>0\), and \(n>1\).
For its four-photon coefficient, the usual massive Dirac continuation
gives a superposition of Dirac discontinuities
with weight \(-b(\lambda)\partial_{M^2}\) at \(M^2=m^2+\lambda\).
Indeed, \(-\partial_{M^2}e^{-tM^2}=te^{-tM^2}\); the statement about
cuts uses the same Feynman boundary value at every mass. At fixed timelike
energy the two-particle cut has bounded mass support. Real local polynomial
terms have zero discontinuity.

Take two incoming real photons of equal circular helicity. Set
\(S=(q_1+q_2)^2>0\), \(z=\cos\theta\),
\(\delta=4M^2/S\), \(\beta=\sqrt{1-\delta}\), and
\(\alpha_{\rm em}=e^2/(4\pi)\). With the outgoing spins summed, the
Born formula for a single unit-charge, colorless Dirac fermion is
\begin{align}
 \frac{d\sigma^D_0}{dz}
 &= \frac{4\pi\alpha_{\rm em}^2}{S}
    \frac{\beta(1-\beta^4)}{(1-\beta^2z^2)^2},
    \qquad 0<\delta<1 .
 \label{eq:ward-bw-angular-uv}
\end{align}
This is the \(J_z=0\) channel of Ref.~\cite{WardBreitWheeler},
Eq.~(2), with the color multiplicity removed and
\(dt/dz=S\beta/2\).
Integration over the full angular interval yields
\(\sigma^D_0=(4\pi\alpha_{\rm em}^2/S)F(\delta)\), where
\begin{align}
 F(\delta)&=\beta(1+\beta^2)
       +(1-\beta^4)\operatorname{artanh}\beta,
       &&0<\delta<1, \nonumber\\
 F'(\delta)&=2\beta^2\operatorname{artanh}\beta
              -2\beta-\beta^{-1}.
 \label{eq:ward-bw-deriv-uv}
\end{align}
Extend \(F\) by zero for \(\delta\ge1\). Since \(F(1)=0\),
differentiation across threshold produces only an integrable square-root
singularity. For \(S>4m^2\) the signed forward cut, in cross-section
normalization and with any common external form factors stripped off, is
\begin{equation}
 \Sigma_0(S)=-\frac{16\pi\alpha_{\rm em}^2}{S^2}
    \int_0^{S/4-m^2} b(\lambda)
       F'\!\left(\frac{4(m^2+\lambda)}{S}\right)\,d\lambda .
 \label{eq:ward-bw-signed-uv}
\end{equation}
At small \(\delta\),
\(F'(\delta)=\log(4/\delta)-3+
O(\delta|\log\delta|)\).
The shell-transform estimate
\(b(\lambda)\le C(1+L^2\lambda)^{-(n+1)/2}\)
makes \(b\) integrable, and \(b(0)=1\).
On \(\delta\le1/2\), the ratio
\(F'(\delta)/\log(S/m^2)\) is uniformly bounded for large \(S\)
and tends to one at each fixed \(\lambda\).
On \(1/2\le\delta\le1\), \(|F'|\le C(1-\delta)^{-1/2}\),
and the integral of \(b|F'|\) there is \(O(S^{-(n-1)/2})\).
Dominated convergence therefore gives
\begin{equation}
 \lim_{S\to\infty}\frac{S^2\Sigma_0(S)}{\log(S/m^2)}
 =-16\pi\alpha_{\rm em}^2\int_0^\infty b(\lambda)\,d\lambda<0 .
 \label{eq:ward-bw-sign-uv}
\end{equation}
For the shell law the integral on the right is
$C_b=n^2/[(n-1)L^2]$ as evaluated in the normal Green function.
For the two-current cut the Dirac threshold factor is
\((1+\delta/2)\sqrt{1-\delta}\), whose negative mass derivative
is nonnegative. The four-current test probes a stronger requirement:
a positive physical Hilbert space and a unitary photon scattering operator
require a nonnegative leading forward absorptive part for this initial
polarization. For the isometric shell readout each real photon has multiplier
$\chi_n(0)=1$. A common external factor
\(|f(q_1)f(q_2)|^2\) preserves the sign wherever it is nonzero.
Thus this scalar-weighted Dirac continuation by itself fails the physical
test. An additional charged channel or a different continuation would
require a new calculation of the complete cut.

\paragraph{A positive-transfer comparison with a complete static source.}
A spatial realization on a chosen time slice isolates an exact cancellation
from the oriented-shell measure. Let $K\ge4$, $L=\sqrt{G_N}$ and $m=L^{-1}$
in carrier units. With $y$ sampled from the transverse A4 unit-ball law
$d\beta$, normalize its projection onto three orthonormal spatial axes to
obtain $a(y)\in S^2$. The zero projection has measure zero. By rotational
invariance $a$ is uniform, so $\langle a_i a_j\rangle=\delta_{ij}/3$.
Let $\pi=-i\nabla-A$, $B=\nabla\times A$, and use the Pauli matrices
$\sigma$ on spin and $\tau$ on two energy bands. Set
\[
 P_a=\pi^2-(\sigma\cdot a)(a\cdot B),\qquad
 E_A=\int (m^2+P_a)^{1/2}\,d\beta .
\]
The two spin blocks of $P_a$ are the expectations of the positive
operator $(\sigma\cdot\pi)^2$ in the constant spin states along $\pm a$.
Thus $E_A\ge m$. To give the free soft particle the Dirac linear term,
let $\chi$ be the radial characteristic function of $d\beta$ and use
\[
\begin{aligned}
 F(s)&=\arctan(s/m)\chi(L|s|)^2,\\
 U_A&=e^{-i\tau_2 F(\sigma\cdot\pi)/2},\\
 h_A&=U_A\tau_3E_AU_A^* .
\end{aligned}
\]
For a smooth static field all operators are spatially gauge covariant.
At $A=0$ the filled negative band has positive excitation energies
$\sqrt{m^2+|k|^2}$. The linear free spatial term is
$\tau_1\sigma\cdot k$. Static occupied energies are independent of $U_A$.

On a flat periodic spatial box with smooth $A$, the complete
field-dependent one-matter-loop vacuum energy of this Hamiltonian is
\[
 \int_0^\infty \frac{e^{-m^2t}\,dt}{2\sqrt\pi\,t^{3/2}}
   \int d\beta\,\operatorname{Tr}(e^{-tP_a(A)}-e^{-tP_a(0)}).
\]
This formula includes all variations of the square root, including direct
quadratic and higher contacts. The constant heat coefficient cancels.
The coefficient of $t$ vanishes by the spin trace. Dividing out the
common spin and free heat factors, the next local coefficient is
\[
 \tfrac12(a\cdot B)^2-\tfrac16|B|^2 .
\]
Indeed, for $P=-(\nabla_A^2+S)$ the flat coefficient~\cite[Eq.~(4.28)]{WardHeatManual} contains
$S^2/2+[\nabla_i,\nabla_j][\nabla_i,\nabla_j]/12+\nabla_A^2S/6$.
The curvature contraction is $-2|B|^2$ and the derivative has zero spin
trace. The shell average therefore cancels the whole coefficient
pointwise. Its heat-trace remainder is $O(t^{3/2})$, making the displayed
integral absolutely convergent at zero for arbitrary smooth static $A$.
Differentiating the smooth-parameter heat expansion gives the same bound
for each fixed number of insertions. At large $t$ positivity and $m>0$
control the trace: writing $t=ns$ with $1\le s\le2$, differentiation of
$n$ contractive semigroup factors costs only a polynomial in $n$ at each
fixed order. Every static one-matter-loop coefficient is therefore finite.
This comparison selects a spatial spin compression. Identifying the
resulting Hamiltonian and its covariant charged current from the native
action, controlling internal dynamical photon momenta, and proving
positivity for that interacting theory require additional arguments.

\subsection{Testing changes in the charged current}
\label{sec:nov-transverse-uv}

A deficit in pair production has two possible sources. The matrix
element that creates a pair can be too large, or the closed trace
can give too little inclusive weight. We first hold the trace fixed
and vary the on-shell matrix element. We then recompute both from
one charged operator. All comparisons retain the post-locking
support $\bin=(s,b_s)$ and the prescribed vectorial pole limits of
that support. The Lorentzian variable $S$ is the invariant squared
mass on its Minkowski readout; the absolute heat estimates use
Euclidean momenta after the sector-plane reconstruction.

\paragraph{An analytic bound on the static magnetic moment.}
Hold the four-dimensional paired-shell determinant fixed. Work in units
$mL=1$ at $K=249$, write $n=248$, $\tau=S/(4m^2)$, and strip the
common transverse photon readout. Let $F_1,F_2$ be Dirac and Pauli
form factors, real analytic off the cut $[1,\infty)$ and bounded
there in the disk coordinate below. Unit electric charge fixes
$F_1(0)=1$. Put $G_M=F_1+F_2$, $G_E=F_1+\tau F_2$ and let $R$ be
the inclusive determinant density divided by its unit Dirac value.
In the pair rest frame the magnetic and electric Born terms have angular
weights \(1+z^2\) and \((1-z^2)/\tau\), with \(z=\cos\theta\).
Their integrals have ratio \(2\tau:1\). The proton analysis in
\cite[Eq.~(2)]{NextSachs} exhibits these weights with a common
final-state Coulomb factor. Here only the spin weights enter the
unit-Dirac normalization. Positivity then gives the pole-pair bound on each rim,
\begin{equation}
 2\tau |G_M|^2+|G_E|^2\le(2\tau+1)R(\tau).
 \label{eq:next-moment-sachs-uv}
\end{equation}
In particular, if $F_1=1$ on the cut, completing the square in
$F_2$ bounds the pair ratio from below by
$2(\tau-1)^2/[(\tau+2)(2\tau+1)]$, attained at $F_2=-3/(\tau+2)$.
Already at $\tau=40$ this lower bound is $169/189$. To bound the
shell ratio there, set $v=39(1-x^2)$, $H=1-39x^2/40$.
With squared mass shifts in units of $m^2$, write
$b(v)=|\chi_n(\sqrt v)|^2$. The ratio is the $H$-weighted
average of $b(v)$ over $0<x<1$.
For the difference of two independent ball coordinates, the
cosine bound through degree four and the single-coordinate moments
$1/(n+2),\,3/[(n+2)(n+4)]$ give
$b(v)\le1-v/250+v^2/125000$. Integrating this polynomial yields
$R(40)\le13056437/14765625<0.89<169/189$.
The first excited circle threshold at trivial holonomy is at
$\tau=1+4\pi^2>40$. The same bound therefore applies to the
charged circle tower at this energy.

For independent $F_1,F_2$ map the cut plane to the unit disk by
$z=(1-\sqrt{1-\tau})/(1+\sqrt{1-\tau})$, so $\tau=4z/(1+z)^2$.
Set $a=\sqrt3+\sqrt2$ and
\[
 V(z)=\frac{(\sqrt8\,G_M,\,(1+z)G_E)}{a+a^{-1}z}.
\]
On the unit circle $\tau=\sec^2(\arg z/2)$ and
$|a+a^{-1}z|^2=8+4/\tau$. Hence $\|V\|^2\le R$.
Project onto the unit vector parallel to $V(0)$ and apply the
logarithmic mean inequality. One may first multiply by
$(1+z)/(1+\eta z)$, $0\le\eta<1$, to give a bounded analytic
function even at the image of infinity. The logarithmic mean of
this multiplier is zero and its value at the origin is one.
The result, with the two lips identified by reality, is
\begin{equation}
 \begin{aligned}
 8|G_M(0)|^2+1&\le a^2\exp\!\int_1^\infty\log R\,d\mu,\\
 d\mu(\tau)&=\frac{d\tau}{\pi\tau\sqrt{\tau-1}} .
 \end{aligned}
 \label{eq:next-moment-mean-uv}
\end{equation}

The shell determines a numerical upper bound without fitting its shape.
In the dimensionless variable $v$, for $\beta=\sqrt{1-1/\tau}$,
\[
 R(\tau)=\frac{3}{4\tau^2\beta(1+1/(2\tau))}
 \int_0^{\tau-1}\frac{(1+v)b(v)}{\sqrt{1-(1+v)/\tau}}\,dv .
\]
For $h_n(y)=n(1-|y|^2)_+^{(n-2)/2}/(2\pi)$ on the normal plane,
planar Plancherel gives
\[
 \begin{aligned}
 M_1:=\int_0^\infty(1+v)b(v)\,dv
 &=4\pi(\|h_n\|_2^2+\|\nabla h_n\|_2^2)\\
 &=\frac{n^2}{n-1}+\frac{n^2(n-2)}{n-3}.
 \end{aligned}
\]
Take $A=256$, $p_A=\mu([\!A,\infty))=2\arcsin(1/\sqrt A)/\pi$
and write $\mathbb E_A$ for conditioning on this interval. Exchanging
the two positive integrals in $\mathbb E_A[\tau^2 R]$ uses, for $u=1+v$,
\[
 \int_{\max(A,u)}^\infty
 \frac{\tau^{-3/2}\,d\tau}{\sqrt{1-u/\tau}}\le\frac{\pi}{\sqrt A}.
\]
This follows by substituting $\sqrt{u/\tau}$; for $u\le A$ use
$\arcsin x\le\pi x/2$. Using $p_A\ge2/(\pi\sqrt A)$ gives
$\mathbb E_A[\tau^2R]\le C_A=3\pi M_1/[8(1-1/A)]$.
Also $p_y/p_A\ge\sqrt{(1-1/A)A/y}$ for $y\ge A$ by the elementary
bounds on $\arcsin$. Integrating this tail with $y=Ae^t$ gives
$\mathbb E_A[\log(\tau/A)]\ge2\sqrt{1-1/A}$. Since $R\le1$,
concavity of the logarithm now gives
\[
 \int\log R\,d\mu\le
 p_A[\log(C_A/A^2)-4\sqrt{1-1/A}]<-0.154 .
\]
Equation~\eqref{eq:next-moment-mean-uv} consequently gives
$|G_M(0)|<0.968$ in units of the charge-one Dirac moment.

\paragraph{An explicit charge-normalized pair amplitude within the bound.}
Consider the bounded analytic functions
\[
 F_1(z)=-F_2(z)=\frac{(1+z)^3}{(1-3z/4)^2(1+3z/4)} .
\]
They have $F_1(0)=1$, $G_M=0$ and $G_E=(1-\tau)F_1$.
On the cut the proposed pole-pair ratio is
\[
 G(\tau)=
 \frac{2^{18}(\tau-1)^2}{(\tau+48)(2\tau+1)(49\tau-48)^2}
 < \min(3/4,\,55/\tau^2) .
\]
The inequality follows from $(49\tau-48)>49(\tau-1)$.
For $0\le v\le15$ the shell moment bound
$\chi_n(\sqrt v)\ge1-v/500$ gives $b(v)\ge c=(1-15/500)^2$.
Thus $R\ge c$ on $1\le\tau\le16$ (with its threshold limit).
For $\tau\ge16$, express $R$ as the ratio of the preceding
$v$ integral to its value at $b=1$. The denominator is at most
$4\tau^2/3$ and the numerator is at least
$c\int_0^{15}(1+v)dv$, giving $R>89/\tau^2$. Therefore
$G<R$ everywhere above threshold. These open matrix elements meet
the two-current inequality of the specified determinant. A common action would also have to supply the intermediate states
responsible for their timelike cut. The original determinant still
has the negative four-photon cut found in the preceding section.

\paragraph{Gauging a magnetic moment in both the determinant and the open source.}
A transverse spin coupling tests a way out of scalar spectral rigidity.
On the selected Euclidean four-plane write
\[
 \begin{aligned}
 K_B&=m+d_B+P_B,\qquad H_B=K_B^\dagger K_B,\\
 P_B&=\frac{\kappa}{4m}\sigma^E_{\mu\nu}F_{\mu\nu}[B],
 \qquad \sigma^E_{\mu\nu}=\frac i2[\gamma_\mu,\gamma_\nu].
 \end{aligned}
\]
The sign of the real magnetic moment is fixed by the convention
$G_M(0)=1+\kappa$ for the corresponding continued Dirac vertex;
reversing the sign of $P_B$ reverses $\kappa$ in that vertex.
The shell-read connection $B$ and the probability law of its
normal separation are as above. Consider its polar completion
$\mathcal D_B=L^{-1} K_BH_B^{-1/2}\exp[-\Phi(H_B)]$ near zero
background. Its open inverse is $L F(H_B)K_B^\dagger$,
$F(x)=e^{\Phi(x)}/\sqrt{x}$. Its parity-even fermionic action is
$\operatorname{Tr}\Phi(H_B)$, up to a source-independent constant.
Thus every insertion of $P_B$ also changes the closed trace,
including its direct two-source contact. A mass atom
$A=\lim_{x\downarrow0}xF(x)>0$ contributes $LAK_B^{-1}$ on the
open line. After external normalization its charged pair has the
full Dirac--Pauli vertex. The continuum contributes no double
mass atom, by the resolvent representation at the two open ends.

The ultraviolet behavior can be tested before continuing the energy.
For smooth compactly supported backgrounds, the field-dependent part
of the first local heat coefficient of $H_B$, traced over spin, is
$\operatorname{tr}P_B^2$. To verify this result, write $H_B$ as
$-\nabla_B^2+[P_B,\gamma_\mu]\nabla_{B,\mu}$ plus its zeroth-order
terms. Complete the square in the shifted connection
$\nabla_{B,\mu}-[P_B,\gamma_\mu]/2$. The ordinary spin-curvature
term, the term linear in $mP_B$, and derivatives linear in $P_B$
have zero spin trace. At quadratic order the traced coefficient is
\[
 -\operatorname{tr}P_B^2
 -\tfrac14\sum_\mu\operatorname{tr}[P_B,\gamma_\mu]^2
 =\operatorname{tr}P_B^2 .
\]
Here $\sum_\mu\gamma_\mu P_B\gamma_\mu=0$ for a Clifford bivector in
four dimensions. In particular
$\operatorname{tr}P_B^2=\kappa^2 F_{\mu\nu}[B]F_{\mu\nu}[B]/(2m^2)$.
The complete trace contains $(4\pi t)^{-2}t\int\operatorname{tr}P_B^2$.
For the two independent shell marks $W(t)\sim C_b t$, so the
quadratic field response contains a positive logarithm,
\begin{equation}
 \frac{C_b}{32\pi^2}
 \int d^4x\,\operatorname{tr}P_B^2
 \int_0^\epsilon\frac{dt}{t}.
 \label{eq:next-pauli-heat-uv}
\end{equation}
This divergence is present for every nonzero Pauli coefficient on a
background with nonzero readout curvature. Fixed nonzero separation
does suppress this short-time integral.

\paragraph{The pair cut of the same modified trace.}
The absorptive test works even when the local ultraviolet coefficient
diverges, since the displayed local term in $B$ has zero discontinuity
in the timelike test after the common shell factors are removed. More generally take the transverse Dirac and magnetic
coefficients at a given timelike $S$ to be $f,k$ after stripping the
common external readout. For the preceding operator they are $1,\kappa$.
The argument permits real analytic photon form factors with the same
boundary values on the tested interval and no independent absorptive
cut there. They act on the external photon, independently of the
charged mass variable.

The identity $K_B^\dagger\gamma^5=\gamma^5K_B$ gives
\[
 H_B+v=(K_B+i\sqrt v\,\gamma^5)^\dagger
                 (K_B+i\sqrt v\,\gamma^5).
\]
Hence the parity-even reference at $v$ has the pair cut of a Dirac
field with the pseudoscalar mass $i\sqrt v\,\gamma^5$ in addition
to $m$. Chiral rotation gives the real mass $M=\sqrt{m^2+v}$,
leaves the vector current intact, and rotates its magnetic tensor by
$\cos\vartheta=m/M$, $\sin\vartheta=\sqrt v/M$ into a magnetic
and an electric dipole tensor. A parity-odd local Jacobian has no
absorptive part in this test. Put $E=\sqrt S/2$, $\tau=E^2/m^2$
and $\beta_v=\sqrt{1-M^2/E^2}$. In the pair rest frame take the
particle momentum along the third axis. After dividing the current
matrix elements by $2E$, the two transverse spin amplitudes have
coefficient $f+k$. The longitudinal triplet has coefficient
$(M/E)f+(E/M)k$, and the longitudinal singlet has magnitude
$|k|(E\beta_v/m)\sin\vartheta$. These follow by inserting the
ordinary rest-frame spinors of mass $M$. The magnetic tensor
interchanges the transverse and longitudinal coefficients $2E,2M$
of the Dirac matrix element, while the tensor times $\gamma^5$
has only the singlet matrix element of magnitude $2E\beta_v$.
The spin sums, angular average and two-body phase space then give
\begin{equation}
 \begin{aligned}
 12\pi^2\rho_v(S)
  &=\tfrac32 |f+k|^2\beta_v\\
  &\quad+\tfrac12(\tau|k|^2-|f|^2)\beta_v^3 .
 \end{aligned}
 \label{eq:next-pauli-twist-uv}
\end{equation}
Set $\rho_v=0$ for $v\ge E^2-m^2$. With $k=0$ and $f=1$ this
reduces to $\beta_v(1+2M^2/S)/(12\pi^2)$. The threshold
singularity in its mass derivative is integrable.

Differentiating at fixed $S,f,k$ gives
\begin{equation}
 \begin{aligned}
 -16\pi^2 E^2\beta_v\,\partial_v\rho_v
 &= (1-\beta_v^2)|f|^2+2\operatorname{Re}(\bar f k)\\
 &\quad+(1+\tau\beta_v^2)|k|^2 .
 \end{aligned}
 \label{eq:next-pauli-mass-derivative-uv}
\end{equation}
The determinant of this two-dimensional Hermitian form is
$\beta_v^2[\tau(1-\beta_v^2)-1]=\beta_v^2v/m^2$.
Both diagonal entries are positive. Thus $-\partial_v\rho_v>0$
throughout $0<v<E^2-m^2$ whenever $(f,k)\ne(0,0)$.
For the paired-shell logarithm the same mass-derivative identity
used for the four-photon cut gives the determinant density
$\int_0^{E^2-m^2} b(v)(-\partial_v\rho_v)\,dv$.
Since $b(v)<1$ for $v>0$, the determinant has strictly less
spectral weight than $\rho_0(S)$, the contribution of its own
open mass-atom pair. The same inequality holds with
$b(v)=\mathbb E J_0(a\sqrt v)$ for any positive law with
$\Pr(a>0)>0$ for which the stated determinant boundary value and
open atom exist. In particular, a fixed nonzero separation still
fails this two-source test.

For a general positive Stieltjes open spectrum with a mass atom
the bounded-phase representation of $\Phi=\log(\sqrt{x}F(x))$
gives $b_F(v)\le1$ as in the scalar rigidity argument. If the
pair inequality is required at unboundedly many energies with
$(f,k)\ne(0,0)$, strict positivity in
\eqref{eq:next-pauli-mass-derivative-uv} forces $b_F=1$ almost
everywhere on the positive axis. It follows that $F(x)=C/x$.
The parity-even determinant is then the ordinary Dirac--Pauli
determinant. For a nonzero local $\kappa$, the same first heat
coefficient now yields a power divergence. At $\kappa=0$ the
ordinary Dirac charge logarithm remains. This calculation rules
out a unitary ultraviolet-finite completion by this particular
magnetic deformation of the normal logarithm. Couplings to
independent charged histories, or form factors with their own
physical cuts, have different source data.

\paragraph{What fixed separation recovers.}
Keeping a fixed normal displacement $a=L$ in this same Dirac--Pauli
model gives a counterterm-free convergence result on the Euclidean
plane for $K>6$. Each charged heat vertex still contains all contacts
of $K_B^\dagger K_B$. Use the rooted constant-connection identity
of the preceding joint graph bound. An undifferentiated transverse
potential obtains one factor of $|q|$ and a heat displacement,
while $F[B]$ already contributes $|q|$ and $\partial F[B]$ at
most $|q|^2$. In the expanded operator every linear mark
differentiates at most one adjacent free kernel. The direct
quadratic terms, including $P_B^2$, have no charged derivative.
The same Gaussian derivative and time-simplex bounds therefore
give, for each finite number of marks, a finite sum of powers of
$t$ times $e^{-m^2t}\prod_j |q_j|(1+L|q_j|)|c_L(q_j)|$,
with constants depending on $m,L$ and the order. Each remaining
inverse power at zero is integrable against
$e^{-L^2/(4t)}dt/t$. Every sewn photon edge is bounded by a
constant times $(1+L|q|)^2|c_L(q)|^2$, bounded at zero and
integrable over $\mathbb R^4$ since $K>6$. Spanning-tree
integration thus controls all charged and photon loop momenta.
The two external factors of $|q|$ and dominated convergence retain
a finite perturbative Maxwell coefficient. For tested open
lines the identity
\[
 (H_B+v)^{-1}K_B^\dagger
   =\tfrac12\sum_{\eta=\pm1}(K_B+\eta i\sqrt v\,\gamma^5)^{-1}
\]
gives the marked-chain estimate with one factor
$(1+L|q_j|)|c_L(q_j)|$ per insertion and $M^{-r-1}$ on a chain
of $r\ge1$ marks. The positive free Stieltjes measure for fixed
separation has $\int M^{-2}\,d\nu(v)<\infty$. This proves the
ultraviolet bound in the same infrared regime used above. The
fixed-separation source still violates the optical bound because
its mass weight is $J_0(L\sqrt v)<1$ for every $v>0$.

\paragraph{Independent spinor parities and charged flavor mixing.}
The single-scalar restriction can be removed while keeping minimal
gauging by the same Euclidean Dirac operator $d_B$ on the selected
sector-plane. Consider $N$ charge-one flavors with open covariance
$G(d_B)$. Let $G(z)$ be an $N\times N$ analytic matrix in the
upper half-plane, with $\operatorname{Im}G(z)<0$ there and
$G(\bar z)=G(z)^*$. These are the negative-imaginary-part conventions for matrix Herglotz functions~\cite{NextHerglotz}. This sign holds, for example, for positive
matrix weights of $(z+M)^{-1}$ and $(z-M)^{-1}$, $M>0$, whenever
the weights together see all flavors. These two terms allow
independently weighted Dirac parities. Flavor matrices at distinct
masses may fail to commute. Assume a gap about zero where $G$ is
analytic and Hermitian on the real axis, with $G(0)$ invertible.

\begin{proposition}[Optical lower bound on the bare charge logarithm]
\label{prop:parity-optical-log-uv}
Suppose the same Gaussian covariance generates the parity-even
charged determinant. Let an atom of rank $r\ge1$ at a signed mass
of magnitude $m>0$ give $r$ orthogonal charge-one particles with
the usual crossed pair interpretation. Use the same nonzero
transverse shell readout in the open current and its inclusive
cut. If the determinant includes the positive contributions of
these pairs at every energy above threshold and its physical
cut is nonnegative elsewhere, its bare quadratic transverse
response at fixed nonzero Euclidean transfer has a logarithmic
ultraviolet divergence. One charged mass atom suffices, even
when the other flavor directions have no such atom.
\end{proposition}
\begin{proof}
For $\operatorname{Im}x>0$ choose $z=-i\sqrt{x}$ with the
principal square root. The scalar symbol of the parity-even
logarithm is
\[
 \Psi(x)=\tfrac12\{\operatorname{tr}\log G(z)
                        +\operatorname{tr}\log G(-z)\}.
\]
The two logarithms use arguments in $(0,\pi)$ and $(-\pi,0)$,
respectively. On $x>0$ they are adjoints. Pairing the two Dirac
eigenvalues gives precisely $\operatorname{Tr}\Psi(-d_B^2)$
for the parity-even fermion effective action. Define the
real phase density $b(v)=-2\operatorname{Im}\Psi(-v+i0)/\pi$,
$v>0$. Each trace of a logarithm is the sum of the logarithms of
the eigenvalues with algebraic multiplicity, so $|b|\le N$.
The boundary value exists almost everywhere because the imaginary
part is bounded harmonic. Near $v=0$, the negative inertias
of $G(\sqrt v)$ and $G(-\sqrt v)$ agree, hence $b=0$ there.
The Poisson formula, with its real constant fixed at $x_*>0$,
gives
\[
 \Psi(x)-\Psi(x_*)=
 \tfrac12\int_0^\infty b(v)
       \left(\frac1{x+v}-\frac1{x_*+v}\right)dv .
\]
Only this real constant drops out in the source derivatives.
The complete Dirac spin heat trace, with its direct contact,
therefore has the weight $W(t)=t\int_0^\infty b(v)e^{-tv}dv$.
Its two-current density, with the identical external factors
divided out, is
\[
 \rho(S)=\frac1{2\pi^2}\int_0^1 h\, b(Sh)\,du,
 \qquad h=u(1-u).
\]
This discontinuity also follows from the resolvent difference
in the preceding display without first asserting that the bare
quadratic coefficient exists.

A signed mass pole contributes its flavor residue times a
minimally gauged Dirac resolvent. The simultaneous pole
coefficient of its first source derivative is the ordinary
Dirac current between two such resolvents. Unit normalization
on the range of the residue cancels the external overlaps.
The two signed mass conventions give the same pair-production
weight. Thus $\rho(S)\ge r\rho_{\rm D}(S)$ for $S>4m^2$, with
$\rho_{\rm D}(S)\to1/(12\pi^2)$. This is a lower bound even
in the presence of additional physical channels.

Fix $v_*>0$ and $S_0>4m^2$. For $R>S_0$, change variables
$v=Sh$ in the logarithmic energy average. Uniform boundedness
of $b$ and integrability of $h|\log h|$ give
\[
 \int_{S_0}^{R}\rho(S)\frac{dS}{S}
 =\frac1{12\pi^2}\int_{v_*}^{R} b(v)\frac{dv}{v}+O(1).
\]
Indeed changing $Rh$ to $R$ costs at most $N|\log h|$, and
changing $S_0h$ to $v_*$ has an analogous bound. The optical
lower bound and the explicit Dirac threshold factor now imply
\[
 \int_{v_*}^{R} b(v)\frac{dv}{v}
       \ \ge\ r\log(R/v_*)-O(1).
\]
At fixed $X=|Q|^2>0$ the bare polarization before the external
readout is the improper integral with integrand
$(2\pi^2)^{-1}\int_0^1h\,e^{-thX}W(t)\,du/t$.
Integrating over $t\ge\epsilon>0$ gives the inner integral
$\int b(v)e^{-\epsilon(v+hX)}\,dv/(v+hX)$.
It differs from $\int_{v_*}^{1/\epsilon} b(v)\,dv/v$
by $O(1)$ uniformly in $u$ as $\epsilon\downarrow0$.
Near zero this uses the gap in $b$. At large $v$ the shift of
the denominator costs $O(X/v^2)$; the exponential tail above
$1/\epsilon$ and the replacement of $e^{-\epsilon v}$ by
one below that point have bounded logarithmic integrals.
The polarization thus grows at least as
$r\log(1/(\epsilon v_*))/(12\pi^2)-O(1)$.
Here $\epsilon$ tests the lower endpoint of the original
proper-time integral. No subtraction is used to assign that
integral a finite value.
\end{proof}

There is also an equality case. If the mass $m$ is the bottom of
the spectrum, assume $G$ is Hermitian analytic on $(-m,m)$.
Its derivative there is negative semidefinite by the Cauchy
representation of a matrix with the stated imaginary-part sign.
Consequently, except at isolated zeros,
\[
 b(v)=n_-\!\left(G(-\sqrt v)\right)
         -n_-\!\left(G(\sqrt v)\right)\le0,\qquad v<m^2 ,
\]
where $n_-$ counts negative eigenvalues. Above $m^2$ the upper
bound is $N$. If $r=N$, the optical lower bound requires equality
in the integral against $h$ for each $S>4m^2$. Changing from
$u$ to $v=Sh$ gives strictly positive weight throughout
$0<v<S/4$. Hence $b=N\,\mathbf1_{\{v>m^2\}}$ almost everywhere
and $\Psi(x)=C-(N/2)\log(x+m^2)$ in fixed units. This fixes the
parity-even determinant to that of $N$ Dirac fields. It does not
identify the individual open flavor propagators.

For example, a single flavor with positive residues $3/4$ and
$1/4$ at opposite signed masses has
\[
 G(z)=\frac{3}{4(z+m)}+\frac{1}{4(z-m)}
     =\frac{z-m/2}{z^2-m^2}.
\]
Both open species have physical mass $m$. Nevertheless the zero of
the numerator contributes a negative Dirac cut with mass $m/2$.
The determinant's density is
$2\rho_{\rm D}(S;m)-\rho_{\rm D}(S;m/2)$, strictly negative on
$m^2<S<4m^2$, where the open source has no pair at all. In
contrast, putting these same two poles on independent flavor
diagonals gives the ordinary positive two-species cut and two
bare Dirac charge logarithms.

\subsection{Positive charged-pair production on a spatial slice}
\label{sec:nov-positive-uv}

An independent transverse pair interaction can change both the
open amplitude and the closed spectral sum. Here a single
Hamiltonian defines them. The GTL delta selects the support, and
the projected displacements
are sampled from the original velocity-shell measure. The additional
spin pairings and spatial readouts specified below provide a
comparison whose physical successes and failures can both be
calculated.

\paragraph{A Gauss charge and a positive pair channel from the same shell.}
Work on a fixed $\bin=(s,b_s)$ and one of its spatial slices. Keep
the ambient GTL pullback and normalized oriented A4 measure. For
$n=K-1\ge3$ let $h(p)=\chi_n(L|p|)^2$ on three orthonormal
readout axes, $L=\sqrt{G_N}$ in carrier units, $m=1/L$. This is
the characteristic function of the projected difference of two
independent shell records. The transverse readout
$h(1-h)=h-h^2$ is the difference between sampling once and twice
with this law. These readouts, the time slice and the charged
pairing below are additional choices.

Use positive-energy particle and antiparticle spin slots $b,d$,
with energy $E_k=\sqrt{|k|^2+m^2}$, and a transverse Maxwell
field of frequency $r=|p|$. Denote the free Fock Hamiltonian
by $H_0$. Define the three triplet pair annihilators $C_i(p)$ by
\[
 C_i(p)=\int_k h(k)h(p-k)\,
             b(k)^{\mathsf T}i\sigma_2\sigma_i d(p-k) ,
 \qquad \int_k=\int\frac{d^3k}{(2\pi)^3}.
\]
Let $Y_i(p)$ be the magnetic photon annihilator with one-photon
matrix element
\[
 y_i(p,\epsilon)=L^{-1}h(p)[1-h(p)]
              \frac{i(p\times\epsilon)_i}{\sqrt{2r}} .
\]
These annihilators are spatial functions of the canonical
transverse potential and electric field at the chosen time. Define
a fourth pair mark $C_{\rm s}$ by replacing $i\sigma_2\sigma_i$ in
$C_i$ by the singlet tensor $i\sigma_2$. In the same spatial
chart let
\[
 Y_2(x)= L\,(h\mathbf E^{(+)})(x)\cdot(h\mathbf B^{(+)})(x).
\]
The superscript denotes the annihilation part of the canonical
operator. Its two photon factors act on separate shell records. In momentum
space $Y_2(p)$ denotes the spatial Fourier transform of $Y_2(x)$. Write $\rho(q)$ for the difference
of the particle and antiparticle number densities in Fourier space:
on a fixed particle sector each summand translates one momentum
by $q$ with the sign of that charge. Specify the action in Hamiltonian
form by the usual charged and Maxwell symplectic terms and
\begin{equation}
\begin{aligned}
 H_e={}&H_0+L^2\int_p\sum_i (C_i+eY_i)^\dagger(C_i+eY_i)\\
 &+L^2\int_p(C_{\rm s}+e^2Y_2)^\dagger(C_{\rm s}+e^2Y_2)\\
 &+\frac{e^2}{2}\int_q\frac{|h(q)|^2}{|q|^2}
                  \rho(q)^\dagger\rho(q).
\end{aligned}
\label{eq:positive-gauss-comparison-uv}
\end{equation}
Here $e$ counts electric insertions and can be set to the chosen
primitive charge. In particular the pair interaction $L^2 C^\dagger C$
already acts at $e=0$. The positive Coulomb square follows by solving
Gauss' law for the electric source $e h(q)\rho(q)$ and includes its
diagonal self-energy. Its zero-transfer charge on one particle is $e$.
Gauge covariance before reduction can use the dressed field
$e^{-ie h\varphi}b$ with opposite phase on $d$, where
$\varphi=\Delta^{-1}\nabla\cdot A$ and $h$ acts by spatial convolution.
Under $A\mapsto A+\nabla\alpha$ the undressed slots transform by
their corresponding $e^{\pm ie h\alpha}$; the invariant temporal
source is $h(A_0-\partial_t\varphi)$. This prescription fixes the
longitudinal dressing. It does not impose a minimal transverse
electric coupling. The positive squares retain both pair interactions,
the one- and two-photon transitions, the quadratic and quartic
photon contacts, and all Gauss contractions. This is a non-Gaussian charged model.

The total charge is conserved. Each vertex changes the auxiliary
count $N=N_b+N_d+2N_\gamma$ by at most two. On an incoming sector
of fixed count the Coulomb form has operator norm at
most $(e^2N^2/2)\int_q |h(q)|^2/|q|^2<\infty$.
For two charged legs with fixed total momentum $p$, the relative
momentum functional in $C_i$ has norm bounded uniformly by a spin
constant times $\|h\|_2\|h\|_\infty$. The photon multiplier $y$ is
bounded. Each leg of $Y_2$ has envelope
$g(k)=\sqrt{|k|}\,|h(k)|\in L^2\cap L^\infty$, so its analogous
relative-momentum norm is at most a constant times
$\|g\|_2\|g\|_\infty$. Summation over finitely many occupied
legs now bounds every interaction block from a fixed sector. The same estimates
hold with absolute values on every individual contraction kernel:
the Coulomb interaction is bounded by convolution with an $L^1$
function, while the pair maps obey the same Cauchy--Schwarz bound in
relative momentum. Consequently each ordered contraction on a finite
time interval, between finite-particle packets, has an absolutely
convergent integral over its spatial momenta and ordered times.
The charge-independent pair interactions preserve the count;
on the finitely many sectors visited at a fixed electric order,
their expansions sum absolutely by the simplex factorial.
The estimates cover intermediate sectors reached by finitely many
external charged fields or transverse probes. They concern the
canonical time-ordered representation. A bound on separately routed
four-frequency integrands does not follow from the time integrals.

For the full Hamiltonian view the three sets of annihilation or
density maps in the squares as operators into spatial $L^2$ of
Hilbert-space vectors, with $h(q)/|q|$ included in the density
map. They are closable: finitely many sectors in the target
belong to the adjoint domain by the same block bounds and the
finite change in count. Sum their closed squared-norm forms
with the free positive form on their dense common domain. This
gives a translation-invariant self-adjoint Hamiltonian with the
vacuum at zero and $H_e\ge H_0\ge|\mathbf P|$.
Couple an external real transverse probe by shifting the canonical
photon field in both squares, with conjugate shifts for the adjoints.
This fixes the one-source current and every direct source contact.
The resulting current on the vacuum includes the charged pair
$C_i^\dagger|0\rangle$, with nonzero matrix element at generic
nonzero momentum and transverse polarization. The same Hamiltonian
evolves that state and the open charge insertions. Hence the
inclusive spectral measure of each tested Hermitian current is its
actual sum of positive squared projections. Finite-time evolution
is unitary. These statements do not require general scattering
wave operators or a covariant identification of asymptotic currents.

Two real photons can create a pair in the singlet channel. For
back-to-back momenta and equal circular helicities their symmetrized
matrix element in $\mathbf E^{(+)}\cdot\mathbf B^{(+)}$ is nonzero
where $h(p)\ne0$: both magnetic polarization vectors are the
same helicity-dependent multiple of their electric vectors.
At order $e^2$, only $C_{\rm s}^\dagger Y_2$ connects this initial
state to a single pair. The triplet terms preserve $N$ while
these two sectors have counts four and two. At zero electromagnetic coupling the neutral
singlet pair block at total momentum zero is
$D_0+L^2|v\rangle\langle v|$ with $D_0=2E_k$, where $v$ is the
singlet form factor of $C_{\rm s}$. Put
$G(z)=\langle v,(D_0-z)^{-1}v\rangle$. The rank-one resolvent
gives the inclusive singlet weight per unit energy proportional to
\[
 \frac{\operatorname{Im}G(E+i0)}
      {\pi\,|1+L^2G(E+i0)|^2}.
\]
The free numerator is positive at generic $E>2m$ since $v(k)$
is proportional to $h(k)^2$ on the energy sphere. This gives an
explicit, nonzero, positive leading two-real-photon absorptive
channel, with the same singlet interaction as in the open states.
It supplies a sign and sewing test in this representation; its
magnitude is specific to this interaction. At the pair threshold
$\operatorname{Im}G(E+i0)$ is a positive constant times
$\sqrt{E-2m}$ to leading order, since $h(0)=1$. The real
threshold value of $G$ is positive and finite. The incoming
shell factor is nonzero at $|p|=m$, since the single-record
characteristic function there averages a positive cosine for
$mL=1$. Thus the leading threshold weight is genuinely nonzero
immediately above $2m$.

\paragraph{Two mass-shell tests in the positive comparison.}
At fixed total momentum $p\ne0$ the space of one transverse
photon and one spin-triplet pair is reducing: the singlet
annihilator vanishes on the latter and $Y_2$ kills both.
The free and Coulomb operators preserve spin, and the triplet
coupling acts within this space. On this subspace let
$T_p$ be the pair annihilation map to its three spin components,
$y_p$ the analogous map from the two photon polarizations, and
$D_p$ the free pair energy plus its complete positive Coulomb
energy. In particular $D_p\ge 2m$. The two one-photon polarization trials have
trial energies at most $r+e^2 L^2\|y_p\|^2$, below $2m$ for small
$r$. The orthogonal pair subspace has Rayleigh quotient at least
$2m$. The min--max principle gives two eigenvalues counting multiplicity, isolated below $2m$ in this
reducing subspace seen by the transverse one-photon probe. For either eigenvalue $z$, the
Schur complement on the photon space is
\begin{equation}
\begin{aligned}
 (z-r)x&=e^2 L^2 y_p^\dagger (I+L^2 M_p(z))^{-1}y_p x ,\\
 M_p(z)&=T_p(D_p-z)^{-1}T_p^\dagger .
\end{aligned}
\label{eq:positive-gauss-speed-uv}
\end{equation}
This follows by eliminating the pair block
$D_p+L^2T_p^\dagger T_p-z$. Uniform boundedness of $T_p$ bounds
$M_p(z)$ uniformly for small $r$, when $z<m$.
Also $y_p^\dagger y_p=L^{-2}r h(p)^2[1-h(p)]^2 I/2$,
and $1-h(p)=L^2r^2/(n+2)+O(L^4r^4)$ by the second shell moment.
Thus for fixed nonzero $e$ and sufficiently small $r$,
\[
 0<c\,e^2L^4r^5\le z-r\le C\,e^2L^4r^5 .
\]
Elimination of the pair block also bounds the norm of the pair
component by $O(r^{5/2})$ times the photon norm at fixed $L,e$.
The positive photon poles therefore have limiting speed one and
their combined relative photon weight tends to one. At finite
nonzero soft momentum their energies exceed $r$.

In the count-one sector of charge $+e$ the conversion square
vanishes and $\rho(q)^\dagger\rho(q)=I$. Hence a single charged
state has the exact energy
\[
 E_{\rm ch}(k)=\sqrt{|k|^2+m^2}+\Delta_{\rm C},\qquad
 \Delta_{\rm C}=\frac{e^2}{2}\int_q\frac{|h(q)|^2}{|q|^2}>0 .
\]
This dispersion disagrees away from rest with a relativistic
mass shell of mass $m+\Delta_{\rm C}$. In particular its kinetic
mass from the small-$k$ curvature remains $m$. This finite
mismatch follows with the self-contraction retained. Together the
two tests exhibit a model with charge, UV convergence, positive
current weights and a gapless photon of unit limiting speed, and
identify where its charged and photon dispersions still disagree
with exact sector-plane kinematics.

\paragraph{Restoring the moving charge's transverse current.}
In \eqref{eq:positive-gauss-comparison-uv} the magnetic transition
has zero transverse matrix element between one-charge states. The
spatial charge motion can instead be gauged by the same shell that
measures Gauss' law. For a particle or hole use the positive kinetic
form
\[
 m\|\psi\|^2+\frac12\int\nu_m(dz)\,
                  \|(U_{\pm A}(z)-I)\psi\|^2 .
\]
Here $U_{\pm A}(z)$ is translation by $z$ with the straight Wilson
link of $\pm e hA$. The temporal readout $\pm e h A_0$ gives the same
Coulomb square. Retain the longitudinal dressings of the pair marks
before Gauss reduction. The mass determines the positive jump measure
\[
 \nu_m(dz)=d^3z\int_0^\infty
 \frac{e^{-m^2t}}{2\sqrt\pi\,t^{3/2}}
       (4\pi t)^{-3/2}e^{-|z|^2/(4t)}\,dt .
\]
Indeed, integrating $1-\cos(k\cdot z)$ gives $E_k-m$. The Wilson
translations are unitary and $\int\min(1,|z|^2)\nu_m(dz)<\infty$.
The difference map to $L^2(\nu_m)$ is closed by almost-everywhere
subsequence passage. Smooth finite-particle tests belong to its
domain by the first displacement estimate on their translation
and the smeared field. Sum the closed positive forms as above.

At order $a\ge1$ the kinetic vertex is, up to charge signs and
photon normalization, the polarization contraction of
\[
 \prod_j h(q_j)\int_{[0,1]^a}
   (\nabla^{a}E)\!\left(k+\sum_j t_jq_j\right)\,d^at .
\]
Each $t_j$ is an insertion position. Differentiating the jump
identity gives all contacts. At $a=1$ combine the opposed jumps
$z,-z$ first: the sine contributes the second power of $|z|$
at zero. For $a\ge2$ the jump integral has that power
already. The $a$th derivative is bounded by $C_a m^{1-a}$
uniformly in $k$. The soft one-photon matrix element is therefore
the transverse projection of $\pm e\,k/E_k$, as required for
the free moving charge. For a dynamical photon restore the factor
$1/\sqrt{2|q_j|}$ at each insertion. The resulting envelope is
square-integrable in spatial momentum, so each new emission,
absorption and contraction satisfies the finite-sector absolute
momentum bound. Only finitely many sectors occur at fixed order. The one-current
triplet creation and the leading two-photon singlet cut persist:
the kinetic vertices conserve matter number and act as the
identity on pair spin, so they cannot convert a triplet into
the singlet at order $e^2$. This change restores the leading
charged convection vertex without field-dependent counterterms.

Radiation by moving charges destroys the reducing one-photon
and triplet-pair subspace. There is nevertheless a full interacting
soft estimate. In a neutral fiber at total momentum $p$, states
without matter have free photon energy at least $r=|p|$.
The remaining matter sectors have energy at least $2m$ from the
gauged kinetic forms alone. These kinetic forms preserve matter
number; the conversion and Gauss forms are nonnegative. It
follows that $H_e\ge\min(r,2m)$. For $0<r<2m$ the bare transverse
photon of unit polarization has expectation at most
$r+C e^2 L^4r^5$ for sufficiently small $r$: the singlet square,
Gauss and matter kinetic forms vanish on that trial state. Hence
its normalized spectral measure $\mu_{p,\epsilon}$ satisfies
\[
 \mu_{p,\epsilon}([(1+\delta)r,\infty))
       \le C e^2 L^4r^4/\delta,\qquad \delta>0 .
\]
Thus $E/r$ concentrates at one in this source as $r\downarrow0$.
An isolated photon pole and a relativistic charged dispersion have
yet to be derived for this kinetic energy. The charge density of
both preparations also has the following exact boost test.

\paragraph{A boost test for the current that couples to Gauss' law.}
The Wilson kinetic term supplies a moving charge with its leading
convection current. The failure of a different matrix element gives
a sharper test of the same theory. In both Hamiltonians the temporal
charge insertion at zero external source is $J^0(q)=e h(q)\rho(q)$
in the reduced variables. It annihilates the vacuum for every $q$.
Time evolution preserves this identity, since the vacuum is an
eigenvector. A spatial transverse current, on the other hand,
creates a triplet pair with nonzero weight for suitable packets.
If this particular $J^\mu$ were a Lorentz four-vector on the
selected sector-plane with an invariant vacuum, boosting the
identity $J^0(f)|0\rangle=0$ along any spatial axis would give
\[
 \big[\cosh\eta\,J^0(f_\eta)
       \mathbin{\pm}\sinh\eta\,J^i(f_\eta)\big]|0\rangle=0
\]
for every test $f$, where $f_\eta$ is its boosted test.
Since the boost maps tests onto tests, this would force
$J^i(f)|0\rangle=0$ for every $i$ and $f$. The measured pair
channel rules that out. No ultraviolet limit enters the argument.

A covariant physical charge density for a representation with this
pair channel must therefore itself have vacuum-to-pair matrix
elements at nonzero spatial momentum. Its integrated charge can
still equal the number charge. For example, a polarization coupled
to the electric field contributes a spatial divergence to the
density and a time derivative to the current. Both matrix elements
and the contacts generated by dynamical electric-field sewing
would have to follow from the ambient action. Changing the
integrated charge, the mass parameter or the pole prescription
alone cannot supply this missing matrix element.

\subsection{Sewing a positive physical current}
\label{sec:nov-ring-uv}

The time-ordered model uses vertex times and spatial momenta. To
identify it with covariant bare O--QED one would also need a
prescription for integrating a virtual four-momentum shared by
two physical current responses. The vacuum graph with two leading
charged polarizations has three independent loop momenta, including
the momentum of the connecting photons. It gives a necessary test
of that prescription at order $e^4$. It uses only a positive
timelike band; a charged-particle pole is unnecessary.

\paragraph{A fourth-order test of the physical photon sewing.}
A cancellation in a single closed polarization can be tested with
two identical charged loops. Work on one reconstructed sector-plane.
Write $X=|p|^2$ for the Euclidean invariant and $K_T(X)=X\Pi(X)$
for the complete transverse current Hessian, including its contacts.
Suppose that, up to a common sign, $\Pi$ has a positive spectral
dispersion with kernel $(S+X)^{-1}$, finitely many subtractions, and
nonzero measure on a compact timelike mass band. Local covariant
transverse polynomial tensors add polynomials to $\Pi$. A separate
mass contact $cP_T$ would give $c/X$ in $\Pi$; it is excluded here.
The hypotheses concern the measured physical current.

Choose $N\ge1$ large enough to remove the polynomial terms and to
difference the spectral integral. For
$\Delta_h f(X)=f(X+h)-f(X)$, set $h=X$ after differencing. Then
\[
 |\Delta_X^N\Pi(X)|
 =\int\frac{N!X^N\,d\rho(S)}
               {\prod_{j=0}^N[S+(j+1)X]}
 \ \ge\ \frac{c}{X}
\]
for all sufficiently large $X$, by restriction to the occupied
band. The square of the finite difference is at most
$C_N\sum_{j=0}^N|\Pi((j+1)X)|^2$. Integration against $X\,dX$
and rescaling thus give
\begin{equation}
 \int_{X_0}^{(N+1)R} X|\Pi(X)|^2\,dX
 \ \ge\ c_N\log(R/X_0),\qquad R\ge X_0 .
 \label{eq:normal-physical-ring-uv}
\end{equation}
With ordinary Maxwell sewing the left side, up to positive
normalization, is the absolute radial integral for two distinct
complete charged loops joined by two photon lines. The symmetry
factor of this connected fourth-order ring is nonzero. Ward contacts
on both charged loops have already been combined. Under a graphwise
finiteness requirement, a different connected topology cannot cancel
its absolute integral.

This argument also tests a common internal geometric record $\xi$
with a normalized positive trace. At a fixed momentum suppose the
same $\xi$ is retained around the ring, the Euclidean Hessian
$K_\xi$ is Hermitian and the positive photon covariance is $D_\xi$.
Its two-polarization ring is the squared Hilbert--Schmidt norm of
$D_\xi^{1/2}K_\xi D_\xi^{1/2}$, averaged over $\xi$. Let $u_\xi$ be
the physical port in a fixed transverse polarization and suppose
\[
 |u_\xi\rangle\langle u_\xi|\ \preceq\ C X D_\xi,\qquad
 \mathbb E_\xi\langle u_\xi,K_\xi u_\xi\rangle=X\Pi(X)
\]
for large $X$, with $C$ uniform. These are conditions on the actual
sewing and readout. The first holds if the positive photon spectral
measure includes a massless residue uniformly dominating this port.
Writing $u_\xi=D_\xi^{1/2}v_\xi$ with $\|v_\xi\|^2\le CX$ gives
\[
 |\Pi(X)|^2\le C^2\mathbb E_\xi
       \|D_\xi^{1/2}K_\xi D_\xi^{1/2}\|_{\rm HS}^2 .
\]
The estimate follows from the quadratic-form inequality and
Cauchy--Schwarz; it also holds by approximation on finite test
projections whenever the traced continuum form exists. Thus
\eqref{eq:normal-physical-ring-uv} tests even a shared geometric
label with signed two-port responses. A proposed geometry that
avoids this conclusion must change at least one of the identified
physical overlaps, the ring contraction, or the spectral hypotheses.

\paragraph{A test with two separately marked charged copies}
A common geometric record and independently integrated records can
enter the same source. Their roles can be read directly from a
mixed derivative. On a finite real transverse photon projection,
condition on a common record $\eta$. The photon $A$ is centered
Gaussian with covariance $D_\eta$. Let two identical charged
copies have private records independent of one another and of $A$
conditional on $\eta$. Suppose their first marked coefficients in
the complete action are
\[
 q_i=c_i+\tfrac12 A^{\mathsf T}K_i A,\qquad K_i=K_i^{\mathsf T},
\]
with finite second moments under the reference law. Each $K_i$
includes its direct two-photon source variation. Couple the copies
with separate counting variables $g_i$ and assume the action has
no mixed $g_1g_2$ term, including the logarithm of any moving
boundary density. Interpret the derivative below as extraction of a
perturbative coefficient. For an ordinary derivative of $\log Z$,
assume that the integral exists near zero and permits the two
differentiations. The marks distinguish copies of the source. Write
$B_i=D_\eta^{1/2}K_iD_\eta^{1/2}$, $\bar B=\mathbb E[B_i\mid\eta]$
and $m=\mathbb E[c_i\mid\eta]+\tfrac12\Tr\bar B$. Then
\[
 \left.\partial_{g_1}\partial_{g_2}\log Z\right|_0
     =\tfrac12\mathbb E\|\bar B\|_{\rm HS}^2
                          +\operatorname{Var}(m).
\]
Indeed the mixed coefficient is $\operatorname{Cov}(q_1,q_2)$.
The covariance conditional on all records is
$\tfrac12\Tr(B_1B_2)$ by Gaussian pairing. The remaining
covariance of the two conditional means reduces to
$\operatorname{Var}(m)$ by independence of the private records.
In particular a scalar vacuum score cannot cancel the mixed
two-photon ring in this preparation. If the same private record
is instead retained around both copies, its ring uses
$\mathbb E\|B_i\|_{\rm HS}^2$; the increase is the nonnegative
$\mathbb E\|B_i-\bar B\|_{\rm HS}^2$, before the symmetry factor.

For a concrete check choose two off-diagonal curvature ports
separated by $d_i=m_0(\eta)+\delta_i$, with private Gaussian
increments of covariance $2u(\eta)I_4$ and $u>0$. In Maxwell
units their normalized off-diagonal entries are
$e^{\ii q\cdot d_i}$ and its conjugate, as obtained by translating
one curvature in a bilinear pairing; complex Fourier slots are
paired with their real-field conjugates. At $X=|q|^2$ the trace of
the product of the two port matrices averages to
$2\mathbb E e^{-2uX}$. The common displacement $m_0$ cancels.
Hence, apart from the angular volume and the symmetry factor, the
whole radial integral is
\[
 \int_0^\infty 2X\,\mathbb E e^{-2uX}\,\dd X
                          =\tfrac12\mathbb E[u^{-2}].
\]
More generally, suppose the two displacements are jointly Gaussian
at fixed $\eta$ with equal conditional means. Write $\Sigma_\Delta$
for the conditional covariance of their difference. The complete
integral of this product trace in the convention $\dd^4q$ is
$2(2\pi)^2\mathbb E[(\det\Sigma_\Delta)^{-1/2}]$, with
infinite cost for a singular covariance. This follows by integrating
$2\exp(-q^{\mathsf T}\Sigma_\Delta q/2)$ before averaging the
positive result. A null direction that changes orientation with
$\eta$ still has infinite conditional cost. If the same increment
crosses both copies the product trace is two at every momentum.
These absolute costs concern the graph after its prescribed Gaussian
integrations. At fixed displacements the product trace is
$2\cos(q\cdot(d_1-d_2))$, and taking its modulus there would
discard the phase cancellation. Thus even a Gaussian marginal for each
photon end leaves the answer undetermined until the gluing rule
specifies its relative increment. This example is a curvature-port
test; it supplies no independently measured charged-pair cut.

Distinct correlated records can also be tested. Suppose the photon
acts in a common transverse port and the two loop responses
$b_X(\xi)$ in that port are glued by the same self-adjoint
nonnegative Markov operator $T_X$ on $L^2(\pi)$ for the normalized
boundary measure $\pi$. With the Maxwell factors included in $b_X$,
let $\int b_Xd\pi=\Pi(X)$. The ring contraction in this port is
\[
 \langle b_X,T_X b_X\rangle
 =|\Pi(X)|^2+
 \langle b_X-\Pi(X),T_X(b_X-\Pi(X))\rangle
 \ \ge\ |\Pi(X)|^2,
\]
because $T_X1=1$. The test covers a decoupled physical port, and
also a matrix of loop entries when the ring sums these same
nonnegative forms over the entries. It applies to normalized
reversible heat transfers, their probability mixtures, and
reflection gluings $T_X=P_X^*P_X$ preserving constants. Dependence
of the propagation time on $X$ leaves this lower bound intact.
A nonnegative transition density alone does not imply operator
nonnegativity: interchanging two equiprobable records has eigenvalue
$-1$ on the orthogonal complement of constants. The condition is
therefore a test of the actual geometric transfer.

One can quantify how much the same overlap would have to decrease.
Replace the fixed port lower bound by
$D_\xi(X)\succeq z(X)|u_\xi\rangle\langle u_\xi|/X$ for positive
nonincreasing $z$. The radial ring integral is bounded below by
$\int Xz(X)^2|\Pi(X)|^2dX$. In the finite-difference proof, multiply
the squared estimate by $Xz((N+1)X)^2$. Monotonicity allows the
weight at each dilated argument $(j+1)X$ to be replaced by the
larger value $z((j+1)X)^2$. Absolute convergence of the ring
therefore requires
\[
 \int_{X_0}^{\infty} z(X)^2\,\frac{dX}{X}<\infty .
\]
For $z(X)\asymp[\log(X/X_*)]^{-\gamma}$ at large $X$ this means
$\gamma>1/2$. The rate is sharp for the isolated ring with one
positive mass atom and scalar sewing $D=z(X)/X$. Such a scalar
sewing with $z(X)\to0$ already fails the positive uniform massless
residue bound. A physical implementation would need to derive the
changing joint overlap together with the on-shell photon residue.

\paragraph{A normalized response to a spatially varying Maxwell probe.}
There is a related test at order $e^2$ with nonzero momentum
transfer. Vary the free Maxwell density by
$t h(x)F_{\mu\nu}F_{\mu\nu}/4$ with a real Schwartz probe $h$,
leaving the matter coupling fixed. In
$\log [Z_e[th]Z_0[0]/(Z_e[0]Z_0[th])]$ the coefficient of
$e^2t^2$ closes the complete matter Hessian on
$D V_h D V_h D$. This follows by expanding the inverse Maxwell
quadratic form $(D^{-1}+tV_h)^{-1}$ on transverse tests. The free
photon normalization cancels from the displayed ratio. With
$p'=p+q$, the matrix in $V_h$ between $p$ and $p'$, before
$\widehat h(q)$, is
\[
 M_{\mu\nu}=(p\cdot p')\delta_{\mu\nu}-p'_\mu p_\nu .
\]
Contraction with $p_\mu$ or $p'_\nu$ vanishes, so transverse
projectors leave it unchanged. In four dimensions
$\|M\|_{\rm HS}^2=p^2p'^2+2(p\cdot p')^2$.
The absolute density after the finite tensor trace, apart from
a fixed normalization, is consequently
\[
 |\widehat h(q)|^2\frac{|K_T(p^2)|}{p^2}
 \left[1+2\frac{(p\cdot p')^2}{p^2p'^2}\right].
\]
The bracket lies between one and three wherever the momenta are
nonzero. Proposition~\ref{prop:carrier-sewn-uv} therefore forces
absolute divergence for every nonzero such probe, including probes
whose Fourier transforms are supported away from zero. This is a
connected response with two prescribed Maxwell kinetic-density
insertions. It tests the same physical Hessian without changing any
matter or boundary parameter.

\paragraph{A light-cone test for positive transverse polarization channels.}
A wider test allows the positive spectrum, its polarization and
its invariant mass distribution to vary with spatial momentum.
On the transverse photon subspace let the complete polarization
at its first electromagnetic order have the matrix representation
\begin{equation}
 T_k(X)=\int_{s_0}^{\infty}\frac{dR_k(s)}{X+s},
 \qquad X=\omega^2+|k|^2,\quad s_0>0 .
 \label{eq:null-positive-port-uv}
\end{equation}
Here $dR_k$ is a measurable nonnegative Hermitian matrix measure
on the two transverse polarizations, with even scalar trace in
$k$. Both the matrices and their mass support may depend on
the direction and magnitude of $k$. Suppose the actual Maxwell
sewing at this order is proportional to
$\int d\omega\,d^3k\,\operatorname{tr}T_k(X)$, and that the
corresponding term in the transverse inverse is $e^2 X T_k(X)$.
Under spacelike commutativity of the measured magnetic field at
this order, every nonzero integrated channel gives a divergent
ring. No bound on the high-mass moments of $R_k$ is required.

Indeed, integrating the positive trace in $\omega$ shows that
finiteness of the ring would make
\[
 dM(k,s)=\frac{d^3k\,\operatorname{tr}dR_k(s)}{2E_s(k)},
 \qquad E_s(k)=\sqrt{|k|^2+s},
\]
a finite positive measure. Put
$W(t,x)=\int e^{-iE_s t+ik\cdot x}\,dM$ and
$J(t,x)=i[W(t,x)-W(-t,-x)]$. Both are bounded and continuous.
For each fixed time the spatial Fourier measure has an integrable
density, so $J(t,x)$ tends to zero at spatial infinity.

The order-$e^2$ correction to the photon covariance is
\[
 -\frac{T_k(X)}{X}
 =\int\frac{dR_k(s)}{s}
          \left(\frac{1}{X+s}-\frac1X\right).
\]
In the trace of the magnetic commutator the two curls supply
$|k|^2$ on each transverse polarization, since the squared
spatial curl is $|k|^2$ times the identity on $k^\perp$.
The even trace measure therefore weights the difference of the
massive and massless sine propagators. Applying the wave
operator kills the massless sine propagator and multiplies the
massive one by $-s$. Thus spacelike commutativity implies
$\Delta_x J=0$ on the spacelike exterior. These identities hold
as distributions: the auxiliary inverse powers of $s$ are
controlled by $s\ge s_0$, and spatial polynomials are tested on
Schwartz functions. Continuity of $J$ allows the spatial identity
to be used at each fixed time on $|x|>|t|$.

Let $\overline J(t,r)$ be the spherical average in $x$. A radial
harmonic function in three spatial dimensions has the form
$A(t)+B(t)/r$ on $r>|t|$. Decay at infinity gives $A(t)=0$.
Consequently $\partial_r[r\overline J(t,r)]=0$ throughout the
exterior. The angular average of a spatial plane wave is
$\sin(r|k|)/(r|k|)$. Therefore
\[
 \partial_r[r\overline J(t,r)]
     =2\int\sin(E_s t)\cos(r|k|)\,dM .
\]
The differentiation is legitimate even without a momentum moment,
and the right side is continuous in $r$ because $M$ is finite.
Setting $r=|t|$ now gives, for every real $t$,
\[
 \int\{\sin[t(E_s+|k|)]+\sin[t(E_s-|k|)]\}\,dM=0 .
\]
This is the sine transform of the sum of two finite positive
pushforward measures. Both are carried by $(0,\infty)$ since
$E_s>|k|$. Uniqueness of Fourier transforms would identify
their sum with its reflection about zero, forcing zero total
mass. This contradicts a nonzero channel.

This proof uses the matrix weights only through their positive
trace and the source's two explicit Maxwell contractions. It
also applies to a non-Gaussian charged action if those are its
actual leading amplitudes. Adding a finite local polynomial to
the traced ring integrand cannot restore absolute convergence:
at each fixed $k$ the Stieltjes term tends to zero as
$|\omega|\to\infty$, whereas a nonzero polynomial in $\omega$
already has a divergent absolute frequency integral. A different
microscopic gluing must establish its own relation between the
measured magnetic response and its complete bare graph.

\subsection{The directed action and its scalar coefficient}
\label{sec:path-action-type-uv}
On one aligned support $\bin=(s,b_s)$, put
$e_{\bin}=\hat p_{0,\bin}$ and $u_{\bin}=\sigma_\epsilon e_{\bin}$.
The real trace and supertrace in the printed curvature monomial
return a scalar coefficient.
The physical action on this support takes values on the carrier line.
If $a_{g,\bin}$ and $a_{m,\bin}$ denote scalar coefficients of the
geometric A4 and matter A3 reductions, respectively, their directed
exports and their inverse coefficient maps are
\[
\begin{aligned}
 \boldsymbol S_{g,\bin}&=u_{\bin}a_{g,\bin},&
 a_{g,\bin}&=u_{\bin}\boldsymbol S_{g,\bin},\\
 \boldsymbol S_{m,\bin}&=\sigma_s u_{\bin}a_{m,\bin},&
 a_{m,\bin}&=\sigma_s u_{\bin}\boldsymbol S_{m,\bin}.
\end{aligned}
\]
The inverse formulas use the one-generator geometric product
$e_{\bin}^2=1$. These are restrictions to the corresponding reductions.
In a metric-origin representative the neutral metric-copy fold stays
inside its geometric coefficient. The A3 sign is tested on the induced
matter action; inserting an independent matter action would change that
representative. Actions on distinct supports are compared only after
the prescribed carrier identification.

The restrictions also give a simple test for a proposed mixed term.
Take dimensionless coordinates $x,y$ in geometry and matter,
respectively. The scalar expression $x^2+y^2+xy$ restricts to $x^2$
and $y^2$ on the two axes. In the sector $\sigma_s=-1$, both
$u_{\bin}(x^2-y^2+xy)$ and $u_{\bin}(x^2-y^2-xy)$ pass those
two directed tests. Their mixed second derivatives differ by $2u_{\bin}$.
The mixed Hessian can be computed once the proposed curvature supplies its carrier-valued pullback or an equivalent assignment of its
coefficient before source differentiation.

For a locked A5 ordering branch use its carrier unit $u_{\bin}$
and total scalar coefficient $a$ with the prescribed source fold.
With the sources suppressed its phase is
$\mathrm{Exp}_{u_{\bin}}(i u_{\bin}a/\hbar)$. For real $a$ its two
idempotent components are $e^{ia/\hbar}$ and $e^{-ia/\hbar}$ with respect to
$(1\pm u_{\bin})/2$. If a carrier unit varies with the raw field, the variation includes
$a\,\delta u$ as well as $u\,\delta a$. The carrier unit enters once in the phase. Each character of the phase has modulus one. For fixed carrier data, derivatives
of $u_{\bin}a$ with respect to the fields have the same scalar
coefficients up to this fixed unit. The directed lift supplies neither
a new displacement in an internal history nor decay in its hard
momentum. The matter pole keeps its separate signed infinitesimal
$i\sigma_s\sigma_\epsilon\varepsilon e_{\bin}$, with the stated
opposite pole displacements and $\varepsilon\downarrow0$.

\subsection{Two ordered paths in the charged source}
\label{sec:path-tests-uv}
How does an intermediate displacement change the current that creates
a charged pair and the loop obtained by closing the same source?
On an Abelian link, going first along $a$ and then along $b$
has the same endpoints as going in the opposite order. Their
difference measures the flux between the paths and vanishes in
zero field. This lets us change the current while leaving the
free mass spectrum fixed. We compute that current. We then
put both paths inside the determinant and compare the weight of
a physical pair inferred in these two ways.

\paragraph{A spin-sensitive rectangular charged insertion.}
On a selected Euclidean four-plane choose a rank-four isometric
readout of the normalized A4 ball in $n=K-1\ge4$ dimensions.
Let $a,b$ be independent displacements from this projection
at length $L$, and let $f(u)=\chi_n(L\sqrt u)$, so
$\mathbb E e^{ip\cdot a}=f(p^2)$. The radial power series defines
$f$ at negative and complex arguments. This is a Euclidean
preparation of the charged comparison. Its analytic continuation
is distinct from a Lorentz-invariant probability law for real,
future-directed displacements.

Put $\Sigma^{\alpha\beta}=a^\alpha b^\beta-b^\alpha a^\beta$.
In the charged bilinear between $x$ and $x+a+b$, test the kernel
\[
 +\frac{i\lambda}{8mL^4}\sigma^E_{\alpha\beta}
 \,\mathbb E\!\left[\Sigma^{\alpha\beta}(U_{ab}-U_{ba})
                  \psi(x+a+b)\right].
\]
Here $\sigma^E_{\alpha\beta}=i[\gamma_\alpha,\gamma_\beta]/2$.
The expectation includes the translated field. The link
$U_\gamma=\exp(ie\int_\gamma A)$ brings the endpoint field back
to the base in the charge convention $\psi\mapsto e^{-ie\alpha}\psi$.
The sign of the real dimensionless insertion marker $\lambda$ fixes
the sign of the continued Pauli vertex; one may set $\lambda=1$
for this comparison. The area weight reverses under exchange of
$a,b$ together with the link difference. Under the change $(a,b)\mapsto(-b,-a)$, each reversed
path is the adjoint link, the area changes sign and $i$ is conjugated.
The displayed translation operator is therefore Hermitian for real
$A$, and commutes with $\gamma^5$. It vanishes at $A=0$.
The two boundary records and the spin contraction are specified
charged pairing data. They need a derivation to be selected
by the ambient curvature action.

Define $I(t)=\int_0^1 e^{iut}du$. At one photon the bare line
difference, omitting the common base point phase and the coupling,
is exactly
\[
 i I(Q\cdot a) I(Q\cdot b)\,Q_\rho\Sigma^{\rho\mu}.
\]
This follows by adding the two straight segment integrals in
each order and using $e^{it}-1=itI(t)$. In particular the Ward
contraction vanishes identically; common endpoints are sufficient.
The second derivative of a link is the product of its two
line integrals times $(ie)^2$. The contact is the difference of
the two such products, with the same area and endpoint field.
For incoming momentum $p$ at the translated endpoint put
\[
 M_{\alpha\beta}
 =-\int_0^1du\,\left.
       \partial_{\ell_\alpha}\partial_{\ell_\beta} f(\ell^2)
                        \right|_{\ell=p+uQ}.
\]
Indeed $M_{\alpha\beta}=\mathbb E[a_\alpha a_\beta
e^{ip\cdot a}I(Q\cdot a)]$. Independence of the two actual
records and antisymmetry of $\sigma^E_{\alpha\beta}$ give
$\sigma^E_{\alpha\beta}\mathbb E[
 e^{ip\cdot(a+b)}I(Q\cdot a)I(Q\cdot b)
 \Sigma^{\alpha\beta}\Sigma^{\rho\mu}]
 =4M^{\rho\alpha}M^{\mu\beta}\sigma^E_{\alpha\beta}$.
The resulting vertex is
\begin{equation}
 \delta\Gamma^\mu(p+Q,p)=
 -\frac{ie\lambda}{2mL^4}\,
 Q_\rho M^{\rho\alpha}M^{\mu\beta}\sigma^E_{\alpha\beta}.
 \label{eq:rect-pair-mark-uv}
\end{equation}
In particular the spin contribution survives the shell average.
At $p=0$, $Q\to0$, $M_{\alpha\beta}\to
L^2\delta_{\alpha\beta}/(n+2)$. The entire expression transforms
tensorially under complexified rotations of the four-plane, since
it is built from the metric, $p,Q$, and derivatives of an entire
scalar. On the continued mass shell this supplies a Lorentz
tensor for this vertex. A boost generator on the interacting
Hilbert space does not follow from this identity.

\paragraph{Timelike pair and spacelike elastic tests of the same mark.}
Assume the open charged source has the ordinary normalized Dirac
pole and the induced insertion at that pole is
\eqref{eq:rect-pair-mark-uv}. This holds, for example, when the
displayed term is added to the minimally gauged Dirac inverse.
The same test applies to the atom of a polar source whose atom
contributes that inverse. On the continued pair shell let
$p^2=(p+Q)^2=-m^2$, $Q^2=-S$ and $E=\sqrt S/2$.
Choose the pair rest frame, $Q_4=2iE$, $p_4=-iE$,
$p_3=\sqrt{E^2-m^2}=:P$, and $Q_{1:3}=p_{1:2}=0$.
For $s_u=-m^2+S u(1-u)$ the matrix $M$ is diagonal with
\begin{equation}
 \begin{aligned}
 M_{44}&=-2f'(-m^2),\\
 M_{11}=M_{22}&=-2\int_0^1 f'(s_u)\,du=:M_\perp,\\
 M_{33}&=M_\perp-4P^2\int_0^1 f''(s_u)\,du .
 \end{aligned}
 \label{eq:rect-pair-rest-uv}
\end{equation}
The first equality follows by integrating
$Q_\rho\partial^\rho\partial_\alpha f$ in $u$; oddness about
$u=1/2$ removes the mixed entries. At threshold $M_{33}=M_\perp>0$
and $M_{44}>0$: $f(-w)$ is the positive hyperbolic-cosine
moment of the longitudinal ball coordinate and $f'(-w)<0$.
The Dirac pair has a nonzero triplet matrix element of
$\sigma^{0i}$ at threshold. At rest the temporal current vanishes,
as conservation requires. A boost along a nonzero spatial
matrix element then gives a nonzero temporal matrix element of
this continued vertex. At zero spatial transfer its extra
density is identically zero by the line difference. Thus the
construction realizes a nonzero vacuum-to-pair density matrix element
while preserving total charge at the one-current level.

At large $S$ one instead finds
\[
 M_\perp\sim \frac{4f(-m^2)}{S},\qquad
 M_{33}\longrightarrow 2f'(-m^2).
\]
For $j=1,2$, substitution of $v=s_u$ writes
$S\int_0^1 f^{(j)}(s_u)du$ as twice the integral of
$f^{(j)}(v)[1-4(v+m^2)/S]^{-1/2}$ on $[-m^2,P^2]$.
The radial Bessel estimate on the positive axis gives
$f^{(j)}(v)=O(v^{-(n+1)/4-j/2})$. The upper half of that
interval contributes $o(1)$ even with its integrable endpoint
singularity; dominated convergence applies on the lower half.
The two limits of the rescaled integrals are $-2f(-m^2)$
and $-2f'(-m^2)$.

With the conventional signs for the Pauli moment let
$\kappa_T=\lambda M_{44}M_\perp/L^4$ and
$\kappa_L=\lambda M_{44}M_{33}/L^4$. In a static elastic
Breit limit the transverse magnetic insertion has
$\kappa_0=4\lambda f'(-m^2)^2/L^4$.
Consequently $\kappa_T=O(S^{-1})$ and $\kappa_L\to-\kappa_0$.
For a unit minimal vertex plus this rectangular insertion the
two transverse triplet amplitudes and single longitudinal triplet
give, with $\tau=S/(4m^2)$, $\beta=\sqrt{1-1/\tau}$,
\begin{equation}
 12\pi^2\rho_{\rm pole}(S)=
 \beta\left[|1+\kappa_T|^2+
           \frac{|1+\tau\kappa_L|^2}{2\tau}\right].
 \label{eq:rect-pair-optical-uv}
\end{equation}
The normalization is that of the unit Dirac density after
removing any common photon readout and $e^2$.
For $\kappa_0\ne0$ the pole-pair requirement grows as
$\kappa_0^2S/(96\pi^2m^2)$. This is a lower bound on a positive
inclusive cut for that same current, not a calculation of the
closed determinant. In particular the earlier paired-shell
density, already bounded above by the Dirac value, cannot be
kept fixed while replacing the open vertex by this one.

A further path test uses a third independent shell displacement
$c$. Prefix both routes by $c/2$ and append the remaining $c/2$
to both. At one mark the common segments cancel in the link
difference; translating the insertion by $c/2$ and the endpoint
by $c$ multiplies \eqref{eq:rect-pair-mark-uv} by exactly
$f((p+Q/2)^2)$. On the pair shell its argument is $S/4-m^2$,
so the rectangular contribution to the pair weight now decays
for $n\ge4$. The static moment retains a nonzero multiplier
$f(-m^2)$. The minimally coupled unit Dirac contribution still
survives at high pair energy.

There is an independent crossed-channel cost. For spacelike
elastic transfer $Q^2=X>0$ the same midpoint argument is
$-m^2-X/4$. In the Breit frame the transverse magnetic
coefficient of the unsplit rectangle contains $-2f'(-m^2)$
times $-2\int_0^1f'(-m^2-Xu(1-u))du$. Its integrand is positive.
Writing $f(-w)=\mathbb E\cosh(\sqrt w\,a_1)$ gives
$-f'(-w)=\mathbb E[a_1\sinh(\sqrt w\,a_1)]/(2\sqrt w)$.
The nonzero probability of $a_1>L/2$ bounds this below by an
exponential divided by $\sqrt w$ at large $w$. Restricting to
$1/4<u<3/4$ establishes exponential growth in $\sqrt X$ for the
original elastic magnetic form factor. The extra midpoint shell
enhances that growth. Neither preparation gives an elastic
magnetic form factor polynomially bounded on that spacelike ray
for a nonzero rectangular coupling. The time orientation and
the vectorial pole prescription choose the physical boundary
value; they do not change these entire multipliers.

\paragraph{Two ordered paths inside the paired normal determinant.}
A direct test keeps the full displacements in the charged operator.
Use the four-dimensional Euclidean readout of the selected sector-plane,
$n=K-1\ge4$, $L=\sqrt{G_N}$ in carrier units, and $m=1/L$.
Let $d\mu(a)$ be the projected law of $a=LT^*y$ for one normalized
A4 ball record. The full record retains its original oriented shell
sign. The rank-four isometry $T$ and two independent records are
choices of this comparison. The projected law is even and has a
bounded density supported in $|a|\le L$. Set
\[
 (V_a[B]\psi)(x)=
 e^{i\int_0^1 a\cdot B(x+ta)\,dt}\psi(x+a).
\]
Put $T_\mu[B]=\int a_\mu V_a[B]d\mu(a)$. Absorb the charge marker
in $B$ when writing covariant translations. Use the connection of
\eqref{eq:upgrade-shell-pair-dirac-uv}. With repeated indices summed,
use the preceding unsplit rectangle with $\lambda=1$. Its operator form is
\begin{equation}
\begin{aligned}
 R_B&=\frac{i}{4L^3}\sigma^E_{\mu\nu}
                    [T_\mu[B],T_\nu[B]],\\
 K_B&=m+d_B+R_B,\qquad H_B=K_B^\dagger K_B,\\
 \mathcal D_B&=L^{-1}K_B H_B^{-1/2}e^{-\Phi(H_B)} .
\end{aligned}
\label{eq:rect-polar-uv}
\end{equation}
Here $\sigma^E_{\mu\nu}=i[\gamma_\mu,\gamma_\nu]/2$ and $\Phi$ is
the two-record normal Green function. The two orderings of the
Wilson translations carry the same endpoint phases. The commutator
weights their difference by the oriented spin area
$\sigma^E_{\mu\nu}a_\mu b_\nu$. For real $B$,
$V_a^\dagger=V_{-a}$ and $T_\mu^\dagger=-T_\mu$.
Thus $R_B$ is Hermitian and commutes with $\gamma^5$.
At zero field the $T_\mu$ are commuting Fourier multipliers.
Consequently $R_0=0$ and the free charged mass atom and its
positive spectral measure are the same as in the paired model.
These properties also give gauge covariance and
$K_B^\dagger=\gamma^5 K_B\gamma^5$. The perturbative phase is constant in the gapped
neighborhood of zero background by this same adjoint symmetry.
The normalization in
\eqref{eq:rect-polar-uv} specifies a charged representation;
the curvature monomial alone has not fixed it.

\paragraph{Ultraviolet test before the short-path expansion.}
Every source-dependent Euclidean graph built from the complete
charged vertices of \eqref{eq:rect-polar-uv} is ultraviolet
finite at each finite electromagnetic order in the rooted
representation used above. Its transverse photon is perturbatively
massless. In particular the nonzero local Pauli heat divergence
in \eqref{eq:next-pauli-heat-uv} does not persist for this finite
bilocal insertion.

To prove the improvement, differentiate the actual translations.
A mark on one link inserts a component of $a$ and an integral
over its position in $[0,1]$. For fixed positions and transfers
each marked $T_\mu$ thus has, up to its transverse photon factors,
symbol $\widehat{a^\alpha\mu}(k+P)$ at a translated charged
momentum. This function is in $L^2\cap L^\infty(\mathbb R^4)$,
as are all its momentum derivatives, by compact support, bounded
density and Plancherel. Each marked $R$ contains the product of
two such factors at possibly different translations. Cauchy--Schwarz
gives an $L^1(d^4k)$ bound uniform in both translations, as well
as a uniform supremum bound. This applies to every distribution of
marks on the two links, including direct contacts at one vertex.

For the relative, electromagnetically marked trace use
\[
 \Phi(H)=\tfrac12\int_0^\infty b(v)(H+v)^{-1}dv ,
 \qquad K_{B,+}=K_B+i\sqrt v\,\gamma^5 .
\]
Writing $D_5=\gamma^5 K_B=D_5^\dagger$ gives
$K_{B,+}=\gamma^5(D_5+i\sqrt v)$ and $D_5^2=H_B$. Thus, as
an operator identity before tracing,
\[
 \operatorname{Re}\!\left(K_{B,+}^{-1}
          \frac{i\gamma^5}{2\sqrt v}\right)
 =\tfrac12(H_B+v)^{-1}.
\]
At zero field the free inverse at $k$ has norm
$(m^2+v+|k|^2)^{-1/2}$.
Subtract the minimally gauged trace solely to estimate their
difference, retaining both pieces in the amplitude. Every word
in the difference has at least one marked $R$. Use its $L^1$
bound for each closed circulation and use supremum bounds for
the other factors. The mass integration is harmless since
$\int b(v)dv/\sqrt v<\infty$. This controls the charged variables
uniformly in all photon transfers.

The diagonal kernel of the traced word sum is gauge invariant at its root.
In its multilinear variations replace each plane wave by its
difference from its value at that root, exactly as in
\eqref{eq:upgrade-paired-root-bound-uv}. Express each difference
as a transfer derivative integrated along the incident ray.
Differentiating a translated link transform preserves its bounds;
differentiating a free inverse inserts another inverse and a
bounded Clifford matrix. The rooted word still contains both
link transforms of at least one $R$. Each transverse incidence
therefore has a bound with $|q|\,|c_L(q)|$. The previously proved
minimal bound, Maxwell spanning-tree integration and dominated
soft limit now apply to the full trace. On open lines use
\[
 (H_B+v)^{-1}K_B^\dagger
 =\tfrac12\sum_{\eta=\pm1}(K_B+\eta i\sqrt v\,\gamma^5)^{-1}.
\]
Every word with at least one marked vertex has at least two
massive inverses, even for a many-photon contact. The positive
free Stieltjes measure integrates $(m^2+v)^{-1}$. This proves
the open-chain hard estimate with the same external infrared
convention. No positive lower bound on a shell displacement was
used. Expanding $V_a[B]$ in charged derivatives before the loop
integral would lose both integrable link transforms.

\paragraph{A threshold test of the very same pair of paths.}
Let $U=S/4-m^2>0$ and let $\rho_v(S)$ be the parity-even
two-current discontinuity of the first-order reference
$K_B+i\sqrt v\,\gamma^5$, with the common external shell factors
removed. Its free charged mass is $M=\sqrt{m^2+v}$. Near the
lowest threshold its nonanalytic term comes from the ordinary
two free propagator poles, with both on-shell vertices evaluated
at those poles. In the terms containing $R$, every link transform
is entire in the incident momenta. In a bounded imaginary-energy
strip its measure acquires only a bounded exponential tilt, with
the same $L^2$ estimates. Route the bubble symmetrically by $k\pm Q/2$.
Off a small ball about $k=0$ the two energy denominators stay
uniformly away from zero on the real contour near the first
threshold, including its boundary values reached from below.
The tilted transforms are holomorphic as $L^2$-valued functions
and their product is locally bounded in $L^1$. The hard part
therefore continues analytically. At a single-propagator contact
all the transfer can be routed through the entire link factors,
so it has no threshold pinch. The two poles inside the ball
give the ordinary two-body phase space by their residues. The pure minimal term has the usual Dirac
two-particle cut. This argument concerns the local discontinuity
on the prescribed matter boundary; it requires no contour
closure through a direction of exponential growth at infinity.

The threshold matrix element is nonzero. At $v=0$ set
$Q=2im e_4$, with incoming and outgoing pole momenta $\mp im e_4$
and spatial polarization $j$. Write
\[
 A=\int a_4\sinh(m a_4)d\mu,\quad
 C=\int a_j^2\int_0^1 e^{(1-2t)ma_4}dt\,d\mu .
\]
The change in the free $T_4$ multiplier is $-2A$ and
$\delta T_j=iC$; the crossed components vanish by reflection.
Both antisymmetric terms in the commutator therefore give the
forward correction $h\sigma^E_{4j}$ for $h=AC/L^3$ and the
reverse correction $-h\sigma^E_{4j}$. The rest projectors are
$P_\pm=(1\pm\gamma_4)/2$. Direct multiplication gives
\[
 P_\pm(i\gamma_j\pm h\sigma^E_{4j})P_\mp
       =(1+h)P_\pm i\gamma_j P_\mp .
\]
Here $h>0$ and $h\le\sinh(1)\cosh(1)/(n+2)^2$ by the ball
second moment and $mL=1$. For the twisted free mass $m+i\sqrt v\,
\gamma^5$ the same unitary matrix
$e^{-i\vartheta\gamma^5/2}$ acts on both sides of the kinetic
operator, with $\tan\vartheta=\sqrt v/m$. It leaves the vector
vertex invariant and turns the spin vertex into its product
with $e^{-i\vartheta\gamma^5}$. These matrices and their
two on-shell link factors are continuous uniformly for all
spatial directions as $U\downarrow0$ and $0\le v\le U$.
Consequently the angular spin trace times two-body phase space is
\[
 \rho_v(S)=[c_*+o(1)]\sqrt{U-v},\qquad
 c_*=\frac{(1+h)^2}{8\pi^2m}>0 ,
\]
uniformly on this interval in the unit Dirac normalization.

The open inverse is $L F(H_B)K_B^\dagger$. Its Stieltjes atom
therefore contributes $Z_\psi K_B^{-1}$. Marking it once gives
$-Z_\psi K_0^{-1}(\delta K)K_0^{-1}$. In the continuum term
the two external residue factors select $v=0$ through
$x/(x+v)$ and $y/(y+v)$ as the two squared-mass distances
$x,y\downarrow0$. Near zero the numerators are bounded, and
away from zero the open-chain mass bound controls the tail.
No continuous mass contributes a double atom. The external
spinor normalizations remove the common $Z_\psi$. Thus
$\rho_0$ is exactly this source's normalized pole-pair
density, including the full one-photon rectangle at both ends.
The analytic external photon readout multiplies the open and
closed tests by the same nonzero square near threshold.

Applying the mass derivative to the complete reference
two-current response and integrating by parts gives
\begin{equation}
 \rho_{\rm cl}(S)=\int_0^U b(v)(-\partial_v\rho_v(S))dv
 =\rho_0(S)+\int_0^U b'(v)\rho_v(S)dv .
 \label{eq:rect-threshold-uv}
\end{equation}
One can also integrate by parts in the analytic polarization
before taking its discontinuity. This avoids differentiating an
individual threshold singularity. The shell second moment gives
$b(0)=1$ and $b'(0)=-L^2/(n+2)$. It follows that
\[
 \frac{\rho_{\rm cl}(S)}{\rho_0(S)}
 =1-\frac{2L^2}{3(n+2)}U+o(L^2U),\qquad U\downarrow0 .
\]
Thus the full paths cure the marked Pauli ultraviolet divergence
and still give a strict optical deficit on an open interval
immediately above threshold. Positivity with the usual physical
charged-pair interpretation requires at least $\rho_0$ on that
interval. The test uses the atom of the specified free endpoint;
it makes no assumption of an isolated interacting charged pole.

More generally, for this mass-derivative trace and a continuous nonzero
rest-pair amplitude, the first relative correction is
$2b'(0)U/3+o(U)$ when $b(0)=1$. For the difference of independent
copies of one bounded normal coordinate $Y$, its weight satisfies
$b(v)=|\mathbb E e^{i\sqrt vY}|^2$ and $b'(0)=-\operatorname{Var}Y$.
Every such nondegenerate coordinate gives a deficit of the same sign.

\paragraph{A geometric normal boundary state with an exact optical interval.}
The strict deficit arbitrarily close to the light threshold can be
removed by changing the charged normal boundary state. Its behavior
at higher energy can be tested in the same determinant. Keep the
paired normal Green operator, the minimally gauged polar charged
inverse, and the four-dimensional transverse shell readout above.
For this comparison embed the dual two-dimensional normal plane
isometrically in a Euclidean four-plane. Average the massless
four-dimensional Green kernel over its sphere of radius $R=m$,
using normalized surface measure $d\omega$. Put $a=R^2$,
$v=|\xi|^2$ and define the boundary wavefunction by its transform
\begin{equation}
 \begin{aligned}
 \widehat\psi_R(\xi)&=\int_{S^3}\frac{R^2}{|\xi-R\omega|^2}\,d\omega
 =\begin{cases}1,&0\le v\le a,\\ a/v,&v>a,\end{cases}\\
 b_R(v)&=|\widehat\psi_R(\xi)|^2 .
 \end{aligned}
 \label{eq:fix-threshold-sphere-uv}
\end{equation}
Rotation invariance and harmonicity give a bounded constant
radial potential inside the sphere and a multiple of $1/v$
outside. Its value at the center and its asymptotic at infinity
fix both constants. The surface singularity is integrable and
the two expressions agree there. The spectrum extends to
arbitrarily high momentum with nonzero weight. The selected
sphere in dual normal momentum and the coherent boundary state
are extra charged representation data. The oriented A4 shell,
its velocity delta and its sector carrier retain their original
roles in the geometric action.

This is an actual square-integrable boundary state in the
two-dimensional normal position space, up to the normalization
of the two boundary insertions. Indeed
$\int_0^\infty b_R(v)dv=2a$, so planar Plancherel applies.
It also has an integrable real, signed amplitude. To see the
tail, the planar distributional Laplacian of $\widehat\psi_R$
is $4R^2|\xi|^{-4}$ outside the disk plus the surface measure
on its boundary with coefficient $-2/R$. The angular Fourier
integral of either term is bounded by $C_R |y|^{-1/2}$ for
$|y|>1/R$, using the bound for $J_0$. Division by $|y|^2$
gives $|\psi_R(y)|\le C_R|y|^{-5/2}$. Local integrability follows
from square integrability. Thus $\int\psi_R=1$ in the stated
Fourier convention. A nonnegative probability preparation with
unit characteristic function on a neighborhood of the origin
would be a point mass. For example, equality in the triangle
inequality makes the Fourier phase constant almost surely for
a countable dense set of those momenta. The coherent state in
\eqref{eq:fix-threshold-sphere-uv} has signed, unbounded spatial
support.

The matrix element between these two normal boundary states gives
$W_R(t)=t\int_0^\infty e^{-tv}b_R(v)dv$ and the following exact
logarithm for real $z>0$:
\begin{equation}
 2\Phi_R(z)=\log(1+a/z)+\frac{a}{z}
               -\frac{a^2}{z^2}\log(1+z/a).
 \label{eq:fix-threshold-log-uv}
\end{equation}
This follows by splitting the resolvent integral at $a$.
Since $0\le b_R\le1$ and $W_R(t)\sim2at$ at short time,
the positive free Stieltjes construction and the joint Euclidean
graph proof \eqref{eq:upgrade-paired-root-bound-uv} apply at
every finite electromagnetic order to the normalized source.
This includes the contact vertices, the tested open lines and the
soft transverse Maxwell limit in that proof. Here
$F_R(z)=e^{\Phi_R(z)}/\sqrt z$ satisfies
$\lim_{z\downarrow0} zF_R(z)=R e^{1/4}$.
Equation~\eqref{eq:fix-threshold-log-uv} gives a meromorphic
continuation of $F_R$ across $-a<z<0$, with the pole at zero
as the only spectral contribution below $a$. The normalized
open residue therefore gives the ordinary light-Dirac matrix
elements of the same gauged source.

Let $\mathcal C_v(S)$ denote either the two-current Dirac
spectral density or the equal-helicity forward two-photon
cut, with charged mass squared $m^2+v$. Strip common external
readout factors and keep the same physical pole boundary.
Both functions vanish for $v\ge U=S/4-m^2$, are positive
for $0\le v<U$, and vanish continuously at the upper endpoint.
For these one-loop cuts the mass-derivative representation gives
\begin{equation}
 \mathcal C_{\rm cl}(S)=\mathcal C_0(S)
 -{\bf1}_{U>a}\int_a^U \frac{2a^2}{v^3}\mathcal C_v(S)\,dv ,
 \qquad U>0 .
 \label{eq:fix-threshold-cuts-uv}
\end{equation}
Integration by parts has no interior delta term because $b_R$ is
continuous at $a$. The complete closed cuts therefore equal the
weight of their normalized open pole pair throughout
$4m^2<S\le8m^2$ for $R=m$. They are strictly below it at every
higher $S$. At large $S$ the equal-helicity cut is negative:
the argument of \eqref{eq:ward-bw-sign-uv} applies with
$\int b_R=2a$ and a $v^{-2}$ tail. These conclusions concern
the two specified channels at their stated perturbative order.
They isolate the range on which this revised representation
repairs the previously immediate threshold mismatch.

\paragraph{Which added states could supply the missing weight?}
The harmonic boundary permits a stronger check than the two specific
cuts above. Define a positive probability measure
$d\lambda(v)=2a^2{\bf1}_{v>a}\,dv/v^3$. For $z>0$ its scalar
resolvent identity is
\begin{equation}
 2\Phi_R(z)=\int_a^\infty
       \log\!\left(\frac{z+v}{z}\right)d\lambda(v).
 \label{eq:fix-relative-spectrum-uv}
\end{equation}
Indeed integration by parts gives
$\int_0^\infty b_R(v)/(z+v)\,dv$ on the right. Functional
calculus gives the same identity for the massive squared Dirac
operator. Thus its parity-even closed loop is a light Dirac log
determinant minus the probability average of the heavier ones.
The spectral identity combines the logarithms \emph{before} taking
an ultraviolet trace. It does not give a finite value to either
isolated bare charge logarithm.

At any finite number of photon insertions in this one charged loop,
fix a physical channel and take its normal two-charged-line cut.
For real incident photons of total invariant $4m^2<S\le8m^2$,
the intermediate states of every heavier term in this cut are
kinematically excluded, since $2\sqrt{m^2+v}>\sqrt S$.
Consequently the cut is exactly the light Dirac cut, with the
common external readouts. The open light atom produces the same
Dirac amplitude: the remainder of its Stieltjes measure starts at
$v=a$ and has no external double mass atom. This statement concerns
the specified normal cut at nonsingular external kinematics, without
making an assertion about anomalous thresholds, other channels,
or graphs containing internal photons.

One might try adding independent positive species with mass law
$d\lambda$ to replace the negative closed contribution. They would
also create the corresponding additional open physical pairs. At
two currents, if a spectator sector has closed weight $c$ equal to
the sum of its positive squared production amplitudes, tensoring it
with the original vacuum adds $c$ to both sides of the optical test.
The original deficit is unchanged. In the explicit harmonic example
the heavy spectral average cancels algebraically in the combined
closed loop, leaving the light Dirac ultraviolet logarithm, while
the new positive open pairs remain. A repair must change the
relation between the measured charged states and the original loop
itself.

\paragraph{What compact coherent shell histories can achieve.}
If the normal boundary amplitude is a finite signed measure of
compact support, its Fourier transform on every complex momentum
line is entire. The autocorrelation on a real line extends as
$A(z)\overline{A(\bar z)}$. Agreement with unity on any interval
forces that function to equal unity everywhere on the line.
For radial data this gives $b\equiv1$, $W\equiv1$ and the
ordinary charge logarithm. More generally an exact two-current
optical equality for the minimal source on $0<U<U_0$, with
$0\le b\le1$, already forces $b=1$ almost everywhere there:
the negative mass derivative of the massive Dirac density is
strictly positive for $0<v<U$. This argument also covers an
autocorrelation formed from signed compact histories. For a
positive mixture $b(v)=\int e^{-sv}d\eta(s)$ with $\eta$
a probability law, equality at even one $v>0$ forces
$\eta=\delta_0$. Positive mixtures of proper-time exponentials
cannot make a finite interval of exact equality with a decaying
tail.

There is a compact realization of arbitrarily high \emph{finite}
contact with the optical threshold. Let $\mu$ be the original
two-dimensional projected record law, $\widetilde\mu(E)=\mu(-E)$,
$K=\mu*\widetilde\mu$ and $g(v)=|\chi_n(L\sqrt v)|^2$.
At a positive integer $p$ superpose finitely many convolutions
of these records with the following normalized signed amplitude:
\[
 \begin{gathered}
 \mu_p=\sum_{j=1}^p(-1)^{j+1}{p\choose j}K^{*j},\\
 A_p(v)=1-(1-g(v))^p,\qquad b_p(v)=A_p(v)^2 .
 \end{gathered}
\]
The shell variables are fully integrated on each history.
The integer counts the number of available convolutions.
Here $0\le b_p\le1$, $b_p\le p^2g^2$ and the same marked
Euclidean bounds hold for each finite $p$. Put
$c=L^2/(n+2)$. The shell second moment gives
$b_p(v)=1-2(cv)^p+O(v^{p+1})$. Since the Dirac
reference density is uniformly a common positive constant
times $\sqrt{U-v}[1+O(U/m^2)]$, its two-current ratio obeys
\[
 \frac{\rho_{\rm cl}}{\rho_0}
 =1-2p\,\mathrm B(p,3/2)(cU)^p+o((cU)^p).
\]
Here $\mathrm B$ is Euler's beta function. At fixed $p$ the
high-energy equal-helicity sign test still applies.
As $p\to\infty$, $b_p\uparrow1$ except at the discrete shell
zeros. Hence $W_p(t)\uparrow1$ at every positive time, and
monotone convergence recovers the divergent short-time Dirac
charge integral. The exact-interval construction has the same
limit if $R\to\infty$ at fixed charged mass. Its charge integral
is $2\Phi_R(m^2)=\log(a/m^2)+1/2+o(1)$ in that limit.

\paragraph{Positive first-order species and the price of exact sewing.}
A different charged representation uses physical first-order
Dirac species of mass $M$ and charge $e(M)$ in orthogonal
sectors with positive multiplicity $d\nu(M)$. At finite
multiplicity the determinant is the product of their ordinary
determinants. Each normal label propagates through a cut, so
the two-current and four-current inclusive weights equal sums
of the corresponding squared pair amplitudes. A positive
direct integral with a specified normal trace has the same
comparison by integration whenever the coefficients exist.
For minimal coupling the common-field coefficient of $F^2$
at short time has a positive multiple of
\[
 \int_0^\epsilon\frac{dt}{t}
         \int |e(M)|^2 e^{-tM^2}d\nu(M).
\]
If a bounded mass interval carries nonzero charge weight, its
contribution alone bounds the inner integral away from zero.
Every nonzero positive charged measure on finite masses has
such an interval. A nonzero external shell readout preserves
this divergence. Negative multiplicities in a formal
superdeterminant may cancel its coefficient, but a negative
resolved normal species cannot be counted as a positive
intermediate particle. This conclusion concerns positive
minimally coupled Dirac species with the stated trace. Other
spins, gauge constraints or genuinely different transverse
charged dynamics require their own sewing and ultraviolet
calculations.

\subsection{A shell measurement in the charged rest frame}
\label{sec:path-rest-tests-uv}
The rectangular spin insertion changes the Pauli current. Its failure at
the pair threshold suggests testing the electric form factor itself.
First hold the closed spectrum fixed and ask which open amplitudes
fit beneath it. Then insert the proposed amplitude in two explicit
off-shell constructions and recalculate their consequences.

\paragraph{A fixed plane and a boosted preparation.}
The projected Euclidean probability used in the determinant and a
probability on real Lorentzian histories answer different questions.
For example, there is no boost-invariant probability carried by nonzero
future timelike displacements on a Minkowski plane. For such a displacement
$a$, the real number $a^0+a^1$ is strictly positive. A boost of rapidity
$\log 2$ along the first axis can double it. Invariance would give all
intervals $[2^jL,2^{j+1}L)$ the same mass. Their union is $(0,\infty)$,
so no probability on that set is possible. A law conditional on a
transforming external four-velocity avoids this argument. The charge's
four-velocity and the carrier's unit geo-momentum have distinct roles.

\paragraph{A covariant A4 readout in the charged rest space.}
There is a different way to use the velocity data at a physical
charged pole. Use Minkowski signature $(+---)$ for this calculation.
For an on-shell charged leg of mass $m>0$ and momentum $p$,
take $u=p/m$. Project three coordinates of the normalized
A4 ball in $n=K-1$ dimensions isometrically into the spacelike
rest space $u^\perp$, endowed with metric $-g|_{u^\perp}$.
The induced law is independent of the spatial triad and the family
of laws transforms covariantly with $u$. This construction uses
the physical charged velocity as additional boundary information.
The ambient unit vector $\hat p_0$ alone does not identify that
rest space. The one-mark characteristic function is
\[
 \chi_n\!\left(L\sqrt{(Q\cdot u)^2-Q^2}\right).
\]
For pair production $Q=p+p'$ and elastic scattering $Q=p'-p$,
with $p^2=p'^2=m^2$, both signs give $(Q\cdot p)^2=(Q^2)^2/4$.
Either charged leg therefore yields the same Dirac and Pauli form factors
for the trial vector current:
\begin{equation}\label{eq:charged-rest-readout-uv}
  F_1(\tau)=H(\tau)=
  \chi_n\!\left(2Lm\sqrt{\tau(\tau-1)}\right).
\end{equation}
Here $F_2(\tau)=0$ and $\tau=Q^2/(4m^2)$. Contraction with $Q$ vanishes between the physical spinors by
their free Dirac equations. Both static Sachs form factors are
one, $H(1)=1$, and $|H|\le1$ on both real physical rays.
For fixed $Lm>0$, the Bessel asymptotic gives
$H=O(|\tau|^{-(n+1)/2})$ at either real end. These are on-shell
statements.

\paragraph{An all-energy two-particle test at the native scale.}
For $n=248$ and $Lm=1$, let $b(v)=|\chi_n(\sqrt v)|^2$. Denote
by $r(\tau)$ the stationary paired-shell density divided by the
unit point-Dirac density. In the cut variable $v$, the ratio is
the average of $b(v)$ on $0\le v\le\tau-1$ with weight
$w_\tau(v)=(1+v)/\sqrt{1-(1+v)/\tau}$. For $\tau>1$ the
charged-pair optical requirement is $|H(\tau)|^2\le r(\tau)$.
It holds on the entire physical ray.

For $1<\tau\le16$, put $k=2\sqrt{\tau(\tau-1)}$. Then
$\sqrt v< k\le\sqrt{960}<31$. The functions $\chi_n$ and
$\chi_{n+2}$ are positive on $[0,31]$. Indeed, their Bessel
orders are $124,125$, and $h(x)=\sqrt x J_\nu(x)$ satisfies
$h''=((\nu^2-1/4)/x^2-1)h$. Starting positive and increasing
near zero, $h$ stays convex and increasing up to $31$.
Since $\chi'_n(x)=-x\chi_{n+2}(x)/(n+2)$, we obtain
$|H(\tau)|^2<b(v)$ throughout the averaging interval.
For $\tau\ge16$, the coordinate variance
$1/(n+2)$ and $\cos x\ge1-x^2/2$ give
$b(v)\ge(1-15/500)^2$ on $[0,15]$. The full denominator satisfies
$\int_0^{\tau-1}w_\tau(v)dv\le4\tau^2/3$, while
$w_\tau(v)\ge1+v$. Thus
\[
  r(\tau)\ge
  \frac{3(1-15/500)^2}{4\tau^2}
  \left(15+\frac{15^2}{2}\right)>\frac{89}{\tau^2}.
\]
The normalized one-coordinate density is
$g_n(x)=g_n(0)(1-x^2)_+^{(n-1)/2}$, with
$g_n(0)=\Gamma(1+n/2)/(\sqrt\pi\,\Gamma((n+1)/2))$.
Its total variation is $2g_n(0)$, hence
$|\chi_n(k)|\le2g_n(0)/k$ by integration by parts.
Log-convexity of $\Gamma$ gives $g_n(0)^2\le(n+1)/(2\pi)$.
Consequently
\[
 |H(\tau)|^2\le
 \frac{249}{2\pi\tau(\tau-1)}<\frac{43}{\tau^2}
 \qquad(\tau\ge16).
\]
At threshold both ratio and squared form factor approach one.
This leaves positive spectral weight available above threshold,
without identifying the states or interactions that would supply it.

\paragraph{The price of the rest-space continuation.}
The even power series for $\chi_n$ makes $H$ entire in $\tau$.
On $\tau=1/2+iy$, however, its argument is imaginary and at the
native scale
\[
 \begin{aligned}
 H(1/2+iy)&=\mathbb E\cosh\!\left(2Y_1\sqrt{1/4+y^2}\right)\\
 &\ge {\mathbb P}(|Y_1|>1/2)\cosh\!\left(\sqrt{1/4+y^2}\right).
 \end{aligned}
\]
Here $Y_1$ is one coordinate of the unit A4 ball and the
probability is strictly positive. The trial current violates
polynomial boundedness on the complex physical sheet; bounds
assuming a bounded analytic form factor do not apply to it.
The real rest-space probability law requires a timelike charged
velocity. Extending the expression $u=p/\sqrt{p^2}$ to arbitrary
internal momenta fails at $p^2=0$ and loses that real probability
interpretation for spacelike $p$. The all-energy inequality
therefore proves a kinematic compatibility result. An open
charged operator, its common closed determinant, and its Ward
contact terms still have to produce the same spectral data.

For example, at $\tau=2$ a direct quadrature gives $|H|^2=0.968505$
and $r=0.997125$. At $\tau=16$ these numbers are $0.020861$ and
$0.953911$. The analytic bounds cover the intervals between these tests.

\paragraph{An off-shell charged vertex with the same physical values.}
One can extend the rest-space mark without dividing by the virtual
charged momentum squared. In Euclidean variables let $Q=k'-k$,
$P=(k'+k)/2$, $X=Q^2$ and put
\begin{equation}
\begin{aligned}
 b^\mu&=(1+X/(4m^2))Q^\mu+(P\cdot Q)P^\mu/m^2,\\
 Z&=Q\cdot b=X+X^2/(4m^2)+(P\cdot Q)^2/m^2,\\
 d(Z)&=[f(Z)-1]/Z,\qquad f(Z)=\chi_n(L\sqrt Z).
\end{aligned}
\label{eq:rest-vertex-symbol-uv}
\end{equation}
Use $d(0)=f'(0)$. Both $f$ and $d$ are entire in their scalar argument.
For a unit-charge free Dirac inverse, strip the common factors $i,e$
from its one-photon vertex and take
\begin{equation}
 \Gamma^\mu=\gamma^\mu+
       d(Z)\big[ Z\gamma^\mu-b^\mu(\gamma\cdot Q)\big].
 \label{eq:rest-ward-vertex-uv}
\end{equation}
Contraction gives $Q_\mu\Gamma^\mu=\gamma\cdot Q$, exactly the
free inverse difference in these units. The vertex is symmetric
under interchange of the charged momenta, regular at $Q=0$ and
equal to $\gamma^\mu$ there. On the two equal mass shells $P\cdot Q=0$.
The continued Dirac equations kill the term containing
$\gamma\cdot Q$ between physical spinors. The on-shell Dirac
form factor is exactly the rest-space $H$, on both physical rays.

This construction is bounded in its \emph{scalar} shell argument
on the real Euclidean contour since $Z\ge0$. It also identifies a
hard direction the scalar mark cannot see. Fix $Q\ne0$ and take
$P$ in its three-dimensional orthogonal complement. Then $Z$ is
independent of $P$ and any transverse polarization samples
$f(X+X^2/(4m^2))\,\gamma\cdot\epsilon$. More generally, in
a slab $|P\cdot Q|\le c$ and for a unit polarization orthogonal
to both $P$ and $Q$, the same contraction equals
$f(Z)\gamma\cdot\epsilon$. On a sufficiently small interval of
nonzero $X$ and sufficiently narrow slab its norm has a strictly
positive lower bound at arbitrarily large $|P|$. Thus the one-mark
Ward identity, exact physical form factor and absence of Euclidean
exponential growth can all be realized together. An absolute bound
for the charged circulation must still use the propagators, higher
contacts and boundary measure of a specified action. The one-mark
calculation alone cannot decide their combined ultraviolet degree.

\paragraph{Gauging this vertex at both endpoints.}
The unseen directions can be tested in a complete charged determinant.
Start with the local Euclidean Dirac inverse with mass $m$.
Write $V^\mu=\Gamma^\mu-\gamma^\mu$. Take its linear coupling
$ie V^\mu(P,Q) A_\mu(Q)$ as a midpoint Fourier symbol. In position
space the midpoint is $(x+y)/2$ and $P$ is conjugate to $y-x$.
Multiply this bilocal kernel by the straight Abelian parallel
transport from $y$ to $x$, and add it to the minimally gauged inverse.
The linear kernel is invariant under a gauge shift since $Q\cdot V=0$;
the link supplies the two endpoint phases. On the real Euclidean
contour this symbol has polynomial growth and defines a tempered
kernel on test fields. Reversal of the endpoints gives the usual
Euclidean Dirac adjoint convention. This specifies the higher
contacts as well as the one-photon vertex; independent terms
quadratic in the field strength would change the action.

Take $Q=q e_4$ with $q>0$ fixed and choose the transverse polarization
$e_1$. Write $t=P\cdot Q$, $r=|\mathbf P|$. The complete linear vertex is
\[
 \Gamma^1=f(Z)\gamma_1-
        \frac{q\,t}{m^2}d(Z)P_1\gamma_4.
\]
Choose $q$ small and keep $t$ in a fixed compact positive interval
where $|d(Z)|$ is bounded below. Such choices exist because $f'(0)\ne0$.
Restrict the spatial directions to $|P_1|>r/2$. For both orientations
of the photon the term proportional to $P_1$ has the same sign.
With the free propagators at $P\pm Q/2$, the leading absolute
spin trace of the two-vertex loop is
\[
 4\left|q\,t\,d(Z)/m^2\right|^2 P_1^2/r^2+O(r^{-1}).
\]
This follows from
$\operatorname{tr}[(\boldsymbol\gamma\cdot\mathbf P)\gamma_4
 (\boldsymbol\gamma\cdot\mathbf P)\gamma_4]=-4r^2$;
the remainder is uniform on the stated region.

The quadratic contact is determined by one derivative of the endpoint
link multiplying the linear kernel. A $1$-polarized link insertion
brings down $(y-x)_1$ and integrates its position along the chord.
In the Fourier symbol this differentiates $V^1$ in $P_1$ at
$P+sQ$, $|s|\le1/2$, plus the exchanged photon assignment.
At fixed $Q=\pm q e_4$, $Z$ and $P\cdot Q$ are independent of $P_1$.
The derivative is therefore bounded uniformly in $r$, also at all
the shifted arguments. The contact trace, with one massive propagator,
is $O(r^{-1})$. The absolute value of the full two-mark integrand
is bounded below by a positive constant for large $r$ in this region.
Its integral out to spatial radius $R$ grows at least as $cR^3$.
This is an exhaustion of a divergent bare coefficient. No finite
momentum scale was inserted in the action. Thus this explicit
gauge completion of the successful pole form factor already fails
the absolute UV test at two currents. The conclusion uses the
specified straight chord and local free charged inverse.

\paragraph{A virtuality-dependent replacement and its small-transfer cost.}
Write $X=Q^2>0$, $P=(k'+k)/2$ and
$\Delta=P^2+X/4+m^2$. In the Euclidean comparison define
\[
 \begin{aligned}
 Z&=X\left[1+\frac{X}{4m^2}+(\Delta/m^2)^2\right],\\
 h(P,Q)&=\chi_n(L\sqrt Z),\\
 \Gamma^\mu&=h\gamma^\mu+\frac{1-h}{X}Q^\mu\slashed Q .
 \end{aligned}
\]
The quotient extends by the entire power series, the Ward contraction
is $\slashed Q$, and $\Delta=0$ at equal charged poles. This gives
the same charged-rest form factor there. Consider the raw Dirac
inverse $D_0=i\slashed P+m$ with the transverse correction linear
in the field strength and gauge completed by its straight midpoint
Wilson chord. Fix $Q=q e_4$, $q>0$, and use polarization $e_1$.
In the convention $D_1=ie h\gamma_1$ the two-mark contact is
\[
 D_2=2ie^2\gamma_1\int_{-1/2}^{1/2}
                \partial_{P_1}h(P+sQ,Q)\,ds .
\]
The chord gives the derivative through its endpoint displacement.
The two equal terms exchange photon labels. The transverse second
variation of $-\log\det D$ is the integral of
$\operatorname{tr}(S_+D_1S_-D_1-SD_2)$ against
$d^4P/(2\pi)^4$, with $S_\pm=D_0(P\pm Q/2)^{-1}$ and
$S=D_0(P)^{-1}$. Both terms are absolutely integrable at each
$q>0$ for $n\ge4$. Indeed the Bessel argument grows as $|P|^2$,
and both $\chi_n(x)$ and $\chi'_n(x)$ are
$O(x^{-(n+1)/2})$ at infinity.

The small-transfer limit does not commute with this integration.
For the indicated two-mark coefficient one obtains
\begin{equation}\label{eq:virtmark-smallq-uv}
 \Pi_{11}(qe_4)=\frac{e^2m^2}{8\pi^2 Lq}
       \int_0^\infty[2\chi_n(t)-\chi_n(t)^2]dt+o(q^{-1}).
\end{equation}
To see the coefficient, set $P=\rho w$ with
$\rho=m/\sqrt{Lq}$ and $r=|w|$. Then
$h(\rho w+sQ,Q)\to j(r)=\chi_n(r^2)$ and its transverse
$w$ derivative converges to that of $j$. After division by
$-e^2\rho^2$, the two spin traces with the rescaled measure
have angular means $2j^2/r^2$ and $2j'/r$, respectively,
before the radial measure $r^3dr/(8\pi^2)$.
Here $\langle w_1^2\rangle=r^2/4$.
For small $q$, the rescaled integrands prior to the radial measure
are bounded by $C(1+r^{-2})$ for $r\le1$, and by
$C r^{-(n+1)}$ for $r\ge1$. These estimates follow from the
positive mass and the same Bessel derivative bound, uniformly along
the bounded chord. Dominated convergence and integration of
$r^2j'(r)$ by parts give \eqref{eq:virtmark-smallq-uv} after
$t=r^2$. If $g_n$ is the normalized one-coordinate ball density,
Fourier inversion and Plancherel give
\[
 \int_0^\infty(2\chi_n-\chi_n^2)dt
 =\pi\!\left(2g_n(0)-\int_{-1}^1g_n(x)^2dx\right)>0,
\]
since $0\le g_n\le g_n(0)$ and $\int g_n=1$.
Thus the virtuality factor cures the fixed-transfer hard cone of
the preceding vertex, yet its complete straight-chord two-mark
response has an inverse-$|Q|$ singularity. Even exact
transversality leaves this failure of the regular Maxwell
small-transfer limit visible. At $n=248$ the coefficient of
$e^2m^2/(Lq)$ is $0.323249$ to six decimal places.
Here $\int g_n^2=g_n(0)^2\mathrm B(1/2,n)$, obtained by direct
integration of the one-coordinate ball density.

Increasing the power of the virtuality delays the same failure.
Replace $(\Delta/m^2)^2$ by $(\Delta/m^2)^{2N}$ with integer $N>1$.
The hard scale is now $\rho=m(Lq)^{-1/(2N)}$. With $\beta=1/N$,
the same rescaling gives
\[
 \Pi_{11}(qe_4)\sim
 \frac{e^2\rho^2}{8\pi^2 N}
 \int_0^\infty t^{\beta-1}(2\chi_n(t)-\chi_n(t)^2)\,dt .
\]
Indeed the rescaled form factor tends to $\chi_n(r^{2N})$;
its radial derivative gains $r^{2N-1}$ while the Bessel bound
gives $O(r^{-N(n+1)})$ at infinity. These estimates justify the
same integration by parts. For independent ball coordinates
$Y,Y'$ the integral is $\Gamma(\beta)\cos(\pi\beta/2)$ times
$2\mathbb E|Y|^{-\beta}-\mathbb E|Y+Y'|^{-\beta}$. This follows
by inserting $e^{-\epsilon t}$ in the cosine integral, using
$|\epsilon-ix|^{-\beta}\le|x|^{-\beta}$, then taking the limit.
The negative moments exist since $\beta<1$ and the densities are
bounded. A centered interval has maximal probability under the even,
decreasing coordinate density. Integrating this inequality over the
positive superlevel sets of $|x|^{-\beta}$ gives
$\mathbb E|Y+y|^{-\beta}\le\mathbb E|Y|^{-\beta}$ for every $y$.
The coefficient is strictly positive. Thus every finite power in
this family leaves a transverse singularity of order $q^{-1/N}$.

\paragraph{A two-record contact that repairs the soft photon response.}
A quadratic charged kernel can keep the positive A4 measure and the
straight endpoint Wilson chord. It exposes exactly which new charged
operator would repair that response. Use $X_i=Q_i^2$, $x=P^2$,
$D=x+m^2$ and let $h_i=h(P,Q_i)=f(Z_i)$ be the original
virtuality-dependent weight, $f(Z)=\chi_{248}(L\sqrt Z)$.
Define a new quadratic charged kernel at its midpoint by
\begin{align}
 \delta D[A,A]&=\frac{e^2}{m}\,
  F_{\mu\nu}(Q_1)F_{\mu\nu}(Q_2)\,\mathcal K,\nonumber\\
 \mathcal K&=\frac{h_1h_2}{12D}
 -\frac{x}{8}\left[
 (1-h_1)\frac{\partial_x h_2}{X_2}
 +(1-h_2)\frac{\partial_x h_1}{X_1}\right].
 \label{eq:virial-contact-uv}
\end{align}
Here $\Delta_i=x+X_i/4+m^2$ as in the one-mark vertex.
The real transverse field-strength contractions use the same midpoint;
momentum integrations and the final endpoint Wilson chord are understood.
The apparent singularities are removable:
$\partial_x h_i/X_i=2\Delta_i f'(Z_i)/m^4$.
The symmetric products require independent records. The derivative is
the virtuality derivative of the shell characteristic function. This
operator changes no free pole or one-mark matrix element. It is
additional action data selected here by a soft matching condition.

For a transverse component at $Q=q e_4$, $q>0$, the new second mark
is $4e^2q^2\mathcal K/m$. The scalar Dirac trace contributes
$-16e^2q^2\int\mathcal K D^{-1}d^4P/(2\pi)^4$. To compute the sum,
first retain the exact $h(P,Q)$ and omit only the $\pm Q/2$ shifts
in the two free propagators and the $sQ$ shifts in the straight
chord. Angular averaging and a radial integration by parts give this
part of the original response exactly as
\[
 \frac{e^2}{8\pi^2}\int_0^\infty
   \left(1-\frac{m^4}{D^2}\right)(2h-h^2)\,dx .
\]
The derivative part of \eqref{eq:virial-contact-uv} contributes
$e^2/(8\pi^2)\int_0^\infty (x^2/D)\partial_x(2h-h^2)\,dx$.
Since $(x^2/D)'=1-m^4/D^2$ and the boundary terms vanish, these
expressions cancel at every positive $q$. This identity accounts for
the full contribution from the moving large-momentum scale in the
unshifted integrand.

In the shifted bubble divided by $q^2$, the remaining coefficient
at fixed $x$, after its $h^2/(12D)$ contact, is
$4m^2/(3D^3)+m^4/(3D^4)$ in units $e^2h^2$.
This follows by expanding the two propagators with the angular
means $\langle P_1^2\rangle=x/4$ and
$\langle P_1^2P_4^2\rangle=x^2/24$.
It has $x\,dx$ integral $13/18$. The discarded part of this
propagator expansion, after division by $q^2$, is bounded by
$Cq^2/D^3$ for $q\ll m$. Thus the first part of the limit is an
absolutely convergent integral. The residual chord term has a
finite contribution from the moving scale. With $u=Lqx/m^2$ its
coefficient in units $e^2/(16\pi^2)$ is
\[
 -\frac43\int_0^\infty\!
 \left[\frac{u}{4}\chi_{248}''(u)
              +\frac{u^2}{12}\chi_{248}'''(u)\right]du
 =-\frac19.
\]
Here one expands $\partial_xh(x+2s q P_4+s^2q^2)$ about $x$ and
uses $\int_{-1/2}^{1/2}s^2ds=1/12$. The linear term odd in $s$
vanishes. The first three derivatives of the ball characteristic
function decay as $O((1+u)^{-249/2})$, giving convergent majorants
for the displayed rescaled limit. Combining the two terms yields
\begin{equation}
 \lim_{q\downarrow0}\frac{\Pi_{11}(q e_4)}{q^2}
       =\frac{11e^2}{288\pi^2}.
 \label{eq:virial-maxwell-uv}
\end{equation}
For $L=m=e=1$, direct integration gives $0.00382971$, $0.00386030$
and $0.00386892$ at $q=0.04,0.01,0.001$, respectively, against the
limit $0.00386991$. This construction preserves the original
candidate's on-shell form factor and gives a finite Maxwell response
at two marks. A derivation of \eqref{eq:virial-contact-uv} from the
GCM boundary dynamics, its higher-order continuation, and positivity
of its own closed spectrum remain separate requirements.

\paragraph{The residual branch and a pole-free continuation.}
A finite Maxwell coefficient alone does not give a mass gap.
To test the next order, integrate the original straight-chord contact
by parts in $P_1$ and translate the integration variable. Its shifted
part now contains $8h[J(P-sQ)-J(P)]$, averaged over
$-1/2\le s\le1/2$, where $J(P)=D^{-1}-2P_1^2D^{-2}$.
At fixed radius the angular mean is
$\langle\partial_4^2J\rangle=-2m^4/D^4$. Odd powers of $s$
drop out. Together with the bubble, the response divided by
$e^2q^2/(16\pi^2)$ becomes
\[
 \int_0^\infty x\,dx\left[
  \left(\frac{4m^2}{3D^3}+\frac{m^4}{3D^4}\right)h^2
               -\frac{2m^4}{3D^4}h\right]+O(q^2/m^2).
\]
The fourth-order propagator and chord Taylor remainders after
division by $q^2$ are bounded by $Cq^2/D^3$. Their coefficients
divided by $q^2$ have integrable limits. Set $a=Lq$. Only the
$4m^2/(3D^3)$ term changes at order $a$. Substitution
$u=aD/m^2$ and dominated convergence, using
$|1-\chi_{248}(t)^2|\le C\min(t^2,1)$ for real $t$, gives
\begin{equation}\label{eq:virial-soft-cusp-uv}
 \frac{16\pi^2\Pi_{11}}{e^2q^2}
   =\frac{11}{18}
     -\frac{2\pi}{3}\mathbb E|Y+Y'|\,Lq+o(Lq).
\end{equation}
Indeed the linear coefficient is
$(4/3)\int_0^\infty (\chi_{248}(u)^2-1)du/u^2$.
Here $Y,Y'$ are independent coordinates of the normalized A4 ball.
For $n=248$ their absolute-sum mean is
$4g_n(0)^2 B(1/2,n+1)/(n+1)\simeq0.07140051$, where $g_n$ is
the normalized coordinate density. The formula follows by averaging
the two signs and integrating the density of
$\max(|Y|,|Y'|)$. The nonzero $\sqrt{Q^2}$ term rules out analyticity
at zero transfer and hence a purely gapped dispersion with ordinary
local subtractions for this two-mark coefficient.

The factor $1/D$ in the scalar contact also needs a pole test in an
open charged amplitude. With equal continued charged poles and
$P=(p+p')/2$, $D=-(p'-p)^2/4$ in the Euclidean conventions.
At forward transfer the factor can diverge for virtual external
photons with a nonzero field-strength contraction. Replace only that
$1/D$ in \eqref{eq:virial-contact-uv} by
\[
 \begin{gathered}
 T(D)+m^2T(D)^2,\\
 T(D)=\frac{1-w(D/m^2)}{D},\qquad w(y)=\chi_{248}(y)^2 .
 \end{gathered}
\]
Evenness and $w(0)=1$ make $T$ entire with $T(0)=0$. At large
positive $D$ the replacement equals $D^{-1}+m^2D^{-2}$ up to terms
with a rapidly decaying factor $w$. This removes the open midpoint
pole and cancels the $D^{-3}$ tail responsible for
\eqref{eq:virial-soft-cusp-uv}.

Writing $w=w(D/m^2)$, the integrand in brackets in the preceding
radial formula becomes
\[
 \frac{m^4}{3D^4}[(h-1)^2-1]
 +\frac43 h^2
 \left[\frac{w}{D^2}+\frac{m^2(2w-w^2)}{D^3}\right].
\]
The first square contributes
$a^2\int_0^\infty(1-\chi_{248}(u))^2du/(3u^3)+o(a^2)$.
Both endpoints are integrable. The $w$ terms admit a second-order
expansion in $q$ under the integral. Together with the convergent
Taylor remainder above, this yields
\begin{equation}
 \frac{16\pi^2\Pi_{11}}{e^2q^2}=C_0+C_2q^2+o(q^2)
 \label{eq:fix-entire-contact-uv}
\end{equation}
for a finite $C_2$.
The constant is
\[
 C_0=-\frac1{18}
 +\frac43\int_1^\infty(y-1)
 \left[\frac{w(y)}{y^2}
       +\frac{2w(y)-w(y)^2}{y^3}\right]dy .
\]
Numerically $C_0\simeq2.63553431$. To see its strict positivity,
observe that on $[1,2]$,
$w\ge(1-4/500)^2$ by the coordinate variance. The first positive
term alone therefore exceeds $(4/3)(1-4/500)^2(\log 2-1/2)>1/18$.
For $L=m=e=1$ direct integration gives $\Pi_{11}/q^2=0.01665162$,
$0.01668016$, and $0.01668912$ at $q=0.08,0.04,0.01$, against
the limit $0.01668972$. This continuation passes the forward-pole
test and has a first derivative in $Q^2$ at zero. Analyticity on a
neighborhood of zero and positivity of its full continued spectrum
require additional arguments.

\paragraph{Putting the same form factor into an off-shell connection.}
The rest-space construction has a simple off-shell comparison on the
Euclidean plane. It helps identify what a successful charged source would
have to do differently. Let $X=|Q_E|^2$. In a normalized A4 ball of
dimension $n\ge5$, use four orthonormal coordinates $Y_{1:4}$ and one
more orthogonal coordinate $Y_5$ from the \emph{same} draw. Rotational
invariance gives the exact identity
\begin{equation}
\begin{aligned}
 h(X)&=\mathbb E\exp\!\left[iL
       \left(Y_{1:4}\cdot Q_E+Y_5\frac{X}{2m}\right)\right]\\
 &=\chi_n\!\left(L\sqrt{X+X^2/(4m^2)}\right).
\end{aligned}
\label{eq:rest-five-coordinates-uv}
\end{equation}
Both terms in the parentheses have momentum units. This choice couples
the fifth mark to virtuality. That coupling is additional operator data.
Analytic continuation $X=-S$ gives exactly $H(S/(4m^2))$. On the
Euclidean ray it is smooth at zero, equals one there, and is
$O(|Q_E|^{-(n+1)})$ at large momentum.

As a direct test, replace the photon connection throughout a charged
functional by $B(Q)=[h(X)P_T(Q)+P_L(Q)]A(Q)$. This map preserves
Abelian gauge shifts. Its apparent singularity at $Q=0$ is removable
since $1-h(X)=O(X)$. The chain rule applies at every order,
including contacts. At two currents the old closed spectral density
and the old squared pole matrix element are \emph{both} multiplied by
$|H|^2$ in the transverse channel. Wherever $H\ne0$ their ratio is
therefore unchanged. Even at a real zero, both vanish, and the strict
deficit on neighboring energies persists. The entire continuation of the
common factor permits this local statement about the cut without a
large-complex-energy dispersion assumption. This calculation explains why
the successful inequality for a trial \emph{charged} form factor does
not transfer to the determinant by simply changing the photon readout.

\subsection{Constructing the charged current on its physical mass shells}
\label{sec:fix-fock-uv}
The rest-space shell can be tested on actual charged states before
choosing how a virtual photon is sewn. Work on the Minkowski readout
of one locked sector-plane. Use the positive Hilbert space of a free
mass-$m$ Dirac field, with both particle and antiparticle states, and
the normalization $mL=1$, $n=K-1=248$ for the explicit checks.
Write $h(s)=H(s/(4m^2))$ for the rest-space form factor in
\eqref{eq:charged-rest-readout-uv}. Let $j_D^\mu$ be the normally
ordered free Dirac current of unit charge. Define an operator-valued
distribution on this Hilbert space by its real Fourier transform:
\begin{equation}
 \widetilde J_h^\mu(Q)=h(Q^2)\widetilde j_D^\mu(Q).
 \label{eq:fix-fock-current-uv}
\end{equation}
The Wilson and Euclidean vertex calculations use independent virtual
energies. Here $Q^0$ is the \emph{actual} difference of energies at a
number-preserving matrix element, and it is their sum at a pair-creation
matrix element. These two operations need not give the same
time-ordered distributions. The scalar $h$ is smooth on the real axis.
Its derivatives as a function of real four-momentum grow at most
polynomially, by the compact shell integral on the two physical rays
and the entire power series on the intervening bounded interval.
It therefore acts on Schwartz tests. Current conservation,
Hermiticity and the Lorentz-vector transformation law follow from the
corresponding properties of $j_D$. The integrated charge still counts
particles minus antiparticles: density conservation kills the pair
matrix element at zero spatial transfer, and $h(0)=1$ on a diagonal
charged state. At nonzero spatial transfer the density has the boosted
pair matrix element. The Hilbert space and the current in this test
are specified together.

\paragraph{All finite orders in a smooth external source.}
Let $a_\mu(t,\mathbf x)$ be a smooth classical probe, compactly
supported in time and rapidly decreasing in space with uniform
Schwartz seminorms. On an equal-clock-time slice of the reconstructed plane, decompose $J_h(a)$ into
the two number-preserving blocks and the pair creation and annihilation
blocks. With spinors divided by $\sqrt{2E_p}$, the spin matrices have
a uniform norm bound. The kernel of each number-preserving block,
even after taking its pointwise absolute value, is bounded by a
constant times $|\widehat a(t,\mathbf p'-\mathbf p)|$, summing the
four probe components. Young's inequality gives a bounded operator
on one-particle momentum space.

The pair-creation block is Hilbert--Schmidt even after taking the
pointwise absolute value. In fact, for fixed $\mathbf k=\mathbf p+
\mathbf p'$ and $s=(E_p+E_{p'})^2-|\mathbf k|^2$,
\begin{equation}
 \int d^3p\,|h(s)|^2
 \le C\int_{4m^2}^\infty \sqrt{s+|\mathbf k|^2}\,|h(s)|^2\,ds .
 \label{eq:fix-fock-HS-uv}
\end{equation}
Insert a delta function for $E_p+E_{p'}$, change total energy
to $s$, and restore the invariant two-body measure
$d^3p\,d^3p'/(4E_pE_{p'})$. Its integral at fixed total
four-momentum is bounded. Also $4E_pE_{p'}\le s+|\mathbf k|^2$.
These two facts give the displayed bound. The right side is at most
$C_{m,L}(1+|\mathbf k|)$ in fixed units by the real-axis decay of $h$.
Integrating it against $|\widehat a(t,\mathbf k)|^2$ proves the claim.

In the vacuum Wick expansion at distinct ordered times, each closed
fermion cycle contains at least one pair creation and one annihilation:
the vacuum is empty of physical particles and holes. More explicitly
the earliest and latest vertices on the cycle have two outgoing and
two incoming physical contractions, respectively. The two
Hilbert--Schmidt norms bound the absolute trace, with ordinary operator
norms for the intervening absolute kernels. Additional pair kernels
are also bounded. Time evolution on the charged lines contributes
phases of modulus one. The time integrals of these bounds are finite.
The same argument with endpoint wave packets bounds open chains.
Thus \emph{each fixed-order canonical matter graph with these classical
probes} has an absolute momentum bound. This uses the geometric
shell weight at the pair transitions. Making the probes themselves
quantum photon fields introduces further integrations, and passing to
a covariant time ordering requires its contacts.

\paragraph{A positive spectral realization of the two-current contact.}
The same Fock current fixes a covariant vacuum-to-pair density
including both spins and the complete two-body phase space:
\[
 w(s)^2=\rho_h(s)=
  \frac{\sqrt{1-4m^2/s}\,(1+2m^2/s)}{12\pi^2}|h(s)|^2,
 \quad s\ge4m^2 .
\]
This is an equality from the charged Hilbert space. It supplies an
explicit target for the quadratic contact. Put
$M_j=\int_{4m^2}^\infty s^j\rho_h(s)ds$.
The moments $M_{-1},M_0,M_1$ are finite and positive at the
stated dimension. For instance $M_0/m^2\simeq0.20579$ and
$M_{-1}\simeq0.01364$ at $mL=1$. The transverse canonical
spatial-current covariance at $Q=(\omega,\mathbf k)$ has
numerator $s$ on a polarization orthogonal to $\mathbf k$. With
$X=\omega^2+|\mathbf k|^2$, its spectral integral satisfies
\[
 \begin{aligned}
 \int\frac{s\rho_h(s)}{X+s}\,ds&=M_0-X\Pi_h(X),\\
 \Pi_h(X)&=\int\frac{\rho_h(s)}{X+s}\,ds .
 \end{aligned}
\]
The $M_0$ is fixed by the same spectral density as the exchange.

One can obtain their required combination from an action. Introduce
a continuum of real vector fields $B_{s,\mu}$ with $s\ge4m^2$. In
the Euclidean quadratic form take the positive density
\begin{equation}
\begin{aligned}
 {\cal L}_2={}&\tfrac14 F_A^2+
 \int_{4m^2}^\infty ds\,
 \left[\tfrac14 F_{B_s}^2+
       \tfrac12|\sqrt{s}B_s-e\,w(s)A|^2\right].
\end{aligned}
\label{eq:fix-fock-square-uv}
\end{equation}
Under $A\mapsto A+\partial\alpha$ shift
$B_s\mapsto B_s+e\,w(s)\partial\alpha/\sqrt{s}$. This shift has
finite species norm by $M_{-1}<\infty$. On each transverse
polarization, Gaussian integration gives the exact inverse
$X[1+e^2\Pi_h(X)]$. The square generates the direct term
$e^2 M_0$ and the exchange $-e^2\int s\rho_h(s)/(X+s)ds$ in one
operation. The source therefore has regular Maxwell behavior at
zero and precisely the specified leading current cut. For the
dynamical photon the transverse mass operator on the photon and
continuum variables is nonnegative by the same square. Its zero
vector is proportional to $(1,e w(s)/\sqrt{s})$ in species space.
The massless photon residue in this quadratic model is exactly $(1+e^2 M_{-1})^{-1}>0$,
and its remaining spectral measure is nonnegative by the spectral
theorem. This quadratic realization matches the charged current's
vacuum two-point data. Its higher Gaussian moments do not represent
the higher fermionic current correlators.

Closing the continuum on the ordinary Maxwell photon now costs an
explicit amount. Per unit volume in four dimensions the relative
transverse half log determinant is
$\tfrac32\int d^4Q/(2\pi)^4\log(1+e^2\Pi_h(X))$ relative to $e=0$, with the source-independent free factors removed. Since $X\Pi_h(X)\to M_0$, the quadratic
term in $e$ grows as $3e^2 M_0 R^2/(32\pi^2)$ on balls
$|Q|\le R$. The fourth-order term grows as
$-3e^4 M_0^2\log(R^2/m^2)/(64\pi^2)$ up to a bounded term.
Every coefficient of order $e^{2j}$ with $j\ge3$ in this relative
determinant is absolutely convergent: $\Pi_h$ is bounded at zero
and is at most $M_0/X$ at infinity. The same leading tail appears
in the connected, vacuum-normalized two-probe Maxwell-density
test in Section~\ref{sec:nov-ring-uv}, whose absolute integrand
is bounded below by a constant times $|\widehat a(q)|^2\Pi_h(p^2)$.
Thus achieving the positive pair cut and the exact massless photon
residue repairs the two-current sector. It also exhibits explicitly
the unsuppressed contraction when these currents are sewn with the
ordinary physical photon.

\paragraph{A spacelike commutator in the charged sector.}
The operator in \eqref{eq:fix-fock-current-uv} lets us test whether
its covariant matrix elements come from a local current. At one
time let $p$ and $q$ denote collinear signed spatial momenta and
$E_k=\sqrt{m^2+k^2}$. For an incoming positive-energy electron,
sum over all intermediate spin states of the free Dirac Hamiltonian.
Writing $\eta=\pm1$ for the energy sign gives
\[
\begin{aligned}
 w_\eta(p,k)&=\tfrac12\left[1+
           \eta\frac{m^2+pk}{E_p E_k}\right],\\
 s_\eta(p,k)&=2m^2+2pk-2\eta E_pE_k,\\
 B(p,k)&=\sum_{\eta=\pm1}w_\eta(p,k)|h(s_\eta(p,k))|^2 .
\end{aligned}
\]
Indeed $w_\eta$ is the spin-averaged product of the incoming Dirac
projector and the intermediate energy projector. At equal time the
density vertex has the identity as its spin matrix between these
projectors. The canonical anticommutation relations give the forward
matrix element of the density commutator at transfers $q,-q$,
with its vacuum value removed, as
\[
 C_p(q)=B(p,p-q)-B(p,p+q).
\]
For bilinears this is the one-body commutator: any scalar from normal
ordering drops out in the vacuum difference. Both signs of the
intermediate energy are essential.

At fixed $p$ all the absolute transfers $|s_\eta(p,p\pm q)|$ tend
to infinity as $|q|\to\infty$, since $E_p>|p|$. Hence $C_p(q)\to0$.
But set $p=q$ and let $p/m\to\infty$. The intermediate momentum
zero contributes zero in the limit. At $k=2p$ the positive-energy
weight tends to one and $s_+(p,2p)\to-m^2/2$. The other weight
tends to zero. Consequently
\[
 \lim_{p/m\to\infty}C_p(p)=-|h(-m^2/2)|^2<0 .
\]
The shell argument at this point is $3/4$ when $mL=1$. The result
is approximately $-0.997753$ for $n=248$. There are therefore
fixed finite momenta $p$ for which $C_p$ is nonzero and decays
along the same real line. A distribution in the equal-time relative
coordinate supported only at zero has a polynomial Fourier transform
on that line; no such polynomial has these properties. The calculation
can equivalently be made with a narrow external momentum packet,
integrating the center coordinate to take the forward matrix element.
The fixed-$p$ asymptotics are uniform on compact sets of packet momenta.
Thus the charge density of this otherwise covariant current fails
spacelike commutativity. Positive vacuum two-point data alone miss
this obstruction.

The same calculation gives a general restriction. Suppose a scalar
Dirac form factor on this free Fock space is continuous, unit
normalized at zero and tends to zero on both real hard rays. A local
equal-time density would force the polynomial $C_p(q)$ to vanish for
every fixed $p$. Now take $q=\alpha p$ with $0<\alpha<1$ and send
$p/m\to\infty$. Both intermediate spatial momenta are positive.
Their positive-energy weights tend to one. Using
$E_p E_{cp}-cp^2\to m^2(c+c^{-1})/2$ for $c>0$ gives
\[
 |h(-m^2 u)|^2=|h(-m^2 v)|^2,\qquad
 u=\frac{\alpha^2}{1-\alpha},\quad
 v=\frac{\alpha^2}{1+\alpha}.
\]
Every $u>0$ occurs exactly once. The resulting continuous map
$u\mapsto v$ has $0<v<u$ and no positive fixed point, so its iterates
decrease to zero. Continuity would give $|h(-m^2u)|=1$ at every
spacelike transfer, contradicting decay. This excludes a local
density for this entire class of scalar modifications of the free
Dirac bilinear. New current operators with further intermediate
states or different spin structure require a new commutator test.

\paragraph{Can the pole form factor be an overlap of boosted boundary states?}
Suppose elastic suppression is due solely to a normalized internal
boundary state of the same particle. If a unitary one-parameter
boost transports that state through relative rapidity $t$, its
overlap $k(t)$ must be a positive-definite function on the boost
line. For the proposed rest-space readout at collinear velocities,
$k(t)=\chi_n(mL\sinh t)$. Three rapidities $-t,0,t$ already decide
the test. The putative squared norm for coefficients $(1,-2,1)$ is
$6-8k(t)+2k(2t)=k^{(4)}(0)t^4+O(t^6)$. The ball-coordinate
moments give, with $d=mL$,
\[
 k^{(4)}(0)=\frac{d^2}{n+2}
                  \left(\frac{3d^2}{n+4}-4\right).
\]
It is strictly negative at the tested native scale. A covariant
family of one-mark probability laws, or a perfectly good operator
form factor, need not define positive inner products between the
states assigned to different velocities.

There is a second test for any normalized, real, even overlap of
this kind. By the spectral theorem for the boost generator it has
the form $k(t)=\int\cos(t\lambda)d\mu(\lambda)$ with $\mu$ a
symmetric probability. Suppose the boost state admits exponential
moments for every imaginary rapidity $0<y<\pi$, and the elastic
continuation is regular between zero and pair threshold, with a
continuous common threshold value for the two physical channels.
In that strip,
\[
 h(4m^2\sin^2(y/2))=k(iy)=\int\cosh(y\lambda)d\mu(\lambda).
\]
The limit at $y=\pi$ is strictly greater than one, possibly
infinite, for any positive weight at $\lambda\ne0$. The
paired-shell optical ceiling tends to one there. Thus a
nonconstant, purely Dirac boundary overlap satisfying these
analytic hypotheses cannot solve the crossing test. The explicit
massive-link comparison below allows a subthreshold singularity.
A genuine local current can also involve several spin amplitudes
and matrix elements beyond a single transported boundary state.

\subsection{Which neutral-vector masses pass the charged-pair test?}
\label{sec:vector-link-uv}
Proposition~\ref{prop:carrier-portal-uv} obtains the measured electric
form factor $b/(b-S)$ from a massive vector link, where $S$ is timelike
squared momentum and $b$ is the squared mass at fixed external
connection. Section~\ref{sec:path-rest-tests-uv} instead obtains an
inclusive ceiling $r(\tau)$ from the paired shell at $mL=1$,
$n=K-1=248$, with $\tau=S/(4m^2)$. The two calculations give a
sharp mass selection when one asks whether the pair from the first
calculation fits beneath the continuum from the second.

Set $s_*=4m^2$. Consider any positive probability law $\nu$ on
$0<b<s_*$ and a purely Dirac form factor
$F(S)=\int b/(b-S)\,d\nu(b)$. Its normalization gives $F(0)=1$.
A finite mixture has an explicit gauge realization: give each
connection $B_j$ the portal action of
Proposition~\ref{prop:carrier-portal-uv}, using the same external
$A$ in every mass term, and couple the charge to
$\sum_j w_j B_j$, where $w_j$ are the probability weights.
A simultaneous gauge shift acts on this connection with unit weight.
Completing each independent gauge-field square gives the stated
form factor at the first order in the electric charge.

The complete pair inequality for this class is equivalent to one
threshold moment:
\begin{equation}
 |F(s_*\tau)|^2\le r(\tau)\ \ (\tau>1)
 \quad\Longleftrightarrow\quad
 \int\frac{b}{s_*-b}\,d\nu(b)\le1.
 \label{eq:vector-link-threshold-uv}
\end{equation}
Indeed, $r(\tau)\to1$ as $\tau\downarrow1$. Monotone convergence
at threshold proves necessity, including laws with infinite moment.
For sufficiency use
$(s_*\tau-b)^{-1}\le[\tau(s_*-b)]^{-1}$ to get
$|F(s_*\tau)|\le1/\tau$. For $1<\tau\le16$, the coordinate
variance $1/250$ of the unit ball and $\cos x\ge1-x^2/2$ give
$r(\tau)\ge[1-(\tau-1)/500]^2\ge\tau^{-2}$, since
$0\le v\le\tau-1$ throughout the defining average of
$|\chi_n(\sqrt v)|^2$. For $\tau\ge16$ the bound
$r(\tau)>89/\tau^2$ from Section~\ref{sec:path-rest-tests-uv}
suffices. For a single link the criterion is precisely $b\le2m^2$.
More generally strict convexity of $b/(s_*-b)$ implies
$\int b\,d\nu\le2m^2$, with equality only for the single atom
$b=2m^2$. Thus the largest possible mean squared mass fixes the
link mass without a second scale. This comparison leaves the
subthreshold neutral states as explicit additions to the spectrum.

There is a further distinction when the probed $A$ also gets a
canonical Maxwell term. For a single link the transverse eigenvalues
at Euclidean $X>0$ are $X$ and $X+2b$. Hence
\[
 D_{BB}=D_{AA}=\tfrac12\left(\frac1X+\frac1{X+2b}\right),
 \qquad D_{BA}=\frac{b}{X(X+2b)} .
\]
The pole of $D_{BA}/D_{AA}=b/(X+b)$ at $S=b$ is now a zero
of the propagator used for amputation. At tree level the two
dynamical vector poles lie at $S=0,2b$; at the extremal mass the
second falls exactly at the charged-pair threshold. Moreover
$D_{BB}=(X+b)^{-1}+[b/(X+b)]^2D_{AA}$. The first term is the
conditional fluctuation of the charged connection. It persists
even when the external transfer is small. The fixed-transfer
charged-loop divergence proved in Proposition~\ref{prop:carrier-portal-uv}
therefore remains a separate test of the native internal displacement.

\subsection{Closing a pair of paths on its boundary}
\label{sec:path-closure-tests-uv}
A second choice concerns which displacements the two sides of a
closed observable share. Proposition~\ref{prop:closure-native-perimeter-uv}
gives a general perimeter obstruction for thin currents. For two ordered
links the exact angular coefficient also tests the higher contacts.

\paragraph{A closed test that retains the same intermediate record.}
The difference of two paths has an ultraviolet cost even when their
net displacement and their duration are fixed. Work on the reconstructed
Euclidean four-plane. For real, noncollinear increments $a,b$ let $C$
be the parallelogram following $ab$ and returning along the reverse of
$ba$, with perimeter $P(C)=2(|a|+|b|)$. Reversal here denotes the
adjoint parallel transport. Define $j_C(q)=\oint_C e^{iq\cdot x}dx$.
The gauge-invariant positive square of the charged link difference is
\[
 |U_{ab}-U_{ba}|^2=2-2\cos\!\left(e\oint_C A\right).
\]
This includes the direct two-source variations of the links. At order
$e^2$ the Maxwell contraction of this quantity, with the charge
coefficient stripped, is
\begin{equation}
 I_C=\int\frac{d^4q}{(2\pi)^4}\frac{|j_C(q)|^2}{|q|^2}.
 \label{eq:paired-perimeter-closed-uv}
\end{equation}
The transverse projector acts as the identity since $q\cdot j_C=0$.
At zero momentum $j_C=O(|q|)$ by closure.

This square can be obtained by integrating explicit charged variables.
For a single record use two Grassmann slots transforming at its
two endpoints, and take the Hermitian matrix
\[
 B[A]=\begin{pmatrix}1&\overline{\kappa w}\\
                  \kappa w&1\end{pmatrix},
 \qquad w=U_{ab}-U_{ba},\quad |\kappa|\le\tfrac14 .
\]
The second slot transforms at the terminal endpoint. Then
$\det B=1-|\kappa|^2|w|^2$ and $B[0]=I$. Here $\kappa$ is a
dimensionless bounded record weight; an exchange-odd weight is
allowed. The coefficient of $-\log\det B$ is
$|\kappa|^2 I_C$ after photon contraction. The same cost occurs in
a normalized open observable for a prepared record:
$(\det B)(B^{-1})_{11}=1$, so its diagonal two-point function
after photon averaging is $1/\langle\det B\rangle_A$ and has
the same positive order-$e^2$ correction. This computation concerns
that finite-dimensional charged source and a common boundary record.

\paragraph{Exact hard-momentum cost of a thin polygon.}
Use the normalized angular average on $S^3$, written $\langle\cdot
\rangle_{\rm ang}$. The following limit identifies which history
contractions can use a fixed rank-four shell readout:
\begin{equation}
 \lim_{r\to\infty} r\,\langle|j_C(rn)|^2\rangle_{\rm ang}=4P(C).
 \label{eq:paired-perimeter-asymptotic-uv}
\end{equation}
Indeed, writing $I(t)=\int_0^1 e^{ist}ds$, the two contributions to
$j_C$ are $aI(q\cdot a)(1-e^{iq\cdot b})$ and
$bI(q\cdot b)(e^{iq\cdot a}-1)$ up to a common phase. For
$t=n\cdot a/|a|$ the angular density is $2\sqrt{1-t^2}/\pi$.
Plancherel gives $\int_{\mathbb R}|I(u)|^2du=2\pi$.
At fixed $t$, the conditional average of $e^{irn\cdot b}$ is
$e^{irt b_\parallel}\operatorname{sinc}
 (r|b_\perp|\sqrt{1-t^2})$, where $b_\parallel,b_\perp$
are the components along and perpendicular to $a$ and
$\operatorname{sinc}(x)=\sin(x)/x$. Rescale $u=r|a|t$.
Since $b_\perp\ne0$, the interference term tends to zero under
the integrable $|I(u)|^2$ majorant. Thus $r$ times the angular
mean of the squared $a$ contribution tends to $8|a|$.
The other square gives $8|b|$. To bound their cross term,
the projection of uniform measure on $S^3$ onto two orthonormal
coordinates is uniform on the unit disk. The joint law of
$(n\cdot a,n\cdot b)$ therefore has bounded compactly supported
density for each noncollinear pair. Using
$|I(t)|\le\min(1,2/|t|)$ in both variables gives
$O(r^{-2}\log^2(2+r))$ for the angular cross term.
As a check, for perpendicular unit sides the values of
$r\langle |j_C|^2\rangle_{\rm ang}$ at $r=20,320,1280$ from the
one-coordinate integral are $14.60834$, $15.97216$ and $15.99952$,
respectively. The predicted limit is $16$.
For collinear increments the two Abelian transports coincide and
$j_C=0$ exactly.

Consequently the improper integral in
\eqref{eq:paired-perimeter-closed-uv}, exhausted by balls of
radius $R$, grows as $P(C)R/(2\pi^2)+o(R)$. The exhaustion
measures a divergence in the bare coefficient. Averaging over
records with positive weight $|\kappa|^2$ cannot remove it.
For bounded marks and integrable perimeter the coefficient of
linear growth averages exactly, with collinear pairs omitted:
the separate angular squares obey the uniform bound
$r\langle |j_C|^2\rangle_{\rm ang}\le C_0 P(C)$ by the same
one-coordinate estimate. Retaining a nonzero probability of
noncollinear records with $\kappa\ne0$ gives a strictly positive
coefficient.

If the actual measured transverse connection has an additional
radial multiplier $c(r)$, continuous and bounded near zero,
the integrand instead contains $|c(r)|^2$. Equation
\eqref{eq:paired-perimeter-asymptotic-uv} gives the sharp test
\[
 I_{C,c}<\infty\quad\Longleftrightarrow\quad
       \int_1^\infty |c(r)|^2dr<\infty ,
\]
assuming local square integrability at finite momenta.
For instance the isometric Euclidean A4 ball readout
$c(r)=\chi_{K-1}(Lr)$ passes. The comparison makes explicit
whether the shell is used only to choose the sides of the
polygon or is also measured by the connection on each side.
Those choices give different closed photon contractions.
If two records are glued by an operator $T_r$ on their normalized
boundary space with $T_r\succeq p(r)I$, $p(r)\ge0$, the same
positive lower bound applies with $p(r)|c(r)|^2$ to the corresponding
angularly averaged quadratic form of the weighted path current.
In particular a hard-momentum-independent diagonal component with
a nonnegative operator remainder retains the unsmeared perimeter
divergence. A purely coherent average before squaring requires
its own contraction.

\paragraph{A finite two-mark test that fails at the next electric order.}
Consider an A4 boundary law for a common translation $\eta$ of a fixed
nondegenerate parallelogram. On the Euclidean plane choose four
orthonormal readout axes ($K\ge5$) and let the characteristic function
of this probability law be $c(q)=\chi_{K-1}(L|q|)$.
Write $\ell_\eta=\oint_{C+\eta}A$. The gauge-invariant, Hermitian
odd part of the relative holonomy
$U_{ba,\eta}^\dagger U_{ab,\eta}$ is $\sin(e\ell_\eta)$.
This compares the same two unreduced routes, using adjoint transport
only in the boundary pairing. Two distinct coherent boundary
quantities with exactly the same first derivative are
\[
 Y=\int\sin(e\ell_\eta)\,d\mu(\eta),\qquad
 Y'=\sin\!\left(e\int\ell_\eta\,d\mu(\eta)\right).
\]
Here $Y$ averages completed holonomies; $Y'$ averages the connection
before exponentiating. They make different predictions for the full
source contacts.

To compute the bare Taylor coefficients, first exhaust photon momenta
by balls of radius $R$. Write $C_R(\eta-\eta')$ for the free Maxwell
covariance of $\ell_\eta,\ell_{\eta'}$ and $V_R=C_R(0)$.
The variance of the average is
\[
 A_R=\iint C_R(\eta-\eta')\,d\mu(\eta)d\mu(\eta')
    =\int_{|q|\le R}\frac{d^4q}{(2\pi)^4}
                \frac{|c(q)j_C(q)|^2}{|q|^2}.
\]
It tends to a finite $A>0$: $c(0)=1$, the soft area response is
nonzero on an open set, and the sharp radial test above controls the
hard tail. For two centered jointly Gaussian variables the product
of their sines is half the difference of the cosines of their
difference and sum. Thus
\begin{equation}
 \langle Y^2\rangle_R=e^{-e^2V_R}\iint
         \sinh\big(e^2 C_R(\eta-\eta')\big)\,d\mu(\eta)d\mu(\eta').
 \label{eq:paired-coherent-Wilson-uv}
\end{equation}
The order-$e^2$ coefficient equals $A_R$. The order-$e^4$
coefficient equals $-V_RA_R$ and diverges linearly as $R\to\infty$
by \eqref{eq:paired-perimeter-asymptotic-uv}. It comes from a
contraction within the third derivative of one Wilson observable
sewn to the other observable's first derivative. The identity
includes the full fourth-order coefficient of this positive square.
It is a bare perturbative obstruction even if a pointwise limit at
fixed $e\ne0$ is taken in a different order. In a two-slot determinant
as above with off-diagonal entry $\kappa Y$ (and a common endpoint
transport if needed), the additional order-$e^4$ term of the
negative logarithm is $3|\kappa|^4 A_R^2/2$ after averaging, which
is finite and cannot cancel the term $-|\kappa|^2 V_RA_R$.
For the second preparation,
$\langle(Y')^2\rangle_R=(1-e^{-2e^2A_R})/2$; every fixed-order
coefficient has a finite limit. Agreement of the one-photon vertices
and of their two-mark test therefore does not fix the ultraviolet
behavior of the contacts of a paired history.

There is also an exact test at fixed nonzero $e$. In the same
ball exhaustion the Euclidean Maxwell scalar kernel away from
the origin is
\[
 D_R(x)=\frac{1-J_0(R|x|)}{4\pi^2 |x|^2},
\]
by the four-dimensional angular integral and $J_0'=-J_1$.
The difference set $C-C$ is a finite union of sets of dimension
at most two. For $\eta-\eta'\notin C-C$ the two compact contours
are separated, hence the double contour integral $C_R(\eta-\eta')$
is bounded uniformly in $R$. The projected A4 ball law has a
density in the four readout coordinates, so this separation
holds almost surely for independent $\eta,\eta'$. The Gaussian
sine product has absolute value at most one. As $V_R\to\infty$,
dominated convergence in \eqref{eq:paired-coherent-Wilson-uv}
now gives $\langle Y^2\rangle_R\to0$ for each fixed $e\ne0$.
Thus taking the fixed-coupling limit erases this charged response
even though its second Taylor coefficient has a finite positive
limit and its fourth one diverges.

\paragraph{An explicit boundary measure that closes the Euclidean photon estimates.}
Use the two charged slots above at each center of a stationary
marked Poisson process of intensity $L^{-4}$ on the reconstructed
Euclidean plane, with independent normalized marks. For example
rotate a fixed noncollinear pair $a,b$ uniformly under $O(4)$ and
assume $|a|+|b|\le M L$. Take a nonzero constant $|\kappa|\le1/4$.
On all links use a connection $B$ whose transverse part is measured
with the isometric A4 characteristic function
$c(q)=\chi_{K-1}(L|q|)$ and whose longitudinal part is unchanged.
Thus $B$ transforms as a connection. The point process and this
per-link measurement are constitutive choices for this comparison.
In a bounded region the Poisson sum of the product of the exact
two-slot determinants gives
\begin{equation}
 -\log Z_{\rm slots}[A]=L^{-4}\int dx\,
   \mathbb E_{\rm mark}|\kappa|^2
              \left[2-2\cos\!\left(e\oint_{x+C} B\right)\right].
 \label{eq:paired-measured-gas-uv}
\end{equation}
Indeed a Poisson law of mean $m$ gives $\mathbb E\,z^N=e^{m(z-1)}$;
conditionally on $N$, centers and marks are independent. The
infinite-space translation-invariant vertices are then defined
per unit volume, stripping the overall momentum delta.

Every photon incidence in a closed vertex carries its own
$c(q)j_C(q)$, including contacts at arbitrary order. Closure
and bounded sides give $|j_C(q)|\le C_M L\min(1,L|q|)$.
Every internal photon is bounded by a multiple of
\[
 h(q)=L^2 |c(q)|^2\min(1,L|q|)^2/|q|^2.
\]
This function is bounded at zero and integrable on $\mathbb R^4$:
$c(q)=O((1+L|q|)^{-K/2})$, $K\ge5$. Bound the factors on
a spanning tree by $\|h\|_\infty$ and integrate the chord
momenta with $\|h\|_1$; a self-loop is a chord. The finite
mark integrals preserve this bound. All source-dependent
closed Euclidean photon graphs of this action are absolutely
convergent at fixed electric order. Each amputated external
leg supplies its own $O(|q|)$ soft factor, so all proper
two-point corrections are $O(|q|^2)$ near zero. This removes
the photon mass term in this Euclidean comparison. The two
charged slots do not furnish a propagating Dirac species.

The same action tests whether these bounds imply a positive
spectral reconstruction. Its leading current Hessian is
$2 e^2 L^{-4}|\kappa|^2 |c(q)|^2\,\mathbb E\,
\overline{j_C^\mu(q)} j_C^\nu(q)$. Isotropy and transversality
give the structure $(q^2\delta_{\mu\nu}-q_\mu q_\nu)p(q^2)$.
The area response makes $p$ strictly positive near positive
$q^2=0$; every positive Bessel zero of $c$ gives a zero of $p$.
A scalar with a nonzero positive spectral measure and finitely
many polynomial subtractions has a strictly signed derivative
at all $X=q^2>0$, at any sufficiently high order $n$ where
the differentiated integral converges: up to the fixed
sign its magnitude is $n!\int (s+X)^{-n-1}d\rho(s)>0$.
Repeated Rolle's theorem excludes infinitely many distinct
positive zeros. A pure polynomial term also cannot have them
unless it vanishes identically. Thus this leading current
cannot have that spectral representation. In particular
the Euclidean graph theorem for this boundary choice does
not establish interacting spectral positivity.

\paragraph{Inclusive propagation of a measured shell displacement.}
The elastic A4 amplitude can be small even while the inclusive
propagation through the shell is of Maxwell size. Work on a single
reconstructed Euclidean four-plane, after its sector-plane pullback.
Let $(Y,\pi)$ be the normalized boundary space and
$P_s=\exp(-sH)$ a conservative reversible Markov semigroup there.
Write $\mathcal E$ for its Dirichlet form and assume the real,
dimensionless readout coordinates $f_1,\ldots,f_4$ have finite
form energy. For $L>0$ set
\[
 \begin{aligned}
 v_Q(y)&=e^{iL Q\cdot f(y)}, & W(Q)u&=v_Q u,\\
 M_{ij}&=\mathcal E(f_i,f_j), & C&=\lambda_{\max}(M).
 \end{aligned}
\]
The squared norm of $v_Q$ is one. Its elastic port in the constant
state is $\int v_Q\,d\pi$. If that port is attenuated, the exact
inclusive sum over an orthonormal boundary basis still equals one
by Parseval. Keeping a dynamical transfer between the two marks
requires a stronger estimate.

For the symmetric stationary measure $\pi(dy)P_s(y,dz)$,
the squared-difference representation of the form and
$|e^{ia}-e^{ib}|\le |a-b|$ give
\[
 \begin{aligned}
 \mathcal E(v_Q,v_Q)
 &=\lim_{s\downarrow0}\frac1{2s}\int
       |v_Q(y)-v_Q(z)|^2\,\pi(dy)P_s(y,dz)\\
 &\le L^2 Q^{\mathsf T} M Q\le L^2 C|Q|^2 .
 \end{aligned}
\]
In particular the modulated state belongs to the form domain.
This argument includes conservative symmetric jump laws;
it does not require smooth diffusion paths. For $X=|Q|^2>0$
the ordinary, simultaneously propagated transverse port obeys
\begin{equation}
 \frac{1}{(1+C)X}\ \le\
  \langle v_Q,(X+L^{-2}H)^{-1}v_Q\rangle_\pi\
 \le\frac1X .
 \label{eq:fix-inclusive-Markov-uv}
\end{equation}
The spectral measure of $H$ in $v_Q$ has mass one and first
moment $\mathcal E(v_Q,v_Q)$. Jensen's inequality for the
positive convex inverse gives the lower bound. Equivalently
the inclusive heat expectation is at least
$\exp(-s\mathcal E(v_Q,v_Q))$ at dimensionless time $s$.
The region $s\asymp(L^2X)^{-1}$ therefore retains finite weight
at arbitrarily hard momentum.

\paragraph{Every simultaneous short-duration face in a reflected history.}
Let $h\ge1$ be fixed and $K_i=K_0+\sum_{j=1}^{i}Q_j$ the cumulative
momenta on successive arcs of one history; the $Q_j$ can have
arbitrary relative sizes. For $\mu>0$ set $d_i=\mu^2+|K_i|^2$.
Sew this history to its adjoint using the complete boundary
resolution after the last arc. The resulting squared norm is
$\|\Psi(t)\|^2$, where
\[
 A_i=e^{-t_i H/(2L^2)},\qquad
 \Psi(t)=A_h W(Q_h)\cdots A_1 W(Q_1)v_{K_0}.
\]
Let $B_h=\int_{(0,\infty)^h}\|\Psi(t)\|^2
\prod_i e^{-t_i d_i}\,dt_i$. The $t_i$ here are the ordinary
positive duration variables of the comparison action; $\mu$ can
be the intrinsic inverse length $L^{-1}$. For
$\delta=[2h^2(1+C)]^{-1}$, one has the uniform two-sided bound
\begin{equation}
 \frac{(1-e^{-\delta})^h}{4}\prod_{i=1}^h d_i^{-1}
 \ \le B_h\le\prod_{i=1}^h d_i^{-1}.
 \label{eq:fix-inclusive-Markov-history-uv}
\end{equation}
Indeed all factors are contractions, and with every $A_i$ replaced
by $I$ each prefix is exactly $v_{K_i}$. Telescoping and the
spectral inequality $(1-e^{-z})^2\le z$ for $z\ge0$ give
\[
 \|\Psi(t)-v_{K_h}\|
 \le \sum_i\|(A_i-I)v_{K_i}\|
 \le \sum_i\sqrt{t_i C|K_i|^2/2}.
\]
Restricting independently to $0\le t_i\le\delta/d_i$ bounds the
last sum by $1/2$. Integration on this box proves the lower
bound; the upper bound follows from $\|\Psi(t)\|\le1$.
The constants do not depend on the number of independent hard
scales present at the fixed order $h$, or on cancellations in the
$K_i$. An exactly coincident mark has $A_i=I$ and contributes
no duration integral; the other arcs obey the same argument.
Thus a positive reflected history of this form has precisely
the same powers of all its arc denominators as one with no
shell propagation. A signed sum of distinct histories, a
different photon port, and a source contact with a different
vertex require their own sewing identities.

\paragraph{Sharp coefficient for the original ball measure.}
On the unit A4 ball in $n=K-1\ge4$ transverse dimensions let
$\pi$ be normalized volume, $H=-\Delta_N$, and let $T$ isometrically
embed the reconstructed four-plane into $\mathbb R^n$. Take
$Q\cdot f(y)=(TQ)\cdot y$. The form identity is now exact,
$\mathcal E(v_Q,v_Q)=L^2 X$, whereas the elastic probability is
$|\chi_n(L|Q|)|^2$. In fact
\begin{equation}
 X\langle v_Q,(X+L^{-2}H)^{-1}v_Q\rangle_\pi
 \longrightarrow \tfrac12\quad(|Q|\to\infty)
 \label{eq:fix-inclusive-ball-uv}
\end{equation}
uniformly in direction. For a smooth interior cutoff $\phi$,
$g_Q=\phi v_Q$ is in the Neumann operator domain and
$\|(L^{-2}H-X)g_Q\|\le C_\phi(1+\sqrt X)$ by differentiation.
The resolvent identity now gives
$\|X(X+L^{-2}H)^{-1}g_Q-g_Q/2\|=O_\phi(X^{-1/2})$.
Since $\|X(X+L^{-2}H)^{-1}\|\le1$, interior cutoffs tending
to one in $L^2(\pi)$ extend the result to $v_Q$; no boundary
condition on the modulated state itself was imposed.

For the closed thin parallelogram of nonzero area, the measured
transverse current is $j_C(Q)v_Q(y)$. It is transverse for
each $y$ since $Q\cdot j_C=0$. Using the inclusive covariance
in \eqref{eq:fix-inclusive-ball-uv} in place of $1/X$ in
its Maxwell variance gives
$P(C)R/(4\pi^2)+o(R)$ under spherical exhaustion. The estimate
uses the uniform limit and the perimeter angular law above.
All intermediate normal excitations are included. This explicit
coefficient detects a divergence of the gauge-invariant
closed current already at second order for this sewn source.
It shows precisely why an exclusively elastic measurement,
or an independent coherent shell average at each incidence,
has different ultraviolet behavior.

\paragraph{Positive boundary dynamics and the physical photon spectrum.}
There is a spectral test even for boundary generators whose readout
coordinates have infinite energy. Suppose the measured transverse
scalar $D(X)$ along a Euclidean ray is to have the usual covariant
positive spectral representation with finitely many polynomial
subtractions. Suppose also $0\le D(X)\le1/X$ for all $X>0$ and
$\lim_{X\downarrow0}XD(X)=1$. Then $D(X)=1/X$ identically.
To see this without assuming a finite spectral mass in advance,
take a difference of order $N$ large enough to remove all
subtraction polynomials, with step $X$. The spectral formula gives
\[
 X|\Delta_X^N D(X)|
 =\int\frac{N!X^{N+1}}{\prod_{j=0}^N[s+(j+1)X]}\,d\rho(s).
\]
The left side is bounded by $\sum_{j=0}^N\binom Nj/(j+1)$.
Fatou's lemma as $X\to\infty$ shows that $\rho$ has finite total
mass. Its ordinary Stieltjes transform therefore exists and has
the same sufficiently high derivative as $D$, so their difference
is a polynomial. Both tend to zero at infinity and the polynomial
vanishes. Taking the soft limit under the finite measure gives
$\rho(\{0\})=1$. This pole alone saturates the upper bound $1/X$,
leaving no positive massive weight.

For the inclusive propagator in \eqref{eq:fix-inclusive-Markov-uv}
the upper bound follows just from $H\ge0$ and the unit phase.
At small $Q$, $v_Q\to1$ in $L^2(\pi)$, and $H1=0$ gives
$XD(X)\to1$. At any $Q$ for which $v_Q$ has spectral weight
outside $\ker H$, the estimate $XD(X)\le1$ is strict. Thus
the measured scalar at such a ray fails this positive covariant
spectral test. The full boundary propagation itself is positive;
allowing the physical readout to vary with the Euclidean time
frequency is precisely the extra ingredient tested here. For
the ball its hard pole coefficient is exactly one half by
\eqref{eq:fix-inclusive-ball-uv}.

\subsection{A fixed Newton constant and a fluctuating geometric area}
\label{sec:path-newton-tests-uv}
A delta function can fix a coupling, or it can tie a duration to the
geometry of a particular history. These operations give different
results when several momenta become hard at once.

\paragraph{A gravitational scale constraint and a common geometric record.}
Let $g>0$ have the units of Newton's constant. At the level of a
source functional, with the delta normalized in $dg$,
\[
 \int_0^\infty dg\,\delta(g-G_N)\,Z_g[A]=Z_{G_N}[A].
\]
The same identity holds for its marked derivatives. This insertion
also leaves the tensor type of the action unchanged. A constraint
$C(g,q)=0$ with a simple positive root instead contributes the
coarea factor $|\partial_g C|^{-1}$ at that root. If the root or
the weight depends on the source, its derivatives belong to the
current and contact terms.

A positive distribution over a \emph{common} $g$ requires a
different test. On normalized sectors write $Z_g[0]=1$,
$W_g[A]=-\log Z_g[A]$, and set
$W[A]=-\log\int e^{-W_g[A]}\,d\mu(g)$ with $\mu$ a probability
measure independent of $A$. Assume vanishing one-point currents
and the finite, differentiable tested moments under consideration.
Use $H_{g,ij}$ and $V_{g,1234}$ for the second and fourth
derivatives of $W_g$ at zero on the indicated probe directions.
Expansion gives the complete derivatives
\begin{equation}
\begin{aligned}
 H_{ij}&=\mathbb E_\mu H_{g,ij},\\
 V_{1234}&=\mathbb E_\mu V_{g,1234}
   -\sum_{(ij)(kl)}
       \operatorname{Cov}_\mu(H_{g,ij},H_{g,kl}).
\end{aligned}
\label{eq:path-g-annealed-uv}
\end{equation}
The sum runs over the three pairings and includes the contacts
within each conditional derivative. For one real probe the extra
term is $-3\operatorname{Var}_\mu(H_g)$. Thus a nontrivial global
scale can change a connected four-current coefficient. Its new
term couples two responses through their shared value of $g$.
For a diagonal ensemble close two polarization insertions with two
ordinary Maxwell lines. With one charged loop per polarization,
the microscopic vacuum graph has three independent loop momenta:
two inside the polarizations and one through their connecting
photons. For leading polarizations of order $e^2$ it occurs at
order $e^4$. For the responses actually inserted the nonnegative
ring integral is proportional to
\[
 \int d\mu(g)\int_{X_0}^{\infty}X|\Pi_g(X)|^2\,dX .
\]
Tonelli's theorem makes this infinite if the fixed-$g$ ring
diverges on a set of positive $\mu$-measure. If independent
constant labels are averaged on the two loops, the test instead
uses $|\mathbb E_\mu\Pi_g|^2$; the finite-band ring bound applies
when the averaged response obeys its stated positive dispersion
hypotheses. A spectral bound on a fully interacting response only
implies an obstruction at a specified bare order if that coefficient
itself obeys the bound. A covariance term in a different topology
does not supply graphwise absolute convergence. Likewise the vacuum norm
of the density in a diagonal mixture is the mean of its
conditional squared norms. If every conditional density
annihilates the vacuum whereas a spatial current creates a pair
in a positive-measure set of sectors, that boost obstruction
survives. An operator changing $g$ would require off-diagonal
matrix elements and a new action.

\paragraph{Testing a length delta built from two forward displacements.}
One can instead constrain a length by the actual history. For an
explicit comparison on a fixed reconstructed plane of signature
$(+---)$ let $L^2=\hbar G_N/c^3$ in ordinary units, with $c$ here
the scalar magnitude of the light speed. Use carrier length units
for all four displacement components. Choose two temporal steps of
length $L$ in this readout and write
\[
 a=L(1,Y_1),\qquad b=L(1,Y_2).
\]
The dimensionless $Y_i$ are independent projections onto three
orthonormal axes of the normalized A4 ball $B^n$, $n=K-1\ge3$.
This fixes a boundary preparation on that readout. Both segments
are future timelike almost surely. Their dimensionless Lorentz
Gram invariant and a possible length constraint are
\begin{equation}
\begin{aligned}
 U&=\frac{(a\cdot b)^2-a^2b^2}{L^4}
   =|Y_2-Y_1|^2-|Y_1\times Y_2|^2,\\
 &\hspace{1cm}d(\ell^2)\,\delta(\ell^2-L^2U).
\end{aligned}
\label{eq:path-g-area-uv}
\end{equation}
The invariant uses only the two projected displacements; the
chosen equal steps still distinguish a frame. The small-area
tail can be computed even near the lightlike edge. Each $Y_i$
has density proportional to $(1-|Y_i|^2)^\alpha$ on $B^3$,
$\alpha=(n-3)/2$. Set $e=(Y_2-Y_1)/|Y_2-Y_1|$ away from the
diagonal and write $Y_1=z+se$, $Y_2=z+te$, $z\perp e$, $s<t$.
With $w=\sqrt{1-|z|^2}$, one has $|s|,|t|<w$ and
$U=w^2(t-s)^2$. Polar integration over $Y_2-Y_1$ gives volume
element $(t-s)^2\,d\Omega_e\,d^2z\,ds\,dt$. Setting
$s=wu$, $t=wv$, the negative moment $\mathbb E U^{-p}$ is,
up to a finite positive constant, the product of
\[
 \begin{gathered}
 \int_0^1 w^{\,2n-1-4p}\,dw,\\
 \int_{-1<u<v<1}(v-u)^{2-2p}
        (1-u^2)^\alpha(1-v^2)^\alpha\,du\,dv .
 \end{gathered}
\]
The second integral diverges at the interior diagonal when
$p\ge3/2$. For $p<3/2$ both are finite, since $n\ge3$ and
$\alpha\ge0$. Thus $U>0$ almost surely, with finite negative
moments precisely for $p<3/2$.

The delta in \eqref{eq:path-g-area-uv} changes a marked graph
only if $\ell$ enters its kernel. For a heat comparison with
one common area, take $X_j=|Q_j|^2$ in Euclidean units,
$p_j>0$ and $P=\sum_jp_j$. Direct integration yields
\begin{equation}
\begin{aligned}
 &\int_{(0,\infty)^k}\!\prod_j(X_j^{p_j-1}\,dX_j)\,
       \mathbb E e^{-L^2U\sum_j X_j}\\
 &\qquad=L^{-2P}\Big(\prod_j\Gamma(p_j)\Big)
                                      \mathbb E U^{-P}.
\end{aligned}
\label{eq:path-g-joint-uv}
\end{equation}
Each marginal with $p_j=1$ is finite. Two such independent hard
momenta sharing the same area give a divergent integral. Giving
each its own independent area would replace the last moment by
the product of the marginal moments. A positive weight $U^h$,
$h\ge0$, in the common preparation changes the criterion to
$P-h<3/2$. These tests determine what an area-dependent heat
factor would accomplish. The actual Wilson mark has its own
oscillatory line integrals, and its joint law and hard bound must
be obtained from the charged action.

\subsection{When the measured displacements span the hard loops}
\label{sec:fix-record-rank-uv}

\paragraph{When geometric records control every independent hard momentum.}
There is a sharp positive-kernel test that allows geometry to be shared
between transfers. Fix a finite connected graph with edge set $E$. Draw
independent positive dimensionless records $u_r$ at the elementary marks
of the proposed source, before choosing any graph momentum coordinates.
Each $u_r$ is a sum of $h_r\ge1$ independent copies of the two-displacement
area $U$ above. Give edge $e$ the positive clock
$s_e=\sum_{r\in I_e}u_r$, where the nonempty set $I_e$ specifies the
marks on which that transfer depends. The same mark may enter many edges.
Use dimensionless Euclidean momentum units, and in this comparison give
edge $e$ the factor $e^{-s_e|q_e|^2}$. For a set $S$ of records put
$E(S)=\{e:I_e\subseteq S\}$. For any edge set $F$, denote the dimension
of the cycle space of its spanning subgraph by $b_1(F)$, including
parallel edges and self-loops. The integral of the expected product over
all loop momenta at fixed finite external momenta is finite precisely if
\begin{equation}
 \sum_{r\in S}\kappa_r>2b_1(E(S)),\qquad
 \kappa_r=3h_r/2 ,
 \label{eq:geo-record-subset-uv}
\end{equation}
for every nonempty set $S$ of records. Thus only cycles wholly supported
on the transfers that can shrink together cost independent four-momentum
volumes. Bridges may carry records but introduce no additional loop.

Here is a proof that includes overlapping hard regions. Write the edge
momenta as $q=Ck+p$ in a fundamental integral cycle basis and put
$s=\operatorname{diag}(s_e)$. Gaussian integration gives a constant times
$\det(C^{\mathsf T}sC)^{-2}$ times an exponential with exponent
$-\min_k\sum_e s_e|q_e|^2$. This exponential lies between a fixed
positive constant and one because the clocks have bounded support and
$p$ is fixed. By Cauchy--Binet,
\[
 \det(C^{\mathsf T}s C)=\sum_T\prod_{e\notin T}s_e ,
\]
where $T$ ranges over spanning trees. Indeed a maximal row minor
with complement $T$ is unimodular when $T$ is a tree and vanishes
otherwise. Normalize the records by a fixed constant to be at most one,
and bin $u_r\asymp 2^{-k_r}$ with $k_r\ge0$. In these bins
$s_e\asymp 2^{-k_e}$ with $k_e=\min_{r\in I_e} k_r$. Set
$S_j=\{r:k_r\ge j\}$. A maximum-weight spanning tree built from large
$k_e$ first contains a maximal forest in each $E(S_j)$. Its omitted
edges there number $b_1(E(S_j))$. Consequently the determinant is
comparable, with constants independent of the bins, to
$2^{-\sum_{j\ge1} b_1(E(S_j))}$.

The area law has finite negative moments of every order below $3/2$.
Independence and the arithmetic--geometric mean inequality yield finite
negative moments of $u_r$ at every order below $\kappa_r$.
Each bin therefore has probability at most
$C_\eta 2^{-(\kappa_r-\eta)k_r}$ for $0<\eta<\kappa_r$.
There is also a lower bound $c 2^{-\kappa_r k_r}$ for all sufficiently
large bins. For a single area, restrict both projected velocities to an
interior ball and their relative displacement to an annulus of radius
comparable to $2^{-k_r/2}$. In the interior the area is comparable to
the squared relative displacement, and a subannulus of fixed relative
width can be chosen inside the target area bin. For $h_r$ areas,
restrict each to a subinterval of the target bin divided by $h_r$.
The joint lower bound follows by independence.

Under \eqref{eq:geo-record-subset-uv}, the finite set of strict
inequalities permits an $\eta>0$ leaving
$\sum_{r\in S}(\kappa_r-2\eta)>2b_1(E(S))$ for every nonempty $S$.
After the determinant bound, the sum of bin contributions is thus
bounded by a constant times $\sum_{(k_r)}2^{-\eta\sum_r k_r}$.
Conversely, put all records in a failing set in bin $k$ and the
others in fixed interior bins. These disjoint regions contribute at
least a constant times $2^{[2b_1(E(S))-\sum_{r\in S}\kappa_r]k}$.
This proves divergence also at equality.

For one independent area per edge, four parallel edges between two
vertices give equality: their three independent loop momenta cost six
powers of the clock, exactly the total small-clock exponent. For two
independent areas per edge, $\kappa_e=3$ and $b_1(F)\le |F|$, so every
finite graph passes this Gaussian test. A single common vertex clock
can instead fail at high valence: $m$ self-loops at that vertex cost
$2m$ powers even if their number of distinct geometric records stays
fixed. Polynomial absolute bounds can also be tested. A factor
$(1+|q_e|)^{d_e}$ costs at most a constant times $s_e^{-d_e/2}$ after
halving its heat exponent. A sufficient strict inequality replaces the
right side of \eqref{eq:geo-record-subset-uv} by
$2b_1(E(S))+\tfrac12\sum_{e\in E(S)}d_e$. Applying this to charged
vertices requires the marked determinant and its physical contacts to
have such a bound. The area delta by itself produces no heat factor.

\paragraph{Realizing the rank test with forward A4 displacements.}
The required momentum bounds can also follow from actual characteristic
functions. Here is a boundary preparation on one fixed readout, with
no Gaussian inserted in momentum space. Take $n=K-1\ge3$ and two
independent normalized A4 ball records $y,y'$. Let $Y$ be the
three-axis spatial projection of $y$ and $t=|y'|$. The radial law is
$n t^{n-1}\,dt$ on $(0,1)$. In units of the geometric length choose
the four-displacement $D=(t,tY)$, which is future timelike almost
surely. Its four-dimensional density is
\[
 f(t,\mathbf x)=n t^{n-4} f_Y(\mathbf x/t),\qquad 0<t<1,
 \qquad \|f\|_2^2=\frac{n^2}{2n-4}\|f_Y\|_2^2 .
\]
Here $f_Y$ is the three-coordinate marginal on the unit ball, bounded
for $n\ge3$. This preparation identifies the duration with a radial
shell record and includes arbitrarily short durations. It uses the
chosen readout frame. Two independent such forward displacements,
paired with opposite Fourier phases, give
$F(q)=|\widehat f(q)|^2$. By Plancherel $F\in L^1(\mathbb R^4)$,
$0\le F\le1$, and $F$ is bounded below by a positive constant near
zero. The difference of the two Fourier marks need not be a
future-directed displacement.

Let $b$ be the number of independent graph momenta $k$. Allow a finite
family of recorded phase directions $c_i\in\mathbb R^b$, determined
by the actual mark incidences and a loop basis. Each mark samples an
independent forward pair as above with duration rescaled by
$\sqrt{u_{r(i)}}$, conditional on the shared geometric areas. Thus
its coefficient is exactly $F(\sqrt{u_{r(i)}}(c_i\cdot k+p_i))$,
where $p_i$ is the routed external momentum. Define
\[
 M=\sum_i u_{r(i)} c_i c_i^{\mathsf T},\quad
 d(S)=b-\operatorname{rank}\{c_i:r(i)\notin S\}.
\]
If the full direction family fails to span, an entire loop direction
is invisible. At zero external momentum its integral is infinite.
For a spanning family and $p_i=0$, the conditional loop integral
of the product of the $F$'s is bounded above and below by positive
constants times $(\det M)^{-2}$, uniformly in the areas. For the
upper bound choose a basis $J$ maximizing
$(\det[c_i]_{i\in J})^2\prod_{i\in J}u_{r(i)}$. Cauchy--Binet makes
this comparable to $\det M$. Keep these factors, bound the others
by one, and use the chosen linear variables to integrate $b$ copies
of $F$. This upper bound allows arbitrary fixed $p_i$. For the
lower bound use $\sum_i u_{r(i)}|c_i\cdot k|^2\le\delta^2$ with
$\delta$ small. Every factor there is bounded below, and the
ellipsoid volume is a constant times $(\det M)^{-2}$.

The dyadic proof of \eqref{eq:geo-record-subset-uv} now applies to
maximal independent subsets of the $c_i$. A basis chosen from largest
clocks first uses precisely $d(S_j)$ directions from the records in
$S_j$. This gives the sharp criterion
\begin{equation}
 \sum_{r\in S}\kappa_r>2d(S)\quad
 \text{for every nonempty record set } S .
 \label{eq:geo-actual-phase-rank-uv}
\end{equation}
Divergence at a failing set persists at finite fixed external momenta:
choose the outside clocks in fixed bins small enough that
$\sqrt{u_{r(i)}}|p_i|<\delta$ and let the failing clocks shrink.
The same ellipsoid supplies the bin lower bounds. In the edge-clock
example there is a separate forward pair for each incidence $(e,r)$,
all pairs with record $r$ sharing only its area. The rank deficiency
is then exactly $b_1(E(S))$. In particular two independent areas per
edge, with one forward pair for each area, control all loop momenta
of the bounded-vertex comparison by genuine geometric phases.
This does not identify those independently sampled marks in the
Taylor derivatives of a Wilson holonomy or a charged determinant.

A common vector displacement is more restrictive than a common
scalar area. For example, attach one displacement $z_v$ to each
vertex and give each oriented edge $e:v\to w$ the phase of
$z_w-z_v$. The sum $\sum_e q_e\cdot(z_w-z_v)$ depends only on the
external injections, by momentum conservation at the vertices.
All loop directions disappear identically, with any law for the
$z_v$. This is a direct test for a proposed native identification
of two histories: the retained relative marks must survive the
cycle projection after endpoint matching and all contacts.

\paragraph{Shared velocities and the short-duration faces.}
Independence of areas on adjacent native steps can itself be tested.
Give the vertices of a geometry graph $J$ independent three-coordinate
A4 velocity records $Y_i$. Each momentum-graph edge $e$ is assigned
one pair $i(e)\ne j(e)$ in $J$ and the clock
$s_e=U(Y_{i(e)},Y_{j(e)})$. Multiple edges may use the same pair.
More generally let $s_e$ be a sum of $h$ such areas from $h$
independent copies of the vertex records. Edges meeting at a vertex
have correlated clocks even within one copy. Write $r_J(H)$ for the
size of a maximal forest of the geometry edges indexed by $H$, and
$b_1(H)$ for the cycle dimension in the momentum graph. For $n=K-1\ge5$
both the Gaussian comparison and the forward-pair comparison with one
independent conditional pair per edge are finite precisely when
\begin{equation}
 (3h/2)\,r_J(H)>2b_1(H)
 \quad\text{for all nonempty sets of transfers }H .
 \label{eq:geo-velocity-forest-uv}
\end{equation}
For $n=3,4$ this condition is necessary, and replacing $3h/2$ by
$h(n+1)/4$ is sufficient.

The velocity-shell rim matters for the upper bound. Fix $y=t e_1$
by spatial rotation and put $Y=(z,w)$ with $w\in\mathbb R^2$ and
$a=1-t^2$. Directly, $U(y,Y)=(z-t)^2+a|w|^2$. Write
$\alpha=(n-3)/2$. On $U\le\delta^2$, $|z-t|\le\delta$ and
$1-z^2\le C(a+\delta)$. Integrating the bounded marginal density
$(1-z^2-|w|^2)^\alpha$ over the allowed disk gives
\[
 \Pr(U(y,Y)\le\delta^2)
 \le C\delta(a+\delta)^\alpha
          \min\{\delta^2/a,\,a+\delta\}
 \le C'\delta^{\min(3,\alpha+2)} .
\]
At $a=0$ use the second disk bound. The last inequality follows by
splitting at $a=\delta$ and holds uniformly in the first velocity.
Set $\kappa=\min(3/2,(n+1)/4)$. On a bin $s_e\asymp 2^{-k_e}$
each constituent area is at most $C2^{-k_e}$. Select a
maximum-weight spanning forest of $J$ with edge weights $k_e$.
Integrating its leaves successively and using independence between
copies bounds the bin probability by
$C'2^{-h\kappa\sum_{j\ge1}r_J(H_j)}$,
$H_j=\{e:k_e\ge j\}$. The momentum determinant contributes at
most $C''2^{2\sum_{j\ge1}b_1(H_j)}$. Strict margins over the finite
set of edge subsets yield a convergent sum of bins.

For necessity restrict to an interior patch where $U(y,y')$ is
comparable to $|y-y'|^2$. In each component of the geometry subgraph
on $H$, keep one velocity free in this patch and put all the others
in balls of radius $c\delta$ at offsets of order $\delta$ from it.
This costs at most $3h\,r_J(H)$ powers of $\delta$ in the sense
that the probability is bounded below by $c'\delta^{3h r_J(H)}$.
The offsets may be chosen so one specified transfer in $H$ has
$s_e\asymp\delta^2$, giving disjoint events along a geometric
subsequence. Every clock in $H$ is $O(\delta^2)$; every cotree
contains at least $b_1(H)$ of these edges, so the inverse squared
determinant is bounded below by a constant times $\delta^{-4b_1(H)}$.
A small initial patch also keeps all external shifts small in the
forward-pair lower bound. This proves divergence at the stated
threshold, including equality.

Two parallel momentum transfers using the same two velocities
already fail at $h=1$: $r_J=1$ and $b_1=1$. With $m$ parallel
transfers on the same pair $b_1=m-1$, so no fixed number of copies
controls arbitrarily high contact multiplicity in this comparison.
Assigning independent velocity pairs to every transfer makes the
geometry graph a matching and gives $r_J(H)=|H|$. These choices
give different predictions with the same one-step A4 marginal.

For independent velocity pairs the sharp forest threshold also
extends to the two lowest dimensions in this comparison, $n=3,4$.
To see the effect of integrating the first velocity as well as
the second, write $a=1-|Y_i|^2$. Its density on $(0,1)$ is
proportional to $a^\alpha(1-a)^{1/2}$, with
$\alpha=(n-3)/2$. For $a\ge\delta$ the conditional estimate above
is at most $C\delta^3 a^{\alpha-1}$; for $a<\delta$ it is at
most $C\delta^{\alpha+2}$. Thus for $0<\delta<1/2$,
\[
 \Pr\{U(Y_i,Y_j)\le\delta^2\}\le
 \begin{cases}
 C\delta^3[1+\log(1/\delta)],& n=3,\\
 C\delta^3,& n=4 .
 \end{cases}
\]
Indeed the first interval costs at most
$C\delta^3\int_\delta^1 a^{2\alpha-1}da$, and the second costs
$C\delta^{2\alpha+3}$. For a geometry matching the clocks on
distinct pairs are independent, so at dyadic depth $k_e$ the
upper bound for $h$ copies has the same exponential power
$2^{-3h k_e/2}$ in both dimensions. The extra polynomial factor
in $k_e$ when $n=3$ is summable under the strict subset margins.
The preceding interior lower bound still proves necessity, including
at equality. Sharing a velocity among several pairs makes the
conditional rim cost recur at that vertex and requires a joint test.

\paragraph{Subdivision and the arbitrarily short physical transfer.}
The distinction between a mark at a physical interaction and an
artificial subdivision has a quantitative consequence. For two adjacent
forward segments with common temporal step $h$, the dimensional clock
$h^2 U$ is bounded by $4h^2$, since $U\le|Y_1-Y_2|^2\le4$. For $N$
such cells with $h=t/N$ the sum of all clocks is at most $4t^2/N$,
regardless of correlations. Their product of heat factors at any fixed
momentum tends to one. Retaining a fixed number of independent areas
per cell only multiplies this bound by that number. Assigning a fresh,
fixed-width record to every artificial split is also inconsistent with
an additive, stationary, independently divisible clock of bounded
nonzero variance. If its unit-time law has support width $D$, each of
its $N$ independent identically distributed increments has support
width $D/N$. Variance additivity bounds the unit-time variance by
$D^2/(4N)$ for every $N$.

Even an unbounded clock with a genuine positive convolution semigroup
faces the duration integral. In dimensionless units let
$\mathbb E e^{-x S_t}=e^{-t\psi(x)}$ for an additive positive clock
with $\psi(0)=0$. The negative log Laplace transform $\psi$ is concave
by H\"older's inequality. Hence $\psi(x)\le x\psi(1)$ for $x\ge1$.
With the usual unweighted resolvent duration, $m>0$, the averaged
transfer is exactly
\begin{equation}
 \int_0^\infty e^{-mt}\mathbb E e^{-xS_t}\,dt
  =\frac{1}{m+\psi(x)}\ \ge\ \frac{1}{m+x\psi(1)},\qquad x\ge1 .
 \label{eq:geo-clock-resolvent-uv}
\end{equation}
For $x=|Q|^2$ the four-dimensional integral of the square of this
positive transfer diverges at least logarithmically. It is the
two-line, constant-vertex comparison. Native numerators, gauge
cancellations and a different boundary law must be tested on their own
source. The subset criterion thus identifies a sufficient incidence
structure and the clock calculation identifies exactly what is lost
if the geometry is assigned to arbitrary time slices or integrated
against an unsuppressed short-duration measure.

\paragraph{Testing the charged operator and its measure in one calculation.}
A boundary law that depends on the external field can change both
comparisons. On a regular chart with a normalized reference measure
$d\mu(\xi)$, write the relative source weight as $e^{-L_\xi[A]}$,
where
\[
 L_\xi[A]=-\log m_\xi[A]
   -\operatorname{Tr}\log(D_\xi[A]D_\xi[0]^{-1}).
\]
Here $m_\xi[0]=1$ includes the boundary density and all Jacobians
of the conditional measure. This identity is applied on finite
tests, or when the displayed marked traces and differentiated
averages exist. Write $V_i=\partial_iD_\xi|_0$,
$V_{ij}=\partial_i\partial_jD_\xi|_0$, $S_\xi=D_\xi[0]^{-1}$,
and $L_i,L_{ij}$ for its relative logarithmic derivatives. For
$W=-\log\mathbb E e^{-L}$ and the normalized open inverse
$G=\mathbb E(e^{-L}D^{-1})/\mathbb E e^{-L}$ one obtains
\begin{equation}
\begin{aligned}
 W_{ij}&=\mathbb E L_{ij}
                  -\operatorname{Cov}(L_i,L_j),\\
 L_{ij}&=\operatorname{Tr}
       (S_\xi V_iS_\xi V_j-S_\xi V_{ij})
                  -\partial_{ij}\log m_\xi|_0,\\
 G_i&=-\mathbb E(S_\xi V_iS_\xi)
                  -\operatorname{Cov}(S_\xi,L_i).
\end{aligned}
\label{eq:fix-native-derivatives-uv}
\end{equation}
These are bilinear connected covariances, taken componentwise for
an open matrix element. They follow by differentiating the single
normalized integral. In particular a nonconstant boundary score
$L_i$ can affect the open pole amplitude as well as the two-current
contact. A factor depending on $A$ alone has zero covariance with
$S_\xi$ and changes only the closed derivative here. Independently
prepared endpoint spinors, if field dependent, have their own
derivatives in $G_i$. In the proposed GCM derivation the data to
compute are $V_i$, $V_{ij}$, those endpoint derivatives and the
boundary derivatives on the same history, prior to any decay
estimate or use of an optical inequality.

\subsection{Deriving the charged propagator and its contacts on one GCM chart}
\label{sec:gcmone-transport-uv}
There are two distinct questions about the curvature principle.
Given a representation, what charged source does its integral produce?
Which representation, if any, does the unrestricted physical principle
select? For the first question we construct the charged boundary map,
integrate its auxiliary constraints and differentiate the result.
The two measurable outputs are a pole-pair amplitude with specified
external states and the discontinuity of the closed determinant. A
positive Hilbert-space interpretation requires the latter to contain
at least the former's squared norm, including every photon contact
at the same order.

\paragraph{A velocity-shell exit law and its charged transport}

One can obtain the normal multiplier in
\eqref{eq:gcmone-combined-action-uv} by transporting a norm.
Fix the same aligned sector plane $\bin=(s,b_s)$,
$L=\sqrt{G_N}$ in the stated carrier units, $m=L^{-1}$ and $a=m^2$.
Write $y$ for the transverse unit-ball coordinate of the full
oriented A4 velocity shell in $n=K-1\ge4$ dimensions. In the
selected rank-four subspace form $\omega=P_4y/|P_4y|$ off the
null set $P_4y=0$. Its direction is uniform on $S^3$ by rotation
invariance, including when $n=4$. Choose a rank-two orthogonal
projector $P_2$ inside this four-space. Follow the ray
$\lambda\mapsto\lambda P_2\omega$ until it exits the disk of
radius $a$ in this squared-mass coordinate. Its exit value is
$V=a/|P_2\omega|$. Normalizing four independent standard Gaussian
coordinates gives a uniform direction. The two complementary sums
of two squares are independent chi-squares with two degrees of
freedom, hence their ratio to their sum gives
$|P_2\omega|^2\sim{\rm Uniform}(0,1)$. In particular
\[
 \mathbb P(V>v)=b_a(v)=\begin{cases}1,&0\le v\le a,\\
                   a^2/v^2,&v>a.\end{cases}
\]
The rank choices and the exit readout specify a history representation
of this chart. The GTL insertion, shell measure and A3/A4 directed
lifts retain their prescribed roles. In particular the geometric carrier is
$\sigma_\epsilon\hat p_{0,\bin}$, the matter carrier is
$\sigma_s\sigma_\epsilon\hat p_{0,\bin}$, and the two matter boundary
values use $\pm i\sigma_{s\epsilon}\varepsilon\hat p_{0,\bin}$ with
$\varepsilon\downarrow0$. The exit parameter is an internal normal
coordinate. No sign of a geometric-time pole is assigned to it.

Let $K_B$ include the two ordered Wilson translations of
\eqref{eq:rect-polar-uv}, and write $G=\gamma_5 K_B$,
$H=K_B^\dagger K_B$. On the real Euclidean connection $G=G^\dagger$:
the Dirac differential is skew and anticommutes with $\gamma_5$,
and the spin-area term is Hermitian and commutes with $\gamma_5$.
Thus $H=G^2$. Initially work on a finite test space with $H\ge hI$
for $h>0$, using a gauge-covariant realization of these relations.
On each shell record give the charged coefficient space the positive
metric $M_\lambda=H+\min(\lambda,V)I$. A parallel section
$w'=-C_\lambda w$ preserves $w^\dagger M_\lambda w$ precisely
when $C_\lambda^\dagger M_\lambda+M_\lambda C_\lambda=M_\lambda'$.
Requiring $C_\lambda$ to be self-adjoint for this metric fixes it
uniquely almost everywhere as $M_\lambda^{-1}M_\lambda'/2$.
The metric is continuous at the exit, so parallel transport is
continuous there. Since the metrics commute, its exact step from
$u$ to $v$ is the positive square root of $M_u M_v^{-1}$.
The self-adjoint choice excludes a freely assigned unitary rotation.

For a partition $0=\lambda_0<\cdots<\lambda_N=R$ put
$A_j=\mathbb E[(M_{\lambda_{j-1}}M_{\lambda_j}^{-1})^{1/2}]$.
All shell records act on the same Berezin fields. The one shell
integral of the action therefore averages each step before the
fields are integrated. Write $J=\operatorname{sgn}G$,
$\chi_0=\Psi$, $d_j=\chi_j-A_j\chi_{j-1}$ and
$\bar d_j=\bar\chi_j-\bar\chi_{j-1}A_j$. A finite charged
action with these geometric step maps is
\begin{equation}
 \sum_{j=1}^N(\bar\eta_j d_j+\bar d_j\eta_j)
 +\tfrac m2(\bar\Psi J\chi_N+\bar\chi_N J\Psi).
 \label{eq:gcmone-exit-chain-uv}
\end{equation}
Each auxiliary pair has its independent Berezin measure.
At fixed external spinors the linear map from the $\chi_j$ to
the $d_j$ is lower triangular with identity diagonal. The adjoint
chain gives its dual. Their combined determinant is a constant
sign, fixed by the common measure normalization. Integration over
the multipliers implements $d_j=\bar d_j=0$ without a
field-dependent determinant. Since $J$ and all $A_j$ commute,
the retained bilinear has coefficient $m J A_N\cdots A_1$ on this real form.
The fixed dual change by $\gamma_5$ gives the charged operator
$m K_B H^{-1/2} A_N\cdots A_1$.
The block rows, their adjoints, and the endpoint pairing in
\eqref{eq:gcmone-exit-chain-uv} are the actual coefficients to
represent in the four-arrow curvature identity
\eqref{eq:g5-arrow-product}. Each row coefficient is used there as an even
endpoint map paired with its two Grassmann entries. The links,
metric transport and polar endpoint $J$ specify those maps.
The construction assumes this represented functional calculus and
the indicated mixed curvature pairing; the physical selection of this chart is open.

The limit can be estimated uniformly in the spectral variable
$z\ge h$. For a step $I_j$ write
$x_j=\tfrac12\int_{I_j}\mathbf1_{v<V}(z+v)^{-1}\,dv$, so its
averaged eigenvalue is $\mathbb E e^{-x_j}$. Jensen's inequality
and $e^{-x}\le1-x+x^2/2$ for $x\ge0$ imply
\[
 0\le\log\mathbb E e^{-x_j}+\mathbb E x_j
       \le\tfrac12\mathbb E x_j^2 .
\]
Writing $|\pi|$ for the maximum step, Cauchy--Schwarz bounds the
sum of the right sides by
$(|\pi|/8)\int_0^R(h+v)^{-2}\,dv\le |\pi|/(8h)$.
For $R\ge a$, the omitted tail of
$\Phi_a(z)=\tfrac12\int_0^\infty b_a(v)/(z+v)\,dv$
is at most $a^2/(4R^2)$. Consequently
\begin{equation}
 0\le \log\!\left[\prod_{j=1}^N A_j(z)\right]+\Phi_a(z)
       \le |\pi|/(8h)+a^2/(4R^2).
 \label{eq:gcmone-exit-error-uv}
\end{equation}
The positive transfer converges in operator norm on this gapped
domain to $e^{-\Phi_a(H)}$. The inverse propagator with its restored polar factor and the
parity-even log determinant are thus the ones of the same
conditional action. The auxiliary approximation concerns the
squared-mass transport coordinate and imposes no tangential
momentum cutoff. Continuum traces or differentiated amplitudes
still require the estimates for that induced charged action.

Keeping one random exit value through an entire chain and averaging
only its final transport gives a different source. Its scalar
factor is $\mathbb E\sqrt{z/(z+V)}$, strictly larger than
$e^{-\Phi_a(z)}$ by Jensen since
$\Phi_a(z)=\tfrac12\mathbb E\log(1+V/z)$.
As $z\downarrow0$ the two factors, divided by $\sqrt{z/a}$,
tend respectively to $4/5$ and $e^{-1/4}$. Indeed
$V=a/\sqrt U$ for uniform $U$,
$\mathbb E U^{1/4}=4/5$, and $\mathbb E\log(V/a)=1/2$.
These different limits change the charged pole normalization. The
inverse transfer gives residues $5/4$ and $e^{1/4}$, respectively,
in carrier units $Lm=1$.
The probe, contacts and determinant must use the same stage of
shell integration throughout the chosen source. Neither transport
reverses the branch's vectorial pole prescription.

\paragraph{Obtaining the positive square roots from a local normal equation}
The transport also has a realization that uses the same charged derivative
on every normal slice. Fix a smooth background in the gapped chart
and first work on finite spectral tests. At depth parameter $v\ge0$,
extend boundary data $f$ into an auxiliary half-line with coordinate
$u$ of length dimension. Use the quadratic form
\[
 \int_0^\infty
  \big(\|\partial_u\phi\|^2+\|K_B\phi\|^2+v\|\phi\|^2\big)\,du,
 \qquad \phi(0)=f .
\]
The third term is the normal gradient energy of a Fourier mode of
squared normal momentum $v$. The decaying solution of the Euler equation
is $e^{-u(H_B+v)^{1/2}}f$. Integration by parts gives the boundary
quadratic form $\langle f,N_v f\rangle$, where the outward normal
derivative is $N_v f=(H_B+v)^{1/2}f$. Thus the fixed-record step between
depths $v_1$ and $v_2$ is $N_{v_2}^{-1}N_{v_1}$ when the outgoing
normal derivative at the first slice is used as the incoming datum at
the second. At the external boundary, unit Neumann data $g$ give
Dirichlet data $N_0^{-1}g$. Measuring their charged derivative gives
$K_BN_0^{-1}g$, exactly the polar factor in the terminal bilinear.
These identities extend by spectral calculus to the natural domains.
They recover both square roots from one normal equation. The choice to
match normal derivatives and to measure the charged derivative at the
endpoint is part of this represented boundary law.

This construction evaluates the boundary map of the normal equation.
A Gaussian integral over additional, independently quantized normal fields
would also contribute a field-dependent determinant. In the stated
representation the normal map is a constraint, with the normalized
unit-diagonal Berezin implementation above. This fixes its contribution
to the charged trace and prevents an uncounted boundary determinant.

\paragraph{The ordered charged kernel}
Take two independent rank-four projected displacements in the charged
Wilson translations of \eqref{eq:rect-polar-uv}, with
$R_B=i\sigma^E_{\mu\nu}[T_\mu[B],T_\nu[B]]/(4L^3)$ and
$K_B=m+d_B+R_B$. The commutator itself is a curvature component:
with internal exterior generators $\theta^\mu$, the comparison
connection $d+\sum_\mu\theta^\mu T_\mu[B]/L^2$ has internal
curvature $[T_\mu,T_\nu]/L^4$. The fixed spin projection gives $R_B$.
The two orders of transport have identical charged endpoints. This is
why the same variation differentiates both orders even when two
photons strike a single charged insertion. The spin projection and its
unit normalization specify the charged chart.
With the common exit transfer, its scalar charged action is
\begin{equation}
\begin{aligned}
 S_{\rm ch}[B]&=\langle\bar\Psi,L^{-1}K_B H_B^{-1/2}
                         e^{-\Phi_a(H_B)}\Psi\rangle,\\
 2\Phi_a(H_B)&=\int_0^\infty b_a(v)(H_B+v)^{-1}\,dv .
\end{aligned}
\label{eq:gcmone-combined-action-uv}
\end{equation}
Let $\mathcal D_B$ denote this bilinear coefficient. At fixed $B$
the normalized finite-test generating functional for Grassmann sources
$\bar j,j$ is
\[
 \frac{\det\mathcal D_B}{\det\mathcal D_0}
             \exp\langle\bar j,\mathcal D_B^{-1}j\rangle .
\]
The photon variable in $B$ is integrated with the same Maxwell
measure when photon lines are sewn. The charge $e$ labels the
perturbative derivatives at the chosen charge normalization.
The conditional geometry introduces no adjustable scale beyond $L$.
At the level of the original curvature integral each constraint
bilinear is one block in the direct sum of the represented
superconnection, with the same external connection and the original
GTL coarea. The stage at which the shell is integrated, the internal
projection, and the boundary gluing are specified data of this
representation. In particular this is a construction on the GCM
action in \eqref{eq:uvpeA-curvature}, and does not follow from its
scalar curvature monomial alone.

The harmonic state in \eqref{eq:fix-threshold-sphere-uv} has the same
resolvent matrix element, $\Phi_a(H_B)=2\pi\langle\psi_m,
(-\Delta_N+H_B)^{-1}\psi_m\rangle_N$. This equality concerns the
measured normal transfer. The state on the dual normal plane and
the first-exit preparation are two constructions of that weight.

\paragraph{What the same Gaussian source predicts on open and closed lines}
Its parity-even effective action is
$\operatorname{Tr}\Phi_a(H_B)$ up to the field-independent part.
Its open covariance is $L F_a(H_B)K_B^\dagger$ with
$F_a(z)=e^{\Phi_a(z)}/\sqrt z$. The zero-field Stieltjes law,
including its normalization, follows from
\eqref{eq:fix-threshold-log-uv}. It has residue $e^{1/4}$ in
the open Dirac propagator and its next spectral support occurs
at squared mass $m^2+a$. At nonzero $B$ the atom contributes
$e^{1/4}K_B^{-1}$. Its marked open amplitudes therefore
contain the actual ordered-path contacts of $K_B$. Dividing
by the external pole normalizations removes $e^{1/4}$.

At each finite electromagnetic order consider connected coefficients
with external source insertions, dividing the generating functional by
its value at zero external sources. Charged endpoints are tested using
the preceding infrared prescription. Absolute convergence
here concerns the gauge-completed rooted integrands, after combining
the contact terms prescribed by the action. For the
purely minimal words $W_a(t)=t\int b_a(v)e^{-tv}dv\sim2at$
at short time. Every additional closed word contains a marked
$R_B$ with two link transforms in $L^2(d^4k)$, uniformly in
their real translations; their product and all differentiated
products are in $L^1(d^4k)$. In the twisted resolvent proof the
mass bound is now explicit:
$\int_0^\infty b_a(v)dv/\sqrt v=8\sqrt a/3$.
At a diagonal root a constant connection is a pure gauge. Replacing
each wave by its difference from its value at the root retains one
$|q|$ per transverse incidence. The same photon spanning-tree
bound controls the hard photon circulations. The dominated soft
limit of the proper closed two-point function gives a finite Maxwell
coefficient at every perturbative order. The positive Stieltjes law integrates
the open bound $(m^2+v)^{-1}$, including many-photon contacts.
These estimates use full translations and integrate over all
hard momentum regions and all allowed shell displacements.

\paragraph{Sewing the complete two-photon contact}
For $K_v=K_B+i\sqrt v\,\gamma^5$ one has
$K_v^\dagger K_v=H_B+v$. Integration by parts in
\eqref{eq:gcmone-combined-action-uv} gives the exact identity
\[
 2\Phi_a(H_B)=-\log H_B+
       \int_a^\infty\frac{2a^2}{v^3}\log(H_B+v)\,dv .
\]
All logarithms use common fixed units; the coefficients of a
constant sum to zero. This operator identity can be applied to
discontinuities or kept together in the original resolvent.
The individual unsmeared minimal determinant traces can have
divergent real local terms. They need no separate assignment in
this proof.

Consider the two-current spectral density and the forward
elastic four-current amplitude for two nonzero real photons,
with a fixed physical polarization state. Define $C_v(S)$ to be
the corresponding pair cut in the parity-even first-order
reference, with mass $M=\sqrt{m^2+v}$, in a common cut normalization.
For a polarization that distinguishes the two pseudoscalar mass
signs, this is the equal average of their separate forward cuts:
$\log|\det K_v|$ averages the two Euclidean log determinants
because $K_v^\dagger=\gamma^5(K_B-i\sqrt v\,\gamma^5)\gamma^5$.
At $v=0$ they coincide.
For the second test let $G_v=K_v[0]^{-1}$ and let $K_j,K_{12}$
denote the first and mixed second photon variations of $K_B$.
The amputated tree operator at either side of this cut is, up
to a common overall sign,
\[
    K_{12}-K_1G_vK_2-K_2G_vK_1 .
\]
Indeed differentiating $K_v^{-1}$ twice puts precisely this
combination between two external propagators. It includes both
photon orderings and the variation of the Wilson transports.

The physical conjugation of this contact can be checked before
invoking a probability interpretation. On real Euclidean
momenta $K_B^\dagger=\gamma^5K_B\gamma^5$ exchanges incoming
and outgoing fermion and conjugates each photon wave. At real
physical momenta complex conjugation of the continued Euclidean
energy equals its reflection. For spatial transverse polarizations
this reflection acts on the vector and tensor marks by
$J=\gamma^5\gamma_4$, since conjugation by $J$ fixes $\gamma_j$,
$j=1,2,3$, and negates $\gamma_4$. The ball records are reflection
invariant. It follows that each reversed continued mark is the
$\gamma_4$-adjoint of the forward mark, with conjugated
polarizations. The free twist also obeys this rule:
adjunction changes its sign and reflection changes the sign of
$\gamma^5$. Equivalently the same chiral rotation acts on the two
sides of $K_v$ to give ordinary free mass $M$.
Every charged mark involving the commutator is entire and has the
same $L^2$ link bounds after bounded imaginary energy translations.
One may isolate the compact pinching region at any finite physical
energy, leaving the hard part analytic in a neighborhood by the
tilted link bounds. The purely minimal contribution has the usual
massive Dirac cut. Two successive massive on-shell propagators
cannot be separated by one nonzero real null photon, since its
scalar product with a timelike on-shell momentum is nonzero.
In forward kinematics the crossed transfer is spacelike or zero.
Thus the physical pair pinch joins the indicated tree and its
adjoint through the two massive residues; a contact lacking the
two pinching propagators contributes no such cut. This establishes
the squared pair weight $C_v\ge0$ for these channels with all
contact signs fixed. The argument uses finite contour deformations
on the prescribed matter boundary.

Writing $U=S/4-m^2$, the complete discontinuity from this very
source is consequently
\begin{equation}
 C_{\rm cl}(S)=C_0(S)-
 {\bf1}_{U>a}\int_a^U \frac{2a^2}{v^3}C_v(S)\,dv,
 \qquad U>0 .
 \label{eq:gcmone-combined-cuts-uv}
\end{equation}
Below and at $S=8m^2$ the physical pair cut, when open, agrees
exactly with the pole-pair contribution of the same open source
in both tests. This includes the full two-photon spin-area
contact. It is a perturbative sewing statement for these
amplitudes. Construction of an interacting Hilbert space is a separate requirement.

\paragraph{The first heavy threshold of the combined action}
The two-current channel gives a strict obstruction immediately
above the exact interval. At a reference threshold use
\[
 \begin{aligned}
 h_M={}&L^{-3}\!\int u_4\sinh(Mu_4)\,d\mu(u)\\
 &{}\times\int u_j^2\!\int_0^1
       e^{(1-2s)M u_4}\,ds\,d\mu(u)>0 ,
 \end{aligned}
\]
where $j$ is any spatial polarization and $d\mu$ is the same
projected four-dimensional displacement law at scale $L$.
The forward and reverse rest vertices before chiral rotation
have tensor corrections $\pm h_M\sigma^E_{4j}$. The rotation
multiplies the tensor by $e^{-i\vartheta\gamma^5}$,
$\cos\vartheta=m/M$. The term
$\sigma^E_{4j}\gamma^5$ commutes with $\gamma_4$, so it drops
out between the opposite rest projectors. The surviving triplet
factor is $1+(m/M)h_M$. In particular there is a nonzero pair
weight for $v$ just below $U$ at every finite $U>a$, so
\eqref{eq:gcmone-combined-cuts-uv} is strictly below $C_0$ in
the two-current channel throughout that range after the common
external readout is stripped. More precisely, with
$M_*=\sqrt{m^2+a}$, $\epsilon=U-a\downarrow0$ and the unit
Dirac normalization of the two-current density,
\[
 \rho_{\rm cl}-\rho_0=
 -\frac{[1+(m/M_*)h_{M_*}]^2}{6\pi^2 a M_*}\,
       \epsilon^{3/2}+o(\epsilon^{3/2}).
\]
This coefficient follows by integrating the leading term of the
uniform reference threshold law,
\[
 \frac{[1+(m/M_*)h_{M_*}]^2}{8\pi^2M_*}\sqrt{U-v}.
\] The harmonic state and full finite paths
therefore realize Euclidean UV regularity and exact sewing on
a common finite interval in one action. The explicit loss of
positive inclusive weight above it precludes the unrestricted
physical interpretation of that action.

The four-current loss is also strict just above $8m^2$ for an
explicit polarization. At $M=M_*$ send the two real photons along
$\pm e_3$ with perpendicular linear polarizations $e_1,e_2$.
At pair threshold the charged endpoints have Euclidean momenta
$\mp iM e_4$ and the intermediate charged momenta are $\pm M e_3$.
The two marks on one link vanish by the independent transverse
reflections applied to $u_\mu u_1u_2$. With one mark on each
link only $\sigma^E_{12}$ remains. The two orders have equal
scalar products by $e_3\mapsto-e_3$ and the equality of the two
transverse second moments. Their commutator cancels. Thus $K_{12}$
vanishes at these kinematics in the specified action.

For a single mark the change in the temporal $T$ multiplier is at
most $L\sinh w/(n+2)$ and the spatial change at most $Lw/(n+2)$,
where $w=ML$. These estimates use
$|\sin(Mu_3)|\le M|u_3|$ and the convexity bound for $\sinh$
on $[0,w]$. The transverse link variation has magnitude at most
$L^2\sinh(w)/[w(n+2)]$: integrate the absolute tilted second
moment along the link using reflection in $u_4$. Consequently the
operator norm of each correction to $i\gamma_j$, before or after
the unitary chiral rotation, is at most
\[
 d=\frac{(\sinh w+w)\sinh w}{2w(n+2)^2}<0.064
 \qquad (w=\sqrt2,\ n\ge4).
\]
Between the opposite rest projectors the two ordinary exchange
terms sum, up to sign, to $i\gamma_1\gamma_2\gamma_3/M$,
of norm $1/M$. Each free intermediate propagator has norm
$1/(\sqrt2 M)$. The combined change from both marked exchange
terms is therefore bounded by $\sqrt2(2d+d^2)/M<0.19/M$.
The full threshold matrix element is nonzero. Continuity and
two-body phase space now give $C_v(S)=[c+o(1)]\sqrt{U-v}$
with $c>0$ for $a\le v\le U$ as $U\downarrow a$.
Equation~\eqref{eq:gcmone-combined-cuts-uv} gives a strictly
negative difference of order $(U-a)^{3/2}$ in this physical
four-current polarization as well.

\paragraph{Positive stopping laws for the same charged action}
Replace the harmonic sphere by a probability law $d\pi(T)$ on
finite positive squared-mass shifts. Assume
$\mathbb E|\log(T/m^2)|<\infty$. Its survival weight
$b(v)=\Pr(T>v)$ gives, by Tonelli on the positive Euclidean axis,
\begin{equation}
 2\Phi_\pi(x)=\mathbb E\log(1+T/x),\qquad
 W_\pi(t)=\mathbb E(1-e^{-tT}),\qquad x,t>0 .
 \label{eq:gcmone-stopping-law-uv}
\end{equation}
The Stieltjes argument applies because $0\le b\le1$.
The mass-atom residue of $L F_\pi(x)=Le^{\Phi_\pi(x)}/\sqrt x$
is $Lm\exp[\tfrac12\mathbb E\log(T/m^2)]>0$.
The complete Euclidean estimate for the same rectangular
$K_B$ even allows infinite $\mathbb E T$:
$\int_0^\infty W_\pi(t)e^{-m^2t}dt/t=\mathbb E\log(1+T/m^2)$
is finite. The planar normal amplitude with Fourier transform
$\sqrt b$ is in $L^2$ exactly when $\int b=\mathbb E T<\infty$.
With infinite mean it defines a generalized boundary preparation
through its finite resolvent quadratic form at $m^2$.
The displayed time integral controls the two-mark minimal trace, and the higher
rooted minimal words have extra short-time powers. In the difference
containing a rectangular mark, the two link transforms control the
charged integral. Each marked twisted resolvent word also retains
at least two free inverses. The remaining mass integral is bounded
by $\int_0^\infty dv/[\sqrt v(m^2+v)]=\pi/m$.
The open and photon estimates above apply with this positive
Stieltjes law and the same infrared convention.

For either one-loop channel of \eqref{eq:gcmone-combined-cuts-uv}
use its locally bounded reference cut $C_v(S)\ge0$ and its usual
two-particle boundary. The logarithm in
\eqref{eq:gcmone-stopping-law-uv} is the difference of the
heavy and light operator logarithms before tracing. Taking the
complete cut at $U=S/4-m^2>0$ gives the precise equality criterion
\[
 C_0(S)-C_{\rm cl}(S)=\int_{(0,U)}C_v(S)\,d\pi(v).
\]
The upper endpoint contributes zero by threshold phase space.
Equality holds exactly when the positive integrand vanishes
$\pi$-almost everywhere below $U$, including when there is no
mass there. If $a_*=\inf\operatorname{supp}\pi$, equality holds
through $U=a_*$. In the rectangular two-current channel the rest
triplet is continuous and nonzero at every finite threshold.
Consequently all $U$ immediately above $a_*$ have a strict deficit.
More generally, if a positive-measure set of shifts has a nonzero
cut on a positive-measure energy set for each such shift, Tonelli
applied first in energy gives a strict deficit on an energy set of
positive measure. A gap in the stopping law can postpone the
mismatch. A probability law on finite shifts cannot restore
all-energy equality for this charged two-current channel.

\paragraph{Changing a charged contact to repair the cut}
There are two restrictions on such a change. First, keep the free
operator and its first photon derivative fixed. In a first-order
reference the difference of two-current loops then has only one
charged propagator and a changed two-photon contact. If the contact
is polynomial or is built from compact Wilson displacements with
the stated tilted bounds, both external momenta can be routed
through the contact. The solitary charged pole is independent of
their transfer. The localized integral has no two-particle pinch,
and an integrably dominated hard part is holomorphic across this
cut. Polynomial local ambiguities also have zero discontinuity.
Under these hypotheses an alteration starting at the second
photon derivative cannot supply the missing two-current weight.
The large-virtuality contact treated below has to be tested
separately because its tilted bounds fail.

Second, even a change of the first photon mark is constrained by
crossing. Consider any stopping-law action with the logarithm
\eqref{eq:gcmone-stopping-law-uv} whose massive reference cuts
are positive squared pair amplitudes. Suppose that for almost every
shift the first-order reference is Lorentz covariant, has a
conserved on-shell current and a nonzero elastic electric charge.
Assume its tree vertex, after removing common photon readout
factors, is holomorphic on a connected complex neighborhood in
the product of the two mass shells joining an open pair region
to forward elastic momenta. This hypothesis is satisfied by
entire finite-path marks with ordinary external Dirac poles when
the stated covariance and charge normalization hold.

If the inclusive cut contained the light-pair weight at every
energy, the positive defect integral above would vanish. Tonelli
on compact energy intervals then gives zero pair weight at almost
every energy, for almost every reference shift. At timelike
transfer the physical spin sums and the three transverse virtual
photon polarizations are positive. Vanishing of their sum of
squares forces every projected pair amplitude to vanish almost
everywhere in two-body phase space. Continuity, covariance and
current conservation give zero for every component on an open real
pair region. In Dirac variables the spin projectors are polynomial
in the complex on-shell momenta; crossing one leg to negative
energy gives the pair spin sum. The projected current therefore
vanishes on a real open set of maximal real dimension in the
product of the complex mass shells. Successive application of
the one-variable identity theorem in local coordinates continues
that zero to the elastic region. Its forward electric matrix
element would vanish, contrary to the assumed charge. Thus a
positive scalar stopping measure, ordinary charged crossing and
a nonzero reference charge exclude an all-energy repair within
this determinant class. The argument makes no assumption of a
physical Hilbert space for a different GCM boundary dynamics.

\subsection{The matched electric contact under the same normal transport}
\label{sec:gcmone-contact-uv}
The virtuality-dependent electric kernel addresses a different physical
test from the compact ordered paths. Use it in place of the ordered
kernel, keeping the boundary transfer and its determinant fixed as
functions of $H_B=K_B^\dagger K_B$. This substitution defines a second
action. Its previously measured soft coefficient has to be recomputed.

\paragraph{A fixed-momentum test of positive contact factorization}
The pole-free contact following \eqref{eq:virial-soft-cusp-uv} is a
real symmetric kernel. For fixed $x=P^2>0$ let $X=Q^2$ and write
$h(X)=\chi_{248}(L\sqrt{Z(x,X)})$, $b(X)=1-h(X)$ and
$c(X)=\partial_x h(X)/X$, with the removable value at $X=0$.
Its scalar coefficient multiplying $F_{\mu\nu}F_{\mu\nu}$ is
\begin{equation}\label{eq:gcmone-contact-Gram-uv}
 K_x(X,Y)=\frac{\mathcal R(x)}{12}h(X)h(Y)
 -\frac{x}{8}\,[b(X)c(Y)+c(X)b(Y)] .
\end{equation}
Here $\mathcal R(x)=T(x+m^2)+m^2T(x+m^2)^2$ with
$w(y)=\chi_{248}(y)^2$. Restricting photon transfers to one
Euclidean axis and their polarization to a common perpendicular
axis multiplies this kernel by nonzero separable kinematic factors.

For each $x>0$ the argument $L\sqrt Z$ increases continuously
and strictly from zero to infinity with $X$. At a positive simple
zero of $\chi_{248}$ the diagonal equals $-xc(X)/4$.
Since $\partial_x Z/X=2\Delta/m^4>0$, the sign of $c$ there
is the sign of $\chi_{248}'$. The characteristic function is a
positive multiple of $t^{-124}J_{124}(t)$, whose slope is negative
at its first positive zero and positive at its second.
The contact diagonal consequently has both signs, independently
of the pole-removing choice of $\mathcal R$. By continuity this
persists on open intervals. A real analytic photon readout normalized at zero
can have isolated zeros without changing this conclusion.

There is a test even in an arbitrarily small positive transfer
interval. The functions $h,b,c$ are real analytic and linearly
independent there. Indeed, any linear identity continues to all
positive $X$. Since $h,c\to0$ and $b\to1$ at large $X$, the
coefficient of $b$ must vanish. At a simple zero of $h$, $c\ne0$,
which also removes the coefficient of $c$. Thus three samples in
the chosen interval give an invertible evaluation matrix. Their
Gram candidate is congruent to the real symmetric matrix with
diagonal entries $(\mathcal R/12,0,0)$ and sole off-diagonal
pair $(2,3),(3,2)$ equal to $-x/8$. Since $\mathcal R>0$, its
inertia is two positive and one negative. At $x=0$ the explicit
factor $x$ removes the mixed term.

In particular this literal contact cannot be the covariance or
Gram kernel of one family of marks in a positive inner-product
space. Independent positive A4 records can still represent its
products as signed bilinear insertions. Positivity of the full
electromagnetic response concerns all insertions and the state
space and does not reduce to positivity of a single contact.

\paragraph{Apply the harmonic normal trace to the complete charged kernel}
In the harmonic normal construction of
\eqref{eq:gcmone-combined-action-uv}, use as the full charged
kernel $K$ the virtuality-dependent Dirac vertex tested in
\eqref{eq:virtmark-smallq-uv}, its straight-chord completion
and the quadratic contact of \eqref{eq:virial-contact-uv}
with the pole-free replacement following \eqref{eq:virial-soft-cusp-uv}.
This choice has no ordered-path spin insertion. On the Euclidean
contour $Q=q e_4$, $X=q^2$, $x=P^2$,
$\Delta=x+X/4+m^2$, $Z=X[1+X/(4m^2)+(\Delta/m^2)^2]$,
and $h=\chi_{248}(L\sqrt Z)$. The stripped one-mark vertex is
$h\gamma^\mu+(1-h)Q^\mu\slashed Q/X$. Keep $m,L$ and all marks
fixed when differentiating the
normal mass parameter $v$. Write $D=x+m^2$, $U=m^2+v$ and
$D_v=x+U$. The harmonic normal state with matching radius $m$ gives
the signed marked trace
\[
 d\zeta(v)=\delta_0(dv)
       -{\bf1}_{v\ge m^2}\frac{2m^4}{v^3}\,dv .
\]
This follows by integrating the mass derivative of the complete
first-order reference $K+i\sqrt v\,\gamma^5$ by parts at fixed
nonzero transverse transfer. The reference tends to zero for
$v\to\infty$ at fixed transfer by the decay of the marks. A chiral
rotation leaves a vector mark invariant and rotates a scalar mark by
$\exp(-i\theta\gamma^5)$, $\cos\theta=m/\sqrt U$. Its one-propagator
spin trace is therefore $4m/D_v$.
Using $4\sqrt U/D_v$ here would change the two-current calculation.

Here $w(y)=\chi_{248}(y)^2$ on real $y$, and the entire
extension is the algebraic square. Put
$T(D)=[1-w(D/m^2)]/D$ with its removable value at zero,
$\mathcal R=T(D)+m^2T(D)^2$ and $V(x)=1-x\mathcal R$.
For each $v$ the unshifted bubble and chord give
$(e^2/8\pi^2)\int(1-U^2/D_v^2)(2h-h^2)dx$.
The derivative part of the scalar contact contributes
$(e^2/8\pi^2)\int (x^2/D_v)\partial_x(2h-h^2)dx$.
Since $(x^2/D_v)'=1-U^2/D_v^2$, cancellation holds for every
normal mass before averaging. For $Q=q e_4$, $q>0$, the remaining
transverse coefficient, divided by $e^2q^2/(16\pi^2)$, is
\begin{align}
 \int_0^\infty dx\bigg[&
 \frac{4V(x)h^2}{3D_v}
 -\frac{4U^2h^2}{3D_v^3}
 +\frac{xU^2(h^2-2h)}{3D_v^4}\bigg]
 +O(q^2/U).
 \label{eq:combined-harmonic-virtuality-uv}
\end{align}
Expand the shifted vector bubble to second order and integrate the
straight chord by parts in $P_1$ before doing the angular average.
Its fourth-order remainder and the bubble remainder, after division
by $q^2$, are bounded in the radial integral by $Cq^2 x/D_v^3$
for $q\ll m$. For the displayed radial terms, $|h|\le1$,
$0<V\le1$ and $\int V(x)/D\,dx<\infty$.
The other absolute integrals are uniform since
$\int_0^\infty U^2 D_v^{-3}dx=1/2$ and
$\int_0^\infty xU^2D_v^{-4}dx=1/6$.
In particular both orders of the soft limit and the finite signed
mass integration are justified at leading order.

At fixed $v$ the dimensionless Maxwell coefficient is
$-13/18+(4/3)\int V(x)/D_v\,dx$. The mass-independent part cancels
in the actual normal trace, which instead gives
\begin{equation}\label{eq:combined-harmonic-Maxwell-uv}
 \begin{aligned}
 C_{\rm pol}&=\frac43\int_0^\infty V(x)\mathcal B(x)\,dx>0,\\
 \mathcal B(x)&=\int\frac{d\zeta(v)}{D_v}\\
 &=\frac{2m^2}{D^2}\left[1-\frac{m^2}{D}
                  \log(1+D/m^2)\right].
 \end{aligned}
\end{equation}
Here $V\ge m^4/D^2$ since $0\le w\le1$. Both factors are strictly
positive. For $w(y)=\chi_{248}(y)^2$, numerical integration gives
$C_{\rm pol}=1.11756608$, hence
$\lim\Pi_{11}/q^2=0.00707707\,e^2$ with common external readout
factors stripped.

\paragraph{A joint short-duration obstruction to a gapped spectrum}
The next term tests normal masses and charged virtualities escaping
together. For $j=3,4$ let
$A_j(x)=\int U^2 D_v^{-j}d\zeta(v)$. As $x/m^2\to\infty$,
\[
 A_j(x)=\frac{m^4}{x^j}\big[-2\log(x/m^2)+O(1)\big].
\]
On $m^2<v<x$ replace $U^2/v^3$ by $1/v+O(m^2/v^2)$ and
$(1+U/x)^{-j}$ by $1+O(U/x)$. Both integrated errors are bounded,
as is the scaled integral on $v>x$. This logarithm is the
truncated second mass moment of the harmonic trace.
With $s=Lq$ and $n=248$, the combined coefficient satisfies
\begin{equation}\label{eq:combined-harmonic-double-log-uv}
 \begin{aligned}
 \frac{16\pi^2\Pi_{11}(q e_4)}{e^2q^2}
 &= C_{\rm pol}-\frac{4s^2}{3(n+2)}\log^2(1/s)\\
 &\quad +O\!\left(s^2\log(1/s)+q^2/m^2\right).
 \end{aligned}
\end{equation}
Indeed the two $A_j$ terms in
\eqref{eq:combined-harmonic-virtuality-uv}, less their $h=1$ values,
give at large $x$ the leading radial difference
\[
 \frac{2m^4\log(x/m^2)}{x^3}
       \left[(h^2-1)+\tfrac23(h-1)\right].
\]
On $m^2<x<m^2/s$ the quadratic term of the characteristic function
is $h-1\sim -s^2x^2/[2(n+2)m^4]$ in the joint region
$m^2\ll x\ll m^2/s$. Its integral gives the coefficient in
\eqref{eq:combined-harmonic-double-log-uv}. To bound the error,
use the fourth-order Taylor bound on this interval; the rescaled
fourth power contributes $O(s^2\log(1/s))$. Above $m^2/s$ the
bounded marks give the same bound by integrating
$\log(x/m^2)/x^3$. The nonlogarithmic parts of $A_j$ contribute
$O(s^2\log(1/s))$, and $V\mathcal B=O(m^6/x^4)$ contributes
$O(s^2)$ to the difference.

Thus this same open and closed candidate has a finite Maxwell
coefficient, but its scalar two-current function has no finite
derivative in $Q^2$ at zero and cannot have a purely gapped
dispersion with finite local subtractions. This is a failure of the
gap test; the sign alone does not exclude a positive gapless
measure. The finite-imaginary-strip proof for compact translations
also cannot be imported into this calculation. For $q=i\eta$,
$0<\eta<m$, the large-$x$ mark contains
$\mathbb E\cosh((L\eta x/m^2+O(1))Y)$, exponentially large by the
positive measure of $|Y|>1/2$. Exchanging real-loop integration
with such continuation requires a separate contour construction.
The Euclidean nonanalyticity already excludes the stated gapped
continuation independently of any such construction.

\subsection{What the boundary measure must determine}
\label{sec:gcmone-boundary-uv}
The preceding deficits have an immediate test in a lossless real-time
description. One can derive the test without guessing a density for the
unobserved boundary marks.

\paragraph{What a source-dependent probability can change in a physical cut}
Use finite one-particle tests and a real compact source pulse $eA$.
Let $P_\pm$ be the free particle and filled-sea subspaces. In the
interaction picture assume each conditional one-particle transfer
$U_\xi(e)$ is unitary and $U_\xi(0)=I$. The same fermionic action
determines its vacuum amplitude and its vacuum-to-pair mark:
\[
 z_\xi(e)=\det(P_-U_\xi(e)P_-|_{P_-}),\qquad
 \beta_\xi(e)=P_+ U_\xi(e)P_- .
\]
The vacuum amplitude follows by taking the exterior product of
all occupied orbitals. Unitarity gives the exact identity
$|z_\xi|^2=\det(I-\beta_\xi^\dagger\beta_\xi)$ on the occupied
subspace. All exclusive occupations are minors of the same
transferred Slater determinant; their squared moduli sum to one by
Cauchy--Binet. Thus the inclusive probability is fixed before
assigning a probability to the boundary record.

Now let $d\pi_e$ be any twice differentiable normalized positive
reweighting of those records, with uniformly integrable derivatives
on the tests. Put $\pi=\pi_0$, $z'_\xi(0)=i\theta_\xi$ with
$\theta_\xi$ real, and $T_\xi=\beta'_\xi(0)$. The coherent vacuum
readout $\bar z(e)=\int z_\xi(e)d\pi_e$ obeys
\begin{align}
 1-|\bar z(e)|^2
 &=\mathbb E_{\pi_e}(1-|z_\xi|^2)
                     +\operatorname{Var}_{\pi_e}z_\xi,\nonumber\\
 [e^2](1-|\bar z(e)|^2)
 &=\|\mathbb E_\pi T_\xi\|_{\rm HS}^2+
   \mathbb E_\pi\|T_\xi-\mathbb E_\pi T_\xi\|_{\rm HS}^2+
   \operatorname{Var}_\pi\theta_\xi .
 \label{eq:gcmone-boundary-optical-uv}
\end{align}
Indeed $1-|z_\xi|^2=e^2\|T_\xi\|_{\rm HS}^2+o(e^2)$ and
$z_\xi-1=ie\theta_\xi+O(e^2)$. This proves the cancellation
of both derivatives of the normalized record density in the
leading absorptive coefficient. If only charged-pair occupations
are counted, the last, boundary-only term is absent. A lossless
rotation on boundary factors alone also cannot change the first
pair transition from the vacuum: its generator has zero matrix
element between those two occupation sectors. These statements
apply to a common zero-pair initial state and this lossless
real-time realization; a Euclidean density derivative need not
by itself possess such a realization.

For a given light-pair amplitude $u$ and a nonnegative target
$r\|u\|^2$ with $r<1$, any additional amplitude $d$ on the same
pair states must satisfy $\|d\|\ge(1-\sqrt r)\|u\|$ if the light
channel is to fit inside that inclusive weight. This follows from
$\|u+d\|^2\le r\|u\|^2$ and the triangle inequality. Additional
orthogonal channels only strengthen the bound. If the closed
weight is negative, even removing the entire light amplitude
cannot give it a positive inclusive interpretation.

\paragraph{Why the choice of boundary gluing matters for ultraviolet power counting}
There is a short independent check at the level of a complete determinant.
Let $T$ average translations of a charged endpoint with the forward A4
law of \eqref{eq:geo-actual-phase-rank-uv}, using its rank-four density
in $L^2$, and let $t(p)$ be its Fourier multiplier. Give a Pauli
curvature insertion the two independent endpoint marks $T^*M_{TF}T$,
where $M_{TF}$ multiplies by the averaged curvature with its spin
contraction. Consecutive insertions put $|t(p)|^2$ on every massive
charged line. Two smoothed curls sewn by a Maxwell line have bound
$C|t(q)|^2$ as well. Both line bounds are integrable and bounded.
Bounding a spanning tree by the suprema and integrating the independent
chords proves the absolute fixed-graph estimate for this curvature
source. If one displacement is held fixed on every charged endpoint
before its final average, it instead conjugates the whole loop by a
translation. Its phases cancel in the trace. In a constant magnetic
field $F_{12}=B$ the spin trace on each momentum sphere then leaves
the positive term
\[
 2(eLB)^2m^2\int\frac{d^4p}{(2\pi)^4(m^2+|p|^2)^2}.
\]
This integral diverges logarithmically. The Pauli source supplies a
transverse vertex of order $|q|$ at soft transfer. These two
realizations of the same velocity marginal already have different
ultraviolet behavior before the stronger requirement of a nonzero
soft electric charge is imposed.

\paragraph{Consequences for the physical GCM action}
The GTL coarea, the velocity-shell delta and the directed carrier
specify an event on the locked plane. The cylinder exit and the
normal boundary equation specify how a charged state moves between
events on the represented chart. Varying the resulting one source
yields \eqref{eq:gcmone-combined-cuts-uv}, including its complete
second photon contact. The deficit above the first heavy threshold
places a physical restriction on this representation even when every
rooted Euclidean estimate holds. Changing the positive law of a
scalar exit time can move the threshold. The sign of the missing
weight follows from the same determinant for each such law. With positive
reference cuts and analytic charged crossing, the only way to annihilate
all the heavy pair amplitudes within this construction would also
annihilate their elastic electric charge.

For the virtuality-dependent electric vertex the pole-free contact
has a different outcome after the normal integral. Its finite
coefficient in \eqref{eq:combined-harmonic-Maxwell-uv} and its
nonanalytic correction in \eqref{eq:combined-harmonic-double-log-uv}
are predictions of one action. The independently matched contact
cannot be carried unchanged into a gapped physical theory with
that boundary law. At the level of real-time probability,
\eqref{eq:gcmone-boundary-optical-uv} requires the transition on
the observed pair states to change if their coherent weight is to
decrease.

A physical choice of $\Omega_\alpha$ would have to determine this
charged transition, its higher contacts and the gluing of the
unseen normal histories. Once specified, the physical cut can be
compared directly with the open pair and the complete marked
source tested on the hard subsets of
\eqref{eq:geo-actual-phase-rank-uv}. The restricted construction
shows which choices suffice for Euclidean convergence and which
fail the physical test. Selection of a different charged
representation by the ambient dynamics and verification of its
cut and hard bounds remain necessary for the unrestricted claim.

\subsection{Restoring the charged measure and testing the resulting photon}
\label{sec:next-physical-join-uv}
An all-energy physical test must use the same charged variables for
both the pair-production matrix element and the determinant. The calculations use the scalar matter coefficient on
$\bin=(s,b_s)$ with directed lift $\sigma_s\sigma_\epsilon
\hat p_{0,\bin}$; the geometric lift is $\sigma_\epsilon\hat p_{0,\bin}$.
The pole displacements are $\pm i\sigma_{s\epsilon}\varepsilon
\hat p_{0,\bin}$, $\varepsilon\downarrow0$. The spectral variable
$S$ belongs to the reconstructed Minkowski readout. The forward
A4 distribution and the original GTL pullback define the geometric
input. There is an exact measure that restores the missing discontinuity in the
stopping-law construction. Its field dependence can be computed
before any ultraviolet estimate. This provides a direct test of
what would happen if the normal transport were simply a change
of variables in the physical action.

\paragraph{Restoring the measure of the charged field}
There is an exact way of restoring the missing pair weight on the
charged chart of \eqref{eq:gcmone-combined-action-uv}. It also
reveals its ultraviolet cost. Use the real Euclidean connection
and the same four-dimensional readout on $\bin=(s,b_s)$. On a
finite gapped test space set $G_B=\gamma_5 K_B=G_B^\dagger$ and
$H_B=G_B^2$. For the stopping law $d\pi(v)$ in
\eqref{eq:gcmone-stopping-law-uv}, define the positive operator
\[
 P_B=\exp\!\left[\frac12\int
       \log\big((H_B+v)/m^2\big)\,d\pi(v)\right],\qquad Lm=1 .
\]
The logarithmic moment makes it finite. The factor of $H_B^{1/2}$
in the exit law cancels its reciprocal in the polar factor. Since
$P_B$ commutes with $G_B$, the charged bilinear is exactly
\begin{equation}
 \mathcal D_B=K_B P_B^{-1},\qquad
 \gamma_5\mathcal D_B=P_B^{-1/2}G_BP_B^{-1/2}.
 \label{eq:join-congruence-uv}
\end{equation}
Thus one can obtain the bilinear by a change of both charged
coordinates in the first-order theory. Put
$\varphi=P_B^{-1/2}\chi$ and
$\bar\varphi=\bar\chi P_B^{-1/2}$ in the Hermitian form.
Berezin integration supplies $\det P_B$ in these coordinates.
Keeping that Jacobian multiplies the normalized photon functional
by $\det P_B/\det P_0$. The total closed determinant becomes
$\det K_B/\det K_0$. With probes coupled to $\chi$, its normalized
open inverse is still $\mathcal D_B^{-1}$. This calculation includes
the field dependence of the change of variables. Differentiation
therefore includes its measure contacts at every order. The
ordinary first-order reference has the light-pair cut $C_0$ in both
one-loop channels of \eqref{eq:gcmone-combined-cuts-uv}; in the
restored determinant this agreement extends to all finite pair
energies. It is an equality of discontinuities in those channels.

The determinant identity survives noncommuting changes of charged
coordinates. For an invertible common endpoint map $V_B$,
$\det(V_B^\dagger G_B V_B)=|\det V_B|^2\det G_B$.
The corresponding change of Berezin measure cancels the factor
$|\det V_B|^2$. This holds even if $V_B$ is a product of
noncommuting path transports. It is a statement about one
invertible coordinate change on the full test space. Projecting
out states or averaging distinct actions before the determinant
requires its own gluing calculation.

\paragraph{The current cut fixes the added mass measure}
Consider adding a determinant weight with the same shifted charged
vertices at each $v$, described on finite tests by a finite signed
measure $d\nu(v)$ in place of $d\pi(v)$ in $\log\det P_B$.
Use the normalized background difference and suppose the marked
integrals exist. Write $U=S/4-m^2$ and suppress the common photon
readout. In the complete two-current channel its cut is
\[
 C_{\nu}(S)=C_0(S)+\int_{[0,U)} C_v(S)\,d(\nu-\pi)(v).
\]
Here $C_v(4(m^2+U))=(U-v)^{1/2}c(U,v)$. On compact triangles
$0\le v\le U$, the rectangular tree vertex gives a continuous
$c$ with two continuous $U$ derivatives up to the diagonal. The
compact Wilson displacements give smooth marks on each bounded
complex energy strip. In the pair rest frame the angular integral
removes odd powers of relative momentum and the remaining
two-body phase space gives the displayed square root. The rest
triplet calculated above gives $c(v,v)>0$ at every finite shift.
These regularity and nonvanishing conditions also suffice for
the following test for other reference cuts.

Suppose $C_\nu=C_0$ at every energy. Set $d\mu=d\nu-d\pi$.
Differentiate the zero cut difference in $U$ in distribution sense.
The new kernel is $(U-v)^{-1/2}d(U,v)$ with
$d=c/2+(U-v)\partial_U c$. Integrating once more against
$(x-U)^{-1/2}$ gives
\[
 \int_{[0,x]}J(x,v)d\mu(v)=0,\qquad
 J(x,v)=\int_0^1\frac{d(v+t(x-v),v)}{\sqrt{t(1-t)}}\,dt .
\]
Fubini applies locally to the total variation by the finite beta
integral. Taking a right limit in $x$ includes any atom at $x$.
In particular the beta integral gives $J(v,v)=\pi c(v,v)/2>0$.
Differentiation as measures gives
\[
 J(x,x)d\mu(x)=-
    \left[\int_{[0,x]}\partial_xJ(x,v)d\mu(v)\right]dx .
\]
On each compact interval the diagonal has a positive lower bound
and $\partial_xJ$ is bounded. Hence the cumulative total variation
$M(x)=|\mu|([0,x])$ obeys $M(y)\le C\int_0^y M(x)dx$.
Gronwall's inequality gives $M=0$. Thus exact light-pair saturation
through a given energy range fixes the added measure below its
upper threshold, including its atoms. Exact saturation on the
whole ray fixes $\nu=\pi$. If other physical channels contribute,
the relevant test is the inclusive inequality; saturation need
not hold.

\paragraph{The Maxwell cost of a positive determinant restoration}
Any nonzero finite \emph{positive} $d\nu$ in this class already
changes the divergent coefficient of the photon two-point function.
For its minimally gauged shifted Dirac factor the short-time
Maxwell density in $-\log\det$ is
\begin{equation}
 \frac{e^2}{48\pi^2} F_{\mu\nu}F_{\mu\nu}
       e^{-m^2t}\int e^{-tv}d\nu(v)\,\frac{dt}{t}.
 \label{eq:join-maxwell-cost-uv}
\end{equation}
Here $F$ is the curvature of the represented Abelian readout with
the charge stripped; the density is per unit Euclidean volume.
For example, in a constant $F_{12}=B$ the spin trace is
$4(4\pi t)^{-2}e^{-(m^2+v)t}(eBt)\coth(eBt)$.
Its quadratic coefficient, with the factor $dt/(2t)$, gives
\eqref{eq:join-maxwell-cost-uv}. The same local term occurs
for smooth varying fields with nonzero curvature. The difference
between this first-order trace and the trace with the full ordered
rectangle has at least one marked commutator. Its two square-integrable
link transforms give an integrable charged momentum bound at fixed
photon transfers. The remaining inverse factors are uniformly
bounded for $v\ge0$ since $m>0$. This also bounds their finite
mass average, including the second photon contact. Thus the
ordered paths leave the displayed coefficient intact. As
$t\downarrow0$ the integral in it tends to the total mass of
$\nu$. The original normalized exit source has a finite Maxwell
coefficient. Adding any positive amount of these ordinary
reference determinants therefore introduces a bare logarithm.

The full Jacobian gives an all-energy repair of the two tested
cuts and restores the logarithm with unit weight. A partial
positive restoration incurs the logarithm as soon as it carries
any weight. This identifies a required physical change: charged
boundary states that restore inclusive probabilities while retaining
the ultraviolet bound must change the charged transition or its
physical sewing beyond this class of added determinants.

\paragraph{The ultraviolet cost of delaying the first missing pair}
Let $\mu$ be the scalar exit probability measure and
$M_{\log}=\int\log(1+T/m^2)\,d\mu(T)$. To match the light-pair
cut throughout $4m^2<S\le S_*$, every charged reference with
a nonzero threshold transition must have
$T\ge S_*/4-m^2$ almost surely: otherwise integrate its
strictly positive threshold contribution in energy and use
Tonelli. Consequently
\[
 M_{\log}\ \ge\ \log\frac{S_*}{4m^2}.
\]
When $S_*/4-m^2$ is an allowed shift, concentrating the exit there
attains equality. A larger imposed lower endpoint on the exits
strengthens the bound. It quantifies the cost of postponing the
deficit to arbitrarily high energy. The logarithmic moment is
the exact weighted two-mark proper-time integral
$\int W_\mu(t)e^{-m^2t}dt/t$.

For the minimally coupled Dirac kernel there is also a directly
observable version. Feynman parametrization of the complete relative
two-current trace gives, per copy and with common readout factors
removed, the Euclidean Maxwell coefficient in the positive screening
convention
\[
 \frac{e^2}{2\pi^2}\int_0^1 dx\,x(1-x)
 \int \log\!\left(1+\frac{T}{m^2+x(1-x)Q^2}\right)d\mu(T).
\]
Here integrating the common routed light-minus-heavy Gaussian in
proper time gives the logarithm; the transverse Dirac numerator
gives $x(1-x)$. At $Q^2=0$ the coefficient is
$e^2 M_{\log}/(12\pi^2)$. The bound also has an approximate
version. At a single $U=S/4-m^2>0$ put
$\delta=(C_0-C_{\rm cl})/C_0$. For $0<\eta<1$ the
minimal reference cut satisfies $C_T/C_0\ge\sqrt{1-\eta}$ on
$T\le\eta U$: its phase-space ratio is $\sqrt{1-T/U}$ and its
Dirac spin-factor ratio is at least one. Hence
\[
 M_{\log}\ge
 \left(1-\frac{\delta}{\sqrt{1-\eta}}\right)_+
                     \log(1+\eta U/m^2).
\]
This estimate uses the measured deficit at one energy. A sequence
of such minimally coupled preparations approaching the complete
Dirac pair weight on expanding energy ranges cannot keep a uniformly
finite bare Maxwell coefficient. The spin-area coupling has a
different spin factor, so the last quantitative estimate is
specific to the minimal kernel.

\subsection{Can mixing normal states preserve charged probabilities?}
\label{sec:next-physical-matrix-uv}
The A4 shell can prepare several normal states with the same electric
charge. Ordered matrix transport retains more information than a
scalar exit depth. We can ask separately what the source remembers
when a charged line is opened and when it is closed.

\paragraph{Mixing the boundary states at successive depths}
A scalar exit time loses information about which internal charged state
reached a particular depth. To retain that information, take $d$ copies
of the same unit-charge carrier and let $E(dT)$ be a positive
operator-valued probability on squared-mass exits $a\le T<\infty$,
with $E([a,\infty))=I_d$. Here $a=m^2$, $Lm=1$. Write
$\mathsf b(v)=E((v,\infty))$ for its survival operator. Assume $\int\log(1+T/a)\,\operatorname{tr}E(dT)<\infty$.
For finite $d$ this implies
$\int_0^\infty \|\mathsf b(v)\|\,dv/(a+v)<\infty$ by positivity.
Only the normal gluing is changed: the bulk $K_B$ and its photon
coupling act identically on the $d$ copies.

Use the positive normal boundary quadratic form on each side of
the gluing. For a fixed record its metric at depth $v$ is
$(H_B+\min(v,T))^{1/2}$, the map from boundary value to outward
normal derivative derived above. Transport preserving this metric,
with self-adjoint metric connection, has eigenvalue
$[(z+\min(u,T))/(z+\min(v,T))]^{1/4}$ on a step $u\to v$.
Average each step against $E$ before integrating the charged fields.
At a fixed positive squared singular value $z$ the product of steps
converges to $A_z(\infty)$, where
\begin{equation}
 \frac{d A_z(v)}{dv}
    =-\frac{\mathsf b(v)}{4(z+v)}A_z(v),\qquad A_z(0)=I_d .
 \label{eq:nextmatrix-gluing-uv}
\end{equation}
Indeed the first-order term of a step is the integral of its
displayed generator, and the remainder on a bounded depth interval
is bounded by a constant times the squared mesh width. The tail is
controlled by the integrable norm of the generator. Pair the two
ends by the adjoint, giving $V(z)=A_z(\infty)^\dagger A_z(\infty)$
for real $z>0$. The resulting charged bilinear is
$m K_B H_B^{-1/2}V(H_B)$ with $H_B=K_B^\dagger K_B$.
It uses the same auxiliary constraint with unit diagonal on both arms.
The adjoint pairs the boundary states on the fixed sector plane.

\paragraph{Closed and open responses of the matrix transport}
The identity $d\log\det A=\operatorname{tr}(A^{-1}dA)$ removes
the ordering from the determinant. One obtains
\begin{equation}
 -\log\det V(z)
 =\tfrac12\int_0^\infty\frac{\operatorname{tr}\mathsf b(v)}{z+v}\,dv
 =\tfrac12\int\log(1+T/z)\,\operatorname{tr}E(dT).
 \label{eq:nextmatrix-trace-uv}
\end{equation}
At finite spectral tests, tracing over the charged carrier now gives
exactly the relative effective action of the scalar stopping source
with mass measure $\operatorname{tr}E$ of total weight $d$. Every
closed photon derivative follows from this identity at nonzero
background. In particular, all closed rooted estimates proved for
that source hold, including simultaneous short-duration limits.

The open amplitude does retain matrix information. Write $A_z^>$
for the transport from $a$ to infinity and put
$C(z)=(A_{\bar z}^>)^\dagger A_z^>$. These functions are analytic
for $|z|<a$, by the norm-convergent ordered series there, uniformly
on compact subdisks. The evolution below $a$ is scalar, so
$V(z)=\sqrt{z/(z+a)}\,C(z)$ on $z>0$. The pole term of the open
covariance has coefficient $R=(\sqrt a/m)C(0)^{-1}>0$. More explicitly,
\[
 \begin{aligned}
 \mathcal D_B^{-1}&=R K_B^{-1}+J(H_B)K_B^\dagger ,\\
 J(z)&=\frac{m^{-1}\sqrt{z+a}\, C(z)^{-1}-R}{z}
               \quad (z\text{ near }0).
 \end{aligned}
\]
The remainder is regular at the light mass. At the one-photon
transition, differentiating a regular matrix function introduces
divided differences regular at both external poles. After pole
amputation and normalization by $R^{-1/2}$ on each charged end,
each copy consequently has the light pair amplitude of $K_B$.

There is a direct hard-momentum estimate on the open source.
In the ordered series for the inverse transport each unmarked
generator is bounded by $\|\mathsf b(v)\|/[4(m^2+v)]$. To differentiate
a marked resolvent put $G=\gamma_5K_B$ on the real Euclidean chart.
Then $H_B=G^2$ and $(H_B+v)^{-1}$ is the difference of
$(G-i\sqrt v)^{-1}$ and $(G+i\sqrt v)^{-1}$ divided by
$2i\sqrt v$. A word with at least one marked vertex has at least
two free inverses, each of norm at most $(m^2+v)^{-1/2}$.
Using the bounded compact-path marks with their transverse readout
factors gives the integrable depth majorant
$C_r/[\sqrt v(m^2+v)]$ for $r\ge1$ fixed, after those factors
are extracted. Extra inverses cost at most $1/m$ apiece.
Summing intervening unmarked ordered words costs at most the
exponential of the integrated unmarked norm on each interval.
The remaining polar factor $H_B^{-1/2}K_B^\dagger$ has norm one
at zero field. Its marked terms have the same integrable majorant using its integral of $(H_B+v)^{-1}K_B^\dagger$ with weight
$dv/(\pi\sqrt v)$ and the mean of the two twisted charged inverses.
This proves the open hard bound at every fixed order, with the
same tested endpoints, photon tree estimate and infrared convention
as the scalar construction. No matrix-valued spectral positivity
assumption enters this Euclidean bound.

If $C_T(S)$ denotes the two-current pair weight of the first-order
mass-shifted reference and $U=S/4-m^2$, the light-pair weight
minus the closed cut is
\[
 d C_0(S)-C_{\rm cl}(S)
       =\int_{[a,U)} C_T(S)\,\operatorname{tr}E(dT).
\]
The integral is zero for $U\le a$. Positivity and the nonzero
reference pair transition give the same obstruction on the full
energy ray as for scalar exits. Retaining matrix order can change
the pole residues while leaving the missing positive measure in
the cut unchanged.

\paragraph{A field-dependent rotation at a charged endpoint}
A separate test identifies a limit of field-dependent mixing.
On finite tests let $\widetilde{\mathcal D}_B=X_B\mathcal D_B Y_B$
with invertible covariant endpoint maps and
$\det X_B\det Y_B$ constant in $B$. The normalized determinant
and all its derivatives agree. Suppose the endpoint maps, their
inverses and first photon variations have regular kernels at the two
external mass shells. The derivative of
$Y_B^{-1}\mathcal D_B^{-1}X_B^{-1}$ that acts on an endpoint
contains only one open covariance. Its double residue at the
two charged poles consequently vanishes. The derivative on the
middle factor has the original double residue transported by the
zero-field endpoint maps. With the induced external states and
their pole normalizations this is the original pair transition.
Thus an invertible regular redefinition of the observed endpoints
cannot supply the missing weight. A prospective repair through
mixing must alter the physical transition inside the propagation,
the measure of the closed trace, or the singular structure of
the charged boundary.

\paragraph{The free spectral test at three exit depths}
For a free Dirac carrier in this matrix construction, positivity of
the vector spectral weight can already decide which gluings survive.
Suppose the exit has just three values $0<a<b<c$. Write the four
successive tails as $I_2,P,Q,0$ with $0\le Q\le P\le I_2$.
The three positive effects are $I_2-P,P-Q,Q$.
They can all be prepared on the uniform A4 shell coordinate by
partitioning it into three equal intervals and assigning three times
the corresponding effect as its density on each interval.
In this two-component model the free spectral measure is positive
if and only if $P$ and $Q$ commute.

To prove the nontrivial direction, remove the common positive
constant $1/m$ from the vector part of the open covariance. With
the principal branches fixed by $z>0$, its matrix multiplier is
\begin{align}
 F(z)&=\frac{\sqrt{z+a}}{z}\,D(z)
       \left(\frac{z+c}{z+b}\right)^{Q/2}D(z),\nonumber\\
 D(z)&=\left(\frac{z+b}{z+a}\right)^{P/4}.
 \label{eq:nextmatrix-free-test-uv}
\end{align}
This is the exact ordered product of the three constant-generator
steps. At $z=-u+i0$ with $a<u<b$, take the eigenbasis of $P$
with eigenvalues $p_i$, and put $\alpha_i=\pi p_i/4$,
$G=[(c-u)/(b-u)]^{Q/2}$. Up to the positive factor
$\sqrt{u-a}/u$ and conjugation by an invertible positive diagonal
matrix, the candidate spectral density $-\operatorname{Im}F$ has
entries
\[
                 G_{ij}\cos(\alpha_i+\alpha_j).
\]
Here the imaginary part is the Hermitian imaginary part.
If $[P,Q]\ne0$, both matrices have simple spectra, and a normalized
top eigenvector $w$ of $Q$ has two nonzero components in this basis.
As $u\uparrow b$, $G/\operatorname{tr}G\to ww^\dagger$. The
determinant of the corresponding normalized candidate tends to
\[
 -|w_1w_2|^2\sin^2(\alpha_1-\alpha_2)<0,
\]
since $\cos(2r)\cos(2s)-\cos^2(r+s)=-\sin^2(r-s)$.
Congruence preserves the sign of the determinant. A negative
spectral direction therefore persists on an open mass interval.
If $P,Q$ commute, their common eigenbasis gives two scalar
positive survival laws. Each open multiplier is Stieltjes by the
scalar argument. Its Dirac scalar weight is $m$ times the vector
weight on masses at least $m$, so the free charged positivity
inequality holds in each component. This proves the equivalence.

For example, take $b=2a$, $c=3a$,
$P=\operatorname{diag}(1,1/2)$ and let $Q$ have all four
entries $1/4$. At $u=3a/2$ the positive diagonal congruence is
the identity. The first diagonal entry of the candidate vanishes
and the off-diagonal entry is
$(3^{1/4}-1)\cos(3\pi/8)/2\ne0$. The resulting negative
determinant occurs despite a positive Euclidean transfer, a
normalized positive exit measure and finite Euclidean graphs.
The three-depth equivalence does not assume or establish a
classification for general continuous exit laws.

\paragraph{A fourth geometric depth restores positive charged propagation}
The three-depth obstruction can be repaired while keeping mixed
charged histories. Choose exits $a,2a,3a,4a$. Their successive
matrix tails are $I_2,P,rI_2,Q,0$, where
\[
 P=\operatorname{diag}(3/4,1/2),\qquad r=1/2,\qquad
 Q=\tfrac18\begin{pmatrix}1&1\\1&1\end{pmatrix}.
\]
Their differences are positive and sum to $I_2$, and $[P,Q]\ne0$.
Four equal intervals of the uniform shell coordinate, with four
times each difference as its preparation density, realize this exit
measure. The extra interval transports all boundary flavors at
the same rate between the two anisotropic steps.

For this construction the vector part of the free Dirac propagator
is matrix Stieltjes. To verify the claim, form the ordered product as
in \eqref{eq:nextmatrix-free-test-uv}. Each interval of constant
tail $H$ from $l$ to $h$ contributes a factor
$[(z+h)/(z+l)]^{H/4}$ to the inverse open transfer on each arm.
For $z=-u+i0$, $u>a$, the scalar prefactor is
$-i\sqrt{u-a}/u$. Only the factor on the interval containing $u$
has a complex phase. All earlier factors act by invertible
congruence, all later factors form a positive middle matrix.
On $a<u<2a$ the middle matrix, up to a positive scalar, is
$G=[(4a-u)/(3a-u)]^{Q/2}$. After removal of the positive diagonal
modulus of the active factor, its spectral sign is determined by
the Hermitian real part of $UGU$, $U=e^{-i\pi P/4}$.
The smallest eigenvalue is at least
\[
 \cos(3\pi/8)-\|G-I_2\|
 \ \ge\ \cos(3\pi/8)-(2^{1/8}-1)>0 .
\]
This follows by perturbing $\operatorname{Re}(U^2)$ and using
unitarity of $U$. On $2a<u<3a$ the active phase is scalar, so the
spectral matrix is a positive congruence times $\cos(\pi r/2)>0$.
On $3a<u<4a$ the active $Q$ phase has no later anisotropic factor
and gives $\cos(\pi Q/2)>0$. Above $4a$ every remaining factor
acts by positive congruence. There is a positive matrix pole at
zero and no continuum for $0<u<a$.

Cauchy's integral formula for this finite product now represents it
by that pole and the computed positive discontinuities. The large
circle vanishes because the multiplier is $O(|z|^{-1/2})$. Small
circles about negative thresholds vanish as well: at each interior
threshold the worst growth exponent is bounded by half the sum
of the two adjacent tail norms, strictly below one. At $-a$ the
factor $\sqrt{z+a}$ supplies an extra positive half-power.
Thus there is no omitted endpoint measure or polynomial.
The scalar spinor weight is $m$ times the positive matrix vector
weight, so the ordinary free charged spectral inequality holds.

This example passes both the free charged positivity test and the
rooted Euclidean ultraviolet bound with genuinely noncommuting
gluings. Its closed pair cut still obeys
\eqref{eq:nextmatrix-trace-uv} and undercounts its light pair
once the first physical heavy reference has support. The distinction
identifies the remaining issue in this example as the charged
current and its inclusive trace even after free positivity is
recovered.

\subsection{A shell preparation with an analytic soft photon response}
\label{sec:nextsoft-uv}
Two integrations cause the double logarithm in
\eqref{eq:combined-harmonic-double-log-uv}. A bounded exit time
can eliminate the logarithm from the normal mass integral alone.
For example, exit from the full four-dimensional sphere gives the
constant shift $v=m^2$ in the same squared-mass coordinates.
For the original virtuality mark the signed moments in that
calculation now obey $A_j(x)= -3m^4x^{-j}+O(x^{-j-1})$ at large
$x$, with fixed $m$. Set $s=Lq$. In the range
$m^2\ll x\ll m^2/s$ the term linear in $h-1$ gives
$-4s^2/[(n+2)x]$ in the radial integrand. The quadratic terms,
the ends of the range and the shifted-propagator remainders are
$O(s^2)$ after integration. Thus the dimensionless scalar response
still has a term $-4s^2\log(1/s)/(n+2)$. The second logarithm in
the original obstruction measures the growth of the charged mark
with virtuality.

\paragraph{A change of the charged mark from the same A4 shell}
Keep the original cylinder exit, its charged normal transport and
its harmonic shift law $d\pi(v)=2m^4{\bf1}_{v\ge m^2}dv/v^3$.
Use an independently sampled projected shell direction for the
one-photon mark. The squared transverse fraction
$u=|P_2\omega|^2$ is uniform on $[0,1]$ as above. For dimensionless virtuality $z$, the loss of normal heat transfer
over a fraction $u$, divided by that fraction, is
$(1-e^{-uz})/u$. Its shell expectation is
\[
 \ell(z)=\mathbb E\frac{1-e^{-uz}}{u}
        =\int_0^1\int_0^z e^{-ut}\,dt\,du .
\]
The removable value at $u=0$ is $z$. This entire shell average has
$\ell(0)=0$, $\ell'(0)=1$ and $\ell'(z)=(1-e^{-z})/z$.
On the positive real axis $\ell(z)=\log z+O(1)$ at large $z$.
Use it as the virtuality term in the transverse charged vertex:
\begin{equation}
\begin{aligned}
 Z&=X[1+X/(4m^2)+\ell(\Delta/m^2)],\\
 h&=\chi_n(L\sqrt Z),\qquad \Delta=x+X/4+m^2 .
\end{aligned}
\label{eq:nextsoft-mark-uv}
\end{equation}
Use $n=248$, $Lm=1$, $x=P^2$ and $X=Q^2$ on the Euclidean
plane. The even series for $\chi_n$ removes the apparent square-root
branch. The vertex is still $h\gamma^\mu+(1-h)Q^\mu\slashed Q/X$.
Its Ward contraction is $\slashed Q$, and the light on-shell form
factor is unchanged because $\Delta=0$ there. Give this vertex the
same straight Wilson chord and evaluate the pole-free scalar contact
\eqref{eq:gcmone-contact-Gram-uv} with the new $h$ and its actual
derivative $\partial_x h/X$. The normal determinant, the new charge
mark and both contact derivatives now belong to one specified action.
The choice of the averaged heat loss as the off-shell coupling is
additional charged representation data. It adds no fitted scale.

\paragraph{Analyticity up to the light threshold}
Consider the photon Hessian at leading charged-loop order, with
$Q=q e_4$ and transverse polarization $e_1$. Write $U=m^2+v$,
$D_U=x+U$ and let $\delta_\pi f=f(m^2)-\int f(m^2+v)d\pi(v)$
act on the mass argument $U$. Its mean shift is
$\int v\,d\pi(v)=2m^2$. The free spin traces after the constant
chiral mass rotation involve
\[
\begin{aligned}
 J_U(P)&=D_U^{-1}-2P_1^2 D_U^{-2},\\
 B_U(P,q)&=\frac{-4(D_U-q^2/4-2P_1^2)}
         {(D_U+q^2/4)^2-q^2 P_4^2}.
\end{aligned}
\]
Define $b_U=B_U(P,q)-B_U(P,0)$ and let $j_U$ be the average of
$J_U(P-tQ)-J_U(P)$ over $-1/2\le t\le1/2$.
With $\mathcal R=T(D)+m^2T(D)^2$, $D=x+m^2$, as in
\eqref{eq:combined-harmonic-Maxwell-uv}, the completed response is
\begin{equation}
\begin{aligned}
 \frac{16\pi^2}{e^2}\Pi_{11}(q e_4)
  =\int_0^\infty x\,dx\,
  \Big\langle &h^2\delta_\pi b_U+8h\delta_\pi j_U\\
   &-\tfrac43 q^2\mathcal R h^2\delta_\pi D_U^{-1}
  \Big\rangle .
\end{aligned}
\label{eq:nextsoft-rooted-uv}
\end{equation}
The brackets denote the normalized four-dimensional angular mean.
Here one initially takes real $q>0$ and forms the signed mass
difference within the integrand. In that difference the unshifted
bubble and chord cancel the derivative contact by
$(x^2/D_U)'=1-U^2/D_U^2$. The boundary term at infinity vanishes
because $x^2\delta_\pi D_U^{-1}$ is bounded by the mean shift
and $h\to0$. Translating the integrated-by-parts chord then gives
$j_U$. These operations can also be performed on a ball followed
by exhaustion. The surface error for the signed chord is $O(|h|)$
at large radius. At fixed nonzero real $q$, the ball transform
and its derivative decay as $(1+L\sqrt Z)^{-(n+1)/2}$,
and $\sqrt Z$ grows proportionally to $\sqrt{\log x}$.
After the mass difference, the original signed bubble has an
integrable radial bound $C|h|^2/x$ at infinity. The chord and
derivative contact are integrable there using
$|\partial_xh|\le C_q x^{-1}(\log x)^{-(n+3)/4}$,
uniformly under the bounded real chord shifts at large $x$.
This argument never integrates the two unsuppressed reference
determinants individually.

For complex $q$ in a compact set, splitting the integral for
$\ell$ at $(1+|\Delta|/m^2)^{-1}$ gives
$|\ell(\Delta/m^2)|\le C[1+\log(1+x/m^2)]$ for large real $x$.
The compact ball characteristic function satisfies
$|\chi_n(\zeta)|\le e^{|\operatorname{Im}\zeta|}$.
Consequently $|h|\le C_\epsilon(1+x/m^2)^\epsilon$ for every
$\epsilon>0$, uniformly on that compact set. At large $x$ the
denominators are uniformly comparable to their unshifted values.
Evenness of $B_U$ and averaging $t$ with $-t$ in $j_U$ give
\[
 |\partial_U b_U|+|\partial_U j_U|
       \le C |q|^2/(x+m^2)^3
\]
by Taylor's formula in $q$ applied to the displayed rational
kernels. Also $|\mathcal R\,\delta_\pi D_U^{-1}|
\le C/(x+m^2)^3$ there. Hence after division by $q^2$ the hard
radial part of \eqref{eq:nextsoft-rooted-uv} has the integrable
majorant $C x(1+x/m^2)^{2\epsilon}/(x+m^2)^3$ for
$0<\epsilon<1/2$. For bounded $x$ no free denominator vanishes
on real loop momenta if $|\operatorname{Im}q|<2m$: each charged
energy is shifted by at most $|\operatorname{Im}q|/2<m$.
On compact subsets of this strip the same Taylor argument at zero
and ordinary boundedness elsewhere apply uniformly in $U\ge m^2$.
The quotient is therefore even and holomorphic throughout that strip.
In particular its Taylor series in $q^2$ converges for $|q^2|<4m^2$.
At $q=0$ the mark is one. Taking that limit in the rooted formula
recovers exactly \eqref{eq:combined-harmonic-Maxwell-uv}, including
$\lim\Pi_{11}/q^2=0.00707707\,e^2$ after stripping the common
external readout. All soft derivatives now exist for this coefficient.

\paragraph{The spectral test of the repaired action}
The same bounds control the hard part locally at every finite
continued energy. All free-pole pinches can be enclosed in a
bounded loop-momentum region there. In that region the marks are
entire. A single-propagator contact permits a local contour
deformation and has no two-particle cut. The bubble has the
usual two massive residues, with both charged marks evaluated at
them. This local argument keeps the hard signed contribution
together. For a finite energy interval only a bounded set of
normal shifts can have an open pinch; two-body phase space
is locally bounded and vanishes at the upper shift endpoint.
On the mass shell of a reference with shift $v$, one has
$X=-S$ and $\Delta=-v$. The chiral rotation preserves the vector
mark; its longitudinal component drops out of the transverse pair
amplitude. In the usual unit Dirac normalization the reference
pair weight for $S>4(m^2+v)$ is
\begin{equation}
\begin{aligned}
 C_v(S)&=\frac{1+2(m^2+v)/S}{12\pi^2}
        \sqrt{1-4(m^2+v)/S}\ \chi_n(k_v)^2,\\
 k_v&=L\sqrt{S[S/(4m^2)-1-\ell(-v/m^2)]}.
\end{aligned}
\label{eq:nextsoft-cut-uv}
\end{equation}
The normal logarithm therefore gives
$C_{\rm cl}(S)= C_0(S)-\int_{m^2}^{S/4-m^2} C_v(S)d\pi(v)$
with zero integral for $S\le8m^2$. For every $S>8m^2$ the heavy
contribution is strictly positive: $k_v$ is strictly increasing
in $v$ in the open channel and the zeros of $\chi_n$ are isolated.
At each positive light form-factor zero $j_{n/2,k}$, the light
pair vanishes. These energies obey
\[
  S_k/m^2=2\big(1+\sqrt{1+j_{n/2,k}^2}\big) .
\]
For $n=248$ the first is $S_1/m^2\simeq268.92969$. At this and
every subsequent zero the closed spectral density is strictly
negative, and by continuity remains negative on an open interval.
Thus the same action passes the complete one-loop soft analyticity
test and fails a positive-spectrum test at infinitely many finite
energies. The calculation establishes neither graphwise finiteness
at arbitrary electromagnetic order for the new vertex nor all-order
control of its higher photon contacts.

\paragraph{Preserving the zeros of the measured light-pair form factor}
The negative intervals in \eqref{eq:nextsoft-cut-uv} suggest a direct
change in the charged action. At a light-pair zero, positivity of the
closed cut in a scalar stopping law requires the heavy amplitudes to
vanish for almost every open shift. Independent shell displacements
allow one of the two marks to be common to every charged mass.

Use another measured normal phase from the same uniform squared
fraction $u$. Its even phase loss gives
\[
 \lambda(z)=\int_0^1\frac{1-\cos(uz)}{u}\,du .
\]
This is entire and even. For real $z$ it is nonnegative, and
$\lambda'(z)=(1-\cos z)/z$ makes it strictly increasing
on $z>0$. On the positive axis
$\lambda(z)=\log z+O(1)$ by integrating the cosine term by parts
away from zero. Replace the single ball mark by the product of two
independent ball marks in the charged kernel:
\begin{equation}
\begin{aligned}
 h_\times(P,Q)&=h_0(X)\,
    \chi_n\!\left(\tfrac{X}{m^2}\sqrt{\lambda(\Delta/m^2)}\right),\\
 h_0(X)&=\chi_n\!\left(L\sqrt{X[1+X/(4m^2)]}\right).
\end{aligned}
\label{eq:nextsoft-product-uv}
\end{equation}
The second factor is entire in its squared argument. For the same
realization use $h_\times$ in the electric vertex, its straight chord
and its differentiated pole-free scalar contact, all before varying
the determinant. This is a different charged vertex with the same
original light on-shell value, normal exit and charge normalization.

The radial identity \eqref{eq:nextsoft-rooted-uv} applies verbatim.
At fixed nonzero real $q$ the second ball argument grows as
$\sqrt{\log x}$, and its $x$ derivative is
$O_q(1/[x\sqrt{\log x}])$. On real loop momentum and bounded complex
$q$, $\operatorname{Im}(\Delta/m^2)$ is bounded; splitting the
integral at $u=(1+|\Delta|/m^2)^{-1}$ again bounds $|\lambda|$
by $C[1+\log(1+x/m^2)]$. This gives the previous subpower
majorant for the signed hard integral. All its soft derivatives
therefore exist, with the same Maxwell coefficient. The compact
pinch calculation gives the complete one-loop transverse density.
Writing $U=S/4-m^2$, let $\rho_v(S)$ be the unit Dirac factor in
\eqref{eq:nextsoft-cut-uv} before multiplication by the ball
transform. For $U>0$ set
\[
 g_v(S)=\chi_n\!\left(\tfrac{S}{m^2}\sqrt{\lambda(v/m^2)}\right),
 \qquad
 D(S)=\int_{m^2}^{U}\rho_v(S)g_v(S)^2\,d\pi(v),
\]
with $D=0$ for $U\le m^2$. The two densities are exactly
\begin{equation}
\begin{aligned}
 C_{\rm pair}(S)&=h_0(-S)^2\rho_0(S),\\
 C_{\rm cl}(S)&=h_0(-S)^2[\rho_0(S)-D(S)] .
\end{aligned}
\label{eq:nextsoft-positive-uv}
\end{equation}
Every $g_v$ has a real ball argument and $|g_v|\le1$. At fixed
$S$, the Dirac factor decreases as the reference mass increases:
the derivative of $(1+2y)\sqrt{1-4y}$ is
$-12y/\sqrt{1-4y}$ for $0<y<1/4$. It follows that
\[
 b_a(U)C_{\rm pair}(S)\le C_{\rm cl}(S)\le C_{\rm pair}(S).
\]
The repaired closed density is nonnegative on the whole ray, including
every light-pair zero. For $U>m^2$ the integral $D$ is strictly
positive because $\lambda$ is strictly increasing on the open shift
interval. Away from the common zeros the optical deficit consequently
persists at every such energy. It becomes small at high energy:
$\lambda(v/m^2)\ge\lambda(1)>0$ uniformly on the heavy support.
The large-argument ball bound gives
$D(S)/\rho_0(S)=O((m^2/S)^{n+1})$ as $S\to\infty$.
This action thus repairs the sign of the measured two-current cut
and the analytic Maxwell limit simultaneously. Exact inclusive
sewing still fails, even though its relative error tends to zero.

\paragraph{Arbitrary photon multiplicity on one charged loop}
A further heat-transfer choice controls every derivative of the
charged vertex at large real loop momentum. In
\eqref{eq:nextsoft-product-uv} replace $\lambda(z)$ by
$\lambda_H(z)=\ell(z^2)/2$, and set $w_H(y)=e^{-y^2}$ in the
pole-free scalar coefficient $T(D)=[1-w_H(D/m^2)]/D$.
Use $\ell$ from the first heat construction and the same mass and length as before. The unit-depth heat
profile in $w_H$ and the projected heat average in $\lambda_H$
specify this comparison. Both functions are entire.
On the real axis $\lambda_H$ is even and nonnegative, is strictly
increasing for positive argument, and grows as $\log z+O(1)$.
The preceding soft and finite-energy cut arguments still apply;
the positive Maxwell coefficient now uses $w_H$ in
\eqref{eq:combined-harmonic-Maxwell-uv}. In particular the
nonnegative density bounds and strict off-node optical deficit
persist for the modified action.

Fix any finite number $N\ge2$ of nonzero real Euclidean momenta
summing to zero and transverse external polarizations. The
connected coefficient of their single charged determinant loop
converges absolutely with its signed normal mass difference kept
in the integrand. The constants can be chosen uniformly for
compact sets of marks bounded away from zero. To see this, put
$p=|P|$ in a hard region larger than all external translations
and use $l=1+\log(1+p/m)$.
For $\alpha=(n+1)/4>1$, at each fixed nonzero transfer the
transverse one-mark multiplier obeys
\[
 |\partial_P^d h_\times|\le C_d p^{-d}l^{-\alpha}
 \qquad(d\ge0).
\]
Indeed $\lambda_H'(z)=(1-e^{-z^2})/z$, whose subsequent real
derivatives have bounds $C_j z^{-j-1}$ at positive infinity.
The ball argument grows as $\sqrt{\log p}$, each positive-order
momentum derivative costs $p^{-1}$ per derivative, and every
fixed derivative of the ball transform has the same large-argument
Bessel bound. The chord translations are bounded. The heat choice
in $T$ gives $|\partial_P^d \mathcal R|\le C_dp^{-2-d}$.
It follows that a vector insertion carrying $r\ge1$ photons
costs at most $C p^{1-r}l^{-\alpha}$. At $r=1$ this is the full
transverse vertex; for $r>1$ the constant part of its correction
is killed by a chord derivative. An insertion from the scalar
contact with $r\ge2$ costs at most $C p^{2-r}l^{-\alpha}$.
Indeed $x\partial_x h/X$ and all its momentum derivatives have
the same logarithmic gain, and each additional straight-chord
mark is a momentum derivative integrated over a bounded interval.

For a word with $k$ reference propagators use
$(K_0+i\tau\gamma^5)^{-1}$ on every charged segment, with the
same twist $\tau$ throughout. The signed word is the word at
zero twist minus the mean of the two words at
$\tau=\pm\sqrt v$. Taylor's formula with the
linear terms canceled bounds this difference by
$C_k v p^{-k-2}$ times the vertex bounds. Each twist derivative
inserts a bounded Clifford matrix and one additional inverse, and
the unmarked inverses are bounded by $C/p$ uniformly in the real
twist. The finite mean of $v$ makes the normal integral absolute
in this paired representation. For $k\ge3$ the power alone is
integrable in four dimensions. For $k=2$ the worst radial bound
is $Cp^{-1}l^{-\alpha}$ and is integrable. With $k=1$ a vector
insertion has $r=N\ge2$ and hence contributes at least one
further inverse power as well as the logarithmic gain. For a
scalar insertion the spin trace of the signed inverse is exactly
$4m\,\delta_\pi D_U^{-1}=O(p^{-4})$, again accompanied by the
logarithmic gain. All bounded momenta are regular because of the
mass gap. This proves the fixed-external-momentum claim for all
photon multiplicities of one and the same determinant. Bounds
uniform on simultaneous internal photon circuits, including the
regions in which some transfers approach zero, are still required
for dynamical Maxwell sewing.

There is one simultaneous limit in this proof for which the dependence
on small transfer can also be controlled. In units $m=L=1$, set
$s_i=|Q_i|^2$. In the hard radial region let $t=\log x$
with $x=P^2$ large, and take bounded soft transfers.
From $x|\lambda_H'(\Delta)|\le C$ and the radial ball bound on
$\chi'_n(y)/y$ one obtains, for some $\gamma>2$,
\[
 \frac{x|\partial_x h_i|}{s_i}
    \le C s_i(1+s_i\sqrt t)^{-\gamma}.
\]
The apparent quotient at $s_i=0$ is removable. In a marginal
word with two scalar contacts the two dangerous derivatives of
this type share a single charged radial integral. Their
logarithmic cost is uniformly bounded since
\[
 s_i s_j\int_0^\infty
 \frac{dt}{(1+s_i\sqrt t)^\gamma(1+s_j\sqrt t)^\gamma}
 \le\int_0^\infty \frac{2y\,dy}{(1+y)^{2\gamma}}.
\]
This is Cauchy--Schwarz with $y=s_i\sqrt t$ and $y=s_j\sqrt t$
in the respective squared norms. Thus these two insertions do
not lose their joint hard-momentum bound when their transfers
become small together. Uniform sewing of arbitrary hard photon
circulations remains a separate requirement.

\subsection{The physical spectral weight required by a charged pole}
\label{sec:next-outer-uv}
A cut is measured on the timelike ray, whereas electric charge is measured
at zero transfer. Analytic continuation ties the two measurements together.
The minimally gauged harmonic source of \eqref{eq:fix-threshold-cuts-uv}
makes this constraint quantitative. In this subsection the spin-area
commutator is absent. Use the Minkowski invariant $S$ on the fixed
sector plane, put $\tau=S/(4m^2)$, and strip the common photon readout.
Let $r(\tau)$ be the ratio of the closed two-current density to that
of a unit Dirac pair. In the units $a=m^2$, write $U=\tau-1$ and
$x=v/m^2$ for the dimensionless squared-mass shift. The two-body
phase space and spin sum give
\begin{equation}
\begin{aligned}
 r(\tau)&=1 &&(1<\tau\le2),\\
 r(\tau)&=1-\int_1^U \frac{2}{x^3}
    \left(1+\frac{x}{2\tau+1}\right)
    \sqrt{1-\frac{x}{U}}\,dx &&(\tau>2).
\end{aligned}
\label{eq:next-outer-ratio-uv}
\end{equation}
The shifted Dirac weight decreases strictly with squared mass at fixed
energy. In particular $U^{-2}\le r<1$ above the heavy threshold. The
lower bound is the probability of a mass shift above $U$.
There is also a useful upper bound $r=O((1+\log U)/U^2)$ at large $U$.
Indeed put $y=x/U$ and $A=U/(2\tau+1)$. Rationalizing
$1-(1+Ay)\sqrt{1-y}$ gives a numerator
$y[(1-2A)+(2A-A^2)y+A^2y^2]$ and a denominator at least one.
Now $1-2A=3/(2\tau+1)$. Integrating the resulting upper bound
against $2\,dx/x^3$ and adding the tail $U^{-2}$ proves the estimate.
It also gives a stable way to evaluate the ratio numerically.

\paragraph{A necessary bound for a regular charged vertex}
Use the disk coordinate $z$ in \eqref{eq:next-moment-mean-uv},
so $z(0)=0$ and its upper and lower rims have
$\tau_\theta=\sec^2(\theta/2)$. Define
\[
 \log\alpha=\frac{1}{2\pi}\int_0^\pi\log r(\tau_\theta)\,d\theta .
\]
The preceding bounds ensure that this integral exists. For unit charge,
$G_E(0)=1$, the logarithmic mean inequality of
\eqref{eq:next-moment-mean-uv} now gives
\begin{equation}
 1+8|G_M(0)|^2\ \le\ (\sqrt3+\sqrt2)^2\alpha^2.
 \label{eq:next-outer-moment-uv}
\end{equation}
The hypotheses are a pole-free charged vertex in the cut plane and the
bounded analytic or Hardy-class growth that permits the disk mean
inequality for its weighted form factors. A direct quadrature gives
$\alpha\simeq0.42025236$, hence $|G_M(0)|\lesssim0.305835$.
The exclusion of the Dirac value $G_M(0)=1$ is also analytic.
For $\tau\ge4$ the shifted-to-light ratio in the integral is at least
$1/\sqrt2$ on $1\le x\le3/2$ and $1/\sqrt3$ on $3/2\le x\le2$.
Those intervals carry weights $5/9$ and $7/36$. Their contributions
sum to more than $1/2$, so $r<1/2$ on the last third of the angular
interval. Thus $\alpha\le2^{-1/6}$, already small enough to exclude
the Dirac value. A change of the static magnetic moment or of the
closed spectral strength has observable consequences even at zero
momentum transfer.

\paragraph{Two explicit solutions of the boundary magnitude problem}
Consider the narrower comparison with one scalar dressing $F$, so
$G_E=G_M=F$. Following the disk formulation of unitarity bounds~\cite{PhysicalGluingAbbas}, for a positive even rim weight $w$ with integrable
logarithm, its zero-free outer function is
\[
 O_w(z)=\exp\!\left\{\frac{1}{4\pi}\int_{-\pi}^{\pi}
    \frac{e^{i\theta}+z}{e^{i\theta}-z}\log w(\theta)\,d\theta\right\}.
\]
The real part of the exponent is the Poisson integral of
$\tfrac12\log w$. Hence $|O_w|^2=w$ almost everywhere on the rim,
and $O_r(0)=\alpha$ with $r$ evaluated at $\tau_\theta$.
For a bounded analytic $F$ with $|F|^2\le r$, the mean inequality
forces $|F(0)|\le\alpha$. Two different changes can restore the
unit charge.

First keep the continuum fixed and allow a pole. With
$B_b(z)=(b-z)/(1-bz)$, the function $F=O_r/B_\alpha$ has unit
charge, the exact modulus $\sqrt r$ on both rims, and one simple
pole at
\[
 S_p=\frac{16m^2\alpha}{(1+\alpha)^2}\simeq3.33349\,m^2 .
\]
If $b$ is the position of a single positive real pole in any such
meromorphic scalar solution and $B_b F$ is bounded analytic, then
$b=|(B_bF)(0)|\le\alpha$. The displayed solution attains this
bound. A physical pole would require an actual subthreshold state
and its one-particle contribution in the full current spectrum.
Multiplying an old photon Hessian by $F^2$ generally gives a double
pole at $S_p$. A positive spectral measure allows an atom there
with a simple pole only.

Alternatively keep the first interval and choose a smooth threshold
onset for additional continuum strength. This is an explicit
comparison law. It is not selected by the original boundary action.
Put $g(\tau)=(1-2/\tau)_+^{3/2}$ and let
$\langle g\rangle=\pi^{-1}\int_0^\pi g(\tau_\theta)\,d\theta$.
The threshold and exponent are fixed in this comparison. Set
\begin{equation}
 \widetilde r=r\,e^{2\beta g},\qquad
 \beta=-\log\alpha/\langle g\rangle .
 \label{eq:next-outer-repair-uv}
\end{equation}
Here $\langle g\rangle=\Gamma(5/4)/(2\sqrt\pi\,\Gamma(7/4))$
since $g(\tau_\theta)=(-\cos\theta)_+^{3/2}$.
The outer function $O_{\widetilde r}$ has $F(0)=1$, no first-sheet pole,
and equal electric and magnetic form factors. It leaves the light
interval's pair probability unchanged. Its continuum
$\rho_D\widetilde r$ is everywhere positive; its unsubtracted
dispersion integral $\int_{4m^2}^\infty
\rho_D(s)\widetilde r(s/(4m^2))/(s+Q^2)\,ds$ converges for $Q^2\ge0$.
Indeed $g$ is bounded and $\rho_D$ tends to a constant.
The extra continuum supplies the positive logarithmic mean required by
charge normalization, with no fitted constant once $g$ has been chosen.

These formulas solve a two-current inverse problem. On the free Dirac
state space one may multiply the current matrix element by the spacelike
value $F(S)$ and use conjugate boundary values on the future and past
pair shells. This gives a Hermitian conserved current on finite-particle
test states and the positive two-point spectral weight just calculated.
In a local theory with a single elastic eigenchannel and time-reversal
invariance its boundary values must also satisfy the elastic phase relation~\cite{PhysicalGluingWatson},
$F_+=S_{\rm el}F_-$, with $S_{\rm el}$ the elastic scattering phase.
Trivial elastic scattering on an open part of the unchanged interval
would make $F$ real there. Since $|F|=1$, it would be constant there.
Reflection across that interval and analytic uniqueness would make
it constant throughout the cut domain, contrary to the high-energy
weight. The outer solution is regular on the interior of this interval.
Thus a local physical implementation of this scalar repair with these
elastic hypotheses needs charged scattering already below $8m^2$,
where the measured pair rates still agree exactly. If other channels
are open, a corresponding coupled-channel unitarity relation must be
derived instead. The positive two-current construction alone supplies
neither this scattering matrix nor the higher photon contacts.

\paragraph{Applying the charge bound to the positive product source}
The product source in \eqref{eq:nextsoft-positive-uv} allows the same
test after stripping its common light form factor $h_0(-S)$.
The relative continuum weight is $r_\times=1-D/\rho_0$ on the open ray,
including the continuous limiting ratio at a common zero. Write
$\alpha_\times$ for the preceding logarithmic mean with $r_\times$.
Since $r_\times\le1$ and is strictly less than one on a set of
positive angular measure, $\alpha_\times<1$. A bounded analytic
relative scalar multiplier with unit charge cannot fit even this
positive spectrum at all energies. Numerically
$\alpha_\times\simeq0.94850219$. At $S=12m^2$ the missing fraction
is $D/\rho_0\simeq0.39817605$, even though the ratio tends back to
one at large energy. In the single-pole scalar construction the
extremizing pole would lie at $S_p\simeq3.99721m^2$, just below
the light threshold. This location is a test of this particular
completion and does not predict a bound state of the physical theory.

Allow separate relative electric and magnetic form factors satisfying
the same bounded analytic hypotheses after the common factor is
removed. Equation~\eqref{eq:next-outer-moment-uv} applies with
$\alpha_\times$. Its numerical upper bound on the magnetic moment is
$0.99409$ in unit Dirac normalization. This test concerns
analyticity of the relative form factors across the common
zeros and their growth. A general charged amplitude with poles in
those ratios must account for them explicitly. A model preserving
the inferred light charge, magnetic moment and crossed analyticity
has to change the closed spectral strength or the singular
structure of the current.

\paragraph{An integer certificate for the logarithmic mean}
The strict inequality $-\log\alpha_\times>0.048$ can also be
checked with rational arithmetic. Use dimensionless $s=S/m^2$.
The angular measure in the defining logarithmic mean is
$ds/[2\pi s\sqrt{s/4-1}]$. Keep the cells $s_i=8+i/5$,
$0\le i<260$, and $v_j=1+j/20$, $0\le j<100$. In each cell put
$b=s_i+1/5$, $u=s_i/4-1$, $v=v_j$, $w=v+1/20$.
Use the integer scale $D=10^{12}$ and
\[
 B(w)=\sum_{k=1}^9\frac{(-1)^{k+1}w^{2k}}{2k(2k)!},
 \qquad Q=\lceil D b^2 B(w)/4\rceil .
\]
The alternating tail gives $\lambda(w)\le B(w)$
on this interval. Discard cells with $w\ge u$ or $4Q>45^2D$.
On the other cells $2\sqrt{Q/D}$ upper-bounds the Bessel
argument: $\lambda'(v)=(1-\cos v)/v\ge0$.
The normalized ball transform $\chi_{248}$ is positive and
decreasing up to 45. Indeed $(xJ_\nu')'=(\nu^2/x-x)J_\nu$
gives positivity of $J_{124}$ and $J_{125}$ there from their
positive leading terms; the derivative of $x^{-124}J_{124}(x)$
is $-x^{-124}J_{125}(x)$.

Start $l=h=D$. For $k=21,20,\ldots,1$ use the simultaneous update
\[
 (l,h)\longleftarrow
 \left(D-\left\lceil\frac{Qh}{Dk(124+k)}\right\rceil,
       D-\left\lfloor\frac{Ql}{Dk(124+k)}\right\rfloor\right).
\]
The endpoints enclose $D$ times the nested partial sum of the
$ {}_0F_1(125;-Q/D)$ series. The neglected tail after odd
degree 21 begins positive and decreases in absolute value.
Consequently $c=\max(0,l)/D$ is a lower bound for the positive
ball transform throughout the retained cell.
For nonnegative rational $p/q$ put
$R(p/q)=\operatorname{isqrt}(\lfloor D^2p/q\rfloor)$,
the integer part of $D\sqrt{p/q}$. Form the integer sum at fixed $i$
\[
 d_i=\sum_j\left\lfloor(v^{-2}-w^{-2})
       (1+2v/(b+2))R(1-w/u)c^2\right\rfloor .
\]
Here the mass density $2/v^3$ was integrated exactly on each cell.
Then $d_i/D$ bounds the heavy-pair deficit from below throughout
the energy cell. Using six positive terms of $-\log(1-d_i/D)$ gives
\[
 \begin{aligned}
 e_i&=\sum_{k=1}^6\left\lfloor
          \frac{d_i^k}{kD^{k-1}}\right\rfloor,\\
 A&=\sum_i\left\lfloor\frac{e_iR(1/(b/4-1))}{5bD}\right\rfloor
     =302551431196 .
 \end{aligned}
\]
The outer weight is decreasing, and all discarded contributions
are nonnegative. Since $\pi<22/7$, the certified lower bound is
$7A/(44D)=529465004593/11000000000000>0.048$.
For the charge test, $\sqrt3<1.7321$ and $\sqrt2<1.4143$
by squaring. Therefore $(\sqrt3+\sqrt2)/3<1.0488
<1+0.048+0.048^2/2<e^{0.048}$.

\subsection{Testing the photon used to sew charged loops}
\label{sec:next-physical-photon-uv}
The propagator used to join two charged responses is itself a
physical observable. Its massive spectral weight supplies a test
even when the charge insertions mix and their signs vary. The
following argument uses Euclidean momenta on one reconstructed
sector plane and the positive spectral representation of its
physical transverse photon.

\paragraph{A test using the spectrum of the photon that joins the loops}
Suppose the normal gluing changes the relative overlap of a current
probe and an internal Maxwell line. The photon read out after that
gluing supplies a further test, even for a signed charged Hessian.
Work on one reconstructed plane at Euclidean $X=|q|^2>0$ and in a
fixed transverse physical polarization. Conditional on a geometric
record $\xi$, let the free covariance on the photon ports be
$D_{0,\xi}=C_\xi/X$ with $C_\xi\succeq0$, and let $K_\xi(X)$
be the complete first quadratic insertion including its contacts.
Write
\[
 B_\xi=D_{0,\xi}^{1/2}K_\xi D_{0,\xi}^{1/2},\qquad
 w_\xi=C_\xi^{1/2}u_\xi ,
\]
where $u_\xi$ is the physical photon readout. Both $C_\xi$ and
$u_\xi$ are fixed as $X$ varies. Allow noncommuting port matrices and
assume $M_4=\mathbb E\|w_\xi\|^4<\infty$. For a conditional
Gaussian this last quantity is proportional to the fourth moment
of the free readout at fixed $X$. Suppose the same record goes
through both occurrences of the insertion in the ring. Its absolute
radial norm per unit volume is proportional to
$\int X A(X)\,dX$ with $A(X)=\mathbb E\|B_\xi(X)\|_{\rm HS}^2$.
These hypotheses specify a gluing rule that can be tested on finite
port projections before any continuum limit.

First hold the normalized record law, its covariance numerator
and the physical readout fixed at this order. Differentiating the photon
inverse then gives $X\delta D(X)=-\mathbb E\langle w_\xi,B_\xi
w_\xi\rangle$. Therefore
\begin{equation}
 |X\delta D(X)|^2\le M_4 A(X).
 \label{eq:sew-dressed-photon-ring-uv}
\end{equation}
Here $\delta D$ is the change in the covariance of the physical
transverse photon. The inequality follows from the operator norm
bound by the Hilbert--Schmidt norm and Cauchy--Schwarz in $\xi$.
No sign or lower bound on $\langle u_\xi,K_\xi u_\xi\rangle$
enters it. Suppose the first change in that photon's spectrum is
\[
 \delta D(X)=\frac{\delta Z}{X}
        +\int_{(0,\infty)}\frac{d\sigma(s)}{X+s},
 \qquad d\sigma\ge0 ,
\]
where the photon remains massless and the integral exists at positive
$X$. This is a hypothesis on the leading massive spectral weight.
It holds for a leading sum of squared photon-to-state matrix
elements under the usual spectral construction when the free photon
has only its massless atom and the pole residue has the indicated
derivative. If $\int^\infty X A(X)dX$ were finite, $X\delta D(X)$
would approach zero along some unbounded sequence. Its spectral
formula is nondecreasing, so monotone convergence gives
$\delta Z+\int d\sigma=0$. For a compact band
$I=[s_0,s_1]\subset(0,\infty)$ this implies
\[
 |X\delta D(X)|=\int\frac{s\,d\sigma(s)}{X+s}
       \ge\frac{s_0\sigma(I)}{X+s_1}.
\]
By \eqref{eq:sew-dressed-photon-ring-uv}, any occupied band
would already force logarithmic divergence of the ring.
The conclusion concerns the photon that actually propagates
between the loops. Its positive pair continuum alone can test the
gluing even if the averaged charged response has no established
spectral representation.

If the source also varies the normalized record law and the
readout, the derivative needs two additional terms. Suppose that
the law has an integrable logarithmic derivative $s_\xi$ of mean
zero and the readout has derivative $\dot u_\xi$, both independent
of $X$, that $C_\xi$ stays fixed, and that differentiating these
expectations is justified. The two contributions to $X\delta D(X)$
sum to the constant
\[
 \kappa_\partial=\mathbb E\!\left[
 s_\xi\|w_\xi\|^2+
 2\operatorname{Re}\langle\dot u_\xi,C_\xi u_\xi\rangle\right].
\]
The inequality then bounds $|X\delta D(X)-\kappa_\partial|^2$.
The same proof tests the massive band with the residue variation
$\delta Z-\kappa_\partial$. A momentum-dependent boundary
response requires its own estimate, derived from the same source.

One way to establish the needed massive photon weight directly
from an action is the Maxwell equation on physical transverse
matrix elements. If its first variation is $\partial_\mu F^{\mu\nu}
=e J^\nu$ to leading electric order for the same readout, a state
of squared mass $s>0$ has transverse photon matrix element of
magnitude $|e\langle0|J_T|n\rangle|/s$ at that order. In the
convention where a transverse current weight is $s\,d\rho_J(s)$,
the coefficient of $e^2$ in the corresponding photon weight is
$d\rho_J(s)/s$. Thus one nonzero light-pair band already supplies
the $d\sigma$ needed in the test. If the boundary variation adds a
current in this equation, its coherent matrix element must be
included before squaring. This identifies which native
Euler derivative is relevant to the obstruction.

There is also an exact version for a quadratic photon completion
with coupling $g>0$. Assume $I+gB_\xi(X)\succeq\delta I$ with
constant $\delta>0$ and use the same physical readout. Put
$C_*=\mathbb E\|w_\xi\|^2$ and
$X D_g(X)=\mathbb E\langle w_\xi,(I+gB_\xi)^{-1}w_\xi\rangle$.
The resolvent identity gives
\[
 |C_*-XD_g(X)|^2\le g^2\delta^{-2} M_4 A(X).
\]
If the ring is finite and $D_g$ has a positive unsubtracted photon
spectral measure, its total mass must be $C_*$ by the same sequence
argument. A massive band would make $C_*-XD_g(X)$ at least of order
$1/X$, contradicting the bound after radial integration. Only
$D_g=C_*/X$ is then possible in this polarization. This exact
conclusion applies to the specified quadratic completion. The
fixed-order test makes no assumption about summing an infinite
perturbative series.

\paragraph{Positive geometric sewing that passes every ring integral}
The hypotheses can be probed directly on the forward A4 shell.
With $U=|P_2\omega|^2$ as in the shell exit construction, take
its logarithmic radial depth $r=-\log U$. It has probability
measure $e^{-r}dr$ on $(0,\infty)$. Allow the amplitude of the
transverse photon in a record to scale by $r/\sqrt2$, so that
$c_r=r^2/2$ is its covariance coefficient and $\mathbb E c_r=1$.
Use the ordinary heat evolution of its curvature for normal
duration $rL^2$, where $L$ is the length already fixed by $G_N$.
In the transverse sector the conditional action is
\[
 \frac1{4c_r}\int F_A^2+
 \frac{g}{4}\int F_{A+b}\cdot
       e^{-rL^2(-\Delta)}F_{A+b}.
\]
Here $b$ probes the same curvature as the dynamical photon in
the insertion. Use the normalized central trace over fixed
records for the loop coefficients and the same prepared probability
to read out the conditional covariances. This is a prescribed
comparison representation. The free readout has Maxwell covariance,
while its fourth moment records the common random amplitude.
A global random $r$ also correlates different separated probes.
No charged pair has been postulated in this quadratic example.

Write $u=L^2X$. The measured scalar insertion after removing $X$
is $\mathbb E e^{-ru}=(1+u)^{-1}$, which has a positive mass
atom at $L^{-2}$. In contrast each closed insertion carries
$c_r e^{-ru}$. For an arbitrary positive integer $p$ the exact
$p$-insertion ring before radial integration is
\begin{equation}
 \mathbb E (c_r e^{-ru})^p
       =\frac{(2p)!}{2^p(1+pu)^{2p+1}} .
 \label{eq:sew-log-depth-rings-uv}
\end{equation}
This follows by integrating $r^{2p}e^{-(1+pu)r}$. The measure
$X\,dX$ is integrable against the result at both ends, including
$p=1$. Every ring coefficient of this quadratic action is
therefore absolutely finite without a momentum cutoff.

Its readout photon nevertheless fails a physical test for every
$g>0$. Exactly,
\[
 XD_g(X)=\mathbb E\frac{c_r}{1+g c_r e^{-ru}},
 \qquad 0<1-XD_g(X)\le\frac{6g}{(1+u)^5}.
\]
Also $\lim_{X\downarrow0}XD_g(X)=\mathbb E[c_r/(1+gc_r)]<1$.
A positive photon measure of total mass one with this smaller
massless residue would have positive mass on some compact
positive-mass interval. Its deficit $1-XD_g(X)$ would then
be bounded below by a positive multiple of $1/X$. The
contradiction uses Euclidean data on the positive axis throughout.
The positive measured atom and even convergence of every ring
can therefore coexist with a failure of spectral positivity
for the same propagating photon.

\paragraph{What a physical gluing law has to supply}
An invertible transport of the charged coordinates has a computable
Jacobian. Keeping it restores the two tested light-pair discontinuities
over the full ray and restores the Maxwell logarithm. Dropping it gives
a different quantum source. A positive matrix exit can change the open
pole residues and can even preserve the free charged spectrum, while
its trace still subtracts physical pair weight. The product vertex in
\eqref{eq:nextsoft-product-uv} changes what happens inside a
charged propagation. It gives an analytic one-loop Maxwell response
and a nonnegative two-current density at every energy; its explicit
deficit in \eqref{eq:nextsoft-positive-uv} still violates the
inclusive bound above $8m^2$ away from the common zeros.

These calculations narrow the missing physical input to a joint
choice. The charged action must give the measured electric
transition and any other states that account for its timelike
discontinuity. The same action must prescribe how two such
responses exchange a dynamical photon. Equations
\eqref{eq:next-outer-moment-uv} and
\eqref{eq:sew-dressed-photon-ring-uv} test these operations with
different data, the static moment and the spectral measure of the
propagating photon. In a theory with ordinary common-record Maxwell
sewing and nonzero leading massive photon weight, the latter test
already forces a logarithm. Proposing new boundary variables only
changes this conclusion if they change an assumption of that
calculation in the actual generating functional. The selection of
such a charged and photon boundary by $\Omega_\alpha$, and the
hard-subgraph estimates for the resulting source, remain open.

\subsection{Which parts of the physical finiteness claim can hold together?}
\label{sec:joint2-map-uv}
A pair of charged particles can be counted in two ways. One can
prepare the pair and square its transition amplitude, or open the
corresponding intermediate states in a closed current response. A
physical photon must also reproduce that response at spacelike momentum
and propagate between two such currents in a larger graph. Each step
uses more information about the action. Three explicit comparisons
separate their requirements.

First, the earlier product vertex has a positive timelike density, yet
its spacelike response vanishes at a sequence of momenta where its own
spectral integral is strictly positive. Bounding the heat increment
at large Euclidean transfer controls all photon loop integrations, but
a strictly positive inclusive deficit survives. Second, joint heat
preparation of the normal and charged states supplies the exact
one-loop Dirac cuts. Its measured current fails a reflection test at
separated points. Third, an interaction in the neutral pair channel
reproduces the full observed two-point weight and a positive photon
spectrum. Resolving the modes in the action shows precisely which
individual vacuum rings remain divergent.

All calculations start on a specified locked plane $\bin=(s,b_s)$.
The real coefficient of the directed curvature action in
\eqref{eq:uvpeA-curvature} is used after projection on its carrier. Its aligned
geometric and matter directions are $u_{g,\bin}=\sigma_\epsilon
\hat p_{0,\bin}$ and $u_{m,\bin}=\sigma_s u_{g,\bin}$. The boundary
displacements on the matter carrier are $\pm i\sigma_{s\epsilon}
\varepsilon\hat p_{0,\bin}$ according to the chosen ordering. The
radial Euclidean variable $X=|q_E|^2$ below enters after the declared
continuation on that plane. The GT-locked coordinate readout and the
Minkowski clock readout remain distinct. In these comparisons the GTL
and velocity-shell distributions and Newton scale are fixed. Each choice
of charged boundary state or photon coupling is identified where it
enters the generating functional.

\subsection{A spacelike test of the positive timelike cut}
\label{sec:joint2-disp-uv}
The product vertex in \eqref{eq:nextsoft-product-uv} repairs the sign
of the one-loop cut. A physical current must also reproduce its
Euclidean response from that spectrum. Work on the same locked plane,
strip the common external photon readout, and use $m=L^{-1}=1$ for
numerical values. Let $X=q^2>0$,
$F_E(X)=\Pi_{11}(q e_4)/(e^2X)$ and write $C=C_{\rm cl}$ for the
density in \eqref{eq:nextsoft-positive-uv}. Its positive Stieltjes
transform is
\[
 F_C(X)=\int_4^\infty\frac{C(s)}{s+X}\,ds .
\]
Here $s$ is timelike squared mass in units of $m^2$. Both $C$ and the
open light-pair density $C_{\rm pair}$ are integrable: the timelike
ball argument is $\sqrt{s(s/4-1)}$ and the squared ball transform
is $O(s^{-249})$ at infinity. Any ordinary dispersion representation
with this density and finitely many local subtractions can therefore be
written as $F_C(X)+p(X)$ for a polynomial $p$.

\paragraph{Zeros on the spacelike axis}
Let $j_{124,k}$ be the $k$th positive zero of $J_{124}$. Unlike
the timelike zeros, the following are at real Euclidean momentum:
\begin{equation}
 X_k=2\big(\sqrt{1+j_{124,k}^2}-1\big),\qquad F_E(X_k)=0 .
 \label{eq:joint2-zero-disp-uv}
\end{equation}
Indeed $h_0(X_k)=0$ identically in the charged momentum. Every term
of the complete rooted Hessian \eqref{eq:nextsoft-rooted-uv} has at
least one factor $h_0$ at positive transfer. This includes the
one-propagator contact. On the other hand $F_C(X_k)>0$, as can
already be seen from $5\le s\le6$, below the first heavy threshold.
There the density is at least
$c_*=(497/500)^2/(9\sqrt5\,\pi^2)$ by the ball-coordinate variance
and the free phase space. Thus $F_C(X_k)\ge c_* /(X_k+6)$.

Since $F_C(X_k)\to0$, any polynomial satisfying
$p(X_k)=-F_C(X_k)$ for every $k$ would have to vanish identically.
There is also a finite test for each fixed subtraction degree. At
$N+1$ distinct such roots the $N$th divided difference of $F_E$
vanishes. If $N>\deg p$, a dispersive response instead obeys
\[
 (-1)^N[X_1,\ldots,X_{N+1}]F_C
   =\int_4^\infty\frac{C(s)\,ds}
             {\prod_{j=1}^{N+1}(s+X_j)}>0 .
\]
Consequently the full Euclidean Hessian cannot represent its own
positive cut with finitely many local contacts. Both functions are
analytic near zero and along the positive axis. If all their
Taylor coefficients beyond some order were equal, their difference
would continue as a polynomial there. Thus infinitely many soft
coefficients differ. These conclusions hold for the phase and heat
versions of the product vertex.

The finite-node test also covers proposed additional positive
channels whose total weight at high energy is unknown. Suppose
$d\rho$ contains nonzero positive mass $M$ on $[a,b]$ and admits
an $N$-subtracted dispersion at $X_*>0$, with subtraction polynomial
of degree less than $N$. The required integrability is only
$\int(1+s)^{-N-1}d\rho(s)<\infty$. For distinct positive spacelike
nodes $Y_0,\ldots,Y_N$ its response obeys
\[
 (-1)^N[Y_0,\ldots,Y_N]F
 =\int\frac{d\rho(s)}{\prod_{j=0}^N(s+Y_j)}
 \ \ge\ \frac{M}{\prod_{j=0}^N(b+Y_j)} .
\]
Indeed the first $N$ Taylor terms of $(s+X)^{-1}$ at $X_*$ are
polynomials killed pointwise by the divided difference; apply that
finite linear combination directly to the integrable subtracted
kernels. Choosing $N+1$ distinct $X_k$ forces a zero left side.
Thus no positive completion with any finite subtraction order and
the same spacelike response can conceal the observed massive band,
even if its total spectral weight is infinite. The charged Euclidean
kernel itself must change if the band and dispersion are retained.

\paragraph{The first two moments}
The discrepancy can be tested before reaching the first spacelike
zero. In the stated units write $C_{\rm pol}=16\pi^2 F_E(0)$.
For either product vertex $h(0,x)=1$ and
$\partial_Xh(0,x)=-1/[2(n+2)]$ independently of $x$.
The second factor of the product starts at order $X^2$.
The strip bounds for the rooted formula justify differentiating
its radial integral. At unit mark the shifted chord difference
integrates to zero at each transfer after signed mass pairing:
the integrand is a translate of $\delta_\pi J_U$ minus itself and
its boundary translation error vanishes at infinity.
The derivative of the two marks in the bubble and scalar contact
therefore contributes $-C_{\rm pol}/(n+2)$ to $16\pi^2F_E'(0)$.
For a single reference mass $U$ the angular average of the $X^2$
coefficient of $B_U(P,q)$ is
\[
 -\frac{30U^3+51U^2x+32Ux^2+8x^3}{24(U+x)^6}.
\]
Its $x\,dx$ integral is $-4/(15U)$. Taking the light-minus-heavy
difference, with $\mathbb E[v/(1+v)]=2(1-\log2)$, gives
\[
 -16\pi^2 F_E'(0)
 =\tfrac8{15}(1-\log2)+C_{\rm pol}/(n+2).
\]
For the phase product, direct quadrature gives $F_E(0)=0.00707707$
and $F_C(0)=0.01078752$. The negative slopes are $0.00106466$
and $0.00112786$, respectively. At $X_1\simeq264.92969$ the
dispersive value is $0.00054840$. For the heat product the
corresponding zero-transfer values are $0.00232698$ and
$0.01133709$, with negative slopes $0.00104566$ and $0.00115760$.
The analytic zero test is independent of these quadratures.

\paragraph{A lower bound from the open charged pair alone}
For any of the pole-free profiles in this contact family with
$0\le w\le1$, \eqref{eq:combined-harmonic-Maxwell-uv} gives
$12\pi^2F_E(0)\le\int_0^\infty\mathcal B(x)dx
=\mathbb E\log(1+v)=1$, since $V\le1$. In contrast an
unsubtracted positive completion containing just the observed light
pair already requires
\begin{equation}
 12\pi^2\int_4^\infty C_{\rm pair}(s)\frac{ds}{s}>1.08 .
 \label{eq:joint2-open-soft-uv}
\end{equation}
For an elementary check, partition $[5,26]$ at $7,10,14,18,22$.
For each interval $[a,b]$ bound its contribution from below by
\[
 \left(1+\frac2b\right)l_a
 \left(1-\frac{b(b-4)}{2000}\right)^2\frac{2(b-a)}{a+b}.
\]
Here $l_a=(44,65,77,84,88,90)/100$ in increasing order, each
bounded by $\sqrt{1-4/a}$. This uses
$\chi_n(y)\ge1-y^2/500$ and
$\log(b/a)\ge2(b-a)/(a+b)$. The six rational terms sum to more
than $1.0814$. Thus even arbitrary retuning of $w$ within its
positive unit range cannot retain the soft coefficient and supply
the positive pair dispersion in this unsubtracted comparison.

\paragraph{What a minimal dispersive repair would cost}
There are two separate changes in the minimal positive current
response with exactly the measured open-pair density. Write
$F_{\rm pair}(X)=\int C_{\rm pair}(s)ds/(s+X)$. Then
\[
 F_{\rm pair}-F_E=(F_{\rm pair}-F_C)+(F_C-F_E).
\]
The first term has the explicitly missing positive heavy-reference
weight from \eqref{eq:nextsoft-positive-uv}. The second term
compensates a discrepancy with the known discontinuity held fixed.
It cannot be supplied by a finite local polynomial. This construction
gives a positive two-current dispersion; its derivation as a
boundary variation of the GCM action remains to be established.

A stronger high-momentum obstruction applies even to marks without
spacelike zeros. Work in the same units and assume the rooted identity
\eqref{eq:nextsoft-rooted-uv} with $h(P,Q)=f(X)g(P,Q)$. The
factors $f,g$ are real and satisfy $|f|,|g|\le1$ on the Euclidean slice. Assume the
native scalar coefficient $\mathcal R$ is nonnegative there, as for
the heat and phase profiles used here. Uniformly for $X\ge1$,
\begin{equation}
\begin{aligned}
 X F_E(X)&\le C(|f|+f^2)\log(2+X),\\
 |F_E(X)|&\le C\left[(|f|+f^2)\frac{\log(2+X)}{X}+f^2\right] .
\end{aligned}
\label{eq:joint2disp-hard-uv}
\end{equation}
To see the logarithm, set $D_\pm=|P\pm Q/2|^2+U$. In the bubble the
absolute numerator before the factor $4$ is bounded by
$\sqrt{D_+D_-}$: reflect the first component of $P+Q/2$ and apply
Cauchy--Schwarz to the vectors augmented by $\sqrt U$. The mass
derivative of $B_U$ is consequently bounded by
$C(D_+^{-2}+D_-^{-2})$, by weighted arithmetic--geometric mean.
The mass derivative of the chord integrand at $P-tQ$ is bounded by
$C/(1+|P-tQ|^2)^2$. Integrating the unsigned mass derivatives of
both rooted differences on $|P|\le2\sqrt X$ thus costs at most
$C\log(2+X)$ by translation of the ball and $\int v\,d\pi=2$.
On $x>4X$ every momentum shift along the transfer segment has length
at most $|P|/4$. Two transfer derivatives of the mass-differentiated
kernels and evenness give the bound $CX/(x+1)^3$ for each rooted
mass difference; its $x\,dx$ integral costs $O(1)$.
Finally $\mathcal R\,\delta_\pi D_U^{-1}$ is nonnegative and
bounded by $C/(1+x)^3$. The scalar contact therefore lowers
$XF_E$ and has absolute magnitude at most $CXf^2$. These three
bounds give \eqref{eq:joint2disp-hard-uv} without any integration
of individual reference determinants.

Positive spectral weight on a compact massive band would imply
$XF_E(X)\ge c>0$ at large $X$ for an unsubtracted Stieltjes
response. The first inequality would then require
$|f(X)|\ge c'/\log(2+X)$ for every sufficiently large $X$.
Even a common mark with no finite zero and arbitrarily weak
algebraic decay fails this test. The same lower bound excludes
every radial mark with $\int^\infty X|f(X)|^2dX<\infty$.
In that class a sequence $X_k\to\infty$ has $f(X_k)\to0$; the
second inequality then forces $F_E(X_k)\to0$. When the candidate
positive spectral transform exists at zero it tends to zero at
infinity by dominated convergence, so no nonzero finite local
polynomial can bridge their difference. For the two product
actions the common factor is $o(X^{-1/2})$ by the ball-transform
bound. The absolute estimate gives $XF_E(X)\to0$.

By monotone convergence the required repair therefore has
$X(F_{\rm pair}-F_E)\to\int C_{\rm pair}(s)ds>0$. In the
stated units the limit is approximately $0.20579$ and
$F_{\rm pair}(0)\simeq0.01364334$. Sewing two copies of this
repair by ordinary transverse Maxwell lines entails the radial
norm $\int^\infty X|F_{\rm pair}-F_E|^2dX$, which diverges
logarithmically. Any further positive spectral channels only raise
the minimal dispersive weight. The same two-current calculation
therefore tests the charged cut, the spacelike response and the
cost of its ordinary photon sewing.

\subsection{A bounded heat increment and fully sewn Euclidean photon graphs}
\label{sec:joint2-graphs-uv}
The heat vertex of \eqref{eq:nextsoft-product-uv} leaves a specific
difficulty at large photon transfer. Repeated derivatives of the
charged midpoint act on a ball transform with an argument proportional
to that transfer squared. A single unit of geometric heat evolution
gives a bounded coefficient for all those derivatives on the
Euclidean plane. It also permits a separate, quantitative test of the
continued vertex.

Use units $m=L=1$ in the estimates. Write $s=|q|^2$,
$x=|P|^2$, $\Delta=1+x+s/4$ and $W(t)=\chi_n(\sqrt t)$,
$n=248$. Keep the harmonic normal mass law, $\lambda_H(z)=\ell(z^2)/2$,
$w_H(y)=e^{-y^2}$ and the transverse photon readout
$c_L(q)=\chi_n(L|q|)$. Change the charged mark to
\begin{equation}
 \begin{aligned}
 A(s)&=1-e^{-s},\\
 h(P,q)&=h_0(s)W\big(A(s)^2\lambda_H(\Delta)\big).
 \end{aligned}
 \label{eq:joint2-cap-uv}
\end{equation}
Thus $A(s)$ is the difference between the initial and final heat
amplitudes across a single normal interval $L^2$ for the transfer
$|q|^2$. The choice to put this increment in the second ball mark
is a prescription for the charged representation. The original
GTL localization, velocity shell, pole boundary and $G_N$ scale are
the same. Insert $h$ in the electric vertex, its straight Wilson
chord and the complete quadratic contact with $\partial_x h/s$,
before integrating the charged determinant. The quotient is regular
at zero since $A(s)^2/s$ is entire with a removable zero there.
The Ward contraction and the light on-shell mark $h_0$ are preserved.

\paragraph{A uniform bound before joining photon lines}
At every fixed perturbative order the connected, source-dependent
Euclidean graphs with closed charged components and external
photons converge absolutely in the complete parity-even determinant
representation. This statement includes the continuous four-momenta
of the massless transverse photons and simultaneous hard subgraphs.
The free vacuum normalization is omitted. For a complete charged
component with $N\ge2$ incidences, strip its external readouts.
On conserved momenta there is a bound of the form
\begin{equation}
 \begin{aligned}
 |V_N(q_1,\ldots,q_N)|&\le C_N
    \prod_{i=1}^N\langle q_i\rangle^{6}\min(1,|q_i|),\\
 \langle q\rangle&=(1+|q|^2)^{1/2}.
 \end{aligned}
 \label{eq:joint2-component-uv}
\end{equation}
The corresponding bound also controls the integral of the absolute
gauge-completed densities used below. The exponent six is sufficient
at every multiplicity; the constants may depend on $N$.

To prove the bound, take the diagonal position-space trace density
of the normal resolvent before imposing momentum conservation at
its root. For $G_B=\gamma^5 K_B$, the operator
$(G_B^2+u)^{-1}$ is the difference of its two twisted first-order
resolvents divided by $2i\sqrt u$. Its integrated weight is
$b(u)du/2$, with $\int b(u)du=2$ in these units. A word of $k$
vertices consequently contains $k+1$ free first-order inverses,
and the difference of the two twists gains one more inverse by
telescoping the product. An inverse at momentum $p_j$ has norm
$(1+u+|p_j|^2)^{-1/2}$. This proves uniformity as $u\downarrow0$
and $u\to\infty$. All shifts $p_j-p$, including midpoint and chord
shifts, are bounded by a fixed multiple of $R=1+\sum_i|q_i|$.
It is enough to treat $|p|>C_N R$ and its complement.

On the exterior put $t=\log |p|^2$, taking $C_N$ large enough
that $t\ge1$. Since $0\le A(s)\le\min(s,1)$, $d$ midpoint
derivatives of a vector mark cost at most $C_d |p|^{-d}$.
Indeed the positive-axis derivatives of $\lambda_H$ of order
$j\ge1$ are bounded by $C_j\Delta^{-j}$, and the derivatives of
$W$ are bounded there by its ball integral or Bessel series.
The same symbol estimate holds for $x\,\partial_x h/s$.
Without derivatives it has the stronger bound
\[
 \left|x\,\frac{\partial_x h}{s}\right|
 \le C A(s)(1+A(s)\sqrt t)^{-\gamma},\qquad \gamma>2 ,
\]
because $x\lambda'_H(\Delta)$ is bounded and
$|W'(y^2)|\le C(1+y)^{-\gamma}$. The other scalar symbol
$\mathcal R(x)$ and its $d$-th midpoint derivatives have order
$|p|^{-2-d}$. Each photon on a Wilson chord differentiates the
corresponding symbol once. A vector insertion with $r$ photons
therefore costs $C_r|p|^{1-r}$, and a scalar insertion costs
$C_r|p|^{2-r}$ times one $|q_i|$ for each of its two curvature
owners. These estimates are uniform in all transfers.

For $N\ge3$ and $k\ge3$ the resolvent difference is already
$O(|p|^{-k-2})$. For $k=1$ a vector vertex has at least two
midpoint derivatives. A scalar vertex has at least one; its odd-twist
spin trace divided by $\sqrt u$ is bounded by
$C m/(D_0D_1)$ where $D_j=m^2+u+|p_j|^2$.
For $k=2$ the only mixed word without a midpoint derivative has
one scalar and one vector. Its corresponding spin numerator is
bounded by $C m\sum_j |p_j|$ over $D_0D_1D_2$: the odd twist leaves
a trace of one vector gamma matrix with two remaining numerators,
and only a term with one scalar mass and one slashed momentum
survives. These statements apply also when the total transfer at
the position-space root is nonzero. All these hard radial integrals
have at least $|p|^{-5}$ decay prior to the four-dimensional measure,
after extracting the possible curvature momenta. The sole marginal
word is $N=4$, with two undifferentiated quadratic scalar vertices.
Either a factor $\mathcal R$ supplies extra decay or two factors of
$x\partial_x h/s$ occur. The resulting radial integral is bounded
uniformly by Cauchy--Schwarz, since
\[
 \int_0^\infty A^2(1+A\sqrt t)^{-2\gamma}\,dt
 \le \int_0^\infty 2y(1+y)^{-2\gamma}\,dy
 \qquad (0\le A\le1).
\]
In this marginal word all four photon momenta already appear as
explicit curls.

For clarity, the small-transfer factors in
\eqref{eq:joint2-component-uv} must hold independently at all
soft incidences. For $N\ge3$ a constant connection at one selected
incidence differentiates the common charged momentum in the lower
diagonal word. This follows by conjugating $K_B$ by a constant
gauge gradient, equally on both sides of every Wilson chord.
The $(N-1)$-marked lower integrand is $O(|p|^{-4})$ or better on
a large sphere, by the same count including the scalar spin trace.
Its outward flux vanishes as the sphere expands. Thus the integrated
diagonal coefficient vanishes when any one of its transfers is zero,
also away from total momentum conservation. Apply a finite difference
between the plane wave and its constant value at the root at every
soft incidence, with its transverse polarization held fixed along the
radial interpolation. This is an exact representation of that coefficient.
The marginal four-curl word vanishes in every newly constant argument
individually and already has all its soft factors. In each other word,
one soft-transfer derivative per soft incidence costs at most a fixed
power of $1+\log |p|$ on the exterior, times the displayed integrable
power. In fact a derivative through a routing shift is a midpoint or
propagator derivative, and a derivative through $A(s)^2\lambda_H$
at a soft owner gives at most $C(1+\log |p|)$. The apparent $1/s$
in the scalar vertex occurs only through the smooth $A(s)^2/s$.
This proves the product of soft factors under the exterior integral.

On $|p|\le C_N R$ the mass gap bounds every free inverse and its
derivatives. The displayed heat symbols and their midpoint
derivatives are uniformly bounded. Soft-transfer derivatives cost
at most a fixed power of $1+\log(2+R)$, which is bounded by $C_N R$.
The four-volume costs $C_N R^4$ and the explicit hard curvatures
cost at most one $\langle q_i\rangle$ per incidence. Since
$R\le C_N\prod_i\langle q_i\rangle$, this gives precisely the
growth permitted by \eqref{eq:joint2-component-uv}.
When $N=2$ use the exact, conserving, gauge-completed Hessian
\eqref{eq:nextsoft-rooted-uv}. Its two-point quotient by $|q|^2$
is bounded near zero, including under its absolute rooted integral.
For hard $q$, the mass-differenced exterior kernel costs
$C|q|^2|p|^{-6}$ and its interior has the preceding volume and
curvature bound, since real Euclidean denominators stay massive.
Vertices with a single photon vanish in the vacuum current by
charge conjugation or by a constant momentum shift at zero transfer.

\paragraph{The joint photon integral}
Each internal photon between closed charged components, including a
self-loop, now has the scalar majorant
\[
 C\,|c_L(q)|^2\langle q\rangle^{12}
       \frac{\min(1,|q|)^2}{|q|^2}.
\]
It is bounded and integrable on $\mathbb R^4$ since
$|c_L(q)|^2=O((1+|q|)^{-249})$. Contract each charged component
to a vertex, choose a spanning tree, bound its line factors by
their suprema and integrate each independent chord by its $L^1$
norm. This proves absolute convergence in the gauge-completed
representation, including overlapping charged and photon hard
regions. Each external soft photon retains its own $|q|$.
The proper transverse two-point correction is consequently
$O(|q|^2)$ at every finite order. The argument excludes a
momentum-independent photon mass term; it does not reconstruct
the interacting photon spectrum. For open charged chains the
free positive Stieltjes measure integrates the uniform bound
$(m^2+u)^{-1}$ from the two or more massive inverses of each
marked chain. The same vertex bounds then give the hard-momentum
estimate with the original external charged infrared prescription.

\paragraph{Open charged lines tested as Euclidean distributions}
The massless photon can also be retained on open charged graphs
when the unamputated, off-shell external momenta are tested with
Schwartz packets. The open Stieltjes bound above is uniform in the
charged momentum at every point of a marked chain. The open
vertices need no soft zero in the following argument. Use the
momentum deltas on each open chain to integrate its outgoing charged
momentum. Its incoming momentum remains integrable against the
packet; no photon conservation constraint is left at that chain.
One can bound a joint Schwartz packet by a product of separately
rapidly decreasing weights on the incoming charged and external
photon momenta, uniformly in the eliminated momenta.

A photon joining two closed vertices has a majorant in
$L^1\cap L^\infty$. One between an open and a closed vertex has a
majorant in $L^1\cap L^2$, since it has at most a $1/|q|$
singularity at zero. One joining two open vertices has an
$L^1$ majorant with at most a $1/|q|^2$ singularity. The shell
readout controls infinity for each of these bounds. Every
open--open line momentum can be integrated directly. On each
connected cluster of the remaining closed vertices choose an
internal spanning tree, bound its edge factors by their suprema,
and use the closed-vertex constraints to eliminate their momenta.
One conservation law remains for the boundary of this cluster.
For at least two boundary edges ending on open chains, integrate
one freely against the translated dependent edge by
Cauchy--Schwarz and their two $L^2$ bounds. The other boundary
edges and all photons outside the internal spanning tree that
join closed vertices are then integrable by $L^1$. This argument
also treats an isolated closed vertex. For a single edge ending
on an open chain its momentum is fixed by the total external
photon transfer to that cluster. Integrating any one of those
external photon packets controls the remaining translated
$L^1$ factor. With no such external photons a singly attached
cluster vanishes: the vectorlike charged determinant is even
under charge conjugation, while contraction of its photon legs
inside the cluster removes an even number of legs. Its evenness
follows by transposing the oppositely oriented charged chords;
the symmetric scalar two-curvature contact is even as well.
Connectedness excludes a closed cluster with no attachment to the open part.
This proves absolute finiteness of these Euclidean coefficients as
tested distributions. Amputation, sharp exceptional momenta and
the charged on-shell infrared limit require their own analysis.

\paragraph{How close the same action comes to its own pair weight}
At every finite continued energy the new mark has the same bounded
complex-transfer estimates as in \eqref{eq:nextsoft-rooted-uv}:
on a fixed compact transfer set $A(s)$ is bounded, so the growth
from the ball transform still costs only $e^{C\sqrt{\log x}}$ on
the hard real charged momentum. The one-loop transverse response
is analytic before the light threshold. Its closed and open cuts
are \eqref{eq:nextsoft-positive-uv} with
\[
 g_v(S)=\chi_{248}\!\left[
   (e^{S/m^2}-1)\sqrt{\lambda_H(v/m^2)}\right].
\]
At a common light zero use the relative deficit $D/\rho_0$ after
stripping the light multiplier. For every $S>8m^2$ the heavy
average is strictly positive because its argument is strictly
increasing with the shift and the Bessel zeros are isolated.
Nevertheless the discrepancy is uniformly tiny. On $v\ge m^2$
and $S\ge8m^2$ the argument exceeds $1875$, since
$\ell(1)\ge 1-1/4+1/18-1/96$. For integer order,
$|J_{124}(z)|\le1$ on the real axis by its angular integral.
Monotonicity of the unit Dirac phase space with mass now gives
\begin{equation}
 0<\frac{D(S)}{\rho_0(S)}
 \le \left(\frac{2^{124}124!}{1875^{124}}\right)^2
 <10^{-322},\qquad S>8m^2 .
 \label{eq:joint2-tiny-uv}
\end{equation}
At $S=12m^2$ the argument exceeds $10^5$ and the same bound
is below $10^{-750}$. Below the heavy threshold the deficit
is exactly zero. Thus one action can have fully sewn,
source-dependent Euclidean photon graphs and nearly saturate
its two-current pair cut over the whole energy ray. Exact
inclusive unitarity still fails above threshold away from the
common zeros. The complex growth of this heat readout is also
uncontrolled at infinity by the finite-strip argument. No
unrestricted physical finiteness theorem follows by rounding
the strictly positive missing weight to zero.

\paragraph{Two finite spacelike transfers test the propagating photon}
The cut accuracy in \eqref{eq:joint2-tiny-uv} does not measure
physical consistency of the photon that joins two charged loops.
There is a direct test using only positive Euclidean squared momenta.
Keep the transverse source $B=c_LP_T A+P_L A$ used in this
construction and let $f(X)=X\delta D_T(X)$ denote the leading
correction to the physical transverse photon covariance. Since
both external probes of the charged Hessian pass through this
linear connection, $f$ contains $c_L(q)^2$, $X=|q|^2$.
The stripped two-current coefficient is finite there. Consequently
$f(X_1)=f(X_2)=0$ at
\[
 X_i=L^{-2}j_{124,i}^2,\qquad i=1,2 .
\]
Here $j_{124,i}$ are the first two positive zeros of $J_{124}$.
At zero, Dyson differentiation of the transverse covariance gives
$f(0)=-\lim_{q\to0}\Pi_{11}(qe_4)/q^2<0$ in the Maxwell
normalization used above.

If the free photon has only its unit massless atom, the leading
spectral law being tested takes the form
$f(X)=\delta Z+\int_{(0,\infty)}X\,d\sigma(s)/(X+s)$,
$d\sigma\ge0$. For any two distinct $0<X_1<X_2$ a nonzero
measure gives a strictly positive difference, and a zero measure
makes $f$ constant. Neither possibility reproduces these three
Euclidean values. For a leading continuum supported above
$s_0=4m^2$ the comparison is quantitative:
\[
 f(X_2)-f(X_1)\ge
 \frac{s_0(X_2-X_1)}{X_1(X_2+s_0)}
                  [f(X_1)-f(0)].
\]
This is the pointwise ratio of the two spectral differences,
bounded by its value at the lowest allowed mass. In the stated
carrier units the multiplier is approximately
$2.28303\times10^{-5}$. The strict two-node contradiction
already follows without any gap estimate. It applies to a
linearly read out transverse photon with the same connection at
both external probes. Replacing that readout by a nonvanishing
profile removes these exact zeros and still leaves the general
sewn-photon spectral test to be satisfied.

\paragraph{Why an increasingly accurate cut need not give an exact physical limit}
A sequence of integer heat depths gives a test of this near agreement.
Set $A_j(s)=(1-e^{-js})/j$ for $j=1,2,\ldots$, with $s=L^2|q|^2$.
This is a sequence of different charged preparations. On $s\ge0$
one has $A_j\le\min(s,1)$ and $0\le A_j'\le1$. The proof of
\eqref{eq:joint2-component-uv} is uniform in $j$. In particular
$|(A_j^2/s)'|\le3$ and there is at most one soft-transfer derivative
at each owner. The hard-momentum estimate needs any number of
midpoint derivatives, and these do not differentiate $A_j$.
The one-loop Maxwell coefficient is the same positive heat-profile
value for every $j$. On the pair cut the smallest heavy ball argument
is $(e^{8j}-1)\sqrt{\lambda_H(1)}/j$. The upper bound in
\eqref{eq:joint2-tiny-uv} therefore tends uniformly to zero,
even after the common light factor is stripped.

Consider now the complete Hessian
\eqref{eq:nextsoft-rooted-uv}. At any fixed real $q\ne0$,
$A_j(q^2)\to0$, so $h\to h_0(q^2)$ throughout its charged integral.
Since the real-axis multiplier is bounded and the rooted integrand
has the uniform integrable mass-difference bound, the Euclidean
limit is obtained by inserting $h_0$ there. This limiting Hessian
is analytic in $q^2$ near zero and retains the nonzero Maxwell
coefficient, by the same rooted formula. Its bubble has the heavy
normal mass cut with marks $h_0(-S)$ independent of mass. The
shifted single-propagator terms have no pair pinch. Its continued
density is consequently
\[
 h_0(-S)^2\left[\rho_0(S)
                  -\int_{m^2}^{S/4-m^2}\rho_v(S)\,d\pi(v)\right]
 \qquad(S>8m^2).
\]
Continuing each finite-$j$ Hessian first gives instead a density
tending to $h_0(-S)^2\rho_0(S)$.

One can detect the failure of this interchange even below threshold.
In carrier units take $q=i\eta$ with $0<\eta<2m$. In the real
charged integration $\Delta=x+m^2-\eta^2/4>0$.
The ball factor is bounded there and tends to zero since
$|A_j(-L^2\eta^2)|\to\infty$ and $\lambda_H(\Delta/m^2)>0$.
The denominators stay away from every pinch on compact subthreshold
intervals. The same signed hard majorant permits dominated
convergence; the continued response at this fixed transfer tends to
zero. Thus the limiting prethreshold values cannot be the analytic
continuation of the spacelike limit with positive Maxwell coefficient.
Uniform Euclidean graph control and arbitrarily accurate
real-energy pair rates do not supply a single analytic limiting source.

\subsection{Joint heat preparation of a charged boundary state}
\label{sec:joint2-charge-uv}
A boundary record can specify a duration for the whole charged
state to evolve in the normal extension. This changes the action
before the fermionic integral. It gives an example in which the
closed determinant counts the full light pair at every finite
energy, the measured free charged propagator is positive, and the
Euclidean loop integrals still converge. The dispersive moments of
the current give a further independent test.

\paragraph{One positive boundary measure}
Use the aligned sector plane and the original oriented A4 measure.
Retain the Abelian connection $B=[c_LP_T+P_L]A$ of
\eqref{eq:upgrade-shell-pair-dirac-uv}, $K_B=m+d_B$,
$H_B=K_B^\dagger K_B=m^2-d_B^2$, and $Lm=1$. On each of
$j=1$ or $2$ independent A4 records let $u_i=|P_2\omega_i|$, the
radius of the two-dimensional projection of its uniform four-space
direction. Since $u_i^2$ is uniform on $[0,1]$, the duration
$\tau=L^2\max_i u_i$ obeys, with $p=2j$ and $a=L^2$,
\begin{equation}
 W_p(t)=\Pr(\tau\le t)=\min\{(t/a)^p,1\},\qquad t>0.
 \label{eq:joint2-duration-uv}
\end{equation}
All positive durations and all real momenta occur. The maximum
selects the time at which all these recorded projections have
arrived. This use of the A4 readout is part of the charged
representation. The sector label and directed pole displacements
are those of the original matter carrier.

On a two-dimensional Euclidean normal plane $N$, give the state the
joint quadratic generator $\mathcal L_B=-\Delta_N+H_B$ of its
normal gradient and charged kinetic energy. Prepare its boundary map
on a charged test spinor $f$ by
$J_{\tau,B}f=e^{-\tau\mathcal L_B/2}(\delta_0\otimes f)$.
The planar diagonal heat kernel is $(4\pi t)^{-1}$. In particular
$\mathbb E\|J_{\tau,B}f\|^2\le
\|f\|^2\mathbb E(1/\tau)/(4\pi)$ and
$\mathbb E(1/\tau)=p/[(p-1)a]$. This is a square-integrable
preparation over the shell. Its measured normal resolvent is
\begin{equation}
\begin{aligned}
 \Phi_p(H_B)&=2\pi\,\mathbb E
         [J_{\tau,B}^\dagger\mathcal L_B^{-1}J_{\tau,B}],\\
 2\Phi_p(z)&=\int_0^\infty W_p(t)e^{-tz}\,\frac{dt}{t}
             =\mathbb E E_1(\tau z),\qquad z>0 .
\end{aligned}
\label{eq:joint2-normal-uv}
\end{equation}
Here $E_1(w)=\int_1^\infty e^{-wv}dv/v$. The first equality on
the second line follows by integrating the joint heat semigroup
from its recorded duration. The preparation depends on the same
$H_B$ as the resolvent. Its photon variations are therefore part
of every determinant contact.

\paragraph{The charged observable and its graph bounds}
Use the polar charged bilinear
$\mathcal D_B=mK_BH_B^{-1/2}e^{-\Phi_p(H_B)}$.
In the Hermitian charged form set $G_B=\gamma_5K_B$ and
$M_B=mH_B^{-1/2}e^{-\Phi_p(H_B)}$. Then $G_B^2=H_B$,
$\gamma_5\mathcal D_B=G_BM_B$, and $M_B>0$. Choose the charged
observable at each external end to be $M_B^{1/2}$ times the
integrated field in this form. Its normalized two-source kernel,
at every background on a finite test, is exactly
\[
 M_B^{1/2}(G_BM_B)^{-1}M_B^{1/2}=G_B^{-1}.
\]
The fixed dual Clifford map converts it to the ordinary Dirac
open kernel $K_B^{-1}$. This is a choice of which charged state
is measured. The integrated Berezin coordinates and their
determinant are those of $\mathcal D_B$. No coordinate Jacobian
has been inserted. Computing the open functional first includes
the derivatives of both measured endpoints. Its free spectral
measure is consequently the positive Dirac pole of mass $m$.

Equivalently, on a finite charged test the normalized electromagnetic
integrand with Grassmann sources is
\[
\begin{aligned}
 \mathcal Z[B,\bar\eta,\eta]&=e^{-S_{\rm M}[A]}
        \frac{\det(G_BM_B)}{\det(G_0M_0)}\,e^{\langle\bar\eta,\mathcal S_B\eta\rangle},\\
 \mathcal S_B&=M_B^{1/2}(G_BM_B)^{-1}M_B^{1/2}.
\end{aligned}
\]
The Maxwell factor is shown explicitly. Photon integration in a fixed
Abelian gauge and division by the free photon normalization follow. It specifies both the determinant and the
open observable, including the dependence of their boundary factors on
the same connection. The directed GCM monomial supplies the shell
measure; assigning its durations to this joint heat generator and its
external states to these two endpoint maps specifies the representation.

The parity-even charged effective action from the same Gaussian
is $\operatorname{Tr}\Phi_p(H_B)$, with the charged bilinear
measured in units of $m$. Proposition~\ref{prop:ward-normal-graphs-uv}
applies to its closed vertices: use the same rooted heat-kernel
identity at positive $t$ and multiply the entire vertex, including
all its contacts, by $W_p(t)dt/(2t)$. For $N\ge2$ insertions the
joint short-time bound is proportional to
$W_p(t)t^{N-3}dt$ after integration over the heat simplex. Every closed
photon incidence retains its $|q c_L(q)|$. Open chains at fixed
external endpoints use the massive Dirac inverse and have no
independent charged circulation. The transverse photon
spanning-tree estimate thus gives absolute convergence of all the
complete rooted Euclidean electromagnetic graphs at each finite
order with the stated charged infrared convention. It also gives
a perturbatively massless transverse photon as before. For a
connected vacuum graph with at least one photon exchange, remove
the overall volume and apply the same chord bound to the
contracted closed vertices. For $p=4$ the free charged
parity-even vacuum density itself is finite: its integral is
proportional to $\int_0^\infty W_4(t)e^{-m^2t}dt/t^3$.
This last observation concerns the charged determinant per
volume in the normalized free Maxwell measure.

\paragraph{Exact light-pair cuts on the entire ray}
The observed amplitudes for one or two photons to create the
light pair are those of $K_B$. For the normalized Dirac
two-current scalar on the represented connection, Feynman
parametrization of its complete one-loop heat trace gives, with
$X=Q^2$ and $b_x=x(1-x)$,
\begin{equation}
 \Pi_p(X)=\frac{e^2}{2\pi^2}\int_0^1 b_x\,dx
     \int_0^\infty W_p(t)e^{-t(m^2+b_x X)}\,\frac{dt}{t}.
 \label{eq:joint2-maxwell-uv}
\end{equation}
Use the convention in which this scalar decreases with $X>0$.
Its finite-energy cut is the full ordinary Dirac density. Indeed
$E_1(w)+\gamma_{\rm E}+\log w$ is entire, and $0<\tau\le a$
with $\mathbb E|\log(\tau/a)|<\infty$. Thus $2\Phi_p(z)$
differs from $-\log(az)$ by an entire function of $z$.
For $N\ge3$ photon marks, compare the complete one-charged-loop
vertex at this source with the minimally coupled Dirac vertex on $B$.
Their difference is supported on $0<t<a$ with weight $W_p-1$.
To see why it is entire in the external momenta, keep all ordered
contacts at each fixed $t>0$ and use the diagonal-root heat estimate.
On a compact complex momentum set the routing shifts $P$ are bounded.
A segment of duration $v$ gives
\[
 |e^{-v(k+P)^2}|=e^{v|\operatorname{Im}P|^2}
                       e^{-v|k+\operatorname{Re}P|^2} .
\]
The segment lengths sum to $t$, and bounded imaginary shifts in the
orbital numerators change only the compact-set constant for $t\le a$.
Constant-connection Ward subtractions on the diagonal hold also for
these analytic vertices. The rooted bound on the difference is thus
$C_{N,K}e^{C_Kt}t^{N-3}dt$, integrable at zero after integration over the heat simplex,
uniformly on compact sets of external momenta. Holomorphic integration on compact sets gives an entire
difference on the momentum-conservation hyperplane. Thus every finite
normal cut of this one-charged-loop amplitude agrees with its Dirac
value. Odd vertices vanish by charge conjugation. For two marks the
preceding logarithm fixes the cut directly. Since the open measured
kernel at background $B$ is the Dirac inverse, the ordinary cut pair
amplitudes carry precisely the same external readouts. This includes
the two-current and forward four-current tests at every finite energy
on the chosen vectorial matter boundary. Intermediate states containing
additional dynamical photons require a separate calculation.

\paragraph{A current can have the right discontinuity and the wrong dispersion}
Equation~\eqref{eq:joint2-maxwell-uv} determines more than its cut.
For integers $k\ge1$ define its alternating derivatives
$M_{p,k}=(-1)^k\Pi_p^{(k)}(0)$. The corresponding dispersive
Dirac values are obtained with $W=1$ by differentiating the
once-subtracted Dirac polarization. Their differences are
\begin{equation}
\begin{aligned}
 M_{{\rm D},k}-M_{p,k}
 &=\frac{e^2}{2\pi^2}\int_0^1 b_x^{k+1}dx\\
 &\quad\times\int_0^a[1-(t/a)^p]t^{k-1}e^{-m^2t}\,dt>0 .
\end{aligned}
\label{eq:joint2-moments-uv}
\end{equation}
For $k=1$ the Dirac value is $e^2/(60\pi^2m^2)$, while its
fraction retained by the joint preparation is
$\mathbb E e^{-m^2\tau}$. It equals $2-4/e$ for $p=2$ and
$24-64/e$ for $p=4$, with $e$ in these two numbers the base of
the natural logarithm. Thus the normalizable two-record example
retains about $45.6\%$ of the first dispersive moment even though
it reproduces the full pair cut. No finite polynomial in $X$ can
correct all these derivatives. More generally, for an integrable
signed short-time change supported in $[0,a]$, matching all
sufficiently high derivatives makes its corresponding signed
moment measure zero: multiply it by a sufficient power of $t$
and by $e^{-m^2t}$ and use density of polynomials on the compact
interval. A change supported exactly at zero would be a local
contact prescription. For ordinary weights this argument forces
$W=1$ almost everywhere if that whole cut is to have its ordinary
finite-subtraction dispersion.

There is a separate test on the actual transverse $A$ readout.
Its Euclidean scalar is $c_L(\sqrt X)^2\Pi_p(X)$, with $c_L$
here denoting the radial ball transform. It vanishes at the
infinitely many positive Bessel zeros and is positive between
them. With a nonzero positive spectral band and finitely many
local subtractions, its $k$th derivative would have a strict
alternating sign on $X>0$ for all sufficiently large $k$, by
differentiating the positive dispersion integral. Repeated
Rolle's theorem on the positive-axis zeros excludes that
possibility. The two-current loss of physical dispersion persists
both before and after the external shell readout. This construction
isolates a source-dependent entire contribution to the measured
current. Agreement of the charged cuts, free positivity, and the
rooted ultraviolet bounds leave that contribution to be tested.

\paragraph{A negative reflection norm for two positive-time density pulses}
There is a direct Euclidean witness for the missing positive-state
interpretation, independent of the number of dispersion subtractions.
Fix a nonzero spatial momentum $k$ and test the temporal component
of the conserved current. Up to its fixed Wick sign its Fourier
kernel is
\[
 C(\omega,k)=|k|^2\Pi_p(\omega^2+|k|^2)
\]
on the represented connection. For the full transverse readout
multiply this by $|c_L(\sqrt{\omega^2+|k|^2})|^2$. Both kernels
are nonnegative and nonzero on an interval. Since $p\ge2$,
$W_p(t)\le(t/a)^{7/4}$. Consequently
\[
 0<\Pi_p(X)\le C X^{-7/4}\int_0^1 [x(1-x)]^{-3/4}dx
 \quad(X>0).
\]
In particular $(1+\omega^2)C(\omega,k)$ is integrable. Its time
transform $K(t,k)$ is twice continuously differentiable and even,
with $K(0,k)>0$ and
$K''(0,k)=-\int \omega^2C(\omega,k)d\omega/(2\pi)<0$.
Take two density pulses at positive times $\epsilon,2\epsilon$
with coefficients $1,-1$. With the sign that gives a positive
diagonal at short time, their vacuum reflection form is
\[
 K(2\epsilon,k)-2K(3\epsilon,k)+K(4\epsilon,k)
       =\epsilon^2K''(0,k)+o(\epsilon^2)<0
\]
for all sufficiently small $\epsilon>0$. The other overall sign
already fails on a single pulse. The strict test survives smoothing
the temporal pulses inside positive time and localizing the spatial
Fourier test near $k$. A local two-current contact is supported
at zero time difference and drops out at all three reflected
separations. Thus the loss of physical current positivity occurs
in the complete Gaussian response at separated Euclidean points
despite exact real-axis pair-cut agreement.

The two-pulse test also identifies a wider incompatibility. For
any nonzero positive weight $W(t)$ in the same minimal Dirac
heat trace, integrate $\omega^2 C(\omega,k)$ first in frequency.
The Gaussian integral, followed by the Feynman-parameter integral,
bounds it by a constant times
\[
 |k|^2\int_0^1 b_x^{-1/2}dx
       \int_0^\infty W(t)e^{-m^2t}\,\frac{dt}{t^{5/2}}.
\]
For $0\le W\le1$ this is finite whenever the bare free charged
vacuum density $\int_0^\infty W(t)e^{-m^2t}dt/t^3$ is finite.
The zeroth frequency moment is then finite as well. The same
negative reflection norm follows. Thus any such positive scalar
heat preparation that makes the unmarked four-dimensional
charged vacuum finite must modify the current's dependence on
the photon beyond this minimal $H_B$ if it is to have a positive
physical density. This restriction permits other charged
operators and source-dependent measures.

\paragraph{An exact relation between the unmarked vacuum and the hard current}
The same calculation gives a quantitative obstruction for an arbitrary
positive scalar duration weight in this minimal charged generator.
Write $\mathcal V_W=\int_0^\infty W(t)e^{-m^2t}dt/t^3$ for the
proper-time integral in the four-dimensional free charged vacuum
density, omitting its fixed spin and heat-kernel factors. For $e\ne0$,
\begin{equation}
 \lim_{X\to\infty}X^2\Pi_W(X)=\frac{e^2}{\pi^2}\mathcal V_W .
 \label{eq:joint2-vac-hard-uv}
\end{equation}
Here $X=Q^2$ and $\Pi_W$ is defined by \eqref{eq:joint2-maxwell-uv}
with $W_p$ replaced by $W$. To prove the identity set
$H(y)=\int_0^1 b_x e^{-y b_x}dx$. By symmetry it suffices to use
$0\le x\le1/2$, where $x/2\le b_x\le x$.
With $v=yx$, the integrand of $y^2H(y)$ on this half interval,
including its factor of two, is bounded by $2v e^{-v/2}$ and converges
to $2v e^{-v}$ as $y\to\infty$. Thus $y^2H(y)$ is uniformly bounded
and tends to $2$. Use this with $y=Xt$ in the nonnegative proper-time
integral. Dominated convergence proves \eqref{eq:joint2-vac-hard-uv}
for finite $\mathcal V_W$; Fatou's lemma gives an infinite limit if
$\mathcal V_W=\infty$. In the two-record case
$\mathcal V_{W_4}/m^4\simeq0.37393308$ for $Lm=1$.

In particular, a finite unmarked vacuum forces $X\Pi_W(X)\to0$.
An unsubtracted positive spectral measure with mass $w>0$ on any
finite interval $[s_1,s_2]$ instead contributes at least
$wX/(X+s_2)$ to $X\Pi(X)$. Thus the same positive scalar weight
cannot give both the finite unmarked vacuum and that physical
current dispersion in this source. The separated-time test above
also covers local contact conventions. Suppression of internal
momenta in a theory with a physical current must distinguish the
relative displacement inside a loop from the external current
momentum that enters this spectral bound.

\paragraph{A finite curvature term that supplies the missing current response}
The current moments determine a specific repair on the represented
connection $B$. Preserve its coefficient at $X=0$ and require the
ordinary once-subtracted dispersion of its own light-pair cut.
The required scalar change is
\begin{equation}
\begin{aligned}
 \Delta\Pi(X)&=\frac{e^2}{2\pi^2}\int_0^1 b_x\,dx
     \int_0^a\frac{dt}{t}(1-W_p(t))e^{-m^2t}
                     (e^{-b_xtX}-1).
\end{aligned}
\label{eq:joint2-complement-uv}
\end{equation}
This is entire and finite for every finite $X$. Its geometric
realization is a single bilocal curvature term in the action.
For each $x,t$ draw a displacement $Y$ from the four-dimensional
heat kernel at duration $b_xt$ (covariance $2b_xt I$). With the
usual $\tfrac14\int F_{\mu\nu}F_{\mu\nu}$ convention the added
action is
\[
 -\frac{e^2}{16\pi^2}\int_0^1 b_x\,dx\!
 \int_0^a\frac{(1-W_p(t))e^{-m^2t}}{t}\,
 \mathbb E_Y\|F_B(\cdot+Y)-F_B\|_2^2\,dt .
\]
Translation invariance of the norm and the Fourier transform of
the heat kernel give \eqref{eq:joint2-complement-uv}. At small
$t$ the expected square is at most
$8b_xt\|\nabla F_B\|_2^2$ for smooth curvature. Every integral
in this bare difference is finite, including the duration endpoint.
The missing fraction of trajectories fixes the weight. Preserving
the zero-transfer coefficient fixes the only subtraction constant
of the specified once-subtracted dispersion. This is a finite
additional boundary coupling selected by a two-current matching
condition. Its appearance in an unrestricted GCM variation would
need a physical derivation.

Adding the curvature term to the same generating functional gives
\[
 \Pi_{\rm tot}(X)=\Pi_p(0)-
    \frac{e^2}{2\pi^2}\int_0^1 b_x\log(1+b_xX/m^2)\,dx .
\]
Indeed the contributions proportional to $W_p$ and $1-W_p$
now have the common factor $e^{-b_xtX}-1$ after taking the
difference with $X=0$. This reproduces every dispersive moment
and the full positive Dirac two-point spectrum at separated
points on $B$. Local polynomial terms are supported at coincidence and do not
enter the reflection test at positive reflected separations. The new term
is quadratic, so the higher contact derivatives of the charged
determinant with respect to $B$ are unchanged.

The cost depends on which photon propagates. When $B$ itself
has ordinary Maxwell covariance, $\Pi_{\rm tot}(X)$ has the
large-$X$ term $-e^2\log(X/m^2)/(12\pi^2)$. The absolute
two-insertion ring then contains a radial integral proportional
to $\int^\infty X|\Pi_{\rm tot}(X)|^2dX$, which diverges.
When the same A4 readout $B=[c_LP_T+P_L]A$ stays in the
dynamical source, the plane-wave curvature marks on this
additional vertex cost $|q c_L(q)|$ per end and at most
$C[1+\log(1+L^2|q|^2)]$ per vertex. The logarithm can be
split equally between its two ends. The corresponding photon
edge bound at large momentum remains integrable and bounded
since $|c_L(q)|^2=O(|q|^{-(n+1)})$, $n\ge4$. At small momentum
the new scalar is $O(|q|^2)$, hence the previous infrared
convention suffices. The same spanning-tree bound and the
one-momentum trace for a closed chain of quadratic insertions
retain the finite-order Euclidean estimates.

It is then $A$ that couples to the measured current. Its density
at fixed $k\ne0$ has scalar multiplier
$|k|^2 c_L(\sqrt X)^2\Pi_{\rm tot}(X)$, $X=\omega^2+|k|^2$.
This real, even, nonzero Fourier kernel and its product with
$|\omega|$ are integrable. Its time transform is differentiable,
vanishes at infinity, and has zero derivative at the time origin.
A positive spectral representation on positive times, allowing a
fixed Wick sign, would have finite total mass by continuity from
the right. Monotone convergence of the first difference quotient
would give zero first energy moment. Only zero-energy weight
could remain, and the large-time limit eliminates that weight.
This is impossible for the nonzero kernel. Restoring the missing
two-current response on the charged connection thus distinguishes
two physical requirements: the high-momentum photon used for
joining loops and the positive-state photon coupled to a probe
still have to be realized by a single physical law.

\subsection{An inverse construction with the full charged-pair weight}
\label{sec:joint2-inverse-uv}
There is a positive invariant-mass model in which the same geometry
gives the open pair amplitude, the photon spectrum and the complete
quadratic trace. The extra datum is an interaction among the pair
modes. It can be determined by an inverse problem. The forward
square in \eqref{eq:fix-fock-square-uv} starts with a fixed free
mass in that sector. Here the prescribed mass and current belong
to the \emph{observed} continuum after mixing.

Fix one transverse component on the reconstructed sector plane. Let
$t_0=4m^2$, let $T$ multiply by $s$ on
$\mathcal H=L^2([t_0,\infty),ds)$ and choose the observed pair row
\[
 \begin{aligned}
 u(s)&=h_0(-s)\sqrt{\rho_D(s)},\\
 \rho(s)&=|u(s)|^2,\qquad M_j=\int s^j\rho(s)\,ds .
 \end{aligned}
\]
Here $h_0$ is the light-shell mark of
\eqref{eq:nextsoft-product-uv} and $\rho_D$ includes the spin
sum and two-body phase space for the charge-one Dirac pair.
Thus $\rho=C_{\rm pair}$ in \eqref{eq:nextsoft-positive-uv}.
This is the cyclic combination of states seen in this channel;
orthogonal spin combinations may be appended unchanged.
The real-axis A4 ball bound at $n=248$ makes $M_{-1},M_0,M_1,M_2,M_3$
finite. Use a dimensionless coupling marker $g$ with
$g^2M_{-1}<1$. Define
\begin{equation}
\begin{aligned}
 T_g&=T-g^2|u\rangle\langle u|,\\
 B_g(c,f)&=g u c+T_g^{1/2} f,\qquad
 H_g=B_g^\dagger B_g .
\end{aligned}
\label{eq:joint2-inverse-square-uv}
\end{equation}
Cauchy--Schwarz with $T^{1/2}f$ gives
$T_g\ge(1-g^2M_{-1})T>0$. The attractive pair interaction is
uniquely fixed within this ansatz by the condition
$B_gB_g^\dagger=T$ with the observed row $g u$ held fixed.
All quantities in the rank-one change are finite. In carrier
units $Lm=1$, even the choice $g=1$ lies within this range.
Indeed for $y=L\sqrt{s(s/(4m^2)-1)}$ one has $ds\le2m^2dy$.
The ball's one-coordinate density $f_n$ obeys
$\int_0^\infty\chi_n(y)^2dy=\pi\int f_n^2\le\pi f_n(0)$.
The gamma ratio for $f_n(0)$ and Cauchy--Schwarz give
$f_n(0)\le\sqrt{(n+1)/(2\pi)}$. Since
$\rho_D\le1/(12\pi^2)$, it follows that
$M_{-1}\le M_0/t_0\le f_n(0)/(24\pi)<1$.

\paragraph{The open channel and the photon pole}
The map $W=B_g^\dagger T^{-1/2}$ is an isometry from
$\mathcal H$ onto $(\ker B_g)^\perp$ and $H_gW=WT$ on the
spectral domain. For the unit bare photon $e_A=(1,0)$ it gives
$W^\dagger e_A=g T^{-1/2}u$. The two physical readouts of the
same source are therefore explicit. The massless photon has
residue $Z=1-g^2M_{-1}$. Its massive spectral density is
$g^2\rho(s)/s$. The measured transverse current obtained from
$H_ge_A/g$ has continuum amplitude $\sqrt{s}\,u(s)$ in this
same basis. Dividing its spin weight by $s$ gives exactly
$C_{\rm pair}(s)$ at every open energy. The dressed photon
obeys the dispersion identity
\begin{equation}
\begin{aligned}
 R(X)&=\int_{t_0}^\infty\frac{\rho(s)}{X+s}\,ds,\\
 D_g(X)&=\frac{1-g^2R(X)}{X}
   =\frac{Z}{X}+g^2\int_{t_0}^\infty\frac{\rho(s)}{s(X+s)}\,ds .
\end{aligned}
\label{eq:joint2-inverse-spectrum-uv}
\end{equation}
For constant $g$ the covariant vector square has the usual Maxwell
kinetic terms for $A$ and the continuum fields $F$ and mass form
$\|g u A+T_g^{1/2}F\|^2/2$ componentwise. Under
$A\mapsto A+d\Lambda$, take $F\mapsto F-gT_g^{-1/2}u\,d\Lambda$.
This verifies the constant-coupling Ward symmetry. The massive
spectrum is carried by neutral spin-one modes representing the
measured pair channel. Producing charged asymptotic spinors and
their crossed or multipair amplitudes requires further dynamics.

\paragraph{The pair determinant changes the ultraviolet conclusion}
The bare susceptibility of the coupled pair modes is
$P_g(X)=\langle u,(X+T_g)^{-1}u\rangle=R(X)/(1-g^2R(X))$
by the rank-one resolvent identity. The Schur factor of the
photon is $1+g^2P_g(X)$. The pair determinant, relative to the
same modes at $g=0$, is $1-g^2R(X)$. Both arise in the native
block trace. Their product is one at each $X>0$, before the
center momentum is integrated. Each Taylor coefficient of the
\emph{complete relative trace} at constant $g$ consequently
vanishes. In particular at fourth order the term from
$P_g=R+g^2R^2+\cdots$ as well as the bare-pair trace must be
retained. The isolated two-insertion Maxwell ring still has
logarithmic cost since $XR(X)\to M_0>0$. The cancellation
relates different contributions from the same block. It does
not make that ring an absolutely convergent graph in isolation.

\paragraph{A nonconstant source tests more than isospectrality}
A constant coupling tests only the spectrum. To probe local
variation, replace $g$ by a smooth scalar source $g(x)$ in
the mass block alone and keep the ordinary center kinetic operator.
This calculation concerns a single scalar transverse comparison
component. A spacetime-dependent local gauge transformation of the
full vector model would also vary its polarization constraints.
Put $H_0=0\oplus T$ and $V(g)=H_g-H_0$. Each $V(g)$ has finite
rank. Two exact trace identities at every center point are
\[
 \operatorname{tr}V(g)=0,\qquad
 \operatorname{tr}V(g)^2=2g^2M_1=-2\operatorname{tr}(H_0V(g)).
\]
Indeed its diagonal blocks are $g^2 M_0$ and
$-g^2|u\rangle\langle u|$ and its off-diagonal vector
is $g T_g^{1/2}u$. These identities give integrability at every
finite multiplicity of the external scalar source for the
complete one-loop relative log determinant.

For the hard estimate write $t=|q|^2$ and
$G_0(q)=(t+H_0)^{-1}$; an infrared mass may be held fixed
throughout the estimate. At fixed external momenta the one-insertion
trace is $g^2M_1/t^2+O(t^{-3}g^2M_2)$. The leading term of the
two-insertion trace in $\operatorname{Tr}\log(I+G_0 V)$ is
$-\operatorname{tr}V(g)^2/(2t^2)$, integrated over the center.
The two terms cancel by the displayed identity, coefficient by
coefficient in the external source. For fixed transfer $p$,
replacing $G_0(q+p)$ by $t^{-1}I$ next to a marked vertex costs
$O_p(|q|^{-3})$ from the center momentum and $O(t^{-2})$ times
the norm with one internal $H_0$. These latter norms are finite
since $M_3<\infty$. For example the square-root resolvent formula
writes $(T_g^{1/2}-T^{1/2})u$ as an integral of
$-g^2 r^{1/2}(T+r)^{-1}u\,R(r)/(\pi[1-g^2R(r)])$.
Its Taylor coefficients converge also in the graph norm of $T$.
Hence the complete two-term remainder is integrable in four
center dimensions. Each word with at least three insertions has
the integrable bound $O(|q|^{-6})$ or better. This proves the
claim at fixed external momenta in the specified infrared
convention, after summing the source-prescribed contacts.

There is an independent check at quadratic source order. In the
relative half log determinant the coefficient of $|\widehat
g(p)|^2$, using center loop measure $d^4q/(2\pi)^4$, is
\begin{equation}
 \frac{1}{32\pi^2}\int_{t_0}^\infty s\rho(s)\,ds
       \int_0^1\log(1+x|p|^2/s)\,dx .
 \label{eq:joint2-inverse-gradient-uv}
\end{equation}
Before the $q$ integration the difference per unit $s\rho(s)/2$ is
$|q|^{-2}[(|q|^2+s)^{-1}-(|q+p|^2+s)^{-1}]$.
Feynman parametrization of this convergent difference gives the
formula. It vanishes for a constant probe and is strictly positive
at nonzero transfer. Its leading small-momentum coefficient is
$M_0|p|^2/(64\pi^2)$. The surviving gradient response is
therefore finite and nontrivial, even though every constant-source
vacuum coefficient cancels exactly.

\paragraph{A Ward-invariant test with spatially varying mixing}
The finite marked comparison also has a gauge-invariant version.
At nonzero Euclidean momentum put $P(q)=I-qq^{\mathsf T}/|q|^2$,
apply the same projector to every vector component of the photon
and pair fields, and use the mass form
$\langle P Y,H_{g(x)}P Y\rangle/2$ together with the free Maxwell
kinetic forms. Independent gradient shifts of $Y$ leave this
action invariant for every scalar profile $g(x)$. The projection
can be written using the curvature and the inverse Laplacian.
It is an explicit nonlocal choice of the comparison action. Use
the transverse functional measure and the same infrared convention
for the zero mode. At constant $g$, the response to conserved
vector probes has precisely the positive mass spectrum in
\eqref{eq:joint2-inverse-spectrum-uv}.

The external source estimate survives this change. A one-insertion
trace carries $\operatorname{tr}P(q)=3$, whereas a two-insertion
trace carries $2+(\hat q\cdot\widehat{q+p})^2$. The latter equals
$3+O_p(|q|^{-2})$ at hard momentum and is bounded everywhere on
the integration domain. All higher products of projectors are
bounded as well. Thus the leading one- and two-insertion terms
still cancel exactly, and the remainder at each finite multiplicity
is integrable. At quadratic order the transverse answer is three
times \eqref{eq:joint2-inverse-gradient-uv} plus
\[
 \frac12\int s\rho(s)\,ds\int\!
   \frac{d^4q}{(2\pi)^4}\,
   \frac{1-(\hat q\cdot\widehat{q+p})^2}
        {|q|^2(|q+p|^2+s)} .
\]
The extra term is finite and nonnegative for real Euclidean
momenta. This calculation supplies the Ward identity for all
external scalar mixing profiles. It concerns the indicated
neutral gauge-invariant probe and the transverse spectrum at
constant coupling. Causal Lorentzian coupling of an arbitrary
mixing profile, genuine charged endpoints and the ultraviolet
cost of quantizing that profile each require additional tests.

\paragraph{Sewing the geometric marks with the same dynamical photon}
The external-source estimate can be made uniform enough to sew
photons. Fix the infrared convention, with a common positive
infrared mass $\mu$ on the Euclidean comparison propagators
restricted to their transverse subspaces. Scale all momenta in
units of $m$. For $N$ scalar mixing marks of momenta $p_i$ put
$P_*=1+\sum_i|p_i|$. After grouping the one- and two-insertion
contacts in the preceding proof, the absolute integral over
the center circulation is at most
\[
 C_{N,\mu}\,[1+\log(1+P_*)].
\]
On $|q|\le2P_*$ the one-insertion internal trace is bounded by
$C/(|q|^2+\mu^2)^2$. For two insertions, Cauchy--Schwarz
bounds the integral of the two translated resolvent envelopes
by $C[1+\log(1+P_*)]$ on this ball. All shifts are bounded by
$P_*$. For three or more insertions, H\"older with exponent equal
to their number gives a uniform bound on the whole space.
On $|q|>2P_*$ the leading one- and two-insertion terms cancel.
Their geometric and momentum-shift remainders are bounded by
$C_N(|q|^{-6}+P_*|q|^{-5})$ there, with constants independent of
the transfers. The same bound applies to the transverse projection
since $\|P(q+p)-P(q)\|\le C|p|/|q|$ in this region.
Integration gives the stated uniform estimate.

For an explicit nonlinear test keep the transverse projection on
both unsmeared slots of the mass block. Average the Maxwell
curvature with the four-coordinate A4 ball displacement on the
Euclidean readout, writing $\mathcal T F_A$ for its convolution
and $\chi_n(L|q|)$ for the multiplier. Take $n=248$ and choose
\begin{equation}
 g(e,A)=e\left[1+
  \tanh\!\left(e^2 L^4\sum_{\mu<\nu}
      (\mathcal T F_{A,\mu\nu})^2\right)\right] .
 \label{eq:joint2-native-mark-uv}
\end{equation}
Use this coefficient pointwise in the \emph{same}
$H_{g(e,A)}$ of \eqref{eq:joint2-inverse-square-uv}. The response
uses the coefficient $e$ at zero background. This bounded function
is a specified representation choice, with no fitted scale.
For $|e|\le1$ one has $|g|\le 2$; the preceding bound at
$n=248$ gives $4M_{-1}<1$. The Euclidean mass form is positive
at each real field. The gradient gauge symmetry of the projected
action is exact. Taylor coefficients in $e$ involve only finitely
many curvature marks at each fixed order.

In the bare expansion there are two distinguished unsmeared vector
slots at each mass vertex. Every extra photon slot from
\eqref{eq:joint2-native-mark-uv} has a curl and a ball multiplier.
Remove every internal line with at least one such marked end. The
remaining lines form paths and closed cycles. Fix the momenta
flowing through the removed lines. A closed unmarked cycle is
exactly an external-source trace of the kind bounded above.
Its contact grouping can be made before sewing: first label the
external mixing marks, take the complete coefficient of the
auxiliary scalar source on the cycle, then substitute its
curvature polynomial at each label and Wick-contract the resulting
photon slots. The multinomial coefficients are those of this
substitution. This includes a contraction between marks on the
same cycle. An unmarked path carries no independent circulation;
its propagators and finite-rank marked coefficients have bounded
operator norms at the fixed infrared scale.
A removed internal line of momentum $k$ has $b=1$ or $2$ marked
ends and scalar envelope
\[
 C\frac{|k|^b |\chi_n(L|k|)|^b}{|k|^2+\mu^2}.
\]
Each envelope is bounded and integrable on $\mathbb R^4$ even after
multiplication by any fixed logarithmic power. This follows from
$|\chi_n(y)|\le C(1+|y|)^{-(n+1)/2}$ on the real axis.
Contract the paths and completed cycles to vertices. Their
logarithmic factors are bounded by products of fixed powers
of $1+\log(1+|k|)$ on the incident removed lines and constants
depending on the external tests. Bounding a spanning tree by
supremum norms and integrating the independent chords proves
absolute convergence at every finite order in $e$ for these
completed Euclidean coefficients. All center and exchanged
momenta range over the whole plane, including overlapping hard
regions.

This is a nonlinear dynamical photon and neutral pair-mode example.
The all-order statement uses the contacts fixed by its common
mass form on each trace. Individual terms before those contacts
are combined can diverge. Its positive light residue and exact
pair continuum at constant mixing are the quadratic, unsewn
spectral data. Positivity of a Euclidean density does not establish
reflection positivity after the four-dimensional curvature
smearing. Interacting physical spectral positivity, canonical
electric charge for individual fermions, and selection of this
gluing by the GCM dynamics require separate demonstrations.

\paragraph{Protecting the same massless direction at every field}
Pointwise singularity of a mass matrix leaves room for a
positive cost when its kernel changes with position. For example,
hold a transverse constant photon $c$ fixed, prescribe
$g(x)=\varepsilon h(x)$ with real smooth $h$ of zero mean,
and minimize the quadratic action over its pair field. The
coefficient of $\varepsilon^2$ in the minimized action is
\[
 \frac12\int\frac{d^4 k}{(2\pi)^4}|\widehat h(k)|^2
 \left[M_0|c|^2-|P(k)c|^2
      \int\frac{s\rho(s)}{|k|^2+s}\,ds\right].
\]
One obtains this formula by completing the massive square with
the transverse resolvent $P(k)(|k|^2+T)^{-1}$ at zero mixing.
The bracket is at least $|c|^2|k|^2R(|k|^2)$, strictly positive
away from zero for $c\ne0$. Constant backgrounds can be defined
in a periodic box before taking the continuum limit. This
prescribed-profile test diagnoses the cost of rotating the light
direction. The dynamical curvature law in \eqref{eq:joint2-native-mark-uv}
allows its first mass correction to be computed explicitly in
\eqref{eq:next3-moving-mass-uv}.

\paragraph{Fluctuations of the moving null direction generate a finite mass term}
For the alternative coupling \eqref{eq:joint2-native-mark-uv}, the
position of the unit null vector of $H_g$ depends on the measured
curvature. The prescribed-profile stiffness can be carried through the
first photon contraction. Work in the exact free basis
$Y=n_e a+W_e b$ at $g=e$, where the quadratic masses are $0,T$.
At this order it suffices to write
\[
 H(x)=L^4\sum_{\mu<\nu}(\mathcal T F_{a,\mu\nu})^2(x),
 \qquad g-e=e^3H+O(e^4).
\]
The remainder allows for the $O(e)$ massive component of the
observed photon in this basis. The interaction depends on $a$ through
curvature at these marked slots, so shifting $a$ by a constant
physical polarization $c$ leaves $g$ unchanged.

Differentiate $H_e n_e=0$ and put $w_e=W_e^\dagger H'_e n_e$,
where a prime differentiates the mass matrix with respect to $g$ at
$e$. The vanishing eigenvalue gives two identities,
\[
 \langle n_e,H'_e n_e\rangle=0,\qquad
 \tfrac12\langle n_e,H''_e n_e\rangle=\langle w_e,T^{-1}w_e\rangle.
\]
At $e=0$, $w_0=T^{1/2}u$ and the last scalar is $M_0$.
Keep the constant null background fixed in a periodic test box and
eliminate the massive field, including its constant component. The
term linear in $c$ starts with $e^3\int H\,c\cdot\langle w_0,b\rangle$.
The term quadratic in $c$ starts with
$e^6 M_0 |c|^2\int H^2/2$. In the negative logarithm, the former
enters with minus one half its variance and the latter with its free
expectation. No other term contributes to the constant quadratic
coefficient at that order: the first derivative of the mass form
has zero null--null entry, and every variation of a curl under a
constant shift vanishes. At leading order $H$ depends only on the
free photon and is independent of $b$.

Write $C_H(k)$ for the Fourier density of the connected covariance of
$H$. The disconnected term $|\mathbb E H|^2$ cancels between the
direct contact and the massive response at $k=0$. This cancellation
is exactly the Schur correction from a constant change of the massive
field. It can equivalently be obtained at small nonzero external
momentum with $P_q c=c$ before taking the limit. The coefficient of
$e^6$ in the remaining quadratic energy density is
\begin{equation}
\begin{aligned}
 \tfrac12\int \frac{d^4 k}{(2\pi)^4}C_H(k)
 \left[M_0|c|^2-|P_kc|^2
                 \int\frac{s\rho(s)}{|k|^2+s}\,ds\right].
\end{aligned}
\label{eq:next3-moving-mass-uv}
\end{equation}
Indeed the bracket is bounded above by $M_0|c|^2$ and below by
$|c|^2 |k|^2R(|k|^2)$ with $R(X)=\int\rho(s)ds/(X+s)$.
For the free A4 photon, the covariance of the filtered curls in
momentum space is a positive matrix bounded by $C|c_L(k)|^2$.
Wick's rule gives $C_H$ as twice the sum of the corresponding
convolutions, including the factor $L^8$. It is nonnegative as a
spectral density, integrable, and strictly positive near zero since
its value there is the integral of a nonzero matrix square.
Consequently \eqref{eq:next3-moving-mass-uv} is finite and strictly
positive for $c\ne0$. This integral extends over the entire center
momentum space.

Equivalently one can compute the quadratic correction after eliminating
$b$ and then take the soft transverse momentum to zero. In the two
first-variation terms with both external legs undifferentiated, the
connecting line has positive mass and the filtered internal photon
lines have a curl at either end. Their integrable bounds give the same
limit by translation continuity. Terms with an external curl vanish
in this limit. At leading order the two first-variation vertices each
have one massive slot and three null slots. Hence the massive slot
must contract between the vertices, and an internal null contraction
with one or both unmarked ends has at worst the locally integrable
$|k|^{-2}$ singularity. Any single massless bridge would require a
nonvanishing third-order null--null two-point coefficient on both
sides, which is absent by the odd number of massive fields at a
first-variation vertex. In the two-vertex graph the massive line
can carry the balancing momentum, so the two photon contractions
have independent momenta, also for a self-contraction. On hard
regions with external curls one uses the complete one- and
two-insertion trace bound above. It is uniform at this order on
soft external momenta: each unsmeared two-line trace contains
a gapped line, and the remaining single massless inverse is
locally integrable. Thus external curl factors are retained after
integration. This gives a mass term for the transverse
inverse after the massive Schur complement at order $e^6$. It is
a property of this moving-kernel action. The protected fixed-kernel
action below removes precisely this contribution by keeping $n_e$
in the kernel at every field value.

\paragraph{A fixed null direction for the nonlinear coupling}
One can remove this cost at the level of an interacting action.
Fix $e$ with $e^2M_{-1}<1$, write $n_e$ for the unit null vector
of $H_e$ chosen with positive photon overlap, and put
$W_e=B_e^\dagger T^{-1/2}$. Choose $v_1,v_2$ by orthonormalizing
$u,Tu$ in that order and let
$J=|v_2\rangle\langle v_1|-|v_1\rangle\langle v_2|$.
The continuum support of $u$ guarantees their independence and
$[T,J]\ne0$. The A4 tail gives $u,Tu\in D(T^2)$ at $n=248$.
In the projected Euclidean mass form use
\begin{equation}
\begin{aligned}
 a_0&=\langle n_e,Y\rangle,\\
 t(x)&=e^2 L^4\sum_{\mu<\nu}(\mathcal T F_{a_0,\mu\nu})^2(x),\\
 U_t&=\exp(tJ),\\
 \widetilde H(x)&=W_e U_{t(x)}^\dagger T U_{t(x)} W_e^\dagger .
\end{aligned}
\label{eq:joint2-fixed-null-uv}
\end{equation}
Here the curvature in the interaction measures the physical
null photon $a_0$. The original optical probe remains $e_A$.
At zero field $\widetilde H=H_e$, so its complete quadratic
spectrum and optical normalization are unchanged. The kinetic
form is the same as above. Every $\widetilde H(x)$ is positive
and annihilates the same $n_e$, and the full action is invariant
under $Y\mapsto Y+n_e c$ for constant $c$. Its local massive
coupling has a nonzero linear term $[T,J]t$. This constitutes
an interaction with two photon curls, fixed by a geometric
smearing and a specified rotation in the measured continuum.

The splitting $Y=n_e a_0+W_e b$ sharpens both estimates. Every
unsmeared interaction slot is now a gapped $b$ slot, with free
resolvent $(|q|^2+T)^{-1}$. Write
$V_t=U_t^\dagger T U_t-T$. The identities
\[
 \operatorname{tr}V_t=0,\qquad
 2\operatorname{tr}(TV_t)+\operatorname{tr}(V_t^2)=0
\]
follow by conjugating $T$ and $T^2$. All differences and all
cyclic permutations used here are finite rank on vectors in
the requisite domains. For instance the range of $U_t-I$ is
spanned by $u,Tu$, which have the required higher moments.
Thus closed $b$ cycles obey the preceding one- and two-insertion
contact cancellation with a strictly positive gap $t_0$ replacing
the infrared comparison mass. The remaining $b$ paths have
uniformly bounded propagators. Every internal $a_0$ line has
a smeared curl at \emph{both} ends. Its full tensor envelope is
$C|\chi_n(L|k|)|^2$, bounded at zero and integrable with any
fixed logarithmic weight in four dimensions. The spanning-tree
argument therefore proves cutoff-free convergence of the
contact-completed coefficients with the true massless internal
photon, also when external momenta range over a compact set.
Vacuum traces throughout are relative to the free kinetic sectors.

In a proper two-point coefficient each external $a_0$ must enter
through its own curl. Amputation consequently leaves two powers
of external momentum with uniformly bounded coefficients; each
internal null line already has its two curls. At nonzero
external momentum tending to zero along a fixed ray the
remaining integrals have finite directional limits by the same
bounds (the projectors are bounded and continuous off a null
set). The symmetry $b\mapsto-b$ makes the mixed two-point
function vanish. The inverse transverse null propagator thus
equals $|q|^2$ times a formal series with tree coefficient one.
This protects the perturbative massless pole to every finite
order, without an infrared mass or a mass counterterm. It
concerns the neutral-mode completion \eqref{eq:joint2-fixed-null-uv}.
For this four-dimensional smearing, Section~\ref{sec:next3-null-uv} computes a negative reflected photon norm at the first interacting order. The quadratic Lorentzian measure by itself still supplies the stated one-pair spectrum.

\paragraph{Using the observed photon in the nonlinear interaction}
The protected direction can also be retained when the curvature that
controls the massive rotation is that of the original optical probe.
In the same representation let $A_{\rm obs}=\langle e_A,Y\rangle$
and set
\[
 t=e^2L^4\sum_{\mu<\nu}(\mathcal T F_{A_{\rm obs},\mu\nu})^2,
 \qquad A_{\rm obs}=\sqrt{Z}\,a_0+
       \langle W_e^\dagger e_A,b\rangle .
\]
Its free two-point function at this fixed mixing has the massless pole
and the measured pair continuum of \eqref{eq:joint2-inverse-spectrum-uv}.
The massive rotation leaves the null vector fixed at every field.
Expand in the number of curvature insertions on this background.
An internal null line still has two marked ends and the bounded
integrable envelope $C|\chi_n(L|k|)|^2$. A massive line with one or
two marked ends has envelope bounded by
$C|k\chi_n(L|k|)|^j/(|k|^2+t_0)$, $j=1,2$. This is bounded at
zero and integrable at infinity even with fixed logarithmic weights.
Deleting all lines with a marked end again leaves only gapped paths
and complete gapped cycles. The trace identities for $V_t$ hold
before substitution of the observed curvature. Consequently the same
absolute convergence argument applies to these coefficients, with
massive external states tested in their Hilbert space. The comparison
may also be expanded at $e=0$: the bounded rank-one change and the
strict gap of $T_e$ give analytic normal vectors and matrix elements
in the required finite graph norms near zero.

The reflection $b\mapsto-b$ alone is no longer a symmetry. Instead
write the proper transverse two-point operator in its null and massive
blocks. A null leg always enters through a curl. The null--null
correction is therefore $O(|q|^2)$ and the mixed corrections are
$O(|q|)$ in operator norm for compact soft external momenta. The
massive--massive correction is bounded, by the gapped-path and
cycle bounds above. In formal perturbation theory its inverse block
stays bounded near zero order by order, since the free inverse on
this block has norm at most $1/t_0$. Eliminating that block adds to
the null inverse a product of two mixed factors and one bounded
inverse, hence another $O(|q|^2)$ term. The two-curl factors and the
integrable edge bounds give finite soft limits along fixed rays.
With unit tree coefficient the perturbative massless mode survives.
Its nonzero overlap $\sqrt Z$ keeps it visible in the original
optical photon as a formal series. This gives a common observed and
interacting photon in a neutral mode model. Section~\ref{sec:next3-null-uv} tests its interacting null photon and finds a negative reflected norm at fourth order. The coupling to individual charged states is tested separately in Section~\ref{sec:next3-car-uv}.

\subsection{What spectral weight must the free photon already carry?}

\paragraph{One neutral massive level and the physical light-pair weight}
There is a positive spectral comparison in which the same transverse
photon has a massless pole, the prescribed leading pair continuum and
finite vacuum rings of every positive length. It uses a massive
geometric mode already at zero electric coupling. To see exactly
what this costs, let
\[
 d\kappa(s)=C_{\rm pair}(s)\,ds/s,\qquad
 k_j=\int s^j\,d\kappa(s).
\]
Use the open-pair density in \eqref{eq:nextsoft-positive-uv} in
unit-charge normalization. Its support begins at $4m^2$. For this
construction assume $k_0,k_1,k_2,k_3$ finite and $k_1>0$. These
conditions hold for the specified ball form factor. Match the
moments $k_0,k_1,k_2$ by one neutral pole and a massless reservoir.
The matching equations have the unique solution
\begin{equation}
 b=\frac{k_2}{k_1},\qquad c_b=\frac{k_1^2}{k_2},\qquad
 c_0=k_0-c_b .
 \label{eq:joint2-radau-moments-uv}
\end{equation}
This choice minimizes the massive residue required of any positive
reservoir of neutral states in this class. Indeed let $d\nu$ be
the decrease of the massive neutral spectral measure per unit
$g$, with an independently adjustable change of the massless atom.
Assume finite third moments. Finiteness of the first absolute
four-dimensional trace requires cancellation of the signed zeroth,
first and second spectral moments: the first nonzero term through
that order gives a relative insertion asymptotic to a nonzero
constant times $X^{-j}$, $j\le2$. In particular
$\int s\,d\nu=k_1$ and $\int s^2d\nu=k_2$, so
$\int d\nu\ge k_1^2/k_2=c_b$. Equality forces $s=b$
almost surely in this massive measure. Also its support must reach
at least $b$, since its $s\,d\nu/k_1$ mean is $b$.
Thus a reservoir restricted below the physical light-pair threshold
cannot supply this repair. Both weights in the displayed solution
are nonnegative by Cauchy--Schwarz. The squared mass
$b$ is the mean of $s$ with respect to $s\,d\kappa/k_1$ and is
strictly above threshold for this continuum. At $n=248$, $Lm=1$,
the dimensionless quantities are $b/m^2\simeq21.98930$,
$c_b\simeq0.00935871$ and $c_0\simeq0.00428463$.
One may reproduce these values with
$y=s/m^2=2+2\sqrt{1+k^2}$, integrating
$y^{j-1}(1+2/y)\sqrt{1-4/y}\,\chi_{248}(k)^2\,dy/(12\pi^2)$.
The decay $\chi_{248}(k)=O(k^{-249/2})$ controls the tail.

Use $g\in[0,1]$ to count the transfer of spectral weight, with $g$
of order $e^2$. Normalize the permanent massless residue to one. Put
\begin{align}
 D_g(X)={}&\frac1X+
 (1-g)\left(\frac{c_0}{X}+\frac{c_b}{X+b}\right)
 +g\int\frac{d\kappa(s)}{X+s},\qquad X>0 .
 \label{eq:joint2-radau-photon-uv}
\end{align}
This is the covariance of the Gaussian transverse action with
curvature multiplier $[XD_g(X)]^{-1}$. It has a positive photon
spectral representation at every $g$. If $D_0$ denotes its free
covariance, set $h=(D_g-D_0)/(gD_0)$, defined by the derivative
at zero. Its sign and absolute ultraviolet power follow from
an exact interpolation remainder:
\begin{equation}
 D_0(X)h(X)=
 -\frac{1}{X(X+b)^2}
   \int\frac{s(s-b)^2}{X+s}\,d\kappa(s).
 \label{eq:joint2-radau-error-uv}
\end{equation}
Indeed subtract from $(X+s)^{-1}$ the quadratic in $s$ with the
same value at zero and the same value and slope at $b$. Their
difference is $-s(s-b)^2/[X(X+b)^2(X+s)]$. The three matched
moments integrate that quadratic exactly. Hence $h\le0$ and
$h=O(X^{-3})$ at large $X$. It stays bounded at zero. The
inverse quadratic action relative to the free one has multiplier
$(1+gh)^{-1}-1$; the permanent pole bounds $1+gh$ away from zero.
Every trace word of positive total order $p$ costs at most a
constant times $\int_0^\infty X|h(X)|^p\,dX$, finite already at
$p=1$. Direct higher inverse-action contacts satisfy the same
bound. The proof integrates the full radial momentum range.

There is a sharp condition even without a third moment. With
$k_0,k_1,k_2$ finite, $XD_0(X)$ lies between two positive
constants. On $X\ge b$, Tonelli and
$X/(X+b)^2\asymp1/X$ show that the first trace is finite
exactly when
\[
 \int (s-b)^2\log(1+s/b)\,d\kappa(s)<\infty .
\]
The elementary inner integral is
$\int_b^\infty dX/[X(X+s)]=\log(1+s/b)/s$.
The normalizations, the neutral mass and the amount of residue
removed were fixed by matching the measured continuum moments.
The preparation of the initial neutral level and its transfer into
that continuum are extra physical hypotheses. The leading change
in the spectrum contains the negative atom $-c_b\delta_b$ as
well as $d\kappa$ and the massless residue change. Therefore the
positive-leading-massive-change hypothesis of
\eqref{eq:sew-dressed-photon-ring-uv} does not hold here.
Equation~\eqref{eq:joint2-radau-photon-uv} is a quadratic
photon model with the desired two-point data. To derive it from
the physical GCM one would need the dynamics of that neutral
level and its coupling to the actual charged-pair states, with
the same higher contacts and functional measure. In particular,
the ordinary source-free Maxwell equation at zero electric
coupling would exclude a massive transverse photon matrix
element there. The required neutral contribution to the photon
Euler equation is a physical signature of this route.

\paragraph{A self-adjoint mass operator that realizes the spectral transfer}
The covariance in \eqref{eq:joint2-radau-photon-uv} can be realized
with one source and a definite charged-pair matrix element. This
gives a test of the degrees of freedom removed when only its photon
quadratic form is kept. Work in a single physical transverse
polarization in the invariant-mass representation. Let $H_0\ge0$
be the squared-mass operator. Give it a permanent massless unit
state $a_0$, a second massless state $n_0$, a neutral eigenstate
$n_b$ at $b$, and the actual free charged two-body continuum in
this polarization. These spaces are mutually orthogonal.
Use the continuum of the measured light-pair amplitude to prepare
a unit wave packet $w$ with spectral measure $d\kappa/k_0$. In a
direct integral over pair spins and invariant mass, use the
normalized vector of pair amplitudes in each mass fiber, multiplied
by the square root of this scalar density with respect to the
two-body measure. Its phase and spin dependence can thus be chosen
to equal those of the given measured transverse amplitude.
Define
\[
 u=(\sqrt{c_0}n_0+\sqrt{c_b}n_b)/\sqrt{k_0},\qquad
 a=a_0+\sqrt{k_0}\,u.
\]
The weights in \eqref{eq:joint2-radau-moments-uv} make $u$ a
unit vector. Both $u$ and $w$ belong to the domain of $H_0$
since the continuum has finite second moment.

Let $T=|w\rangle\langle u|-|u\rangle\langle w|$,
$U_\theta=e^{\theta T}$, and set
$H_\theta=U_\theta^\dagger H_0U_\theta$ with $g=\sin^2\theta$.
The rotation fixes the orthogonal complement of these two vectors.
In particular $U_\theta a=a_0+\sqrt{k_0}(\cos\theta\,u+
\sin\theta\,w)$. By orthogonality of their spectral supports
as subspaces, the following two predictions are exact:
\[
 \begin{aligned}
 \langle a,(X+H_\theta)^{-1}a\rangle&=D_g(X),\\
 \|U_\theta^\dagger E_{\rm pair}(I)U_\theta a\|^2
                  &=g\int_I d\kappa(s).
 \end{aligned}
\]
Here $E_{\rm pair}(I)$ selects the original charged continuum in
the mass interval $I$. The second line measures the corresponding
exact eigenchannels of $H_\theta$. Thus the photon resolvent and
the specified one-pair spectral amplitude follow from the same
self-adjoint quadratic dynamics. The second neutral massless
component and the massive level are physical inputs to this
comparison, not consequences of a choice of pole sign.

There is an independent ultraviolet test on the enlarged action.
Its first change $K=[H_0,T]$ is bounded, finite rank and self-adjoint
because $H_0u,H_0w$ have finite norms. It is nonzero: the
charged continuum packet is not concentrated at one eigenvalue.
With $R_0=(X+H_0)^{-1}$, the microscopic two-insertion vacuum word
has positive factor $\|R_0^{1/2}KR_0^{1/2}\|_{\rm HS}^2$.
Strong convergence $\sqrt X R_0^{1/2}\to I$ and finite rank give
\begin{equation}
 \lim_{X\to\infty}X^2\|R_0^{1/2}KR_0^{1/2}\|_{\rm HS}^2
   =\|K\|_{\rm HS}^2>0 .
 \label{eq:joint2-rotated-ring-uv}
\end{equation}
The radial measure $X\,dX$ therefore produces a logarithmic divergence
in this individual bare graph. Its coefficient can be read directly
from the target density. The first two moments of $H_0$ on $u$ and
$w$ agree by construction. Their centered images are orthogonal to
$u,w$ and to each other, each with squared norm
$k_2/k_0-(k_1/k_0)^2$. Expanding $[H_0,T]$ in these four orthogonal
directions gives $\|K\|_{\rm HS}^2=4[k_2/k_0-(k_1/k_0)^2]$.
The total change of vacuum determinant
vanishes, as follows on finite tests when the mass matrix is conjugated.
In the continuum the corresponding relative Fredholm determinant is
also one: $(X+H_\theta)(X+H_0)^{-1}$ factors as $U_\theta^\dagger$
times $(X+H_0)U_\theta(X+H_0)^{-1}$. Both are the identity plus
finite rank and their determinants are the reciprocal unitary
determinants. For the second factor this follows from the finite
rank determinant formula on the two transported columns in the
domain of $H_0$. The vanishing total change includes the contact
from differentiating the same $H_\theta$ twice. Consequently
finiteness of the projected photon rings and cancellation in the
full vacuum coefficient can coexist with an ultraviolet-divergent
primitive ring. A physical charged action must also supply amplitudes
and contacts outside this one-pair quadratic representation.

\subsection{A photon continuum prepared from the measured charged pair}
\label{sec:joint2-jacobi-uv}
The photon-ring obstruction in \eqref{eq:sew-dressed-photon-ring-uv}
assumes a free transverse photon with only a massless pole. A different
physical starting point is a massless pole accompanied by neutral
massive states that can mix with the charged pair. This admits a direct
test on the same sector plane. We first compute the complete photon
response and its rings, then resolve the normal states in its determinant.
The two computations use the same measured photon.

Write $s_0=4m^2$. Let $C_{\rm pair}(s)$ be the unit-charge light-pair
weight of \eqref{eq:nextsoft-positive-uv}. In the spectral convention
of \eqref{eq:sew-dressed-photon-ring-uv} set
\[
 d\kappa(s)= C_{\rm pair}(s)\,ds/s,\quad
 k_0=\int d\kappa(s)>0,\qquad d\mu=d\kappa/k_0 .
\]
This is the massive photon weight required by the transverse Maxwell
matrix element on these pair states. It has finite moments through at
least degree six: the argument of $\chi_{248}$ grows linearly in $s/m^2$
at high energy and its square decays as the inverse 249th power of
that argument. In $L^2(d\mu)$ let $M$ be multiplication by $s$,
$u_0=1$, and write $\langle\cdot\rangle$ for the $\mu$ expectation.
Two orthogonal normal modes and their next link are fixed by the measure:
\[
\begin{aligned}
 a_0&=\langle s\rangle,& b_0^2&=\langle(s-a_0)^2\rangle,\\
 u_1&=(s-a_0)/b_0,&a_1&=\langle s u_1^2\rangle,\\
 b_1u_2&=(s-a_1)u_1-b_0u_0,&\|u_2\|&=1 .
\end{aligned}
\]
Both $b_j$ are positive because the measured mass law has a continuum.
Let $P$ project onto $u_0,u_1$. The compression $PMP=J$ has diagonal
$a_0,a_1$ and off-diagonal $b_0$. The only coupling to the orthogonal
complement is $V=b_1(|u_1\rangle\langle u_2|+
|u_2\rangle\langle u_1|)$. Define
\[
 M_0=M-V=J\oplus M_{\rm tail},\qquad M_g=M_0+gV .
\]
The tail is an operator on the whole complement. All these operators
are self-adjoint on the domain of $M$ by bounded perturbation.
Both compressions of $M$ are bounded below by $s_0$, hence
$M_g\ge s_0 I$ for $0\le g\le1$ by convexity. Here $g$ turns on
one bond for the comparison. Full coupling means $g=1$.

\paragraph{The closed photon and its measured continuum}
For $X=|q_E|^2>0$ define
\begin{equation}
 \begin{aligned}
 R_g(X)&=\langle u_0,(X+M_g)^{-1}u_0\rangle,\\
 D_g(X)&=Z/X+k_0 R_g(X),\qquad Z>0 .
 \end{aligned}
 \label{eq:joint2-jacobi-photon-uv}
\end{equation}
$Z$ is the massless field normalization and may be set equal to one.
Every $D_g$ has an unsubtracted positive transverse spectral measure,
with precisely this massless residue. At $g=1$ its massive measure
is exactly $d\kappa$. The two free neutral masses squared at $g=0$
are the eigenvalues of $J$. This is the two-node Gauss rule for the
measure of $u_0$~\cite{JointGluingGolubWelsch}. In units $m=1$ the
nodes are approximately $10.14321$ and $36.16724$, with normalized
weights $0.810158$ and $0.189842$. These numbers follow from the
A4 light-pair measure; they are properties of the assumed free
preparation. They are not additional predicted particle masses.
The continuum of independent modes has been retained at zero coupling
and is then invisible to this particular probe.

The opening of the normal line can be checked explicitly. On the
upper half-plane put $t(z)=\langle u_2,(M_{\rm tail}-z)^{-1}u_2\rangle$
and $p(z)=(a_0-z)(a_1-z)-b_0^2$. Resolvent elimination gives
\[
 R_g(-z)=
 \frac{a_1-z-g^2b_1^2t(z)}
 {p(z)-(a_0-z)g^2b_1^2t(z)} .
\]
On the continuous part of the real spectrum where the boundary values
exist and the denominator is nonzero, its imaginary part is
$|g b_0 b_1/[p(s)-(a_0-s)g^2 b_1^2t(s+i0)]|^2
\operatorname{Im}t(s+i0)$. This is the amplitude to reach the tail,
times its adjoint, with the tail's own state density. Any isolated
eigenvalue must be included separately. At full coupling the spectral
theorem gives the whole measure, including any exceptional points,
directly from multiplication by $s$. The identification with charged
pairs is at the level of the cyclic, neutral spin-one channel spanned
by the measured pair vector. This construction does not specify the
propagation and interactions of its two charged constituents.

\paragraph{Every ring of the projected transverse action}
The action $\tfrac14\int F_{\mu\nu}[X D_g(X)]^{-1}F_{\mu\nu}$
has covariance $D_g$ on a physical polarization. It respects the
linear Ward identity. In $R_g$ the first four moments, of orders zero
through three, agree with those of $R_0$. A path of matrix elements
from $u_0$ to the tail and back has at least four steps. Equivalently
$M_gu_0=a_0u_0+b_0u_1$ and
\[
 \|M_g^2u_0\|^2-\|M_0^2u_0\|^2=g^2 b_0^2 b_1^2 .
\]
For a positive measure with finite fourth moment the resolvent
expansion through that moment follows by polynomial division and
dominated convergence on $X>0$. Therefore $R_g-R_0=
g^2 b_0^2b_1^2 X^{-5}+o(X^{-5})$. The relative quadratic insertion
$D_0/D_g-1$ is $O(X^{-4})$ at infinity and $O(X)$ at the origin.
Every fixed positive power is thus absolutely integrable against
the four-dimensional radial measure $X\,dX$, including the
one-insertion ring. Expanding in $g$ gives the same conclusion
coefficient by coefficient: in the Schur formula its successive
even coefficients at large $X$ carry additional factors
$b_1^2t(-X)\langle u_1,(X+J)^{-1}u_1\rangle=O(X^{-2})$.
This is a positive photon with a massive continuum and a finite
relative vacuum determinant in the projected quadratic action.
No momentum has been excluded.

\paragraph{What happens to the same bond in the resolved determinant}
If the normal coordinates are independent Gaussian fields, their
determinant must be varied too. Put $a(X)=\langle u_1,
(X+J)^{-1}u_1\rangle$. Per transverse component the full
relative determinant is exactly
\begin{equation}
 \det\big[(X+M_g)(X+M_0)^{-1}\big]
       =1-g^2 b_1^2 a(X)t(-X).
 \label{eq:joint2-jacobi-full-uv}
\end{equation}
This is a finite-rank determinant, even on the infinite normal space.
Both $Xa(X)$ and $X t(-X)$ tend to one. For a real Gaussian degree
of freedom the resulting quadratic vacuum coefficient per four-volume
has hard radial density $-g^2 b_1^2\,dX/(32\pi^2 X)$ to leading order.
Thus the visible photon passes the ring test while the resolved
transition between $u_1$ and $u_2$ already yields a bare logarithm.

One might try to suppress precisely this last transition by replacing
$g$ on the bond with $g f(X)$, where $f$ is real, bounded and the
resulting positive Euclidean inverses exist. The same determinant
can be finite in the ultraviolet at fixed $g\ne0$ only if
$\int^\infty |f(X)|^2\,dX/X$ is finite. This forces $f$ to tend
to zero along an unbounded sequence. The photon on that sequence
gives a decisive test of this repair. Indeed Schur elimination
yields the exact identity
\begin{equation}
 p(-X)^2[R_{gf}(X)-R_0(X)]
 =\frac{g^2 b_0^2 b_1^2 f(X)^2t(-X)}
          {1-g^2 b_1^2 f(X)^2a(X)t(-X)} .
 \label{eq:joint2-jacobi-identity-uv}
\end{equation}
If this massive remainder has a positive spectral measure $d\nu$,
its first four moments of orders zero through three are fixed by
the $O(X^{-5})$ bound to agree with the two-node reference.
This needs no moment assumption on $d\nu$: successive divisions
by $X+s$ followed by monotone convergence recover each next
moment from the positive remainder of the preceding division.
Polynomial division in $s$ then identifies the left side of
\eqref{eq:joint2-jacobi-identity-uv} with
$\int p(s)^2\,d\nu(s)/(X+s)$, because $p$ vanishes at both reference
nodes. Multiply by $X$ and pass to the sequence on which $f\to0$.
Monotone convergence gives $\int p(s)^2d\nu(s)=0$. A positive
measure must therefore be supported on the two reference nodes,
and its known first moments fix its weights. This rules out every
new positive pair band in this damped-bond construction, including
when the protected massless photon is retained separately.
In this class the missing physical ingredient is a change to the
resolved normal dynamics or its measure beyond a single attenuated
bond. The success of the projected quadratic action does not decide
whether the native GCM boundary implements that change.

\subsection{Can correlations in the photon measure supply the missing weight?}

\paragraph{A fourth-moment test for a non-Gaussian photon boundary}
There is a direct test that permits a general positive joint law
of the photon and its geometric records. In a finite Fourier
test projection use independent real coordinates as a choice of
basis, without assuming statistical independence. For each
selected physical polarization let $a_j$ be a real normalized
Fourier coefficient with $X_j=|q_j|^2>0$, and write
$Z_j=X_ja_j^2$. Let $V$ be the complete real first variation of
the negative log weight of a normalized source $e^{-gV+O(g^2)}$.
Assume square-integrable modal intensities and insertion, and
keep the measured field coordinates fixed as $g$ varies. For a
centered free photon the first covariance change is
$d_j=X_j\delta D_j=-\operatorname{Cov}(Z_j,V)$. This includes
variations of the action and boundary law whenever that score exists.
For the positive matrix $G_{ij}=\operatorname{Cov}(Z_i,Z_j)$,
\begin{equation}
 d\in\operatorname{ran}G,\qquad
 d^{\mathsf T}G^+d\le\operatorname{Var}V,\qquad
 \|d\|^2\le\|G\|\operatorname{Var}V .
 \label{eq:joint2-intensity-bound-uv}
\end{equation}
Here $G^+$ is its Moore--Penrose inverse. Project
$V-\mathbb EV$ orthogonally onto the span of the centered
intensities in $L^2$ of the joint law. The squared norm of the
projection is precisely the middle expression on the left.
This proof uses neither Wick factorization nor a Gaussian shift.
The connected two-insertion coefficient is the variance when the
two copies of $V$ are marked independently. A source containing
a second derivative of the action has that additional contact
in its second derivative with respect to a single coupling.

Suppose the variance is bounded per unit four-volume in the
large-box limit and that the norm of $G$ is bounded uniformly
there and over ultraviolet test projections. Dividing the last
inequality by volume, then taking the infinite-volume limit on
each fixed annulus, bounds $\int X|X\delta D(X)|^2dX$ for an
isotropic covariance change. A nonzero positive massive band
of the leading photon spectral change, with the free spectrum
purely massless, contradicts this bound by the monotonicity
argument of \eqref{eq:sew-dressed-photon-ring-uv}. In fact the
product of $\|G\|$ and the variance density must grow at least
logarithmically on exhausting four-momentum annuli whenever the
finite-volume derivative reproduces that continuum response.
Thus a non-Gaussian gluing with a regular score needs correlations
of the intensities strong enough to violate this uniform bound,
or a genuine extra variation of the physical readout.

\subsection{What changes when the charged pair is opened a second time?}
\label{sec:next3-map-uv}
Two moves that look innocuous at the level of a two-point function
change the ultraviolet and physical-state tests. The first is to
replace a measured pair of fermions by an elementary neutral mode of
the same invariant mass. Two such modes can propagate independently;
two pairs of fermions can also exchange partners. The second is to
suppress the spatial center momentum of a pair in the clock frame
selected for a calculation. A pair can reach arbitrarily large center
momentum by a boost while its invariant mass and relative threshold
stay fixed. The shell distribution in the relative variable alone
then leaves an entire hard integration uncontrolled.

The simultaneous-duration criterion in \eqref{eq:next3-joint-record-uv}
shows why assigning three charged pairs a common A4 clock can lose
ultraviolet convergence even after fermionic exchange.

The following calculations start on the same aligned sector plane
$\bin=(s,b_s)$, with geometric carrier $\sigma_\epsilon\hat p_{0,\bin}$
and matter carrier $\sigma_s\sigma_\epsilon\hat p_{0,\bin}$. Positive
Hamiltonian time below refers to the chosen clock readout after locking.
The spatial averages introduced in a chosen frame do not identify that
frame with an observer-independent geometric-time direction. In a
spectral calculation the symbol $s$ used as an integration variable
denotes squared mass; the sector label stays fixed. Real-energy cuts
use the declared vectorial boundary on that matter carrier.

Section~\ref{sec:next3-car-uv} tests whether the measured pair can
be opened at actual charged endpoints. Section~\ref{sec:next3-dipole-uv}
constructs a Hamiltonian in which the same electric coupling fixes
both that transition and the photon contact, then computes the cost of
boosting a pair in a fixed invariant-mass band. The subsequent current
and vector tests identify which gauge contacts and which physical states
have to enter a covariant construction. A spatially prepared neutral
interaction and a charged velocity-transport determinant provide two
further tests directly on the geometric histories and their measures.

\subsection{What the neutral pair mode fixes on a charged Fock space}
\label{sec:next3-car-uv}
An observed pair already contains two fermions. Opening two copies
of its line can therefore exchange one constituent between them.
This gives a test of the inverse construction
\eqref{eq:joint2-inverse-square-uv} before a dynamical photon is sewn.
On the reconstructed plane let $\mathfrak h_+,\mathfrak h_-$ be the
physical one-particle and antiparticle spaces. A normalizable pair
packet $K$ is a Hilbert--Schmidt map from the conjugate antiparticle
space to the particle space. In orthonormal bases set
$B^\dagger(K)=\sum_{ij}K_{ij}a_i^\dagger d_j^\dagger$ and
$B(K)=B^\dagger(K)^\dagger$. Particle and antiparticle operators
obey the canonical anticommutation relations. Both endpoints carry
the original charged sector label. Smearing the transverse pair row
of Section~\ref{sec:fix-fock-uv} gives such packets.

\paragraph{An exchange test with arbitrarily many pairs}
The overlap of two exponential pair sources is exactly
\begin{equation}
 \langle e^{B^\dagger(K)}0,e^{B^\dagger(L)}0\rangle
       =\det(1+K^\dagger L).
 \label{eq:next3-car-det-uv}
\end{equation}
Here and below $0$ inside a state denotes the vacuum. At degree $r$
the amplitudes are the $r$-rowed minors of the pair matrix. Summing
their products gives the determinant by Cauchy--Binet, first at
finite rank, then by Hilbert--Schmidt approximation and trace-norm
continuity of the Fredholm determinant. In particular, if $p_j$
are the squared singular values of $K$, then
\[
 \|(B^\dagger(K))^r|0\rangle\|^2=(r!)^2 e_r(p_1,p_2,\ldots),
\]
where $e_r$ is the elementary symmetric polynomial. At $r=2$ the
answer is $2[(\operatorname{tr}K^\dagger K)^2-
\operatorname{tr}((K^\dagger K)^2)]$. The second term is absent
for an elementary Gaussian pair field with the same vacuum
two-point function. It is strictly nonzero for every nonzero packet
and cancels the entire two-pair norm for a rank-one packet. The
spectral density at one pair supplies a squared norm on mass fibers.
It does not determine these four-constituent exchange terms.

\paragraph{A sharp simultaneous-duration test with physical Pauli exchange}
There is a boundary norm that can be evaluated for an arbitrary joint
law of the durations. On a spatial three-torus in the chosen frame,
take one spin slot of each charge, opposite momenta, and the positive
boundary heat amplitude $K_t(p,-p)=\exp[-t(m^2+|p|^2)]$. Here $p$
ranges over the entire reciprocal lattice. Each $t>0$ gives an actual
normalizable pair packet. The diagonal heat map and the choice to take
a positive mixture of its squared norms specify this test. Let
$\nu_r$ be the joint law of $r$ recorded durations $0<t_i\le L^2$.
For $N_r=\|B^\dagger(K_{t_1})\cdots B^\dagger(K_{t_r})|0\rangle\|^2$
there are positive constants at fixed volume, $mL$ and $r$ such that
\begin{equation}
 c_r\prod_{i=1}^r (t_i/L^2)^{-3/2}\ \le\ N_r\ \le\
 C_r\prod_{i=1}^r (t_i/L^2)^{-3/2} .
 \label{eq:next3-joint-record-uv}
\end{equation}
For the upper bound, the amplitude for any set of distinct occupied
momenta is the sum over assignments of these momenta to the $r$
recorded pairs. Pair creators in different slots commute. All
assignment amplitudes have the same sign after both species are
ordered. Cauchy--Schwarz on at most $r!$ assignments, followed by
summing the squared amplitudes, gives at most $r!\prod_i\sum_p
\exp[-2t_i(m^2+|p|^2)]$. For the lower bound split the lattice into
$r$ cosets by the first momentum index modulo $r$. Require the pair
at slot $i$ to use the $i$th coset and retain that assignment. Its
squared norm is the product of the restricted sums. Other assignments
in the same final state can only increase the norm. Each coset heat
sum is comparable to $(t/L^2)^{-3/2}$ throughout $0<t\le L^2$.
Indeed the small-$t$ Gaussian Riemann sum on a coset has $1/r$ of
the full positive leading volume coefficient. Positivity and continuity
on the remaining compact interval give uniform lower and upper bounds.
This proves both inequalities without discarding the exchange terms.

Consequently the prepared $r$-pair norm is finite precisely when the
joint negative moment on the right of \eqref{eq:next3-joint-record-uv}
is integrable against $d\nu_r$. For example, use the two-record A4
duration of \eqref{eq:joint2-duration-uv}, whose lower-tail power
is four. If the same duration is reused at every pair, the integral
near zero is proportional to $\int_0^1 u^{3-3r/2}du$. It is finite
for one and two pairs and diverges already for three. With independent
copies of exactly the same marginal it is finite for every fixed $r$.
The difference comes entirely from which intermediate displacements
the measure retains. In particular the one-pair spectrum and Pauli
blocking do not by themselves control simultaneous short durations.
This is an exact charged-boundary test with zero total pair momenta.
Applying it to a proposed GCM representation requires deriving its
actual amplitude and joint law, including any coherent interferences
between different records, from the source before taking a norm.

\paragraph{Which mass rotations lift to charged one-body transformations?}
Consider a constant-profile copy of $U_t=\exp(tJ)$ in
\eqref{eq:joint2-fixed-null-uv}, acting on the relative pair state
at each total spatial momentum, with the other pair channels fixed.
Take the same bounded generator in the fibers. This translation-
invariant rotation is norm-continuous at zero and mixes different
invariant masses since $[T,J]\ne0$. A field-linear CAR transformation
that preserves the empty vacuum and the charge acts on one pair as
$V_+\otimes V_-$: a creator component of a transformed annihilator
would create a state from the vacuum. Charge conservation prevents
the two one-body species from mixing. Commutation of the tensor
product with translations implies
\[
 V_\pm^\dagger e^{i a\cdot\mathbf p}V_\pm
       =e^{\pm i a\cdot d}e^{i a\cdot\mathbf p} .
\]
Indeed two unitary tensor factors with product the identity must be
reciprocal scalars, and their continuous characters in $a$ determine
the vector $d$. A nonzero shift moves a packet of sufficiently narrow
momentum support to an orthogonal packet. Its tensor product action
is then at least $\sqrt{2}$ from the identity in operator norm.
For small $t$ this is impossible. With $d=0$ both momenta are
preserved and the free pair mass, a scalar function of these momenta,
is preserved too. Thus no such field-linear lift realizes this
mass rotation. The argument allows arbitrary spin rotations at
each constituent momentum.

There is an explicit lift with a two-body generator. In the sector
with $n_+$ particles and $n_-$ antiparticles use
$\mathbb J=\sum_{i=1}^{n_+}\sum_{j=1}^{n_-}J_{ij}$, with $J_{ij}$
acting on that pair of constituents, and restrict to the separately
antisymmetric spaces. The sum respects both antisymmetries and is
skew-adjoint, of norm at most $n_+n_-\|J\|$ in that sector. Its
sectorwise exponential defines a unitary on the full Fock space.
It gives exactly $U_t$ on one pair and fixes the single-charge
sectors. On two pairs it acts on all four possible choices of
partners. Conjugation by this unitary preserves the CAR, but acts
nonlinearly on individual charged operators. Conjugating a free
Hamiltonian and every probe together is a change of representation.
A new physical interaction requires a change in their relative action.

\paragraph{Two positive charged sources with identical optical pair data}
An explicit Hamiltonian comparison shows how much physical information
is still missing. Take the tensor product of the free transverse photon
and the massive charged Fock space. Use as optical probe the sum of
$\sqrt Z$ times the free photon and a pair operator plus its adjoint.
On pair mass $s$ take its vacuum matrix element to be $g h_0(-s)/s$
times the transverse Dirac current matrix element. Only the timelike
pair block occurs in this term, so $s\ge4m^2$. Its dispersion is
that of an observable evolved by the Hamiltonian. The spin sum of
the Dirac transverse current in the pair rest frame contributes
$s\rho_D(s)$, giving photon weight $g^2\rho(s)ds/s$ with
$\rho=|h_0|^2\rho_D$. With $Z=1-g^2M_{-1}>0$ as above, this is
precisely the complete positive quadratic photon law of
\eqref{eq:joint2-inverse-spectrum-uv}. Charged endpoints are the
actual CAR creators and annihilators. The comparison has a positive
physical Hilbert space on the chosen clock slice.

One may now keep the free Hamiltonian or add the following positive
term on every charged sector. Sum over unordered quartets $I$ of two
particles and two antiparticles, and in their equal-clock-time position
variables use
\begin{equation}
 V_4=m\sum_I\int_0^1
   \exp\!\left[-\frac{u}{L^2}\sum_{a<b\,\in I}
             |\mathbf x_a-\mathbf x_b|^2\right]du .
 \label{eq:next3-car-quartet-uv}
\end{equation}
The uniform $u$ is the A4 projection fraction. Coupling it to a
quartet separation in this manner is an explicit choice. This
bounded perturbation in each sector gives a positive self-adjoint
direct sum with the free Hamiltonian. It preserves translations,
rotations and charge. The quartet term tends to zero upon separating
its constituents into distant clusters. The vacuum, each single-charge
propagator and the entire optical two-point function are identical
for the two choices. Nevertheless take a smooth finite-energy pair
packet for which $Y=(B^\dagger(K))^2|0\rangle\ne0$. In the
pair-number-resolved four-source test the difference of the negative
right derivatives at $\tau=0$ of $\langle Y,e^{-\tau H}Y\rangle$
is exactly $\langle Y,V_4Y\rangle>0$. This is a measurable difference
at four insertions among sources that agree on the full one-pair band.

Every quartet factor is an average of normalized Gaussian momentum
transfers (one independent Gaussian per pair difference at fixed $u$).
These measures have full relative support and total mass one. Thus
the absolute transfer operator on each fixed charged sector is bounded
by the number of its quartets times $m$. At fixed order a smooth
external pair insertion or removal has absolute-kernel norm bounded
by $\sqrt{(n_++1)(n_-+1)}\,\|K\|_{\rm HS}$ for creation and the
adjoint bound for removal. These bounds follow by antisymmetrizing
the tensor product with the packet. Before antisymmetrization the
absolute permutation sum gives the same bound by the triangle
inequality. The free Euclidean propagation factors are bounded by
one in momentum space. Only finitely many charged sectors are
visited at fixed order, so each permutation contraction on a finite
time interval with smooth probe packets and open endpoints has an
absolute momentum bound. Positive transfer holds for the exact two
Hamiltonians. This example leaves boost covariance of the quartet,
local electromagnetic coupling of a single charge, and dynamical
photon sewing open. It shows that even positive reconstruction and
knowledge of every one-pair mass do not select the physical
four-charge part of the GCM boundary law.

\subsection{A charged electric polarization with its canonical photon contact}
\label{sec:next3-dipole-uv}
A transverse photon can couple to a true fermion pair through an electric
polarization. In Hamiltonian variables the same coupling shifts the
canonical electric momentum. This fixes a quadratic contact as well as
the measured current and permits an interacting positive-state test.
Work in one chosen carrier-time frame on the locked sector plane. Start
with the free Maxwell field and the actual particle and antiparticle CAR
spaces of a Dirac field of mass $m$. Use the sector orientation of the
matter carrier. Write $E_p=(m^2+|p|^2)^{1/2}$ and use unit-normalized
spinors in the one-particle spatial momentum measure. Set
$h(k)=\chi_{248}(L|k|)$, $Lm=1$, and keep the invariant pair mark
$h_0(-s)$ of \eqref{eq:nextsoft-product-uv}. For $r=k-p$, put
$w=E_p+E_r$, $s=w^2-|k|^2$ and
$j_i(p,r)=u(p)^\dagger\alpha_i v(r)$, including both spin indices.
Define the pair-creation part of a polarization by
\begin{equation}
 P_i^+(k)=h(k)\int_p \frac{h_0(-s)}{i w}\,
                j_i(p,k-p)b_p^\dagger d_{k-p}^\dagger .
 \label{eq:next3-dipole-uv}
\end{equation}
Set $P_i(k)=P_i^+(k)+P_i^+(-k)^\dagger$. The time-local pair operator
uses the energies of its two constituent modes. The inverse $w$ converts
a polarization into a current. It is bounded since $w\ge2m$.
The center mark $h$ is the three-spatial-coordinate readout of one A4
record. Its placement and the electric polarization are representation
data. The CAR spaces retain individual charged endpoints and all
fermionic exchange signs at higher multiplicity.

Couple the polarization to the full electric field and use the
free-charge density $\rho_f(k)=h(k)\rho(k)$, where $\rho$ counts
particles minus antiparticles as in \eqref{eq:positive-gauss-comparison-uv}.
In the longitudinally dressed charged coordinates used there, the
instantaneous scalar source generates this free charge. Eliminating
Gauss' law gives the reduced Hamiltonian
\begin{equation}
\begin{aligned}
 H(e)={}&H_F+H_\gamma+e\langle\Pi_T,P_T\rangle
                 +\tfrac{e^2}{2}\|P_T\|^2\\
 &+\tfrac{e^2}{2}\|(-\Delta)^{-1/2}
                    (\rho_f-\nabla\cdot P)\|^2 .
\end{aligned}
\label{eq:next3-dipole-ham-uv}
\end{equation}
Inner products integrate the spatial components and each operator square
is the squared norm on a state. The mixed transverse term is Hermitian
because the matter and photon time-zero operators commute. The first
three photon terms are the free Maxwell form with its electric momentum
shifted from $\Pi_T$ to $\Pi_T+eP_T$, relative to the free photon vacuum
energy. The last square contains the free-charge self-energy. Integration
of a polarization divergence gives zero, so the zero-transfer charge
on a single particle is $e$. The longitudinal dressing specifies how
to lift the reduced source under Abelian gradient shifts. An elastic
minimal Dirac transverse vertex would be additional dynamics.

At first electric order the current that creates a pair is $\dot P_i$.
Indeed a classical transverse source enters the action as an electric
field, and integration by parts differentiates $P_i$. The free vacuum to
pair matrix element of its time derivative is $iw$ times
\eqref{eq:next3-dipole-uv}. Thus the measured transverse Born
amplitude is exactly $e h(k)h_0(-s)j_i$. The bound-charge density
$-\nabla\cdot P$ gives its conserved temporal component by the free
Dirac pair identity. In the infinite-volume free phase-space integral
the transverse pair weight is therefore $e^2|h(k)|^2 C_{\rm pair}(s)$
in the normalization of \eqref{eq:nextsoft-positive-uv}. In the pair
rest frame this is the full prescribed weight at every finite mass.
The order-$e^2$ absorptive part of the response to the same probe has
precisely these intermediate fermion states. The fixed canonical
contact contributes no finite-energy discontinuity. At nonzero total
spatial momentum the factor $|h(k)|^2$ changes under boosts at fixed
$s$. This gives an explicit failure of invariant optical normalization,
even when the rest-frame spectrum agrees exactly.

For an independent ultraviolet and positivity result put the system
on a fixed spatial torus, omit the free harmonic photon mode and use
the zero-mean Coulomb inverse. One may restrict to total charge zero
or use a compensating homogeneous charge for the zero mode in charged
sectors. Spatial momenta still have unbounded range. For fixed $k$,
sum the absolute pair-creation matrix entries. On
$|p|\le2(|k|+m)$ their sum before $h(k)$ is $O(\langle k\rangle^3)$
with momenta in carrier units. Outside this ball $s\ge c|p|^2$ and
the invariant A4 ball decays as $O((1+s/m^2)^{-249/2})$.
Thus $C|h(k)|\langle k\rangle^3$ bounds both the absolute sum and
the operator norm, by the bounded CAR generators. This envelope
remains summable after multiplication by $1+\sqrt{|k|}$. The linear
photon term is therefore infinitesimally form bounded relative to
$H_\gamma$, by the creation and annihilation estimates. The longitudinal
square is a closed nonnegative form, using the closure of its charge
map on finite-particle tests. Closability follows by testing its
adjoint against finitely many spatial modes and finite-particle vectors.
The sum defines a semibounded Hamiltonian with compact resolvent since
its form dominates $\tfrac12(H_F+H_\gamma)-C_e$. Vacuum correlations
in any ground state have a positive time-reflection form by the
semigroup inner-product identity. This assertion concerns the physical
Hilbert-space observables with the Hamiltonian adjoint under reflection.

Every ordered contraction on a finite time slab also has an absolute
spatial momentum bound. At fixed perturbative order and finitely many
external field insertions only finitely many particle counts occur.
The pair bound above applies with individual matrix entries replaced
by their absolute values. Each free-charge insertion is a sum of at
most the particle count of unit momentum translations multiplied by
$|h(k)|$. The Coulomb and transverse squares then have summable transfer
kernels, including their equal-time contractions. A photon field supplies
at most $C\sqrt{|k|}$ and a bounded occupation factor per marked end.
Free Euclidean propagation between ordered times has norm at most one.
Iteration yields an absolute bound for each individual contraction,
including closed fermion circulations with no external charged endpoints.
This argument uses the canonical ordered representation and free-vacuum
normalization. A Lorentz-invariant infinite-volume interacting charged
source, including its elastic current and the boost test above, remains
to be obtained from the physical GCM.

\paragraph{The complete bare vacuum coefficient under boosts of a fixed pair}
There is a quantitative test of the cost of restoring the invariant
pair normalization in this Hamiltonian. Allow any common real bounded
radial center mark $h$ on the pair, keep the measured invariant density
$\rho(s)=C_{\rm pair}(s)$, and write $r=|k|$, $w=\sqrt{r^2+s}$.
In the continuum free spectral normalization $ds/(2w)$ the spatial
current tensor on a mass shell is $(s\delta_{ij}+k_i k_j)\rho(s)$.
Dividing by $w^2$ for the two polarization factors, the vacuum squared
norm per $ds$ with the common $|h(k)|^2$ stripped is
$s\rho(s)/(2w^3)$ for each transverse polarization and
$\rho(s)/(2w)$ for the longitudinal one. The latter contributes
directly to the Coulomb square since $\rho_f(k)$ annihilates the
free vacuum. The transverse contact in each mode has coefficient
$e^2/2$. The two linear insertions give the second-order connected vacuum
coefficient with opposite sign and magnitude
$e^2 r/[2(w+r)]$, since the created state has one photon and one pair.
They combine to the positive coefficient $e^2w/[2(w+r)]$ before any
momentum integration. Per unit volume the canonical order-$e^2$ bare
vacuum-energy coefficient, including its factor of $e^2$, is therefore
\begin{equation}
 \frac{e^2}{4}\int\!\frac{d^3k}{(2\pi)^3}|h(k)|^2
  \int_{4m^2}^\infty\!ds\,\rho(s)
       \left[\frac{1}{w}+\frac{2s}{w^2(w+r)}\right].
 \label{eq:next3-dipole-energy-uv}
\end{equation}
One can obtain this coefficient from the connected finite-time vacuum
expansion by dividing by its duration and spatial volume. For a created
state of energy $w+r>0$, the ordered time integral per unit duration
tends to $1/(w+r)$. Finite test sums followed by exhaustion on the
nonnegative combined kernel give the displayed integral. The calculation
requires only this perturbative coefficient and makes no assumption of
analyticity of the exact ground-state energy.
For $M_j=\int s^j\rho(s)ds$ the A4 invariant mark makes
$0<M_0,M_1<\infty$. The first inner term is asymptotic to $M_0/r$
and the second to $M_1/r^3$, by dominated convergence. Hence the
coefficient is ultraviolet finite exactly when
$\int^\infty r|h(r)|^2dr<\infty$. Keeping $h=1$ restores the desired
boost behavior of the leading pair current but gives a quadratic
ultraviolet divergence. Even its completed transverse contribution
alone is logarithmic. A compact invariant mass band of positive weight
already supplies both lower bounds: arbitrarily large momenta occur
by boosting pairs of that same invariant mass. This test concerns the
bare vacuum energy. A vacuum-normalized connected functional can omit
that energy; it does not make the underlying coefficient finite.

\paragraph{Coupling each open charge to the same transverse photon}
There is a minimal kinetic completion of the individual charged
endpoints in \eqref{eq:next3-dipole-ham-uv}. For the particle and
antiparticle signs $a=\pm1$ replace the free one-body energies by
\[
 K_a(A)=\sqrt{m^2+(-i\nabla-ae A_h)^2},\qquad
       \widehat A_h(k)=h(k)\widehat A(k) .
\]
Sum this operator over occupied spin slots. Its square root is the
positive spectral square root at each equal-time photon configuration.
Prior to transverse reduction it is covariant under multiplication by
$\exp(iae h*\alpha)$ for $A\mapsto A+\nabla\alpha$, with $h*$ here
meaning convolution by the spatial readout. Use the same longitudinal
dressing for the pair operator as above. In the number-preserving
sector the linear kinetic vertex is
\[
 -ae h(k)\frac{p'+p}{E_{p'}+E_p},\qquad p'=p+k .
\]
Its contraction with $k$ gives $-ae h(k)(E_{p'}-E_p)$ as required by
the free-charge density. Its zero-transfer limit is the charge times
the relativistic velocity with the minimal-coupling sign. This
spin-independent kinetic action leaves the finite-transfer magnetic
spin current to be specified. Its contacts are fixed by $K_a(A)$.

They admit an order-independent routing estimate. With
$G_t(p)=(m^2+|p|^2+t)^{-1}$, expand the positive-square-root
resolvent formula prior to photon contraction. A term with $l$
linear kinetic insertions and $d$ direct quadratic insertions
contains $l+d+1$ resolvents and weight $t^{1/2}dt/\pi$. Each
linear spatial numerator satisfies
\[
 (|p'|+|p|)G_t(p')^{1/2}G_t(p)^{1/2}
        \le 2(m^2+t)^{-1/2}.
\]
After allocating half resolvents to each adjoining insertion,
the term of photon degree $N=l+2d$ is bounded by
$C_N\int_0^\infty t^{1/2}(m^2+t)^{-1-N/2}dt$ times the
external readouts. This is finite for every $N\ge2$; the displayed
one-photon divided difference treats $N=1$. All bounds are uniform in
the charged momenta. Each canonical photon brings its summable
$|h(k)|/\sqrt{2|k|}$ on the nonzero torus modes. An elastic insertion
preserves charged count and translates an occupied momentum. Thus
the ordered graph bounds include all these new contacts and their
internal equal-time contractions. Every pure kinetic interaction
annihilates the free charged vacuum. By charged-count orthogonality,
its cross term with the single pair vertex also vanishes in the
order-$e^2$ vacuum coefficient. Equation~\eqref{eq:next3-dipole-energy-uv}
therefore continues to hold.

The compact-volume positivity statement survives this completion.
The kinetic sum is at least $mN_{\rm ch}$ on each photon configuration,
where $N_{\rm ch}$ counts particles and antiparticles. The full form
dominates that sum plus $H_\gamma/2-C_e$. To see compactness, truncate
first the charged count at $J$ in the compactness argument. The omitted
squared norm on a bounded-energy set is $O(1/J)$ by this lower bound.
On the remaining sectors the ordinary free kinetic form is controlled:
at fixed configuration $p^2\le2(p-ae A_h)^2+2e^2\|A_h\|_\infty^2$
as forms, and the scalar supremum commutes with the one-body operator.
Monotonicity of the square root bounds each $|p|$ by a constant
times $K_a(A)+|e|\|A_h\|_\infty$. The latter spatial supremum has
expectation bounded by $C\langle H_\gamma+1\rangle^{1/2}$ on unit
vectors, by summing the absolute readout Fourier quadratures.
Both the kinetic sum and $H_\gamma$ preserve the count sectors.
The free form has compact resolvent there, which proves compactness
of the full form embedding. The count truncation is removable and
introduces no restriction on the charged Hilbert space.

\paragraph{A Gauss-law test independent of the chosen polarization}
Part of the preceding obstruction does not depend on the form of
$P_i$. Suppose that in a fixed carrier frame the weak-charge source
obeys the ordinary Maxwell constraint and uses its positive Coulomb
energy. Suppose its measured leading conserved current has a
Lorentz-covariant positive vacuum spectrum with scalar mass measure
$d\nu$ of weight $w_*>0$ on a finite interval $[s_1,s_2]\subset(0,\infty)$.
These hypotheses concern the complete measured current, including its
actual boundary states. Restrict that current, applied to the vacuum,
to the stated invariant-mass interval by its orthogonal spectral
projection. At spatial momentum $k$ its density norm squared per
unit volume is at least
\[
 |k|^2\int_{s_1}^{s_2}\frac{d\nu(s)}{2\sqrt{|k|^2+s}}
 \ge\frac{|k|^2w_*}{2\sqrt{|k|^2+s_2}} .
\]
This is the $00$ component of the transverse spectral tensor. The
Coulomb kernel cancels $|k|^2$. Hence the bare order-$e^2$ expectation
of the Gauss square per unit volume is at least
\[
 \frac{e^2w_*}{4}\int\frac{d^3k}{(2\pi)^3\sqrt{|k|^2+s_2}}
       =+\infty .
\]
If the density exists only as a distribution in time, test it on a
unit-integral smooth time packet approaching a delta. Its spectral
norm gains the nonnegative factor $|\widehat a(\delta\sqrt{|k|^2+s})|^2$,
which tends pointwise to one. Fatou's lemma gives the same divergent
lower limit. Extra positive sectors and the sign of a transition
amplitude cannot remove this individual positive contribution. To
apply the test one must verify that this Gauss contraction is indeed
part of the bare physical source. Other action-prescribed terms can
affect its sum with transverse graphs. Equation~\eqref{eq:next3-dipole-energy-uv}
computes that sum in the electric-polarization realization and shows
there that it still diverges when the center mark is removed.

\subsection{The elastic Dirac current and the algebra of the charge density}
\label{sec:next3-gauss-car-uv}
The charge counted at zero momentum fixes one number. A local Abelian
generator must also give the correct composition of two independent
spatial gauge variations. The distinction can be tested with the same
constituent spinors used to open the observed pair.

\paragraph{A spin-sensitive elastic vertex}
In the positive charged Hamiltonian above one can use the Pauli square
root $[m^2+(\boldsymbol\sigma\cdot(-i\nabla-ae A_h))^2]^{1/2}$ on
the two spin components of each sign $a$. The squared Pauli operator
is positive. Its one-photon matrix element, for polarization $\epsilon$,
is $-ae h(k)$ times
\[
 \frac{(p'+p)\cdot\epsilon+
             i\boldsymbol\sigma\cdot[(p'-p)\mathbin{\times}\epsilon]}
                {E_{p'}+E_p},\qquad p'=p+k .
\]
When $E_{p'}=E_p$ this is the exact spatial elastic Dirac current
in the unit-normalized positive-energy spinors, with the indicated
charge and readout. Indeed the upper and lower spinor components
then give $(\boldsymbol\sigma\cdot p')\boldsymbol\sigma\cdot\epsilon+
\boldsymbol\sigma\cdot\epsilon\,\boldsymbol\sigma\cdot p$ over $2E_p$.
This includes the elastic Breit frame. The low-momentum Zeeman energy
is $-ae\boldsymbol\sigma\cdot B_h/(2m)$, the Dirac magnetic moment.
The absolute resolvent estimates for the scalar square root carry over
because $\|\boldsymbol\sigma\cdot p\|=|p|$. The same norm identity
bounds the ordinary kinetic form by the magnetic one plus the smoothed
field supremum on each count sector. These statements concern the
spatial current. The time component represented by the number-density
translation still differs from $u(p')^\dagger u(p)$ at finite transfer.

\paragraph{Two opposite gauge transfers give an operator test}
Suppose one attempts to restore the complete Dirac elastic density
while retaining the A4 weight only on the pair transition. Use the
free Dirac one-body space before passing to the filled-sea Fock
representation, with energy projections
$P_\pm(p)=[I\pm(\boldsymbol\alpha\cdot p+\beta m)/E_p]/2$.
Let $R_k$ shift $p$ to $p'=p+k$. Its matrix from band $j$ to band $i$
is $P_i(p')P_j(p)$ for $i=j$ and $a(p',p)P_i(p')P_j(p)$ otherwise,
where
\[
 a(p',p)=h_0(-s),\qquad s=(E_{p'}+E_p)^2-|p'-p|^2 .
\]
The common real spatial mark $h(k)$ is stripped from both bands.
This is precisely the time component of the proposed pair current
combined with the ordinary elastic density. It has $R_k^\dagger=R_{-k}$
and the correct unit charge at $k=0$. For an external positive-energy
spinor define the loss in the intermediate negative-energy channel by
\[
 d(p',p)=\frac{1-a(p',p)^2}{2}
       \left(1-\frac{m^2+p'\cdot p}{E_{p'}E_p}\right).
\]
The exact positive-band compression of $R_{-k}R_k$ at $p$ is
$[1-d(p+k,p)]P_+(p)$. Indeed the intermediate projections sum
to the identity at unit weight. The Clifford anticommutator gives
$P_+(p)P_-(p')P_+(p)=[1-(m^2+p'\cdot p)/(E_{p'}E_p)]P_+(p)/2$.
In the reverse order the answer contains $d(p-k,p)$. In particular
\begin{equation}
 P_+[R_{-k},R_k]P_+
       =[d(p-k,p)-d(p+k,p)]P_+(p).
 \label{eq:next3-charged-algebra-uv}
\end{equation}
Put $p=k$ and let $|p|/m$ grow. Both invariant masses in the two
weights tend to infinity, so both $a$'s tend to zero. The directional
spin factor tends to zero for $p'=2p$ and to $1/2$ for $p'=0$.
The commutator consequently approaches $P_+(p)/2$ in the sense of
its scalar coefficient. Restoring the common $|h(k)|^2$ leaves a
nonzero commutator at arbitrarily large finite transfers away from
its zeros. Continuity gives the same strict test on wave packets.
At $p=0$ the two losses agree. The failure is thus momentum-dependent;
a central term in a normal-ordered density algebra cannot cancel it.
Real sine and cosine gauge parameters test the same commutator.

\paragraph{The accompanying longitudinal contact is calculable}
The obstruction applies when these densities alone are the matter
parts of field-independent local Gauss generators. In a polarization
Hamiltonian the measured bound-charge density and the generator of
gauge transformations have different roles. One can display the
additional contact that a field-dependent generator would need.
On finite smooth gauge tests let $R_\alpha$ be the smeared Hermitian
density and $\chi_A=\Delta_{\rm sp}^{-1}\nabla\cdot A$ on the nonzero
spatial modes. A shift $A\mapsto A+\nabla\alpha$ sends $\chi_A$ to
$\chi_A+\alpha$. Conjugate $-i\partial_\alpha$ on this gauge orbit
by $U[A]=\exp(-ie R_{\chi_A})$. The exact generators commute by
conjugacy and have expansion
\begin{equation}
 U(-i\partial_\alpha)U^\dagger=-i\partial_\alpha+eR_\alpha
           -\frac{i e^2}{2}[R_{\chi_A},R_\alpha]+O(e^3).
 \label{eq:next3-charged-connection-uv}
\end{equation}
The two derivatives of the displayed contact cancel the order-$e^2$
commutator in \eqref{eq:next3-charged-algebra-uv}. This construction
solves the one-body longitudinal consistency problem on its test
domain. It gives an explicit constraint on a proposed lift to the
charged determinant and its boundary measure. As a conjugation on the
gauge orbit it supplies no transverse two-photon response, and does
not alter the bare boost integral in \eqref{eq:next3-dipole-energy-uv}.

\subsection{Which vector states propagate inside the neutral pair loop?}
\label{sec:next3-proca-uv}
The fixed null direction in \eqref{eq:joint2-fixed-null-uv} protects
the photon pole in perturbation theory. The preceding positive spectrum
at constant mixing fixes the propagator between conserved probes. A
space-dependent mass rotation also probes the longitudinal component of
a massive internal vector. This distinction can be tested at the first
nonconstant mixing coefficient, before any all-orders argument.

\paragraph{A discontinuity below the first allowed pair threshold}
Retain the notation $T,u,J$ of \eqref{eq:joint2-fixed-null-uv}, let
$T_t=e^{-tJ}Te^{tJ}$, and put $K=[T,J]$, $D=[[T,J],J]$. The massive
squared masses $s,r$ range over $[t_0,\infty)$, $t_0=4m^2$.
Write $C_s(q)=P(q)/(|q|^2+s)$ for the internally projected vector
covariance. At two separated scalar mixing probes the connected
correlation of $O=\langle b,Kb\rangle/2$ has the bubble
\[
 \frac12\int |K(s,r)|^2\,ds\,dr
   \int\frac{d^4q}{(2\pi)^4}
       \operatorname{tr}\{C_s(q)C_r(q+p)\} .
\]
The scalar contact containing $D$ changes the coincident response. The
complete twice-varied relative trace can equivalently be taken first,
so no individual divergent real part is assigned to this expression.
The tensor identity with $x=|q|^2$ is
\begin{equation}
 C_s(q)=\frac{I+qq^{\mathsf T}/s}{x+s}
                    -\frac{qq^{\mathsf T}}{sx} .
 \label{eq:next3-projector-pole-uv}
\end{equation}
The last term drops out between conserved external vectors. In a mass
vertex with a nonzero momentum transfer the two internal vectors have
different momenta, so this term can contribute at both ends. On the
continued invariant $0<S<t_0$ only the two zero-mass denominators
can pinch. On their common shell the contracted squared numerator is
$S^2/(4sr)$. With spectral density normalized as $\pi^{-1}$ times
the imaginary part at $p_E^2=-S-i0$, the massless scalar bubble has
constant density $1/(16\pi^2)$. This also follows by continuing its
Feynman-parameter logarithm. The cut of the complete response in the
connected-current sign convention is therefore
\begin{equation}
 \rho_{OO}(S)=\frac{S^2}{128\pi^2}\,\mathcal I,
 \qquad \mathcal I=\int\frac{|K(s,r)|^2}{sr}\,ds\,dr>0 .
 \label{eq:next3-subthreshold-uv}
\end{equation}
This calculation uses the usual two-line discontinuity on the selected
plane~\cite{Cutkosky1960}. Polynomial contacts have zero discontinuity.
If the two propagating particles were drawn only from the declared
massive modes, their squared threshold would be $4t_0$. At zero
mixing the scalar probe acts only on these massive modes and has zero
matrix element into physical null photons. Thus the strictly positive
weight in \eqref{eq:next3-subthreshold-uv} cannot be assigned to the
stated intermediate states. When $t$ is a squared photon curvature,
the same channel is present in the leading four-photon response for
polarizations with nonzero curvature contraction.

Its size is fixed by the original optical row. With $M_j=\int
s^j|u(s)|^2ds$, the normalized rotation of $u,Tu$ gives
\[
 K(s,r)=\frac{(s-r)^2u(s)\overline{u(r)}}
                     {\sqrt{M_0M_2-M_1^2}} .
\]
Hence $\mathcal I=(2M_3M_{-1}-8M_2M_0+6M_1^2)/(M_0M_2-M_1^2)$.
All moments used here exist for the A4 readout. Direct integration gives
$\mathcal I\simeq5.10733$ for $Lm=1$, $n=248$. Positivity follows
from the integral of a square and the nontrivial mass continuum.

\paragraph{Restoring the physical massive vector and computing its cost}
Use the local Proca quadratic form on each massive mode and keep the
same mass rotation and its full second derivative. The inverse is now
$C_s^{\rm P}=(I+qq^{\mathsf T}/s)/(x+s)$, which has only the
massive pole. The connected cut then starts at the physical threshold.
There is a separate Euclidean cost. In the relative half log
determinant the coefficient of $|\widehat t(p)|^2$ is $1/4$ times
the integral of $|K(s,r)|^2$ against the difference of the two
tensor bubbles at transfers zero and $p$. The zero-transfer term is
the actual $D$ contact: at the same momentum
$C_s^{\rm P}-C_r^{\rm P}=(r-s)C_s^{\rm P}C_r^{\rm P}$, and the
diagonal matrix elements of $D$ give precisely this difference of
resolvents. There is no adjustable subtraction. For
$y=|q+p|^2$, $g_s(x)=(x+s)^{-1}$, the tensor contraction is exactly
\begin{equation}
\begin{aligned}
 &\operatorname{tr}\{C_s^{\rm P}(q)
                [C_r^{\rm P}(q)-C_r^{\rm P}(q+p)]\} \\
 &=3g_s(x)[g_r(x)-g_r(y)]
   +\frac{x |p|^2-(q\cdot p)^2}{sr(x+s)(y+r)} .
\end{aligned}
\label{eq:next3-proca-cost-uv}
\end{equation}
One can verify it by resolving each vector into the longitudinal and
three transverse directions at its own momentum. The first term
falls as $|q|^{-5}$ at fixed $p$. The second is nonnegative. On a
hard angular sector of nonzero measure it has order $|p|^2/|q|^2$
for any $p\ne0$ and any bounded mass interval carrying nonzero
$|K|^2$. Its leading angular average, including the measure
$d^4q/(2\pi)^4$ and the factor $1/4$, is
$3\mathcal I |p|^2\,dx/(256\pi^2)$. Thus even the complete
contact and bubble leave a quadratic ultraviolet divergence in the
local massive-vector action. The spacetime-constant rotation is an
exception because its contact cancels the full relative trace.

\paragraph{Transporting the kinetic term as well as the mass}
There is a precise test of whether a connection fixed by the same
rotation supplies the missing contact. In the Proca curl replace
$\partial_\mu$ by $\partial_\mu+cJ\partial_\mu t$ for a real comparison
constant $c$. For a prescribed smooth scalar $t$ the unitary field
change $b'=e^{ctJ}b$ brings that curl exactly to the free curl and
leaves the mass $e^{-(1-c)tJ}Te^{(1-c)tJ}$. The rotation acts on two
normal directions and has unit determinant at each point. Equivalently
the two quadratic operators are exactly conjugate, including their
measures on finite tests with transported kinetic forms. The hard
coefficient in \eqref{eq:next3-proca-cost-uv} is therefore multiplied
by $(1-c)^2$. At $c=1$ the prescribed scalar disappears from the
whole massive quadratic form expressed in $b'$, and its normalized
determinant is independent of $t$. In this class a flat internal
connection can remove the divergence only by removing that bulk
response. This reasoning also applies to the null-curvature version
of the model conditional on the null field. If the observed curvature
itself contains $b$, a field-dependent coordinate change requires its
own Jacobian; the prescribed-profile result alone does not evaluate it.

\subsection{The first interacting reflection test on the protected photon}
\label{sec:next3-null-uv}
The fixed-null construction in \eqref{eq:joint2-fixed-null-uv}
protects the location of the photon pole. Its first nonlinear
correction decides a different question about the physical states.
Work on its reconstructed Euclidean plane with the same four-dimensional
A4 convolution $\mathcal T$, the true massless field $a_0$, and the
transverse Gaussian massive field $b$. Write $G_s(q)=(|q|^2+s)^{-1}$,
$P_q=I-qq^{\mathsf T}/|q|^2$, and $V=[T,J]$, where $J$ generates
the specified rotation. Its finite-rank mass kernel satisfies
$V(s,s')=(s-s')J(s,s')$ and is nonzero.

First prescribe a scalar rotation profile $t(x)$. The linear
coefficient of the relative massive determinant vanishes since the
diagonal mass kernel of $V$ vanishes. Write the quadratic coefficient
of the half log determinant as $\int |\widehat t(p)|^2 K(p)\,d^4p/(2\pi)^4$.
The double variation of the same mass square gives
\begin{equation}
\begin{aligned}
 4K(p)&=\iint |V(s,s')|^2 ds\,ds'\,
      \int\frac{d^4q}{(2\pi)^4}G_s(q)\\[-2pt]
 &\quad\times\big[3G_{s'}(q)-\operatorname{tr}(P_qP_{q+p})
                                      G_{s'}(q+p)\big].
\end{aligned}
\label{eq:next3-null-k-uv}
\end{equation}
With $C=[[T,J],J]$, symmetrizing its one-insertion trace gives
$\operatorname{tr}[(|q|^2+T)^{-1}C]=
\iint |V(s,s')|^2G_s(q)G_{s'}(q)\,ds\,ds'$.
This is the first term of \eqref{eq:next3-null-k-uv}. The second
is the ordinary two-insertion term. Their difference is absolutely
integrable for each $p$ by the gapped trace estimates. To determine
its sign, split off the part with $3-\operatorname{tr}(P_qP_{q+p})
=1-(\hat q\cdot\widehat{q+p})^2\ge0$. The remaining difference of
scalar bubbles, per unit spin factor, is
\[
 \frac{1}{16\pi^2}\int_0^1\log\!\left(
 1+\frac{x(1-x)|p|^2}{xs+(1-x)s'}\right)dx .
\]
This follows by integrating the Feynman-parameter difference, with
its common leading ultraviolet term canceled at the integrand level.
Consequently $K(0)=0$, and $K(p)>0$ at every $p\ne0$. The preceding
trace bound gives $K(p)\le C[1+\log(1+L|p|)]$.

Now use the actual nonlinear interaction
$t=e^2L^4I$, with $I=\sum_{\mu<\nu}(\mathcal T F_{a_0,\mu\nu})^2$.
The first photon interaction is the positive quartic form just computed,
at order $e^4$. For an external photon of momentum $Q$ and transverse
unit polarization $\epsilon$, denote its order-$e^4$ proper Hessian,
with the factor $e^4$ stripped, by $H_4(Q,\epsilon)$. There are two
types of photon contraction. Pairing the two fluctuations in a single
$I$ fixes its total momentum to zero and gives $K(0)=0$. Pairing
one fluctuation in each $I$ gives, up to a positive overall constant,
\begin{equation}
\begin{aligned}
 H_4(Q,\epsilon)&\propto L^8 |c_L(Q)|^2 \int\frac{d^4 k}{(2\pi)^4}
     K(Q+k)|c_L(k)|^2\\
 &\qquad\times\frac{|(k\cdot Q)\epsilon-(k\cdot\epsilon)Q|^2}{|k|^2} .
\end{aligned}
\label{eq:next3-null-h-uv}
\end{equation}
The numerator is already transverse to $k$. All cross contractions
have the same sign because they give the variance of the linear term
in $I$ in a background field. Thus $H_4$ is strictly positive for
$Q\ne0$ away from the zeros of $c_L(Q)$. It is bounded by
$C |Q|^2 |c_L(Q)|^2[1+\log(1+L|Q|)]$. All integrals here use the
full geometric form factor over the whole momentum space.

Choose a nonzero spatial momentum $\mathbf k$ and a spatial unit
polarization perpendicular to it. For $Q=(\omega,\mathbf k)$,
$H_4$ is even and integrable in $\omega$. Write $A_j(\mathbf k)=\int \omega^j H_4(\omega,\mathbf k)\,d\omega/(2\pi)$
for $j=0,2$. Both numbers are finite and positive for the A4 profile.
The time transform obeys $\widetilde H_4(t,\mathbf k)\ge
A_0-t^2 A_2/2$ by the elementary cosine bound. Take a smooth
spatially transverse positive-time probe $f$ and put $h=-\Delta f$.
It is also a gauge-invariant curvature test, since integration by
parts gives $a_0(h)=\tfrac12\int F_{a_0,\mu\nu}(\partial_\mu f_\nu-\partial_\nu f_\mu)$
for divergence-free $f$. The free reflection norm of $a_0(h)$ vanishes because applying these
two Laplacians to its free Maxwell covariance leaves only a contact.
For the first interacting coefficient the propagator varies as
$-D_0H_4D_0$. The same two Laplacians therefore leave $-H_4$ on
the chosen polarization. Take a nonzero, nonnegative time pulse
supported in $0<t<\tfrac12\sqrt{A_0/A_2}$. Every reflected
separation then gives $\widetilde H_4\ge A_0/2$. The resulting coefficient
of $e^4$ in the reflection norm is strictly negative. Transverse
spatial wave packets near $\mathbf k$ give ordinary smooth tests
by continuity. Any local covariance contact vanishes on these
separated reflected supports. This rules out a reflection-positive
completion with the stated small-coupling expansion of this
observable: it would have a negative leading norm on a free null state.

The same conclusion holds for the full observable algebra when the
rotation is driven by $F_{A_{\rm obs}}$. In the free splitting at each
$e$, $A_{\rm obs}=\sqrt Z\,a_0+\langle W_e^\dagger e_A,b\rangle$,
where $Z=1+O(e^2)$ and $W_e^\dagger e_A=O(e)$. At any single
first-variation vertex, two external null legs exhaust its two curl
slots and the remaining massive trace vanishes. The double variation
and the two-vertex contractions first contribute at order $e^4$,
where the curl slots reduce to those of $a_0$. Thus the witness above
is still a physical-null-mode test in that algebra. A theory whose
accessible algebra contains only $A_{\rm obs}$ requires a separate
test on that restriction. The fixed pole and the positive quadratic
pair spectrum both survive this calculation. It is the interacting
four-dimensional readout that fails the separated-time norm test.

The sign test uses no special zero of the A4 transform. It applies
to any bounded real radial multiplier $c(|q|)$ that is nonzero
on an open interval and satisfies
\[
 \int_0^\infty r^3 |c(r)|^2[1+\log(1+Lr)]\,dr<\infty .
\]
This is the same type of absolute photon-line estimate used to
close the geometric graphs. Since $\log(1+L|Q+k|)$ is bounded
by the sum of the two individual logarithms, the estimate on
$H_4(Q,\epsilon)$ still holds with $c$ in place of $c_L$.
At fixed nonzero spatial momentum the change of variables
$r=(\omega^2+|\mathbf k|^2)^{1/2}$ has bounded Jacobian
$d\omega/dr$ for $r>2|\mathbf k|$. The displayed $r^3$ moment
therefore controls the $r^2$ weight in the frequency integral of
$H_4$. Near the lower endpoint the Jacobian has only an integrable
inverse-square-root singularity. Choosing $|\mathbf k|$ below
some interval where $c$ is nonzero gives the same strict
reflection test. Hence varying the radial decay profile cannot
repair this interaction while it damps both ends of the physical
null photon strongly enough for this photon-line bound. Geometric
suppression carried on genuinely internal relative variables is
outside the scope of this obstruction.

\subsection{Positive interacting time evolution on a locked spatial slice}
\label{sec:next3-os-uv}
The distinction between a spatial displacement and a clock displacement
allows an exact test of reflection positivity. Use the normal square
of \eqref{eq:joint2-inverse-spectrum-uv} at fixed $e^2M_{-1}<1$.
Write $a$ for its massless spatially transverse photon and $b_i$ for
three independent real massive spatial components with mass-squared
operator $T$. These components have ordinary scalar kinetic operators
on the continuum of masses. On divergence-free spatial tests the
observed field is
$A=\sqrt Z\,a+\langle eT^{-1/2}u,b\rangle$, with the exact free
spectrum in \eqref{eq:joint2-inverse-spectrum-uv}. The corresponding
transverse projection on the massive optical term may be inserted
explicitly. Its curl and all the stated optical tests give the same
answer. All component fields are even under Euclidean time reflection.

Let $S$ average over the spatial marginal of the centered A4 ball at
length $L$, so its real multiplier is $c(\mathbf k)=\chi_{248}(L|\mathbf k|)$.
Put $D=I-S$, $v=u/\|u\|$, and define on each clock slice
\begin{equation}
 \mathcal B=SD^2\nabla\!\times A,
 \qquad \mathcal Y=SD\langle v,b\rangle .
 \label{eq:next3-os-marks-uv}
\end{equation}
Each $D$ measures a field relative to its shell average. Using two
such differences on the magnetic field and one on the massive packet
will also control the long-distance photon. Consider the interaction
density, with $mL=1$,
\[
 e^2m^4[1-\chi_{248}(eL^2|\mathcal B|)]
       [1-\chi_{248}(eL|\mathcal Y|)] .
\]
Each bracket is an average of $1-\cos$ of a real displacement phase.
The density lies between zero and $4e^2m^4$ and its first derivatives
in the two marked vectors are bounded, by the characteristic-function
formula for the ball. The interaction starts with a nonzero mixed
quartic vertex. This is a choice of boundary interaction and a
preferred spatial slicing.

\paragraph{The exact reflection test and the primitive graphs}
On a spatial torus of side $R$ omit the spatially constant free photon
mode as an infrared convention. All spatial momenta above it and all
normal masses are present. The marked fields at a spatial point exist
as equal-time Gaussian variables. Their variances are bounded by
sums of $|c|^2$ times at most two spatial momenta divided by the
positive free frequency. In the free field configuration representation
the spatial integral of the density is thus a bounded self-adjoint
multiplication operator. The free Fock Hamiltonian plus this potential
is self-adjoint and bounded below. At fixed $R$, if
$4e^2m^4R^3<\min(2\pi/R,\sqrt{t_0})/2$, the isolated simple free
vacuum persists as an isolated ground state by bounded perturbation.
Its Euclidean time correlations obey reflection positivity since a
reflected ordered expectation is the squared norm of a vector obtained
with the corresponding positive-time semigroup. For arbitrary allowed
$e$ there is a direct periodic Euclidean construction. Weight the free
Gaussian on a time circle by the exponential of minus the spacetime
integral of the density. The free reflected kernel is positive by its
positive-frequency decomposition, also for the spatial transverse
projector. The weight factors as $W_+\Theta W_+$ with $W_+$ depending
only on the positive half circle. All bounded positive-time observables
then obey the interacting reflection inequality.

Every primitive bare graph at fixed order is ultraviolet finite in
this model. The Taylor coefficients contain finitely many ordinary
marked fields per order. A null line has a magnetic curl at both ends.
With $k=|\mathbf k|$ it has absolute envelope
$Ck^2|c|^2/(\omega^2+k^2)$. Massive lines have envelopes
$C(1+k)^2|c|^2/(\omega^2+k^2+t_0)$ by the normal gap and the bounded
normal packet norms. The factors $1-c$ are bounded. Both envelopes
are bounded and integrable over all four momenta, including the
origin, since integrating frequency costs only $\pi/k$ or
$\pi/\sqrt{k^2+t_0}$ before the numerator. Bound a spanning tree by
supremum norms and integrate each independent chord in $L^1$. This
works for self-contractions as well. No coincident-field subtraction
or cancellation between primitive rings is needed. The perturbative
claim concerns the Gaussian expansion of connected coefficients;
convergence of the full perturbation series is a separate question.

\paragraph{Infinite-volume states and the massless scaling limit}
For an optical or normal test $f$ let $C_0$ denote the free covariance
on all components and $K$ the pair of marks in
\eqref{eq:next3-os-marks-uv}. Translation of the Gaussian by its
Cameron--Martin vector gives the uniform estimate
\begin{equation}
 \left|\log\mathbb E_{\rm int} e^{h\phi(f)}
          -\tfrac12 h^2 C_0(f,f)\right|
       \le C_e |h|\,\|K C_0 f\|_{L^1},\qquad h\in\mathbb R .
 \label{eq:next3-os-tilt-uv}
\end{equation}
Indeed the difference of the two shifted interactions is bounded
pointwise by the integral of their bounded gradients times the
shift of the marked fields. The bound is uniform on expanding
periodic four-boxes. For the massless covariance the spatial curl
removes its transverse projector. The kernel $SD^2\nabla G_0$,
$G_0(x)=1/(4\pi^2|x|^2)$, is locally integrable and has integrable
tail: the centered compact shell differences cancel spatial Taylor
terms through degree three, giving $O(|x|^{-7})$. Periodization
bounds its $L^1$ norm on each box by the continuum norm. The
omitted photon modes have spatial momentum zero and vanish in the
curl. Massive Green kernels have $L^1$ norm $1/s\le1/t_0$; a
spatial derivative of such a kernel has finite $L^1$ norm uniformly
for $s\ge t_0$ by the free heat representation. The optical probe is
tested with spatially transverse vectors, and the massive fields
enter the unprojected normal marks. These choices prevent a
spurious spatial projector tail in this bound.

For smooth periodized tests their free variances converge. The
bounds above also give uniform continuity of the characteristic
functionals at zero in the Schwartz topology: by applying the
exponential moment bound to both signs and rescaling $f$, the second
moment is bounded by a constant times
$C_0(f,f)+\|KC_0 f\|_1^2$. One can first use compactly supported
smooth tests on each fixed spacetime ball and a countable dense
family of normal packets. Tightness of the test marginals, diagonal
extraction and this continuity give a translation-invariant
subsequential state on the test distributions. Periodic reflection
passes to the limit on bounded cylinder observables. The exponential
bounds extend it to their polynomial moments. This proves joint
interacting reflection positivity at every allowed real coupling
in this sense; uniqueness of the infinite-volume state is open here.

There is also a local identity that distinguishes the limiting state
from the free one. Let $h$ be a compactly supported smooth shift of
the fields, spatially transverse on the null component and in the
domain of $T$ on its finitely many normal profiles. For a bounded continuous
cylinder observable $F$ Gaussian translation gives
\[
\begin{aligned}
 \mathbb E F(\phi+h)=\mathbb E\big[ F(\phi)
 \exp\{&\phi(C_0^{-1}h)-\tfrac12\langle h,C_0^{-1}h\rangle\\
       &+V(\phi)-V(\phi-h)\}\big] .
\end{aligned}
\]
Here $V$ denotes the integrated interaction, and only its difference
occurs. That difference has compact support in spacetime because
$Kh$ does. The identity passes to each thermodynamic limit. To see
this, smooth the marks in a small spacetime neighborhood before
using finite test marginals. Their pointwise free variances are
finite, and their free mean-square approximation errors vanish in
the large-box limit followed by removal of this smoothing.
The corresponding $L^1$ term in \eqref{eq:next3-os-tilt-uv} is
the mollification error of $K C_0 K^\dagger$, also tending to zero.
Indeed the compact spatial kernels of $K^\dagger$, including the
differentiated A4 marginal, have finite total variation. The
preceding $L^1$ estimates apply to their convolutions and to their
periodizations. Thus the interacting marks converge in mean square
on each compact region under this approximation. Bounded gradients
of the density control the error in the local interaction difference,
and the uniform exponential moments control the remaining Gaussian
translation factor. These approximations justify the limiting identity.

Under the free measure the translation identity has no interaction
difference. Equality of the two states would therefore force that
difference to vanish almost surely for every such $h$, by testing
all bounded continuous cylinder observables. Choose $h$ with both marked
vectors nonzero on a common open region. At the zero configuration
its difference equals $-V(-h)<0$ for $e\ne0$, and remains negative
on a small $L^2$ neighborhood of the zero marks over the compact
region. A centered Gaussian with finite local second moment assigns
positive probability to every such ball, as follows by finite-rank
approximation of its covariance. Thus the limiting joint state is
actually changed by the interaction. This conclusion neither
requires nor establishes uniqueness of the thermodynamic state.

For a smooth compactly supported spatially transverse optical test
use $f_\lambda(x)=\lambda^{-3}f(x/\lambda)$ on the full plane.
By the second-order Taylor formula each $D$ costs two spatial
derivatives in $L^1$. The massless contribution to $\|KC_0
f_\lambda\|_1$ is $O(\lambda^{-2})$ by scaling five derivatives
of $G_0*f$, which are integrable. For the gapped components move
the derivatives onto $f_\lambda$ and use the massive $L^1$ bound.
Even the single $D$ in $\mathcal Y$ gives $O(\lambda^{-1})$.
The free massive optical variance is $O(\lambda^{-2})$ by the gap.
Equation~\eqref{eq:next3-os-tilt-uv} consequently squeezes every
joint optical generating function at large distances to that of
the free transverse massless field with residue $Z$. This is an
exact Gaussian scaling limit in each of these states. It gives no
boost-invariant mass-shell assertion at finite momentum. The
instantaneous magnetic interaction has changed the covariant
four-dimensional readout, and its normal fields describe a neutral
pair channel. An unrestricted charged GCM dynamics must still
supply its fermionic states and a physical transformation law for
the chosen slice.

\subsection{What a charged pair can measure in its own rest plane}
\label{sec:next3-rest-uv}
There is a covariant way to reproduce the observed A4 pair amplitude
while sampling only spatial displacements at the time of the pair.
It exposes a difficulty in using the same average off shell. Work on
one locked sector plane with Minkowski signature $(+---)$. For two
on-shell particles of equal mass set $Q=p_++p_-$, $P=(p_+-p_-)/2$,
$S=Q^2\ge4m^2$, and $v=Q/\sqrt S$. Project the normalized A4 ball
of dimension $n=248$ onto three axes in $v^\perp$. Dilate the
projected displacement by $L\sqrt S/m$ and pair it with $P$. This
prescription for the boundary state is invariant under changes of
spatial axes in the pair rest frame. A boost transports both the
state and the hyperplane. Its amplitude is
\begin{equation}
 h(G)=\chi_n\!\left(\frac Lm\sqrt G\right),\qquad
 G=(P\cdot Q)^2-P^2Q^2 .
 \label{eq:next3-rest-gram-uv}
\end{equation}
Indeed $P\cdot Q=0$ and $G=S(S/4-m^2)$. The factor $\sqrt S/m$
in the length is essential to the precise agreement with the
charged-leg rest-space measurement in \eqref{eq:charged-rest-readout-uv}.
This is spatial averaging on a physical two-body intermediate state.
It gives the full open mark, including its sign between Bessel zeros.
The unit $v$ is a velocity on the observed sector plane; the
vectorial pole displacement remains the one on the original carrier.

\paragraph{A spacelike transfer supplies no rest-space probability}
Suppose instead one tries to select a real future timelike unit
normal by a normalized positive law conditional on $Q$ alone,
equivariant under proper orthochronous Lorentz transformations.
In a frame with $Q=(0,0,0,q)$ the stabilizer contains the boost
of rapidity $\log2$ along the first spatial axis. For a future
unit normal $v$ the positive variable $v^0+v^1$ can be doubled by
this boost. The disjoint dyadic intervals of this variable would
all have the same probability, which precludes normalization.
A stronger statement applies directly to a real displacement $Y$.
Invariance of a probability gives $Y^0+Y^1=Y^0-Y^1=0$ almost
surely, by the same argument on each sign and with inverse boosts.
The boost along the second spatial axis then forces that component
to vanish as well. Thus $Y$ belongs to the line spanned by $Q$ almost
surely. If the displacement is required to be orthogonal to $Q$,
it vanishes. Conditional data from another measured worldline can
change this conclusion. Averaging over an independent normalized
Lorentz-invariant timelike normal in the vacuum cannot supply it.

\paragraph{The unique analytic extension using only the bivector norm}
At the two physical poles $G$ ranges over the whole nonnegative
ray. If a holomorphic off-shell mark depends solely on $G$ and is
regular on a connected continuation domain containing both rays,
its values there determine it uniquely by the identity theorem.
The entire even Bessel series gives precisely
\eqref{eq:next3-rest-gram-uv}. Under a joint Wick rotation
$G=(P_E\cdot Q_E)^2-|P_E|^2|Q_E|^2$. For $Q_E=q e_4$, $q>0$,
write $r=|\mathbf P_E|$ and $y=Lqr/m$. The continued mark is
$\chi_n(iy)=\mathbb E\exp(y U_1)$, where $U_1$ is one coordinate
of the normalized ball. For any $0<\delta<1$ it is bounded below by
$e^{\delta y}\Pr(U_1>\delta)$. Moreover
$|\partial_y\chi_n(iy)|\le\chi_n(iy)$, where the derivative means
that of the real function $y\mapsto\chi_n(iy)$. The enhancement at
real Euclidean loop momentum follows directly from positivity of
the displacement measure.

One can test the complete determinant contact for this continuation
first on bounded momentum regions, then by exhaustion.
Use the local massive Dirac inverse and the straight Wilson chord on
the midpoint kernel as in \eqref{eq:rest-ward-vertex-uv}. In Euclidean
notation set $b=(P\cdot Q)P-P^2Q$, $d(G)=[h(G)-1]/G$, with its
removable value at zero. The transverse correction to the ordinary
one-photon vertex is $d(G)[G\gamma^\mu-b^\mu\slashed Q]$. It obeys
the Ward identity since $Q\cdot b=G$. Complete the term linear in
curvature by its endpoint Wilson chord, with no separate quadratic
curvature coupling. In polarization $1$, on the integration region
$|P_4|,|P_1|<m$ and $r\to\infty$, the full vertex is
$h(G)\gamma^1+O_q(h(G)/r^2)$. The spin trace of the two free
propagators with unit $\gamma^1$ vertices has magnitude
$(4+o(1))/r^2$, uniformly on this region. Consequently the bubble
has magnitude at least $c h(G)^2/r^2$ for large $r$. The chord
contact is an integral of a first midpoint derivative of the linear
kernel at $P+\theta Q$, $|\theta|\le1/2$, times one free propagator.
Those shifts leave $r$ fixed. The derivative estimate above and the
polynomial factors in $b$ bound the contact by a polynomial in $r$
times $h(G)$. The bubble therefore dominates the absolute value of
the complete integrand at large $r$. The remaining two transverse
components range over annuli of unbounded area. The two-current
coefficient already fails absolute ultraviolet convergence at fixed
nonzero spacelike transfer.

This failure concerns the holomorphic one-invariant extension and
the specified chord action. A second invariant measuring how far the
charged lines lie from their poles can distinguish the spacelike
integration region while leaving the measured cut fixed. Its coupling,
its higher variations and its measure still have to be derived from
the physical source. A different choice of continuation of the square
root in \eqref{eq:next3-rest-gram-uv} cannot change its even series.

\paragraph{A second invariant repairs the Euclidean sign}
The zeros of the inverse charged propagators provide a direct test
of a modification. In Lorentzian variables define
\[
 D_\pm=m^2-(P\pm Q/2)^2,\qquad Z=G(1-D_+D_-/m^4).
\]
Use $h=\chi_n((L/m)\sqrt Z)$ in the bare charged vertex. The two
physical poles give $Z=G$. On the real Euclidean contour put
$Q=q e_4$, $u=P_4$, $r=|\mathbf P|$ and
$D_\pm=m^2+r^2+(u\pm q/2)^2$. Then
$Z=q^2 r^2(D_+D_-/m^4-1)\ge0$.
The same analytic mark now reproduces the timelike amplitude and
is bounded at all real Euclidean momenta. Two possible curvature
couplings with precisely this scalar mark have different ultraviolet
behavior.

First, replacing $G$ and $b$ in the polynomial vertex above by
$Z$ and $(1-D_+D_-/m^4)b$ preserves their ratio. For polarization
$1$ the vertex is
$h\gamma^1+(h-1)uP_1\gamma^4/r^2$.
At fixed $q>0$, take $|P_1|<m$, $r<u<2r$ and $r\to\infty$.
The mark and its first midpoint derivatives decay faster than any
inverse power needed in this estimate at $n=248$. The two vertices
are $O(r^{-1})$, and their bubble is $O(r^{-4})$. In the Wilson
contact the derivative of the limiting $\gamma^4$ coefficient is
$-u(1-2P_1^2/r^2)/r^2$, uniformly up to the bounded chord shifts
in $u$. The trace of $\gamma^4$ with the single propagator is
proportional to $u/(m^2+r^2+u^2)$, with a nonzero constant.
Hence the absolute contact is bounded below by $c/r^2$ on this
region. The integral of the combined graph diverges. Even the
bounded Euclidean scalar mark does not suppress a transverse
component left by this particular Ward completion.

For a second explicit action put the transverse correction in the
projected connection, using $\Gamma^\mu=h\gamma^\mu+
(1-h)Q^\mu\slashed Q/Q^2$ at $Q\ne0$, with a straight endpoint
Wilson chord on its midpoint kernel. The zero mode takes the
original charge normalization. This is a separate nonlocal
field-strength coupling. Its transverse vertex is exactly
$h\gamma^1$. At fixed $q>0$ and $|P|$ large its ball argument
is comparable to $q r|P|^2 L/m^3$. Its first transverse midpoint
derivative costs at most $C_q(1+|P|^2)$ before the decay of the
ball derivative. In four-dimensional polar coordinates write
$P=R\omega$ and $a=\sqrt{1-\omega_4^2}$. The absolute contact,
including its propagator, is bounded in the exterior by
$C_q R(1+c_qR^3a)^{-(n+1)/2}$; the bubble has a smaller bound.
The radial integral of the contact there is bounded by $C_q a^{-5/3}$,
integrable on the three-sphere, also at the two collinear directions.
The complete bare Hessian therefore exists at each real nonzero
transfer.

Its soft limit gives a sharper test of the repair. Translate the
chord after integration by parts in $P_1$, as for
\eqref{eq:virtmark-smallq-uv}. With the spin kernels $B_U,J_U$ of
\eqref{eq:nextsoft-rooted-uv}, at $U=m^2$ the exact integrand
of $\Pi_{11}/e^2$ against $d^4P/(2\pi)^4$ is
$h^2 B_U(P,q)+8h\int_{-1/2}^{1/2}J_U(P-\theta Q)d\theta$.
Put $\rho=(m^3/(Lq))^{1/3}$ and rescale $P=\rho w$.
The limiting mark is $j=\chi_n(|w|^3a)$, and the rescaled
integrand after extracting $\rho^2$ tends to
$4(1-2\omega_1^2)(2j-j^2)/|w|^2$. For $q\le m$ one has
$D_+D_--m^4\ge |P|^4$. Both shifted spin kernels are bounded
by $C/(m^2+|P|^2)$, uniformly over the chord. Thus
$C|w|^{-2}(1+|w|^3a)^{-(n+1)/2}$ is an integrable majorant.
Dominated convergence and radial substitution yield
\begin{equation}
 \lim_{q\downarrow0}\frac{\Pi_{11}(q e_4)}{e^2\rho^2}=C_n>0,
 \label{eq:next3-gram-soft-uv}
\end{equation}
where, writing $d\omega$ for surface area on the unit three-sphere,
\[
\begin{aligned}
 C_n&=\frac{4}{3(2\pi)^4}\int_{S^3}\!
   (1-2\omega_1^2)a^{-2/3}d\omega\ I_n,\\
 I_n&=\int_0^\infty t^{-1/3}
     [2\chi_n(t)-\chi_n(t)^2]dt .
\end{aligned}
\]
At fixed $\omega_4$ the angular mean of the first factor in
parentheses is $1-2a^2/3>0$. For independent one-coordinate
ball samples $Y,Y'$ the damped cosine Mellin integral with exponent
$2/3$ gives $I_n=\Gamma(2/3)\cos(\pi/3)
[2\mathbb E|Y|^{-2/3}-\mathbb E|Y+Y'|^{-2/3}]$.
The even decreasing density of $Y$ puts the most mass in a centered
interval, so conditioning on $Y'$ bounds the second expectation
by the first. Both moments are finite. At $n=248$ direct integration
gives $C_n\simeq0.13362$. Thus the complete transverse coefficient
grows as $q^{-2/3}$ as the external momentum approaches zero. The
virtuality product repairs the Euclidean sign, but in the projected
source the momenta $|P|\sim\rho$ obstruct a regular Maxwell limit.
This result concerns the stated minimal chord action. Additional
curvature contacts would have to supply their own compensating
contribution and pass the physical spectral test.

Raising the virtuality product to an integer power cannot cure this
limit. Take $Z=G[1-(D_+D_-/m^4)^j]$, $j\ge1$, in the same
projected action, and put $d=2j+1$ and
$\rho_j=(m^{2j+1}/(Lq))^{1/d}$. The rescaled mark tends to
$\chi_n(|w|^d a)$. For $q\le m$ the bound
$(D_+D_-/m^4)^j-1\ge[(D_+D_-/m^4)-1]^j$ gives the same
integrable majorant with power $d$ in its argument. The limit of
$\Pi_{11}/(e^2\rho_j^2)$ has prefactor $4/[d(2\pi)^4]$,
angular weight $a^{-2/d}$ and Mellin exponent $2/d$. Since
$0<2/d<1$, the same centered-interval argument makes its coefficient
strictly positive. Every such finite power leaves a singularity
of order $q^{-2/(2j+1)}$ on the two-current test.

\paragraph{A soft obstruction for every nonnegative Euclidean mark in this source}
The minimal projected coupling admits a more general test that does
not need a power law. Let $h(P,q)$ be rotationally invariant in the
three directions orthogonal to $Q=q e_4$, with $0\le h\le1$ on real
Euclidean momenta. Suppose $h(P,q)\to1$ at each fixed $P$ as
$q\downarrow0$, as required by a regular pointwise limit of this
vertex with its free Dirac normalization. Retain the preceding
straight-chord action and its convergent translated contact identity.
Then a finite regular transverse response at zero momentum is
impossible, however fast its geometric width grows with virtuality.

Indeed let $D=m^2+r^2+u^2$ and put
$J_*=(D-2r^2/3)/D^2\ge1/(3D)$. Angular averaging over the spatial
two-sphere replaces $P_1^2$ by $r^2/3$. At small $q/m$ the resulting
bubble obeys $\overline B_U=-4J_*[1+O(q/m)]$, and each translated
chord kernel satisfies $\overline J_U(P-\theta Q)=J_*[1+O(q/m)]$
uniformly in $r,u$ and $|\theta|\le1/2$. For the first estimate
use the exact numerator $-4(D-q^2/4-2r^2/3)$ and the denominator
$D_+D_-$; for the second, shifting $u$ by at most $q/2$ changes
both $D$ and $D-2r^2/3$ by a relative $O(q/m)$ because both
dominate $m^2+u^2$. Hence the complete angular integrand is bounded
below by $cJ_*h$ for all sufficiently small $q$, using $h^2\le h$.
Fatou's lemma and $\int d^4P/D=\infty$ then give
$\liminf_{q\downarrow0}\Pi_{11}(q e_4)=+\infty$ for nonzero charge.
No assumption on the momentum scale that dominates the limit is
needed. A sign-changing mark can evade this particular inequality;
the explicit A4 virtuality-product family still has the divergent
Mellin coefficient computed above. Independent curvature contacts
or a different charged determinant change the identity being bounded.

\subsection{Retaining the shell velocity in the charged trace}
\label{sec:next3-velocity-uv}
There is a charged transport in which a contractible loop first acquires
an electromagnetic response when its actual path encloses an area.
The number of changes of direction then controls the short-duration
end of the determinant. Start with the normalized ball in
\eqref{eq:monomial-owner-uv} on a fixed aligned plane. In velocity units
let $y$ be its rank-four orthogonal projection, with law $d\nu(y)$.
For the uniform $n$-ball, $n=248$, this projected law has bounded
density on the unit four-ball, positive in its interior. On the retained
velocity space $L^2(\nu)$ put $P=|1\rangle\langle1|$. A stationary
Markov rule that completely refreshes the direction at rate $\lambda>0$
has generator $\lambda(P-I)$. On the Euclidean readout choose the
charged transport operator
\begin{equation}
 D_B=m+y\cdot(\partial+i e B)+\lambda(I-P),\qquad m>0.
 \label{eq:next3-flight-uv}
\end{equation}
The mass and refresh rate have inverse length units. This velocity
transition law is an additional specification of the charged curvature
representation; the shell delta determines its marginal. We specify its relative Grassmann determinant below by a cyclic
proper-time trace. The same operator gives, with the normalized
constant velocity state at each endpoint, the open inverse
$P D_B^{-1}P$. The field here has a scalar transport index. An electric
Dirac spin coupling and a physical charged-state reconstruction would
require further input.

\paragraph{Which closed paths actually contribute?}
Use the cyclic proper-time prescription
$-\int_0^\infty dt\,\operatorname{Tr}_{\rm cyc}
(e^{-tD_B}-e^{-tD_0})/t$. At fixed $t>0$ the trace first closes the
base path under its return density, then contracts the velocity
indices. This order is part of the path trace on this comparison
chart. With a continuous velocity a nonzero decomposable streaming
operator on the velocity fibers need not be compact even in a finite
base region. No ordinary unprojected relative Fredholm determinant
is assumed. Expanding $e^{-tD_B}$ at $k$ refreshes
gives the Poisson weight $e^{-\lambda t}(\lambda t)^k/k!$, the common
factor $e^{-mt}$, and the Wilson phase of a polygon. Under the cyclic
trace, for $k\ge1$ its $k$ velocities are independent with law $\nu$.
The initial and final intervals in the ordered integral join into one
arc. If $h_j$ are the resulting arc lengths in the transport parameter,
$\sum_j h_j=t$, and returning to the base point inserts
$\delta^4(\sum_j h_j y_j)$. Integrating the longest arc velocity first
bounds this return density by $\|d\nu/d^4y\|_\infty(k/t)^4$ at every
point of the time simplex. A closed straight path, or a path with only
two sides that exactly retrace each other, has unit Abelian holonomy.
Thus every field derivative is zero on the no-refresh streaming sector
and on the one- and two-refresh terms. For the streaming sector this
cancellation can be tested without assigning a value to a velocity
delta on its diagonal: on closing the base path its relative kernel
is proportional to $\delta^4(ty)(e^{-ie\int B\cdot dz}-1)$, which
is identically zero as a distribution in the velocity, since the
phase on $y=0$ is one. This vanishing holds against smooth
velocity test kernels before their trace is taken. Three refreshes admit nondegenerate triangular returns
on a set of positive return measure.

At a plane-wave photon insertion the Wilson integral on a closed path
is $J(q)=\oint e^{iq\cdot z}dz$. Closure gives $q\cdot J=0$ and,
by subtracting its phase at the base point, $|J(q)|\le
\min(t,|q|t^2)$. In inverse-length units fixed by $L$, this is bounded
by $t(1+t)\min(1,|q|)$. All $N$ derivatives of the Wilson exponential
therefore retain their separate soft factors, including coincident
variations. Near $t=0$ the Poisson average of
$k^4\mathbf1_{k\ge3}$ is $O(t^3)$. Including the return density,
$t^N$ and $dt/t$ gives $O(t^{N-2})dt$, integrable for $N\ge2$.
At large $t$ the Poisson fourth moment grows polynomially and
$e^{-mt}$ controls every fixed multiplicity. Odd vacuum vertices vanish
under reversal. These bounds apply under the absolute integrals before
performing the photon integrations. They concern source-dependent
coefficients per overall volume.

\paragraph{The photon may be averaged at a single clock reading}
Take $\mathcal T$ to be convolution with the three-spatial-coordinate
pushforward of the same ball, of multiplier $\chi_n(L|\mathbf q|)$.
Couple $B=\mathcal T A$ on the loop. For ordinary gauge transformation
of the open endpoints use $B=\mathcal T A+\partial\phi_A$, with
$\phi_A=\Delta_{\rm sp}^{-1}(I-\mathcal T)\nabla_{\rm sp}\cdot
\mathbf A$ on decaying tests. The pure gradient integrates to zero on
every closed path, and $B$ shifts by $\partial\alpha$ when $A$ does.
Each internal massless Maxwell line between two completed closed
vertices has the envelope
\[
 C|\chi_n(L|\mathbf q|)|^2
       \frac{\min(1,|q|)^2}{|q|^2},\qquad q=(q_4,\mathbf q).
\]
It is bounded on $\mathbb R^4$ and integrable there: use
$\int dq_4/(q_4^2+|\mathbf q|^2)=\pi/|\mathbf q|$ and the ball
transform decay at spatial infinity. A spanning tree bound now proves
absolute convergence of every connected graph of these closed charged
vertices and dynamical photon lines at each finite order, with no
frequency damping or momentum cutoff. A thin point coupling instead
retains the perimeter divergence of
Proposition~\ref{prop:closure-native-perimeter-uv} already on the
triangular returns. The spatial preparation selects a clock slicing.

\paragraph{A positive spectral test on the open transport}
At zero field the open scalar inverse is isotropic. For $X=|q|^2$,
$M=m+\lambda$ and $v=y_1$ set
$a(X)=\int M/(M^2+Xv^2)\,d\nu$. The rank-one inverse gives
$F(X)=a(X)/(1-\lambda a(X))$. Equivalently
\[
 F(X)=\langle1,(M-\lambda P+Xv^2/M)^{-1}1\rangle .
\]
Since $M-\lambda P\ge mI$, the spectral theorem gives a positive
Stieltjes measure for $F$ by conjugating $v^2/M$ by
$(M-\lambda P)^{-1/2}$. There is no constant term since $v^2$ has
trivial kernel. The unswitched continuum starts at $M^2$. On
$0<s<M^2$, $a(-s)$ increases strictly from $1/M$ to
$n/[M(n-1)]$, by the one-coordinate ball density. Hence the retained
field has a simple positive-residue pole strictly below that continuum
exactly when $\lambda>(n-1)m$. For example $\lambda=nm$ fixes this
ratio from the shell dimension as a representation choice. This is
positivity for the free scalar endpoint observable. The transport
calculation makes no identification of that pole with a Dirac particle
and leaves positivity of the interacting current to its own test.

The complete Gaussian already gives a stricter test on its even
algebra. Couple a real scalar probe $h(x)$ to the density of the
retained charged state by the bilinear $hP$. In the same cyclic relative
trace, Duhamel variation opens the two arcs at their rank-one source
insertions and gives the two projected resolvents. Their connected
kernel at distinct points is
\[
 C_h(x,y)=-F(x-y)F(y-x).
\]
Here $F(x)$ denotes the inverse Fourier transform of the projected
free inverse. The minus sign is the closed Grassmann contraction. The
positive massive Stieltjes representation makes $F(x)$ real, even
and strictly positive away from zero: each massive Green kernel is
an integral of positive heat kernels there. The spectral integral
converges on separated compact supports by its Stieltjes bound.
For any nonzero nonnegative smooth test supported strictly in positive
time, the connected reflection form of this centered scalar density
is therefore strictly negative. A local contact contributes zero at
these separated points. This excludes an interpretation of this
bilinear as a reflection-even Hermitian physical scalar. In a Dirac
representation the spin numerators and reflected dual change this
contraction. The missing spin and charge adjoint can thus be tested
before adding a second charged loop.

\paragraph{A Clifford transition repairs free charged-state positivity}
The geometry of the return proof permits a different spin amplitude at
an actual change of velocity. Take Hermitian Euclidean matrices with
$\{\gamma_i,\gamma_j\}=2\delta_{ij}$. Replace $P$ at each refresh by
$J\Xi$ on $L^2(\nu)\otimes\mathbb C^4$, where
\[
 (J\xi)(y)=\xi,\qquad
 \Xi f=\int(\gamma\cdot y) f(y)\,d\nu(y).
\]
Use the single charged operator $D_B^{\rm sp}=M+y\cdot(\partial+i e B)
-\lambda J\Xi$ with $M=m+\lambda$, and use $\Xi$ and $J$ as its two
open Berezin source maps. The spin transition $J\Xi$ is an amplitude;
it is not a Markov projection. The time weights in the expansion of
its streaming semigroup are still Poisson weights times $e^{-mt}$.
Each cyclic spin weight has magnitude at most four, since it is the
trace of a product of $\gamma\cdot y_j$ of norm at most one. Thus the
same complete closed-graph bounds hold, including those for dynamical
photons with the spatial A4 readout. In this case the triangular spin
trace vanishes too. Four-sided paths are the first that can have both
nontrivial electromagnetic holonomy and a nonzero spin trace.

Here the projected free Grassmann covariance has a positive spinor
spectral representation. Let $X=|q|^2$, $v=y_1$, and define
\[
 b(X)=\int\frac{d\nu(y)}{X+M^2/v^2},\qquad
 G(q)=-i\gamma\cdot q\,b(X).
\]
Indeed $G=\Xi(M+i y\cdot q)^{-1}J$ by reflection and rotational
invariance. Resolvent elimination gives $S=(I-\lambda G)^{-1}G$
for the inverse of the full action between its two source maps.
At $q=i z e_4$ its two $\gamma_4$ components are $f(z)$ and $f(-z)$,
with
\[
 a(z)=z b(-z^2),\qquad f(z)=a(z)/(1-\lambda a(z)).
\]
Let $d\kappa$ be the symmetric probability law giving equal weight
to $r=\pm M/|v|$. In $L^2(\kappa)$ let $T$ multiply by $r$. Partial
fractions give $a(z)=\langle1,(T-z)^{-1}1\rangle$. Consequently $f$
is the resolvent matrix element on the same unit vector of the
self-adjoint operator $T-\lambda|1\rangle\langle1|$. Denote its
positive spectral law by $d\kappa_\lambda$. Rotating back to real
Euclidean $q$ gives the exact formula
\begin{equation}
 S(q)=\int_{\mathbb R}\frac{r-i\gamma\cdot q}{X+r^2}
                    \,d\kappa_\lambda(r).
 \label{eq:next3-spin-spectrum-uv}
\end{equation}
Thus the vector spectral measure is the pushforward under $r\mapsto
r^2$ and the scalar measure is the pushforward of $r\,d\kappa_\lambda$.
They obey $|d\rho_{\rm scalar}|\le\sqrt{s}\,d\rho_{\rm vector}$.
Each integrand is a free Dirac covariance of signed mass $r$ with the
usual reflected spin dual; either mass sign gives a positive CAR
space. Taking the direct integral of those spaces and its fermionic
Fock space realizes the free Wick functions of the projected field.
There is a charged continuum with positive norm. An isolated charged
pole is a further test: in the gap $(-M,M)$ the function $a$ is
increasing and $a(M-)=1/[M(n-1)]$. Since $0<\lambda<M$ and $n=248$,
there is no gap pole for this particular spin transition. The scalar
pole of the previous comparison cannot be transferred to this one.

Gauge covariance of $D_B^{\rm sp}$ gives the usual endpoint phases
and the open longitudinal Ward identity with the same electric
charge $e$. The free spin measure alone does not determine the
gauged determinant. For example at $\lambda=0$ the closed field
response vanishes in the cyclic prescription, whereas a pure-gauge
variation of the open kernel gives the difference of its two endpoint
phases. This limit exposes the independent role of the transition in
the physical current. At nonzero refresh rate, equality of the full
current cut to the actual charged-pair probabilities, interacting
reflection positivity and covariance of the spatial photon averaging
under boosts remain independent requirements.

\paragraph{A separated three-point test for the gauged spinor}
There is a direct spectral obstruction even for the repaired free
spinor. Keep the charged endpoints at $(0,0)$ and $(T,\mathbf x)$
and vary a transverse electromagnetic source at $(t,\mathbf z)$,
with $0<t<T$. At first electric order the open kernel is the sum of
paths with one current mark. The zero-refresh term is particularly
simple. Its free spin kernel at $x=(T,\mathbf x)$ is
$(\gamma\cdot x)w(x)$ with $w(x)>0$: substituting $y=x/u$ in
the streaming inverse gives
$w(x)=\int_{|x|}^\infty e^{-Mu}u^{-5}(d\nu/d^4y)(x/u)\,du$.
Its gauge factor is the Wilson line on this one straight segment.
Put $\vartheta=t/T$. After stripping the common electric phase, the
$i$th spatial current mark is therefore
\[
 \frac{x_i}{T}(\gamma\cdot x)w(x)\,
              \delta^3(\mathbf z-\vartheta\mathbf x).
\]
For the other terms the joint displacement density in $(x,z)$ is
absolutely continuous. At fixed positive arc lengths and an interior
mark on one arc, the two displacement vectors are linear combinations
of the independent four-velocities, with two linearly independent
rows whenever there is at least one refresh. The velocity densities
then push forward to an ordinary eight-dimensional density. End times
of arcs have measure zero. The absolute total mass of all these marked
measures is at most $C\int_0^\infty u e^{-mu}du$, since the Clifford
weights have bounded products. Thus for almost every ordered $T,t$
the refreshed density is $L^1$ in $(\mathbf x,\mathbf z)$.

Take $\mathbf k$ parallel to a fixed axis perpendicular to $e_i$ and
Fourier-transform these two spatial variables at
$(-\vartheta\mathbf k,\mathbf k)$. After taking the spin trace against
$\gamma_i$, the straight contribution is the strictly positive constant
$4T^{-1}\int x_i^2 w(x)d^3\mathbf x$. The remaining transform tends to
zero by the Riemann--Lebesgue lemma. The gradient term in the open
connection vanishes for a spatially transverse probe. Hence the actual
spatially averaged photon vertex has a nonzero algebraic tail along
the successive maxima of $|\chi_n(L|\mathbf k|)|$. This conclusion also
holds under smooth tests in sufficiently small balls of fixed radius
about the indicated spatial momenta: the straight contribution is
continuous in their deviation from the displayed relation, the
refreshed transform vanishes uniformly at infinity, and a smooth
spatial transverse polarization can be used throughout each ball.

By contrast, in a tempered Lorentzian three-point function satisfying
the forward-cone spectrum condition the ordered Euclidean time gaps
$t,T-t$ give Fourier factors $e^{-tE_1-(T-t)E_2}$. At these spatial
momenta the two intermediate spatial norms are
$(1-\vartheta)|\mathbf k|$ and $\vartheta|\mathbf k|$, up to bounded
errors. Each energy is at least its corresponding spatial norm. On
the translated smooth tests temperedness thus bounds the result by a
polynomial times $e^{-c|\mathbf k|}$ for some $c>0$. This contradicts
the algebraic lower bound. The leading marked response cannot itself
be the continuation of such a current matrix element. If it is to be
obtained by differentiating an interacting Lorentzian theory, the
comparison assumes that differentiation is valid as a tempered
distribution and commutes with continuation at the separated times.
The obstruction holds at every finite refresh rate. A charged velocity
law with such a surviving straight component must change its temporal
current sampling or its path law to pass this test.

\paragraph{Why gauging the continuum as independent charged particles fails}
Equation~\eqref{eq:next3-spin-spectrum-uv} normalizes one observed
spinor. Treating every spectral mass as an independently charged Dirac
species would also put those mass labels in the physical current trace.
Choose a compact interval $I\subset(M,\infty)$ of positive continuous
spectral weight. Such intervals exist because the original mass law has
positive smooth density there, and the imaginary part of its resolvent
is multiplied by $|1-\lambda a|^{-2}$ in $f$. With standard diagonal
minimal coupling the pair creation kernel of a spacetime-smeared
current acts diagonally in $\mathcal H_I=L^2(I,d\kappa_\lambda)$.
Choose a timelike probe above $2\sup I$, with nonzero transverse pair
amplitude. The free Dirac pair phase space is strictly positive
throughout this interval. On a subset of positive measure one common
unit spatial-spin packet has fiber image of norm bounded below by a
positive constant. This follows by choosing from a countable dense
set of packets and shrinking the mass subset. Split the subset into
infinitely many disjoint parts, normalize their indicators, and
multiply by that packet. The resulting orthonormal inputs each have
output norm bounded away from zero. The pair kernel is therefore not
Hilbert--Schmidt, whereas its Hilbert--Schmidt norm squared is exactly
the squared norm of the pair created from the CAR vacuum. This is a
divergence on a finite invariant-mass band, independent of the decay
of the observed endpoint at large mass.

There is a crossing test for a more general, possibly mass-changing
vertex. Suppose an elastic spin projection and a crossed pair projection
are values of one bounded operator $T(z)$ on $\mathcal H_I$ after
removing kinematic factors. Assume norm holomorphy on a connected
crossing domain containing the elastic point $z=0$, norm continuity
up to an open ordinary timelike rim, and unit elastic charge
$T(0)=I$. On that rim assume finite inclusive pair probability controls
the squared $L^2(I\times I)$ norm of the kernel of this projection.
This last condition specifies what the measured pair channel must
control; it holds for a bounded spin projection when its two-body
phase-space weights are bounded below there. Then every rim value is
Hilbert--Schmidt and hence compact. Map $T$ to the quotient of bounded
operators by compact operators. Every continuous linear functional
of this quotient gives a scalar holomorphic function with zero
boundary values on the rim. Schwarz reflection across an open real
subinterval and the identity theorem force it to vanish throughout
the domain. Since linear functionals separate points, $T(0)$ would be
compact, contrary to the identity on this infinite-dimensional space.
The argument uses norm analyticity of the specified projection; separate
analyticity on each fixed charged wave packet is a weaker condition.
In particular the charge-one limit can be a strong limit of compact
transitions without being their norm limit. This test distinguishes a
continuous mass spectrum of one measured field from a continuum of
independent particles whose charges all enter the same measured current.

The strength of the topology assumption can be checked with a shell
coordinate. Parameterize the mass interval by its spectral quantile
$u\in[0,1]$, distributed uniformly just as $|P_2\omega|^2$ on the A4
four-direction. On this interval put $H=1-\partial_u^2$ with Neumann
boundary conditions and use dimensionless squared transfer $z=L^2Q^2$.
The transition $T(z)=(I+\sqrt z\,H)^{-1}$ on the principal sheet
approaches the identity strongly as $z\downarrow0$. For $z>0$ its
distance from the identity in norm is one. On the timelike rim
$z=-s+i0$, $s>0$, its squared Hilbert--Schmidt norm is
\[
 \sum_{j=0}^\infty\frac{1}{1+s(1+\pi^2j^2)^2}<\infty .
\]
Thus this concrete mass-changing kernel can retain charge on each
fixed wave packet and create finite inclusive pair norm on the tested
mass interval. The square root places a branch at zero transfer; the
operator has no norm-analytic continuation across the charge point.
It illustrates the price of evading the crossing test. Assigning that
branch to real gapless degrees of freedom and obtaining this vertex
from an electromagnetic action would be separate physical derivations.

\paragraph{Continuous turning, the full velocity trace and charged poles}
The straight component can be removed by propagating the observed A4
four-direction $v\in S^3$ at unit speed while its direction evolves by
spherical Brownian motion. Its stationary law is the normalized
angular law of the four-coordinate shell projection. In the same
Euclidean units consider instead
\begin{equation}
 D_B=m-\kappa\Delta_{S^3}+v\cdot(\partial+i e B),\qquad m,\kappa>0 .
 \label{eq:next3-angular-flight-uv}
\end{equation}
Here $\Delta_{S^3}$ has nonpositive spectrum. One may set $m=\kappa=L^{-1}$
in this comparison. The choice of angular diffusion is additional
dynamics. In particular its paths have no interval of unchanged
velocity almost surely. The action has an ordinary smooth heat kernel
at every positive transport duration. Its diagonal now permits an
estimate on the unmarked charged vacuum itself.

Write $p_t$ for the mass-free transition density of $(x_t,v_t)$,
where $\dot x_t=v_t$ and $v_t$ has generator $\kappa\Delta_{S^3}$.
For $0<t\le1$ one has a crude bound $\sup p_t\le C t^{-N}$ for a
finite $N$. Indeed the rotations of the sphere span its tangent
plane. Their brackets with $v\cdot\partial_x$ span the three base
directions perpendicular to $v$. Bracketing a second time includes
the missing longitudinal base direction. This gives the uniform
finite-step parabolic H\"ormander condition, so the short-time density
bound from Malliavin integration by parts applies~\cite{Next3Kusuoka}.
One may first periodize the base on a torus of side greater than
four and work in finite sphere charts. Finite speed and translation
invariance give the same short-time estimate on the plane. The
transition density of the adjoint process reverses the base drift
and preserves the spherical diffusion.

In fact for $a=(0,v)$ the diagonal satisfies
\begin{equation}
 p_t(a,a)\le C t^{-N}e^{-c/t}\quad(0<t\le1).
 \label{eq:next3-angular-return-uv}
\end{equation}
Put $h=t/2$ and split the midpoint integral of
$p_h(a,z)p_h(z,a)$ at $v\cdot x_z=h/2$. On the lower part a forward
half-path must turn until its velocity has projection at most $1/2$
on $v$. This costs at most $C e^{-c/h}$: in ambient coordinates the
velocity has bounded diffusion matrix and bounded drift, so the
exponential martingale maximal inequality applied to each coordinate
bounds any fixed-size excursion on a short interval. Bound the other
half-kernel by its polynomial supremum. On the upper part use the
adjoint half-path started at $a$: its initial base drift is $-v$, so
reaching $v\cdot x_z>h/2$ requires the same large angular excursion.
This proves the estimate. At large time $\sup p_t\le\sup p_1$ by the
semigroup law and invariance of phase-space volume. Consequently
$\int_0^\infty e^{-mt}dt\int p_t((0,v),(0,v))d\nu(v)/t$ is finite.
It defines the absolute proper-time vacuum trace per base volume on
this generator with its chosen normalization. On a closed unit-speed
path the same bounds $|J(q)|\le\min(t,|q|t^2)$ apply. Angular
parallel propagation without an additional spin potential acts as the
identity in the four spin slots. Thus all its complete charged loops,
including their Wilson contacts, obey the same massless-photon sewing
proof with the spatial A4 connection.

Use again $J\xi=\xi$, $\Xi f=\int\gamma\cdot v\,f(v)d\nu(v)$ for
the charged endpoints. Set $H=m-\kappa\Delta_{S^3}$ and
$C=H^{-1/2}v_1H^{-1/2}$. This compact self-adjoint operator has
trivial kernel, since $v_1$ vanishes only on a null set. Its spectral
law $d\mu(c)$ on the unit constant function is a symmetric
probability by angular parity. With $X=|q|^2$ the scalar matrix
element $F=\langle1,(H+i|q|v_1)^{-1}1\rangle$ satisfies
$mF(X)=\int (1+Xc^2)^{-1}d\mu(c)$. Integrating the angular resolvent
equation against one gives the projected spinor
\begin{equation}
 S(q)=-i\gamma\cdot q\int\frac{d\mu(c)}{X+c^{-2}} .
 \label{eq:next3-angular-spin-uv}
\end{equation}
This is a positive Dirac spectral sum with equal weights for the two
mass signs at each $|c|^{-1}$. Compactness makes the occupied nonzero
$|c|$ discrete, with accumulation possible only at zero. In particular
the largest occupied $|c|$ gives an isolated positive-residue charged
mass level. Any infinitely many further masses can accumulate only at
infinity. The projected free charged Wick functions have a positive
CAR realization. This proof supplies both the true angular vacuum
trace and charged mass poles in one represented operator. The
separated three-point test must now be recomputed for its continuously
turning histories; elimination of the exact straight component alone
does not establish cone support, interacting positivity or charged
cut saturation.

One can measure just the positive-mass Dirac residue at the lowest
level. Let $m_*$ be this derived mass. Apply the covariant differential
map $(m_*-\gamma\cdot\nabla_B)/(2m_*)$ at the observed spinor end
and its Dirac dual at the other end, before source differentiation.
At zero field the two Fourier factors are both
$(m_*-i\gamma\cdot q)/(2m_*)$. On a signed-mass Wightman component
the field amplitude consequently changes by $(m_*+r)/(2m_*)$.
Equivalently put $s=c^{-2}$ in \eqref{eq:next3-angular-spin-uv}.
The filtered Euclidean inverse is
\[
 \frac{1}{4m_*^2}\left[
 i\gamma\cdot q-2m_*
 +\int\frac{2m_*s-i\gamma\cdot q\,(m_*^2+s)}{X+s}\,d\mu(c)
 \right].
\]
The first two terms are the coincident distribution dictated by the
two derivatives; the integral exists without additional moments.
At $s=m_*^2$ its pole has precisely the numerator
$m_*-i\gamma\cdot q$, up to a positive residue. At every higher
level the two signed residues remain nonnegative in their respective
Dirac spaces because the field amplitudes were multiplied before
squaring. This identifies a single-sign charged pole with the
original determinant and no adjustable residue subtraction. Derivatives
of both filtered endpoints are part of any electromagnetic transition
for this measured particle.

\paragraph{The charged mass and its threshold transition on the unit-speed branch}
Both can be determined from the same angular operator. In the zonal
subspace on $S^3$ use the Chebyshev polynomials $U_j(v_4)$ of the
second kind, orthonormal for the normalized semicircle marginal.
Their spherical Laplacian eigenvalues are $-j(j+2)$, as follows from
the zonal spherical Laplacian $(1-u^2)\partial_u^2-3u\partial_u$.
Multiplication by $v_4$ has neighboring coefficients $1/2$,
with $U_{-1}=0$. At $m=\kappa=L^{-1}$ the resulting matrix $C/L$
has zero diagonal and off-diagonals
$d_j=1/[2(j+1)(j+2)]$, $j=0,1,\ldots$. The constant vector is
cyclic on this subspace. The inverse largest positive eigenvalue is
$m_*L$. In particular
\[
                     3.78984<m_*L<3.78985 .
\]
For a reproducible rational bracket, keep ten rows. At a trial
$x$ the diagonal elimination pivots are $u_0=x$ and
$u_{j+1}=x-d_j^2/u_j$ through $j=8$. For $x=1/3.78985$ a pivot
is negative. For $x=1/3.78984$ they are all positive and
$u_9-d_9^2/(x-2d_{10})>0.24$. The omitted tail has norm at most
$2d_{10}$ by row sums, which proves the upper bound on the largest
eigenvalue by its Schur complement. This mass belongs to the stated
angular dynamics. It is not a determination of the physical electron
mass from Newton's constant.

For a double-pole current measurement let $\varphi_+$ be the angular
ground state at the first crossing of $H-Ev_4$, $H=m-\kappa\Delta_{S^3}$,
and let $\varphi_-(v)=\varphi_+(-v)$. The operator is strictly positive
before the crossing. At the crossing its elliptic ground state on
the connected sphere is simple and strictly positive, hence invariant
under spatial $SO(3)$ and visible to the constant endpoint. Conversely
the first crossing must be zonal by this same argument, so it is the
mass just calculated. Its slope is nonzero since
$E\langle v_4\rangle=\langle H\rangle>0$ in $\varphi_+$. The
minimal transverse current vertex is $v_i$ times the identity on the
spin indices. At the opposite rest poles its double residue contains
\[
              \langle\varphi_+,v_i\varphi_-\rangle=0,
                       \qquad i=1,2,3 .
\]
It transforms as a spatial vector between two rotation-invariant
angular states. The observed spinor maps multiply this matrix
element at the two ends. A photon derivative of a differential
endpoint filter supplies only one charged resolvent and cannot
change the double residue. Bounded perturbation by $i\mathbf p\cdot
\mathbf v$ and the isolated angular gap make the double residue
$O(|\mathbf p|)$ for small relative momentum. With ordinary two-body
phase space its pair weight is consequently $O(\beta^3)$ or smaller,
where $\beta=(1-4m_*^2/S)^{1/2}$. The target Dirac transverse
current has a nonzero spin-triplet matrix element at rest and the
A4 mark equals one there; its weight begins at order $\beta$.
Thus the charged pole and the finite unrestricted-duration loop
integrals coexist in one represented operator, yet its minimally
gauged angular motion misses a specific spin transition already at
the first physical threshold.

\paragraph{The radial A4 measure and the clock for changes of direction}
The unit direction of a four-coordinate projection and the magnitude
of that projection have different effects on a return. Write the
displaced velocity in dimensionless units as $y=R u$, with $u\in S^3$
and $0<R\le1$. For the rotationally invariant ball marginal the
conditional law of $u$ is uniform. Its radial law has density
proportional to $R^3(1-R^2)^{(n-4)/2}$, $n=248$. If only velocity is
multiplied by $R$ and the angular diffusion clock is kept fixed, the
base return density per four-volume at a given cycle duration is
multiplied by $R^{-4}$. The radial integral of this factor diverges.
Thus replacing the velocity by its unit direction changes a genuine
boundary-measure test, even though $R$ vanishes with probability zero.

There is a second, intrinsic choice of angular clock. Let turning occur
at a fixed rate per unit distance traveled in the Euclidean readout.
Conditional on the recorded $R$, multiply both the streaming generator
and the angular diffusion generator by $R$. In particular the scalar
charged operator on this branch has the form
\[
 D_{B,R}=m+R[ u\cdot(\partial+ieB)-\kappa\Delta_{S^3} ].
\]
The same time change acts in each of the four spin slots of
\eqref{eq:next3-angular-flight-uv}. The mass factor remains $e^{-mt}$. During a
coordinate interval $t$, the path has geometric length $s=Rt$ and
$R\int_0^t u_{R\tau}\,d\tau=\int_0^{Rt}u_v\,dv$. Consequently its
mass-free diagonal return kernel is precisely the unit-speed kernel
at duration $Rt$. No extra radial Jacobian appears in this identity.
If the unit-speed diagonal is bounded by $Cs^{-N}e^{-c/s}$ near zero,
then it is bounded by $C'e^{-c/(2t)}$ uniformly over $0<R\le1$ for
small $t$. At later times boundedness follows from the unit-speed
short-time estimate and the semigroup bound after any fixed positive
duration. Closed Wilson insertions use the actual path length $Rt\le t$.
The unmarked vacuum and all the fixed-order, source-completed bounds
proved with the unit clock consequently admit the full radial average.
A common $R$ shared across several cycles is harmless for these
uniform estimates. This radius is a classical geometric record, with
its prescribed probability trace; declaring an independent Berezin
field at every $R$ would require a separate trace over those fields.

For each $R>0$, its projected free inverse is $R^{-1}$ times the
unit-speed projected inverse with mass coefficient $m/R$. The
positive-mass free spinor reconstruction on this angular branch
therefore passes to a positive mixture over $R$. The positive factor
$R^{-1}$ is harmless on Euclidean tests since the accretive resolvent
bound before factoring $R$ is $1/m$. With coordinate mass $m$ the averaged measure in fact has no mass
atoms. For a prescribed $M>0$, a pole at $M$ requires $-m/R$ to be
an eigenvalue of $-\kappa\Delta_{S^3}-Mv_1$ (or its parity conjugate).
This elliptic operator has discrete spectrum on the compact sphere,
so only countably many radii can contribute to that atom. The A4
radial measure assigns those radii zero weight.

A mass accumulated per unit geometric arclength gives a different
physical preparation. Let the full generator at $R$ be $R D_B^{\rm unit}$,
where $D_B^{\rm unit}$ is \eqref{eq:next3-angular-flight-uv} with
$m=\kappa=L^{-1}$. At every $R>0$ the substitution $s=Rt$ preserves
$dt/t=ds/s$. Both the unmarked loop trace and the marked Wilson
amplitudes integrated over $s$ coincide with those at unit speed,
uniformly as completed integrals in $R$. The free projected inverse
at $R$ is $R^{-1}S(q)$. Its radial average is finite since
$\mathbb E[R^{-1}]<\infty$. Choosing the measured endpoint amplitude
$\sqrt R$ at each end gives exactly $S(q)$ in the full radial average,
with its discrete pole and its threshold selection rule. This choice
of mass and source weighting reparametrizes the same unit-speed paths.
It supplies no additional independent displacements for simultaneous
hard circulations. Both radial constructions use the stated probability
on the classical shell record. An independently quantized radial
multiplicity would require its own determinant and trace.

\subsection{Which measurement can distinguish two finite charged sources?}
\label{sec:four-map-uv}
A finite vacuum graph can still assign the wrong weight to a physically
allowed intermediate state. To locate the discrepancy, keep the same
aligned sector plane, directed action, GTL constraints and velocity
shell as in Section~\ref{sec:next3-velocity-uv}. Work first with
the intrinsic unit-arclength angular generator. The polarizations,
Euclidean continuation and masses below refer to the chosen plane.
In particular a Euclidean frequency in a rest-frame test is a component
of the continued four-momentum. Its sign does not replace the sector's
vectorial pole prescription.

Three measurements expose different parts of a charged source. The
two sides of its first pair cut compare the determinant with an open
charged transition. Its magnetic moment measures the elastic current
near zero transfer. Finally the frequency dependence of its full
two-current response tests whether a positive physical spectrum could
produce the same Euclidean data. The last test also detects changes
that pass both of the first two.

\subsection{The signed light-pair cut of the angular determinant}
\label{sec:fourcone-uv}
The first threshold test of \eqref{eq:next3-angular-flight-uv}
concerns an observed open pair. The closed determinant provides an
independent test of the same transition. For the minimally gauged
angular operator its first nonzero pair weight has the wrong sign.
Work in the pair rest frame on the declared matter boundary. The
spatial photon readout then equals one. Write $M=m_*$,
$H=m-\kappa\Delta_{S^3}$, and $R(k)=(H+i k\cdot v)^{-1}$ on one
angular spin slot. The four spin components propagate identically.

\paragraph{The poles of the full angular propagator}
The poles of $R(k)$ at complex four-momentum can be identified
without assuming that angular eigenfunctions carry a bounded
representation of complex boosts. The compact operator
$H^{-1/2}(k\cdot v)H^{-1/2}$ is Hilbert--Schmidt since
$H^{-1/2}$ has summable fourth powers of singular values on $S^3$.
The regularized determinant $\det_2(I+iH^{-1/2}(k\cdot v)H^{-1/2})$
is an entire function of the four components of $k$. Every homogeneous
Taylor term is a rotation-invariant polynomial, hence a multiple of
a power of $k\cdot k$ or zero. This follows by restricting first to
real spheres and then using polynomial continuation. On the fourth
axis the zeros therefore fix the complete polar set:
$k\cdot k=-c^{-2}$ for the nonzero eigenvalues $c$ of
$H^{-1/2}v_4 H^{-1/2}$. Compactness isolates the lowest mass, and the
elliptic ground state is simple in the angular space at each of its
two rest poles. None of these facts determines the sign of a joined
current matrix element.

\paragraph{The sign after joining the two residues}
Let $z=v_4$ and normalize the real positive ground function $f(z)$ of
$H-Mz$ to have squared norm one. Put $a=\int z f(z)^2 d\nu>0$.
The positive and negative rest residues of $R$ with respect to
$k^2+M^2$ are $Z_* |f(\pm z)\rangle\langle f(\pm z)|$ with
$Z_*=2M/a$. This follows by differentiating the simple eigenvalue
of $H-Ez$ at $E=M$. For a small common spatial momentum $p$ at the
pinch, $w=(M^2+|p|^2)^{1/2}$. The angular functions at the poles
$k_4=\pm iw$ are
\[
 f_\pm(v)=f\big((\pm w z-i p\cdot\mathbf v)/M\big).
\]
Indeed these are small complex rotations of the two real solutions.
The zonal ODE has a regular analytic solution at both poles of the
sphere, so this continuation is legitimate near threshold. The
resolvent is symmetric under transpose in the real angular basis.
The residues and their normalization therefore use bilinear
continuation of the real projectors. At the two conjugate physical
momenta $f_-(v)=\overline{f_+(-v)}$.

For a spatial polarization $i$ the angular vertex is multiplication
by $v_i$. Its two matrix elements inside the trace are the same
bilinear number $L_i=\int f_+(v)v_i f_-(v)d\nu(v)$, by transpose
symmetry. Reversing $v$ in its complex conjugate gives
$\overline{L_i}=-L_i$. In fact
\[
 \begin{aligned}
 L_i&=-\frac{i p_i}{M}b+O(|p|^3),\\
 b&=\int v_i^2[f'(z)f(-z)+f(z)f'(-z)]\,d\nu>0 .
 \end{aligned}
\]
To see strict positivity, the derivative of
$(1-z^2)^{3/2}f'(z)$ is $(1-z^2)^{1/2}(m-Mz)f(z)/\kappa$.
The primitive vanishes at both endpoints, and its derivative changes
sign exactly once, from positive to negative. Thus $f'>0$ inside.
Rotational symmetry gives the same $b$ for each spatial component.

Use the scalar-current convention in which an ordinary charge-one
Dirac loop has density $(1+2M^2/S)\beta/(12\pi^2)$, where
$\beta=(1-4M^2/S)^{1/2}$. With two propagator poles and vertices
$i e v_i$, the same loop sign and two-body phase space multiply the
spin numerator $4Z_*^2 L_i^2$ by $\beta/(16\pi^2 S)$ and average
it over relative directions. For calibration, the Dirac threshold
spin numerator is $8M^2$, giving $\beta/(8\pi^2)$ to leading order.
Here $\langle p_i^2\rangle=S\beta^2/12$, giving the closed,
unit-charge density
\begin{equation}
 \rho_{\rm ang}(S)=-\frac{b^2}{12\pi^2 a^2}\,\beta^3
                         +O(\beta^5).
 \label{eq:fourcone-signed-uv}
\end{equation}
At $m=\kappa=L^{-1}$ the zonal Jacobi calculation gives
$a\simeq0.45634$, $b\simeq0.33372$, and a leading coefficient
approximately $-0.00451550$. The sign follows from the squares and
the strict inequalities above. The result concerns the actual closed
Grassmann trace. Filtering the external charged residues changes the
open observable alone. No positive inclusive sum of such observed
states can equal this negative density.

\paragraph{Why the hard angular modes cannot cancel the threshold}
The cut can be taken in a bounded neighborhood of the two poles.
Here is a bound justifying that localization for this nonlocal
charged response. For real $k$ rotate $k\cdot v$ to $|k|z$ and put
$u=R(k)g$, $\|g\|\le1$. The real part of the angular energy identity
gives $\|u\|_{H^1}\le C$. Test its imaginary part against $u$ times
a smooth approximation to $\operatorname{sgn}z$, with derivative
bounded by $C/\delta$, equal to the sign for $|z|\ge\delta$ and
of the same sign elsewhere. This gives
$\int_{|z|>\delta}|u|^2\le C/(|k|\delta^2)$. In the complementary
strip the $H^1$ to $L^6$ sphere bound gives $C\delta^{2/3}$.
Taking $\delta=|k|^{-3/8}$ proves $\|R(k)\|\le C|k|^{-1/8}$ at
large real momentum. Also $\|R(k)\|_{\rm HS}\le\|H^{-1/2}\|_4^2$
by the spectral theorem in the middle angular factor. Both bounds
persist for bounded complex shifts there by a Neumann series.

Take sufficiently many derivatives in the external rest frequency.
Routing it through one propagator gives one additional $R$ per
derivative; two Hilbert--Schmidt factors bound the trace. More than
32 additional factors make the hard four-momentum integral locally
holomorphic on every compact continuation range. To justify passage
from the cyclic heat trace to this differentiated bubble one can
first smooth the base return delta. Each frequency derivative gains
at most one power of the total duration from the displacement between
the two marks. The small-time majorant is $Ct^{N+1-N_0}dt$ for $N$
derivatives if the mass-free density supremum is $Ct^{-N_0}$. It is
integrable for sufficiently large $N$, uniformly under normalized
base smoothing. Positive-time kernels then give the differentiated
resolvent formula by convolution, and the smoothing can be removed.
For the hard contribution the smoothings can equivalently be removed
at the undifferentiated level in its sufficiently subtracted Taylor
remainder. The local integration constants are analytic and do not
supply a threshold discontinuity. In the remaining bounded region
the isolated simple angular poles give the stated pinch by the usual
two-line continuation~\cite{Cutkosky1960}. This argument establishes local threshold
behavior while leaving possible complex growth at infinity unrestricted.

\subsection{Spin transport on the sphere of velocity directions}
\label{sec:fourspin-berry-uv}
The zero at the first pair threshold in
\eqref{eq:next3-angular-flight-uv} arose because the two angular
states were spatial scalars and the spin index was a spectator.
There is a minimal coupling on the same sphere for which both
threshold matrix elements and the closed-loop spin trace can be
computed. In a Hermitian Euclidean Clifford representation put
$P(v)=(1+\gamma\cdot v)/2$, $v\in S^3$. The range of $P(v)$ is a
two-dimensional complex spin space. On its sections use the induced
unitary connection $\nabla^P=P\,d$ and the operator
\begin{equation}
 \mathcal D_B=m+\kappa(\nabla^P)^*\nabla^P
                     +v\cdot(\partial+ieB).
 \label{eq:fourspin-berry-action-uv}
\end{equation}
The projector transforms with the spin representation under every
oriented Euclidean rotation. Its choice is additional charged
representation data in the directed action. Unitary parallel transport
on the angular Brownian path replaces the previous spin multiplicity
by a closed holonomy of absolute trace at most two. It changes neither
the path, the Abelian Wilson factor nor its differentiated contacts.
The exponential return bound \eqref{eq:next3-angular-return-uv}
therefore controls the absolute unmarked charged trace per volume.
With the instantaneous spatial A4 photon preparation the same joint
photon-edge estimate proves convergence of the source-completed closed
graphs at every finite electromagnetic order. The intrinsic-arclength
choice also gives the previous uniform estimates over the projected
radius. The spectral calculation below uses unit speed, or the common
spectrum with the intrinsic mass and endpoint normalization.

\paragraph{Two radial angular equations}
Set $v=(\sin\theta\,n,\cos\theta)$, $|n|=1$, and let
$g=\cos\theta+i\sin\theta\,n\cdot\boldsymbol\sigma$.
Take $\gamma\cdot v$ to have off-diagonal blocks $g,g^\dagger$.
A local orthonormal frame for $P(v)$ is the two-block column
$W=(g^{1/2},g^{-1/2})^{\mathsf T}/\sqrt2$, where the superscript
indicates block stacking. With $b=\sin^2(\theta/2)$, its connection
is $W^\dagger dW=i b(n\mathbin{\times}dn)\cdot\boldsymbol\sigma$.
The covariant two-sphere Laplacian consequently contributes
$[\mathbf L^2+2b\,\boldsymbol\sigma\cdot\mathbf L+2b^2]/\sin^2\theta$
to the angular kinetic energy, with $\mathbf L=-i n\times\nabla_n$.
The polar part is $-\partial_\theta^2-2\cot\theta\,\partial_\theta$.
For orbital $l$ the spin--orbit eigenvalues are $l$ and $-(l+1)$,
the latter available only for $l\ge1$. The lowest potential in each
family occurs at $l=0$ and $l=1$, respectively. They are
\[
 V_0(\theta)=\tfrac12\tan^2(\theta/2),\qquad
 V_1(\theta)=\tfrac12\cot^2(\theta/2).
\]
Their angular spinors are $\xi$ and $i(n\cdot\boldsymbol\sigma)\xi$.
The endpoint conditions are those of the covariant kinetic form on
the smooth sphere, equivalently the regular radial Friedrichs forms.

At positive rest energy $E$, the zero of the inverse is tested in
$m+\kappa(\nabla^P)^*\nabla^P-E\cos\theta$. Its lowest radial
state lies in $V_0$. To compare the two minima, take a positive radial
ground state of $V_1$ with the tilt $-E\cos\theta$. At each pair of
reflected angles place the larger of its values on the northern half.
This preserves the norm and radial derivative integral. For $V_0$
and the tilt, both lower potentials now coincide with the larger
value; for $V_1$ and the same tilt their lower values were on opposite
sides. One of the two inequalities is strict on a set of positive
measure, since the $V_1$ ground state cannot be symmetric when the
odd part of its potential is nonzero. The other partial waves lie
above the two compared minima by the displayed angular eigenvalues.
At $E=0$ the $V_0$ ground state is $\sqrt{1+\cos\theta}$ in normalized
spherical measure, with covariant eigenvalue $3/4$. The lowest
$V_0$ eigenvalue decreases strictly thereafter and crosses zero once:
its derivative is $-a(E)$, where $a(E)=\langle\cos\theta\rangle$ is
strictly increasing from $1/4$. The reduced-resolvent second derivative
is strictly negative because multiplication by $\cos\theta$ has a
component orthogonal to the radial ground state. A test in a northern
cap gives a negative eigenvalue for sufficiently large $E$.

\paragraph{The measured pair and its closed spin trace}
Write $E_*$ for the first crossing, $u=\cos\theta$, and $f(u)>0$ for
the normalized northern radial ground state there. At the opposite
rest pole the angular state is $f(-u)i(n\cdot\boldsymbol\sigma)\xi$.
With $d\nu_0(u)=2\sqrt{1-u^2}\,du/\pi$, define
\[
 \begin{aligned}
 a&=\int u f(u)^2\,d\nu_0(u),\\
 b_*&={1\over3}\int \sqrt{1-u^2}\,f(u)f(-u)\,d\nu_0(u).
 \end{aligned}
\]
Between these two doublets the minimally gauged spatial velocity is
$i b_*\sigma_i$, up to a common phase convention, since
$\int n_i n_j\,d\Omega/(4\pi)=\delta_{ij}/3$. Its reverse is its
adjoint. Constant spinors projected by $P(v)$ see both poles: angular
integration of $W$ gives the northern combination of chirality blocks
$(I,I)$ and the southern combination $(I,-I)$, with equal nonzero
absolute overlap. The rest propagator factors are the corresponding
orthogonal projections divided by $a(E_*+ip_4)$ and $a(E_*-ip_4)$.
The temporal diagonal matrix element is $\pm a$, so pole normalization
gives unit charge. The normalized transverse transition has magnitude
$b_*/a$ relative to a unit Dirac vertex. In the closed bubble the two
resolvents have exactly the same pole signs as the free Dirac pair.
Both vertices carry $ie$ and the cyclic spatial spin trace is
$2b_*^2\delta_{ij}$. Thus the light-pair cut and the squared observed
transition have the same positive leading weight. It begins at order
$\beta$ and has fraction $(b_*/a)^2$ of the specified Dirac threshold
coefficient. A single-resolvent contact has no such double pole.

This fraction is strictly between zero and one for every $m,\kappa>0$.
Indeed Cauchy--Schwarz applied to $\sqrt{1-u}\,f(u)$ and
$\sqrt{1+u}\,f(-u)$ gives $3b_*\le1-a$, whereas $a>1/4$ at the
crossing. For $m=\kappa=L^{-1}$ one finds $E_*L\simeq4.27635$ and
$(b_*/a)^2\simeq0.04856$. To reproduce the numbers factor
$\sqrt{1+u}$ from the northern radial state. The remaining orthogonal
polynomials have Jacobi weight $(1-u)^{1/2}(1+u)^{3/2}$ and kinetic
eigenvalues $j(j+3)+3/4$. Diagonalize the tridiagonal multiplication
by $u$ between inverse square roots of these shifted energies including
$m/\kappa$. The reciprocal largest positive eigenvalue gives $E_*/\kappa$.
The inequalities above are independent of truncation. Intrinsic angular
spin transport restores the threshold power while fixing a deficient
normalization. Neither its correct charge nor a rescaling of external
spinor packets supplies the missing transverse amplitude.

\subsection{What a positive charged endpoint measures}
\label{sec:four-endpoint-uv}
The Berry generator admits a positive, canonically normalized free
spinor readout on its entire mass spectrum. Coupling this readout to
a varying connection tests an additional part of the boundary law.
Let $J\xi(v)=P(v)\xi$, with normalized spherical measure, and put
\[
 H=m+\kappa(\nabla^P)^*\nabla^P,\quad R(p)=(H+i v\cdot p)^{-1},
 \quad F(p)=J^\dagger R(p)J.
\]
Here $J^\dagger J=I/2$ and $HJ=hJ$, where $h=m+3\kappa/4$.
The second identity follows either from the two radial equations
or from the projected connection on a constant ambient spinor.
Set $A_0=1/(2h)$, so $F(0)=A_0 I$.
All the following momenta are Euclidean variables on this one plane.
Continuation to its mass shell uses the declared vectorial boundary.

\paragraph{A positive spectral readout on every free angular level}
Take $p=i z e_4$ and a unit constant spinor with $\gamma_4\xi=\xi$.
Let $d\mu(c)$ be the spectral measure of the compact self-adjoint
operator $C=H^{-1/2}v_4H^{-1/2}$ in $H^{-1/2}J\xi$.
There is no kernel, since $v_4$ vanishes only on a null set.
Rotation symmetry makes the measure independent of the spin
polarization. Reflection of the three spatial components is implemented
on spinors by $\gamma_4$, and commutes with the fourth-axis resolvent.
Rotations taking $e_4$ to $-e_4$ interchange its two spin blocks.
Consequently
\[
 F_+(z)=\int\frac{d\mu(c)}{1-zc},\qquad
 F_-(z)=F_+(-z),\qquad \int d\mu=A_0 .
\]
A positive measure here is a measure in inverse velocity-weighted
energy $c$. In a spinor resolvent the coefficient of
$(c^{-1}-z)^{-1}$ in $F_+$ is $c^{-1}\mu(\{c\})$.
Both signs of $c$ occur in this observed block. Indeed the lowest
northern-channel eigenvalue of $H+zv_4$ starts positive and becomes
negative on a trial function in the southern cap for large positive
$z$. Its first zero has a positive radial ground state with nonzero
overlap with the $J\xi$ radial factor $\cos(\theta/2)$. The same
argument at the opposite tilt gives a positive $c$ pole. A negative
coefficient of $(r-z)^{-1}$ is incompatible with a positive spinor
spectral weight in this block, including when the signed mass $r$ is
negative. A polynomial contact leaves that residue unchanged.

There is an exact free repair. With $\slashed p=\gamma\cdot p$
(where the displayed slash denotes Clifford contraction), define
\begin{equation}
 \begin{aligned}
 S_0(p)&=(-i\slashed p)^{-1}\frac{F(p)-A_0 I}{A_0}\\
       &=\int_{\mathbb R}\frac{r-i\slashed p}{p^2+r^2}\,d\nu(r),\\
 \nu&=(c\mapsto c^{-1})_*(\mu/A_0).
 \end{aligned}
 \label{eq:four-endpoint-spectral-uv}
\end{equation}
On the positive fourth-axis spin block this is just
$(F_+(z)-A_0)/(A_0z)=\int(r-z)^{-1}d\nu(r)$. The other block
and real four-momenta follow from reflection, rotation, and analytic
continuation from the origin. The right side removes the apparent
singularity at $p=0$. Directly, $J^\dagger v_i J=\gamma_i/8$
and $S_0(0)=I/(4h)$. The measure is a probability with
$|r|\ge1/\|C\|>0$. Each integrand is a free positive Dirac
covariance of signed mass $r$ (a change of mass sign preserves the
positive Hilbert norm). Thus the free spinor has the canonical
equal-time normalization and a positive spectral representation.
This operation uses the angular generator itself on both mass
branches. It does not alter its closed determinant.

\paragraph{A Ward identity leaves a transverse singularity}
A natural candidate for covariantizing the first expression in
\eqref{eq:four-endpoint-spectral-uv} is
\[
 \begin{aligned}
 -d_B S_B&=(F_B-A_0 I)/A_0,\\
 d_B&=\gamma\cdot(\partial+ieB),\qquad
 F_B=J^\dagger\mathcal D_B^{-1}J .
 \end{aligned}
\]
An actual gapped source must be regular also when just one leg has
zero Euclidean momentum. Vary once with a probe of polarization
$i$ taking $p$ into $p'$. At $p'=0$, the numerator for division
by $-i\slashed p'$ is
\[
 N_i(p)=ie\left[\gamma_i S_0(p)
         -\frac{1}{h A_0}J^\dagger v_i R(p)J\right].
\]
The longitudinal contraction vanishes by the defining free identity and
$i p_iJ^\dagger v_iR(p)J=I/2-hF(p)$, summed over $i$.
If every $N_i$ vanished for all real $p$, Clifford contraction would
give $F(p)=4h A_0S_0(p)=2S_0(p)$, since
$\sum_i\gamma_i J^\dagger v_iR(p)J=F(p)$ on the projected bundle.
The defining identity would then force
$S_0(p)=(2/A_0+i\slashed p)^{-1}$, a single positive mass.
This contradicts the spectral support just proved. Thus there exists
a transverse Euclidean probe and a real $p$ for which division
produces a $1/|p'|$ endpoint singularity. Gauge covariance alone has
passed its longitudinal test and failed to keep this massive source
regular.

\paragraph{A regular Wilson-coupled boundary on real Euclidean momenta}
There is a different covariant extension with an explicit contact at
each order. Set $M(p)=S_0(p)F(p)^{-1}$. For real $p$ the middle
resolvent in $R(p)$ has Hermitian part
$(I+(H^{-1/2}(v\cdot p)H^{-1/2})^2)^{-1}$. In particular
\[
 \mathop{\rm Re}F(p)\ \ge\
   \frac{A_0}{1+(|p|/m)^2}\,I .
\]
Thus $M$ is smooth on the real four-plane. Its derivatives of every
fixed order have polynomial bounds, by the resolvent derivative
identity, $\|R(p)\|\le1/m$, and the displayed inverse bound.
Transport its translation-invariant kernel along a straight chord,
\[
 M_B(x,y)=\exp\!\left(-ie\int_y^x B\cdot d\ell\right)
                 \check M(x-y),\qquad S^{\rm chord}_B=M_B F_B .
\]
For a smooth Schwartz connection the product with the chord phase
is defined distributionally. Both factors transform with the same
endpoint gauge rule and $S^{\rm chord}_0=S_0$. For plane-wave probes
the linear coefficient of $M_B$ is $e\int_0^1
\partial_{p_i}M(p+t(p'-p))\,dt$ in the convention $-i\partial+eB$.
Higher chord variations integrate higher derivatives of $M$ over
unit cubes, as follows by differentiating the Abelian exponential.
These formulas remain regular for bounded real momenta, including
$p'=0$. They supply a particular measured dual endpoint and all its
contacts. The charged determinant still uses $\mathcal D_B$. Regularity
on the real axis and the positive free mass mixture leave two
physical tests outstanding: reflection of this extended boundary at
nonzero coupling and agreement of its open cuts with the determinant
at higher thresholds. The next estimate controls all differentiated chords at real momenta.

\paragraph{The Wilson endpoint has bounded differentiated marks}
There is a uniform estimate even when several photon transfers are
large. For a real unit direction $n$ extend $F_n(z)=F(zn)$ into a
sector $\Re z\ge\delta|z|>0$. In the spectral calculus of
$H^{-1/2}(v\cdot n)H^{-1/2}$ one has
\[
 \Re\frac{z}{1+izc}=\frac{\Re z}{|1+izc|^2},\qquad
 \Re(zF_n(z))\ge\frac{A_0\Re z}{(1+|z|/m)^2}\,I .
\]
The denominator has no zero there. This gives
$\|F_n(z)^{-1}\|\le(1+|z|/m)^2/(A_0\delta)$.
For $|z|\ge1$ the endpoint formula
$M(zn)=(-iz\gamma\cdot n)^{-1}(I/A_0-F_n(z)^{-1})$
therefore gives $\|M(zn)\|\le C_\delta(1+|z|)$.
Cauchy's estimate on a disk of radius proportional to $r$ about each
$r\ge1$ bounds its $k$th radial derivative by $C_k r^{1-k}$.
Reflection and rotation covariance give
$M(rn)=u(r)I+w(r)\gamma\cdot n$. Angular derivatives in Cartesian
momentum accordingly cost one inverse power of $r$ each. Smoothness
at the origin completes the symbol bound
\[
       \|\partial_p^\alpha M(p)\|
            \le C_\alpha(1+|p|)^{1-|\alpha|}.
\]
Every chord term differentiated at least once is thus uniformly
bounded, regardless of the momenta sampled in its unit-cube integral.
In an open chain $M_B F_B$ the undifferentiated $M$ acts at the
observed outgoing momentum, so it costs at most its first moment
against an external packet. Inside $F_B$ every real minimal angular
resolvent has norm at most $1/m$, and each minimal mark is bounded.
These estimates eliminate ultraviolet growth of the marked open
vertex at fixed external momenta. An open endpoint can still have a
nonvanishing soft mark; infrared sewing requires its own momentum
routing and estimates.
\paragraph{Absolute sewing on charged Schwartz packets}
For the minimal Berry action with the chord readout one can also do
the soft routing. Consider a connected graph at fixed order whose
external legs are charged Schwartz packets, with at least one open
chain and no prescribed external photon transfers. At each open
chain use its momentum delta to eliminate the outgoing momentum.
Its packet absorbs the possible linear endpoint factor there,
uniformly in all transfers. The incoming packet variable remains
absolutely integrable. Joint Schwartz packets obey the same bound
by a product of sufficiently high inverse powers.

Let $h_j(Q)=|\chi_{248}(L|\mathbf q|)|^2
\min(1,|Q|)^{2-j}/|Q|^2$ when $j$ of the two ends of a sewn
photon lie on open chains. Each closed Wilson mark supplies one soft
zero. Its duration cost has finite moments by the diagonal return
bound. In four dimensions $h_0\in L^1\cap L^\infty$,
$h_1\in L^1\cap L^2$ and $h_2\in L^1$. Integrate the unconstrained
open--open edges directly. In each connected component formed only
by closed vertices and their mutual edges, use its internal tree to
eliminate all but one momentum delta; the eliminated factors cost at
most $\|h_0\|_\infty$. Integrate its remaining internal edges in $L^1$.
The last delta sets the sum of its $k$ incident open--closed transfers
to zero. For $k\ge2$ their integral is bounded by
$\|h_1\|_2^2\|h_1\|_1^{k-2}$ by convolution.

The case $k=1$ vanishes before sewing. The closed path law is invariant
under time reversal together with $v\mapsto-v$; the Abelian line
integral changes sign. A constant chirality matrix identifies the
antipodal Berry projectors and their induced connections. Reversed
parallel transport gives inverse holonomy, with the same trace:
the connection takes values in $\mathfrak{su}(2)$ in the quaternionic
frames above. Thus every odd closed minimal Wilson vertex is zero.
A component incident to a single open--closed edge has odd total
valence. This proves fixed-order absolute convergence of the stated
packet coefficients, including differentiated endpoint contacts.
It does not address sharp exceptional external momenta, on-shell
limits, or reflection positivity of the interacting source.

\paragraph{A finite Wilson circulation supplies an electric spin vertex}
A curvature transition can be introduced without an unbounded photon
frequency at its vertex. This is useful when the photon is averaged on
one clock slice. Work with the continuous angular generator and the
same Abelian connection $B$ as in \eqref{eq:next3-angular-flight-uv}.
For an oriented orthonormal two-frame $(a,b)$ in the Euclidean
sector-plane readout let $C_{ab}(x)$ be the positively oriented circle
of radius $L$ centered at $x$ in its plane. Average the two-frame with
normalized rotation-invariant measure. With Hermitian Euclidean
Clifford matrices define the multiplication operator
\begin{equation}
 \mathcal P_B(x)=\frac{e}{L}\int \gamma(a)\gamma(b)
                    \oint_{C_{ab}(x)} B\cdot dy\ d\mu(a,b).
 \label{eq:four-circulation-uv}
\end{equation}
Use $D_B+\mathcal P_B$ in both the Grassmann determinant and the
open inverse. Reversing the circle orientation changes both factors
in the integral by a minus sign. A gauge gradient has zero circulation.
For real $B$ the added spin matrix is anti-Hermitian. Its ordered
exponential along any real history is therefore unitary. The coefficient
$1/L$ specifies a geometric strength; its choice still requires a
charged representation of the GCM curvature. In particular the
following threshold calculation does not fix the overall physical
strength from the angular diffusion alone.

For $Q=q e_4$ and spatial polarization $i$, Stokes' formula on the
planar disk gives the linear mark
\[
  \frac{e}{L}\,i q\,\gamma_4\gamma_i\,
        \frac{\pi L^2}{6}\Phi(Lq),\qquad
        \Phi(z)=8J_2(z)/z^2,\quad \Phi(0)=1.
\]
Indeed $z_0=|P_{ab}e_4|^2$ is uniform on the unit interval and
$\sum_i(a_4 b_i-b_4 a_i)^2=z_0$. At fixed $z_0$ each spatial
direction contributes $z_0/3$. The normalized disk transform is
$2J_1(Lq\sqrt{z_0})/(Lq\sqrt{z_0})$. Integrating against $z_0/3$
yields the expression above. For $q=iE$, $E>0$, the area coefficient
is $4\pi L^2 I_2(LE)/[3(LE)^2]>0$. The rest vertex at $E=2m_*$
is independent of the charged angular state. Its double pole thus
has the strictly positive angular overlap
$\int_{S^3}\varphi_+(v)\varphi_-(v)d\nu(v)$. On the two spin spaces
selected by the single-sign pole filter, $\gamma_4\gamma_i$ maps
one eigenspace of $\gamma_4$ onto the opposite eigenspace. Hence a
transverse rest matrix element is nonzero, and the squared open
transition begins at order $\beta$ under ordinary two-body phase
space. Endpoint derivatives with only one resolvent still have zero
double residue. Obtaining the target coefficient and accounting for
all the states of the closed determinant are distinct tests.

The unmarked angular heat trace is exactly the previous one. Every
spin mark on a plane wave obeys the perimeter bound
$C\min(1,L|Q|)$ with a fixed prefactor in carrier units: subtract
the phase at the circle center for the small-$Q$ estimate. Central
Wilson insertions obey the closed-path estimate used for the angular
vacuum. At fixed photon multiplicity the ordered spin insertions cost
at most a polynomial in the proper duration. Absolute values are
bounded before the angular bridge and duration integrations by the
same diagonal return estimate. With the spatial A4 photon readout
each sewn massless photon again has an $L^1\cap L^\infty$ envelope
$C|\chi_{248}(L|\mathbf q|)|^2\min(1,L|Q|)^2/|Q|^2$.
Thus the unmarked charged vacuum and the complete source-dependent
closed Euclidean graphs remain finite at every fixed order. In
particular the electric-spin transition uses neither an extra frequency
cutoff nor a fourth-coordinate multiplier on the Maxwell propagator.
The circles themselves extend in Euclidean time. A separated-time
physical-state test and the high-energy open versus closed spectral
weights do not follow from these ultraviolet and threshold results.

\paragraph{Matching the first transverse pair channel}
There is a sign and multiplicity check on the circulation. Write
$\mathcal B[B]=\int\gamma(a)\gamma(b)\oint_{C_{ab}}B\cdot dy\,d\mu$
and $A(q)=\pi L^2\Phi(Lq)/6$. Let $E_0=2m_*$ and normalize the real
positive angular modes $\varphi_\pm$ of $H\mp m_*v_4$ to norm one.
Set $a_*=\langle\varphi_+,v_4\varphi_+\rangle>0$ and
$\eta=\langle\varphi_+,\varphi_-\rangle>0$. Time reversal sends
$\varphi_+(v)$ to $\varphi_-(v)=\varphi_+(-v)$.
The temporal pole kernels in the two time directions have weights
$|\varphi_\pm\rangle\langle\varphi_\pm|\otimes I_4/a_*$.
Indeed, the eigenvalue slopes at the two simple zeros have opposite
signs, as do the orientations of the respective Fourier contours.
For the previously filtered endpoints, $b=\int\varphi_+d\nu$ and
$d=\int v_4\varphi_+d\nu$ give the positive rest residue $bd/a_*$
on $P_\pm=(1\pm\gamma_4)/2$: the factor $\Xi\varphi_\pm$ equals
$\pm d\gamma_4$ and $\varphi_\pm^*J=bI$. Two charged-pole residues
followed by external-leg normalization consequently give an angular
factor $\eta/a_*$ multiplying a spin mark independent of $v$.
The determinant has that factor squared for the same two poles.

For a real anti-Hermitian circulation coupling $e g\mathcal B$,
$g\in\mathbb R$, the two marks at $q=\pm iE_0$ have opposite real
coefficients times $\gamma_4\gamma_i$. Their product has positive
spin trace, since $(\gamma_4\gamma_i)^2=-I$. The corresponding
Euclidean Dirac trace, with the same determinant and contour signs,
is $\operatorname{tr}[P_+(ie\gamma_i)P_-(ie\gamma_i)]=-2e^2$.
Multiplying the circulation by $i$ corrects the sign but still counts
both spin transitions in the unfiltered closed line.

An explicit, different constitutive choice for the charged kernel is
\begin{align*}
 D_B^{\rm match}&=D_B- i e g\,
 \frac12\left(1-\frac{\gamma\cdot\partial}{E_0}\right)
                \mathcal B[B],\\
 g&=\frac{a_*}{E_0 A(iE_0)\eta}.
\end{align*}
The derivative acts only on the averaged circulation. Its value on a
plane wave is a bounded function of the real photon momentum: the
isotropic transverse form factor is $i A(|Q|)\gamma(Q\wedge B)$
with $A(r)=O(r^{-5/2})$ for large real $r$, where
$\gamma(Q\wedge B)=[\gamma(Q),\gamma(B)]/2$. Thus the two pieces
scale there as $O(r^{-3/2})$ and $O(r^{-1/2})$, respectively; the
small-momentum zero is preserved. The same fixed-order marked-loop
bounds apply to the new matrix vertices. Arbitrary-background spin
transport need no longer be unitary.

At $q=\pm iE_0$ the differential factor becomes the left projector
$P_\pm$. Since $P_\pm\gamma_4=\pm P_\pm$, the two marks become
$+i e g E_0 A(iE_0)P_\pm\gamma_i$. The normalized open vertex is
therefore $P_+(ie\gamma_i)P_-$, and the closed lowest-pair trace is
exactly $-2e^2$ including $(\eta/a_*)^2$. Relative to a Dirac pair
of the same mass, this matches the coefficient of the leading
square-root phase space in this isolated channel. The minimal $v_i$
mark has zero rest matrix element by spatial symmetry. Endpoint
contacts with a single propagator have zero double-pole residue.
This is also the full one-loop leading threshold coefficient for this
specified action. To apply the hard-region argument of
Section~\ref{sec:fourcone-uv}, factor each entire function of the
external frequency outside the loop. The remaining bubble insertions
are the bounded angular matrices $v_i$ or $I$. The differentiated
resolvent bound and the short-duration smoothing argument apply to
each pair of these matrices, leaving the hard region analytic across
the first threshold. The spin trace of a mixed minimal and circulation
pair vanishes. The two minimal vertices first contribute at order
$\beta^3$. This construction specifies the interaction through a
matching rule using the angular mass and eigenfunctions. Derivation
from the physical GCM action and positivity or agreement away from the
first threshold require additional arguments. The differential factor
is a projector at the two chosen frequencies.

\paragraph{Matching the lowest Berry doublet with a bounded circulation}
The two-dimensional pole spaces in \eqref{eq:fourspin-berry-action-uv}
permit a sharper test of \eqref{eq:four-circulation-uv}. Insert
$i c P(v)\mathcal P_B(x)P(v)$ in the same bulk operator for the inverse
and determinant, where $c$ is real. The factor $i$ gives a Hermitian
Euclidean spin coupling for real $B$. Its phase is measurable in the
closed cut. Let $\Gamma_i=\gamma_4\gamma_i$. In the local frame above,
\[
 W^\dagger\Gamma_i W=-i\sin\theta\,
            (n\times e_i)\cdot\boldsymbol\sigma .
\]
Multiplying on the right by the southern spin factor
$i n\cdot\boldsymbol\sigma$ and averaging $n$ gives
$2i\sin\theta\,\sigma_i/3$. The full north-to-south angular
matrix element in the convention
$\langle {+}|\Gamma_i|{-}\rangle$ is $2i b_*\sigma_i$.
The reversed matrix is also $2i b_*\sigma_i$ by anti-Hermiticity.

Put $\mathcal A(q)=\pi L^2\Phi(Lq)/6$ and
$d_*=2E_*\mathcal A(2i E_*)/L>0$. After removing the common
$ie$ the circulation contributes $ic q\mathcal A(q)\Gamma_i/L$.
At incoming photon energy $q=2i E_*$ it has coefficient $-c d_*$,
and the reverse vertex has coefficient $+c d_*$. The resulting
transverse matrices are $i b_*(1-2c d_*)\sigma_i$ and its Hermitian
adjoint. The diagonal charge stays at unity because the circulation
vanishes at zero transfer. The choice
\begin{equation}
 c_*=\frac{1-a/b_*}{2d_*}
 \label{eq:four-berry-circulation-match-uv}
\end{equation}
therefore gives the unit Dirac rest transition after division by $a$.
The closed spin trace is $2a^2\delta_{ij}$ before the two residue
factors; the same two pole spaces saturate the leading two-body cut.
For $m=\kappa=L^{-1}$, $d_*\simeq274.32219$ and
$c_*\simeq-0.00644825$. This is an exact matching of the leading
threshold coefficient. Its input is a physical normalization condition
on a new bulk coupling.

The free angular return is unchanged. At any fixed source order the
projected circulation is bounded by its perimeter estimate, since
$\|P(v)\|=1$ and the intervening free spin transports are unitary.
This proves the same fixed-order Euclidean closed-graph estimates
for the Hermitian insertion. It makes no assumption of unitarity
of transport in a fixed external field. The equalities at threshold
do not determine the cut away from threshold or the signs of the
higher angular channels.

\paragraph{Magnetic moment of the same doublet}
The signed threshold transition and the elastic magnetic moment
probe the same bulk operator at different transfers. To compute the
second quantity, let $H_*=m+\kappa(\nabla^P)^*\nabla^P-E_*v_4$,
$P_N$ be its northern ground projection, and $R$ its inverse on the
orthogonal complement. For static soft fields put
$\pi_i=-i\partial_i+e B_i$. The spatial term in the inverse is
$i v_i\pi_i$. Eliminating the gapped complement gives the leading
terms on $P_N$,
\[
 a(E_*-E)+\sum_{ij} G_{ij}\pi_i\pi_j,\qquad
 G_{ij}=P_Nv_i R v_jP_N .
\]
This is a derivative expansion at zero spatial momentum and to first
order in the magnetic field strength.
The angular vector $v_i f= h(u)n_i$, with
$h(u)=\sqrt{1-u^2}f(u)$, lies entirely in $l=1$. Let $C_1$ and
$C_3$ be the radial inverse quadratic forms of $H_*$ on $h$ for
total spin $1/2$ and $3/2$, respectively. The spin-$1/2$
projector on three spinor-vector components is $\sigma_i\sigma_j/3$,
and the spherical average of $n_i^2$ is $1/3$. These identities give
\[
 \begin{aligned}
 G_{ij}&=d\delta_{ij}+i b_M\epsilon_{ijk}\sigma_k,\\
 d&=(C_1+2C_3)/9,\qquad b_M=(C_1-C_3)/9 .
 \end{aligned}
\]
The $3/2$ potential exceeds the $1/2$ potential by
$6\kappa\sin^2(\theta/2)/\sin^2\theta$. Both radial operators are
strictly positive at the northern crossing. The variational formula
for their inverse quadratic forms, tested at the $3/2$ maximizer,
gives $C_1>C_3>0$ because the potential difference is strictly
positive in the interior. Rotational covariance of the isolated free
pole gives $d/a=1/(2E_*)$ by expanding its dispersion to second order.
Using $[\pi_i,\pi_j]=-ie F_{ij}$, the coefficient of
$e\boldsymbol\sigma\cdot\mathbf B_{\rm mag}$ relative to that of
$\boldsymbol\pi^2$ is the tree gyromagnetic ratio divided by two:
\begin{equation}
  0<\frac{g_{\min}}2=\frac{C_1-C_3}{C_1+2C_3}<1.
  \label{eq:fourspin-soft-uv}
\end{equation}
At $m=\kappa=L^{-1}$, radial Jacobi inversion gives
$C_1/L\simeq0.295714$, $C_3/L\simeq0.130555$ and
$g_{\min}\simeq0.593215$. For reproduction the two $l=1$ radial
bases factor out $\sqrt{1-u}$ and $\sqrt{1-u}(1+u)$; the remaining
polynomials are orthogonal in the corresponding squared weights
times $d\nu_0$. Their dimensionless kinetic eigenvalues are
$3/4+j(j+3)$ and $19/4+j(j+5)$, respectively.

For a local Hermitian insertion
$ie t\sum_{\mu<\nu}P\gamma_\mu\gamma_\nu P F_{\mu\nu}$, the
average of the spatial Clifford rotation in the northern ground state
is $i s_B\epsilon_{ijk}\sigma_k$, where $s_B=(1+2a)/3$.
Indeed conjugation in the two chiral blocks rotates the Pauli vector
through opposite angles $\theta$ about $n$. Their half-sum leaves
$u\sigma_k+(1-u)n_k(n\cdot\boldsymbol\sigma)$. Its angular average
gives the stated factor. Together with the crossed electric matrix
element above this yields the signed magnetic and threshold measurements
\[
 \begin{aligned}
 g/2&=g_{\min}/2-2E_*t s_B/a,\\
 F_{\rm thr}&=(b_*/a)(1-4E_*t).
 \end{aligned}
\]
In canonical units the first measurement requires
$t_{\rm soft}/L\simeq-0.0634288$ for $g=2$; the second requires
$t_{\rm thr}/L\simeq-0.2068234$ for $F_{\rm thr}=1$.
This comparison uses the tree Dirac values specified for matching.

For a circle of radius $r$, the ratio of its threshold to its
zero-transfer curvature coefficient is
$\Phi(2iE_*r)=8 I_2(2E_*r)/(2E_*r)^2$. At $r=L$ it is about
$61.2575$, whereas the required ratio is $3.26072$.
The normalized modified-Bessel series increases strictly from one
to infinity for positive radius. It therefore fixes one radius
$r/L\simeq0.471132$ for this family. Setting the overall circulation
strength to give $t_{\rm soft}$ then matches both measurements while
preserving the bounded perimeter and the unmarked angular return.
The fitted dimensionless radius and strength use two independent
physical conditions. Their values do not follow from the original
velocity-shell measure alone.

\paragraph{A negative higher-channel cut in the minimal Berry determinant}
The sign problem in the minimally gauged Berry operator can also be
localized to identifiable pole spaces. In the same $l=0$ radial
potential $V_0$, let $f$ and $g$ be the positive ground states
at the first zeros for tilts $-M u$ and $+M' u$, respectively.
Here $M=E_*$, and $M'>M$ by the strict reflected comparison above
and a southern cap trial at large energy. The two zero equations
and self-adjointness give $\int ufg\,d\nu_0=0$.
For the radial Wronskian
$W_0=(1-u^2)^{3/2}(f'g-fg')$ those equations give
\[
 W_0'=-(M+M')u\sqrt{1-u^2}\,fg/\kappa .
\]
The regular boundary conditions give $W_0(-1)=W_0(1)=0$. Hence $W_0$ is
strictly positive in the interior.

Take the spatial polarization parallel to a small common line momentum
$p e_i$ at the two-line pinch in the pair rest frame.
The infinitesimal complex rotations to these poles shift the northern
and southern radial arguments by $-ip v_i/M$ and $+ip v_i/M'$,
respectively. In the frame $W(v)$ of the angular bundle the spin
part of the rotation is proportional to
$(n\times e_i)\cdot\boldsymbol\sigma$.
Indeed the projected Clifford rotation and the variation of
the orthonormal frame both have this form, directly from
the spin-frame block $\cos(\theta/2)+i\sin(\theta/2)n\cdot\boldsymbol\sigma$
in the frame formula. Its product with the velocity $v_i$ has
zero angular average. The matrix element is thus
$-ip I\,\mathbf1_2+O(p^2)$, up to the common sign convention, where
\[
 I=\frac13\int(1-u^2)
     \left(\frac{f'g}{M}-\frac{fg'}{M'}\right)d\nu_0>0 .
\]
To see the strict inequality, integration by parts gives
$\int(1-u^2)(fg)'d\nu_0/3=\int ufg\,d\nu_0=0$; the remaining
weighted difference is a positive integral of $W_0$.

Spatial inversion acts by $K\psi(\mathbf v,u)=
\gamma_4\psi(-\mathbf v,u)$. It is even on both $l=0$ doublets and
sends $v_i$ to $-v_i$. In particular $K$ conjugates the angular
pencil at real $\mathbf p$ and imaginary rest frequency to its
adjoint. Continuing its two rest projectors by the small complex
rotations gives positive $K$ norm on both continued subspaces,
since this norm equals one at rest and the rotations preserve $K$.
The reverse velocity matrix in the closed trace is consequently
minus the adjoint of the forward one. The two positive scalar pole
residues multiply a strictly negative spin trace on a set of
relative directions of positive area. This mixed pair contributes
negative weight at order $|p|^3$, including its two-body phase
space. The result concerns this exclusive pair of minimally gauged
levels; other channels can already be open at $M+M'$. A circulation
changes the transition matrix.

The same exclusive test reaches both fitted circulations. For a real
infinitesimal rotation that shifts $u$ by $v_i$ the spin-frame
generator on a radial function $f$ is
\[
    v_i f'+i\frac{b}{\sin\theta}
                  (n\times e_i)\cdot\boldsymbol\sigma f .
\]
The spin coefficient follows by adding the projected ambient
$-\gamma_4\gamma_i/2$ to $W^\dagger dW$ along
$dn=-u(e_i-n_i n)/\sin\theta$. Inserting the projected
$\gamma_4\gamma_i=-i\sin\theta(n\times e_i)\cdot\boldsymbol\sigma$
in the frame, the two spin variations add. Since
$\langle |n\times e_i|^2\rangle=2/3$ and
$\int bfg\,d\nu_0=\int fg\,d\nu_0/2$, the crossed matrix for this
Clifford insertion is $ip(1/M+1/M')\int fg\,d\nu_0/3$ at first
order. Its current coefficient is $-(M+M')t$ at this transfer,
where $t$ includes the circle form factor. Integrating $W_0$ by parts
against $u$ reduces the minimal velocity term to the same expression
times a moment of $u^2$. Cancellation occurs precisely at
\begin{equation}
 t= -\frac{\int u^2 f g\,d\nu_0}{2\kappa\int f g\,d\nu_0},
 \qquad -\frac1{2\kappa}<t<0 .
 \label{eq:fourspin-mixed-zero-uv}
\end{equation}
For $m=\kappa=L^{-1}$ the opposite extremal eigenvalue of the same
northern Jacobi pencil gives $M'L\simeq10.14254$.
The cancellation value is about $-0.09031L$. At the actual transfer
the circle gives $t/L\simeq-21.81$ for the threshold fit at $r=L$
and $-1.1151$ for the simultaneous fit at $r/L\simeq0.471132$.
Both exceed the possible cancellation magnitude $1/(2\kappa)$.
Thus even the circulation calibrated to both light-pair measurements
retains a negative exclusive mixed-pair contribution.

Varying the scalar form factor freely still leaves a second polarization.
Put $Z=M+M'$, $S=\int fg\,d\nu_0$, $J=Z\int u^2fg\,d\nu_0/(\kappa S)$
and $\delta=(M'-M)/Z$. For momentum $p e_j$ and polarization
$e_i$ perpendicular to it, the full current coefficient of
$\epsilon_{ijk}\sigma_k$ to first order is
\[
 -\frac{p S}{6}(M^{-1}+M'^{-1})
             \big[1+Zt(J+\delta)\big].
\]
For example, setting $X_i=(n\times e_i)\cdot\boldsymbol\sigma$,
the angular identities are
$\langle n_iX_j\rangle=\epsilon_{ijk}\sigma_k/3$ and
$\langle X_iX_j\rangle=i\epsilon_{ijk}\sigma_k/3$ for $i\ne j$.
The spin variations give the minimal term and the part proportional
to $\delta$. The derivatives of the radial waves multiplying the
Clifford insertion give the part proportional to $J$ by the same
Wronskian integral. Equation~\eqref{eq:fourspin-mixed-zero-uv}
already requires $Zt=-J/2$. Thus a single form factor can annul
both transitions only if $J(J+\delta)=2$. In canonical units
$J\simeq2.604$, $\delta\simeq0.407$, and the product is about
$7.84$. Even an independently fitted scalar circulation at this higher
transfer leaves an exclusive negative channel in this example.

The strict mismatch can be checked with eight northern Jacobi modes.
In units $L=1$ write $C=H^{-1/2}uH^{-1/2}$, with
$H_{jj}=7/4+j(j+3)$ starting at $j=0$. In its first eight rows the
extreme roots bracket as $(-0.09860,-0.09859)$ and $(0.23384,0.23385)$,
and the other six belong to $(-0.030,0.047)$. These enclosures follow
by the tridiagonal determinant recurrence. The discarded diagonal
block has norm at most $1/H_{88}<0.012$; its single connecting
entry is less than $0.0063$. The extreme normalized eight-mode
vectors have full residuals below $3\times10^{-7}$, found by
multiplying their last components by this connecting entry.
The spectral gaps are consequently greater than $0.06$ and $0.17$
by the min--max bound. Applying $H^{-1/2}$ and normalizing changes
each approximate radial vector by less than $10^{-4}$, since the
two eight-mode norms before normalization exceed $0.47$.
Their scalar overlap lies in $(0.4272,0.4274)$ and their $u^2$
overlap in $(0.0771,0.0773)$, computed by two tridiagonal
multiplications in the full basis. The error in each overlap is
below $0.0002$, as $\|u\|\le1$. These bounds already imply
$J>2.5$ and $\delta>0$.

An independent check uses a connection displaced along an oriented
Clifford transport. For $m=\kappa=L^{-1}$ a two-orientation
mixture on a single chiral carrier can match both the magnetic moment
and the first transverse pair coefficient. Its normalized perpendicular
amplitude has slope between $-17.55$ and $-16.99$ in
$x=|\mathbf p|^2/m_*^2$ at threshold. The point Dirac value is $1/2$.
The detailed angular calculation is given in the companion monograph.
Thus even a pair of on-shell matching conditions leaves the continuum
shape independently measurable.

\subsection{A global Euclidean test after threshold matching}
\label{sec:four-angular-superconvergence-uv}
There is a common test of the continuously turning closed actions. It
uses the full determinant response, including coincident Wilson
derivatives, and is independent of the values matched at threshold.
Take $q=\omega e_4$, $X=\omega^2$. The spatial photon preparation
equals unity here. In the rest-current convention of the spacelike
dispersion test, write $F(X)=\Pi_{ii}(\omega e_4)/(e^2X)$ for one
fixed transverse direction. For the minimal angular action, its Berry
version, and either of them with the bounded circulation marks above
that vanish at large real $|\omega|$,
\begin{equation}
             \lim_{X\to\infty} X F(X)=0.
 \label{eq:four-angular-superconvergence-uv}
\end{equation}
To prove this statement, condition on a closed path of duration $t$
at a base and velocity point $a$ with $p_t(a,a)>0$. For $0<h<t$
the bridge density at $z$ is
$p_h(a,z)p_{t-h}(z,a)/p_t(a,a)$. It is absolutely continuous.
Since the equator $v_4=0$ has zero spherical measure, Fubini gives
zero time on this equator for almost every such bridge. Now
$x_4(h)$ is continuously differentiable with derivative $v_4(h)$.
The open set where this derivative is nonzero is a countable union
of intervals of monotonicity. On each, ordinary change of variable
pushes $v_i(h)dh$ forward to an absolutely continuous signed measure
in $x_4$. Their summed total variations are bounded by $t$. This
remains true however often the path turns in clock time. Hence
\[
 J_i(\omega)=\int_0^t v_i(h)e^{i\omega x_4(h)}dh
                      \longrightarrow 0
\]
by the Riemann--Lebesgue lemma. For a minimal Wilson weight its two
derivatives give $J_i(\omega)J_i(-\omega)$ times the free spin
holonomy. The modulus is bounded by the spin multiplicity times
$t^2$. Any additional term with a decaying circulation mark tends
to zero uniformly on each fixed-duration path, with the same form
of duration bound. This includes the bounded derivative of the
averaged circle; a point curvature coupling growing with frequency
is outside the argument. Integrability of $t^2$ against the
mass-weighted diagonal proper-time trace follows from
\eqref{eq:next3-angular-return-uv}. Dominated convergence proves
\eqref{eq:four-angular-superconvergence-uv}.

Suppose this conserved-current scalar has a positive spectral
measure with weight $M>0$ on a finite band $[s_0,s_1]$ and admits
a dispersion relation with finitely many local subtractions. Choose
$N$ large enough to remove its subtraction polynomial and make its
subtracted kernel integrable. At the nodes $X_j=(j+1)R$,
$0\le j\le N$, the divided difference satisfies
\[
 (-1)^N[X_0,\ldots,X_N]F
 \ \ge\ \frac{M}{\prod_{j=0}^N(s_1+X_j)} .
\]
Indeed the divided difference of the subtracted spectral kernel is
$(-1)^N/\prod_j(s+X_j)$, so every spectral channel contributes with
the same sign. But the interpolation formula expresses the left
side as $R^{-N}$ times a fixed linear combination of $F((j+1)R)$.
Equation~\eqref{eq:four-angular-superconvergence-uv} makes it
$o(R^{-N-1})$, contradicting the positive band. This excludes a
tempered positive conserved-current reconstruction with the stated
dispersion and the same Euclidean Hessian, even when upper-channel
weights or the subtraction degree are unknown. It leaves the free
charged propagator test intact and identifies the missing physical
response beyond any fitted first threshold. A modified source must
supply a high-frequency contribution that survives this test.

One can see the conflict without assuming rotation invariance of the
physical current. At zero spatial momentum $\Pi_{ii}(\omega e_4)$ is
the Fourier transform of a finite temporal measure $\Lambda$, possibly
signed or complex. The two Wilson marks push forward measures of
total variation at most $t^2$. Each undifferentiated circle mark is
itself a finite measure by its perimeter bound. A differentiated
circle adds a multiplier proportional to $\omega^2\mathcal A(\omega)$,
hence to $J_2(L\omega)$; this too is the Fourier transform of a finite
measure since
\[
 J_2(z)=-\int_{-1}^1 e^{izu}\,
             \frac{2u^2-1}{\pi\sqrt{1-u^2}}\,du .
\]
Translations and ordered multiplication of such marks preserve the
stated variation bound up to fixed constants. Every mark vanishes at
zero transfer on a closed path, so $\Lambda(\mathbb R)=0$.
Averaging its Fourier transform over $[-R,R]$ and taking $R\to\infty$
gives $\Lambda(\{0\})=0$ by \eqref{eq:four-angular-superconvergence-uv}
and dominated convergence applied to $\sin(R\tau)/(R\tau)$.
In a positive stationary representation the separated-time autocorrelation
of one Hermitian current is a Laplace transform of its nonnegative
spectral measure in $|\tau|$, up to the common Euclidean current sign.
A nonzero positive band gives a strict sign for all $|\tau|>0$.
This is incompatible with the two measure identities even when a
different contact is allowed at $\tau=0$. The test concerns the
complete current measured by the source Hessian.

\paragraph{An instantaneous spatial displacement supplies the missing atom}
Consider an oriented rectangular Wilson loop whose two spatial sides
have length $\ell$, directions $n,-n$, and clock separation $\tau$.
At zero spatial momentum its transverse mark, up to a translation
phase and choice of orientation, is $\ell n_i(e^{i\omega\tau}-1)$.
Take $n$ uniform on $S^2$, and $\tau$ independently exponential with
rate $E>0$. The expected squared mark is
\[
 \frac{2\ell^2}{3}\int_0^\infty
     E e^{-E\tau}(1-\cos\omega\tau)\,d\tau
       =\frac{2\ell^2}{3}\frac{X}{X+E^2}.
\]
This is the quadratic response of an averaged gauge-invariant term
$1-\cos(e\oint B\cdot dy)$ after removing $e^2$. Both source
derivatives act on the same loop. Its temporal measure, in units of
$2\ell^2/3$, is $\delta_0-(E/2)e^{-E|\tau|}d\tau$.
Thus a spatial displacement made at one clock time provides precisely
the contact required by a positive massive band consisting of one
atom. The rectangle perimeter has every finite moment under this
law. This is a rest-frame source comparison; deriving such a contact
from the GCM boundary and reproducing the existing charged transition
are separate requirements.

\subsection{A closed geometric excursion with an exact spectral response}
\label{sec:four-excursion-uv}
Can a charged history supply the missing high-frequency term while its
short closed returns stay suppressed? Merely adding ordinary
translational diffusion to its streaming coordinate loses the original
vacuum estimate. For example, add an independent four-dimensional
diffusion of strength $D>0$ per unit angular duration $t$. Conditional
on the angular path, the added four-displacement has Gaussian
density at the required return of at least
$(4\pi Dt)^{-2}\exp[-t/(4D)]$, since the original displacement has
length at most $t$. The scalar angular heat trace grows as a positive
constant times $t^{-3/2}$ at zero. The unmarked proper-time trace
with measure $dt/t$ now diverges. Any independent positive law on $D$
giving positive weight to a compact subinterval of $(0,\infty)$
has the same obstruction.

A closed excursion can change the current without changing the
displacement of the angular center. We give an exactly soluble
construction on the spin-spectator branch of
\eqref{eq:next3-angular-flight-uv}. At the base point $x$ draw an
independent pinned four-dimensional Brownian loop $b$ from zero to
zero of heat duration $T$ (generator $\Delta$, so each free component
has variance $2T$). Its duration has the gamma law of shape $1/2$
and scale $\theta>0$, of dimension length squared. Equivalently it
is the square of a real centered Gaussian of variance $\theta/2$.
Define the gauge-invariant multiplication operator
\begin{align*}
 w_B(x)&=\mathbb E\exp\!\left(ie\oint B(x+b)\circ db\right),\\
 D_B^{\rm exc}&=D_B+\lambda(1-w_B)I .
\end{align*}
The circle on the stochastic differential denotes Stratonovich
integration, so a gradient has zero circulation. Orientation reversal
makes $w_B$ real for real $B$, gives $w_0=1$, and kills its first
variation. The signed coefficient $\lambda$ has inverse-length units.
This is an explicit new curvature coupling in the charged inverse and
determinant. At zero field the unmarked heat trace is exactly the
original one.

\paragraph{The contact computed from its paths}
For $Q=\omega e_4$, $X=\omega^2$, let
$J_i=\oint e^{i\omega b_4}\circ db_i$, $i\ne4$.
Given the fourth component of the bridge, its transverse component
is an independent pinned Gaussian. Conditional variance gives
$2T-2|\int_0^T e^{i\omega b_4(u)}du|^2/T$; the projected increment
of the fourth bridge has variance $2|u-v|(1-|u-v|/T)$.
The same pinned increment covariance occurs in worldline representations
of one-loop responses~\cite{ExcursionSchubert}. Here the heat duration
is drawn anew at each inserted excursion. The width at fixed $T$ is
\[
 \mathbb E(|J_i|^2\mid T)=
 2T\left[1-\int_0^1 e^{-TX y(1-y)}dy\right].
\]
In passing from the double time integral to $dy$ one uses symmetry
under $y\mapsto1-y$. Weighting the gamma law by $T/\mathbb ET$
raises its shape to $3/2$. The identity
$\int_0^1[1+a y(1-y)]^{-3/2}dy=(1+a/4)^{-1}$ follows by setting
$2y-1=z$. With $s=4/\theta$ the exact result is
\begin{equation}
 K_s(X)=\mathbb E|J_i|^2=\frac4s\,\frac{X}{X+s}.
 \label{eq:four-excursion-response-uv}
\end{equation}
Rotational invariance gives the same transverse kernel at every
real Euclidean four-momentum before any photon readout. The large
frequency limit is $4/s$. The constant part and the exponential
separated-time part are fixed by the same pinned-path measure.

The sign in the determinant matters. Let $k_0>0$ denote the internal
spin and angular trace of the free inverse at coincident base point,
per unit Euclidean volume. It is finite by the angular return bound;
its scalar heat kernel has positive weight on closed turns.
For $\Gamma=-\operatorname{Tr}\log D_B^{\rm exc}$ the added
quadratic transverse Hessian is $-\lambda e^2 k_0 K_s(X)$. With
$c=-4\lambda k_0/s$ the corresponding variation of a unit Maxwell
covariance at order $e^2$, before photon preparation, is
\[
 \delta D_T(X)=\frac{e^2 c}{s}
       \left[\frac{1}{X+s}-\frac{1}{X}\right].
\]
A positive massive residue requires $\lambda<0$. The positive Poisson
rate that inserts independent Wilson excursions in the original
Grassmann heat trace gives the opposite sign. The negative coefficient
is a bounded curvature potential for real probes and is a separate
choice of action. For example $2|\lambda|<m$ suffices for a uniform
positive real-part lower bound on real backgrounds. Each fixed-order
calculation at zero field already uses the original free generator.

There is also a positive-rate realization of the required sign in a
different charged representation. Enlarge the path record by a
two-element group, let each excursion flip its element, and integrate
only the odd character in the charged fiber. A jump then transports
$-1$ times the Wilson phase. Its generator adds
$\lambda_+(1+w_B)$ with $\lambda_+>0$. At zero field the free mass
becomes $m+2\lambda_+$, and the Hessian has the positive massive
sign just derived, with $k_0$ evaluated at that shifted mass.
The curvature potential is nonnegative on real backgrounds. This
introduces a definite discrete holonomy in the action. The shifted
charged poles and all of their transitions are to be compared at
the new mass. The original shell measure alone does not specify the
odd character or this jump rate.

A distribution of $s$ can realize an entire positive finite band
in this Hessian. For a desired scalar measure $d\rho(s)$ supported
above $s_0>0$, take the signed coupling measure
$d\lambda(s)=-s\,d\rho(s)/(4k_0)$. The total coupling magnitude
is finite when $\int s\,d\rho(s)<\infty$. The added
$\delta\Gamma_T/(e^2X)$ is then $\int d\rho(s)/(X+s)$.
This inverse prescription specifies a spectral law as input.
In particular it cannot derive that law from the original GTL and
shell constraints. Adding the entire desired spectrum to a determinant
with a nonzero prior cut would count an additional response. Matching
requires the difference between the measured spectrum and the actual
prior determinant, including its sign.

\paragraph{The rate and the charged mass in the odd representation}
Prescribing a massive band in the positive-rate, odd-character
representation fixes a nonlinear normalization. Let $m>0$ be its
streaming mass before the flips, $\ell=\int d\lambda_+(s)$ their
total rate and $\mu=m+2\ell$ the shifted streaming mass. Write
$h(t)$ for the spin-spectator diagonal return trace per unit base
volume before the mass term, so that
$k_0(\mu)=\int_0^\infty e^{-\mu t}h(t)dt$. Take an externally
prescribed positive band $d\rho$ on $s\ge s_0>0$ and put
$M_1=\int s\,d\rho(s)<\infty$. Matching the additional scalar
Hessian of this particular excursion gives
\[
 d\lambda_+(s)=\frac{s\,d\rho(s)}{4k_0(\mu)},\qquad
 M_1=2(\mu-m)k_0(\mu).
\]
Thus the attainable first moments form a finite interval. Its upper
endpoint is the maximum of the right-hand side over $\mu\ge m$.
The function is continuous, strictly positive for $\mu>m$, and
vanishes at both ends. To check the limit at infinity, put
$z=(\mu-m)t$; the integrand is bounded by
$2e^{-mt}h(t)/(\mathrm e t)$ since $ze^{-z}\le1/\mathrm e$.
This majorant is integrable by the previously obtained return bound
and the positive mass. Pointwise convergence also proves the limit.
In particular every target satisfies
\[
 M_1\le\frac{2}{\mathrm e}\int_0^\infty e^{-mt}h(t)\,\frac{dt}{t}.
\]
Every positive moment below the maximum has at least two solutions,
one on either side of a maximizing value; the derivative at $m$ is
$2k_0(m)>0$, giving a unique nearby small-rate branch. Zero moment
for a positive gapped target forces zero rate at finite mass.
These statements hold for a fixed target and the additional two-source
contact. A target computed from the shifted charged poles must itself
be evaluated at the same $\mu$.

\paragraph{Sewing this contact and opening its source}
For a real transverse unit polarization $\epsilon$, the two bridge
components along $Q$ and $\epsilon$ are independent. The same conditional
Gaussian argument and bridge moments give, for every finite $p\ge2$,
\[
 \|J_\epsilon(Q)\|_{L^p(\cdot\mid T)}
       \le C_p\min(\sqrt T,|Q|T).
\]
For the second bound subtract one from $e^{iQ\cdot b}$ and use
$|e^{iQ\cdot b}-1|\le |Q|\sup|b|$. Longitudinal marks are
closed gradients. H\"older's inequality bounds mixed insertions even
when their momenta differ and they occur on the same excursion.
All fixed moments of $T$ exist, uniformly in the above mixture for
$s\ge s_0$. Thus in fixed carrier units the per-end real-momentum
bound is $C_p\min(1,|Q|)$. Insertions in the angular heat trace
cost only powers of its total duration at each fixed electromagnetic
order. The original short-return estimate and the mass at long
duration still control their joint proper-time integrals.

With the spatial A4 readout at each photon end the sewn envelope is
bounded and integrable in the four-momentum, exactly as for the
minimal closed graphs. This proves fixed-order absolute convergence
for the source-completed closed angular graphs with these new contacts.
The temporal-axis kernel above is unchanged. At nonzero spatial
momentum the two readout factors also appear, so invariance of the
prepared photon spectrum does not follow. With an unsmeared Maxwell
line the single contact already costs an integral proportional to
$\int d^4Q\,K_s(Q^2)/Q^2$, quadratically divergent for a nonzero
single coupling. The angular return factor here is a finite constant
independent of the photon circulating through that contact.

Finally, $w'_0=0$. The added two-current Hessian has a pole even
though the added one-current transition between charged states is
identically zero. This difference is observable by opening the
charged line. Closing the unrestricted theorem would require
physical intermediate states producing the same residue, together
with their vertices and a photon sewing that bounds their hard
momenta in every clock frame.

The same excursion has a measurable forward two-current response
on an open charged line. Take two plane-wave probes with momenta
$Q,-Q$ and the same real transverse unit polarization. Compare at
the same free streaming mass $\mu$, before spatial photon preparation.
Write $\lambda_{\rm eff}=\lambda$ for the original signed insertion
and $\lambda_{\rm eff}=-\lambda_+$ in the odd representation.
Translation cancels the two plane-wave factors at the base point.
On the full angular and spin space the extra second derivative of
the charged operator is consequently
$\lambda_{\rm eff}e^2 K_s(X)I$. Its first derivative vanishes.
Let $G(p;\mu)$ be the free charged inverse or its compression by
fixed external maps. Differentiating the inverse at both sources
and subtracting the comparison at the same $\mu$ gives exactly
\[
 \Delta G_{Q,-Q}(p)
   =\lambda_{\rm eff}e^2 K_s(X)\,
     \partial_\mu G(p;\mu).
\]
Indeed $\partial_\mu D_0=I$ and
$\partial_\mu D_0^{-1}=-D_0^{-2}$. The mass derivative holds
the endpoint maps fixed, including their chosen normalization.
If the readout itself depends on the probe, its derivatives common
to both comparisons cancel in this difference. At the simple
angular pole of mass $M$, let $d=\langle f,v_4 f\rangle>0$ in
its normalized northern ground state. Differentiating its zero
eigenvalue gives $dM/d\mu=1/d$. The extra double pole in the
invariant denominator $(p^2+M^2)^2$ therefore has coefficient
$-2M\lambda_{\rm eff} e^2K_s(X)/d$ times the original residue.
The closed contact and this open two-current coefficient measure
the same insertion. The additional on-shell one-current matrix
element remains zero.

\paragraph{A lower bound on the bare vacuum graph for a gapped conserved current}
The excursion also illustrates a more general consequence for ordinary
physical photon sewing. Suppose the measured current is covariant on
its plane, has a gapped, nonzero positive spectral band, and its
transverse scalar $F$ admits finitely many dispersion subtractions.
Allow arbitrary local gauge-invariant quadratic contact polynomials.
Their transverse tensors have the form
$(Q^2\delta_{\mu\nu}-Q_\mu Q_\nu)P(Q^2)$ and hence add a polynomial
to $F$. A massless pole of the current itself is excluded here.
Put $K(X)=X F(X)$. Taking $N$ above the subtraction degree, the
band lower bound on the divided difference in
Section~\ref{sec:four-angular-superconvergence-uv} gives
$\sum_{j=1}^{N+1}|F(jR)|\ge c_N/R$ for all sufficiently large $R$.
Multiply this inequality by $R$ and integrate from $Y$ to $2Y$.
Changing variables $X=jR$ gives
\[
 \int_Y^{2(N+1)Y}|K(X)|\,dX\ \ge\ c'_N Y .
\]
The four-dimensional radial measure is proportional to $X\,dX$.
Thus sewing this very Hessian with $D_T(X)\ge Z/X$, $Z>0$, costs
at least order $\Lambda^2$ in the absolute integral through Euclidean
radius $\Lambda$. The growing integration domain diagnoses convergence;
it is absent from the action. This strengthens
Proposition~\ref{prop:carrier-sewn-uv} under the stated absence of a
massless current pole and the transverse local-contact hypothesis.
With arbitrary polynomial additions directly to $K$, that proposition
still gives its logarithmic obstruction. The distinction matters: a
constant in $K$ is a $1/X$ term in $F$. A finite band already suffices,
whatever the ultraviolet tail of the positive spectral measure.
A proposal to evade this estimate through correlated normal records
must compute the actual joint contraction and also explain why closing
the observed physical photon does not reproduce the factorized graph.

\subsection{Can the electric and transverse currents share the boost cost?}
\label{sec:four-boost-uv}
The boost obstruction in \eqref{eq:next3-dipole-energy-uv} used a
purely electric polarization. A more general linear source can excite the
same pair through both the electric momentum and the transverse vector
potential. Their relative sign matters. Consider a free Maxwell oscillator
of spatial momentum $k$, $r=|k|>0$, and a pair of energy
$w=\sqrt{r^2+s}$. For one transverse polarization write $j$ for the
prescribed charge-one current amplitude on that final-state fiber.
Include the full spin and any geometric records in the Hilbert norm on
the fiber. If $P$ is the electric polarization amplitude and $J_A$ the
additional direct current, the measured amplitude is
\[
 j=iwP+J_A .
\]
Indeed the source action contains $-e\dot A_TP+e A_T J_A$ and the
external-current variation differentiates $P$ in time. Keep the canonical
contact generated by $(\Pi_T+eP_T)^2/2$. A direct transverse coupling
$-e A_T J_A$ has no quadratic contact in the comparison below. Write
$x=iP$, so $J_A=j-wx$. Creation of a photon from the vacuum supplies
$i\sqrt{r/2}$ in $\Pi_T$ and $1/\sqrt{2r}$ in $A_T$. The sum of the
absolute contact and virtual intermediate-state contributions to the
bare vacuum coefficient is therefore
\begin{equation}
 \frac{e^2}{2}\left(\|x\|^2+
      \frac{\|j-(w+r)x\|^2}{r(w+r)}\right)
 \ \ge\ \frac{e^2\|j\|^2}{2(w+r)(w+2r)} .
 \label{eq:four-boost-floor-uv}
\end{equation}
This is a positive Hilbert-space square completion with equality at
$x=j/(w+2r)$. It holds even when the electric and direct channels have
unobserved components. On an invariant positive mass band, each
transverse Dirac current has fiber norm squared $s\rho(s)ds/(2w)$ in
the preceding convention. The lower bound then costs a positive constant
times $dr/r$ at large $r$ after the center integration. This divergence
persists if the longitudinal energy is removed altogether. Thus changing
the split between two linear current couplings cannot make these two
individual transverse graphs absolutely finite with an ordinary Maxwell
oscillator and the specified positive band.

\paragraph{Which part can be repaired at the complete coefficient}
In the signed vacuum coefficient the second term of the left side of
\eqref{eq:four-boost-floor-uv} enters with a minus sign. One can solve
for an exact cancellation at each momentum while retaining the original
measured amplitude. For example take
\[
 iwP=\bigl(1-\sqrt{r/(w+r)}\bigr)j,\qquad
 J_A=\sqrt{r/(w+r)}\,j .
\]
Substitution gives $\|x\|^2=\|j-(w+r)x\|^2/[r(w+r)]$. These are
specified transverse matrix elements of the time-local Hamiltonian on
finite tests. The dependence on $w$ is evaluated on its constituent
pair, with the Hermitian adjoint used for annihilation. The direct
current couples only to $A_T$. Longitudinally retain $iwP_L=j_L$ for
the desired conserved pair charge. Choose the gauge-invariant electric
kinetic form $\|E_T\|^2/2+\langle E_L,g^{-1}E_L\rangle/2$, with
ordinary Maxwell magnetic energy and
\begin{equation}
 g(r)=4\int_0^1 u^3\,|\chi_{248}(Lru)|^2du .
 \label{eq:four-boost-long-uv}
\end{equation}
This is the positive squared readout averaged over the earlier shell
radius with distribution function $u^4$. It is strictly positive at
every finite $r$ and $g(0)=1$. Changing variables to $v=Lru$ gives
$g(r)\sim 4(Lr)^{-4}\int_0^\infty v^3|\chi_{248}(v)|^2dv$.
The conjugate longitudinal displacement now gives the positive Gauss
energy with kernel $g(r)/r^2$. Consequently the complete order-$e^2$
coefficient per four-volume in this comparison is simply
\[
 \frac{e^2}{4}\int\frac{d^3k}{(2\pi)^3}\,g(r)
                  \int_{4m^2}^\infty \frac{\rho(s)}{w}\,ds <\infty .
\]
This calculation combines the transverse integrands before integration.
It preserves the transverse Maxwell dispersion, its free residue and
the invariant measured Born pair weight. It adds no divergent constant.
The noncovariance has moved into a test on the field of an electric
charge: its static exchange is $e^2g(r)/r^2$. The unit massless pole
alone in a covariant positive spin-one exchange between conserved probes
contributes $e^2/r^2$ to the static exchange. Additional positive
massive weights increase it, whereas $g(r)<1$ for $r>0$. This excludes
that interpretation of the modified static field with the same physical
probe and no extra contact. Both the electric splitting and the spatial
stiffness are representation choices. Their finite combined vacuum
coefficient is compatible with the lower bound on its separate terms;
neither calculation establishes a full interacting Hamiltonian from
this unsmeared continuum source.

\subsection{A transverse photon converted into a charged pair}
\label{sec:four-excite-uv}
There is a way to realize the invariant optical band in the actual
charged Fock space without its vacuum boost integral. It changes the
boundary Hamiltonian. Take the two transverse Maxwell polarizations
and the positive-energy particle and antiparticle spaces of the Dirac
field. Let $Q=N_\gamma+(N_++N_-)/2$. A conversion of one photon into
one particle and one antiparticle preserves both $Q$ and electric
charge. The translation-invariant conversion acts at fixed total
spatial momentum $k$, $r=|k|>0$. Resolve the opposite-charge pair by
invariant mass squared $s\ge4m^2$, with energy $w=\sqrt{s+r^2}$.
For each transverse polarization use the signed spin and angular row
$u(s,k)$ of the marked Dirac pair, with its two-body Jacobian absorbed
in the fiber norm. Its squared norm at fixed $s$ is
$\rho(s)=C_{\rm pair}(s)$ from \eqref{eq:nextsoft-positive-uv}.
The two transverse rows are orthogonal after the angular and spin
sum. With these conventions the measured current row is
$\sqrt{s/(2w)}\,u(s,k)$.

Write $M_{-1}=\int\rho(s)ds/s$ and suppose
$Z=1-e^2M_{-1}>0$. In each polarization define
\begin{align}
 q_r(s)&=e\sqrt{r/(ws)}\,u(s,k),&
 d_r&=Z+\|q_r\|^2, \nonumber\\
 v_r&=q_r\,\frac{\arctan(\|q_r\|/\sqrt Z)}{\|q_r\|}.
 \label{eq:four-excite-angle-uv}
\end{align}
The quotient has its continuous value at zero. The skew rotation with
off-diagonal row $v_r$ sends the photon to
$(\sqrt Z,q_r)/\sqrt{d_r}$ in the photon-plus-pair space.
On the full Fock space use the skew generator $J_e$ whose forward
term annihilates one photon of momentum $k$ and creates its constituent
fermions with row $v_r$, integrated over $k$ and summed over the two
polarizations. Subtract its adjoint. Before (anti)symmetrization the
map on one incident photon is the direct integral of bounded maps
$f(k)\mapsto v_r f(k)$. In a fixed $Q$ sector the choices of photon,
particle and antiparticle give a finite combinatorial norm factor.
Thus $U_e=\exp J_e$ is unitary sectorwise, preserves momentum and
charge, and fixes the vacuum and single charged particles.

Use the positive geometric quartet $V_4$ of
\eqref{eq:next3-car-quartet-uv}, with the normalization there, and
let $H_0$ be the sum of free one-particle energies. Take
\[
 H_e=U_e^\dagger(H_0+e^2V_4)U_e,\qquad
 A_{\rm obs}(k)=\sqrt{d_r/(2r)}\,[a(k)+a^\dagger(-k)]
\]
on transverse tests. The positive self-adjoint direct sum obeys
$H_e\ge|\mathbf P|$ because $H_0\ge|\mathbf P|$ and $U_e$
commutes with translations. Its vacuum expectation satisfies time
reflection positivity by the Hilbert norm of the time-ordered
positive-time vectors, using adjunction and order reversal on reflection.
The quartet vanishes on $Q=1$. The exact transverse two-point
function at fixed spatial momentum is therefore
\begin{equation}
 C(\tau,r)=\frac{Z}{2r}e^{-r|\tau|}
 +e^2\int\frac{\rho(s)}{s}\frac{e^{-w|\tau|}}{2w}\,ds .
 \label{eq:four-excite-optical-uv}
\end{equation}
Fourier transformation in $\tau$ gives precisely
\eqref{eq:joint2-inverse-spectrum-uv} with $X=k_0^2+r^2$ and $g=e$.
The exact continuum state in this representation is $U_e^\dagger$
applied to the free pair. The transverse Maxwell wave operator on its matrix
element of $A_{\rm obs}$ supplies $s$ up to the overall current sign.
Dividing by $e$ recovers the measured signed current, including its
spin row, in every momentum fiber. The scalar spectral law on transverse
tests depends only on $X$. The massless residue remains $Z$
at every allowed coupling. The factor $d_r$ is determined by the
equal-time sum: $2rC(0,r)=Z+e^2\int (r/w)\rho(s)ds/s$.
It increases strictly from $Z$ to one for a nonzero massive band.
A bare electromagnetic coordinate creating a unit photon from this
invariant vacuum has constant $2rC(0,r)=1$; the measured coordinate
must have the stated normalization in this representation.

Current derivatives also distinguish a time-ordered response from its
continuum cut. For $R(X)=\int\rho(s)ds/(X+s)$ and
$M_0=\int\rho(s)ds$, the same covariance obeys
\begin{align*}
 XD_e(X)&=1-e^2R(X),\\
 X^2D_e(X)&=X-e^2 M_0+e^2\int\frac{s\rho(s)}{X+s}\,ds .
\end{align*}
The polynomial terms are the contact distributions of the differentiated
two-point function. Their coefficients are fixed and finite. Removing
them when identifying the positive continuum discontinuity leaves
their contribution to the full time-ordered response intact.

The cutoff-free perturbation estimate includes open charged endpoints.
Geometric pair suppression gives $\int\rho(s)ds<\infty$.
Uniformly in $r$, the squared conversion norm is bounded using
$r/(ws)\le1/s$. Commutation with $H_0$ multiplies its row by $w-r$,
and $(w-r)^2r/(ws)\le1$. The same two bounds hold for the absolute
momentum kernels. The scalar functions multiplying $q_r$ in
\eqref{eq:four-excite-angle-uv} have convergent expansions near
$e=0$ uniform in $r$, since $\int(r/w)\rho(s)ds/s\le M_{-1}$.
In expanding the unitary conjugation at fixed order, each commutator
with the free energy therefore supplies one bounded absolute conversion
row. Each remaining conversion and each quartet transfer also has a
bounded absolute kernel on a fixed sector, by the Gaussian measure
bound following \eqref{eq:next3-car-quartet-uv}. Spectator permutations
have finite multiplicity and free Euclidean propagation is contractive.
Composing these bounds and integrating the ordered times over a finite
interval proves absolute momentum convergence of each such term with
smooth finite-particle packets at fixed order. Every interaction term
annihilates the vacuum. No subtraction of vacuum diagrams is required
in this Hamiltonian expansion.

This construction implements a dynamical photon source, Pauli exchange,
and a nonzero interaction among four charged constituents. The
Hamiltonian preserves $Q$, whereas the measured field changes $Q$ by
one. Its current therefore has zero elastic matrix element between
single-charge packets, even for an off-shell transverse probe and a
moving charge. The quartet uses a chosen time slice. The invariant
scalar transverse law and the forward energy cone consequently settle
neither covariant higher-point functions nor minimal coupling of a
single charge to its Gauss field. Deriving those vertices and this
boundary conversion from the same physical GCM action remains necessary
for the unrestricted theorem.

\paragraph{An interacting representation of the boosts}
\label{par:four-excite-mass-uv}
The quartet also admits a realization in the invariant mass. In each
free constituent sector containing two positive and two negative charges,
use total spatial momentum $\mathbf P$ and rest-frame relative momenta,
absorbing the change of density unitarily in the internal wave functions.
Sum the quartet potential of \eqref{eq:next3-car-quartet-uv} in the
coordinates conjugate to the internal rest momenta. The common position
drops out of its pair differences. The resulting $V_{\rm rest}$ is
positive, independent of $\mathbf P$, rotationally scalar, and bounded
in each fixed sector. Its Fourier transfers preserve zero total rest
momentum and have bounded absolute mass. Set it to zero in sectors
without a quartet. In the affected sectors $M_0=\sum_i E_i^{\rm rest}
\ge4m$. Use
\[
 M_e=M_0+e^2V_{\rm rest},\qquad
 H_{\rm ref}=\sqrt{|\mathbf P|^2+M_e^2}.
\]
Since the internal rotation generators commute with $M_e$, its spectral
and spin decomposition specifies a direct integral of positive-energy
massive Wigner representations, with the usual $\mathbf P$ translations
and the displayed energy. Use the free representations on the remaining
sectors, and conjugate the entire representation by $U_e$. This
constructs boosts satisfying the group law and leaves
\eqref{eq:four-excite-optical-uv} unchanged.

For completeness let $f_P(x)=\sqrt{P^2+x^2}$, $P=|\mathbf P|$.
On $|z-x|\le x/4$, $x\ge4m$, one has
$\operatorname{Re}(P^2+z^2)\ge P^2+x^2/2$ and
$|f'_P(z)|\le 2$. Cauchy's inequality and the simplex formula for
divided differences give, for $n\ge1$ and all $x_i\ge4m$,
\[
 |f_P[x_0,\ldots,x_n]|
 \le \frac{2}{n}\left(\frac{4}{4m}\right)^{n-1}.
\]
In the free mass representation the $n$th coefficient in the mass
perturbation consists of this divided difference times the $n$ successive
$V_{\rm rest}$ kernels, integrated over their intermediate variables.
One obtains the formula first for step-function masses of bounded
range by the resolvent expansion, then by spectral approximation on
compactly supported packets and the displayed uniform absolute bound.
The coefficient is dominated in operator norm by the bound times
$\||V_{\rm rest}|\|^n$, where the bars inside this norm indicate the
absolute transfer kernel. Indeed the absolute transfer product on $|\psi|$ gives an $L^2$
majorant for each coefficient on a packet $\psi$. Truncating the free
masses from above therefore gives strong coefficient limits by dominated
convergence. The uniform geometric bounds sum near $e=0$ in each sector,
and strong resolvent convergence identifies their sum with $H_{\rm ref}-H_0$. These spectral
approximations prove the identity; they are absent from the action and
the resulting coefficients. The previous absolute convergence proof
therefore extends to this choice of interacting boosts. A representation
of spacetime symmetry does not itself give cluster separability for
widely separated subsystems or a locally transforming electromagnetic
current. Those remain additional tests on this construction~\cite{CarrierClusters}.

\subsection{Testing the excursion in a magnetic background}
\label{sec:fifth-mag-uv}
A pole measured in the vacuum can be tested by switching on a
static magnetic field and repeating the same two current variations.
This leaves the physical clock translation symmetry available for a
positive spectral test. Use the same sector plane and Euclidean
continuation as in Section~\ref{sec:four-excursion-uv}. Choose the
background in the $12$ plane, put $b=|eF_{12}|>0$, and probe the
third component with transfer $Q=\omega e_4$, $X=\omega^2$.
A normalized spatial readout preserves a constant field and is unity
at this probe's spatial momentum. The calculation first concerns the
excursion in the charged determinant. The contribution from the rest
of that determinant is identified below.

\paragraph{A positive vacuum pole becomes a double pole}
The planar Brownian area average on a pinned loop of heat duration
$T$ is $bT/\sinh(bT)$. Indeed the covariant scalar Laplacian in this
plane has levels $(2n+1)b$ of density $b/(2\pi)$ per unit area.
Dividing its diagonal heat kernel by $(4\pi T)^{-1}$ gives this
ratio~\cite{ExcursionSchubert}. The coordinates $3,4$ of the same
pinned loop are independent of its first two coordinates. Therefore
the factor multiplies the earlier conditional width at fixed $T$
before averaging the heat duration. With $\mu=s/4$ and
$S_n=s+4(2n+1)b$, the exact contact is
\begin{align}
 \frac{K_{s,b}(X)}{X}
  &=3b\sqrt\mu\sum_{n\ge0}(S_n/4)^{-5/2}\nonumber\\
  &\quad\times\left[\frac1{X+S_n}
           +\frac{2S_n}{3(X+S_n)^2}\right].
 \label{eq:fifth-magnetic-contact-uv}
\end{align}
To obtain it expand $1/\sinh(bT)$ into its positive exponential
series. Weighting the original gamma density by $T^2$ gives shape
$5/2$. The remaining chord integral is
\[
 \int_0^1[1+4z y(1-y)]^{-5/2}dy
       =\frac{1+z/3}{(1+z)^2},\qquad z=X/S_n .
\]
It also follows by differentiating the shape-$3/2$ identity with
respect to the mass parameter. The positive expansion for real $X$
and local uniform convergence away from the poles justify the
meromorphic continuation of the displayed response.

A positive stationary current at fixed spatial momentum can have a
simple pole from a discrete state. A locally finite positive spectral
measure cannot produce this double pole: near an isolated spectral atom
its singular term is its mass divided by $X+S_n$. The same result
follows by integrating the discontinuity against a smooth test
function; a double pole gives a derivative of a delta distribution.
Finite local contacts leave the pole order unchanged. This is a
failure already in the response of a single excursion law whose
zero-field two-current test gave exactly one positive simple pole.

One can expose it in the complete spin-spectator angular determinant.
In a constant field the excursion potential is constant in $x$.
Reflection of the third bridge component gives zero first excursion
variation for this probe. Its second variation thus adds exactly the
kernel above times $-\lambda e^2 k_b$ to the transverse Hessian,
where $k_b$ is the coincident angular inverse at the shifted mass
$m_b=m+\lambda(1-w_{B_{\rm bg}})$. The rest is the minimally gauged
angular Hessian at this field and mass. For sufficiently weak real
backgrounds $k_b$ stays nonzero by dominated convergence from $k_0>0$.
The odd-character version uses its own shifted mass and the opposite
excursion sign. On a closed central path of angular duration $t$,
the clock separation of two marks is bounded by $t/2$: either arc
connects the same pair of points at speed at most one. The minimal
Wilson derivatives at complex $\omega$ consequently cost at most
$t^2\exp(|\operatorname{Im}\omega|t/2)$. The real magnetic phase
has modulus one. The original return estimate at short duration and
$m_b>0$ at long duration give a holomorphic minimal contribution in
$|\operatorname{Im}\omega|<2m_b$. Whenever $S_0<4m_b^2$, the first
double pole is inside this strip and survives in the complete Hessian.
This gives a class of actions for which no cancellation by the minimal
part is possible. Outside the strip the isolated-contact test still
holds, and the other part must be continued separately.

\paragraph{Which superpositions can keep a positive magnetic spectrum?}
First let $\rho(s)\,ds$ be the zero-field scalar spectral weight assigned
to the excursion contact, with support in $s\ge s_0>0$. Normalize
each component by $s/4$ so its zero-field kernel is $(X+s)^{-1}$.
The physical determinant may also multiply the result by a common
background-dependent factor $k_b/k_0$. Write $g(s)=s^{3/2}\rho(s)$,
extended by zero below its support, and $d_n=4(2n+1)b$. The spectral
distribution for the normalized contact of the mixture is
\begin{equation}
 \rho_b(S)=\frac{8b}{S^{3/2}}\sum_{n\ge0} g'(S-d_n).
 \label{eq:fifth-mag-mixture-uv}
\end{equation}
The derivative includes threshold jumps. In fact the coefficient of the
simple kernel after $S=s+d_n$ is $12b\,g(s)/S^{5/2}$. Integrating
the double kernel by parts adds the derivative of
$8b\,g(s)/S^{3/2}$, cancelling the term with $S^{-5/2}$.
These identities hold distributionally on compact spectral intervals.
For integrable $g$ the transform converges at infinity. The
individual delta-mass case follows directly from
\eqref{eq:fifth-magnetic-contact-uv}.

\paragraph{A sharp half-moment obstruction, including all finite-rate mixtures}
The preceding formula also holds for a locally finite spectral
measure $d\rho$, upon setting $dg_0(s)=s^{3/2}d\rho(s)$ and
interpreting $g'$ as the distributional derivative of $dg_0$.
Let $G(z)=\int e^{-zs}\,dg_0(s)$, $\Re z>0$. Suppose
\begin{equation}
       \int_{s_0}^\infty \sqrt{s}\,d\rho(s)<\infty .
       \label{eq:fifth-half-moment-uv}
\end{equation}
Then $G$ is holomorphic and $xG(x)\to0$ as $x\downarrow0$:
the integrand relative to $\sqrt{s}\,d\rho(s)$ is
$xs e^{-xs}$, bounded and pointwise vanishing.

At a fixed $b$, positivity of the contact spectrum in
\eqref{eq:fifth-mag-mixture-uv} would make the derivative of the
one-sided sum of translates of $dg_0$ by $nD$, $D=8b$, a positive
measure. Multiplication by the smooth positive $S^{3/2}$ and the
common translation by $4b$ preserve the sign. The Laplace transform
of that derivative is $H_D(z)=zG(z)/(1-e^{-Dz})$. These distributional
identities follow by testing at $\Re z>0$, where the series converges
absolutely. For a positive measure $|H_D(x+i\omega)|\le H_D(x)$.
At the first harmonic $\omega=2\pi/D$ the denominators coincide,
so
\begin{equation}
 |G(x+i\omega)|\le\frac{xG(x)}{|x+i\omega|}.
 \label{eq:fifth-field-harmonic-uv}
\end{equation}
If positivity held on any open interval of magnetic strengths,
its first harmonics would fill an open interval of nonzero frequencies.
The right side goes to zero uniformly on every compact subinterval.
Thus $G$ extends continuously by zero to that imaginary-axis segment.
Schwarz reflection and the identity theorem force $G$ to vanish
in the whole right half-plane. Since $dg_0$ is positive, this is
possible only for the zero law.

Consequently a nonzero mixture with the stated positivity throughout
an interval of fields must have an infinite spectral half-moment.
This excludes every finite-rate positive target in the preceding
gamma-clock prescription since $s\ge s_0>0$ and its rate is
proportional to $\int s\,d\rho$. The argument covers atoms and
unbounded tails. It tests the added term as an independently
observable positive current. Other variations of the complete
generator can change the measured spectrum, as the strip test above
makes explicit.

\paragraph{Infinitely many short excursions can pass this test}
There is a law at the sharp moment boundary. Take
$g(s)=C\,\mathbf1_{[s_0,\infty)}(s)$, $C>0$. The vacuum spectral
density is integrable and the magnetic spectral measure is a positive
sum of atoms of weights $8bC(s_0+d_n)^{-3/2}$ at $s_0+d_n$.
As $b$ decreases to zero these midpoint Riemann sums recover the
vacuum density $C s^{-3/2}\mathbf1_{[s_0,\infty)}ds$ on continuous
compactly supported tests. The first spectral moment is infinite,
so the finite-rate mixture cannot implement this law.

Averaging the gamma clocks before coupling to the angular operator
gives the required magnitude of intensity in heat duration,
\begin{equation}
 d\eta(T)=\frac{C}{2k_0\sqrt\pi}\,
              T^{-3/2}e^{-s_0 T/4}\,dT .
 \label{eq:fifth-excursion-levy-uv}
\end{equation}
Indeed the gamma density is
$\sqrt{s}\,e^{-sT/4}/(2\sqrt{\pi T})$ and its coupling magnitude
is $s\rho(s)ds/(4k_0)$. The resulting tempered half-stable
L\'evy measure has infinite total mass but finite
$\int(1\wedge T)d\eta(T)$. For the massive positive spectral
sign, the added potential is
$-\int(1-w_{B,T})d\eta(T)$, where $w_{B,T}$ denotes the normalized
Wilson excursion at fixed $T$. For smooth bounded field strength
the integral exists. For example, using radial gauge about the
excursion base bounds the loop's small-$T$ flux in $L^2$ by
$C_B T$ for a smooth field with bounded derivatives. Orientation
reversal and $1-\cos u\le u^2/2$ give an integrable $O(T^2)$
bound. Equivalently its perturbative variations at zero field can
be defined directly by the following joint inequalities.

At least two photon marks occur on each nontrivial excursion. For
$r\ge2$, H\"older's inequality and the conditional bridge estimate in
Section~\ref{sec:four-excursion-uv} bound their mixed moment by
$C_r\prod_j\min(\sqrt T,|Q_j|T)$. At $T\le1$ this is at most
$C_r T^{r/2}\prod_j\min(1,|Q_j|)$ in fixed length units.
Its integral against $d\eta$ is finite even at $r=2$. At $T\ge1$
exponential decay controls every required duration moment. This gives
the same per-end bounds for all fixed-order closed Euclidean graphs
with spatial A4 photon sewing, including graphs with several marks
on one excursion. At zero field the potential vanishes exactly and
the original unmarked vacuum is unchanged. This construction takes
its difference inside the geometric Wilson observable at each $T$;
no divergent constant is assigned a value.

A positive rate of odd-character jumps would pay $2\int d\eta$ in
the free angular mass. More generally a unitary jump with mean real
character $u\in[-1,1]$ at that heat scale pays $1-u$ per jump,
and $1-u\ge2(-u)_+$. A divergent net intensity of the negative
character required for this sign cannot be realized at finite
unmarked cost by adding ordinary positive-rate channels of this
type. The signed curvature potential and a positive jump measure
therefore have different physical content. At large constant $b$
the signed potential itself decreases in proportion to $-\sqrt b$,
by scaling $T\mapsto T/b$ in
\eqref{eq:fifth-excursion-levy-uv}. Its uniform positive-mass
majorant can then fail. The positive spectral weights above apply to the current
parallel to the background field. The orthogonal current is tested in
Section~\ref{sec:six-mag-uv}.
It supplies neither a positive charged opening nor unrestricted
Maxwell sewing of the resulting spectrum.

\paragraph{Opening the pinned excursion at its two geometric ends}
The Brownian loop also defines a covariant two-end kernel. Work on
the Euclidean locked four-plane, before the spatial photon readout.
Let $\mu_0=1/\theta>0$ have inverse-area units and put
$H_B=-(\nabla^B)^2$ for the scalar magnetic Laplacian with Wilson
phase $\exp(ie\int B)$. For a smooth real connection its heat kernel
on the diagonal is the free return density $(4\pi T)^{-2}$ times
the pinned Wilson expectation at duration $T$. The gamma density is
$\sqrt{\mu_0/\pi}\,T^{-1/2}e^{-\mu_0 T}dT$. Dividing by the free
return density before integrating gives exactly
\[
 w_B(x)=12\pi^2\sqrt{\mu_0}\,
       [(H_B+\mu_0)^{-5/2}](x,x).
\]
Indeed the remaining power of $T$ is $3/2$ and
$\Gamma(5/2)=3\sqrt\pi/4$. At $B=0$ the free diagonal is
$1/(12\pi^2\sqrt\mu_0)$ before multiplication by the displayed
constant. Replacing the diagonal by $(x,y)$ in this same heat
integral opens a charged Wilson path with its two endpoint phases.
The free Fourier multiplier of this particular opening is
proportional to $(p^2+\mu_0)^{-5/2}$.

There is an explicit obstruction to treating that multiplier as the
two-point function of a positive scalar charged field. Omit its
positive constant, fix spatial momentum $k$, and set
$a=(|k|^2+\mu_0)^{1/2}$. Its time transform $h$ has
\[
 h(0)=\frac{2}{3\pi a^4},\qquad h'(0)=0,\qquad
 h''(0)=-\frac{1}{3\pi a^2}.
\]
These follow by integrating $(\omega^2+a^2)^{-5/2}$ and its product
with $-\omega^2$ against $d\omega/(2\pi)$. Both integrals converge
absolutely. Positive-time probes at $t$ and $2t$ have reflected
determinant
$h(2t)h(4t)-h(3t)^2=h(0)h''(0)t^2+o(t^2)<0$.
Smooth time pulses and a small spatial momentum packet preserve the
strict sign. Equivalently, the multiplier has no nonzero positive
Stieltjes measure since its product with $p^2$ tends to zero.
This test concerns the opened scalar kernel; the closed excursion
has the distinct two-source pole at $s=4\mu_0$ computed above.

One cannot assign that pole to an ordinary independent massive
particle and antiparticle merely by matching its location to the
threshold. If each constituent has mass $M=\sqrt\mu_0$, in the pair
rest frame $s=4(M^2+|p|^2)$. The spectral projection onto the
singleton $s=4M^2$ multiplies the relative-momentum wavefunction
by the indicator of $p=0$, a null set for the two-body measure.
Every square-integrable spin and momentum vertex gives zero mass to
this singleton, even if its continuum density is unbounded at the
edge. A nonzero pole requires a separate normalizable state or
different interacting spectral data.

There is also a test of the duration measure at the action level.
A positive direct integral of ordinary four-dimensional scalar
determinants, all with the same Wilson coupling and determinant
sign, weights their pinned loops per unit volume by
$T^{-3}\int e^{-sT}d\sigma(s)\,dT$ up to a common constant.
The power includes the proper-time trace and free return density.
To reproduce the prescribed gamma duration law literally would
require the positive mass Laplace transform to be proportional to
$T^{5/2}e^{-\mu_0 T}$. At $T\downarrow0$ every nonzero positive
measure on $s\ge0$ has a strictly positive or infinite limiting
Laplace mass, whereas this function tends to zero. Signed
combinations, additional marked vertices and interacting auxiliary
fields have other duration measures. Their charged endpoints,
determinants and physical cuts require separate constructions.

A different readout leaves the free charged endpoints positive and
makes the price of closing them explicit. Write $L_B=H_B+\mu_0$.
Applying this local covariant operator at each endpoint gives, as an
operator identity on smooth tests,
\[
 L_B\,L_B^{-5/2}\,L_B=L_B^{-1/2}
 =\frac1\pi\int_0^\infty
       \frac{du}{\sqrt{u}}\,(L_B+u)^{-1}.
\]
The integral follows from the scalar spectral theorem. At zero field
this is a positive continuum of ordinary charged scalar masses,
tested as a generalized free field. Reclosing with the original
geometric diagonal requires a factor $L_B^{-1}$ at each port,
or $L_B^{-2}$ in the trace. In the free four-plane the ordinary
diagonal integral of $L_0^{-1/2}$ diverges, whereas that of
$L_0^{-5/2}$ is finite. With smooth endpoint tests one recovers
the original kernel exactly before taking the finite pinned-diagonal
limit. In a varying field the two reclosure factors also contribute
derivatives fixed by $\delta L_B^{-1}=-L_B^{-1}(\delta L_B)L_B^{-1}$.
This supplies a concrete covariant sewing with a positive free
opening and the original finite diagonal. Its sewing is
source dependent and extends in clock time. Its multi-current cuts
and its relation to ordinary completeness in the measured charged
Hilbert space still have to be established; the negative reflected
determinant above applies to the original scalar opening.

For the infinite-activity law \eqref{eq:fifth-excursion-levy-uv},
the same heat-kernel calculation gives a diagonal \emph{difference}
proportional to
$(H_B+s_0/4)^{-3/2}-(H_0+s_0/4)^{-3/2}$. The power is now $3/2$
because $T^2$ from the inverse free return multiplies $T^{-3/2}$
from the intensity. The individual diagonals have a linear
four-momentum divergence. The difference exists with the two
Wilson terms combined at the same duration, as proved above.
The covariant off-diagonal kernel before taking that difference
decays as $|p|^{-3}$ in the free theory. A nonzero scalar propagator
with a positive spectral measure has
$\liminf_{X\to\infty}X D(X)>0$ at $X=p^2$ by monotone convergence.
Thus the elementary scalar opening also fails for the law that
passed the magnetic spectral test.

\subsection{Which normal records can change a physical photon closure?}
\label{sec:fifth-joint-completeness-uv}
Fix one reconstructed sector plane and its free physical transverse
Maxwell field. Three different operations on a normal record can be
tested separately: conditioning the vacuum, inserting a filter
between photon operators, and selecting an outcome jointly with a
charged pair. A positive probability or a Ward identity by itself
does not identify these operations.

\paragraph{Conditioning a pure Maxwell vacuum}
Suppose the zero-charge state restricts to the pure Fock vacuum
$\omega_\gamma$ on the full transverse photon Weyl algebra, including
both the field and its conjugate momentum. For a normal record effect
$0\le E\le I$ commuting with this algebra define
$\phi_E(a)=\omega(aE)$. On positive $a$,
$0\le\phi_E(a)\le\omega_\gamma(a)$. The two positive functionals
$\phi_E$ and $\phi_{I-E}$ decompose the pure marginal. Normalizing
their nonzero summands therefore gives the same vacuum in each, so
\[
             \omega(aE)=\omega_\gamma(a)\,\omega(E).
\]
The argument also applies to any positive conditional states whose
average is this pure marginal. Applying it to record events and
bounded photon tests, then to integrable coefficients by approximation,
factors the measured two-port kernel from the photon contraction.
No translation invariance of individual record realizations is
needed. Knowledge of the probability law of Euclidean field
coordinates alone is weaker than purity on the stated algebra.

\paragraph{A filter between the two photon ends}
For a bounded effect that can act on the photon, write $a(f)$ for
annihilation in a physical packet $f$ and normalize the packet inner
product by the free one-photon norm. Its two-port weight obeys
\begin{align*}
 \langle a^\dagger(f)\Omega,E a^\dagger(g)\Omega\rangle
    ={}&\langle f,g\rangle\,\omega(E)\\
       &+\omega([[a(f),E],a^\dagger(g)]) .
\end{align*}
Commuting the annihilation operator to the right proves the identity,
understood on finite packets as a weak matrix-element identity.
Thus the difference has a specified one-photon matrix element. For
example the vacuum projector has unit vacuum expectation and zero
one-photon compression.

If the fixed or inclusively averaged filter commutes with translations
and the proper orthochronous Lorentz group, its compression to the
observed one-photon subspace commutes with the massless helicity
representation. On each helicity it is a constant $z_\lambda I$ by
irreducibility, with $0\le z_\lambda\le1$. With several identical
copies it is a fixed operator on their multiplicity space. In
particular the separated-time covariance in one observed helicity
is $z_\lambda e^{-r|\tau|}/(2r)$ at $r=|\mathbf k|>0$.
A nonzero overlap retains Maxwell order at every momentum.
For a complete family of outcomes $\sum_a E_a=I$ the constants
sum to one. This test permits noncommuting normal records and applies
to a filter on the photon alone. It does not assume that an arbitrary
nonlocal action is represented by such a filter.

\paragraph{A positive joint filter that attenuates the boost integral}
A real photon $k$ and a massive pair $P$ have a relative invariant
even when their individual masses and polarizations are fixed. Use
the mostly-minus convention for these future physical momenta only.
On their joint space consider the scalar effect
\[
 E(k,P)=|\chi_{248}(L^2 k\cdot P)|^2 ,
\]
where $L$ is the geometric length and $\chi_{248}$ is the normalized
one-coordinate A4 characteristic function. In the vacuum intermediate
state of Section~\ref{sec:four-boost-uv}, put
$k=(r,\mathbf k)$, $P=(w,-\mathbf k)$, $w=\sqrt{r^2+s}$.
Then $k\cdot P=r(w+r)\ge2r^2$. The squared ball transform is
$O((1+|x|)^{-249})$. For a compact occupied mass band this weight
makes the transverse direct-current boost integral finite: its bare
radial integrand is $O(1/r)$ at infinity from the virtual denominator
$r(w+r)$ and pair density $s\,d\rho(s)/(2w)$.
The estimate tests a conditional transition with this assigned joint
effect. A complete measurement with the second effect $I-E$ restores
the full norm. Indeed
$v\mapsto(\sqrt E\,v,\sqrt{I-E}\,v)$ is an isometry. Retaining only
the first outcome needs a physical selection rule to identify it as
the entire sector entering the charged determinant.

There is a separate calibration with the pair and photon moving
in the same spatial direction with equal spatial momenta. Then
$k\cdot P=r(w-r)$ takes every value in $(0,s/2)$ as $r$ grows.
These energy-mismatched matrix elements can be tested with a
time-dependent source. If a normalized positive record law with
characteristic function $h$ is required to obey $|h|=1$ throughout
that interval of arguments, its record is deterministic in that
direction. In fact equality makes the phase of the record constant
almost surely at each tested frequency. Two frequencies with
irrational ratio force the difference of two independent records
to vanish almost surely. The same $h$ then has unit modulus on the
oppositely moving configurations as well. This conditional calibration
is extra physical input; fixing zero-transfer charge only gives the
single value $h(0)$.

\paragraph{How accurately must the joint channel be opened?}
Let $0\le E\le I$ on photon and charged-record packet spaces. Suppose
the loss $\langle f\otimes g,(I-E)f\otimes g\rangle$ is at most
$\epsilon$ on every normalized product packet. For a normalized
opening $\psi=\sum_j c_j f_j\otimes g_j$ in Schmidt form, the triangle
inequality for $(I-E)^{1/2}\psi$ gives
\[
 1-\langle\psi,E\psi\rangle
       \le\epsilon\Big(\sum_j c_j\Big)^2
       \le \epsilon\,\operatorname{rank}_{\rm Schmidt}(\psi)
\]
at finite rank, with $c_j\ge0$. Exact equality with unit probability
on all product packets forces $E=I$ by density. The rank dependence
at nonzero error is sharp. For two $d$-dimensional record spaces put
$\Psi_d=d^{-1/2}\sum_{j=1}^d e_j\otimes e_j$ and
$E=I-|\Psi_d\rangle\langle\Psi_d|$. Every product packet loses at
most $1/d$ by Cauchy--Schwarz, yet the observed weight of $\Psi_d$
is zero. Thus high-rank coherent records can require an actual joint
opening even when independent normal and photon tests are accurate.
These criteria specify which missing conditional matrix elements a
proposed GCM gluing law has to derive.

\paragraph{Paying for every normal excitation in the intermediate energy}
A complete sum can still attenuate an energy denominator if the
complementary states cost physical energy. The A4 marginal gives a
quantitative example. For $d=248$ let $w_d(y)\,dy$ be the normalized
one-coordinate ball law on $[-1,1]$, with
$w_d(y)\propto(1-y^2)^{(d-1)/2}$. On $L^2(w_d\,dy)$ specify the
nonnegative operator $H$ of the form $\int |f'|^2w_d\,dy$ with
natural boundary, and ground state $1$. The characteristic function
of $v_t(y)=e^{ity}$ in this state is exactly $\chi_d(t)$. Its
\emph{entire} spectral law $\mu_t$ for $H$ satisfies
\begin{align*}
 \int(\lambda-t^2)^2\,d\mu_t(\lambda)&=t^2 I_d,\\
 I_d=\int (w_d'/w_d)^2w_d\,dy&=\frac{d(d-1)}{d-3}.
\end{align*}
Indeed $Hv_t=(t^2-it\,w_d'/w_d)v_t$ with vanishing endpoint flux;
the logarithmic derivative is square integrable for $d>3$.
The beta integral for $y^2$, of parameters $1/2,(d+1)/2$, gives
the displayed constant. In particular the probability of
$\lambda<t^2/2$ is at most $4I_d/t^2$.

Specify the mass in the intermediate charged channel to be
$m_{s,\lambda}=\sqrt{s}+\sqrt\lambda/L$. In the opposite-momentum
vacuum graph put $t=L^2r(w+r)$ as above. The isometric normal
transition uses $v_t$ itself, with no projection onto its elastic
component. With direct transverse coupling and its initial current
row of squared norm $s\,d\rho(s)/(2w)$, its fully summed energy
denominator is
\[
 B(r,s)=\int\frac{d\mu_t(\lambda)}
                   {r+\sqrt{r^2+m_{s,\lambda}^2}}
 \ \le\ \frac{4I_d}{t^2(r+w)}+\frac{\sqrt2 L}{t}.
\]
The first term bounds the low normal energies; above $t^2/2$ the
mass itself is at least $t/(\sqrt2 L)$. For $r\to\infty$,
$t\ge2L^2r^2$, while at small $r$ one uses $B\le1/(r+w)$.
Thus the full radial integral, up to fixed positive factors,
\[
   \int_0^\infty dr\int \frac{r s}{w} B(r,s)\,d\rho(s)
\]
is finite for a gapped positive band with $\int s\,d\rho(s)<\infty$.
Without the excitation mass its compact-band tail is a positive
multiple of $dr/r$. On each such compact band $B(r,s)\sim L/t$:
$\lambda/t^2$ concentrates at one by the variance identity, and
the exceptional lower tail contributes at most
$4I_d/[Lt(r+w)]$ to $(t/L)B$. The high-energy suppression here
comes from recoil into accessible states, even though their total
transition norm is exactly one.

This assigns a definite normal kinetic form and excited mass law
to the joint intermediate source. The calculation concerns the
direct transverse vacuum coefficient. It does not change an
instantaneous Gauss contact by an energy denominator. In particular,
the input row is indexed by the base mass $\sqrt{s}$, whereas a
state opened after recoil has mass $m_{s,\lambda}$. Its measured
current, boosts, contacts and higher transitions have to result
from the same action before this comparison can represent physical
photon sewing.

One can test the proposed identification of the row with an ordinary
matter current directly. If the photon at $\mathbf k$ is opened
and removed, use the remaining row and the displayed excited energies
as that current's spectrum at $-\mathbf k$. At $r=0$ its normal law
is $\mu_0=\delta_0$, so every occupied initial band survives. At
large $r$, a final squared mass at most $S$ requires $s\le S$ and
$\lambda\le L^2(\sqrt S-\sqrt s)^2$. The probability of this region
is $O_S(r^{-4})$ by the same low-energy bound. After multiplication
by $2r$ the whole transverse weight below $S$ is consequently at
most $O_S(r^{-4})\int_{s\le S}s\,d\rho(s)$. A covariant conserved
matter current on a nonzero fixed transverse band has a strictly
positive limit there, since $r/\sqrt{r^2+s}\to1$. Thus this
reassignment of the joint row to the matter current fails a measured
boost test. If the transition is intrinsically joint with the photon,
its definition must instead explain why removal of that photon is
not the physical current opening.

\paragraph{Using the measured recoil energy inside the normal phase}
The last construction uses $w=\sqrt{r^2+s}$ when it forms
the phase $t=L^2r(w+r)$. One might try to determine a single
phase from the energy of its own positive recoil distribution.
That proposed prescription already fails a fixed-point test. Put
\[
 w(t)=\int\sqrt{r^2+m_{s,\lambda}^2}\,d\mu_t(\lambda),
 \qquad u=Lr .
\]
This integral is finite since $\int\lambda\,d\mu_t=t^2$ by the
quadratic form of $H$. For $t>0$, the variance bound above gives
\[
 \int |\sqrt\lambda-t|\,d\mu_t
 \le t^{-1}\int|\lambda-t^2|\,d\mu_t\le\sqrt{I_d}.
\]
In particular $w(t)\ge(t-\sqrt{I_d})/L$. The self-consistent equation
$t=L^2r(r+w(t))$ would require
$(1-u)t\ge u^2-u\sqrt{I_d}$. For $u>\max(1,\sqrt{I_d})$ its two
sides have opposite signs. No positive phase can satisfy it.
The conclusion concerns this single-phase, mean-energy prescription.
A joint transfer law that resolves each outgoing energy separately
would have to derive both its normalization and its inclusive measure.

\paragraph{Opening each recoil energy at its own phase}
There is a more direct prescription to test before introducing a
new joint law. Let $\Pi_\lambda$ be the individual spectral
projections of $H$, use the actual outgoing energy in each phase,
and define
\[
 t_\lambda=L^2r\bigl(r+\sqrt{r^2+m_{s,\lambda}^2}\bigr),
 \qquad
 \Psi_{r,s}=\sum_\lambda \Pi_\lambda v_{t_\lambda}.
\]
The norm of this row need not equal that of any single phase. In
fact $\|\Psi_{r,s}\|^2=O((Lr)^{-3})$, uniformly in the base mass
as $Lr\to\infty$. Write $u=Lr\ge2$. Since
$t_\lambda\ge u(u+\sqrt\lambda)$, the exact variance identity
separately gives
\[
 \|\Pi_\lambda v_{t_\lambda}\|^2
 \le\frac{I_d}{(1-u^{-2})^2u^2(u+\sqrt\lambda)^2}.
\]
To sum this bound, conjugate $H$ by multiplication by
$\zeta=\sqrt{w_d}$. The resulting interval operator has kinetic
term $-\partial_y^2$ and potential $\zeta''/\zeta$. At $d=248$
this potential is bounded below and tends to positive infinity at
both ends. On the smooth weighted core the transformed function
vanishes there and integration by parts gives the Schr\"odinger
form. Its closure is consequently contained in $H^1_0(-1,1)$.
The interval min--max principle gives compact resolvent and
$\lambda_j\ge c j^2-C$, with eigenvalues counted with multiplicity.
Therefore $\sum_j(u+\sqrt{\lambda_j})^{-2}=O(u^{-1})$.
Orthogonality of the output levels proves the assertion, including
the entire tower. Normalizing a nonzero row of this form would cost
at least a factor proportional to $u^{3/2}$ in amplitude. Thus
using each outgoing energy in the same unaltered A4 phase ceases
to implement an isometry. A conserved joint transfer would require
different matrix elements or a different physical normalization.

\paragraph{Integrating the recoil phase along a unitary transit}
One can use each instantaneous recoil energy and still preserve the
full transition norm by integrating a generator along the normal
transit. This changes the normal dynamics. Enlarge the normal
coordinate to $y\in\mathbb R$ and prepare only its entry state by
the original A4 law, $\psi(y)=\sqrt{w_d(y)}$ on $[-1,1]$ and zero
outside. Its conjugate momentum is $p=-i\partial_y$, with
$\|p\psi\|^2=I_d/4$. Give a channel the squared mass $s+p^2/L^2$.
For $u=Lr>0$ and $\sigma=L^2s$ put
\[
 g(p)=u\bigl(u+\sqrt{u^2+\sigma+p^2}\bigr),\qquad
 A=\tfrac12\{y,g(p)\}.
\]
The rate is precisely $L^2r(r+\sqrt{r^2+s+p^2/L^2})$.
For constant rate $g$, $e^{igy}$ would give the original A4 mark.
Use $U=e^{iA}$ for one dimensionless transit. In the momentum
representation $A=i(g\partial_p+g'/2)$. The coordinate
$z=\int_0^p dp'/g(p')$ ranges over the whole line; multiplying
the wavefunction by $\sqrt{g(p)}$ in this coordinate conjugates
$A$ to $i\partial_z$. Thus $U$ is unitary, and its momentum law
is the pushforward of $|\widehat\psi(p_0)|^2dp_0$ by
$\dot p=g(p)$ at transit time one. This gives the complete normal
measure, including its Jacobian. At $r=0$ the transit is the identity.

For $|p_0|\le u^2/2$, $g\ge2u^2$ implies $p(1/2)\ge u^2/2$,
after which $g(p)\ge u p$ implies $p(1)\ge u^2 e^{u/2}/2$.
The complementary event has probability at most $I_d/u^4$.
Writing $p_1=p(1)$ therefore gives the inclusive estimate
\[
 \mathbb E\frac1{r+\sqrt{r^2+s+p_1^2/L^2}}
 \ \le\ \frac{2L}{u^2}e^{-u/2}+\frac{I_d}{r u^4}.
\]
The ordinary $1/(r+w)$ bound covers small $r$. The preceding direct
transverse graph is again absolutely finite for the stated positive
mass law. The first-order generator defines an internal unitary
transit; its spectrum is not bounded below and it is not a
positive-time physical Hamiltonian. Only the entry distribution is
confined to the A4 interval, and its conjugate momentum already
gives a spread of masses at $r=0$. For fixed $s$ a finite output
mass band containing a neighborhood of $p=0$ has positive weight
there. At large $r$ all bounded output mass bands together below
any fixed mass ceiling have weight at most $I_d/u^4$. Thus this
explicit isometry resolves the normal-transition problem while
still failing the same-current boost test if its photon is removed.
It does not supply the longitudinal Gauss contact or a common
physical source for all the sewings.

\subsection{Elastic photon emission with the same invariant pair mark}
\label{sec:fifth-elastic-uv}
The missing single-charge vertex in Section~\ref{sec:four-excite-uv}
can be restored to first order without losing finite bare
coefficients at higher orders. The issue is whether a spacelike
invariant readout can control a real photon of arbitrarily large
spatial momentum in a boosted frame. The transverse Dirac current
supplies an angular factor needed for the estimate.

On the same locked plane take positive-energy on-shell massive
momenta $p,p'$, write $k=\mathbf p-\mathbf p'$, $r=|k|$,
and put $D=r^2-(E_p-E_{p'})^2\ge0$. Use the very same entire
$F(D)=h_0(D)$ that gives $h_0(-s)$ on the open pair in
\eqref{eq:nextsoft-product-uv}. Normalize spinors to unit Euclidean
norm at each spatial momentum and write
$j^\mu=\overline u(p')\gamma^\mu u(p)$ in that normalization.
The one-body transfer $R_\lambda(k)$ sends $p$ to $p'$ with matrix
$F(D)\,\boldsymbol\epsilon_\lambda(k)\cdot\mathbf j$. Sum these
transfers over positive and negative charges with their respective
elastic Dirac rows and signs. Choose real orthonormal transverse
polarizations equal at $k$ and $-k$ on momentum patches, so
$R_\lambda(k)^\dagger=R_\lambda(-k)$.

\paragraph{A bound uniform in the incident boost}
For $\ell=0,1$ and a fixed incident spin,
\begin{equation}
 \sup_{\mathbf p}\sum_{\lambda,\sigma'}\int
 \frac{d^3k}{r^{1+\ell}}\,|F(D)|^2
 |\boldsymbol\epsilon_\lambda\cdot\mathbf j|^2 <\infty .
 \label{eq:fifth-elastic-bound-uv}
\end{equation}
Indeed the on-shell Dirac identity $k\cdot\mathbf j=(E_p-E_{p'})j^0$
gives $|\mathbf j_T|^2=-j^\mu\overline j_\mu+(D/r^2)|j^0|^2$ in
the mostly-minus convention for this calculation. Summing final spin,
$\sum(-j^\mu\overline j_\mu)=(p\cdot p'-2m^2)/(E_p E_{p'})$,
by $\gamma_\mu\slashed p'\gamma^\mu=-2\slashed p'$.
Also $\sum|j^0|^2\le1$ by orthogonal projection. Here the
four-vector contractions denote ordinary on-shell mass invariants.

In the incident rest frame set $t=D/(2m^2)$, $a=\sqrt{t(t+2)}$,
$c=1+t$. Let $u$ be the cosine of the outgoing direction with the
boost axis, $v=|\mathbf p|/E_p$ and $\gamma=E_p/m$. Then
\[
 \begin{aligned}
 d^3p'/E_{p'}&=m^2a\,dt\,d\Omega,\\
 E_{p'}&=\gamma m(c+vau),\\
 r^2&=m^2[2t+\gamma^2(t+vau)^2].
 \end{aligned}
\]
For $\ell=0$ the invariant contraction, bounded in absolute value,
costs at most $C m^2 a c\,dt/\sqrt{2t}$ after angles, before $|F|^2$.
For the density term one must integrate $E_{p'}D/r^3$. If $v\ge1/2$,
substitute $y=\gamma(t+vau)$; the integral over $u$ is at most
$2(c+a)/(va)$ since
$\int_{\mathbb R}2t(2t+y^2)^{-3/2}dy=2$. For $v<1/2$ it is at most
$C(c+a)/\sqrt{2t}$. Combining with $m^2a\,dt$ bounds both parts
by $C m^2(1+t)^{3/2}dt$. The additional $1/r$ for $\ell=1$ costs
at most $1/(m\sqrt{2t})$. Both resulting integrals against
$|h_0(2m^2t)|^2$ converge at the origin and at infinity by the A4
ball bound. Taking absolute values of the final spin and polarization
entries changes only finite matrix factors. The estimate holds on the
antiparticle band as well.

\paragraph{A positive completion with its induced contact}
Write $\int_k=\int d^3k/(2\pi)^3$, use photon oscillators with this
measure, and set on each fixed charged-constituent sector
\begin{align}
 H_{\rm ref}&=H_{\rm ch}+e^2 V_4+
       \sum_\lambda\int_k r\,T_\lambda(k)^\dagger T_\lambda(k),
       \nonumber\\
 T_\lambda(k)&=a_\lambda(k)+
             \frac{e R_\lambda(k)}{\sqrt{2}\,r^{3/2}},
 \qquad H=U_e^\dagger H_{\rm ref}U_e .
 \label{eq:fifth-elastic-square-uv}
\end{align}
Here $H_{\rm ch}$ is the free constituent energy, $V_4$ the previous
positive quartet, and $U_e$ the original pair conversion. The
square fixes its contact $e^2\sum_\lambda\int_k R_\lambda^\dagger
R_\lambda/(2r^2)$. With $n$ charged constituents this contact has
absolute kernel norm at most $C e^2 n^2$ by the $\ell=1$ estimate
and the triangle inequality for the constituent sum. The linear term
is bounded in form by $C|e|n\|\psi\|\|H_\gamma^{1/2}\psi\|$.
The closed positive sector forms thus define an unrestricted
self-adjoint direct sum. The vacuum has zero energy. In an ordered
expansion the $\ell=0$ bound controls absolute creation from $N$
photons with factor $C n\sqrt{N+1}$, and absorption by adjunction.
Only finitely many number sectors occur at a fixed order from
finite-particle packets. Together with the contact and conversion
bounds in Section~\ref{sec:four-excite-uv} this proves absolute
momentum convergence of every resulting finite-time ordered term,
including intermediate emissions from a charge and their reabsorption.
No photon momentum cutoff, mass correction or vacuum subtraction is
used.

At first order the linear coupling is the product of the canonical
photon coordinate and the elastic transfer. The oscillator equation
therefore supplies $eR_\lambda$ up to the overall Maxwell-source sign;
the pair conversion supplies its previous timelike row. Thus the same
measured transverse field has the full Dirac elastic and open-pair
Born matrix elements with the analytic mark $h_0$, including the
moving-charge, zero-transfer current. The new interaction annihilates
the zero-constituent sector. More precisely, let
$f_k(s)=u(s,k)/\sqrt{2ws}$ in the notation of
\eqref{eq:four-excite-angle-uv}. For $\tau\ge0$ the exact two-point
function at fixed $k$ is
\begin{equation}
 C(\tau,k)=\frac{Z e^{-r\tau}}{2r}
 +e^2\langle f_k,e^{-\tau H_{\rm ref}^{(1,1)}(k)}f_k\rangle .
 \label{eq:fifth-elastic-optical-uv}
\end{equation}
The second Hamiltonian acts on one particle, one antiparticle and any
number of photons at that total momentum. It is bounded below by
$2m$. This preserves a massless photon of weight $Z$ at all momenta
and identifies its entire pole for $0<r<2m$. The remaining spectral
measure is positive. Its first perturbative change relative to
\eqref{eq:four-excite-optical-uv} has order $e^4$, since a single
linear insertion changes photon number. Its exact invariant density
requires a separate calculation.

There is also a sharp infrared price for attempting to keep the
original isolated single-charge state by a differentiable Fock-space
rotation. At first order the required one-photon displacement divides
the emission row by $E_{p-k}+r-E_p=r(1-\mathbf v\cdot\hat k)+O(r^2)$.
Since $F(0)=1$, its squared norm for a moving charge has a soft
radial factor $dr/r$, with positive angular factor
\[
 \int d\Omega\,
 \frac{|\mathbf v|^2-(\mathbf v\cdot\hat k)^2}
      {(1-\mathbf v\cdot\hat k)^2}
 =4\pi\left[\frac{1}{v}\log\frac{1+v}{1-v}-2\right],
 \qquad 0<v<1 .
\]
This rules out that perturbative sharp-state rotation, including
recoil. It does not impair the finite-time UV estimates. The
Hamiltonian completion is a transverse boundary choice. Boost
covariance of the new interacting continuum, a local Gauss generator
and derivation of the same square and contacts from the physical
charged determinant remain necessary.

\paragraph{A variational test of the forward energy cone}
The positivity of \eqref{eq:fifth-elastic-square-uv} can be tested
against boosts without assuming that a charged mass-shell eigenvector
exists. In the one-charge fiber of total momentum $\mathbf p$ consider
a no-photon vector of fixed spin. Write $W_T(p)$ for its transverse
contact energy divided by $e^2$. Define $W_L(p)$ by the same integral
with $\sum|j^0|^2$ in place of $\sum|\mathbf j_T|^2$, the usual
Coulomb contact for the elastic density alone. Let $G(p)$ be the
one-photon second-order gain, given by the $W_T$ integrand multiplied
by $r/(r+E_{p'}-E_p)$. All three integrals are finite. As $|\mathbf p|\longrightarrow\infty$, they satisfy
\begin{equation}
 \begin{gathered}
  W_T(p)\to W_*,\quad W_L(p)\to2W_*,\quad G(p)\to2W_*,\\
  W_*=\frac{\pi^2}{2(2\pi)^3}
       \int_0^\infty |F(\kappa^2)|^2d\kappa>0 .
 \end{gathered}
 \label{eq:fifth-elastic-boost-uv}
\end{equation}
For example, at fixed $k$, write $\mu=\hat k\cdot\hat p$.
The limiting density and transverse spin sums are $1$ and $1-\mu^2$,
and the gain multiplier tends to $1/(1-\mu)$. Changing variables
to $\kappa=r\sqrt{1-\mu^2}$ gives the angular integrals
$\int_{-1}^1\sqrt{1-\mu^2}\,d\mu=\pi/2$ for $W_T$ and
$\int_{-1}^1(1-\mu^2)^{-1/2}d\mu=\pi$ for $W_L$.
For $G$ the latter integrand is multiplied by $1+\mu$, with the
same integral.

To check that no weight escapes to an increasing lab momentum, use
the rest variables above with $y=\gamma(t+vau)$ and
$B=\sqrt{2t+y^2}$. The exact density spin sum is
$1-t/[2\gamma(\gamma+y)]$, from the positive-energy projection.
Apart from the common prefactor $2\pi m|F(2m^2t)|^2dt/[2(2\pi)^3v]$,
the $y$ integrands for $W_L,W_T$ are respectively
\[
 \frac{b}{B^2},\qquad
 \frac{t-1}{\gamma^2B^2}+\frac{2t b}{B^4},
 \qquad b=1+\frac{y}{\gamma}-\frac{t}{2\gamma^2}.
\]
The interval has endpoints $\gamma(t\pm va)$ and expands to the real
line at each $t>0$. There $0\le b\le c+a$.
For $G$ multiply the second integrand by $B/(B+y)\le B^2/t$.
The term containing $t-1$ then has absolute integral
$\le2va|t-1|/(\gamma t)$ and tends to zero. The other terms are
dominated by $2(c+a)/B^2$. Their $y$ integrals are bounded by
$C(c+a)/\sqrt t$, integrable in $t$ against the A4 weight.
These bounds justify the stated limits, including the collinear region.

In the exact variational test use a no-photon vector minus $e$ times
the one-photon emission row divided by $r+E_{p'}-E_p$, restricted
in the trial vector to a compact interval of $r>0$ and a closed
angular subinterval of $-1<\mu<1$. Its one-photon norm is uniformly
bounded for large $|\mathbf p|$. Choose this fixed window so its
limiting gain exceeds $3W_*/2$; positivity and the angular calculation
guarantee such a window. In the expectation of $H_{\rm ref}-E_p$ on
the unnormalized trial vector, the free recoil and the linear term
together give minus $e^2$ times this window's gain. The contact gives
$e^2 W_T(p)+O(e^4)$, uniformly by \eqref{eq:fifth-elastic-bound-uv}.
The quartet vanishes. Normalization divides by $1+O(e^2)$. Since
$E_p-|\mathbf p|\to0$, the variational energy lies strictly below
$|\mathbf p|$ for sufficiently small fixed $e\ne0$ and sufficiently
large $|\mathbf p|$. The unitary $U_e$ preserves total momentum, so
the same violation holds for $H$. Thus the positive transverse square
fails the relativistic spectral condition.

The ordinary linear transverse current plus the Coulomb energy of the
elastic density, with no extra transverse square, passes this particular
second-order check: its limiting contact $2W_*$ equals the limiting
recoil gain. Adding both squares would leave $W_*$ in the diagonal
second-order expression. This test selects a contact combination to
derive from the action. It still requires the vacuum-to-pair Gauss
matrix elements and their other action-prescribed terms, which the
elastic number density alone cannot supply.

\paragraph{The recoil logarithm after the leading cancellation}
For that Coulomb choice one can go a step further. In the same $y$
variables the integrand of $W_L-G$ simplifies to
\[
 \frac{b y}{B^3}-\frac{t-1}{2t\gamma^2}(1-y/B).
\]
A primitive is
\[
 -\frac{1-t/(2\gamma^2)}{B}
 +\frac{\operatorname{arsinh}(y/\sqrt{2t})-y/B}{\gamma}
 -\frac{t-1}{2t\gamma^2}(y-B).
\]
Set $t=\cosh\eta-1$, $\gamma=\cosh\zeta_{\rm r}$.
Then $y_\pm/\sqrt{2t}=\sinh(\eta/2\pm\zeta_{\rm r})$: the inverse-hyperbolic-sine
term contributes exactly $2\zeta_{\rm r}/\gamma$ at every transfer.
Indeed $\gamma(B_-^{-1}-B_+^{-1})=
\gamma\sinh\zeta_{\rm r}/(\gamma^2+t/2)\le1$ and
$[(y-B)]_-^+=2(a-t)\sinh\zeta_{\rm r}$.
Thus all remaining terms after multiplication by $\gamma$ are bounded
by $C(1+t+t^{-1/2})$, integrable against the mark. Consequently
\begin{equation}
 \gamma[W_L(p)-G(p)]
 =\frac{m}{8\pi^2v}\left[
 2\zeta_{\rm r}\int_0^\infty |F(2m^2t)|^2dt+O(1)\right].
 \label{eq:fifth-elastic-recoil-uv}
\end{equation}
Nor does retuning constant weights of these two instantaneous contacts
remove this logarithm. In fact $\gamma(W_L-2W_*)$ and
$\gamma(W_T-W_*)$ are bounded. For the leading $1/B^2$ integrand
the omitted $y$ tails are bounded by
$4e^{-\zeta_{\rm r}}\cosh(\eta/2)/\sqrt{2t}$, using
$y=\sqrt{2t}\sinh z$ and $\operatorname{sech}z\le2e^{-|z|}$.
The $2t/B^4$ tails are smaller. The corrections proportional to
$y/\gamma$ cost at most $\eta/\gamma$ in the first integral, since
$|\log(B_+/B_-)|\le\eta$, and $2/\gamma$ in the second.
The remaining $\gamma^{-2}$ terms are integrable by the previous
bounds. Any constant combination cancelling the leading gain therefore
has the same growing term in its diagonal energy coefficient times
$\gamma$. A Lorentz-invariant displacement of a sharp charged
mass shell would instead give constant $\gamma\,\delta E$ at this
order. This comparison identifies a missing contribution even after
the leading Coulomb balance. It does not require a charged pole for
the earlier variational conclusion.

\paragraph{The complementary recoil denominator on spacelike transfer}
The full defect can be isolated, including its finite part. Write
$d=E_{p'}-E_p$, $G_+=G$, and define $G_-$ by using $r-d$ in place
of $r+d$ in the same emission integral. Both denominators are
positive at nonzero transfer. Their equal average supplies
$[1/(r+d)+1/(r-d)]/2=r/D$. Using the transverse Ward identity inside
the Coulomb term gives exactly
\begin{align}
 W_L-\tfrac12(G_++G_-)
 &=\tfrac12\int_k\frac{|F(D)|^2}{D}
          \sum_{\sigma'}j^\mu\overline j_\mu,\nonumber\\
 \gamma[W_L-\tfrac12(G_++G_-)]
 &=\frac{m}{8\pi^2}\int_0^\infty
       \frac{a(1-t)}{t}|F(2m^2t)|^2dt .
 \label{eq:fifth-elastic-symmetric-uv}
\end{align}
The last line uses the invariant measure $d^3p'/E_{p'}$ and is finite
at both endpoints. The complete comparison is the covariant exchange
between the conserved on-shell rows on spacelike transfers, where
there is no photon pole. It gives exactly constant $\gamma\,\delta E$
at every boost. The excess in the forward-only energy is
$(G_--G_+)/2=\int_k |F(D)|^2 d\sum|\mathbf j_T|^2/(2rD)$.
Thus a specific current contraction identifies the missing term,
including its sign and coefficient, without fitting a function.
The physical intermediate state consisting of the recoiling charge
and a positive-energy photon has energy $E_{p'}+r$, and generates
only $G_+$. Assigning $G_-$ to the same state would change its
energy denominator. A derivation of the covariant answer needs the
other time orderings or contacts of the physical charged action
and an independent check of their positive-state sewing.

\subsection{A physical curvature test before choosing how to sew a path}
\label{sec:six-physical-uv}
A proposed gluing law can be tested on the photon that it actually
produces. This gives an obstruction even when the unobserved normal
records are correlated with the photon. Work on one reconstructed
plane, with its declared continuation and its physical normalization
of the Abelian field strength. Let $X=Q^2$ be the Euclidean momentum
invariant, and assume that, on conserved probes, the observed photon
has transverse scalar covariance $D(X)$ with positive spectral measure
$d\nu(s)$. Here $s$ under a mass integral denotes invariant mass
squared, independently of the sector index in $\bin=(s,b_s)$.
Allow the local contacts of a tempered spectral representation.
Suppose there is a finite interval $[s_0,s_1]\subset(0,\infty)$
of nonzero finite weight. A photon that obeys the usual Maxwell
equation with a nonzero charged-current band has this property.
Indeed, on an intermediate state of mass squared $s$, its transverse
matrix elements obey $s A_T=\pm e j_T$, up to the common current
sign. In the convention with current spectral tensor
$(q_\mu q_\nu-g_{\mu\nu}q^2)d\rho_J(s)$ this gives
$s\,d\nu(s)=e^2d\rho_J(s)$ on that band. If the physical
Maxwell equation contains additional polarization currents, the
band test uses their \emph{total} matrix element.

\paragraph{The separated curvature response}
For a nonzero Euclidean separation $y$, define the connected
point-split response in the vacuum by
\[
 {\cal C}(y)=\tfrac14\sum_{\mu,\nu=1}^4
  \langle F_{\mu\nu}(y) F_{\mu\nu}(0)\rangle_c .
\]
This is a contraction of two observed field strengths, with the
Euclidean time components continued in the same convention as the
spatial components. It is distinct from the Hamiltonian energy
density. Before taking the inverse Fourier transform the contraction
is $3X D(X)/2$ up to local polynomials: the longitudinal part drops
out of each field strength. On the separated test,
$X/(X+s)=1-s/(X+s)$ leaves only the second term. With $r=|y|$ and
$K_1$ the modified Bessel function this gives
\begin{align}
 {\cal C}(y)&=-\frac32\int_{(0,\infty)}
                      s\,\Delta_s(r)\,d\nu(s),\nonumber\\
 \Delta_s(r)&=\frac{\sqrt{s}}{4\pi^2r}K_1(r\sqrt s),\qquad
 \lim_{r\downarrow0}r^2[-{\cal C}(y)]
    =\frac{3}{8\pi^2}\int_{(0,\infty)}s\,d\nu(s).
 \label{eq:six-curvature-floor-uv}
\end{align}
The limit allows $+\infty$. The kernel follows by integrating
$(4\pi t)^{-2}\exp[-r^2/(4t)-st]$ over $t>0$.
The function $zK_1(z)$ decreases from one to zero since its
derivative is $-zK_0(z)$. Monotone convergence gives the limit.
At positive separation the exponential decay also allows a tempered
high-mass tail and finitely many local dispersion subtractions. All
polynomials in $Q$ have support at $y=0$ after Fourier inversion.
The massless photon by itself gives zero in this particular separated
trace. The strict sign comes from the occupied massive interval,
regardless of the photon pole residue or the sign of any unobserved
geometric determinant.

There is a perturbative consequence with explicit hypotheses. Suppose
the order-$e^2$ coefficient of the actual photon two-point function
on separated probes contains a nonzero positive massive measure
$d\nu_2$, together with a possible correction to the massless
residue and local contacts. For example this follows for an ordinary
Maxwell port on a Born pair channel when its cut exhausts the massive
weight at that order. Apply \eqref{eq:six-curvature-floor-uv}
to $d\nu_2$. If the \emph{local} coincident quadratic Maxwell
insertion belongs to the claimed complete bare graphs, and separating
its two field legs multiplies the same integrands by
$e^{iQ\cdot y}$, absolute convergence of those integrands bounds the
response uniformly down to $y=0$. This follows by dominated
convergence, with an absolutely summable majorant if the source has
infinitely many resolved pieces at that order. The spectral result
instead gives a divergence at least proportional to $1/r^2$.
This proves incompatibility on that branch at the first occupied
perturbative order. It needs no factorization of the hidden variables,
choice of momentum cutoff, or interchange of a nonperturbative
limit with a coupling derivative. A spectral band known only in the
exact theory proves the exact separated statement; assignment to a
particular perturbative coefficient requires the coefficient hypothesis.

For an independent free Maxwell reference there is a direct consequence
for an ordinary connected vacuum graph. Work first on finite transverse
test modes in a translation-invariant vacuum, with zero one-current
expectation. Suppose the charged source couples through the same $eA$
and is integrated at fixed $A$. Write its normalized two-port term as
$e^2 A^{\mathsf T}M A/2$, including every contact of that source, and
let $D$ be the bare transverse Maxwell covariance. Wick contraction
gives the vacuum coefficient $\operatorname{Tr}(MD)/2$ and the photon
covariance coefficient $G_2=DMD$. Inserting the two curvatures on
that line, with separation $y$, gives per unit volume
\[
 {\cal C}_2(y)=\tfrac12\operatorname{Tr}
       [e^{iQ\cdot y}D^{-1}G_2]
     =\tfrac12\operatorname{Tr}[e^{iQ\cdot y}MD],
\]
where the trace includes the momentum density. Thus the separated
test marks the photon in the ordinary order-$e^2$ vacuum contraction.
Absolute convergence of that contraction would bound the marked
answer uniformly in $y$. An occupied positive Born band contradicts
this bound. The identity uses the complete source before closing its
ports. If the photon and the charged record have a joint reference
law, the mixed derivatives in Section~\ref{sec:six-joint-uv} replace
this Wick argument.

A fixed separated measurement remains well defined. The inference
to a bare graph specifically requires closing these two physical
curvatures at coincidence. If the native pairing
$\operatorname{Str}[\mathbb F\mathbb P_*\mathbb F]_{\rm top}$
keeps a relative displacement even when the two marked curvatures
are closed, its insertion is a different measured kernel. The
following calculation tests that possibility quantitatively.

\paragraph{What a retained relative displacement has to do}
Let a positive, atomless geometric record assign a bounded
displacement $Y$ between the two marked curvatures. Consider first
the separated part of their response. If
$0<\int s\,d\nu(s)<\infty$, its absolute geometric average is
finite exactly when $\mathbb E|Y|^{-2}<\infty$.
The upper bound follows from $zK_1(z)\le1$, and the lower bound
near zero follows from any one of the occupied finite mass intervals.
With infinite first moment the inverse-square test remains necessary.
Atoms at coincidence require the full distributional contact and a
separate source derivation. Even for an atomless law, a contact
$\delta(Y)$ samples the density at zero when that density is smooth;
one cannot discard it by observing that $\Pr(Y=0)=0$.

For instance take the actual $n$-ball coordinate of the A4 shell in
\eqref{eq:monomial-owner-uv}, $n=248$. If a curvature pair samples
$Y=L P u$ with $P$ an orthogonal rank-four projection and $u$ uniform
in the unit $n$-ball, $|Pu|^2$ has beta parameters $2,(n-2)/2$.
Thus $\mathbb E |Y|^{-2}=n/(2L^2)$. This is a test of a specified
map from the shell to the relative coordinate. Setting both
curvatures at the same projected $u$ gives $Y=0$ instead.
The shell integral by itself makes no choice between the maps.

Several returns test this requirement simultaneously. Let
$B_1,\ldots,B_m$ be fixed real linear maps of the same ball coordinate
and let $B_S$ denote their vertical stack for indices in $S$. For
positive exponents $a_i$ the exact criterion is
\begin{equation}
 \begin{gathered}
  \int_{B^n}\prod_i|B_i u|^{-a_i}du<\infty\\
  \Longleftrightarrow\quad
       \sum_{i\in S}a_i<\operatorname{rank}B_S\\
  \text{for every }\varnothing\ne S\subseteq\{1,\ldots,m\}.
 \end{gathered}
 \label{eq:six-relative-ranks-uv}
\end{equation}
A zero map causes immediate divergence. For the other maps write
each inverse power as the gamma integral of
$t_i^{a_i/2-1}e^{-t_i|B_i u|^2}$. The $u$ integral is bounded above
and below by fixed positive multiples of
$\det(I+\sum_i t_i B_i^{\mathsf T}B_i)^{-1/2}$. For the upper
bound insert $e^{1-|u|^2}$ and extend to all space. For the lower
bound integrate on the unit ellipsoid of that matrix. Factors with
$t_i\le1$ cost finite constants and change the determinant by at most
a fixed factor. Order the remaining $t_i$ from largest to smallest
and divide each by the next, using one for the final denominator.
In the Cauchy--Binet expansion of the determinant, greedily taking
independent rows from the largest weights gives a monomial with degree
$\operatorname{rank} B_S$ in each successive prefix ratio. Add
identity rows to complete the basis. This monomial dominates the
others up to fixed constants. The ratio integral on $[1,\infty)$
then converges precisely at the strict inequality for its prefix.
Every subset is a prefix in some ordering, and equality already
diverges logarithmically.

For $a_i=2$ a single 248-dimensional record cannot control a product
of 124 such coincidences, even if every individual projection has
rank four. Independent fresh rank-four records at each incidence
do pass the displayed test at every finite multiplicity. The same
criterion holds on a finite product of ball domains by comparison
with inscribed and circumscribed balls. It is an exact positive
source-integral result, including simultaneous collapses. A connected
graph of the native theory would have to yield the stated product
or an appropriate bound after its spin trace and contact terms.
The criterion does not assume that every graph contains a product
of 124 curvature insertions. If the maps vary with duration or with
other boundary records, their determinant constants must also be
integrated on the physical measure.

\paragraph{Opening the curvature with the kinetic norm of the same paths}
An attenuation in a curvature matrix element also enters the photon
covariance when the same transported curvature defines the kinetic
energy. This issue can be settled for the following class of bare
quadratic actions. Fix one transverse polarization, $X=Q^2>0$, and a
positive law on its normal records. Let $T_Q$ map potential amplitudes
to the space of square-integrable transported amplitudes. Give a mode
$a$ the kinetic form $X\|T_Qa\|^2/2$, after quotienting its null
directions. Read the photon with the same map by
$B_Q(a)=\langle v,T_Qa\rangle$, where the record state $v$ has unit
norm. In the completion of the kinetic form the variable $T_Qa$
ranges over $\mathcal R_Q=\overline{\operatorname{ran}T_Q}$.
The variance of this linear readout is its squared dual norm,
\[
 X D_B(X)=\|\Pi_{\mathcal R_Q}v\|^2\le1 .
\]
This proof also covers nonclosed range. For one scalar potential with
record multiplier $h_\gamma(Q)$ and $v=1$ the squared norm is
$|\mathbb E h_\gamma|^2/\mathbb E|h_\gamma|^2$ whenever the
denominator is nonzero. Keeping a phase only in the numerator would
miss the kinetic normalization.

Suppose this same readout has a positive, unsubtracted transverse
spectral representation and its massless pole has residue one in the
above units. Then $X D_B(X)=1+\int_{s>0}X\,d\nu(s)/(X+s)\ge1$.
Both inequalities must be equalities at each $X>0$. The measured
photon is pure Maxwell in this quadratic realization, and
$v\in\mathcal R_Q$ for each such momentum. A residue below one leaves
a finite allowance for massive weight. A nonlinear charged vacuum
can change the quadratic law. Its measured covariance, kinetic
variation and mixed source must then be computed together.

\subsection{Independent photon and normal sewings through an ordered transit}
\label{sec:six-recoil-uv}
A photon may return to a different vertex from the normal excitation
created with it. This gives a multiloop test in which each field has
its own complete intermediate-state sum. Work in one physical
transverse polarization and its chosen clock. Write $v_r(p)$ for the
normalized outgoing momentum wavefunction of the ordered
transit generated by $A=\{y,g(p)\}/2$ at $u=Lr$ (one may fix $\sigma=1$). Introduce a free normal
oscillator for each $p\in\mathbb R$ with energy
$\epsilon(p)=\sqrt{1+p^2}/L$. It is an independent field. The
ordinary Maxwell oscillator still has frequency $r=|\mathbf k|$.

At a photon incidence use the real normal coordinate
$\mathcal Y(v_r)=b^\dagger(v_r)+b(v_r)$. The wavefunctions $v_r$
can be taken real because the initial A4 amplitude is real and even
and its momentum-space transit is a real change of variables. Set
\[
 b(r)=\int\frac{|v_r(p)|^2}{\epsilon(p)}\,dp .
\]
The transit bound gives $b(r)\le L$ and, for $u\ge1$,
\[
 b(r)\le 2L e^{-u/2}/u^2+L I_d/u^4 .
\]
Indeed, outside the exceptional event the final $|p|$ is at least
$u^2e^{u/2}/2$; on that event $\epsilon^{-1}\le L$.
Consequently both $\int d^3k\,b(r)/r$ and
$\int d^3k\,b(r)/r^2$ are finite. These are complete normal sums,
including the low-energy component.

\paragraph{A convergent determinant with overlapping hard momenta}
Place an emitter at a specified spatial point. Couple the Maxwell
coordinate linearly to $\mathcal Y(v_r)$ there, using the same
incidence for external and internal fields. In canonical coordinate
variables the interaction is the real bilinear of the photon
coordinate and the normal coordinate with kernel
$e\sqrt{2\epsilon(p)}\,v_r(p)$, up to the spatial translation phase.
Let $d k=d^3\mathbf k/(2\pi)^3$. After multiplication by the square
roots of the two free covariances at clock frequency $\nu$ the
off-diagonal kernel without $e$ is
\[
 C_\nu(\mathbf k,p)=
 \frac{\sqrt{2\epsilon(p)}\,v_r(p)}
 {\sqrt{(\nu^2+r^2)(\nu^2+\epsilon(p)^2)}} .
\]
Writing $A(\nu)=\|C_\nu\|_{\rm HS}^2$, one has
\begin{align*}
  A(\nu)&\le A(0)=2\int dk\, b(r)/r^2,\\
  \int\frac{d\nu}{2\pi}A(\nu)
    &=\int\frac{dk}{r}\int
              \frac{|v_r(p)|^2}{r+\epsilon(p)}\,dp <\infty .
\end{align*}
The second identity integrates the two scalar frequency denominators
before the remaining positive integrals. If $e^2 A(0)<1$, the
instantaneous coupled potential form is positive by its free-energy
square-root bound at $\nu=0$. Ordinary canonical kinetic terms
therefore give a nonnegative oscillator dynamics. Free normalization
fixes the relative vacuum log determinant per unit time to
\[
 \frac12\int\frac{d\nu}{2\pi}
            \log\det(I-e^2 C_\nu^\dagger C_\nu).
\]
This follows by block elimination on finite test spaces and by the
trace-class expansion of $C_\nu^\dagger C_\nu$ in the limit.
The absolute integral of the order-$2n$ ring without $e^{2n}$ is
at most $A(0)^{n-1}\int A(\nu)d\nu/(4\pi n)$, even with absolute
values taken on every incidence kernel before sewing. Indeed replacing
$v_r$ by $|v_r|$ leaves its Hilbert--Schmidt bound unchanged, and
the trace of the $n$th power is bounded by $A(\nu)^n$.
The fourth-order trace is strictly positive and integrates two
independent photon momenta and two independent normal momenta. The
argument controls their simultaneous hard regions at every order.
The normalization removes the free determinant; every displayed
interaction coefficient is already absolutely integrable.

\paragraph{Time ordering with different photon and normal pairings}
The bound also applies on a finite clock slab of length $T$ when
each linear incidence includes a bounded matter operator $J(\mathbf
k)$, with $J(-\mathbf k)=J(\mathbf k)^\dagger$ and uniformly bounded
norm. Between vertices use a nonnegative free matter Hamiltonian
on normalized packets. For example, a spatial translation phase
$e^{i\mathbf k\cdot X}$ at the position operator of a mobile emitter
has unit norm. Finite internal spin matrices can be included.
In each vacuum pairing of $2n$ linear interactions, bound a photon
edge by $1/(2r)$. A different pairing joins the normal fields.
For a normal edge between incidences of momenta $r,r'$ separated
by time $t$, its absolute kernel is bounded by the integral of
$|v_r v_{r'}|e^{-\epsilon |t|}$ in $p$. The two endpoint times
integrated over the full slab cost at most
\[
 2T\int \frac{|v_r(p)v_{r'}(p)|}{\epsilon(p)}\,dp
     \le 2T\sqrt{b(r)b(r')}.
\]
Drop time ordering after bounding the matter packet matrix element by
$\sup\|J\|^{2n}$. Every vertex is the end of exactly one normal
edge, so these independent time integrations supply one full
$b(r)$ per photon. The remaining bound is proportional to
$[T\int dk\,b(r)/r]^n$. It includes crossed pairings and all
simultaneously short lifetimes with their full normal measures.
For finite spin spaces absolute internal entry sums change only
the fixed-order constant. This proof concerns linear incidences.
It does not ascribe a duration integral to a coincident two-photon
contact.

\paragraph{An exact contact test on a shared unitary record}
Consider instead a spinless charge of mass $M$ on a spatial torus
with transverse kinetic energy $(P-e\mathcal A_T)^2/(2M)$.
In $\mathcal A_T$ let emission at momentum $\mathbf k$ act by
$a^\dagger(\mathbf k)e^{-i\mathbf k\cdot X}U_r/\sqrt{2r}$
times its transverse unit polarization, and let absorption be its
adjoint. Here $U_r$ is the unitary transit on the same normal
space at both incidences. At charge momentum zero and in the photon
vacuum the entire term linear in $e$ annihilates the initial
vector: its only momentum factor is proportional to
$\mathbf k\cdot\boldsymbol\epsilon_\lambda(\mathbf k)=0$.
This remains true for every freely propagated normal packet.
The complete order-$e^2$ transverse coefficient in the energy
(or per unit slab time in the corresponding diagonal amplitude,
with its Euclidean sign) is therefore the seagull alone,
\[
 \frac{1}{2M}\int\frac{dk}{2r}
              \sum_\lambda\|U_r\psi\|^2
\]
in the continuum normalization for a unit packet. The torus has the
corresponding volume-normalized sum over nonzero modes. Its ultraviolet
tail diverges quadratically. The two opposing normal transports have
multiplied to the identity with zero intervening duration. The independent oscillator
mark has the same equal-time obstruction if its linear coupling is
put into this kinetic square, since
$\|\mathcal Y(v_r)\Omega_{\rm n}\|^2=1$. No virtual graph at
this order can cancel the contact at rest.

The same conclusion holds for a spectral kinetic energy
$f((P-e\mathcal A_T)^2)$ whenever its first two perturbative
variations exist and $f'(0)>0$. Put
$W_1=-(P\cdot\mathcal A_T+\mathcal A_T\cdot P)$ and
$W_2=\mathcal A_T^2$. The first variation on the rest vector equals
$f^{[1]}(P^2,0)W_1$ on that vector and vanishes. Both ends of
the two-$W_1$ term in the diagonal second variation vanish for the
same reason. The $W_2$ term costs precisely $f'(0)$ times its
diagonal expectation. For $f(x)=\sqrt{M^2+x}-M$ this is again
$1/(2M)$, so relativistic square-root dispersion leaves the full
transverse divergence intact. These identities follow on finite mode
tests before taking the positive mode sum. A first-order Dirac
action has different kinetic variations.

The independent-oscillator model uses a prepared oscillator at an
emitter and only transverse coupling. Its observable current creates
a normal excitation when its photon is removed. The recoil depletes
bounded normal energies at high spatial momentum; the analysis
does not identify a covariant charged-pair band in that current.
Conversely the spinless contact shows precisely where a geometric
bound earned on separated incidences is lost under ordinary minimal
closure. A common GCM source must determine which normal propagation
and which coincident contacts its actual charged operator requires.

\paragraph{The complete magnetic correction at rest}
The positive Pauli kinetic energy
$f([\boldsymbol\sigma\cdot(P-e\mathcal A_T)]^2)$ supplies a
nonzero virtual spin transition. Let the normal preparation be a
ground state with energy shifted to zero, $N\ge0$ the intermediate
normal energy, and $v$ its normalized state after emission. This
includes $v_r$ created from the independent oscillator vacuum.
For a single transverse polarization set
$\Delta=f(r^2)-f(0)>0$ and $d_r=\Delta/r^2$. Before applying $f$
the linear mark on a rest spinor has spin matrix, up to sign,
$(\boldsymbol\sigma\cdot\mathbf k)
 (\boldsymbol\sigma\cdot\boldsymbol\epsilon)$.
Its squared norm is $r^2$ on each unit incident spin.
Removing the common photon factor $1/(2r)$, the direct second
coefficient of the kinetic energy is
$f'(0)+r^2 f[0,r^2,0]=d_r$, by the repeated-node divided-difference
formula. The linear mark of $f$ has the additional factor $d_r$.
The complete signed coefficient per mode is therefore
\[
 \frac{d_r}{2r}\left\langle v,
     \frac{r+N}{r+\Delta+N}v\right\rangle
 \ \ge\ \frac{d_r}{2(r+\Delta)} .
\]
Here the virtual term is $-r^2d_r^2$ times the expectation of
$(r+\Delta+N)^{-1}$ before restoring $1/(2r)$.
After the spatial radial measure, the lower bound is proportional to
$\Delta/(r+\Delta)\,dr$. It diverges linearly for both the
Schr\"odinger and relativistic square-root dispersions. Even a
monotone dispersion approaching a positive plateau above its rest
value gives a logarithm. On the ordered oscillator output the
fraction in the expectation tends to one for each of the first
two dispersions: $N$ exceeds $u^2e^{u/2}/(2L)$ outside an event
of probability $O(u^{-4})$. Suppression of the virtual spin
transition consequently leaves the direct coefficient at leading
order. This calculation concerns the positive Pauli kinetic
operator and its own contact, not the determinant of a first-order
relativistic spinor.

\paragraph{A canonical normal coordinate controls the coincident contacts}
The ordered transit has a stronger tail than the preceding second
moment estimate requires. If $p_0\ge-u^2e^{u/4}/4$, the differential
inequality $\dot p\ge-up$ while $p\le0$ gives
$p(1/4)\ge-u^2/4$. From $g\ge2u^2$ and then $g\ge up$ one obtains
$p(1/2)\ge u^2/4$ and $p(1)\ge u^2e^{u/2}/4$. The excluded entry
probability is at most $4I_d e^{-u/2}/u^4$. Consequently
\[
 b(r)\le4L e^{-u/2}(u^{-2}+I_d u^{-4}),\qquad u=Lr\ge1 .
\]
All spatial polynomial moments of the hard envelope now exist.

Couple to the canonical normal coordinate itself, whose vacuum
matrix element contains $1/\sqrt{2\epsilon(p)}$. Normalize its
soft strength by $b(0)>0$. On the same normal oscillator put
\begin{align*}
 m_r&=v_r/\sqrt{\epsilon b(0)},\qquad
 h_r=\langle m_0,m_r\rangle,\\
 n_r&=m_r-h_r m_0,\qquad
 \mathcal N_r=h_r I+b^\dagger(n_r)+b(n_r).
\end{align*}
The specified boundary reads the $m_0$ component elastically and
retains its orthogonal complement as actual oscillator excitations.
The data are real, $\|m_0\|=1$ and $\mathcal N_0=I$. The complete
equal-time weight is exactly
$\|\mathcal N_r\Omega_{\rm n}\|^2=|h_r|^2+\|n_r\|^2=b(r)/b(0)$.
The normal transit is still unitary. The change is the physical
coordinate measured at its exit and its specified elastic channel.

For a spin-$1/2$ particle at $X$ use the free transverse Maxwell
coordinate sampled with $e^{i\mathbf k\cdot X}\mathcal N_r$ at
each incidence; denote this connection by $B_T(X)$. The same
factor multiplies an external transverse potential. With unit
charge write $\rho(\mathbf k)=e^{i\mathbf k\cdot X}\mathcal N_r$.
Take the Pauli kinetic form
$\|\boldsymbol\sigma\cdot(P-eB_T)\Psi\|^2/(2M)$, the free
Maxwell and normal energies, and the full positive Coulomb form
\[
     \frac{e^2}{2}\int \frac{dk}{r^2}\,\|\rho(\mathbf k)\Psi\|^2 .
\]
These forms have a common dense finite-particle core with smooth
momentum packets. The transverse field and its fixed spatial
derivatives are defined there by the finite moments just proved;
the real covariant derivative is symmetric and closable. For the
Coulomb form the adjoint of the map into the weighted $k$-space
contains finite sums of smooth compactly supported test functions
times core vectors, so this map is closable too. The closure of the
sum defines a nonnegative Hamiltonian. A finite number of charges
uses the sum of their signed densities.

On every finite clock slab, take fixed-order Euclidean perturbative
coefficients between charged Schwartz packets and the free photon
and normal vacuum, with a fixed finite number of charged particles.
Include all kinetic and Coulomb contacts before integration. In
each normal Wick pairing a line between marked endpoints is at most
$\|m_r\|\|m_{r'}\|$, even at zero time, and an elastic endpoint
costs at most $\|m_r\|$. Absolute values on the normal momentum
wavefunctions satisfy the same bound. Every transverse photon
therefore costs at most a constant times $b(r)/r$, and each
instantaneous Coulomb line costs at most a constant times $b(r)/r^2$.
At fixed order all covariant-derivative numerators are polynomials
in the finitely many transfer momenta and the incoming charged
momenta. Free propagation is contractive, the time simplex has finite
volume, and Schwartz packets absorb the external polynomial.
Integration of each independent internal momentum converges both at
zero and at infinity; no time separation of a seagull was used.
This proves absolute convergence of every such fixed-order graph.

One may also use the positive spectral kinetic energy
$\sqrt{M^2+[\boldsymbol\sigma\cdot(P-eB_T)]^2}-M$, with the square
defined by its closed form. At positive perturbative order the
resolvent integral for the square root gives terms with $j\ge1$
kinetic insertions separated by $j+1$ free resolvents. Their
auxiliary integral is bounded, apart from the same momentum
polynomials, by $\pi^{-1}\int_0^\infty
z^{1/2}(M^2+z)^{-j-1}dz$. Thus its fixed-order coefficients obey
the same estimates.

The densities commute at equal times because $n_r$ are real.
For real smooth gauge tests put
$\Theta_\chi=\int dk\,\widehat\chi(\mathbf k)\rho(\mathbf k)$
with the Fourier convention
$\chi(X)=\int dk\,\widehat\chi(\mathbf k)e^{i\mathbf k\cdot X}$.
Then $G_\chi=e^{ie\Theta_\chi}$ represents the additive gauge
group. A longitudinal external field $\nabla\chi$ acts by conjugating
the entire Coulomb-gauge Hamiltonian by $G_\chi$, with scalar
potential coupled by $e\Theta_{A_0}$. This gives the usual
time-dependent gauge law because
$i\dot G_\chi G_\chi^\dagger=-e\Theta_{\dot\chi}$. In particular
the normal energy is conjugated as well; the gauge shift has finite
normal form energy since $\|\sqrt\epsilon m_r\|=1/\sqrt{b(0)}$.
This construction retains the total soft charge, an ordinary bare
Maxwell field, and finite instantaneous Gauss and magnetic terms.
It uses the chosen spatial clock and a fixed charged-particle
sector. Its action supplies no vacuum charged-pair continuum or
covariant identification of the high-boost current, and its normal
coordinate readout has not been derived from the GCM action.

\subsection{The orthogonal polarization rules out the half-stable excursion}
\label{sec:six-mag-uv}
Keep the constant background \(eF_{12}\) and the same infinite-activity
duration law as in Eq.~\eqref{eq:fifth-excursion-levy-uv}. The positive
atomic answer in Section~\ref{sec:fifth-mag-uv} used the current along
the magnetic field. Rotate the current into direction 1, keeping the
transfer \(Q=\omega e_4\), \(X=\omega^2\). Both are spatially uniform
transverse probes of the same action. The second probe connects neighboring
Landau levels and gives a different spectral test.

Put \(b=|eF_{12}|>0\), \(\mu=s_0/4\),
\(A_n=\mu+(2n+1)b\), and \(r_b=k_b/k_0\).
Write \(K_1(X)\) for the added transverse determinant Hessian
with \(e^2\) removed. Reflection of both in-plane bridge coordinates
about the excursion base leaves the constant magnetic flux and the
frequency phase invariant and changes the sign of the linear
direction-1 mark. The first excursion variation thus vanishes at
each base point. The coincident angular trace supplies the common
factor \(k_b\). The determinant sign is the same as for the parallel
polarization. With the constant \(C\) of the chosen law, the result is
\begin{align}
 K_1(X)&=3Cr_b b^2\sum_{n=0}^\infty(n+1)
                         [I_n(0)-I_n(X)],\nonumber\\
 I_n(X)&=\int_0^1
  \frac{du}{[(1-u)A_n+uA_{n+1}+u(1-u)X]^{5/2}}.
 \label{eq:six-mag-transverse-uv}
\end{align}

To obtain this formula, divide the covariant heat diagonal by the
free return density, as in the opening calculation. Multiplying by
the excursion intensity gives the heat-trace weight
\(C(4\pi)^2 T^{1/2}e^{-\mu T}/(2k_0\sqrt\pi)\) per \(dT\).
In the second variation, the local quadratic vertex is independent of
\(X\). The two separated linear vertices have the opposite sign.
Their Duhamel time integral is
\(T^2\int_0^1du\). With the charge removed each vertex is twice the
covariant momentum in direction 1. Its squared matrix element between
levels \(n,n+1\) is \(2b(n+1)\), from the oscillator ladder operators.
The reverse transition contributes equally after changing \(u\) to
\(1-u\). Multiply by the Landau degeneracy \(b/(2\pi)\) and by the
free Gaussian integral \(e^{-Tu(1-u)X}/(4\pi T)\) in directions 3 and 4.
Integrating the remaining \(T^{3/2}\) gives
\(\Gamma(5/2)=3\sqrt\pi/4\) and the coefficient in
Eq.~\eqref{eq:six-mag-transverse-uv}.
A constant perturbing potential has zero circulation on every loop, so
\(K_1(0)=0\) fixes the local quadratic contribution. Each series for
the two terms converges since \((n+1)I_n(0)=O(n^{-3/2})\).

In fact the two-mark term is elementary. Set
\(a=\sqrt{A_n}\), \(h=\sqrt{A_{n+1}}\), and \(S=(a+h)^2\).
With \(D(u)=(1-u)A_n+uA_{n+1}+u(1-u)X\), direct integration gives
\(\int_0^1D(u)^{-3/2}du=2(a^{-1}+h^{-1})/(X+S)\).
Applying \(-(2/3)\partial_\mu\) at fixed field gives
\[
 I_n(X)=\frac{2(a^{-3}+h^{-3})}{3(X+S)}
      +\frac{4(a^{-1}+h^{-1})S}{3ah(X+S)^2}.
\]
The sum continues locally uniformly away from its isolated poles.
The double-pole coefficient of \(K_1(X)/X\) at \(X=-S\) is exactly
\[
 4Cr_b b^2(n+1)\frac{a+h}{a^2 h^2}.
\]
It is nonzero for weak fields where \(r_b>0\); all other levels are
regular there. A positive spectral measure at fixed spatial momentum
can assign an atom to this energy, yielding a simple pole. It cannot
supply the derivative of an atom required here. A local contact does
not change this conclusion. Thus the duration law that passes the
longitudinal-to-field test fails as an independently positive current
when the probe is rotated.

There is also a direct separated-time check. The Fourier transform in
\(\omega\) of \((\omega^2+E^2)^{-2}\) is
\((1+E|t|)e^{-E|t|}/(4E^3)\).
For the normalized excursion response \(K_1/X\) the lowest transition
therefore gives \(h(t)=(A t+B)e^{-Et}\) for \(t>0\), up to
exponentially faster terms, with \(E=\sqrt{\mu+b}+\sqrt{\mu+3b}\)
and \(A\ne0\). The higher transition energies grow as \(\sqrt n\)
and their coefficients grow at most polynomially, which bounds their
tail at positive time. Two reflected probes at \(t,2t\) consequently
give \(h(2t)h(4t)-h(3t)^2\sim-A^2t^2e^{-6Et}<0\).
The same determinant sign is obtained from the separated-time part of
\(K_1\), whose nonzero double-pole coefficient differs by the factor
\(-E^2\). Narrow smooth time pulses preserve the strict sign.

For the spin-spectator determinant the minimal angular contribution is
holomorphic for \(|\operatorname{Im}\omega|<2m_b\) by the bounded
central-path separation, including for direction 1.
Whenever \(\sqrt{\mu+b}+\sqrt{\mu+3b}<2m_b\), the first double pole
therefore survives in the complete Hessian. In particular, if
\(\mu<m^2\) at zero field, sufficiently small nonzero fields meet this
strict inequality by continuity. Above the proved strip there is still
the excursion obstruction; deciding the full response there requires
the continuation of the other charged variations. The two-polarization
test leaves the fixed-order Euclidean convergence estimates and the
sector-plane vectorial pole prescription untouched.

\paragraph{Positivity for unbounded field strengths has a cost at zero field}
There is a necessary condition for the whole duration-mixture class.
Consider the excursion current itself, with the common determinant factor
removed, a nonzero nonnegative vacuum measure \(d\rho(s)\) supported
in \(s\ge s_0>0\), and
\(\int(1+s)^{-1}d\rho(s)<\infty\).
Assume the contact has a positive stationary spectrum for both polarizations
at a sequence of field strengths \(b\to\infty\). This stronger hypothesis
concerns the excursion functional on those backgrounds. An angular
determinant known to exist only on a weak-field interval need not meet it.

First apply Eq.~\eqref{eq:fifth-mag-mixture-uv} to the field-parallel
probe. On any fixed compact set of vacuum mass coordinates \(s\), translate
the output variable by \(4b\). For sufficiently large \(b\) none of the
higher translates reaches that set. Positivity implies that the distributional
derivative of \(s^{3/2}d\rho(s)\), extended by zero below threshold, is
positive. A distribution with positive derivative is represented by a
nondecreasing function: integrate the positive derivative and observe that
the difference has zero derivative. Thus \(d\rho(s)=s^{-3/2}g(s)\,ds\)
with a nonnegative nondecreasing \(g\). Write
\(g(s)=\int_{\mu\le s/4}d\nu(\mu)\) almost everywhere, with
\(d\nu\ge0\) supported in \(\mu\ge s_0/4\). This expresses the law as a
positive mixture of the preceding step laws.

Use the direction-1 probe. On any compact interval of \(\mu\) the first
transition \(S=(a+h)^2\), \(a=\sqrt{\mu+b}\), \(h=\sqrt{\mu+3b}\),
is separated from the first transition of the next level for all
sufficiently large \(b\). Indeed the two squared thresholds divided by
\(b\) tend respectively to \((1+\sqrt3)^2\) and
\((\sqrt3+\sqrt5)^2\). The first is strictly increasing in \(\mu\).
Let \(P_b,Q_b\) be its simple- and double-pole coefficients in
\(K_1/X\) for a unit step with \(r_b=1\). On this isolated interval,
pulling the spectral distribution back to a mass-coordinate test
\(\psi(\mu)\) gives
\[
 \int P_b\psi\,d\nu-\int\frac{Q_b}{\partial_\mu S}
                         \psi'\,d\nu\ \ge0
 \qquad(\psi\ge0).
\]
The second term is the derivative of the spectral atom in the double
pole. Higher thresholds lie outside the test. One may mix the
partial-fraction responses first on compact mass sets. The remaining
large-mass tail is locally analytic in the spectral variable, bounded
there by a constant times \(\int_{M}^\infty\mu^{-3/2}d\nu\);
this moment is finite by the vacuum hypothesis and the step representation.
Indeed the subtracted series divided by \(X\) is bounded there by
\(C_b\sum_n(n+1)[\mu+(2n+1)b]^{-7/2}=O_b(\mu^{-3/2})\).

The explicit coefficients in Eq.~\eqref{eq:six-mag-transverse-uv} obey,
uniformly for \(\mu\) in any fixed compact interval,
\[
 \begin{gathered}
  P_b=O(b^{-1/2}),\qquad
  \frac{Q_b}{\partial_\mu S}=\frac{4b^2}{ah(a+h)},\\
  \frac{Q_b}{\sqrt b\,\partial_\mu S}
       \longrightarrow \frac4{\sqrt3(1+\sqrt3)}.
 \end{gathered}
\]
Dividing the test inequality by \(\sqrt b\) and taking the limit gives
\(-\int\psi' d\nu\ge0\). Hence \(d\nu\) itself has a nonnegative
nondecreasing density. For a nonzero law that density is bounded below
by a positive constant on some upper ray. In particular \(g(s)\)
grows at least linearly there and \(d\rho(s)/ds\ge c s^{-1/2}\)
for all sufficiently large \(s\), up to null sets.

This necessary tail already has a consequence in the vacuum. Integrating
over \(s\in[X,2X]\) gives
\(K(X)=X\int d\rho(s)/(X+s)\ge c'\sqrt X\) at large \(X\).
A continuous normalized spatial A4 readout has modulus bounded below on
a ball about zero spatial momentum. On that ball, direct Maxwell sewing
of the same contact has absolute integrand bounded below by a positive
constant times \(1/|\omega|\) at large clock frequency. Its vacuum
contraction therefore diverges at least logarithmically even with that
spatial readout. This rules out simultaneous independent contact positivity
on the stated unbounded fields and absolute vacuum sewing in this
mixture class. It does not substitute that contact for a complete
sourced Hessian with further terms or different gluing.

\subsection{Opening a genuinely joint photon and charged source}
\label{sec:six-joint-uv}
A shared geometric record can correlate the photon with the charged
response even in the reference measure. There is then no justification
for replacing the contraction by a product of two marginal covariances.
Two source tests can be made before specifying how that correlation
is generated. Work on one locked sector plane, initially on finitely
many real transverse test modes. Let $A$ be the dynamical photon,
$b$ a classical connection probe in the same modes, and $J$ an
independent source coupled linearly to $A$. At zero charge let
$(A,\xi)$ have any joint reference law with $\mathbb EA=0$ and
$\mathbb EAA^{\mathsf T}=D$. The record $\xi$ may contain the entire
charged path. Suppose the coupling enters through $e a$, $a=A+b$.
After including variations of the charged weight, determinant, and
boundary measure, write its first two variations as
\[
 W_e(a)=1+e\ell_i a_i
              +\tfrac12 e^2 M_{ij}a_i a_j+O(e^3).
\]
The symmetric matrix $M$ and vector $\ell$ can depend on the whole
reference configuration. They are derivatives of the weight, so $M$
includes products of linear action variations. Take $\mathbb E\ell=0$
and use the connected source $F=\log\mathbb E[e^{J\cdot A}W_e(A+b)]$.
Assume the indicated finite-test derivatives exist.

\paragraph{The mixed photon--current response counts the closed ports}
Evaluate all derivatives at the origin. Put
\[
 V=\partial_e^2F,\quad
 K_{ij}=\tfrac12\partial_{b_i}\partial_{b_j}\partial_e^2F,\quad
 H_{ij}=\tfrac12\partial_{b_i}\partial_{J_j}\partial_e^2F.
\]
Then, with $B=\mathbb E(\ell\cdot A)$ and
$T_{ij}=\mathbb E(\ell_i A_j)$, one has
\begin{equation}
 \begin{aligned}
 K&=\mathbb EM,\qquad V=\mathbb E(A^{\mathsf T}MA)-B^2,\\
 H_{ij}&=\mathbb E[(MA)_iA_j]-B T_{ij},\qquad
 \operatorname{tr}H=V .
 \end{aligned}
 \label{eq:six-joint-cross-uv}
\end{equation}
Indeed the second coupling derivative at a fixed tilt $J$ is
$\mathbb E_J(a^{\mathsf T}Ma)-[\mathbb E_J(\ell\cdot a)]^2$,
by two differentiations of the logarithm. The source derivatives give the
formulas, using $\operatorname{tr}T=B$. The trace identity counts
the two ways to open one of the two connection ports. It uses
neither Gaussianity nor conditional independence. For an ordinary
Gaussian Maxwell source independent of the charged coefficients and
a current of zero mean, $H=KD$. A joint law can change $V$ while
preserving $K$ only by changing the trace of this mixed response
relative to the ordinary opening. This is measurable with one photon
source and one connection probe. A coupling-dependent term that does
not depend on $e(A+b)$ has its own derivative and must be included
separately.

There is also a four-source check. Set
$N_{ij;kl}=\partial_{J_k}\partial_{J_l}K_{ij}(J)|_{J=0}$, where
$K_{ij}(J)$ is the same two-$b$ derivative at tilt $J$. Direct
differentiation gives
\[
 V-\operatorname{tr}(DK)
 =\sum_{ij}N_{ij;ij}+\operatorname{tr}(T^2).
\]
In fact $N_{ij;kl}$ is the covariance of $M_{ij}$ with $A_kA_l$
minus the sum $T_{ik}T_{jl}+T_{il}T_{jk}$. Thus the ordinary absence of
both an order-$e$ mixed current--photon term and an order-$e^2$
two-probe, two-photon connected term already restores the usual
vacuum contraction. All these identities survive the removal of
test projections whenever the corresponding contractions converge
absolutely in the declared source assembly.

\paragraph{A positive conditional tensor that defeats two diagonal tests}
Here is an explicit comparison in one block of $d\ge2$ transverse
real photon coordinates. Whiten the Maxwell field as $Y=D^{-1/2}A$,
with marginal $N(0,I_d)$, and write $r=|Y|^2$. The reference
record may contain $Y$ itself. Put $P_Y=YY^{\mathsf T}/r$ off the
null event $r=0$. For $c>0$ and $R>d+3$, choose
\[
 \begin{aligned}
 h(r)&=\frac{2dc\,\mathbf1_{r>R}}
                  {\mathbb E[(r-d)r\,\mathbf1_{r>R}]},\qquad
 t=\frac{dc-\mathbb Eh}{d-1},\\
 L_Y&=h(r)P_Y+t(I-P_Y),\qquad P=Y^{\mathsf T}L_Y Y=r h(r).
 \end{aligned}
\]
Here $t>0$, since $\mathbb Eh\le2dc/[R(R-d)]$. Give this block
the actual weight
$\exp[-e^2(Y+D^{-1/2}b)^{\mathsf T}L_Y(Y+D^{-1/2}b)/2]$.
Its conditional quadratic tensor is positive. Isotropy gives
$\mathbb E L_Y=cI_d$. The normalized physical photon covariance
also has coefficient $[e^2]\operatorname{Cov}_e(Y)=-cI_d$, since
$\mathbb E[(r-d)P]=2dc$. Yet $\mathbb EP\le2dc/(R-d)$.
In these whitened coordinates the magnitude of the mixed opening is
$(\mathbb EP/d)I_d$, whereas independent sewing would give $cI_d$.

For instance set $D=I_d/X$ and $c(X)=\alpha s/(X+s)$ with
$\alpha,s>0$. The photon variation is exactly
$\alpha[(X+s)^{-1}-X^{-1}]I_d$: it adds a positive massive
atom and shifts the massless residue. Taking independent blocks
with $R>d+3+(L^2X)^2$ gives a UV-integrable order-$e^2$ vacuum
density under the four-dimensional radial measure $X\,dX$ while
retaining both specified diagonal responses. This is an imposed
conditional coupling of transverse modes. It changes the mixed
opening, and its rapidly growing rare-event weight has observable
higher-order consequences.

\paragraph{A fourth-order lower bound without factorized sewing}
Suppose the observed photon has the free Gaussian marginal at zero
charge. After any conditional charged integrations, suppose the
actual source-free weight has the even expansion
$\exp[-e^2P/2-e^4Q/24+O(e^6)]$ for real $P,Q$ on the reference
measure. All native fourth-order terms belong in $Q$.
For any finite orthonormal family of whitened Maxwell modes,
put $G=[e^2]\operatorname{Cov}_e(Y)$. Then
\begin{equation}
 G=-\tfrac12\mathbb E[(YY^{\mathsf T}-I)P],\qquad
 \operatorname{Var}(P)\ge2\|G\|_{\rm HS}^2 .
 \label{eq:six-joint-fourth-uv}
\end{equation}
To prove the inequality, contract with an arbitrary real symmetric
matrix $U$. The centered Gaussian quadratic
$Y^{\mathsf T}UY-\operatorname{tr}U$ has squared norm
$2\operatorname{tr}(U^2)$. Cauchy--Schwarz with $P-\mathbb EP$,
and then $U=G$, gives the result. The other photon modes and all
charged variables may be correlated with $P$. A fourth-order
connected vacuum coefficient, on each finite realization, is
$\operatorname{Var}(P)/8-\mathbb E Q/24$.

For a translation-invariant transverse polarization with
$D(X)=1/X$, suppose its physical covariance variation has a massless
residue change and a positive massive spectral measure. Its whitened
eigenvalue is
$g(X)=z_1+\int X\,d\eta(s)/(X+s)$. A nonzero compact band at
positive mass gives $g(2X)-g(X)\ge c_1/X$ at large $X$, independently
of $z_1$. Squaring, using $|u-v|^2\le2(|u|^2+|v|^2)$, and changing
variables in the dilated term forces
$\int^{2R} X|g(X)|^2dX\ge c_2\log R-O(1)$.
In modes per unit volume this is precisely a lower bound on the
squared Hilbert--Schmidt norm in \eqref{eq:six-joint-fourth-uv}.
Thus a purely quadratic joint weight still fails the fourth-order
bare vacuum test even when its second-order vacuum is finite, as
in the conditional tensor example. A complete action could correlate
a fourth-order contact with this variance. For the combined
coefficient to converge the contact must carry a contribution of
the same divergent order and sign in $\mathbb E Q$; an absolutely
finite fourth-order contact cannot do so. This requirement follows
from the measured photon covariance and the reference marginal
before any decision about gluing independent current kernels.

\subsection{Which charged endpoint changes have a determinant?}
\label{sec:sixnative-endpoint-uv}
A background-dependent change of both Grassmann coordinates cancels
its Berezin Jacobian on each finite test. In infinite dimension
even a connection gauge equivalent to zero can prevent the two
endpoint determinants from existing separately. Take a flat
four-torus on one locked plane and a positive translation-invariant
filter $T_0$. Along the Fourier modes $n e_1$ let its eigenvalues
$t_n>0$ tend to zero for $n\to+\infty$. The smooth periodic gauge
transformation $U(x)=\exp(2\pi i x_1/\mathcal L)$ for torus side
$\mathcal L$ shifts the integer momentum by one. Gauge covariance
gives $T_{0^U}=UT_0U^\dagger$. The normalized endpoint
$T_{0^U}T_0^{-1}$, initially on finite Fourier sums, has entries
$r_n=t_{n-1}/t_n$ along this ray. If it extended to the identity
plus a trace-class operator, then $\sum_{n\ge1}|r_n-1|<\infty$.
The positive product would converge to a finite nonzero number.
But $\prod_{n=1}^N r_n=t_0/t_N\to\infty$. This obstruction
requires no particular rate of damping. A two-sided cancellation
between opposite rays cannot define either ordinary Fredholm
determinant. A complete gauge-covariant relative determinant can
still exist. The comparison identifies a difficulty with separating
its factors and says nothing by itself about a gauge anomaly.

\paragraph{A determinant-class change constructed from the A4 record}
Here is an explicitly different possibility. Let $a_L(z)$ be the
four-coordinate displacement density of the projected A4 ball at its
existing length $L$, normalized in $L^1$. At $n=248$ it and its
first derivatives are square integrable. On a torus of side larger
than $2L$, use its unique supported displacement and the straight
Wilson phase $W_B[x,y]$ to define
\[
 A_B(x,y)=a_L(x-y)W_B[x,y],\qquad
 X_B=A_B^\dagger M_{1-w_B}A_B .
\]
Here $M$ denotes multiplication and $w_B$ is the Wilson expectation
of the closed gamma excursion. For periodic probes it can be
computed on their periodic lift with the same pinned Euclidean
paths. For real smooth probes $X_B$ is nonnegative and trace class.
Indeed $A_B$ is Hilbert--Schmidt, $|W_B|=1$,
$-1\le w_B\le1$, and the exact trace is
\[
 \operatorname{Tr}X_B=c_L\int (1-w_B(x))d^4x,\qquad
 c_L=\int|a_L(z)|^2d^4z .
\]
The Wilson phases cancel within this trace. They remain in the
open kernel and give covariant endpoint phases. The dimensions
are fixed: $c_L$ is an inverse four-volume.

Use the existing angular charged generator $D_B$ of the
spin-spectator branch. Lift $X_B$ to the normalized constant
angular slot by its rank-one orthogonal projection and use the
identity on the retained spin space. Denote the lift by $\widehat X_B$
and its spin multiplicity by $d_{\rm sp}$. Define a new quadratic
action in the original charged variables and with their original
flat measure by
\[
 E_B=\exp(-\widehat X_B/2),\qquad
 \widetilde D_B=E_B D_B E_B .
\]
Its inverse is $E_B^{-1}D_B^{-1}E_B^{-1}$. The free operator is
unchanged because $X_0=0$. Its closed effective action differs from
the original by exactly $d_{\rm sp}\operatorname{Tr}X_B$, at every
source order. In finite dimension this follows by multiplication
of determinants. On a smooth compact, massive test the relative
factor $\widetilde D_BD_B^{-1}=E_B(D_BE_BD_B^{-1})$ has its
ordinary Fredholm determinant. In fact $D_B\widehat X_B$ is trace
class: the single spatial derivative lands on the outer
Hilbert--Schmidt kernel and the angular projection has smooth
finite rank. Thus $D_B\widehat X_BD_B^{-1}$ is trace class with
the same trace as $\widehat X_B$ by cyclicity using bounded
$D_B^{-1}$. Exponentiation gives the second determinant factor.
The sign of the added two-source spectral weight is positive.
The construction fixes its magnitude by $d_{\rm sp}c_L$ and the
excursion measure. Changing variables back in this new integral
retains the compensating Berezin factor and gives exactly the same
additional term.

In the infinite-volume connected expansion use densities or smear
the total translation coordinate. The extra closed vertices are
precisely the single pinned-loop vertices, without any powers of
the containing angular duration. At fixed order their transverse
ends obey the earlier bridge moment bounds, hence the same spatial
A4 photon readout and spanning-tree bounds prove absolute
ultraviolet convergence for all connected closed graphs of this
comparison, including graphs containing both types of vertex.
A field-independent unmarked trace remains that of the original
angular action. No determinant-class claim for an extensive
translation-invariant change on the whole of infinite space is
needed in this graph argument.

This active modification has a measurable limitation. Since
$X'_0=0$, it leaves the entire open one-current vertex unchanged.
At forward two-current order the change of the open inverse is
\[
 \tfrac12\bigl(\widehat X''_0D_0^{-1}
                 +D_0^{-1}\widehat X''_0\bigr).
\]
For fixed Euclidean transfer $Q,-Q$ its spatial multiplier on the
constant angular slot is proportional to
$K_s(Q^2)\widehat a_L(p)\widehat a_L(-p)$.
The compact displacement makes the two transforms entire at
finite charged momentum $p$. This contribution has at most one
charged pole. The earlier additive insertion into $D_B$ instead
had a double pole. More generally the derivative of any regular
invertible congruence $E_0=I$ adds terms $E'_0D_0+D_0E'_0$ to
the amputated first vertex. Sandwiching between the two charged
propagators cancels one outer inverse in each added term. Thus it
cannot change the double charged-pole residue at independently
continued external momenta. Its active measure can still change
the closed spectrum. The new finite action implements both
readouts and demonstrates exactly why its new vacuum pole
requires additional physical intermediate states.

\paragraph{Can a genuine charged pair supply the strength of the new pole?}
A positive interaction can create a normalizable threshold state
in the original particle--antiparticle space. Its current residue
obeys a further constraint. On one transverse polarization take the
existing geometrically marked, charge-one Dirac pair row $u(s)$,
including its spin, angular variables and two-body Jacobian.
Write $\|u(s)\|^2=\rho(s)$ and $M_j=\int s^j\rho(s)ds$.
In the invariant mass normalization of the earlier pair calculation
the current vector is $j(s)=\sqrt{s}\,u(s)$. Thus a scalar spectral
atom of weight $b$ at $s_0$ costs $s_0 b$ of the squared norm
$\|j\|^2=M_1$. This statement uses the conserved-current scalar
before any spatial photon readout. At zero spatial momentum the
latter is one.

Here is a positive quadratic form on the actual two-constituent
space. Fix $s_0=4m^2$, let $T$ multiply by $s$, $Y=T-s_0$, and
put $\phi=u/\sqrt{M_0}$. Project the relative kinetic form
orthogonally away from the single direction $Y^{1/2}\phi$. The
resulting mass-squared operator is
\[
 T_c=T-\frac{|Y\phi\rangle\langle Y\phi|}
                    {\langle\phi,Y\phi\rangle},
 \qquad T_c\ge s_0 I .
\]
The subtraction in this formula is the finite projector term of
the specified interaction. Its coefficient is fixed by the row
$u$. The A4 mark gives $M_2<\infty$, so $Y\phi$ is in the
pair Hilbert space and the perturbation is bounded and rank one.
Positivity follows by applying Cauchy--Schwarz in the relative
kinetic form. Equality holds precisely at multiples of $\phi$,
since $Y$ has no zero eigenvector in the original two-particle
continuum. Consequently $T_c\phi=s_0\phi$, and this genuine
threshold state carries the current strength
\[
 |\langle j,\phi\rangle|^2=M_{1/2}^2/M_0>0 .
\]
There is still essential spectrum on $[s_0,\infty)$ under a finite
rank perturbation. The strict Cauchy--Schwarz inequality for this
nontrivial band gives $M_{1/2}^2/M_0<M_1$. Relative-momentum
contractions at any fixed number of these vertices converge:
the absolute kernel has rank one with row $|Y\phi|$ of finite
squared norm. Bounded free Euclidean mass propagation preserves
this estimate; smooth external pair packets close its two ends.
This is a two-constituent form on one spin channel. It does not
give the multicharge action or its photon sewing.

There is an absolute normalization test even if the mass
interaction is chosen differently. Take $n=248$, $mL=1$ and the
signed row of the previously derived Dirac pair with its single
A4 mark, so $\rho(s)=\chi_n(y)^2\rho_D(s)$,
$y=L\sqrt{s(s/(4m^2)-1)}$. Here
$0\le\rho_D\le1/(12\pi^2)$, $s=2m^2(1+\sqrt{1+y^2})$ and
$ds/dy\le2m^2$. Hence
\[
 M_1\le \frac{m^4}{3\pi^2}
              \int_0^\infty(2+y)\chi_n(y)^2dy<9m^4 .
\]
For the last bound let $f_n(x)=c_n(1-x^2)_+^{(n-1)/2}$ be
the one-coordinate ball density. Parseval gives
$\int_0^\infty\chi_n^2=\pi\|f_n\|_2^2\le\pi c_n$ and
$\int_0^\infty y^2\chi_n^2=\pi\|f'_n\|_2^2$. Beta integration
gives $\|f'_n\|_2^2/\|f_n\|_2^2=(n-1)(n-1/2)/(2(n-2))$.
Cauchy--Schwarz and $c_n\le\sqrt{(n+1)/(2\pi)}$ now prove
the number quoted.

By contrast, the four-coordinate A4 density used in the Wilson
sandwich above gives
\[
 c_L=\frac{(124\cdot123)^2}{\pi^2\,245\cdot246\,L^4}.
\]
This follows by squaring its normalized density, proportional to
$(1-|z|^2/L^2)_+^{122}$, and integrating on the four-ball.
At $\theta=m^{-2}$ its added excursion alone would require
$s_0 b=4d_{\rm sp}c_L>1564m^4$, already for a single spin slot.
Every spectral projection of any self-adjoint pair mass operator
has squared overlap with this fixed current bounded by $M_1$.
Thus even allowing a bound state at threshold does not make the
specified charge-one pair row large enough to supply the new pole.
The conclusion concerns the normalization predicted by this
particular determinant-class action and this fixed current. A
different charge-carrying source, additional physical states, or
an overlapping contribution of the opposite sign in the total
closed coefficient requires its own calculation.

\subsection{The opposite-energy Dirac pole and the charged recoil test}
\label{sec:six-dirac-block-uv}
Does the missing denominator in \eqref{eq:fifth-elastic-symmetric-uv}
come from the antiparticle of the same Dirac field? This can be checked
without assigning a second time ordering to an elastic intermediate
state. In the ordinary equal-time CAR source include all four bilinears
of the Dirac current, with the same invariant mark $F=h_0$ in each
band. Couple the transverse photon linearly to that current; let the
density in the Coulomb energy be its temporal component. The following
calculation concerns the connected coefficient of one occupied charge.
It does not assign a finite value to the independent bare vacuum
coefficient. The chosen plane and its vectorial pole prescription are
those of the preceding charged comparison.

Take an occupied particle $p$ of fixed spin, and let $l$ denote a
positive-energy antiparticle momentum. Put $E=E_p$, $w=E+E_l$,
$k=\mathbf p+\mathbf l$, $r=|k|$, $s=w^2-r^2$. Use unit-normalized
spinors as in Section~\ref{sec:fifth-elastic-uv} and set
$j^\mu_{\rm b}=\overline u(p)\gamma^\mu v(l)$, where the subscript
indicates the pair transition blocked by the occupied charge.
Define $S_0=\sum_{\sigma_l}|j^0_{\rm b}|^2$ and
$S_T=\sum_{\sigma_l}|\mathbf j_{{\rm b},T}|^2$.
Alongside the elastic terms of \eqref{eq:fifth-elastic-boost-uv}
the complete connected coefficient divided by $e^2$ is
\begin{align}
 B_{\rm CAR}(p)&=W_L(p)-G_+(p)+G_{\rm b}(p)-W_{\rm b}(p),
 \nonumber\\
 G_{\rm b}(p)&=\tfrac12\int_l
 \frac{|F(-s)|^2 S_T}{r(r+w)},\nonumber\\
 W_{\rm b}(p)&=\tfrac12\int_l
 \frac{|F(-s)|^2 S_0}{r^2} .
 \label{eq:six-dirac-block-uv}
\end{align}
Here $\int_l=\int d^3l/(2\pi)^3$. To see both signs, expand the
pair creation term $b^\dagger d^\dagger$ on the occupied state.
Compared to the vacuum with a spectator particle, the CAR deletes
precisely the amplitude creating the occupied mode. This removes a
negative virtual contribution and a positive instantaneous density
contribution. All other spectator terms cancel in the connected
coefficient. Equivalently these are its connected exchange
contractions. The new pair plus photon costs $w+r$, even though the
internal fermion pole is on the opposite frequency branch.
The finite-mode CAR calculation gives this identity before any
continuum integral; the following bounds give each connected term
its own continuum limit.

The crossed marks agree more strongly than their asymptotic powers
suggest. Write $p\cdot l=m^2c$, $c=1+t=\cosh\eta$, $a=\sinh\eta$.
The elastic spacelike mark at the same relative rapidity and the pair
timelike mark obey
\begin{equation}
 F(2m^2t)=F(-2m^2(t+2))=\chi_{248}(Lm\sinh\eta).
 \label{eq:six-dirac-crossmark-uv}
\end{equation}
This follows directly from $X[1+X/(4m^2)]$ in
\eqref{eq:nextsoft-product-uv}. The two bands therefore involve
exactly the same geometric weight.

\paragraph{What the blocked pair contributes at large boost}
Let $\gamma=E/m$, $v=|\mathbf p|/E$, and use the incident rest frame,
$u$ the cosine with the boost axis. Then
$E_l=\gamma m(c+vau)$, $s=2m^2(1+c)$ and
$d^3l/E_l=m^2a\,dt\,d\Omega$.
The positive-energy spin sum gives
$S_0=1-s/(4E E_l)$. Pair conservation gives
$w j^0_{\rm b}=k\cdot\mathbf j_{\rm b}$, hence
$S_0/r^2\le C/w^2$ even near $k=0$. These formulas yield
\begin{align}
 \lim_{\gamma\to\infty}\gamma W_{\rm b}
 &=\frac{m}{8\pi^2}\int_0^\infty |F(2m^2t)|^2
          [\eta-\tanh(\eta/2)]\,dt,\nonumber\\
 \lim_{\gamma\to\infty}\gamma G_{\rm b}&=0 .
 \label{eq:six-dirac-block-limit-uv}
\end{align}
For the first limit, $\gamma E_l/w^2\le1/(4m)$ is an integrable
dominator after including the radial A4 weight and $a\,dt$. At fixed
$t$, $r/w\to1$, $S_0\to1$, and the angular integral times $a$ is
$\int_{-1}^1 a(c+au)/(1+c+au)^2\,du
=\eta-a/(1+c)$.
For the second limit, on $r\ge w/2$ the same dominator works.
There $S_T\to0$ at fixed $t$ since
$S_T=(p\cdot l+2m^2)/(EE_l)-s S_0/r^2$.
On $r<w/2$ one has $s\ge3w^2/4$, $w\ge E$ and $w>2r$.
The A4 decay bounds the integral on this region by a constant times
$\int_0^\infty r\,dr\,(1+(E+r)/m)^{-6}/(E+r)$, which is
$o(1/\gamma)$ as $E\to\infty$. This argument also controls the
threshold $r=0$ of the blocked graph.

With $I_F=\int_0^\infty |F(2m^2t)|^2dt$, Eq.~\eqref{eq:fifth-elastic-recoil-uv} now gives
\[
 \gamma B_{\rm CAR}(p)=\frac{m}{8\pi^2v}
 \left[2 I_F\operatorname{arcosh}(\gamma)+O(1)\right].
\]
Thus restoring the actual antiparticle, with its crossed amplitude
and its fixed Coulomb contraction, leaves the growing rapidity term.
The calculation distinguishes the two time orderings in a spacelike
exchange between specified current rows from the energy denominators
of an occupied-charge self-energy. If a sharp charged mass shell is
claimed perturbatively, its Lorentz-invariant displacement must
give constant $\gamma\,\delta E$ at this order. A derivation from the
physical action must therefore exhibit additional connected contacts
or intermediate states, or derive the appropriate charged spectral
observable in a representation without that sharp pole. Merely
reinstating the negative-frequency Dirac pole cannot supply the
required term.

\paragraph{Covariant exchange and a real photon give different tests}
The joint preparation of Section~\ref{sec:joint2-charge-uv} provides
another calculation on one specified action: after varying both
charged endpoints the open kernel at fixed connection is exactly
$K_B^{-1}$. Write $f(X)$ for the transverse multiplier of $B$.
The following applies both to the radial A4 connection there and to
the choice $f(X)=h_0(X)$ with the same longitudinal connection.
A transverse virtual photon between two insertions is weighted by
$f(X)^2$ at its integrated Euclidean virtuality $X$. For example,
the one-loop rest mass coefficient of this open Dirac kernel, in the
usual positive mass-shift convention, is
\begin{equation}
 \frac{\delta m}{e^2}=2m\int\frac{d^4q}{(2\pi)^4}
 \frac{f(X)^2(X+2q_4^2)}{X(X^2+4m^2q_4^2)},\qquad X=q_E^2 .
 \label{eq:six-open-rest-mass-uv}
\end{equation}
Use transverse Euclidean gauge. Between on-shell spinors the
longitudinal numerator divided by the intermediate Dirac denominator
reduces, by the free Ward identity, to an odd term proportional to
$p\cdot q$; its symmetric integral vanishes. At rest the remaining
spin numerator and denominator, before taking the even part in $q_4$,
give $2(m+iq_4)/(X+2imq_4)$ in addition to $f(X)^2/X$.
This proves \eqref{eq:six-open-rest-mass-uv} as a Euclidean
on-shell boundary value. Near $q=0$ the radial measure of the even
integrand is bounded by $C\,d|q|$ after the angles: use
$\int_{-1}^1 dz/(|q|^2+4m^2z^2)\le C/(m|q|)$.
At infinity the A4 tail controls the logarithmic radial power.
Euclidean rotational covariance fixes the scalar mass coefficient
by continuation in $p_E^2$ from above $-m^2$.
This virtual result uses the frequency-dependent connection; it
does not convert an antiparticle energy denominator into $r-d$.

For an emitted photon one instead has exactly $f(0)=1$ on the whole
null ray. Test the positive-energy charged endpoint at total spatial
momentum zero. Into an electron of momentum $-k$ and one physical
photon of momentum $k$, the free positive-energy projection on the
incoming open Dirac kernel gives $1/(E_r+r-m)$ with
$E_r=\sqrt{m^2+r^2}$. Its signed transverse Dirac row has summed
squared norm $(E_r-m)/E_r$. Thus the positive one-photon cut has
radial weight
\begin{equation}
 \frac{dN_{e\gamma}}{dr}
 =\frac{e^2}{4\pi^2}
   \frac{r(E_r-m)}{E_r(E_r+r-m)^2}
 \sim\frac{e^2}{16\pi^2 r}.
 \label{eq:six-open-radiative-uv}
\end{equation}
The equivalent is for $r\to\infty$; the density is $O(r)$ at zero. This hard logarithm already occurs at rest.
If the exact charged endpoint has a positive spectral measure of
canonical total weight one, and all but $O(e^2)$ of that weight stays
in a fixed low-energy window containing the free state, a positive
continuum above that window must have bounded total weight after division
by $e^2$. Applying Fatou's lemma on increasing finite energy windows
to the proposed leading cut contradicts
\eqref{eq:six-open-radiative-uv}. This test specifies the
normalization and weak-coupling assumption it uses. Finite-time
smeared coefficients can still exist.

\paragraph{A curvature contribution from the two charged endpoint operators}
There is an explicit change of the measured boundary state that
repairs this positive cut on the same action. On the chosen plane
let $\pi=i\partial-eB$, $H_{\rm s}=\pi^2-m^2$ and
$H_{\rm D}=(\gamma\cdot\pi)^2-m^2$ in the on-shell convention.
Use the original ball average $g(z)=\chi_{248}(L^2z)$ and multiply
the measured charged endpoint by
\[
 O_B=I+g(H_{\rm s})-g(H_{\rm D}),
\]
with the corresponding adjoint at the other end. These covariant
operators have exactly the same free value, so $O_0=I$ on every
free momentum. The difference measures their spin curvature; its
first divided difference is analytic even at equal arguments.
The underlying charged determinant is unaltered.

Let $p'$ be the on-shell outgoing charge, $k$ the real photon,
$P=p'+k$, and $z=P^2-m^2=2p'\cdot k>0$. For $k\cdot\epsilon=0$
the linear variation of $H_{\rm s}-H_{\rm D}$ has matrix
$e\slashed k\slashed\epsilon$. Indeed the two squared momenta
give $-e(P+p')\cdot\epsilon$ and
$-e(\slashed P\slashed\epsilon+\slashed\epsilon\slashed p')$.
The variation of $O_B$ supplies the same matrix times
$[g(z)-1]/z$. The ordinary open row is
$e(2p'\cdot\epsilon+\slashed k\slashed\epsilon)u(p')/z$.
Choose the two transverse polarizations in the rest frame of $P$;
then $p'\cdot\epsilon=0$. The two terms sum to the original
row times $g(z)$. Consequently \eqref{eq:six-open-radiative-uv}
gains the squared geometric mark $|g((E_r+r)^2-m^2)|^2$. It is
now integrable at both endpoints. In fact the relative coefficient is
fixed within this pair of endpoint operators. Multiplying their
difference by $b$ leaves the hard multiplier $1-b+b g(z)$; its
squared integral against $dz/z$ is finite precisely for $b=1$.
The additional first-order
contact lacks an incoming charged pole and preserves the original
double charged-pole residues of the one-photon vertex.
This specifies a different measured observable on the same
background action; it neither moves the Hamiltonian mass spectrum
nor changes the photon determinant. Positivity of its complete
multi-insertion measure and the coincident bounds for all its
higher boundary variations still require derivation.

\paragraph{The two-photon contact and a boundary state for real radiation}
There is an instructive test of $O_B$ at the next order. On an open
Dirac wave with its outgoing charged leg on shell, the background
equation gives $H_{\rm D}\psi_B=0$. Thus $O_B\psi_B=g(H_{\rm s})\psi_B$.
Take two nonparallel real photons at total spatial momentum zero and
choose both spatial polarizations normal to the plane of their
momenta. All scalar one-photon insertions vanish: each scalar
momentum lies in that plane. The $e^2 B^2$ term in $H_{\rm s}$
survives. Consequently its exact two-photon tree row is
\[
 g(z)\mathcal M_2^{\rm D}
 +2e^2(\epsilon_1\cdot\epsilon_2)\frac{g(z)-1}{z}u(p'),
\]
where $z=(p'+k_1+k_2)^2-m^2$. Here $\mathcal M_2^{\rm D}$ includes both Dirac orderings with unit
spinor normalization. On an angular cone where all pairwise directions
are separated and all three energies are comparable to $R$, every
intermediate charged denominator is of order $R^2$, each numerator
is $O(R)$, and $z$ is of order $R^2$. Thus
$\mathcal M_2^{\rm D}=O(e^2/R^2)$. The positive rest spin projection
of $u(p')$ has norm bounded below for a suitable spin. The contact
therefore leaves a row of magnitude at least $ce^2/R^2$ on this cone
at large $R$. The positive two-photon measure
$d^3k_1d^3k_2/(r_1r_2)$ on a dyadic shell grows as $R^4$, giving a
logarithm in the total order-$e^4$ cut. The single repaired
radiative row does not pay for this compulsory boundary contact.

The spatial Gauss primitive already used in
\eqref{eq:next3-charged-connection-uv} gives a second boundary
choice on the same action. Define $\chi_B=\Delta_{\rm sp}^{-1}
\nabla\cdot\mathbf B$ on nonzero spatial modes, with the inherited
homogeneous gauge convention. Let $H_{\rm flat}=(i\partial-e\,d\chi_B)^2-m^2$
and replace $H_{\rm s}$ by $H_{\rm flat}$ in $O_B$. This is a flat
connection with the usual Abelian gauge variation. For any collection
of transverse radiation waves on the chosen time slice $\chi_B=0$.
Since $H_{\rm D}\psi_B=0$, the new endpoint on that wave is exactly
$g(-\partial^2-m^2)\psi_B$. At every tree multiplicity the complete
radiation row therefore satisfies
\begin{equation}
 \mathcal M_n^{\rm flat}(P;p',k_1,\ldots,k_n)
 = g(P^2-m^2)\,\mathcal M_n^{\rm D}(P;p',k_1,\ldots,k_n).
 \label{eq:six-open-radiative-all-uv}
\end{equation}
In the total rest frame $W=P^0=E_{p'}+\sum r_i$ controls every hard
photon simultaneously, including nearly collinear ones.

A direct bound also gives the scope of this statement. Use infrared
test factors $r_i/(m+r_i)$ in the amplitude, which leave the hard
power unchanged. Each ordered Dirac numerator is $O(W)$; the charged
denominator after a nonempty subset $I$ satisfies
$(p'+\sum_{i\in I}k_i)^2-m^2\ge (m^2/W)\sum_{i\in I}r_i$.
The amplitude before the endpoint filter is therefore bounded by
$C_n(W/m)^{2n}\prod_i r_i^{-1}$. The remaining product measure in
its squared norm, including the tests, is bounded by
$C_n\prod_i r_i\,dr_i/(m+r_i)^2$ after angles.
Its mass up to energy $W$ grows at most logarithmically in each
variable. The original ball gives $|g(W^2-m^2)|^2=O(W^{-498})$
at large $W$, sufficient by this bound for all $n\le124$.
There is also a same-measure endpoint preparation for arbitrary fixed
$n$: weight the uniform unit A4 ball in the boundary matrix element
by the squared normalized state proportional to
$\exp[-1/(1-|u|^2)]$ in its interior. Its extension by zero is smooth
and compactly supported, and its characteristic function decreases
faster than every power by integration by parts. Using this function
as $g$ in \eqref{eq:six-open-radiative-all-uv} proves convergence
for every finite $n$ with the stated infrared tests. This choice
changes the measured boundary state. Neither the photon determinant
nor the charged Hamiltonian has been modified. The result covers
exclusive tree radiation on this slice. It leaves local boost
covariance, positive reconstruction of the complete boundary
observable, and interactively sewn loops to be tested.

\subsection{The exact photon record at a cut of a path integral}
\label{sec:sixth-photon-seam-uv}
Fix one labelled, post--GTL sector plane and its chosen Euclidean
clock. Take the transverse Maxwell field with the spatial A4
preparation $h(\mathbf k)=\chi_{248}(L|\mathbf k|)$ of the angular
comparison. The following calculation integrates that field exactly.
It identifies the information that has to pass across a clock slice
when amplitudes are composed. Write $k=|\mathbf k|$,
$d\pi(\mathbf k)=d^3\mathbf k/(2\pi)^3$, and let $P(\mathbf k)$
be the transverse projector. The free transverse covariance before
the two incidence factors is $P(\mathbf k)e^{-k|t-s|}/(2k)$, the
inverse clock Fourier transform of $P/(\omega^2+k^2)$.

For a real spatial current $j=j_-+j_+$ with records on either side
of $t=0$, write $V(j)$ for its Maxwell variance, including $h$ on
both ends. The Gaussian influence is $I_e(j)=\exp[-e^2V(j)/2]$.
For each side form the transverse vector
\begin{equation}
 b_\pm(\mathbf k)=\frac{h(\mathbf k)}{\sqrt{2k}}
 \int_{\mathbb R_\pm}e^{-k|t|}P(\mathbf k)
                          j_\pm(dt,\mathbf k).
 \label{eq:sixth-seam-data-uv}
\end{equation}
The inner product below uses $d\pi$ and its real part on real
currents. Since $|t-s|=|t|+|s|$ on opposite sides, the mixed
factor is exactly
\[
 \frac{I_e(j_-+j_+)}{I_e(j_-)I_e(j_+)}
       =\exp[-e^2\operatorname{Re}\langle b_-,b_+\rangle].
\]
Equivalently condition on the free photon at the cut. Its contribution
to a half-current is an isonormal real Gaussian $X(b_\pm)$, with
covariance $\operatorname{Re}\langle b_-,b_+\rangle$. Within each half,
the remaining covariance has the scalar factor
\[
      \frac{e^{-k|t-s|}-e^{-k(|t|+|s|)}}{2k}.
\]
This is the Dirichlet covariance, obtained by subtracting the
conditional mean contribution. The two Dirichlet fields are
independent. Multiplying their conditional characteristic functions
and integrating their common $X$ reconstructs $I_e$ exactly.
Conjugating and clock-reflecting a half-current consequently yields
a squared conditional norm in this free photon sector.

The vector in \eqref{eq:sixth-seam-data-uv} is necessary to
distinguish pasts by all future transverse probes, once their
individual self weights have been recorded. Indeed, for $e\ne0$
equality of all mixed factors gives orthogonality of the difference
to every attainable $b_+$. Choose a nonnegative smooth time pulse
of nonzero integral supported at positive times. Its Laplace
transform is strictly positive on $k>0$. Multiplying by arbitrary
smooth transverse momentum packets away from $k=0$ and from the
isolated radial zeros of $h$ gives a dense set in the transverse
one-photon space. Disjoint packets in a fixed such annulus already
give an infinite-rank mixed form. No finite list of linear current
moments specifies it.

There is a path-level witness. Put an oriented spatial circle of
radius $r$ at clock time $-a$, $a>0$, and its reflected copy at $a$.
At momentum $k e_1$ in the circle plane the $e_2$ component of its
line current has modulus $2\pi r|J_1(kr)|$, by direct angular
integration. Hence its boundary vector has strictly positive norm.
Reversing the past circle leaves its own quadratic weight and
length unchanged while reversing that vector. The two influences
with the same future circle have ratio $\exp[-2e^2\|b_-\|^2]$,
with the equally oriented case in the numerator. One may append
the circles as past probes before the same terminal charge and
velocity record. When a path turns back in clock time, one groups
its current by actual clock half-space. Splitting it once at an
arbitrary value of its path parameter need not produce these
two half-currents.

For an oriented path of finite spatial variation $\ell$ the full
phase variance satisfies
\[
  0\le V(j)\le C_h\ell^2,\qquad
  C_h=\int\frac{|h(\mathbf k)|^2}{2k}\,d\pi(\mathbf k)<\infty .
\]
The estimate follows by absolute integration of the two spatial
line elements; the angular projector has norm one. The small-$k$
radial measure is proportional to $k\,dk$, and the real A4 tail
controls infinity. The same absolute majorant works before the
clock frequency integral because
$\int d\omega/[2\pi(\omega^2+k^2)]=1/(2k)$. Thus every finite
collection of exchanged transverse photons on fixed finite-variation
histories is integrable over all its photon momenta, including
self contractions. At order $2n$ the integrated coefficient of one
cycle is bounded in absolute value by
\[
 \frac{C_h^n}{2^n n!}\int\ell^{2n}\,d|\mu|,
\]
where $\mu$ is the complete charge-cycle measure before photon
integration. For the continuous-turning heat trace per unit volume,
$\ell\le t$ and \eqref{eq:next3-angular-return-uv} supplies all
these moments at the origin; the positive mass supplies them at
large duration. This derives finite transverse self sewing directly
for that charged generator. To apply the estimate to unrestricted
physical GCM histories one must derive the same incidence factors
and the necessary variation moments for its determinant, contacts,
boundary measure, and joint short-duration faces. Positivity of the
conditioned photon factor by itself imposes no such moment bound
on the charged measure.

\paragraph{The infinite excursion potential under the fluctuating photon field}
Brownian paths have infinite variation. Their Wilson coupling under
the Gaussian photon can nevertheless be constructed without
restricting momentum or deleting any duration. Use the positive
intensity $d\eta$ of \eqref{eq:fifth-excursion-levy-uv}, and keep
the sign of the potential separate. For a pinned Euclidean Brownian
excursion $b$ of heat duration $T$ at $x$, let $u_{x,b}$ be its
transverse current in the real Gaussian one-photon space. In four-
momentum its representative is the spatially projected Fourier
line current multiplied by $h(\mathbf k)/(\omega^2+k^2)^{1/2}$, with
measure $d\omega\,d\pi/(2\pi)$. Its squared norm is the conditional
Wilson-phase variance before multiplication by $e^2$.

In each spatial polarization perpendicular to $\mathbf k$, the
stochastic line integral is a centered Gaussian conditional on the
longitudinal spatial and clock bridge coordinates. This follows
from independence of the remaining transverse bridge component.
The deterministic integrand in this conditional integral has
modulus one. The covariance of the pinned increments subtracts the
constant-time projection, giving variance at most $2T$ in each of
the two directions. Thus
\[
  \mathbb E_b\|u_{x,b}\|^2\le 4C_hT .
\]
In particular $u_{x,b}$ exists in the Gaussian Hilbert space for
almost every bridge. The isonormal Gaussian evaluation $X(u_{x,b})$
is defined in the joint path and field measure by mean-square
approximation. For any complete orthonormal basis, Parseval's identity and the
variance bound give convergence of the partial isonormal sums in
the joint mean-square norm.

For the real fluctuating photon, orientation reversal gives the
nonnegative potential magnitude
\[
 Y_A(x)=\int d\eta(T)\,
       \mathbb E_b[1-\cos(eX(u_{x,b}))\mid A] .
\]
The notation denotes conditional averaging over an independent
excursion with $A$ held fixed. Since $1-\cos v\le v^2/2$,
Tonelli gives $\mathbb E_A Y_A(x)\le2e^2 C_h\int T\,d\eta(T)$,
which is finite. Hence the signed potential $-Y_A(x)$ exists
almost surely at each fixed $x$.

More is true on every finite central history $x_s$, $0\le s\le t$,
independent of the photon before weighting. Condition on that history.
For any $z\ge0$ its signed-potential exponential has a finite
Maxwell average, uniformly in the choice of history. Indeed split
the duration integral at $\delta>0$. The intensity above $\delta$
is finite and $1-\cos v\le2$. Below $\delta$ form the positive
trace-class operator
\begin{align*}
 M_\delta
   &=\int_0^t ds\int_0^\delta d\eta(T)\,
       \mathbb E_b\,|u_{x_s,b}\rangle\langle u_{x_s,b}|,\\
 \operatorname{Tr}M_\delta
   &\le 4t C_h\int_0^\delta T\,d\eta(T).
\end{align*}
The last bound tends to zero at fixed $t$. Choose $\delta$ so that
$ze^2\operatorname{Tr}M_\delta<1$. Applying the quadratic cosine
bound only to this part gives
\begin{equation}
 \mathbb E_A e^{z\int_0^tY_A(x_s)ds}
 \le
 \frac{\exp[2zt\,\eta([\delta,\infty))]}
      {\det(I-ze^2M_\delta)^{1/2}} .
 \label{eq:sixth-excursion-exponential-uv}
\end{equation}
To see the denominator, diagonalize $M_\delta$. The dominating
Gaussian quadratic form is a sum of its nonnegative eigenvalues
times independent squared standard normals. The product of their
Gaussian integrals converges because the eigenvalues are summable
and $ze^2\|M_\delta\|<1$. Both ranges of the original intensity enter
the bound.

For $t$ in a bounded interval one choice of $\delta$ serves the
entire interval. The determinant bound is uniform there:
$-\log(1-v)\le v/(1-v)$ controls its logarithm by the trace bound.
The minimal real Wilson phase of the central path has modulus one.
Consequently the angular small-return estimate
\eqref{eq:next3-angular-return-uv} still controls the absolute
vacuum trace at short central duration after the fluctuating field
and the signed infinite excursion potential are included, at every
fixed real coupling. This is an actual short-duration path-integral
estimate for the stated angular generator. Growth at arbitrarily
large central duration, other charged sectors and a positive
physical cut require additional information.

\subsection{What groupoid convolution retains between two curvatures}
\label{sec:seven-groupoid-uv}
Work on one locked Euclidean sector-plane and on the positive even
part of the represented curvature pairing. For the regular groupoid
trace the coefficient of $F^\ast\!\circledast F$ at the unit of an object $x$ is
the integral of $F(\gamma)^\dagger F(\gamma)$ over arrows with
source $x$, in the corresponding Haar-fiber measure. Thus replacing
each arrow kernel by $e^{iQ\cdot d(\gamma)}F(\gamma)$ leaves its
inclusive quadratic coefficient unchanged. This calculation uses
the same raw adjoint at the return; it applies even if $d$ is not
additive. A physical graded pairing or a projection mixing arrows
requires its actual contracted product.

There is a further restriction on an additive displacement. On a
transitive chart choose arrows $g_x:o\to x$ and write every arrow
$x\to y$ as $g_y h g_x^{-1}$ with $h$ an isotropy loop at $o$.
A continuous additive displacement maps compact isotropy trivially
into $\mathbb R^4$. More generally, with orthogonal transport $R$
suppose $d(gh)=d(g)+R(g)d(h)$. For compact isotropy its normalized
Haar mean $v=\int d(h)dh$ obeys $d(h)=v-R(h)v$ by left invariance.
Putting $b(x)=d(g_x)+R(g_x)v$ gives
$d(\gamma)=b(y)-R(\gamma)b(x)$. With incoming momentum
$R(\gamma)^{\mathsf T}Q$ the phase is endpoint conjugation by
$e^{iQ\cdot b}$. A cyclic trace eliminates a common transported
twist on a closed momentum cycle. Distinct routings require
their own calculation. A Gibbs state with diagonal
endpoint density still eliminates this conjugation. A noncommuting
density can retain it but changes cyclic gluing: for an invertible
finite density $H$, moving $A$ across
$\operatorname{Tr}(HAB)$ inserts $HAH^{-1}$ after $B$.
The modular transport must then belong to the physical sewing law.

For example, for the heat transit $e^{t\Delta}$ on a flat
four-plane, the Gibbs functional with momentum density
$e^{-\beta |k|^2}$, per unit volume, gives the exact ratio
\[
 \frac{\int e^{-\beta |k|^2-2t|k-Q|^2}d^4k}
      {\int e^{-(\beta+2t)|k|^2}d^4k}
 =\exp[-2t\beta X/(\beta+2t)],\quad X=|Q|^2 .
\]
This preserves a separation in the tested amplitude. If the ratio
is the entire curvature kinetic coefficient, its local potential
inverse grows exponentially at large real Euclidean momentum and
has no tempered spectral reconstruction. The Gibbs parameter here
belongs to the internal heat preparation. A covariantly moving
density would also contribute its own source derivatives.

A noncompact winding subgroup is a different case. For isotropy
$\mathbb Z$ with displacement $n\ell$, $\ell\ne0$, take a finite
word sum with coefficients $a_n$. In a single character of angle
$\vartheta$ the squared amplitude is
$|\sum_n a_ne^{in(\vartheta+Q\cdot\ell)}|^2$. It is periodic in
$Q\cdot\ell$ with mean $\sum_n|a_n|^2$. Integration over all
characters gives this constant at each $Q$. Thus neither the
regular representation nor this fixed irreducible quotient yields
uniform hard decay for a nonzero finite winding sum.

\paragraph{Two placements of the curvature in the same heat representation}
At zero background curvature, a single nonlocal pairing between
local curvature variations still gives
$\operatorname{Tr}(M_f\mathbb P_*M_f)=\operatorname{Tr}(\mathbb P_*M_{f^2})$
on its trace domain. Separation requires an off-diagonal curvature
variation or additional boundary transport. Take the pair-groupoid convolution on the Euclidean
four-plane and a positive self-adjoint transit
$A=(m^2-\Delta)^{-p}$, $m>0$, $1<p<2$. This comparison specifies
a positive superposition of heat transits with density
$t^{p-1}e^{-m^2t}/\Gamma(p)$; it changes no GTL or shell measure.
Let $a(y)$ be its convolution kernel and $c=\int |a(y)|^2d^4y$.
For smooth real compactly supported curvature $f$, the raw
adjoint square $\operatorname{Tr}[(M_f A)^\dagger M_f A]$
is just $c\int f^2$. Instead placing one curvature on each visit
gives the positive square
\[
 \operatorname{Tr}[(A^{1/2}M_f A^{1/2})^2]
 =\int f(x) f(x+y)|a(y)|^2d^4x\,d^4y .
\]
The marked operator on the left is Hilbert--Schmidt: in momentum
space the convolution of two copies of $(m^2+|k|^2)^{-p}$ is
bounded by its squared $L^2$ norm. The identity follows by
convolution and approximation. After normalizing
at zero momentum, the transverse curvature multiplier is
$h_p(X)=c^{-1}\widehat{|a|^2}(Q)$, $X=|Q|^2$.
Covariantizing $A$ by the magnetic Laplacian leaves this flat
Hessian unchanged; higher vertices include its variations.
The resulting local potential covariance is $D_p(X)=1/[Xh_p(X)]$.

At $p=3/2$ the heat integral gives
$a(y)=e^{-m|y|}/(4\pi^2|y|)$ and $c=1/(32\pi^2m^2)$.
Angular Fourier integration of $|a|^2$ gives
$h_p(X)=8m^2[1-2m/\sqrt{4m^2+X}]/X$. Hence
\[
 D_p(X)=\frac1X+\frac1{8m^2}
      +\frac{1}{4m(\sqrt{4m^2+X}+2m)} .
\]
The last term has positive Stieltjes density
$\sqrt{s-4m^2}/(4\pi m s)$ on $s>4m^2$, as follows by setting
$s=4m^2+u^2$ in its dispersion integral. The massless pole and
local contact thus come from the same kinetic action as the
continuum. Its separated-return cost
$c^{-1}\int |a(y)|^2|y|^{-2}d^4y$ still diverges logarithmically.
The contact in $D_p$ also survives unsmeared photon sewing.

The tradeoff is sharp within the stated fractional family.
The heat integral gives $a(y)\asymp |y|^{2p-4}$ at zero and
exponential decay at infinity. Thus the separated cost is finite
exactly for $p>3/2$. In that range $h_p(X)\sim C X^{2-2p}$:
rescaling its momentum convolution by $\sqrt X$ leaves an integrable
majorant $|z|^{-2p}|z+\hat Q|^{-2p}$ in four dimensions.
Hence $D_p(X)\sim C^{-1}X^\alpha$ with $0<\alpha=2p-3<1$.
For any $N\ge1$ the signed divided difference
$(-1)^N[1,2,\ldots,N+1]x^\alpha$ is strictly negative by the
sign of its $N$th derivative. Scaling the nodes by $R$ gives
the same strict sign for $D_p$ at large $R$.
A positive spectral representation with finitely many local
subtractions gives a nonnegative signed divided difference when
$N$ exceeds the subtraction degree. Thus the improved coincidence
moment in this range is incompatible with that local photon
spectrum. These conclusions apply to the stated curvature placements and transits.

There is an instructive exception outside that range. The pinned
gamma excursion itself provided $p=5/2$ as a charged transit.
Used here between curvatures it has $a(y)\propto e^{-m|y|}$.
The normalized displacement density has value $4m^4/(3\pi^2)$
at zero and inverse-square moment $2m^2/3$. With $b=4m^2$ its
Fourier transform is $h(X)=(1+X/b)^{-5/2}$. The ensuing
$D(X)=(1+X/b)^{5/2}/X$ has a unit massless pole and an exactly
positive twice-subtracted continuum:
\[
 \begin{aligned}
 D(X)&=X^{-1}+5/(2b)+15X/(8b^2)\\
 &\quad+X^2\int_b^\infty\frac{(s/b-1)^{5/2}}
                          {\pi s^3(X+s)}\,ds .
 \end{aligned}
\]
The density follows from the cut of the displayed fractional power,
and two subtractions suffice by its growth. This gives a positive
separated-time spectrum and a finite cost for each bounded spectral
band in the inverse-square test. The full spectral weight has
divergent first moment. The density at zero handles an undifferentiated
delta only; its radial cusp cannot test a Laplacian of a delta.
Indeed the full momentum contraction has $h(X)D(X)=1/X$. Sewing two
curvatures with their own kinetic transit contributes $Xh(X)D(X)=1$
per transverse mode, before four-dimensional momentum integration.

\paragraph{Two charged rest spaces and the cost of a small deflection}
On a locked plane use \eqref{eq:uvpeP-groupoid} with $K=R=I$ on three coordinates. The underlying Gaussian records
$Z_0,Z_a$, normalized to unit marginal variance at two visits
separated by $a>0$, have cross covariance $\rho I$, $\rho=e^{-a}$.
Condition on future unit charged velocities $u,v$ with
$u\cdot v=\cosh\eta$ in the mostly-minus convention. Embed the
records at scale $L$ into their respective rest spaces, identifying
triads by the boost in the $u,v$ plane. This triad-independent boundary map is a representation choice.

In the $u$ rest frame take $v=(\cosh\eta,\sinh\eta,0,0)$ and write
$x=(Q_1,Q_\perp)$ and $y=(\cosh\eta\,Q_1-\sinh\eta\,Q_0,Q_\perp)$.
The phase of the difference of the physical displacements has
characteristic function $h(Q)=\exp[-L^2(|x|^2+|y|^2-2\rho x\cdot y)/2]$.
For real four-momentum the exact squared integral is
\begin{equation}
 \int_{\mathbb R^4}|h(Q)|^2\,d^4Q
 =\frac{\pi^2}{2L^4\sinh\eta\,(1-\rho)\sqrt{1-\rho^2}} .
 \label{eq:seven-deflection-rank-uv}
\end{equation}
Indeed $(Q_0,Q_1)\mapsto(x_1,y_1)$ has absolute Jacobian
$\sinh\eta$. The longitudinal Gaussian has determinant $1-\rho^2$,
and each of the two transverse directions has coefficient $2(1-\rho)$.
Two distinct rest spaces thus resolve four real momentum components;
equal velocities leave one component unresolved even for independent
records. Here the four Fourier components are real; continuation to a complex
pole uses the declared boundary branch on the sector plane.

The action also permits a residual triad rotation $O$ in the transport,
replacing $x\cdot y$ by $x\cdot Oy$. Let $S$ be the symmetric
part of its compression to the common two-dimensional transverse
plane. In \eqref{eq:seven-deflection-rank-uv} replace $1-\rho$
by $\sqrt{\det(I-\rho S)}$. To verify this, integrate the
six-dimensional Gaussian in dimensionless $x,y$ against
$\delta^{(2)}(x_\perp-y_\perp)$. Its full integral is
$\pi^3/(1-\rho^2)^{3/2}$, and the transverse difference of the
normalized Gaussian has covariance $(I-\rho S)/(1-\rho^2)$.
Since $-I\le S\le I$, aligned transport gives the largest cost.
No such rotation removes $1/\sinh\eta$ at fixed positive $a$.

A covariant conditional law permits an exact joint short-duration
test. As a comparison let the velocity diffuse for the same $a$ on
the future unit hyperboloid with generator $\Delta_{H^3}$, starting
at the given $u$, independently of the Gaussian normal process.
The transition density at distance $\eta$, against
$4\pi\sinh^2\eta\,d\eta$ after angles, is
\[
 (4\pi a)^{-3/2}e^{-a-\eta^2/(4a)}\frac{\eta}{\sinh\eta}.
\]
Conjugating the radial Laplacian by $\sinh\eta$ verifies the heat equation. Thus
$\mathbb E_a[1/\sinh\eta]=e^{-a}/\sqrt{\pi a}$ by integrating
$\eta e^{-\eta^2/(4a)}$. The mean of
\eqref{eq:seven-deflection-rank-uv} grows precisely as $C a^{-2}$
at zero. With the dimensionless A4 duration $a=\tau/L^2$ of
\eqref{eq:joint2-duration-uv}, its mean is finite precisely for
$p>2$. One record ($p=2$) gives a logarithm and two ($p=4$)
give a finite integral.

With fresh deflections, clocks and Gaussian pairs the integrals
factor for arbitrarily many independent circulations. Sharing just
the clock among $b$ otherwise independent copies requires $p>2b$.
Sharing the same deflection as well costs
$\mathbb E_a[(\sinh\eta)^{-b}]$: at every fixed $a>0$ the radial
probability is proportional to $\eta^2d\eta$ near zero, so three
circulations already diverge logarithmically. 

There is a different rank loss when the actual adjacent visits of
one path share their intervening endpoint. For $b$ consecutive
differences the map sends $X_j\in u_j^\perp$, $0\le j\le b$, to
$(X_1-X_0,\ldots,X_b-X_{b-1})$. Its kernel is the common displacement
in $\bigcap_j u_j^\perp$. Writing $r=\dim\operatorname{span}\{u_j\}$,
its rank is $3b-1+r$, so there are $b+1-r$ unresolved Fourier
directions. Explicitly, each relation $\sum_{i=0}^b\lambda_i u_i=0$
gives the invisible routing $Q_j=\sum_{i<j}\lambda_i u_i$,
$1\le j\le b$, by telescoping the phase. Even with nonparallel
consecutive velocities, four
unrestricted four-momenta cannot all be controlled this way.
For routed graph momenta the obstruction applies to the intersection
of the routing image with this nullspace.

A fourth normal coordinate can supply a fluctuating displacement along
the chosen carrier. It is internal to this readout and does not identify
a geometric-time coordinate with a physical clock reading. In the same quadratic
action use four-dimensional records $Z_j=\rho_j Z_{j-1}
+\sqrt{1-\rho_j^2}\,\xi_j$, $\rho_j=e^{-a_j}$, with independent
standard Gaussian innovations. Put $X_j=LB_j Z_j$ with proper
Lorentz boosts $B_j$ to the visited velocities. Conditional on $Z_0$,
the map from innovations to differences is block triangular with
diagonal $LB_j\sqrt{1-\rho_j^2}$. Thus for $b$ circulations
\[
 \int |h(Q_1,\ldots,Q_b)|^2d^{4b}Q
 \le (\pi^2/L^4)^b\prod_j(1-\rho_j^2)^{-2}.
\]
Conditioning gives equality before averaging the initial phase.
A future timelike boundary readout is also possible: replace
$(z_0,\mathbf z)$ before each boost by
$(|z_0|+|\mathbf z|,\mathbf z)$. Almost surely $X_j=T_jV_j$ with $T_j>0$ and $V_j^2=1$, $V_j^0>0$.
The map has unit absolute Jacobian almost everywhere and multiplicity two. Plancherel and
Cauchy--Schwarz on inverse images multiply the bound by at most
$2^b$. With fresh $p=4$ durations both readouts work at every finite
$b$, even with shared intermediate endpoints. Their physical
selection remains to be derived.

For the three-coordinate test at a physical pair opening
$Q=m(u+v)$ one has $x=-y$ with length $m\sinh\eta$. For aligned transport the
squared phase average gives
$\exp[-2L^2m^2(1+\rho)\sinh^2\eta]$, strictly positive at each
fixed pair mass and invariant under a common boost. It leaves the
positive Gauss boost integral on a bounded occupied mass band
divergent. Nor can the photon alone select a substitute rest
space: for fixed null $k$, the null rotations fixing it translate
the transverse components of a displacement by $(k\cdot X)b$.
On $k\cdot X\ne0$ the transverse components divided by $k\cdot X$
would have a translation-invariant probability law. Thus $k\cdot X=0$
almost surely. The same rotations then shift the component along $k$
by $b\cdot X_\perp$, forcing $X_\perp=0$ by the same argument.
A covariant single-null-momentum displacement law is supported along
$k$ itself and has $e^{ik\cdot X}=1$. Real charged velocities can
provide the additional directions; the sewing law still has to
retain them at the physical insertions.

\subsection{A curvature contact on the actual shell and its radial groupoid}
\label{sec:seven-radial-uv}
The coarea calculation in \eqref{eq:monomial-owner-uv} supplies a
uniform point $u$ of the unit $n$-ball, $n=248$. The subgroup of
rotations preserving its distinguished axis acts on each sphere
$|u|=r$ and fixes $r$. Write $d\pi(r)=n r^{n-1}dr$. Consider a
particular representation on one Euclidean locked plane: project four
normal axes onto its four curvature coordinates with geometric length
$L$, and average the direction by normalized rotation Haar measure
at each curvature mark. This prescription keeps the radius shared
between the two ends of a quadratic insertion. The two angular
averages are taken inside the observable, before it is squared. A four-coordinate
projection is part of this comparison. Averaging only spatial axes
would leave the independent clock frequency in a local contact.

At radius $r$ the displacement density is
$a_r(y)=(Lr)^{-4}a(y/(Lr))$, where
$a(z)=c_n(1-|z|^2)_+^{(n-6)/2}$ and $\int a=1$.
This follows by projecting surface measure on $S^{n-1}$ and
integrating the remaining $n-4$ coordinates. In particular
$a\in H^1(\mathbb R^4)$. For a physical photon with an unsubtracted
positive transverse spectral measure $d\nu(s)$ consider the complete
smeared curvature trace of Section~\ref{sec:six-physical-uv}. It is
\[
 \frac32\int d\nu(s)\int\frac{d^4Q}{(2\pi)^4}
       |\widehat a_r(Q)|^2\frac{X}{X+s},\qquad X=|Q|^2 .
\]
This expression keeps the contact at zero displacement. At fixed
$r>0$ the inner integral times $1+s$ is bounded above and below by
strictly positive constants in fixed mass units. For the upper bound
use $(1+s)X/(X+s)\le1+X$ and the $H^1$ norm. For the lower bound
integrate on a finite annulus where $\widehat a_r$ is nonzero. Thus
the exact test is $\int(1+s)^{-1}d\nu(s)<\infty$. It can succeed
with infinite first spectral moment. Extra polynomial contacts must
be evaluated on their own test derivatives.

For a free Maxwell field the single component
$\overline F_{12,r}(x)=(a_r*F_{12})(x)$ is a real Gaussian with
variance $\|a\|_2^2/[2(Lr)^4]$. Indeed its momentum tensor is
$(Q_1^2+Q_2^2)/X$, with angular mean $1/2$. Reusing one radius
through every insertion has an exact cost. The characteristic function
of the dimensionless smeared flux is
\begin{align*}
 \mathbb E_{r,A}\exp(ieL^2\overline F_{12,r})
   &=p\int_1^\infty t^{-p-1}e^{-g t}dt,\\
 p&=n/4=62,\qquad g=e^2\|a\|_2^2/4 .
\end{align*}
For $j<p$ its $j$th Taylor coefficient in $g$ is
$(-1)^j p/[j!(p-j)]$. The $p$th derivative at $g>0$ is
$(-1)^p p\int_1^\infty e^{-gt}dt/t$, which grows in magnitude
as $p\log(1/g)$. The logarithm of the characteristic function has
the same first nonanalytic term since its constant is one and its
lower derivatives are finite. This is a complete gauge-invariant
Gaussian source with an explicitly connected cumulant. A bounded
source at every real coupling can have a divergent perturbative
coefficient. Here each repeated quadratic contact costs four
inverse powers of the radius; the inverse-square criterion in
\eqref{eq:six-relative-ranks-uv} concerned the separated massive
part. No coincident delta can be dropped on the grounds that the
two displacements differ almost surely.

\paragraph{When radial propagation suffices in a cyclic determinant}
The radius can instead propagate between successive visits. The
following calculation distinguishes a genuine radial quantum state
from the trace over all angular states. Take the Neumann Laplacian
$H$ on the same unit ball and impose the rotation singlet constraint
on the auxiliary state space. Its radial Hilbert space has measure
$\pi$. Set $W(r)=r^{-4}$. For $0<t\le1$, its radial heat kernel
on the diagonal, in this measure, satisfies
\[
 p_t^{\rm rad}(r,r)\asymp
 t^{-n/2}(1+r^2/t)^{-(n-1)/2}.
\]
Indeed the Neumann kernel on a ball has two-sided Gaussian bounds
on this time range~\cite{SevenNeumann}. Average them over the
angular separation $\theta$: the volume element is proportional
to $\sin^{n-2}\theta\,d\theta$ and the squared chord length is
$2r^2(1-\cos\theta)$. For $r>\sqrt t$ a cap of angular radius
$\sqrt t/r$ gives both the order of the lower bound and, by
Gaussian integration, that of the upper bound. For $r\le\sqrt t$
the angular integral is of order one. Substitution $r=\sqrt t\,z$
now gives
\[
 \int W(r)p_t^{\rm rad}(r,r)d\pi(r)\asymp t^{-2};
\]
the remaining integrand behaves as $z^{n-5}$ at zero and as $z^{-4}$
at infinity. If all angular modes run in the trace, the corresponding
power is $t^{-n/2}$ instead. Rotation invariance of the operator
alone does not project a quantum trace to singlets.

Allow four-dimensional propagation of the auxiliary field as well.
For definiteness take a flat four-torus of side $\mathcal L>2L$;
its Fourier momenta run over the full lattice. Write
\[
 K_d=(1+H-L^2\Delta_x)^{-d},\qquad d>0 .
\]
For integral $d$ this is the composition of $d$ identical resolvents.
Each is one ordinary geometric diffusion stopped at a unit-rate
exponential time. The total duration has density
$t^{d-1}e^{-t}/\Gamma(d)$, including every positive short duration.
The four-torus heat trace adds a factor of order $t^{-2}$ at short
time. By the positive heat integral, the trace of
$K_d^{1/2} W K_d^{1/2}$ exists exactly when
\begin{equation}
 \begin{aligned}
 d&>4 &&\text{on radial singlets},\\
 d&>(n+4)/2 &&\text{on all angular states}.
 \end{aligned}
 \label{eq:seven-radial-critical-uv}
\end{equation}
At equality the trace diverges logarithmically. These traces can be
defined first with bounded $W$ and evaluated by monotone integration.
At large $t$ the exponential law and $\int Wd\pi<\infty$ suffice.

Here is an action and a complete graph calculation for that bound.
Give a complex neutral auxiliary field the free covariance $K_d$,
independent of the free Maxwell photon, and add its positive quadratic
interaction with multiplication by $g_0 L^4|\overline F_{12,r}(x)|^2$,
$g_0\ge0$. The reference fields are normalized Gaussians. Integrating
the auxiliary field gives the inverse Fredholm determinant of
$I+g_0 L^4 K_d^{1/2}M_{|\overline F|^2}K_d^{1/2}$. Its full relative
photon average lies in $(0,1]$ for $d>4$ on singlets. In fact the
expected trace of the positive perturbation before multiplication
by $g_0$ is finite by \eqref{eq:seven-radial-critical-uv}. On the
cubic torus the variance is exactly
$(\|a_r\|_2^2-\mathcal L^{-4})/2$, by permutation symmetry and
Parseval on the periodized density. The constant mode has zero
curvature. Thus the first interacting vacuum coefficient already
diverges for $d\le4$ on singlets; the bounded term in this variance
cannot cancel its $r^{-4}$ singularity near the origin.

Every individual interaction graph is absolutely integrable at
every fixed order when $d>4$ on singlets. To see the simultaneous
momentum bound, write one positive heat integral per auxiliary
segment in a closed cycle. For fixed segment durations $t_i$ and
routed photon transfers, the auxiliary momenta are $q+P_i$.
Completing the square in $\sum t_i|q+P_i|^2$ bounds the sum over
$q$ by the unshifted torus heat trace at $\sum t_i$, uniformly in
all transfers. A shifted Gaussian lattice sum is maximal at the
unshifted lattice by Poisson summation. The radial semigroup kernels
are nonnegative. Dropping the remaining Kronecker constraints
increases the absolute sum. On each photon edge the curvature
multiplier is bounded by one. The sum of the absolute endpoint
factors divided by the box volume is at most
$\|a_r\|_2\|a_{r'}\|_2$ by Cauchy--Schwarz. There are two ends at
each quadratic interaction. They cost a constant times $W(r)$
per vertex, including when the two ends are joined at that vertex.
The remaining absolute integral for a cycle of $j$ vertices is
bounded up to fixed constants by $\operatorname{Tr}(B^j)$, with
$B=W^{1/2}K_d W^{1/2}$. Its trace agrees with that in
\eqref{eq:seven-radial-critical-uv} by Hilbert--Schmidt factorization;
$\operatorname{Tr}(B^j)\le(\operatorname{Tr}B)^j$. This covers
crossed photon pairings between cycles and all joint short durations.
Five transits therefore suffice in this positive Euclidean singlet
model. The same proof applies to each of the finitely many curvature
components if their squares are summed.

\paragraph{Which cyclic orders fail at the borderline?}
On the singlet space with the four-torus base the precise answer for
the positive majorant is
\begin{equation}
 \operatorname{Tr}(B^j)<\infty
 \quad\Longleftrightarrow\quad j(d-2)>2
 \qquad(d>2,\ j=1,2,\ldots).
 \label{eq:seven-radial-cycle-uv}
\end{equation}
Indeed, in a base momentum slot $q$ put $\lambda=1+L^2|q|^2$ and
$B(q)=r^{-2}(\lambda+H)^{-d}r^{-2}$.
The positive heat integral and its radial $t^{-2}$ estimate give
$\operatorname{Tr} B(q)\asymp\lambda^{2-d}$ as $|q|\to\infty$.
Consequently $\operatorname{Tr} B(q)^j\le(\operatorname{Tr}B(q))^j$.
For the matching lower bound choose a normalized smooth radial bump
$h_q$ obtained by dilating a fixed annular bump to
$\lambda^{-1/2}<r<2\lambda^{-1/2}$. Then
$\|r^{-2}h_q\|^2\asymp\lambda^2$ and the normalized vector
$g_q=r^{-2}h_q/\|r^{-2}h_q\|$ has kinetic energy
$\langle g_q,Hg_q\rangle=O(\lambda)$.
Convexity of $x\mapsto(\lambda+x)^{-d}$ in the spectral measure
of $g_q$ gives
$\langle h_q,B(q)h_q\rangle\gtrsim\lambda^{2-d}$.
A second application of Jensen gives the lower bound of order
$\lambda^{j(2-d)}$ for $\operatorname{Tr}B(q)^j$.
The sum over the four-dimensional momentum lattice proves the
criterion, with logarithmic divergence at equality.
In particular, when $d=4$ all cycles containing at least two marks
have a finite positive majorant. The actual one-mark Maxwell
contraction still diverges. Thus the first bare coefficient fixes
the threshold even when every longer single cycle is harmless.

\paragraph{A harmonic transit that retains every angular state}
There is a positive realization on the full angular space at $d=5$.
At direction $\omega\in S^{n-1}$ draw a fresh A4 radius $v$ with
law $\pi$. From the midpoint $v\omega/2$ of the radial segment, run
ordinary ball Brownian motion until its first exit. Read the exit
direction $\omega'$. The resulting rotationally invariant Markov
operator $J$ respects reversal with the uniform angular measure.
The harmonic extension of a spherical harmonic of degree $j$ is
the same harmonic times the $j$th power of the radius. Hence its
transition eigenvalue is
\[
 \eta_j=\mathbb E(v/2)^j=2^{-j}\frac{n}{n+j},
 \qquad j=0,1,\ldots .
\]
All angular multiplets remain in the trace. Their dimensions grow
polynomially in $j$, so $\sum_j (\dim\mathcal H_j)\eta_j<\infty$.
The stopping time permits arbitrarily short excursions. The midpoint
launch and the separate angular transit specify a geometric
representation of the auxiliary field.

Use $K_d J$ as its covariance, keeping the same radial curvature
observable at each vertex. The two operators commute. In angular
slot $j$ the radial heat generator gains the nonnegative potential
$j(j+n-2)/r^2$, with the regular radial form at the origin.
Its heat kernel is dominated by the singlet heat kernel by positive
semigroup comparison (equivalently kill the radial diffusion at that
rate). The upper trace bound is therefore at most the singlet bound
times $\sum_j(\dim\mathcal H_j)\eta_j$; the singlet itself supplies
the lower bound since $\eta_0=1$. The first trace threshold is again
exactly $d>4$. The angular kernels and their compositions are
nonnegative. Alternatively $j$ stays fixed around each auxiliary
cycle because its photon vertices are angular scalars, and at $m$
vertices the same convergent angular sum has $\eta_j^m$. The proof
of absolute convergence thus includes the complete angular trace.
This example uses a harmonic exit on the geometry with no momentum
restriction. The A4 shell measure alone does not put this transit
in the inverse of a charged kinetic action.

The distinction between this auxiliary field and a physical charged
particle is observable. Its radial zero mode at $d=5$ has a fifth-order
free pole, incompatible with a positive one-particle spectral measure.
There is also a test of the most direct attempt to charge it. Replace
the base Laplacian in the inverse by the ordinary covariant Laplacian,
keeping the action as the fifth power of that operator. On each
spectral mode its logarithmic determinant is five times the ordinary
charged scalar logarithm. A constant nonzero magnetic field on the
four-plane gives the scalar heat trace per volume the factor
$z/\sinh z=1-z^2/6+O(z^4)$, $z=|eF_{12}|L^2t$, relative to the
free factor $(4\pi L^2 t)^{-2}$. The relative logarithm in the
constant radial slot alone consequently contains a nonzero multiple
of $\int_0 dt/t$. The higher radial slots have the same sign.
The charge-one connection and its gauge-required contacts would
therefore undo the benefit of the high power in the propagator.
The successful cyclic bound belongs to the specified neutral
curvature coupling. Deriving a physical charged action would require
the GCM trace to generate its measured vertices as well as its
radial transits.

\subsection{A groupoid charge density that opens into a physical pair}
\label{sec:seven-pair-uv}
The canonical normal-coordinate construction admits a density that
creates opposite electric charges and still generates an Abelian gauge
action. Work on one reconstructed plane in its chosen clock. The
following representation of the momentum pair groupoid tests the
effect of keeping a complementary slot in a charged propagation.

\paragraph{Charge, current and the complementary slot}
Use Hermitian Dirac Hamiltonian matrices $\beta,\alpha_i$.
Put $m=L^{-1}$, $E_p=(m^2+|p|^2)^{1/2}$, and define
\begin{align*}
 S(p)&=\exp[-(L/2)\beta\boldsymbol\alpha\cdot p\,
                                  e^{-L^2|p|^2}],\\
 h(p)&=S(p)\beta E_p S(p)^\dagger,\qquad
 W(p',p)=S(p')^\dagger S(p).
\end{align*}
$S$ is unitary and $W(p'',p')W(p',p)=W(p'',p)$.
For $k=p'-p$ the spatial current kernel is
\[
 J_i(p',p)=S(p')^\dagger
       \int_0^1 (\partial_i h)(p+t k)\,dt\ S(p).
\]
It satisfies $J_i(p,p')=J_i(p',p)^\dagger$ and
$k_iJ_i=\beta E_{p'}W-W\beta E_p$ with the spatial index summed.
In the block $P_+J_iP_-$ the sum of the two positive energies multiplies
the density. At $p=p'=0$ the chosen generator gives $J_i=\alpha_i$.

Use the filled negative slot, equivalently the empty particle--hole
Fock vacuum in the diagonal energy basis. Normal-order the bilinears
with kernels $W$ and $J_i$ in that basis and call them $R(k)$ and
$J_i(k)$, with transfer $k$. Their off-diagonal blocks create or
annihilate an actual particle and hole. On a torus the off-diagonal
entries of $S$ are square summable. Its momentum-by-momentum
Bogoliubov rotation has a common Fock implementer: the product of
the opposite-momentum pair rotations converges on the vacuum since
their squared mixing amplitudes sum to a finite number, and then
on finite excitations. The $R(k)$ are the conjugates of the
commuting onsite charge translations. The normal-ordering constant
at $k\ne0$ vanishes by momentum conservation; at zero transfer
the rotation preserves net charge. Thus the smeared densities
commute and $R(0)=N_+-N_-$.

If one inserts $P_+=(1+\beta)/2$ at every intermediate arrow,
the two charge paths with small transfers $k,l$ in opposite
orders disagree already at rest. The leading difference of the
path with $k$ followed by $l$ and the one with $l$ followed by $k$
is $P_+[T_k,T_l]P_+$, where $T_k=-(L/2)\beta\boldsymbol\alpha
\cdot k$. In the two-component slot it equals
$-iL^2\boldsymbol\sigma\cdot(k\times l)/2$.
Indeed the omitted product on each path is $P_+W P_- W P_+$;
the two complete products agree exactly. The same off-diagonal
arrows permit the density to create a physical pair.

\paragraph{A positive Hamiltonian with its electric contacts}
Keep the real canonical-coordinate mark of Section~\ref{sec:six-recoil-uv}, $\mathcal N_r$
for $r=|k|$. Put $D(k)=\mathcal N_r R(k)$ and
$V_\lambda(k)=\mathcal N_r\boldsymbol\epsilon_\lambda(k)
\cdot\mathbf J(-k)$, with transverse real polarizations chosen
equally at $k$ and $-k$ on momentum patches. In a box of volume
$V$ omit the fixed Maxwell zero mode. With $H_f$ the positive free
fermion energy and $H_{\rm n}$ the normal oscillator energy, take
\begin{align*}
 H_e={}&H_f+H_{\rm n}+\sum_{k,\lambda}r\,T_\lambda(k)^\dagger
 T_\lambda(k)\\
 &+\frac{e^2}{2V}\sum_{k\ne0}\frac{D(k)^\dagger D(k)}{r^2},\\
 T_\lambda(k)={}&a_\lambda(k)-
            \frac{e V_\lambda(k)}{\sqrt{2V}\,r^{3/2}} .
\end{align*}
The first square supplies the dynamical transverse interaction
and its contact. The last term is the complete instantaneous Gauss
energy. These positive forms are closable on finite-particle smooth
packets: the weighted row maps have adjoints on finite mode families
of core vectors, and the high-mode sums on the core converge by the
bounds below.

The commuting $D(k)$ specify the scalar coupling and gauge generator
$\Theta_\chi=\sum_k\widehat\chi(k)D(k)$ in the box Fourier-series
convention. Use the transverse probe in the same linear current as
the dynamical photon. A longitudinal probe $\nabla\chi$ conjugates
the whole transverse-slice Hamiltonian by $\exp(ie\Theta_\chi)$,
with scalar term $e\Theta_{A_0}$. The time-dependent gauge law
follows as for the canonical normal coordinate. At zero field
$k\cdot J_{\rm full}(k)=[H_f+H_{\rm n},D(k)]$ in this convention.
The contraction includes $[H_{\rm n},\mathcal N_r]R(k)$ in
addition to the matter current. Thus a normal excitation contributes
its actual energy to the longitudinal Ward identity.

\paragraph{A closed vacuum graph and its opening}
At second order resolve $D$ and $V_\lambda$ on the vacuum into a
pair and zero or one normal quantum. Write $w>0$ for the energy
of this intermediate matter and normal state and $d,j_\lambda$
for the respective rows at the appropriate transfer. After
removing $e^2$, the complete coefficient per unit time in the
large-slab vacuum energy expansion is
\[
 \frac1{2V}\sum_{k\ne0}\frac1{r^2}\sum_{\rm states}
 \left[|d|^2+\frac{w}{r+w}\sum_\lambda|j_\lambda|^2\right].
\]
The transverse contact supplies $|j_\lambda|^2$ and a virtual
photon plus pair subtracts $r|j_\lambda|^2/(r+w)$ in the same
units. The result is strictly positive. At rest $\mathcal N_0=I$
and the first transverse pair matrix tends to $P_+\alpha_i P_-$
at threshold. In infinite volume the pair hence has a continuum
with the ordinary leading Dirac threshold weight. The density
also creates pairs for small nonzero $k$ because $W(p+k,p)$
mixes the two slots.

The one-body matrices $W,J_i$ are uniformly bounded. Their
off-diagonal entries are bounded by a constant times
\[
 G(p)+G(p+k)+\int_0^1G(p+t k)\,dt ,
\]
with $G(p)=e^{-cL^2|p|^2}$ for some $c>0$. Subtract the
identity from $S$ and $\beta\partial_i E$ from $\partial_i h$;
the remaining polynomial factors are absorbed by the Gaussian.
The relative momentum sum of this bound divided by $V$ is uniform
in $k$ and in large boxes. In the empty particle--hole vacuum
every closed fermion cycle needs a pair-changing vertex, including
with the indicated order inside a square. Use one such vertex
to sum its independent relative momentum absolutely and bound
the other matrices uniformly. Open chains are tested on smooth
packets. The normal vacuum contractions, including coincident
ones, cost at most $\|m_r\|$ per end. Thus transverse photons cost
$C b(r)/r$, and instantaneous transverse and Gauss transfers cost
$C b(r)/r^2$. On a finite torus the nonzero mode sums at each
fixed order converge absolutely on a finite slab, including
overlaps and contacts. In the displayed second-order coefficient
per volume the continuum integrals are also finite, since the
relative row is integrable and $b(r)/r^2$ is locally integrable.

The summed squared transverse pair norm per volume is at most
$C b(r)$, including the excited normal output. It decays
exponentially even after multiplication by $2r$ at large spatial
momentum. A covariant conserved current with an occupied compact
transverse mass band would retain a positive lower limit there.
Thus the same source gives a conserved charge, a vacuum pair
continuum and finite bare graphs in the indicated box, while
failing that continuum's boost test. Exact composition of these
comparison arrows is insufficient to identify them with the
physical GCM source.

\paragraph{A shared normal record fails a volume test}
The volume restriction is visible at fourth order. Take two
instantaneous Gauss insertions at distinct slab times, close the
two charged densities at each insertion on their own pair cycle,
and join the insertions by two normal Wick contractions. Write
$F_V(k)=\|R(k)\Omega_f\|^2$ and
$C_{rr'}(t)=\langle n_r,e^{-|t|\epsilon} n_{r'}\rangle$.
Up to positive fixed factors this connected graph has integrand
\[
 V^{-2}\sum_{k,l\ne0}
 \frac{F_V(k)F_V(l)}{|k|^2|l|^2}
              C_{|k|,|l|}(t)^2 .
\]
The square is real and nonnegative. On a fixed small open momentum
bin away from zero, $F_V(k)/V$ is bounded below for large boxes by
the nonzero continuous pair row. For any $r>0$ the transit strictly
increases the normal momentum on each trajectory, so $v_r$ cannot
equal a scalar multiple of its normalized entry wavefunction.
Thus $n_r\ne0$. At equal radii $C_{rr}(t)>0$ for finite $t$;
continuity gives a positive lower bound on a smaller bin and a
compact time interval. The graph grows at least as $V^2$. Its
absolute weight per volume diverges even though all its momentum
sums converge at fixed volume. This is a graphwise obstruction;
the argument isolates two specified Gauss contacts.

There is a spatial resolution of this obstruction. Give the normal
oscillators a spatial momentum $u$ and positive energy
$\sqrt{|u|^2+\epsilon(p)^2}$. Replace the fluctuating part
of $D(k)$ by
\[
 V^{-1/2}\sum_u \eta(u)
       [b^\dagger(u,n_r)+b(-u,n_r)]\,R(k-u),
\]
leaving $h_rR(k)$ as its elastic component. Use the same spatial
convolution at transverse current incidences. One may take $\eta$
proportional to the real even A4 profile $\chi_{248}(L|u|)$ with
unit continuum $L^2(d^3u/(2\pi)^3)$ norm. The commuting density
algebra and $D(0)=R(0)$ survive. Each charged cycle in the preceding
graph now requires its two normal transfers to sum to zero. Four
normal incidences cost $V^{-2}$ and only one normal momentum sum
remains. This replaces the $V^2$ bound by $O(V)$: the additional
continuum integral is bounded by $\int|\eta(u)|^4d^3u$, and the
shifted relative-momentum estimates remain uniform. This repairs
the specified connected graph. An all-orders volume bound would
also have to control the soft photon transfers of every connected
topology.

\subsection{What a shadow field does to a physical cut}
\label{sec:seven-shadow-uv}
A field eliminated by its classical equation can generate vertices
without adding an asymptotic particle. This is one possible meaning
of a shadow field. In the nonlocal construction of Kleppe and
Woodard the eliminated field appears only on trees of auxiliary
lines, and the gauge measure is a separate ingredient
\cite{SevenKleppeWoodard}. The following test applies that elimination
rule to a geometric transfer on one fixed sector plane. It concerns
an additional action prescription. Forgetting the carrier or support
when comparing with a Minkowski chart introduces no such field.

In the Euclidean scalar comparison let $D=k_E^2+m^2$, $m>0$, and
give two fields covariances $C_d=F(D)/D$, $C_h=(1-F(D))/D$. Use a
local quartic interaction $V(\phi+h)$ and keep external particles in
the $\phi$ channel. Eliminating $h$ by $h=-C_h V'(\phi+h)$ generates
the auxiliary trees order by order. Consequently every cycle in the
two-color representation contains a $d$ line. A loop-momentum
variation vanishing on the $d$ lines would be a conserved circulation
entirely on a forest; removal of successive leaves sets it to zero.
For fixed positive depth and an exponentially damped $C_d$, this
gives an integrable joint Euclidean momentum envelope for each such
massive graph with polynomial vertices. Indeed the collection of its
damped linear routings is injective on the loop space, and its sum
of norms bounds the loop norm below up to a fixed external shift.
This controls the interaction graphs. The free determinant, source
measure, and any additional induced vertices still belong to their
own action. Integrating $h$ as an independent Gaussian instead
restores the missing cycles: internal lines of $\phi+h$ then have
covariance $C_d+C_h=1/D$. If both independent fields were positive
physical scalars with unsubtracted spectral representations, uniqueness
of the Stieltjes transform would force their two positive spectral
measures to sum to the single atom at $m^2$. Giving that whole residue
to $C_d$ leaves zero measure for $C_h$. Thus a nonzero complement
with this free sum cannot itself be an independent positive physical
scalar. Physical cross-correlations between nonindependent fields
would require a different matrix source calculation.

\paragraph{An auxiliary continuum beneath the forward amplitude}
Even matching the residue of each charged or massive particle pole
does not settle this test. Take the scalar transfer
\[
 F(D)=\exp[-a(\sqrt{D+M^2}-M)],\qquad a>0,\quad M>0 .
\]
Here the argument includes the physical mass shift; $F(0)=1$.
The principal square root is understood in the coefficient of the
declared vectorial boundary prescription. The shadow has no pole
at $D=0$, but has a cut beginning at $s=\mu^2=m^2+M^2$. Its signed
spectral representation is
\begin{equation}
 \begin{aligned}
 C_h(k_E)&=\int_{\mu^2}^\infty\frac{r_h(s)}{k_E^2+s}\,ds,\\
 r_h(s)&=\frac{e^{aM}\sin(a\sqrt{s-\mu^2})}
                    {\pi(s-m^2)} .
 \end{aligned}
 \label{eq:seven-shadow-density-uv}
\end{equation}
Taking $-\operatorname{Im}/\pi$ on the upper rim gives this sign.
The Cauchy contour closes because $F$ is bounded on the principal
slit plane and $C_h=O(|D|^{-1})$ at infinity. In particular $r_h$
is strictly positive just above its first threshold.

The free vacuum also has a direct source test if $\phi$ is integrated
with the ordinary field measure. Differentiating
$\tfrac12\operatorname{Tr}\log(C_d^{-1})$ with respect to $m^2$
at fixed $a,M$ gives, per volume,
\[
 \frac12\int\frac{d^4k_E}{(2\pi)^4}
   \left(\frac1D+\frac{a}{2\sqrt{D+M^2}}\right).
\]
Both terms are positive and their sum is ultraviolet divergent. Forest
coverage of interaction graphs never differentiates this determinant.
Normalizing the free measure separately at every $m$ removes this
response by prescription and does not derive it from the vacuum trace.

In the forward quartic bubble the retained color sum is the ordinary
two-line product minus $C_h C_h$. Up to the common positive
vertex and symmetry factor, the bubble density that has been
\emph{removed} is
\[
 \int_{\sqrt{s_1}+\sqrt{s_2}\le\sqrt u}
 r_h(s_1)r_h(s_2)\,
 \frac{\sqrt{(u-s_+)(u-s_-)}}{16\pi^2 u}\,ds_1ds_2,
\]
where $s_\pm=(\sqrt{s_1}\pm\sqrt{s_2})^2$ and $s_i\ge\mu^2$. To take this discontinuity in the decomposition
of the finite bubble, differentiate once in the Euclidean invariant
of the external pair. The Feynman-parameter denominator at positive
invariant is at least $2\sqrt{x(1-x)s_1s_2}$ and
$\int |r_h(s)|\,ds/\sqrt{s}<\infty$. The differentiated spectral
integrations are therefore absolute. A constant of integration
has zero discontinuity. No subtraction is performed on the original
finite colored integral.

For $u=4\mu^2+\delta$ with small $\delta>0$, both densities in
the removed term are positive. With $s_i=\mu^2+y_i$ each contributes
a positive multiple of $\sqrt{y_i}$, and phase space contributes a
square root at $2(y_1+y_2)=\delta+O(\delta^2)$. Rescaling
$y_i=\delta z_i$ gives a strictly positive leading term proportional
to $\delta^{7/2}$. For $M=0$ the density instead starts as
$a/(\pi\sqrt{s-m^2})$ and the deficit starts as $\delta^{3/2}$
at the ordinary pair threshold. The physical tree amplitude is the
original quartic coupling and its two-particle contribution already
accounts for the entire ordinary cut at this order. The retained
bubble is smaller. Additional positive intermediate states can only
increase the required inclusive weight. A compensating source or
measure contribution would have to be derived in the same action.
The deficit is specific to the square-root transfer; an entire
shadow propagator has no such finite cut.

\paragraph{Graded determinants and physical sources}
A quantized partner of opposite statistics instead produces a
determinant ratio. For identical operators and boundary conditions
the finite-mode factor $\det L[A]/\det L[A]$ equals one at every
background. Its current and all contact derivatives consequently
vanish along the same diagonal source. Differentiating the numerator
alone can recover the original current, while its derivative is then
a different opening from the one used in the vacuum. A BRST-exact
state has zero transition from a closed vacuum under a physical
BRST-closed current, by moving the conserved BRST charge through
that current. This requires the physical pairing, invariant measure
and boundary domain. A negative weight in a graded trace is not an
extra negative-probability physical channel.

On a closed real-time contour $\operatorname{Tr}(U[J]\rho
U[J]^\dagger)=1$ for equal branch sources. This statement permits
nonzero responses to independent branch probes. It sums contour
assignments and presupposes the indicated evolution and state; it
supplies no integrable majorant for the absolute value of an individual
bare Euclidean vacuum graph. Thus neither attaching a graded sign
nor setting the contour sources equal resolves the explicit
square-root cut deficit while preserving its original source opening.

\paragraph{A dynamical multiplet can cancel the vacuum with a nonzero current}
A graded fluctuation cancellation need not erase the response to a
physical probe. A finite-mode supersymmetric comparison shows the
additional action terms that make this possible. Use bosonic
coordinates $z=(A,x)$ and their momenta, two positive-CAR modes
$f$, and
\[
 M_g=\begin{pmatrix}\omega&g\\g&m\end{pmatrix},\qquad
 H_g=\tfrac12(p^{\mathsf T}p+z^{\mathsf T}M_g^2z)
       +f^\dagger M_g f-\tfrac12\operatorname{tr}M_g .
\]
Take $\omega,m>0$, $g^2<\omega m$. The scalar constant follows
from the supersymmetric oscillator ordering:
$H_g=\{Q,Q^\dagger\}/2$ for $Q=f^\dagger\cdot(p-iM_g z)$.
The two positive bosonic frequencies are the eigenvalues of $M_g$.
Their zero-point sum and the displayed constant cancel in the
ordinary ground state for every such $g$. This is a ground-state
calculation in a positive Hilbert space. At nonzero temperature an
ordinary thermal trace has different fermionic and bosonic boundary
conditions and is not the index.

Couple a probe $b$ by replacing $A$ with $A+b$ throughout the
interaction $H_g-H_0$. Its bosonic part is
$g(\omega+m)(A+b)x+g^2[(A+b)^2+x^2]/2$.
There is also the fermionic mixing fixed by $M_g$.
The measured linear current, after division by $g$ at zero coupling,
is $(\omega+m)x$ and its one-quantum transition has norm squared
$(\omega+m)^2/(2m)>0$. The two-$b$ action Hessian at order $g^2$,
including the contact, is $1-(\omega+m)^2D_m(\nu)$ with
$D_t(\nu)=(\nu^2+t^2)^{-1}$. Sewing only these two ports gives
$-(4m)^{-1}$ after integration against $D_\omega d\nu/(4\pi)$.
The additional $g^2x^2/2$ contact contributes $+(4m)^{-1}$.
At zero temperature the complete cancellation is already pointwise
in the real frequency trace. Indeed the total bosonic coefficient is
\[
 \tfrac12[D_\omega+D_m-(\omega+m)^2D_\omega D_m]
       =(\nu^2-\omega m)D_\omega D_m .
\]
The fermionic coefficient is its negative, from
$\operatorname{Re}[(i\nu+\omega)^{-1}(i\nu+m)^{-1}]$.
Equivalently $|\det(i\nu+M_g)|^2=\det(\nu^2+M_g^2)$.
The same probe is used at both bosonic ports. The native partner
terms alter the full vacuum graph sum while the source still opens
a physical state. This is an explicit comparison for a transverse
oscillator, with no charged Gauss law or relativistic spectrum
asserted. A corresponding cancellation in the shell action would
require deriving such partner couplings and the ordering from that
action, and testing the absolute bounds demanded for each retained
graph or complete source coefficient.

\paragraph{What must be derived at the physical charged boundary}
The tests give an order in which to verify a proposed representation.
Compute the two derivatives of its curvature pairing on one common
configuration, including the shell measure and any source-dependent
boundary state. In the resulting cyclic product identify the relative
displacement sampled by the measured photon. Test the charged opening
of the same product on a positive invariant mass band, including its
boosts and the mandatory contact. If a shadow contributes, varying its
action specifies which of these sources it opens and which new cuts it
creates. Only after these identifications do the all-order joint
momentum estimates of the chosen representation apply to physical
bare O--QED.

\paragraph{Reading these tests as constraints on $\Omega_\alpha$}
At the two-current level there is now a sharper requirement on that
postulate. Test the high-frequency current on the same plane as the
measured low-mass band. Equations~\eqref{eq:four-angular-superconvergence-uv}
and~\eqref{eq:four-excursion-response-uv} show how the time structure
of the history can remove or supply the required contact. Changing a
spin normalization or a distribution over Newton's constant does not
select which charged states saturate that contact. An action-derived
rule must identify them in the opened source and in the determinant.

The GTL distributions, the directed velocity shell and the scalar
coefficient along each unit carrier have been kept in the underlying
curvature action. The Hilbert space on which a charged history ends,
the ordering of its velocity changes and the charge density that acts
there are further data in each example. These choices are visible in
an observable. They change, respectively, the norm of a two-pair state,
the first nonzero closed Wilson path and the commutator of two gauge
variations. None of these quantities is fixed by the invariant mass
density of a single measured pair.

For a proposed native action three low-order computations now precede
any general power count. First open it twice into the same charged
one-particle spaces and compare with the exchange determinant in
\eqref{eq:next3-car-det-uv}. Next take its two electric variations,
including its Gauss and transverse contacts, and boost a fixed mass
band before integrating the center momentum. Finally reflect the
first interacting correction on a free Maxwell null test. Each uses the
same physical source and its own boundary measure. If an additional
physical normal or gauge mode cancels a term in any of them, its
norm, its transition amplitude and its determinant have to appear
in this calculation. The geometric scale determines the widths in
the examples. Picking a clock-frame smearing or a mass rotation by
hand does not derive that additional dynamics from the GCM action.

\paragraph{What remains to be derived from the directed action}
These calculations give an operational test for a proposed physical
realization of $\Omega_\alpha$. Differentiate its normalized source
first at two electric probes and then with that response joined to
another on the same propagating photon. The first derivative pair must
pass both the open-pair cut and the separated-point current test. At
the next step the determinant and the contact from varying the measure
must be included before estimating each independent hard circulation.
If a contact cancels a ring, the proof must identify the action-defined
complete graph to which both contributions belong and exhibit an
absolutely integrable integrand for it. A finite sum of numbers
assigned separately to divergent integrals cannot establish this.

In particular, specifying an additional distribution for Newton's
constant concentrated on $G_N$ would reproduce the original integral
for every observable by the delta identity. It leaves this measured
current and its joined graph exactly as before. The datum still needed
is the transformation of the charged boundary states, the neutral
pair modes and their common integration measure under a variation of
the physical connection. The constrained GCM action alone, with the
representations considered here, has yet to fix that datum. The
unrestricted all-orders theorem therefore remains open. These tests
can be applied at one and two source insertions before attempting the
full overlapping-subgraph bound of the next section.

\subsection{Geometric moments in simultaneous hard subgraphs}
\label{sec:nativefull-subgraphs-uv}

\paragraph{Several marginal photon momenta in the same boundary state.}
The operator and boundary state in the curvature action might
retain correlated displacements on one velocity shell. There is a
test which uses their joint measure directly. Fix a locked
sector-plane and an individual complete Euclidean contribution.
Suppose its charged marks and their physical pairing give
\[
 C(k)=\int e^{i\sum_{j=1}^{b}k_j\cdot d_j}\,d\pi(d),
 \qquad C(k)\ge0 .
\]
Here $b\ge1$, each $k_j\in\mathbb R^4$, and $d_j$ is the
difference of the two routed displacements across the pairing.
The finite measure $\pi$ is the pushforward of the complete
marked source, including the GTL delta integral, oriented shell,
determinant, contacts and boundary weights of this contribution.
In a squared amplitude the measure can be signed or complex even
when $C$ is nonnegative. For a nonpositive weight in momentum
space this argument requires a separate positive majorant.
Assume that all residual hard factors are bounded by a constant
times $\prod_j(\mu^2+|k_j|^2)^{-2}$ outside the infrared regions,
with the analogous bounds on partially hard regions and a specified
infrared treatment. Here $\mu>0$ sets comparison units and the
integrals extend over the full momentum space.

For nonnegative $\pi$ the following equality holds, including
the value $+\infty$:
\begin{equation}
\begin{aligned}
 &\int C(k)\prod_{j=1}^b
      \frac{d^4k_j}{(2\pi)^4(\mu^2+|k_j|^2)^2}\\
 &\qquad =\int\prod_{j=1}^b B_\mu(d_j)\,d\pi(d),\\
 B_\mu(x)&=\frac{1}{(4\pi)^2}\int_0^\infty
   e^{-\mu^2 t-|x|^2/(4t)}\,\frac{dt}{t}
   =\frac{K_0(\mu |x|)}{8\pi^2}.
\end{aligned}
\label{eq:nativefull-log-energy-uv}
\end{equation}
Here $K_0$ is the modified Bessel function of the second kind.
For a finite signed or complex measure the identity holds when
the integral on the right against $d|\pi|$ is finite. Indeed the
Schwinger representation of each denominator and $C\ge0$ permit
Tonelli's theorem in $(k,t)$. At fixed positive proper times the
Gaussian Fourier integrals can be taken against the finite measure.
The resulting kernels are positive, so Tonelli applies again for
nonnegative $\pi$, and absolute integration suffices in the other
case. The proper times here represent the residual bounds and are
distinct from any durations integrated inside the charged source.

In particular $B_\mu$ grows only as $\log(1/(\mu |x|))/(8\pi^2)$
at the origin and is bounded away from it. To test simultaneous
short durations, disintegrate $|\pi|$ over the actual remaining
geometrical variables $\lambda$, using conditional probability
measures and marginal mass $d\nu(\lambda)$. Suppose each displacement
satisfies
\[
 \Pr_\lambda\{|d_j|<r\}\le
  \min\{1,C_j(r/a_j(\lambda))^{\delta_j}\},\qquad
 a_j(\lambda)>0,\quad \delta_j>0 ,
\]
with uniform constants $C_j\ge1$. Integrating the tail of
$\log^+(a_j/|d_j|)$ gives a bounded $b$th moment; here
$\log^+x=\max(0,\log x)$. Conditional H\"older with $b$ factors
then proves convergence whenever
\[
 \int\prod_{j=1}^b
   \big[1+\log^+(1/(\mu a_j(\lambda)))\big]\,d\nu(\lambda)
 <\infty .
\]
This controls correlated durations up to their joint boundary.
On a regular constrained chart, for instance, a conditional
density bounded along a dimensionless shell or depth coordinate
and a monotone component of $d_j$ with derivative of magnitude
at least $a_j$ give the required first-power bound by a change of
variable. All charts and their weights must be included. A
direction removed by the GTL or boundary constraints is unavailable
for this calculation.

For a sharp example of the role of the trace, take a coherent
amplitude integrated against $dr$ on $[0,1]$ in a dimensionless
depth coordinate. Let every phase read $r\,\ell n_j$, with $\ell>0$
and fixed Euclidean unit directions $n_j$. Its squared amplitude
has $d_j=(r-r')\ell n_j$ for all $j$, with $dr\,dr'$ on the
unit square. The entire marked list uses one scalar coordinate.
Nevertheless the near-diagonal cost in
\eqref{eq:nativefull-log-energy-uv} is bounded by a constant
times $1+|\log|r-r'||^b$, which is integrable for every finite
$b$. A boundary pairing that imposes $r=r'$ instead gives
$C(k)=1$; its marginal loop integrals diverge. This calculation
gives a sufficient mechanism for the stated logarithmic residual
powers even with shared normal data. Its application to the
physical action requires deriving both the residual power on
each hard region and the displacement measure after the actual
charged source has been sewn.

\paragraph{Displacement retained by overlapping hard subgraphs.}
The geometric test also works when the line momenta are dependent.
Use fixed reference units. Write the residual line momenta as
$a_e\cdot k+p_e$, with $a_e\in\mathbb R^b$ spanning the loop
index space and $p_e$ fixed. Suppose a positive majorant of the
complete contribution is $C(k)$ times
$\prod_e(m_e^2+|a_e\cdot k+p_e|^2)^{-\nu_e}$, with $C(k)$
the nonnegative characteristic integral above, $\nu_e>0$,
$m_e>0$. A massless amplitude needs its stated infrared bounds.
For nonempty sets of lines $S$ write
\[
 \begin{aligned}
 V_S&=\{v:\ a_e\cdot v=0\text{ for }e\notin S\},\\
 \delta_S&=\sum_{e\in S}\nu_e-2\dim V_S,\qquad
 x_S(d)=\|(\operatorname{proj}_{V_S}\otimes I_4)d\|.
 \end{aligned}
\]
Thus $x_S$ measures the surviving phase displacement precisely
in directions invisible to the lines outside $S$. Put
$H_\delta(x)=1$ for $\delta>0$,
$H_0(x)=1+\log^+(1/x)$ and
$H_\delta(x)=1+x^{-2|\delta|}$ for $\delta<0$,
with the corresponding infinite values at zero.

\begin{proposition}[A geometric moment test on each nested set of hard lines]
\label{prop:nativefull-sewn-flags-uv}
The bare absolute momentum integral is finite if for every ordering
of the lines, with successive nonempty prefixes $S_i$, the native
displacement measure satisfies
\[
 \int \prod_i H_{\delta_{S_i}}(x_{S_i}(d))\,d|\pi|(d)<\infty .
\]
The moments are joint moments of the complete boundary pairing.
\end{proposition}
\begin{proof}
Introduce the positive heat parameters $t_e$ of the residual
denominators. Their loop matrix is $M=\sum_e t_e a_e a_e^{\mathsf T}$.
As in \eqref{eq:nativefull-log-energy-uv}, nonnegativity allows
the heat and momentum integrals to be exchanged first. Fourier
integration at fixed $t_e$ contributes, in modulus, at most
\[
 {\rm const}\,(\det M)^{-2}
 \exp[-d^{\mathsf T}(M^{-1}\otimes I_4)d/4] .
\]
Completing the square with the real external shifts adds a phase
and a nonpositive real exponent. Order the parameters at most one
from smallest to largest. The others have integrable large-time
tails from their masses. For the determinant lower bound those
others may be set equal to one. Write $t_i=r_i\cdots r_h$ in the
ordered small parameters, $0<r_i\le1$. A maximal basis minor
obtained by scanning columns from largest width to smallest
uses exactly $\dim V_{S_i}$ columns in $S_i$. Cauchy--Binet
therefore bounds the heat density including $(\det M)^{-2}$ by
a constant times $\prod_i r_i^{\delta_{S_i}}\,dr_i/r_i$.
For vectors in $V_{S_i}$ the quadratic form of $M$ is at most
$C t_i$ times their squared length, since only lines in $S_i$
contribute. The variational formula for the inverse quadratic
form gives
\[
 d^{\mathsf T}(M^{-1}\otimes I_4)d\ \ge\
          x_{S_i}(d)^2/(C t_i).
\]
Divide this estimate among the finitely many indices with
$\delta_{S_i}\le0$ and use $t_i\le r_i$. Their scale integrals
are bounded by
$\int_0^1r^{\delta-1}e^{-c x^2/r}dr\le C_\delta H_\delta(x)$;
for $\delta>0$ omit the exponential. Completing any ordering
of the small parameters to an ordering of all the lines causes
no loss since $H_\delta\ge1$. Absolute integration in $|\pi|$
finishes the proof.
\end{proof}

In particular, when all faces have $\delta_S\ge0$, one scalar
component of the projected displacement with finite logarithmic
moments of all finite degrees on each equality face suffices by
H\"older's inequality, even for perfectly correlated visits.
Only moments up to the number of equality faces of a tested
ordering are needed there. A zero-net-winding return can qualify
if its two histories still differ along these directions after
gluing. The bound is conditional on the charged operator actually
yielding the stated majorant; positivity of a single propagator
does not identify the determinant, contact or measure of that
source.

\paragraph{The simultaneous duration cost of a pinned normal network}
For the Gaussian normal action one can compute a sharp joint cost
even when many short segments meet at contacts. Take a finite
connected multigraph with held vertices $B\ne\varnothing$, interior
$I$, normal dimension $n_\perp$, fixed $D>0$, and one kernel
$(4\pi Dt_e)^{-n_\perp/2}\exp[-|y_i-y_j|^2/(4Dt_e)]$ per edge.
Hold all boundary positions equal, integrate the interior positions,
and use reference units $0<t_e<1$. Orient the edges and let $C_I$
be their incidence matrix restricted to interior rows. Cauchy--Binet
in the normal Gaussian integral gives the factor
\[
 (4\pi D)^{-(|E|-|I|)n_\perp/2}\,\Psi_B(t)^{-n_\perp/2},
 \qquad \Psi_B(t)=\sum_F\prod_{e\notin F} t_e ,
\]
where $F$ runs over spanning forests with one held root per component.
Self-loops have zero incidence and stay in every complementary product.
If the positive duration density is comparable to
$\prod_e t_e^{\eta_e-1}dt_e$, then its complete short-duration
integral is finite precisely when, for every nonempty edge set $S$,
\[
 \sum_{e\in S}\eta_e>
 \tfrac{n_\perp}{2}\bigl(|S|-\operatorname{rank}C_{I,S}\bigr).
\]
Indeed order the durations $t_1\le\cdots\le t_h$ and set
$t_i=r_i\cdots r_h$. Greedily select independent incidence
columns from the smallest duration upward. The resulting basis
forest contains exactly $\operatorname{rank}C_{I,S_i}$ edges
of each prefix $S_i=\{1,\ldots,i\}$. Its complementary monomial
dominates every other one since all $r_i\le1$. The positive
forest sum is bounded between this monomial and the number of
forests times it. The scale integrals therefore have measure
$dr_i/r_i$ and the exponents displayed in the strict inequalities.
Every subset is a prefix of some ordering. Equality on any prefix
already gives a logarithm with the other ratios held in a compact
subinterval of $(0,1)$.

For example, two serial segments joining two equal held endpoints
give $(t_1+t_2)^{-n_\perp/2}$ up to its fixed normalization.
Each separate endpoint limit at a fixed positive opposite duration
has finite kernel, yet the common corner requires
$\eta_1+\eta_2>n_\perp/2$ in addition to $\eta_i>0$.
This extra power counts the path between the held vertices even
though the unpinned graph has zero cycle rank.

At fixed equal held positions the same upper bound controls ordinary
source derivatives of the full normal kernel if for every retained
jet $J$ the durations obey $|\partial_J t_e|\le C_Jt_e$, uniformly
on the comparison chart. The product rule on each positive
complementary monomial gives $|\partial_J\Psi_B|\le C'_J\Psi_B$;
differentiation of its negative power preserves the same majorant.
Moving endpoints, a varying boundary density, and charged or ghost
determinants contribute their own jets. Separated held endpoints
supply an additional heat exponential, so the necessity statement
here applies to the equal-position test.

\paragraph{A nonidentity test on regular analytic branches.}
There is a way to establish the logarithmic moments from the
restored integration domain. Consider finitely many constrained
charts supporting the complete paired density; on each, assume a
compact coordinate region, a density in $L^p$ for some $p>1$, and
real-analytic displacement maps on connected neighborhoods of that
region. For each logarithmic face $S$ require that at least one
component $f$ of $(\operatorname{proj}_{V_S}\otimes I_4)d$ is
not identically zero on each occupied chart. Then all its required
logarithmic moments exist. The component may be stationary at
individual configurations, and the different faces may use the
same geometric variable.

Indeed, at each zero of $f$ a finite derivative is nonzero by
analyticity. There is a direction in which a derivative of some
order $m\ge1$ is bounded away from zero on a small box. For a
one-dimensional slice of this box, a set of length $s$ on which
$|f|<r$ contains $m+1$ points separated by at least $s/(2m+2)$.
The divided difference at the points is bounded by $C_m r/s^m$
and equals $f^{(m)}(\xi)/m!$ for an intermediate point. Thus
$s\le C r^{1/m}$. Fubini and a finite cover give a positive
power bound on sublevel volume. H\"older against the complete
$L^p$ density preserves a positive power. The logarithmic moments
follow by integrating that bound over the logarithmic scale.
In particular this resolves all the marginal faces in
Proposition~\ref{prop:nativefull-sewn-flags-uv} under its residual
power bounds. If the physical velocity shell or duration domain
requires infinitely many charts, the weighted costs must remain
summable over their union. These hypotheses concern the full
paired measure after imposing the delta functions, including any
singular part concentrated on the displacement zero set.

Smoothness alone cannot replace the analytic nonidentity test.
For a coherent amplitude let $r$ be uniform on $(0,1)$, read the
displacement $g(r)n$ along a fixed unit direction $n$, and take
$g(r)=e^{-1/r^2}$. In the squared amplitude the difference
$g(r)-g(r')$ has no atom. For $T\ge1$, however, both variables
lie below $T^{-1/2}$ with probability $1/T$. On this event
$|g(r)-g(r')|\le e^{-T}$. Its logarithmic moment is infinite, so
the positive one-loop integral in
\eqref{eq:nativefull-log-energy-uv} diverges. This is the
distinction between the retention of an intermediate mark and
quantitatively resolving it on the actual preparation.

\subsection{The internal trace and the measured displacement}
\label{sec:monomial-trace-uv}

\paragraph{Which word traces can retain an intermediate displacement?}
Fix one sector-plane $\bin=(s,b_s)$ and the positive Hilbert space of
its internal word records. The following tests concern the positive
even part of a candidate path trace. A graded supertrace on all
fields need not be positive. A Euclidean displacement $d$ below
is a variable dual to a routed hard momentum on the locked chart.
Its four components play no role in specifying the Lorentzian pole.

\begin{proposition}[Finite cyclic trace and exact prefix records]
\label{prop:monomial-prefix-uv}
Suppose $P_v$ are nonzero orthogonal vertex projections in a unital
operator algebra with a faithful positive finite trace $\tau$.
A directed edge $e:v\to u$ has a record operator $S_e$ with
$S_e^*S_e=P_u$ and $S_eS_e^*\leq P_v$.
At each $v$ suppose the ranges of its outgoing edges are orthogonal.
On a finite graph let $M_{vu}$ count its edges from $v$ to $u$.
Then $w_v=\tau(P_v)>0$ satisfies $Mw\leq w$, so the spectral
radius of $M$ is at most one. The unital corner of a single vertex
with two such outgoing loops admits no normalized positive tracial state.
\end{proposition}
\begin{proof}
Sum the pairwise orthogonal range projections at $v$ and take the
trace. Cyclicity assigns each range its initial weight $w_u$. The
matrix with entries $w_v^{-1}M_{vu}w_u$ has row sums at most one
and is similar to $M$. For two loops at one vertex the trace
inequality already reads $2w_v\leq w_v$.
\end{proof}

There is a tracial way to distinguish words if orthogonality is
needed only between the prepared vectors. Use the left regular
unitaries $U_i$ of the free group on the oriented step labels, with
$U_g\delta_h=\delta_{gh}$ and trace
$\tau(T)=\langle\delta_1,T\delta_1\rangle$. On group polynomials
this functional extracts the identity coefficient. It is positive
and cyclic since $gh=1$ exactly when $hg=1$. Distinct positive
words give orthogonal vectors $U_w\delta_1$. Each unitary has the
whole space as its range. For a physical outgoing step $a$ and a
physical return step $b$, the record $ba$ is distinct from the
empty word when $b$ is a separate generator. This models
unreduced histories in an enlarged word algebra. Using
$U_b=U_a^*$ instead erases the backtrack. One must specify which
operation represents path reversal and which represents the
physical adjoint before declaring this an action-derived register.

This distinction is measurable at the source. For a finite set
of distinct recorded words and real displacements $d_w$, put
$J(k)=\sum_w a_w e^{ik\cdot d_w}U_w$. Then
\begin{equation}
 \tau\bigl(J(k)^*J(k)\bigr)=\sum_w|a_w|^2 .
 \label{eq:monomial-word-inclusive-uv}
\end{equation}
The identity coefficient $\tau(U_w)$ vanishes for a nonempty
positive word, while the inclusive two-mark coefficient
$\tau(U_w^*U_w)$ equals one. Thus this parameter-free word
representation distinguishes a geometrically retraced trip and
preserves a positive trace. Its inclusive source loses the entire
phase modulus upon sewing.

\paragraph{An atomic displacement at any depth.}
The preceding hard-channel test extends beyond perfectly
orthogonal records. Let $\mu$ be a finite-variation, Hilbert-valued
Borel measure on $\mathbb R^4$ for the complete marked amplitude
in a specified positive channel, and put
$V(k)=\int e^{ik\cdot d}\,d\mu(d)$. For the hard cubes
$D_R=[R,2R]^4$,
\begin{equation}
 \lim_{R\to\infty}\frac1{R^4}\int_{D_R}\!\|V(k)\|^2\,d^4k
   =\sum_{d\in\mathbb R^4}\|\mu(\{d\})\|^2 .
 \label{eq:monomial-word-atom-uv}
\end{equation}
To see this, write $d\mu=h\,d|\mu|$ with $\|h\|=1$ almost
everywhere for the variation measure. In the squared norm the
averaged phase is the product of four interval averages for
$d_j-d'_j$. Its modulus is at most one. It tends to zero away
from $d=d'$ and equals one on the diagonal. Dominated convergence
leaves just the atoms of $|\mu|\otimes|\mu|$ on that diagonal,
which give the displayed sum. If the remaining positive sewing
factor on these cubes is bounded below by $c|k|^{-p}$ for large
$|k|$ and $p\leq4$, one surviving atom forces a divergent
integral because each of the disjoint dyadic cube interiors
contributes at least a constant times $R^{4-p}$. At $p=4$ the sum already
diverges. A nonzero displacement atom counts just as a zero
displacement atom does. Amplitudes at the same displacement are
combined before its norm is squared. An atomless measure passes
this necessary test, while the rate required for convergence
still depends on its Fourier decay. Applying the test to the
physical theory requires deriving this positive channel, its
measure and its lower sewing weight from the complete source.

\paragraph{A one-scale source with an inclusive Gaussian bound.}
A positive boundary state can retain a displacement with finite
resolution. Use locked units with $\hbar=1$ and unit invariant
speed magnitude. Set $\ell=\sqrt{G_N}$ for this comparison and define
$\gamma_\ell(y)=(2\pi\ell^2)^{-2}e^{-|y|^2/(2\ell^2)}$ on the
Euclidean chart. A displacement $d$ of law $\gamma_\ell(d)\,d^4d$
prepares the unit vector $\psi_d(x)=\sqrt{\gamma_\ell(x-d)}$ in
$L^2(\mathbb R^4,d^4x)$. Thus the two densities, for displacement
and for a reading conditional on it, use the same width. The
one-mark source vector is
\[
 V(k)=\int \gamma_\ell(d)e^{ik\cdot d}\psi_d\,d^4d .
\]
Its kernel has the quadratic Euclidean action
$(2|d|^2+|x-d|^2)/(4\ell^2)$. Summing over \emph{all} final
readings gives $\langle\psi_{d'},\psi_d\rangle
=\exp[-|d-d'|^2/(8\ell^2)]$. For two independent integration
variables of law $\gamma_\ell$, their difference has covariance
$2\ell^2I$. Its Gaussian integral therefore gives
\begin{equation}
 \|V(k)\|^2=\left(\frac23\right)^2
          \exp\!\left(-\frac23\ell^2|k|^2\right).
 \label{eq:monomial-resolution-uv}
\end{equation}
The suppression comes from the nonzero spread of the relative
displacement after the same final reading is paired across both
amplitudes. Each conditional pointer probability is normalized.

For independent tensor copies of this specified source, the
normalized squared norm of the product at marks $j$ is exactly
$\exp[-(2\ell^2/3)\sum_j |(Lk+b)_j|^2]$, where $L$ routes all
independent hard momenta to the marks and $b$ is fixed. If $L$
is injective, its least singular value is positive in this
finite graph. The exponential then integrates every fixed
polynomial in $k$ at infinity. Independence here is a property
of the specified comparison measure. The equality of the widths,
the tensor composition and the absence of vanishing widths on
joint short-duration faces would all need to follow from the
physical action. The calculation supplies an inclusive estimate
of this Euclidean source, leaving its identification with the
determinant and the Lorentzian physical current open.

\subsection{Closing the same physical photon source}
\label{sec:monomial-joint-uv}

\paragraph{What a homogeneous random kernel can change at photon sewing.}
Fix a reconstructed plane in infinite Euclidean volume. Let $A$ be
the centered transverse boundary field and $\xi$ the remaining
data at zero electric coupling. Use a jointly translation-invariant
probability law. Let $T_\xi[u,v]$ be the complete symmetric two-port
Taylor weight on real test connections, with signs from the
determinant, measure and direct source variations included.
Assume that the external and sewn ports use this same weight and
that the connected coefficient has zero first-order vacuum term.
Thus the measured bilinear form is $K=\mathbb E T_\xi$, and closing
the two ports on a finite test projection $P$ costs
$\tfrac12\mathbb E T_\xi[PA,PA]$, up to the common action sign.
These identifications must be checked from the source.

\begin{proposition}[Translation-invariant conditional kernels]
\label{prop:monomial-ergodic-sewing-uv}
Suppose that translations are ergodic on the marginal law of $A$:
every translation-invariant event of that field has probability zero
or one. Assume $T_\xi$ is a measurable continuous distributional
kernel, with integrable test coefficients, and is translation
covariant as a random kernel,
$T_{\tau_x\xi}=U_xT_\xi U_x^{-1}$, where $U_x$ translates a test
field. If almost every realization is itself translation invariant,
then $T_\xi$ is independent of $A$. In particular every finite
projection gives the direct contraction of $K$ with the photon
two-point covariance. This holds for signed kernel entries.
\end{proposition}
\begin{proof}
The invariant-kernel hypothesis makes $T_{\tau_x\xi}=T_\xi$ almost
surely for each $x$. For any event $B$ generated by its test
coefficients, joint stationarity makes $\mathbb E[\mathbf1_B\mid A]$
an invariant function of $A$, modulo null sets. Ergodicity makes this
function the constant $\Pr(B)$. A countable determining family of
test pairs suffices to recover the distributional kernel. Independence
now factors each finite sum of terms $\mathbb E[T_{ij}a_i a_j]$,
where $a_i$ are the real coordinates of $PA$. The result is
$\operatorname{Tr}(D_P K_P)$ before the symmetry factor, with
$D_P$ their covariance matrix.
\end{proof}

The infinite-volume free transverse Maxwell Gaussian has this
ergodicity property. For Schwartz tests the translated cross
covariance is the Fourier transform of an integrable function:
the transverse projector is bounded and $|q|^{-2}$ is locally
integrable in four dimensions. The cross covariance tends to zero
at large separation. Joint Gaussian characteristic functions of
finite test families then factor in the limit, and density of
cylinder observables gives mixing and hence ergodicity. This
argument uses no momentum cutoff. Combined with the positive-band
test above it excludes a change of the hard integral produced solely
by random translation-invariant multipliers, even when the conditional
photon law is non-Gaussian. Stationarity of the averaged kernel is
a weaker condition. Realizations that mix different momenta can
still have a translation-invariant average.

\paragraph{A computable test for momentum mixing.}
Take real coordinates $a_i$ on a finite photon test projection with
centered Gaussian marginal of covariance $D$. For square-integrable
$T_{ij}$ put $H_{ij}(a)=\mathbb E[T_{ij}\mid a]$. Relative to
direct sewing, the change before the symmetry factor is exactly
\[
 \sum_{ij}\mathbb E\!\left[(a_i a_j-D_{ij})H_{ij}(a)\right].
\]
Only the orthogonal projection of $H_{ij}$ onto the span of the
centered quadratics $a_k a_l-D_{kl}$ enters. When $H_{ij}$ is
twice weakly differentiable with the Gaussian integrability needed
for integration by parts, the same expression is
\[
 \sum_{i,j,k,l}D_{ik}D_{jl}\,
                 \mathbb E[\partial_k\partial_l H_{ij}(a)] .
\]
The first identity follows by conditioning. The second follows by
two integrations by parts in the Gaussian coordinates. Both are
finite-projection identities. Inferring an absolutely convergent
physical graph requires control as the test spaces exhaust the
continuum, including the geometric variables in the native
order of integration.

For example, take independent standard Gaussian vectors $a,z$ in
$\mathbb R^d$, $d\ge2$, and write $b=a+\sigma z$. The squared norm of the exterior product, $\|u\wedge b\|^2
=|u|^2|b|^2-(u\cdot b)^2$, tests the marked connection $u$.
The mean port matrix is $(d-1)(1+\sigma^2)I$, so independent
sewing gives $d(d-1)(1+\sigma^2)$. Actual sewing gives
$\mathbb E\|a\wedge b\|^2=\sigma^2d(d-1)$ because
$a\wedge b=\sigma a\wedge z$. The centered-quadratic term accounts
for the entire difference. At $\sigma=0$ the background probe
still has a nonzero Hessian while the closed bare interaction
vanishes identically. This finite test exhibits the need to
identify the propagating charged current and its cut using the same
source that performs the sewing.

\paragraph{Marking a photon pair in a single action monomial.}
A further test allows the kernel to mix momenta. Fix a perturbative
coefficient of the native source, with the counting variable
independent of the number of photon slots. Suppose that after
the native variables and their boundary preparation have been
integrated with their full background dependence, the coefficient
in question is a finite sum of photon polynomials
$V_r(a+h)=L_r[(a+h)^{\otimes r}]$. Here $a$ is a centered Maxwell
Gaussian, $h$ is an external connection, and $L_r$ is a symmetric
multilinear kernel on real smooth test fields, including the
signed weights of its diagrams. Assume that the coefficient of the connected
functional is $\mathbb E\sum_r V_r(a+h)$, with all lower-order
products needed for connectedness already accounted for. Thus the
background shifts every occurrence of the photon in this expression.
A single microscopic curvature monomial need not have this finite
polynomial form after the charged determinant is evaluated.

\begin{proposition}[Marked-pair identity for a common photon source]
\label{prop:monomial-marked-sewing-uv}
For each finite test projection let $D$ be its covariance and put
$K_r=\mathbb E\,\nabla^2V_r(a)$. Then
\[
       \operatorname{Tr}(D K_r)=r\,\mathbb E V_r(a).
\]
In the continuum, absolute convergence of every complete Wick
pairing of the symmetric kernels implies absolute convergence of
the traced sewing of the total measured Hessian. Consequently the
positive-band obstruction excludes this class if that Hessian and
the photon covariance satisfy the spectral and contact hypotheses
of the direct-sewing test.
\end{proposition}
\begin{proof}
Gaussian integration by parts gives
$\operatorname{Tr}(D\,\mathbb E\nabla^2V_r)
=\mathbb E[a\cdot\nabla V_r]=r\,\mathbb E V_r$
by homogeneity. Equivalently, in a pairing of $r=2m$ fields one can
mark and orient one of the $m$ pairs in exactly $2m$ ways. Deleting
the marked covariance produces the Hessian integrand with two
specified ports; sewing restores that covariance. Odd terms have
zero Gaussian expectation on both sides. In continuum momentum
space the absolute integral of a marked contraction, evaluated by
first integrating the unmarked variables and then taking the
absolute value, is bounded by the full absolute integral of its
parent pairing, including the finite polarization contraction in the
same order. Sum over the finitely many marks, pairings and degrees.
For a stationary vacuum use the coefficient per unit volume, with
the overall momentum delta removed. This argument uses fixed smooth
test projections for the algebra and absolute integration of the
actual continuum contraction kernels for the estimate.
\end{proof}

If graphwise convergence is claimed only for sums of diagrams
related by Ward identities, the same bound holds when every marked
sum is a marginal of an absolutely integrable completed sum in
a common momentum routing, with its signs and pair multiplicities
preserved. Absolute convergence obtained by cancellation between
different photon degrees alone gives no such estimate: marking
weights the degree-$r$ contribution by $r$. A polynomial at each
of infinitely many perturbative orders only gives this test
coefficient by coefficient.

For the exterior-product example the common shift changes both
slots to $a+h$ and $a+h+\sigma z$. Their squared area is exactly
$\sigma^2\|(a+h)\wedge z\|^2$. Integrating $z$ gives
$\sigma^2(d-1)|a+h|^2$, a quadratic to which the marked-pair
identity applies. At $\sigma=0$ its measured Hessian vanishes.
This distinguishes a probe that reads only one slot from the
physical background variation of a common action.

This argument also applies to infinitely many degrees at a fixed
physical order if $\sum_r r B_r<\infty$, where $B_r$ sums the
absolute integrals of the complete pairings at degree $r$. The
factor $r$ counts the allowed markings. If two covariant tempered
tensor responses on the reconstructed plane agree on separated-time
tests, their difference is supported on the equal-time hyperplane.
Its support is Lorentz invariant. The only point lying on every
boost of that hyperplane is the origin. A distribution supported at
one point is a finite sum of derivatives of a delta distribution,
whose Fourier transform is polynomial. Thus a proposed nonpolynomial
contact must change a separated-time response, or one of these
covariance or distributional hypotheses.

\paragraph{A joint contraction can pass at second order and fail at fourth.}
On one locked plane let a two-valued internal normal label $j=0,1$
be carried around a loop, with equal trace weights. Take two real
vector fields on each block. In the same natural units put
$\mu^2=G_N^{-1}>0$, $X=Q_E^2$, and on each transverse polarization use
the quadratic bare Hessian
\begin{equation}
 H_j(z;X)=
 \begin{pmatrix}(1+2jz^2)X+2j\mu^2&zX\\ zX&X+\mu^2\end{pmatrix}.
 \label{eq:monomial-doublet-uv}
\end{equation}
Here $z$ marks a curvature mixing. Its kinetic form is the sum of
$(F_B+zF_A)^2$ and $[1+(2j-1)z^2]F_A^2$, in the same Maxwell
normalization. The two mass terms are positive semidefinite and can
be written with Abelian Stueckelberg fields. For $|z|<1$ the kinetic
matrix is positive definite on both blocks. Conjugating the mass
matrix by its inverse kinetic square root exhibits nonnegative
squared masses and nonnegative spectral residues for transverse
observables. One zero mass persists on $j=0$. The label $j$ belongs
to the internal trace on this plane, independently of its N/A label.

Eliminating $B$ gives the inverse $A$ covariance
$D_j^{-1}+z^2K_j$, where
\[
 D_j=(X+2j\mu^2)^{-1},\qquad K_j=2jX-\frac{X^2}{X+\mu^2}.
\]
The first term is the contact from differentiating the same quadratic
form. The second comes from exchange of the massive vector.
After extracting the transverse Maxwell factor $X$, the mean scalar
response is $\mu^2/(X+\mu^2)$, with positive spectral weight $\mu^2$.
The Euclidean Hessian itself has residue $-\mu^4$ at $X=-\mu^2$
because the exchange enters with the Schur-complement sign. The matrix
element producing the pole is a neutral vector transition. It
supplies a positive-mass spectral atom, with no electric pair cut
in this example. The sewn blocks instead give
$D_0 K_0=-X/(X+\mu^2)$ and $D_1 K_1=X/(X+\mu^2)$.
Thus the traced second-order integrand vanishes pointwise, even though
both the mean external kernel and the mean covariance are positive.
The same internal label must be used at the two ends of each
contraction to obtain this equality.

For the normalized internal trace of the relative log determinant
(the vacuum effective action has the additional factor $1/2$), the
two blocks now give exactly
\[
 \frac12\sum_{j=0}^1\log\frac{\det H_j(z;X)}{\det H_j(0;X)}
   =\frac12\log\left[1-z^4\left(\frac{X}{X+\mu^2}\right)^2\right].
\]
The fourth-order integrand tends to $-1/2$ at large $X$.
The fourth-order ring therefore diverges even after this internal
trace and both native contacts. Each identity is per mode before
the momentum integral. Replacing the internal trace by a random
macroscopic choice between two vacua would define a different
generating functional.

\subsection{Photon sewing in the preceding comparison sources}
\paragraph{Photon sewing after the charged transfer intervals are retained.}
We first keep the carrier durations in the two-band determinant
and follow them through the dynamical Maxwell field. For this
calculation hold the normal coordinate fixed, or prescribe fixed
widths $a,\sigma$. Use zero-temperature transfer on the spatial continuum. Smear every electromagnetic
vertex on its spatial slice by the harmonic readout
$w(p)=e^{-a|p|}$, $a>0$, and keep the hopping width $\sigma>0$
fixed. Use Coulomb gauge for the photon. The instantaneous Coulomb
energy is the square of the charge density normal ordered with
respect to the onsite filled vacuum, with its usual factor $1/2$.
The square includes the self-energy of an individual charge. The
coupling $q$ is the electromagnetic expansion parameter; $g$ is the
off-diagonal hopping amplitude in $t(k)=g e^{-\sigma^2|k|^2}$.

\begin{proposition}[An all-orders soft bound for the ordered two-band source]
\label{prop:closing-uv-ordered-Maxwell}
For $2|g|<m$ the matter logarithm has an absolutely summable
expansion in closed hopping paths at each fixed positive order in
$q$, per unit spacetime volume with the external momentum delta
removed. Mark each complete charged path by the current of its
spatial and temporal links, then sew the marked paths with dynamical
transverse and instantaneous Coulomb photons. For each connected vacuum or two-point graph with no internal photon
bridge, the integrals over path displacements, charged durations and
remaining independent photon four-momenta converge absolutely. In particular, for external
transverse polarizations the proper correction at each electric order
is $O(\omega^2+|p|^2)$ as $(\omega,p)\to0$.
\end{proposition}
\begin{proof}
At zero background the onsite inverse on the two slots propagates
on oppositely oriented carrier intervals. Its kernel at carrier
separation $u$ is bounded by $e^{-m|u|}$ times a spatial delta.
A real Euclidean scalar connection attaches the usual unitary Wilson
phase to each such interval: this is the inverse of the onsite
first-order operator with covariant carrier derivative. The onsite
vacuum has zero normal-ordered density and contributes no scalar
response on its own. Each hopping mark is instantaneous and has
kernel bounded by $|g|G_\sigma(z)$ with the normalized spatial
heat density $G_\sigma$ and displacement $z$. Its Wilson phase is
integrated on the straight spatial segment. The log expansion around
the two onsite bands therefore consists of closed paths alternating
these horizontal and vertical segments, with trace weight bounded
by $2/n$ for length $n$ before the $n$ hopping factors. The linear
term vanishes because one hop changes the slot. One way to see the
determinant identity before taking the translation-invariant
density is to restrict the hopping to a finite time interval on a
spatial torus and impose the onsite vacuum endpoint state. The two
off-diagonal products of an onsite inverse and a hopping kernel
are Hilbert--Schmidt there, by the Gaussian and the one-sided time
exponential. Their product is trace class, and the ordinary
determinant after eliminating one slot generates precisely the even
closed words. The occupied onsite reference propagates with its
vacuum scalar phase fixed by the indicated charge operator.
For the continuum coefficient base the path at one point and impose
$\sum_j z_j=0$, $\sum_j u_j=0$. Integrate one closing density using
each delta. The other $n-1$ hops integrate to one each, and their
carrier intervals cost at most $(2/m)^{n-1}$. The resulting
summable bound contains $|g|^n$. It also permits any fixed polynomial
in $n$ and in the total spacetime length
$\ell=\sum_j(|u_j|+|z_j|)$: choose a sufficiently small $\eta>0$
and use $\exp(\eta\ell)$ in the same argument. The tilted Gaussian
is bounded and integrable, its mass tends to one as $\eta\downarrow0$,
and the tilted time cost tends to $2/m$. This proves the claimed
convergence of the hopping expansion with any fixed set of marks.

For one complete path a transverse mark inserts its line current
$J_T$ on the horizontal segments. The constant vector potential
sums to zero around the spatial projection of the path. With the
basepoint subtracted, the difference of the two plane waves costs
at most $(|\omega|+|p|)\ell$. Thus, in fixed reference units,
\[
 |J_T(\omega,p)|\le C(1+\ell)^2\min(1,\sqrt{\omega^2+|p|^2}).
\]
For a density mark on a vertical segment the spatial position $x_j$
is fixed. Since $\oint e^{i\omega t}\,\mathrm dt=0$ even for
$\omega=0$, one can subtract the same spatial phase at an anchor
$x_*$ from every segment. With $t_j$ its starting time this gives
\[
 J_0(\omega,p)=\sum_j(e^{ip\cdot x_j}-e^{ip\cdot x_*})
                  \int_{t_j}^{t_j+u_j}e^{i\omega t}\,\mathrm dt .
\]
The differences are at most $\min(2,|p|\ell)$ and each time integral
is at most $\min(|u_j|,2/|\omega|)$ when $\omega\ne0$. It follows
that
\[
 |J_0(\omega,p)|\le C n(1+\ell)^2\,
       \frac{\min(1,|p|)}{1+|\omega|}.
\]
The full spatial current before transverse projection also obeys
$\omega J_0+p\cdot\mathbf J=0$, by integrating the total differential
of the plane wave around the closed spacetime path. The common spatial
readout preserves this identity. Differentiating the Abelian phase
gives a product of these currents;
all marks can occupy the same link. The Coulomb sewing follows
directly from the Hamiltonian density square. Between hops the
signed densities are stationary in space. Integrating the Gaussian
scalar field over those time intervals gives the exponential of
minus one half the time integral of their Coulomb quadratic form.
Self-contractions are included. Equivalently the scalar covariance
is independent of frequency, with spatial multiplier $|p|^{-2}$
before the two readouts. This explains why the temporal line
integrals must be retained even for an instantaneous interaction.

After the polynomial length costs have been absorbed, put
$Q^2=\omega^2+|p|^2$. The two possible internal lines are bounded by
\[
 \begin{aligned}
 &e^{-2a|p|}\frac{\min(1,Q)^2}{Q^2},
 \qquad
 e^{-2a|p|}\frac{\min(1,|p|)^2}{|p|^2(1+|\omega|)^2}.
 \end{aligned}
\]
Both functions are bounded and integrable in
$\mathrm d\omega\,\mathrm d^3p$. For the first use
$\int_\mathbb R(\omega^2+|p|^2)^{-1}\,\mathrm d\omega=\pi/|p|$.
For the second the two integrals factor. Choose independent photon
momenta on chords of a spanning tree of the connected graph of
charged paths. Each chord uses its integrable envelope and each tree
edge its supremum. In a proper graph every internal photon belongs to
a cycle, so zero spatial momentum of an internal photon occurs only
on a null set of chord variables. The two external transverse marks
supply the two remaining soft factors. There are finitely many sewn
topologies at any fixed positive electric order after the charged
paths have been summed, and a connected vacuum normalization removes
empty unmarked components.
\end{proof}

The ordered-source bound suppresses a Euclidean transverse mass term
at each finite order. To keep the charged frequencies as integration
variables, differentiate the vacuum-normalized determinant with its
onsite reference intact. Already for two marks omitting that reference
produces the divergent single trace in \eqref{eq:closing-uv-density-hard}.
The positive Hamiltonian
uses a preferred carrier, a prescribed spatial heat hopping and a
harmonic charge profile. A physical chiral GCM derivation must
identify its own temporal links, Coulomb self-energy and boundary
charged slots.

\paragraph{The temporal density in the determinant of the prepared vacuum.}
The hard-frequency test in \eqref{eq:closing-uv-density-hard} concerns
one charged trace. The two-band vacuum functional in
Proposition~\ref{prop:closing-uv-ordered-Maxwell} also specifies
what happens if all hopping is switched off while the scalar
background is kept on. Evaluating that functional first determines
which trace enters the response of this particular action.

Work at fixed $a,\sigma,m>0$ in the time-local two-band source.
Write $D_\pm[\phi]=\partial_\tau+iqw\phi\pm m$, where $w\phi$
denotes the common spatial readout. Let $S_\pm[\phi]$ be their
onsite vacuum inverses (oppositely directed carrier propagators).
The off-diagonal hopping from one slot to the other is $T[A]$,
with its conjugate in the opposite block. On the symmetric vacuum
slab, restrict the hopping to the slab as in the preceding path
construction. Normal ordering the charge with respect to the filled
onsite vacuum gives
$\rho(x)\lvert0\rangle=0$ in the sense of spatial smearing.
The evolution at zero hopping therefore has vacuum amplitude one
for every smooth scalar background, after the constant onsite
vacuum energy is fixed. In particular that amplitude is known
before any photon derivatives are taken.

The Schur formula for the matter factor on this slab reads
\begin{equation}
 \mathcal M_g[A,\phi]=
 \det\!\left( I-S_-[\phi]T[A]^*S_+[\phi]T[A]\right).
 \label{eq:native-relative-schur-uv}
\end{equation}
The determinant here is ordinary: the product of the two
Hilbert--Schmidt hopping blocks is trace class. One can verify the
formula by block elimination in finite subspaces and convergence
of the product in trace norm, or by the even trace expansion.
The factor which would otherwise multiply it is the onsite
amplitude just evaluated. A division by $\mathcal M_g[0,0]$
only fixes a field-independent normalization. Formula
\eqref{eq:native-relative-schur-uv} defines the charged part
without assigning an infinite determinant to either diagonal
differential operator.

\begin{proposition}[Joint four-momentum bound for a complete two-band matter loop]
\label{prop:native-relative-density-uv}
In the translation-invariant vacuum matter coefficients of this
action, every fixed positive number of external electric marks
(including spatial Wilson contacts) admits an absolutely integrable
representation in the original loop variables $(\omega,k)$. The
bound is uniform in the external energy and momentum shifts at
fixed mark order and widths. Here the temporal vertices use the
common onsite normalization in \eqref{eq:native-relative-schur-uv}.
\end{proposition}
\begin{proof}
Between marks let $\xi_j=\omega+b_j$, $k_j=k+c_j$, with shifts
given by the external momenta in their cyclic order. Put
\[
 \begin{split}
 R_\lambda(\xi,k)
   &=(i\xi+m\sigma_3+\lambda t(k)\sigma_1)^{-1},\\
 t(k)&=g e^{-\sigma^2|k|^2},\qquad 0\le\lambda\le1 .
 \end{split}
\]
For an all-scalar word with $n$ marks the completed trace,
up to its common charge, cyclic factor and spatial readouts, is
\begin{equation}
 \operatorname{tr}\!\left[
    \prod_{j=1}^n R_1(\xi_j,k_j)
    -\prod_{j=1}^n R_0(\xi_j,k_j)\right].
 \label{eq:native-relative-word-uv}
\end{equation}
To see where the second product comes from, turn on the physical
hopping as $\lambda T$ in the Schur determinant. At $\lambda=0$
its logarithm and all its scalar derivatives vanish. The
$\lambda$ derivative of the logged full block is the trace of
its inverse with one hopping insertion. After the $n$ scalar
derivatives, cyclicity puts this insertion at each of the $n$
resolvents of the scalar word, with the same derivative and
cyclic factor as in $\partial_\lambda\operatorname{tr}\prod R_\lambda$.
Integrate in $\lambda$. Equivalently, the even Schur series removes
exactly the pure onsite word. In translation-invariant coefficients
the inverses in this computation are the gapped vacuum inverses
displayed above. They are obtained from the filled negative band.
On a fixed torus the onsite sea has nonzero overlap with this
vacuum: the sum of squared band mixing angles is finite by the
Gaussian decay of $t$, and no single-mode overlap vanishes.
The bounds below also justify the differentiated interpolation
directly in infinite-volume momentum variables.

With $F(\xi)=(m^2+\xi^2)^{-1/2}$ one has
$\|R_\lambda(\xi,k)\|\le F(\xi)$ and
$\|\partial_\lambda R_\lambda(\xi,k)\|
\le |t(k)|F(\xi)^2$. In each differentiated summand the $k$
integral therefore costs at most $\|t\|_1$ times constant
matrix factors. Keep the squared $F$ at that insertion and
bound the others by $m^{-1}$. The $\omega$ integral of $F^2$
is $\pi/m$, independent of its translation. This proves absolute
convergence of the difference in \eqref{eq:native-relative-word-uv}
for arbitrary external shifts, including simultaneous hard
spatial momenta and carrier frequencies.

For example put $\xi'=\omega+\nu$, $k'=k+p$,
$t'=t(k')$, $t=t(k)$, and
$L(\xi,v)=\xi^2+m^2+v^2$. With common external factors stripped
off, the complete two-mark trace is
\begin{equation}
 \begin{aligned}
 &\frac{2(m^2+tt'-\omega\xi')}
        {L(\omega,t)L(\xi',t')}
 -\frac{2(m^2-\omega\xi')}
        {L(\omega,0)L(\xi',0)} .
 \end{aligned}
 \label{eq:native-relative-bubble-uv}
\end{equation}
The trace at zero hopping has zero frequency integral at each
spatial momentum. Each of its two scalar fractions before taking
the slot trace is a product of two same-band propagators, whose
poles lie in the same half-plane and whose product is
$O(|\omega|^{-2})$. Its absolute integral over all $(\omega,k)$
diverges. The Schur identity prescribes the difference at common
loop momentum before either four-dimensional integral is formed.
As a check, write $E=(m^2+t^2)^{1/2}$, $E'=(m^2+t'^2)^{1/2}$,
and $D=E+E'$. Resolving the two propagators into their spectral
projectors and integrating $\mathrm d\omega/(2\pi)$ gives
\[
 -\left(1-\frac{m^2+tt'}{EE'}\right)\frac{D}{D^2+\nu^2}.
\]
Only transitions between opposite bands contribute. With the
conventional fermion loop sign restored, this is the positive
occupied-to-empty response obtained from the same Hamiltonian.

For a word with spatial marks, each block of $d\ge1$ marks on a
single Wilson hop has the two-slot matrix factor $\sigma_1$ and a
spatial multiplier which, up to fixed charges and phases, is an
average of $d$ directional derivatives of $t$ at translated
arguments. This follows by differentiating the straight segment
phase before Fourier transformation. The multiplier has $L^1$
and $L^\infty$ bounds uniform in the translations, by the
corresponding bounds on Gaussian derivatives. For at least two
propagators, use one marked hop in $L^1$, every remaining vertex
in $L^\infty$, and any two translated $F$ factors in $L^2$ for
the frequency integral. With a single contact and just one
propagator, the only slot trace is
$\operatorname{tr}(R_1(\omega,k)\sigma_1)
=2t(k)/L(\omega,t(k))$. Its frequency integral is absolutely
bounded as well. The onsite reference has no spatial derivative.
There are finitely many orderings and contact partitions at the
specified order.
\end{proof}

The prepared onsite sea fixes this result. Its scalar-dependent
inverse is part of the original Schur factor and its vacuum
amplitude is fixed for every field. No adjustable photon term
enters \eqref{eq:native-relative-bubble-uv}. Replacing an
unspecified physical chiral determinant by this Schur factor would
be an additional choice of theory. For that determinant one must
derive the reference amplitude and its charged boundary state from
the native action.

\paragraph{A Coulomb contraction tests a second independent hard frequency.}
The two-density estimate at fixed external transfer has a precise
limitation. At fixed widths take $qg\ne0$ and write
$R_g(v,k)=(iv+h(k))^{-1}$ and $R_0(v)=(iv+m\sigma_3)^{-1}$.
Set $\eta=v+\Omega$. With the charge and readout factors removed, put
\[
 \begin{aligned}
 B(\eta,v;k,p)
 ={}&\operatorname{tr}\big[R_g(\eta,k+p)R_g(v,k)\big]\\
    &-\operatorname{tr}\big[R_0(\eta)R_0(v)\big] .
 \end{aligned}
\]
Here $v$ is the charged frequency and $\Omega$ the frequency of an
instantaneous Coulomb line joining the two density marks. For fixed
external $(\Omega,p)$ the relative density derivative is absolutely
integrable in $(v,k)$. Sewing it in the displayed routing also
integrates over $\Omega$. Since $i\eta R_g(\eta,k+p)\to I$,
\begin{equation}
 \lim_{|\eta|\to\infty}\eta B(\eta,v;k,p)
 =\frac{2v\,t(k)^2}{(v^2+m^2)(v^2+m^2+t(k)^2)}.
 \label{eq:aud-coulomb-sewing-uv}
\end{equation}
The limit is uniform on compact sets of $(v,k,p)$. Choose a positive
interval for $v$ and spatial balls with $p$ away from zero and with
nonzero readout. On suitable compact subsets the right side has a
strictly positive minimum. The instantaneous Coulomb multiplier
$q^2|w(p)|^2/|p|^2$ is bounded below there as well. Hence the absolute
integral of this routed coefficient diverges at least as
$\int^\infty\mathrm d\Omega/\Omega$. The same conclusion holds for
the spatial torus by fixing two admissible spatial momenta.
A single invertible affine change of the loop variables preserves this
absolute divergence.

The hard region here has one bounded charged frequency and one
unbounded Coulomb frequency. It differs from the common charged
frequency limit at fixed external transfer. The canonical ordered
vacuum amplitude specifies how the two current insertions occur on
a complete charged path. In a momentum proof, the permutations of
marks on each hopping interval must be retained before taking
absolute values. Such a proof concerns the graph generated by that
specified grouping. It cannot establish absolute integrability of
$B$ with the Coulomb line in the routing displayed above.

The value of the same Coulomb graph has a direct positive expression
at every real hopping strength. Let $E_k=(m^2+t(k)^2)^{1/2}$,
$n(k)=(t(k),0,m)/E_k$, and define the transition weight
$W(k,p)=|n(k+p)-n(k)|^2/4$. This is the trace of the lower-band
projector at $k$ times the upper-band projector at $k+p$.
For $\ell_E(v)=2E/(v^2+E^2)$ the convolution identity reads
\[
 \int \frac{\mathrm dv}{2\pi}\,
 \ell_{E_k}(v)\ell_{E_{k+p}}(v+\Omega)
 =\ell_{E_k+E_{k+p}}(\Omega).
\]
It follows by Fourier transforming $e^{-E|t|}$.
The right side is the sum of the two Euclidean time orderings of
the physical transition. The positive frequency representation of
the Coulomb vacuum energy per spatial volume at order $q^2$ is
therefore
\begin{equation}
 \begin{aligned}
 \mathcal E_{\rm C}^{(2)}
 =\frac{q^2}{2}\int
 &\frac{\mathrm d^3p\,\mathrm d^3k}{(2\pi)^6}
       \frac{|w(p)|^2}{|p|^2} W(k,p)\\
 &\times\int\frac{\mathrm dv\,\mathrm d\Omega}{(2\pi)^2}
           \ell_{E_k}(v)\ell_{E_{k+p}}(v+\Omega).
 \end{aligned}
 \label{eq:aud-positive-coulomb-uv}
\end{equation}
The inner integral equals one by positivity and normalization.
Since $\int W(k,p)\,\mathrm d^3 k\le
|p|^2\|\nabla n\|_2^2/4$, the spatial integral converges at both
ends for the harmonic readout $w(p)=e^{-a|p|}$.
This coefficient also follows by evaluating the positive Coulomb
density square in the filled hopping vacuum, including the
diagonal self-energy. It is strictly positive for $qg\ne0$.
Equation~\eqref{eq:aud-positive-coulomb-uv} is an absolute
representation obtained from the spectral projectors of this
Hamiltonian. It leaves intact the hard-frequency counterexample
for the separately routed trace $B$.

\paragraph{Absolute convergence with the charged four-momenta retained.}
The ordered estimate in Proposition~\ref{prop:closing-uv-ordered-Maxwell}
can be strengthened. Work on one selected plane at fixed positive
$m,a,\sigma$, with the neutral onsite preparation of the two-band
action. Use the ordinary determinant for the even hopping words of
that preparation. Each active charged loop then contains at least
two spatial heat hops. The complete scalar vertex on a carrier
interval includes every ordering of its insertions. This last fact
matters for a photon exchanged instantaneously along that same interval.

\begin{proposition}[Joint four-momentum estimate for the complete hopping vertices]
\label{prop:physical-joint-hopping-uv}
Suppose $2|g|<m$. At each finite electric order, each connected graph
of complete hopping loops and transverse or Coulomb photons, with no
internal photon bridge, has an absolutely integrable momentum
representation on the spatial continuum. It contains one independent
charged four-momentum per loop and every independent photon
four-momentum. The estimate is summable in the number of hops and
holds per unit spacetime volume after the overall delta is removed.
The same ultraviolet conclusion holds on a spatial torus with the
stated Gauss zero-mode convention. No charged frequency must be
integrated first.
\end{proposition}
\begin{proof}
For a loop with $n$ hops keep the spatial closure of its displacements
$z_j$ and the temporal closure of its oriented intervals $u_j$ as
Fourier integrals in $k\in\mathbb R^3$ and $\xi\in\mathbb R$.
The onsite factors, in the Fourier convention $e^{isu}$, are
$R_\pm(s)=(\pm m-is)^{-1}$. A scalar insertion with carrier
frequency $\nu$ on an interval gives
$(e^{i\nu u}-1)/(i\nu)$, with its starting phase accounted for
separately. For several such insertions on that interval the exact
multiplier, before integrating $\xi$, is
\[
 \left(\prod_j \frac{T_{\nu_j}-I}{i\nu_j}\right)R_\pm(\xi+b),
 \qquad (T_\nu f)(s)=f(s+\nu).
\]
Here $b$ is a real shift supplied by the other phases.
At zero frequency use the derivative limit. A single divided
difference is the product of the two adjacent resolvents.
For several marks the expression is their sum over all orders:
partition the product of interval integrals into ordered simplices,
or apply the single-mark identity successively. For example a
Coulomb pair of frequencies $\nu,-\nu$ on one interval gives
$[2R(s)-R(s+\nu)-R(s-\nu)]/\nu^2$.
Estimating each ordered summand separately can lose a power of
the hard frequency.

Choose reference units for the next bounds. For $|\nu|>1$ the
two translations in a divided difference cost $2/|\nu|$.
For $|\nu|\le1$ use the integral of a translated derivative.
Mixed differences commute, even when a partial sum of the hard
frequencies vanishes. Thus every scalar end contributes a factor
bounded by $C/(1+|\nu|)$ times translates of resolvents and their
derivatives. On any loop two distinct interval factors can be
estimated in $L^2(\mathrm d\xi)$ and the others in $L^\infty$.
These bounds are uniform in all translations. For an unmarked
interval the norms are $(\pi/m)^{1/2}$ and $m^{-1}$ respectively.
The derivative bounds cost constants depending on $m$ and the number
of marks. A horizontal Wilson mark is an integral over the unit
segment of a spatial phase times a component of $z_j$. Fourier
transformation therefore gives a derivative of a translated
$e^{-\sigma^2|k|^2}$. One heat factor on each loop is in
$L^1(\mathrm d^3k)$, including all its fixed-order derivatives;
the others are bounded, and an undifferentiated heat factor is
bounded by one. The same uniform translated Gaussian sum applies
in a spatial box. This already controls the charged variables with
all photon momenta kept independent.

To cover arbitrarily soft photons in the continuum, choose the
following routing of each complete conserved current before Fourier
transformation. At a scalar end with $|p|\le1$, subtract the phase
at one fixed spatial vertex from the spatial phase on every vertical
segment. The removed current is proportional to
$\oint e^{i\nu t}\,\mathrm dt=0$ for every closed carrier path.
Each resulting spatial difference supplies $|p|$. At a transverse
end with $Q=(\nu^2+|p|^2)^{1/2}\le1$, subtract the common basepoint
phase on every horizontal segment. The removed vector is proportional
to $\sum_j z_j=0$. Each difference now supplies $Q$. Explicitly,
$e^{iy}-1=iy\int_0^1e^{i\theta y}\,\mathrm d\theta$; after the
momentum factor is taken out, $y$ contributes a linear combination
of signed displacements and, for the transverse mark, carrier
intervals. Their coordinates differentiate the corresponding
Gaussian or resolvent. The remaining phases translate its momentum.
There are only finitely many additional derivatives at a given
electric order. These identities hold on each complete path with
its closure constraints in place. They specify the Fourier
representative used at soft momenta. In a fixed box the spatial
zero mode is omitted and these soft subtractions are unnecessary for
the ultraviolet bound.

For clarity, each expansion term now factorizes in the charged
variables into translated differentiated Gaussian and interval
factors. Fourier inversion of the closure convolution is justified by
the same $L^1$ spatial and $L^2$ times $L^2$ temporal bounds.
Its absolute integral in each loop's four-momentum is at most
$C_M n^{d_M}m^{1-n}$ times its marked photon factors, where $M$ is
the fixed electric order. The polynomial counts mark positions and
the choice of a coordinate in each basepoint difference.
The factor $|g|^n$ from the hopping makes the sum over $n$
convergent. At fixed positive $M$ only finitely many charged loops
can occur in a connected graph. A zero-order vacuum loop is
estimated by the same two unmarked interval norms.

The remaining transverse and scalar lines have the respective
envelopes
\[
 e^{-2a|p|}\frac{\min(1,Q)^2}{Q^2},
 \qquad
 e^{-2a|p|}\frac{\min(1,|p|)^2}
                     {|p|^2(1+|\nu|)^2}.
\]
Both are bounded and integrable in four dimensions. Use the photon
momenta off a spanning tree of the graph of charged loops as independent
variables. Integrate the envelope for each of these chords and use the
supremum for each tree edge. Self-edges are chords. Every internal
line in the stated graphs varies with at least one independent
momentum, so the values at the isolated zero denominators do not
affect the result. The estimates bound the absolute joint integrand
also in regions in which charged and Coulomb frequencies grow
together. All widths here occur in the original source.
\end{proof}

The determinant and vertex prescription in this proposition are
specific. The onsite vacuum already has unit scalar amplitude before
hopping is turned on. The raw temporal bubble of
\eqref{eq:closing-uv-density-hard} uses a different determinant
coefficient and still fails its absolute-integrability test.
Even within a hopping loop the proposition takes the absolute value
after the vertex has included its allowed orders. A proposed
physical chiral realization must derive its own vacuum functional,
charged block and complete vertices for this proof to apply.

\paragraph{Charged endpoints, including paths with zero hops.}
At fixed spatial volume the same action also controls the photon
exchange between open charged lines. Use Coulomb gauge, the common
readout and the zero spatial mode convention of
Proposition~\ref{prop:cl-gauss-charge-uv}. Charged intermediate
sectors then have the uniform compensating background on the torus.
Take external elementary charged fields in oppositely charged
pairs, with fixed Euclidean momenta or smooth wave packets. Photon
momenta below refer to amputated external probes.

\begin{corollary}[Finite-volume charged and photon functions at each bare order]
\label{cor:physical-joint-open-uv}
At fixed positive widths on the spatial torus and $2|g|<m$, every
connected graph of the prepared two-band source at any finite
electric order has an absolutely summable and integrable
representation in all its independent charged and photon
four-momenta, for any fixed number of external charged endpoints
and photon probes. This includes one-photon-reducible graphs and
open lines on which the charge never hops. The constants may depend
on the volume.
\end{corollary}
\begin{proof}
The Gaussian fermion integral with Grassmann sources multiplies
the Schur determinant \eqref{eq:native-relative-schur-uv} by
the exponential of the source bilinear in the inverse of the same
two-slot operator. Complete the Grassmann square, or differentiate
the vacuum generating functional in the endpoint sources. Writing
$S=\operatorname{diag}(S_+,S_-)$ and $V$ for the off-diagonal
Wilson hopping block, the open inverse is
\[
 (I+SV)^{-1}S=\sum_{n\ge0}(-SV)^n S .
\]
For real background fields the absolute onsite kernel is dominated
in carrier time by a one-sided exponential of mass $m$.
The spatial kernel of $V$ is dominated by $|g|$ times a normalized
periodized heat density. Hence $\|S\|\le m^{-1}$ and
$\|V\|\le |g|$ in the slab or on the time line. Expanding the
inverse gives $n$ spatial hops and $n+1$ carrier intervals. At
fixed endpoints there is no closure constraint on their total
displacement; the external Fourier factors specify the momentum
entering one end of each open line. Momentum conservation at the
other end fixes the sum of the photon transfers on that line.
There is no further independent charged four-momentum on it.

Apply the divided-difference identity of
Proposition~\ref{prop:physical-joint-hopping-uv} on every interval.
Each scalar end supplies $(1+|\nu|)^{-1}$ times bounded translates
or derivatives of its onsite resolvent. This holds also if the same
interval contains both ends of an instantaneous exchange. For an
open line use the supremum in its frequency and spatial arguments
for every factor. With $M$ marks, the resulting bound on the line,
apart from its photon factors, is at most
$C_M(1+n)^{d_M}|g|^n m^{-(n+1)}$ for some $d_M$. In particular
$n=0$ has a finite bound and still includes all scalar orders on
the one remaining interval. A horizontal mark requires a hop and
cannot occur at $n=0$. Closed lines are bounded in their own
independent four-momenta by the preceding proposition.

Momentum balance on each open or closed line now leaves one photon
four-momentum per chord of a spanning tree of the connected graph
of charged lines. Self-edges are chords. At fixed $R$, for
$p\in R^{-1}\mathbb Z^3\setminus\{0\}$ the line envelopes
\[
 \frac{e^{-2a|p|}}{\nu^2+|p|^2},
 \qquad
 \frac{e^{-2a|p|}}{|p|^2(1+|\nu|)^2}
\]
are bounded and integrable with respect to frequency and the
spatial lattice sum. Use their integrals on chords and their
suprema on tree edges, including photon bridges whose momenta are
fixed by the external data. A line with excluded spatial momentum
is absent. No assumption on soft cancellation at an open endpoint
is required. At fixed $M$ the mark counts are polynomial in each
hop number, their geometric sums converge, and the finite number
of endpoint pairings and contact partitions leaves the bounds
finite. With momentum packets one may finally integrate the fixed
external data against their absolute weights. A free external
two-point factor, when present alone, is bounded at the specified
Euclidean momentum with the same zero-mode convention.
\end{proof}

This corollary uses the canonical endpoint field of the specified
Hamiltonian. It makes no assumption about a charged Lorentz pole.
A volume-uniform infrared estimate for charged endpoints would
require additional control of the soft Coulomb field; the
finite-volume ultraviolet bound uses the entire unbounded momentum
lattice.

\paragraph{A common geometric depth in the full electric vacuum functional.}
The estimates also allow the electric source to reach arbitrarily
short normal distances. On a spatial torus let $a=r$ and
$\sigma^2=\kappa r$, $0<r\le1$, $\kappa>0$. Draw one $r$ for the
whole vacuum preparation with law $\mu$. A sufficient condition is
$\int e^{c/r^3}\,\mu(\mathrm dr)<\infty$ for some $c>0$.
For example the unperturbed radial ground state in the nonzero
angular sector of Proposition~\ref{prop:cl-gauss-charge-uv}
has this property. Use a finite symmetric vacuum slab, with hopping
multiplied by a smooth even function $h$ supported in its interior.
A convenient preparation is
\[
 h(t)=\frac{\int \varphi(s+t)\varphi(s)\,\mathrm ds}
            {\int\varphi(s)^2\,\mathrm ds},
 \qquad 0\le\varphi\in C^\infty_c(\mathbb R),\quad\varphi\ne0 .
\]
Then $0\le h\le1$, and the Fourier density of $h$, with inverse
Fourier normalization, is a probability density. Its frequency
transfers are therefore integrable with mass one even in words
containing arbitrarily many hoppings.

\begin{corollary}[Common-depth electric vacuum normalization and marked graphs]
\label{cor:physical-electric-common-depth-uv}
If $2|g|<m$, the complete Maxwell--Coulomb vacuum partition
function on this slab is positive and finite for every real electric
charge. At each finite electric order the absolute joint graph
estimates are integrable also in the shared depth. These statements
use one normalization after the depth integral.
\end{corollary}
\begin{proof}
At fixed $r$ the unexpanded logarithm of the matter vacuum in real
transverse and auxiliary scalar fields is an absolutely convergent
sum over the closed hopping paths above. Every phase has modulus
one. Pin the first hopping in the slab and use a single periodized
heat density in the spatial closure. Its supremum is at most
$C(1+r^{-3/2})$, by its Fourier series. All other spatial densities
have mass one. The charged time integrals of a word of length $n$
are bounded by $(2/m)^{n-1}$ times the slab duration. All switching
factors have modulus at most one. The resulting bound for the
absolute logarithm is $C_{V,h,m,g,\kappa}(1+r^{-3/2})$.
Conjugate words have conjugate phases, as in
Theorem~\ref{thm:closing-uv-positive-susceptibility}. Hence the
vacuum factor is real and strictly positive; the same upper bound
on its exponential holds after the Maxwell and Coulomb integrations.
The exponential moment of $\mu$ dominates this growth at the tip
and proves finiteness of the joint normalization.

For a marked graph expand each $h$ by its Fourier density. The new
arguments only translate the charged frequencies in
Proposition~\ref{prop:physical-joint-hopping-uv}. For a connected
term the overall time constraint eliminates one switching frequency
using the bounded density; all remaining switching frequencies
integrate with mass one, after solving the photon tree constraints.
All spatial Gaussian derivatives cost bounded powers of
$\sqrt{\kappa r}$ at small $r$. Each charged loop's spatial sum
costs at most $C(1+r^{-3/2})$. Integrating a transverse photon
chord in frequency costs $\pi/|p|$; summing it with its two readouts
costs $C(1+r^{-2})$. For a Coulomb chord the integrable frequency
envelope leaves the spatial cost $C(1+r^{-1})$. Tree-edge suprema
are bounded on the mean-zero box. Thus every fixed-order connected
term costs at most $C_M(1+r^{-1})^{B_M}$ after summing the hops,
for some finite $B_M$. Products of components obey the same type of
bound. The exponentiated zero-charge determinant multiplies these
coefficients by at most $\exp[C(1+r^{-3/2})]$, also when several
charged loops visit the tip together. Integrating in $r$ is therefore
absolute. The positive constant term of the joint normalization
permits its formal quotient at every finite order.
\end{proof}

Here the common depth is constant during each charged history and
is integrated only once. A dynamical radial trajectory in carrier
time changes the interval operators; the proof above uses the
constant-depth direct integral.

For the unexpanded normalization alone the small-hopping condition
can also be removed on a slab of duration $T$ with constant
Hamiltonian between its two vacuum endpoints. At fixed $r$ let $v$
be the normalized photon and onsite matter vacuum. The pair-form
bound in Proposition~\ref{prop:cl-gauss-charge-uv} and the
nonnegative Coulomb energy give
$H(r)\ge-Cg^2(1+r^{-3/2})/m$ on the fixed box, uniformly in the
real electric charge. The vacuum has zero mean energy in the form
sense, using the same vacuum energy convention as above. Jensen's
inequality for its spectral measure therefore gives
\[
 1\le \langle v,e^{-T H(r)}v\rangle
 \le \exp\!\bigl[C T g^2(1+r^{-3/2})/m\bigr].
\]
The same $\mu$ normalizes this joint vacuum for all real $g,q$.
This last bound concerns the full transfer at finite $T$; the
preceding absolute sum of marked hopping words uses $2|g|<m$.

\paragraph{Removing the hopping restriction at fixed infrared volume.}
The carrier frequency of an instantaneous Coulomb photon has no decay
from its free covariance. Gauge covariance of the complete source
allows that frequency to be controlled before either loop integral is
formed. It also extends the joint estimate to elementary open charged
lines and to finitely many internal spinor slots.

\begin{theorem}[Bare charged and photon graphs at arbitrary hopping strength]
\label{thm:ct-uv-gauged-graphs}
Fix a spatial torus of period $2\pi R$, omit the spatial harmonic photon
and Coulomb modes, and use the neutral onsite preparation. Take the same
spatial readout $w(p)=e^{-a|p|}$, $a>0$, on every electric component.
On a finite-dimensional ordinary fermion space of equal electric charge
let $\beta=\beta^*=\beta^{-1}$ and let
\[
 \begin{gathered}
 h(k)=m\beta+T(k),\qquad m>0,\\
 T(k)=T(k)^*,\qquad \beta T(k)+T(k)\beta=0 .
 \end{gathered}
\]
Assume $T$ is a finite matrix sum of Schwartz symbols, with each spatial
hop gauged by the straight Wilson transport of this common readout.
Use the local carrier derivative and the vacuum boundary determinant of
the onsite filled negative slots. At every finite electric order, the
connected bare vacuum, photon and elementary charged Green functions
at fixed external Euclidean momenta admit complete graph integrals
absolutely convergent jointly in all independent charged and photon
momenta. This holds for every finite hopping amplitude. No subtraction
of an adjustable photon term enters the result.
\end{theorem}
\begin{proof}
First fix the scalar field $\phi$ and the transverse potential. Choose
an even smooth real function $\theta(\nu)$ which is zero for
$|\nu|\le\nu_*$ and one for $|\nu|\ge2\nu_*$, with an arbitrary
$\nu_*>0$. A gauge change of the charged variables can cancel the
part $\theta\widehat\phi$ of the temporal potential. Take the gauge
parameter to be its carrier-time primitive, with the sign fixed by
the temporal connection. Its measured value in the charge phase is
$w\chi$, since the spatial readout commutes with both derivatives.
The potential in the transformed scalar vertex is
$(1-\theta)\widehat\phi$. A spatial Wilson vertex also sees a
longitudinal potential proportional to
$p\,\theta(\nu)\widehat\phi/\nu$. An elementary charged endpoint
receives its compensating phase; differentiating it gives a factor
proportional to $w(p)\theta(\nu)/\nu$. This is a change of the
matter variables for each photon field. The Maxwell and Coulomb
Gaussian laws still integrate every original mode.

For smooth sources the transformation follows directly by conjugating
the carrier derivative and the Wilson links. Their endpoint factors
cancel at each internal sewing. A time-slab onsite vacuum is invariant
under these charge phases: their multiplication operators commute with
the constant onsite occupied projection and act on the particle and
hole Fock factors. They fix the vacuum and introduce no measure
factor depending on the electric source.
Alternatively, take test fields smooth in frequency with the transformed
part away from zero; their time primitives vanish at both infinities.
Differentiation at zero source gives the transformed complete vertices.
Open fields must retain the endpoint factors in this identity.

For $T_\lambda=\lambda T$, $0\le\lambda\le1$, one has
$h_\lambda^2=m^2+\lambda^2 T^2$. Thus all the dressed charged
inverses obey the bound $F(\xi)=(\xi^2+m^2)^{-1/2}$. At fixed
volume the filled negative slots of $h$ indeed give the vacuum reached
by long onsite preparations. In each off-diagonal singular-vector pair
of $T(k)$ the band rotation has angle of magnitude
$\tfrac12\arctan(t/m)$, where $t$ is that singular value.
Unpaired slots do not rotate. The sum of the squared angles is finite
by Schwartz decay, the product of their cosines is nonzero, and
the sum of the magnitudes of the sea energy shifts is finite by
$\sqrt{m^2+t^2}-m\le t^2/(2m)$. The excitation gap is
at least $m$. This gives the stationary vacuum representation for
each finite hopping amplitude.

The relative determinant calculation in
Proposition~\ref{prop:native-relative-density-uv} applies to these
slots. For an all-scalar closed word the onsite-normalized difference
is the integral of its $\lambda$ derivative. One differentiated
inverse supplies $F^2\|T(k+d)\|$ for a shifted momentum $k+d$.
This gives a finite frequency integral and a finite spatial lattice
sum uniformly in all shifts. Whenever at least one spatial Wilson
vertex is present, its multiplier is an average of derivatives of the
Schwartz hop with translated arguments. One such vertex supplies the
uniform spatial sum. For a loop of two or more propagators keep two
translated $F$ factors in $L^2$ and all others in $L^\infty$.
For a single Wilson contact $V$ the term proportional to $i\xi$ in
$\operatorname{tr}[(i\xi+h)^{-1}V]$ vanishes: $(\xi^2+h^2)^{-1}$
is block diagonal in $\beta$ and $V$ is off diagonal. The remaining
trace is bounded by a fixed slot multiplicity times
$\|h\|\,\|V\|/(\xi^2+m^2)$. Consequently the charged four-momentum
of each closed component is integrable uniformly in all photon shifts.
The free closed vacuum energy was included above. On an open line the
external momentum fixes the initial fermion momentum. After the photon
momenta are chosen, each internal charged inverse costs at most $1/m$.
All spatial contact multipliers are bounded there.

It remains to bound the photon momenta. Taking absolute values only
after the gauge completion just described, the possible factors at
one scalar end, apart from $w(p)$ and the bounded matter vertices,
are at most
\[
 (1-\theta(\nu))+
       (1+|p|)\,|\theta(\nu)/\nu|
 \qquad (0\le\theta\le1).
\]
The extra one in the second term includes elementary endpoints.
The squared expression times $w(p)^2/|p|^2$ bounds a Coulomb line.
It is bounded and integrable in
$\mathrm d\nu$ times the spatial lattice sum over $p\ne0$:
the low-frequency part has compact frequency support, the high part
is $O(|\nu|^{-2})$, the spatial sum has exponential decay, and the
smallest allowed $|p|$ is $1/R$. The transverse line
$w(p)^2/(\nu^2+|p|^2)$ has the same two properties.
Contract each charged open or closed component to a vertex, choose
a spanning tree of the connected graph, and integrate the line
envelopes on its chords. Each tree line uses its supremum. A tree
line forced into the omitted spatial mode contributes zero. Momentum
conservation along an open chain fixes all shifts from its endpoint
and its incident photons. At finite electric order there are finitely
many contact partitions and transformed endpoint marks. These bounds
control the joint absolute integral, including coincident hard
frequencies. The arbitrary $\nu_*$ only chooses a representative
of the gauge identity.
\end{proof}

For fixed external order the constants here use finitely many
Schwartz seminorms of $T$, a spatial photon sum, and inverse powers
of $m$. With $m$ bounded below, $a=r$, and $T$ a finite sum of
polynomials times $e^{-c r|k|^2}$, $c>0$, whose coefficients grow
at most polynomially in $1/r$, these constants grow at most as an
inverse power of $r$ for $0<r\le1$ at fixed $R$. This follows by
Gaussian scaling, the uniform sum of translated Gaussians, and the
corresponding exponential shell sums for $w$. A geometry-dependent
vacuum normalization must still be included in any integration over
$r$. The fixed-box theorem asserts ultraviolet convergence; uniform
thermodynamic or soft bounds require control of the spatial modes
approaching zero.

There is also a background-field inverse estimate in this family.
On the full carrier line take real smooth bounded measured potentials;
the spatial hopping may be switched on only in a slab to implement
the neutral preparation. Its Hermitian Wilson kernel remains bounded
by the integrable absolute hopping kernel and anticommutes with
$\beta$. With the sign convention
$D=\partial_\tau+iqB_0+m\beta+T[B]$, one has
\[
 \beta D=m+S,\qquad
 S=\beta\partial_\tau+iq\beta B_0+\beta T[B].
\]
Here $S$ is skew-adjoint on the domain of the carrier derivative:
$\beta\partial_\tau$ is skew-adjoint and both other terms are
bounded skew-adjoint perturbations. The spectral theorem gives
$\|D^{-1}\|\le1/m$ for every such background and hopping strength.
When the hopping has finite carrier support in a finite spatial box,
the two off-diagonal onsite-inverse times hopping blocks are
Hilbert--Schmidt. The estimate excludes zeros of their ordinary Schur
determinant there. It does not provide a uniform bound on the
derivatives of that determinant after the photon field is integrated.

\paragraph{A continuum soft estimate for the same dressed charged loop.}
For a closed charged component one can also control the spatial photon
momenta approaching zero at arbitrary hopping strength. The point at
which the hopping is varied supplies a center for a gauge choice.
This gives the requisite soft factors while the independent charged
four-momentum is still present.

\begin{proposition}[Continuum bound at arbitrary hopping for closed components]
\label{prop:ct-uv-rooted-soft}
Under the fixed-width spinor and preparation assumptions of
Theorem~\ref{thm:ct-uv-gauged-graphs}, work in the zero-temperature
spatial continuum and remove the overall momentum delta. Every
connected bare vacuum or photon graph of complete matter vertices
with no internal photon bridge admits a joint absolutely integrable
representation in the charged and photon four-momenta at each finite
electric order, for every finite hopping strength. Each external
transverse mark also admits a factor
$\min(1,(\nu^2+|p|^2)^{1/2})$ at soft Euclidean momentum $(\nu,p)$.
\end{proposition}
\begin{proof}
Scale the covariant hopping to $\lambda T[B]$,
$0\le\lambda\le1$, for the measured Abelian connection $B$.
The relative matter logarithm is zero at $\lambda=0$ for each scalar
source. Its derivative in $\lambda$ has density
\[
 \operatorname{tr}_{\rm slots}
        (D_\lambda[B]^{-1}T[B])(X,X)
\]
at its base point $X$. A background-independent vacuum normalization
drops out of positive electric orders. Both kernels in this expression
transform by the endpoint charge phases. The density is invariant
at its given $X$ before the base point is integrated. The equal-time
jump of the inverse is proportional to the slot identity times the
spatial delta, whose trace against the off-diagonal hopping vanishes.
The marked version of this identity also follows by differentiating
the vacuum inverse with compactly supported test sources. At zero
source its kernels are the dressed resolvents used above.

For each incident mode transform the measured connection prior to
differentiation. First use the temporal primitive of the preceding
proof for its high-frequency scalar component. Leave the resulting
spatial mark, of size $w(p)|p|\theta(\nu)/|\nu|$, in that gauge.
For a remaining scalar mode with $|p|\le1$, or a transverse mode
with $Q=(\nu^2+|p|^2)^{1/2}\le1$, put that component alone into
radial gauge about $X$. Explicitly the connection appearing at $Y$
then has components
\[
 B^{(X)}_\mu(Y)=(Y-X)^\alpha\int_0^1 s\,
 F_{\alpha\mu}(X+s(Y-X))\,\mathrm ds .
\]
Subtracting the derivative in $Y$ of the segment integral of $B$
from $X$ to $Y$ gives the displayed formula. A scalar plane wave
in the temporal component has field strength of magnitude bounded
by $|p|$ times its amplitude. A transverse plane wave has field
strength bounded componentwise by a constant times $Q$ times its
amplitude. The factor $w(p)$ of the original plane wave persists.
The gauge changes for all selected modes can be superposed since
the connection is Abelian. They change the charged variables only;
the external measure need not realize the radial connection as a
separate photon configuration. Components outside the indicated
regions are left in their original gauge.

For precision consider a fixed coefficient of the rooted density.
Differentiation of $D^{-1}$ produces finite words of unforced
resolvents and temporal or Wilson vertices, with a hopping at
the root even if every electric mark is temporal. A radial-gauge
mark adds one coordinate of $Y-X$ to its field strength. Along a
rooted word, this coordinate is a sum of coordinates of the successive
edge displacements, with a fraction of the last displacement if
$Y$ lies inside a Wilson segment. The remaining plane waves translate
their Fourier variables, also when the parameter $s$ is integrated.
Thus every coordinate factor differentiates a hopping symbol, a
marked spatial symbol, or a dressed resolvent. Temporal displacement
of an instantaneous hop is zero. This construction uses the rooted
kernel without shifting a difference of two divergent loop traces.

Here are bounds that also justify these kernel manipulations. The
unforced carrier kernel at positive or negative time separation $u$
is, up to its one-sided sign, the corresponding spectral projector of
$h_\lambda(k)$ times $e^{-|u|\sqrt{m^2+\lambda^2T(k)^2}}$. Its
onsite spatial delta part remains at $k=\infty$. The remainder is
Schwartz in $k$, with every fixed seminorm bounded by a polynomial
in $|u|$ times $e^{-m|u|}$ times a constant depending on the fixed
hop. Indeed the projectors are
$(I\pm h_\lambda(m^2+\lambda^2T^2)^{-1/2})/2$; the gap bounds their
derivatives. Spatial derivatives of the exponential can be taken by
the product integral for its derivative, with the two factors of
the exponential bounded by their respective gapped intervals.
The zeroth difference from the onsite exponential follows by the
same interpolation. Thus the position kernels, including the
onsite delta as a measure, have the polynomial displacement moments
needed at each order. The root hop is spatially smooth and enforces
one integrable spatial factor at closure. These estimates justify
the radial identity also for plane waves: one can first multiply the
gauge functions by smooth compactly supported functions equal to one
on growing balls about $X$, then pass to their jets in the marked
kernel by these moment bounds. Terms from derivatives of the spatial
or temporal localization are bounded by the same polynomial moments
and disappear away from the center in this limit.

In momentum variables all derivatives of the root hop have translated
$L^1(\mathrm d^3k)$ bounds; other spatial vertices use their suprema.
Differentiating a resolvent yields products of resolvents and bounded
derivatives of $h_\lambda$, so every fixed joint derivative is
bounded by a constant times $F(\xi)$ uniformly in spatial momentum.
With two or more resolvents use the two translated $L^2$ frequency
bounds from the theorem. With one propagator, all marks lie on the
root hop. Its trace against any resulting off-diagonal vertex has
frequency decay $O((m^2+\xi^2)^{-1})$ as established there.
A differentiated propagator either has at least two inverse factors
or retains the same cancellation. All constants are uniform in
$\lambda$ and in the shifts supplied by the other photons. Integration
over the base points restores the momentum deltas between charged
components.

The transverse and scalar photon lines therefore have the respective
bounded, integrable envelopes
\[
 \begin{aligned}
 &w(p)^2\min(1,Q)^2/Q^2,\\
 &w(p)^2\frac{(1+|p|)^2\min(1,|p|)^2}
                   {|p|^2(1+|\nu|)^2},
 \end{aligned}
\]
up to fixed constants.
The factor $1+|p|$ covers the high-frequency scalar mode transferred
to a longitudinal spatial mark. Its square is absorbed by the spatial
exponential. These estimates apply simultaneously to all marks.
Spanning-tree integration now works in the spatial continuum. On a
graph with no internal photon bridge each internal line varies with
a chord momentum, so its value at the single omitted spatial mode
does not change the integral. This proves absolute convergence on
the full momentum domain and the stated external soft suppression.
\end{proof}

The base-point cancellation uses a closed charged trace. The finite-box
estimate for externally charged lines includes endpoint phases, and
does not by itself give an infinite-volume charged scattering state.

\paragraph{Charged endpoints in the spatial continuum.}
An open charged line can respond to a constant potential through its
endpoints. Its soft behavior can nevertheless be estimated with all
charged frequencies present. Impose the assumptions of
Theorem~\ref{thm:ct-uv-gauged-graphs} in the spatial continuum,
at zero temperature with the filled onsite reference. Use elementary
charged fields in Coulomb gauge and amputate external photon legs.
In a fixed connected term contract each charged component to a
vertex and keep its incident photon edges. Mark the vertices that
carry external sources. Let $b_v\in\mathbb R^3$ be the net external
spatial momentum entering a marked vertex, with $\sum b_v=0$
after removal of the overall momentum delta. Call the transfers
nonexceptional for this term when no nonempty proper set of marked
vertices has vanishing total $b_v$. This condition is empty for one
marked vertex. A photon bridge separating a component without any
external source has spatial momentum identically zero. Such a term
is absent with the specified harmonic mode removed; below one
can equivalently restrict to graphs without such a bridge.

\begin{theorem}[Continuum bare coefficients with charged endpoints]
\label{thm:seam-uv-open}
For every finite hopping strength and every finite electric order
the connected bare graphs with open charged lines have jointly
absolutely convergent Euclidean charged and photon momentum
representations at nonexceptional transfers. They also act on
joint Schwartz packets in all independent external momenta by
absolutely convergent integrals through the exceptional transfer
sets. The fixed-transfer constants can be chosen uniformly on
compact subsets of the nonexceptional region for a fixed graph.
\end{theorem}
\begin{proof}
Use the dressed open inverses and endpoint factors of
Theorem~\ref{thm:ct-uv-gauged-graphs}, and the rooted, gauge-invariant
density of Proposition~\ref{prop:ct-uv-rooted-soft} on each closed
component. These are identities for the respective functionals of
one background before differentiating them. They can therefore be
used on the two ends of a photon even if one component is open and
the other closed. On an open line the initial external momentum
and all photon transfers determine the intermediate momenta. The
gap bounds its inverses uniformly. The integrals of the absolute
charged loop kernels on closed components are bounded uniformly in
all their photon transfers as in the rooted proof.

For an internal photon let $b\in\{0,1,2\}$ count its ends on
open components, with both ends counted on a self-edge.
In fixed reference units and up to constants, a common majorant
after the charged estimates is
\begin{equation}
 W_b(\nu,p)=\frac{e^{-c|p|}}{(1+|\nu|)^2}
              \big(1+|p|^{-b}\big),\qquad 0<c<2a .
 \label{eq:seam-uv-open-envelope}
\end{equation}
Indeed a transverse line supplies $w(p)^2/Q^2$ and a factor
$\min(1,Q)$ at each closed end, where $Q^2=\nu^2+|p|^2$.
For an instantaneous scalar line the time-primitive identity
supplies $(1+|\nu|)^{-1}$ per end, up to a spatial polynomial,
and the spatial radial-gauge estimate also gives
$\min(1,|p|)$ per closed end. Any fixed polynomial at large
$|p|$ is absorbed in the exponential. In four-momentum measure
all three majorants are in $L^1$. Moreover $W_0$ is bounded and
$W_1\in L^2$, since $\int_{|p|<1}|p|^{-2}\,\mathrm d^3p<\infty$.

Fix nonexceptional transfers. If there are at least two marked
vertices, take $\epsilon>0$ so small that the number of internal
edges times $\epsilon$ is smaller than the minimum norm of the
proper, nonempty marked subset sums. For a routing, keep all
edges with $|p|\ge\epsilon$. All marked vertices lie in one
connected component of these edges: otherwise the balance on one
component containing a proper set of marks would express its
injection as a sum of at most that many boundary transfers of
norm below $\epsilon$. Thus there exists a tree of kept edges
connecting all marks, possibly through other vertices. Sum upper
bounds over its finitely many possible choices. With one mark,
use just that vertex. On the region for each choice, every tree
factor is bounded uniformly. Solve its momentum constraints and
contract it to one vertex, called the root. The tree momenta are
linear combinations of the remaining momenta and the external
injections, with unit Jacobian in the usual cycle coordinates.
After replacing the tree factors by their bounds the other
variables can range over their full domains.

An edge from the root to itself now integrates independently and
costs at most $\|W_b\|_1$. Remove the root and consider one of the
remaining connected components. All its internal edges have two
closed ends. Choose a spanning tree inside it, bound its factors
by $\|W_0\|_\infty$, and integrate the remaining internal chords
at cost $\|W_0\|_1$ each. Conditional on its $d$ signed transfers
from the root, the only unsolved constraint is their total. Each
root edge has at least one closed end, so its majorant is bounded
by a fixed multiple of $W_1$. For $d\ge2$ the convolution at
any prescribed total costs at most
\[
 C^d\|W_1\|_2^2\|W_1\|_1^{d-2}
\]
by Cauchy--Schwarz and Young's inequality. A component with $d=1$
is precisely a source-free bridge case excluded above. This proves
the fixed-transfer assertion, including simultaneous soft transfers
and arbitrarily hard carrier frequencies.

For the packet assertion choose one external four-momentum at each
marked vertex as its transfer coordinate, subtracting the signed
sum of the other external momenta there. After the overall delta
there is one fewer independent transfer than marked vertices.
This linear change of variables has unit absolute Jacobian. Bound
the packet by a rapidly decreasing integrable envelope in the
remaining relative external momenta, uniformly in the transfers.
Integrate the separate vertex balances to eliminate the free
transfer coordinates. This identifies all marked vertices with one
root in the remaining momentum constraints without division by any
external transfer. The root-edge and convolution bounds just
proved apply. In particular this estimate does not assign a finite
point value to a Coulomb exchange at a forced zero external
transfer. It defines its action on the packet.
\end{proof}

For the charged two-point function with no external photons, there
is a single marked vertex. Its Euclidean coefficient is therefore
finite at every real external four-momentum in this representation.
For a four-point function the transfer restriction has a physical
origin. Set $T=0$ and take two distinguishable charged particle slots.
At first Coulomb exchange their amputated direct interaction is
proportional to $q_1q_2|w(p)|^2/|p|^2$. With nonzero charges it
diverges as the spatial transfer $p$ tends to zero, while its
integral against a smooth compactly supported packet is finite.

For a different observable one can see the endpoint cancellation
directly. Couple a neutral two-point probe to
$D_B^{-1}(X,Y)U_B(Y,X)$, where the return transport follows the
straight spacetime segment with the same measured Abelian
connection $B$. The product is invariant under background gauge
changes. At fixed $X$ put each soft component into radial gauge
about $X$, as in Proposition~\ref{prop:ct-uv-rooted-soft}.
Its line integral on the straight return then vanishes, since
$(Z-X)^\mu B_\mu^{(X)}(Z)=0$. Each such mark on the remaining
inverse gains the same field-strength factor as on a closed
component. For a high-frequency scalar mode the temporal primitive
moves the coupling to a longitudinal spatial potential on both
the line and the return; the two endpoint phases cancel. A mark
on the return contributes a component of $Y-X$. In Fourier
space it differentiates the open kernel or a vertex along the
chain, with continuously translated momenta along the segment.
The gap and the spatial Schwartz bounds control every fixed
number of these derivatives; the position kernel has the required
polynomial displacement moments by the rooted-kernel estimate.
Thus complete insertions of this specified neutral probe have the
bounded $W_0$ type of soft envelope even at their open endpoints.
This construction attaches the return to the observable. Theorem~
\ref{thm:seam-uv-open} uses the elementary charged fields and
requires no return attached to them. Both statements concern
Euclidean coefficients away from a charged mass-shell continuation.

\paragraph{The limit of the spatial boxes at fixed electric order.}
The continuum integrals just obtained can be recovered from the
harmonic-mode-free boxes themselves. Periodize a single fixed
Schwartz hopping kernel and attach its straight Wilson transport to
each lifted displacement before summing the lifts. Sample the same
harmonic multiplier $w$ on the reciprocal lattices. Keep the mass,
width, and hopping fixed as the period $2\pi R$ tends to infinity.
These hypotheses specify how the box actions are compared.

\begin{corollary}[Thermodynamic limit of complete Euclidean coefficients]
\label{cor:seam-uv-open-boxlimit}
In Theorem~\ref{thm:seam-uv-open}, take lattice external spatial
momenta tending to fixed nonexceptional values. Remove the overall
translation delta with its usual volume normalization. At each
finite electric order the closed-component-completed graph
coefficient converges to its continuum value. The result includes
the elementary charged two-point coefficients for all real external
momenta. The same convergence holds against joint Schwartz
external packets, with the spatial packet integrals replaced by
normalized lattice sums in each box. External transverse photon
polarizations, where specified, are taken continuous near their
fixed nonzero spatial momenta.
\end{corollary}
\begin{proof}
First compare one closed charged vertex at the allowed input photon
momenta. Use the hopping-interpolation trace in
Proposition~\ref{prop:ct-uv-rooted-soft} and the temporal primitive
of Theorem~\ref{thm:ct-uv-gauged-graphs}; retain the original
spatial gauge in this comparison. The root supplies a marked or
unmarked Schwartz hop. For periodic photon modes the multiplier
of a Wilson mark in the box is precisely the continuum chord
multiplier sampled at the lattice fermion momentum: summing over
all lifts unfolds the spatial integral. With Fourier convention
$e^{ikx}$ the loop measure in the box is $(2\pi R)^{-3}\sum_k$,
and in space it is $(2\pi)^{-3}\int\mathrm d^3 k$.

After extracting the readouts and temporal-primitive factors, every
fixed number of spatial derivatives of the word has integrable
absolute spatial and carrier-frequency norm, uniformly under photon
momentum translations. One derivative of a resolvent inserts a
bounded derivative of $h$ between two inverses. Derivatives of the
root and marked hops retain their Schwartz bounds. Use the two
translated $L^2$ carrier-frequency factors, or the off-diagonal trace
cancellation in a one-propagator word, exactly as in
Theorem~\ref{thm:ct-uv-gauged-graphs}. In the latter case the
canceling diagonal trace is identically zero in the spatial
momentum and its spatial derivatives vanish as well. The charged
loop at zero electric order has the same property after its free
vacuum reference is included. Fourier integration by parts and
Poisson summation in the spatial loop variable therefore bound the
sum of its nonzero image corrections by $C_N R^{-N}$ for every
$N>3$, times the extracted factors. The constants do not require
the photon transfers to remain bounded. The open-line multipliers
are sampled directly because their initial momenta are fixed.
This argument avoids applying a nonperiodic radial gauge
transformation on a finite spatial torus.

Subtract products of vertices one vertex at a time. To bound the
resulting errors, drop all soft Ward gains and use $W_2$ from
\eqref{eq:seam-uv-open-envelope} on every internal photon. Its
normalized lattice sum and frequency integral are uniformly finite
when $p=0$ is omitted. Its supremum is $O(1+R^2)$. Choosing a
spanning tree bounds the total error by $C_N R^{-N}(1+R^2)^{d}$
with $d$ no larger than the number of internal edges. Choose $N>2d+3$.
This includes winding corrections of the closed charged vertices.
There remain the continuum vertices sampled on the photon grid.

For these sampled vertices the proof of the theorem applies to
normalized lattice counting as well as to Lebesgue measure.
Indeed for any $0<\eta<1/2$, shell counting in three spatial
dimensions, with the single zero mode omitted, gives uniform
$L^1$ bounds on $W_2^{1+\eta}$, $L^2$ bounds on $W_1^{1+\eta}$
and an $L^1$ and supremum bound on $W_0^{1+\eta}$. Here
$W_1^{1+\eta}$ also has a uniform $L^1$ bound. Cauchy--Schwarz and
convolution use the same normalized measure as each independent
chord sum. The tree and root argument consequently gives a
uniform $L^{1+\eta}$ estimate for the photon integrand constructed
from the integrated continuum vertices. Multiplying every envelope
by $(1+|p|+|\nu|)^\gamma$ for sufficiently small $\gamma>0$
still gives the required bounds, hence uniformly vanishing tails
in the independent momenta.

On a compact set of these variables, a nonbridge photon spatial
momentum varies with coefficient $1$ or $-1$ in at least one
independent chord. The region where its norm is below $\epsilon$
has normalized grid measure at most $C(\epsilon+R^{-1})^3$.
The uniform $L^{1+\eta}$ bound and H\"older's inequality make
all such neighborhoods negligible. Under the nonexceptional
hypothesis any retained bridge stays away from zero. On the
complement the integrated vertices are continuous in their
momenta by the translated kernel bounds, and ordinary Riemann
summation gives the asserted limit. For external packets use their
free transfer coordinates as in the theorem. Any bridge depending
on external transfers then has the same small-neighborhood
estimate, and rapid decrease controls the remaining external
coordinates. This proves coefficientwise convergence. Constructing
an interacting limiting state or continuing a charged line to its
mass shell requires further information.
\end{proof}

\subsection{Which reference determinant can the physical vacuum remove?}
\label{sec:reference-test-uv}
The fixed-width calculation raises a question for a physical chiral
domain. Can its filled reference sea be divided out in the presence
of an electric source? The answer can be read from a positive
particle--hole matrix element before any ultraviolet estimate.
Throughout this argument the spatial charge profile and the
physical adjoint are the ones used in preparing the sea.

\begin{proposition}[Scalar response of a filled reference sea]
\label{prop:reference-test-uv}
Consider a gapped quasifree reference Hamiltonian $h_0$ with filled
negative-energy projection $P$. Use its occupied vacuum, and let a
real scalar source act by a bounded self-adjoint one-body operator
$Q$ and its vacuum-normal-ordered second quantization. The normalized
vacuum functional is identically one for every time profile of every
allowed scalar source if $[P,Q]=0$ for each source. Conversely,
where the vacuum response exists, this identity requires
$(1-P)QP=0$ for every such source. The assertion applies first to
finite orbital spaces. In the continuum the same necessity holds
on spectral windows where the particle--hole matrix elements and
the smeared response exist.
\end{proposition}
\begin{proof}
Let $J(Q)$ denote the normally ordered density. The component of
$J(Q)|0\rangle$ in the one-particle, one-hole sector is the matrix
$(1-P)QP$. All number-preserving terms annihilate the vacuum.
If this matrix vanishes, $J(Q)|0\rangle=0$; time evolution in
the reference action with any sequence of scalar insertions
therefore leaves the normalized vacuum amplitude equal to one.
Conversely at a positive time $t$ the two insertions of the same
real probe have connected correlation
\[
 \big\|e^{-t(H_0-E_0)/2}J(Q)|0\rangle\big\|^2 .
\]
A vanishing scalar functional derivative forces this norm to zero.
The exponential has trivial kernel and the vacuum one-point
function vanishes. A nonzero matrix element on a bounded spectral
window already gives a strictly positive contribution. This also
explains the window formulation for a continuum sea.
\end{proof}

For a translation-invariant reference with conventional electric
density, a Fourier component of the charge shifts $k$ to $k+p$ with
coefficient $q w(p)I$. Commutation with all real sine and cosine
probes near $p=0$, where $q w(p)\ne0$, requires
$P(k+p)=P(k)$ almost everywhere for those transfers. A measurable
matrix invariant under all translations in a neighborhood is
constant almost everywhere (convolve each entry with a smooth
approximate identity). The onsite sea used above meets this test.
It can still support a nontrivial scalar response once the charged
hopping has been switched on.

The usual four-component massive Dirac sea provides a direct check.
In units with limiting speed one let $h_D(k)=\boldsymbol\alpha\cdot k+
m\beta$ and $E_k=(m^2+|k|^2)^{1/2}$. Write $P_\pm=(I\pm h_D/E_k)/2$.
The spin sum in its electric response is
\[
 \operatorname{tr}[P_+(k+p)P_-(k)]
   =1-\frac{m^2+(k+p)\cdot k}{E_{k+p}E_k}.
\]
At zero external frequency the two orders of the density marks
multiply this by $2q^2|w(p)|^2/(E_{k+p}+E_k)$ under the
$\mathrm d^3k/(2\pi)^3$ integral. On any closed angular cone where
$|p|^2-(\widehat k\cdot p)^2$ is bounded below by a positive
constant, the trace is
$(|p|^2-(\widehat k\cdot p)^2)/(2|k|^2)+O(|k|^{-3})$ for fixed
$p\ne0$. Thus the positive static response has a logarithmic
ultraviolet tail whenever $q w(p)\ne0$. Dividing out this
source-dependent Dirac determinant would discard a physical
response. The calculation explains why the choice of reference in
the chiral theory has to follow from its quantization and boundary
state. It also provides a test for applying a Schur reduction: the
reference may carry its own response, in which case that response
must be kept in the complete coefficient.

\paragraph{A charged spinor with the Dirac spatial derivative at small momentum.}
\label{par:reference-spin-uv}
The finite-slot construction admits a spinor whose leading
dispersion and electric current on the sector plane agree with the
massive Dirac ones. This gives a more discriminating test of how
the charged geo-phase enters the action. At a fixed geometric
depth $r>0$ take four charged slots, an onsite mass $m\beta$, and
the spatial heat displacement of lapse $\kappa r$, $m,\kappa>0$.
Contract its first derivative with the Hermitian Dirac matrices
$\alpha_i$, with $\{\alpha_i,\beta\}=0$ and
$\{\alpha_i,\alpha_j\}=2\delta_{ij}I$. Transport each displacement
by the same spatial Wilson link as the charged spinor. Its
zero-field Fourier symbol, in carrier-length units, is
\begin{equation}
 \begin{aligned}
 h_r(k)&=m\beta+v\,\boldsymbol\alpha\cdot k\,e^{-\kappa r|k|^2},\\
 E_r(k)^2&=m^2+v^2|k|^2e^{-2\kappa r|k|^2}.
 \end{aligned}
 \label{eq:reference-spin-uv}
\end{equation}
Here $v$ fixes the low-momentum slope in these units. The exponential
is the characteristic function of the action's spatial heat step:
integration by parts in its displacement gives the additional $k_i$.
The readout $e^{-r|p|}$ is the harmonic boundary map used on every
component of the electric connection. In particular, a constant
spatial connection translates $k$ in $h_r$ and gives current
$q\partial_{k_i}h_r$. The leading derivative is $qv\alpha_i$ at
$k=0$ and the squared dispersion agrees with $m^2+v^2|k|^2$
through order $|k|^2$. The carrier and geometric depth break
boost symmetry in the higher powers. This construction fixes a
comparison representation; identifying $v$ or the physical
curvature pairing on a PE-CGG-HOUQFT branch requires the native
stationary equations.

There is an explicit realization of the spinor bilinear by the
four-arrow GCM curvature product in~\eqref{eq:g5-arrow-product}.
Its component trace is $\operatorname{Re}(\bar\zeta VU\xi)$.
In the spatial Hamiltonian take $\xi=\psi(t,y)$,
$\bar\zeta=\psi^\dagger(t,x)$, $U$ the oriented spatial Wilson
transport from $y$ to $x$, and $V=-iv\sum_i\alpha_i\partial_{z_i}
G_{\kappa r}(z)$ at $z=x-y$, where $G$ is the normalized heat
density. Integrate over the lifted displacements. Exchanging endpoints
takes the term to its adjoint, so the real pairing gives the
Hermitian hopping. A diagonal copy with $V=m\beta,U=I$ gives the
onsite Hamiltonian. On the real carrier-time action the local jet
$iD_t\psi$ in the column gives the canonical time term, with the
spatial terms entering with the Hamiltonian sign. The integrated
time jet uses the usual adjoint boundary pairing. Vacuum transfer
then gives the Euclidean inverse $\partial_\tau+h_r$ at zero
field. The Abelian Maxwell block and the spatially common charge
readout give the transverse field and the Gauss square already
derived above. The copies in the curvature product are evaluations
of one spinor, so its Fock boundary contributes one charged
determinant. This specifies a represented curvature action and
its boundary preparation. The argument neither derives the choice
of spatial curvature contraction nor selects this vacuum from the
unrestricted physical branch.

\begin{corollary}[Bare spinor coefficients at fixed and shared depth]
\label{cor:reference-spin-uv}
On a fixed spatial torus with its harmonic gauge mode fixed,
the stationary neutral-reference electromagnetic coefficients of
\eqref{eq:reference-spin-uv} admit jointly absolutely convergent
bare momentum representations at every finite electric order,
for every finite real $v$. This includes arbitrary fixed external
charged endpoints as specified in the preceding estimates.
For $0<r\le1$ each fixed graph bound costs at most an inverse
power of $r$. A diagonal preparation with one shared geometric
depth and all inverse moments therefore integrates these bounds.
\end{corollary}
\begin{proof}
The matrix hopping is off diagonal for the onsite splitting by
$\beta$. Its kernel is a first derivative of the Gaussian and
remains Hermitian after Wilson transport because the reversed
chord has the conjugate phase. Theorem~\ref{thm:ct-uv-gauged-graphs} applies with $t_i(k)=v k_i e^{-\kappa r|k|^2}$. To track the
depth constants, each spatial mark differentiates the hopping
symbol or translates its argument under a bounded chord parameter
integral. At every fixed order the resulting Gaussian polynomials
have translated lattice $L^1$ and supremum norms bounded by powers
of $r^{-1}$. Here lattice $L^1$ denotes the momentum sum with its
usual box normalization; comparison with Gaussian shells is
uniform under translations. The gap $E_r\ge m$ controls the carrier
resolvents. The supremum of $\|h_r\|$ grows at most as
$m+C|v|(\kappa r)^{-1/2}$. The harmonic envelopes on the photon
lines give further inverse powers by spatial shell counting with
the zero mode removed. Only finitely many of these costs occur in
a fixed electric graph. Integrating their product uses one inverse
moment of the shared $r$ and needs no independence of the visits.
\end{proof}

A separate estimate controls the actual vacuum amplitude when the
coefficient is sewn with a neutral finite-time preparation. Let
$K_{r,A}$ denote the off-diagonal spinor hopping on the torus in
the transverse Maxwell field. The modulus of each Wilson phase
is one. In three dimensions the derivative of the heat kernel is
bounded by $C(\kappa r)^{-1/2}$ times a heat kernel at twice the
lapse, by its explicit Gaussian formula. Periodization and
Parseval thus give
\[
 \|K_{r,A}\|_{\rm HS}^2
      \le C_{V,v,\kappa}(1+r^{-5/2}) .
\]
On particle--hole Fock space its pair term is bounded in form by
$mN/2+C_m\|K_{r,A}\|_{\rm HS}^2$, by the fermion estimate of
Proposition~\ref{prop:cl-gauss-charge-uv}. The positive Coulomb
square includes the one-charge self-energy. Keeping that square
and the photon energy gives $H(r)\ge-C(1+r^{-5/2})$. At fixed
$r$ the same form argument defines a self-adjoint Hamiltonian.
For a normalized neutral preparation its exact slab vacuum
amplitude of duration $T$ is consequently bounded by
$\exp[CT(1+r^{-5/2})]$ at any real electric charge. A boundary
law with finite $\int e^{c/r^3}\,\mathrm d\mu(r)$ for some $c>0$
absorbs this cost times every fixed inverse power at finite $T$.
The exponential is derived here from the Hamiltonian of the same
charged spinor. A general normalization of a candidate boundary
density would not provide this estimate. This slab bound concerns the unexpanded vacuum amplitude. The
stationary coefficient estimates use fixed $r$ throughout each graph. Promoting $r$
to a local field that changes along a path also changes the
operator and its determinant.

\paragraph{Unexpanded susceptibility of the positive two-band transfer.}
The same charged paths give a real-coupling bound for the photon.
Keep $m,\sigma,a>0$, $2|g|<m$, and a spatial torus of period
$2\pi R$, $R\ge1$. Prepare the neutral onsite fermion vacuum and
the free photon vacuum at the ends of a symmetric Euclidean slab
$[-L,L]$. Switch the hopping on in that slab. Remove the spatial
photon and Coulomb zero modes consistently with the zero-mean
Gauss inverse. Write $C_T=(-\partial_\tau^2-\Delta_x)^{-1}$ for
the transverse free photon covariance on the full time line.
All probes below are physical transverse tests; the scalar Gaussian
variable used in the proof represents the Coulomb energy.

\begin{theorem}[Two-sided susceptibility with a pair-producing positive transfer]
\label{thm:closing-uv-positive-susceptibility}
There exists $c=c(m,g,\sigma,a)$, independent of $R,L$, such that
the exact connected photon covariance obeys, as a quadratic form,
\[
 \frac{1}{1+cq^2}\,C_T
 \ \preceq\ \operatorname{Cov}_{q,g;R,L}(A_T)
 \ \preceq\ \frac{1}{1-cq^2}\,C_T \qquad(cq^2<1).
\]
The lower bound holds for every finite real $q$. In any stationary
positive limit obtained from these neutral preparations, the
inequalities persist for limiting smeared covariances.
\end{theorem}
\begin{proof}
First take smooth real backgrounds for both the transverse field and
the auxiliary temporal potential. The normalized matter vacuum in
these backgrounds is the exponential of the convergent closed-path
logarithm in Proposition~\ref{prop:closing-uv-ordered-Maxwell}.
This is a positive real number. Indeed each nonzero word has even
length, with equal numbers of forward upper-slot propagators and
backward lower-slot propagators. Reverse the polygon, interchange
the two slot labels and use the symmetry of the spatial heat density.
The ordered onsite signs and trace coefficient are unchanged, all
segment phases are conjugated, and the restrictions to the slab
are invariant. Winding $w$ in space is paired with $-w$.
The onsite vacuum is neutral, so there is no extra scalar background
weight at zero hops. At fixed $R,L$ the absolute log bound is at most
$C(2L)(2\pi R)^3$, independently of either background.

Vary the two readouts along real test potentials $b_T,b_0$ on the
full time line, initially smooth. In a closed word let $(\tau,x)$
be its basepoint and $\ell$ its total length. For a horizontal mark,
subtract $b_T(\tau,x)$ from each segment. Translation in the basepoint
variables gives
$\|b_T(\cdot+d)-b_T\|_2\le |d|\|\nabla_{\tau,x}b_T\|_2$.
The subtracted integrals cost at most
$\ell^2\|\nabla_{\tau,x}b_T\|_2$ in this norm. The constant on a
winding polygon contributes $w\cdot b_T(\tau,x)$. If $w\ne0$,
$2\pi R\le |w|\le\ell$ and
$\|b_T\|_2\le R\|\nabla_x b_T\|_2$ by zero spatial mean at each
time. Thus this term obeys the same length bound up to a constant.
For the vertical marks subtract $b_0(s,x)$ at each actual segment
time $s$. Its signed time integral around every closed path is
zero, since there is no time winding with vacuum endpoints.
Only spatial translations remain, with norm cost
$\ell^2\|\nabla_x b_0\|_2$. No time derivative of the scalar potential
is needed. These estimates also bound the basepoint integral
restricted to the slab, by extending the positive majorant to all
times.

The two derivatives of each Abelian word bring the product of two
complete currents. Cauchy--Schwarz in its basepoint bounds them by
$Cq^2\ell^4(\|\nabla_{\tau,x}b_T\|_2^2+\|\nabla_x b_0\|_2^2)$.
The length-tilted time and spatial integrals in the preceding proof
still converge; one closing tilted spatial Gaussian has uniformly
bounded periodization for $R\ge1$. Thus this is a uniform Hessian
bound for the charged log weight after the length moment is summed.
Each spatial readout is a contraction in the displayed norms. In the
free joint metric $Q=\operatorname{diag}(C_T^{-1},-\Delta_x)$,
the absolute Hessian of the log weight is consequently at most $cq^2 Q$.

Finite Gaussian projections now have a real negative log density $S$
with Hessian between $(1-cq^2)Q$ and $(1+cq^2)Q$. In coordinates
where the free covariance is identity, integration by parts gives
$\operatorname{Cov}(u\cdot X,u\cdot\nabla S)=|u|^2$ and
$\mathbb E[(u\cdot\nabla S)^2]\le(1+cq^2)|u|^2$.
Cauchy--Schwarz proves the lower bound. Brascamp--Lieb variance
comparison gives the upper bound in the strictly convex regime.
Restriction to tests in the transverse block retains both bounds.

To remove the projections, integrate a mode approximation against each
path current first. The smeared temporal source of one vertical
segment is an interval indicator times its spatial profile, with
finite norm in the instantaneous Coulomb covariance. Its integral
is a Gaussian increment; the horizontal smeared field at each fixed
time also has finite free variance. Projections converging in these
Gaussian norms give convergence in mean square of the corresponding
line integrals for each finite polygon. Their phases have modulus
one. Dominated convergence in the convergent path series and the
fixed-volume log bound yield convergence of the full weight and of
its tested Maxwell moments. Equivalently, conditioning on the
occupations between successive hops, integration of the scalar
Gaussian produces precisely $\exp[-\frac12\int \rho C_{\rm C}\rho]$
with both self and cross terms; $C_{\rm C}=(-\Delta_x)^{-1}$ on
each spatial slice and $\rho$ includes the charge and the common
spatial profile. This is the diagonal Coulomb evolution in the
canonical time ordering. The transverse Gaussian is the vacuum
oscillator evolution between these same hops. Thus the positive
Hamiltonian preparation supplies the magnetic correlator appearing
in the bound. In particular its slab magnetic observables satisfy
time-reflection positivity: reflecting a product of positive-time
insertions in the symmetric preparation gives the squared norm of
its half-slab evolved vector, divided by the vacuum amplitude.
\end{proof}

This result bounds the full susceptibility using a model with
charged pairs and an instantaneous Coulomb self-energy. Spatial
transfer controls the ultraviolet cost while the fermions propagate
in carrier time. The estimate assumes no carrier-time heat factor on
the measured photon. Uniform comparison of the Euclidean response
allows a gapless spectral threshold in the bulk and does not specify
a sharp Lorentz-invariant photon pole or a unique bulk state.

There is a quantitative particle-spectrum consequence in a stationary
positive vacuum limit. Suppose a nonzero spatial transverse mode
of length $k$ has a positive squared-energy measure of total mass
one in $G_k(x)=\int(x+s)^{-1}\rho_k(\mathrm ds)$, $x=\omega^2$.
At finite volume for a stationary ground vector the total mass follows
from $\langle[A,[H,A]]\rangle=1$ for a real unit-normalized photon
mode: the double commutator is the free canonical one, since the
hopping and Coulomb terms commute with $A$. This calculation uses
the ground-state form domain of the mode and its energy moment, or
the corresponding sum rule obtained by smearing the mode. For a bulk
limit the normalization can be tested there directly. Push $\rho_k$
to the dimensionless variable $u=s/k^2$ and put $\delta=cq^2<1$.
The positive function
\[
 \frac{(u-1)^2}{(u+1)(u+2)}=1+\frac4{u+1}-\frac9{u+2}
\]
vanishes at the Maxwell energy. Evaluate the susceptibility bounds
at $x=k^2,2k^2$. They give the following bound, uniform in $k$,
for each $\eta>0$:
\begin{equation}
 \rho_k\!\left(\left\{s:\left|s/k^2-1\right|\ge\eta\right\}\right)
 \le \frac{(2+\eta)(3+\eta)}{\eta^2}
        \frac{5\delta-\delta^2}{1-\delta^2}.
 \label{eq:closing-uv-photon-band}
\end{equation}
Indeed the integral of the displayed function is at most
$1+2/(1-\delta)-3/(1+\delta)$. Off the indicated central interval
the function is at least $\eta^2/[(2+\eta)(3+\eta)]$, by monotonicity
on either side of one. Thus at sufficiently weak charge a positive
fraction of the measured spectral weight lies in any prescribed
relative band about $E=k$, uniformly down to soft $k$. This excludes
a uniform positive photon gap when such modes approach zero.

If the canonical normalization has not been passed through a bulk
limit, the susceptibility bound still yields
$(1+\delta)^{-1}\le\int\rho_k\le(1-\delta)^{-1}$ by monotone
convergence at large $x$, assuming the positive representation.
The same absolute-weight estimate then uses
$6\delta/(1-\delta^2)$ in place of the last factor in
\eqref{eq:closing-uv-photon-band}. Neither estimate requires an atom
in the indicated interval.

\begin{corollary}[Existence of stationary positive photon limits in the two-band model]
\label{cor:closing-uv-pair-bulk}
In the weak-charge regime of the preceding theorem, the neutral vacuum
slabs have subsequential stationary photon limits after first sending
$L\to\infty$ at fixed $R$ and then $R\to\infty$. Their photon
covariances satisfy the same Maxwell comparisons on transverse
test fields. The magnetic observables satisfy time-reflection
positivity. At almost every nonzero spatial momentum, in each real
transverse polarization, these limits have a positive squared-energy
measure and obey the unnormalized version of
\eqref{eq:closing-uv-photon-band}.
\end{corollary}
\begin{proof}
The joint slab weight is invariant under simultaneous reversal of the
two gauge fields, by the conjugation and reality of its matter
factor. Hence linear fields have mean zero. A real linear tilt leaves
the Hessian bound unchanged. Integrating the resulting variance bound
twice in the tilt strength gives, for every real transverse smearing
$f$,
\[
 \mathbb E e^{tA_T(f)}
 \le \exp\!\left(\frac{t^2}{2(1-\delta)}
                         \langle f,C_T f\rangle\right).
\]
On a fixed spatial box this is uniform in $L$. For bulk Schwartz
tests periodized in space the free quadratic form is bounded
uniformly in $R\ge1$ by a continuous Schwartz seminorm squared:
integrating the time frequency of the Maxwell majorant costs at
most $\pi/|k|$ times a rapidly decreasing spatial bound. The
normalized spatial sums of $|k|^{-1}$ near the origin are uniformly
integrable in three dimensions, with the origin omitted. On the
complement ordinary rapid decay controls the sums and gives their
continuum limits. These bounds give tight finite-dimensional laws
and uniform integrability of all fixed moments. Extract on a
countable dense set of tests at each of the two limits. The uniform
Schwartz bound makes the limiting characteristic functional continuous
at zero; the characteristic functional theorem on Schwartz space~\cite{BulkLevyContinuity}
then gives a probability on transverse tempered fields. This
construction retains the covariance comparisons and spatial
translations. Spatial inversion also holds since the hopping and
readout are even, and the vector potential transforms as a vector.

For time stationarity fix $R$ during the first limit. Write $v$ for
the neutral preparation, $E_v$ for the infimum of the spectral
support of $v$ under the lower-bounded Hamiltonian, and
$v_t=e^{-tH}v/\|e^{-tH}v\|$. Its exponentially tilted energy law
concentrates at $E_v$. Consequently for fixed $h$ one has
$\|e^{-(t+h)H}v\|/\|e^{-tH}v\|\to e^{-hE_v}$ and
$\|v_{t+h}-v_t\|\to0$. These limits follow by the spectral theorem;
for positive $h$ the exponential test functions are bounded on the
support, and negative shifts follow by using $t+h$ as base time.
For bounded time-zero Weyl insertions at
$s_1<\cdots<s_n$ the normalized slab expectation is the matrix
element between $v_{L-s_n}$ and $v_{L+s_1}$ of the bounded ordered
product with intervening transfers. Its scalar factor is
$\|e^{-(L-s_n)H}v\|\,\|e^{-(L+s_1)H}v\|/\|e^{-LH}v\|^2$.
It tends to $e^{(s_n-s_1)E_v}$. Simultaneously shifting the
$s_j$ by a fixed amount changes neither the limiting factor nor
the limiting matrix elements. The argument does not require a
ground eigenvector at $E_v$.

Such equal-time tests and their finite-time smearings determine the
photon law. Indeed the upper Maxwell bound controls the squared
increment of a spatially smooth mode between two times by a constant
times their separation: in the free covariance use
$(1-e^{-|k||t-s|})/|k|\le |t-s|$. This justifies mollifying time
evaluations or taking Riemann sums in the time smearings, uniformly
in the slab width. Stationarity and the reflected squared-norm
inequalities for the bounded Weyl observables therefore persist in
the subsequences. Differentiation at zero source is justified by the
uniform moment bounds.

Finally the invariant covariance is a spectral multiplier in
four-momentum. In any real transverse polarization its Fourier
measure is absolutely continuous and bounded above and below by
the stated multiples of the Maxwell density, by testing these form
inequalities on localized functions. At almost every nonzero $k$ the
time-frequency density is integrable. Its time transform $F_k(t)$
is bounded, continuous, and even for paired spatial modes $k,-k$
by reality and spatial inversion. Time-reflection positivity gives
positive matrices $[F_k(t_i+t_j)]$ for $t_i>0$: localize the
spatial test functions around the paired momenta, first use rational
times and rational coefficients, then use continuity. The Laplace
representation theorem for exponentially convex functions gives
$F_k(t)=\int e^{-Et}\mu_k(\mathrm dE)$ for $t>0$ with
$\mu_k\ge0$~\cite{Widder1941}. Boundedness at large time excludes
$E<0$, and continuity at zero gives finite total mass. There is no
constant-time component: it would put an atom at zero frequency in
a Fourier measure dominated by the Maxwell density at this $k$.
Taking the time transform identifies $\rho_k$ with the pushforward
of $2E\,\mu_k$ by $E\mapsto E^2$. The covariance bounds hold for
this Stieltjes transform at every positive squared frequency by
continuity. The band estimate follows as above, without passing a
canonical commutator through either limit.
\end{proof}

For soft $k$ these positive states exclude a uniform spectral gap
at sufficiently weak charge. The result establishes existence of
subsequential states. Uniqueness, an isolated photon pole, and a
relativistic chiral completion would require further dynamics.

\subsection{A massless photon and the spectral threshold of a moving charge}
\label{subsec:uv-rigid-charge-threshold}

A photon can have a sharp light cone even when a charged probe has
zero overlap with any state at the bottom of its spectrum. Both
spectral measures follow from the Hamiltonian of a rigid current coupled linearly
to the transverse photon field, a source of van Hove type~\cite{SpectrumDerezinski2003}. Fix a reconstructed sector plane
$\bb=(s,b_s)$, use its single physical time and three spatial
coordinates, and use units with $\hbar=1$ and unit limiting speed
on this carrier. The ambient geometric-velocity vector retains its
selected direction. Prescribe a charge $e$ moving at
constant velocity $\mathbf v$, with $u=|\mathbf v|<1$, coupled to
the Maxwell action through its conserved current. Define its spatial
profile by harmonic extension through a half-space of depth $a>0$.
Minimizing the Dirichlet energy at fixed boundary value gives, at
spatial frequency $\mathbf k$, the decaying solution
$e^{-a|\mathbf k|}$ times that value. The associated unit-mass
boundary density and the current are
\begin{equation}
\begin{aligned}
 p_a(\mathbf x)&=\frac{a}{\pi^2(a^2+|\mathbf x|^2)^2},
 &\widehat p_a(\mathbf k)&=e^{-a|\mathbf k|},\\
 (j^0,\mathbf j)(t,\mathbf x)&=
 e(1,\mathbf v)p_a(\mathbf x-\mathbf vt).
\end{aligned}
\label{eq:uv-rigid-geometric-current}
\end{equation}
Here the Fourier transform has no prefactor in the forward direction.
The current satisfies $\partial_t j^0+\nabla\!\cdot\!\mathbf j=0$.
The depth is part of the source definition. Its spatial Fourier
multiplier occurs at the current vertex, while every photon momentum
in $\mathbb R^3$ remains in the Hilbert space.

Use Coulomb gauge. Write $r=|\mathbf k|$, $\mathbf n=\mathbf k/r$,
$d=1-\mathbf v\cdot\mathbf n$, and
$\omega=rd$ for the photon frequency in the generator of time
translations in the moving coordinates. In the absence of the source
this generator is $K_0=H_0-\mathbf v\cdot\mathbf P=
\mathrm d\Gamma(\omega)$ on transverse photon Fock space. The
one-photon inner product uses $\sum_{\lambda=1}^{2}\int
\mathrm d^3\mathbf k$, with unit transverse polarizations
$\boldsymbol\epsilon_\lambda(\mathbf n)$. In the charged sector
Gauss' law fixes the longitudinal electric field. Including its
energy gives
\begin{equation}
\begin{aligned}
 K_1={}& E_{\rm c}+E_{\rm C}+\mathrm d\Gamma(\omega)
                      +a^*(g)+a(g),\\
 g_\lambda(\mathbf k)={}&
 -\frac{e\,\mathbf v\cdot\boldsymbol\epsilon_\lambda(\mathbf n)}
             {(2\pi)^{3/2}\sqrt{2r}}\,e^{-ar},\\
 E_{\rm C}={}&\frac{e^2}{2(2\pi)^3}
  \int_{\mathbb R^3}\frac{e^{-2ar}}{r^2}\,\mathrm d^3\mathbf k
  =\frac{e^2}{8\pi^2 a}.
\end{aligned}
\label{eq:uv-rigid-sector-H}
\end{equation}
$E_{\rm c}>0$ is the prescribed core energy in this generator.
The coupling is $-\int\mathbf j\cdot\mathbf A_T$, with the
source stationary in the moving coordinates. Matter recoil and
virtual pair creation are absent from this specified action.

\begin{proposition}[The charge threshold and its spectral weight]
\label{prop:uv-rigid-spectral-atom}
The bottom $E_*$ of the charged spectrum is strictly positive for
every real $e$. In the neutral vacuum the transverse magnetic
two-point function has a positive one-photon atom with laboratory
energy $r$ at each fixed $\mathbf k\ne0$. In the charged sector,
take the probe that prepares the prescribed Gauss field and the
transverse Fock vacuum $|0\rangle$. For $eu\ne0$ its energy law is
absolutely continuous on $[E_*,\infty)$. With $x=E-E_*\downarrow0$
its density is asymptotic to $C x^{\alpha-1}/\Gamma(\alpha)$,
where
\begin{equation}
\begin{aligned}
 E_*&= E_{\rm c}+\frac{e^2}{8\pi^2a}
       (1-u^2)\frac{\operatorname{arctanh}u}{u},\\
 w(\mathbf n)&=\frac{e^2}{2(2\pi)^3}
       \frac{u^2-(\mathbf v\cdot\mathbf n)^2}{d^2},\\
 \alpha&=\int_{S^2} w\,\mathrm d\Omega
  =\frac{e^2}{2\pi^2}
       \left(\frac{\operatorname{arctanh}u}{u}-1\right),\\
 C&=(2a)^\alpha\exp\!\left[-\int_{S^2} w\log d\,\mathrm d\Omega\right].
\end{aligned}
\label{eq:uv-rigid-edge-constants}
\end{equation}
At $u=0$ use $\operatorname{arctanh}(u)/u=1$.
If $eu=0$, this probe gives a unit atom at $E_*$. In either case
neutral photons of arbitrarily small positive energy are available.
\end{proposition}

\begin{proof}
The transverse magnetic covariance at $\mathbf k\ne0$ and
positive Euclidean time $\tau$ is
$r(\delta_{ij}-n_i n_j)e^{-\tau rd}/2$, with the spatial
momentum-conservation delta suppressed. This gives its positive
matrix weight at moving energy $rd$ and hence at laboratory energy
$r$. The static source's transverse state is also free when $u=0$.

For the charged block, $g$ and $g/\sqrt\omega$ are square
integrable, so the linear term defines an infinitesimally
form-bounded perturbation of $\mathrm d\Gamma(\omega)$.
Completing the square gives the lower bound
$E_{\rm c}+E_{\rm C}-\|g/\sqrt\omega\|^2$. Coherent vectors
with displacement $-g/\omega$ restricted to $r>\epsilon$
approach this bound as $\epsilon\downarrow0$. This variational
sequence is used only to identify the infimum of the unmodified
Hamiltonian. After angular symmetrization under
$\mathbf n\mapsto-\mathbf n$, the integrand of
$E_{\rm C}-\|g/\sqrt\omega\|^2$ relative to that of $E_{\rm C}$
is $(1-u^2)/(1-(\mathbf v\cdot\mathbf n)^2)$. Integrating gives
\eqref{eq:uv-rigid-edge-constants}. The finite shift is the
energy of the transverse displacement in the same Hamiltonian.

For the spectral law, square completion for finitely truncated
displacements and the vacuum coherent-state identity give
\begin{equation}
\begin{aligned}
 L(\tau)&=\langle0|e^{-\tau(K_1-E_*)}|0\rangle\\
 &=\exp\!\left[-\int_{S^2} w(\mathbf n)
    \log\!\left(1+\frac{\tau d}{2a}\right)\mathrm d\Omega\right].
\end{aligned}
\label{eq:uv-rigid-exact-laplace}
\end{equation}
Indeed the exponent before radial integration is
$-\int w\int_0^\infty (1-e^{-\tau rd})e^{-2ar}\,\mathrm dr/r\,
\mathrm d\Omega$. The integral is finite at both endpoints.
Convergence of the truncated couplings in the one-photon norm and
in the $\omega^{-1/2}$ norm gives convergence of their semigroups
by quadratic forms. Thus the identity holds on the full space.
The angular integration of $w$ gives $\alpha$ as stated.
When $eu\ne0$, the divergent soft norm
$\|g/\omega\|^2=\int w\int_0^\infty e^{-2ar}\,\mathrm dr/r\,
\mathrm d\Omega$ explains why the charged atom disappears.
The vanishing of $L(\tau)$ at large $\tau$ measures its weight
directly. At each finite $\tau$ the exponential in
\eqref{eq:uv-rigid-exact-laplace} has a convergent power series in
$e^2$ for every coupling. For $eu\ne0$ its coefficients at positive
orders grow as powers of $\log\tau$ at large time.

For the density, let $y>0$ denote emitted energy, write
$b(\mathbf n)=2a/d$ and $b_{\max}=2a/(1-u)$, and decompose the
positive measure in the Laplace exponent as
\[
 \frac{\mathrm dy}{y}\int_{S^2}w e^{-b(\mathbf n)y}\,\mathrm d\Omega
   =\alpha e^{-b_{\max}y}\frac{\mathrm dy}{y}
       +h(y)\,\mathrm dy .
\]
Here $h\ge0$ is bounded at zero and has finite integral
$M=\int w\log(b_{\max}/b)\,\mathrm d\Omega$.
The first term generates a gamma law with shape $\alpha$ and rate
$b_{\max}$. The second generates the positive compound Poisson law
$e^{-M}\delta_0+e^{-M}\sum_{j\ge1}(h\,\mathrm dy)^{*j}/j!$,
whose density part is bounded near zero by convolution of the
bounded density $h$ with finite measures. Their convolution gives
the asserted absolutely continuous law. Its density is strictly
positive for $x>0$ from the gamma contribution. The density part
of the second law contributes $O(x^\alpha)$ at the edge, and
$e^{-M}b_{\max}^{\alpha}=C$. For $eu=0$ the logarithmic exponent
vanishes identically.
\end{proof}

\paragraph{Which charged probe could have a pole?}
For a transverse coherent probe of displacement $f\in L^2$,
the threshold atom has weight
$\exp[-\|f+g/\omega\|^2]$, interpreted as zero for an infinite
norm. This follows from the same coherent-state formula with
$|f+g/\omega|^2(1-e^{-\tau\omega})$ in the exponent and monotone
convergence as $\tau\to\infty$. When $eu\ne0$, all such overlaps
vanish. Coherent vectors span a dense subspace, so the charged Fock
Hamiltonian has no eigenvector at $E_*$. A non-Fock infrared
dressing changes the representation and requires its own state and
probe. If an actual current/dressing pair instead has a residual
displacement of order $r^{-3/2+\delta}$ at zero, bounded in angular
$L^2$ with $\delta>0$, its soft norm is finite. A nonzero leading
$r^{-3/2}$ angular coefficient gives logarithmic divergence.
The half-space profile tends to one as $r\downarrow0$ and leaves
this distinction intact. To infer a mass shell for dynamical matter
one must derive its physical charged probe, its soft current and its
recoil spectrum from the same action. Equations
\eqref{eq:uv-rigid-sector-H}--\eqref{eq:uv-rigid-exact-laplace}
are fixed-velocity Hilbert-space spectral statements; the vectorial
boundary prescription is applied separately to a reconstructed
resolvent on the chosen sector plane.

\subsection{Conditional crossing and locality}

Local holomorphy in Theorem~\ref{thm:holomorphy} is a statement
about a graph representation. Physical crossing relates boundary
values in different scattering channels. To invoke a standard
crossing result one needs the applicable physical hypotheses:
locality of reconstructed observables, spectral support,
appropriate stable asymptotic particles, LSZ limits, and a
connected analytic domain joining the channels with controlled
boundary growth. Kernel and vertex singularities must be
included in that domain analysis.

For two-particle scattering with the required stable particles and a
strictly positive physical mass gap, Bros, Epstein, and Glaser construct
a connected on-shell continuation domain \cite{BrosEpsteinGlaser1965}.
Their gap concerns all nonvacuum physical states, not just the external
particles \cite[March preprint, p.~6]{BrosEpsteinGlaser1965}. The same paper explicitly
separates the additional identification of on-shell boundary values
with the physical amplitudes and postpones its proof
\cite[March preprint, p.~2]{BrosEpsteinGlaser1965}. Here that identification is part
of the separately assumed scattering and boundary-limit hypotheses;
it is not inferred from existence of the analytic domain alone.
A massless physical channel lies outside this invocation. Neither the
cited domain result nor formal permutation symmetry of a Euclidean
integrand proves crossing for an arbitrary damped completion.

The implication justified here is therefore
\[
 \begin{gathered}
 \text{completed physical reconstruction and LSZ hypotheses}\\
 \text{an admissible channel domain}\\
 \text{and identification of its physical boundary values}\\
 \Longrightarrow\quad
 \text{the corresponding physical crossing relation}.
 \end{gathered}
\]
This statement is separate from the fixed-graph UV theorem and
does not assert a massless scattering construction.

\section{Kernel families, coordinate forms, and worked examples}
\label{sec:examples}

\subsection{Admissible Euclidean kernel families}
\label{sec:kernel-families}

The families below use the dimensionless variable \(z=a^2q_E^2\).
Their suitability for Euclidean UV estimates does not automatically
supply the independent continuation or physical hypotheses.

\paragraph{A single heat depth}
For \(\tau_*>0\),
\begin{equation}
 \Ph(z)=e^{-\tau_*z},\qquad\rho_{\bb}=\delta^{\rm heat}_{\bb,\tau_*}.
 \label{eq:heat-family}
\end{equation}
This satisfies the gap route, and all inverse moments are finite.
It is entire, but its timelike continuation has the growth in
\eqref{eq:heat-continuation}.

\paragraph{Power kernels}
For dimensionless \(b>0\) and \(\nu>0\),
\begin{equation}
 \Ph(z)=(1+z/b)^{-\nu},\qquad
 \rho_{\bb}(\dd\tau)=\frac{b^\nu}{\Gamma(\nu)}
                 \tau^{\nu-1}e^{-b\tau}\dd\tau .
 \label{eq:power-family}
\end{equation}
For \(0<\beta<\nu\),
\begin{equation}
 M_\beta=b^\beta\frac{\Gamma(\nu-\beta)}{\Gamma(\nu)},
 \label{eq:gamma-inverse-moments}
\end{equation}
and the inverse moment diverges for \(\beta\geq\nu\).
Thus \(\nu>1+D/2\) suffices for the global theorem.
Weaker graph-specific conditions may suffice, as the bubble
example shows. The singularity at \(z=-b\) must be retained in a
Lorentzian analysis.

\paragraph{Finite positive exponential mixtures}
For \(w_j>0\), \(\sum_jw_j=1\), and \(\tau_j>0\),
\begin{equation}
 \Ph(z)=\sum_{j=1}^Nw_je^{-\tau_jz},\qquad
 \rho_{\bb}=\sum_{j=1}^Nw_j\delta^{\rm heat}_{\bb,\tau_j}.
 \label{eq:finite-mixture}
\end{equation}
The gap is \(\min_j\tau_j\). These are exponential mixtures,
not rational functions or Pad\'e approximants.

\paragraph{Stretched exponentials}
For \(0<r<1\),
\begin{equation}
 \Ph(z)=e^{-z^r}
 \label{eq:stable-family}
\end{equation}
has a positive stable representing measure. It has no gap, but
every inverse moment is finite:
\begin{align}
 M_\beta
 &=\frac1{\Gamma(\beta)}
     \int_0^\infty z^{\beta-1}\Ph(z)\dd z
   =\frac{\Gamma(\beta/r)}{r\Gamma(\beta)},
 \qquad \beta>0.
 \label{eq:stable-inverse-moments}
\end{align}
The first equality follows from
\(\tau^{-\beta}
 =\Gamma(\beta)^{-1}\int_0^\infty z^{\beta-1}e^{-\tau z}\dd z\)
and Tonelli. The restriction to orders below \(r\) concerns
positive moments of the stable variable, not inverse moments.
At \(r=1\), the measure reduces to a point mass.

For the monograph's canonical \(r=\tfrac12\) specialization,
\begin{equation}
 \rho_{\bb}(\dd\tau)=\frac1{2\sqrt\pi}\,
     \tau^{-3/2}e^{-1/(4\tau)}\dd\tau,\qquad
 M_\beta=\frac{4^\beta\Gamma(\beta+\tfrac12)}{\sqrt\pi}.
 \label{eq:stable-half-density}
\end{equation}
The line multiplier is \(e^{-a|q_E|}\), with
superpolynomial large-momentum decay. The large-\(\tau\) tail
of \eqref{eq:stable-half-density} excludes any positive
exponential moment. Its UV usefulness and its more delicate
continuation therefore coexist.
The O--QED application uses the separately specified normalized massive
variant in Section~\ref{sec:OQED-moments}; its positive auxiliary scale
changes the large-parameter tail and local analytic hypotheses without
creating a positive Laplace support gap.

An ungapped measure does not acquire a gap by mixing it with a
positive atom away from zero: its original near-zero support
remains. Translating the entire measure by \(\tau_*>0\), or
truncating and renormalizing it below that point, does create a
gap. Translation corresponds to
\(\Ph(z)\mapsto e^{-\tau_*z}\Ph(z)\).

The distinction between inverse moments and exponential moments is
summarized in Table~\ref{tab:kernel-comparison}. Write
\[
 \mathcal E_{\rho,\bb}=\left\{\epsilon>0:
      \int e^{\epsilon\tau}\rho_{\bb}(\dd\tau)<\infty\right\}.
\]
Inverse moments enter the UV estimate; exponential moments give a
sufficient local continuation condition. The continuation of an
amplitude also uses the graph-level hypotheses of
Section~\ref{sec:analytic}.

\begin{table*}[t]
\centering
\small
\setlength{\tabcolsep}{3pt}
\renewcommand{\arraystretch}{1.2}
\caption{Representative kernel properties. Here
\(\tau_*,b,\nu>0\), \(\beta>0\), \(0<r<1\), and \(\mu>0\)
in the normalized massive row.}
\label{tab:kernel-comparison}
\begin{tabularx}{\textwidth}{@{}>{\raggedright\arraybackslash}p{0.20\textwidth}
 >{\raggedright\arraybackslash}p{0.28\textwidth}
 >{\raggedright\arraybackslash}p{0.14\textwidth}
 >{\raggedright\arraybackslash}X@{}}
\toprule
Kernel \(\Ph(z)\) & Inverse moment \(M_\beta\) &
\(\mathcal E_{\rho,\bb}\) & Complex continuation\\
\midrule
\(e^{-\tau_*z}\) & \(\tau_*^{-\beta}\), all \(\beta>0\) &
\((0,\infty)\) & Entire\\
\((1+z/b)^{-\nu}\) &
\(\dfrac{b^\beta\Gamma(\nu-\beta)}{\Gamma(\nu)}\),
\(0<\beta<\nu\) & \((0,b)\) &
Pole or branch point at \(z=-b\)\\
\(e^{-z^r}\) &
\(\dfrac{\Gamma(\beta/r)}{r\Gamma(\beta)}\), all \(\beta>0\) &
\(\varnothing\) & Branch point at \(z=0\)\\
\(e^{\mu-\sqrt{z+\mu^2}}\) &
\eqref{eq:OQED-inverse-moments}, all \(\beta>0\) &
\((0,\mu^2]\) & Branch point at \(z=-\mu^2\); no support gap\\
\((2e^z-1)^{-1}\) &
\(\sum_{n\geq1}2^{-n}n^{-\beta}\), all \(\beta>0\) &
\((0,\log2)\) & Poles at \(-\log2+2\pi\ii n\),
\(n\in\mathbb Z\)\\
\bottomrule
\end{tabularx}
\end{table*}

\subsection{Coordinate-space covariance and coincidence limit}

For a finite positive line measure and \(m\geq0\), the bounded
multiplier defines the following Fourier transform initially as a
tempered distribution. Under the coincidence hypotheses stated below,
it is an ordinary continuous covariance function:
\begin{align}
 C_{\Ph,E}(x-y)
 &=\int\rho_{\bb}(\dd\tau)
     \int\frac{\dd^4q_{E,\bb}}{(2\pi)^4}\,
       \frac{e^{\ii q_E\cdot(x-y)}
                 e^{-a^2\tau q_E^2}}{q_E^2+m^2}
 \label{eq:coordinate-covariance}\\
 &=\int\rho_{\bb}(\dd\tau)\int_0^\infty\dd\alpha\,
     \frac{e^{-\alpha m^2}}
          {[4\pi(\alpha+a^2\tau)]^2}\nonumber\\
 &\qquad\times\exp\!\left[-\frac{|x-y|^2}
                         {4(\alpha+a^2\tau)}\right].
 \label{eq:coordinate-heat-form}
\end{align}
The second equality uses the elementary Schwinger representation
\cite{Schwinger1951}. Each positive heat parameter spreads the field before the covariance is
sampled. The formula averages those heat-smoothed Euclidean covariances;
it does not average over different sector-plane identities. At coincidence, for \(\tau>0\),
\[
 \int_0^\infty
 \frac{e^{-\alpha m^2}}{(\alpha+a^2\tau)^2}\dd\alpha
       \leq\frac1{a^2\tau}.
\]
Therefore
\begin{equation}
 C_{\Ph,E}(0)\leq
       \frac{M_1[\rho_{\bb}]}{16\pi^2a^2},
 \label{eq:coincidence-bound}
\end{equation}
provided there is no atom at zero and \(M_1<\infty\).
For \(m=0\), equality holds. In the stable-half case,
\(M_1=2\), giving
\(C_{\Ph,E}(0)=1/(8\pi^2a^2)\) for the massless scalar
comparison covariance. This finite coincidence limit is a
Euclidean statement, not a reflection-positivity proof.

\subsection{The literal two-line scalar bubble}
\label{sec:literal-bubble}

Fix one support and \(m>0\). The following calculation is a
Euclidean coefficient benchmark with the explicitly fixed virtual
carrier convention. Its displayed bubble function omits the
propagator numerators, vertices, and overall Dyson coefficient.
These data are restored below. For a running-unit Lorentzian
bubble with momenta \(k\) and \(P-k\), the two line factors
\(\sgn(k_\parallel)\) and
\(\sgn(P_\parallel-k_\parallel)\) must instead remain in the
integral. The scalar radial computation does not evaluate that
different integrand.

With one multiplier on each uncut line, the Euclidean scalar bubble is
\begin{equation}
 B_{\Ph,E}(P_E)=
 \int\frac{\dd^4k_{E,\bb}}{(2\pi)^4}\,
 \frac{\Ph(a^2k_E^2)\Ph(a^2(k_E+P_E)^2)}
      {(k_E^2+m^2)((k_E+P_E)^2+m^2)} .
 \label{eq:two-line-bubble}
\end{equation}
Introduce independent \(\tau_1,\tau_2\), independent
\(\alpha_1,\alpha_2\), and \(w_i=\alpha_i+a^2\tau_i\).
Completing the square gives
\[
\begin{aligned}
 &w_1k_E^2+w_2(k_E+P_E)^2\\
 &\quad=(w_1+w_2)\left(k_E+\frac{w_2}{w_1+w_2}P_E\right)^2
       +\frac{w_1w_2}{w_1+w_2}P_E^2.
\end{aligned}
\]
The exact parameter formula is
\begin{align}
 B_{\Ph,E}(P_E)
 =\frac1{16\pi^2}
 &\int\rho_{\bb}(\dd\tau_1)\rho_{\bb}(\dd\tau_2)
       \int_0^\infty\dd\alpha_1\dd\alpha_2 \nonumber\\
 &\times\frac{
 e^{-m^2(\alpha_1+\alpha_2)}
 e^{-\frac{w_1w_2}{w_1+w_2}P_E^2}
 }{(w_1+w_2)^2}.
 \label{eq:two-line-parameters}
\end{align}
This is the actual linewise rule. A single factor attached after
combining and shifting denominators is a different effective
model unless a separate identity establishes equivalence.

In the normal ascending scalar reference convention
\(\sigma_s=\sigma_\epsilon=+1\), use a real \(\phi^4\)
coupling, \(S=1+\ii T\), and tree coefficient
\(\mathcal M_{\rm tree}=-\lambda\). The one-loop
coefficient in this benchmark is
\begin{equation}
\begin{aligned}
 &\mathcal M^{\rm coeff}_{+,\bb}(\sMand,\tMand,\uMand)\\
 &\quad=\mathcal P^{(4)}_{\bb\mid{\rm ext}}
 \left[-\lambda+\frac{\lambda^2}{2}
    \sum_{c=1}^{3}B^{\rm coeff}_{\Ph,+}(v_c)\right]+O(\lambda^3).
\end{aligned}
 \label{eq:scalar-amplitude}
\end{equation}
Here \(v_1=\sMand\), \(v_2=\tMand\), and \(v_3=\uMand\);
the three channel contributions are counted with their multiplicity.
The function \(B^{\rm coeff}_{\Ph,+}\) is the extracted boundary
of the declared continuation of \eqref{eq:two-line-bubble}, as in
\eqref{eq:analytic-function-directed-interface}. This extraction is
performed on the admitted path for each channel; it is not an
independent complex displacement of three constrained invariants. This is a specified scalar benchmark,
not the full gauge completion. The factor \(1/2\) belongs to
the graph symmetry factor and must not be inserted into the
definition of \(B_\Ph\).

Let \(\mathcal M^{\rm coeff}_{-,\bb}\) denote the complete
opposite-branch reference with its own boundary values. For this
real-coupling scalar benchmark with the reflected endpoint and
vertex conventions,
\(\mathcal M^{\rm coeff}_{-,\bb}
=\overline{\mathcal M^{\rm coeff}_{+,\bb}}\). Hence
\begin{align}
 \vec{\mathcal M}_{\bb}
 &=\mathsf P_+\mathcal M^{\rm coeff}_{+,\bb}
      +\mathsf P_-\mathcal M^{\rm coeff}_{-,\bb}
   =\embed(\mathcal M^{\rm coeff}_{+,\bb}),\nonumber\\
 \vec{\mathcal A}_{\bb}
 &=\mathsf P_+\ii\mathcal M^{\rm coeff}_{+,\bb}
      -\mathsf P_-\ii\mathcal M^{\rm coeff}_{-,\bb}
   =\je\vec{\mathcal M}_{\bb}.
 \label{eq:scalar-benchmark-Dyson-reconstruction}
\end{align}
Here \(\vec{\mathcal A}\) is the Dyson transition coefficient;
its opposite branch includes the opposite Dyson sign. Thus a real
tree value \(-\lambda\) gives
\(\vec{\mathcal M}_{\rm tree}=-\lambda\) and
\(\vec{\mathcal A}_{\rm tree}=-\lambda\je\).
If \(\mathcal M_+=A+\ii B\), then
\(\vec{\mathcal M}=A+\je B\) but
\(\vec{\mathcal A}=-B+\je A\). The two quantities have
different parity decompositions. Outside this explicitly reflected
benchmark, use the two complete representatives in
\eqref{eq:complete-graph-character-reconstruction} directly.

\subsection{Exact Gaussian bubble and its discontinuity}
\label{sec:gaussian-bubble}

Choose \(\Ph(z)=e^{-\tau_*z}\) and put \(t=a^2\tau_*>0\).
Then \(w_i=\alpha_i+t\).
Set \(W=w_1+w_2\) and \(x=w_2/W\). The integration region is
\[
\begin{gathered}
 0<x<1,\qquad W\geq W_{\min}(x):=\frac{t}{\min(x,1-x)},\\
 \dd w_1\dd w_2=W\,\dd W\dd x .
\end{gathered}
\]
Equation~\eqref{eq:two-line-parameters} becomes
\begin{equation}
 B_{t,E}(P_E)
   =\frac{e^{2tm^2}}{16\pi^2}
        \int_0^1\dd x\int_{W_{\min}(x)}^\infty
        \frac{\dd W}{W}
             e^{-W[m^2+x(1-x)P_E^2]} .
 \label{eq:gaussian-W}
\end{equation}
Initially for \(\Re z>0\), set
\(E_1(z)=\int_1^\infty e^{-zu}\dd u/u\), and choose its
continuation from the Euclidean region, with the principal branch
and its standard logarithmic jump \cite{DLMF6}. Define
\begin{equation}
\begin{aligned}
 \mathsf S_{\eta,\bb}&:=\sMand+\ii\eta\epsvec,\\
 \Delta^{\rm dir}_{\eta,\bb}(x)
       &:=m^2-x(1-x)\mathsf S_{\eta,\bb},\\
 \EoneDir&:=\embed\circ E_1\circ\embed^{-1}.
\end{aligned}
 \label{eq:Gaussian-directed-channel}
\end{equation}
The directed boundary bubble and its extracted coefficient are
\begin{align}
 \mathcal B^{\rm dir}_{t,\eta,\bb}(\sMand)
 &=\lim_{\varepsilon\downarrow0}
   \frac{e^{2tm^2}}{16\pi^2}
      \int_0^1\EoneDir\!\left(
          W_{\min}(x)\Delta^{\rm dir}_{\eta,\bb}(x)
                                  \right)\dd x,
          \nonumber\\
 B^{\rm coeff}_{t,\eta}(\sMand)
 &:=\embed^{-1}\!\left[
          \mathcal B^{\rm dir}_{t,\eta,\bb}(\sMand)\right]
   =\lim_{\varepsilon\downarrow0}
                    B_t(\sMand+\ii\eta\varepsilon),
 \label{eq:Gaussian-directed-boundary-and-coefficient}
\end{align}
where the ordinary analytic sheet function is
\begin{equation}
\begin{gathered}
 B_t(z)=\frac{e^{2tm^2}}{16\pi^2}
    \int_0^1 E_1\!\left(
       W_{\min}(x)[m^2-x(1-x)z]\right)\dd x,\\
 z\in\CC\setminus[4m^2,\infty).
\end{gathered}
 \label{eq:gaussian-E1}
\end{equation}
This representation is locally uniformly convergent on the
stated cut plane. Indeed,
\(W_{\min}(x)x(1-x)=t\max(x,1-x)\leq t\).
For \(z\) in a compact set and \(x\) near either endpoint,
the real part of the \(E_1\) argument therefore tends uniformly
to \(+\infty\), and its positive-half-plane integral gives
exponential endpoint control. On the remaining compact
\(x\)-interval, the argument avoids the \(E_1\) cut whenever
\(z\notin[4m^2,\infty)\). The physical boundary has only
integrable logarithms at the roots of
\(m^2-x(1-x)\sMand\), including the double root at threshold.
Thus the directed boundary values follow from this evaluated
Euclidean function. This argument does not assume a contour
rotation of the growing unparametrized Lorentzian graph.

For \(\sMand>4m^2\), let
\(\beta(\sMand)=\sqrt{1-4m^2/\sMand}\) denote the positive
real phase-space root. The directed branch identity is
\[
\begin{aligned}
 \lim_{\varepsilon\downarrow0}\Bigl[&\EoneDir\!\left(W[-d-\ii\eta c\epsvec]\right)\\
       &-\EoneDir\!\left(W[-d+\ii\eta c\epsvec]\right)\Bigr]
       =2\pi\eta\je,
\end{aligned}
\]
with \(W,d,c>0\).
Here \(W\) has inverse-momentum-squared units, \(d\) has
momentum-squared units, and \(c\) is dimensionless; the
\(E_1\) arguments are dimensionless. On the open cut interval
use \(W=W_{\min}(x)\), \(c=x(1-x)\), and
\(d=c\sMand-m^2>0\). Consequently
\begin{align}
 \Disc^{\rm dir}_\eta\mathcal B_t
 &:=\mathcal B^{\rm dir}_{t,\eta,\bb}
       -\mathcal B^{\rm dir}_{t,-\eta,\bb}
   =\eta\je\,\frac{e^{2tm^2}}{8\pi}
                \beta(\sMand)\theta(\sMand-4m^2),
        \nonumber\\
 \Disc^{\rm coeff}_\eta B_t
 &:=B^{\rm coeff}_{t,\eta}-B^{\rm coeff}_{t,-\eta}
   =\eta\ii\,\frac{e^{2tm^2}}{8\pi}
                \beta(\sMand)\theta(\sMand-4m^2).
 \label{eq:gaussian-disc}
\end{align}
The ordinary absorptive comparison coefficient is recovered as
\begin{equation}
 \frac{\Disc^{\rm dir}_\eta\mathcal B_t}{2\eta\je}
   =\frac{\Disc^{\rm coeff}_\eta B_t}{2\eta\ii}
   =\frac{e^{2tm^2}}{16\pi}
              \beta(\sMand)\theta(\sMand-4m^2).
 \label{eq:Gaussian-directed-absorptive-coefficient}
\end{equation}
This jump concerns the numerator-stripped bubble. Its vertices,
symmetry factor and Dyson coefficient are restored through
\eqref{eq:scalar-benchmark-Dyson-reconstruction}.

With the standard on-shell one-particle measures, for
future-directed \(P\), \(\eta=+1\), and
\(P^2=\sMand>4m^2\), the two-body
phase-space integral is
\begin{equation}
\begin{aligned}
 \int\dd\Phi_{2,\bb}(P)
 &=\int\prod_{i=1}^2\frac{\dd^3\mathbf q_i}{(2\pi)^3\,2\omega_i}\,
        (2\pi)^4\delta^{(4)}_{\bb}(P-q_1-q_2)\\
 &=\frac{\beta(\sMand)}{8\pi},\qquad
 \omega_i=\sqrt{\mathbf q_i^2+m^2}.
\end{aligned}
 \label{eq:two-body-phase-space}
\end{equation}
The \(\omega_i>0\) here are intrinsic on-shell momentum
coefficients in the scalar comparison measure, not replacements
for carrier-valued geo-momentum charges.
For the opposite orientation the invariant measure is
\begin{equation}
 \dd\Pi_{1,\eta,\bb}(q)
  :=\frac{\dd^4q_{\bb}}{(2\pi)^4}\,2\pi\delta^{\rm sh}_{\bb}(q^2-m^2)
       \theta(\eta\ee\!\cdot q),
 \label{eq:oriented-one-particle-measure}
\end{equation}
and the two-body integral with \(P=q_1+q_2\) is
\(\beta(\sMand)\theta(\eta\ee\!\cdot P)
\theta(\sMand-4m^2)/(8\pi)\).
An outgoing cut uses \(q_1=k\) and \(q_2=P-k\), so both
selectors have the same declared physical orientation. Evolution-branch reversal changes the routed covectors together with
the boundary prescription, as in \eqref{eq:fw-cut-routing}; the physical
state measure remains positive and unchanged. An A final particle is
integrated as an outgoing physical state even when in the ascending-GT
ket. The scalar comparisons that follow use $\eta=+1$ in this routed
coefficient convention.

Consequently, a prescribed undressed unit-vertex bubble cut would
give \(\ii\beta/(8\pi)\), whereas
\eqref{eq:gaussian-disc} gives
\(e^{2tm^2}\ii\beta/(8\pi)\).
They differ even though \(\Ph(0)=1\), the kernel is entire, and
there is a strict Laplace-depth gap. The missing factor must be
accounted for by the actual physical matching and completion if
this operational cut is to represent an analytic discontinuity.

The familiar positive-boundary scalar benchmark normalizations are
\[
 \Disc B_{\rm std}=\frac{\ii\beta}{8\pi},\qquad
 \Im B_{\rm std}=\frac{\beta}{16\pi},\qquad
 \Im\mathcal M_{\rm bubble,std}
      =\frac{\lambda^2\beta}{32\pi}.
\]
The last expression includes the \(1/2\) symmetry factor.
For the Gaussian benchmark it acquires the factor \(e^{2tm^2}\)
in the same convention. The tree term \(-\lambda\) is real;
the factor \(\ii\) in the Dyson coefficient \(\ii\mathcal M\)
is not a tree-level absorptive part.

\subsubsection*{Finite shell normalization and the remaining amplitude}

The next proposition concerns the ordinary analytic sheet functions
\(B_t(z)\), \(b_t(z)\), \(b_{\rm std}(z)\), and \(R_t(z)\).
Their physical directed boundary coefficients are obtained by the
explicit transport below; these coefficient identities do not
replace the full graph reconstruction.

The Gaussian example also permits an explicit normalization test.
Multiplying each scalar line by \(e^{-tm^2}\) makes its massive
pole coefficient unity. For the bubble this gives the comparison
normalization
\begin{equation}
 \widehat B_t(\sMand)=e^{-2tm^2}B_t(\sMand),\qquad
 b_t(\sMand)=\widehat B_t(\sMand)-\widehat B_t(0).
 \label{eq:normalized-bubble}
\end{equation}
Equation~\eqref{eq:scalar-amplitude} continues to use the original
\(B_t\). Define the finite ordinary comparison function by
\begin{equation}
 b_{\rm std}(\sMand)
 =\frac1{16\pi^2}\int_0^1\dd x\,
   \log\!\frac{m^2}{m^2-x(1-x)\sMand},
 \label{eq:subtracted-standard-bubble}
\end{equation}
with the logarithm continued from \(\sMand<4m^2\).

\begin{proposition}[Entire remainder after finite normalization]
\label{prop:entire-bubble-remainder}
For \(m>0\) and \(t>0\), on
\(\CC\setminus[4m^2,\infty)\),
\begin{equation}
 b_t(\sMand)=b_{\rm std}(\sMand)+R_t(\sMand),
 \label{eq:bubble-entire-decomposition}
\end{equation}
where the remainder is entire and has the representation
\begin{equation}
 R_t(\sMand)
 =\frac1{16\pi^2}\int_0^1\dd x
   \int_0^{W_{\min}(x)}\frac{\dd W}{W}\,
       e^{-m^2W}\bigl[1-e^{Wx(1-x)\sMand}\bigr].
 \label{eq:entire-bubble-remainder}
\end{equation}
It vanishes at zero, is real on the real axis, and satisfies
\begin{equation}
 |R_t(\sMand)|
   \leq\frac{3t|\sMand|}{64\pi^2}e^{t|\sMand|}.
 \label{eq:entire-remainder-bound}
\end{equation}
Thus \(b_t\to b_{\rm std}\) uniformly on compact subsets of the
cut plane as \(t\downarrow0\).
\end{proposition}

\begin{proof}
For \(\Re\sMand<4m^2\), the proper-time form of
\eqref{eq:subtracted-standard-bubble} has lower limit \(W=0\).
The normalized difference in \eqref{eq:normalized-bubble} has lower
limit \(W_{\min}(x)\). Subtracting the two convergent integrals
gives \eqref{eq:entire-bubble-remainder}.

Set \(c(x)=x(1-x)\). Throughout the remaining integration region,
\(0\leq Wc(x)\leq t\max(x,1-x)\leq t\). Consequently,
\[
 \frac{|1-e^{Wc(x)\sMand}|}{W}
       \leq c(x)|\sMand|e^{t|\sMand|}.
\]
This is integrable on the region because
\(\int_0^1c(x)W_{\min}(x)\dd x=3t/4\).
It supplies a common majorant on every compact \(\sMand\)-set,
so the remainder is entire by Theorem~\ref{thm:holomorphy} and
satisfies \eqref{eq:entire-remainder-bound}.
Analytic continuation proves \eqref{eq:bubble-entire-decomposition}
on the cut plane; the same bound gives compact-uniform convergence.
\end{proof}

For the directed argument \(\mathsf S_{\eta,\bb}\), set
\(\LogDir:=\embed\circ\log\circ\embed^{-1}\), with the
branch continued from the real subthreshold domain. The same
identity reads, before taking its boundary limit,
\begin{align}
 b_{\rm std}^{\rm dir}(\mathsf S_{\eta,\bb})
 &=\frac1{16\pi^2}\int_0^1
    \LogDir\!\left[
       \frac{m^2}{m^2-x(1-x)\mathsf S_{\eta,\bb}}
             \right]\dd x,\nonumber\\
 R_t^{\rm dir}(\mathsf S_{\eta,\bb})
 &=\frac1{16\pi^2}\int_0^1\dd x
       \int_0^{W_{\min}(x)}\frac{\dd W}{W}\,e^{-m^2W}\nonumber\\
 &\quad\times\left[1-\ExpDir\!\left(
          Wx(1-x)\mathsf S_{\eta,\bb}\right)\right],
           \nonumber\\
 b_t^{\rm dir}(\mathsf S_{\eta,\bb})
 &=b_{\rm std}^{\rm dir}(\mathsf S_{\eta,\bb})
                         +R_t^{\rm dir}(\mathsf S_{\eta,\bb}).
 \label{eq:normalized-bubble-directed-transport}
\end{align}
Each superscript \({\rm dir}\) means transport by \(\embed\);
applying \(\embed^{-1}\) gives the preceding coefficient
identities. Since \(R_t\) is entire and real on the real axis,
its real-axis directed boundary is carrier-even and has no jump.
Its contribution to a complete Dyson graph still carries that
graph's vertex and Dyson factors.

The finite normalization removes the extra shell factor from the
discontinuity, while the entire term retains nontrivial momentum
dependence: \(R_t(\sMand)<0\) for \(0<\sMand<4m^2\).
It also changes the real part above threshold. Matching the
discontinuity therefore does not determine the full amplitude.
At every \(t>0\), the subtraction in
\eqref{eq:normalized-bubble} removes an already finite reference
value. The controlled limit concerns that normalized difference;
the unsubtracted value in \eqref{eq:gaussian-zero} still diverges
as the damping is removed.

\begin{corollary}[The Gaussian remainder is not polynomially bounded]
\label{cor:Gaussian-growth-obstruction}
For real \(s>0\), the entire remainder of
Proposition~\ref{prop:entire-bubble-remainder} satisfies
\begin{equation}
 -R_t(s)\geq
 \frac{\log2}{48\pi^2}e^{-2tm^2}
       \left(e^{2ts/9}-1\right).
 \label{eq:Gaussian-remainder-growth-lower}
\end{equation}
Thus no addition of a finite polynomial makes \(R_t\), or the
normalized bubble boundary \(b_t\), polynomially bounded along
the positive real axis at fixed \(t>0\).
\end{corollary}

\begin{proof}
For positive \(s\), the integrand of \(-R_t(s)\) in
\eqref{eq:entire-bubble-remainder} is nonnegative. Restrict to the
rectangle \(1/3\leq x\leq2/3\), \(t\leq W\leq2t\),
which lies inside \(W\leq W_{\min}(x)\).
There \(x(1-x)\geq2/9\) and \(e^{-m^2W}\geq e^{-2tm^2}\).
Its \(x\)-length is \(1/3\) and
\(\int_t^{2t}\dd W/W=\log2\), proving the bound.
For the ordinary comparison, direct integration of
\eqref{eq:subtracted-standard-bubble} gives on the upper rim,
for \(s>4m^2\),
\[
 16\pi^2 b_{\rm std}(s+\ii0)
 =2-\beta(s)\log\frac{1+\beta(s)}{1-\beta(s)}
       +\ii\pi\beta(s).
\]
It grows at most logarithmically as \(s\to+\infty\).
It cannot cancel the real negative exponential lower bound, nor
can a finite polynomial.
\end{proof}

The isolated Gaussian benchmark therefore still fails the growth
assumptions of a finite-subtraction dispersion argument after its jump
has been matched. For a completed physical amplitude, all contributions
must be analyzed together: additional nonpolynomial dynamics could
supply a cancellation. The compact-uniform limit \(t\downarrow0\)
remains valid, but it cannot be interchanged with this large-\(s\) limit.

\subsubsection*{A reproducible Euclidean numerical check}

At \(P_E=0\), radial integration gives
\begin{align}
 B_{t,E}(0)
 &=\frac1{16\pi^2}
     \int_0^\infty
       \frac{u\,e^{-2tu}}{(u+m^2)^2}\dd u \nonumber\\
 &=\frac{(1+x)e^xE_1(x)-1}{16\pi^2},
 \qquad x=2tm^2 .
 \label{eq:gaussian-zero}
\end{align}
For \(\tau_*=1\) and \(am=1\),
\begin{equation}
 B_{t,E}(0)=
 0.0005318466124095473\ldots .
 \label{eq:numerical-bubble}
\end{equation}
The same value follows from the independent dimensionless
proper-time integral
\[
 \frac1{16\pi^2}\int_0^\infty
       \frac{v\,e^{-v}}{(v+2)^2}\dd v .
\]
This number specifies the kernel, scales, external momentum,
measure, and absence of coupling or symmetry factors. It is an
off-shell Euclidean check, not a predicted scattering measurement.
\subsection{A massive stable-half bubble need not be analytic at zero}
\label{sec:stable-bubble-origin}

The absence of a positive exponential moment for the canonical
stable-half law cannot in general be bypassed by citing positive
internal masses. Consider the same literal, fixed-carrier scalar
coefficient model as \eqref{eq:two-line-bubble}, now with
\[
 C_{a,m}(q)=\frac{e^{-a|q|}}{q^2+m^2},\qquad a>0,\quad m>0.
\]
Let \(\mathcal C_{a,m}(x)\) be its inverse Fourier transform
with measure \(\dd^4q_{\bb}/(2\pi)^4\), and write
\[
 B_{a,m}(P)=\int\frac{\dd^4q_{\bb}}{(2\pi)^4}
                   C_{a,m}(q)C_{a,m}(q+P).
\]

\begin{proposition}[A zero-momentum nonanalyticity of the stable-half bubble]
\label{prop:stable-bubble-nonanalytic}
The function \(B_{a,m}\) is finite for all real \(P\) and is
\(C^5\). It is not \(C^6\) at \(P=0\); in particular, no
function holomorphic in \(P^2\) on a neighbourhood of zero can
restrict to this bubble there. Its coordinate covariance has the
positive large-distance tail
\begin{equation}
 \mathcal C_{a,m}(x)=
   \frac{3a}{4\pi^2m^2}|x|^{-5}+O(|x|^{-6}).
 \label{eq:stable-covariance-long-tail}
\end{equation}
For $b(h)=B_{a,m}(h e_1)$ and real $h\to0$, the first nonanalytic
term is
\begin{equation}
 \begin{aligned}
 b(h)={}&b(0)+\frac{b''(0)}{2}h^2+\frac{b^{(4)}(0)}{24}h^4\\
        &-\frac{a^2}{8192\pi^2m^4}h^6\log\frac{m}{|h|}
          +O(h^6).
 \end{aligned}
 \label{eq:stable-bubble-leading-log}
\end{equation}
\end{proposition}

\begin{proof}
The complete Fourier and finite-difference argument is given in
\ref{sec:stable-bubble-proof}. In brief, the inverse
transform of \(e^{-a|q|}\) is the positive four-dimensional
Poisson kernel
\(3a/[4\pi^2(a^2+|x|^2)^{5/2}]\). Convolution with the
positive massive free Green function, whose total integral is
\(m^{-2}\), gives \eqref{eq:stable-covariance-long-tail}.
The bubble is the Fourier transform of
\(\mathcal C_{a,m}(x)^2\). This nonnegative function has all
absolute moments of order less than six. At order six, the radial
moment grows logarithmically. This gives
\eqref{eq:stable-bubble-leading-log} and rules out a sixth-order
Taylor expansion at zero.
\end{proof}

This is an infrared regularity obstruction, not an ultraviolet
divergence: all inverse Laplace moments and the fixed-graph UV
estimates remain valid. A positive mass in the denominator does
not remove the long-distance tail introduced by the multiplier's
cusp at \(q=0\). Finite polynomial additions cannot restore the
missing differentiability. A physical completed amplitude may
contain different vertices and compensating contributions; it
must establish any claimed analyticity separately. The proposition
does not identify this one auxiliary bubble with a physical
observable or prove failure of every completion.

The result tests the monograph's canonical kernel directly through
its stated Euclidean line rule. It is not a
restatement of the monograph's conditional holomorphy theorem.
The effective single-block model in the next subsection has a
different kernel placement and is not subject to this conclusion
merely because it uses the same scalar profile.

\subsection{An effective radial model}
\label{sec:effective-block}

\begin{remark}[Linewise prescription and channel boundary path]
\label{rem:linewise-versus-channel-boundary}
For real Lorentzian momenta \(k,P\) and finite
\(\varepsilon>0\), the common directed prescription combines
exactly as
\begin{align}
 &\frac1{[k^2-m^2+\ii\eta\epsvec]
             [(P-k)^2-m^2+\ii\eta\epsvec]}
       \nonumber\\
 &\qquad=\int_0^1
       \frac{\dd x}{[(k-xP)^2-\Delta(x)
                                      +\ii\eta\epsvec]^2},
 \nonumber\\
 &\Delta(x)=m^2-x(1-x)P^2.
 \label{eq:exact-directed-Feynman-combination}
% Monograph source: eq:feynman_param_final.
\end{align}
In each character the affine denominator has imaginary part
\(\kappa\eta\varepsilon\), so it has no zero for
\(0\leq x\leq1\). For complex external momenta the same
real-parameter representation is used only where both character
segments remain nonzero, or through a separately justified
analytic continuation.
The linewise infinitesimal remains independent of \(x\).
With \(\sMand=P^2\), continuation of an already evaluated
invariant function through
\(P_E^2\mapsto-\mathsf S_{\eta,\bb}\) gives
\begin{equation}
 \Delta^{\rm dir}_{\eta,\bb}(x)
   =m^2-x(1-x)\mathsf S_{\eta,\bb}
   =\Delta(x)-\ii\eta x(1-x)\epsvec.
 \label{eq:channel-directed-effective-denominator}
% Monograph source: eq:effective-kernel-channel-vectorial-denominator.
\end{equation}
These are distinct finite-\(\varepsilon\) formulas. Their
boundary values can be identified only after the selected
threshold-excluding continuation and its non-tangential limit
have been justified. In particular, the Gaussian closed form is
defined by its Euclidean continuation; no equality with a raw
growing Lorentzian loop integral is assumed.
\end{remark}

For a running-unit graph the real shift \(\ell=k-xP\)
also transforms its numerator signs:
\begin{equation}
 \sgn(k_\parallel)\sgn(P_\parallel-k_\parallel)
 =\sgn(\ell_\parallel+xP_\parallel)
   \sgn((1-x)P_\parallel-\ell_\parallel).
 \label{eq:running-signs-after-parameter-shift}
% Monograph source: eq:raw-kernel-under-feynman-translation; eq:running-unit-versus-fixed-virtual-carrier.
\end{equation}
Consequently the denominator combination alone does not turn
that graph into a radial fixed-carrier integral.

For comparison, one may declare an effective single virtual block
after combining denominators:
\begin{equation}
\begin{aligned}
 B_{{\rm eff},E}(P_E)
   &=\int_0^1\dd x\int\frac{\dd^4\ell_{E,\bb}}{(2\pi)^4}
       \frac{\Ph_{\rm eff}(a^2\ell_E^2)}
            {[\ell_E^2+\Delta_E(x)]^2},\\
 \Delta_E(x)&=m^2+x(1-x)P_E^2.
\end{aligned}
 \label{eq:effective-model}
\end{equation}
The measure \(\rho_{\rm eff,\bb}\) belongs to this model. It is not
automatically the measure for a product of the two line kernels
in \eqref{eq:two-line-bubble}.

For the explicit choice \(\Ph_{\rm eff}(z)=e^{-\sqrt z}\),
define
\begin{equation}
 J(u;a)=\int_0^\infty r^3e^{-ur^2-ar}\dd r,
 \qquad u>0,\quad a>0.
 \label{eq:J-radial}
\end{equation}
Here $r$ is a radial momentum variable, not the normal propagation
depth used in Section~\ref{sec:update-geometric-source}. The function
$J(u;a)$ is a scalar calculus identity; the fixed-copy momentum measure
has already been reduced to its radial Jacobian. A Schwinger parameter
for the squared denominator gives
\begin{equation}
 B_{{\rm eff},E}(P_E)
   =\frac1{8\pi^2}\int_0^1\dd x
        \int_0^\infty\dd u\,
              u\,e^{-u\Delta_E(x)}J(u;a).
 \label{eq:effective-J}
\end{equation}
The corresponding directed external-invariant continuation,
with \(\Delta^{\rm coeff}_{\eta,\bb}(x)
:=\embed^{-1}[\Delta^{\rm dir}_{\eta,\bb}(x)]\),
is initially defined where
\(\Re\Delta^{\rm coeff}_{\eta,\bb}(x)>0\):
\begin{equation}
 \mathcal B^{\rm dir}_{{\rm eff},\eta,\bb}(\sMand;\varepsilon)
  =\frac1{8\pi^2}\int_0^1\dd x\int_0^\infty\dd u\,
       u J(u;a)\ExpDir\!\left[
                  -u\Delta^{\rm dir}_{\eta,\bb}(x)\right].
 \label{eq:effective-directed-J}
\end{equation}
Its ordinary coefficient is obtained by \(\embed^{-1}\);
the physical boundary additionally takes \(\varepsilon\downarrow0\)
when that limit exists on the prescribed domain.
Beyond that directly convergent domain it denotes the continuation
of the established function, whenever the required domain and
majorant exist, not the same positive-real proper-time integral.
The roots of positive real \(u\) in \(J(u;a)\) remain ordinary
real roots; the directed argument is the continued invariant
\(\Delta^{\rm dir}_{\eta,\bb}\).

The exact radial function is
\begin{equation}
 \boxed{\;
 J(u;a)=\frac{
 1+v^2-\sqrt\pi\,v(v^2+\tfrac32)e^{v^2}\erfc(v)}
 {2u^2},
 \qquad v=\frac{a}{2\sqrt u}.
 \;}
 \label{eq:J-closed}
\end{equation}
It follows by applying \(-\partial_a^3\) to
\[
 \int_0^\infty e^{-ur^2-ar}\dd r
 =\frac{\sqrt\pi}{2\sqrt u}\,
      e^{a^2/(4u)}\erfc\!\left(\frac{a}{2\sqrt u}\right).
\]
The minus sign comes from differentiating \(e^{-ar}\) three times.
The useful endpoint checks are
\begin{equation}
 J(u;0)=\frac1{2u^2},\qquad
 J(0;a)=\frac6{a^4}.
 \label{eq:J-limits}
\end{equation}
The second is the limit \(u\downarrow0\) at fixed \(a>0\);
it makes the small-\(u\) integrand of
\eqref{eq:effective-J} integrable.
For numerical evaluation one may use
\(\erfcx(v)=e^{v^2}\erfc(v)\), with the integral
\eqref{eq:J-radial} or a controlled endpoint expansion when
subtraction causes loss of precision.

The alternative
\[
 J_M(u;a,m)
    =\int_0^\infty r^3e^{-ur^2-a\sqrt{r^2+m^2}}\dd r
\]
defines a different massive profile rather than a change of variables
in \eqref{eq:J-radial}. Both integrals are used through their stated
representations; no Tricomi or Appell reconstruction is assumed.
The original linewise kernel, the effective block, and the massive
radial profile therefore retain distinct model labels.

\subsection{Gauge and ghost covariance coefficients}

For a flat massless Euclidean gauge benchmark with covariant
parameter \(\zeta>0\), the quadratic coefficients may be chosen as
\begin{equation}
 D_{\mu\nu,E}^{\Ph}(q)
   =\frac{\Ph(a^2q^2)}{q^2}
         \bigl(P_{T,\mu\nu}+\zeta P_{L,\mu\nu}\bigr),
 \qquad
 C_{{\rm gh},E}^{\Ph}(q)=\frac{\Ph(a^2q^2)}{q^2},
 \label{eq:gauge-covariances}
\end{equation}
where \(P_{L,\mu\nu}=q_\mu q_\nu/q^2\) and
\(P_T=\delta-P_L\).
The action and ghost ordering that realize these coefficients
must still be specified, as in
Section~\ref{sec:brst}. In particular, transverse damping cannot
be deduced from BRST-exact gauge fixing alone.

For Lorentzian coefficients put
\begin{equation}
 Z_{\eta,\bb}(k):=k^2+\ii\eta\epsvec,
 \qquad
 F^{\rm dir}_{\eta,\bb}(k):=
              \Ph_{\ee}(-a^2Z_{\eta,\bb}(k)).
 \label{eq:gauge-directed-kinetic-data}
\end{equation}
The following displays use the monograph's resolved scalar-pole
covariance prescription at finite \(\varepsilon>0\); their
admitted boundary distributions take the common limit
\(\varepsilon\downarrow0^+\). They are not universal
finite-\(\varepsilon\) identities for the inverse of an
arbitrary regulated matrix Hessian. For the Lorentzian Maxwell coefficient $-F_{\mu\nu}F^{\mu\nu}/4$
and the corresponding covariant gauge fixing in signature $(+---)$,
the directed gauge and ghost templates are
\begin{align}
 \vec D^{\Ph}_{\mu\nu;\eta,\bb}(k)
  &=-\eta\ii\widehat u_{k,\bb}\,
      F^{\rm dir}_{\eta,\bb}(k)
       \left[\frac{g_{\mu\nu}}{Z_{\eta,\bb}(k)}
        -(1-\zeta)\frac{k_\mu k_\nu}
                           {Z_{\eta,\bb}(k)^2}\right],
                   \nonumber\\
 \vec C^{\Ph}_{{\rm gh};\eta,\bb}(k)
  &=\eta\ii\widehat u_{k,\bb}\,
      \frac{F^{\rm dir}_{\eta,\bb}(k)}
           {Z_{\eta,\bb}(k)}.
 \label{eq:physical-directed-gauge-ghost-lines}
% Monograph source: eq:general-directed-gauge-coefficient; eq:general-directed-ghost-coefficient.
\end{align}
The gauge tensor is symmetric in its two Lorentz indices before
multiplication by the running unit. Consequently its unrestricted
running-unit version has the same exchange obstruction as the real
scalar: it is not an inverse-Hessian covariance for a single real
commuting gauge field without an additional compatible endpoint
realization. The fixed-carrier benchmark does not have this oddness
obstruction. Ghost/antighost and Dirac fields instead require their
complete graded, distinct-field source matrices.

The tensor and ghost coefficients inside the oriented numerator
are extracted by \(\embed^{-1}\) on their directed scalar
entries. Their full lines are reconstructed as in
\eqref{eq:scalar-line-coefficient}, or character by character
when they form part of a mixed carrier graph. The longitudinal
double pole for \(\zeta\ne1\) is already present before the
kernel is inserted.
The minus sign in the photon template is the Maxwell sign also used in
\eqref{eq:OQED-directed-lines}. It is not an antiverse sign and is not
removed by taking an absolute Euclidean bound. The ghost sign additionally
uses the ordered ghost/antighost source convention. A different quadratic
sign would be a different coefficient action, not the same physical
photon covariance.

For reference, the corresponding massive tensor coefficient is
\begin{equation}
 \mathcal R^{\rm dir}_{\mu\nu;\eta,\bb}(k)
 =\frac{F^{\rm dir}_{\eta,\bb}(k)}
           {Z_{\eta,\bb}(k)-m^2}
    \left[g_{\mu\nu}-(1-\zeta)
       \frac{k_\mu k_\nu}{Z_{\eta,\bb}(k)-\zeta m^2}\right].
 \label{eq:massive-directed-gauge-coefficient}
\end{equation}
The ghost denominator is
\(Z_{\eta,\bb}(k)-\zeta m^2\). These coefficients specify
only the quadratic part; a massive gauge theory still requires
its compatible Higgs/Goldstone, ghost and interaction completion.

For a Lorentzian Dirac line use
\(\{\gamma^\mu,\gamma^\nu\}=2g^{\mu\nu}\one\) and
\(\slashed k=\gamma^\mu k_\mu\). In the monograph's
entrywise coefficient convention, define
\begin{equation}
\begin{aligned}
 R^{D,\Ph}_\eta(k;\varepsilon)
   &:=(\slashed k+m)R^\Ph_\eta(k;\varepsilon),\\
 \vec S^{\Ph}_{\eta,\bb}(k;\varepsilon)
   &:=\eta\ii\widehat u_{k,\bb}\,
                     \embed(R^{D,\Ph}_\eta(k;\varepsilon)).
\end{aligned}
 \label{eq:physical-directed-Dirac-line}
\end{equation}
This rationalized Dirac coefficient prescribes the scalar pole
boundary. For example, without the kernel and with
\(D=k^2-m^2\),
\begin{equation}
 (\slashed k-m)
     \frac{\slashed k+m}{D+\ii\eta\varepsilon}
   =\frac{D}{D+\ii\eta\varepsilon}\,\one.
 \label{eq:Dirac-finite-boundary-check}
\end{equation}
Thus it is not the inverse of the unmodified Dirac matrix at
finite \(\varepsilon\). When an exact regulated matrix
inverse or a Ward identity at finite regulator is required, the
completed Hessian and its boundary datum must be specified and
inverted together in each character. Its identification with a
resolved scalar-pole boundary requires the corresponding limit
argument.
% Monograph source: eq:Omni-CC-Hessian-Green-vectorial-boundary-firewall; eq:background-commuting-virtual-covariance.
The complete spinor coefficient, including its matrix entries,
is embedded consistently. If a different central matrix
realization is retained, use its two character representatives
and \eqref{eq:complete-graph-character-reconstruction}; do not
apply \(\embed^{-1}\) to a mixed-algebra matrix.

\subsection{QED vacuum polarization and the Ward identity}
\label{sec:QED}

Use canonically normalized Euclidean matrices
\(\{\gamma^\mu,\gamma^\nu\}=2\delta^{\mu\nu}\), and set
\begin{equation}
\begin{aligned}
 S_0(k)&=\frac{-\ii\slashed k+m}{k^2+m^2},&
 S_0^{-1}(k)&=\ii\slashed k+m,\\
 \Gamma_0^\mu&=\ii\gamma^\mu,&
 F_k&=\Ph(a^2k^2).
\end{aligned}
 \label{eq:QED-conventions}
\end{equation}
The two-line graph with bare vertices is
\begin{equation}
 \Pi_{\rm bv}^{\mu\nu}(Q)
 =-e^2\int\frac{\dd^4k_{\bb}}{(2\pi)^4}\,
 F_kF_{k+Q}
 \Tr\!\left[\Gamma_0^\mu S_0(k)
                 \Gamma_0^\nu S_0(k+Q)\right].
 \label{eq:bare-QED-loop}
\end{equation}
An overall plane binder
\(\mathcal P^{(2)}_{\bb\mid{\rm ext}}\) is understood.
The factors \(\ii\) are included in \(\Gamma_0^\mu\).
Under an otherwise identical canonical complex identification of
the mass, gamma matrices and binder, this tensor is the negative
of the monograph's displayed \(\gamma\)-vertex loop with the
same prefactor \(-e^2\): the two vertices contribute
\(\ii^2=-1\) \cite{ChiapponeGHOUQFT2026}.
% Monograph source: eq:Pi-Omni-full.
The Ward contraction and the benchmark below use
\eqref{eq:QED-conventions} throughout. Comparison with a physical
current correlator also retains its declared tensor and
Wick-rotation normalizations.

The calculation uses Euclidean coefficients and a fixed virtual
carrier. Running-unit Lorentzian graphs retain their mode signs
and require a separate evaluation.

Contract with \(Q_\mu\), using
\(Q_\mu\Gamma_0^\mu=S_0^{-1}(k+Q)-S_0^{-1}(k)\).
Cyclicity of the trace gives
\[
 Q_\mu\Pi_{\rm bv}^{\mu\nu}
 =-e^2\int\frac{\dd^4k_{\bb}}{(2\pi)^4}F_kF_{k+Q}
  \Tr\{\Gamma_0^\nu[S_0(k)-S_0(k+Q)]\}.
\]
Assume that the original tensor integrand is absolutely integrable
componentwise at the fixed real Euclidean \(Q\). Its contraction
is initially kept under one integral: the two summands need not
be separately integrable. For the finite positive completely
monotone kernels considered here, Lemma~\ref{lem:Ward-common-cutoff}
justifies the translation on a common bounded domain, proves that
the boundary error vanishes, and proves absolute integrability of
the final remainder. The lemma treats \(m>0\) and also covers
the unchanged locally integrable massless free expressions. It yields
\begin{equation}
 Q_\mu\Pi_{\rm bv}^{\mu\nu}(Q)
 =-e^2\int\frac{\dd^4k_{\bb}}{(2\pi)^4}
 F_k\,[F_{k+Q}-F_{k-Q}]
           \Tr[\Gamma_0^\nu S_0(k)].
 \label{eq:Ward-remainder}
\end{equation}
This remainder is generally nonzero.

\begin{proposition}[Nontransversality under an exact one-loop integrability test]
\label{prop:QED-general-obstruction}
Let \(F(u)=\Ph(a^2u)\) for a nonzero finite positive completely
monotone law, let \(m\geq0\), and assume
\begin{equation}
 \mathcal Q_F:=\int_0^\infty F(u)^2\dd u<\infty.
 \label{eq:QED-Ltwo-condition}
\end{equation}
Then \eqref{eq:bare-QED-loop} is absolutely convergent for every
real \(Q\), is continuous in \(Q\), and obeys
\begin{equation}
\begin{aligned}
 \Pi_{\rm bv}^{\mu\nu}(0)&=\kappa_F\delta^{\mu\nu},\\
 Q_\mu\Pi_{\rm bv}^{\mu\nu}(Q)
      &=\kappa_F Q^\nu+o(|Q|),\\
 \kappa_F
 &=\frac{e^2}{8\pi^2}\int_0^\infty
     \frac{u(u+2m^2)}{(u+m^2)^2}F(u)^2\dd u.
\end{aligned}
 \label{eq:general-Ward-obstruction}
\end{equation}
Here \(\kappa_F>0\) when \(e\ne0\).
For \(m=0\), the rational factor in the integral is one for
\(u>0\). Conversely, absolute convergence of the original
nontrivial tensor at any fixed real \(Q\) implies
\eqref{eq:QED-Ltwo-condition}, by the radial-tail argument in
Lemma~\ref{lem:Ward-common-cutoff}. No differentiability at
\(u=0\) is assumed.
\end{proposition}

\begin{proof}
Choose Hermitian Euclidean gamma matrices and put
\(S_F(k)=F(k^2)S_0(k)\). Then
\[
 \operatorname{tr}[S_F(k)^\dagger S_F(k)]
       =\frac{4F(k^2)^2}{k^2+m^2}.
\]
Equation~\eqref{eq:QED-Ltwo-condition} makes all matrix entries
of \(S_F\) square-integrable, including at \(k=0\) when
\(m=0\). Cauchy--Schwarz bounds the tensor product at every
translate \(Q\); continuity follows from continuity of
translations in \(L^2\). Direct evaluation of the trace at zero
and spherical averaging give
\[
 \Pi_{\rm bv}^{\mu\nu}(0)
 =2e^2\delta^{\mu\nu}\int\frac{\dd^4k_{\bb}}{(2\pi)^4}
          \frac{k^2+2m^2}{(k^2+m^2)^2}F(k^2)^2,
\]
which is \eqref{eq:general-Ward-obstruction}. Its integrand is
positive. The stated small-\(Q\) relation uses only continuity.
For the converse, the trace and monotonicity estimates in
\eqref{eq:Ward-radial-tail} imply
\(\int^\infty rF(r^2)^2\dd r<\infty\). Boundedness of
\(F\) near zero completes \eqref{eq:QED-Ltwo-condition}.
\end{proof}

The exact one-loop condition also has a direct measure test. Tonelli's
theorem applied to the two independent line parameters gives
\begin{equation}
 \mathcal Q_F=\frac1{a^2}
   \int_{[0,\infty)^2}
       \frac{\rho_{\bb}(\dd\tau)\rho_{\bb}(\dd\sigma)}{\tau+\sigma}.
 \label{eq:QED-convolution-moment}
\end{equation}
The equality holds in \([0,\infty]\), with the integrand infinite
at \((0,0)\). If finite, it excludes an atom of \(\rho_{\bb}\) at
zero and equals \(a^{-2}M_1[\rho_{\bb}*\rho_{\bb}]\), where \(\rho_{\bb}*\rho_{\bb}\)
is the convolution measure. It need not require \(M_1[\rho_{\bb}]<\infty\).
Indeed, integrating \(e^{-a^2(\tau+\sigma)u}\) over \(u>0\)
proves the formula directly. This test is for the two-line QED graph;
it is not an all-orders replacement for the moment criteria in
Section~\ref{sec:uv}.

\paragraph{How much mass can the two heat parameters place near zero?}
The convolution test can be read directly from the lower tail of
the two durations. Let $\nu$ be a finite positive joint heat law
and $H_+(t)=\nu\{(\tau,\sigma):\tau+\sigma\le t\}$. In fixed units,
\[
 \int\frac{\dd\nu(\tau,\sigma)}{\tau+\sigma}
       =\int_0^\infty H_+(t)\frac{\dd t}{t^2}.
\]
The equality includes a possible infinite value at the common zero.
It follows by integrating $\mathbf1_{\tau+\sigma\le t}/t^2$
first in $t$. At zero transfer, with the same bare Dirac vertices
and the product of the two conditional heat factors, this is the
exact ultraviolet test. Indeed the radial trace coefficient in
\eqref{eq:general-Ward-obstruction} contains the joint transform at
$a^2u$, multiplied by a bounded rational factor tending to one.
Its finite-$u$ part is integrable. No separate test of either heat
time detects the whole joint cost.

For independent copies of $\rho_{\bb}$ put
$H(t)=\rho_{\bb}([0,t])$. The two inclusions for nonnegative
durations give $H(t/2)^2\le H_+(t)\le H(t)^2$.
The exact test is therefore finiteness of
$\int_0^1 H(t)^2\dd t/t^2$. For two copies of the same random
duration $H_+(t)=H(t/2)$, so the test uses $H$ to the first power.
For instance, suppose near zero $H(t)$ is bounded above and below
by positive multiples of $t^\beta[\log(1/t)]^{-\gamma}$ with
$\beta>0$. The independent case converges for $\beta>1/2$, or
for $\beta=1/2$ and $\gamma>1/2$. The shared case converges for
$\beta>1$, or for $\beta=1$ and $\gamma>1$. All remaining cases
diverge. These thresholds concern the unmodified two-vertex
comparison at zero transfer.

An integration by parts also gives
\begin{equation}
 \kappa_F=-\frac{e^2}{4\pi^2}\int_0^\infty
                 \frac{u^2F(u)F'(u)}{u+m^2}\dd u.
 \label{eq:general-Ward-obstruction-derivative}
\end{equation}
Indeed, complete monotonicity gives \(F'\leq0\) for \(u>0\),
and the boundary term
\(u^2F(u)^2/(u+m^2)\) vanishes at both endpoints.
At infinity this follows from integrability and monotonicity of
\(F^2\); near zero it follows from boundedness of \(F\).
One first integrates on a compact positive interval and then uses
the nonnegative integrands to pass to the endpoints. This proves
convergence of the derivative formula even for a cusp at zero.

For example, in the massless canonical stable-half comparison,
\begin{equation}
 F(u)=e^{-a\sqrt u},\qquad
 \kappa_F=\frac{e^2}{16\pi^2a^2}.
 \label{eq:stable-half-Ward-coefficient}
\end{equation}
For \(F(u)=(1+a^2u)^{-\nu}\) the condition is exactly
\(\nu>1/2\), and at \(m=0\),
\(\kappa_F=e^2/[8\pi^2a^2(2\nu-1)]\).
These are finite nontransverse coefficients of the declared bare
comparison tensor, not physical photon masses. The general remainder
is \(o(|Q|)\); a cubic error requires additional smoothness.

For a sharper remainder, suppose \(k\mapsto F(k^2)\) is \(C^3\).
Require the first through third \(Q\)-derivatives of the complete
weighted expression
\[
 F_k F((k+Q)^2)\Tr[\Gamma_0^\nu S_0(k)]
\]
to have a common \(L^1(\dd^4k_{\bb})\) majorant for \(Q\) near zero.
This includes the numerator weight and any infrared behaviour, not
just derivatives of the scalar profile alone. Then use
\(\Tr[\Gamma_0^\nu S_0(k)]=4k^\nu/(k^2+m^2)\) and
\[
 F((k+Q)^2)-F((k-Q)^2)
       =4(k\cdot Q)F'(k^2)+O(|Q|^3).
\]
Rotational symmetry then gives
\begin{equation}
 Q_\mu\Pi_{\rm bv}^{\mu\nu}(Q)
 =-4e^2Q^\nu
   \int\frac{\dd^4k_{\bb}}{(2\pi)^4}
       \frac{k^2F(k^2)F'(k^2)}{k^2+m^2}
       +O(|Q|^3).
 \label{eq:Ward-smallQ}
\end{equation}
For the Gaussian kernel \(F(k^2)=e^{-tk^2}\), with \(m>0\),
put \(x_t=2tm^2\). The linear coefficient has the closed form
\begin{align}
 \kappa_t
 &:=-4e^2\int\frac{\dd^4k_{\bb}}{(2\pi)^4}
       \frac{k^2F(k^2)F'(k^2)}{k^2+m^2} \nonumber\\
 &=\frac{e^2t}{4\pi^2}\int_0^\infty
       \frac{u^2e^{-2tu}}{u+m^2}\dd u \nonumber\\
 &=\frac{e^2}{16\pi^2t}
       \left[1-x_t+x_t^2e^{x_t}E_1(x_t)\right].
 \label{eq:gaussian-Ward-coefficient}
\end{align}
The last equality uses
\(u^2/(u+m^2)=u-m^2+m^4/(u+m^2)\).
The positive integral proves \(\kappa_t>0\) for \(e\neq0\), so
\begin{equation}
 Q_\mu\Pi_{\rm bv}^{\mu\nu}(Q)
       =\kappa_tQ^\nu+O(|Q|^3),\qquad
 \Pi_{\rm bv}^{\mu\nu}(0)=\kappa_t\delta^{\mu\nu}.
 \label{eq:gaussian-Ward-violation}
\end{equation}
The second equality follows by rotational covariance and the smooth
expansion at \(Q=0\), or directly by integrating the zero-momentum
trace by parts. This finite nontransverse contribution is a specific
term that the completed functional must cancel; it does not determine
a physical photon mass.

\subsubsection*{A compatible longitudinal vertex}

Let the completed inverse coefficient propagator be
\[
 S_\Ph^{-1}(p)=\ii A(p^2)\slashed p+B(p^2).
\]
For \(S_\Ph=F S_0\), one has \(A=F^{-1}\) and
\(B=mF^{-1}\).
Define the symmetric average and divided differences by
\[
\begin{aligned}
 \overline A(p',p)&=\tfrac12[A(p'^2)+A(p^2)],\\
 A^{[1]}(p',p)&=\frac{A(p'^2)-A(p^2)}{p'^2-p^2},\\
 B^{[1]}(p',p)&=\frac{B(p'^2)-B(p^2)}{p'^2-p^2}.
\end{aligned}
\]
A longitudinal vertex of Ball--Chiu type \cite{BallChiu1980} is
\begin{equation}
\begin{aligned}
 \Gamma_L^\mu(p',p)
 ={}&\ii\overline A(p',p)\gamma^\mu\\
 &+(p'+p)^\mu\Bigl[
     \tfrac{\ii}{2}A^{[1]}(p',p)(\slashed p'+\slashed p)\\
 &\hspace{6.4em}+B^{[1]}(p',p)\Bigr].
\end{aligned}
 \label{eq:longitudinal-vertex}
\end{equation}
Contracting with \(q=p'-p\) gives exactly
\[
 q_\mu\Gamma_L^\mu
       =S_\Ph^{-1}(p')-S_\Ph^{-1}(p).
\]
When the functions are differentiable, equal-invariant difference
quotients are interpreted by their derivative limits.
An arbitrary transverse addition remains possible.

Equation~\eqref{eq:longitudinal-vertex} repairs this one Ward
identity. It does not specify the entire theory. Contact terms,
higher photon vertices, and the non-Abelian generalization must
come from the completed action, and their growth must be included
in the UV estimates. In particular, the inverse \(F^{-1}\) in
this vertex may counteract line suppression.

\begin{proposition}[An exact norm test for a Ward-compatible auxiliary vertex]
\label{prop:Ward-vertex-growth}
Let \(S_\Ph^{-1}(p)=F(p^2)^{-1}(\ii\slashed p+m)\),
with \(F>0\), \(F(0)>0\), \(m\geq0\), and Hermitian Euclidean
gamma matrices. Write \(A_R=F(R^2)^{-1}\), \(A_0=F(0)^{-1}\).
Any vertex satisfying \eqref{eq:dressed-WTI} obeys the exact identity
\begin{equation}
 \left\|n_\mu\Gamma_\Ph^\mu(Rn,0)\right\|^2
       =A_R^2+\frac{m^2}{R^2}(A_R-A_0)^2,
 \label{eq:Ward-vertex-exact-norm}
\end{equation}
for a unit vector \(n\) and \(R>0\). In particular,
\begin{equation}
 \left\|n_\mu\Gamma_\Ph^\mu(Rn,0)\right\|
 \geq\frac1R\left[
       \frac{\sqrt{R^2+m^2}}{F(R^2)}-\frac{m}{F(0)}
                      \right].
 \label{eq:Ward-vertex-growth-lower}
\end{equation}
The norm is the matrix operator norm. For a superpolynomially
decreasing \(F(R^2)\), no globally polynomially bounded vertex
can satisfy this identity in the stated auxiliary field coordinates.
\end{proposition}

\begin{proof}
Set \(p'=Rn\), \(p=0\) in the Ward identity and divide by
\(R\). The contraction is
\(\ii A_R\slashed n+m(A_R-A_0)\one/R\).
Since \(\slashed n\) is Hermitian and \((\slashed n)^2=\one\),
its adjoint times itself is
\([A_R^2+m^2(A_R-A_0)^2/R^2]\one\). This proves the exact
norm. The reverse triangle inequality also gives
\eqref{eq:Ward-vertex-growth-lower}. A transverse addition leaves
this contraction unchanged.
\end{proof}

For the Gaussian and stable-half kernels the lower bound grows
respectively like \(e^{tR^2}\) and \(e^{aR}\). This first test
uses a growing photon momentum \(q=Rn\); by itself it is not a
large-loop-momentum estimate at fixed external photon momentum.
The following diagonal test addresses that distinct regime.

\paragraph{Soft photon, large fermion momentum}
Suppose additionally that \(F\) is differentiable for positive
arguments and the vertex is continuous at \((p',p)=(p,p)\)
for each \(p\ne0\). Put \(A(u)=F(u)^{-1}\). Dividing the
Ward identity with \(p'=p+hn\) by \(h\), then taking
\(h\to0\) for every unit \(n\), fixes
\begin{equation}
\begin{aligned}
 \Gamma_\Ph^\mu(p,p)
 &=\frac{\partial S_\Ph^{-1}(p)}{\partial p_\mu}\\
 &=\ii A(p^2)\gamma^\mu
    +2p^\mu A'(p^2)(\ii\slashed p+m).
\end{aligned}
 \label{eq:soft-Ward-vertex}
\end{equation}
A continuous transverse addition has zero diagonal limit, by the
same directional argument. At \(p=Rn\), write \(A'_R=A'(R^2)\).
For any unit \(v\) orthogonal to \(n\), Hermiticity of the
Euclidean gamma matrices gives
\begin{equation}
\begin{aligned}
 \|n_\mu\Gamma_\Ph^\mu(p,p)\|^2
   &=(A_R+2R^2A'_R)^2+4m^2R^2(A'_R)^2,\\
 \|v_\mu\Gamma_\Ph^\mu(p,p)\|&=A_R.
\end{aligned}
 \label{eq:soft-Ward-norms}
\end{equation}
This proves nonpolynomial vertex growth for the same kernels with
photon momentum fixed at zero and fermion momentum unbounded,
provided the stated diagonal regularity holds. It is still a
vertex test, not a lower bound on a complete loop integral.
The full product of vertices and covariances must be estimated;
individual vertex growth does not prove graph divergence or exclude
cancellation by the required contact graphs. Cancellations or changed
field coordinates require their own estimates.
Proposition~\ref{prop:residual-coercivity} gives explicit sufficient
conditions when nonpolynomial vertex growth leaves a positive damping
margin. This obstruction to a particular sufficient hypothesis is not a proof
that every completed graph diverges. Together with
Proposition~\ref{prop:pure-redefinition-cancellation}, it gives two
explicit tests of the interaction completion instead of assuming
that source normalization or a single Ward identity supplies it.

\paragraph{The stable-half vertex at zero momentum}
The equal-invariant derivative prescription in
\eqref{eq:longitudinal-vertex} only applies where those derivatives
exist. For \(F(p^2)=e^{-a|p|}\), \(a>0\), and \(m>0\),
no Ward-compatible auxiliary vertex can be jointly continuous at
\((p',p)=(0,0)\). Indeed, for every unit \(n\), the Ward identity
forces
\begin{equation}
 \lim_{R\downarrow0}n_\mu\Gamma_\Ph^\mu(Rn,0)
       =\ii\slashed n+am\one.
 \label{eq:stable-half-vertex-origin-limit}
\end{equation}
A continuous vertex would make the left side
\(n_\mu\Gamma_\Ph^\mu(0,0)\), which is odd in \(n\).
The nonzero scalar term \(am\one\) is even in \(n\), a
contradiction. The same obstruction follows from the trace alone.
Thus the difference quotients do not silently supply a regular
zero-momentum completion for this massive cusp profile. This is a
bounded directional discontinuity test, not an assertion of an
infinite vertex norm: the squared contracted norm in
\eqref{eq:Ward-vertex-exact-norm} tends to \(1+a^2m^2\).
For \(m=0\) this particular continuity obstruction
is absent; no higher-regularity conclusion is inferred. These
statements concern the specified auxiliary inverse propagator, not
a separately defined physical current vertex.

The preceding equations are ordinary Euclidean coefficient
calculations. For the full carrier graph, let
\(S_\kappa\), \(\Gamma_\kappa^\mu\), and
\(\Pi_\kappa^{\mu\nu}\) be the complete character
representatives supplied consistently by the completed functional,
with all factors in each character using its complete branch data. Their reconstructions are
\begin{equation}
 \vec S=\sum_{\kappa=\pm1}\mathsf P_{\kappa}S_\kappa,
 \qquad
 \vec\Gamma^\mu=\sum_{\kappa=\pm1}\mathsf P_{\kappa}
                                  \Gamma_\kappa^\mu,
 \qquad
 \vec\Pi^{\mu\nu}=\sum_{\kappa=\pm1}\mathsf P_{\kappa}
                                  \Pi_\kappa^{\mu\nu}.
 \label{eq:QED-character-reconstruction}
\end{equation}
One character is used throughout each traced loop product.
If the inverse propagator exists in both characters, then
\(\vec S^{-1}=\sum_\kappa\mathsf P_{\kappa}S_\kappa^{-1}\).
Consequently, the pair of coefficient Ward identities
\begin{equation}
 q_\mu\Gamma_\kappa^\mu(p',p)
        =S_\kappa^{-1}(p')-S_\kappa^{-1}(p)
 \quad(\kappa=\pm1)
 \label{eq:QED-character-Ward-identities}
\end{equation}
reconstructs exactly the carrier-algebra identity
\(q_\mu\vec\Gamma^\mu
=\vec S^{-1}(p')-\vec S^{-1}(p)\).
Likewise, transversality of the completed tensor must hold in
each character. A scalar coefficient check with omitted running
line signs does not establish this identity for the unrestricted
Lorentzian graph.

For the full completed tensor, write
\begin{equation}
 \Pi_{\rm full}^{\mu\nu}
      =\Pi_{\rm bv}^{\mu\nu}+\Pi_{\rm comp}^{\mu\nu}.
 \label{eq:Pi-complete}
\end{equation}
Here \(\Pi_{\rm comp}\) includes the changes in vertices and all
additional graphs required by the completed functional.
Its Ward identity gives
\(Q_\mu\Pi_{\rm full}^{\mu\nu}=0\).
For \(Q\ne0\), if the completed plane also has the stated
Euclidean covariance and no additional background tensors, then
\begin{equation}
 \Pi_{\rm full}^{\mu\nu}(Q)
    =(Q^2\delta^{\mu\nu}-Q^\mu Q^\nu)
                       \Pi_{\rm full}(Q^2).
 \label{eq:Pi-transverse}
\end{equation}
The scalar function in this equation cannot be computed from
\eqref{eq:bare-QED-loop} alone. The decomposition is not a proof of
regularity at \(Q=0\): a transverse scalar form factor may have a
singularity there. If the tensor is continuous at zero, transversality
forces \(\Pi_{\rm full}^{\mu\nu}(0)=0\). If it is analytic
there, the scalar coefficient has the corresponding regular extension.
These properties require the completed graph sum and its limits.

\subsubsection*{Parameter and cut bookkeeping}

In the two-line Schwinger representation, define
\[
 w_i=\alpha_i+a^2\tau_i,\qquad
 W=w_1+w_2,\qquad x=\frac{w_2}{W}.
\]
The Gaussian exponent after translating the loop variable is
\[
 -m^2(\alpha_1+\alpha_2)-x(1-x)WQ^2.
\]
If \(\alpha_1,\alpha_2,\tau_1,\tau_2\) are still independent
integration variables, \(x\) is this function of them, not an
additional independent Feynman parameter. Changing variables to
\((W,x)\) requires the corresponding Jacobian and lower limits.
The tensor numerator and completion terms must be retained through
that change of variables.

As a reference for a conventional unit-vertex cut with physical
normalization, ordinary QED gives
\begin{equation}
 \Im\Pi^{\rm coeff}_{{\rm std},+}(\sMand)
  =\frac{e^2}{12\pi}
    \left(1+\frac{2m^2}{\sMand}\right)
       \sqrt{1-\frac{4m^2}{\sMand}}\,
       \theta(\sMand-4m^2).
 \label{eq:QED-reference-spectrum}
\end{equation}
Let \(\Pi_{\rm std}^{\rm sheet}\) denote the ordinary
Lorentzian current-correlation scalar in the stated normalization,
and set
\begin{equation}
\begin{aligned}
 \Pi^{\rm dir}_{{\rm std},\eta,\bb}(\sMand)
  &:=\lim_{\varepsilon\downarrow0}
    (\embed\circ\Pi^{\rm sheet}_{\rm std}
                   \circ\embed^{-1})(\mathsf S_{\eta,\bb}),\\
 \Pi^{\rm coeff}_{{\rm std},\eta}
       &:=\embed^{-1}(\Pi^{\rm dir}_{{\rm std},\eta,\bb}).
\end{aligned}
 \label{eq:QED-directed-spectrum-interface}
\end{equation}
For the reflected ordinary reference,
\begin{equation}
 \Disc^{\rm dir}_\eta\Pi_{\rm std}
 =2\eta\je\,
                  \Im\Pi^{\rm coeff}_{{\rm std},+}(\sMand).
 \label{eq:QED-directed-reference-jump}
\end{equation}
The square root in the spectral density is the nonnegative real
phase-space root. This reference scalar omits the full graph's
Dyson factors; the completed polarization and scattering graphs
use \eqref{eq:QED-character-reconstruction}.

This is a normalization benchmark for the ordinary current
matrix element. The completed damped theory has its own
on-shell subamplitudes, residues, and possible kernel branch
contributions. Equality with \eqref{eq:QED-reference-spectrum}
requires their matching; it does not follow from deleting
\(\Ph\) on the cut lines. Resummation of
\(\Pi_{\rm full}\) likewise requires analyzing its actual zeros
and singularities rather than presuming an unchanged pole set.

\section{The standalone raw O--QED action and its fixed-depth tests}
\label{sec:OQED-application}

In the standalone construction of Ref.~\cite{ChiapponeIgnition}
one writes the support-diagonal local matter, Maxwell and gauge-fixing
terms directly. Equation~\eqref{eq:OQED-source-action} gives this raw
reference action. The additional prescription in
\eqref{eq:OQED-normalized-kernel} assigns a positive-depth multiplier
to selected virtual lines. The following calculations test that
prescription with each indicated choice of vertices, contacts and
measure. Section~\ref{sec:pe-actions} defines a second electromagnetic
theory by retaining the photon and charged sources in a GCM quantum
pushforward. Matching a quadratic covariance is only the first of
the source identities needed to carry results between these constructions.

\subsection{Source action, normalization, and Abelian BRST}
\label{sec:OQED-action}

Fix one support $\bb=(s,b_s)$, one carrier, and one ordering branch.
All fields and independent variational slots in the reference density
are on that copy. Let $e_{\rm src}$ denote its unrescaled charge and
write
\begin{equation}
\begin{aligned}
 D_{\mu,\bb}&=\partial_\mu+\ii e_{\rm src}A_{\mu,\bb},\\
 \bar D_{\mu,\bb}&=\partial_\mu-\ii e_{\rm src}A_{\mu,\bb},\\
 F_{\mu\nu,\bb}&=\partial_\mu A_{\nu,\bb}-\partial_\nu A_{\mu,\bb},\\
 \mathcal L_{{\rm mat},\bb}
  =\tfrac12\bigl[&\bar\Psi_{\bb}(\ii\gamma^\mu D_{\mu,\bb}-m_e)\Psi_{\bb}\\
  &+\bar\Psi_{\bb}^{\,*}
       (-\ii\gamma^{\mu *}\bar D_{\mu,\bb}-m_e)\Psi_{\bb}^*\bigr].
\end{aligned}
\label{eq:OQED-source-matter}
\end{equation}
The gamma matrices refer to the fixed intrinsic spin frame. The starred
slots are the literal componentwise conjugates of the corresponding
unstarred slots; $\bar\Psi_{\bb}^{\,*}=(\bar\Psi_{\bb})^*$ is not a
silently substituted Dirac adjoint of $\Psi_{\bb}^*$ in a different
representation. Independent variation precedes an admitted conjugation
restriction. No charge-conjugate spinor is substituted for these data.

With $\bb_N=(N,b_N)$ and $\bb_A=(A,b_A)$, the full gauge-fixed coefficient
and its directed intrinsic assembly are
\begin{equation}
\begin{aligned}
 \mathcal L_{{\rm ref},\bb}
  ={}&\mathcal L_{{\rm mat},\bb}
       -\tfrac14F_{\mu\nu,\bb}F^{\mu\nu}_{\bb}\\
    &-\frac{(\partial_\mu A^\mu_{\bb})^2}{2\alpha_{\rm gf}}
                  +\bar c_{\bb}(-\partial^2)c_{\bb},\\
 \vec{\mathcal L}^{(4)}_{{\rm ref},\epsilon,\bb}
  &:=\sigma_s\sigma_\epsilon\hat p_{0,\bb}
                       Z_{\rm GT}\mathcal L_{{\rm ref},\bb},\\
 \vec S_{{\rm ref},\epsilon}^{\rm lock}
  &=\sum_{s\in\{N,A\}}\int\dd\mu_{\bb_s}\,
                    \vec{\mathcal L}^{(4)}_{{\rm ref},\epsilon,\bb_s}.
\end{aligned}
\label{eq:OQED-source-action}
\end{equation}
Its ambient form is \eqref{eq:fw-ambient-directed-action}, with one
$\Delta_{{\rm GTL},\bb_s}$ at each event before pullback. At a common
carrier and $\sigma_\epsilon=+1$ this gives the source's
$\hat p_0\sum_s\sigma_s\int\dd\mu_{\bb_s}Z_{\rm GT}\mathcal L_{{\rm ref},\bb_s}$.
The measure contains no orientation sign. The free Abelian ghost can be
omitted from a decoupled physical-amplitude display, not from the stated
BRST functional. The gauge parameter is nonzero; Landau gauge is a
controlled limit, not substitution of zero into its denominator.

This standalone O--QED sector fold is not the HOUQFT metric-copy fold.
Nor does its two-sector sum introduce a mixing vertex: every displayed
interaction is support diagonal. The ultraviolet estimates below are
fibrewise and do not cancel divergent coefficients by subtracting N from A.
The photon and ghost blocks occur once per support, not once per fermion
species.

The locked normalization $Z_{\rm GT}=\ell_{\rm GT}^{K-1}$ is not the
virtual depth or a plasma response length. The canonical coefficient
variables used in the loop tests are
\begin{equation}
\begin{aligned}
 \Psi_{c,\bb}&=Z_{\rm GT}^{1/2}\Psi_{\bb},&
 A_{c,\mu,\bb}&=Z_{\rm GT}^{1/2}A_{\mu,\bb},\\
 e&=Z_{\rm GT}^{-1/2}e_{\rm src},&
 J_{c,\bb}&=Z_{\rm GT}^{-1/2}J_{{\rm src},\bb}.
\end{aligned}
\label{eq:fw-oqed-rescaling}
\end{equation}
The source relation assumes no extra $Z_{\rm GT}$ in the pairing:
$J_{c,\bb}\Psi_{c,\bb}=J_{{\rm src},\bb}\Psi_{\bb}$.
For any other convention, transport the complete pairing and its actual
prefactors. Ghosts, antighosts and auxiliary fields scale by
$Z_{\rm GT}^{1/2}$ in the corresponding quadratic block; the mass does
not change. The later coefficient loop formulas suppress $c$ and use
$e$ for the matched charge.

The doubled conjugation display does not double a physical Dirac loop.
Species multiplicities and integration slots belong to the declared
measure. For $N_f$ diagonal species the answer is a sum of closed species
traces, each with its own charge, mass, kernel and contacts. It is not a
factor $N_f$ at each vertex.

Before gauge fixing,
$A_{\mu,\bb}\mapsto A_{\mu,\bb}-\partial_\mu\alpha_{\bb}$ and
$\Psi_{\bb}\mapsto e^{\ii e_{\rm src}\alpha_{\bb}}\Psi_{\bb}$.
A compatible left odd differential is
\begin{equation}
\begin{aligned}
 \BR A_{\mu,\bb}&=-\partial_\mu c_{\bb},& \BR c_{\bb}&=0,\\
 \BR\Psi_{\bb}&=\ii e_{\rm src}c_{\bb}\Psi_{\bb},&
 \BR\bar\Psi_{\bb}&=\ii e_{\rm src}\bar\Psi_{\bb}c_{\bb},\\
 \BR\Psi_{\bb}^*&=-\ii e_{\rm src}c_{\bb}\Psi_{\bb}^*,&
 \BR\bar\Psi_{\bb}^{\,*}&=-\ii e_{\rm src}\bar\Psi_{\bb}^{\,*}c_{\bb},\\
 \BR\bar c_{\bb}&=B_{\bb},& \BR B_{\bb}&=0.
\end{aligned}
\label{eq:OQED-Abelian-BRST}
\end{equation}
The right-ordered ghost in the barred slot follows from the graded
Leibniz rule. It gives $\BR(\bar\Psi_{\bb}\Psi_{\bb})=0$,
$\BR(D_{\mu,\bb}\Psi_{\bb})=\ii e_{\rm src}c_{\bb}D_{\mu,\bb}\Psi_{\bb}$,
$\BR F_{\mu\nu,\bb}=0$ and $\BR^2=0$.
For example, in canonical Euclidean variables take the field-independent
operator $H$, $\zeta>0$, and $\mathcal F_{\bb}=-\partial\cdot A_{\bb}$:
\begin{equation}
\begin{aligned}
 \Psi_{{\rm gf},\bb}
  &=\int\dd\mu_{E,\bb}\,\bar c_{\bb}H
                         (\mathcal F_{\bb}-\zeta B_{\bb}/2),\\
 \BR\Psi_{{\rm gf},\bb}
  &=\int\dd\mu_{E,\bb}\,[B_{\bb}H\mathcal F_{\bb}
          -\tfrac\zeta2B_{\bb}HB_{\bb}
          +\bar c_{\bb}H(-\partial_E^2)c_{\bb}].
\end{aligned}
\label{eq:fw-oqed-gaugefermion}
\end{equation}
The admitted auxiliary-field contour, not a positive real $B$ Gaussian,
yields positive Euclidean gauge fixing. Abelian ghosts remain free in
this linear gauge. These statements verify the reference differential;
they do not construct a quantum anomaly-free nonlocal master action.
The trial below specifies its own covariant vertices and contacts, whose
growth is tested separately.

\subsubsection*{A covariant trial that generates the longitudinal vertex and contacts}
\label{sec:polish-covariant-trial}
A concrete coefficient action makes the relation between a dressed line,
its Ward vertex and its contact terms testable. It is a specified
comparison, not a derivation of the electromagnetic channel of A4.
Use one canonically normalized support, $a,M,m_e>0$. In the following
operator calculation $\mathcal A=\mathcal A_{\bb}$ and
$\mathcal D=\mathcal D_{\bb}$; all domains and source slots are on
that fixed copy. Define
\begin{equation}
 \begin{aligned}
 D_\mu[\mathcal A]&=\partial_\mu+\ii e\mathcal A_\mu,
       & L_{\mathcal A}&=-D_\mu[\mathcal A]D_\mu[\mathcal A],\\
 D_0[\mathcal A]&=\gamma_E^\mu D_\mu[\mathcal A]+m_e,
       &f(u)&=e^{a(\sqrt{u+M^2}-M)},\\
 \mathcal D_{\rm sym}[\mathcal A]
       &=\tfrac12\{f(L_{\mathcal A}),D_0[\mathcal A]\}.
 \end{aligned}
 \label{eq:polish-covariant-dirac}
\end{equation}
The braces denote an anticommutator. The matter coefficient action is
$\int\dd\mu_{E,\bb}\,\bar\psi_{\bb}\mathcal D_{{\rm sym},\bb}[\mathcal A_{\bb}]\psi_{\bb}$,
with the independent source and conjugation slots specified as above.
It may be combined with the reference Maxwell, gauge-fixing and ghost
coefficients for this trial. These remain scalar Euclidean coefficients;
returning to the directed theory retains the action unit, source maps,
$\widehat u_{k,\bb}$, and $\ii\eta\epsvec$ of the fixed copy.

Work first with finite spectral matrices, or on a common invariant
continuum domain on which every displayed composition and variation
exists. The unbounded exponential $f(L_{\mathcal A})$ is not defined on
all smooth fields merely by writing a formal series. In a continuum
claim the domain and its gauge transport must be exhibited. Under an
admitted multiplication $g=e^{\ii e\alpha}$, with
$\mathcal A' =\mathcal A-\dd\alpha$, spectral covariance gives
$L_{\mathcal A'}=gL_{\mathcal A}g^{-1}$ and
$\mathcal D_{\rm sym}[\mathcal A']=g\mathcal D_{\rm sym}[\mathcal A]g^{-1}$.
An arbitrary finite-mode cutoff need not preserve these gauge maps;
the finite matrices establish the variation identities, not a gauge
invariant continuum measure.

At zero background, the two free factors commute and
$\mathcal D_{\rm sym}[0](k)=f(k^2)(\ii\slashed k+m_e)$.
Its inverse is exactly $F_{a,M}(k^2)S_0(k)$.
The first variation, with the charge removed from the vertex, is
\begin{equation}
 \begin{aligned}
 \Gamma_{\rm sym}^\mu(k',k)
  ={}&\frac{f(k'^2)+f(k^2)}2\,\ii\gamma_E^\mu\\
  &+(k'+k)^\mu f^{[1]}(k'^2,k^2)
       \left[\frac\ii2(\slashed k'+\slashed k)+m_e\right].
 \end{aligned}
 \label{eq:polish-trial-vertex}
\end{equation}
Here $f^{[1]}(u,v)=(f(u)-f(v))/(u-v)$ with $f^{[1]}(u,u)=f'(u)$.
Indeed $e^{-1}\delta L=(k'+k)\cdot\delta\mathcal A$ and
$e^{-1}\delta D_0=\ii\gamma_E\cdot\delta\mathcal A$ in the free
Fourier basis. The matrix derivative of $f(L)$ is its divided difference
multiplying $\delta L$. Substitution into the anticommutator proves
\eqref{eq:polish-trial-vertex}. With $P=k'-k$, contraction gives
\[
 P_\mu\Gamma_{\rm sym}^\mu(k',k)
 =f(k'^2)(\ii\slashed k'+m_e)-f(k^2)(\ii\slashed k+m_e),
\]
so this is the Ball--Chiu longitudinal choice
\eqref{eq:longitudinal-vertex}, not a separately fitted vertex.

The second insertion is fixed by the same action. For independent
background variations $\delta_1,\delta_2$, put $R_z=(z-L_{\mathcal A})^{-1}$
and let $\mathscr C$ surround the finite spectrum inside the analytic
domain of $f$. Then
\begin{equation}
 \begin{aligned}
 \delta_1\delta_2 f(L_{\mathcal A})
  =\frac1{2\pi\ii}\oint_{\mathscr C} f(z)\bigl[&R_z\delta_1\delta_2L\,R_z\\
    &+R_z\delta_1L\,R_z\delta_2L\,R_z\\
    &+R_z\delta_2L\,R_z\delta_1L\,R_z\bigr]\dd z,\\
 \delta_1\delta_2L
    &=2e^2\delta_1\mathcal A_\mu\delta_2\mathcal A_\mu.
 \end{aligned}
 \label{eq:polish-trial-second-variation}
\end{equation}
Consequently
\begin{equation}
 \begin{aligned}
 \delta_1\delta_2\mathcal D_{\rm sym}
  ={}&\tfrac12\{\delta_1\delta_2 f,D_0\}
    +\tfrac12\{\delta_1 f,\delta_2D_0\}\\
    &+\tfrac12\{\delta_2 f,\delta_1D_0\},
       \qquad\delta_1\delta_2D_0=0.
 \end{aligned}
 \label{eq:polish-trial-contact}
\end{equation}
Differentiating $(z-L)^{-1}$ gives $R_z(\delta L)R_z$ and hence
\eqref{eq:polish-trial-second-variation}. In particular, omitting
$\delta_1\delta_2L$ would omit an action-generated contact.

For a constant background on the noncompact flat chart, the free
Fourier symbol is shifted by $k\mapsto k+e\mathcal A$. Thus this trial
has exactly the constant-background bubble/contact coefficient of
Proposition~\ref{prop:OQED-constant-contact}. Its nonintegrable trace
therefore applies to this particular covariant trial as well. Classical
covariance, a correct Ward contraction and a damped free inverse do not
establish complete-graph UV convergence or physical unitarity. This
construction supplies a concrete test action; it does not change the
conditional scope of Theorem~\ref{thm:OQED-fixed-graph} or assert a
normalizable continuum fermion determinant. At $M=0$ the zero-spectral
cusp requires separate domain and differentiability analysis.

\subsection{The specified kernel passes every finite inverse-moment test}
\label{sec:OQED-moments}

The following formulas refer to the assigned fixed depth of the
O--QED test theory. A selector average defines a different source
functional. Its relevant finite-graph bounds are
\eqref{eq:update-joint-moment-budget}--\eqref{eq:update-depth-budget},
not an independent substitution of one marginal kernel on every line.


Use the following dictionary, without equating independent scales:
\[
 a:=a_{\rm ker}^{(p)}>0,\qquad M:=M_{\rm ker}\geq0,\qquad
 \ell_\Omega=\hbar a,\qquad \mu:=aM.
\]
The electron pole parameter \(m_e\) remains in its denominator;
\(M\) is the auxiliary kernel scale. With a dimensionful nonnegative
kinetic variable \(u\), the source's normalized representative is
\begin{equation}
\begin{aligned}
 F_{a,M}(u)&:=e^{-a h_M(u)},&
 h_M(u)&:=\sqrt{u+M^2}-M,\\
 \Ph_\mu(z)&:=e^{\mu-\sqrt{z+\mu^2}},&
 F_{a,M}(u)&=\Ph_\mu(a^2u).
\end{aligned}
\label{eq:OQED-normalized-kernel}
\end{equation}
In particular \(F_{a,M}(0)=1\). For \(M>0\) this is not the same
profile as either \(e^{-a\sqrt u}\) or the unnormalized
\(e^{-a\sqrt{u+M^2}}\).

\begin{proposition}[Explicit O--QED kernel certificate]
\label{prop:OQED-moment-certificate}
The dimensionless Laplace measure of \(\Ph_\mu\) is the probability
measure
\begin{equation}
 \rho_{\mu,\bb}(\dd\tau)=\frac{e^\mu}{2\sqrt\pi}\,
 \tau^{-3/2}\exp\!\left(-\mu^2\tau-\frac1{4\tau}\right)\dd\tau.
\label{eq:OQED-dimensionless-measure}
\end{equation}
It has no atom at zero and no positive support gap, including when
\(\mu>0\). Every inverse moment is finite. For \(r>0\),
\begin{equation}
 M_r[\rho_{\mu,\bb}]=
 \begin{cases}
 \displaystyle\frac{e^\mu}{\sqrt\pi}(2\mu)^{r+1/2}
                      K_{r+1/2}(\mu),&\mu>0,\\[2pt]
 \displaystyle\frac{4^r\Gamma(r+1/2)}{\sqrt\pi},&\mu=0.
 \end{cases}
\label{eq:OQED-inverse-moments}
\end{equation}
Here \(K_\nu\) is the modified Bessel function, not the number of GT
directions. In particular,
\begin{equation}
\begin{aligned}
 M_1&=2(1+\mu),\\
 M_2&=4(\mu^2+3\mu+3),\\
 M_3&=8(\mu^3+6\mu^2+15\mu+15).
\end{aligned}
\label{eq:OQED-first-moments}
\end{equation}
For \(\mu>0\) its positive exponential moments exist precisely for
\(0<b\leq\mu^2\), with
\begin{equation}
 \int e^{b\tau}\rho_{\mu,\bb}(\dd\tau)
       =\exp\!\left(\mu-\sqrt{\mu^2-b}\right).
\label{eq:OQED-exponential-moment}
\end{equation}
There is no positive exponential moment when \(\mu=0\).
\end{proposition}

\begin{proof}
Rescale the dimensionful heat parameter as $\lambda=a^2\tau$ in the
positive subordination law for \eqref{eq:OQED-normalized-kernel}.
The stable-half integral identity gives
\(\int e^{-z\tau}\rho_{\mu,\bb}(\dd\tau)=\Ph_\mu(z)\), and evaluation
at zero proves normalization. The density is strictly positive on
\((0,\infty)\); tempering its large-\(\tau\) tail creates no gap
at small \(\tau\). The factor \(e^{-1/(4\tau)}\) makes every inverse
moment integrable at zero, and the finite measure tail controls infinity.
Using the standard integral representation of \(K_\nu\)
\cite[Eq.~10.32.10]{DLMF10}, or its elementary change of variable,
\[
 \int_0^\infty \tau^{-r-3/2}e^{-\mu^2\tau-1/(4\tau)}\dd\tau
       =2(2\mu)^{r+1/2}K_{r+1/2}(\mu).
\]
For \(\mu=0\), substitute \(v=1/(4\tau)\). The three integer
moments follow by half-integer evaluation, or by integrating
\(z^{r-1}\Ph_\mu(z)/\Gamma(r)\) after setting
\(y=\sqrt{z+\mu^2}-\mu\). Finally, replacing \(\mu^2\) by
\(\mu^2-b\) in the density proves
\eqref{eq:OQED-exponential-moment}. At \(b=\mu^2\) the tail is
\(\tau^{-3/2}\), still integrable; beyond it the integral diverges.
At that endpoint its first differentiated moment diverges, so finiteness
there does not extend holomorphy across the branch point.
\end{proof}

The dimensionful inverse moment is
\begin{equation}
 \int\lambda^{-r}\rho^{\rm heat}_{a,M,\bb}(\dd\lambda)
       =a^{-2r}M_r[\rho_{aM,\bb}].
\label{eq:OQED-dimensional-moment}
\end{equation}
Thus the actual condition \(M_{1+D/2}<\infty\) is verified for every
finite numerator degree, without inferring it from normalization or
from absence of an atom. For \(a\geq a_->0\) and \(aM\leq\mu_+\), the density also gives
\[
 a^{-2r}M_r[\rho_{aM,\bb}]
 \leq a_-^{-2r}e^{\mu_+}\frac{4^r\Gamma(r+1/2)}{\sqrt\pi}.
\]
This supplies explicit uniform kernel constants on the corresponding
compact parameter sets. No uniform damping-removal limit follows.
A useful separate bound is
\begin{equation}
 e^{-a\sqrt u}\leq F_{a,M}(u)
       \leq\min\{1,e^{aM}e^{-a\sqrt u}\},\qquad u\geq0.
\label{eq:OQED-linear-exponential-bounds}
\end{equation}
It follows from \(\sqrt u-M\leq h_M(u)\leq\sqrt u\).
Increasing \(M\) at fixed \(a,u>0\) weakens the suppression; it does
not improve the large-momentum exponential rate.


\subsection{Positive-depth stability and singular depth limits}
\label{sec:OQED-depth-limits}

The UV theorem holds at fixed positive depth. The next results address
three distinct questions: controlled differentiation at positive depth,
large-depth localization, and the singular removal of damping. Each
concerns the positive Euclidean scalar line or bubble comparison; a
completed polarization tensor requires its own calculation.

\begin{proposition}[Depth derivatives and one-line asymptotics]
\label{prop:OQED-depth-stability}
For $F_{a,M}$ of \eqref{eq:OQED-normalized-kernel} and every
integer $n\geq0$,
\begin{equation}
 \partial_a^nF_{a,M}(u)=(-h_M(u))^nF_{a,M}(u),\qquad
 h_M(u)=\sqrt{u+M^2}-M.
 \label{eq:OQED-depth-derivatives}
\end{equation}
For $a\in[a_-,a_+]$, $0<a_-\leq a_+$, and $0\leq M\leq M_+$,
\begin{equation}
 |\partial_a^nF_{a,M}(q^2)|
 \leq e^{a_+M_+}|q|^n e^{-a_-|q|}.
 \label{eq:OQED-depth-uniform-envelope}
\end{equation}
In particular, such derivatives preserve an established complete
graph bound when the additional polynomial numerator degree and
infrared factors are controlled. For $u>0$, $\partial_MF_{a,M}(u)>0$, with the right derivative at
$M=0$.

Define a Fourier-normalized one-line comparison for $d\geq0$:
\begin{equation}
 \mathscr Q_d(a,M,m)
 :=\int_{\RR^4}\frac{\dd^4q_{\bb}}{(2\pi)^4}
   \frac{|q|^dF_{a,M}(q^2)}{q^2+m^2},
 \quad m\geq0.
 \label{eq:OQED-one-line-depth}
\end{equation}
For every fixed $m,M\geq0$, it is finite at $a>0$, and
\begin{equation}
 \lim_{a\downarrow0}a^{d+2}\mathscr Q_d(a,M,m)
 =\frac{\Gamma(d+2)}{8\pi^2}.
 \label{eq:OQED-zero-depth-one-line}
\end{equation}
The convergence is uniform for $(m,M)$ in bounded nonnegative sets.
For fixed $m>0$, the two different strong-depth limits are
\begin{equation}
 \begin{aligned}
 \lim_{a\to\infty}a^{d+4}\mathscr Q_d(a,0,m)
   &=\frac{\Gamma(d+4)}{8\pi^2m^2},\\
 \lim_{a\to\infty}a^{(d+4)/2}\mathscr Q_d(a,M,m)
   &=\frac{(2M)^{(d+4)/2}\Gamma((d+4)/2)}
           {16\pi^2m^2},\quad M>0.
 \end{aligned}
 \label{eq:OQED-strong-depth-one-line}
\end{equation}
These asymptotics hold in their stated regimes, not uniformly over
arbitrary joint limits of $a$ and $M$.
\end{proposition}
\begin{proof}
The derivatives follow by differentiating the exponential at fixed
$M$. The inequalities
$0\leq h_M(q^2)\leq|q|$ and
$F_{a,M}(q^2)\leq e^{a_+M_+}e^{-a_-|q|}$ prove the uniform bound.
Also $\partial_Mh_M(u)=M/\sqrt{u+M^2}-1<0$ for $u>0$.

Four-dimensional radial integration gives
\[
 \mathscr Q_d=\frac1{8\pi^2}\int_0^\infty
 \frac{r^{d+3}e^{-a(\sqrt{r^2+M^2}-M)}}{r^2+m^2}\dd r.
\]
For the zero-depth limit, set $r=v/a$ and multiply by $a^{d+2}$.
The resulting denominator is $v^2+a^2m^2$ and the exponent is
$-\sqrt{v^2+a^2M^2}+aM$. It converges pointwise to $-v$, with
an integrable common upper bound proportional to
$v^{d+1}e^{-v}$ for $a$ sufficiently small and bounded masses.
Dominated convergence gives \eqref{eq:OQED-zero-depth-one-line},
including $m=0$. For $M=0$ and $a\to\infty$, the same $r=v/a$
substitution, now multiplied by $a^{d+4}$, leaves the limiting
massive denominator $m^2$ and the moment $\Gamma(d+4)$.
For fixed $M>0$, set $r=v/\sqrt a$; then
$a(\sqrt{M^2+v^2/a}-M)\to v^2/(2M)$.
The radial Gaussian integral gives the second limit.
For the last limit, the rescaled exponent is
$v^2/[\sqrt{M^2+v^2/a}+M]$. For all $a\geq1$ it is
at least $v^2/(2M+v)$, which gives an integrable dominator
proportional to $v^{d+3}e^{-v^2/(2M+v)}$ after bounding the
massive denominator by $m^{-2}$. This also justifies the earlier
changes of limit.
\end{proof}


\begin{corollary}[An explicit O--QED exponential momentum tail]
\label{cor:OQED-exponential-line-tail}
For $a>0$, $M,m\geq0$, $d\geq0$ and $R>0$, the
Fourier-normalized line integral of
\eqref{eq:OQED-one-line-depth} obeys
\begin{equation}
 \begin{aligned}
 &\int_{|q|>R}\frac{\dd^4q_{\bb}}{(2\pi)^4}
     \frac{|q|^dF_{a,M}(q^2)}{q^2+m^2}\\
 &\qquad\leq \frac{e^{aM}}{8\pi^2a^{d+2}}
                     \Gamma(d+2,aR),
 \end{aligned}
 \label{eq:OQED-exponential-line-tail}
\end{equation}
where $\Gamma(s,x)=\int_x^\infty
 t^{s-1}e^{-t}\dd t$ is the upper incomplete gamma
function. For each fixed positive depth the right side is
exponentially small as $R\to\infty$. Selected-line graph tails
inherit finite products and union bounds only when the
\emph{complete} anisotropic numerator and unselected-line
estimate in Theorem~\ref{thm:optimized-selected-lines}
has been verified. These are Euclidean coefficient bounds,
not contour-arc or physical-cut estimates.
\end{corollary}
\begin{proof}
Use $F_{a,M}(r^2)\leq e^{aM}e^{-ar}$ and
$(r^2+m^2)^{-1}\leq r^{-2}$ for $r>0$. Four-dimensional
radial integration bounds the tail by
$e^{aM}(8\pi^2)^{-1}\int_R^\infty r^{d+1}e^{-ar}\dd r$,
which is the stated incomplete-gamma expression.
\end{proof}

For every $M\geq0$, the literal two-line scalar comparison
used below is already defined for real Euclidean $P$ by
\begin{equation}
 B_{a,M,m}(P):=\int\frac{\dd^4q_{\bb}}{(2\pi)^4}
 \frac{F_{a,M}(q^2)F_{a,M}((q+P)^2)}
 {(q^2+m^2)((q+P)^2+m^2)},\quad m>0.
 \label{eq:OQED-literal-depth-bubble}
\end{equation}
This real Euclidean definition also covers $M=0$; the stronger
holomorphic-tube statement below assumes $M>0$.


\begin{proposition}[Depth monotonicity of the literal scalar bubble]
\label{prop:OQED-bubble-depth-monotone}
Fix $m>0$, $M\geq0$, and real Euclidean $P$. Put
$H_{M,P}(q):=h_M(q^2)+h_M((q+P)^2)$. For each $a>0$,
$a\mapsto B_{a,M,m}(P)$ is completely monotone:
\begin{equation}
 \begin{aligned}
 &(-1)^n\partial_a^n B_{a,M,m}(P)\\
 &\quad=\int\frac{\dd^4q_{\bb}}{(2\pi)^4}
 \frac{H_{M,P}(q)^n e^{-aH_{M,P}(q)}}
 {(q^2+m^2)((q+P)^2+m^2)}>0
 \end{aligned}
 \label{eq:OQED-bubble-depth-derivatives}
\end{equation}
for every integer $n\geq0$. For $a\geq a_->0$,
$0\leq M\leq M_+$, and any real $P$,
\begin{equation}
 |\partial_a^nB_{a,M,m}(P)|\leq
 c_{n,a_-}B_{a_-/2,M_+,m}(0),
 \label{eq:OQED-bubble-depth-uniform}
\end{equation}
where $c_{0,a_-}=1$ and
$c_{n,a_-}=(2n/(\mathrm e\,a_-))^n$ for $n\geq1$.
The bubble is holomorphic in the \emph{depth} $a$ on
$\Re a>0$ at fixed real $P$. For fixed $a>0$ it is
nondecreasing in the auxiliary mass $M$ (strictly for
$P$ in the stated real domain). None of these statements
implies analyticity in complex external momentum when $M=0$,
or reflection positivity of a physical two-point function.
\end{proposition}
\begin{proof}
The two line multipliers equal $e^{-aH_{M,P}}$ with
$H_{M,P}\geq0$ and $H_{M,P}>0$ almost everywhere. On
$\Re a\geq a_->0$, derivatives are dominated because
$|H^n e^{-aH}|=H^n e^{-(\Re a)H}
\leq c_{n,a_-}e^{-a_-H/2}$.
The resulting positive integrand is integrable; its integral
is bounded by $c_{n,a_-}B_{a_-/2,M,m}(P)$.
Cauchy--Schwarz for the real even line covariance gives
$B_{a_-/2,M,m}(P)\leq B_{a_-/2,M,m}(0)$.
Since $\partial_Mh_M(u)<0$ for $u>0$, the latter is at
most $B_{a_-/2,M_+,m}(0)$. The same majorant proves
holomorphy in the stated depth half-plane. The depth derivatives have the alternating signs in
\eqref{eq:OQED-bubble-depth-derivatives}. Strict increase in $M$
follows by integrating the strict pointwise inequality away from
the zero momenta.
\end{proof}

\begin{theorem}[Controlled crossover for $aM$ fixed]
\label{thm:OQED-strong-depth-crossover}
Fix $m>0$ and $d\geq0$, and consider $M=\mu/a$ with fixed
$\mu\geq0$ as $a\to\infty$. Define
\begin{equation}
 \mathcal H_d(\mu):=\int_0^\infty v^{d+3}
                  e^{-\sqrt{v^2+\mu^2}+\mu}\dd v.
 \label{eq:OQED-crossover-H}
\end{equation}
Then, uniformly for $\mu$ in compact subsets of $[0,\infty)$,
\begin{equation}
 \begin{aligned}
 \lim_{a\to\infty}a^{d+4}\mathscr Q_d(a,\mu/a,m)
    &=\frac{\mathcal H_d(\mu)}{8\pi^2m^2},\\
 \lim_{a\to\infty}a^4B_{a,\mu/a,m}(0)
    &=\frac{3+6\mu+4\mu^2}{64\pi^2m^4}.
 \end{aligned}
 \label{eq:OQED-strong-depth-crossover}
\end{equation}
Here $B_{a,M,m}$ is the literal real Euclidean comparison
\eqref{eq:OQED-literal-depth-bubble}. Furthermore $\mathcal H_0(\mu)=6+6\mu+2\mu^2$,
$\mathcal H_d(0)=\Gamma(d+4)$, and
\begin{equation}
 \mathcal H_d(\mu)\sim\tfrac12(2\mu)^{(d+4)/2}
          \Gamma((d+4)/2)\qquad(\mu\to\infty).
 \label{eq:OQED-large-mu-crossover}
\end{equation}
The last limit concerns the crossover coefficient; it is not a
uniform joint $a,M$ asymptotic over unbounded $aM$.
\end{theorem}
\begin{proof}
Set $r=v/a$ and $M=\mu/a$. The exponent becomes
$\sqrt{v^2+\mu^2}-\mu$, and the denominator tends to $m^2$,
respectively $m^4$. For $\mu$ bounded above by $\mu_+$,
$e^{\mu-\sqrt{v^2+\mu^2}}\leq e^{\mu_+}e^{-v}$,
which supplies a common integrable radial majorant.
The limiting bubble integral is
\[
 \frac1{8\pi^2m^4}\int_0^\infty v^3
  e^{-2(\sqrt{v^2+\mu^2}-\mu)}\dd v.
\]
With $t=\sqrt{v^2+\mu^2}-\mu$ one has
$v^3\dd v=(t^3+3\mu t^2+2\mu^2t)\dd t$;
its integral equals $3/8+3\mu/4+\mu^2/2$.
The same substitution with $e^{-t}$ gives $\mathcal H_0$.
Finally rescale $v=\sqrt{2\mu}\,z$ in $\mathcal H_d$;
the exponential tends to $e^{-z^2}$ and is bounded by a function
proportional to $e^{-c\min(z^2,z)}$ for $\mu\ge1$. Indeed the positive exponent magnitude is
$2z^2/(1+\sqrt{1+2z^2/\mu})$.
This yields \eqref{eq:OQED-large-mu-crossover}.
\end{proof}

\subsection{Fixed-graph application and the remaining vertex condition}
\label{sec:OQED-graph-application}

In the flat coefficient prototype, with this representative assigned to
the designated photon and fermion lines, use
\begin{equation}
\begin{aligned}
 S_{a,M}(q)&=F_{a,M}(q^2)
                 \frac{-\ii\slashed q+m_e}{q^2+m_e^2},\\
 D_{a,M,E}^{\mu\nu}(q)&=
           \frac{F_{a,M}(q^2)}{q^2}\delta^{\mu\nu}
                 \quad(\alpha_{\rm gf}=1).
\end{aligned}
\label{eq:OQED-Euclidean-lines}
\end{equation}
These are auxiliary uncut coefficients, not a redefinition of the
source-normalized physical two-point functions. There is one kernel
per designated uncut line. If dressed lines are used, no additional
copy of their product is inserted into the graph. The same rule applies
to contacts and to graphs with no internal photon. The Lorentzian
source convention, at finite boundary magnitude before coefficient
extraction and the admitted common limit, retains
\begin{equation}
\begin{aligned}
 z_{\eta,\bb}(q)&=q^2+\ii\eta\epsvec,\\
 F^{\rm dir}_{a,M}(q)&=
    \ExpDir\!\left[aM-a\SqrtDir(M^2-z_{\eta,\bb}(q))\right],\\
 \vec S(q)&=\eta\ii\widehat u_q\,
    \embed\!\left[(\slashed q+m_e)
      \embed^{-1}\!\left(\frac{F^{\rm dir}_{a,M}(q)}
                               {z_{\eta,\bb}(q)-m_e^2}\right)\right],\\
 \vec D^{\mu\nu}(q)&=-\eta\ii\widehat u_q\,
       \frac{F^{\rm dir}_{a,M}(q)}{z_{\eta,\bb}(q)}g^{\mu\nu}.
\end{aligned}
\label{eq:OQED-directed-lines}
\end{equation}
The photon minus sign follows the Maxwell coefficient in
\eqref{eq:OQED-source-action}; it is not erased by an absolute
Euclidean estimate.
The spinor embedding is the entrywise convention of
\eqref{eq:physical-directed-Dirac-line}. These remain resolved
scalar-pole templates, with the same finite-regulator and source-exchange
qualifications as Sections~\ref{sec:carrier} and~\ref{sec:LSZ}.

\begin{theorem}[O--QED application with explicit kernel hypotheses]
\label{thm:OQED-fixed-graph}
Fix the support, branch, source/endpoint conventions, and positive
line depths of the O--QED record. Suppose a finite completed loop graph
has full-rank flat routing and the infrared hypotheses of
Theorem~\ref{thm:global-uv}. On a full-rank set of loop lines use
\(F_{a_i,M_i}\), with \(a_i>0\), \(M_i\geq0\), and bound all
other multipliers as in Corollary~\ref{cor:damped-basis}.
If the full completed numerator has a finite polynomial bound of degree
\(D_G\), the graph is absolutely Euclidean convergent, uniformly on
the admitted compact external set. The kernel part of this assertion
is established by \eqref{eq:OQED-inverse-moments}, not left as an
unevaluated assumption.

Alternatively, assume bounded covariance envelopes and a complete
numerator bound
\begin{equation}
 |N_G(k,p)|\leq C(1+\norm{k}^{D_G})
          \exp\!\left(\sum_i v_i|q_i|\right),\qquad v_i\geq0.
\label{eq:OQED-vertex-budget}
\end{equation}
In this second alternative, every exponentially bounded line has its
specified \(a_i>0\); for a merely bounded unselected multiplier set
\(a_i=v_i=0\) in the bound, without changing its model definition.
It is sufficient that the signed residuals \(h_i=a_i-v_i\) obey
\eqref{eq:signed-linear-coercivity}; a positive right side in
\eqref{eq:signed-linear-basis-certificate} is one explicit certificate.
The simpler nonnegative-residual, full-loop-rank test remains sufficient.
These bounds concern the product of all vertices and numerators in the
graph, including contact vertices. The constants \(e^{a_iM_i}\) belong
to the envelope, not to a renormalization of the kernel measure.

If either condition, or the verified forest alternative, holds for
all graphs and insertions of one fixed completed action, and the
independent anomaly-free BRST hypotheses hold for the same action in
both characters, its admitted coefficients are graphwise UV finite and
satisfy the identities of Section~\ref{sec:brst} at every finite order.
This last implication does not certify an unspecified completion.
\end{theorem}

\begin{proof}
For each selected line the required moment is finite:
\[
 a_i^{-2-D_G}M_{1+D_G/2}[\rho_{a_iM_i,\bb}]<\infty.
\]
Corollary~\ref{cor:damped-basis} therefore applies directly. In the second case
\eqref{eq:OQED-linear-exponential-bounds} leaves
\(\exp[-\sum_i(a_i-v_i)|q_i|]\) times a fixed constant
\(\exp(\sum_i a_iM_i)\). The linear residual-coercivity proof then
supplies an integrable polynomial times \(e^{-c\norm{k}}\).
The character norm transfers the two coefficient estimates to the
carrier graph. The finite-order and BRST conclusions use their respective
corollary and theorem with the same complete graph data.
\end{proof}

For the conventional local trilinear comparison vertices, the fermion
numerators alone give \(D_G\leq I_f\), with \(I_f\) the number
of internal fermion lines. Additional polynomial vertex derivatives
increase that bound by their total degree. The polarization bubble,
electron self-energy, electron--photon vertex, and four-photon fermion
box have respectively the sufficient benchmark degrees
\(2,1,2,4\), hence moment orders \(2,3/2,2,3\).
All four required moments are available. Together with the local
infrared analysis below, this proves convergence of these specified
comparison integrands; it does not claim that unchanged local vertices
satisfy the dressed Ward identity. The massless photon prevents a
direct invocation of the globally bounded-covariance hypothesis until
that separate infrared analysis is supplied.

\paragraph{Massless photons and infrared control}
The physical photon denominator in \eqref{eq:OQED-Euclidean-lines}
is massless. A positive-mass comparison is not substituted for it inside
a Ward identity. In the vacuum one-loop self-energy and vertex graphs
with \(m_e>0\), choose the photon momentum as the loop variable.
At its origin the only photon singularity is \(|q|^{-2}\), locally
integrable with the four-dimensional measure; the remaining massive
fermion factors are bounded there on compact real external sets.
The polarization and fermion box have only massive internal fermions.
Their polynomial-vertex benchmark integrals are therefore finite
without adding a photon mass. This off-shell Euclidean statement does
not remove on-shell soft-photon singularities or prove uniform
infrared bounds at higher orders.

A useful extension avoids assuming individually bounded denominators.
Let \((\Xi,\nu_{\bb})\) be the positive measure space of all loop or spectral
variables, and let \(R:\Xi\to[0,\infty)\) be a measurable radial size
in reference momentum units. Let the nonnegative measurable function
\(W_G\) contain the rational, matrix-norm, and polynomial factors left
after extracting a residual exponential. Suppose, with constants
\(C,C_0,c>0\), \(N\geq0\) uniform on the external domain,
\begin{equation}
\begin{aligned}
 |I_G(\xi,p)|&\leq C e^{-cR(\xi)}W_G(\xi,p),\quad c>0,\\
 \int_{R(\xi)\leq R_0}W_G(\xi,p)\dd\nu_{\bb}(\xi)
    &\leq C_0(1+R_0)^N,\qquad R_0\geq1.
\end{aligned}
\label{eq:OQED-spectral-volume-certificate}
\end{equation}
Then the graph is absolutely convergent: splitting into
\(n\leq R<n+1\) bounds it by
\(CC_0\sum_{n\geq0}e^{-cn}(n+2)^N<\infty\).
The series beginning at \(n=\lfloor R_*\rfloor\) also bounds the
numerical tail beyond \(R_*\). Local infrared integrability is part of the
second inequality, not inferred from exponential UV suppression.
This also gives an explicit sufficient extension to a specified
thermal integration/sum or Landau-level representation, provided its
multiplicities, transition matrix elements, and completed vertex growth
are included in \(W_G\). It does not infer those estimates from the
flat vacuum propagator.

\subsection{An evaluated one-loop certificate and its Ward limitation}
\label{sec:OQED-one-loop}

The exact one-loop integrability condition can be evaluated without
special functions. For every \(p>0\), the substitution
\(y=\sqrt{u+M^2}-M\), \(u=y^2+2My\), gives
\begin{equation}
 \int_0^\infty F_{a,M}(u)^p\dd u
   =2\int_0^\infty(y+M)e^{-pa y}\dd y
   =\frac{2}{p^2a^2}+\frac{2M}{pa}.
\label{eq:OQED-kernel-Lp}
\end{equation}
Consequently \(\mathcal Q_{F_{a,M}}=1/(2a^2)+M/a\) and
Proposition~\ref{prop:QED-general-obstruction} applies to the source's
vacuum prototype with bare coefficient vertices. In particular,
\begin{equation}
\begin{aligned}
 \kappa_{a,M,m_e}
 &=\frac{e^2}{8\pi^2}\int_0^\infty
  \frac{u(u+2m_e^2)}{(u+m_e^2)^2}
  e^{-2a(\sqrt{u+M^2}-M)}\dd u,\\
 0<\kappa_{a,M,m_e}
 &\leq\frac{e^2}{16\pi^2a^2}(1+2aM)\quad(e\ne0).
\end{aligned}
\label{eq:OQED-explicit-Ward-obstruction}
\end{equation}
Equality holds for \(m_e=0\) and is strict for \(m_e>0\).
Indeed, the rational factor is
\(1-m_e^4/(u+m_e^2)^2\), between zero and one. Thus
\(\Pi_{\rm bv}^{\mu\nu}(0)=\kappa_{a,M,m_e}\delta^{\mu\nu}\)
and \(P_\mu\Pi_{\rm bv}^{\mu\nu}(P)
=\kappa_{a,M,m_e}P^\nu+o(|P|)\). The finite bare comparison
is not transverse. Neither normalization at \(u=0\) nor choosing
\(M>0\) removes this obstruction.

\begin{center}
\small
\begin{tabular}{@{}crr@{}}
\toprule
\(aM\) & \(a^2\mathcal Q_F\) & \(\kappa/e^2\), \(a=m_e=1\)\\
\midrule
0 & 0.5 & 0.003774745738\\
0.5 & 1.0 & 0.008196730399\\
1 & 1.5 & 0.013357748315\\
2 & 2.5 & 0.024602484262\\
\bottomrule
\end{tabular}
\end{center}
These are dimensionless coefficient benchmarks in one reference
momentum unit, not calibrated plasma inputs or physical photon masses.
The tabulated values were checked using both the positive integral and
its integration-by-parts form
\eqref{eq:general-Ward-obstruction-derivative}. Values with \(M\leq m_e\)
are valid Euclidean comparisons, not certificates of an isolated
Lorentzian shell below the kernel branch point.

The normalized massive kernel also makes the distinction between
regularity and ultraviolet growth concrete. For
\(A(u)=F_{a,M}(u)^{-1}\),
\begin{equation}
 A'(u)=\frac{a}{2\sqrt{u+M^2}}A(u),\qquad M>0.
\label{eq:OQED-inverse-derivative}
\end{equation}
Hence the Ball--Chiu-type vertex of
\eqref{eq:longitudinal-vertex}, with \(B=m_eA\), has regular
divided differences at the origin and diagonal value
\(\Gamma^\mu(0,0)=\ii\gamma^\mu\). This removes the specific
massive zero-momentum cusp obstruction of the \(M=0\) kernel.
It does not bound the vertex polynomially at large momentum:
\(A(R^2)\sim e^{a(R-M)}\).

\begin{proposition}[A Ward-compatible two-vertex bubble fails the graphwise test]
\label{prop:OQED-Ward-bubble-obstruction}
Use the flat vacuum coefficient \(S_{a,M}=F_{a,M}S_0\),
\(a>0\), \(m_e>0\), \(M\geq0\), and \(e\ne0\).
Assume the vertex satisfies the dressed Ward identity and is continuous
at \((p',p)=(k,k)\) for every \(k\ne0\). Consider precisely the
two-vertex term of the source's polarization skeleton at \(P=0\),
without its contact contribution. Its trace over external Euclidean
indices is not absolutely integrable. More precisely, if
\[
 I_W^{\mu\nu}(k):=-e^2\Tr\bigl[
       \Gamma^\mu(k,k)S_{a,M}(k)
       \Gamma^\nu(k,k)S_{a,M}(k)\bigr],
\]
then, for \(u=k^2>0\), \(\Omega=\sqrt{u+M^2}\),
\begin{equation}
\begin{aligned}
 \sum_{\mu=1}^4I_W^{\mu\mu}(k)=e^2\biggl[
 &\frac{8(u+2m_e^2)}{(u+m_e^2)^2}
 -\frac{8a u}{\Omega(u+m_e^2)}\\
 &-\frac{4a^2u}{u+M^2}\biggr]
 \longrightarrow-4e^2a^2 .
\end{aligned}
\label{eq:OQED-Ward-bubble-trace}
\end{equation}
Thus this particular two-vertex graph cannot meet a graphwise
absolute-convergence theorem on an external set containing zero.
\end{proposition}

\begin{proof}
The soft Ward identity \eqref{eq:soft-Ward-vertex} fixes the diagonal
vertex; a continuous transverse addition has zero diagonal limit.
Multiplying by the specified propagator gives the exact cancellation
\begin{equation}
 \Gamma^\mu(k,k)S_{a,M}(k)
       =\ii\gamma^\mu S_0(k)+\frac{a k^\mu}{\Omega}\one.
\label{eq:OQED-Ward-line-cancellation}
\end{equation}
Let \(G_0^\mu=\ii\gamma^\mu S_0\). The Euclidean Clifford trace
identities give
\[
 \sum_\mu\Tr[G_0^\mu G_0^\mu]
     =-\frac{8(u+2m_e^2)}{(u+m_e^2)^2},\qquad
 \Tr G_0^\mu=\frac{4k^\mu}{u+m_e^2}.
\]
The cross terms in the squared product contribute
\(8a u/[\Omega(u+m_e^2)]\), and its scalar-square terms contribute
\(4a^2u/(u+M^2)\). Multiplication by \(-e^2\) proves
\eqref{eq:OQED-Ward-bubble-trace}. Its nonzero asymptote implies
failure of absolute integrability of at least one tensor component.
For an explicit radial exhaustion,
\begin{equation}
 \int_{|k|<R}\frac{\dd^4k_{\bb}}{(2\pi)^4}
       \sum_\mu I_W^{\mu\mu}(k)
       =-\frac{e^2a^2}{8\pi^2}R^4+O(R^3).
\label{eq:OQED-Ward-bubble-exhaustion}
\end{equation}
This radius only displays the asymptotic failure; it is not a new
physical cutoff. The locally integrable \(M=0\) endpoint does not
affect the ultraviolet conclusion.
\end{proof}

\begin{corollary}[The obstruction persists at nonzero transfer]
\label{cor:OQED-fixed-transfer-obstruction}
Retain \(S_{a,M}=F_{a,M}S_0\), \(a>0\), \(m_e>0\),
\(M\geq0\), and \(e\ne0\), but no longer assume diagonal
continuity. Fix real \(P\ne0\), let \(k'=k+P\), and require both
vertices of the two-vertex bubble to satisfy
\(P\cdot\Gamma=S_{a,M}^{-1}(k')-S_{a,M}^{-1}(k)\) in its common
external orientation. Define its complete two-vertex coefficient by
\[
 \begin{aligned}
 I_W^{\mu\nu}(k,P):=-e^2\Tr\bigl[&\Gamma^\mu(k',k)S_{a,M}(k)\\
                    &\Gamma^\nu(k,k')S_{a,M}(k')\bigr].
 \end{aligned}
\]
This graph is not componentwise absolutely integrable. Specifically,
\begin{equation}
\begin{aligned}
 P_\mu P_\nu I_W^{\mu\nu}(k,P)
 =-e^2\Tr\bigl[&S_{a,M}^{-1}(k')S_{a,M}(k)\\
               &+S_{a,M}^{-1}(k)S_{a,M}(k')-2\one\bigr],
\end{aligned}
\label{eq:OQED-double-Ward-contraction}
\end{equation}
and along \(k=rn\), \(|n|=1\), \(r\to\infty\),
\begin{equation}
 P_\mu P_\nu I_W^{\mu\nu}(rn,P)
       \longrightarrow-8e^2\bigl[\cosh(a n\cdot P)-1\bigr].
\label{eq:OQED-fixed-transfer-obstruction}
\end{equation}
\end{corollary}

\begin{proof}
Use the two vertices \(\Gamma^\mu(k',k)\) and
\(\Gamma^\nu(k,k')\). Their Ward identities have the same sign when
both are contracted with \(P\); the second follows by reversing its
own transfer. Expanding the product of the two inverse differences
and using the ordered inverse identities gives
\eqref{eq:OQED-double-Ward-contraction}.
For a finite-momentum check put \(u=k^2\), \(u'=k'^2\),
\(b=k\cdot k'+m_e^2\), and
\(d_\Omega=\sqrt{u'+M^2}-\sqrt{u+M^2}\). Direct Clifford traces give
\begin{equation}
\begin{split}
 P_\mu P_\nu I_W^{\mu\nu}(k,P)
 =-4e^2\biggl[ b\biggl(&\frac{e^{a d_\Omega}}{u+m_e^2}
                  +\frac{e^{-a d_\Omega}}{u'+m_e^2}\biggr)-2\biggr].
\end{split}
\label{eq:OQED-double-Ward-finite-trace}
\end{equation}
No factors \(F\) are discarded in this identity. Along \(k=rn\),
\(d_\Omega=n\cdot P+O(r^{-1})\) uniformly on the unit sphere,
and \(b/(u+m_e^2)=1+O(r^{-1})\), as does its primed partner.
This proves the stated limit with a uniform \(O(r^{-1})\) error. It is strictly negative
on an angular set of nonzero measure, so its absolute radial integral
diverges. Absolute integrability of every tensor component would imply
integrability of this fixed contraction, a contradiction. Transverse
vertex additions do not change either Ward contraction.
\end{proof}

The fixed-transfer contraction proves failure of absolute convergence
for the specified vacuum prototype's two-vertex term, beyond the earlier
vertex-growth test. The Ward contractions alone still do not establish
that the vertices are functional derivatives of one completed action.
This test does not compute the contact contribution, prove divergence
of the bubble-plus-contact sum, or exclude a physical completion with
a different admitted source and field realization. If an exact source
pushforward generates a single two-photon vertex from these terms,
its complete unregularized integrand must be formed and supplied with
its own absolute majorant.
When the original bare expansion counts the bubble and the contact as
different graphs, convergence of their sum certifies only their summed
response. The bubble's graphwise obstruction still applies on that
expansion. Nor does the Ward identity determine every
transverse or contact term. The source already leaves those objects
in its completion record; the present calculation identifies a specific
requirement they must address.

The choice of expansion matters for a geo-phase as well. An exact
centered Gaussian average gives $e^{-g^2V(k)/2}$ for a phase of
variance $g^2V(k)$. At nonzero real $g$ this can damp a hard loop,
while its coefficient at order $g^0$ is one. If $g$ is counted as a
bare interaction, its original coefficients therefore need their own
estimates. For each graph, suppose $I_G(k;g)$ is holomorphic on
a neighborhood of $|g|\le r$, $r>0$. On the circle assume
$|I_G(k;g)|\le M_G(k)$ with $M_G$ integrable and independent of
any approximation.
Cauchy's formula gives $|[g^n]I_G(k;g)|\le r^{-n}M_G(k)$.
An action that puts the normal phase in its free measure defines a
different expansion and retains the phase when counting photon vertices.

\paragraph{A radial profile cannot remove the diagonal Ward obstruction}
There is an estimate independent of the chosen square-root kernel.
Set $S_F(k)=F(k^2)S_0(k)$ with $F$ positive and differentiable on
$(0,\infty)$. Use a diagonal-continuous coefficient vertex obeying
the dressed Ward identity. For $k\ne0$ and any unit vector $v\perp k$,
\eqref{eq:soft-Ward-vertex} gives
$v_\mu\Gamma^\mu(k,k)S_F(k)=\ii\slashed v S_0(k)$. Since
$\slashed v\slashed k\slashed v=-\slashed k$ and
$\Tr\one=4$, the two-vertex tensor obeys
\begin{equation}
 v_\mu v_\nu I_W^{\mu\nu}(k,0)
     =\frac{4e^2}{k^2+m_e^2}.
 \label{eq:uvpeB-ward-tangential}
\end{equation}
Taking a norm of the tensor at each $k$ allows $v$ to depend on $k$.
For $e\ne0$ its integral therefore grows at least quadratically with
the radius of a four-dimensional momentum ball, for every such $F$.
This is a bound on the absolute two-vertex graph, independent of any
contact and of the sign of the Euclidean index trace.

\paragraph{A finite-wavelength test including the local Dirac action}
At fixed $P\ne0$ put $k'=k+P$, $r=|k|$, $r'=|k'|$,
$t(r)=\sqrt{r^2+m_e^2}$, and introduce an arbitrary positive reference
$t_*$ of momentum dimension so that $U(r)=\log[t(r)/(t_*F(r^2))]$
is dimensionless. For the two-vertex Ward contractions of
Corollary~\ref{cor:OQED-fixed-transfer-obstruction}, the exact counterpart
of \eqref{eq:OQED-double-Ward-finite-trace} is
\begin{equation}
 \begin{aligned}
 P_\mu P_\nu I_W^{\mu\nu}(k,P)
   &=8e^2[1-c(k,k')\cosh(U(r')-U(r))],\\
 c(k,k')&=\frac{k\cdot k'+m_e^2}{t(r)t(r')}.
 \end{aligned}
 \label{eq:uvpeB-ward-cosine}
\end{equation}
This follows by tracing the two ordered inverse--propagator products.
The factor $c$ is the cosine between the Euclidean vectors
$(k,m_e)$ and $(k',m_e)$. Suppose $|rU'(r)|\le C$ at all sufficiently
large $r$. This covers $F\equiv1$ and, for example, every
$F(u)=(1+u/b)^{-\nu}$, $b>0$, $\nu>0$.
For $k=rn$, uniformly in the unit direction $n$ at fixed $P$,
\[
 \begin{aligned}
 r'-r&=n\cdot P+O(r^{-1}),\\
 1-c(k,k')&=\frac{|P|^2-(n\cdot P)^2}{2r^2}+O(r^{-3}).
 \end{aligned}
\]
To verify the second line, use
$t(r)^2t(r')^2-(k\cdot k'+m_e^2)^2
=t(r)^2|P|^2-(k\cdot P)^2$ and rationalize $1-c$.
On $|n\cdot P|<\delta|P|$ the mean-value bound on $U$ gives
$|U(r')-U(r)|\le 2C(\delta|P|+O(r^{-1}))/r$. The positive
$1-c$ term then dominates $c[\cosh(U(r')-U(r))-1]$ for all
large $r$ if $\delta>0$ is chosen sufficiently small. Thus
$P_\mu P_\nu I_W^{\mu\nu}\ge c_0 e^2|P|^2/r^2$ on a cone of
positive solid angle for some $c_0>0$. Its absolute integral again
grows at least quadratically. The square-root exponential has
unbounded logarithmic slope and instead obeys the stronger asymptote
\eqref{eq:OQED-fixed-transfer-obstruction}.
For $F=1$ this proves failure of absolute integrability directly
for the minimally coupled local Dirac bubble at nonzero external
momentum. A Ward identity imposed after a shift of a divergent
integral does not change that result. For a projected GCM charged
current the two-vertex test applies only if its actual inverse,
source insertions, and routing have this form.

\subsection{Contact completion must be tested before integration}
\label{sec:OQED-contact-test}

The response skeleton in Ref.~\cite{ChiapponeIgnition} explicitly
retains a contact sector but does not evaluate it. We first state the mathematical
requirement for combining it with the bubble. A vanishing integrated
Ward contraction is not, by itself, a bound on the absolute integrand.

\begin{lemma}[Finite common-routing sums]
\label{lem:OQED-grouped-convergence}
At a fixed perturbative order, suppose a finite set of action-generated
terms has been placed on one measure space \((\Xi,\nu_{\bb})\), with all
routing Jacobians, signs, and multiplicities retained. Let
\[
 J_{\mathcal B}(\xi,p)=\sum_{j\in\mathcal B}c_jJ_j(\xi,p)
\]
be their pointwise sum. Suppose the chosen common finite-domain
prescription represents that action coefficient, and there is an
integrable envelope \(M_{\mathcal B}\) such that
\[
 |J_{\mathcal B}(\xi,p)|\leq M_{\mathcal B}(\xi),\qquad p\in\mathcal K.
\]
Then its common-domain limit exists absolutely and is independent of
exhaustion. If the same property holds locally on a complex domain
and the summed integrand is holomorphic there almost everywhere, the
summed coefficient is holomorphic. These conclusions do not assert
absolute convergence of the separate terms.
\end{lemma}
\begin{proof}
Integrate the sum before taking the exhaustion limit. Dominated
convergence gives the absolute limit and its uniform tail bound; the
complex conclusion follows from Theorem~\ref{thm:holomorphy}. No exchange
with separately divergent integrals is used.
\end{proof}

This is a criterion for a declared finite sum, not an application of
Theorem~\ref{thm:global-uv} to a divergent constituent. Separate changes
of loop routing inside divergent terms are not licensed: any boundary
terms induced by such changes must be controlled in the same common
domain. All subintegration regions, not only simultaneous large
momenta, must satisfy the envelope at higher loops. At one loop a
local integrable envelope together with
\(\lVert J_{\mathcal B}(k,p)\rVert\leq C(1+|k|)^{-4-\delta}\),
\(\delta>0\), outside a fixed ball is one simple sufficient condition.
Cancellation after angular averaging alone is not an absolute bound.

For example, if the contact contribution to the double contraction
has a continuous uniform angular limit at \(k=rn\), absolute
integrability of the complete one-loop sum forces that limit to be
\[
 \lim_{r\to\infty}P_\mu P_\nu I_{\rm contact}^{\mu\nu}(rn,P)
       =8e^2[\cosh(a n\cdot P)-1].
\]
Otherwise a nonzero limiting sum on an open angular set would reproduce
the divergent radial integral in
Corollary~\ref{cor:OQED-fixed-transfer-obstruction}. This cancellation
of the leading term is necessary under the stated limit assumption,
not sufficient: slower-decaying remainders can still be nonintegrable.

\paragraph{A separately specified contact test}
To examine that distinction explicitly, impose an additional assumption
not fixed by the response source alone: in a flat vacuum coefficient
calculation with
field-independent measure, a constant Abelian background enters the
quadratic fermion symbol only by the shift
\[
 \mathcal D[\mathcal A](k)=\mathcal D(k+e\mathcal A),\qquad
 \mathcal D(k)=F_{a,M}(k^2)^{-1}(\ii\slashed k+m_e).
\]
Here \(\mathcal A\) is the constant background on the noncompact flat
chart, and the two external insertions carry zero momentum. This specifies the constant-background response. The covariant trial
in Section~\ref{sec:polish-covariant-trial} has this limit for a constant
background and fixes its higher vertices at nonconstant fields through
the action. Those vertices are not inferred from the Ward identity alone. Fix \(a,M,m_e>0\) and \(e\ne0\).
Using the same bubble sign as \eqref{eq:bare-QED-loop}, the corresponding
contact contribution from a second variation of the trace logarithm is
\(e^2\Tr[\mathcal D_{,\mu\nu}S]\), with \(S=\mathcal D^{-1}\)
and commas denoting momentum derivatives. No second multiplier is
inserted on this already dressed line. This definition fixes the
relative bubble/contact sign for the following auxiliary test; it is
not inferred from the Ward identity alone.

\begin{proposition}[A contact cancellation that is not absolute convergence]
\label{prop:OQED-constant-contact}
In this constant-background realization the combined coefficient is
\begin{equation}
\begin{aligned}
 J_{\rm cb}^{\mu\nu}(k)
 &=e^2\Tr[\mathcal D_{,\mu\nu}S
           -\mathcal D_{,\mu}S\mathcal D_{,\nu}S]\\
 &=e^2\partial_\mu\partial_\nu\log\det\mathcal D(k).
\end{aligned}
\label{eq:OQED-contact-Hessian}
\end{equation}
It has the exact tensor form, with \(u=k^2\), \(\Omega=\sqrt{u+M^2}\),
\begin{equation}
\begin{split}
 \frac{J_{\rm cb}^{\mu\nu}}{e^2}
 =4a\left(\frac{\delta^{\mu\nu}}{\Omega}
               -\frac{k^\mu k^\nu}{\Omega^3}\right)
   +\frac{4\delta^{\mu\nu}}{u+m_e^2}
               -\frac{8k^\mu k^\nu}{(u+m_e^2)^2}.
\end{split}
\label{eq:OQED-contact-explicit}
\end{equation}
The leading constant of the two-vertex term cancels, but the summed
tensor is still not componentwise absolutely integrable. In fact,
\begin{equation}
\begin{split}
 \sum_\mu J_{\rm cb}^{\mu\mu}(k)
 =e^2\left[
 \frac{4a(3u+4M^2)}{(u+M^2)^{3/2}}
       +\frac{8(u+2m_e^2)}{(u+m_e^2)^2}\right]>0,
\end{split}
\label{eq:OQED-contact-trace}
\end{equation}
and its radial exhaustion is exactly
\begin{equation}
 \int_{|k|<R}\frac{\dd^4k_{\bb}}{(2\pi)^4}
            \sum_\mu J_{\rm cb}^{\mu\mu}(k)
  =\frac{e^2}{2\pi^2}\left[
        \frac{aR^4}{\sqrt{R^2+M^2}}+\frac{R^4}{R^2+m_e^2}\right].
\label{eq:OQED-contact-exhaustion}
\end{equation}
\end{proposition}
\begin{proof}
Differentiate \(S\mathcal D=\one\), so
\(\partial_\nu S=-S\mathcal D_{,\nu}S\), and use cyclicity of the
finite matrix trace. This proves \eqref{eq:OQED-contact-Hessian}.
In four Dirac components,
\[
 \log\det\mathcal D(k)
   =4a(\sqrt{u+M^2}-M)+2\log[(u+m_e^2)/p_{\rm ref}^2]
\]
up to a constant. The fixed momentum unit \(p_{\rm ref}>0\) only makes
the logarithm dimensionless; all derivatives are independent of it.
Twice differentiating gives \eqref{eq:OQED-contact-explicit}, whose
trace is \eqref{eq:OQED-contact-trace}. The divergence theorem on a
finite ball, with sphere area \(2\pi^2R^3\), gives
\eqref{eq:OQED-contact-exhaustion}. Its leading term is
\(e^2aR^3/(2\pi^2)\), proving the asserted failure. The determinant
here is only a finite matrix identity at fixed momentum; no determinant
of an unbounded field operator is being defined by it.
\end{proof}

Subtracting the separately typed ordinary symbol
\(\mathcal D_{\rm ref}=\ii\slashed k+m_e\) in this same constant-background trial
removes the second term of \eqref{eq:OQED-contact-trace}, not the
nonintegrable first term. Nor may the integral be set to zero merely
because \eqref{eq:OQED-contact-Hessian} is a total derivative: its
surface term is displayed explicitly. A symmetry-preserving prescription
for an otherwise divergent functional is a different construction from
an absolutely convergent Lebesgue integral. This calculation neither
assigns a physical photon mass nor changes the fixed-depth kernel rule.
It tests the additional constant-background, trace-log, and measure
assumptions just stated. Additional action terms, measure determinants,
or source maps need their own derivation and common-integrand estimate;
they are not silently excluded or supplied by this test.

\paragraph{The longitudinal contact at nonzero transfer}
The same radial vacuum symbol gives a further test with nonconstant
probes. Assume gauge covariance on the differentiated Fourier
domain, $\mathcal D[A+\dd\lambda]
=e^{-\ii e\lambda}\mathcal D[A]e^{\ii e\lambda}$, and retain the
field-independent measure and trace-log sign above. Fix real $P\ne0$.
Use $k,k+P$ for the bubble momenta and $k$ for the contact momentum,
on one common ball in $k$ when testing an exhaustion. Polarizing with
$\lambda_\pm(x)=e^{\pm\ii P\cdot x}$ tests the double longitudinal
contraction, since the factors $\ii$ and $-\ii$ in their gradients
multiply to one.

The mixed derivative of the kinetic operator is
$-e^2[\lambda_+,[\lambda_-,\mathcal D]]$, where the bracket uses
multiplication by these Fourier modes. Its diagonal at momentum $k$
is $e^2[\mathcal D(k+P)+\mathcal D(k-P)-2\mathcal D(k)]$.
Adding the doubly contracted bubble
\eqref{eq:OQED-double-Ward-contraction} gives the pointwise coefficient
\begin{equation}
 P_\mu P_\nu J^{\mu\nu}(k,P)=
 e^2\Tr\big[\mathcal D(k-P)S(k)-\mathcal D(k)S(k+P)\big].
 \label{eq:finaluv-finite-transfer-contact}
\end{equation}
All products here are finite matrices at the indicated momenta.

For $k=rn$ with $|n|=1$, set $u=n\cdot P$ and
$v=|P|^2-u^2$. The square-root damper obeys
$\Omega(k\pm P)-\Omega(k)=\pm u+v/(2r)+O(r^{-2})$,
where $\Omega(k)=\sqrt{|k|^2+M^2}$. Both spin trace ratios,
after removing the exponential and the spin factor four, are
$1-u/r+O(r^{-2})$. The first exponential in
\eqref{eq:finaluv-finite-transfer-contact} is
$e^{-au}[1+av/(2r)+O(r^{-2})]$, and the second is
$e^{-au}[1-av/(2r)+O(r^{-2})]$. Therefore
\[
 P_\mu P_\nu J^{\mu\nu}(rn,P)
   =4e^2 a e^{-au}\frac{v}{r}+O(r^{-2}),
\]
uniformly on the sphere. Since $v>0$ on a set of positive solid
angle, this contraction is not absolutely integrable in the
declared routing. It has a nonzero cubic radial divergence.

Equation~\eqref{eq:finaluv-finite-transfer-contact} has the form
$g(k)-g(k+P)$. Shifting the second integral on a finite ball
moves its domain and gives the boundary contribution measured by
the cubic divergence. A regrouping or trace prescription that
changes that boundary term needs an independent integrability
argument. This calculation tests the longitudinal coefficient
of this trace-log trial in the designated common routing.

\begin{corollary}[Constant-background contact obstruction for a general radial damper]
\label{cor:OQED-general-contact-obstruction}
Retain the constant-background shift, trace-log sign, one-species
normalization, and field-independent measure of
Proposition~\ref{prop:OQED-constant-contact}. Replace \(F_{a,M}\) by
any positive radial multiplier \(F(u)\) that is \(C^2\) for
\(u>u_0>0\) and satisfies \(F(u)\to0\) as \(u\to\infty\).
Set \(h_F=-\log F\), keep \(m_e>0\) and \(e\ne0\), and subtract
the ordinary symbol's contact-plus-bubble coefficient \emph{before}
integration. Then the trial-minus-local difference is not componentwise absolutely
integrable on the exterior of a ball. No monotonicity or rate of decay
of \(F\) is needed for this assertion.
\end{corollary}
\begin{proof}
At fixed momentum, \(\mathcal D_F=F^{-1}\mathcal D_{\rm ref}\)
has four Dirac components. Subtracting its finite matrix logarithmic
derivatives therefore gives exactly
\begin{equation}
 \Delta J_F^{\mu\nu}(k)
     =4e^2\partial_\mu\partial_\nu h_F(k^2).
 \label{eq:OQED-general-contact-difference}
\end{equation}
Choose \(R_0^2>u_0\). On each finite annulus the divergence theorem
and the retained Fourier normalization give
\begin{equation}
\begin{split}
 \int_{R_0<|k|<R}\frac{\dd^4k_{\bb}}{(2\pi)^4}
             \sum_\mu\Delta J_F^{\mu\mu}(k)
 =\frac{e^2}{\pi^2}\bigl[&R^4h_F'(R^2)\\
                      &-R_0^4h_F'(R_0^2)\bigr].
\end{split}
 \label{eq:OQED-general-contact-flux}
\end{equation}
Indeed, the outward radial derivative of \(4e^2h_F(k^2)\) is
\(8e^2R h_F'(R^2)\), and the sphere area is \(2\pi^2R^3\).
If every tensor component were absolutely integrable outside \(R_0\),
the trace integrals on the left would be bounded independently of \(R\).
Hence \(|u^2h_F'(u)|\leq C\) for all sufficiently large \(u\).
It would follow that \(\int^\infty |h_F'(u)|\dd u<\infty\), so
\(h_F\) had a finite limit. This contradicts
\(h_F(u)=-\log F(u)\to+\infty\).
\end{proof}

The distinction between trace cancellation and absolute tensor control
can also be read off pointwise. In the Euclidean spectral matrix norm,
define, for real \(u>u_0\),
\[
 \mathcal H_F(u):=\max\{|h_F'(u)|,
                         |h_F'(u)+2u h_F''(u)|\}.
\]
The radial and three tangential eigenvalues of
\(\Delta J_F=8e^2h_F'I+16e^2h_F''kk^T\) give the exact identity
\begin{equation}
 \int_{|k|>R_0}\frac{\dd^4k_{\bb}}{(2\pi)^4}
            \|\Delta J_F(k)\|_2
   =\frac{e^2}{2\pi^2}\int_{R_0^2}^\infty u\mathcal H_F(u)\dd u.
 \label{eq:OQED-contact-norm-test}
\end{equation}
This identity holds with either side infinite. Finite-dimensional norm
equivalence makes finiteness equivalent to componentwise absolute
integrability. Since \(\mathcal H_F\geq|h_F'|\), it also gives a
direct alternative proof of the obstruction. Conversely, a nonzero
limit of \(F\) alone is not sufficient: for
\(h_F(u)=1-(1+u)^{-1}\), the multiplier tends to \(e^{-1}\)
and the flux in \eqref{eq:OQED-general-contact-flux} stays bounded,
but \(u|h_F'(u)|\sim u^{-1}\) is not integrable. The trace test
must therefore not be reversed into a tensor-convergence criterion.

This result concerns only the constant-background realization of the covariant trial above,
including its ordinary-symbol subtraction. It is not a statement about
every contact completion, a renormalized response difference, or an
infinite-dimensional determinant. It shows why changing the radial
profile alone cannot supply the missing common-integrand bound while
those assumptions are kept. Field-dependent measure contributions,
different endpoint realizations, or additional action terms must be
constructed and estimated, not presumed absent or presumed to cancel it.

\subsection{Local analytic gain and the directed pole test}
\label{sec:OQED-analytic}

For \(M>0\), the positive exponential moments in
\eqref{eq:OQED-exponential-moment} satisfy the additional kernel
hypothesis of Lemma~\ref{lem:massive-neighbourhood}. Therefore every
flat positive-mass graph with the lemma's polynomial holomorphic
numerator has a holomorphic neighbourhood of each compact real
Euclidean external set. In particular, for \(m_e>0\) the bare
polarization and the literal two-line scalar benchmark with mass \(m_e\)
are locally holomorphic at zero. For the latter, Euclidean rotational
invariance makes its local convergent expansion a function of \(P_E^2\).
The \(M=0\) nonanalyticity of Proposition~\ref{prop:stable-bubble-nonanalytic}
is consequently not a theorem about this \(M>0\) profile. No uniform
analytic neighbourhood as \(M\downarrow0\) is asserted. Massless
photon graphs, nonpolynomial completed vertices, and a medium require
separate common-majorant arguments.

\begin{proposition}[An explicit holomorphic tube for the scalar comparison]
\label{prop:OQED-bubble-tube}
For \(a,M,m>0\), let
\(C(q)=F_{a,M}(q^2)/(q^2+m^2)\), with the principal square-root branch, and set
\[
 \mathcal T_*=\{P\in\CC^4:|\Im P|<2m_*\},\qquad
 m_*:=\min\{m,M\}.
\]
The literal scalar comparison bubble has the holomorphic continuation
\begin{equation}
 B(P)=\int_{\RR^4}\frac{\dd^4k_{\bb}}{(2\pi)^4}
                     C(k-P/2)C(k+P/2),\qquad P\in\mathcal T_*.
\label{eq:OQED-bubble-tube}
\end{equation}
For every \(0\leq d<2m_*\), its integrand obeys the common bound
\begin{equation}
 |C(k-P/2)C(k+P/2)|
 \leq\frac{e^{2aM}e^{-2a|k|}}{(m^2-d^2/4)^2},
 \qquad |\Im P|\leq d.
\label{eq:OQED-bubble-tube-majorant}
\end{equation}
Near zero write \(B(P)=b(P^2)\), with \(P^2=\sum_\mu P_\mu^2\).
Then \(b(w)\) is holomorphic for \(|w|<4m_*^2\). For
\(0<r<\mathcal R<4m_*^2\), its Taylor coefficients \(b_n\) satisfy
\begin{equation}
\begin{aligned}
 |b_n|&\leq B_{\mathcal R}\mathcal R^{-n},\\
 \sup_{|w|\leq r}\left|b(w)-\sum_{n=0}^N b_n w^n\right|
   &\leq B_{\mathcal R}
          \frac{(r/\mathcal R)^{N+1}}{1-r/\mathcal R},\\
 B_{\mathcal R}&:=\frac{3e^{2aM}}
           {64\pi^2a^4(m^2-\mathcal R/4)^2}.
\end{aligned}
\label{eq:OQED-bubble-Taylor-bound}
\end{equation}
\end{proposition}
\begin{proof}
For real \(P\), a real translation of the absolutely convergent
integral gives the symmetric routing in \eqref{eq:OQED-bubble-tube}.
For complex \(P=x+\ii y\), put \(q_\pm=k\pm x/2\pm\ii y/2\).
For \(\Lambda=m\) or \(M\),
\[
 \Re(q_\pm^2+\Lambda^2)
   =|k\pm x/2|^2+\Lambda^2-|y|^2/4>0.
\]
Both denominators are nonzero and both principal square roots are
holomorphic. For \(\Re z>0\),
\(\Re\sqrt z=\sqrt{(|z|+\Re z)/2}\geq\sqrt{\Re z}\).
Thus \(|F_{a,M}(q_\pm^2)|\leq e^{aM}e^{-a|k\pm x/2|}\),
and \(|q_\pm^2+m^2|\geq m^2-d^2/4\). The triangle inequality
\(|k-x/2|+|k+x/2|\geq2|k|\) proves the envelope. Its integral is
\[
 \int\frac{\dd^4k_{\bb}}{(2\pi)^4}e^{-2a|k|}
                =\frac{3}{64\pi^2a^4}.
\]
Theorem~\ref{thm:holomorphy} proves the tube statement. Real rotational
invariance makes the local power series a series in \(P^2\).
Equivalently, the even holomorphic function \(B(te_1)\) for
\(|t|<2m_*\) descends to \(b(w)\), \(w=t^2\), throughout the
stated disk. On \(|w|=\mathcal R\), choose either square root of
\(w\); the bound applies with \(d=\sqrt{\mathcal R}\).
Cauchy's estimate and the geometric-series tail give
\eqref{eq:OQED-bubble-Taylor-bound}. No complex loop-momentum
translation or Lorentzian contour rotation is used.
\end{proof}

This is a sufficient domain and an explicit remainder for the literal
scalar comparison, not a maximal physical sheet. The same two-line
routing with a polynomial numerator admits this tube by an additional
compact-set polynomial majorant, with the stated positive masses.
The result does not cover the nonintegrable Ward-vertex bubble, the
constant-background contact test, or the unevaluated medium response.
It demonstrates concretely how a positive \(M\) changes the local
analytic problem; the guaranteed tube shrinks to zero as \(M\downarrow0\).

\paragraph{Certified scalar truncation and the first Taylor coefficient}
For the same symmetric routing, let \(B^{[R]}(P)\) be
\eqref{eq:OQED-bubble-tube} restricted to \(|k|\leq R\), \(R\geq0\).
The pointwise envelope gives a uniform absolute tail bound throughout
\(|\Im P|\leq d<2m_*\):
\begin{equation}
\begin{aligned}
 |B(P)-B^{[R]}(P)|
 &\leq\frac{e^{2aM-2aR}}{8\pi^2(m^2-d^2/4)^2}\,T_a(R),\\
 T_a(R)&:=\frac{R^3}{2a}+\frac{3R^2}{(2a)^2}
                   +\frac{6R}{(2a)^3}+\frac6{(2a)^4}.
\end{aligned}
 \label{eq:OQED-bubble-radial-tail}
\end{equation}
Four-dimensional polar integration reduces the omitted envelope to
\(\int_R^\infty r^3e^{-2ar}\dd r=e^{-2aR}T_a(R)\).
The radius is a numerical integration parameter, not a new theory scale.
This bound may be added to a separately certified quadrature error on
the retained ball; a floating-point comparison alone is not such a
certificate. Both errors refer to the same routing and complex domain.

The first two coefficients in \(b(w)=b_0+b_1w+\cdots\) also have
one-dimensional representations. Put
\(f(u)=F_{a,M}(u)/(u+m^2)\). Then
\begin{equation}
\begin{aligned}
 b_0&=\frac1{16\pi^2}\int_0^\infty u f(u)^2\dd u,\\
 b_1&=-\frac1{32\pi^2}\int_0^\infty u^2 f'(u)^2\dd u<0,\\
 f'(u)&=-f(u)\left[\frac{a}{2\sqrt{u+M^2}}+\frac1{u+m^2}\right].
\end{aligned}
 \label{eq:OQED-bubble-first-coefficients}
\end{equation}
To prove the derivative identity, differentiate the real bubble at zero
and integrate by parts, using smoothness at the origin and exponential
decay at infinity. Rotational invariance gives
\[
 \partial_{P_i}\partial_{P_j}B(0)
 =-\int\frac{\dd^4k_{\bb}}{(2\pi)^4}\partial_i C(k)\partial_j C(k)
 =2\delta_{ij}b_1.
\]
Since \(\partial_i C=2k_i f'(k^2)\), spherical averaging yields
\eqref{eq:OQED-bubble-first-coefficients}; strict negativity follows
from \(f'\ne0\). This is the Euclidean invariant coefficient. Under
\(w=-s\), its first invariant derivative changes sign. Neither sign
is a physical damping rate or an evaluation of the completed O--QED
polarization. Together with \eqref{eq:OQED-bubble-Taylor-bound}, these
formulas provide controlled approximation data for the stated scalar
comparison without interchanging an unbounded Taylor series with a loop
integral.

\begin{proposition}[Signed scalar approximation bounds]
\label{prop:OQED-scalar-CM-enclosure}
For the scalar comparison with \(a,M,m>0\), let \(w=P_E^2\)
on the real Euclidean axis. The invariant function \(b(w)\)
admits a finite positive Laplace measure \(\nu_{\bb}\) on
\((0,\infty)\), with representation
\begin{equation}
 \begin{aligned}
  b(w)&=\int_0^\infty e^{-\lambda w}\nu_{\bb}(\dd\lambda),
       \qquad w\geq0,\\
  \int_0^\infty e^{s\lambda}\nu_{\bb}(\dd\lambda)&<\infty,
       \qquad 0<s<4m_*^2.
 \end{aligned}
 \label{eq:OQED-scalar-positive-measure}
\end{equation}
Thus \(b\) is holomorphic for \(\Re w>-4m_*^2\) and, on the
real half-line, \((-1)^n b^{(n)}(w)>0\) for every \(n\geq0\).
Its coefficients \(b_n=b^{(n)}(0)/n!\) alternate in sign. Every
finite truncation \(p_N(w)=\sum_{n=0}^N b_nw^n\) satisfies
\begin{equation}
 0\leq(-1)^{N+1}[b(w)-p_N(w)]
       \leq |b_{N+1}|w^{N+1},\qquad w\geq0.
 \label{eq:OQED-scalar-signed-enclosure}
\end{equation}
In particular, odd truncations are lower bounds and even truncations
are upper bounds on the positive Euclidean invariant axis.
\end{proposition}
\begin{proof}
Use the two independent line parameters in
\eqref{eq:two-line-parameters}, with \(w_i=\alpha_i+a^2\tau_i\)
and \(\rho_{\bb}=\rho_{aM,\bb}\). Push forward the positive measure
\[
 \frac{e^{-m^2(\alpha_1+\alpha_2)}}{16\pi^2(w_1+w_2)^2}
 \dd\alpha_1\dd\alpha_2\rho_{\bb}(\dd\tau_1)\rho_{\bb}(\dd\tau_2)
\]
by \(\lambda=w_1w_2/(w_1+w_2)>0\). Its mass is the finite
number \(b_0\), and Tonelli proves the first equality for \(w\geq0\).
Since \(\lambda\leq(w_1+w_2)/4\), multiplication by
\(e^{s\lambda}\) is bounded by reducing the Schwinger mass square
to \(m^2-s/4>0\) and tilting each line measure by
\(e^{sa^2\tau/4}\). For \(s<4m_*^2\), the tilt is strictly
below \((aM)^2\); its mass and required inverse moments remain
finite. The existing Gaussian estimate therefore proves the second
line of \eqref{eq:OQED-scalar-positive-measure}. It supplies all
moments of \(\nu_{\bb}\), a common holomorphic majorant on compact
subsets of the half-plane, and
\[
 (-1)^n b^{(n)}(w)=\int\lambda^n e^{-\lambda w}\nu_{\bb}(\dd\lambda).
\]
Strict positivity follows because the measure is nonzero and supported
on positive \(\lambda\). Finally the integral form of the Taylor
remainder of \(e^{-x}\), for \(x\geq0\), has sign
\((-1)^{N+1}\) and modulus at most \(x^{N+1}/(N+1)!\).
Integrating that finite-order inequality proves
\eqref{eq:OQED-scalar-signed-enclosure} without exchanging an
infinite series and an integral.
\end{proof}

This positive measure is a proper-time measure for the auxiliary
\emph{bubble}; it is not a positive physical mass-spectral measure or
an OS reconstruction. Its half-plane continuation agrees with
\eqref{eq:OQED-bubble-tube} wherever the latter is used, by uniqueness
from the real Euclidean region. The signed inequalities apply at each
finite order even outside the guaranteed Taylor disk; they do not
assert convergence of that series there. They supplement, rather than
replace, the complex-disk estimate \eqref{eq:OQED-bubble-Taylor-bound}.
For a useful independent coefficient check, integration by parts and
four-dimensional rotational averaging also give
\begin{equation}
 \begin{aligned}
 b_2&=\frac1{192\pi^2}\int_0^\infty
             u[2f'(u)+uf''(u)]^2\dd u>0,\\
 b_1^2&\leq2b_0b_2 .
 \end{aligned}
 \label{eq:OQED-scalar-second-coefficient}
\end{equation}
Indeed, \(\partial_{P_1}^4B(0)=\int(\partial_1^2C)^2\dd^4k_{\bb}/(2\pi)^4
=24b_2\). For a radial \(C\), its integrated value is one eighth
of \(\int(\Delta C)^2\dd^4k_{\bb}/(2\pi)^4\); here
\(\Delta C=4[2f'+uf'']\). Boundary terms vanish under the same
positive-mass smoothness and exponential bounds. The last inequality
is Cauchy--Schwarz for \(1,\lambda\) in \(L^2(\nu)\).

At the line level the branch in \eqref{eq:OQED-directed-lines} starts
at the timelike invariant \(q^2=M^2\). If \(0<m_e^2<M^2\), the
numerator is holomorphic and nonzero near the electron shell, but its
coefficient there is
\begin{equation}
 f_{a,M}(m_e^2)
    =e^{a(M-\sqrt{M^2-m_e^2})},
\label{eq:OQED-electron-shell-weight}
\end{equation}
not one. The photon shell at zero instead has unit numerator for
\(M>0\). These are local auxiliary pole statements; physical residues
are those of the completed source matrix. If \(m_e^2=M^2\), the pole
coincides with the kernel branch point and the isolated-meromorphic
pole argument is unavailable. For \(m_e^2>M^2\) the two numerator
boundaries are \(e^{aM}e^{\pm\ii a\sqrt{m_e^2-M^2}}\), and the
cut--discontinuity test retains their branch contribution. The source's
kernel-only inequality \(m_e^2<M^2\) does not exclude medium cuts,
vertex singularities, or competing physical poles.

The internal vectorial datum \(\ii\eta\epsvec\) remains separate
from the retarded response limit \(z=\omega+\ii0^+\) after the
physical readout. Positive Laplace weights do not establish that
readout, reflection positivity, or undressed operational-cut matching.
The spectral obstruction in \eqref{eq:spectral-obstruction} continues
to apply to a raw covariance with this exponentially decaying profile.


\subsection{A finite zero-depth limit after common scalar subtraction}
\label{sec:OQED-scalar-subtraction}

The positive-depth graph theorem does \emph{not} extend the
unsubtracted scalar bubble uniformly to $a=0$. Nevertheless, a
particular \emph{difference formed before integration} has a
finite limit. This provides a useful comparison with ordinary
subtracted scalar field theory, not a kernel-off renormalization
of the Ward-completed gauge polarization.

\begin{theorem}[Subtracted literal bubble at vanishing depth]
\label{thm:OQED-subtracted-zero-depth}
Fix $m>0$ and $M\geq0$. For real Euclidean $P$, $w=P^2\geq0$,
define $C_{a,M,m}(q):=F_{a,M}(q^2)/(q^2+m^2)$ and
\begin{equation}
 D_{a,M,m}(w):=B_{a,M,m}(0)-B_{a,M,m}(P).
 \label{eq:OQED-subtracted-bubble-def}
\end{equation}
At every $a>0$ this is an absolutely convergent \emph{common}
integral and satisfies
\begin{equation}
 0\leq D_{a,M,m}(w)\leq\frac{w}{64\pi^2m^2}.
 \label{eq:OQED-subtraction-uniform-bound}
\end{equation}
As $a\downarrow0$ with $M$ fixed,
\begin{equation}
 \boxed{\displaystyle
 \lim_{a\downarrow0}D_{a,M,m}(w)
   =\frac1{16\pi^2}\int_0^1
      \log\left(1+\frac{x(1-x)w}{m^2}\right)\dd x.}
 \label{eq:OQED-subtraction-limit}
\end{equation}
Convergence is uniform when $w$ and $M$ range over compact
subsets of $[0,\infty)$ and $m$ ranges over a compact subset of
$(0,\infty)$. In addition,
\begin{equation}
 \lim_{a\downarrow0}D_{a,M,m}'(0)
  =\frac1{96\pi^2m^2}.
 \label{eq:OQED-subtraction-slope-limit}
\end{equation}
In contrast, the \emph{unsubtracted} comparison has the
universal finite-part asymptotic
\begin{equation}
 B_{a,M,m}(0)=\frac1{8\pi^2}
  \left[\log\frac1{2am}-\gamma_{\mathrm E}-\frac12\right]
  +o(1),\quad a\downarrow0,
 \label{eq:OQED-unsubtracted-log-depth}
\end{equation}
where $\gamma_{\mathrm E}$ is Euler's constant. In this
normalization the constant inside the logarithm is specified,
not absorbed into an arbitrary reference momentum scale.
The comparison makes no assertion about a jointly completed
O--QED kernel-off gauge functional or about uniform Lorentzian
boundary values.
\end{theorem}
\begin{proof}
For $a>0$, $C_{a,M,m}$ is real, even, and square-integrable.
A real translation gives the manifestly positive \emph{single}
representation
\begin{equation}
 D_{a,M,m}(P^2)=\frac12\int\frac{\dd^4k_{\bb}}{(2\pi)^4}
 [C_{a,M,m}(k-P/2)-C_{a,M,m}(k+P/2)]^2.
 \label{eq:OQED-subtraction-difference-square}
\end{equation}
This is an identity between convergent terms at $a>0$; its
right-hand side remains the preferred representation at the limit.

For the bound and the limit, use the positive dimensionful heat-depth
probability law $\rho^{\rm heat}_{a,M,\bb}$ for $F_{a,M}$, defined
as the push-forward of $\rho_{aM,\bb}$ under $\tau\mapsto a^2\tau$,
so $F_{a,M}(u)=\int e^{-\lambda u}\rho^{\rm heat}_{a,M,\bb}
(\dd\lambda)$. Use it in \eqref{eq:two-line-parameters}. Let
$w_i=\alpha_i+\lambda_i$ and $S=\alpha_1+\alpha_2$.
The proper-time coefficient for $D$ has the nonnegative factor
\[
 \frac{e^{-m^2S}}{(w_1+w_2)^2}
 \left[1-e^{-\frac{w_1w_2}{w_1+w_2}w}\right].
\]
Since $\lambda:=w_1w_2/(w_1+w_2)\leq(w_1+w_2)/4$ and
$1-e^{-\lambda w}\leq\lambda w$, this factor is at most
$w e^{-m^2S}/[4(w_1+w_2)]\leq w e^{-m^2S}/(4S)$.
The integral of $e^{-m^2S}/S$ over the positive $\alpha$ quadrant
is $m^{-2}$. Including $1/(16\pi^2)$ proves the bound.

For each fixed $u\geq0$, $F_{a,M}(u)\to1$ as $a\downarrow0$.
The normalized positive heat-depth measures therefore concentrate at
$\lambda_i=0$: for any $s,\delta>0$,
\[
 \rho^{\rm heat}_{a,M,\bb}([\delta,\infty))
 \leq\frac{1-F_{a,M}(s)}{1-e^{-s\delta}}\longrightarrow0.
\]
The preceding $\alpha$-integrable upper bound permits first taking
the heat-measure limit for $\alpha_1,\alpha_2>0$ and then using
dominated convergence in those variables. With both heat depths set
to zero, put $\alpha_1=xS$, $\alpha_2=(1-x)S$. The resulting
positive integral is
\[
 \frac1{16\pi^2}\int_0^1\dd x\int_0^\infty
  \frac{e^{-m^2S}}S
       [1-e^{-Sx(1-x)w}]\dd S,
\]
which evaluates to \eqref{eq:OQED-subtraction-limit}. The same
uniform bound applied to finite differences, or the bounded first
$w$-derivative of the proper-time integrand, gives equicontinuity
and hence compact-uniform convergence. The derivative at $w=0$
converges by the identical domination and equals
$[16\pi^2m^2]^{-1}\int_0^1x(1-x)\dd x$.

The argument is uniform on the stated compact $M,m$ sets:
$1-F_{a,M}(s)\leq a\sqrt{s}$ controls the heat-law
concentration independently of $M$, and $m\geq m_->0$ gives
one common integrable Schwinger envelope. The same domination
controls the first $w$-derivative. The explicit finite part
in \eqref{eq:OQED-unsubtracted-log-depth} is proved in
Proposition~\ref{prop:OQED-zero-depth-finite-part} below.
Thus finiteness is achieved only after the \emph{common}
subtraction; the unsubtracted integral remains divergent.
\end{proof}


\begin{proposition}[Universal logarithmic finite part at zero depth]
\label{prop:OQED-zero-depth-finite-part}
For the literal scalar bubble at $P=0$, with $m>0$, $M\geq0$,
put $s=2am>0$ and
$L_{a,m}:=[\log(1/s)-\gamma_{\mathrm E}-1/2]/(8\pi^2)$.
Then
\begin{equation}
 |B_{a,0,m}(0)-L_{a,m}|
 \leq \frac{(2+\pi/4)s}{8\pi^2},
 \label{eq:OQED-unsubtracted-finitepart-error}
\end{equation}
and
\begin{equation}
 0\leq B_{a,M,m}(0)-B_{a,0,m}(0)
 \leq(e^{2aM}-1)B_{a,0,m}(0).
 \label{eq:OQED-kernel-mass-finitepart-error}
\end{equation}
Consequently \eqref{eq:OQED-unsubtracted-log-depth} holds
uniformly for $M$ in bounded nonnegative sets and $m$ in
compact positive intervals. Its divergent coefficient and
finite constant do not depend on $M$ under this limiting
prescription. It is a coefficient-level statement about the
unsubtracted \emph{auxiliary scalar bubble}, not a physical
counterterm or an interacting QED amplitude.
\end{proposition}
\begin{proof}
At $M=0$, radial substitution $r=mx$ gives
\begin{equation}
 8\pi^2B_{a,0,m}(0)
 =\int_0^\infty\frac{x^3e^{-sx}}{(1+x^2)^2}\dd x
 =I(s)-K(s),
 \label{eq:OQED-massless-kernel-log-integrals}
\end{equation}
where
$I(s)=\int_0^\infty e^{-sx}x/(1+x^2)\dd x$ and
$K(s)=\int_0^\infty e^{-sx}x/(1+x^2)^2\dd x$.
Let
$g(x)=x/(1+x^2)-\mathbf1_{[1,\infty)}(x)/x$.
One checks directly that
$\int_0^\infty g(x)\dd x=0$ and
$\int_0^\infty x|g(x)|\dd x=1$. Hence
\[
 I(s)=E_1(s)+\int_0^\infty e^{-sx}g(x)\dd x,
 \qquad
 |I(s)+\gamma_{\mathrm E}+\log s|\leq2s,
\]
using $|E_1(s)+\gamma_{\mathrm E}+\log s|\leq s$.
Also $K(0)=1/2$ and
$|K(s)-1/2|\leq s\int_0^\infty
 x^2(1+x^2)^{-2}\dd x=\pi s/4$.
Together these inequalities establish
\eqref{eq:OQED-unsubtracted-finitepart-error}.

For general $M$, let
$\delta_M(r):=r+M-\sqrt{r^2+M^2}\in[0,M]$.
Then
$e^{-2ah_M(r^2)}=e^{-2ar}e^{2a\delta_M(r)}$.
Multiplying by the positive massive radial weight and
integrating gives \eqref{eq:OQED-kernel-mass-finitepart-error}.
For bounded $M$, its right side is
$O(aM\log(1/a))$ as $a\downarrow0$, uniformly when
$m$ stays in a compact positive interval. This establishes
the stated universal finite part with no subtraction
of two separately undefined integrals.
\end{proof}

\subsection{The kernel-off comparison is not a new UV theorem}
\label{sec:OQED-reference}

Following the matched kernel-on/kernel-off construction in
Ref.~\cite{ChiapponeIgnition}, retain the
same state, gauge, source normalization, action-completion prescription,
and renormalization/matching record for
\begin{equation}
 \Delta\Pi_{R,\Phi}(a)
       :=\Pi_{R,\Phi}^{\rm matched}(a)-\Pi_{R,\rm ref}^{\rm ren},
       \qquad \Phi_{\rm off}\equiv1.
\label{eq:OQED-matched-increment}
\end{equation}
The reference is separately defined. Finiteness of the fixed-depth
term does not render an unrenormalized kernel-off QED integral finite.
The shared record fixes physical normalization conditions and the
reference prescription; it does not mean that identical divergent bare
counterterms must be added to an already finite graph. Where the two
completed renormalized tensors and their analytic continuations exist,
their Ward identities imply transversality of their difference. This
linearity argument determines neither its sign nor a spatial power.

The optional real interpolation illustrates the boundary of the UV
hypotheses. To avoid confusing its parameter with the sector label,
write
\begin{equation}
\begin{aligned}
 F_t&=(1-t)+tF_{a,M},\qquad 0\leq t\leq1,\\
 \rho_{t,\bb}&=(1-t)\delta^{\rm heat}_{\bb,0}+t\rho_{aM,\bb}.
\end{aligned}
\label{eq:OQED-interpolation-atom}
\end{equation}
For every \(t<1\), its measure has a nonzero atom at zero and
\(F_t(u)\to1-t\) as \(u\to\infty\). Thus the moment theorem
cannot be applied along the whole path merely because it applies at
\(t=1\). For the bare two-line zero-momentum prototype,
\begin{equation}
 \left.\partial_tF_t(u)^2\right|_{t=0}=2(F_{a,M}(u)-1),\qquad
 \partial_t^2F_t(u)^2=2(F_{a,M}(u)-1)^2.
\label{eq:OQED-interpolation-derivatives}
\end{equation}
The two factors tend to \(-2\) and \(2\), respectively. Multiplying
the undamped QED tensor, whose absolute radial tail is proportional
to \(r\,\dd r\) on an angular set of nonzero measure, makes neither
derivative integrable. The unrenormalized difference factor
\(F_{a,M}^2-1\to-1\) has the same obstruction. These examples do
not disprove differentiability of a separately completed, consistently
renormalized response family. They show why the uniform derivative
majorants assumed in the source study's optional interpolation audit are
additional hypotheses, not consequences of fixed-depth UV convergence.

In contrast, derivatives with respect to the depth on a compact interval
\(0<a_-\leq a\leq a_+\) satisfy
\[
 |\partial_a^nF_{a,M}(u)|=h_M(u)^nF_{a,M}(u)
              \leq u^{n/2}F_{a_-,M}(u).
\]
Every fixed finite number of line derivatives therefore has a polynomial
cost covered by the verified moments, when the remaining completed
vertex derivatives and infrared bounds are supplied. This is a
positive-depth parameter estimate, not a physical zero-depth limit.

\subsection{A labelled elastic benchmark and its precise limitations}
\label{sec:fw-elastic-benchmark}

The graph estimates and Ward tests above are not cross sections. A
restricted two-species comparison makes explicit how the same-copy source
pairing and the GT endpoint convention would enter such a calculation.
Use canonically normalized electron and muon fields on $\bb$, with
positive masses $m_e,m_\mu$. In the increasing-$\gamma_{\bb}$ coefficient
chart put $p_\gamma=\kappa_{\mathbf p}>0$ and
$p_\parallel=\sigma_s p_\gamma$; the spin frame is transported with that
change of chart. The plane's geometry is still $\RGTL$, not an
unlabelled replacement by $\RMink$.

\paragraph{An action that supplies the photon-only tree rule}
Let $H_F=F_{a,M}(-\Delta_E)^{-1}$ on its spectral domain. As a distinct
Abelian trial, take ordinary massive Dirac coefficients and the Euclidean
photon/gauge-fixing form
\begin{equation}
\begin{aligned}
 S_{E,\bb}^{\gamma F}
  =\int\dd\mu_{E,\bb}\bigg[&
    \sum_{f=e,\mu}\bar\psi_{f,\bb}
       \bigl(\gamma_E^\nu(\partial_\nu+\ii eA_{\nu,\bb})+m_f\bigr)
                  \psi_{f,\bb}\\
  &+\tfrac14 F_{\mu\nu,\bb}H_FF_{\mu\nu,\bb}
    +\frac1{2\zeta}(\partial\cdot A_{\bb})H_F(\partial\cdot A_{\bb})\\
  &+\bar c_{\bb}H_F(-\Delta_E)c_{\bb}\bigg],\qquad\zeta>0.
\end{aligned}
\label{eq:fw-photon-only-action}
\end{equation}
This is a declared coefficient action on a common domain, not the original
local reference Lagrangian with a line edited afterwards. Its photon
Hessian and inverse at nonzero Euclidean momentum are
\begin{equation}
 K_{A,E}=q^2H_F(q^2)(P_T+\zeta^{-1}P_L),\qquad
 K_{A,E}^{-1}=\frac{F_{a,M}(q^2)}{q^2}(P_T+\zeta P_L).
\label{eq:fw-photon-hessian}
\end{equation}
The real-space operator domain and Lorentzian continuation remain
hypotheses. The gauge-fixing/ghost pair is generated as in
\eqref{eq:fw-oqed-gaugefermion}. The photon occurs once; it is not
repeated for each species. This trial leaves the fermion propagators
and vertices ordinary. Its fermion vacuum-polarization loop consequently
has no damped internal photon and is not covered by a claim of universal
internal-line suppression. It is not the fermion-dressed trial
\eqref{eq:polish-covariant-dirac}.

\paragraph{Why zero carrier transfer needs a source prescription}
For elastic electron--muon scattering in its center-of-spatial-momentum
frame, $q_{\rm tr}=p'_e-p_e$ has $q_{{\rm tr},\gamma}=0$, hence
$q_{{\rm tr},\parallel}=0$, although $q_{\rm tr}^2<0$ away from the
forward direction. The normalized running GT unit is then undefined.
It cannot be replaced by zero or silently by the fixed pole unit.
For $q_\parallel\ne0$ let $r_q=\sgn(q_\parallel)$. On the fixed directed
line, the raw Maxwell tag and canonical coefficient representative obey
\begin{equation}
\begin{aligned}
 \vec D_{\eta,\bb}^{\rm tag}(q)&=r_q\vec D_{\eta,\bb}^{\rm can}(q),\\
 \vec D_{\eta,\bb}^{\rm can}(q)
   &:=-\eta\je\,
       \frac{F^{\rm dir}_{\eta,\bb}(q)}
            {q^2+\ii\eta\epsvec}\,g^{\mu\nu}.
\end{aligned}
\label{eq:fw-tagged-canonical-pairing}
\end{equation}
Indices $\mu\nu$ on both sides are understood. A declared dual attachment
$r_q$ recovers the canonical kernel away from that hyperplane. Its
spacelike coefficient then has a continuous zero-carrier extension in
this photon-only action. The raw tag has opposite one-sided limits where
the even coefficient is nonzero. This source pairing must be part of the
model before graph integration; equality almost everywhere does not
define an external evaluation on $q_\parallel=0$. An unrestricted
running-unit theory without that pairing has not supplied this benchmark.

\paragraph{Amplitude, spin trace, and positive flux}
Let $p,p'$ be electron and $k,k'$ muon momenta in the $\gamma_{\bb}$
coefficient chart, and put
$\sMand=(p+k)^2$, $\tMand=(p'-p)^2$, $\uMand=(p-k')^2$.
These are scattering invariants, not sector labels or the clock readout.
The two assignments of the distinguishable species cancel the Dyson
factor $1/2!$. Each vertex has its own event measure and ambient support
delta before pullback. After coefficient extraction and removal of the
single overall conservation distribution, the tree coefficient is
\begin{equation}
\begin{aligned}
 \mathcal M_{\bb}^{\rm toy}
   &=\frac{e^2F_{a,M}(-\tMand)}{\tMand}
       [\bar u_e(p')\gamma^\mu u_e(p)]
       [\bar u_\mu(k')\gamma_\mu u_\mu(k)],\\
 \overline{|\mathcal M_{\bb}^{\rm toy}|^2}
   &=\frac{2e^4F_{a,M}(-\tMand)^2}{\tMand^2}\,
                  \mathcal B(\sMand,\tMand,\uMand),\\
 \mathcal B&=(\sMand-m_e^2-m_\mu^2)^2
            +(\uMand-m_e^2-m_\mu^2)^2\\
   &\quad+2\tMand(m_e^2+m_\mu^2).
\end{aligned}
\label{eq:fw-elastic-amplitude}
\end{equation}
All spinors here belong to $\bb$. Use
$\sum_hu_f(p,h)\bar u_f(p,h)=\slashed p+m_f$ and average over the four
physically prepared initial spins. The electron trace is
$4[p'^\mu p^\nu+p'^\nu p^\mu+(m_e^2-p'\cdot p)g^{\mu\nu}]$;
contracting it with the muon trace gives the displayed $\mathcal B$.
The on-shell equations give $q_{\rm tr}\cdot j_e=q_{\rm tr}\cdot j_\mu=0$,
so the longitudinal photon coefficient drops from this tree calculation.
This is not the loop Ward test of a dressed fermion action.

With the positive normalization \eqref{eq:fw-positive-state-measure},
define $\lambda_K(x,y,z)=x^2+y^2+z^2-2xy-2xz-2yz$. The incident flux,
outgoing phase space and differential cross-section coefficient are
\begin{equation}
\begin{aligned}
 \mathcal F_{\rm in}&=2\sqrt{\lambda_K(\sMand,m_e^2,m_\mu^2)},\\
 p_*&=\frac{\sqrt{\lambda_K(\sMand,m_e^2,m_\mu^2)}}{2\sqrt{\sMand}},
 &\dd\Phi_{2,\bb}&=\frac{p_*}{16\pi^2\sqrt{\sMand}}\dd\Omega,\\
 \frac{\dd\sigma_{\bb}^{\rm toy}}{\dd\tMand}
  &=\frac{\overline{|\mathcal M_{\bb}^{\rm toy}|^2}}
            {16\pi\lambda_K(\sMand,m_e^2,m_\mu^2)}.
\end{aligned}
\label{eq:fw-elastic-phase-space}
\end{equation}
These coefficient conventions agree with the usual derivation
\cite{TongQFTFramework}. For N with ascending GT order, the final pair
occupies the bra. For A with ascending GT order, the final pair occupies
the ket, while the spin average and outgoing phase-space integral still
follow physical preparation and detection. Reverse GT order reverses
those slots as in Table~\ref{tab:fw-endpoints}; the complete adjoint,
spinor duals and pole branch must be transported together. Subject to
those physical identifications the modulus is unchanged, not the
unextracted complex amplitude by fiat.

At fixed $\tMand<0$ this trial yields the ratio $F_{a,M}(-\tMand)^2$.
The forward Coulomb singularity persists because $F_{a,M}(0)=1$.
No finite total cross section, loop-corrected infrared-safe observable,
positive Lorentzian reconstruction or all-orders gauge-theory theorem
follows from this tree comparison. Those require the same complete
source/state/action construction tested throughout the paper.

\subsection{From the UV certificate to the medium response}
\label{sec:OQED-medium}

The vacuum line \eqref{eq:OQED-Euclidean-lines} carries no density,
temperature, magnetic spectrum, or distribution function. The source
instead specifies a state \(\mathcal S\), its integration/summation
measure \(\dd\mu_{E,\bb}^{\mathcal S}\), and a nonnegative kinetic
operator/eigenvalue \(\mathcal Q_{E,\mathcal S}\). In the scalar
spectral realization used there,
\begin{equation}
 S_{E,\Phi}(Q;\mathcal S)
   =F_{a,M}(\mathcal Q_{E,\mathcal S}(Q))
                     S_{E,\rm ref}(Q;\mathcal S).
\label{eq:OQED-medium-line}
\end{equation}
A noncommuting background realization instead requires its ordered
operator covariance and common domain. One cannot infer scalar
commutation with spin or magnetic projectors from
\eqref{eq:OQED-medium-line}.
With all arguments and index contractions provided by that realization,
the source's one-loop skeleton is
\begin{equation}
\begin{aligned}
 \Pi_{E,\Phi}^{\mu\nu}(P)
  ={}&-e^2\int\dd\mu_{E,\bb}^{\mathcal S}(Q)\,
    \Tr[\Gamma_\Phi^\mu S_{E,\Phi}(Q)
              \Gamma_\Phi^\nu S_{E,\Phi}(Q+P)]\\
   &+\Pi_{E,\rm contact,\Phi}^{\mu\nu}(P).
\end{aligned}
\label{eq:OQED-completed-skeleton}
\end{equation}
For a magnetic spectral basis, the compact \(Q+P\) notation includes
the actual transition kernels and sums, not a scalar substitution for
Landau-level matrix elements. The UV proof applies after those factors
and the contact terms have been expanded at a specified perturbative
order or bounded as a declared skeleton. A skeleton built from dressed
objects does not by itself have one strict perturbative loop order.

Equation~\eqref{eq:OQED-spectral-volume-certificate} states a concrete
sufficient estimate for the resulting measure, including all discrete
multiplicities. KMS or a separately justified Euclidean state is required
for this route. A genuinely nonequilibrium response needs the real-time
functional stated in the source, with its own estimates; the vacuum
Gaussian proof is not automatically a theorem for that functional.
Similarly, the common holomorphic domain and retarded boundary topology
must be proved for the completed increment, not inferred from a
pointwise expansion of the kernel. At even order that expansion must
retain odd--odd products with vertices, propagators, and state factors.

\paragraph{A quantitative transfer bound}
The ultraviolet analysis can supply convergent representations and
controlled numerical tails. To use those bounds in the source's RT
pole calculation, the response error must be controlled on the same
complex continuation domain. The following estimate makes that
requirement quantitative.
Let \(X=\Delta\Pi_{R,\Phi}\), \(D_0=D_{\rm EM,0}\),
\(U=V_{\rm RT,EM}\), and \(V=V_{\rm EM,RT}\). First work with
bounded operators on fixed electromagnetic and retained RT Hilbert
spaces: \(D_0\) is boundedly invertible, \(X,Y\) act on the
electromagnetic space, and \(U,V\) have the indicated opposite channel
directions. Assume \(r_X=\|D_0X\|<1\). The source's exact Schur
correction is
\begin{equation}
 \delta K_X=-U(R_X-D_0)V,\qquad
 R_X=(D_0^{-1}-X)^{-1},
\label{eq:OQED-Schur-increment}
\end{equation}
with \(K_0=D_{\rm RT,0}^{-1}\). For another admissible response
\(Y\), with \(r_Y=\|D_0Y\|<1\), the resolvent identity gives
\begin{equation}
\begin{aligned}
 \delta K_X-\delta K_Y&=-U R_X(X-Y)R_Y V,\\
 \|\delta K_X-\delta K_Y\|
 &\leq\frac{\|U\|\|V\|\|D_0\|^2}
                {(1-r_X)(1-r_Y)}\|X-Y\|.
\end{aligned}
\label{eq:OQED-Schur-error-bound}
\end{equation}
No commutation is used: multiplication of both inverse identities
proves the ordered resolvent formula. An unbounded electromagnetic
inverse can also be admitted on a fixed graph-norm domain
\(\mathcal D_{\rm EM}\). Require both \(D_0^{-1}\) and
\(D_0^{-1}-Y\) to be bounded isomorphisms from
\(\mathcal D_{\rm EM}\) to \(\mathcal H_{\rm EM}\) in these norms, and use
\(X-Y:\mathcal D_{\rm EM}\to\mathcal H_{\rm EM}\),
\(R_Y:\mathcal H_{\rm EM}\to\mathcal D_{\rm EM}\),
\(U:\mathcal D_{\rm EM}\to\mathcal H_{\rm RT}\), and
\(V:\mathcal H_{\rm RT}\to\mathcal H_{\rm EM}\).
Every inverse and product below then uses those fixed domain/codomain
norms; in particular \(R_Y(X-Y)\) is bounded on
\(\mathcal D_{\rm EM}\). A merely formal common dense domain is not
a substitute for these mapping bounds. The retained RT operator
estimates below remain bounded-space statements.

For the pole calculation use a finite retained pencil, or a separately
admitted analytic Fredholm simple-pole realization. Assume the operators
are holomorphic on the same neighbourhood of the algebraically simple
baseline pole \(z_*\), and use the source's \emph{pole} bookkeeping family
\begin{equation}
 K_{X,\lambda_{\rm p}}(z)
       :=K_0(z)+\lambda_{\rm p}\delta K_X(z),\qquad
 z_X(0)=z_*,\quad
 \dot z_X:=\left.\partial_{\lambda_{\rm p}}z_X\right|_0.
 \label{eq:OQED-pole-bookkeeping-family}
\end{equation}
Here \(X\) and its full Schur correction are held fixed while
\(\lambda_{\rm p}\) is varied. This is not the family obtained by
replacing \(X\) by \(\lambda_{\rm p}X\) inside the electromagnetic
resolvent: the latter has first variation \(-UD_0XD_0V\), rather
than \(\delta K_X\). Neither bookkeeping parameter is a new coupling.
Let \(v\) and \(w^\dagger\) be the baseline right/left null vectors
with \(w^\dagger K_0'(z_*)v\ne0\). The pole-family derivatives and
rate coefficients satisfy
\begin{equation}
\begin{aligned}
 \dot z_X&=-\frac{w^\dagger\delta K_X(z_*)v}
                      {w^\dagger K_0'(z_*)v},\qquad
 \Gamma_X^{(1)}=-2\Im\dot z_X,\\
 |\Gamma_X^{(1)}-\Gamma_Y^{(1)}|
 &\leq2\Lambda_*\|\delta K_X(z_*)-\delta K_Y(z_*)\|,\\
 \Lambda_*&=\frac{\|w\|\|v\|}{|w^\dagger K_0'(z_*)v|}.
\end{aligned}
\label{eq:OQED-pole-error-bound}
\end{equation}
These follow by differentiating the null-vector equation and by the
operator norm inequality. The adjoint is the physical represented
adjoint, not the carrier reality operation. The result propagates a
certified response error into the exact-Schur, first-order pole rate;
it does not bound the finite pole-shift remainder. That remainder and
its projected rate accuracy remain separate requirements of the response model.
A numerical error proved only on real Euclidean momenta cannot be
inserted into \eqref{eq:OQED-Schur-error-bound} at a complex RT pole
without a controlled continuation of the error itself.

\begin{corollary}[Certification from an approximate response]
\label{cor:OQED-aposteriori-Schur}
Suppose only an approximate response \(Y\) and a certified bound
\(\|X-Y\|\leq\epsilon_\Pi\) are available on the common complex domain.
Assume \(R_Y=(D_0^{-1}-Y)^{-1}\) exists there and
\(\|R_Y\|\leq B_Y\), with \(q:=B_Y\epsilon_\Pi<1\), uniformly on
the domain under consideration. Then the exact inverse exists and
\begin{equation}
\begin{aligned}
 \|R_X\|&\leq\frac{B_Y}{1-q},\\
 \|\delta K_X-\delta K_Y\|
       &\leq E_K:=\frac{\|U\|\|V\|B_Y^2\epsilon_\Pi}{1-q},\\
 |\Gamma_X^{(1)}-\Gamma_Y^{(1)}|&\leq2\Lambda_*E_K
                 \quad\text{at the common baseline pole}.
\end{aligned}
\label{eq:OQED-aposteriori-Schur}
\end{equation}
The norms of \(U,V\) are their uniform bounds on that domain, or
their pointwise bounds when only a pointwise statement is made.
For the graph-domain realization, \(B_Y\) is the
\(\mathcal H_{\rm EM}\to\mathcal D_{\rm EM}\) bound and
\(\epsilon_\Pi\) the \(\mathcal D_{\rm EM}\to\mathcal H_{\rm EM}\)
bound. Analyticity of \(X,Y,D_0,U,V\) in the indicated bounded-operator
spaces gives holomorphic inverses and Schur corrections; the norm
inequalities alone do not establish that analytic hypothesis.
\end{corollary}
\begin{proof}
Write \(E=X-Y\). The ordered factorization
\[
 D_0^{-1}-X=(D_0^{-1}-Y)(\one-R_YE)
\]
and the Neumann series give
\(R_X=(\one-R_YE)^{-1}R_Y\). The exact resolvent difference
\(R_X-R_Y=R_X E R_Y\) gives the second estimate, and
\eqref{eq:OQED-pole-error-bound} gives the third. This argument does
not assume a bound on the unknown \(\|D_0X\|\). In particular,
\(B_Y=\|D_0\|/(1-r_Y)\) is available when \(r_Y<1\).
\end{proof}

The strict margin \(\Gamma_Y^{(1)}>2\Lambda_*E_K\) certifies a
positive exact first-order transferred rate on the same represented
branch. It does not certify a positive finite shifted-root rate:
that additional statement needs the independently controlled pole
remainder. Near a null rate the absolute error bound remains meaningful,
whereas a relative error ratio need not. The bound uses the physical
pole functional and is invariant under rescaling either null vector.

\begin{proposition}[A finite-pole enclosure on a retained matrix space]
\label{prop:OQED-finite-pole-enclosure}
On a finite-dimensional retained RT space, define the full pencils
\(K_X=K_0+\delta K_X\) and \(K_Y=K_0+\delta K_Y\).
Suppose both are holomorphic on a neighbourhood of the closed disk
\(\overline{\mathbb D}(z_Y,r)\), \(r>0\), and that
\(\det K_Y\) has exactly one zero there, the simple zero at \(z_Y\),
with no zero on the boundary. Assume a certified bound
\(\|K_X(z)-K_Y(z)\|\leq E_K(z)\) on that boundary and
\begin{equation}
 \sup_{|z-z_Y|=r}\|K_Y(z)^{-1}\|E_K(z)<1.
 \label{eq:OQED-finite-pole-contour-test}
\end{equation}
Then \(\det K_X\) has exactly one zero in the disk, also simple,
and its location \(z_X\) satisfies \(|z_X-z_Y|<r\).
For \(Z\in\{X,Y\}\), use these enclosed endpoint roots to define
the rate-comparison quantities relative to one fixed \(z_*\):
\begin{equation}
\begin{aligned}
 \Gamma_Z^{\rm fin}&:=-2\Im(z_Z-z_*),\\
 |\Gamma_X^{\rm fin}-\Gamma_Y^{\rm fin}|&<2r.
\end{aligned}
 \label{eq:OQED-finite-pole-rate-enclosure}
\end{equation}
Thus \(\Gamma_Y^{\rm fin}\geq2r\) certifies a strictly positive
\(\Gamma_X^{\rm fin}\). Identifying this enclosed root with the
selected baseline mode additionally requires the contour condition
below.
\end{proposition}
\begin{proof}
On the boundary, every matrix
\(K_Y+t(K_X-K_Y)=K_Y[\one+tK_Y^{-1}(K_X-K_Y)]\),
\(0\leq t\leq1\), is invertible by the Neumann estimate.
The argument principle applied to its determinant preserves the number
of interior zeros counted with algebraic multiplicity along this
homotopy. That number is one at \(t=0\), so it is one at \(t=1\).
The unique zero is simple and lies strictly inside the disk. The rate
bound follows from \(|\Im(z_X-z_Y)|\leq|z_X-z_Y|\).
\end{proof}

Here \(E_K\) may be supplied by
Corollary~\ref{cor:OQED-aposteriori-Schur}, provided its common-domain
inverse hypotheses exclude electromagnetic poles throughout the disk.
The approximate pencil's root and its isolation must themselves be
certified; an estimated slope or a small residual at one point is not
the contour test. This proposition bounds the \emph{finite} root
difference without a first-order replacement. It is distinct from
\eqref{eq:OQED-pole-error-bound} and does not assert an
infinite-dimensional determinant theorem. Root clusters, branch singularities
intersecting the disk, and unproved continuum reductions require
separate spectral arguments.

\paragraph{Endpoint matching is not baseline-mode tracking}
The proposition matches roots through \(K_Y+t(K_X-K_Y)\) at the
physical endpoint \(\lambda_{\rm p}=1\). This is not the pole family
\eqref{eq:OQED-pole-bookkeeping-family}. A concrete counterexample,
in fixed reference units, is
\[
 \begin{gathered}
 K_0(z)=z^2-1,\qquad K_Z(z)=z^2-q_Z,\\
 q_Y=-1-\ii\delta_{\rm ex},\qquad q_X=-1+\ii\delta_{\rm ex}.
 \end{gathered}
\]
Take \(\delta_{\rm ex}=0.01\), and track each affine family
\(z^2-[1+\lambda_{\rm p}(q_Z-1)]\) from the baseline root \(+1\).
The two scalar radicands stay nonzero along their respective affine
paths. Starting at (+1) therefore selects the principal square
root along each path. The endpoints are \(z_Y=\sqrt{q_Y}\) and \(z_X^{\rm tracked}=\sqrt{q_X}\),
near \(-\ii\) and \(+\ii\), respectively. Nevertheless, on the
circle \(|z-z_Y|=r=0.05\),
\[
 \frac{\|K_X-K_Y\|}{|K_Y|}
 \leq\frac{2\delta_{\rm ex}}{r(2|z_Y|-r)}<0.206<1.
\]
The enclosed \(X\) root is therefore \(-\sqrt{q_X}\), not
\(z_X^{\rm tracked}\). With baseline \(+1\), the matched endpoint
has a positive rate-comparison quantity while the tracked \(X\) branch
has the opposite sign. This is an auxiliary matrix-pencil example, not an O--QED model. It
does not contradict the root-count theorem; it shows why physical mode
identification is an additional input.

One sufficient identification certificate uses a fixed positively
oriented simple contour \(\Gamma_0\). Its interior contains the endpoint
disk and exactly one baseline zero, the simple root \(z_*\). Assume
all pencils are holomorphic on a neighbourhood of the closed interior
and \(K_0\) is invertible on the contour. Require
\begin{equation}
 \sup_{z\in\Gamma_0}\|K_0(z)^{-1}\|
       \max\{\|\delta K_X(z)\|,\|\delta K_Y(z)\|\}<1.
 \label{eq:OQED-common-baseline-contour}
\end{equation}
The boundary then stays invertible for the entire two-parameter family
\(K_0+\lambda_{\rm p}[(1-t)\delta K_Y+t\delta K_X]\),
\((\lambda_{\rm p},t)\in[0,1]^2\). The argument principle gives
one simple interior root throughout. Uniqueness identifies both endpoint
roots with the continuations of the same baseline mode. A less restrictive
certified contour family may be used, but it must establish that same
identification, not only endpoint root counts.

\begin{corollary}[A contour certificate for the finite pole remainder]
\label{cor:OQED-pole-Taylor-remainder}
On the same finite retained matrix space, let \(K_0\) and \(\delta K\)
be holomorphic near \(\overline{\mathbb D}(z_*,R_{\rm p})\),
with \(R_{\rm p}>0\).
Suppose \(\det K_0\) has precisely one zero in this disk, the simple
zero \(z_*\), and none on its boundary \(\Gamma_*\). Assume a
certified constant \(0<\vartheta_{\rm p}<1\) satisfies
\begin{equation}
 \sup_{z\in\Gamma_*}\norm{K_0(z)^{-1}\delta K(z)}
                    \leq\vartheta_{\rm p}.
 \label{eq:OQED-complex-pole-margin}
\end{equation}
The pole family \(K_{\lambda_{\rm p}}=K_0+\lambda_{\rm p}\delta K\)
has one simple interior root \(z(\lambda_{\rm p})\), holomorphic for
\(|\lambda_{\rm p}|<\vartheta_{\rm p}^{-1}\), with
\(z(0)=z_*\). Writing \(z(\lambda_{\rm p})-z_*
=\sum_{n\geq1}c_n\lambda_{\rm p}^n\), one has
\begin{equation}
 \begin{aligned}
  |c_n|&\leq R_{\rm p}\vartheta_{\rm p}^{\,n},\\
  \left|z(1)-z_*-\sum_{n=1}^{N}c_n\right|
     &\leq\frac{R_{\rm p}\vartheta_{\rm p}^{\,N+1}}
                    {1-\vartheta_{\rm p}},\\
  |\Gamma^{\rm fin}-\Gamma^{(1)}|
     &\leq\frac{2R_{\rm p}\vartheta_{\rm p}^{\,2}}
                    {1-\vartheta_{\rm p}}.
 \end{aligned}
 \label{eq:OQED-pole-remainder-certificate}
\end{equation}
Here \(N\geq1\), \(c_1=\dot z\), and the rates use the same
baseline and pole bookkeeping parameter as
\eqref{eq:OQED-pole-bookkeeping-family}. The radius and margin are
certification data, not additional theory parameters.
\end{corollary}
\begin{proof}
For every \(|\lambda_{\rm p}|<\vartheta_{\rm p}^{-1}\), the
Neumann estimate keeps the boundary inverse well defined. Homotopy
from zero preserves the number of interior zeros with multiplicity,
so there is exactly one, and it is simple. The argument principle
also gives the root itself:
\begin{equation}
 z(\lambda_{\rm p})=\frac1{2\pi\ii}\oint_{\Gamma_*}
 z\,\Tr\!\left[K_{\lambda_{\rm p}}(z)^{-1}
                    \partial_zK_{\lambda_{\rm p}}(z)\right]\dd z.
 \label{eq:OQED-pole-contour-moment}
\end{equation}
The contour integrand is locally uniformly holomorphic in
\(\lambda_{\rm p}\), so the root is holomorphic throughout that
disk and remains within distance \(R_{\rm p}\) of \(z_*\).
Cauchy's coefficient estimate on \(|\lambda_{\rm p}|=r_{\lambda}<\vartheta_{\rm p}^{-1}\) gives
\(|c_n|\leq R_{\rm p}r_{\lambda}^{-n}\). Letting
\(r_{\lambda}\uparrow\vartheta_{\rm p}^{-1}\) and summing the
geometric tail at \(\lambda_{\rm p}=1\) proves the first two
bounds. The rate bound follows by taking the imaginary part.
No inference from a real-parameter estimate to complex analyticity is
used: the affine pencil and boundary inverse were controlled for
complex \(\lambda_{\rm p}\) first.
\end{proof}

For \(\delta K=\delta K_X\), the unknown exact correction may be
bounded on \(\Gamma_*\) using
\[
 \norm{K_0^{-1}\delta K_X}
 \leq\norm{K_0^{-1}}\bigl(\norm{\delta K_Y}+E_K\bigr).
\]
If this proves \eqref{eq:OQED-complex-pole-margin}, and the same
response certificate controls the baseline point, then
\begin{equation}
 |\Gamma_X^{\rm fin}-\Gamma_Y^{(1)}|
 \leq2\Lambda_*E_K(z_*)+
       \frac{2R_{\rm p}\vartheta_{\rm p}^{\,2}}
                     {1-\vartheta_{\rm p}}.
 \label{eq:OQED-combined-finite-rate-error}
\end{equation}
This is one explicit, possibly conservative route to the previously
separate pole-remainder requirement. A margin below one guarantees a
bound, not necessarily a small bound. In the \(e^{-\ii zt}\)
response convention, for an unstable baseline,
positive \(\Gamma^{\rm fin}\) means reduced growth; strict decay of
this mode requires \(\Gamma^{\rm fin}>2\Im z_*\), not merely
\(\Gamma^{\rm fin}>0\).

\paragraph{A counted root and its observable residue}
A simple determinant zero is a pole of the matrix inverse, not of every
source/readout component. Let \(K(z_0)v=0\),
\(w^\dagger K(z_0)=0\), and
\(d_{\rm p}:=w^\dagger K'(z_0)v\ne0\) at a simple zero.
The local resolvent form is \cite[Theorem~2.4]{Beyn2012}
\begin{equation}
 K(z)^{-1}=\frac{v w^\dagger}{d_{\rm p}(z-z_0)}+H(z),
 \qquad H\text{ holomorphic near }z_0.
 \label{eq:OQED-readout-resolvent}
\end{equation}
For completeness, the adjugate formula gives a rank-one pole
coefficient proportional to \(v w^\dagger\); Jacobi's formula for
\(\partial_z\det K\) fixes its coefficient to \(d_{\rm p}^{-1}\).
For holomorphic input and output maps on that neighbourhood,
\begin{equation}
 \begin{aligned}
 \mathcal O(z)&=L_{\rm out}(z)K(z)^{-1}L_{\rm in}(z),\\
 \mathop{\rm Res}_{z=z_0}\mathcal O
   &=\frac{[L_{\rm out}(z_0)v]
                   [w^\dagger L_{\rm in}(z_0)]}{d_{\rm p}}.
 \end{aligned}
 \label{eq:OQED-observable-pole-residue}
\end{equation}
Both factors must be nonzero for this readout to contain the pole.
For example, \(K=\diag(z-z_0,1)\),
\(L_{\rm out}=(0,1)\), and \(L_{\rm in}=(0,1)^T\) give
\(\mathcal O=1\), despite the simple determinant zero. Thus root
location, continuation of the baseline mode, and its nonzero overlap
with the selected source and readout are distinct certificates. This
qualification changes neither the internal vectorial prescription nor
the pole of the full retained inverse. The maps are physical response
maps, not the carrier characters. The residue in
\eqref{eq:OQED-readout-resolvent} is not in general an idempotent:
for \(K(z)=2(z-z_0)\) it is \(1/2\). It must not be confused
with a normalized spectral projector.

\begin{proposition}[Observable residue error]
\label{prop:OQED-observable-residue-error}
Suppose the finite matrix pencils \(K_X,K_Y\) are holomorphic near
and inside a positively oriented rectifiable simple contour \(\Gamma\),
and each has exactly one interior determinant zero, both simple.
Use the same holomorphic input and output maps in
\(\mathcal O_Z=L_{\rm out}K_Z^{-1}L_{\rm in}\), \(Z=X,Y\),
with no other interior singularities. In compatible operator norms,
assume the following certified uniform boundary bounds
\[
 \begin{gathered}
  \|K_Y^{-1}\|\leq B_\Gamma,\qquad
  \|K_X-K_Y\|\leq\epsilon_K,\qquad
  q_\Gamma:=B_\Gamma\epsilon_K<1,\\
  \|L_{\rm out}\|\leq L_{{\rm out},\Gamma},\qquad
  \|L_{\rm in}\|\leq L_{{\rm in},\Gamma}.
 \end{gathered}
\]
Let \(\ell(\Gamma)\) be the contour length and
\(\mathcal R_Z^{\rm obs}\) the residue of \(\mathcal O_Z\) at
its enclosed zero, allowing this residue to vanish. Then
\begin{equation}
 \begin{aligned}
  \|\mathcal R_X^{\rm obs}-\mathcal R_Y^{\rm obs}\|
      &\leq E_{\rm obs},\\
  E_{\rm obs}&:=\frac{\ell(\Gamma)}{2\pi}
       L_{{\rm out},\Gamma}L_{{\rm in},\Gamma}
       \frac{B_\Gamma^2\epsilon_K}{1-q_\Gamma}.
 \end{aligned}
 \label{eq:OQED-observable-residue-error}
\end{equation}
A certified lower bound
\(\|\mathcal R_Y^{\rm obs}\|\geq\underline r_Y>E_{\rm obs}\)
therefore proves that the exact readout contains its enclosed pole.
\end{proposition}
\begin{proof}
Put \(E=K_X-K_Y\). The boundary Neumann estimate gives
\(\|K_X^{-1}\|\leq B_\Gamma/(1-q_\Gamma)\), and
\[
 K_X^{-1}-K_Y^{-1}=-K_X^{-1}E K_Y^{-1}.
\]
Each observable has at most one simple pole inside \(\Gamma\), so
\(\mathcal R_Z^{\rm obs}=(2\pi\ii)^{-1}
\oint_\Gamma\mathcal O_Z(z)\dd z\). Integrating the ordered
resolvent difference proves \eqref{eq:OQED-observable-residue-error};
the reverse triangle inequality proves the nonvanishing assertion.
The two enclosed pole locations need not coincide.
\end{proof}

The lower bound must include any error in evaluating the approximate
residue. When the input or output maps themselves are approximated,
their errors must also be propagated; the displayed bound assumes common
exact maps. Root counts and baseline-mode identification remain separate
requirements. A nonzero residue certifies visibility in this response
channel, not positivity of that residue, a damping sign, or a completed
physical reconstruction.

\paragraph{Certifying a contour bound from finitely many evaluations}
For a practical sufficient test, use the spectral norm on the finite
matrix space and a circle \(\Gamma\) with nodes \(z_j\). Let every
point of \(\Gamma\) lie within arc length \(h\) of a node. Suppose
\(\underline\sigma_j\) are certified lower bounds for
\(\sigma_{\min}(K_Y(z_j))\), including evaluation and arithmetic
errors, and suppose a bound \(L\geq\sup_\Gamma\|K_Y'(z)\|_2\)
has been proved. For a certified uniform boundary error \(E_*\), put
\begin{equation}
 \begin{aligned}
  d_*&:=\min_j\underline\sigma_j-Lh,\\
  E_*&\geq\sup_\Gamma\|K_X-K_Y\|_2.
 \end{aligned}
 \label{eq:OQED-mesh-contour-margin}
\end{equation}
If \(d_*>E_*\geq0\), then the approximate pencil is nonsingular
on the entire circle and
\(\sup_\Gamma\|K_Y^{-1}\|_2\|K_X-K_Y\|_2\leq E_*/d_*<1\).
Indeed, integration of \(K_Y'\) along the intervening arc bounds
the matrix difference by \(Lh\), and the variational definition of
the smallest singular value gives the same Lipschitz bound for it.
This certifies the boundary inequality, not the interior root count,
which remains a separate hypothesis. Finite sampling without the
between-node bound is insufficient: \(1+(z/r)^n\) equals \(2\) at
\(z_j=r e^{2\pi\ii j/n}\) but vanishes between those nodes.

The remaining root count can, for example, be certified independently by
\begin{equation}
 N_\Gamma=\frac1{2\pi\ii}\oint_\Gamma
          \Tr[K_Y(z)^{-1}K_Y'(z)]\dd z.
 \label{eq:OQED-certified-root-count}
\end{equation}
With holomorphy inside and nonsingularity on the contour, this is the
nonnegative integer number of zeros of \(\det K_Y\), counted with
multiplicity. A certified quadrature value \(\widetilde N\) with
error \(\epsilon_N\) establishes \(N_\Gamma=1\) whenever
\(|\widetilde N-1|+\epsilon_N<1/2\). An uncertified near-integer
floating-point value is not a root count. Nor does a count of one make
a numerical centre an exact zero: if a disk of radius \(r\) about
\(z_c\) encloses \(z_X\), while only
\(|z_Y-z_c|\leq\epsilon_z\) is certified, then the resulting
root-distance bound is \(r+\epsilon_z\), not \(r\). The rate
bound inherits the same additional factor \(2\epsilon_z\).

\paragraph{Error from eliminating unretained modes}
All preceding finite-matrix certificates concern the declared exact
retained pencil. A numerical compression of a larger theory requires
an additional reduction estimate. Write the larger inverse-response
pencil as
\[
 \mathscr K(z)=\begin{pmatrix}A(z)&B(z)\\C(z)&D(z)\end{pmatrix}.
\]
Assume \(D(z)^{-1}\) exists on the common domain. The exact retained
pencil is then \(A-BD^{-1}C\), not merely \(A\). An approximation \(\widetilde A\) which discards that complement
must therefore include, in compatible operator norms,
\begin{equation}
 \|A-BD^{-1}C-\widetilde A\|
 \leq\|A-\widetilde A\|+\|B\|\|D^{-1}\|\|C\|
 \label{eq:OQED-retained-complement-error}
\end{equation}
in its total error budget. The blocks and complementary inverse must
be holomorphic on the same disk neighbourhood, with the required
uniform bounds on its contour. Uniform norm bounds alone do not
supply holomorphy.
Near a singular complement, this reduction and its finite-pole
certificate are unavailable until a larger retained space or another
proved reduction is used. No continuum error estimate is supplied by
matrix quadrature alone.

The same elimination must transport the source and readout maps.
For the compatible block input and output
\(J=(J_r,J_c)^T\) and \(L=(L_r,L_c)\), define
\(S=A-BD^{-1}C\). Wherever the inverses exist, block factorization
gives the exact response identity
\begin{equation}
 \begin{aligned}
  L\mathscr K^{-1}J
     &=L_{\rm eff}S^{-1}J_{\rm eff}+L_cD^{-1}J_c,\\
  L_{\rm eff}&=L_r-L_cD^{-1}C,\\
  J_{\rm eff}&=J_r-BD^{-1}J_c.
 \end{aligned}
 \label{eq:OQED-Schur-source-readout}
\end{equation}
Indeed, solving the complementary row gives
\(x_c=D^{-1}(J_c-Cx_r)\), followed by
\(Sx_r=J_{\rm eff}\); inserting both into the output proves the
formula, with all products in their displayed order. If the blocks,
complementary inverse, sources, and readouts are holomorphic, the last
term is holomorphic and the pole residue comes from the two effective
endpoint maps. Keeping only \(L_r,J_r\) is valid only under conditions
that remove the omitted terms, such as \(L_c=J_c=0\); controlling the
retained pencil alone does not justify that omission. The identity is a
block-response statement, not an additional microscopic O--QED action.

The strict \(H_{k^2}\) status, the low-wavenumber transfer test, and the
finite-band \(H_{T,k^2}\) match in Ref.~\cite{ChiapponeIgnition} are
not consequences of UV finiteness. Neither is incremental passivity or
the sign of the transferred rate. Only after those separate tests can
its reduced law \(\Gamma_{\rm OQED}^{(0)}=e^2k^2
\mathcal K_{r_{\rm resp}}(\ell_{\rm dyn};x,\beta)\) be used.
The laboratory coordinate \(\ell_{\rm dyn}\) is not \(a\),
\(\ell_\Omega\), or \(\ell_{\rm GT}\). No fusion-yield or calibrated
attenuation number is inferred from the kernel certificate.

\begin{table*}[t]
\centering
\small
\setlength{\tabcolsep}{5pt}
\renewcommand{\arraystretch}{1.15}
\caption{Fixed-depth tests of the standalone O--QED prescription. The isolated Ward bubble and the grouped
coefficient in the specified contact test have nonintegrable integrands
under their stated hypotheses. The table records the available estimates
and the checks that remain for a completed interacting theory.}
\label{tab:OQED-certification-status}
\begin{tabularx}{\textwidth}{@{}>{\raggedright\arraybackslash}p{0.21\textwidth}
 >{\raggedright\arraybackslash}p{0.37\textwidth}
 >{\raggedright\arraybackslash}X@{}}
\toprule
Object & Established in this article & Scope or outstanding requirement\\
\midrule
Normalized square-root kernel & All finite inverse moments; positive-depth derivative and explicit selected-line tail estimates & Properties of a specified virtual kernel, not of completed vertices\\
Selected Euclidean graph & Absolute convergence when the full routing, vertex envelope, line moments, and IR bounds meet the stated theorem & These hypotheses are not derived for every graph of the O--QED action\\
Dressed lines with local bare vertices & A finite tensor with a nonzero small-momentum Ward remainder & Not the Ward-completed polarization or a physical photon mass\\
Ward-compatible auxiliary vertices & The indicated vacuum bubble has a nonintegrable contraction at fixed nonzero transfer & Its separately defined contact-completed sum requires a new common-integrand estimate\\
Constant-background contact test & Leading bubble growth cancels, but the declared trace-log Hessian is still not absolutely integrable & Tests an additional auxiliary realization, not every possible completed contact sector\\
Physical amplitude and response & Conditional source, residue, contour, and error bounds & Positive physical reconstruction, cut matching, and a completed medium response remain separate\\
\bottomrule
\end{tabularx}
\end{table*}

\subsection{Scope of the fixed-depth tests}
\label{sec:OQED-status}

For the supplied O--QED profile, normalization, complete monotonicity,
all inverse moments, fixed positive-depth derivatives, their singular
zero-depth and strong-depth limits, and the bare one-loop integrability
test are explicit calculations. The optimized selected-line theorem
also certifies appropriate massless selected denominators provided
the remaining complete graph factors obey its anisotropic bound. For
\(M>0\), the additional local analytic kernel hypothesis is also
verified. The Abelian BRST differential and the source action's sector
and carrier conventions are matched to this article.

A complete standalone theory still requires vertices and contacts that
meet the graphwise or grouped estimates, together with a physical source
realization and state-dependent reconstruction. The Ward-compatible
vacuum bubble in
Proposition~\ref{prop:OQED-Ward-bubble-obstruction} and
Corollary~\ref{cor:OQED-fixed-transfer-obstruction} fails the
individual-graph test, despite the verified kernel moments. The additional
constant-background contact test cancels its leading zero-transfer
term but fails the grouped absolute bound as well. Its raw reference
difference remains nonintegrable for every positive radial multiplier
tending to zero under the same shift, trace-log, and measure assumptions.
This is a specified auxiliary realization, not an evaluation of the
source manuscript's full completion. The scalar tube, its signed real-invariant and
complex Taylor bounds, and its finite zero-depth \emph{subtracted}
comparison are independent results. The first-order and finite-pole
error certificates are conditional bounds on the explicitly stated
analytic and operator domains. Endpoint root enclosure
requires baseline-branch identification; the complex pole-family
margin also supplies an explicit finite-remainder bound. A selected
source/readout channel must have nonzero residue to contain that pole;
\eqref{eq:OQED-observable-residue-error} supplies a sufficient quantitative
test when its contour data are available. These results constrain completions of the standalone fixed-depth
prescription. Applying any one of them to the GCM-induced theory
requires deriving its actual vertices, quantum density and boundary
response and identifying the graph and routing to which the estimate
applies. Table~\ref{tab:OQED-certification-status} records the standalone
tests. Sections~\ref{sec:pe-graphs} and~\ref{sec:pe-curvature} address the
corresponding action-derived estimates.

\section{Consequences for PE-CGG-HOUQFT and comparison with QED}
\label{sec:interpretation}

\subsection{The meaning of graphwise ultraviolet finiteness}
\label{sec:meaning-UV-finite}

The convergence claim concerns each finite Euclidean loop graph whose
complete integrand satisfies a stated criterion. Suppose the hypotheses of
Theorem~\ref{thm:global-uv} hold and the integrand of
\eqref{eq:graph-majorant} bounds \(\norm{\mathscr I_G}\) pointwise.
Then the tensor integrand \(\mathscr I_G\) satisfies
\[
 \sup_{p\in\mathcal K}
 \int_{\RR^{4L}}\frac{\dd^{4L}k_{\bb}}{(2\pi)^{4L}}\,
       \norm{\mathscr I_G(k,p)}<\infty .
\]
The support, branch, positive kernel depths, and infrared prescription
are fixed. The integral covers all loop momenta. Absolute convergence
defines the graph without a UV-divergent subtraction. The integrable
majorant also justifies the stated exchanges of integrations over loop
momenta and parameters. Uniformity refers to the specified compact set
of Euclidean external data. Tree exchange poles, physical thresholds,
and on-shell infrared limits require their own analysis.

The estimate depends on where the damping acts and how much of it
survives the vertices. Theorem~\ref{thm:optimized-selected-lines}
allows separate moment budgets for lines that span the loop directions,
including admitted massless denominators. Proposition~\ref{prop:residual-coercivity}
allows growing vertices when enough damping remains in the complete
integrand. The forest criterion in Theorem~\ref{thm:forest} also
checks regions where only a subset of loop momenta becomes large.
These criteria apply to the actual numerator, denominators, and
insertions of the calculation in question.

For one declared completion, applying the estimates to every graph
gives graphwise finiteness at each finite loop order as in
Corollary~\ref{cor:all-orders}. Each finite collection then has a
common bound on its admitted compact external set. Summing infinitely
many terms at one order, integrating over families of supports, and
summing over orders require additional bounds. The Lean proof inventory
in \ref{sec:lean-verification} records the checked implications and
their premises.

\subsection{Applying the standalone O--QED tests to an induced source}
\label{sec:interpretation-OQED}

The normalized square-root kernel has every positive inverse moment
and its representing measure has no atom at zero, as shown in
Proposition~\ref{prop:OQED-moment-certificate}. With conventional local
trilinear comparison vertices, the one-loop polarization, electron
self-energy, electron--photon vertex, and fermion box all have enough
ultraviolet damping. The massive fermion and the local integrability
of the massless photon denominator supply the infrared checks for
these off-shell Euclidean examples. The explicit scalar two-loop
comparison in \eqref{eq:two-loop-selected-example} independently
tests the selected-line bound and its numerical error estimates.

The polarization example shows why the vertices matter. With bare
vertices, the damped tensor is finite but has the nonzero Ward remainder
in \eqref{eq:OQED-explicit-Ward-obstruction}. With the two vertices
required to obey the dressed Ward identities, inverse propagator
factors grow and offset the damping. The isolated two-vertex term
is nonintegrable at zero transfer under the continuity hypothesis of
Proposition~\ref{prop:OQED-Ward-bubble-obstruction}. The failure
persists at fixed nonzero transfer under the hypotheses of
Corollary~\ref{cor:OQED-fixed-transfer-obstruction}.

The specified constant-background contact term cancels the leading
zero-transfer growth and leaves a nonintegrable remainder.
Corollary~\ref{cor:OQED-general-contact-obstruction} extends this
obstruction for the raw reference difference to positive radial
multipliers that are \(C^2\) at large momentum and tend to zero.
It keeps the same shift, trace-log, species, and measure assumptions.
This result identifies an obstruction that any proposal retaining
that realization must address. A different completion has its own
action, measure, and integrability tests.

Lemma~\ref{lem:OQED-grouped-convergence} gives a criterion for
the combined coefficient. Derive a finite group of terms from one
completed action and express them on a common integration domain,
retaining their signs, multiplicities, and routing Jacobians.
The pointwise sum must have an integrable bound for the full tensor.
The Ward identity provides a separate check on its contraction.
Table~\ref{tab:OQED-certification-status} summarizes which of these
checks have been carried out.

\paragraph{Transporting a graph estimate between the two actions}
There is a precise comparison when the two sources have been reduced
to the same fixed support, finite external data, graph, regulator
$\Lambda$ and loop variables $k$. Include the Jacobian of any change of
variables in each density. Write $I^{\rm ind}_\Lambda$ for the
GCM-induced tensor integrand, $I^{\rm test}_\Lambda$ for the specified
comparison and $E_\Lambda=I^{\rm ind}_\Lambda-I^{\rm test}_\Lambda$.
If $I^{\rm test}_\Lambda$ and $E_\Lambda$ converge locally in $L^1$
and each has uniformly vanishing $L^1$ tails as $|k|\to\infty$,
then $I^{\rm ind}_\Lambda$ has a global $L^1$ limit. Indeed the
triangle inequality on a ball gives the local limit, and on its
complement it bounds the induced tail by the sum of the two tails.
This criterion requires control of the complete difference, including
the determinant and contact terms generated by the electromagnetic
sources.

There is also a converse obstruction at one hard scale. Let
$D_j=\{2^j\le |k|<2^{j+1}\}$, with a fixed infrared prescription.
If $\int_{D_j}\|I^{\rm test}\|\ge c>0$ for every sufficiently
large $j$ and $\int_{D_j}\|E\|\to0$, then
$\int_{D_j}\|I^{\rm ind}\|\ge c/2$ eventually by the reverse
triangle inequality. The induced graph is then divergent. Thus a
UV-integrable correction to a genuinely divergent graph cannot
repair it in this common representation. When the two actions
have different elementary graphs, the comparison must first identify
a coefficient of the same source derivative on a common integration
domain. The resulting statement concerns that coefficient. Finiteness
obtained by canceling multiple independently defined graphs supplies
no individual bound on those graphs.

\subsection{Normal data governing the ultraviolet cost}
\label{sec:update-interpretation}
The normal extension gives the transfer by solving its declared local
boundary problem. With a deterministic conditional measure it then
fixes the complete mixed-source functional; the two-ended covariance
has a one-sided kernel form only on a commuting block. GTL acts on
endpoint coordinates, not on an already assigned positive depth.
Neither late localization nor a support-only action supplies the
normal operator, interaction layer, and selector by itself.

A common fluctuating depth is a field-theory datum, not independent
random damping on every virtual line. Its joint inverse-lapse budget
and exact marginal tail are established in
Section~\ref{sec:update-joint-depth}. The positive winding law with
$B_*>0$ controls every finite polynomial moment order, but actual
induced vertices, contacts, and infrared data still need a common
bound. A saddle approximation at fixed momentum does not determine
the exact ultraviolet rate of the mixed measure.

The physical comparisons are equally specific. The constrained
isotropic free multiplet fails its joint GT-reflection test, and its
unsubtracted quartic gives a negative boundary test at weak positive
coupling. A positive spectral cross completion must carry additional
diagonal weight; for the specified scalar contractions that weight
restores a divergent bubble. The local positive bulk retains normal
fluctuations and a changed boundary spectrum. These successful and
unsuccessful properties belong to different actions; they cannot be
combined by retaining only their preferred propagator entries.
The GT-contour identity, vectorial boundary values, and full source
maps therefore remain tests of one completed physical theory.

\subsection{Regularization, normalization, and the zero-depth limit}
\label{sec:interpretation-renormalization}

The kernel and its positive depth \(a\) are part of the virtual
covariance. For an admitted graph, absolute integrability gives the
same limit for every measurable exhaustion of the integration domain
at that fixed depth. Changing \(a\) changes the integrand. In standard
perturbative renormalization, momentum-space Taylor subtractions can
give an integrable graph coefficient. Zimmermann's forest formula
organizes these subtractions for nested and overlapping divergences
\cite{Zimmermann1969}. The convergence test above applies to the
unsubtracted integrand of the specified covariance.

Finiteness leaves the choice of physical normalization. Fix the
charge through a matching condition on a physical amplitude.
Where the physical interface admits a pole, its position and residue
can receive finite loop corrections and must be matched using the
appropriate source maps. Soft and threshold regions also require
their own estimates.

Switching off the kernel replaces its Laplace--Stieltjes measure
by \(\delta^{\rm heat}_{\bb,0}\). The interpolation in
\eqref{eq:OQED-interpolation-atom} also has a zero atom whenever
the undamped term has nonzero weight. The no-atom inverse-moment
criterion therefore ceases to apply along this path to the ordinary
reference. Equation~\eqref{eq:OQED-matched-increment} uses an
independently defined renormalized QED reference.

Theorem~\ref{thm:OQED-subtracted-zero-depth} proves a specific
zero-depth limit after the common subtraction of the scalar
bubble. The unsubtracted comparison has the logarithmic behavior
in Proposition~\ref{prop:OQED-zero-depth-finite-part}. These two
results show where the finite scalar limit comes from. A limit
of the completed gauge theory would have to control its vertices,
contact terms, insertions, and matching conditions together.

\subsection{Requirements for a comparison with QED}
\label{sec:interpretation-comparison}

For the curvature sector of Section~\ref{sec:pe-curvature}, the
determinant controls the complete depth-integrated coefficients.
The heat-link sector has a direct full-momentum construction at
fixed readout in Theorem~\ref{thm:s6-heatlink-uv} and with one fluctuating
readout in Proposition~\ref{prop:dualrev-uvheat-depth} and
Theorem~\ref{thm:dualrev-uvheat-real}. The fixed-depth O--QED comparison
gives quantitative integration tails under its own hypotheses. Comparing the electromagnetic sector
of PE-CGG-HOUQFT with QED requires a physically normalized current
observable of the induced action. Its measure and normalization must
come from the same source as its interaction vertices.

Equation~\eqref{eq:OQED-matched-increment} specifies a comparison
with separately renormalized QED. Use common physical masses, charge,
and source normalization. Evaluate the same observable for the same
state and external kinematics at the same perturbative order.
For scattering, also match the infrared prescription. If both
completed polarization tensors exist and satisfy their Ward identities,
their difference is transverse. Its size and sign still have to be
calculated.

A first target is a vacuum response to a conserved external
current. Fix the physical normalization at a reference point
and calculate the response at other specified transfers using the
same probe in both models. To connect the result to an experiment,
identify the measured rate or response, include the source
normalization and any soft contributions it requires, and propagate
the calculation and parameter errors. A claimed discrepancy must
be resolved at that accuracy.

Proposition~\ref{prop:pure-redefinition-cancellation} gives a
useful check on the comparison. In its finite-mode construction,
an invertible change of field variables can introduce damping
into the quadratic term while its induced vertices, sources,
and Jacobian restore the original result. To find a physical
difference one must calculate with the full action and normalized
sources.

\subsection{The remaining calculation on the physical GCM branch}
\label{sec:interpretation-priorities}

Postulate~$\Omega_\alpha$ requires the physical graph expansion to come
from one normalized source on the chosen stationary branch. Its
superconnection, curvature pairing, cyclic history trace, BV measure
and BFV boundary state determine both the charged inverse and the photon
derivatives of the closed determinant. The explicit comparison sources
give tests for this derivation.

The heat-link source illustrates the need to keep both parts of the
calculation. A physical projection that kept its complete source would
inherit its ultraviolet estimates and its nonreal charged poles.
If the chiral projection changes the charged covariance or removes
the spectral pairing that makes the Euclidean determinant positive,
its inverse, determinant phase and photon contacts must be evaluated
again with the physical boundary state. Two returns can each suppress
a hard momentum yet lose that suppression on a shared opposite-momentum
channel, as the two-return calculation on p.~\pageref{dualrev-uvb-shared} shows. Matching
the classical density or the free photon transfer fixes only the
corresponding terms of this quantum source.

\paragraph{Determine the weight of identity and return histories}
Compute the charged determinant in the sector with zero visits and
its electromagnetic derivatives. For each unreduced return of positive traveled length, including a
contractible return with zero net winding, evaluate the normal
integrations with the actual vertices and endpoint pairing in place.
The outer signs and the vectorial pole belong to the fixed sector-plane
and ordering branch. The positive normal length carries neither of
those signs. Evaluate also the complete
normalization as a function of the shared depth. A determinant from
the fluctuations can suppress short visits, or other fields can change
its sign in the logarithm. Dividing it out separately at each depth
changes the quantum source when that weight varies with depth.

\paragraph{Resolve the complete one-loop polarization}
Start with the vacuum response at nonzero external transfer.
Compute the action's bubble, contact, and any measure contributions
with one routing. For the dressed bubble of
Corollary~\ref{cor:OQED-fixed-transfer-obstruction}, an absolutely
integrable sum requires the other double contractions to cancel the
profile in \eqref{eq:OQED-fixed-transfer-obstruction} if their sum has
a continuous limit uniform in direction. Then bound the full
remainder. In four loop dimensions, a local integrable
envelope together with a uniform \(O((1+|k|)^{-4-\delta})\)
tail, for some \(\delta>0\), is one sufficient test. Check the Ward
identity for the resulting finite tensor. If the native theory contains
a gauge-root block of the type in Section~\ref{sec:pe-curvature},
calculate its magnetic two-point coefficient, including the normal
polarization and ghost, on the same background. This tests a gauge
loop that a visit confined to charged matter need never encounter.

\paragraph{Extend the estimates and reconstruct physical amplitudes}
At higher orders, treat all ultraviolet subregions, the relevant
infrared limits, and the insertions required by the Slavnov--Taylor
and Nielsen identities. The same action must also meet the
physical reconstruction requirements in Section~\ref{sec:analytic}.
These include reflection positivity and the other OS hypotheses for
the full physical correlations, the external source and pole conditions,
and the matching of prescribed cuts to analytic discontinuities.
Assumptions~\ref{ass:OS-completion}
and~\ref{ass:cut-matching} specify separate checks.

\paragraph{Evaluate a matched observable}
Evaluate the difference for the same physical channel in both
theories. Where the observable depends on a retained pole, the
response and residue estimates bound its errors under their stated
source and contour hypotheses. A medium calculation also requires
its own state, transition kernels, and retarded response as in
Section~\ref{sec:OQED-medium}.


\section{Conclusions}
\label{sec:conclusions}

On the surface selected by the free-slope variational problem,
the curvature action fixes the leading quartic coefficient of the normal
metric. It thereby fixes, for each prepared $j\ne0$, the coefficient of
$r^{-3}$ in the exponential decay of the ground-state density. In the
homogeneous displacement representation the heat lapse is proportional
to proper radial distance. The charged log determinant then grows as
$r^{-2}$ times a cubic logarithm. The difference of these rates pays
for the complete charged vacuum and every finite set of contact and
photon insertions (Corollary~\ref{cor:geo-free-slope-uv}). Smoothness
of the unexpanded partition uses a further coefficient estimate.

For exact power profiles, the ultraviolet question reduces to
an integral of one specified source over the full reciprocal lattice
and a common radial state. Each charged loop supplies a Gaussian
momentum estimate from a marked Wilson vertex. Each internal photon
carries the harmonic readout twice. At fixed volume, for a radial
profile \(r^d\), heat lapse \(r^\alpha\), and nonzero angular number
\(j\), the determinant-weighted free state normalizes exactly when
\(d-1>2\alpha\). It then has all inverse radial moments. Every fixed-order
bare graph is absolutely finite, including in mixed hard-momentum and
short-duration regions. The cubic logarithm in the charged vacuum
excludes equality at this threshold.

The real-coupling partition integrates the ordinary charged Fredholm
determinant against both the Maxwell probability measure and the radial
boundary density. The field-independent mass reference leaves the full
depth-dependent vacuum in this integral. If \(d-1>3\alpha\), the boundary
state controls every finite derivative in the real charge and the
declared smooth sources. At $d-1=3\alpha$ the corresponding
coefficients must be compared. At exact coefficient equality the
ground-state prefactor still controls the vacuum and the responses
to smooth transverse shifts. The first two charge derivatives exist for
every admissible heat exponent there. The Maxwell average gives all finite
charge derivatives there for $\alpha>6/7$
(Proposition~\ref{prop:uvheat-critical-charge}). This includes the
quartic chart $(d,\alpha)=(4,1)$. At $\alpha=6/7$ all bare coefficients
still agree with the asymptotic expansion of the exact source at zero
charge (Corollary~\ref{cor:uv-geom-charge-origin}). The restriction to finite carrier
volume matters: with a single radial preparation held fixed as the
torus expands, the logarithm of the free normalization grows faster
than the volume.

The test of the photon soft limit distinguishes two of the finite
representations. The first-order Dirac realization has a finite
transverse constant at order $e^2$
(\eqref{eq:g5M-soft-mass}). On the fixed-distance continuum vacuum chart of the heat-link
realization, the source contacts and exchange terms together give
a soft factor on each photon leg (\eqref{eq:s6-eachleg}). The charged covariance has its own spectral obstruction. Its nonreal
free poles (\eqref{eq:s6-freepole}) preclude the standard positive
spinor spectral representation for this covariance. Averaging over the
vacuum-weighted radial state also leaves the open vector coefficient
too small to have a nonzero positive mass-spectral measure
(\eqref{eq:geom-uv-open-vector-rate}).

A raw local Dirac loop has the logarithm in
\eqref{eq:uvpeA-raw-log}. An unmarked gauge-root loop also retains a
magnetic logarithm after combining the vector, normal and ghost
contributions (Section~\ref{sec:pe-curvature}). A finite fixed-depth
kernel or a nonzero traveled distance by itself gives only partial
information about these terms. The GCM source must determine which
intermediate visits remain in the determinant and which momentum
directions they control after sewing. The Gaussian criteria in
Section~\ref{sec:pe-graphs} make those questions quantitative.

The scalar normal action illustrates how to check every short-duration
face without imposing a common scaling. Its charged Green kernels
confine differences of radial depths. Its neutral kernels respond to
the depths shared by successive visits. Contracting the charged
intervals gives the subset inequalities in
Theorem~\ref{thm:uv-normal-word}. They prove finiteness at each
positive order for curvature insertions with zero angular mean.
When instead the source varies the flat angular connection, its
neutral two-source contact is singular and the complete angular
sum retains the singularity. This locates the physical question in
the source derivative and boundary measure that the GCM action
assigns to the measured photon. For the scalar model the complete
relative determinant can also be tested before any expansion in the
normal holonomy. At every fixed nonintegral holonomy the short-time
trace excess is positive and proportional at leading order to
$t^{-3/8}$ (Proposition~\ref{prop:quartic-holonomy-uv}). The test
therefore detects sensitivity of the determinant itself, in addition
to failure of the contact derivative. Keeping the normal holonomy
fixed still leaves the horizontal magnetic test. For a locally
minimally coupled scalar the quadratic coefficient per normal
path is independent of its horizontal diffusion time. A winding
on the collapsing circle can carry that coefficient at arbitrarily
short durations. For the specified marked winding trace the sharp
threshold is $p>3$ for a radial insertion $r^p$
(Proposition~\ref{prop:horizontal-magnetic-word-uv} and the
calculation following it).

The relative Wilson source keeps the ordinary free Dirac pole while
each closed electromagnetic graph retains a marked displacement. Its
interacting log determinant is bounded above by $C r^{-2}$ on the
minimizing surface. The radial density decays as
$r^4\exp(-A_* r^{-3})$, so every real coupling is covered at fixed
volume. The continuum one-loop transverse constant changes sign with
the mass. At least one mass makes it vanish after averaging the
specified mass-independent geometric depth. At fourth order the variance of the
conditional response survives inversion of the full photon covariance.
The large-volume lowest-mode test detects it even at a matching mass.
The determinant-averaged charged line also fails the positive energy
measure test at nonzero charge. These conclusions concern the specific
common-depth action and its retained field. A tree-level charge measurement also
fixes where a proposed normalization must enter. Scaling the entire
interaction to normalize each charged vertex leaves a nonzero one-loop
constant, while calibrating the coupling inside the transverse Wilson
phase preserves the matching masses.

At fixed geometric depth the other Wilson source gives positive upper
and lower bounds of Maxwell order on the Euclidean susceptibility.
This result holds for the interacting measure in its weak-charge regime
and includes the first nonzero mode on every torus of the stated size.
A positive transfer would turn an invariant Maxwell singularity into
a mass-zero atom. The charged two-point function and the current cut
each require a separate spectral test on the proposed observable.
In the time-local polarization model the two positive spectral weights
and the gapless photon frequency can all be calculated exactly. In
the rigid-current model a moving charge instead has continuous
spectral weight down to threshold. A physical charged spectrum need
not have an isolated one-particle pole.

An eliminated charged block can carry a divergent closed loop even
when the retained open propagator is fixed
(Proposition~\ref{prop:visit-cycle-uv}). The quantum boundary overlap
tests which loop phases survive sewing
(Proposition~\ref{prop:phys-uv-boundary-overlap}). A physical cut
also sums over every accessible excitation. In the oscillator example
each individual overlap decays exponentially, yet the inclusive spectral
weight can approach a positive constant. Both calculations require the
boundary state and the charged current selected by the same action.

A remaining physical input can be located in the quadratic action.
The variation of the normal flux with squared carrier frequency
measures exactly the temporal curvature weight sampled by the
readout. Static measurements cannot fix this weight. For the
fixed-width heat link, positive Maxwell scaling persists at weak
charge, but the magnetic obstruction persists at every positive mass in every
bulk limit at fixed weak charge as well. A positive energy spectrum
and a finite bare electric loop impose distinct requirements on a
spatially extended source.

The charged-pair examples make the distinction explicit. A spatial
magnetic interaction gives positive interacting photon and charged
spectral weights with absolutely convergent bare graphs.
A Wilson hopping Hamiltonian with the Gauss energy included admits
positive charged transfer and an interacting radial ground state with
all inverse moments in its nonzero angular sector. The scalar density
loop has finite ordered photon sewings at fixed widths in the small-hopping
regime. At weak electric charge these bounds give stationary positive
photon states with weight near the Maxwell energy even at very small
momenta (Corollary~\ref{cor:closing-uv-pair-bulk}).
The occupied onsite vacuum cancels the hard tail of the single
two-density trace at the integrand level. For the physical chiral
source, the charged operator, its temporal and contact derivatives,
and the boundary determinant must likewise be derived together.

\paragraph{What the physical source must decide.}
The charged heat step fixes a displacement width $\sqrt{\tau(r)}$.
The transverse measurement fixes a second length, obtained by
solving its normal curvature boundary problem. Their ratio can
be measured in the variance of a single charged Wilson chord.
The closing charge can also propagate in different angular states:
its relevant energy is the quadratic form on the pair after the
boundary Gauss constraint and physical chiral projection have
acted. Two metrics can have the same volume element and the same
single-charge barrier at every depth, and nevertheless give
different pair energies
(Proposition~\ref{prop:seam-selection-uv}).

On the ordinary neutral quartic chart with a regular transverse
Maxwell form, Proposition~\ref{prop:seam-closed-uv} and
\eqref{eq:seam-regular-optical-uv} give a divergent connected
photon coefficient. Both external legs probe a transverse field
of nonzero spatial momentum. Giving each occupied pair a nonzero
normal angular charge makes this entire two-hop response finite at
every electric order (Proposition~\ref{prop:seam-switched-pair-uv}).
Two pairs can nevertheless exchange partners during three short
intervals and enclose a spatial area that a transverse probe resolves.
Their connected coefficient at $g^4e^{120}$ diverges
(Proposition~\ref{prop:seam-uv-two-pair}). The powers of $g$ here
refer to the joint expansion in hopping and charge.
In the optical representation of
Theorem~\ref{thm:seam-calibrated-joint-uv}, the related
geometric lengths, the charged determinant and the onsite boundary
state give absolute convergence at each finite electric order
even when the depth moves in the neutral angular sector.
Corollary~\ref{cor:seam-optical-shifts-uv} proves the same result when
each creation and annihilation changes the angular state. The
four-hop exchange is included at arbitrary finite hopping strength
in this realization.
No ultraviolet cutoff or counterterm enters either calculation.
The singular normal coefficients in the explicit optical
realization are defined by its weighted Dirichlet form. A
regular, uniformly elliptic constitutive tensor at the same tip
has the different boundary response just calculated.

The original onsite action therefore serves as a controlled
UV-finite comparison. For a charged action with a dispersive
reference, Theorem~\ref{thm:carrier-flight-dressed-uv} instead
retains all four-momentum integrations at arbitrary finite pair
hopping and every finite electric order on the stated slab and
torus. Its kinetic jump samples the gauge field on an unreduced
normal excursion. At zero electric charge, Pauli blocking gives
the neutral matter threshold of Proposition~\ref{prop:carrier-blocked-uv}
in the moving geometry. The independently sampled mass tower in
Proposition~\ref{prop:carrier-mass-visits-uv} places its normal
energy in the mass shell and restores a finite zero-charge vacuum
logarithm per unit volume. Its bounded equal-time density weight
still fails the boosted-current test even with all channels present.
Assigning a spatially local massive normal
field instead of a common radius also gives extensive connected
coefficients in the separately specified sources of
Proposition~\ref{prop:replacement-local-normal-uv}. These derivations
use different geometric preparations and must be compared at the
level of the complete generating functional.

For unrestricted PE-CGG-HOUQFT the decisive further question is how
the physical photon closes an observed charged loop.
Proposition~\ref{prop:carrier-sewn-uv} gives a negative answer for
a whole class of conventional sewings even after all polynomial
contacts are included. The massless physical photon, a positive
timelike current band and an ordinary Maxwell contraction of that
same Lorentz-covariant response cannot together yield an absolutely
finite bare vacuum graph under the stated dispersion hypotheses.
The curvature coupling of Proposition~\ref{prop:repair-pair-sewing-uv}
illustrates the effect by retaining both polarization contacts in
a concrete action. To go beyond this result one must exhibit an
actual modification of the current or of the complete gluing law
and check the physical cuts of the modified source. A finite
external form factor, a choice of vectorial boundary values, or
a finite amplitude in an independent mass channel does not supply
that modification.

The reference action for the physical result is the full ambient
curvature integral in \eqref{eq:uvpeA-curvature}. Its coarea and velocity
root can be evaluated exactly. For an explicitly chosen transfer
on that very shell, \eqref{eq:nativefull-shell-power-uv}
isolates the cost of a simultaneous return and a short segment.
Inclusive electric sewing recovers the unmarked trace in the
separable example, even though a selected transition is finite.
In the shell-averaged Dirac comparison an angular integration
before the determinant protects the photon momenta and gives an
all-orders result for nonzero winding on every charged loop.
The same calculation leaves the zero-winding magnetic divergence.
Giving both charged incidences their own shell readout in
\eqref{eq:nativefull-three-compressions-uv} controls all charged
circulations and photon momenta of the specified transverse
Euclidean source. Its finite one-loop constant in the photon
inverse propagator fails the masslessness test. Gauging each entire shell displacement supplies the additional
vertex terms. The Wilson determinant then preserves a massless
photon at every finite order, including with the full charged
circle tower, but its free charged endpoint has opposite-sign
residues at two positive energies (Section~\ref{sec:ward-hop-uv}).
The normal Green function in Section~\ref{sec:ward-normal-uv}
gives positive free charged weights together with a perturbatively
massless photon. The gauge-completed heat vertices converge even
when their durations shrink together. Yet the open source predicts
more charged-pair production than is present in the determinant
cut. The equal-helicity forward photon cut independently becomes
negative at high energy. The two-current test also constrains an entire class:
for a single Gaussian charged field with a positive mass atom
and covariance $(m-d_B)F(m^2-d_B^2)$, a positive Stieltjes
spectrum and the ordinary inclusive pair interpretation force
$F(x)=C/x$. The corresponding local Dirac trace restores
the logarithmic bare charge divergence.
The local magnetic deformation of that same determinant has a
positive divergent heat coefficient, and finite minimally gauged
flavor mixing also gives a divergent bare polarization as soon as
one charged pole contributes its full physical pair weight.
Section~\ref{sec:nov-positive-uv} gives a different result for
a non-Gaussian source on a spatial slice. It admits UV-finite
canonical perturbation theory and positive one- and two-photon
production from the same charged pair sectors. Its tests of charge
transport and propagation still depend on the chosen slice.
For a covariant current with positive spectral weight even in one
finite band, Section~\ref{sec:nov-ring-uv} derives a logarithmically
divergent ring for ordinary Maxwell sewing. A native replacement
must give both the continued photon residue and the overlap that
appears when two such loops are closed. Even momentum-dependent
positive normal spectra fail the separate causal test in
\eqref{eq:null-positive-port-uv} when the same polarization
is sewn into the Maxwell trace and continued magnetic response.
The ordered-path calculation in Section~\ref{sec:path-tests-uv}
improves on the local magnetic coupling. Its charged determinant and
contacts give UV-finite Euclidean vertices at every finite order, with
a perturbatively massless photon. The near-threshold cut still falls short
of its own pole-pair weight. Section~\ref{sec:path-rest-tests-uv}
constructs a shell-dependent electric vertex satisfying the real
two-current inequality and gives an explicit gauge completion whose
combined two-mark integral diverges. A second completion makes that
integral finite at nonzero transfer. Its transverse photon coefficient
diverges as the inverse transfer magnitude near zero. For the boundary measure,
Section~\ref{sec:path-closure-tests-uv} traces a divergent fourth-order
coefficient to a contraction within one completed holonomy.
Section~\ref{sec:path-newton-tests-uv} computes how much simultaneous
hard momentum a shared area can suppress.
The normal boundary and the charged photon contacts can now be
differentiated together. Section~\ref{sec:gcmone-transport-uv}
derives their common transport from the shell exit geometry and a
specified gluing law. For the ordered Wilson paths the determinant
and the charged pole give identical pair cuts up to the first heavy
threshold. Above it the deficit is strict in the two-current channel
at every energy. For the virtuality-dependent electric vertex, the
same normal integration leaves a finite Maxwell coefficient and
produces the double logarithm in
\eqref{eq:combined-harmonic-double-log-uv}. Neither tested source
satisfies all of the physical requirements. The cut identities and
the rank test in \eqref{eq:geo-actual-phase-rank-uv} apply to
different stages of the calculation: the former constrain charged
transition probabilities, while the latter measures the independent
hard momenta suppressed by actual relative displacements.

A heat version of the product vertex also gives an absolute bound
on a single charged determinant loop with any finite number of
nonzero external transverse probes. The proof keeps its signed normal
mass difference together and includes the higher Wilson contacts.

The charged measure can be computed explicitly on the exit chart.
Equation~\eqref{eq:join-congruence-uv} gives its exact change of
variables. Restoring the measure reproduces the all-energy light-pair
cuts and exposes the bare Maxwell divergence in
\eqref{eq:join-maxwell-cost-uv}. Matrix-valued exits retain the
trace identity \eqref{eq:nextmatrix-trace-uv}. A separate product of
shell marks gives a positive one-loop density and an analytic soft
response. Its missing inclusive weight in \eqref{eq:nextsoft-positive-uv}
is a strict obstruction even when it becomes small at high energy.
The spectrum of the propagating photon gives an independent hard
constraint in \eqref{eq:sew-dressed-photon-ring-uv}.

The joint charged preparation in \eqref{eq:joint2-normal-uv} yields
finite complete Euclidean graphs at each electromagnetic order and
reproduces all finite normal Dirac cuts at one charged loop. Its
separated-time density test nevertheless fails. A two-current curvature
term restores the dispersion on the represented connection and exposes
the cost of Maxwell sewing. In the neutral channel,
\eqref{eq:joint2-inverse-square-uv} gives an exact pair transition,
positive photon spectrum and finite complete relative response to
nonconstant mixing. An individual resolved two-insertion ring still
has a logarithm. These are distinct tests of the same proposed physical
generator. No one of the constructions here satisfies them all for the
unrestricted charged theory.

There is now a charged comparison with actual fermion endpoints, the correct
rest-frame Born pair weight and a Hamiltonian inner product. On a fixed
torus its ordered bare coefficients have absolute ultraviolet bounds at
each fixed order. The same electric coupling fixes its quadratic photon
contact. After that contact is included, boosting a positive finite-mass
band gives \eqref{eq:next3-dipole-energy-uv}. An independent failure of
local charge composition for marked pair and unmarked elastic Dirac
transitions is quantified by \eqref{eq:next3-charged-algebra-uv}.
In the neutral model the previous question about interacting positivity
of the null field now has a negative answer at fourth order, even though
its perturbative pole is protected. A time-local interaction on a
specified spatial slice passes the full reflection test.

For the angular source these tests now give a definite answer.
The scalar angular determinant has a negative first cut.
Parallel transport of a two-component spin state restores its sign
and threshold power, while the unmarked return and fixed-order
closed-graph bounds survive. A matched curvature strength and radius
can also reproduce the tree Dirac coefficient and magnetic moment.
The full current still obeys
\eqref{eq:four-angular-superconvergence-uv}, which excludes the
required positive spectral band for that determinant. A fluctuating
closed excursion supplies a computable high-frequency contact; its
first charged transition is zero. The Hamiltonian comparison in
Section~\ref{sec:four-excite-uv} has a positive transverse photon
spectral law on the charged Fock space. Its measured elastic current
on a single charge is zero.

The magnetic test places a further restriction on the excursion law.
A finite spectral half-moment is incompatible with a positive
excursion spectrum throughout any interval of static fields.
Equation~\eqref{eq:fifth-excursion-levy-uv} gives a law at the
limiting power that passes the field-parallel test while retaining fixed-order
closed-graph bounds with the spatial A4 readout. The elementary
scalar openings of both laws fail positivity. For the field-perpendicular
current the same limiting law gives double spectral poles. When the
first pole lies in the analytic strip of the other variations it also
obstructs positivity of the complete determinant Hessian
(Section~\ref{sec:six-mag-uv}). Covariant endpoint
factors recover a positive free opening with a source-dependent
sewing metric. The
Hamiltonian of \eqref{eq:fifth-elastic-square-uv} does produce
both measured Born transitions. At sufficiently large boost its
charged energy falls below the magnitude of its momentum. The
second recoil ordering in \eqref{eq:fifth-elastic-symmetric-uv}
identifies a specific missing contribution to covariant exchange.

The physical photon supplies a test independent of the proposed path
representation. When its first occupied massive band is positive at
Born order, the separated trace of its measured field strengths grows
in magnitude at least as the inverse square of distance. A bare
coincident Maxwell insertion built by closing the same legs fails
absolute convergence at that order. If the action pairs curvatures
at a relative displacement, Eq.~\eqref{eq:six-relative-ranks-uv}
specifies the exact joint inverse-moment conditions in a linear
shell model. This requires the relative records to survive both
action variations. In the explicit charged Pauli comparison, coupling
to the canonical normal coordinate preserves a decaying factor even
in the instantaneous kinetic and Coulomb terms. On finite clock
intervals this gives absolute convergence at every fixed order in
that model. Its charged Hilbert space and its normal coordinate were
specified in addition to the underlying curvature monomial.

The new cyclic determinant shows that retaining a relative
displacement can suffice even at a coincident curvature contact.
For the specified rotation-singlet representation the exact first
trace threshold is $d>4$ for the shape parameter of the geometric diffusion
(Section~\ref{sec:seven-radial-uv}); an integer $d$ counts
consecutive transits. An additional harmonic exit from the midpoint of
a radial segment gives the same result with every angular multiplet in
the trace. The proof includes all photon contractions
between different auxiliary cycles. Charging the same auxiliary
inverse by ordinary minimal substitution gives a bare magnetic
logarithm already on the constant radial mode. The independent
particle--hole construction in Section~\ref{sec:seven-pair-uv}
has a measurable charged threshold and a conserved charge, but its
high-boost transverse weight tends to zero on every fixed mass band.
These are two distinct actions and their conclusions cannot be
combined into a single charged theorem.

There are also different roles for a shadow. An auxiliary field
eliminated on trees, a positive intermediate charged state, and a
graded determinant partner have different source derivatives. On the
tested square-root branch the tree prescription loses a positive part
of the forward cut. A dynamical partner can cancel a vacuum coefficient
while leaving a measured transition, as the explicit supersymmetric
oscillator illustrates. Its partner vertices and their boundary
conditions would have to follow from the same physical action.

The quartic normal surface gives a concrete first calculation for
the physical theory. Differentiate the charged operator twice with
respect to the measured photon and resolve the result in angular
modes. A closed word that stays away from angular zero pays the
trace cost derived in Section~\ref{subsec:g5N-normal}. A contact
between two zero modes samples the five-dimensional radial volume
at the tip, where an inverse fifth power already produces a
logarithm. If other terms change this conclusion, their weights and
relative signs must come from that same derivative.

For unrestricted induced O--QED, the remaining calculation starts
with a single charged operator and its sewn boundary state under
Postulate~$\Omega_\alpha$. Differentiate their common generating
functional with respect to the measured photon. When the resulting
legs are closed, the determinant, the direct contacts and the
contractions between distinct charged cycles must satisfy the joint
hard-region bounds of Proposition~\ref{prop:nativefull-sewn-flags-uv}.
When those same legs are opened, the current must produce the charged
pair with its measured energy and transverse spectral weight. The
Slavnov--Taylor, contour and reflection conditions apply on the same
sector plane $\bb=(s,b_s)$ with its prescribed pole direction.
The graphwise results above establish convergence for their specified
models. Deriving both the closed bounds and the open spectrum from
the physical action remains the condition for extending them to the
unrestricted bare theory.

\section*{Funding}
This research received no funding.

\section*{Conflicts of interest}
The author declares no conflict of interest.

\section*{Acknowledgements and use of generative AI}
The author wrote the original research manuscripts underlying this work without using agentic AI. In the summer of 2024, before using AI in this research, the author discovered the geo-phase ultraviolet damping mechanism at the same time that he first fully understood geometric-time locking.

\emph{Author's declaration accompanying the earlier preprint.}

The scientific ideas, calculations, and research program underlying this
manuscript were entirely developed by the author.  The geometric-time coordinate
concept, along with the Multidualverse concept and related mathematics, was formalized in 2017.  The multi-geometric-time coordinate concept,
together with the Omniverse concept and related mathematics, was formalized
from 2020 to 2022. The geometric-time locking concept was not fully understood until 2024, still before any AI-assisted research took place. The author wrote more than a thousand pages on these topics before the emergence of modern large language models.

The author used ChatGPT (OpenAI) for proofreading and language refinement,
chiefly condensing and reorganizing the author's earlier research monograph
(initially released in July 2025) and later work, including the monograph~\cite{ChiapponeGHOUQFT2026}. None of the figures were AI-generated
or AI-designed; AI only polished the author's figures for clarity.
The tool did not originate scientific results or ideas and served neither
as an author nor as an independent authority.

The use of AI should not be interpreted as an endorsement of any company, its products, or its business practices. The tool was used under human oversight to help minimize errors in the manuscript, primarily out of consideration for readers and editors, while also accommodating the author's medical circumstances. The author carefully reviewed and edited the AI-assisted output as necessary and assumes full responsibility for the scientific content and conclusions of this work.

\emph{Scope of assistance in the present working revision.}
ChatGPT also assisted with drafting the PE-CGG-HOUQFT sections,
developing and checking the mathematical arguments, and constructing
numerical checks of the restricted models. The attribution and review
statement above describes the earlier preprint.

\appendix
% Keep theorem numbers compact; elsarticle retains its native appendix headings.
\renewcommand{\thetheorem}{\Alph{section}.\arabic{theorem}}

\section{Details of the Gaussian and determinant estimates}
\label{sec:appendix-gaussian}

\subsection{Weighted-basis identity for the affine shift}

Let \(C\) be an \(I\times L\) real matrix of full column rank,
\(W=\diag(w_i)\) with \(w_i>0\), and
\(A=C^TWC\). For one component of the routing vector,
\[
 \mu=A^{-1}C^TWR.
\]
Cauchy--Binet gives
\begin{equation}
 \det A=\sum_{|T|=L}(\det C_T)^2\prod_{i\in T}w_i.
 \label{eq:CB-appendix}
\end{equation}
Terms with singular \(C_T\) vanish.

For component \(j\) of \(\mu\), let \(C^{(j)}\) be \(C\)
with column \(j\) replaced by \(R\).
Cramer's rule and a second Cauchy--Binet expansion give
\begin{align}
 \mu_j
 &=\frac{\det(C^TWC^{(j)})}{\det A}\nonumber\\
 &=\frac{
 \sum_{|T|=L}\det C_T\det C_T^{(j)}
                          \prod_{i\in T}w_i}
 {\det A}.
 \label{eq:Cramer-shift}
\end{align}
For invertible \(C_T\), Cramer's rule in that square subsystem
says
\[
 \det C_T^{(j)}
     =\det C_T\,(C_T^{-1}R_T)_j.
\]
Substitution yields \eqref{eq:shift-mixture}. Since there are
only finitely many such subsystems, compact external data imply
a uniform bound independently of relative widths. The finite convex
representation supplies the uniformity needed when some widths
vanish much faster than others.

\subsection{Gaussian moments with the affine shift retained}

Writing \(n=4L\), let
\(y=(A^{1/2}\otimes\one_4)(k+\bmu)\).
For \(D=0\), the numerator is bounded and the zeroth Gaussian
integral suffices. For \(D>0\), let
\(c_D=\max\{1,2^{D-1}\}\) and
\(B=\sup_{w_i>0,\,p\in\mathcal K}\norm{\bmu(w,p)}<\infty\).
The inequality \((x+y)^D\leq c_D(x^D+y^D)\) gives
\[
 1+\norm{k}^D\leq
 1+c_DB^D+c_D\norm{A^{-1}}^{D/2}\norm{y}^D.
\]
For every real \(D\geq0\), the radial Gaussian moment is
\begin{equation}
 \int_{\RR^n}|y|^D e^{-|y|^2}\dd^ny_{\bb}
   =\frac{\pi^{n/2}\Gamma((n+D)/2)}{\Gamma(n/2)}.
 \label{eq:Gaussian-moment-exact}
\end{equation}
It follows by integrating over the unit sphere and substituting
\(v=|y|^2\). Multiplying by the Jacobian
\((\det A)^{-2}\) proves \eqref{eq:gaussian-bound}.
For \(D=0\), only the zeroth moment is required; the displayed
constant-term convention remains harmless.

\subsection{A direct momentum-space proof of the coverage criterion}
\label{sec:direct-loop-proof}

For Corollary~\ref{cor:damped-basis}, write the bounded multiplier
on an unselected line as \(\Psi_i\), with
\(\sup_{z\geq0}|\Psi_i(z)|\leq B_i\). Its whole line factor is
bounded by $B_i\mathcal I_i$ for any fixed positive scale in the
multiplier argument.
Use the selected momenta as independent variables:
\[
 k=(C_S^{-1}\otimes\one_4)(q_S-R_S(p)),\qquad
 \dd^{4L}k_{\bb}=|\det C_S|^{-4}\dd^{4L}q_{S,\bb}.
\]
The bounded external shift gives
\(1+\norm{k}^D\leq C_{G,\mathcal K,S}
(1+\sum_{j\in S}|q_j|^D)\) for \(D>0\).
For \(D=0\), a constant bound suffices.

For \(0\leq d\leq D\), define the one-line integral
\begin{align}
 Q_{d,j}
 &:=\int_{\RR^4}|q|^d\Ph_j(a_j^2q^2)C_j(q)\dd^4q_{\bb} \nonumber\\
 &=\pi^2\Gamma(2+d/2)
   \int\rho_{j,\bb}(\dd\tau)\int_0^\infty
   \frac{E_j(\alpha)\dd\alpha}
        {(\alpha+a_j^2\tau)^{2+d/2}}.
 \label{eq:one-line-moment}
\end{align}
Here \(|q|^0\) means the constant function one. Tonelli's theorem
and the radial Gaussian integral justify the second equality.
Writing \(M_{0,j}=\rho_{j,\bb}([0,\infty))\),
\eqref{eq:alpha-envelope-bound} gives
\begin{align}
 Q_{d,j}\leq\pi^2\Gamma(2+d/2)
 \bigg[&\frac{\mathcal E_j a_j^{-2-d}}{1+d/2}
           M_{1+d/2}[\rho_{j,\bb}] \nonumber\\
       &+\alpha_0^{-2-d/2}\mathcal I_j M_{0,j}\bigg].
 \label{eq:one-line-moment-bound}
\end{align}
Finite measure mass and the assumed moment make both
\(Q_{0,j}\) and \(Q_{D,j}\) finite. Apart from fixed routing and
Fourier constants, the transformed graph majorant is bounded by
\[
 \left(\prod_{i\notin S}B_i\mathcal I_i\right)
 \left[\prod_{j\in S}Q_{0,j}
   +\sum_{j\in S}Q_{D,j}
        \prod_{h\in S\setminus\{j\}}Q_{0,h}\right].
\]
For \(D=0\), the first product alone is enough. The bound holds
on all loop-momentum space, so it controls partial and simultaneous
UV limits with constants uniform on \(\mathcal K\).
For a positive mass, \(\mathcal I_i=m_i^{-2}\).
The large-radius factor \(r^{D+1}\Ph_j(a_j^2r^2)\) explains the
inverse-moment order \(1+D/2\) in this sufficient estimate.

\section{Normalization and limiting checks for the examples}
\label{sec:appendix-checks}

\subsection{The two-body phase-space factor}

In the future rest frame \(P=(\sqrt{\sMand},\mathbf0)\),
the spatial delta function in \eqref{eq:two-body-phase-space}
sets \(\mathbf q_2=-\mathbf q_1\).
For equal masses, \(\omega=\sqrt{|\mathbf q_1|^2+m^2}\);
the remaining integral is
\[
 \frac1{4(2\pi)^2}
     \int\frac{4\pi r^2\dd r}{\omega^2}
         \delta^{\rm sh}_{\bb}(\sqrt{\sMand}-2\omega).
\]
Since \(\dd(2\omega)/\dd r=2r/\omega\), this equals
\[
 \frac{r_*}{8\pi\omega_*}
       =\frac1{8\pi}\sqrt{1-\frac{4m^2}{\sMand}}.
\]
The imaginary part of a Schwarz-reflecting scalar bubble is
half its discontinuity divided by \(\ii\).
A separate graph symmetry factor must then be applied at the
amplitude level.

\subsection{Zero-momentum Gaussian bubble}

Set \(b=2t\) and \(x=bm^2\). Radial integration of
\eqref{eq:two-line-bubble} at zero external momentum gives
\[
 16\pi^2B_{t,E}(0)
   =\int_0^\infty
          \frac{u\,e^{-bu}}{(u+m^2)^2}\dd u.
\]
Let
\[
 I_1=\int_0^\infty\frac{e^{-bu}}{u+m^2}\dd u
       =e^xE_1(x).
\]
Differentiating with respect to \(m^2\),
\[
 I_2=\int_0^\infty\frac{e^{-bu}}{(u+m^2)^2}\dd u
       =\frac1{m^2}-bI_1.
\]
Using \(u/(u+m^2)^2=1/(u+m^2)-m^2/(u+m^2)^2\)
yields
\[
 16\pi^2B_{t,E}(0)
   =I_1-m^2I_2=(1+x)e^xE_1(x)-1.
\]
As \(x\to\infty\), the bracket is \(x^{-2}+O(x^{-3})\).
As \(x\downarrow0\), it diverges logarithmically.
The finite-\(a\) theorem thus does not authorize interchanging
integration with the formal removal of damping.

\subsection{Radial kernel: endpoint and numerical stability}

Expanding \(e^{-ur^2}\) to a fixed finite order in
\eqref{eq:J-radial} gives, as \(u\downarrow0\) at fixed
\(a>0\),
\begin{equation}
 J(u;a)
 =\frac6{a^4}-\frac{120u}{a^6}
        +\frac{2520u^2}{a^8}
        -\frac{60480u^3}{a^{10}}+O(u^4).
 \label{eq:J-endpoint-expansion}
\end{equation}
Each coefficient follows from
\(\int_0^\infty r^n e^{-ar}\dd r=n!/a^{n+1}\).
More generally, set
\[
 P_N(u;a)=\sum_{n=0}^N
       \frac{(-u)^n(2n+3)!}{n!\,a^{2n+4}}.
\]
Taylor's remainder for \(e^{-ur^2}\), integrated against the positive
weight \(r^3e^{-ar}\), gives the explicit enclosure
\begin{equation}
\begin{gathered}
 0\leq(-1)^{N+1}\bigl[J(u;a)-P_N(u;a)\bigr]\\
 \leq\frac{(2N+5)!}{(N+1)!\,a^{2N+6}}u^{N+1},
 \qquad a>0,\quad u\geq0.
\end{gathered}
 \label{eq:J-truncation-bound}
\end{equation}
Odd truncations are lower bounds and even truncations are upper
bounds. In particular,
\[
 0\leq J(u;a)-P_3(u;a)
       \leq\frac{1663200\,u^4}{a^{12}}.
\]
These estimates apply to each finite truncation. The scaling identity
\[
 J(u;a)=a^{-4}j(u/a^2),\qquad
 j(\xi)=\int_0^\infty x^3e^{-x-\xi x^2}\dd x
\]
identifies \(u/a^2\) as the dimensionless endpoint parameter.
The estimates become useful for approximation when that parameter is
small; no interchange with the infinite asymptotic series is needed.

Two dimensionless numerical checks of
\eqref{eq:J-closed} against the defining radial integral are
\begin{center}
\begin{tabular}{@{}ccr@{}}
\toprule
\(u\)&\(a\)&\(J(u;a)\), both evaluations\\
\midrule
\(1\)&\(1\)&\(0.147563809330584\)\\
\(2\)&\(0.4\)&\(0.086614103683222\)\\
\bottomrule
\end{tabular}
\end{center}
The representation by \(\erfcx\) controls exponential overflow,
but subtractive cancellation for large \(v\) should still be
handled with the integral or an appropriate finite endpoint
expansion.

\subsection{A common-domain justification of the Ward remainder}
\label{sec:Ward-cutoff-proof}

\begin{lemma}[Common-domain Ward translation]
\label{lem:Ward-common-cutoff}
Fix \(m\geq0\), \(a>0\), and \(Q\in\RR^4\), and use
the canonical free symbols of \eqref{eq:QED-conventions}. Let
\[
\begin{gathered}
 F(u)=\Ph(a^2u)=\int_{[0,\infty)}e^{-a^2\tau u}\rho_{\bb}(\dd\tau),\\
 \rho_{\bb}\geq0,\qquad \rho_{\bb}([0,\infty))<\infty.
\end{gathered}
\]
Assume that the original tensor integrand
\[
 H_Q^{\mu\nu}(k)
 =F_kF_{k+Q}
   \Tr[\Gamma_0^\mu S_0(k)\Gamma_0^\nu S_0(k+Q)]
\]
is in \(L^1(\RR^4)\) componentwise. Then
\[
 \mathcal R_Q^\nu(k)
 =F_k[F_{k+Q}-F_{k-Q}]\,
       \Tr[\Gamma_0^\nu S_0(k)]
\]
is also in \(L^1(\RR^4)\), and
\eqref{eq:Ward-remainder} follows without assuming that the two
terms in the contracted difference are separately integrable.
\end{lemma}

\begin{proof}
Put \(T^\nu(k)=4k^\nu/(k^2+m^2)\).
First, absolute integrability of \(H_Q\) forces the radial tail
condition
\begin{equation}
 \int_{r_0}^{\infty}rF(r^2)^2\dd r<\infty
 \quad\hbox{for some }r_0>0.
 \label{eq:Ward-radial-tail}
\end{equation}
Indeed, write \(k=rn\), \(|n|=1\), and fix a diagonal component
\(\mu\), with no sum over \(\mu\). The trace has the uniform
large-\(r\) expansion
\[
 \Tr[\Gamma_0^\mu S_0(k)\Gamma_0^\mu S_0(k+Q)]
 =\frac{4(2n_\mu^2-1)}{r^2}+O_{Q,m}(r^{-3}).
\]
On the fixed angular region \(n_\mu^2\leq1/4\), its modulus is at
least \(c/r^2\) for sufficiently large \(r\).
Since \(F\) is nonnegative and nonincreasing,
\[
 F(r^2)F(|k+Q|^2)\geq F((r+|Q|)^2)^2.
\]
Integration on this angular region, followed by a bounded radial
shift, proves \eqref{eq:Ward-radial-tail}.
For the nonincreasing function \(f(r)=F(r^2)^2\),
\[
 \frac38R^2f(R)
 \leq\int_{R/2}^{R}rf(r)\dd r\longrightarrow0.
\]
Thus \(R^2F(R^2)^2\to0\).

The radius \(R\) below only exhausts the integration domain; it
is not a new physical regulator or kernel parameter.
Now let \(B_R=\{k:|k|<R\}\) and set
\begin{align*}
 I_R^\nu
 &=\int_{B_R}\frac{\dd^4k_{\bb}}{(2\pi)^4}
       F_kF_{k+Q}[T^\nu(k)-T^\nu(k+Q)],\\
 J_R^\nu
 &=\int_{B_R}\frac{\dd^4k_{\bb}}{(2\pi)^4}
       F_k[F_{k+Q}-F_{k-Q}]T^\nu(k),\\
 C_Q^\nu(k)&=F_kF_{k-Q}T^\nu(k).
\end{align*}
All these finite-domain integrals exist. Changing variables on the finite integration domains gives the exact
identity
\begin{align}
 I_R^\nu&=J_R^\nu+E_R^\nu,\nonumber\\
 E_R^\nu
 &=\int_{B_R}\frac{\dd^4k_{\bb}}{(2\pi)^4}C_Q^\nu(k)
   -\int_{B_R+Q}\frac{\dd^4k_{\bb}}{(2\pi)^4}C_Q^\nu(k).
 \label{eq:Ward-finite-cutoff}
\end{align}
For sufficiently large \(R>4|Q|\), the symmetric difference of the balls lies in
a shell of volume \(O_Q(R^3)\), and \(|T^\nu(k)|\leq C/R\)
there. Both \(|k|\) and \(|k-Q|\) are at least
\(R-2|Q|\). Monotonicity therefore gives
\begin{equation}
 |E_R^\nu|
 \leq C_QR^2F((R-2|Q|)^2)^2
 \longrightarrow0.
 \label{eq:Ward-boundary-error}
\end{equation}

It remains to establish the absolute integrability of the final
remainder. For \(u>0\), the positive measure representation gives
\begin{equation}
 -uF'(u)
 =\int_{[0,\infty)}(a^2\tau u)e^{-a^2\tau u}\rho_{\bb}(\dd\tau)
 \leq\frac2{\mathrm e}F(u/2),
 \label{eq:Ward-CM-derivative}
\end{equation}
because \(x e^{-x/2}\leq2/\mathrm e\).
The mean-value theorem along \(k+tQ\), \(-1\leq t\leq1\),
then yields, for sufficiently large \(r=|k|>2|Q|\),
\[
 |F_{k+Q}-F_{k-Q}|
 \leq\frac{C_Q}{r}F(r^2/8).
\]
Using \(|T^\nu(k)|\leq4/r\), the absolute radial tail of
\(\mathcal R_Q^\nu\) is bounded by a constant times
\[
 \int^\infty rF(r^2)F(r^2/8)\dd r.
\]
This is finite by Cauchy--Schwarz, \eqref{eq:Ward-radial-tail},
and a constant rescaling of the radial variable in the second
factor. On compact momentum sets, \(F(0)<\infty\) and
\(|T^\nu(k)|\leq 4/|k|\) give local integrability even at \(m=0\).
Hence \(J_R^\nu\) converges absolutely.
The algebraic Ward contraction identifies
\(-e^2\lim_{R\to\infty}I_R^\nu\) with
\(Q_\mu\Pi_{\rm bv}^{\mu\nu}\). Taking the limit in
\eqref{eq:Ward-finite-cutoff} proves
\eqref{eq:Ward-remainder}.
\end{proof}

At \(m=0\), the original free tensor is bounded locally by
\(C/(|k|\,|k+Q|)\), with the coincident case \(Q=0\) of order
\(|k|^{-2}\). These singularities are locally integrable in four
Euclidean dimensions. With a different infrared prescription,
the Ward contraction and the finite-domain changes of variables
both require verification for the modified symbols.

The distinction between the original tensor and its split terms is
material. For \(F(u)=(1+a^2u)^{-\beta}\), \(m>0\), and
\(1/2<\beta\leq3/4\), the tensor has an integrable absolute
radial tail \(r^{1-4\beta}\), while each split term has a
nonintegrable absolute tail \(r^{2-4\beta}\).
The common-cutoff proof includes this range.

\paragraph{Relation to the monograph}
% Monograph source: eq:Pi-Omni-full; lem:ST-stability-internal-only-early.
The monograph distinguishes the bare-vertex tensor from the
completed transverse tensor and derives the latter's Ward
identity from the completed A5/Q1 functional
\cite{ChiapponeGHOUQFT2026}. The lemma above adds the explicit
finite-cutoff and \(L^1\) justification for the bare-vertex
remainder; it is not a replacement for that action-level argument.
It uses precisely \(\Gamma_0^\mu=\ii\gamma^\mu\) and
\(S_0^{-1}=\ii\slashed k+m\) from
\eqref{eq:QED-conventions}.

\subsection{Stable-half covariance tail and bubble differentiability}
\label{sec:stable-bubble-proof}

Use the inverse Fourier convention of
Section~\ref{sec:stable-bubble-origin}. Subordination by the
stable-half density gives
\begin{align}
 P_a(x)
 &=\int_0^\infty
    \frac{a\,e^{-a^2/(4v)}}{2\sqrt\pi\,v^{3/2}}
    \frac{e^{-|x|^2/(4v)}}{(4\pi v)^2}\dd v\nonumber\\
 &=\frac{3a}{4\pi^2(a^2+|x|^2)^{5/2}}.
 \label{eq:four-dimensional-Poisson-kernel}
\end{align}
The second line follows by substituting
\(y=(a^2+|x|^2)/(4v)\) and using
\(\Gamma(5/2)=3\sqrt\pi/4\). This derives the coefficient
rather than importing a different Fourier normalization.
The positive massive free Green function is
\[
\begin{aligned}
 G_m(y)&=\int_0^\infty e^{-m^2v}
       (4\pi v)^{-2}e^{-|y|^2/(4v)}\dd v,\\
 \int_{\RR^4}G_m(y)\dd^4y_{E,\bb}&=m^{-2}.
\end{aligned}
\]
It has an exponentially decaying integral tail. For example,
\(c|y|\leq |y|^2/(8v)+2c^2v\) implies
\[
 \int e^{c|y|}G_m(y)\dd^4y_{E,\bb}
       \leq\frac{4}{m^2-2c^2}<\infty,
       \qquad 0<c<m/\sqrt2.
\]
Fourier multiplication therefore gives the positive convolution
\(\mathcal C_{a,m}=G_m*P_a\).

Put \(r=|x|\). On \(|y|\leq r/2\) the derivative of
\(P_a\) along the segment from \(x\) to \(x-y\) is
\(O(r^{-6})\) as \(r\to\infty\). The mean-value estimate gives
\(|P_a(x-y)-P_a(x)|\leq Cr^{-6}|y|\), which integrates against
\(G_m\) because the first moment is finite. On \(|y|>r/2\),
the bound \(2\|P_a\|_\infty\) and the exponential integral tail
of \(G_m\) give an error smaller than \(O(r^{-6})\). Thus
\(\mathcal C_{a,m}(x)=m^{-2}P_a(x)+O(r^{-6})\),
which proves \eqref{eq:stable-covariance-long-tail}. Also
\(\|\mathcal C_{a,m}\|_\infty\leq m^{-2}\|P_a\|_\infty\),
so there is no local coordinate singularity.

The momentum covariance is real, even and square-integrable.
Fourier convolution, with the stated \((2\pi)^{-4}\) convention,
gives the exact identity
\begin{equation}
 B_{a,m}(P)=\int_{\RR^4}e^{\ii P\cdot x}
                         \mathcal C_{a,m}(x)^2\dd^4x_{E,\bb}.
 \label{eq:stable-bubble-coordinate-square}
\end{equation}
Since \(\mathcal C_{a,m}(x)^2\asymp r^{-10}\) at infinity,
all its absolute moments of orders \(0\leq j<6\) are finite:
the radial tail is \(r^{j-7}\dd r\). Differentiation under
the integral proves \(C^5\) regularity. On an angular set with
\(|x_1|\geq r/2\), however, the sixth moment has the divergent
tail \(\int^\infty\dd r/r\). Thus
\(\int x_1^6\mathcal C_{a,m}(x)^2\dd^4x_{E,\bb}=\infty\).

For the leading coefficient, set $c=3a/(4\pi^2m^2)$ and
subtract $c^2r^{-10}\mathbf1_{r>m^{-1}}$ from
$\mathcal C_{a,m}(x)^2$. The difference has a finite absolute
sixth moment. Its cosine integral after the Taylor terms through
order four has an $O(h^6)$ remainder. For the subtracted tail
itself, split the radial integral at $r=|h|^{-1}$. Inside that
radius the next cosine term is $-h^6 x_1^6/720$, with error bounded
by $C|h|^8 r^8$. The radial integral of this error against the
tail is $O(h^6)$. Outside that radius the cosine minus its first
three Taylor terms is bounded by $C(1+|h|^4r^4)$, which also
integrates to $O(h^6)$. The sixth term on the intermediate
annulus contributes
\[
 -\frac{c^2 h^6}{720}\log\frac{m}{|h|}
                  \int_{S^3}n_1^6\dd\Omega.
\]
The spherical moment is $2\pi^2\,[1\cdot3\cdot5/(4\cdot6\cdot8)]
=5\pi^2/32$, which gives \eqref{eq:stable-bubble-leading-log}.

One can independently exclude \(C^6\) directly from the divergent
sixth moment. Use \(b(h)=B_{a,m}(h e_1)\) in the symmetric sixth
difference
\[
 d_h^6b:=\sum_{j=0}^6(-1)^j\binom{6}{j}b((3-j)h).
\]
Equation~\eqref{eq:stable-bubble-coordinate-square} yields
\begin{equation}
 -\frac{d_h^6b}{h^6}
   =\int_{\RR^4}\frac{64\sin^6(hx_1/2)}{h^6}
                  \mathcal C_{a,m}(x)^2\dd^4x_{E,\bb}.
 \label{eq:stable-sixth-difference}
\end{equation}
The integrand is nonnegative and converges pointwise to
\(x_1^6\mathcal C_{a,m}(x)^2\). Fatou's lemma makes the
lower limit of the right side infinite. If \(b\) were \(C^6\)
near zero, its left side would tend to the finite number
\(-b^{(6)}(0)\), a contradiction. Holomorphy in \(P^2\)
around zero would in particular imply analyticity of \(b(h)\),
and is therefore also excluded. This completes the proof of
Proposition~\ref{prop:stable-bubble-nonanalytic}.

\section{Logical dependence of the principal results}
\label{sec:dependence}

Sections~\ref{sec:pe-principle}--\ref{sec:pe-curvature} identify the
objects to be tested on the physical GCM branch. The entries below
record hypotheses for individual steps.
In particular, the convergence and symmetry rows must refer to the
same completed action and the same graph family.
The adjoint and reconstruction assumptions concern physical source
channels and full correlation functions, not just their auxiliary
carrier coefficients.

\begin{description}[style=nextline]
\item[Common microscopic source.] Postulate~$\Omega_\alpha$ requires
one assignment of raw data, histories and preparation to a quantum
source. The five recoveries test its classical action, gauge functional,
physical state, fixed-plane kernel and groupoid extension. Carrying out
one kernel test leaves the higher source derivatives to be computed.

\item[Moving charged sources.] For the optical spinor,
the heat lapse fixes the charged displacement width and the normal
Maxwell form fixes the photon sampling length. Theorem~\ref{thm:seam-calibrated-joint-uv}
and Corollary~\ref{cor:seam-optical-shifts-uv} use their relative
scaling and the full moving normal measure. On the regular chart
a two-pair exchange gives the separate obstruction in
Proposition~\ref{prop:seam-uv-two-pair}.

\item[GCM electromagnetic reduction.] Retain photon and charged probes
on the declared physical chart. The normalization, determinants,
constraint and boundary variations contribute to derivatives of this
same functional. Agreement with the raw O--QED density at classical
or quadratic order determines only the terms so compared.

\item[Restricted curvature proof.] Given the three specified curvature
blocks, positive pairing, harmonic normal propagation and common depth
law of Section~\ref{sec:pe-curvature}, the determinant and the
interaction share the same transferred fields. This provides a
cutoff-uniform absolute bound for every fixed-order graph on that
chart. The physical even and odd GCM fields can have different
statistical determinant powers and source couplings.

\item[Transported charged heat displacement.] The fixed positive heat
lapse, charged interval representation, harmonic readout and transverse
Maxwell law give Theorem~\ref{thm:s6-heatlink-uv}. Gauge covariance
of the diagonal resolvent and weighted bounds on its free inverse
give \eqref{eq:s6-eachleg} on the fixed-distance continuum chart.
For a fluctuating readout, Proposition~\ref{prop:dualrev-uvheat-depth}
includes the full common vacuum weight in each graph. The
real-coupling partition and the coefficient-matched derivatives are
treated in Theorem~\ref{thm:dualrev-uvheat-real} and
Proposition~\ref{prop:uvheat-critical-charge}. The spectral test on
the free open charged line is \eqref{eq:s6-freepole}. For the surface selected by the curvature
action, Corollary~\ref{cor:geo-free-slope-uv} gives the direct geometric
graph theorem, with all shared radial limits included.

\item[Plane binding.] \textit{Inputs:} Fixed labelled axial fibre, independent A2/GTP momentum binder,
and support-diagonal action; admitted canonical or source representation.
\textit{Conclusion:} Selection rule on the chosen support; positivity and scattering
unitarity require the physical completion.

\item[Carrier reconstruction.] \textit{Inputs:} Fixed pole carrier,
declared vectorial boundary limits, both character projections of
the same complete functional, and compatible functional-calculus branches.
\textit{Conclusion:} Carrier-valued lines and functionals; one complex coefficient suffices
only under the stated reflection or directed-subalgebra conditions.

\item[Reality and adjoints.] \textit{Inputs:} The declared carrier
reality \(\ddagger\); a separately specified physical adjoint
\(\dagger\). \textit{Conclusion:} Reality relations for coefficients
do not establish physical positivity. A positive unital representation
preserving \(\ddagger\) on the entire carrier algebra is excluded by
Proposition~\ref{prop:carrier-reality-adjoint}.

\item[Source covariance.] \textit{Inputs:} Graded source derivatives at fixed branch, source pairing, and
translation-invariant field components.
\textit{Conclusion:} Exchange symmetry of the complete source matrix. An odd oriented
running line is not an identical commuting scalar covariance.

\item[Geometric source map.] \textit{Inputs:} Declared normal operator,
boundary domain, normalized deterministic fiber, labelled source pairing,
and positive retained covariance. \textit{Conclusion:} The full three-block
source law. Its canonical one-sided multiplier additionally requires
spectral commutation; the construction is not forced by coordinate GTL.

\item[Correlated-depth convergence.] \textit{Inputs:} Joint heat/depth law,
full-rank selected routing, the complete uniform envelope, and the joint
moments in \eqref{eq:update-joint-moment-budget}.
\textit{Conclusion:} Uniform finite-graph absolute convergence including
the depth integral. The independent-line theorem cannot be applied to
marginals of a common-depth law without this extra justification.

\item[Depth statistics and limiting rate.] \textit{Inputs:} The specified
collective measure and its positive coefficients.
\textit{Conclusion:} Exact moments and marginal tail, with common-depth
cumulants retained. The fixed-depth exponential requires a point mass;
its limiting ultraviolet rate is not inferred from concentration at each
fixed spectral value.


\item[Global graph convergence.] \textit{Inputs:} Full-rank routing, completed numerator bound, gap or sufficient
inverse moments, and the stated IR control.
\textit{Conclusion:} Each finite graph at fixed \(a>0\) and admitted IR data;
singular parameter limits need separate bounds.

\item[Optimized line coverage.] \textit{Inputs:} Full-rank
selected routing, line-specific inverse moments, and an anisotropic
bound covering the complete numerator and unselected lines.
\textit{Conclusion:} Absolute convergence even with selected massless
$1/q^2$ denominators; stronger moments give explicit selected-variable
UV tails. Unselected massless singularities are not excused.

\item[Residual damping.]
\textit{Inputs:} A bound on the full nonpolynomial vertex product and
a strictly positive quadratic or linear damping margin, with the
stated bounded infrared covariances.
\textit{Conclusion:} Uniform absolute Euclidean convergence of the
complete graph. Failure of this sufficient test alone proves no
divergence and no continuation statement.

\item[O--QED kernel application.] \textit{Inputs:} The normalized kernel
\(F_{a,M}\), fixed \(a>0\), \(M\geq0\), and the same completed
vertices and infrared/source conventions as in the claimed amplitude.
\textit{Conclusion:} All inverse moments are verified. Combining
those damped lines with local bare vertices gives finite comparison
integrals. The Ward-compatible two-vertex trial fails absolute integrability
before contacts. The GCM-induced vertices require their own estimates.
The declared constant-background contact trial also fails absolute
integrability. Its difference from the local radial reference
has the obstruction in Corollary~\ref{cor:OQED-general-contact-obstruction}.
Other completions remain to be specified. A grouped result requires
its own common-routing envelope.
Medium sums require an explicit spectral/infrared bound, and the
kernel-off reference is separately renormalized and matched.

\item[Forest convergence.] \textit{Inputs:} Uniform simultaneous sector envelopes, every active UV scale
controlled, and IR/angle bounds.
\textit{Conclusion:} All verified sector scales; each damping gain is assigned through
a simultaneous routing estimate.

\item[Slavnov--Taylor identity.] \textit{Inputs:} Completed closed BRST/BV action in both characters, compatible
source/antifield maps and source covariances, invariant anomaly-free
measure, and admitted insertions.
\textit{Conclusion:} Formal perturbative identity for the completed functional;
its vertex growth is checked independently.

\item[Nielsen identity.] \textit{Inputs:} Extended BRST functional, admissible gauge-parameter insertion,
fixed kernel/source maps, and compatible physical normalization.
\textit{Conclusion:} Gauge independence for admitted fixed or Nielsen-transported physical
classes; off-shell functions can depend on the gauge parameter.

\item[Local holomorphy.] \textit{Inputs:} Holomorphic graph integrand and a common pointwise complex
\(L^1\) majorant.
\textit{Conclusion:} The declared complex domain; not a holomorphic neighbourhood of
the origin for the literal massive stable-half bubble. Boundary
growth and channel continuation require their own estimates.

\item[Depth limits and scalar subtraction.] \textit{Inputs:} The normalized
square-root kernel at positive depth and the literal scalar bubble,
with $m>0$ for the subtraction and strong-depth limits.
\textit{Conclusion:} Explicit depth derivatives, distinct strong-depth
regimes and crossover, a nonuniform undamped line limit, and a finite
$a\downarrow0$ limit for the \emph{already subtracted} scalar bubble.
The literal bubble is completely monotone in positive depth,
with an explicit uniform derivative bound and an exponentially
decaying O--QED line tail. Its unsubtracted zero-depth limit has
a universal logarithmic finite part with a quantitative error
bound. These are not interacting gauge or retarded-response results.

\item[Scalar approximation.] \textit{Inputs:} A literal scalar bubble
with \(a,M,m>0\) and its positive proper-time measure. \textit{Conclusion:} Holomorphy for
\(\Re w>-4\min(m,M)^2\), signed finite-order bounds for real
\(w\geq0\), and the separately stated complex Taylor and radial-tail
bounds. This is not physical spectral positivity or a completed gauge
response.

\item[Physical comparison obstructions.] \textit{Inputs:} The stated
source algebra and isotropic transfer, or an ordinary positive matrix
spectral representation. \textit{Conclusion:} The joint and boundary
reflection counterexamples, the cross-spectrum restriction, and the
specified positive-completion bubble divergence. These hypotheses are
not imposed on a ghost or merely diagrammatic kernel.

\item[GT contour.] \textit{Inputs:} One physical state, the declared GT
orderers, and common source domains. \textit{Conclusion:} The two-orderer
identity and the explicit nonlocal-transfer correction. The square-root
virtual line and an undressed Wightman sum generally require additional
action-derived terms before they can share that contour.


\item[Physical reconstruction.] \textit{Inputs:} Completed even-observable
full correlation functions satisfying Assumption~\ref{ass:OS-completion},
including \eqref{eq:OS-hierarchy-growth} uniformly in insertion
number. Odd physical fields require the corresponding graded
reconstruction. Finiteness graph by graph is not this hierarchy
condition or an order-by-order positivity claim.
\textit{Conclusion:} A positive physical reconstruction on the admitted
observable algebra; an individual virtual-line kernel need not be
a physical spectral covariance.

\item[Pole normalization.] \textit{Inputs:} An isolated pole of the
full inverse in a compatible physical source channel; the actual tree
inverse and loop correction; positive one-particle spectral overlaps.
\textit{Conclusion:} Finite normalization on the nonzero residue
subspace, with an invertible field map and its dual source map when
all field coordinates are retained. Neither \(\Ph(0)=1\) nor an
arbitrary nonzero complex Laurent coefficient establishes positivity.

\item[Cut matching and optical relation.] \textit{Inputs:} Normalized physical states and completed subamplitudes; cut--jump
matching. The optical relation also requires unitary scattering
and a complete state resolution, or an equivalent Hermitian-analytic
cut-factorization argument.
\textit{Conclusion:} Completed amplitudes with their actual residues and branch terms;
positive norms and support diagonality alone do not imply the
optical identity.

\item[Gaussian shell normalization.] \textit{Inputs:} The exact fixed-carrier Gaussian bubble and its finite reference
subtraction.
\textit{Conclusion:} The usual normalized jump plus an entire remainder; the remainder
has nonpolynomial positive-real-axis growth.

\item[Response and pole errors.] \textit{Inputs:} A certified complex-domain
response error, compatible operator/domain norms, and an invertible
approximate electromagnetic response with a strict Neumann margin.
The endpoint root enclosure additionally requires the holomorphic
retained matrix pencils and contour isolation in
Proposition~\ref{prop:OQED-finite-pole-enclosure}. Interpreting its root
as the selected physical pole requires baseline-branch identification;
numerical contour samples require between-node bounds, and a compression
of a larger theory requires control of the discarded complement.
\textit{Conclusion:} Schur and first-order pole-family error bounds;
under the respective additional hypotheses, an endpoint root enclosure
and a bound on the selected finite rate. A complex pole-family margin
supplies the explicit Taylor remainder in
\eqref{eq:OQED-pole-remainder-certificate}; an observed response pole
also needs the nonzero input/output residue in
\eqref{eq:OQED-observable-pole-residue}. Its nonvanishing can be
certified by \eqref{eq:OQED-observable-residue-error}; eliminating a
complement also transports the endpoint maps as in
\eqref{eq:OQED-Schur-source-readout}. A real Euclidean error alone
supplies neither conclusion at an unverified complex continuation point.

\item[Physical crossing.] \textit{Inputs:} Local physical reconstruction,
suitable asymptotic poles and LSZ limits, and controlled continuation
between the channels, with the physical boundary-value identification
established independently of the domain result. Invocation of the cited
Bros--Epstein--Glaser domain theorem also requires its stability and
strictly positive physical mass-gap hypotheses.
\textit{Conclusion:} Crossing for the admitted scattering channels
and asymptotic states; the continuation must account for kernel
and vertex singularities.
\end{description}

\section{Relation to the foundational constructions}
\label{sec:monograph-differences}

The cited monograph revision and GTL Parts~I and~II give the
underlying sector-plane conventions
\cite{ChiapponeGHOUQFT2026,ChiapponePartI2026,ChiapponePartII2026}.
The subsequent path-enriched formulation introduces a candidate GCM
quantum source and Postulate~$\Omega_\alpha$. Its action-level
reduction and UV tests are set out independently in
Sections~\ref{sec:pe-principle}--\ref{sec:g5-induction}.
The locked-plane geometry, directed vectorial boundary prescription,
two-character carrier reconstruction, and original Gaussian routing proof
follow that framework. The complete mixed-source identities, commutation
criterion, depth cumulants and measure laws, GT-ordering tests, and
spectral-completion costs are restated here with the hypotheses needed
for this article's convergence problem. The joint-heat theorem is the
explicit integration of those source-derived correlations into the
selected-line proof. These results do not establish that the stated choices are necessary
consequences of the unchanged A3/A4 axioms. The present article states
its hypotheses independently and develops
additional sufficient Euclidean convergence criteria: the endpoint
inverse-moment condition, line-specific full-rank coverage with
massless selected denominators, residual coercivity for some
growing vertices, and explicit tail and complex-domain estimates.
The two-loop comparison following
Theorem~\ref{thm:optimized-selected-lines}, together with
Proposition~\ref{prop:two-loop-complex-certificate}, is a
mathematical test of those criteria, not a graph attributed to the
physical action. Its bounds refer to the explicit nonnegative
selected-line integrand, not to an independently completed
O--QED Feynman graph.

Several calculations delimit the use of the auxiliary covariances.
Source exchange symmetry prevents an unrestricted odd running-unit
line from being identified with the covariance of one identical
commuting real field without a compatible endpoint/source matrix.
Carrier reality $\ddagger$ is not the positive physical adjoint.
The Gaussian bubble, massive stable-half bubble, and reflection-
positivity counterexample separately test discontinuities, low-
momentum analyticity, and physical reconstruction. None changes
the original vectorial pole datum or identifies Euclidean damping
with a physical spectral representation.

The action-level BRST/BV and Nielsen identities require a compatible
completed action, source measures, and insertions. The common-domain
bare QED Ward argument in \ref{sec:Ward-cutoff-proof}
adds an explicit absolute-integration justification; it does not
replace the monograph's completed-action Ward argument.
The Euclidean $\Gamma_0^\mu=\ii\gamma^\mu$ convention used in
\eqref{eq:QED-conventions} is stated separately from the
$\gamma^\mu$-vertex display in the reference, so their overall
prefactors must be converted before comparison.

For the standalone O--QED example, the action and kernel input come
from the supplied revision of Ref.~\cite{ChiapponeIgnition}: its
sector-resolved reference action, directed hyper-propagators and
positive fixed-depth kernel. These are stated explicitly in the
corresponding local formulas of Section~\ref{sec:OQED-application}.
The inverse moments, depth-limit formulas, scalar comparison bounds,
and explicit Ward-bubble obstruction are calculated here. The
constant-background contact test additionally imposes a shift
covariantization, a trace-log realization, and a field-independent
measure; it is not an evaluation of the source's unspecified full
contact completion. The nonintegrability of these particular
terms neither proves divergence of every possible completion nor
supplies a missing combined-graph estimate.

The pole and readout certificates use the stated simple-root,
Schur-complement, and common-contour assumptions; the local inverse
residue invokes the standard result of Ref.~\cite{Beyn2012}.
Reflection positivity and crossing retain their separate OS and
physical-boundary hypotheses \cite{OS1975,BrosEpsteinGlaser1965}.
The restricted curvature theory has a complete fixed-order graph proof
under the hypotheses of Section~\ref{sec:pe-curvature}. The physical
GCM branch must meet the joint graph, gauge, and reconstruction
requirements with its own source, including the chiral fields and
boundary state.

\section{Lean verification}
\label{sec:lean-verification}

\paragraph{Version boundary of the formal inventory}
The formal claims in this section refer to the repository revision
cited below and its unchanged 45-claim inventory.
The geometric-source, joint-heat/depth, induced-transfer, GT-contour,
reflection, and spectral-completion results, including the fixed-depth
Wick-ordered witness and the covariant O--QED trial added here, are
analytical extensions with derivations in this article; they are
\emph{not} claimed to have been added to that pinned Lean development.
The endpoint/arrow maps, explicit ambient action assembly and elastic
benchmark added here are also outside that pinned proof inventory.
A check of the archived \texttt{paper/UV.tex} hash verifies that archived
snapshot, not byte identity with the present expanded manuscript. Extending
the formal coverage requires a new explicit proof inventory and build.


The proposed unified postulate, the electromagnetic pushforward, the
history-dependent graph estimates, the curvature-sector theorem
and the gauge-root obstruction in Sections~\ref{sec:pe-principle},
\ref{sec:pe-actions}, \ref{sec:pe-graphs} and~\ref{sec:pe-curvature}
also lie outside the pinned inventory. This includes the induced
first-order and Wilson graph theorems and the boundary-state results
in Section~\ref{sec:g5-induction}. Their derivations are given in
this article.

The pinned claim inventory in Ref.~\cite{UVLean} records 45 named
mathematical results as verified under their documented hypotheses.
These are the statements of its archived source snapshot, not a blanket
formal-verification claim for every proposition or later refinement
in this manuscript. The
\href{https://github.com/GiovanniC137z/uv-lean/blob/24ca145d8241664896d09751488b90d2fe5077ba/docs/claims.md}
{claim inventory} records the checked scope of each result. The
\href{https://github.com/GiovanniC137z/uv-lean/blob/24ca145d8241664896d09751488b90d2fe5077ba/claims.json}
{machine-readable record} maps the proofs to their source labels
and records the two physical assumptions and one definition.

\subsection{Scope of the formal proofs}

\paragraph{Graphwise ultraviolet convergence}
The Lean proofs cover the archived global and selected-line estimates
corresponding to Theorems~\ref{thm:global-uv} and~\ref{thm:optimized-selected-lines}.
The global proof constructs an invertible routing block and
bounds the joint integral by selected-line integration. The
checked results also include the selected-line tail, residual
coercivity, and the forest criterion under its simultaneous
sector bounds.

\paragraph{Kernel and analytic certificates}
The proofs construct the O--QED probability and heat measures
and derive the square-root kernel's inverse moments, the scalar
bubble bounds and the depth limits. They cover the common-majorant
holomorphy and derivative estimates, the Ward-vertex and contact
obstructions, and the counterexamples to reflection positivity
and sixth-order regularity for the stated heat and stable-half
models, respectively.

\paragraph{Source normalization and response}
The proof of Proposition~\ref{prop:free-source-normalization}
constructs a common reducing graph core from the supplied spectral
resolutions. The response proofs establish the inverse and Schur
bounds, finite-matrix pole enclosures, and the Taylor remainder
for the contour-defined holomorphic root.

\paragraph{Continuum Slavnov--Taylor identity}
For Theorem~\ref{thm:ST}, the formalization derives the partition
and connected Ward identities, the Legendre Slavnov--Taylor
identity, and the nonlinear formal loop hierarchy. When the same
completed graphs and insertions meet the ultraviolet bounds, their
coefficients are finite in both carrier characters. Existence of
the invariant anomaly-free measure, a common functional domain,
and a specified interacting completion remain hypotheses.
Background covariance and the extended Slavnov--Taylor and Nielsen
identities retain their additional premises from
Section~\ref{sec:brst}.

\paragraph{Universal rectifiable Jordan contours}
For Proposition~\ref{prop:OQED-observable-residue-error}, the
independent Jordan companion constructs a filling of the simple
closed contour using the formalized Jordan--Schoenflies theorem
of Ref.~\cite{SchoenfliesLean}. An arc-length representative
preserves the order of traversal and the original geometric
length. Cauchy's residue formula expresses each observable residue
as a contour integral. Integrating the Neumann bound on the
difference of the integrands gives the factor
\(\ell(\Gamma)/(2\pi)\) in
\eqref{eq:OQED-observable-residue-error}. The proof uses the
original compatible operator norms and allows normed input spaces
and Banach output spaces. The enclosed pole locations may differ,
and a readout residue may vanish. The strict residue margin ensures
visibility of the exact pole in the chosen readout.

\subsection{Pinned projects and logical audit}

The development comprises two separately compiled Lean~4 projects
that use mathlib~\cite{Mathlib4}.
\begin{description}[style=nextline]
\raggedright
\item[Root project.]
Lean~4.24.0 and mathlib revision
\href{https://github.com/leanprover-community/mathlib4/tree/f897ebcf72cd16f89ab4577d0c826cd14afaafc7}
{\texttt{f897ebcf72cd}}.
\item[Jordan companion.]
Lean~4.32.2, mathlib revision
\href{https://github.com/leanprover-community/mathlib4/tree/905b95818eb32af7874a58b427f50c1711a5e96c}
{\texttt{905b95818eb3}}, and \texttt{schoenflies-lean} revision
\href{https://github.com/alonamaloh/schoenflies-lean/tree/05a43d29cde026618777db3d4e4316204ccca237}
{\texttt{05a43d29cde0}}.
\end{description}
Each project rebuilds its proof chain using its recorded toolchain
and locked dependency manifest.

Each audit checks transitive axiom dependencies for every local
theorem and every compiled declaration in the \texttt{UV} namespace.
The companion also checks imported Schoenflies declarations.
The audits permit Lean's standard \texttt{propext},
\texttt{Classical.choice}, and \texttt{Quot.sound}, and reject
proof placeholders and project-specific axioms. The mathematical
admission hypotheses enter as premises of the formal theorems.
Both builds and audits passed in the
\href{https://github.com/GiovanniC137z/uv-lean/actions/runs/42544180224}
{continuous-integration run} for the cited revision.

\subsection{Reproduction, access, and interpretation}

With Git, elan, and Python~3.11 or newer installed, clone the public
repository and select the exact revision in Ref.~\cite{UVLean}:
\begingroup
\footnotesize
\begin{verbatim}
git clone https://github.com/GiovanniC137z/uv-lean.git
cd uv-lean
REV=24ca145d8241664896d09751488b90d2fe5077ba
git checkout --detach "$REV"
\end{verbatim}
\endgroup
Then run the following from that repository root.
\begingroup
\footnotesize
\begin{verbatim}
lake exe cache get
lake build
python3 scripts/audit.py
cd jordan
MATHLIB_NO_CACHE_ON_UPDATE=1 lake update
python3 scripts/fetch_cache.py
lake build
python3 scripts/audit.py
cd ..
python3 scripts/claims.py
\end{verbatim}
\endgroup
Cache downloads are optional. The final command checks the claim
inventory, the companion pins, and the SHA-256 digest of the
source snapshot archived as \texttt{paper/UV.tex} before this revision.
The repository README lists cache alternatives.

The verification repository and pinned mathematical dependencies are
public at the time of this revision. The cited continuous-integration
record reports success for the exact commit. This report is distinct
from rerunning the Lean projects during an editorial or numerical
review, and from verifying the added analytical results.

Assumptions~\ref{ass:OS-completion} and~\ref{ass:cut-matching}
still govern physical reconstruction. For an interacting theory,
the vertices, contact terms, insertions, infrared bounds, and
physical source maps must meet the stated conditions.
Convergence of the infinite perturbation series and unrestricted
sums over supports require the separate estimates described in
\ref{sec:dependence}.

\begingroup
\raggedright
\begin{thebibliography}{99}

\bibitem{Next3Kusuoka}
S.~Kusuoka and D.~Stroock, Applications of the Malliavin calculus. II,
J. Fac. Sci. Univ. Tokyo Sect. IA Math. \textbf{32} (1985), 1--76.
\href{https://doi.org/10.15083/00039520}{doi:10.15083/00039520}.

\bibitem{JointGluingGolubWelsch}
G.~H.~Golub and J.~H.~Welsch, Calculation of Gauss quadrature rules,
Math. Comp. \textbf{23} (1969), 221--230.
\href{https://doi.org/10.1090/S0025-5718-69-99647-1}{doi:10.1090/S0025-5718-69-99647-1}.

\bibitem{PhysicalGluingAbbas}
G.~Abbas, B.~Ananthanarayan, I.~Caprini, I.~Sentitemsu Imsong and
S.~Ramanan, Theory of unitarity bounds and low energy form factors,
European Physical Journal A \textbf{45} (2010), 389--399.
\href{https://doi.org/10.1140/epja/i2010-11010-5}{doi:10.1140/epja/i2010-11010-5}.
\href{https://arxiv.org/abs/1004.4257}{arXiv:1004.4257}.

\bibitem{PhysicalGluingWatson}
K.~M.~Watson, The effect of final state interactions on reaction
cross sections, Phys. Rev. \textbf{88}, 1163--1171 (1952).
\href{https://doi.org/10.1103/PhysRev.88.1163}{doi:10.1103/PhysRev.88.1163}.

\bibitem{NextSachs}
E.~Tomasi-Gustafsson, A.~Bianconi, and S.~Pacetti,
New fit of timelike proton electromagnetic formfactors from
$e^+e^-$ colliders, Phys.\ Rev.\ C \textbf{103} (2021), 035203.
\href{https://doi.org/10.1103/PhysRevC.103.035203}{doi:10.1103/PhysRevC.103.035203}.

\bibitem{NextHerglotz}
F.~Gesztesy and E.~Tsekanovskii, On matrix-valued Herglotz functions,
Math.\ Nachr.\ \textbf{218}(1) (2000), 61--138.
\href{https://doi.org/10.1002/1522-2616(200010)218:1\%3C61::AID-MANA61\%3E3.0.CO;2-D}{Journal publication}.
\href{https://arxiv.org/abs/funct-an/9712004}{arXiv:funct-an/9712004}.

\bibitem{WardBreitWheeler}
G.~Jikia and A.~Tkabladze,
QCD corrections to heavy quark pair production in polarized
\(\gamma\gamma\) collisions and the intermediate-mass Higgs boson signal,
Phys.\ Rev.\ D \textbf{54} (1996), 2030--2042,
\href{https://doi.org/10.1103/PhysRevD.54.2030}{doi:10.1103/PhysRevD.54.2030}.

\bibitem{WardHeatManual}
D.~V.~Vassilevich, Heat kernel expansion: user's manual,
Phys.\ Rep.\ \textbf{388} (2003), 279--360.
\href{https://doi.org/10.1016/j.physrep.2003.09.002}{doi:10.1016/j.physrep.2003.09.002}.

\bibitem{SchillingBernstein2012}
R.~L.~Schilling, R.~Song, and Z.~Vondra\v{c}ek,
\textit{Bernstein Functions: Theory and Applications},
2nd rev.\ and ext.\ ed., De Gruyter Studies in Mathematics, Vol.~37
(Walter de Gruyter, Berlin/Boston, 2012).
\href{https://doi.org/10.1515/9783110269338}{doi:10.1515/9783110269338}.

\bibitem{ChiapponeGHOUQFT2026}
G.~J.~Chiappone,
\emph{The Path-Enriched Chirally Graded Groupoid Extension of the
Hyper-Omniverse Unified Quantum Field Theory: A Hidden Ambient Geometric
Multiverse for Single-Parameter UV-Finite Particle Physics and
Vacuum-Creation Cosmology}, Zenodo preprint (2026).
\href{https://doi.org/10.5281/zenodo.24407696}{doi:\nolinkurl{10.5281/zenodo.24407696}}.

\bibitem{ChiapponePartI2026}
G.~J.~Chiappone,
\emph{Geometric-Time Locking and Propagating Spacetime Planes I:
Exact Relativistic Fibers and Noether Geo-Momentum}, Zenodo preprint (2026).
\href{https://doi.org/10.5281/zenodo.24407675}{doi:\nolinkurl{10.5281/zenodo.24407675}}.

\bibitem{ChiapponePartII2026}
G.~J.~Chiappone,
\emph{Geometric-Time Locking and Propagating Spacetime Planes II:
Geometric-Time Quantum Interfaces and the Normal-Transfer Program}, Zenodo preprint (2026).
\href{https://doi.org/10.5281/zenodo.24405252}{doi:\nolinkurl{10.5281/zenodo.24405252}}.

\bibitem{ChiapponePartXXVI2026}
G.~J.~Chiappone,
\emph{Geometric-Time Locking and Propagating Spacetime Planes XXVI:
Metastable Holonomy and Inverse-Margin Ultraviolet Criteria},
Zenodo preprint (2026).
\href{https://doi.org/10.5281/zenodo.24405653}{doi:\nolinkurl{10.5281/zenodo.24405653}}.

\bibitem{ChiapponePartXXIX2026}
G.~J.~Chiappone,
\emph{Geometric-Time Locking and Propagating Spacetime Planes XXIX:
The Path-Enriched Chiral-Cone Master Action and Geo-Phase Transfer},
Zenodo preprint (2026).
\href{https://doi.org/10.5281/zenodo.24405697}{doi:\nolinkurl{10.5281/zenodo.24405697}}.

\bibitem{ChiapponeIgnition}
G.~J.~Chiappone,
\emph{Omniverse QED geo--phase damping in the hydrodynamic window:
a first--principles, hypothesis--scoped fixed--depth response framework
for Rayleigh--Taylor damping in MagLIF},
author-supplied LaTeX revision of 12 November 2026,
\texttt{Ignition14(20261112-230918).tex}
(the front matter retains the original date 8 August 2026).
The 8 August Zenodo edition with this title is archived at
\href{https://doi.org/10.5281/zenodo.21856264}{doi:\nolinkurl{10.5281/zenodo.21856264}}.
The cited content is the O--QED action, hyper-propagators, fixed-depth
kernel, contact-containing response skeleton and kernel-off matching;
equation numbers are version dependent.
Related Zenodo preprint, 24 May 2026,
\emph{Omniverse QED geo--phase damping in the hydrodynamic window:
a first--principles route to Rayleigh--Taylor suppression for MagLIF},
\href{https://doi.org/10.5281/zenodo.20371104}{doi:\nolinkurl{10.5281/zenodo.20371104}}.

\bibitem{UVLean}
G.~J.~Chiappone,
\emph{UV Lean verification},
Lean~4 formalization with a proof inventory and CI audit,
GitHub (2026), commit
\href{https://github.com/GiovanniC137z/uv-lean/tree/24ca145d8241664896d09751488b90d2fe5077ba}
{\texttt{24ca145d8241}}.
\href{https://github.com/GiovanniC137z/uv-lean}
{\texttt{GiovanniC137z/uv-lean}}.
See \ref{sec:lean-verification} for access and dependency pins.

\bibitem{TongQFTFramework}
D.~Tong, \emph{Quantum Field Theory}, University of Cambridge lecture
notes, Michaelmas Terms 2006 and 2007, Secs.~3.6--3.7 and 6.3--6.6.
\href{https://www.damtp.cam.ac.uk/user/tong/qft/qft.pdf}{Lecture notes},
\href{https://www.damtp.cam.ac.uk/user/tong/qft.htm}{course page},
\href{https://www.damtp.cam.ac.uk/user/tong/qft/qfthtml/S3.html}{interacting fields},
\href{https://www.damtp.cam.ac.uk/user/tong/qft/qfthtml/S6.html}{quantum electrodynamics}.

\bibitem{StingaTorrea2010}
P.~R.~Stinga and J.~L.~Torrea,
Extension problem and Harnack's inequality for some fractional operators,
\emph{Commun. Partial Differential Equations} \textbf{35} (2010) 2092--2122.
\href{https://doi.org/10.1080/03605301003735680}{doi:\nolinkurl{10.1080/03605301003735680}}.

\bibitem{DLMF10}
\emph{NIST Digital Library of Mathematical Functions},
Release 1.2.8 (15 September 2026), F.~W.~J. Olver \emph{et al.} (eds.),
\S10.32, ``Integral Representations,'' especially
\href{https://dlmf.nist.gov/10.32.E10}{Eq.~10.32.10};
\href{https://dlmf.nist.gov/10.30.E2}{Eq.~10.30.2} and
\href{https://dlmf.nist.gov/10.40.E2}{Eq.~10.40.2} for the small- and
large-argument bounds used here (accessed 19 November 2026).

\bibitem{Widder1941}
D.~V.~Widder,
\emph{The Laplace Transform}, Princeton Mathematical Series, vol.~6,
Princeton University Press, 1941. Princeton Legacy Library
electronic reissue, 2016.
\href{https://doi.org/10.1515/9781400876457}{doi:\nolinkurl{10.1515/9781400876457}}.

\bibitem{Weinberg1960}
S.~Weinberg,
High-energy behavior in quantum field theory,
\emph{Phys.\ Rev.}\ \textbf{118} (1960) 838--849.
\href{https://doi.org/10.1103/PhysRev.118.838}{doi:\nolinkurl{10.1103/PhysRev.118.838}}.

\bibitem{Zimmermann1969}
W.~Zimmermann, Convergence of Bogoliubov's method of renormalization
in momentum space, \emph{Commun. Math. Phys.} \textbf{15}(3),
208--234 (1969).
\href{https://doi.org/10.1007/BF01645676}{doi:10.1007/BF01645676}.

\bibitem{Abbott1981}
L.~F.~Abbott,
The background field method beyond one loop,
\emph{Nucl.\ Phys.\ B} \textbf{185} (1981) 189--203.
\href{https://doi.org/10.1016/0550-3213(81)90371-0}{doi:\nolinkurl{10.1016/0550-3213(81)90371-0}}.

\bibitem{BRST1976}
C.~Becchi, A.~Rouet, and R.~Stora,
Renormalization of gauge theories,
\emph{Ann.\ Phys.}\ \textbf{98} (1976) 287--321.
\href{https://doi.org/10.1016/0003-4916(76)90156-1}{doi:\nolinkurl{10.1016/0003-4916(76)90156-1}}.

\bibitem{Tyutin1975}
I.~V.~Tyutin,
\emph{Gauge Invariance in Field Theory and Statistical Physics
in Operator Formalism}, P.~N.~Lebedev Physical Institute preprint
No.~39 (1975); arXiv reproduction (cited version v2,
15 December 2008),
\href{https://arxiv.org/abs/0812.0580v2}{arXiv:0812.0580v2}.

\bibitem{PiguetSorella}
O.~Piguet and S.~P.~Sorella,
\emph{Algebraic Renormalization: Perturbative Renormalization,
Symmetries and Anomalies}, Lecture Notes in Physics Monographs,
vol.~28, Springer, Berlin/Heidelberg (original edition 1995;
eBook published 11 September 2008).
\href{https://doi.org/10.1007/978-3-540-49192-7}{doi:\nolinkurl{10.1007/978-3-540-49192-7}}.

\bibitem{Slavnov1972}
A.~A.~Slavnov,
Ward identities in gauge theories,
\emph{Theor.\ Math.\ Phys.}\ \textbf{10} (1972) 99--104
(English translation of \emph{Teor.\ Mat.\ Fiz.} \textbf{10}
(1972) 153--161).
\href{https://doi.org/10.1007/BF01090719}{doi:\nolinkurl{10.1007/BF01090719}}.

\bibitem{Taylor1971}
J.~C.~Taylor,
Ward identities and charge renormalization of the Yang--Mills field,
\emph{Nucl.\ Phys.\ B} \textbf{33} (1971) 436--444.
\href{https://doi.org/10.1016/0550-3213(71)90297-5}{doi:\nolinkurl{10.1016/0550-3213(71)90297-5}}.

\bibitem{ZinnJustin1975}
J.~Zinn-Justin,
Renormalization of gauge theories,
in H.~Rollnik and K.~Dietz (eds.),
\emph{Trends in Elementary Particle Theory},
Lecture Notes in Physics, vol.~37, Springer, Berlin/Heidelberg
(1975), pp.~1--39.
\href{https://doi.org/10.1007/3-540-07160-1_1}{doi:\nolinkurl{10.1007/3-540-07160-1_1}}.

\bibitem{Nielsen1975}
N.~K.~Nielsen,
On the gauge dependence of spontaneous symmetry breaking in gauge
theories, \emph{Nucl.\ Phys.\ B} \textbf{101} (1975) 173--188.
\href{https://doi.org/10.1016/0550-3213(75)90301-6}{doi:\nolinkurl{10.1016/0550-3213(75)90301-6}}.

\bibitem{Bernstein1929}
S.~Bernstein,
Sur les fonctions absolument monotones,
\emph{Acta Math.}\ \textbf{52} (1929) 1--66.
\href{https://doi.org/10.1007/BF02592679}{doi:\nolinkurl{10.1007/BF02592679}}.

\bibitem{Landau1959}
L.~D.~Landau,
On analytic properties of vertex parts in quantum field theory,
\emph{Nucl.\ Phys.}\ \textbf{13} (1959) 181--192.
\href{https://doi.org/10.1016/0029-5582(59)90154-3}{doi:\nolinkurl{10.1016/0029-5582(59)90154-3}}.

\bibitem{OS1973}
K.~Osterwalder and R.~Schrader,
Axioms for Euclidean Green's functions,
\emph{Commun.\ Math.\ Phys.}\ \textbf{31} (1973) 83--112.
\href{https://doi.org/10.1007/BF01645738}{doi:\nolinkurl{10.1007/BF01645738}}.

\bibitem{OS1975}
K.~Osterwalder and R.~Schrader,
Axioms for Euclidean Green's functions II
(with an appendix by Stephen Summers),
\emph{Commun.\ Math.\ Phys.}\ \textbf{42} (1975) 281--305.
\href{https://doi.org/10.1007/BF01608978}{doi:\nolinkurl{10.1007/BF01608978}}.

\bibitem{Arici2018}
F.~Arici, D.~Becker, C.~Ripken, F.~Saueressig,
and W.~D.~van Suijlekom,
Reflection positivity in higher derivative scalar theories,
\emph{J.\ Math.\ Phys.}\ \textbf{59} (2018) 082302.
\href{https://doi.org/10.1063/1.5027231}{doi:\nolinkurl{10.1063/1.5027231}};
\href{https://arxiv.org/abs/1712.04308v1}{arXiv:1712.04308v1}
(original submission, 12 December 2017; the journal DOI identifies the published article).

\bibitem{StreaterWightman}
R.~F.~Streater and A.~S.~Wightman,
\emph{PCT, Spin and Statistics, and All That},
original edition, W.~A. Benjamin (1964); additions and corrections
(1978); Princeton Landmarks in Mathematics and Physics,
Princeton University Press corrected paperback (2000).
Electronic edition, 2016.
\href{https://doi.org/10.1515/9781400884230}{doi:\nolinkurl{10.1515/9781400884230}}.

\bibitem{DappiaggiEtAl2022}
C.~Dappiaggi, B.~A.~Ju\'arez-Aubry, and A.~Marta,
Ground state for the Klein--Gordon field in anti--de Sitter spacetime
with dynamical Wentzell boundary conditions,
\emph{Phys. Rev. D} \textbf{105} (2022) 105017.
\href{https://doi.org/10.1103/PhysRevD.105.105017}{doi:\nolinkurl{10.1103/PhysRevD.105.105017}}.

\bibitem{LSZ1955}
H.~Lehmann, K.~Symanzik, and W.~Zimmermann,
Zur Formulierung quantisierter Feldtheorien,
\emph{Il Nuovo Cimento} \textbf{1} (1955) 205--225.
\href{https://doi.org/10.1007/BF02731765}{doi:\nolinkurl{10.1007/BF02731765}}.

\bibitem{GaoGloriosoLiu2020}
P.~Gao, P.~Glorioso, and H.~Liu,
Ghostbusters: unitarity and causality of non-equilibrium effective field
theories, \emph{JHEP} \textbf{03} (2020) 040.
\href{https://doi.org/10.1007/JHEP03(2020)040}{doi:\nolinkurl{10.1007/JHEP03(2020)040}}.

\bibitem{BrisceseEtAl2024}
F.~Briscese, G.~Calcagni, L.~Modesto, and G.~Nardelli,
Form factors, spectral and K\"all\'en--Lehmann representation in nonlocal
quantum gravity, \emph{JHEP} \textbf{08} (2024) 204.
\href{https://doi.org/10.1007/JHEP08(2024)204}{doi:\nolinkurl{10.1007/JHEP08(2024)204}}.

\bibitem{Cutkosky1960}
R.~E.~Cutkosky,
Singularities and discontinuities of Feynman amplitudes,
\emph{J.\ Math.\ Phys.}\ \textbf{1} (1960) 429--433.
\href{https://doi.org/10.1063/1.1703676}{doi:\nolinkurl{10.1063/1.1703676}}.

\bibitem{BrosEpsteinGlaser1965}
J.~Bros, H.~Epstein, and V.~Glaser,
A proof of the crossing property for two-particle amplitudes in
general quantum field theory,
\emph{Commun.\ Math.\ Phys.}\ \textbf{1}(3) (1965), 240--264.
\href{https://doi.org/10.1007/BF01646307}{doi:\nolinkurl{10.1007/BF01646307}}.
\href{https://repo-archives.ihes.fr/FONDS_IHES/J_Seminaires_Conferences/2J_Seminaires/2J_RUEL_I/2J_RUEL_I_1/2J_RUEL_I_1.pdf}{Authors' March 1965 preprint (IH\'ES archive)}.

\bibitem{TerningS6}
J.~Terning,
Gauging nonlocal Lagrangians,
\emph{Phys. Rev. D} \textbf{44} (1991) 887--897.
\href{https://doi.org/10.1103/PhysRevD.44.887}{doi:\nolinkurl{10.1103/PhysRevD.44.887}}.

\bibitem{DLMF4W}
\emph{NIST Digital Library of Mathematical Functions},
Release 1.2.8 (15 September 2026), F.~W.~J. Olver \emph{et al.} (eds.),
\href{https://dlmf.nist.gov/4.13}{\S4.13}, Lambert $W$-Function
(accessed 19 November 2026).

\bibitem{ExcursionSchubert}
C.~Schubert, Perturbative quantum field theory in the string-inspired
formalism, \emph{Phys. Rep.} \textbf{355}(2--3) (2001), 73--234.
\href{https://doi.org/10.1016/S0370-1573(01)00013-8}{doi:10.1016/S0370-1573(01)00013-8}.

\bibitem{CarrierClusters}
B.~D.~Keister and W.~N.~Polyzou,
Model tests of cluster separability in relativistic quantum mechanics,
\emph{Physical Review C} \textbf{86} (2012), 014002.
\href{https://doi.org/10.1103/PhysRevC.86.014002}{doi:10.1103/PhysRevC.86.014002}.

\bibitem{NativeFullNeumann}
M.~Choulli, L.~Kayser and E.~M.~Ouhabaz, Observations on Gaussian
upper bounds for Neumann heat kernels, \emph{Bull. Aust. Math. Soc.}
\textbf{92} (2015), 429--439.
\href{https://doi.org/10.1017/S0004972715000611}{doi:10.1017/S0004972715000611}.

\bibitem{MonomialSteklov}
A.~Girouard and I.~Polterovich, Spectral geometry of the Steklov
problem (survey article), \emph{J. Spectr. Theory} \textbf{7}
(2017), 321--359. \href{https://doi.org/10.4171/JST/164}{doi:10.4171/JST/164}.

\bibitem{SeamMaximalFourier}
R.~Oberlin, A.~Seeger, T.~Tao, C.~Thiele and J.~Wright,
A variation norm Carleson theorem,
\emph{Journal of the European Mathematical Society}
\textbf{14} (2012), 421--464.
\href{https://doi.org/10.4171/JEMS/307}{doi:10.4171/JEMS/307}.

\bibitem{BulkLevyContinuity}
H.~Bierm\'e, O.~Durieu and Y.~Wang, Generalized random fields and
L\'evy's continuity theorem on the space of tempered distributions,
\emph{Communications on Stochastic Analysis} \textbf{12}, no.~4
(2018), Article~4.
\href{https://doi.org/10.31390/cosa.12.4.04}{doi:10.31390/cosa.12.4.04}.

\bibitem{ClosureDotsenkoVergeles}
V.~S.~Dotsenko and S.~N.~Vergeles, Renormalizability of phase factors
in non-abelian gauge theory, \emph{Nuclear Physics B} \textbf{169}
(1980), 527--546.
\href{https://doi.org/10.1016/0550-3213(80)90103-0}{doi:10.1016/0550-3213(80)90103-0}.

\bibitem{SpectrumLaineVuorinen}
M.~Laine and A.~Vuorinen, \emph{Basics of Thermal Field Theory:
A Tutorial on Perturbative Computations}, Lecture Notes in Physics
\textbf{925} (Springer, Cham, 2016).
\href{https://doi.org/10.1007/978-3-319-31933-9}{doi:10.1007/978-3-319-31933-9}.

\bibitem{SpectrumHopfield}
J.~J.~Hopfield, Theory of the contribution of excitons to the complex
dielectric constant of crystals, \emph{Phys. Rev.} \textbf{112}
(1958), 1555--1567.
\href{https://doi.org/10.1103/PhysRev.112.1555}{doi:10.1103/PhysRev.112.1555}.

\bibitem{SpectrumBrascampLieb}
H.~J.~Brascamp and E.~H.~Lieb, On extensions of the Brunn--Minkowski
and Pr\'ekopa--Leindler theorems, including inequalities for log
concave functions, and with an application to the diffusion equation,
\emph{J. Funct. Anal.} \textbf{22} (1976), 366--389.
\href{https://doi.org/10.1016/0022-1236(76)90004-5}{doi:10.1016/0022-1236(76)90004-5}.

\bibitem{SpectrumDerezinski2003}
J.~Derezi\'nski, \emph{Van Hove Hamiltonians -- Exactly Solvable Models of the
Infrared and Ultraviolet Problem}, Annales Henri Poincar\'e
\textbf{4} (2003), 713--738.
\href{https://doi.org/10.1007/s00023-003-0145-5}{doi:10.1007/s00023-003-0145-5}.

\bibitem{HiroshimaPF2007}
F.~Hiroshima, Fiber Hamiltonians in non-relativistic quantum electrodynamics,
\emph{J. Funct. Anal.} \textbf{252} (2007), 314--355.
\href{https://doi.org/10.1016/j.jfa.2007.06.006}{doi:10.1016/j.jfa.2007.06.006}.

\bibitem{PhysicalDenigSalme}
A.~Denig and G.~Salm\`e, ``Nucleon electromagnetic form factors in the
timelike region,'' \emph{Progress in Particle and Nuclear Physics}
\textbf{68} (2013), 113--157.
\href{https://doi.org/10.1016/j.ppnp.2012.09.005}{doi:10.1016/j.ppnp.2012.09.005}.

\bibitem{DualRevBogachev1998}
V.~I.~Bogachev, \emph{Gaussian Measures},
Mathematical Surveys and Monographs, vol.~62,
American Mathematical Society, Providence, RI, 1998.
\href{https://doi.org/10.1090/surv/062}{doi:10.1090/surv/062}.

\bibitem{SimonG5Normal}
B.~Simon, \emph{Functional Integration and Quantum Physics},
2nd ed., AMS Chelsea Publishing, 2005.
\href{https://doi.org/10.1090/chel/351}{doi:\nolinkurl{10.1090/chel/351}}.

\bibitem{SimonUVpeDet}
B.~Simon, Notes on infinite determinants of Hilbert space operators,
\emph{Adv. Math.} \textbf{24} (1977) 244--273.
\href{https://doi.org/10.1016/0001-8708(77)90057-3}{doi:\nolinkurl{10.1016/0001-8708(77)90057-3}}.

\bibitem{Schwinger1951}
J.~Schwinger,
On gauge invariance and vacuum polarization,
\emph{Phys.\ Rev.}\ \textbf{82} (1951) 664--679.
\href{https://doi.org/10.1103/PhysRev.82.664}{doi:\nolinkurl{10.1103/PhysRev.82.664}}.

\bibitem{DLMF6}
\emph{NIST Digital Library of Mathematical Functions},
Release 1.2.8 (15 September 2026), F.~W.~J. Olver \emph{et al.} (eds.),
Chapter~6 (exponential integrals: definitions, continuation,
and asymptotics); see
\href{https://dlmf.nist.gov/6.2}{\S6.2},
\href{https://dlmf.nist.gov/6.4}{\S6.4},
and \href{https://dlmf.nist.gov/6.12}{\S6.12}
(accessed 19 November 2026).

\bibitem{BallChiu1980}
J.~S.~Ball and T.-W.~Chiu,
Analytic properties of the vertex function in gauge theories.~I,
\emph{Phys.\ Rev.\ D}\ \textbf{22} (1980) 2542--2549.
\href{https://doi.org/10.1103/PhysRevD.22.2542}{doi:\nolinkurl{10.1103/PhysRevD.22.2542}}.

\bibitem{Beyn2012}
W.-J.~Beyn,
An integral method for solving nonlinear eigenvalue problems,
\emph{Linear Algebra Appl.}\ \textbf{436} (2012) 3839--3863.
\href{https://doi.org/10.1016/j.laa.2011.03.030}{doi:\nolinkurl{10.1016/j.laa.2011.03.030}};
\href{https://arxiv.org/abs/1003.1580v1}{arXiv:1003.1580v1}
(original submission, 8 March 2010; the journal DOI identifies the published article).

\bibitem{SchoenfliesLean}
\emph{The Jordan--Sch\"onflies Theorem in Lean}
(\texttt{alonamaloh/schoenflies-lean}), GitHub, commit
\href{https://github.com/alonamaloh/schoenflies-lean/tree/05a43d29cde026618777db3d4e4316204ccca237}
{\texttt{05a43d29cde0}}, Apache~2.0.
\href{https://github.com/alonamaloh/schoenflies-lean}
{Project repository}.

\bibitem{Mathlib4}
The mathlib Community,
\emph{mathlib4}, the mathematical library of Lean~4, GitHub.
\href{https://github.com/leanprover-community/mathlib4}
{Project repository}.
The revisions used here are given in \ref{sec:lean-verification}.

\bibitem{SevenKleppeWoodard}
G.~Kleppe and R.~P.~Woodard, Nonlocal Yang--Mills,
\emph{Nuclear Physics B} \textbf{388} (1992), 81--112.
\href{https://doi.org/10.1016/0550-3213(92)90546-N}{doi:10.1016/0550-3213(92)90546-N}.

\bibitem{SevenNeumann}
P.~Gyrya and L.~Saloff-Coste,
\emph{Neumann and Dirichlet Heat Kernels in Inner Uniform Domains},
Ast\'erisque \textbf{336} (2011).
\href{https://doi.org/10.24033/ast.899}{doi:10.24033/ast.899}.

\end{thebibliography}
\endgroup
\end{document}
